% EDITED 2026-10-29 scripts/patch_continuity_tower_premise_2026-10-29T1500.py
% EDITED 2026-09-15 scripts/patch_value_continuum_declared_cencov_2026-09-15T1540.py
% EDITED 2026-09-15 scripts/patch_continuity_connected_state_space_2026-09-15T1430.py
% =====================================================================
%  SECTION: Continuity Is Not an Axiom
%  Drafted 2026-07-04 from continuity_limit_realisability_2026-07-04.md
%  (now in project documents) and the May record
%  (dynamics_companion_fanout_unitarity_2026-05-20T1430.md).
%
%  PLACEMENT in main.tex: a standalone section between self-reference
%  and division-algebras --
%      % EDITED 2026-10-29 scripts/patch_continuity_tower_premise_2026-10-29T1500.py
% EDITED 2026-09-20 scripts/patch_born_cost_ledger_v2_2026-09-20T0900.py
% =====================================================================
%  SECTION: Self-Reference: One Principle
%  Drafted 2026-07-04. The paper's foundational section; resolves the
%  largest dangling label in the manuscript (sec:self-reference).
%
%  PLACEMENT in main.tex: FIRST of the foundations chain --
%      % EDITED 2026-10-29 scripts/patch_continuity_tower_premise_2026-10-29T1500.py
% EDITED 2026-09-20 scripts/patch_born_cost_ledger_v2_2026-09-20T0900.py
% =====================================================================
%  SECTION: Self-Reference: One Principle
%  Drafted 2026-07-04. The paper's foundational section; resolves the
%  largest dangling label in the manuscript (sec:self-reference).
%
%  PLACEMENT in main.tex: FIRST of the foundations chain --
%      % EDITED 2026-10-29 scripts/patch_continuity_tower_premise_2026-10-29T1500.py
% EDITED 2026-09-20 scripts/patch_born_cost_ledger_v2_2026-09-20T0900.py
% =====================================================================
%  SECTION: Self-Reference: One Principle
%  Drafted 2026-07-04. The paper's foundational section; resolves the
%  largest dangling label in the manuscript (sec:self-reference).
%
%  PLACEMENT in main.tex: FIRST of the foundations chain --
%      \input{self_reference_section_2026-07-04}
%      \input{continuity_section_2026-07-04}
%      \input{division_algebras_section_2026-07-03}
%      \input{hopf_justification_section_2026-06-08}
%
%  SOURCES (re-read 2026-07-04): lawvere_inside_outside_fixed_point_
%  theorem_2026-05-04T1806.md (the spine: the two-principles reduction,
%  the theorem, the sigma candidate, the open items -> op minted per
%  D9); research_summary_neg_prob.md (negativity route; lamp (1+i)/2);
%  dynamics_companion_fanout_unitarity (legs 1a/1c; anti-events; the
%  Brukner-Zeilinger rebuttal per D6); Pure_states_and_mixed_states
%  (local tomography derived; QBism contrast per D7).
%
%  DECISIONS ENCODED (JS 2026-07-04): D1 title as below; D2 Wheeler1990
%  cited here as the title anchor; D3 lineage Popper -> Breuer ->
%  Svozil -> Szangolies (Dalla Chiara omitted); D4 Frauchiger_2018
%  cited here, dissolved in sec:wigners-friend; D5 Kirkwood-Dirac
%  heritage via the Arvidsson-Shukur 2024 review as one-stop; D6
%  Brukner-Zeilinger + the framework rebuttal IN ("Shannon is
%  post-FANOUT"); D7 QBism contrast IN; D8 Aaronson private
%  communication OMITTED (published island paper only); D9 one op
%  minted (op:sigma-functor), cross-linked to op:hopf-frobenius.
%
%  CITATIONS -- existing keys: Wheeler1990, small_2006, Godel1931,
%  Lawvere_1969, Yanofsky2003, CoeckeKissinger2017, Rovelli1996a,
%  Popper1950a, Popper1950b, Breuer1995, svozil_k_2018,
%  szangolies_j_2018, Frauchiger_2018, dirac1942, Feynman_r_82,
%  Feynman1987, Youssef1994 [SWAP-FLAG: intended paper is "Quantum
%  mechanics as complex probability theory", Mod. Phys. Lett. A 9
%  (1994) 2571 -- if that lives under youssef1994quantum, swap],
%  caticha_a_2006, goyal_knuth_skilling_2010, Aaronson2004h, Hardy2001a,
%  muller_p_2011, Thomson_1954, benacerraf1962tasks, cantor1891diagonal,
%  Turing_1937.
%  NEW keys (RIS in self_reference_refs_2026-07-04.ris, all verified
%  this session): tikochinsky1988feynman, renou2021real, chen2022ruling,
%  wu2022refutation, hoffreumonwoods2026real, arvidssonshukur2024kirkwood,
%  bruknerzeilinger2001shannon, fuchsschack2013qbism,
%  shannon1948mathematical.
%
%  CROSS-REFS: forward (fine under two-pass LaTeX): sec:continuity,
%  sec:division-algebras, sec:hopf-justification, prop:f1-dictionary,
%  sec:wigners-friend, op:limit-realisability, op:hopf-frobenius.
%  Still dangling after this section: sec:fanout-boundary (and its
%  sibling label sec:fanout-direction) -- both DRAFTED 2026-07-06.
%
%  AUDITED 2026-07-06 (section-by-section final pass, first stop):
%  proof of prop:inside-outside re-verified by hand; attribution chain
%  clean (Burgin/Khrennikov correctly absent); all header-declared
%  citations confirmed used; lamp arithmetic and antecedents verified;
%  FANOUT smallcaps consistent; no pre-definition use of "frame".
%  Fixes: stray Unicode em-dash (QBism paragraph, the file's only
%  non-ASCII character) -> house ---; trace-blind two-way stitch added
%  to the Born subsection ("one privacy, two duties");
%  op:sigma-functor gains its first standing constraint (F1
%  inheritance via the op:hopf-frobenius discharge) and its previously
%  one-directional local-tomography cargo made explicit.
% =====================================================================

\section{Self-Reference: One Principle}
\label{sec:self-reference-2}

Wheeler asked for it from bit: for physics derived from information, ``every it --- every particle,
every field of force'' obtained from binary answers to yes-or-no questions \autocite{Wheeler1990}.
This paper's route to that programme runs through G\"{o}del: the single principle from which
everything else here is derived, deferred, or priced is \emph{self-reference} --- the fact that any
complete physical description, made from inside the world it describes, must describe the describer.
This section states the principle, proves that the framework's second historical commitment follows
from it as a theorem rather than standing beside it as an axiom, and then extracts the price of the
principle in four instalments: probabilities whose parts can cancel, amplitudes valued in
$\mathbb{C}$, the Born rule, and the local tomography of composite systems. What the section cannot
yet derive is stated as precisely as what it can, in one numbered problem and in the ledger at the
end.

\subsection{Events, not objects}

The framework's ontology is relational: the fundamental entities are events and the relations
between them, and objects are the bookkeeping of relational invariants --- what remains stable
across many relations. The kinship with relational quantum mechanics \autocite{Rovelli1996a} is
genuine and acknowledged: facts are facts relative to the systems that register them. The
difference, developed over the rest of this paper, is that the present framework grounds the
relational structure in a single logical principle and then carries it, through the division
algebras and the Hopf fibrations, all the way to the gauge group and the particle spectrum ---
relations are not only the ontology but the derivation engine. On this reading a measurement is not
the revelation of a pre-existing property; it is the creation of a relation, and the later sections
of this paper read that creation geometrically (\S\ref{sec:hopf-justification}) and even
gravitationally. Here we need only the logical core: whatever describes the world is an event in
the world it describes.

\subsection{One principle, and a theorem}

The framework historically carried two foundational commitments: self-reference, and the
\emph{Inside/Outside Equivalence Principle} --- that the laws of physics must be identical whether
formulated by an embedded participant or by an external observer. Two principles are one too many,
and this subsection reduces them.

The reduction corrects the framework's own origin, and we make the correction in print. The 2005
form of the argument \autocite{small_2006} ran a G\"{o}delian chain: a self-referential process,
being complete, must by the incompleteness theorem \autocite{Godel1931} be inconsistent; an
inconsistent system can negate itself; and the negation of self-reference is the production of
causal sequence. The chain captures the right intuition and is retained here as motivation, but it
is not sound as a proof: classical inconsistency yields \emph{ex falso quodlibet} --- every
proposition provable, every causal sequence an output, no constraint at all. The load-bearing
argument is different, and better: a fixed-point theorem.

Let $\mathcal{C}$ be the category of physical descriptions and admissible processes --- concretely,
the dagger-compact-closed setting with classical structures in which categorical quantum mechanics,
and the whole of this paper, works \autocite{CoeckeKissinger2017}. Lawvere showed that in a
sufficiently rich such setting, any weakly point-surjective self-representation forces every
endomap to have a fixed point, and Yanofsky showed that one diagonal argument of this shape
underlies Cantor's, G\"{o}del's, Turing's and Tarski's theorems alike
\autocite{Lawvere_1969,Yanofsky2003,cantor1891diagonal}. The framework's single commitment is that
$\mathcal{C}$ contains a \emph{self-representing object} --- an object capable of encoding the
processes on itself. That is what ``self-reference is foundational'' means, stated categorically,
and it is the surjectivity hypothesis of the diagonal lemma. Now let
$\sigma \colon \mathcal{C} \to \mathcal{C}$ be the \emph{inside/outside involution}: the operation
exchanging the two complementary presentations of any description --- the inside presentation, as
participant-indexed total space in which the observer is part of the system, and the outside
presentation, as a base-times-fibre decomposition over a classical parameter space. That
$\sigma^{2} = \mathrm{id}$ is not an axiom but a coherence requirement: the outside view of the
inside view of a thing is the thing. The candidate concrete realisation of $\sigma$ is the
Hopf-duality swap --- total space against base-times-fibre for the three fibrations --- which gives
the involution computable content and ties this subsection to \S\ref{sec:hopf-justification}.

\begin{proposition}[Inside/outside fixed point]\label{prop:inside-outside}
Let $\mathcal{C}$ satisfy the richness condition above and contain a self-representing object, and
let $\sigma$ be an involution on $\mathcal{C}$. Then the fixed-point locus
$\mathrm{Fix}(\sigma)$ is non-empty: there exist descriptions invariant under the exchange of
inside and outside perspective.
\end{proposition}

\begin{proof}
Lawvere's construction, applied to $\sigma$ in place of negation: given a weakly point-surjective
$\varphi \colon A \to Y^{A}$, form $g(a) = \sigma\!\left(\varphi(a)(a)\right)$; point-surjectivity
supplies $a^{*}$ with $\varphi(a^{*}) = g$, and then $\varphi(a^{*})(a^{*})$ is a fixed point of
$\sigma$ \autocite{Lawvere_1969,Yanofsky2003}. Where the same construction applied to Boolean
negation yields paradox --- negation has no fixed point on a two-element type, so the conclusion
explodes the premises --- the inside/outside involution is an honest involution on a rich type, and
the conclusion is a constructive fixed point rather than a contradiction.
\end{proof}

\textbf{Definition.} \emph{Physical law is the structure carried by} $\mathrm{Fix}(\sigma)$.

The Inside/Outside Equivalence Principle is thereby demoted from axiom to theorem: physical law is
invariant under the exchange of perspectives \emph{because it is the fixed-point set of that
exchange}, and the fixed-point set is forced to be non-empty by self-reference alone. One
commitment remains foundational; everything the second commitment asserted is now derived. The
reader who accepts categorical quantum mechanics as an existing programme has already accepted the
machinery; the framework adds one premise and collects the rest. And the fixed-point locus is not
empty of structure: it inherits the Frobenius-algebra apparatus of $\mathcal{C}$, so that the
Bell, GHZ and W primitives --- and with them the particle classification of the later sections ---
live inside $\mathrm{Fix}(\sigma)$ from the start. What is proved above is non-emptiness given the
setting; what is assumed is the setting itself and the existence of $\sigma$ as a genuine
involution; what is open is the rigour of both, and it is priced now rather than hidden:

\begin{openproblem}[The involution made functorial]\label{op:sigma-functor}
Fix the category $\mathcal{C}$ rigorously --- the dagger-compact-closed setting with classical
structures is the candidate --- and define $\sigma$ on morphisms, not only on objects, the natural
route being the induced fibre--base exchange on spider diagrams; establish that $\sigma$ is an (at
worst pseudo-)involutive endofunctor; and prove or refute the identification of
$\mathrm{Fix}(\sigma)$ with the three Hopf cross-sections, under which Adams's theorem would become
a corollary of the fixed-point classification. This problem is cross-linked with Open
Problem~\ref{op:hopf-frobenius}: under an affirmative resolution of both, the parallelisable
terminus, the fixed-point classification, and compositional completeness are three statements of
one fact, and progress on either problem constrains the other. One constraint is already in place:
at the complex rung, Proposition~\ref{prop:f1-dictionary} supplies the dictionary between
classical structures and the fibration's lifted skeletons --- the clause of
Open Problem~\ref{op:hopf-frobenius} it discharges --- and any resolution here must reproduce it.
The problem also carries a piece of cargo recorded later in this section: the formalisation, within
the fixed category, of the perspective-accounting that derives local tomography.
\end{openproblem}

\subsection{The lineage}

The obstruction this section builds on is old, and the framework claims its completion, not its
discovery. Popper argued in 1950 that a predictor embedded in the world has irreducibly inferior
knowledge of itself, because it cannot impart information about itself without, by that very act,
rendering the information out of date \autocite{Popper1950a,Popper1950b} --- the passage is quoted
where the framework makes the obstruction structural (\S\ref{sec:wigners-friend}). Breuer proved
the impossibility of accurate state self-measurement within quantum theory \autocite{Breuer1995};
Svozil's monograph surveys the whole territory where physics meets self-reference, undecidability
and Chaitin-style algorithmic limits \autocite{svozil_k_2018}; Szangolies sharpened the modern form
as epistemic horizons \autocite{szangolies_j_2018}; and the Frauchiger--Renner theorem is the
contemporary crisis-form --- nested observers deriving contradictory records when each treats the
others' quantum descriptions as universally valid \autocite{Frauchiger_2018} --- which
\S\ref{sec:wigners-friend} dissolves relationally rather than paradoxically. This framework's
addition to the lineage is twofold. Structurally, Proposition~\ref{prop:f1-dictionary} converts
Popper's temporal obsolescence into algebraic cancellation: the datum a copier cannot put into its
own record is exactly its fibre commitment, which does not arrive stale --- it cancels. And
foundationally, the present section runs the lineage in reverse: where the tradition treats
self-reference as a limit on knowledge \emph{within} physics, the framework takes it as the
generator of physics, and derives the quantum formalism from it.

\subsection{The price is negative probability}\label{sec:complex-amplitudes-2}

What does it cost a probabilistic description to contain its own describer? The first instalment is
the sign. In the regime before any record exists, every event has an anti-event --- the description
must be able to undo what it has provisionally done, or the self-description could never be
revised into consistency --- and a probability calculus closed under undoing cannot be confined to
$[0,1]$: it is a signed calculus, in which the parts of a probability can oppose and cancel. The
framework did not discover negativity as bookkeeping; Dirac admitted negative probabilities as
formal intermediaries \autocite{dirac1942}, Feynman defended them as consistent whenever the final,
recorded probabilities stay positive \autocite{Feynman_r_82,Feynman1987}, and the Kirkwood--Dirac
distribution has carried complex quasiprobability from 1933 to the present, as a recent
comprehensive review documents \autocite{arvidssonshukur2024kirkwood}. What the framework adds is
the forcing: negativity is not a permissible convenience but the first price of self-referential
closure, and it is charged exactly and only in the pre-record regime --- which is why every
negativity in the tradition lives in intermediate, unrecorded quantities and never in a final
record.

The second instalment is complexity. A signed information measure must extend Shannon's
$S = -\ln p$ off the positive axis, and the analytic continuation of the logarithm carries $p$ into
the complex plane: amplitudes, with the phase as the continued information. Here too the extension
has a prior tradition, and the attributions are owed precisely. Tikochinsky derived Feynman's rules
for complex amplitudes from consistency of composition \autocite{tikochinsky1988feynman}; Youssef
reformulated quantum mechanics as complex probability theory outright \autocite{Youssef1994};
Caticha carried objective amplitudes to Bayesian probabilities \autocite{caticha_a_2006}; and
Goyal, Knuth and Skilling derived Feynman's rules from a pair-postulate without assuming
analyticity \autocite{goyal_knuth_skilling_2010}. The extension is prior art. The framework's
claim is the compulsion: the complex extension is not a modelling choice among exotic probability
theories but the unique continuation forced by self-referential closure --- the same forcing that
Aaronson glimpsed when he argued quantum mechanics is what negative probability inevitably becomes
\autocite{Aaronson2004h}.

Why does the continuation stop at $\mathbb{C}$? Downward, the real field is not available: real
quantum theory is experimentally distinguishable from complex quantum theory in network scenarios
and has been disfavoured by experiment
\autocite{renou2021real,chen2022ruling,wu2022refutation}; the recent objection that the exclusion
rests on locally tomographic composition \autocite{hoffreumonwoods2026real} is, for this framework,
not an escape but a confirmation of where the weight sits --- local tomography is \emph{derived}
below, not assumed, so the framework closes from its own resources exactly the door the objection
holds open. Upward, the quaternions and octonions fail as a \emph{base} for the information
continuation: the continued measure is a logarithm, and beyond $\mathbb{C}$ the logarithm of a
product is no longer the sum of logarithms --- non-commutativity severs
$\log(ab) = \log a + \log b$ --- so the additivity that makes it an information measure is lost.
Aaronson's parameter-counting and signalling arguments against a quaternionic base point the same
way \autocite{Aaronson2004h}. The higher algebras are not banished; they return as the
\emph{internal} algebras of the register, which is the business of \S\ref{sec:division-algebras}.
The base is $\mathbb{C}$, for informational reasons, and stays $\mathbb{C}$.

\subsection{The Born rule from fibre privacy}\label{sec:born-rule}

The third instalment is the quadratic rule. An amplitude carries one datum more than any record of
it can: its overall phase. This was asserted throughout the framework's development as the
observer's private, self-referential ignorance; it is now a proved statement, in exactly the form
needed. By Proposition~\ref{prop:f1-dictionary}, a classical structure --- the copy operation any
record requires --- is committed to a specific lift of its skeleton, and that commitment is absent
from every classical record the structure induces: the copier is fibre-committed, the record is
fibre-blind, and the fibre freedom is gauge for every observer. (The computation extends verbatim
from one qubit to the full register, where Proposition~\ref{prop:trace-blind} makes it the ground
of hypercharge's tracelessness --- one privacy, two duties.) Popper's obsolescence, made
cancellation. It follows that anything recordable is invariant under the fibre --- observables and
their statistics must be $U(1)$-invariant functions of the amplitude. The lowest-degree polynomial invariant
of a complex amplitude is the quadratic $|\psi|^{2}$ --- the non-polynomial invariant $|\psi|$,
of degree one, is excluded by the analyticity declared at \S\ref{sec:complex-amplitudes} --- and
normalisation fixes the coefficient:
probability is the minimal record-compatible function of the amplitude,
\[
p \;=\; |\psi|^{2},
\]
the Born rule as the cheapest invariant the fibre permits. The same square is reached from the
cost side in \S\ref{sec:born-cost-ledger}, where a record is the closed loop of a history and its
return leg and the exponent is the number of legs; neither route uses the lattice of projections
that Gleason's theorem needs, and both hold for a single qubit. The derivation is the framework's
own;
its premise --- that the fibre datum is structurally excluded from records --- is
Proposition~\ref{prop:f1-dictionary}, and its geometric reading --- probability as base-point data
under the Hopf projection --- is developed in \S\ref{sec:hopf-justification}.

One standing objection to information-theoretic foundations is answered here, because this
framework needs the answer. Brukner and Zeilinger argued that Shannon information is conceptually
inadequate for quantum measurements: it is basis-dependent, and its operational meaning presupposes
ignorance of a definite but unknown fact, which the quantum case does not supply
\autocite{bruknerzeilinger2001shannon}. The framework's reply is structural, and it is a
demarcation rather than a rebuttal. Shannon's precondition fails exactly where they say it fails
--- in the pre-record regime, where there is no definite unknown to be ignorant of --- and there
the operative measure is the quadratic invariant just derived. Shannon's axioms hold exactly where
the precondition holds --- on records, after \textsc{fanout}, where outcomes are definite and
ignorance is classical --- and there the uniqueness theorem forces $-\log p$
\autocite{shannon1948mathematical}. The basis-dependence dissolves with the same stroke: a basis
\emph{is} selected precisely at \textsc{fanout}, so the basis-dependent measure applies only where
a basis exists. Shannon is not inadequate; Shannon is post-\textsc{fanout}. The framework thereby
keeps both measures, each on its own side of the boundary, and uses the additivity of the
post-record measure --- the very property the previous subsection used to pin the base field ---
with a clear conscience.

\subsection{Thomson's lamp, resolved}

The section owes \S\ref{sec:continuity} a debt incurred there: the lamp's amplitude-side limit.
Thomson's lamp --- toggled at times $1 - 2^{-n}$, its state at $t = 1$ demanded --- has no
classical answer: the task description fixes every finite stage and thereby fixes no limit at all
\autocite{Thomson_1954,benacerraf1962tasks}, which \S\ref{sec:continuity} read as the proof that
the record algebra is not limit-closed. On the amplitude side the limit exists, and the calculus of
this section computes it. The toggle exchanges the two amplitudes; the self-consistent state of the
completed task is the fixed point of the exchange in the complex calculus, and solving it under the
signed continuation gives
\[
P(\mathrm{on}) = \tfrac{1 + i}{2}, \qquad P(\mathrm{off}) = \tfrac{1 - i}{2},
\]
a superposed fixed point whose imaginary part carries the unrecordable datum --- the phase of the
endless alternation, when the lamp \emph{would have} stopped. Two checks certify that the
assignment is calculus rather than decoration. The pair is normalised,
$P(\mathrm{on}) + P(\mathrm{off}) = 1$; and the Born rule of the previous subsection, applied to
the fixed point, returns $|P(\mathrm{on})|^{2} = \tfrac{1}{2}$ --- exactly the symmetric classical
ignorance that is the only sane record-level answer, recovered rather than imposed. The lamp,
which entered the framework at its very beginning as the place where classical computation meets
its boundary, thus exits it as a worked example: undetermined as a record, determinate as an
amplitude, the two sides split precisely along \textsc{fanout} --- and the diagonal construction it
physically instantiates \autocite{cantor1891diagonal,Turing_1937} is the same one that powers
Proposition~\ref{prop:inside-outside}. The status flag is owed here as everywhere: the
undetermination is theorem-level; the specific fixed-point value is the framework's complex
calculus, structural, offered with its two consistency checks shown.

\subsection{Local tomography, derived}\label{sec:local-tomography}

The fourth instalment concerns composites, and it discharges two promises at once. The
Inside/Outside Equivalence Principle --- now a theorem --- requires that an external observer's
characterisation of a composite system be recoverable by the embedded participants, and the
embedded participants have access only to local measurements and their correlations. Equivalence
of the two perspectives therefore entails that the state of a composite is determined by the
statistics of local measurements: \emph{local tomography}, obtained as a consequence rather than
imposed as an axiom. This converts a postulate of the reconstruction programmes
\autocite{Hardy2001a,mullerp2011} into a theorem-shaped consequence of the one principle ---
discharging the caveat recorded in \S\ref{sec:continuity} --- and it closes, from the framework's
own resources, the escape route in the real-number debate noted above: the assumption the
objection of \autocite{hoffreumonwoods2026real} contests is here not an assumption. The status is
flagged plainly: the derivation is structural, one paragraph of perspective-accounting, and its
formalisation within the categorical setting belongs to Open Problem~\ref{op:sigma-functor}.

\subsection{Contrast: QBism, and the ledger}

The nearest neighbour in spirit deserves a paragraph of honest contrast. QBism
\autocite{fuchsschack2013qbism} holds, as this framework does, that quantum probabilities are an
agent's --- that the theory is a calculus of expectation, not a catalogue of properties. What QBism
does not say is why the agent's expectations must be complex-valued and quadratically squared; the
Born rule enters as a normative addition to Bayesian coherence, and the complex amplitude as an
unexplained representation. This section is, in effect, the missing why: the agent's ignorance is
self-referential --- the agent is inside the description --- and self-referential ignorance is
complex, with the quadratic rule as the cheapest record-compatible invariant of it. Where QBism
stops at the agent, the framework derives the agent's calculus.

The ledger, then. Delivered in this section: the reduction of two principles to one
(Proposition~\ref{prop:inside-outside}); the pre-record regime with its anti-events, handed to
\S\ref{sec:continuity} as the five-face regime whose continuity is priced there
(Open Problem~\ref{op:limit-realisability}); the complex base, forced downward by experiment and
upward by the additivity of information; the Born rule from fibre privacy; the lamp's fixed point;
and local tomography. Owed, and stated where they are owed: the functorial rigour of the involution
and the fixed-point classification (Open Problem~\ref{op:sigma-functor}); the remaining premise of continuity (Open Problem~\ref{op:limit-realisability}, in
\S\ref{sec:continuity}); and the
rung-by-rung compulsion of the algebra ladder, posed as Open
Problem~\ref{op:ladder-compulsion} in \S\ref{sec:division-algebras}. The
hand-on is the paper's spine: the $\mathbb{C}$ this section derives natively and the $\mathbb{C}$
that \S\ref{sec:continuity} underwrites through Hardy's theorem are the same field, the floor of
the ladder that \S\ref{sec:division-algebras} climbs and \S\ref{sec:hopf-justification}
geometrises. One principle; everything else is theorem, derivation, or a debt with a stated
invoice.

%      % EDITED 2026-10-29 scripts/patch_continuity_tower_premise_2026-10-29T1500.py
% EDITED 2026-09-15 scripts/patch_value_continuum_declared_cencov_2026-09-15T1540.py
% EDITED 2026-09-15 scripts/patch_continuity_connected_state_space_2026-09-15T1430.py
% =====================================================================
%  SECTION: Continuity Is Not an Axiom
%  Drafted 2026-07-04 from continuity_limit_realisability_2026-07-04.md
%  (now in project documents) and the May record
%  (dynamics_companion_fanout_unitarity_2026-05-20T1430.md).
%
%  PLACEMENT in main.tex: a standalone section between self-reference
%  and division-algebras --
%      \input{self_reference_section_2026-07-04}
%      \input{continuity_section_2026-07-04}
%      \input{division_algebras_section_2026-07-03}
%  (Architecture call SETTLED by the sync manifest: standalone.)
%
%  CITATIONS -- existing keys (as supplied by JS 2026-07-04):
%    Hardy2001a, Aaronson2004h, Thomson_1954, Lawvere_1969,
%    Turing_1937, muller_p_2011, Brukner2009.
%  NEW keys (RIS in continuity_refs_2026-07-04.ris):
%    cantor1891diagonal, benacerraf1962tasks, broeckertomdieck1985compact.
%  Deliberately excluded per JS: Chiribella--D'Ariano--Perinotti;
%  Bargmann/Wigner's theorem. Yanofsky2003 now paired on the
%  Lawvere-shaped clause (the pending-key excuse was void: the key
%  exists and sec:self-reference uses it).
%
%  CROSS-REFS: sec:self-reference now DRAFTED (2026-07-04, later the
%  same day) -- no longer dangling; only sec:fanout-boundary (and its
%  sibling sec:fanout-direction) -- DRAFTED 2026-07-06.
%
%  AUDITED 2026-07-06 (final pass, second stop): both proofs
%  re-verified by hand (Lagrange divisibility sentence exact; the
%  Baire countable-escape chain sound); caveat-discharge with
%  sec:self-reference confirmed two-way (already present); file pure
%  ASCII; no pre-definition "frame"; all citations used as declared.
%  Fixes: Yanofsky2003 paired on the Lawvere clause; the op's two
%  candidate routes anchored to printed propositions (f1-dictionary
%  for the divisibility premise's referent, inside-outside for the
%  closedness demand) -- statuses unchanged, referents sharpened;
%  header staleness corrected. Deliberate FORWARD references --
%  prop:f1-dictionary and sec:hopf-justification (defined later in the
%  paper; fine under two-pass LaTeX and flagged in prose),
%  sec:division-algebras, sec:space-of-computations.
%
%  DECISIONS ENCODED: op:limit-realisability minted (the owed arrow,
%  priced: exactly two derivations, divisibility and closedness);
%  Lagrange used twice (finite exclusion; component-group triviality);
%  the (1+i)/2 lamp value NOT stated here -- "a superposed fixed point"
%  with the complex-probability treatment deferred to
%  sec:self-reference; the two-branch over-determination stated to
%  pre-empt a circularity reading; Thomson spelling throughout.
% =====================================================================

\section{Continuity Is Not an Axiom}
\label{sec:continuity}

Hardy derives quantum theory from five axioms \autocite{Hardy2001a}. Four of them are satisfied
by classical probability theory; the fifth --- that a \emph{continuous} reversible
transformation connects any two pure states --- is the entire difference, the single postulate
that turns the classical simplex into the complex Hilbert space. Every reconstruction programme
pays an equivalent toll: continuous reversibility in Masanes and M\"{u}ller
\autocite{mullerp2011}, the corresponding postulate in Daki\'{c} and Brukner
\autocite{Brukner2009}. The axiom is load-bearing, and in every programme it is unexplained: it
is where the quantumness is put in by hand. The claim of this section is that in the present
framework it is not an axiom. Continuity is the fifth face of the pre-measurement regime that self-reference forces, and what
stands between the framework and a full derivation is not a postulate about values but one
premise about transformations --- that every transformation the inside performs has a stationary
outside presentation --- derived in outline below from the inside/outside involution of
\S\ref{sec:self-reference} and stated as an open problem at the end of the section with nothing
hidden in it. The continuum itself, the completion that every probability theory adopts, is
named for what it is: the real line, shared with classical probability and carrying no quantum
content.

One structural remark prevents a misreading. The framework reaches complex amplitudes natively in
\S\ref{sec:self-reference}, by the negative-probability route; the present section reaches the same
destination by an independent branch --- the regime's continuity, then Hardy's theorem. The two
branches converging on one structure is over-determination, not circularity, and the convergence is
itself evidence that the destination is forced.

\subsection{Five faces of one regime}

The regime before any measurement --- before any \textsc{fanout} --- can be characterised five
equivalent ways: negative probability is available; anti-events exist, every event being undoable;
no classical record exists; no preferred basis is selected; and a continuous reversible
transformation connects any two pure states. These are not five premises to be combined. They are
one regime under five descriptions, and \textsc{fanout} (\S\ref{sec:fanout-boundary}) is the single
event that fails all five at once: it writes the record, selects the basis, kills the anti-events,
discretises the transformation group, and switches off the negativity. The equivalence of the first
two faces is the anti-event identity of \S\ref{sec:self-reference}; the third is the boundary's
definition. The work of this section is the equivalence of the last two faces, and then the route
from the regime's existence to its continuity. The fourth face has, since
Proposition~\ref{prop:f1-dictionary} (\S\ref{sec:hopf-justification} --- a deliberate forward
reference), an exact form: a basis-selection is the commitment of a classical structure, a lifted
antipodal skeleton, so ``no preferred basis'' means precisely that no copy structure has fired.

\subsection{Basis selection is the collapse of continuity}

\begin{proposition}[Continuity forbids a selected basis]\label{prop:basis-discreteness}
If a finite set of pure states is privileged --- a preferred basis --- then every reversible
transformation respecting the privilege carries the set to itself and acts on it as a permutation;
the respecting transformations therefore realise only a finite group of motions among the
distinguished states, and no continuous family of them connects distinct pure states of the set.
Contrapositively, continuous reversibility entails that no basis is selected.
\end{proposition}

\begin{proof}
A privileged finite set is by definition carried to itself by every structure-preserving reversible
map, whose induced action on the set is a permutation of finitely many points. A continuous path of
such maps induces a continuous path in a finite discrete set, which is constant. Hence a continuous
reversible family moving one distinguished pure state to another respects no such privilege, and if
every pair of pure states is so connected, no finite set is privileged.
\end{proof}

The converse is where the historical difficulty lies: the absence of a privileged finite set does
not by itself force continuity, for the transformation group could in principle be infinite and
discrete. Two structural properties close every such escape at once.

\begin{proposition}[Divisible and closed implies connected]\label{prop:divisible-closed}
Let $G$ be the group of reversible transformations of the pre-measurement regime, acting on the
state body. If $G$ is divisible --- every element has an $n$-th root for every $n$ --- and
topologically closed, then $G$ is a compact connected Lie group; and if $G$ moreover acts
transitively on the pure states, then a continuous reversible transformation connects any two of
them: Hardy's fifth axiom.
\end{proposition}

\begin{proof}
The state body is compact and convex with nonempty interior, so its symmetry group is compact, and
a closed subgroup of a compact matrix group is a compact Lie group with finitely many connected
components \autocite{broeckertomdieck1985compact}. The component group $G/G_{0}$ is a quotient of a
divisible group and hence divisible; but a nontrivial finite group is never divisible --- with
$\lvert H\rvert = m$ every element satisfies $h^{m} = e$, so no element other than the identity
possesses an $m$-th root --- and $G/G_{0}$ is finite, so it is trivial and $G = G_{0}$ is connected.
A connected Lie group is path-connected, and transitivity then supplies the continuous reversible
path between any two pure states.
\end{proof}

The same one-line fact of Lagrange closes the remaining escapes. Finite groups are excluded
directly, being nontrivial and never divisible. Countable groups are excluded because a
countable closed subgroup of a compact group is compact, has an isolated point by the Baire
category theorem, and is therefore discrete and finite. Uncountable disconnected groups are
excluded by the component-group argument of the proposition. Every route to a discrete
transformation group therefore passes through a failure of divisibility or a failure of
closedness. Divisibility, however, asks for more than the conclusion needs, and the next
proposition removes it.

\begin{proposition}[Continuity from a connected state space]\label{prop:connected-transitive}
Let the pure states of a finite-dimensional system form a connected space $P$, and let the group
$G$ of reversible transformations of the pre-measurement regime be topologically closed and act
transitively on $P$. Then the identity component $G_0$ already acts transitively on $P$, and a
continuous reversible transformation connects any two pure states: Hardy's fifth axiom holds,
with no divisibility assumed.
\end{proposition}

\begin{proof}
As in Proposition~\ref{prop:divisible-closed}, closedness makes $G$ a compact Lie group with
finitely many connected components \autocite{broeckertomdieck1985compact}, so $G_0$ is a closed
normal subgroup of finite index and is path-connected. Each $G_0$-orbit in $P$ is the continuous
image of a compact group, hence compact, hence closed. Since $G$ is transitive and $G_0$ is
normal, the $G_0$-orbit of $gx$ is $g$ applied to the $G_0$-orbit of $x$, so the cosets of $G_0$
permute the orbits transitively and there are finitely many of them. A finite partition of $P$
into closed sets is a partition into open sets; $P$ is connected; so there is one orbit and
$G_0$ is transitive. Given pure states $x$ and $y$, choose $g \in G_0$ with $gx = y$ and a
continuous path $g_t$ in $G_0$ from the identity to $g$; then $g_t x$ is a continuous reversible
transformation from $x$ to $y$.
\end{proof}

The proposition relocates the premise. Divisibility was asked for because a transformation group
could be infinite and discrete; but a group acting transitively on a continuum cannot be
countable, and a closed group acting transitively on a connected space cannot be disconnected in
the way that matters, so what the argument needs from the regime is not the availability of
$n$-th roots but the connectedness of the state space itself. That is a statement about values:
that the pure states are all the unit vectors of $\mathbb{C}^{K}$ up to phase, the space
$\mathbb{CP}^{K-1}$, connected for $K \ge 2$, and hence that the amplitudes of
\S\ref{sec:self-reference} fill their continuum. The two standing foils show the premise doing
its work. The stabiliser fragment of quantum theory has complex amplitudes and is contextual,
yet it has finitely many pure states, its outcome probabilities take only the values $0$,
$\tfrac12$ and $1$, and its reversible transformations form a finite group, the Clifford group,
so Hardy's fifth axiom fails in it; Spekkens's toy theory \autocite{spekkens2007toy} is the same
with real amplitudes, its ignorance rationed by a knowledge-balance principle. Neither has a
connected state space, and neither is a model of the regime described here, whose credences take
every value. Divisibility, which both also lack --- the phase gate has no square root in the Clifford group
--- is a symptom of the discreteness, not its cause. What both lack can be said exactly, and the
next subsection says it: each contains a transformation that fixes no state it admits.

\subsection{The premises, and what they mean}

Three premises carry Proposition~\ref{prop:connected-transitive}. Transitivity is the regime's
homogeneity: inside/outside equivalence leaves no pure state privileged, and a proper closed
invariant set of pure states would define a superselection observable, a classical record
predating any \textsc{fanout} (\S\ref{sec:fanout-boundary}). Connectedness of the state space is two statements, and only one of them is quantum. The first
is that the pure states are infinite and dense, with a fixed point for every transformation; the
second is that the state space is complete. Quantum theory over the algebraic numbers ---
amplitudes in the algebraic closure of $\mathbb{Q}$ within $\mathbb{C}$ --- has the first
without the second: it has $K = N^{2}$, the Born rule, superposition, entanglement and Bell
violation at the Tsirelson value, every unitary has an algebraic eigenvector, and it is
countable and totally disconnected. So the whole quantum content of Hardy's axiom lies in the
density, and the completion is the real-number convention of every probability theory. Hardy
adopts it before his fifth axiom is stated, in the first, which makes probabilities real numbers
and gives his classical theory a real simplex of mixtures; this framework adopts it at the
Shannon step, where credences are declared to range over a connected domain so that the
continuation of the logarithm is unique (\S\ref{sec:complex-amplitudes}). Given that convention,
density of the pure states is connectedness of their closure. Divisibility, which an earlier
version of this section asked for, is not needed: it follows from the conclusion, every element
of a compact connected Lie group having roots of every order.

Closedness is the completion, limit-realisability: a Cauchy sequence of available
transformations has its limit available. It cannot be derived, and the reason should be stated
rather than hidden. Any collection of first-order principles about states, transformations and
probabilities that has an infinite model has a countable one, and a countable subset of
$\mathbb{CP}^{K-1}$ is totally disconnected; the founding non-computability does not escape
this, since the arithmetical complex numbers form a countable algebraically closed field
containing non-computable numbers and closed under every limit it can itself describe. Hardy's
axiom and closedness are second-order, quantifying over paths and over sequences, and a
second-order principle is grounded as Dedekind completeness is grounded or not at all. The
grounds are these. The completion of a dense set under the finite-precision metric is unique and
depends on no choice, whereas any countable alternative is a choice of definability level, an
extra law with no ground and no test; and the framework's own result is that no observer's
descriptive level captures the regime, since by \S\ref{sec:self-reference} every observer's
state exceeds its own level and the inside/outside principle makes that a law rather than a
perspective. The burden therefore falls on any restriction, not on the completion, which is the
least committal option and carries no quantum content. What limit-realisability adds to the line
is still worth stating exactly, because it is the corrected form of a natural intuition.
Cantor's diagonal \autocite{cantor1891diagonal} does not build the continuum; its output is an
infinite decimal --- a Cauchy sequence of rationals --- which exists \emph{as a point} only
because the line is complete. The diagonal presupposes limit-realisability and then shows that
the completed line outruns every internal listing. And the character of the completion is the
framework's own: the computable reals are countable \autocite{Turing_1937}, hence of measure
zero, so once the transformation group contains its limits almost all of them are
non-computable. Non-computability is a property of the completed continuum, not the reason for
it. The constructive reading, on which only describable limits exist, remains consistent; it is
the countable theory above, and the two readings differ in nothing an experiment can certify.

Density is where the section's title is earned, and it is derived from the paper's own premise.
Write $(\mathrm{I})$ for the statement that every admissible transformation of the regime is a
phase of some admissible classical structure: diagonal, that is, in the record basis of some
observer the regime contains, a phased spider of that basis in the calculus of
\S\ref{sec:hopf-justification} \autocite{coeckeduncan2011interacting} (VERIFY-SUPPORT). It says
of transformations what the global phase of \S\ref{sec:self-reference} says of states: there are
no absolute ones. It follows from three things the paper already holds. The inside/outside
involution is a genuine involution, so the outside presentation of the inside is a presentation
of the same object with nothing lost: the outside can have complete knowledge of the inside, a
potential record, the joint state lying in a structure the outside carries. Time is operation
count, internal to the system that operates (\S\ref{sec:space-of-computations}), so the outside
does not share the inside's count. And before any \textsc{fanout} that count is not classically
recorded, so the outside's description is pure. Let $g$ be the transformation the inside
performs per operation. The outside's complete knowledge cannot depend on the inside's count,
and it is pure, so it is a pure state of the inside's clock and system that is stationary under
``advance the clock and apply $g$''; the stationary states of that map are clock Fourier modes
tensored with eigenvectors of $g$ and nothing else \autocite{pagewootters1983evolution}
(VERIFY-SUPPORT). So the outside holds a structure in which $g$ acts by phases, which is
$(\mathrm{I})$; and the relative phase in that structure is the outside's ignorance of the
inside's computational count, the same ignorance the global phase records one level down. The
consequences are a theorem. A phase of a structure fixes the structure's copyable states, so
every admissible transformation fixes an admissible pure state; a finite transitive action has
an element moving every point (Jordan), so the pure states are infinite; the transformation
carrying a stabiliser state to the fixed point of the composite of a $Z$-phase and an $X$-phase
turns through $\arccos(1/\sqrt{3})$, an irrational fraction of a turn, so the group has an
element of infinite projective order and its orbits are dense. With the completion,
Proposition~\ref{prop:connected-transitive} fires and Hardy's fifth axiom holds. The two foils
fail $(\mathrm{I})$ exactly where they fail continuity: fourteen of the twenty-four Clifford
operations are phases of no Pauli basis and fix no stabiliser state, their fixed points being
the magic states; the toy theory's three-cycles are phases of no epistemic basis; and in each
the outside's knowledge of the inside under the order-three protocol is either count-dependent
or a mixture over the count, so each makes the inside's time classical before any record exists.
The derivation of $(\mathrm{I})$ is an outline, not a proof, and it is what the open problem
below asks for.

\subsection{Why the cut matters physically}

The regime's continuity is not decorative. Reversibility alone, over complex amplitudes, licenses
the full invertible group $GL(n, \mathbb{C})$, and Aaronson showed what that buys: superluminal
signalling, and computational power at the level of \textsf{PP} \autocite{Aaronson2004h}. The cut
from $GL(n, \mathbb{C})$ down to the unitary group is the conserved quadratic form, and continuity
is what entitles the theory to make it: with the fifth axiom in hand, Hardy's theorem delivers the
complex Hilbert space, the quadratic norm and the unitary group in one stroke \autocite{Hardy2001a}
--- cited here, not re-proved. One honest caveat travels with the citation. Hardy's reconstruction
uses his remaining axioms besides --- composition and subspace structure --- and these must be
derived rather than imported: the composition ledger of \S\ref{sec:division-algebras} and the
locality-of-tomography argument of \S\ref{sec:self-reference} carry that duty, and it is discharged
there, not silently assumed here.

\subsection{Super-tasks at the boundary}

The premise just priced has a fifty-year-old stress test in the philosophy of mathematics, and the
framework passes it in an unexpected way. A super-task --- infinitely many operations completed in
finite time \autocite{Thomson_1954} --- presupposes its terminal moment as an actual point of the
timeline: the accumulation point of a generic sequence of step-times exists only in the completed
line, and the defining question, the state \emph{after all the steps}, is meaningful only under
limit-realisability of time. Super-tasks are downstream consumers of exactly the completeness at
issue in this section. Thomson's lamp then exhibits both sides of the \textsc{fanout} boundary at
once. On the record side, the lamp's description fixes its state at every finite stage and thereby
fixes no limit at all: any terminal value is consistent with the whole task
\autocite{Thomson_1954,benacerraf1962tasks}. That is not a puzzle about lamps; it is the proof that
the classical record algebra is not limit-closed, and cannot be. On the amplitude side the limit
exists --- a superposed fixed point, whose complex-probability treatment belongs to
\S\ref{sec:self-reference} --- because that side \emph{is} limit-closed, which is the fifth face
itself. The boundary is thereby the edge of limit-realisability: complete on the reversible side,
incomplete on the record side, with the lamp standing on both sides at once as the boundary's oldest
exhibit. A final clause dissolves rather than answers the classic question. In this framework time
is operation count (\S\ref{sec:space-of-computations}), and ``after infinitely many operations''
indexes no operational moment: the lamp packs its endless \textsc{fanout}s into the emergent
continuum, and asking for its record at the accumulation point asks the discrete contents for a
value that only the completed container possesses. On or off is not undecided; it is unasked.

\subsection{What remains, priced}

\begin{openproblem}[Every transformation a phase of an inside structure]\label{op:limit-realisability}
Proposition~\ref{prop:connected-transitive} reduces Hardy's fifth axiom to three premises:
connectedness of the pure states, transitivity, and closedness. Transitivity is argued above
from the absence of records. Closedness is the completion, adopted on the grounds stated above
and carrying no quantum content. Connectedness is density plus that completion, and density
follows, by Jordan's theorem and the infinite-order argument, from $(\mathrm{I})$: that every
admissible transformation of the regime is a phase of some admissible classical structure. What
remains open is the rigour of the derivation of $(\mathrm{I})$ from the paper's premise. Derive,
inside the fixed-point locus $\mathrm{Fix}(\sigma)$ of Proposition~\ref{prop:inside-outside},
that every endomap is a phase of a classical structure in $\mathrm{Fix}(\sigma)$, from the
self-representing object, the involution $\sigma$ as a genuine involution, and the reading of
the inside presentation as a sequence indexed by operation count with no record before
\textsc{fanout}. The outline above does this with the stationary-state construction of Page and
Wootters; a derivation passes if it makes that outline a theorem of the calculus, with the
outside's count-independent complete description shown to exist and to be pure from the
involution alone. It fails before any detail is checked if every premise it uses holds in the
stabiliser fragment of quantum theory, in Spekkens's toy theory \autocite{spekkens2007toy}, in
the real qubit, or in quantum theory over the Gaussian rationals, all of which violate
$(\mathrm{I})$ outright and all of which make the inside's operation count classical before any
record exists. It also fails if it imports the completion as a premise about transformations:
the completion is the real line and has no derivation, and a derivation of $(\mathrm{I})$ must
hold in the countable algebraic theory, where $(\mathrm{I})$ is true and continuity is not.
\end{openproblem}

The accounting, in one sentence: the axiom that every reconstruction programme pays is, in this
framework, one derivation posed, of Lawvere's shape --- that every endomap of the fixed-point
locus is a phase of a classical structure in it --- and one convention named, the completion,
which classical probability pays too. And
the hand-on is already visible: the complex amplitudes that this section's continuity underwrites
through Hardy's theorem, and the complex amplitudes that \S\ref{sec:self-reference} derives natively,
are one and the same $\mathbb{C}$ --- the floor of the ladder that \S\ref{sec:division-algebras}
climbs.

%      % =====================================================================
%  SECTION: The Division-Algebra Ladder
%  Drafted 2026-07-03. One of the three foundations sections
%  (sec:self-reference, sec:division-algebras, sec:fanout-boundary);
%  this one resolves the dangling \S\ref{sec:division-algebras} in the
%  Hopf-justification section, whose header names it "the
%  Cayley--Dickson / Hurwitz step".
%
%  PLACEMENT in main.tex (REORDERED per JS 2026-07-06): position 6,
%  opening Act II ("why stop at C?") right before its first consumer --
%      \input{self_reference_section_2026-07-04}
%      \input{continuity_section_2026-07-04}
%      \input{L2_metric_global_phase_section_2026_06_07}
%      \input{wigners_friend_in_the_hopf_picture_2026-06-13}
%      \input{derivation_of_the_quantum_postulates}
%      \input{division_algebras_section_2026-07-03}   <-- moved here
%      \input{spacetime_from_non-computability}       <-- first consumer
%                                                          (S^7 exclusion)
%      \input{hopf_justification_section_2026-06-08}
%  The section is position-portable: all adjacency phrasing was
%  retargeted to explicit \S\ref on 2026-07-06. The S^7 hand-off
%  sentence to sec:space-of-computations WRITTEN 2026-07-06 (spacetime
%  file read; the closing now hands first to its immediate consumer,
%  then to hopf-justification). Upload verifications COMPLETE: L2 and
%  quantum-postulates do not consume the ladder; WF's references are
%  forward-voiced or self-carrying.
%
%  DIVISION OF LABOUR (deliberate, to avoid overlap):
%   - sec:self-reference (undrafted): self-reference -> complex
%     amplitudes and the Born rule; why not a real base. This section
%     TAKES that as given and starts at C.
%   - THIS section: the algebra -- the doubling construction, the
%     ledger of losses-as-physics, the internal-algebras-not-a-new-QM
%     scoping (Adler pre-empt), the two appearances of the ladder
%     (register / measurement climb, stated lightly), and Hurwitz as
%     the algebraic termination with the committed zero-divisor
%     reading.
%   - sec:hopf-justification: the geometry (Adams, parallelisability,
%     the fibrations), the join of the two classifications, the
%     entanglement-sensitivity literature, the C-direction honest debt,
%     and the epistemic reading of termination.
%
%  COMMITTED READINGS USED (sources re-read 2026-07-03):
%   - Loss ledger and level dictionary: cayley_dickson_qubits_for_review
%     sec.5 ("each step loses a symmetry, and each lost symmetry is
%     identified with a physical content").
%   - Termination reading: grounding_argument_gauge_group ("the
%     sedenions have zero divisors and cannot support consistent
%     information flow around loops") -- leads; norm-composition-as-
%     probability given as its amplitude-level restatement.
%   - Internal-not-base scoping: cayley_dickson_qubits_for_review Q2 and
%     sec.3.2 caveat (Adler objection; C-qubit treated geometrically).
%
%  HONEST DEBT: ladder proved closed at the top (Hurwitz), derived at
%  the floor. The rung-compulsion revisit-trigger FIRED 2026-07-06:
%  sec:self-reference and op:sigma-functor now exist, so the thesis is
%  minted as op:ladder-compulsion (with the CD-conjugation-as-sigma
%  candidate route), inheriting op:sigma-functor's status.
%
%  AUDITED 2026-07-06 (final pass, third stop): all four citations
%  used; all cross-refs live; CD facts, dims, alternativity, Hurwitz
%  statement verified; Adler subsection delivers everything the
%  register's compressed pointer promises; file pure ASCII; no bare
%  "frame". Fixes: opening rewritten for both predecessors (continuity
%  was inserted between after drafting); "associator's non-vanishing is
%  mass itself" recalibrated to the register's debt-against-selection
%  defusal; the O->S ledger row and Hurwitz subsection reconciled with
%  the fourth-qubit forward remark ("the terminus closes the ladder of
%  forces; what survives past it is, at most, flavour").
%
%  CROSS-REFS: sec:self-reference [sibling dangle -- resolves when that
%  section is drafted], sec:hopf-justification, sec:hopf-information-bound,
%  sec:register, sec:entanglement-classes, sec:generations, sec:masses,
%  sec:wigners-friend. All others exist.
%  CITATIONS: existing keys baez2002octonions, hurwitz1898composition,
%  furey2018three. ONE new key: adler1995quaternionic (RIS in
%  division_algebras_refs_2026-07-03.ris).
%  British English; \autocite; no bullets in prose; the ledger is a
%  table because the comparison is the argument.
% =====================================================================

\section{The Division-Algebra Ladder}
\label{sec:division-algebras}

The foundations behind us end at amplitudes twice over: natively, where the self-referential
argument forced the description of a system into the system's own state space and priced
consistency at a number system whose parts can cancel --- complex amplitudes, and the Born rule
that squares them (\S\ref{sec:self-reference}) --- and independently, where the regime's
continuity delivers the same field through Hardy's theorem (\S\ref{sec:continuity}), the
convergence of the two branches being evidence rather than accident. This section asks the question
both routes leave open. Why
should the number system stop at $\mathbb{C}$? Ordinary quantum mechanics stops there by postulate;
a framework that has forbidden itself postulates must answer, and the answer is that the climb
neither stops at $\mathbb{C}$ nor continues forever. There are exactly two more rungs, each
purchased by surrendering an algebraic law and each paying for itself in physics, and the ladder is
closed at the top by a theorem of 1898. This section sets out the algebra of that ladder: the single
construction that climbs it, the ledger of what each rung costs and buys, and the theorem that ends
it. The geometry of the same ladder --- that each rung acts on a sphere, that exactly those spheres
are parallelisable, and that the algebraic and geometric classifications are one --- is the business
of the geometry section (\S\ref{sec:hopf-justification}).

\subsection{One construction climbs the whole ladder}

The rungs are not a miscellany; a single construction generates them all. Given an algebra with a
conjugation, the Cayley--Dickson construction builds its double on ordered pairs, with a twisted
multiplication in which the conjugate enters and a conjugation $(a,b)^{*} = (a^{*}, -b)$
\autocite{baez2002octonions}. (Conventions differ in the placement of the conjugations; all yield
isomorphic algebras.) Applied repeatedly it produces the sequence
$\mathbb{R} \to \mathbb{C} \to \mathbb{H} \to \mathbb{O} \to \mathbb{S}$, of real dimensions
$1, 2, 4, 8, 16$: the reals, the complex numbers, the quaternions, the octonions, the sedenions.
The ground floor is classical: $\mathbb{R}$ is ordered and positive-definite, the arithmetic of
classical probability. The first doubling is therefore the quantum move itself --- the passage from
probabilities to amplitudes that can interfere --- and it is not made here by choice: it is the
derivation of \S\ref{sec:self-reference}, already behind us. What this section must account for is
the rest of the climb.

\subsection{The ledger}

In mathematics the doublings are usually told as a story of degeneration: each algebra is worse
behaved than the last, and the sequence is studied for where it breaks. This framework reads the
same ledger in the opposite direction. Each rung is bought by surrendering a law, and each
surrendered law is exactly a piece of physics --- the losses are the assets, and the comparison is
sharp enough to be the argument:

\begin{center}
\begin{tabular}{llll}
\hline
doubling & dimension & what is lost & what the loss becomes \\
\hline
$\mathbb{R} \to \mathbb{C}$ & $1 \to 2$ & the order of the reals & phase and interference; amplitudes \\
 & & & that cancel; the Born rule (\S\ref{sec:self-reference}) \\
$\mathbb{C} \to \mathbb{H}$ & $2 \to 4$ & commutativity & non-commuting actions: weak isospin \\
 & & & $SU(2)$, the $B$ register (\S\ref{sec:register}) \\
$\mathbb{H} \to \mathbb{O}$ & $4 \to 8$ & associativity & the associator: colour, chirality, the \\
 & & & generations, and mass as debt \\
 & & & (\S\ref{sec:generations}, \S\ref{sec:masses}) \\
$\mathbb{O} \to \mathbb{S}$ & $8 \to 16$ & division itself & no further force: zero divisors appear \\
 & & (zero divisors) & and the ladder of sectors ends (below) \\
\hline
\end{tabular}
\end{center}

Walking the rows: the order lost at the first doubling is the possibility of interference --- parts
of a probability that can oppose one another --- which is where the derivation of \S\ref{sec:self-reference}
lives. The commutativity lost at the second is the existence of actions whose order matters; the
unit quaternions form the group $SU(2)$, and the framework houses weak isospin there, as the
internal algebra of the $B$ register whose order-sensitive, non-commuting dynamics the interaction sections develop
(\S\ref{sec:register}, \S\ref{sec:entanglement-classes}). The associativity lost at the third ---
the octonions retain the weaker property of alternativity \autocite{baez2002octonions} --- creates
the associator $[a,b,c] = (ab)c - a(bc)$, and the associator is where most of this paper's physics
is transacted: the stabiliser of a preferred complex direction inside the octonions' automorphisms
is the colour group (\S\ref{sec:hopf-justification}); the triality structure and its $K_3$
obstruction generate the three generations and the sterile channel (\S\ref{sec:generations}); and
the associator is the currency of mass --- the debt against the selected direction that
\S\ref{sec:masses} prices, owed by whatever action moves the selection and not by residence at
the rung, which is the defusal \S\ref{sec:register} makes explicit. The algebraic facts
in the third column are theorems; the identifications in the fourth are established in the sections
cited, each carrying its own epistemic flag there, and are not re-litigated here. What this table
asserts on its own account is only the correspondence's shape: one law surrendered per rung, one
sector of physics gained per law.

\subsection{Internal algebras, not a new quantum mechanics}

A reader versed in the quaternionic-quantum-mechanics literature will at this point raise a fair
objection: quantum theory over $\mathbb{H}$ as a foundation is known to be troubled --- the
tensor-product and composite-system difficulties set out by Adler and others
\autocite{adler1995quaternionic} --- and quantum theory over $\mathbb{O}$ founders earlier still,
since a non-associative algebra does not furnish a Hilbert space at all. The framework's position
must therefore be stated exactly, because it does not make the claim the objection targets. The base
of the physical Hilbert space here is $\mathbb{C}$ and remains $\mathbb{C}$; the higher algebras
enter as the \emph{internal} algebras of the register --- $\mathbb{H}$ for the $B$ qubit,
$\mathbb{O}$ for the $C$ qubit (\S\ref{sec:register}) --- in a tensor structure whose outer theory is
ordinary complex quantum mechanics. This is the same move by which the division algebras enter the
constructions of Furey and the related programmes without any proposal to rebase quantum mechanics
itself \autocite{furey2018three}. For the octonionic register the point is sharper still:
$\mathbb{O}^{2}$ is not a Hilbert space, and the framework accordingly treats the $C$ qubit not in
the Dirac--von Neumann manner but geometrically, as the sphere $S^{15}$ with its Hopf structure ---
which is precisely the form in which \S\ref{sec:hopf-justification} needs it. Nothing in this paper therefore
stands or falls with quaternionic or octonionic quantum mechanics as foundations; the ladder is
climbed by the registers' internal algebras, not by the theory's base field.

\subsection{The same ladder, twice}

The ladder appears in this framework in two guises, and it will help the reader to have them
distinguished now. Statically, it is the anatomy of a single relation: the three qubits of the
register take their values one rung apart --- $A$ over $\mathbb{C}$, $B$ over $\mathbb{H}$, $C$ over
$\mathbb{O}$ --- so that every particle carries the whole ladder as its internal structure
(\S\ref{sec:register}). Dynamically, it is the shape of measurement: a single qubit's pure states
fill $S^{3}$, an entangling measurement carries the joint state into $S^{7}$, and a third party into
$S^{15}$ --- the correspondence between qubit entanglement and the Hopf fibrations that
\S\ref{sec:hopf-justification} takes from a mature literature --- so that each act of measurement ascends one rung --- the climb the reader has already made,
narrated in \S\ref{sec:wigners-friend} --- and every completed relation is a terminated climb, which is
the reading \S\ref{sec:hopf-information-bound} gives the tower's endpoint. One ladder, seen twice:
as the anatomy of a relation and as the dynamics of relating. The framework regards this double
appearance not as a coincidence but as the same structure viewed statically and in motion, and the
sections cited carry that identification's weight.

\subsection{Hurwitz closes the ladder}

The construction itself would double forever; what stops it is a classification theorem. By
Hurwitz's theorem, the normed division algebras --- those in which the norm composes,
$\lvert ab\rvert = \lvert a\rvert\,\lvert b\rvert$, and no non-zero elements multiply to zero ---
are exactly $\mathbb{R}$, $\mathbb{C}$, $\mathbb{H}$ and $\mathbb{O}$
\autocite{hurwitz1898composition,baez2002octonions}. The next doubling, the sedenions, has zero
divisors: non-zero elements whose product vanishes. The framework's reading of this failure is the
one that matters physically. Zero divisors cannot support consistent information flow around loops
--- a cycle of compositions can silently annihilate what it carries --- and a framework in which
closed computational paths are load-bearing (\S\ref{sec:hopf-justification},
\S\ref{sec:space-of-computations}) cannot be built on an algebra that permits it. At the level of
amplitudes the same failure reads: the composition law
$\lvert ab\rvert = \lvert a\rvert\,\lvert b\rvert$, which carries probability multiplicatively
through every rung of the ladder, fails beyond the octonions, so that two non-vanishing amplitudes
could compose to nothing and probability would be destroyed by mere composition. The last algebra
fit to carry amplitudes is $\mathbb{O}$. One rung past the terminus is not entirely silent,
however: the sedenions' internal structure --- their automorphism group and their octonion
subalgebras --- is the recorded candidate arena for tagging the generations
(\S\ref{sec:interactions}), and it enters exactly as this section's scoping permits, internally
and never as a base. The terminus closes the ladder of forces; what survives past it is, at most,
flavour. This is the algebraic half of the termination; the
geometric half --- that no sphere above $S^{7}$ is parallelisable --- and the theorem joining the
two halves are the opening business of \S\ref{sec:hopf-justification}, and the epistemic reading of the same endpoint,
as a bound on the information a single relation can carry, is given in
\S\ref{sec:hopf-information-bound}.

\subsection{The honest debt}

As in the geometry section, we close by stating plainly what is proved, what is established, and what is
owed. Proved: the construction, the losses, and the termination are theorems --- the Cayley--Dickson
doubling and its ledger of surrendered laws are standard algebra, and Hurwitz's classification is
among the durable results of the subject. Established: the physical cashing of each rung --- the
Born rule at $\mathbb{C}$, weak isospin at $\mathbb{H}$, colour, generations and mass at
$\mathbb{O}$ --- is carried out in the sections cited above, each under its own epistemic flag. Owed,
and stated without disguise: the ladder as presented here is closed from above and derived at its
floor, but it is not yet derived rung by rung from below. Hurwitz forbids a further rung; the
argument of \S\ref{sec:self-reference} forces the first doubling, from the ordered arithmetic of
classical probability to amplitudes that interfere. That the same self-referential closure, applied
to a system whose amplitudes already take values at one rung, compels the next --- that relational
closure at one level forces precisely the Cayley--Dickson double of that level --- is the
framework's reading of the climb, and it is a thesis, not a theorem --- but the formal apparatus
against which it can be posed now exists, and we pose it.

\begin{openproblem}[The compulsion of the intermediate rungs]\label{op:ladder-compulsion}
Derive, within the categorical apparatus of Open Problem~\ref{op:sigma-functor}, that
self-referential closure at one rung compels the Cayley--Dickson double of that rung, so that the
intermediate steps $\mathbb{C} \to \mathbb{H} \to \mathbb{O}$ follow from the same principle
that forces the floor. The candidate route is the construction's own shape: the Cayley--Dickson
pair with its conjugation, $(a, b)^{*} = (a^{*}, -b)$, is an involution-equipped doubling, and the
route reads it as the algebraic shadow of adjoining the describer --- closure of level-$A$
descriptions under the inside/outside exchange requiring exactly the twisted pair. The problem
inherits the status of Open Problem~\ref{op:sigma-functor}, whose apparatus it presupposes.
\end{openproblem}

If the derivation lands, the whole ladder becomes a consequence of the same self-reference that
opened it. Until then, the honest statement is that the ladder's necessity is established at its
ceiling and at its floor, while the compulsion of the intermediate rungs is inherited from the
physics they are shown to carry rather than derived in advance. What this section hands on is, accordingly, an algebraic ladder with a floor, a ceiling and a
ledger. Its first consumer is immediate: \S\ref{sec:space-of-computations} spends the terminal
fact --- the non-associativity of $S^{7}$'s own multiplication --- as geometry, excluding the
seven-sphere as a ruler and fixing the dimension of space at three.
\S\ref{sec:hopf-justification} then supplies the geometry of the fibrations themselves, and the
join.

%      % EDITED 2026-09-09 scripts/patch_spekkens_citekey_2026-09-09T1810.py
% =====================================================================
%  SECTION: Why Hopf Fibrations Belong in Particle Physics
%  (sec:hopf-justification). Drafted 2026-06-08; F1 arc added
%  prop:f1-dictionary and op:hopf-frobenius's discharge clause.
%  AUDITED 2026-07-06 (final pass, eighth stop).
%
%  PLACEMENT (reordered chain, JS 2026-07-06): position 8 --
%      \input{spacetime_from_non-computability}
%      \input{hopf_justification_section_2026-06-08}
%      \input{entanglement_classes_and_gauge_groups_2026-06-08}
%  No adjacency phrasing in the prose: position-portable. All
%  cross-refs live -- sec:fanout-direction DRAFTED 2026-07-06 (formerly the one undrafted
%  section, by design).
%
%  AUDIT RECORD: prop:gsm, prop:termination and prop:f1-dictionary
%  verified against every quoting consumer (statements canonical);
%  honest-debt subsection names both the selection assumption and the
%  mereology debt; the epistemic two-climbs disambiguation and the
%  compositional rung-evidence verified; sec:canonical-orbits live;
%  file pure ASCII. Fixes: op:hopf-frobenius's cross-link with
%  op:sigma-functor made TWO-WAY (the "four ways" line now names the
%  op); both remaining functor conditions ANCHORED --- the quaternionic
%  to the established fibre physics of sec:wf-reconciliation (residue:
%  clutching map vs strong complementarity), the octonionic to the
%  proved canonical-orbit images (umbrella label app:ghz-w-proofs
%  minted at the appendix for the purpose). Statuses unchanged.
%  British English; \autocite; no bullets in prose.
% =====================================================================

\section{Why Hopf Fibrations Belong in Particle Physics}
\label{sec:hopf-justification}

Invoking Hopf fibrations in a discussion of the Standard Model invites an immediate and fair
objection: that the move is decorative, an analogy between two pieces of attractive mathematics
with no load-bearing connection between them. The purpose of this section is to dispose of that
objection before the construction is put to work. The connection is not an analogy. It is
carried by two independent and independently published theorems, and the structures it relies
upon are not selected for convenience but compelled by a pair of classification results. We take
as already established the correspondence between qubit entanglement and the Hopf fibrations
themselves --- the subject of a mature literature
\autocite{mosseri2001hopf,bernevigchen2003geometry} --- and do not reproduce it here. What this
section addresses is the prior question of why that correspondence is entitled to carry physical
content at all.

\subsection{Two theorems, joined end to end}

The bridge from qubit entanglement to the Standard Model is built from two distinct results,
neither of which is original to the present framework and neither of which depends on the other.

The first is geometric. The unit sphere in the Hilbert space of $n$ qubits is the total space of
a Hopf fibration: $S^3 \to S^2$ for one qubit, $S^7 \to S^4$ for two, and $S^{15} \to S^8$ for
three. For two and three qubits the fibration is \emph{entanglement-sensitive} --- separability
is precisely the condition under which the base degenerates to a product of single-qubit Bloch
spheres and the additional dimensions associated with the higher imaginary directions collapse
\autocite{mosseri2001hopf,bernevigchen2003geometry}. This is a statement about the geometry of
state spaces and requires no physical interpretation whatever.

The second is algebraic. Singling out a preferred complex structure on the octonions ---
equivalently, fixing the split $\mathbb{O} \cong \mathbb{C} \oplus \mathbb{C}^3$ --- leaves a
residual symmetry, and that residual symmetry is exactly the gauge group of the Standard Model.
This was established through the exceptional Jordan algebra by Dubois-Violette and Todorov
\autocite{duboisviolette2016exceptional,duboisviolettetodorov2019exceptional,todorovduboisviolette2018deducing}
and given its sharpest spinorial form by Krasnov \autocite{krasnov2021so9}, who showed that the
relevant group is the stabiliser of a complex structure inside $\mathrm{Spin}(9)$. We record the
result in the form in which it enters the present work.

\begin{proposition}[Dubois-Violette--Todorov; Krasnov]\label{prop:gsm} Let $J_R$ be the complex
structure on the space $\mathbb{O}^2$ of $\mathrm{Spin}(9)$ spinors determined by a choice of
unit imaginary octonion, equivalently the split
$\mathbb{O} \cong \mathbb{C} \oplus \mathbb{C}^3$. The subgroup of $\mathrm{Spin}(9)$ commuting
with $J_R$ is \[ G_{\mathrm{SM}} \;=\; \frac{SU(3) \times SU(2) \times U(1)}{\mathbb{Z}_6}, \]
acting on the singled-out $\mathbb{C}$ summand as a colour singlet of hypercharge $-1$ and on
the $\mathbb{C}^3$ summand as a colour triplet of hypercharge $\tfrac{1}{3}$, both as $SU(2)$
doublets --- the quantum numbers of one left-handed generation.
\end{proposition}

The two results are independent, and the join between them is the only place a physical theory
must do work. Each result is a theorem; what neither supplies is a reason why nature should fix
one complex direction within $\mathbb{O}$ rather than treating all directions symmetrically. We
return to this --- the genuine load-bearing assumption --- in \S\ref{sec:hopf-honest-debt}, and
resolve it, in candidate form, in \S\ref{sec:fanout-direction}.

\subsection{The ladder is forced, not chosen}

It might still be objected that selecting the octonions, or the Hopf fibrations, is itself a
modelling choice carrying as much latitude as the choice of a grand-unified group. It is not,
and the reason is a pair of classification theorems that close off the alternatives entirely.

By Hurwitz's theorem the only normed division algebras are $\mathbb{R}$, $\mathbb{C}$,
$\mathbb{H}$ and $\mathbb{O}$ \autocite{hurwitz1898composition,baez2002octonions}. By Adams's
resolution of the Hopf-invariant-one problem --- together with the parallelisability results of
Bott--Milnor and Kervaire --- the only parallelisable spheres are $S^0$, $S^1$, $S^3$ and $S^7$,
so that there exist exactly three nontrivial Hopf fibrations, in bijection with $\mathbb{C}$,
$\mathbb{H}$ and $\mathbb{O}$
\autocite{adams1960hopf,bottmilnor1958parallelizability,kervaire1958nonparallelizability}. These
are not heuristics; they are among the more durable theorems in topology and algebra. Their
force here is that, once a framework is committed to amplitudes valued in a normed division
algebra --- a commitment that in the present construction is made far upstream, by the
self-referential argument of \S\ref{sec:self-reference} and \S\ref{sec:division-algebras}, and
not by fiat at this stage --- the ladder of Hopf fibrations is not selected but compelled, and
it terminates of its own accord.

\begin{proposition}[Termination]\label{prop:termination} The Cayley--Dickson tower of normed
division algebras and the tower of parallelisable spheres both terminate at real dimension
eight. Consequently there are exactly three nontrivial Hopf fibrations, the largest being the
octonionic $S^{15} \to S^8$, and the construction admits no extension to a fourth case.
\end{proposition}

The termination is worth dwelling on, because it is unusual for a model of this kind to predict
its own boundary. The next rung of the Cayley--Dickson sequence, the sedenions, possesses zero
divisors; the Hopf map can no longer be globally defined; there is no fourth fibration. This is
not a limitation imposed from outside but a structural fact, and it carries a physical
consequence that we develop later: there is no room in the construction for a fourth gauge force
or a fifth interaction, so that the discovery of either would count as evidence against the
framework rather than as an invitation to extend it. The endpoint is a theorem, not a
stipulation.

\subsection{Independent arrival as corroboration}

The most persuasive evidence that the destination is not an artefact of the route is that it is
reached by routes sharing none of the same machinery. Szangolies
\autocite{szangolies2025standardmodel} reaches $G_{\mathrm{SM}}$ from the emergent-gravity and
Kaluza--Klein programmes, reading the singled-out complex direction as a dimensional reduction
of a $9+1$-dimensional spacetime; the present framework reaches it from self-reference
(\S\ref{sec:self-reference}); and Furey and Gillard--Gresnigt reach the same gauge group, and
the same three-generation structure, through the minimal left ideals of the Dixon and Clifford
algebras, without invoking Hopf fibrations at all
\autocite{furey2018three,gillardgresnigt2019three}; the relation between their minimal ideals
and the entanglement classes used here is taken up in \S\ref{sec:canonical-orbits}. The
octonionic route to particle physics is in any case an old one, going back to the use of
$\mathbb{O}$ to model the strong force \autocite{gunaydingursey1974quark}.

That programmes with disjoint starting points converge on
$SU(3) \times SU(2) \times U(1)/\mathbb{Z}_6$ and on the same quark--lepton division is the
coincidence Szangolies himself remarks upon: the characteristic structure of the Standard Model,
routinely regarded as arbitrary, turns up independently in unrelated settings, and its appearing
twice would then be doubly odd unless the two settings are in fact related. We read the
convergence the same way --- as evidence that the division-algebraic description is singled out
by the mathematics rather than imposed by the modeller, and therefore that
Proposition~\ref{prop:gsm} is structural rather than contrived.

\subsection{The honest debt}
\label{sec:hopf-honest-debt}

It would be misleading to leave the impression that this route is free of assumptions, so we
state plainly the one assumption that all of these programmes share. Proposition~\ref{prop:gsm}
requires that \emph{some} complex direction in $\mathbb{O}$ be singled out, and the uniqueness
of that choice is not supplied by the theorem. Krasnov parametrises the complex structure $J_R$
by an arbitrary unit imaginary octonion, and the residual group is the same for every such
choice because $G_2 = \mathrm{Aut}(\mathbb{O})$ acts transitively on the unit imaginary
octonions. The theorem therefore guarantees the gauge group \emph{given} that a direction is
fixed, while remaining silent on \emph{why} one is fixed. Szangolies treats this as a free
choice and flags it, in his own conclusions, as an open element of the construction --- one to
which several inequivalent answers are available.

This, and not the appeal to Hopf fibrations, is the genuine load-bearing assumption of the whole
approach. The Hopf fibrations are forced (Proposition~\ref{prop:termination}); the singling-out
is not. The present framework offers a candidate resolution that the prior literature lacks: the
act of measurement, modelled as \textsc{fanout}, selects the direction physically, and the
transitivity of $G_2$ then guarantees that the gauge group is independent of which direction the
measurement happens to select. This argument is given in \S\ref{sec:fanout-direction}; we flag
here, at the appropriate epistemic status, that it is structural rather than fully proved, and
that the residual question --- why a measurement selects \emph{exactly one} imaginary direction
rather than merely \emph{permitting} a selection --- remains open. It is now priced as Open
Problem~\ref{op:fanout-selection}.

A second debt is shared with every programme of this character and is not closed by anything
said here: why the abstract quantum state should admit a tensor-product decomposition into
qubits in the first place. This is the quantum factorisation, or mereology, problem. The
entanglement-sensitivity on which the construction turns is meaningful only once such a
decomposition is fixed, and we treat the decomposition as a necessary input rather than a
derived result. Naming it is preferable to letting it pass unremarked.

\subsection{The termination read epistemically}
\label{sec:hopf-information-bound}

The termination of the tower (Proposition~\ref{prop:termination}) has so far been stated
algebraically: the Cayley--Dickson sequence acquires zero divisors at the sedenions and the Hopf
map can no longer be globally defined. There is a second reading of the same boundary, epistemic
rather than algebraic, which we offer as a motivated conjecture rather than a settled result,
because it connects the tower's endpoint to a bound that arises independently in the foundations
of measurement.

A particle, in this framework, is a reified relation between events, and its internal content is
carried by the octonionic $A,B,C$ register --- the $S^{7}$ sector serving as internal degrees of
freedom in the channel that connects the events (\S\ref{sec:space-of-computations}). The
termination of the tower is then a statement about that content: a single such relation can
support genuine entanglement among at most three parties, because there is no fourth fibration
to carry a fourth independent party within the same irreducible relation. Every particle is, on
this reading, a \emph{completed} octonionic-level relation --- topped out at three parties not
by circumstance but because the tower admits no further rung. The information that can be
encoded in one relation-between-events, and hence extracted by one measurement of it, is
correspondingly bounded: it is the information of a completed three-party register, and no
measurement of a single particle can yield more.

This is exactly the kind of bound that the epistemic-restriction programmes reach from the other
direction. Spekkens' toy theory is built on a knowledge-balance principle that caps the
information obtainable about a system \autocite{spekkens2007toy}, and Szangolies derives such an
epistemic horizon not by positing it but from the same self-referential non-computability on
which the present tower rests \autocite{szangolies_j_2018}. That the two accounts --- the
algebraic termination of the Hopf tower and the epistemic bound on measurement information ---
should coincide is, we suggest, not a numerical accident but a shared root: both are
consequences of the self-referential non-computability of \S\ref{sec:self-reference}, wearing an
algebraic face in the one case and an epistemic face in the other. The Hurwitz--Adams ceiling
and the knowledge-balance ceiling would then be one limit seen twice, and Szangolies' programme,
which reaches $G_{\mathrm{SM}}$ by an independent route (\S\ref{sec:hopf-honest-debt}) and the
epistemic horizon from the same non-computable source, shares with this framework not merely a
conclusion but its origin.

The reading must be kept distinct from a different climb, with which it would otherwise be
conflated. The Wigner's-friend regress of \S\ref{sec:wigners-friend} ascends without end: each
observer may be measured by a further observer, each carrying a fresh gauge-private circle, and
that chain has no top --- it is the bump in the carpet, unbounded in length. The termination
bounds something else: not the length of the chain of observers, but the internal structure of
each single relation, each particle, that the chain is built from. The regress is unbounded
precisely because each of its links is a bounded, completed three-party object;
termination-of-each-link and non-termination-of-the-chain are not in tension but complementary,
the former individuating the links of the latter. What the bound caps is the information in one
particle; what the regress leaves uncapped is the number of particles and observers that can be
chained. The precise identification of ``the information of a completed three-party register''
with the formal measures of the epistemic-restriction literature is the technical step that
would turn this conjecture into a theorem, and we flag it as owed; the structural claim --- that
the tower's endpoint and the measurement bound share the non-computable root --- is what we
advance here.

\subsection{The termination read compositionally}
\label{sec:hopf-compositional}

There is a third termination theorem, proved in a community that has so far had no contact with
the two joined above. Coecke and Kissinger showed that the GHZ and W classes --- realised as the
commutative special and the anti-special Frobenius algebras of categorical quantum mechanics ---
are \emph{compositionally complete}: every $N$-qubit entanglement class under SLOCC is generated
by composing these tripartite primitives with bipartite entanglement and product states
\autocite{coeckekissinger2010compositional}. Above three parties nothing new appears; the
entanglement alphabet closes at exactly the point where the fibration tower terminates. The
mutual silence of the two literatures is itself striking: Coecke and Kissinger make no reference
to Hopf fibrations or division algebras, and the fibration literature on entanglement
\autocite{mosseri2001hopf,bernevigchen2003geometry} nowhere cites compositional completeness.
Two classifications, in different mathematics, draw the same line in the same place, unaware of
one another. That silence has, in the months this paper was in draft, just begun to break from
the categorical side, with a projective, base-level reading of the same calculus
\autocite{burtonanwar2026meromorphic}; we engage it below.

The framework's conjecture is that this is one termination wearing a third face: the alphabet
has the letters it has --- Bell, GHZ, W --- because the parallelisable spheres have the
dimensions they have. There is partial evidence at every rung. At one qubit, the commutative
special Frobenius algebras on $\mathbb{C}^{2}$ are classified by orthonormal bases --- points of
$S^{2}$, the base of the first fibration --- so that the copied basis plays the role of a
section of the bundle and the spider theorem expresses its idempotence. At two qubits, the
$SU(2)$ fibre of $S^{7} \to S^{4}$ carries exactly the entanglement that survives local
operations, and L\'{e}vay's result that the bundle's natural connection is the $SU(2)$
Yang--Mills instanton \autocite{levay2004geometry} stands on the geometric side where the
interaction of two Frobenius structures stands on the categorical one. At three qubits, the
split of D\"{u}r, Vidal and Cirac \autocite{durvidalcirac2000three} mirrors the algebra exactly:
GHZ is genuinely tripartite, demanding the full non-associative octonion, while W is
pairwise-robust and housed in the associative subalgebra $\mathbb{H} \subset \mathbb{O}$ --- the
special/anti-special Frobenius distinction in geometric dress.

The first of the four conditions the correspondence requires can be discharged at once, and we
do so here, because a problem whose first step has been walked is a different kind of invitation
from one merely posed.

\begin{proposition}[The one-qubit dictionary]\label{prop:f1-dictionary} The $\dagger$-special
commutative Frobenius algebras on $\mathbb{C}^{2}$ correspond bijectively to
\emph{lifted antipodal skeletons} of the complex Hopf fibration $\pi \colon S^{3} \to S^{2}$:
unordered antipodal pairs $P \subset S^{2}$ together with a section of $\pi$ over $P$. The
copyable states of a structure are exactly the section's image, so the lifts are recorded
faithfully --- the projectively trivial centre of $U(2)$ acts non-trivially on structures ---
and the correspondence is $U(2)$-equivariant, the ordered structures forming a single free and
transitive $U(2)$-orbit. The channel a structure induces on density matrices is dephasing in its
skeleton and is independent of the section: the fibre freedom, a torsor for
$\mathrm{Map}(P, U(1))$, is exactly the data absent from the classical record a structure
induces.
\end{proposition}

\begin{proof}
By the classification of Coecke, Pavlovic and Vicary, such structures correspond to orthonormal
bases, recovered as the sets of copyable states \autocite{coeckepavlovicvicary2013orthogonal}.
An orthonormal basis of $\mathbb{C}^{2}$ is two orthogonal rays --- an antipodal pair on the
Bloch sphere --- together with a unit representative of each, which is precisely a section of
$\pi$ over the pair; and $\delta(e^{i\theta}x) = e^{i\theta}\, x \otimes x$ against
$(e^{i\theta}x) \otimes (e^{i\theta}x) = e^{2i\theta}\, x \otimes x$ pins the phases.
Equivariance, the action of the centre, the orbit structure and the induced channel are direct
computations; an executable verification of every finite check accompanies the code release.
\end{proof}

Two remarks fix the proposition's scope. It corrects a natural first reading: a classical
structure does not project to the base but selects lifts into the total space --- classicality
is section data, and what the copy operation broadcasts is a lift, not a ray. And it claims
nothing topological about cloning: the finiteness of the copyable set is linear, structural
no-cloning residing in non-naturality \autocite{abramskyheunen2012hstar}, so the first place the
bundle's global topology can bite is the quaternionic rung, where the degree-one clutching map
stands ready to be compared with the strong-complementarity law governing an interacting pair of
structures. The dictionary is, finally, the complement of the projective treatment cited above:
Burton and Anwar pass functorially to the base and deliberately discard the fibre
\autocite{burtonanwar2026meromorphic}, and the classical structures turn out to live in exactly
the data their projectivisation forgets --- the two readings compose, and the resulting triangle
is a consistency condition available to constrain the remaining conditions of the
correspondence.

\begin{openproblem}[The Hopf-fibration--Frobenius correspondence]\label{op:hopf-frobenius} Let
$\mathbf{Hopf}$ denote the category of the three fibrations
$S^{2^{k}-1} \hookrightarrow S^{2^{k+1}-1} \to S^{2^{k}}$, $k = 1, 2, 3$, with bundle maps as
morphisms, and let $\mathbf{Frob}$ denote the Frobenius algebras on $\mathbb{C}^{2}$ in
$\mathbf{FdHilb}$, with the commutative-special (GHZ) and anti-special (W) classes
distinguished. Construct a functor $F \colon \mathbf{Hopf} \to \mathbf{Frob}$ --- or prove that
none exists --- under which the complex fibration yields the single-qubit classical structures
--- Proposition~\ref{prop:f1-dictionary} identifies these as the lifted antipodal skeletons,
section data over two base points rather than the points alone; the quaternionic fibration
yields the interaction of two such structures, with the $SU(2)$ fibre carrying the entanglement
that survives the classical quotient; the octonionic fibration yields the GHZ/W pair, GHZ on the
full non-associative fibre and W on the associative $\mathbb{H} \subset \mathbb{O}$; and the
composition by which Coecke--Kissinger generate the $N$-qubit classes corresponds to a
well-defined geometric operation on the bundles. If such an $F$ exists, compositional
completeness follows from Adams's theorem as a topological corollary --- there is no fourth
fibration to supply a fourth primitive --- and if $F$ is faithful enough that the implication
also runs back, the termination of the tower and the completeness of the alphabet are equivalent
expressions of a single classification. The first condition is discharged by
Proposition~\ref{prop:f1-dictionary}; what remains open is the quaternionic and octonionic
conditions and, above all, the correspondence of compositions. Both open conditions are already
constrained: the quaternionic fibre's physical content --- the recoverable relative phase and
its purity ledger --- is established (\S\ref{sec:wf-reconciliation}), leaving the
clutching-map-to-strong-complementarity comparison as the structural residue; and any candidate
functor must reproduce the proved canonical-orbit images, GHZ to the colour triplet and W to the
singlet (\S\ref{sec:canonical-orbits}, Appendix~\ref{app:ghz-w-proofs}). A proof that no such
functor exists, or a correspondence at the level of the classifications without functoriality on
compositions, would each also resolve the problem: the first would make the agreement of two
independent classifications a fact demanding its own explanation; the second would establish the
weaker reading.
\end{openproblem}

A resolution in the affirmative would do three things. It would explain compositional
completeness rather than merely prove it --- a corollary of the division-algebra structure of
quantum mechanics rather than an independent algebraic fact. It would give the diagrammatic
language of this paper's own interaction circuits a topological warrant: the GHZ/W spider
calculus would be the unique language the topology permits, not a convenient notation. And it
would open a two-way bridge, translating problems on compositions of Frobenius structures into
operations on bundles and back. The problem sits at the intersection of categorical quantum
mechanics, the quantum-information geometry of the fibrations, and the division-algebra
approaches to particle physics, and it is posed here so that any of the three communities can
take it up in its own language. It is perhaps worth noticing that both classifications already
answer to the same name --- the fibrations are Heinz Hopf's, and so, through the bialgebra
structure of the diagrammatic calculus, are the Hopf algebras --- a coincidence of eponymy that
the conjecture proposes to promote to a theorem. With the epistemic reading above, an
affirmative resolution would leave the tower's endpoint as one limit seen four ways: algebraic,
geometric, compositional, epistemic --- the last being the fixed-point classification of Open
Problem~\ref{op:sigma-functor}, with which this problem is cross-linked, each constraining the
other.

\subsection{What this sets up}

The entanglement-sensitivity of the Hopf maps is not a side feature but the hinge of the
physical reading that follows. Because separability is exactly the collapse of the base to a
product of Bloch spheres, the structure of the entanglement --- the SLOCC class of a multi-qubit
state, and within it the canonical maximally entangled orbit --- is carried by precisely the
degrees of freedom on which $G_{\mathrm{SM}}$ acts in Proposition~\ref{prop:gsm}. This is what
makes it possible, in \S\ref{sec:entanglement-classes}, to read the two genuinely tripartite
entanglement classes onto the $\mathbb{C} \oplus \mathbb{C}^3$ split through their canonical
representatives: the GHZ orbit confined to the colour-triplet $\mathbb{C}^3$ sector, the W orbit
reaching the $\mathbb{C}$ singlet, recovering the quark--lepton division not as an assignment
but as a consequence of the entanglement structure. The scope of those statements --- canonical
orbits rather than whole classes, and the identification of the canonical orbits with the
physical particles --- is made precise there. The Hopf fibrations are the bridge; the
entanglement classes, treated next, are the traffic that crosses it.

%
%  SOURCES (re-read 2026-07-04): lawvere_inside_outside_fixed_point_
%  theorem_2026-05-04T1806.md (the spine: the two-principles reduction,
%  the theorem, the sigma candidate, the open items -> op minted per
%  D9); research_summary_neg_prob.md (negativity route; lamp (1+i)/2);
%  dynamics_companion_fanout_unitarity (legs 1a/1c; anti-events; the
%  Brukner-Zeilinger rebuttal per D6); Pure_states_and_mixed_states
%  (local tomography derived; QBism contrast per D7).
%
%  DECISIONS ENCODED (JS 2026-07-04): D1 title as below; D2 Wheeler1990
%  cited here as the title anchor; D3 lineage Popper -> Breuer ->
%  Svozil -> Szangolies (Dalla Chiara omitted); D4 Frauchiger_2018
%  cited here, dissolved in sec:wigners-friend; D5 Kirkwood-Dirac
%  heritage via the Arvidsson-Shukur 2024 review as one-stop; D6
%  Brukner-Zeilinger + the framework rebuttal IN ("Shannon is
%  post-FANOUT"); D7 QBism contrast IN; D8 Aaronson private
%  communication OMITTED (published island paper only); D9 one op
%  minted (op:sigma-functor), cross-linked to op:hopf-frobenius.
%
%  CITATIONS -- existing keys: Wheeler1990, small_2006, Godel1931,
%  Lawvere_1969, Yanofsky2003, CoeckeKissinger2017, Rovelli1996a,
%  Popper1950a, Popper1950b, Breuer1995, svozil_k_2018,
%  szangolies_j_2018, Frauchiger_2018, dirac1942, Feynman_r_82,
%  Feynman1987, Youssef1994 [SWAP-FLAG: intended paper is "Quantum
%  mechanics as complex probability theory", Mod. Phys. Lett. A 9
%  (1994) 2571 -- if that lives under youssef1994quantum, swap],
%  caticha_a_2006, goyal_knuth_skilling_2010, Aaronson2004h, Hardy2001a,
%  muller_p_2011, Thomson_1954, benacerraf1962tasks, cantor1891diagonal,
%  Turing_1937.
%  NEW keys (RIS in self_reference_refs_2026-07-04.ris, all verified
%  this session): tikochinsky1988feynman, renou2021real, chen2022ruling,
%  wu2022refutation, hoffreumonwoods2026real, arvidssonshukur2024kirkwood,
%  bruknerzeilinger2001shannon, fuchsschack2013qbism,
%  shannon1948mathematical.
%
%  CROSS-REFS: forward (fine under two-pass LaTeX): sec:continuity,
%  sec:division-algebras, sec:hopf-justification, prop:f1-dictionary,
%  sec:wigners-friend, op:limit-realisability, op:hopf-frobenius.
%  Still dangling after this section: sec:fanout-boundary (and its
%  sibling label sec:fanout-direction) -- both DRAFTED 2026-07-06.
%
%  AUDITED 2026-07-06 (section-by-section final pass, first stop):
%  proof of prop:inside-outside re-verified by hand; attribution chain
%  clean (Burgin/Khrennikov correctly absent); all header-declared
%  citations confirmed used; lamp arithmetic and antecedents verified;
%  FANOUT smallcaps consistent; no pre-definition use of "frame".
%  Fixes: stray Unicode em-dash (QBism paragraph, the file's only
%  non-ASCII character) -> house ---; trace-blind two-way stitch added
%  to the Born subsection ("one privacy, two duties");
%  op:sigma-functor gains its first standing constraint (F1
%  inheritance via the op:hopf-frobenius discharge) and its previously
%  one-directional local-tomography cargo made explicit.
% =====================================================================

\section{Self-Reference: One Principle}
\label{sec:self-reference-2}

Wheeler asked for it from bit: for physics derived from information, ``every it --- every particle,
every field of force'' obtained from binary answers to yes-or-no questions \autocite{Wheeler1990}.
This paper's route to that programme runs through G\"{o}del: the single principle from which
everything else here is derived, deferred, or priced is \emph{self-reference} --- the fact that any
complete physical description, made from inside the world it describes, must describe the describer.
This section states the principle, proves that the framework's second historical commitment follows
from it as a theorem rather than standing beside it as an axiom, and then extracts the price of the
principle in four instalments: probabilities whose parts can cancel, amplitudes valued in
$\mathbb{C}$, the Born rule, and the local tomography of composite systems. What the section cannot
yet derive is stated as precisely as what it can, in one numbered problem and in the ledger at the
end.

\subsection{Events, not objects}

The framework's ontology is relational: the fundamental entities are events and the relations
between them, and objects are the bookkeeping of relational invariants --- what remains stable
across many relations. The kinship with relational quantum mechanics \autocite{Rovelli1996a} is
genuine and acknowledged: facts are facts relative to the systems that register them. The
difference, developed over the rest of this paper, is that the present framework grounds the
relational structure in a single logical principle and then carries it, through the division
algebras and the Hopf fibrations, all the way to the gauge group and the particle spectrum ---
relations are not only the ontology but the derivation engine. On this reading a measurement is not
the revelation of a pre-existing property; it is the creation of a relation, and the later sections
of this paper read that creation geometrically (\S\ref{sec:hopf-justification}) and even
gravitationally. Here we need only the logical core: whatever describes the world is an event in
the world it describes.

\subsection{One principle, and a theorem}

The framework historically carried two foundational commitments: self-reference, and the
\emph{Inside/Outside Equivalence Principle} --- that the laws of physics must be identical whether
formulated by an embedded participant or by an external observer. Two principles are one too many,
and this subsection reduces them.

The reduction corrects the framework's own origin, and we make the correction in print. The 2005
form of the argument \autocite{small_2006} ran a G\"{o}delian chain: a self-referential process,
being complete, must by the incompleteness theorem \autocite{Godel1931} be inconsistent; an
inconsistent system can negate itself; and the negation of self-reference is the production of
causal sequence. The chain captures the right intuition and is retained here as motivation, but it
is not sound as a proof: classical inconsistency yields \emph{ex falso quodlibet} --- every
proposition provable, every causal sequence an output, no constraint at all. The load-bearing
argument is different, and better: a fixed-point theorem.

Let $\mathcal{C}$ be the category of physical descriptions and admissible processes --- concretely,
the dagger-compact-closed setting with classical structures in which categorical quantum mechanics,
and the whole of this paper, works \autocite{CoeckeKissinger2017}. Lawvere showed that in a
sufficiently rich such setting, any weakly point-surjective self-representation forces every
endomap to have a fixed point, and Yanofsky showed that one diagonal argument of this shape
underlies Cantor's, G\"{o}del's, Turing's and Tarski's theorems alike
\autocite{Lawvere_1969,Yanofsky2003,cantor1891diagonal}. The framework's single commitment is that
$\mathcal{C}$ contains a \emph{self-representing object} --- an object capable of encoding the
processes on itself. That is what ``self-reference is foundational'' means, stated categorically,
and it is the surjectivity hypothesis of the diagonal lemma. Now let
$\sigma \colon \mathcal{C} \to \mathcal{C}$ be the \emph{inside/outside involution}: the operation
exchanging the two complementary presentations of any description --- the inside presentation, as
participant-indexed total space in which the observer is part of the system, and the outside
presentation, as a base-times-fibre decomposition over a classical parameter space. That
$\sigma^{2} = \mathrm{id}$ is not an axiom but a coherence requirement: the outside view of the
inside view of a thing is the thing. The candidate concrete realisation of $\sigma$ is the
Hopf-duality swap --- total space against base-times-fibre for the three fibrations --- which gives
the involution computable content and ties this subsection to \S\ref{sec:hopf-justification}.

\begin{proposition}[Inside/outside fixed point]\label{prop:inside-outside}
Let $\mathcal{C}$ satisfy the richness condition above and contain a self-representing object, and
let $\sigma$ be an involution on $\mathcal{C}$. Then the fixed-point locus
$\mathrm{Fix}(\sigma)$ is non-empty: there exist descriptions invariant under the exchange of
inside and outside perspective.
\end{proposition}

\begin{proof}
Lawvere's construction, applied to $\sigma$ in place of negation: given a weakly point-surjective
$\varphi \colon A \to Y^{A}$, form $g(a) = \sigma\!\left(\varphi(a)(a)\right)$; point-surjectivity
supplies $a^{*}$ with $\varphi(a^{*}) = g$, and then $\varphi(a^{*})(a^{*})$ is a fixed point of
$\sigma$ \autocite{Lawvere_1969,Yanofsky2003}. Where the same construction applied to Boolean
negation yields paradox --- negation has no fixed point on a two-element type, so the conclusion
explodes the premises --- the inside/outside involution is an honest involution on a rich type, and
the conclusion is a constructive fixed point rather than a contradiction.
\end{proof}

\textbf{Definition.} \emph{Physical law is the structure carried by} $\mathrm{Fix}(\sigma)$.

The Inside/Outside Equivalence Principle is thereby demoted from axiom to theorem: physical law is
invariant under the exchange of perspectives \emph{because it is the fixed-point set of that
exchange}, and the fixed-point set is forced to be non-empty by self-reference alone. One
commitment remains foundational; everything the second commitment asserted is now derived. The
reader who accepts categorical quantum mechanics as an existing programme has already accepted the
machinery; the framework adds one premise and collects the rest. And the fixed-point locus is not
empty of structure: it inherits the Frobenius-algebra apparatus of $\mathcal{C}$, so that the
Bell, GHZ and W primitives --- and with them the particle classification of the later sections ---
live inside $\mathrm{Fix}(\sigma)$ from the start. What is proved above is non-emptiness given the
setting; what is assumed is the setting itself and the existence of $\sigma$ as a genuine
involution; what is open is the rigour of both, and it is priced now rather than hidden:

\begin{openproblem}[The involution made functorial]\label{op:sigma-functor}
Fix the category $\mathcal{C}$ rigorously --- the dagger-compact-closed setting with classical
structures is the candidate --- and define $\sigma$ on morphisms, not only on objects, the natural
route being the induced fibre--base exchange on spider diagrams; establish that $\sigma$ is an (at
worst pseudo-)involutive endofunctor; and prove or refute the identification of
$\mathrm{Fix}(\sigma)$ with the three Hopf cross-sections, under which Adams's theorem would become
a corollary of the fixed-point classification. This problem is cross-linked with Open
Problem~\ref{op:hopf-frobenius}: under an affirmative resolution of both, the parallelisable
terminus, the fixed-point classification, and compositional completeness are three statements of
one fact, and progress on either problem constrains the other. One constraint is already in place:
at the complex rung, Proposition~\ref{prop:f1-dictionary} supplies the dictionary between
classical structures and the fibration's lifted skeletons --- the clause of
Open Problem~\ref{op:hopf-frobenius} it discharges --- and any resolution here must reproduce it.
The problem also carries a piece of cargo recorded later in this section: the formalisation, within
the fixed category, of the perspective-accounting that derives local tomography.
\end{openproblem}

\subsection{The lineage}

The obstruction this section builds on is old, and the framework claims its completion, not its
discovery. Popper argued in 1950 that a predictor embedded in the world has irreducibly inferior
knowledge of itself, because it cannot impart information about itself without, by that very act,
rendering the information out of date \autocite{Popper1950a,Popper1950b} --- the passage is quoted
where the framework makes the obstruction structural (\S\ref{sec:wigners-friend}). Breuer proved
the impossibility of accurate state self-measurement within quantum theory \autocite{Breuer1995};
Svozil's monograph surveys the whole territory where physics meets self-reference, undecidability
and Chaitin-style algorithmic limits \autocite{svozil_k_2018}; Szangolies sharpened the modern form
as epistemic horizons \autocite{szangolies_j_2018}; and the Frauchiger--Renner theorem is the
contemporary crisis-form --- nested observers deriving contradictory records when each treats the
others' quantum descriptions as universally valid \autocite{Frauchiger_2018} --- which
\S\ref{sec:wigners-friend} dissolves relationally rather than paradoxically. This framework's
addition to the lineage is twofold. Structurally, Proposition~\ref{prop:f1-dictionary} converts
Popper's temporal obsolescence into algebraic cancellation: the datum a copier cannot put into its
own record is exactly its fibre commitment, which does not arrive stale --- it cancels. And
foundationally, the present section runs the lineage in reverse: where the tradition treats
self-reference as a limit on knowledge \emph{within} physics, the framework takes it as the
generator of physics, and derives the quantum formalism from it.

\subsection{The price is negative probability}\label{sec:complex-amplitudes-2}

What does it cost a probabilistic description to contain its own describer? The first instalment is
the sign. In the regime before any record exists, every event has an anti-event --- the description
must be able to undo what it has provisionally done, or the self-description could never be
revised into consistency --- and a probability calculus closed under undoing cannot be confined to
$[0,1]$: it is a signed calculus, in which the parts of a probability can oppose and cancel. The
framework did not discover negativity as bookkeeping; Dirac admitted negative probabilities as
formal intermediaries \autocite{dirac1942}, Feynman defended them as consistent whenever the final,
recorded probabilities stay positive \autocite{Feynman_r_82,Feynman1987}, and the Kirkwood--Dirac
distribution has carried complex quasiprobability from 1933 to the present, as a recent
comprehensive review documents \autocite{arvidssonshukur2024kirkwood}. What the framework adds is
the forcing: negativity is not a permissible convenience but the first price of self-referential
closure, and it is charged exactly and only in the pre-record regime --- which is why every
negativity in the tradition lives in intermediate, unrecorded quantities and never in a final
record.

The second instalment is complexity. A signed information measure must extend Shannon's
$S = -\ln p$ off the positive axis, and the analytic continuation of the logarithm carries $p$ into
the complex plane: amplitudes, with the phase as the continued information. Here too the extension
has a prior tradition, and the attributions are owed precisely. Tikochinsky derived Feynman's rules
for complex amplitudes from consistency of composition \autocite{tikochinsky1988feynman}; Youssef
reformulated quantum mechanics as complex probability theory outright \autocite{Youssef1994};
Caticha carried objective amplitudes to Bayesian probabilities \autocite{caticha_a_2006}; and
Goyal, Knuth and Skilling derived Feynman's rules from a pair-postulate without assuming
analyticity \autocite{goyal_knuth_skilling_2010}. The extension is prior art. The framework's
claim is the compulsion: the complex extension is not a modelling choice among exotic probability
theories but the unique continuation forced by self-referential closure --- the same forcing that
Aaronson glimpsed when he argued quantum mechanics is what negative probability inevitably becomes
\autocite{Aaronson2004h}.

Why does the continuation stop at $\mathbb{C}$? Downward, the real field is not available: real
quantum theory is experimentally distinguishable from complex quantum theory in network scenarios
and has been disfavoured by experiment
\autocite{renou2021real,chen2022ruling,wu2022refutation}; the recent objection that the exclusion
rests on locally tomographic composition \autocite{hoffreumonwoods2026real} is, for this framework,
not an escape but a confirmation of where the weight sits --- local tomography is \emph{derived}
below, not assumed, so the framework closes from its own resources exactly the door the objection
holds open. Upward, the quaternions and octonions fail as a \emph{base} for the information
continuation: the continued measure is a logarithm, and beyond $\mathbb{C}$ the logarithm of a
product is no longer the sum of logarithms --- non-commutativity severs
$\log(ab) = \log a + \log b$ --- so the additivity that makes it an information measure is lost.
Aaronson's parameter-counting and signalling arguments against a quaternionic base point the same
way \autocite{Aaronson2004h}. The higher algebras are not banished; they return as the
\emph{internal} algebras of the register, which is the business of \S\ref{sec:division-algebras}.
The base is $\mathbb{C}$, for informational reasons, and stays $\mathbb{C}$.

\subsection{The Born rule from fibre privacy}\label{sec:born-rule}

The third instalment is the quadratic rule. An amplitude carries one datum more than any record of
it can: its overall phase. This was asserted throughout the framework's development as the
observer's private, self-referential ignorance; it is now a proved statement, in exactly the form
needed. By Proposition~\ref{prop:f1-dictionary}, a classical structure --- the copy operation any
record requires --- is committed to a specific lift of its skeleton, and that commitment is absent
from every classical record the structure induces: the copier is fibre-committed, the record is
fibre-blind, and the fibre freedom is gauge for every observer. (The computation extends verbatim
from one qubit to the full register, where Proposition~\ref{prop:trace-blind} makes it the ground
of hypercharge's tracelessness --- one privacy, two duties.) Popper's obsolescence, made
cancellation. It follows that anything recordable is invariant under the fibre --- observables and
their statistics must be $U(1)$-invariant functions of the amplitude. The lowest-degree polynomial invariant
of a complex amplitude is the quadratic $|\psi|^{2}$ --- the non-polynomial invariant $|\psi|$,
of degree one, is excluded by the analyticity declared at \S\ref{sec:complex-amplitudes} --- and
normalisation fixes the coefficient:
probability is the minimal record-compatible function of the amplitude,
\[
p \;=\; |\psi|^{2},
\]
the Born rule as the cheapest invariant the fibre permits. The same square is reached from the
cost side in \S\ref{sec:born-cost-ledger}, where a record is the closed loop of a history and its
return leg and the exponent is the number of legs; neither route uses the lattice of projections
that Gleason's theorem needs, and both hold for a single qubit. The derivation is the framework's
own;
its premise --- that the fibre datum is structurally excluded from records --- is
Proposition~\ref{prop:f1-dictionary}, and its geometric reading --- probability as base-point data
under the Hopf projection --- is developed in \S\ref{sec:hopf-justification}.

One standing objection to information-theoretic foundations is answered here, because this
framework needs the answer. Brukner and Zeilinger argued that Shannon information is conceptually
inadequate for quantum measurements: it is basis-dependent, and its operational meaning presupposes
ignorance of a definite but unknown fact, which the quantum case does not supply
\autocite{bruknerzeilinger2001shannon}. The framework's reply is structural, and it is a
demarcation rather than a rebuttal. Shannon's precondition fails exactly where they say it fails
--- in the pre-record regime, where there is no definite unknown to be ignorant of --- and there
the operative measure is the quadratic invariant just derived. Shannon's axioms hold exactly where
the precondition holds --- on records, after \textsc{fanout}, where outcomes are definite and
ignorance is classical --- and there the uniqueness theorem forces $-\log p$
\autocite{shannon1948mathematical}. The basis-dependence dissolves with the same stroke: a basis
\emph{is} selected precisely at \textsc{fanout}, so the basis-dependent measure applies only where
a basis exists. Shannon is not inadequate; Shannon is post-\textsc{fanout}. The framework thereby
keeps both measures, each on its own side of the boundary, and uses the additivity of the
post-record measure --- the very property the previous subsection used to pin the base field ---
with a clear conscience.

\subsection{Thomson's lamp, resolved}

The section owes \S\ref{sec:continuity} a debt incurred there: the lamp's amplitude-side limit.
Thomson's lamp --- toggled at times $1 - 2^{-n}$, its state at $t = 1$ demanded --- has no
classical answer: the task description fixes every finite stage and thereby fixes no limit at all
\autocite{Thomson_1954,benacerraf1962tasks}, which \S\ref{sec:continuity} read as the proof that
the record algebra is not limit-closed. On the amplitude side the limit exists, and the calculus of
this section computes it. The toggle exchanges the two amplitudes; the self-consistent state of the
completed task is the fixed point of the exchange in the complex calculus, and solving it under the
signed continuation gives
\[
P(\mathrm{on}) = \tfrac{1 + i}{2}, \qquad P(\mathrm{off}) = \tfrac{1 - i}{2},
\]
a superposed fixed point whose imaginary part carries the unrecordable datum --- the phase of the
endless alternation, when the lamp \emph{would have} stopped. Two checks certify that the
assignment is calculus rather than decoration. The pair is normalised,
$P(\mathrm{on}) + P(\mathrm{off}) = 1$; and the Born rule of the previous subsection, applied to
the fixed point, returns $|P(\mathrm{on})|^{2} = \tfrac{1}{2}$ --- exactly the symmetric classical
ignorance that is the only sane record-level answer, recovered rather than imposed. The lamp,
which entered the framework at its very beginning as the place where classical computation meets
its boundary, thus exits it as a worked example: undetermined as a record, determinate as an
amplitude, the two sides split precisely along \textsc{fanout} --- and the diagonal construction it
physically instantiates \autocite{cantor1891diagonal,Turing_1937} is the same one that powers
Proposition~\ref{prop:inside-outside}. The status flag is owed here as everywhere: the
undetermination is theorem-level; the specific fixed-point value is the framework's complex
calculus, structural, offered with its two consistency checks shown.

\subsection{Local tomography, derived}\label{sec:local-tomography}

The fourth instalment concerns composites, and it discharges two promises at once. The
Inside/Outside Equivalence Principle --- now a theorem --- requires that an external observer's
characterisation of a composite system be recoverable by the embedded participants, and the
embedded participants have access only to local measurements and their correlations. Equivalence
of the two perspectives therefore entails that the state of a composite is determined by the
statistics of local measurements: \emph{local tomography}, obtained as a consequence rather than
imposed as an axiom. This converts a postulate of the reconstruction programmes
\autocite{Hardy2001a,mullerp2011} into a theorem-shaped consequence of the one principle ---
discharging the caveat recorded in \S\ref{sec:continuity} --- and it closes, from the framework's
own resources, the escape route in the real-number debate noted above: the assumption the
objection of \autocite{hoffreumonwoods2026real} contests is here not an assumption. The status is
flagged plainly: the derivation is structural, one paragraph of perspective-accounting, and its
formalisation within the categorical setting belongs to Open Problem~\ref{op:sigma-functor}.

\subsection{Contrast: QBism, and the ledger}

The nearest neighbour in spirit deserves a paragraph of honest contrast. QBism
\autocite{fuchsschack2013qbism} holds, as this framework does, that quantum probabilities are an
agent's --- that the theory is a calculus of expectation, not a catalogue of properties. What QBism
does not say is why the agent's expectations must be complex-valued and quadratically squared; the
Born rule enters as a normative addition to Bayesian coherence, and the complex amplitude as an
unexplained representation. This section is, in effect, the missing why: the agent's ignorance is
self-referential --- the agent is inside the description --- and self-referential ignorance is
complex, with the quadratic rule as the cheapest record-compatible invariant of it. Where QBism
stops at the agent, the framework derives the agent's calculus.

The ledger, then. Delivered in this section: the reduction of two principles to one
(Proposition~\ref{prop:inside-outside}); the pre-record regime with its anti-events, handed to
\S\ref{sec:continuity} as the five-face regime whose continuity is priced there
(Open Problem~\ref{op:limit-realisability}); the complex base, forced downward by experiment and
upward by the additivity of information; the Born rule from fibre privacy; the lamp's fixed point;
and local tomography. Owed, and stated where they are owed: the functorial rigour of the involution
and the fixed-point classification (Open Problem~\ref{op:sigma-functor}); the remaining premise of continuity (Open Problem~\ref{op:limit-realisability}, in
\S\ref{sec:continuity}); and the
rung-by-rung compulsion of the algebra ladder, posed as Open
Problem~\ref{op:ladder-compulsion} in \S\ref{sec:division-algebras}. The
hand-on is the paper's spine: the $\mathbb{C}$ this section derives natively and the $\mathbb{C}$
that \S\ref{sec:continuity} underwrites through Hardy's theorem are the same field, the floor of
the ladder that \S\ref{sec:division-algebras} climbs and \S\ref{sec:hopf-justification}
geometrises. One principle; everything else is theorem, derivation, or a debt with a stated
invoice.

%      % EDITED 2026-10-29 scripts/patch_continuity_tower_premise_2026-10-29T1500.py
% EDITED 2026-09-15 scripts/patch_value_continuum_declared_cencov_2026-09-15T1540.py
% EDITED 2026-09-15 scripts/patch_continuity_connected_state_space_2026-09-15T1430.py
% =====================================================================
%  SECTION: Continuity Is Not an Axiom
%  Drafted 2026-07-04 from continuity_limit_realisability_2026-07-04.md
%  (now in project documents) and the May record
%  (dynamics_companion_fanout_unitarity_2026-05-20T1430.md).
%
%  PLACEMENT in main.tex: a standalone section between self-reference
%  and division-algebras --
%      % EDITED 2026-10-29 scripts/patch_continuity_tower_premise_2026-10-29T1500.py
% EDITED 2026-09-20 scripts/patch_born_cost_ledger_v2_2026-09-20T0900.py
% =====================================================================
%  SECTION: Self-Reference: One Principle
%  Drafted 2026-07-04. The paper's foundational section; resolves the
%  largest dangling label in the manuscript (sec:self-reference).
%
%  PLACEMENT in main.tex: FIRST of the foundations chain --
%      \input{self_reference_section_2026-07-04}
%      \input{continuity_section_2026-07-04}
%      \input{division_algebras_section_2026-07-03}
%      \input{hopf_justification_section_2026-06-08}
%
%  SOURCES (re-read 2026-07-04): lawvere_inside_outside_fixed_point_
%  theorem_2026-05-04T1806.md (the spine: the two-principles reduction,
%  the theorem, the sigma candidate, the open items -> op minted per
%  D9); research_summary_neg_prob.md (negativity route; lamp (1+i)/2);
%  dynamics_companion_fanout_unitarity (legs 1a/1c; anti-events; the
%  Brukner-Zeilinger rebuttal per D6); Pure_states_and_mixed_states
%  (local tomography derived; QBism contrast per D7).
%
%  DECISIONS ENCODED (JS 2026-07-04): D1 title as below; D2 Wheeler1990
%  cited here as the title anchor; D3 lineage Popper -> Breuer ->
%  Svozil -> Szangolies (Dalla Chiara omitted); D4 Frauchiger_2018
%  cited here, dissolved in sec:wigners-friend; D5 Kirkwood-Dirac
%  heritage via the Arvidsson-Shukur 2024 review as one-stop; D6
%  Brukner-Zeilinger + the framework rebuttal IN ("Shannon is
%  post-FANOUT"); D7 QBism contrast IN; D8 Aaronson private
%  communication OMITTED (published island paper only); D9 one op
%  minted (op:sigma-functor), cross-linked to op:hopf-frobenius.
%
%  CITATIONS -- existing keys: Wheeler1990, small_2006, Godel1931,
%  Lawvere_1969, Yanofsky2003, CoeckeKissinger2017, Rovelli1996a,
%  Popper1950a, Popper1950b, Breuer1995, svozil_k_2018,
%  szangolies_j_2018, Frauchiger_2018, dirac1942, Feynman_r_82,
%  Feynman1987, Youssef1994 [SWAP-FLAG: intended paper is "Quantum
%  mechanics as complex probability theory", Mod. Phys. Lett. A 9
%  (1994) 2571 -- if that lives under youssef1994quantum, swap],
%  caticha_a_2006, goyal_knuth_skilling_2010, Aaronson2004h, Hardy2001a,
%  muller_p_2011, Thomson_1954, benacerraf1962tasks, cantor1891diagonal,
%  Turing_1937.
%  NEW keys (RIS in self_reference_refs_2026-07-04.ris, all verified
%  this session): tikochinsky1988feynman, renou2021real, chen2022ruling,
%  wu2022refutation, hoffreumonwoods2026real, arvidssonshukur2024kirkwood,
%  bruknerzeilinger2001shannon, fuchsschack2013qbism,
%  shannon1948mathematical.
%
%  CROSS-REFS: forward (fine under two-pass LaTeX): sec:continuity,
%  sec:division-algebras, sec:hopf-justification, prop:f1-dictionary,
%  sec:wigners-friend, op:limit-realisability, op:hopf-frobenius.
%  Still dangling after this section: sec:fanout-boundary (and its
%  sibling label sec:fanout-direction) -- both DRAFTED 2026-07-06.
%
%  AUDITED 2026-07-06 (section-by-section final pass, first stop):
%  proof of prop:inside-outside re-verified by hand; attribution chain
%  clean (Burgin/Khrennikov correctly absent); all header-declared
%  citations confirmed used; lamp arithmetic and antecedents verified;
%  FANOUT smallcaps consistent; no pre-definition use of "frame".
%  Fixes: stray Unicode em-dash (QBism paragraph, the file's only
%  non-ASCII character) -> house ---; trace-blind two-way stitch added
%  to the Born subsection ("one privacy, two duties");
%  op:sigma-functor gains its first standing constraint (F1
%  inheritance via the op:hopf-frobenius discharge) and its previously
%  one-directional local-tomography cargo made explicit.
% =====================================================================

\section{Self-Reference: One Principle}
\label{sec:self-reference-2}

Wheeler asked for it from bit: for physics derived from information, ``every it --- every particle,
every field of force'' obtained from binary answers to yes-or-no questions \autocite{Wheeler1990}.
This paper's route to that programme runs through G\"{o}del: the single principle from which
everything else here is derived, deferred, or priced is \emph{self-reference} --- the fact that any
complete physical description, made from inside the world it describes, must describe the describer.
This section states the principle, proves that the framework's second historical commitment follows
from it as a theorem rather than standing beside it as an axiom, and then extracts the price of the
principle in four instalments: probabilities whose parts can cancel, amplitudes valued in
$\mathbb{C}$, the Born rule, and the local tomography of composite systems. What the section cannot
yet derive is stated as precisely as what it can, in one numbered problem and in the ledger at the
end.

\subsection{Events, not objects}

The framework's ontology is relational: the fundamental entities are events and the relations
between them, and objects are the bookkeeping of relational invariants --- what remains stable
across many relations. The kinship with relational quantum mechanics \autocite{Rovelli1996a} is
genuine and acknowledged: facts are facts relative to the systems that register them. The
difference, developed over the rest of this paper, is that the present framework grounds the
relational structure in a single logical principle and then carries it, through the division
algebras and the Hopf fibrations, all the way to the gauge group and the particle spectrum ---
relations are not only the ontology but the derivation engine. On this reading a measurement is not
the revelation of a pre-existing property; it is the creation of a relation, and the later sections
of this paper read that creation geometrically (\S\ref{sec:hopf-justification}) and even
gravitationally. Here we need only the logical core: whatever describes the world is an event in
the world it describes.

\subsection{One principle, and a theorem}

The framework historically carried two foundational commitments: self-reference, and the
\emph{Inside/Outside Equivalence Principle} --- that the laws of physics must be identical whether
formulated by an embedded participant or by an external observer. Two principles are one too many,
and this subsection reduces them.

The reduction corrects the framework's own origin, and we make the correction in print. The 2005
form of the argument \autocite{small_2006} ran a G\"{o}delian chain: a self-referential process,
being complete, must by the incompleteness theorem \autocite{Godel1931} be inconsistent; an
inconsistent system can negate itself; and the negation of self-reference is the production of
causal sequence. The chain captures the right intuition and is retained here as motivation, but it
is not sound as a proof: classical inconsistency yields \emph{ex falso quodlibet} --- every
proposition provable, every causal sequence an output, no constraint at all. The load-bearing
argument is different, and better: a fixed-point theorem.

Let $\mathcal{C}$ be the category of physical descriptions and admissible processes --- concretely,
the dagger-compact-closed setting with classical structures in which categorical quantum mechanics,
and the whole of this paper, works \autocite{CoeckeKissinger2017}. Lawvere showed that in a
sufficiently rich such setting, any weakly point-surjective self-representation forces every
endomap to have a fixed point, and Yanofsky showed that one diagonal argument of this shape
underlies Cantor's, G\"{o}del's, Turing's and Tarski's theorems alike
\autocite{Lawvere_1969,Yanofsky2003,cantor1891diagonal}. The framework's single commitment is that
$\mathcal{C}$ contains a \emph{self-representing object} --- an object capable of encoding the
processes on itself. That is what ``self-reference is foundational'' means, stated categorically,
and it is the surjectivity hypothesis of the diagonal lemma. Now let
$\sigma \colon \mathcal{C} \to \mathcal{C}$ be the \emph{inside/outside involution}: the operation
exchanging the two complementary presentations of any description --- the inside presentation, as
participant-indexed total space in which the observer is part of the system, and the outside
presentation, as a base-times-fibre decomposition over a classical parameter space. That
$\sigma^{2} = \mathrm{id}$ is not an axiom but a coherence requirement: the outside view of the
inside view of a thing is the thing. The candidate concrete realisation of $\sigma$ is the
Hopf-duality swap --- total space against base-times-fibre for the three fibrations --- which gives
the involution computable content and ties this subsection to \S\ref{sec:hopf-justification}.

\begin{proposition}[Inside/outside fixed point]\label{prop:inside-outside}
Let $\mathcal{C}$ satisfy the richness condition above and contain a self-representing object, and
let $\sigma$ be an involution on $\mathcal{C}$. Then the fixed-point locus
$\mathrm{Fix}(\sigma)$ is non-empty: there exist descriptions invariant under the exchange of
inside and outside perspective.
\end{proposition}

\begin{proof}
Lawvere's construction, applied to $\sigma$ in place of negation: given a weakly point-surjective
$\varphi \colon A \to Y^{A}$, form $g(a) = \sigma\!\left(\varphi(a)(a)\right)$; point-surjectivity
supplies $a^{*}$ with $\varphi(a^{*}) = g$, and then $\varphi(a^{*})(a^{*})$ is a fixed point of
$\sigma$ \autocite{Lawvere_1969,Yanofsky2003}. Where the same construction applied to Boolean
negation yields paradox --- negation has no fixed point on a two-element type, so the conclusion
explodes the premises --- the inside/outside involution is an honest involution on a rich type, and
the conclusion is a constructive fixed point rather than a contradiction.
\end{proof}

\textbf{Definition.} \emph{Physical law is the structure carried by} $\mathrm{Fix}(\sigma)$.

The Inside/Outside Equivalence Principle is thereby demoted from axiom to theorem: physical law is
invariant under the exchange of perspectives \emph{because it is the fixed-point set of that
exchange}, and the fixed-point set is forced to be non-empty by self-reference alone. One
commitment remains foundational; everything the second commitment asserted is now derived. The
reader who accepts categorical quantum mechanics as an existing programme has already accepted the
machinery; the framework adds one premise and collects the rest. And the fixed-point locus is not
empty of structure: it inherits the Frobenius-algebra apparatus of $\mathcal{C}$, so that the
Bell, GHZ and W primitives --- and with them the particle classification of the later sections ---
live inside $\mathrm{Fix}(\sigma)$ from the start. What is proved above is non-emptiness given the
setting; what is assumed is the setting itself and the existence of $\sigma$ as a genuine
involution; what is open is the rigour of both, and it is priced now rather than hidden:

\begin{openproblem}[The involution made functorial]\label{op:sigma-functor}
Fix the category $\mathcal{C}$ rigorously --- the dagger-compact-closed setting with classical
structures is the candidate --- and define $\sigma$ on morphisms, not only on objects, the natural
route being the induced fibre--base exchange on spider diagrams; establish that $\sigma$ is an (at
worst pseudo-)involutive endofunctor; and prove or refute the identification of
$\mathrm{Fix}(\sigma)$ with the three Hopf cross-sections, under which Adams's theorem would become
a corollary of the fixed-point classification. This problem is cross-linked with Open
Problem~\ref{op:hopf-frobenius}: under an affirmative resolution of both, the parallelisable
terminus, the fixed-point classification, and compositional completeness are three statements of
one fact, and progress on either problem constrains the other. One constraint is already in place:
at the complex rung, Proposition~\ref{prop:f1-dictionary} supplies the dictionary between
classical structures and the fibration's lifted skeletons --- the clause of
Open Problem~\ref{op:hopf-frobenius} it discharges --- and any resolution here must reproduce it.
The problem also carries a piece of cargo recorded later in this section: the formalisation, within
the fixed category, of the perspective-accounting that derives local tomography.
\end{openproblem}

\subsection{The lineage}

The obstruction this section builds on is old, and the framework claims its completion, not its
discovery. Popper argued in 1950 that a predictor embedded in the world has irreducibly inferior
knowledge of itself, because it cannot impart information about itself without, by that very act,
rendering the information out of date \autocite{Popper1950a,Popper1950b} --- the passage is quoted
where the framework makes the obstruction structural (\S\ref{sec:wigners-friend}). Breuer proved
the impossibility of accurate state self-measurement within quantum theory \autocite{Breuer1995};
Svozil's monograph surveys the whole territory where physics meets self-reference, undecidability
and Chaitin-style algorithmic limits \autocite{svozil_k_2018}; Szangolies sharpened the modern form
as epistemic horizons \autocite{szangolies_j_2018}; and the Frauchiger--Renner theorem is the
contemporary crisis-form --- nested observers deriving contradictory records when each treats the
others' quantum descriptions as universally valid \autocite{Frauchiger_2018} --- which
\S\ref{sec:wigners-friend} dissolves relationally rather than paradoxically. This framework's
addition to the lineage is twofold. Structurally, Proposition~\ref{prop:f1-dictionary} converts
Popper's temporal obsolescence into algebraic cancellation: the datum a copier cannot put into its
own record is exactly its fibre commitment, which does not arrive stale --- it cancels. And
foundationally, the present section runs the lineage in reverse: where the tradition treats
self-reference as a limit on knowledge \emph{within} physics, the framework takes it as the
generator of physics, and derives the quantum formalism from it.

\subsection{The price is negative probability}\label{sec:complex-amplitudes-2}

What does it cost a probabilistic description to contain its own describer? The first instalment is
the sign. In the regime before any record exists, every event has an anti-event --- the description
must be able to undo what it has provisionally done, or the self-description could never be
revised into consistency --- and a probability calculus closed under undoing cannot be confined to
$[0,1]$: it is a signed calculus, in which the parts of a probability can oppose and cancel. The
framework did not discover negativity as bookkeeping; Dirac admitted negative probabilities as
formal intermediaries \autocite{dirac1942}, Feynman defended them as consistent whenever the final,
recorded probabilities stay positive \autocite{Feynman_r_82,Feynman1987}, and the Kirkwood--Dirac
distribution has carried complex quasiprobability from 1933 to the present, as a recent
comprehensive review documents \autocite{arvidssonshukur2024kirkwood}. What the framework adds is
the forcing: negativity is not a permissible convenience but the first price of self-referential
closure, and it is charged exactly and only in the pre-record regime --- which is why every
negativity in the tradition lives in intermediate, unrecorded quantities and never in a final
record.

The second instalment is complexity. A signed information measure must extend Shannon's
$S = -\ln p$ off the positive axis, and the analytic continuation of the logarithm carries $p$ into
the complex plane: amplitudes, with the phase as the continued information. Here too the extension
has a prior tradition, and the attributions are owed precisely. Tikochinsky derived Feynman's rules
for complex amplitudes from consistency of composition \autocite{tikochinsky1988feynman}; Youssef
reformulated quantum mechanics as complex probability theory outright \autocite{Youssef1994};
Caticha carried objective amplitudes to Bayesian probabilities \autocite{caticha_a_2006}; and
Goyal, Knuth and Skilling derived Feynman's rules from a pair-postulate without assuming
analyticity \autocite{goyal_knuth_skilling_2010}. The extension is prior art. The framework's
claim is the compulsion: the complex extension is not a modelling choice among exotic probability
theories but the unique continuation forced by self-referential closure --- the same forcing that
Aaronson glimpsed when he argued quantum mechanics is what negative probability inevitably becomes
\autocite{Aaronson2004h}.

Why does the continuation stop at $\mathbb{C}$? Downward, the real field is not available: real
quantum theory is experimentally distinguishable from complex quantum theory in network scenarios
and has been disfavoured by experiment
\autocite{renou2021real,chen2022ruling,wu2022refutation}; the recent objection that the exclusion
rests on locally tomographic composition \autocite{hoffreumonwoods2026real} is, for this framework,
not an escape but a confirmation of where the weight sits --- local tomography is \emph{derived}
below, not assumed, so the framework closes from its own resources exactly the door the objection
holds open. Upward, the quaternions and octonions fail as a \emph{base} for the information
continuation: the continued measure is a logarithm, and beyond $\mathbb{C}$ the logarithm of a
product is no longer the sum of logarithms --- non-commutativity severs
$\log(ab) = \log a + \log b$ --- so the additivity that makes it an information measure is lost.
Aaronson's parameter-counting and signalling arguments against a quaternionic base point the same
way \autocite{Aaronson2004h}. The higher algebras are not banished; they return as the
\emph{internal} algebras of the register, which is the business of \S\ref{sec:division-algebras}.
The base is $\mathbb{C}$, for informational reasons, and stays $\mathbb{C}$.

\subsection{The Born rule from fibre privacy}\label{sec:born-rule}

The third instalment is the quadratic rule. An amplitude carries one datum more than any record of
it can: its overall phase. This was asserted throughout the framework's development as the
observer's private, self-referential ignorance; it is now a proved statement, in exactly the form
needed. By Proposition~\ref{prop:f1-dictionary}, a classical structure --- the copy operation any
record requires --- is committed to a specific lift of its skeleton, and that commitment is absent
from every classical record the structure induces: the copier is fibre-committed, the record is
fibre-blind, and the fibre freedom is gauge for every observer. (The computation extends verbatim
from one qubit to the full register, where Proposition~\ref{prop:trace-blind} makes it the ground
of hypercharge's tracelessness --- one privacy, two duties.) Popper's obsolescence, made
cancellation. It follows that anything recordable is invariant under the fibre --- observables and
their statistics must be $U(1)$-invariant functions of the amplitude. The lowest-degree polynomial invariant
of a complex amplitude is the quadratic $|\psi|^{2}$ --- the non-polynomial invariant $|\psi|$,
of degree one, is excluded by the analyticity declared at \S\ref{sec:complex-amplitudes} --- and
normalisation fixes the coefficient:
probability is the minimal record-compatible function of the amplitude,
\[
p \;=\; |\psi|^{2},
\]
the Born rule as the cheapest invariant the fibre permits. The same square is reached from the
cost side in \S\ref{sec:born-cost-ledger}, where a record is the closed loop of a history and its
return leg and the exponent is the number of legs; neither route uses the lattice of projections
that Gleason's theorem needs, and both hold for a single qubit. The derivation is the framework's
own;
its premise --- that the fibre datum is structurally excluded from records --- is
Proposition~\ref{prop:f1-dictionary}, and its geometric reading --- probability as base-point data
under the Hopf projection --- is developed in \S\ref{sec:hopf-justification}.

One standing objection to information-theoretic foundations is answered here, because this
framework needs the answer. Brukner and Zeilinger argued that Shannon information is conceptually
inadequate for quantum measurements: it is basis-dependent, and its operational meaning presupposes
ignorance of a definite but unknown fact, which the quantum case does not supply
\autocite{bruknerzeilinger2001shannon}. The framework's reply is structural, and it is a
demarcation rather than a rebuttal. Shannon's precondition fails exactly where they say it fails
--- in the pre-record regime, where there is no definite unknown to be ignorant of --- and there
the operative measure is the quadratic invariant just derived. Shannon's axioms hold exactly where
the precondition holds --- on records, after \textsc{fanout}, where outcomes are definite and
ignorance is classical --- and there the uniqueness theorem forces $-\log p$
\autocite{shannon1948mathematical}. The basis-dependence dissolves with the same stroke: a basis
\emph{is} selected precisely at \textsc{fanout}, so the basis-dependent measure applies only where
a basis exists. Shannon is not inadequate; Shannon is post-\textsc{fanout}. The framework thereby
keeps both measures, each on its own side of the boundary, and uses the additivity of the
post-record measure --- the very property the previous subsection used to pin the base field ---
with a clear conscience.

\subsection{Thomson's lamp, resolved}

The section owes \S\ref{sec:continuity} a debt incurred there: the lamp's amplitude-side limit.
Thomson's lamp --- toggled at times $1 - 2^{-n}$, its state at $t = 1$ demanded --- has no
classical answer: the task description fixes every finite stage and thereby fixes no limit at all
\autocite{Thomson_1954,benacerraf1962tasks}, which \S\ref{sec:continuity} read as the proof that
the record algebra is not limit-closed. On the amplitude side the limit exists, and the calculus of
this section computes it. The toggle exchanges the two amplitudes; the self-consistent state of the
completed task is the fixed point of the exchange in the complex calculus, and solving it under the
signed continuation gives
\[
P(\mathrm{on}) = \tfrac{1 + i}{2}, \qquad P(\mathrm{off}) = \tfrac{1 - i}{2},
\]
a superposed fixed point whose imaginary part carries the unrecordable datum --- the phase of the
endless alternation, when the lamp \emph{would have} stopped. Two checks certify that the
assignment is calculus rather than decoration. The pair is normalised,
$P(\mathrm{on}) + P(\mathrm{off}) = 1$; and the Born rule of the previous subsection, applied to
the fixed point, returns $|P(\mathrm{on})|^{2} = \tfrac{1}{2}$ --- exactly the symmetric classical
ignorance that is the only sane record-level answer, recovered rather than imposed. The lamp,
which entered the framework at its very beginning as the place where classical computation meets
its boundary, thus exits it as a worked example: undetermined as a record, determinate as an
amplitude, the two sides split precisely along \textsc{fanout} --- and the diagonal construction it
physically instantiates \autocite{cantor1891diagonal,Turing_1937} is the same one that powers
Proposition~\ref{prop:inside-outside}. The status flag is owed here as everywhere: the
undetermination is theorem-level; the specific fixed-point value is the framework's complex
calculus, structural, offered with its two consistency checks shown.

\subsection{Local tomography, derived}\label{sec:local-tomography}

The fourth instalment concerns composites, and it discharges two promises at once. The
Inside/Outside Equivalence Principle --- now a theorem --- requires that an external observer's
characterisation of a composite system be recoverable by the embedded participants, and the
embedded participants have access only to local measurements and their correlations. Equivalence
of the two perspectives therefore entails that the state of a composite is determined by the
statistics of local measurements: \emph{local tomography}, obtained as a consequence rather than
imposed as an axiom. This converts a postulate of the reconstruction programmes
\autocite{Hardy2001a,mullerp2011} into a theorem-shaped consequence of the one principle ---
discharging the caveat recorded in \S\ref{sec:continuity} --- and it closes, from the framework's
own resources, the escape route in the real-number debate noted above: the assumption the
objection of \autocite{hoffreumonwoods2026real} contests is here not an assumption. The status is
flagged plainly: the derivation is structural, one paragraph of perspective-accounting, and its
formalisation within the categorical setting belongs to Open Problem~\ref{op:sigma-functor}.

\subsection{Contrast: QBism, and the ledger}

The nearest neighbour in spirit deserves a paragraph of honest contrast. QBism
\autocite{fuchsschack2013qbism} holds, as this framework does, that quantum probabilities are an
agent's --- that the theory is a calculus of expectation, not a catalogue of properties. What QBism
does not say is why the agent's expectations must be complex-valued and quadratically squared; the
Born rule enters as a normative addition to Bayesian coherence, and the complex amplitude as an
unexplained representation. This section is, in effect, the missing why: the agent's ignorance is
self-referential --- the agent is inside the description --- and self-referential ignorance is
complex, with the quadratic rule as the cheapest record-compatible invariant of it. Where QBism
stops at the agent, the framework derives the agent's calculus.

The ledger, then. Delivered in this section: the reduction of two principles to one
(Proposition~\ref{prop:inside-outside}); the pre-record regime with its anti-events, handed to
\S\ref{sec:continuity} as the five-face regime whose continuity is priced there
(Open Problem~\ref{op:limit-realisability}); the complex base, forced downward by experiment and
upward by the additivity of information; the Born rule from fibre privacy; the lamp's fixed point;
and local tomography. Owed, and stated where they are owed: the functorial rigour of the involution
and the fixed-point classification (Open Problem~\ref{op:sigma-functor}); the remaining premise of continuity (Open Problem~\ref{op:limit-realisability}, in
\S\ref{sec:continuity}); and the
rung-by-rung compulsion of the algebra ladder, posed as Open
Problem~\ref{op:ladder-compulsion} in \S\ref{sec:division-algebras}. The
hand-on is the paper's spine: the $\mathbb{C}$ this section derives natively and the $\mathbb{C}$
that \S\ref{sec:continuity} underwrites through Hardy's theorem are the same field, the floor of
the ladder that \S\ref{sec:division-algebras} climbs and \S\ref{sec:hopf-justification}
geometrises. One principle; everything else is theorem, derivation, or a debt with a stated
invoice.

%      % EDITED 2026-10-29 scripts/patch_continuity_tower_premise_2026-10-29T1500.py
% EDITED 2026-09-15 scripts/patch_value_continuum_declared_cencov_2026-09-15T1540.py
% EDITED 2026-09-15 scripts/patch_continuity_connected_state_space_2026-09-15T1430.py
% =====================================================================
%  SECTION: Continuity Is Not an Axiom
%  Drafted 2026-07-04 from continuity_limit_realisability_2026-07-04.md
%  (now in project documents) and the May record
%  (dynamics_companion_fanout_unitarity_2026-05-20T1430.md).
%
%  PLACEMENT in main.tex: a standalone section between self-reference
%  and division-algebras --
%      \input{self_reference_section_2026-07-04}
%      \input{continuity_section_2026-07-04}
%      \input{division_algebras_section_2026-07-03}
%  (Architecture call SETTLED by the sync manifest: standalone.)
%
%  CITATIONS -- existing keys (as supplied by JS 2026-07-04):
%    Hardy2001a, Aaronson2004h, Thomson_1954, Lawvere_1969,
%    Turing_1937, muller_p_2011, Brukner2009.
%  NEW keys (RIS in continuity_refs_2026-07-04.ris):
%    cantor1891diagonal, benacerraf1962tasks, broeckertomdieck1985compact.
%  Deliberately excluded per JS: Chiribella--D'Ariano--Perinotti;
%  Bargmann/Wigner's theorem. Yanofsky2003 now paired on the
%  Lawvere-shaped clause (the pending-key excuse was void: the key
%  exists and sec:self-reference uses it).
%
%  CROSS-REFS: sec:self-reference now DRAFTED (2026-07-04, later the
%  same day) -- no longer dangling; only sec:fanout-boundary (and its
%  sibling sec:fanout-direction) -- DRAFTED 2026-07-06.
%
%  AUDITED 2026-07-06 (final pass, second stop): both proofs
%  re-verified by hand (Lagrange divisibility sentence exact; the
%  Baire countable-escape chain sound); caveat-discharge with
%  sec:self-reference confirmed two-way (already present); file pure
%  ASCII; no pre-definition "frame"; all citations used as declared.
%  Fixes: Yanofsky2003 paired on the Lawvere clause; the op's two
%  candidate routes anchored to printed propositions (f1-dictionary
%  for the divisibility premise's referent, inside-outside for the
%  closedness demand) -- statuses unchanged, referents sharpened;
%  header staleness corrected. Deliberate FORWARD references --
%  prop:f1-dictionary and sec:hopf-justification (defined later in the
%  paper; fine under two-pass LaTeX and flagged in prose),
%  sec:division-algebras, sec:space-of-computations.
%
%  DECISIONS ENCODED: op:limit-realisability minted (the owed arrow,
%  priced: exactly two derivations, divisibility and closedness);
%  Lagrange used twice (finite exclusion; component-group triviality);
%  the (1+i)/2 lamp value NOT stated here -- "a superposed fixed point"
%  with the complex-probability treatment deferred to
%  sec:self-reference; the two-branch over-determination stated to
%  pre-empt a circularity reading; Thomson spelling throughout.
% =====================================================================

\section{Continuity Is Not an Axiom}
\label{sec:continuity}

Hardy derives quantum theory from five axioms \autocite{Hardy2001a}. Four of them are satisfied
by classical probability theory; the fifth --- that a \emph{continuous} reversible
transformation connects any two pure states --- is the entire difference, the single postulate
that turns the classical simplex into the complex Hilbert space. Every reconstruction programme
pays an equivalent toll: continuous reversibility in Masanes and M\"{u}ller
\autocite{mullerp2011}, the corresponding postulate in Daki\'{c} and Brukner
\autocite{Brukner2009}. The axiom is load-bearing, and in every programme it is unexplained: it
is where the quantumness is put in by hand. The claim of this section is that in the present
framework it is not an axiom. Continuity is the fifth face of the pre-measurement regime that self-reference forces, and what
stands between the framework and a full derivation is not a postulate about values but one
premise about transformations --- that every transformation the inside performs has a stationary
outside presentation --- derived in outline below from the inside/outside involution of
\S\ref{sec:self-reference} and stated as an open problem at the end of the section with nothing
hidden in it. The continuum itself, the completion that every probability theory adopts, is
named for what it is: the real line, shared with classical probability and carrying no quantum
content.

One structural remark prevents a misreading. The framework reaches complex amplitudes natively in
\S\ref{sec:self-reference}, by the negative-probability route; the present section reaches the same
destination by an independent branch --- the regime's continuity, then Hardy's theorem. The two
branches converging on one structure is over-determination, not circularity, and the convergence is
itself evidence that the destination is forced.

\subsection{Five faces of one regime}

The regime before any measurement --- before any \textsc{fanout} --- can be characterised five
equivalent ways: negative probability is available; anti-events exist, every event being undoable;
no classical record exists; no preferred basis is selected; and a continuous reversible
transformation connects any two pure states. These are not five premises to be combined. They are
one regime under five descriptions, and \textsc{fanout} (\S\ref{sec:fanout-boundary}) is the single
event that fails all five at once: it writes the record, selects the basis, kills the anti-events,
discretises the transformation group, and switches off the negativity. The equivalence of the first
two faces is the anti-event identity of \S\ref{sec:self-reference}; the third is the boundary's
definition. The work of this section is the equivalence of the last two faces, and then the route
from the regime's existence to its continuity. The fourth face has, since
Proposition~\ref{prop:f1-dictionary} (\S\ref{sec:hopf-justification} --- a deliberate forward
reference), an exact form: a basis-selection is the commitment of a classical structure, a lifted
antipodal skeleton, so ``no preferred basis'' means precisely that no copy structure has fired.

\subsection{Basis selection is the collapse of continuity}

\begin{proposition}[Continuity forbids a selected basis]\label{prop:basis-discreteness}
If a finite set of pure states is privileged --- a preferred basis --- then every reversible
transformation respecting the privilege carries the set to itself and acts on it as a permutation;
the respecting transformations therefore realise only a finite group of motions among the
distinguished states, and no continuous family of them connects distinct pure states of the set.
Contrapositively, continuous reversibility entails that no basis is selected.
\end{proposition}

\begin{proof}
A privileged finite set is by definition carried to itself by every structure-preserving reversible
map, whose induced action on the set is a permutation of finitely many points. A continuous path of
such maps induces a continuous path in a finite discrete set, which is constant. Hence a continuous
reversible family moving one distinguished pure state to another respects no such privilege, and if
every pair of pure states is so connected, no finite set is privileged.
\end{proof}

The converse is where the historical difficulty lies: the absence of a privileged finite set does
not by itself force continuity, for the transformation group could in principle be infinite and
discrete. Two structural properties close every such escape at once.

\begin{proposition}[Divisible and closed implies connected]\label{prop:divisible-closed}
Let $G$ be the group of reversible transformations of the pre-measurement regime, acting on the
state body. If $G$ is divisible --- every element has an $n$-th root for every $n$ --- and
topologically closed, then $G$ is a compact connected Lie group; and if $G$ moreover acts
transitively on the pure states, then a continuous reversible transformation connects any two of
them: Hardy's fifth axiom.
\end{proposition}

\begin{proof}
The state body is compact and convex with nonempty interior, so its symmetry group is compact, and
a closed subgroup of a compact matrix group is a compact Lie group with finitely many connected
components \autocite{broeckertomdieck1985compact}. The component group $G/G_{0}$ is a quotient of a
divisible group and hence divisible; but a nontrivial finite group is never divisible --- with
$\lvert H\rvert = m$ every element satisfies $h^{m} = e$, so no element other than the identity
possesses an $m$-th root --- and $G/G_{0}$ is finite, so it is trivial and $G = G_{0}$ is connected.
A connected Lie group is path-connected, and transitivity then supplies the continuous reversible
path between any two pure states.
\end{proof}

The same one-line fact of Lagrange closes the remaining escapes. Finite groups are excluded
directly, being nontrivial and never divisible. Countable groups are excluded because a
countable closed subgroup of a compact group is compact, has an isolated point by the Baire
category theorem, and is therefore discrete and finite. Uncountable disconnected groups are
excluded by the component-group argument of the proposition. Every route to a discrete
transformation group therefore passes through a failure of divisibility or a failure of
closedness. Divisibility, however, asks for more than the conclusion needs, and the next
proposition removes it.

\begin{proposition}[Continuity from a connected state space]\label{prop:connected-transitive}
Let the pure states of a finite-dimensional system form a connected space $P$, and let the group
$G$ of reversible transformations of the pre-measurement regime be topologically closed and act
transitively on $P$. Then the identity component $G_0$ already acts transitively on $P$, and a
continuous reversible transformation connects any two pure states: Hardy's fifth axiom holds,
with no divisibility assumed.
\end{proposition}

\begin{proof}
As in Proposition~\ref{prop:divisible-closed}, closedness makes $G$ a compact Lie group with
finitely many connected components \autocite{broeckertomdieck1985compact}, so $G_0$ is a closed
normal subgroup of finite index and is path-connected. Each $G_0$-orbit in $P$ is the continuous
image of a compact group, hence compact, hence closed. Since $G$ is transitive and $G_0$ is
normal, the $G_0$-orbit of $gx$ is $g$ applied to the $G_0$-orbit of $x$, so the cosets of $G_0$
permute the orbits transitively and there are finitely many of them. A finite partition of $P$
into closed sets is a partition into open sets; $P$ is connected; so there is one orbit and
$G_0$ is transitive. Given pure states $x$ and $y$, choose $g \in G_0$ with $gx = y$ and a
continuous path $g_t$ in $G_0$ from the identity to $g$; then $g_t x$ is a continuous reversible
transformation from $x$ to $y$.
\end{proof}

The proposition relocates the premise. Divisibility was asked for because a transformation group
could be infinite and discrete; but a group acting transitively on a continuum cannot be
countable, and a closed group acting transitively on a connected space cannot be disconnected in
the way that matters, so what the argument needs from the regime is not the availability of
$n$-th roots but the connectedness of the state space itself. That is a statement about values:
that the pure states are all the unit vectors of $\mathbb{C}^{K}$ up to phase, the space
$\mathbb{CP}^{K-1}$, connected for $K \ge 2$, and hence that the amplitudes of
\S\ref{sec:self-reference} fill their continuum. The two standing foils show the premise doing
its work. The stabiliser fragment of quantum theory has complex amplitudes and is contextual,
yet it has finitely many pure states, its outcome probabilities take only the values $0$,
$\tfrac12$ and $1$, and its reversible transformations form a finite group, the Clifford group,
so Hardy's fifth axiom fails in it; Spekkens's toy theory \autocite{spekkens2007toy} is the same
with real amplitudes, its ignorance rationed by a knowledge-balance principle. Neither has a
connected state space, and neither is a model of the regime described here, whose credences take
every value. Divisibility, which both also lack --- the phase gate has no square root in the Clifford group
--- is a symptom of the discreteness, not its cause. What both lack can be said exactly, and the
next subsection says it: each contains a transformation that fixes no state it admits.

\subsection{The premises, and what they mean}

Three premises carry Proposition~\ref{prop:connected-transitive}. Transitivity is the regime's
homogeneity: inside/outside equivalence leaves no pure state privileged, and a proper closed
invariant set of pure states would define a superselection observable, a classical record
predating any \textsc{fanout} (\S\ref{sec:fanout-boundary}). Connectedness of the state space is two statements, and only one of them is quantum. The first
is that the pure states are infinite and dense, with a fixed point for every transformation; the
second is that the state space is complete. Quantum theory over the algebraic numbers ---
amplitudes in the algebraic closure of $\mathbb{Q}$ within $\mathbb{C}$ --- has the first
without the second: it has $K = N^{2}$, the Born rule, superposition, entanglement and Bell
violation at the Tsirelson value, every unitary has an algebraic eigenvector, and it is
countable and totally disconnected. So the whole quantum content of Hardy's axiom lies in the
density, and the completion is the real-number convention of every probability theory. Hardy
adopts it before his fifth axiom is stated, in the first, which makes probabilities real numbers
and gives his classical theory a real simplex of mixtures; this framework adopts it at the
Shannon step, where credences are declared to range over a connected domain so that the
continuation of the logarithm is unique (\S\ref{sec:complex-amplitudes}). Given that convention,
density of the pure states is connectedness of their closure. Divisibility, which an earlier
version of this section asked for, is not needed: it follows from the conclusion, every element
of a compact connected Lie group having roots of every order.

Closedness is the completion, limit-realisability: a Cauchy sequence of available
transformations has its limit available. It cannot be derived, and the reason should be stated
rather than hidden. Any collection of first-order principles about states, transformations and
probabilities that has an infinite model has a countable one, and a countable subset of
$\mathbb{CP}^{K-1}$ is totally disconnected; the founding non-computability does not escape
this, since the arithmetical complex numbers form a countable algebraically closed field
containing non-computable numbers and closed under every limit it can itself describe. Hardy's
axiom and closedness are second-order, quantifying over paths and over sequences, and a
second-order principle is grounded as Dedekind completeness is grounded or not at all. The
grounds are these. The completion of a dense set under the finite-precision metric is unique and
depends on no choice, whereas any countable alternative is a choice of definability level, an
extra law with no ground and no test; and the framework's own result is that no observer's
descriptive level captures the regime, since by \S\ref{sec:self-reference} every observer's
state exceeds its own level and the inside/outside principle makes that a law rather than a
perspective. The burden therefore falls on any restriction, not on the completion, which is the
least committal option and carries no quantum content. What limit-realisability adds to the line
is still worth stating exactly, because it is the corrected form of a natural intuition.
Cantor's diagonal \autocite{cantor1891diagonal} does not build the continuum; its output is an
infinite decimal --- a Cauchy sequence of rationals --- which exists \emph{as a point} only
because the line is complete. The diagonal presupposes limit-realisability and then shows that
the completed line outruns every internal listing. And the character of the completion is the
framework's own: the computable reals are countable \autocite{Turing_1937}, hence of measure
zero, so once the transformation group contains its limits almost all of them are
non-computable. Non-computability is a property of the completed continuum, not the reason for
it. The constructive reading, on which only describable limits exist, remains consistent; it is
the countable theory above, and the two readings differ in nothing an experiment can certify.

Density is where the section's title is earned, and it is derived from the paper's own premise.
Write $(\mathrm{I})$ for the statement that every admissible transformation of the regime is a
phase of some admissible classical structure: diagonal, that is, in the record basis of some
observer the regime contains, a phased spider of that basis in the calculus of
\S\ref{sec:hopf-justification} \autocite{coeckeduncan2011interacting} (VERIFY-SUPPORT). It says
of transformations what the global phase of \S\ref{sec:self-reference} says of states: there are
no absolute ones. It follows from three things the paper already holds. The inside/outside
involution is a genuine involution, so the outside presentation of the inside is a presentation
of the same object with nothing lost: the outside can have complete knowledge of the inside, a
potential record, the joint state lying in a structure the outside carries. Time is operation
count, internal to the system that operates (\S\ref{sec:space-of-computations}), so the outside
does not share the inside's count. And before any \textsc{fanout} that count is not classically
recorded, so the outside's description is pure. Let $g$ be the transformation the inside
performs per operation. The outside's complete knowledge cannot depend on the inside's count,
and it is pure, so it is a pure state of the inside's clock and system that is stationary under
``advance the clock and apply $g$''; the stationary states of that map are clock Fourier modes
tensored with eigenvectors of $g$ and nothing else \autocite{pagewootters1983evolution}
(VERIFY-SUPPORT). So the outside holds a structure in which $g$ acts by phases, which is
$(\mathrm{I})$; and the relative phase in that structure is the outside's ignorance of the
inside's computational count, the same ignorance the global phase records one level down. The
consequences are a theorem. A phase of a structure fixes the structure's copyable states, so
every admissible transformation fixes an admissible pure state; a finite transitive action has
an element moving every point (Jordan), so the pure states are infinite; the transformation
carrying a stabiliser state to the fixed point of the composite of a $Z$-phase and an $X$-phase
turns through $\arccos(1/\sqrt{3})$, an irrational fraction of a turn, so the group has an
element of infinite projective order and its orbits are dense. With the completion,
Proposition~\ref{prop:connected-transitive} fires and Hardy's fifth axiom holds. The two foils
fail $(\mathrm{I})$ exactly where they fail continuity: fourteen of the twenty-four Clifford
operations are phases of no Pauli basis and fix no stabiliser state, their fixed points being
the magic states; the toy theory's three-cycles are phases of no epistemic basis; and in each
the outside's knowledge of the inside under the order-three protocol is either count-dependent
or a mixture over the count, so each makes the inside's time classical before any record exists.
The derivation of $(\mathrm{I})$ is an outline, not a proof, and it is what the open problem
below asks for.

\subsection{Why the cut matters physically}

The regime's continuity is not decorative. Reversibility alone, over complex amplitudes, licenses
the full invertible group $GL(n, \mathbb{C})$, and Aaronson showed what that buys: superluminal
signalling, and computational power at the level of \textsf{PP} \autocite{Aaronson2004h}. The cut
from $GL(n, \mathbb{C})$ down to the unitary group is the conserved quadratic form, and continuity
is what entitles the theory to make it: with the fifth axiom in hand, Hardy's theorem delivers the
complex Hilbert space, the quadratic norm and the unitary group in one stroke \autocite{Hardy2001a}
--- cited here, not re-proved. One honest caveat travels with the citation. Hardy's reconstruction
uses his remaining axioms besides --- composition and subspace structure --- and these must be
derived rather than imported: the composition ledger of \S\ref{sec:division-algebras} and the
locality-of-tomography argument of \S\ref{sec:self-reference} carry that duty, and it is discharged
there, not silently assumed here.

\subsection{Super-tasks at the boundary}

The premise just priced has a fifty-year-old stress test in the philosophy of mathematics, and the
framework passes it in an unexpected way. A super-task --- infinitely many operations completed in
finite time \autocite{Thomson_1954} --- presupposes its terminal moment as an actual point of the
timeline: the accumulation point of a generic sequence of step-times exists only in the completed
line, and the defining question, the state \emph{after all the steps}, is meaningful only under
limit-realisability of time. Super-tasks are downstream consumers of exactly the completeness at
issue in this section. Thomson's lamp then exhibits both sides of the \textsc{fanout} boundary at
once. On the record side, the lamp's description fixes its state at every finite stage and thereby
fixes no limit at all: any terminal value is consistent with the whole task
\autocite{Thomson_1954,benacerraf1962tasks}. That is not a puzzle about lamps; it is the proof that
the classical record algebra is not limit-closed, and cannot be. On the amplitude side the limit
exists --- a superposed fixed point, whose complex-probability treatment belongs to
\S\ref{sec:self-reference} --- because that side \emph{is} limit-closed, which is the fifth face
itself. The boundary is thereby the edge of limit-realisability: complete on the reversible side,
incomplete on the record side, with the lamp standing on both sides at once as the boundary's oldest
exhibit. A final clause dissolves rather than answers the classic question. In this framework time
is operation count (\S\ref{sec:space-of-computations}), and ``after infinitely many operations''
indexes no operational moment: the lamp packs its endless \textsc{fanout}s into the emergent
continuum, and asking for its record at the accumulation point asks the discrete contents for a
value that only the completed container possesses. On or off is not undecided; it is unasked.

\subsection{What remains, priced}

\begin{openproblem}[Every transformation a phase of an inside structure]\label{op:limit-realisability}
Proposition~\ref{prop:connected-transitive} reduces Hardy's fifth axiom to three premises:
connectedness of the pure states, transitivity, and closedness. Transitivity is argued above
from the absence of records. Closedness is the completion, adopted on the grounds stated above
and carrying no quantum content. Connectedness is density plus that completion, and density
follows, by Jordan's theorem and the infinite-order argument, from $(\mathrm{I})$: that every
admissible transformation of the regime is a phase of some admissible classical structure. What
remains open is the rigour of the derivation of $(\mathrm{I})$ from the paper's premise. Derive,
inside the fixed-point locus $\mathrm{Fix}(\sigma)$ of Proposition~\ref{prop:inside-outside},
that every endomap is a phase of a classical structure in $\mathrm{Fix}(\sigma)$, from the
self-representing object, the involution $\sigma$ as a genuine involution, and the reading of
the inside presentation as a sequence indexed by operation count with no record before
\textsc{fanout}. The outline above does this with the stationary-state construction of Page and
Wootters; a derivation passes if it makes that outline a theorem of the calculus, with the
outside's count-independent complete description shown to exist and to be pure from the
involution alone. It fails before any detail is checked if every premise it uses holds in the
stabiliser fragment of quantum theory, in Spekkens's toy theory \autocite{spekkens2007toy}, in
the real qubit, or in quantum theory over the Gaussian rationals, all of which violate
$(\mathrm{I})$ outright and all of which make the inside's operation count classical before any
record exists. It also fails if it imports the completion as a premise about transformations:
the completion is the real line and has no derivation, and a derivation of $(\mathrm{I})$ must
hold in the countable algebraic theory, where $(\mathrm{I})$ is true and continuity is not.
\end{openproblem}

The accounting, in one sentence: the axiom that every reconstruction programme pays is, in this
framework, one derivation posed, of Lawvere's shape --- that every endomap of the fixed-point
locus is a phase of a classical structure in it --- and one convention named, the completion,
which classical probability pays too. And
the hand-on is already visible: the complex amplitudes that this section's continuity underwrites
through Hardy's theorem, and the complex amplitudes that \S\ref{sec:self-reference} derives natively,
are one and the same $\mathbb{C}$ --- the floor of the ladder that \S\ref{sec:division-algebras}
climbs.

%      % =====================================================================
%  SECTION: The Division-Algebra Ladder
%  Drafted 2026-07-03. One of the three foundations sections
%  (sec:self-reference, sec:division-algebras, sec:fanout-boundary);
%  this one resolves the dangling \S\ref{sec:division-algebras} in the
%  Hopf-justification section, whose header names it "the
%  Cayley--Dickson / Hurwitz step".
%
%  PLACEMENT in main.tex (REORDERED per JS 2026-07-06): position 6,
%  opening Act II ("why stop at C?") right before its first consumer --
%      \input{self_reference_section_2026-07-04}
%      \input{continuity_section_2026-07-04}
%      \input{L2_metric_global_phase_section_2026_06_07}
%      \input{wigners_friend_in_the_hopf_picture_2026-06-13}
%      \input{derivation_of_the_quantum_postulates}
%      \input{division_algebras_section_2026-07-03}   <-- moved here
%      \input{spacetime_from_non-computability}       <-- first consumer
%                                                          (S^7 exclusion)
%      \input{hopf_justification_section_2026-06-08}
%  The section is position-portable: all adjacency phrasing was
%  retargeted to explicit \S\ref on 2026-07-06. The S^7 hand-off
%  sentence to sec:space-of-computations WRITTEN 2026-07-06 (spacetime
%  file read; the closing now hands first to its immediate consumer,
%  then to hopf-justification). Upload verifications COMPLETE: L2 and
%  quantum-postulates do not consume the ladder; WF's references are
%  forward-voiced or self-carrying.
%
%  DIVISION OF LABOUR (deliberate, to avoid overlap):
%   - sec:self-reference (undrafted): self-reference -> complex
%     amplitudes and the Born rule; why not a real base. This section
%     TAKES that as given and starts at C.
%   - THIS section: the algebra -- the doubling construction, the
%     ledger of losses-as-physics, the internal-algebras-not-a-new-QM
%     scoping (Adler pre-empt), the two appearances of the ladder
%     (register / measurement climb, stated lightly), and Hurwitz as
%     the algebraic termination with the committed zero-divisor
%     reading.
%   - sec:hopf-justification: the geometry (Adams, parallelisability,
%     the fibrations), the join of the two classifications, the
%     entanglement-sensitivity literature, the C-direction honest debt,
%     and the epistemic reading of termination.
%
%  COMMITTED READINGS USED (sources re-read 2026-07-03):
%   - Loss ledger and level dictionary: cayley_dickson_qubits_for_review
%     sec.5 ("each step loses a symmetry, and each lost symmetry is
%     identified with a physical content").
%   - Termination reading: grounding_argument_gauge_group ("the
%     sedenions have zero divisors and cannot support consistent
%     information flow around loops") -- leads; norm-composition-as-
%     probability given as its amplitude-level restatement.
%   - Internal-not-base scoping: cayley_dickson_qubits_for_review Q2 and
%     sec.3.2 caveat (Adler objection; C-qubit treated geometrically).
%
%  HONEST DEBT: ladder proved closed at the top (Hurwitz), derived at
%  the floor. The rung-compulsion revisit-trigger FIRED 2026-07-06:
%  sec:self-reference and op:sigma-functor now exist, so the thesis is
%  minted as op:ladder-compulsion (with the CD-conjugation-as-sigma
%  candidate route), inheriting op:sigma-functor's status.
%
%  AUDITED 2026-07-06 (final pass, third stop): all four citations
%  used; all cross-refs live; CD facts, dims, alternativity, Hurwitz
%  statement verified; Adler subsection delivers everything the
%  register's compressed pointer promises; file pure ASCII; no bare
%  "frame". Fixes: opening rewritten for both predecessors (continuity
%  was inserted between after drafting); "associator's non-vanishing is
%  mass itself" recalibrated to the register's debt-against-selection
%  defusal; the O->S ledger row and Hurwitz subsection reconciled with
%  the fourth-qubit forward remark ("the terminus closes the ladder of
%  forces; what survives past it is, at most, flavour").
%
%  CROSS-REFS: sec:self-reference [sibling dangle -- resolves when that
%  section is drafted], sec:hopf-justification, sec:hopf-information-bound,
%  sec:register, sec:entanglement-classes, sec:generations, sec:masses,
%  sec:wigners-friend. All others exist.
%  CITATIONS: existing keys baez2002octonions, hurwitz1898composition,
%  furey2018three. ONE new key: adler1995quaternionic (RIS in
%  division_algebras_refs_2026-07-03.ris).
%  British English; \autocite; no bullets in prose; the ledger is a
%  table because the comparison is the argument.
% =====================================================================

\section{The Division-Algebra Ladder}
\label{sec:division-algebras}

The foundations behind us end at amplitudes twice over: natively, where the self-referential
argument forced the description of a system into the system's own state space and priced
consistency at a number system whose parts can cancel --- complex amplitudes, and the Born rule
that squares them (\S\ref{sec:self-reference}) --- and independently, where the regime's
continuity delivers the same field through Hardy's theorem (\S\ref{sec:continuity}), the
convergence of the two branches being evidence rather than accident. This section asks the question
both routes leave open. Why
should the number system stop at $\mathbb{C}$? Ordinary quantum mechanics stops there by postulate;
a framework that has forbidden itself postulates must answer, and the answer is that the climb
neither stops at $\mathbb{C}$ nor continues forever. There are exactly two more rungs, each
purchased by surrendering an algebraic law and each paying for itself in physics, and the ladder is
closed at the top by a theorem of 1898. This section sets out the algebra of that ladder: the single
construction that climbs it, the ledger of what each rung costs and buys, and the theorem that ends
it. The geometry of the same ladder --- that each rung acts on a sphere, that exactly those spheres
are parallelisable, and that the algebraic and geometric classifications are one --- is the business
of the geometry section (\S\ref{sec:hopf-justification}).

\subsection{One construction climbs the whole ladder}

The rungs are not a miscellany; a single construction generates them all. Given an algebra with a
conjugation, the Cayley--Dickson construction builds its double on ordered pairs, with a twisted
multiplication in which the conjugate enters and a conjugation $(a,b)^{*} = (a^{*}, -b)$
\autocite{baez2002octonions}. (Conventions differ in the placement of the conjugations; all yield
isomorphic algebras.) Applied repeatedly it produces the sequence
$\mathbb{R} \to \mathbb{C} \to \mathbb{H} \to \mathbb{O} \to \mathbb{S}$, of real dimensions
$1, 2, 4, 8, 16$: the reals, the complex numbers, the quaternions, the octonions, the sedenions.
The ground floor is classical: $\mathbb{R}$ is ordered and positive-definite, the arithmetic of
classical probability. The first doubling is therefore the quantum move itself --- the passage from
probabilities to amplitudes that can interfere --- and it is not made here by choice: it is the
derivation of \S\ref{sec:self-reference}, already behind us. What this section must account for is
the rest of the climb.

\subsection{The ledger}

In mathematics the doublings are usually told as a story of degeneration: each algebra is worse
behaved than the last, and the sequence is studied for where it breaks. This framework reads the
same ledger in the opposite direction. Each rung is bought by surrendering a law, and each
surrendered law is exactly a piece of physics --- the losses are the assets, and the comparison is
sharp enough to be the argument:

\begin{center}
\begin{tabular}{llll}
\hline
doubling & dimension & what is lost & what the loss becomes \\
\hline
$\mathbb{R} \to \mathbb{C}$ & $1 \to 2$ & the order of the reals & phase and interference; amplitudes \\
 & & & that cancel; the Born rule (\S\ref{sec:self-reference}) \\
$\mathbb{C} \to \mathbb{H}$ & $2 \to 4$ & commutativity & non-commuting actions: weak isospin \\
 & & & $SU(2)$, the $B$ register (\S\ref{sec:register}) \\
$\mathbb{H} \to \mathbb{O}$ & $4 \to 8$ & associativity & the associator: colour, chirality, the \\
 & & & generations, and mass as debt \\
 & & & (\S\ref{sec:generations}, \S\ref{sec:masses}) \\
$\mathbb{O} \to \mathbb{S}$ & $8 \to 16$ & division itself & no further force: zero divisors appear \\
 & & (zero divisors) & and the ladder of sectors ends (below) \\
\hline
\end{tabular}
\end{center}

Walking the rows: the order lost at the first doubling is the possibility of interference --- parts
of a probability that can oppose one another --- which is where the derivation of \S\ref{sec:self-reference}
lives. The commutativity lost at the second is the existence of actions whose order matters; the
unit quaternions form the group $SU(2)$, and the framework houses weak isospin there, as the
internal algebra of the $B$ register whose order-sensitive, non-commuting dynamics the interaction sections develop
(\S\ref{sec:register}, \S\ref{sec:entanglement-classes}). The associativity lost at the third ---
the octonions retain the weaker property of alternativity \autocite{baez2002octonions} --- creates
the associator $[a,b,c] = (ab)c - a(bc)$, and the associator is where most of this paper's physics
is transacted: the stabiliser of a preferred complex direction inside the octonions' automorphisms
is the colour group (\S\ref{sec:hopf-justification}); the triality structure and its $K_3$
obstruction generate the three generations and the sterile channel (\S\ref{sec:generations}); and
the associator is the currency of mass --- the debt against the selected direction that
\S\ref{sec:masses} prices, owed by whatever action moves the selection and not by residence at
the rung, which is the defusal \S\ref{sec:register} makes explicit. The algebraic facts
in the third column are theorems; the identifications in the fourth are established in the sections
cited, each carrying its own epistemic flag there, and are not re-litigated here. What this table
asserts on its own account is only the correspondence's shape: one law surrendered per rung, one
sector of physics gained per law.

\subsection{Internal algebras, not a new quantum mechanics}

A reader versed in the quaternionic-quantum-mechanics literature will at this point raise a fair
objection: quantum theory over $\mathbb{H}$ as a foundation is known to be troubled --- the
tensor-product and composite-system difficulties set out by Adler and others
\autocite{adler1995quaternionic} --- and quantum theory over $\mathbb{O}$ founders earlier still,
since a non-associative algebra does not furnish a Hilbert space at all. The framework's position
must therefore be stated exactly, because it does not make the claim the objection targets. The base
of the physical Hilbert space here is $\mathbb{C}$ and remains $\mathbb{C}$; the higher algebras
enter as the \emph{internal} algebras of the register --- $\mathbb{H}$ for the $B$ qubit,
$\mathbb{O}$ for the $C$ qubit (\S\ref{sec:register}) --- in a tensor structure whose outer theory is
ordinary complex quantum mechanics. This is the same move by which the division algebras enter the
constructions of Furey and the related programmes without any proposal to rebase quantum mechanics
itself \autocite{furey2018three}. For the octonionic register the point is sharper still:
$\mathbb{O}^{2}$ is not a Hilbert space, and the framework accordingly treats the $C$ qubit not in
the Dirac--von Neumann manner but geometrically, as the sphere $S^{15}$ with its Hopf structure ---
which is precisely the form in which \S\ref{sec:hopf-justification} needs it. Nothing in this paper therefore
stands or falls with quaternionic or octonionic quantum mechanics as foundations; the ladder is
climbed by the registers' internal algebras, not by the theory's base field.

\subsection{The same ladder, twice}

The ladder appears in this framework in two guises, and it will help the reader to have them
distinguished now. Statically, it is the anatomy of a single relation: the three qubits of the
register take their values one rung apart --- $A$ over $\mathbb{C}$, $B$ over $\mathbb{H}$, $C$ over
$\mathbb{O}$ --- so that every particle carries the whole ladder as its internal structure
(\S\ref{sec:register}). Dynamically, it is the shape of measurement: a single qubit's pure states
fill $S^{3}$, an entangling measurement carries the joint state into $S^{7}$, and a third party into
$S^{15}$ --- the correspondence between qubit entanglement and the Hopf fibrations that
\S\ref{sec:hopf-justification} takes from a mature literature --- so that each act of measurement ascends one rung --- the climb the reader has already made,
narrated in \S\ref{sec:wigners-friend} --- and every completed relation is a terminated climb, which is
the reading \S\ref{sec:hopf-information-bound} gives the tower's endpoint. One ladder, seen twice:
as the anatomy of a relation and as the dynamics of relating. The framework regards this double
appearance not as a coincidence but as the same structure viewed statically and in motion, and the
sections cited carry that identification's weight.

\subsection{Hurwitz closes the ladder}

The construction itself would double forever; what stops it is a classification theorem. By
Hurwitz's theorem, the normed division algebras --- those in which the norm composes,
$\lvert ab\rvert = \lvert a\rvert\,\lvert b\rvert$, and no non-zero elements multiply to zero ---
are exactly $\mathbb{R}$, $\mathbb{C}$, $\mathbb{H}$ and $\mathbb{O}$
\autocite{hurwitz1898composition,baez2002octonions}. The next doubling, the sedenions, has zero
divisors: non-zero elements whose product vanishes. The framework's reading of this failure is the
one that matters physically. Zero divisors cannot support consistent information flow around loops
--- a cycle of compositions can silently annihilate what it carries --- and a framework in which
closed computational paths are load-bearing (\S\ref{sec:hopf-justification},
\S\ref{sec:space-of-computations}) cannot be built on an algebra that permits it. At the level of
amplitudes the same failure reads: the composition law
$\lvert ab\rvert = \lvert a\rvert\,\lvert b\rvert$, which carries probability multiplicatively
through every rung of the ladder, fails beyond the octonions, so that two non-vanishing amplitudes
could compose to nothing and probability would be destroyed by mere composition. The last algebra
fit to carry amplitudes is $\mathbb{O}$. One rung past the terminus is not entirely silent,
however: the sedenions' internal structure --- their automorphism group and their octonion
subalgebras --- is the recorded candidate arena for tagging the generations
(\S\ref{sec:interactions}), and it enters exactly as this section's scoping permits, internally
and never as a base. The terminus closes the ladder of forces; what survives past it is, at most,
flavour. This is the algebraic half of the termination; the
geometric half --- that no sphere above $S^{7}$ is parallelisable --- and the theorem joining the
two halves are the opening business of \S\ref{sec:hopf-justification}, and the epistemic reading of the same endpoint,
as a bound on the information a single relation can carry, is given in
\S\ref{sec:hopf-information-bound}.

\subsection{The honest debt}

As in the geometry section, we close by stating plainly what is proved, what is established, and what is
owed. Proved: the construction, the losses, and the termination are theorems --- the Cayley--Dickson
doubling and its ledger of surrendered laws are standard algebra, and Hurwitz's classification is
among the durable results of the subject. Established: the physical cashing of each rung --- the
Born rule at $\mathbb{C}$, weak isospin at $\mathbb{H}$, colour, generations and mass at
$\mathbb{O}$ --- is carried out in the sections cited above, each under its own epistemic flag. Owed,
and stated without disguise: the ladder as presented here is closed from above and derived at its
floor, but it is not yet derived rung by rung from below. Hurwitz forbids a further rung; the
argument of \S\ref{sec:self-reference} forces the first doubling, from the ordered arithmetic of
classical probability to amplitudes that interfere. That the same self-referential closure, applied
to a system whose amplitudes already take values at one rung, compels the next --- that relational
closure at one level forces precisely the Cayley--Dickson double of that level --- is the
framework's reading of the climb, and it is a thesis, not a theorem --- but the formal apparatus
against which it can be posed now exists, and we pose it.

\begin{openproblem}[The compulsion of the intermediate rungs]\label{op:ladder-compulsion}
Derive, within the categorical apparatus of Open Problem~\ref{op:sigma-functor}, that
self-referential closure at one rung compels the Cayley--Dickson double of that rung, so that the
intermediate steps $\mathbb{C} \to \mathbb{H} \to \mathbb{O}$ follow from the same principle
that forces the floor. The candidate route is the construction's own shape: the Cayley--Dickson
pair with its conjugation, $(a, b)^{*} = (a^{*}, -b)$, is an involution-equipped doubling, and the
route reads it as the algebraic shadow of adjoining the describer --- closure of level-$A$
descriptions under the inside/outside exchange requiring exactly the twisted pair. The problem
inherits the status of Open Problem~\ref{op:sigma-functor}, whose apparatus it presupposes.
\end{openproblem}

If the derivation lands, the whole ladder becomes a consequence of the same self-reference that
opened it. Until then, the honest statement is that the ladder's necessity is established at its
ceiling and at its floor, while the compulsion of the intermediate rungs is inherited from the
physics they are shown to carry rather than derived in advance. What this section hands on is, accordingly, an algebraic ladder with a floor, a ceiling and a
ledger. Its first consumer is immediate: \S\ref{sec:space-of-computations} spends the terminal
fact --- the non-associativity of $S^{7}$'s own multiplication --- as geometry, excluding the
seven-sphere as a ruler and fixing the dimension of space at three.
\S\ref{sec:hopf-justification} then supplies the geometry of the fibrations themselves, and the
join.

%  (Architecture call SETTLED by the sync manifest: standalone.)
%
%  CITATIONS -- existing keys (as supplied by JS 2026-07-04):
%    Hardy2001a, Aaronson2004h, Thomson_1954, Lawvere_1969,
%    Turing_1937, muller_p_2011, Brukner2009.
%  NEW keys (RIS in continuity_refs_2026-07-04.ris):
%    cantor1891diagonal, benacerraf1962tasks, broeckertomdieck1985compact.
%  Deliberately excluded per JS: Chiribella--D'Ariano--Perinotti;
%  Bargmann/Wigner's theorem. Yanofsky2003 now paired on the
%  Lawvere-shaped clause (the pending-key excuse was void: the key
%  exists and sec:self-reference uses it).
%
%  CROSS-REFS: sec:self-reference now DRAFTED (2026-07-04, later the
%  same day) -- no longer dangling; only sec:fanout-boundary (and its
%  sibling sec:fanout-direction) -- DRAFTED 2026-07-06.
%
%  AUDITED 2026-07-06 (final pass, second stop): both proofs
%  re-verified by hand (Lagrange divisibility sentence exact; the
%  Baire countable-escape chain sound); caveat-discharge with
%  sec:self-reference confirmed two-way (already present); file pure
%  ASCII; no pre-definition "frame"; all citations used as declared.
%  Fixes: Yanofsky2003 paired on the Lawvere clause; the op's two
%  candidate routes anchored to printed propositions (f1-dictionary
%  for the divisibility premise's referent, inside-outside for the
%  closedness demand) -- statuses unchanged, referents sharpened;
%  header staleness corrected. Deliberate FORWARD references --
%  prop:f1-dictionary and sec:hopf-justification (defined later in the
%  paper; fine under two-pass LaTeX and flagged in prose),
%  sec:division-algebras, sec:space-of-computations.
%
%  DECISIONS ENCODED: op:limit-realisability minted (the owed arrow,
%  priced: exactly two derivations, divisibility and closedness);
%  Lagrange used twice (finite exclusion; component-group triviality);
%  the (1+i)/2 lamp value NOT stated here -- "a superposed fixed point"
%  with the complex-probability treatment deferred to
%  sec:self-reference; the two-branch over-determination stated to
%  pre-empt a circularity reading; Thomson spelling throughout.
% =====================================================================

\section{Continuity Is Not an Axiom}
\label{sec:continuity}

Hardy derives quantum theory from five axioms \autocite{Hardy2001a}. Four of them are satisfied
by classical probability theory; the fifth --- that a \emph{continuous} reversible
transformation connects any two pure states --- is the entire difference, the single postulate
that turns the classical simplex into the complex Hilbert space. Every reconstruction programme
pays an equivalent toll: continuous reversibility in Masanes and M\"{u}ller
\autocite{mullerp2011}, the corresponding postulate in Daki\'{c} and Brukner
\autocite{Brukner2009}. The axiom is load-bearing, and in every programme it is unexplained: it
is where the quantumness is put in by hand. The claim of this section is that in the present
framework it is not an axiom. Continuity is the fifth face of the pre-measurement regime that self-reference forces, and what
stands between the framework and a full derivation is not a postulate about values but one
premise about transformations --- that every transformation the inside performs has a stationary
outside presentation --- derived in outline below from the inside/outside involution of
\S\ref{sec:self-reference} and stated as an open problem at the end of the section with nothing
hidden in it. The continuum itself, the completion that every probability theory adopts, is
named for what it is: the real line, shared with classical probability and carrying no quantum
content.

One structural remark prevents a misreading. The framework reaches complex amplitudes natively in
\S\ref{sec:self-reference}, by the negative-probability route; the present section reaches the same
destination by an independent branch --- the regime's continuity, then Hardy's theorem. The two
branches converging on one structure is over-determination, not circularity, and the convergence is
itself evidence that the destination is forced.

\subsection{Five faces of one regime}

The regime before any measurement --- before any \textsc{fanout} --- can be characterised five
equivalent ways: negative probability is available; anti-events exist, every event being undoable;
no classical record exists; no preferred basis is selected; and a continuous reversible
transformation connects any two pure states. These are not five premises to be combined. They are
one regime under five descriptions, and \textsc{fanout} (\S\ref{sec:fanout-boundary}) is the single
event that fails all five at once: it writes the record, selects the basis, kills the anti-events,
discretises the transformation group, and switches off the negativity. The equivalence of the first
two faces is the anti-event identity of \S\ref{sec:self-reference}; the third is the boundary's
definition. The work of this section is the equivalence of the last two faces, and then the route
from the regime's existence to its continuity. The fourth face has, since
Proposition~\ref{prop:f1-dictionary} (\S\ref{sec:hopf-justification} --- a deliberate forward
reference), an exact form: a basis-selection is the commitment of a classical structure, a lifted
antipodal skeleton, so ``no preferred basis'' means precisely that no copy structure has fired.

\subsection{Basis selection is the collapse of continuity}

\begin{proposition}[Continuity forbids a selected basis]\label{prop:basis-discreteness}
If a finite set of pure states is privileged --- a preferred basis --- then every reversible
transformation respecting the privilege carries the set to itself and acts on it as a permutation;
the respecting transformations therefore realise only a finite group of motions among the
distinguished states, and no continuous family of them connects distinct pure states of the set.
Contrapositively, continuous reversibility entails that no basis is selected.
\end{proposition}

\begin{proof}
A privileged finite set is by definition carried to itself by every structure-preserving reversible
map, whose induced action on the set is a permutation of finitely many points. A continuous path of
such maps induces a continuous path in a finite discrete set, which is constant. Hence a continuous
reversible family moving one distinguished pure state to another respects no such privilege, and if
every pair of pure states is so connected, no finite set is privileged.
\end{proof}

The converse is where the historical difficulty lies: the absence of a privileged finite set does
not by itself force continuity, for the transformation group could in principle be infinite and
discrete. Two structural properties close every such escape at once.

\begin{proposition}[Divisible and closed implies connected]\label{prop:divisible-closed}
Let $G$ be the group of reversible transformations of the pre-measurement regime, acting on the
state body. If $G$ is divisible --- every element has an $n$-th root for every $n$ --- and
topologically closed, then $G$ is a compact connected Lie group; and if $G$ moreover acts
transitively on the pure states, then a continuous reversible transformation connects any two of
them: Hardy's fifth axiom.
\end{proposition}

\begin{proof}
The state body is compact and convex with nonempty interior, so its symmetry group is compact, and
a closed subgroup of a compact matrix group is a compact Lie group with finitely many connected
components \autocite{broeckertomdieck1985compact}. The component group $G/G_{0}$ is a quotient of a
divisible group and hence divisible; but a nontrivial finite group is never divisible --- with
$\lvert H\rvert = m$ every element satisfies $h^{m} = e$, so no element other than the identity
possesses an $m$-th root --- and $G/G_{0}$ is finite, so it is trivial and $G = G_{0}$ is connected.
A connected Lie group is path-connected, and transitivity then supplies the continuous reversible
path between any two pure states.
\end{proof}

The same one-line fact of Lagrange closes the remaining escapes. Finite groups are excluded
directly, being nontrivial and never divisible. Countable groups are excluded because a
countable closed subgroup of a compact group is compact, has an isolated point by the Baire
category theorem, and is therefore discrete and finite. Uncountable disconnected groups are
excluded by the component-group argument of the proposition. Every route to a discrete
transformation group therefore passes through a failure of divisibility or a failure of
closedness. Divisibility, however, asks for more than the conclusion needs, and the next
proposition removes it.

\begin{proposition}[Continuity from a connected state space]\label{prop:connected-transitive}
Let the pure states of a finite-dimensional system form a connected space $P$, and let the group
$G$ of reversible transformations of the pre-measurement regime be topologically closed and act
transitively on $P$. Then the identity component $G_0$ already acts transitively on $P$, and a
continuous reversible transformation connects any two pure states: Hardy's fifth axiom holds,
with no divisibility assumed.
\end{proposition}

\begin{proof}
As in Proposition~\ref{prop:divisible-closed}, closedness makes $G$ a compact Lie group with
finitely many connected components \autocite{broeckertomdieck1985compact}, so $G_0$ is a closed
normal subgroup of finite index and is path-connected. Each $G_0$-orbit in $P$ is the continuous
image of a compact group, hence compact, hence closed. Since $G$ is transitive and $G_0$ is
normal, the $G_0$-orbit of $gx$ is $g$ applied to the $G_0$-orbit of $x$, so the cosets of $G_0$
permute the orbits transitively and there are finitely many of them. A finite partition of $P$
into closed sets is a partition into open sets; $P$ is connected; so there is one orbit and
$G_0$ is transitive. Given pure states $x$ and $y$, choose $g \in G_0$ with $gx = y$ and a
continuous path $g_t$ in $G_0$ from the identity to $g$; then $g_t x$ is a continuous reversible
transformation from $x$ to $y$.
\end{proof}

The proposition relocates the premise. Divisibility was asked for because a transformation group
could be infinite and discrete; but a group acting transitively on a continuum cannot be
countable, and a closed group acting transitively on a connected space cannot be disconnected in
the way that matters, so what the argument needs from the regime is not the availability of
$n$-th roots but the connectedness of the state space itself. That is a statement about values:
that the pure states are all the unit vectors of $\mathbb{C}^{K}$ up to phase, the space
$\mathbb{CP}^{K-1}$, connected for $K \ge 2$, and hence that the amplitudes of
\S\ref{sec:self-reference} fill their continuum. The two standing foils show the premise doing
its work. The stabiliser fragment of quantum theory has complex amplitudes and is contextual,
yet it has finitely many pure states, its outcome probabilities take only the values $0$,
$\tfrac12$ and $1$, and its reversible transformations form a finite group, the Clifford group,
so Hardy's fifth axiom fails in it; Spekkens's toy theory \autocite{spekkens2007toy} is the same
with real amplitudes, its ignorance rationed by a knowledge-balance principle. Neither has a
connected state space, and neither is a model of the regime described here, whose credences take
every value. Divisibility, which both also lack --- the phase gate has no square root in the Clifford group
--- is a symptom of the discreteness, not its cause. What both lack can be said exactly, and the
next subsection says it: each contains a transformation that fixes no state it admits.

\subsection{The premises, and what they mean}

Three premises carry Proposition~\ref{prop:connected-transitive}. Transitivity is the regime's
homogeneity: inside/outside equivalence leaves no pure state privileged, and a proper closed
invariant set of pure states would define a superselection observable, a classical record
predating any \textsc{fanout} (\S\ref{sec:fanout-boundary}). Connectedness of the state space is two statements, and only one of them is quantum. The first
is that the pure states are infinite and dense, with a fixed point for every transformation; the
second is that the state space is complete. Quantum theory over the algebraic numbers ---
amplitudes in the algebraic closure of $\mathbb{Q}$ within $\mathbb{C}$ --- has the first
without the second: it has $K = N^{2}$, the Born rule, superposition, entanglement and Bell
violation at the Tsirelson value, every unitary has an algebraic eigenvector, and it is
countable and totally disconnected. So the whole quantum content of Hardy's axiom lies in the
density, and the completion is the real-number convention of every probability theory. Hardy
adopts it before his fifth axiom is stated, in the first, which makes probabilities real numbers
and gives his classical theory a real simplex of mixtures; this framework adopts it at the
Shannon step, where credences are declared to range over a connected domain so that the
continuation of the logarithm is unique (\S\ref{sec:complex-amplitudes}). Given that convention,
density of the pure states is connectedness of their closure. Divisibility, which an earlier
version of this section asked for, is not needed: it follows from the conclusion, every element
of a compact connected Lie group having roots of every order.

Closedness is the completion, limit-realisability: a Cauchy sequence of available
transformations has its limit available. It cannot be derived, and the reason should be stated
rather than hidden. Any collection of first-order principles about states, transformations and
probabilities that has an infinite model has a countable one, and a countable subset of
$\mathbb{CP}^{K-1}$ is totally disconnected; the founding non-computability does not escape
this, since the arithmetical complex numbers form a countable algebraically closed field
containing non-computable numbers and closed under every limit it can itself describe. Hardy's
axiom and closedness are second-order, quantifying over paths and over sequences, and a
second-order principle is grounded as Dedekind completeness is grounded or not at all. The
grounds are these. The completion of a dense set under the finite-precision metric is unique and
depends on no choice, whereas any countable alternative is a choice of definability level, an
extra law with no ground and no test; and the framework's own result is that no observer's
descriptive level captures the regime, since by \S\ref{sec:self-reference} every observer's
state exceeds its own level and the inside/outside principle makes that a law rather than a
perspective. The burden therefore falls on any restriction, not on the completion, which is the
least committal option and carries no quantum content. What limit-realisability adds to the line
is still worth stating exactly, because it is the corrected form of a natural intuition.
Cantor's diagonal \autocite{cantor1891diagonal} does not build the continuum; its output is an
infinite decimal --- a Cauchy sequence of rationals --- which exists \emph{as a point} only
because the line is complete. The diagonal presupposes limit-realisability and then shows that
the completed line outruns every internal listing. And the character of the completion is the
framework's own: the computable reals are countable \autocite{Turing_1937}, hence of measure
zero, so once the transformation group contains its limits almost all of them are
non-computable. Non-computability is a property of the completed continuum, not the reason for
it. The constructive reading, on which only describable limits exist, remains consistent; it is
the countable theory above, and the two readings differ in nothing an experiment can certify.

Density is where the section's title is earned, and it is derived from the paper's own premise.
Write $(\mathrm{I})$ for the statement that every admissible transformation of the regime is a
phase of some admissible classical structure: diagonal, that is, in the record basis of some
observer the regime contains, a phased spider of that basis in the calculus of
\S\ref{sec:hopf-justification} \autocite{coeckeduncan2011interacting} (VERIFY-SUPPORT). It says
of transformations what the global phase of \S\ref{sec:self-reference} says of states: there are
no absolute ones. It follows from three things the paper already holds. The inside/outside
involution is a genuine involution, so the outside presentation of the inside is a presentation
of the same object with nothing lost: the outside can have complete knowledge of the inside, a
potential record, the joint state lying in a structure the outside carries. Time is operation
count, internal to the system that operates (\S\ref{sec:space-of-computations}), so the outside
does not share the inside's count. And before any \textsc{fanout} that count is not classically
recorded, so the outside's description is pure. Let $g$ be the transformation the inside
performs per operation. The outside's complete knowledge cannot depend on the inside's count,
and it is pure, so it is a pure state of the inside's clock and system that is stationary under
``advance the clock and apply $g$''; the stationary states of that map are clock Fourier modes
tensored with eigenvectors of $g$ and nothing else \autocite{pagewootters1983evolution}
(VERIFY-SUPPORT). So the outside holds a structure in which $g$ acts by phases, which is
$(\mathrm{I})$; and the relative phase in that structure is the outside's ignorance of the
inside's computational count, the same ignorance the global phase records one level down. The
consequences are a theorem. A phase of a structure fixes the structure's copyable states, so
every admissible transformation fixes an admissible pure state; a finite transitive action has
an element moving every point (Jordan), so the pure states are infinite; the transformation
carrying a stabiliser state to the fixed point of the composite of a $Z$-phase and an $X$-phase
turns through $\arccos(1/\sqrt{3})$, an irrational fraction of a turn, so the group has an
element of infinite projective order and its orbits are dense. With the completion,
Proposition~\ref{prop:connected-transitive} fires and Hardy's fifth axiom holds. The two foils
fail $(\mathrm{I})$ exactly where they fail continuity: fourteen of the twenty-four Clifford
operations are phases of no Pauli basis and fix no stabiliser state, their fixed points being
the magic states; the toy theory's three-cycles are phases of no epistemic basis; and in each
the outside's knowledge of the inside under the order-three protocol is either count-dependent
or a mixture over the count, so each makes the inside's time classical before any record exists.
The derivation of $(\mathrm{I})$ is an outline, not a proof, and it is what the open problem
below asks for.

\subsection{Why the cut matters physically}

The regime's continuity is not decorative. Reversibility alone, over complex amplitudes, licenses
the full invertible group $GL(n, \mathbb{C})$, and Aaronson showed what that buys: superluminal
signalling, and computational power at the level of \textsf{PP} \autocite{Aaronson2004h}. The cut
from $GL(n, \mathbb{C})$ down to the unitary group is the conserved quadratic form, and continuity
is what entitles the theory to make it: with the fifth axiom in hand, Hardy's theorem delivers the
complex Hilbert space, the quadratic norm and the unitary group in one stroke \autocite{Hardy2001a}
--- cited here, not re-proved. One honest caveat travels with the citation. Hardy's reconstruction
uses his remaining axioms besides --- composition and subspace structure --- and these must be
derived rather than imported: the composition ledger of \S\ref{sec:division-algebras} and the
locality-of-tomography argument of \S\ref{sec:self-reference} carry that duty, and it is discharged
there, not silently assumed here.

\subsection{Super-tasks at the boundary}

The premise just priced has a fifty-year-old stress test in the philosophy of mathematics, and the
framework passes it in an unexpected way. A super-task --- infinitely many operations completed in
finite time \autocite{Thomson_1954} --- presupposes its terminal moment as an actual point of the
timeline: the accumulation point of a generic sequence of step-times exists only in the completed
line, and the defining question, the state \emph{after all the steps}, is meaningful only under
limit-realisability of time. Super-tasks are downstream consumers of exactly the completeness at
issue in this section. Thomson's lamp then exhibits both sides of the \textsc{fanout} boundary at
once. On the record side, the lamp's description fixes its state at every finite stage and thereby
fixes no limit at all: any terminal value is consistent with the whole task
\autocite{Thomson_1954,benacerraf1962tasks}. That is not a puzzle about lamps; it is the proof that
the classical record algebra is not limit-closed, and cannot be. On the amplitude side the limit
exists --- a superposed fixed point, whose complex-probability treatment belongs to
\S\ref{sec:self-reference} --- because that side \emph{is} limit-closed, which is the fifth face
itself. The boundary is thereby the edge of limit-realisability: complete on the reversible side,
incomplete on the record side, with the lamp standing on both sides at once as the boundary's oldest
exhibit. A final clause dissolves rather than answers the classic question. In this framework time
is operation count (\S\ref{sec:space-of-computations}), and ``after infinitely many operations''
indexes no operational moment: the lamp packs its endless \textsc{fanout}s into the emergent
continuum, and asking for its record at the accumulation point asks the discrete contents for a
value that only the completed container possesses. On or off is not undecided; it is unasked.

\subsection{What remains, priced}

\begin{openproblem}[Every transformation a phase of an inside structure]\label{op:limit-realisability}
Proposition~\ref{prop:connected-transitive} reduces Hardy's fifth axiom to three premises:
connectedness of the pure states, transitivity, and closedness. Transitivity is argued above
from the absence of records. Closedness is the completion, adopted on the grounds stated above
and carrying no quantum content. Connectedness is density plus that completion, and density
follows, by Jordan's theorem and the infinite-order argument, from $(\mathrm{I})$: that every
admissible transformation of the regime is a phase of some admissible classical structure. What
remains open is the rigour of the derivation of $(\mathrm{I})$ from the paper's premise. Derive,
inside the fixed-point locus $\mathrm{Fix}(\sigma)$ of Proposition~\ref{prop:inside-outside},
that every endomap is a phase of a classical structure in $\mathrm{Fix}(\sigma)$, from the
self-representing object, the involution $\sigma$ as a genuine involution, and the reading of
the inside presentation as a sequence indexed by operation count with no record before
\textsc{fanout}. The outline above does this with the stationary-state construction of Page and
Wootters; a derivation passes if it makes that outline a theorem of the calculus, with the
outside's count-independent complete description shown to exist and to be pure from the
involution alone. It fails before any detail is checked if every premise it uses holds in the
stabiliser fragment of quantum theory, in Spekkens's toy theory \autocite{spekkens2007toy}, in
the real qubit, or in quantum theory over the Gaussian rationals, all of which violate
$(\mathrm{I})$ outright and all of which make the inside's operation count classical before any
record exists. It also fails if it imports the completion as a premise about transformations:
the completion is the real line and has no derivation, and a derivation of $(\mathrm{I})$ must
hold in the countable algebraic theory, where $(\mathrm{I})$ is true and continuity is not.
\end{openproblem}

The accounting, in one sentence: the axiom that every reconstruction programme pays is, in this
framework, one derivation posed, of Lawvere's shape --- that every endomap of the fixed-point
locus is a phase of a classical structure in it --- and one convention named, the completion,
which classical probability pays too. And
the hand-on is already visible: the complex amplitudes that this section's continuity underwrites
through Hardy's theorem, and the complex amplitudes that \S\ref{sec:self-reference} derives natively,
are one and the same $\mathbb{C}$ --- the floor of the ladder that \S\ref{sec:division-algebras}
climbs.

%      % =====================================================================
%  SECTION: The Division-Algebra Ladder
%  Drafted 2026-07-03. One of the three foundations sections
%  (sec:self-reference, sec:division-algebras, sec:fanout-boundary);
%  this one resolves the dangling \S\ref{sec:division-algebras} in the
%  Hopf-justification section, whose header names it "the
%  Cayley--Dickson / Hurwitz step".
%
%  PLACEMENT in main.tex (REORDERED per JS 2026-07-06): position 6,
%  opening Act II ("why stop at C?") right before its first consumer --
%      % EDITED 2026-10-29 scripts/patch_continuity_tower_premise_2026-10-29T1500.py
% EDITED 2026-09-20 scripts/patch_born_cost_ledger_v2_2026-09-20T0900.py
% =====================================================================
%  SECTION: Self-Reference: One Principle
%  Drafted 2026-07-04. The paper's foundational section; resolves the
%  largest dangling label in the manuscript (sec:self-reference).
%
%  PLACEMENT in main.tex: FIRST of the foundations chain --
%      \input{self_reference_section_2026-07-04}
%      \input{continuity_section_2026-07-04}
%      \input{division_algebras_section_2026-07-03}
%      \input{hopf_justification_section_2026-06-08}
%
%  SOURCES (re-read 2026-07-04): lawvere_inside_outside_fixed_point_
%  theorem_2026-05-04T1806.md (the spine: the two-principles reduction,
%  the theorem, the sigma candidate, the open items -> op minted per
%  D9); research_summary_neg_prob.md (negativity route; lamp (1+i)/2);
%  dynamics_companion_fanout_unitarity (legs 1a/1c; anti-events; the
%  Brukner-Zeilinger rebuttal per D6); Pure_states_and_mixed_states
%  (local tomography derived; QBism contrast per D7).
%
%  DECISIONS ENCODED (JS 2026-07-04): D1 title as below; D2 Wheeler1990
%  cited here as the title anchor; D3 lineage Popper -> Breuer ->
%  Svozil -> Szangolies (Dalla Chiara omitted); D4 Frauchiger_2018
%  cited here, dissolved in sec:wigners-friend; D5 Kirkwood-Dirac
%  heritage via the Arvidsson-Shukur 2024 review as one-stop; D6
%  Brukner-Zeilinger + the framework rebuttal IN ("Shannon is
%  post-FANOUT"); D7 QBism contrast IN; D8 Aaronson private
%  communication OMITTED (published island paper only); D9 one op
%  minted (op:sigma-functor), cross-linked to op:hopf-frobenius.
%
%  CITATIONS -- existing keys: Wheeler1990, small_2006, Godel1931,
%  Lawvere_1969, Yanofsky2003, CoeckeKissinger2017, Rovelli1996a,
%  Popper1950a, Popper1950b, Breuer1995, svozil_k_2018,
%  szangolies_j_2018, Frauchiger_2018, dirac1942, Feynman_r_82,
%  Feynman1987, Youssef1994 [SWAP-FLAG: intended paper is "Quantum
%  mechanics as complex probability theory", Mod. Phys. Lett. A 9
%  (1994) 2571 -- if that lives under youssef1994quantum, swap],
%  caticha_a_2006, goyal_knuth_skilling_2010, Aaronson2004h, Hardy2001a,
%  muller_p_2011, Thomson_1954, benacerraf1962tasks, cantor1891diagonal,
%  Turing_1937.
%  NEW keys (RIS in self_reference_refs_2026-07-04.ris, all verified
%  this session): tikochinsky1988feynman, renou2021real, chen2022ruling,
%  wu2022refutation, hoffreumonwoods2026real, arvidssonshukur2024kirkwood,
%  bruknerzeilinger2001shannon, fuchsschack2013qbism,
%  shannon1948mathematical.
%
%  CROSS-REFS: forward (fine under two-pass LaTeX): sec:continuity,
%  sec:division-algebras, sec:hopf-justification, prop:f1-dictionary,
%  sec:wigners-friend, op:limit-realisability, op:hopf-frobenius.
%  Still dangling after this section: sec:fanout-boundary (and its
%  sibling label sec:fanout-direction) -- both DRAFTED 2026-07-06.
%
%  AUDITED 2026-07-06 (section-by-section final pass, first stop):
%  proof of prop:inside-outside re-verified by hand; attribution chain
%  clean (Burgin/Khrennikov correctly absent); all header-declared
%  citations confirmed used; lamp arithmetic and antecedents verified;
%  FANOUT smallcaps consistent; no pre-definition use of "frame".
%  Fixes: stray Unicode em-dash (QBism paragraph, the file's only
%  non-ASCII character) -> house ---; trace-blind two-way stitch added
%  to the Born subsection ("one privacy, two duties");
%  op:sigma-functor gains its first standing constraint (F1
%  inheritance via the op:hopf-frobenius discharge) and its previously
%  one-directional local-tomography cargo made explicit.
% =====================================================================

\section{Self-Reference: One Principle}
\label{sec:self-reference-2}

Wheeler asked for it from bit: for physics derived from information, ``every it --- every particle,
every field of force'' obtained from binary answers to yes-or-no questions \autocite{Wheeler1990}.
This paper's route to that programme runs through G\"{o}del: the single principle from which
everything else here is derived, deferred, or priced is \emph{self-reference} --- the fact that any
complete physical description, made from inside the world it describes, must describe the describer.
This section states the principle, proves that the framework's second historical commitment follows
from it as a theorem rather than standing beside it as an axiom, and then extracts the price of the
principle in four instalments: probabilities whose parts can cancel, amplitudes valued in
$\mathbb{C}$, the Born rule, and the local tomography of composite systems. What the section cannot
yet derive is stated as precisely as what it can, in one numbered problem and in the ledger at the
end.

\subsection{Events, not objects}

The framework's ontology is relational: the fundamental entities are events and the relations
between them, and objects are the bookkeeping of relational invariants --- what remains stable
across many relations. The kinship with relational quantum mechanics \autocite{Rovelli1996a} is
genuine and acknowledged: facts are facts relative to the systems that register them. The
difference, developed over the rest of this paper, is that the present framework grounds the
relational structure in a single logical principle and then carries it, through the division
algebras and the Hopf fibrations, all the way to the gauge group and the particle spectrum ---
relations are not only the ontology but the derivation engine. On this reading a measurement is not
the revelation of a pre-existing property; it is the creation of a relation, and the later sections
of this paper read that creation geometrically (\S\ref{sec:hopf-justification}) and even
gravitationally. Here we need only the logical core: whatever describes the world is an event in
the world it describes.

\subsection{One principle, and a theorem}

The framework historically carried two foundational commitments: self-reference, and the
\emph{Inside/Outside Equivalence Principle} --- that the laws of physics must be identical whether
formulated by an embedded participant or by an external observer. Two principles are one too many,
and this subsection reduces them.

The reduction corrects the framework's own origin, and we make the correction in print. The 2005
form of the argument \autocite{small_2006} ran a G\"{o}delian chain: a self-referential process,
being complete, must by the incompleteness theorem \autocite{Godel1931} be inconsistent; an
inconsistent system can negate itself; and the negation of self-reference is the production of
causal sequence. The chain captures the right intuition and is retained here as motivation, but it
is not sound as a proof: classical inconsistency yields \emph{ex falso quodlibet} --- every
proposition provable, every causal sequence an output, no constraint at all. The load-bearing
argument is different, and better: a fixed-point theorem.

Let $\mathcal{C}$ be the category of physical descriptions and admissible processes --- concretely,
the dagger-compact-closed setting with classical structures in which categorical quantum mechanics,
and the whole of this paper, works \autocite{CoeckeKissinger2017}. Lawvere showed that in a
sufficiently rich such setting, any weakly point-surjective self-representation forces every
endomap to have a fixed point, and Yanofsky showed that one diagonal argument of this shape
underlies Cantor's, G\"{o}del's, Turing's and Tarski's theorems alike
\autocite{Lawvere_1969,Yanofsky2003,cantor1891diagonal}. The framework's single commitment is that
$\mathcal{C}$ contains a \emph{self-representing object} --- an object capable of encoding the
processes on itself. That is what ``self-reference is foundational'' means, stated categorically,
and it is the surjectivity hypothesis of the diagonal lemma. Now let
$\sigma \colon \mathcal{C} \to \mathcal{C}$ be the \emph{inside/outside involution}: the operation
exchanging the two complementary presentations of any description --- the inside presentation, as
participant-indexed total space in which the observer is part of the system, and the outside
presentation, as a base-times-fibre decomposition over a classical parameter space. That
$\sigma^{2} = \mathrm{id}$ is not an axiom but a coherence requirement: the outside view of the
inside view of a thing is the thing. The candidate concrete realisation of $\sigma$ is the
Hopf-duality swap --- total space against base-times-fibre for the three fibrations --- which gives
the involution computable content and ties this subsection to \S\ref{sec:hopf-justification}.

\begin{proposition}[Inside/outside fixed point]\label{prop:inside-outside}
Let $\mathcal{C}$ satisfy the richness condition above and contain a self-representing object, and
let $\sigma$ be an involution on $\mathcal{C}$. Then the fixed-point locus
$\mathrm{Fix}(\sigma)$ is non-empty: there exist descriptions invariant under the exchange of
inside and outside perspective.
\end{proposition}

\begin{proof}
Lawvere's construction, applied to $\sigma$ in place of negation: given a weakly point-surjective
$\varphi \colon A \to Y^{A}$, form $g(a) = \sigma\!\left(\varphi(a)(a)\right)$; point-surjectivity
supplies $a^{*}$ with $\varphi(a^{*}) = g$, and then $\varphi(a^{*})(a^{*})$ is a fixed point of
$\sigma$ \autocite{Lawvere_1969,Yanofsky2003}. Where the same construction applied to Boolean
negation yields paradox --- negation has no fixed point on a two-element type, so the conclusion
explodes the premises --- the inside/outside involution is an honest involution on a rich type, and
the conclusion is a constructive fixed point rather than a contradiction.
\end{proof}

\textbf{Definition.} \emph{Physical law is the structure carried by} $\mathrm{Fix}(\sigma)$.

The Inside/Outside Equivalence Principle is thereby demoted from axiom to theorem: physical law is
invariant under the exchange of perspectives \emph{because it is the fixed-point set of that
exchange}, and the fixed-point set is forced to be non-empty by self-reference alone. One
commitment remains foundational; everything the second commitment asserted is now derived. The
reader who accepts categorical quantum mechanics as an existing programme has already accepted the
machinery; the framework adds one premise and collects the rest. And the fixed-point locus is not
empty of structure: it inherits the Frobenius-algebra apparatus of $\mathcal{C}$, so that the
Bell, GHZ and W primitives --- and with them the particle classification of the later sections ---
live inside $\mathrm{Fix}(\sigma)$ from the start. What is proved above is non-emptiness given the
setting; what is assumed is the setting itself and the existence of $\sigma$ as a genuine
involution; what is open is the rigour of both, and it is priced now rather than hidden:

\begin{openproblem}[The involution made functorial]\label{op:sigma-functor}
Fix the category $\mathcal{C}$ rigorously --- the dagger-compact-closed setting with classical
structures is the candidate --- and define $\sigma$ on morphisms, not only on objects, the natural
route being the induced fibre--base exchange on spider diagrams; establish that $\sigma$ is an (at
worst pseudo-)involutive endofunctor; and prove or refute the identification of
$\mathrm{Fix}(\sigma)$ with the three Hopf cross-sections, under which Adams's theorem would become
a corollary of the fixed-point classification. This problem is cross-linked with Open
Problem~\ref{op:hopf-frobenius}: under an affirmative resolution of both, the parallelisable
terminus, the fixed-point classification, and compositional completeness are three statements of
one fact, and progress on either problem constrains the other. One constraint is already in place:
at the complex rung, Proposition~\ref{prop:f1-dictionary} supplies the dictionary between
classical structures and the fibration's lifted skeletons --- the clause of
Open Problem~\ref{op:hopf-frobenius} it discharges --- and any resolution here must reproduce it.
The problem also carries a piece of cargo recorded later in this section: the formalisation, within
the fixed category, of the perspective-accounting that derives local tomography.
\end{openproblem}

\subsection{The lineage}

The obstruction this section builds on is old, and the framework claims its completion, not its
discovery. Popper argued in 1950 that a predictor embedded in the world has irreducibly inferior
knowledge of itself, because it cannot impart information about itself without, by that very act,
rendering the information out of date \autocite{Popper1950a,Popper1950b} --- the passage is quoted
where the framework makes the obstruction structural (\S\ref{sec:wigners-friend}). Breuer proved
the impossibility of accurate state self-measurement within quantum theory \autocite{Breuer1995};
Svozil's monograph surveys the whole territory where physics meets self-reference, undecidability
and Chaitin-style algorithmic limits \autocite{svozil_k_2018}; Szangolies sharpened the modern form
as epistemic horizons \autocite{szangolies_j_2018}; and the Frauchiger--Renner theorem is the
contemporary crisis-form --- nested observers deriving contradictory records when each treats the
others' quantum descriptions as universally valid \autocite{Frauchiger_2018} --- which
\S\ref{sec:wigners-friend} dissolves relationally rather than paradoxically. This framework's
addition to the lineage is twofold. Structurally, Proposition~\ref{prop:f1-dictionary} converts
Popper's temporal obsolescence into algebraic cancellation: the datum a copier cannot put into its
own record is exactly its fibre commitment, which does not arrive stale --- it cancels. And
foundationally, the present section runs the lineage in reverse: where the tradition treats
self-reference as a limit on knowledge \emph{within} physics, the framework takes it as the
generator of physics, and derives the quantum formalism from it.

\subsection{The price is negative probability}\label{sec:complex-amplitudes-2}

What does it cost a probabilistic description to contain its own describer? The first instalment is
the sign. In the regime before any record exists, every event has an anti-event --- the description
must be able to undo what it has provisionally done, or the self-description could never be
revised into consistency --- and a probability calculus closed under undoing cannot be confined to
$[0,1]$: it is a signed calculus, in which the parts of a probability can oppose and cancel. The
framework did not discover negativity as bookkeeping; Dirac admitted negative probabilities as
formal intermediaries \autocite{dirac1942}, Feynman defended them as consistent whenever the final,
recorded probabilities stay positive \autocite{Feynman_r_82,Feynman1987}, and the Kirkwood--Dirac
distribution has carried complex quasiprobability from 1933 to the present, as a recent
comprehensive review documents \autocite{arvidssonshukur2024kirkwood}. What the framework adds is
the forcing: negativity is not a permissible convenience but the first price of self-referential
closure, and it is charged exactly and only in the pre-record regime --- which is why every
negativity in the tradition lives in intermediate, unrecorded quantities and never in a final
record.

The second instalment is complexity. A signed information measure must extend Shannon's
$S = -\ln p$ off the positive axis, and the analytic continuation of the logarithm carries $p$ into
the complex plane: amplitudes, with the phase as the continued information. Here too the extension
has a prior tradition, and the attributions are owed precisely. Tikochinsky derived Feynman's rules
for complex amplitudes from consistency of composition \autocite{tikochinsky1988feynman}; Youssef
reformulated quantum mechanics as complex probability theory outright \autocite{Youssef1994};
Caticha carried objective amplitudes to Bayesian probabilities \autocite{caticha_a_2006}; and
Goyal, Knuth and Skilling derived Feynman's rules from a pair-postulate without assuming
analyticity \autocite{goyal_knuth_skilling_2010}. The extension is prior art. The framework's
claim is the compulsion: the complex extension is not a modelling choice among exotic probability
theories but the unique continuation forced by self-referential closure --- the same forcing that
Aaronson glimpsed when he argued quantum mechanics is what negative probability inevitably becomes
\autocite{Aaronson2004h}.

Why does the continuation stop at $\mathbb{C}$? Downward, the real field is not available: real
quantum theory is experimentally distinguishable from complex quantum theory in network scenarios
and has been disfavoured by experiment
\autocite{renou2021real,chen2022ruling,wu2022refutation}; the recent objection that the exclusion
rests on locally tomographic composition \autocite{hoffreumonwoods2026real} is, for this framework,
not an escape but a confirmation of where the weight sits --- local tomography is \emph{derived}
below, not assumed, so the framework closes from its own resources exactly the door the objection
holds open. Upward, the quaternions and octonions fail as a \emph{base} for the information
continuation: the continued measure is a logarithm, and beyond $\mathbb{C}$ the logarithm of a
product is no longer the sum of logarithms --- non-commutativity severs
$\log(ab) = \log a + \log b$ --- so the additivity that makes it an information measure is lost.
Aaronson's parameter-counting and signalling arguments against a quaternionic base point the same
way \autocite{Aaronson2004h}. The higher algebras are not banished; they return as the
\emph{internal} algebras of the register, which is the business of \S\ref{sec:division-algebras}.
The base is $\mathbb{C}$, for informational reasons, and stays $\mathbb{C}$.

\subsection{The Born rule from fibre privacy}\label{sec:born-rule}

The third instalment is the quadratic rule. An amplitude carries one datum more than any record of
it can: its overall phase. This was asserted throughout the framework's development as the
observer's private, self-referential ignorance; it is now a proved statement, in exactly the form
needed. By Proposition~\ref{prop:f1-dictionary}, a classical structure --- the copy operation any
record requires --- is committed to a specific lift of its skeleton, and that commitment is absent
from every classical record the structure induces: the copier is fibre-committed, the record is
fibre-blind, and the fibre freedom is gauge for every observer. (The computation extends verbatim
from one qubit to the full register, where Proposition~\ref{prop:trace-blind} makes it the ground
of hypercharge's tracelessness --- one privacy, two duties.) Popper's obsolescence, made
cancellation. It follows that anything recordable is invariant under the fibre --- observables and
their statistics must be $U(1)$-invariant functions of the amplitude. The lowest-degree polynomial invariant
of a complex amplitude is the quadratic $|\psi|^{2}$ --- the non-polynomial invariant $|\psi|$,
of degree one, is excluded by the analyticity declared at \S\ref{sec:complex-amplitudes} --- and
normalisation fixes the coefficient:
probability is the minimal record-compatible function of the amplitude,
\[
p \;=\; |\psi|^{2},
\]
the Born rule as the cheapest invariant the fibre permits. The same square is reached from the
cost side in \S\ref{sec:born-cost-ledger}, where a record is the closed loop of a history and its
return leg and the exponent is the number of legs; neither route uses the lattice of projections
that Gleason's theorem needs, and both hold for a single qubit. The derivation is the framework's
own;
its premise --- that the fibre datum is structurally excluded from records --- is
Proposition~\ref{prop:f1-dictionary}, and its geometric reading --- probability as base-point data
under the Hopf projection --- is developed in \S\ref{sec:hopf-justification}.

One standing objection to information-theoretic foundations is answered here, because this
framework needs the answer. Brukner and Zeilinger argued that Shannon information is conceptually
inadequate for quantum measurements: it is basis-dependent, and its operational meaning presupposes
ignorance of a definite but unknown fact, which the quantum case does not supply
\autocite{bruknerzeilinger2001shannon}. The framework's reply is structural, and it is a
demarcation rather than a rebuttal. Shannon's precondition fails exactly where they say it fails
--- in the pre-record regime, where there is no definite unknown to be ignorant of --- and there
the operative measure is the quadratic invariant just derived. Shannon's axioms hold exactly where
the precondition holds --- on records, after \textsc{fanout}, where outcomes are definite and
ignorance is classical --- and there the uniqueness theorem forces $-\log p$
\autocite{shannon1948mathematical}. The basis-dependence dissolves with the same stroke: a basis
\emph{is} selected precisely at \textsc{fanout}, so the basis-dependent measure applies only where
a basis exists. Shannon is not inadequate; Shannon is post-\textsc{fanout}. The framework thereby
keeps both measures, each on its own side of the boundary, and uses the additivity of the
post-record measure --- the very property the previous subsection used to pin the base field ---
with a clear conscience.

\subsection{Thomson's lamp, resolved}

The section owes \S\ref{sec:continuity} a debt incurred there: the lamp's amplitude-side limit.
Thomson's lamp --- toggled at times $1 - 2^{-n}$, its state at $t = 1$ demanded --- has no
classical answer: the task description fixes every finite stage and thereby fixes no limit at all
\autocite{Thomson_1954,benacerraf1962tasks}, which \S\ref{sec:continuity} read as the proof that
the record algebra is not limit-closed. On the amplitude side the limit exists, and the calculus of
this section computes it. The toggle exchanges the two amplitudes; the self-consistent state of the
completed task is the fixed point of the exchange in the complex calculus, and solving it under the
signed continuation gives
\[
P(\mathrm{on}) = \tfrac{1 + i}{2}, \qquad P(\mathrm{off}) = \tfrac{1 - i}{2},
\]
a superposed fixed point whose imaginary part carries the unrecordable datum --- the phase of the
endless alternation, when the lamp \emph{would have} stopped. Two checks certify that the
assignment is calculus rather than decoration. The pair is normalised,
$P(\mathrm{on}) + P(\mathrm{off}) = 1$; and the Born rule of the previous subsection, applied to
the fixed point, returns $|P(\mathrm{on})|^{2} = \tfrac{1}{2}$ --- exactly the symmetric classical
ignorance that is the only sane record-level answer, recovered rather than imposed. The lamp,
which entered the framework at its very beginning as the place where classical computation meets
its boundary, thus exits it as a worked example: undetermined as a record, determinate as an
amplitude, the two sides split precisely along \textsc{fanout} --- and the diagonal construction it
physically instantiates \autocite{cantor1891diagonal,Turing_1937} is the same one that powers
Proposition~\ref{prop:inside-outside}. The status flag is owed here as everywhere: the
undetermination is theorem-level; the specific fixed-point value is the framework's complex
calculus, structural, offered with its two consistency checks shown.

\subsection{Local tomography, derived}\label{sec:local-tomography}

The fourth instalment concerns composites, and it discharges two promises at once. The
Inside/Outside Equivalence Principle --- now a theorem --- requires that an external observer's
characterisation of a composite system be recoverable by the embedded participants, and the
embedded participants have access only to local measurements and their correlations. Equivalence
of the two perspectives therefore entails that the state of a composite is determined by the
statistics of local measurements: \emph{local tomography}, obtained as a consequence rather than
imposed as an axiom. This converts a postulate of the reconstruction programmes
\autocite{Hardy2001a,mullerp2011} into a theorem-shaped consequence of the one principle ---
discharging the caveat recorded in \S\ref{sec:continuity} --- and it closes, from the framework's
own resources, the escape route in the real-number debate noted above: the assumption the
objection of \autocite{hoffreumonwoods2026real} contests is here not an assumption. The status is
flagged plainly: the derivation is structural, one paragraph of perspective-accounting, and its
formalisation within the categorical setting belongs to Open Problem~\ref{op:sigma-functor}.

\subsection{Contrast: QBism, and the ledger}

The nearest neighbour in spirit deserves a paragraph of honest contrast. QBism
\autocite{fuchsschack2013qbism} holds, as this framework does, that quantum probabilities are an
agent's --- that the theory is a calculus of expectation, not a catalogue of properties. What QBism
does not say is why the agent's expectations must be complex-valued and quadratically squared; the
Born rule enters as a normative addition to Bayesian coherence, and the complex amplitude as an
unexplained representation. This section is, in effect, the missing why: the agent's ignorance is
self-referential --- the agent is inside the description --- and self-referential ignorance is
complex, with the quadratic rule as the cheapest record-compatible invariant of it. Where QBism
stops at the agent, the framework derives the agent's calculus.

The ledger, then. Delivered in this section: the reduction of two principles to one
(Proposition~\ref{prop:inside-outside}); the pre-record regime with its anti-events, handed to
\S\ref{sec:continuity} as the five-face regime whose continuity is priced there
(Open Problem~\ref{op:limit-realisability}); the complex base, forced downward by experiment and
upward by the additivity of information; the Born rule from fibre privacy; the lamp's fixed point;
and local tomography. Owed, and stated where they are owed: the functorial rigour of the involution
and the fixed-point classification (Open Problem~\ref{op:sigma-functor}); the remaining premise of continuity (Open Problem~\ref{op:limit-realisability}, in
\S\ref{sec:continuity}); and the
rung-by-rung compulsion of the algebra ladder, posed as Open
Problem~\ref{op:ladder-compulsion} in \S\ref{sec:division-algebras}. The
hand-on is the paper's spine: the $\mathbb{C}$ this section derives natively and the $\mathbb{C}$
that \S\ref{sec:continuity} underwrites through Hardy's theorem are the same field, the floor of
the ladder that \S\ref{sec:division-algebras} climbs and \S\ref{sec:hopf-justification}
geometrises. One principle; everything else is theorem, derivation, or a debt with a stated
invoice.

%      % EDITED 2026-10-29 scripts/patch_continuity_tower_premise_2026-10-29T1500.py
% EDITED 2026-09-15 scripts/patch_value_continuum_declared_cencov_2026-09-15T1540.py
% EDITED 2026-09-15 scripts/patch_continuity_connected_state_space_2026-09-15T1430.py
% =====================================================================
%  SECTION: Continuity Is Not an Axiom
%  Drafted 2026-07-04 from continuity_limit_realisability_2026-07-04.md
%  (now in project documents) and the May record
%  (dynamics_companion_fanout_unitarity_2026-05-20T1430.md).
%
%  PLACEMENT in main.tex: a standalone section between self-reference
%  and division-algebras --
%      \input{self_reference_section_2026-07-04}
%      \input{continuity_section_2026-07-04}
%      \input{division_algebras_section_2026-07-03}
%  (Architecture call SETTLED by the sync manifest: standalone.)
%
%  CITATIONS -- existing keys (as supplied by JS 2026-07-04):
%    Hardy2001a, Aaronson2004h, Thomson_1954, Lawvere_1969,
%    Turing_1937, muller_p_2011, Brukner2009.
%  NEW keys (RIS in continuity_refs_2026-07-04.ris):
%    cantor1891diagonal, benacerraf1962tasks, broeckertomdieck1985compact.
%  Deliberately excluded per JS: Chiribella--D'Ariano--Perinotti;
%  Bargmann/Wigner's theorem. Yanofsky2003 now paired on the
%  Lawvere-shaped clause (the pending-key excuse was void: the key
%  exists and sec:self-reference uses it).
%
%  CROSS-REFS: sec:self-reference now DRAFTED (2026-07-04, later the
%  same day) -- no longer dangling; only sec:fanout-boundary (and its
%  sibling sec:fanout-direction) -- DRAFTED 2026-07-06.
%
%  AUDITED 2026-07-06 (final pass, second stop): both proofs
%  re-verified by hand (Lagrange divisibility sentence exact; the
%  Baire countable-escape chain sound); caveat-discharge with
%  sec:self-reference confirmed two-way (already present); file pure
%  ASCII; no pre-definition "frame"; all citations used as declared.
%  Fixes: Yanofsky2003 paired on the Lawvere clause; the op's two
%  candidate routes anchored to printed propositions (f1-dictionary
%  for the divisibility premise's referent, inside-outside for the
%  closedness demand) -- statuses unchanged, referents sharpened;
%  header staleness corrected. Deliberate FORWARD references --
%  prop:f1-dictionary and sec:hopf-justification (defined later in the
%  paper; fine under two-pass LaTeX and flagged in prose),
%  sec:division-algebras, sec:space-of-computations.
%
%  DECISIONS ENCODED: op:limit-realisability minted (the owed arrow,
%  priced: exactly two derivations, divisibility and closedness);
%  Lagrange used twice (finite exclusion; component-group triviality);
%  the (1+i)/2 lamp value NOT stated here -- "a superposed fixed point"
%  with the complex-probability treatment deferred to
%  sec:self-reference; the two-branch over-determination stated to
%  pre-empt a circularity reading; Thomson spelling throughout.
% =====================================================================

\section{Continuity Is Not an Axiom}
\label{sec:continuity}

Hardy derives quantum theory from five axioms \autocite{Hardy2001a}. Four of them are satisfied
by classical probability theory; the fifth --- that a \emph{continuous} reversible
transformation connects any two pure states --- is the entire difference, the single postulate
that turns the classical simplex into the complex Hilbert space. Every reconstruction programme
pays an equivalent toll: continuous reversibility in Masanes and M\"{u}ller
\autocite{mullerp2011}, the corresponding postulate in Daki\'{c} and Brukner
\autocite{Brukner2009}. The axiom is load-bearing, and in every programme it is unexplained: it
is where the quantumness is put in by hand. The claim of this section is that in the present
framework it is not an axiom. Continuity is the fifth face of the pre-measurement regime that self-reference forces, and what
stands between the framework and a full derivation is not a postulate about values but one
premise about transformations --- that every transformation the inside performs has a stationary
outside presentation --- derived in outline below from the inside/outside involution of
\S\ref{sec:self-reference} and stated as an open problem at the end of the section with nothing
hidden in it. The continuum itself, the completion that every probability theory adopts, is
named for what it is: the real line, shared with classical probability and carrying no quantum
content.

One structural remark prevents a misreading. The framework reaches complex amplitudes natively in
\S\ref{sec:self-reference}, by the negative-probability route; the present section reaches the same
destination by an independent branch --- the regime's continuity, then Hardy's theorem. The two
branches converging on one structure is over-determination, not circularity, and the convergence is
itself evidence that the destination is forced.

\subsection{Five faces of one regime}

The regime before any measurement --- before any \textsc{fanout} --- can be characterised five
equivalent ways: negative probability is available; anti-events exist, every event being undoable;
no classical record exists; no preferred basis is selected; and a continuous reversible
transformation connects any two pure states. These are not five premises to be combined. They are
one regime under five descriptions, and \textsc{fanout} (\S\ref{sec:fanout-boundary}) is the single
event that fails all five at once: it writes the record, selects the basis, kills the anti-events,
discretises the transformation group, and switches off the negativity. The equivalence of the first
two faces is the anti-event identity of \S\ref{sec:self-reference}; the third is the boundary's
definition. The work of this section is the equivalence of the last two faces, and then the route
from the regime's existence to its continuity. The fourth face has, since
Proposition~\ref{prop:f1-dictionary} (\S\ref{sec:hopf-justification} --- a deliberate forward
reference), an exact form: a basis-selection is the commitment of a classical structure, a lifted
antipodal skeleton, so ``no preferred basis'' means precisely that no copy structure has fired.

\subsection{Basis selection is the collapse of continuity}

\begin{proposition}[Continuity forbids a selected basis]\label{prop:basis-discreteness}
If a finite set of pure states is privileged --- a preferred basis --- then every reversible
transformation respecting the privilege carries the set to itself and acts on it as a permutation;
the respecting transformations therefore realise only a finite group of motions among the
distinguished states, and no continuous family of them connects distinct pure states of the set.
Contrapositively, continuous reversibility entails that no basis is selected.
\end{proposition}

\begin{proof}
A privileged finite set is by definition carried to itself by every structure-preserving reversible
map, whose induced action on the set is a permutation of finitely many points. A continuous path of
such maps induces a continuous path in a finite discrete set, which is constant. Hence a continuous
reversible family moving one distinguished pure state to another respects no such privilege, and if
every pair of pure states is so connected, no finite set is privileged.
\end{proof}

The converse is where the historical difficulty lies: the absence of a privileged finite set does
not by itself force continuity, for the transformation group could in principle be infinite and
discrete. Two structural properties close every such escape at once.

\begin{proposition}[Divisible and closed implies connected]\label{prop:divisible-closed}
Let $G$ be the group of reversible transformations of the pre-measurement regime, acting on the
state body. If $G$ is divisible --- every element has an $n$-th root for every $n$ --- and
topologically closed, then $G$ is a compact connected Lie group; and if $G$ moreover acts
transitively on the pure states, then a continuous reversible transformation connects any two of
them: Hardy's fifth axiom.
\end{proposition}

\begin{proof}
The state body is compact and convex with nonempty interior, so its symmetry group is compact, and
a closed subgroup of a compact matrix group is a compact Lie group with finitely many connected
components \autocite{broeckertomdieck1985compact}. The component group $G/G_{0}$ is a quotient of a
divisible group and hence divisible; but a nontrivial finite group is never divisible --- with
$\lvert H\rvert = m$ every element satisfies $h^{m} = e$, so no element other than the identity
possesses an $m$-th root --- and $G/G_{0}$ is finite, so it is trivial and $G = G_{0}$ is connected.
A connected Lie group is path-connected, and transitivity then supplies the continuous reversible
path between any two pure states.
\end{proof}

The same one-line fact of Lagrange closes the remaining escapes. Finite groups are excluded
directly, being nontrivial and never divisible. Countable groups are excluded because a
countable closed subgroup of a compact group is compact, has an isolated point by the Baire
category theorem, and is therefore discrete and finite. Uncountable disconnected groups are
excluded by the component-group argument of the proposition. Every route to a discrete
transformation group therefore passes through a failure of divisibility or a failure of
closedness. Divisibility, however, asks for more than the conclusion needs, and the next
proposition removes it.

\begin{proposition}[Continuity from a connected state space]\label{prop:connected-transitive}
Let the pure states of a finite-dimensional system form a connected space $P$, and let the group
$G$ of reversible transformations of the pre-measurement regime be topologically closed and act
transitively on $P$. Then the identity component $G_0$ already acts transitively on $P$, and a
continuous reversible transformation connects any two pure states: Hardy's fifth axiom holds,
with no divisibility assumed.
\end{proposition}

\begin{proof}
As in Proposition~\ref{prop:divisible-closed}, closedness makes $G$ a compact Lie group with
finitely many connected components \autocite{broeckertomdieck1985compact}, so $G_0$ is a closed
normal subgroup of finite index and is path-connected. Each $G_0$-orbit in $P$ is the continuous
image of a compact group, hence compact, hence closed. Since $G$ is transitive and $G_0$ is
normal, the $G_0$-orbit of $gx$ is $g$ applied to the $G_0$-orbit of $x$, so the cosets of $G_0$
permute the orbits transitively and there are finitely many of them. A finite partition of $P$
into closed sets is a partition into open sets; $P$ is connected; so there is one orbit and
$G_0$ is transitive. Given pure states $x$ and $y$, choose $g \in G_0$ with $gx = y$ and a
continuous path $g_t$ in $G_0$ from the identity to $g$; then $g_t x$ is a continuous reversible
transformation from $x$ to $y$.
\end{proof}

The proposition relocates the premise. Divisibility was asked for because a transformation group
could be infinite and discrete; but a group acting transitively on a continuum cannot be
countable, and a closed group acting transitively on a connected space cannot be disconnected in
the way that matters, so what the argument needs from the regime is not the availability of
$n$-th roots but the connectedness of the state space itself. That is a statement about values:
that the pure states are all the unit vectors of $\mathbb{C}^{K}$ up to phase, the space
$\mathbb{CP}^{K-1}$, connected for $K \ge 2$, and hence that the amplitudes of
\S\ref{sec:self-reference} fill their continuum. The two standing foils show the premise doing
its work. The stabiliser fragment of quantum theory has complex amplitudes and is contextual,
yet it has finitely many pure states, its outcome probabilities take only the values $0$,
$\tfrac12$ and $1$, and its reversible transformations form a finite group, the Clifford group,
so Hardy's fifth axiom fails in it; Spekkens's toy theory \autocite{spekkens2007toy} is the same
with real amplitudes, its ignorance rationed by a knowledge-balance principle. Neither has a
connected state space, and neither is a model of the regime described here, whose credences take
every value. Divisibility, which both also lack --- the phase gate has no square root in the Clifford group
--- is a symptom of the discreteness, not its cause. What both lack can be said exactly, and the
next subsection says it: each contains a transformation that fixes no state it admits.

\subsection{The premises, and what they mean}

Three premises carry Proposition~\ref{prop:connected-transitive}. Transitivity is the regime's
homogeneity: inside/outside equivalence leaves no pure state privileged, and a proper closed
invariant set of pure states would define a superselection observable, a classical record
predating any \textsc{fanout} (\S\ref{sec:fanout-boundary}). Connectedness of the state space is two statements, and only one of them is quantum. The first
is that the pure states are infinite and dense, with a fixed point for every transformation; the
second is that the state space is complete. Quantum theory over the algebraic numbers ---
amplitudes in the algebraic closure of $\mathbb{Q}$ within $\mathbb{C}$ --- has the first
without the second: it has $K = N^{2}$, the Born rule, superposition, entanglement and Bell
violation at the Tsirelson value, every unitary has an algebraic eigenvector, and it is
countable and totally disconnected. So the whole quantum content of Hardy's axiom lies in the
density, and the completion is the real-number convention of every probability theory. Hardy
adopts it before his fifth axiom is stated, in the first, which makes probabilities real numbers
and gives his classical theory a real simplex of mixtures; this framework adopts it at the
Shannon step, where credences are declared to range over a connected domain so that the
continuation of the logarithm is unique (\S\ref{sec:complex-amplitudes}). Given that convention,
density of the pure states is connectedness of their closure. Divisibility, which an earlier
version of this section asked for, is not needed: it follows from the conclusion, every element
of a compact connected Lie group having roots of every order.

Closedness is the completion, limit-realisability: a Cauchy sequence of available
transformations has its limit available. It cannot be derived, and the reason should be stated
rather than hidden. Any collection of first-order principles about states, transformations and
probabilities that has an infinite model has a countable one, and a countable subset of
$\mathbb{CP}^{K-1}$ is totally disconnected; the founding non-computability does not escape
this, since the arithmetical complex numbers form a countable algebraically closed field
containing non-computable numbers and closed under every limit it can itself describe. Hardy's
axiom and closedness are second-order, quantifying over paths and over sequences, and a
second-order principle is grounded as Dedekind completeness is grounded or not at all. The
grounds are these. The completion of a dense set under the finite-precision metric is unique and
depends on no choice, whereas any countable alternative is a choice of definability level, an
extra law with no ground and no test; and the framework's own result is that no observer's
descriptive level captures the regime, since by \S\ref{sec:self-reference} every observer's
state exceeds its own level and the inside/outside principle makes that a law rather than a
perspective. The burden therefore falls on any restriction, not on the completion, which is the
least committal option and carries no quantum content. What limit-realisability adds to the line
is still worth stating exactly, because it is the corrected form of a natural intuition.
Cantor's diagonal \autocite{cantor1891diagonal} does not build the continuum; its output is an
infinite decimal --- a Cauchy sequence of rationals --- which exists \emph{as a point} only
because the line is complete. The diagonal presupposes limit-realisability and then shows that
the completed line outruns every internal listing. And the character of the completion is the
framework's own: the computable reals are countable \autocite{Turing_1937}, hence of measure
zero, so once the transformation group contains its limits almost all of them are
non-computable. Non-computability is a property of the completed continuum, not the reason for
it. The constructive reading, on which only describable limits exist, remains consistent; it is
the countable theory above, and the two readings differ in nothing an experiment can certify.

Density is where the section's title is earned, and it is derived from the paper's own premise.
Write $(\mathrm{I})$ for the statement that every admissible transformation of the regime is a
phase of some admissible classical structure: diagonal, that is, in the record basis of some
observer the regime contains, a phased spider of that basis in the calculus of
\S\ref{sec:hopf-justification} \autocite{coeckeduncan2011interacting} (VERIFY-SUPPORT). It says
of transformations what the global phase of \S\ref{sec:self-reference} says of states: there are
no absolute ones. It follows from three things the paper already holds. The inside/outside
involution is a genuine involution, so the outside presentation of the inside is a presentation
of the same object with nothing lost: the outside can have complete knowledge of the inside, a
potential record, the joint state lying in a structure the outside carries. Time is operation
count, internal to the system that operates (\S\ref{sec:space-of-computations}), so the outside
does not share the inside's count. And before any \textsc{fanout} that count is not classically
recorded, so the outside's description is pure. Let $g$ be the transformation the inside
performs per operation. The outside's complete knowledge cannot depend on the inside's count,
and it is pure, so it is a pure state of the inside's clock and system that is stationary under
``advance the clock and apply $g$''; the stationary states of that map are clock Fourier modes
tensored with eigenvectors of $g$ and nothing else \autocite{pagewootters1983evolution}
(VERIFY-SUPPORT). So the outside holds a structure in which $g$ acts by phases, which is
$(\mathrm{I})$; and the relative phase in that structure is the outside's ignorance of the
inside's computational count, the same ignorance the global phase records one level down. The
consequences are a theorem. A phase of a structure fixes the structure's copyable states, so
every admissible transformation fixes an admissible pure state; a finite transitive action has
an element moving every point (Jordan), so the pure states are infinite; the transformation
carrying a stabiliser state to the fixed point of the composite of a $Z$-phase and an $X$-phase
turns through $\arccos(1/\sqrt{3})$, an irrational fraction of a turn, so the group has an
element of infinite projective order and its orbits are dense. With the completion,
Proposition~\ref{prop:connected-transitive} fires and Hardy's fifth axiom holds. The two foils
fail $(\mathrm{I})$ exactly where they fail continuity: fourteen of the twenty-four Clifford
operations are phases of no Pauli basis and fix no stabiliser state, their fixed points being
the magic states; the toy theory's three-cycles are phases of no epistemic basis; and in each
the outside's knowledge of the inside under the order-three protocol is either count-dependent
or a mixture over the count, so each makes the inside's time classical before any record exists.
The derivation of $(\mathrm{I})$ is an outline, not a proof, and it is what the open problem
below asks for.

\subsection{Why the cut matters physically}

The regime's continuity is not decorative. Reversibility alone, over complex amplitudes, licenses
the full invertible group $GL(n, \mathbb{C})$, and Aaronson showed what that buys: superluminal
signalling, and computational power at the level of \textsf{PP} \autocite{Aaronson2004h}. The cut
from $GL(n, \mathbb{C})$ down to the unitary group is the conserved quadratic form, and continuity
is what entitles the theory to make it: with the fifth axiom in hand, Hardy's theorem delivers the
complex Hilbert space, the quadratic norm and the unitary group in one stroke \autocite{Hardy2001a}
--- cited here, not re-proved. One honest caveat travels with the citation. Hardy's reconstruction
uses his remaining axioms besides --- composition and subspace structure --- and these must be
derived rather than imported: the composition ledger of \S\ref{sec:division-algebras} and the
locality-of-tomography argument of \S\ref{sec:self-reference} carry that duty, and it is discharged
there, not silently assumed here.

\subsection{Super-tasks at the boundary}

The premise just priced has a fifty-year-old stress test in the philosophy of mathematics, and the
framework passes it in an unexpected way. A super-task --- infinitely many operations completed in
finite time \autocite{Thomson_1954} --- presupposes its terminal moment as an actual point of the
timeline: the accumulation point of a generic sequence of step-times exists only in the completed
line, and the defining question, the state \emph{after all the steps}, is meaningful only under
limit-realisability of time. Super-tasks are downstream consumers of exactly the completeness at
issue in this section. Thomson's lamp then exhibits both sides of the \textsc{fanout} boundary at
once. On the record side, the lamp's description fixes its state at every finite stage and thereby
fixes no limit at all: any terminal value is consistent with the whole task
\autocite{Thomson_1954,benacerraf1962tasks}. That is not a puzzle about lamps; it is the proof that
the classical record algebra is not limit-closed, and cannot be. On the amplitude side the limit
exists --- a superposed fixed point, whose complex-probability treatment belongs to
\S\ref{sec:self-reference} --- because that side \emph{is} limit-closed, which is the fifth face
itself. The boundary is thereby the edge of limit-realisability: complete on the reversible side,
incomplete on the record side, with the lamp standing on both sides at once as the boundary's oldest
exhibit. A final clause dissolves rather than answers the classic question. In this framework time
is operation count (\S\ref{sec:space-of-computations}), and ``after infinitely many operations''
indexes no operational moment: the lamp packs its endless \textsc{fanout}s into the emergent
continuum, and asking for its record at the accumulation point asks the discrete contents for a
value that only the completed container possesses. On or off is not undecided; it is unasked.

\subsection{What remains, priced}

\begin{openproblem}[Every transformation a phase of an inside structure]\label{op:limit-realisability}
Proposition~\ref{prop:connected-transitive} reduces Hardy's fifth axiom to three premises:
connectedness of the pure states, transitivity, and closedness. Transitivity is argued above
from the absence of records. Closedness is the completion, adopted on the grounds stated above
and carrying no quantum content. Connectedness is density plus that completion, and density
follows, by Jordan's theorem and the infinite-order argument, from $(\mathrm{I})$: that every
admissible transformation of the regime is a phase of some admissible classical structure. What
remains open is the rigour of the derivation of $(\mathrm{I})$ from the paper's premise. Derive,
inside the fixed-point locus $\mathrm{Fix}(\sigma)$ of Proposition~\ref{prop:inside-outside},
that every endomap is a phase of a classical structure in $\mathrm{Fix}(\sigma)$, from the
self-representing object, the involution $\sigma$ as a genuine involution, and the reading of
the inside presentation as a sequence indexed by operation count with no record before
\textsc{fanout}. The outline above does this with the stationary-state construction of Page and
Wootters; a derivation passes if it makes that outline a theorem of the calculus, with the
outside's count-independent complete description shown to exist and to be pure from the
involution alone. It fails before any detail is checked if every premise it uses holds in the
stabiliser fragment of quantum theory, in Spekkens's toy theory \autocite{spekkens2007toy}, in
the real qubit, or in quantum theory over the Gaussian rationals, all of which violate
$(\mathrm{I})$ outright and all of which make the inside's operation count classical before any
record exists. It also fails if it imports the completion as a premise about transformations:
the completion is the real line and has no derivation, and a derivation of $(\mathrm{I})$ must
hold in the countable algebraic theory, where $(\mathrm{I})$ is true and continuity is not.
\end{openproblem}

The accounting, in one sentence: the axiom that every reconstruction programme pays is, in this
framework, one derivation posed, of Lawvere's shape --- that every endomap of the fixed-point
locus is a phase of a classical structure in it --- and one convention named, the completion,
which classical probability pays too. And
the hand-on is already visible: the complex amplitudes that this section's continuity underwrites
through Hardy's theorem, and the complex amplitudes that \S\ref{sec:self-reference} derives natively,
are one and the same $\mathbb{C}$ --- the floor of the ladder that \S\ref{sec:division-algebras}
climbs.

%      % EDITED 2026-10-29 scripts/patch_continuity_tower_premise_2026-10-29T1500.py
% EDITED 2026-09-15 scripts/patch_value_continuum_declared_cencov_2026-09-15T1540.py
% EDITED 2026-09-15 scripts/patch_continuity_connected_state_space_2026-09-15T1430.py
% =====================================================================
%  SECTION: The Euclidean metric on the state sphere, and the global
%  phase as self-referential ignorance
%  Origin: drafted in an earlier conversation (2026-06-07), edited by
%  JS; audited into the current architecture 2026-07-06.
%
%  PLACEMENT in main.tex (reordered chain, JS 2026-07-06): position 3 --
%      \input{self_reference_section_2026-07-04}
%      \input{continuity_section_2026-07-04}
%      \input{L2_metric_global_phase_section_2026_06_07}   <-- here
%      \input{wigners_friend_in_the_hopf_picture_2026-06-13}
%  (WF backward-references sec:l2-metric-fibre: order-consistent.)
%
%  AUDITED 2026-07-06 (final pass): Lamperti/Banach classification,
%  Hopf map components, submersion facts, Hardy K=N^r reading all
%  verified; the recoverable-vs-gauge-private fibre distinction
%  confirmed canonical and untouched; file pure ASCII. Fixes: opening
%  rewritten (Popper C+/C/E forward-voiced to sec:wigners-friend;
%  continuity claim made CONDITIONAL per JS -- "assuming those premises
%  can be proved", op:limit-realisability cited); Aaronson paragraph
%  de-duplicated against sec:continuity by stitch; sec:born-rule was
%  dangling paper-wide -- label minted at its canonical home (the Born
%  subsection of sec:self-reference) and the claim here recalibrated
%  ("forced, not chosen" -> lowest-degree-invariant, geometric face);
%  record-absence upgraded from interpretive to theorem via
%  prop:f1-dictionary while the self-ignorance identification stays
%  interpretive; grammar and dash punctuation repaired; field-not-fixed
%  bound pointed at sec:self-reference.
%
%  CITATIONS: existing keys Hardy2001a, Aaronson2004h, mosseri2001hopf.
%  NEW keys (RIS in L2_metric_refs_2026-07-06.ris): Banach1932,
%  Lamperti1958, Hopf1931, Urbantke2003, GluckWarnerZiller1986,
%  wooters_1981 (key spelling preserved as used; note the single-t).
%  British English; \autocite; no bullets in prose.
% =====================================================================

\section{The Euclidean metric on the state sphere, and the global phase as self-referential ignorance}
\label{sec:l2-metric-fibre}
Popper's systems $C^{+}$, $C$ and $E$ --- the predictor, its target and their shared
environment, whose argument \S\ref{sec:wigners-friend} anchors and quotes --- are, in modern
dress, Wigner, his friend and a qubit $S$ to be measured. The natural arena for analysing them
is the Hopf fibration and the Bloch sphere; but even to draw the Bloch sphere one must first
justify the $L^{2}$ metric, and this section earns that right. The principle doing the work is
the reversible continuity of \S\ref{sec:continuity} --- derived there from self-reference
conditionally, on the premises priced as Open Problem~\ref{op:limit-realisability}, of which one, that every
transformation of the regime is a phase of an inside structure, remains open; the continuum
itself is the completion the Shannon step declares (\S\ref{sec:complex-amplitudes}), the
classical convention, and \S\ref{sec:cencov-route} below reaches the same metric by a
second route that consumes it and nothing else that is open. Assuming that premise can be proved,
everything below inherits the derivation; we carry that status, and do not upgrade it.

We have established that the amplitudes of a single qubit form a complex vector space
$\mathbb{C}^{2}\cong\mathbb{R}^{4}$. We have not yet fixed which norm on this space singles out
the physical state manifold. That choice is not free: the principle of
\emph{reversible continuity} forces the Euclidean ($L^{2}$) norm, and hence the round
three-sphere, before any appeal to the Born rule. The quadratic structure usually \emph{posited}
as the Born rule is therefore already present in the metric, and the global phase emerges as a
geometrically forced fibre rather than an added convention.

\subsection{Reversible continuity selects the 2-norm}
\label{sec:reversible-continuity-l2}

The principle is the one isolated by Hardy as the single axiom distinguishing quantum from
classical probability \autocite{Hardy2001a}:
\emph{there exists a continuous reversible transformation between any two pure states.}
Operationally a reversible transformation of states is a symmetry of the state manifold, hence
norm-preserving; the demand is that such symmetries form a \emph{connected} group acting
\emph{transitively} on the pure states. Linearity is inherited from the complex vector-space
structure already in hand, so the admissible transformations are the complex-linear isometries
of whatever norm we place on $\mathbb{C}^{2}$.

Take the candidate norms to be the $\ell^{p}$ family,
$\lVert x\rVert_{p}=\bigl(\sum_{k}\lvert x_{k}\rvert^{p}\bigr)^{1/p}$, with the state manifold
the unit $\ell^{p}$ sphere. Their linear isometries are completely classified.


\begin{proposition}[Lamperti; after Banach]
\label{prop:lamperti}
For $1\le p<\infty$, $p\neq 2$, every linear isometry of $L^{p}$ (and hence of $\ell^{p}$) is a
generalised permutation: a weighted composition operator which, in the finite-dimensional case,
acts as $(x_{k})\mapsto(\omega_{k}\,x_{\sigma(k)})$ with $\lvert\omega_{k}\rvert=1$ and $\sigma$
a permutation of the coordinates.
\end{proposition}

\noindent Banach gave this for $L^{p}[0,1]$; the general $\sigma$-finite statement is Lamperti's
\autocite{Banach1932, Lamperti1958}. Its consequence for our four real dimensions is decisive.
For every $p\neq 2$ the isometry group is, in the real case, the finite group of signed
permutations (the hyperoctahedral group $B_{4}$), and in the complex case the group
$\mathbb{T}^{n}\rtimes S_{n}$ of diagonal phases composed with permutations. Even the complex
group, though it contains the continuous torus $\mathbb{T}^{n}$, fails to act transitively on
the sphere: the orbit of a basis state $e_{1}$ is only the phased basis vectors
$\{\omega\,e_{k}\}$, a measure-zero union of circles. No continuous reversible transformation
carries a basis state into a generic superposition such as $(e_{1}+e_{2})/\sqrt{2}$. Reversible
continuity therefore \emph{fails} for every $p\neq 2$.

At $p=2$ the isometry group is $O(4)$ in the real picture and $U(2)$ in the complex one; the
identity component acts transitively on the unit sphere, and any two unit vectors are joined by
a continuous one-parameter subgroup. Among the $\ell^{p}$ geometries, reversible continuity is
satisfiable at $p=2$ alone.

Aaronson states the same selection in the quantum-foundations setting and sharpens it into a
physical, rather than merely aesthetic, necessity \autocite{Aaronson2004h}: the $2$-norm is the
unique $p$-norm admitting a nontrivial continuous group of norm-preserving linear maps, every
other $p$ yielding only generalised permutations; and the alternative of keeping linearity and
reversibility while \emph{dropping} norm preservation lands one in the full general linear group
$\mathrm{GL}(n,\mathbb{C})$, which permits superluminal signalling and collapses the complexity
class $\mathsf{PP}$ into polynomial time --- the same fact \S\ref{sec:continuity} uses to price
the classical--quantum cut. Bare reversibility on the amplitude space buys only
$\mathrm{GL}(n,\mathbb{C})$; what carves the unitary group $U(n)$ out of it is preservation of
the quadratic ($\ell^{2}$) form, and reversibility of the resulting maps is then a corollary
rather than a premise.

The conclusion is that the physical state manifold of a single qubit is the unit sphere of the
Euclidean norm,
\begin{equation}
  S^{3}=\Bigl\{\,x\in\mathbb{R}^{4}\;:\;\textstyle\sum_{i=1}^{4}x_{i}^{2}=1\,\Bigr\}
  =\bigl\{\,\psi\in\mathbb{C}^{2}\;:\;\lvert z_{1}\rvert^{2}+\lvert z_{2}\rvert^{2}=1\,\bigr\},
\end{equation}
with $z_{j}=x_{2j-1}+\mathrm{i}\,x_{2j}$, carrying the round metric induced from the ambient
inner product,
\begin{equation}
  \mathrm{d}s^{2}=\sum_{i=1}^{4}\mathrm{d}x_{i}^{2}\,\Big|_{S^{3}}
  =\operatorname{Re}\langle\mathrm{d}\psi,\mathrm{d}\psi\rangle,
  \qquad
  \langle\phi,\psi\rangle=\bar\phi_{1}\psi_{1}+\bar\phi_{2}\psi_{2}.
\end{equation}
One bound deserves explicit statement, in the interest of claiming no more than is proved:
reversible continuity fixes the \emph{metric} as $L^{2}$, but not the underlying field. Real and
quaternionic quantum theory equally satisfy reversible continuity and equally carry Euclidean
state-space metrics; the restriction to $\mathbb{C}$ is pinned elsewhere in the framework ---
the analytic continuation of the information measure and the additivity of phase information,
\S\ref{sec:self-reference} --- not here.

\subsection{Hardy's independent route}
\label{sec:hardy-route}

The same conclusion is reached by Hardy along a route that never invokes
Proposition~\ref{prop:lamperti} \autocite{Hardy2001a}. His reconstruction turns on the counting
relation $K=N^{r}$, where $N$ is the number of states distinguishable in a single measurement
and $K$ the number of fiducial probabilities needed to specify a state. The continuity axiom
excludes the classical case $r=1$: there the pure states are the $N$ vertices of a simplex and
the reversible maps merely permute them --- a discrete group of motions --- so there is no
continuous path between distinct pure states, which is the simplex face of the same $\ell^{1}$
isometry fact. A simplicity assumption then selects $r=2$, yielding complex quantum theory with
its quadratic norm and trace rule. The quadratic structure is not extracted from continuity as a
separate step; it is built into Hardy's first axiom (probabilities linear in the fiducial set)
together with $K=N^{2}$.

We thus have two genuinely independent demonstrations that reversible continuity forces the
$L^{2}$ metric: Hardy's operational $K=N^{2}$ counting, and the Lamperti--Aaronson
classification of the isometries. They converge on the same step by different mathematics, and
should be cited as two roads rather than one road traversed twice. Both roads start from
continuity. A third, in the next subsection, starts one step earlier, from the value continuum
that continuity itself rests on, and reaches the same metric without passing through the
transformation group at all.

\subsection{The statistical route: \v{C}encov's theorem and the every-basis lemma}
\label{sec:cencov-route}

The value continuum that \S\ref{sec:continuity} isolates as the one open premise reaches the
$L^{2}$ metric on its own, by a route that uses no transformation group. \v{C}encov's theorem
states that the only Riemannian metric on the simplex of probability distributions invariant
under coarse-graining --- under the Markov maps that merge or relabel outcomes --- is the
Fisher--Rao metric \autocite{bengtsson2006geometry}, and in the coordinates $x_{k} =
\sqrt{p_{k}}$ the Fisher--Rao metric is the round metric on the positive orthant of the unit
sphere. The premise the theorem consumes is that the simplex is a manifold, which is the value
continuum declared in \S\ref{sec:complex-amplitudes}, and the principle it uses, that
distinguishability does not depend on how outcomes are grouped, is one the Shannon bookkeeper of
that section already grants. What it does not supply is the passage from one orthant to the
whole projective space; that is the following lemma.

\begin{lemma}[The every-basis lemma]\label{lem:every-basis}
Let $g$ be a Riemannian metric on $\mathbb{CP}^{K-1}$ whose restriction to the positive real
orthant of every orthonormal basis --- the states $\sum_{k}\sqrt{p_{k}}\,e_{k}$ with $p$ in the
simplex --- is the Fisher--Rao metric. Then $g$ is the Fubini--Study metric, whose linear lift
is the $L^{2}$ inner product.
\end{lemma}

\begin{proof}
A tangent vector to $\mathbb{CP}^{K-1}$ at a state $\psi$ is represented by a vector $\phi$
orthogonal to $\psi$, and the curve $t \mapsto \cos t\,\psi + \sin t\,\phi$ with real $t$ lies
in the positive orthant of any orthonormal basis that contains $\psi$ and $\phi$. Along that
curve the Fisher--Rao metric is the round metric of the sphere in the coordinates $\sqrt{p}$,
which is the Fubini--Study metric restricted to real amplitudes; so $g$ and the Fubini--Study
metric agree on every tangent vector at every point, the direction $i\phi$ being reached by the
basis that contains $i\phi$ in place of $\phi$.
\end{proof}

The lemma's premise, that every orthonormal basis is a legitimate measurement and none is
privileged, is the fourth face of the pre-measurement regime seen from the measurement side
(\S\ref{sec:continuity}); in the stabiliser fragment of quantum theory it fails in the same way
everything else does, only the Pauli bases being measurements there. So the $L^{2}$ metric now
rests on two legs that share exactly one premise. Through reversible continuity it uses the
value continuum, closure and homogeneity (Proposition~\ref{prop:lamperti}); through \v{C}encov
it uses the value continuum, coarse-graining invariance and the equal standing of bases
(Lemma~\ref{lem:every-basis}). Neither leg needs anything else that is open, and nothing could
avoid the continuum, since a metric is a statement about one. That is the status of everything
downstream of this section, the Hopf fibrations included: derived from the completion every physicist grants, which carries no quantum content, and from
one premise about transformations that the framework derives in outline and has not yet proved
(Open Problem~\ref{op:limit-realisability}).

\subsection{The Hopf fibration and the Bloch sphere}
\label{sec:hopf-bloch}

The round metric is exactly what makes the global phase a clean symmetry. The phase action
$\psi\mapsto\mathrm{e}^{\mathrm{i}\theta}\psi$ is a one-parameter subgroup of
$U(2)\subset SO(4)$ and so acts on $S^{3}$ by \emph{isometries}, a statement that holds only
because the metric is the $L^{2}$ one just derived. Its orbits are great circles $S^{1}$, all of
circumference $2\pi$, and the quotient is the complex projective line,
\begin{equation}
  S^{3}\big/U(1)=\mathbb{CP}^{1}\cong S^{2},
\end{equation}
the Bloch sphere. This quotient is the complex Hopf fibration
\autocite{Hopf1931, Urbantke2003, mosseri2001hopf}
\begin{equation}
  S^{1}\hookrightarrow S^{3}\xrightarrow{\;h\;}S^{2},
\end{equation}
whose projection sends an amplitude to its Bloch vector,
\begin{equation}
  h(\psi)=\bigl(\langle\psi\rvert\sigma_{x}\lvert\psi\rangle,\;
                 \langle\psi\rvert\sigma_{y}\lvert\psi\rangle,\;
                 \langle\psi\rvert\sigma_{z}\lvert\psi\rangle\bigr)
        =\bigl(2\operatorname{Re}\bar z_{1}z_{2},\;
               2\operatorname{Im}\bar z_{1}z_{2},\;
               \lvert z_{1}\rvert^{2}-\lvert z_{2}\rvert^{2}\bigr).
\end{equation}
The Hopf map is quadratic in the amplitudes, and this is the geometric face of the Born rule:
the third component is $p_{0}-p_{1}$ with $p_{k}=\lvert z_{k}\rvert^{2}$, and the squaring is
the lowest-degree $U(1)$-invariant available on the fibre --- the record-theoretic derivation of
the same rule is \S\ref{sec:born-rule}'s. Because the metric is round, the projection is a
Riemannian submersion with totally geodesic fibres of constant length
\autocite{GluckWarnerZiller1986}: the same amount of phase is hidden at every point of the base,
so the phase is hidden \emph{uniformly}, not in a position-dependent way that the metric could
partially betray.

\subsection{The fibre as self-referential ignorance}
\label{sec:fibre-self-ignorance}

The interpretation the framework places on this geometry --- and here we are reading physics
into the mathematics, so the claim is interpretive rather than proved --- is that the $S^{1}$
fibre is the observer's own un-narratable self-ignorance. The global phase is the imaginary part
of the complex information measure $S=-\ln\lvert p\rvert-\mathrm{i}\theta$: information the
observer lacks about their own state. Lacking it, they cannot communicate it; there is no
\textsc{fanout} of information one does not possess. Uncommunicable, it cannot enter any
measurement record --- and this last step is no longer interpretive:
Proposition~\ref{prop:f1-dictionary} proves the global phase absent from every classical record
a copying structure induces --- so every observable is $U(1)$-invariant; and a $U(1)$-invariant
quantity is a function on the quotient. The physical state is therefore the orbit, not the
representative point---the base $S^{2}$---and the $S^{1}$ fibre carries exactly the inaccessible
self-referential phase. It is a null direction of statistical distinguishability: two global
phases are indistinguishable by any measurement whatever \autocite{wooters_1981}, so the fibre
is a gauge direction without physical content. This phase is irrecoverably discarded, and must
remain so: it is the relative phase, living in the $S^{3}$ fibre of the quaternionic Hopf
fibration at two qubits, that is genuinely recoverable, whereas the single-qubit global phase of
the complex fibration is gauge-private without remainder. The self-reference that opened the
derivation thus reappears, at its end, as the circle the observer can never see around.

%      % EDITED 2026-09-12 scripts/patch_revert_sedenion_repoint_2026-09-12T1600.py
% EDITED 2026-09-11 scripts/patch_undefined_section_refs_2026-09-11T2030.py
% =====================================================================
%  SECTION: Wigner's friend in the Hopf picture (v2; supersedes the
%  2026-06-07 draft). AUDITED 2026-07-06 (final pass, fifth stop).
%
%  PLACEMENT (reordered chain, JS 2026-07-06): position 4 --
%      \input{L2_metric_global_phase_section_2026_06_07}
%      \input{wigners_friend_in_the_hopf_picture_2026-06-13}
%      \input{derivation_of_the_quantum_postulates}
%  sec:l2-metric-fibre is REAL (backward, order-consistent);
%  sec:fanout-boundary DRAFTED 2026-07-06 (formerly the one undrafted section);
%  the old sec:monogamy placeholder note was VOID -- no such \ref in
%  the body -- and is removed. prop:f1-dictionary and
%  sec:hopf-justification are FORWARD references (position 8): first
%  F1 use now carries the deliberate-forward flag.
%
%  DEPENDENCY: \input{monogamy_aside_2026-06-07} -- received from JS
%  and audited 2026-07-06; the \input sits inside
%  sec:wf-reconciliation, immediately after op:tsirelson-geometry
%  (relocated 2026-07-06 from end-of-file, where the aside's opening
%  antecedent had been orphaned). A bare replace-all during the earlier
%  patch briefly mangled this header block, leaving a live \input in
%  it; repaired by index-splice, and the lesson recorded: anchors that
%  also occur in comments must be uniqueness-asserted.
%
%  AUDIT RECORD: CHSH arithmetic (2sqrt(1+C^2), PR => r^2 = -2) and the
%  C^2+r^2=1 chain verified; F1 wording canonical; the
%  obsolescence-made-cancellation echo two-way with
%  sec:self-reference; redundancy/Darwinism/DiBiagio-Rovelli/Brukner
%  positioning sound; Landauer placed on amplification, correctly.
%  Fixes: first-person voice removed; \textquote replaced with
%  portable quotes; IOEP given its prop ref; stale Tsirelson
%  mini-header (wrong key name Tsirelson1980) retired;
%  op:tsirelson-geometry MINTED from the honesty paragraph's unpriced
%  owed. Popper page-pin DISCHARGED 2026-07-07: the zip bundles carry
%  OCR .txt sidecars (the earlier "no text layer" verdict was wrong --
%  ZIP-of-JPEGs defeated the PDF tools, not the text); the quoted
%  phrase sits at Part II, p. 187, now pinned in the cite. Bonus from
%  the same dig: "C+" is Popper's own notation (Part II, scan pp.
%  15-17), so the L2 opening's attribution stands as written; "E" not
%  sighted in the sampled pages -- the inline gloss covers it either
%  way. Keys unique to this file (presumed in the .bib,
%  since the paper compiles): VerstraeteWolf2002,
%  Cirelson1980QuantumInequality, PopescuRohrlich1994, Held2021,
%  Brukner2021, zurek_2009, Landauer_1961, wigner_1961.
%  British English; \autocite; no bullets in prose.
% =====================================================================

\section{Wigner's friend in the Hopf picture: one state, two charts, and the redundancy ladder}
\label{sec:wigners-friend}

The framework dissolves the Wigner's-friend paradox \autocite{wigner_1961} not by
adjudicating between the two accounts but by recognising that they are written in
two charts on a single manifold, selected as the problem requires, much as one
moves between cylindrical and spherical coordinates without either being the true
one. The epistemological chart (private, non-clonable, phase-carrying) and the
ontological chart (clonable, public, record-bearing) are related by the
\textsc{fanout} boundary (Section~\ref{sec:fanout-boundary}), and that boundary
is not a sharp wall but a graded one, parametrised by how redundantly a piece of
relational content has been copied. Read this way, the fact that a \textsc{fanout}
has occurred for the friend but not for Wigner is not a contradiction to be
removed; it is the boundary behaving exactly as defined.

Let the system be a qubit $\lvert\psi\rangle_{S}=\alpha\lvert 0\rangle
+\beta\lvert 1\rangle$ and let the friend $F$, idealised, with deliberate
economy, as a single qubit, measure it in the computational basis. The
interaction is the controlled coupling that copies $S$'s value into $F$'s
register: the \textsc{copy} comultiplication of the diagrammatic calculus
\autocite{CoeckeKissinger2017}, and in the circuit picture the very gate that, in
the two-qubit Mosseri--Dandoloff geometry \autocite{mosseri2001hopf},
activates the quaternionic Hopf fibration $S^{3}\to S^{7}\to S^{4}$ and carries
the joint state from the surface of the entanglement ball to its interior. The
friend and Wigner read the two ends of that plunge.

Two idealisations in this set-up should be named at once, because the argument
later depends on seeing past them. The first is the friend as a single qubit,
which we return to as the redundancy-one endpoint of a graded boundary
(Section~\ref{sec:wf-redundancy}). The second, less obvious and more consequential,
is the system $S$ as a pristine qubit with no prior entanglement --- a relation
with no history. No real measured particle is of this kind. A particle is itself a
reified relation between prior events, and by the time the friend couples to it, it
already carries the completed internal structure such a relation possesses: in the
framework's terms, the full octonionic $A,B,C$ register that the $S^{7}$ sector
supplies as internal degrees of freedom (\S\ref{sec:space-of-computations}). The
pristine qubit is a schematic convenience, drawn structureless so that the logic of
a single link is visible; the ``system, friend, Wigner'' triple is likewise the
schematic of \emph{one link} in the chain of observers, not its base. The chain
does not bottom out in a structureless qubit --- there is no such thing --- and we
draw on this when we ask, later, what bounds the information a single such link can
carry.

\subsection{The friend's chart: a \textsc{fanout} of redundancy one}
\label{sec:wf-friend}

Relative to $F$, the coupling is a \textsc{fanout}: $F$ acquires a fact. What is
copied is gauge-invariant, the relative phase, the correlation between $F$ and
$S$, which is a genuine observable. The global phase is never the thing copied:
there is nothing observable there to duplicate, so the $S^{1}$ fibre of the
complex Hopf fibration $S^{1}\to S^{3}\to S^{2}$ of
Section~\ref{sec:l2-metric-fibre} remains $F$'s gauge-private self-ignorance,
untouched. In Hopf terms the division is clean, but it lives at two levels, which
Proposition~\ref{prop:f1-dictionary} (\S\ref{sec:hopf-justification} --- a deliberate forward
reference) separates. The copy operation itself is fibre-committed:
copying is faithful to phase, so what $F$'s \textsc{fanout} duplicates is a lift --- a chosen
fibre-point over the outcome --- never a bare point of the base. Yet everything the operation can
broadcast is base data alone: the record it induces is independent of the lift, so the $S^{1}$
freedom, though pinned in the copier, is absent from every classical record it yields. $F$'s classical record is a
base-point; the structure that produced it committed to a fibre-point; and the gap between the two
is precisely the gauge character of the fibre. The observation has a fifty-year-old
anchor: Popper argued that a predictor's inferior self-knowledge follows from the logical
impossibility of a system's imparting information about itself without thereby rendering it out of
date --- ``making this information obsolete'', in his words
\autocites{Popper1950a}[Part~II, p.~187]{Popper1950b}. Proposition~\ref{prop:f1-dictionary}
makes the impossibility structural rather than temporal, and localises it: the datum the copier
cannot put into its own record is not its state wholesale but exactly its fibre commitment, which
does not arrive stale --- it cancels. Conditioned on the outcome it reads, $S$ is for
$F$ a definite point on the Bloch sphere,
\begin{equation}
  \rho_{S}^{(F)}=\lvert 0\rangle\langle 0\rvert
  \quad\text{(given outcome }0\text{)},
\end{equation}
a pure state of Bloch radius $r=1$ on the surface of the ball, concurrence $C=0$.

This is a \textsc{fanout}, but the minimal one. The fact $F$ holds is a
\emph{relative} fact in the sense of relational quantum mechanics
\autocite{Rovelli1996a}: $S$ definite relative to $F$, not an observer-independent
record. It is one copy, redundancy $R=1$, and as such it remains reversible by an
agent who controls the joint system. The single-qubit friend is therefore not
outside \textsc{fanout}, as a careless reading might suggest; it is the $R=1$
endpoint of the \textsc{fanout} continuum---a \textsc{fanout} that has happened,
and a fact that is real for $F$, while being maximally still reversible from
outside.

\subsection{Wigner's chart: no copy has reached him}
\label{sec:wf-wigner}

Wigner treats the laboratory as closed and no copy of the relational content has
reached him. For him the coupling is the unitary
\begin{equation}
  \lvert\Psi\rangle
  =\alpha\lvert 0\rangle\lvert F_{0}\rangle
  +\beta\lvert 1\rangle\lvert F_{1}\rangle,
\end{equation}
the move up to the larger Hilbert space, $\mathbb{C}^{2}\to\mathbb{C}^{4}$,
$S^{3}\to S^{7}$, complex Hopf to quaternionic Hopf, with the fibre grown from
$S^{1}$ to $S^{3}$. What $F$ underwent as the acquisition of a fact, Wigner
describes as correlation deposited into a joint pure state on $S^{7}$. When he
asks specifically about $S$, he traces $F$ out, and the geometry does the
bookkeeping:
\begin{equation}
  \rho_{S}^{(W)}=\operatorname{Tr}_{F}\lvert\Psi\rangle\langle\Psi\rvert
  =\lvert\alpha\rvert^{2}\lvert 0\rangle\langle 0\rvert
  +\lvert\beta\rvert^{2}\lvert 1\rangle\langle 1\rvert,
  \qquad
  r=\bigl\lvert\,\lvert\alpha\rvert^{2}-\lvert\beta\rvert^{2}\bigr\rvert,
  \quad
  C=2\lvert\alpha\beta\rvert .
\end{equation}
Tracing out discards the $S^{3}$ fibre and carries the pure point on $S^{7}$ to a
mixed point in the interior of the Bloch ball at radius $r=\sqrt{1-C^{2}}$
\autocite{mosseri2001hopf}. The inward displacement is governed by the
concurrence: it is the relational content that has migrated into the $S$--$F$
correlation, stored in the recoverable $S^{3}$ fibre. Nothing has been destroyed;
it has been relocated into the fibre, which is why Wigner can in principle
interfere the branches and reverse the interaction, the content sits at low
redundancy and no further copy has fixed it.

\subsection{One state, two charts: \texorpdfstring{$C^{2}+r^{2}=1$}{C2+r2=1}}
\label{sec:wf-reconciliation}

The system is pure for the friend ($r=1$, on the surface) and mixed for Wigner
($r<1$, in the interior). These are not rival claims about a single fact; they are
the friend at the surface and Wigner at the interior of the \emph{same} state on
$S^{7}$, bound by the Mosseri--Dandoloff constraint
\begin{equation}
  C^{2}+r^{2}=1,
\end{equation}
in which the discrepancy in Bloch radius is exactly the concurrence Wigner
retains and the friend has resolved away. This is the Inside/Outside Equivalence
Principle (Proposition~\ref{prop:inside-outside}) made concrete, the same law written by the embedded participant and by
the external observer, and the relativity of facts of relational quantum
mechanics \autocite{Rovelli1996a} realised in geometry rather than asserted: the
\textsc{fanout} boundary lies \emph{between} $F$ and $W$, so a fact exists
relative to $F$ and not relative to $W$. Neither chart is privileged; one reads in
the epistemological or the ontological chart as the question requires, and their
mutual incompleteness is a theorem of the construction, not a defect of either.

\medskip
\noindent\emph{Aside: the Tsirelson bound and the absence of PR boxes.}
\label{sec:tsirelson}
A ball has a unique centre and no point beyond it. Maximal entanglement sits at
that centre \autocite{mosseri2001hopf}, $r=0$, $C=1$, and because the
Bloch radius cannot be negative, the concurrence cannot exceed unity. For a
two-qubit pure state the maximal CHSH correlator is
\begin{equation}
  \langle\mathrm{CHSH}\rangle_{\max}=2\sqrt{1+C^{2}}=2\sqrt{2-r^{2}}  
\end{equation}
\autocite{VerstraeteWolf2002},
which the centre maximises to $2\sqrt{2}$, the Tsirelson bound
\autocite{Cirelson1980QuantumInequality}. A Popescu--Rohrlich box \autocite{PopescuRohrlich1994},
with CHSH value $4$, would require $C^{2}=3$ and hence $r^{2}=-2$: a point two
units past the centre, which the ball does not contain. The super-quantum
correlation is excluded by the very positivity that places maximal entanglement
at the centre, one cannot trace out to a more-than-maximally-mixed reduced
state. The Tsirelson bound is in this sense the image, under the CHSH map, of the
boundedness of the $L^{2}$ state ball derived in
Section~\ref{sec:l2-metric-fibre}, and the prohibition on PR boxes shares its
origin with the roundness of the sphere: both trace to the reversible-continuity
metric, not to a separate principle.

\noindent This is the two-qubit, pure-state face of the bound, and its status is
worth stating plainly. It exhibits PR correlations as having no representing state
in the Hilbert geometry, and it recovers the value $2\sqrt{2}$ at the centre; but
forcing \emph{exactly} $2\sqrt{2}$ as a theorem of the geometry, rather than
reading it from the known concurrence--CHSH relation, and lifting the argument
from two qubits to the general Tsirelson bound, remain to be made rigorous. The
closest cognate in the literature is Held's derivation of the same bound from the
geometric algebra $G^{3}\cong\mathbb{H}$ \autocite{Held2021}, the same
quaternionic structure in a different language.

\begin{openproblem}[The Tsirelson bound as geometry]\label{op:tsirelson-geometry}
Derive $2\sqrt{2}$ as a theorem of the round state geometry rather than reading it off the
concurrence--CHSH relation, and lift the argument from the two-qubit pure case to the general
Tsirelson bound --- so that the exclusion of super-quantum correlation becomes, like the exclusion
exhibited above, a consequence of the $L^{2}$ ball's positivity (\S\ref{sec:l2-metric-fibre}) and
of nothing else. Held's quaternionic derivation \autocite{Held2021} is the nearest cognate and the
natural first tool.
\end{openproblem}

\input{monogamy_aside_2026-06-07}

\subsection{The redundancy ladder: from relative fact to objective fact}
\label{sec:wf-redundancy}

The friend's fact becomes a fact for Wigner only when Wigner copies it when he
measures $F$. It becomes a fact for a third party only when that party copies it
from Wigner, and so on up the von Neumann chain. The point at which the content
stops being negotiable is therefore not a point but a limit: once it has been
copied into enough mutually independent fragments that many parties can read it
without disturbing it, no one can practically un-copy it, and it is operationally
an objective fact. This is exactly quantum Darwinism \autocite{zurek_2009}:
objectivity is redundancy. It is not an external graft, because the framework's
account of decoherence already \emph{is} this picture---many fragments each take
their own copy of the clonable content, while their private phases, each
fragment's own self-ignorance, fail to add coherently and wash out.

This supplies the grading parameter that ``graded \textsc{fanout}'' had left
unnamed. \textsc{fanout} is graded not because the operation is fuzzy but because
redundancy $R$ runs continuously: $R=1$ is the single relational copy---the
friend's relative fact, reversible by Wigner---and $R\gg 1$ is the objective
fact, effectively irreversible because un-copying it would mean recovering every
fragment. The single-qubit friend is the $R=1$ endpoint; the macroscopic friend
is the same operation at large $R$. One axis, two ends. The relation to
Di~Biagio and Rovelli's distinction \autocite{DiBiagioRovelli2021} is then exact:
their \emph{stable} facts are the high-$R$ limit and their \emph{relative} facts
the low-$R$ regime, and the standard conditional-probability law they use to
define stability is the operational signature of redundancy having grown large.
Brukner's no-go theorem \autocite{Brukner2021}, that qubits are not observers,
sits on the same axis without conflict: it bites only against calling an $R=1$
relative fact an objective one, which the framework never does---the redundancy
ladder is precisely the structure that separates the relational definiteness a
single qubit does carry from the objectivity it does not.

The thermodynamics tracks the same axis. The single relational copy is unitary and
free; the $kT\ln 2$ cost \autocite{Landauer_1961} is incurred not by the friend's
act but by the amplification and stabilisation that drive $R$ upward, when
fragments are irreversibly committed. The relative fact is thermodynamically free;
the objective fact is what costs $kT\ln 2$.

\subsection{Phase promotion, and the bump in the carpet}
\label{sec:wf-phase-promotion}

One point keeps the load-bearing irrecoverability intact while explaining what
Wigner can do that the friend cannot. What Wigner recovers is \emph{not} the
$S^{1}$ self-ignorance; the global phase remains gauge for every observer and that
circle is spent for good. The quantity the friend resolves away is the relative
phase between the $\lvert 0\rangle$ and $\lvert 1\rangle$ components, and moving up
to $\mathbb{C}^{4}$ \emph{promotes} it: the single-qubit relative phase becomes the
relative phase between $\lvert 0\rangle\lvert F_{0}\rangle$ and
$\lvert 1\rangle\lvert F_{1}\rangle$, which lives in the recoverable $S^{3}$ fibre
(itself a circle bundle over $S^{2}$, so the promotion adds a two-sphere of
relative-phase directions to the lone gauge circle). Gauge phase at one level
becomes physical, interferable phase one level up---which is why Wigner can run
the interference experiment the friend cannot. But Wigner now carries his own
fresh global-phase $S^{1}$ at the top of the joint state: his own irreducible
self-ignorance. Ascending the tower does not eliminate the self-reference, only
displaces it one level higher, and the regress up the chain is precisely the
redundancy ladder of the previous subsection. This is the framework's bump in the
carpet: each observer sits atop a Hopf tower with a single gauge circle he can
never see around, and a measurement relocates the next observer's blind circle
into his own recoverable fibre without ever closing the regress.

Two cautions close the section. First, the recoverability Wigner
enjoys is the low-redundancy limit: for a macroscopic, high-$R$ friend the
reversal is exponentially hard and the fact is objective for all practical
purposes, so Wigner's ``no record'' is the $R\!\approx\!1$ idealisation of a
graded boundary, not a free lunch. Second, naming redundancy as the grading
parameter of \textsc{fanout}, and the phase-promotion reading itself, are
syntheses of the framework's parts: each rests on established pieces (the
clonability criterion, quantum Darwinism, the tensor-product promotion of phase)
but is assembled here, and is offered as the cleanest statement of why the two
observers never disagree about anything both can check.
% =====================================================================
%  SPLICE: "Where the chain goes next" -- forward-looking coda for
%  sec:wigners-friend, pointing to sec:sedenions-generations.
%  Created 2026-08-27 per JS.
%
%  THE POINT (JS, 2026-08-27). The Wigner's-friend charts stop at
%  Wigner because the Hopf tower runs out at the octonions, while the
%  text says the chain of observers obviously continues. The sedenions
%  show what the continuation adds: the observer of Wigner (fourth
%  wire, C (x) S) brings a NEW piece of structure -- the finite S_3 of
%  Brown's theorem, the generation label -- and the observer of the
%  observer of Wigner (fifth wire, the trigintaduonions) brings nothing
%  new in kind. This coda records that, so the section's termination at
%  Wigner reads as a fact about the tower rather than an economy of the
%  exposition.
%
%  WHAT THE ALGEBRA ACTUALLY SAYS (verified 2026-08-27, notebook T --
%  notebook_T_trigintaduonion_automorphisms_2026-08-27.py):
%    - Aut(S) = G_2 x S_3 (Brown 1967), DIRECT product; the new S_3 is
%      generated by the doubling involution and Brown's explicit
%      order-three map psi, which fixes 1 and the doubling unit and
%      rotates every (a, a u) plane, a imaginary, by a third of a turn.
%    - psi permutes the three octonion HALVINGS O, psi O, psi^2 O of S
%      (each doubled by the same unit; meeting only in R). It does NOT
%      permute the three copies through the quaternion core -- those are
%      cycled by a lifted G_2 element (the colour-line three-cycle); see
%      the sedenion section's two-triples remark and the corrections
%      ledger of the same date.
%    - Brown: Aut(A_{n-1}) x S_3 = Aut(A_n) with EQUALITY for n = 4, 5, 6
%      (as quoted in gresnigtgourlayvarma2023three, sec. 4.3); Eakin and
%      Sathaye prove Brown's conjecture in general (VERIFY-CITE for the
%      exact general form -- their step is Aut(A_{n+1}) = Aut(A_n) x G,
%      G in {Z/2, S_3}; over R with the standard parameters the branch is
%      S_3 at every rung >= 4, which is what notebook T checks directly
%      at rungs 4, 5, 6). Derivation algebra g_2 from the octonions on
%      (Schafer 1954; notebook T: dim Der = 14 at n = 3, 4, 5).
%  So "nothing new at the fifth wire" is stated PRECISELY: the fifth
%  doubling adds another finite S_3 and nothing else -- no connected
%  symmetry, no fibration; the increment is the same at every rung
%  thereafter. Not "silent": REPETITIVE. The prose says so.
%
%  STATUS: the algebra is ESTABLISHED (Brown; Eakin--Sathaye; Schafer)
%  and directly checked (notebook T). The reading of the repeated finite
%  factor as what the redundancy ladder consumes is a RHYME, flagged.
%  The reading of the fourth rung's discrete orientation as the
%  generation label is the sedenion section's CANDIDATE identification,
%  pointed to, not re-argued here.
%
%  PLACEMENT: a new \subsection at the END of sec:wigners-friend, after
%  the "Two cautions" paragraph. Because that paragraph reads "Two
%  cautions close the section.", the patch also changes that one word:
%  "close the section" -> "close the argument". Both anchors
%  uniqueness-asserted (the July lesson: anchors that also occur in
%  comments must be asserted unique).
%
%  CROSS-REFS OUT (all exist in the synced source except the two
%  MAP-LABELs):
%    sec:hopf-justification   -- the "two climbs" paragraph (length vs
%                                termination); this coda sharpens it
%    sec:wf-redundancy         -- the redundancy ladder (this file)
%    sec:wf-phase-promotion    -- phase promotion (this file)
%    sec:sedenions-generations -- the sedenion section (2026-08-27 rev.)
%    prop:termination          -- Hurwitz/Adams termination
%  MAP-LABELs:
%    app:notebook-record       -- the notebook-record appendix
%                                (appendices/appendix_notebook_record_
%                                2026-07-11.tex); substitute its label.
%  NEW KEYS (RIS delivered: cayley_dickson_automorphism_refs_2026-08-27.ris):
%    eakinsathaye1990automorphisms, schafer1954cayleydickson
%  IN-LIBRARY: brown1967generalized.
%
%  OPTIONAL CONSUMER-POINTER (not patched; one clause for the "two
%  climbs" paragraph of hopf_justification_section_2026-06-08.tex, after
%  "that chain has no top --- it is the bump in the carpet, unbounded in
%  length."):
%      "(unbounded in length, and --- as \S\ref{sec:wf-chain-type}
%      records --- bounded in type: after the fourth observer nothing
%      new in kind arrives)"
%
%  EDITED 2026-08-30 (per JS): one sentence added to the fourth-wire
%  paragraph anchoring the flavour sector's irremovable phase (the
%  CP-phase recognition) on this rung, with the development at
%  sec:sedenions-mixing (2026-08-30 revision). No new keys here.
%  British English; pure ASCII; \autocite; no bullets; spaced dashes.
% =====================================================================
% SPLICE-BEGIN wf-chain-type 2026-08-27

\subsection{Where the chain goes next: one more rung, then copies}
\label{sec:wf-chain-type}

The regress just described is unbounded in length, and nothing here shortens
it: the observer of Wigner carries a gauge circle of his own, and so does the
observer of that observer, as \S\ref{sec:hopf-justification} insists when it
separates the endlessness of the chain from the termination of the tower. But
the algebra says how much of the endless chain adds structure of a new kind,
and the answer is exact: one more step, and then none.

The count runs as follows. The triple system, friend, Wigner is three wires,
and three wires are the three Hopf rungs $\mathbb{C}$, $\mathbb{H}$,
$\mathbb{O}$ --- the completed link of the earlier caveat --- which is why the
charts above stop at Wigner: the chart of whoever looks in on the triple,
$S^{15}\to S^{8}$, is the last fibration the division algebras supply
(Proposition~\ref{prop:termination}). The observer of Wigner is a fourth wire,
and the fourth wire is the sedenions,
$(\mathbb{C}^{2})^{\otimes 4} = \mathbb{C}\otimes\mathbb{S}$
(\S\ref{sec:sedenions-generations}). No fibration comes with it --- the
sedenions have zero divisors and admit no map $S^{31}\to S^{16}$ --- so the
fourth observer gains no recoverable fibre of the tower's kind: the continuous
promotion of relative phase into a larger fibre, the mechanism of the previous
subsection, has no fourth instance. What the doubling adds instead is finite.
Brown's theorem gives $\mathrm{Aut}(\mathbb{S}) = G_{2}\times S_{3}$
\autocite{brown1967generalized}: the identity component is unchanged from the
octonions, and the entire gain is a component group, generated by an
order-three automorphism that fixes the new doubling direction and rotates
every imaginary direction of the completed link into its copy by a third of a
turn. The relative orientation of the link and its copy, which at the Hopf
rungs was a continuous phase promoted into a fibre, survives at the fourth rung
only as a discrete choice among three, and that choice is the generation label
(\S\ref{sec:sedenions-generations}). The fourth observer, then, sees no new
force and no new fibre; he sees the same link three ways, and cannot gauge the
difference, because a component group has no Lie algebra. His rung is also where
the flavour sector's one irremovable phase finds its anchor: three generations
are the threshold at which the mixing matrix first carries a phase no convention
removes, and \S\ref{sec:sedenions-mixing} reads that phase --- through the
paper's negativity--contextuality chain --- as this wire's contextual cargo,
real only from one rung above.

The observer of the observer of Wigner is a fifth wire, the trigintaduonions,
and here the tower has nothing new to say. Brown proved that the next two
doublings each add exactly one further $S_{3}$ and nothing else, Eakin and
Sathaye that Brown's pattern continues at every rung, and the derivation
algebra --- the Lie algebra of the identity component --- is $\mathfrak{g}_{2}$
from the octonions onward
\autocite{brown1967generalized,eakinsathaye1990automorphisms,schafer1954cayleydickson};
the fifth and sixth rungs are also checked directly in the notebook record
(Appendix~\ref{app:notebook-record}). So the fifth observer's doubling repeats
the fourth's gift --- another third-of-a-turn symmetry between the four-wire
structure and its copy, commuting with everything below --- and adds no
connected symmetry and no fibration. Nothing new in kind arrives at the fifth
wire or at any wire after it. The increments of the tower are, in order, a
force, a force, a label, and then the same finite permutation of copies at
every step. That is the precise sense in which the chain ``obviously
continues'': it continues, and after the fourth observer it has nothing further
to announce. The framework reads the repeated factor as what the redundancy
ladder of \S\ref{sec:wf-redundancy} consumes --- a permutation symmetry among
copies, which is the freedom a record has and a type does not --- and records
that reading as a structural resonance rather than a derivation.

The two climbs of \S\ref{sec:hopf-justification} are thereby sharpened. The
chain of observers is unbounded in length and bounded in type: three forces by
the third wire, one family label by the fourth, copies thereafter. The observer
of Wigner is the last observer whose arrival changes what there is.

% SPLICE-END wf-chain-type 2026-08-27

% =====================================================================
%  FIGURE: the Wigner's-friend chain mapped onto the Cayley--Dickson
%  ladder. Created 2026-08-30 per JS: the Hopf-fibration view of
%  Wigner's friend has been load-bearing for months without being drawn;
%  the fibrations themselves are DEFERRED to Szangolies's introduction
%  (szangolies2025standardmodel, in library), and this figure carries
%  only the mapping.
%
%  CONTENT (left to right, one panel per chart, staircase rising):
%    Panel 1: F's chart of S            -- S^3 -> S^2, fibre S^1,  C, one wire
%    Panel 2: W's chart of {S,F}        -- S^7 -> S^4, fibre S^3,  H, two wires
%    Panel 3: O_W's chart of {S,F,W}    -- S^15 -> S^8, fibre S^7, O, three wires
%             (the last fibration)
%    Panel 4: the next observer's chart of {S,F,W,O_W} -- C^16 = C (x) S,
%             no fibration (S^31 -/-> S^16); the doubling's S_3 badge:
%             the generation label, a third of a turn.
%  Each describer perches on his sealed box with a dashed circle over
%  his head (his own gauge circle -- the self-referential ignorance he
%  cannot see around); dashed promotion arrows carry each circle into
%  the next chart's fibre ("gauge here -> interferable one rung up",
%  sec:wf-phase-promotion); the third arrow terminates at "no fourth
%  fibration", with what survives being the discrete third-of-a-turn
%  (sec:wf-chain-type, sec:sedenions-generations). The regress
%  continues to the right: copies only.
%
%  PREAMBLE REQUIREMENT (confirm in main.tex; add if absent):
%      \usepackage{tikz}
%      \usetikzlibrary{arrows.meta}
%  PLACEMENT: \input by patch_wigner_hopf_ladder_fig_2026-08-30.py,
%  after the two-idealisations paragraph of sec:wigners-friend, with
%  the accompanying deferral paragraph (in the patch).
%  CROSS-REFS OUT (all real): sec:wf-phase-promotion,
%  sec:wf-chain-type, sec:sedenions-generations.
%  KEY: szangolies2025standardmodel (in library).
%  British English; pure ASCII; \autocite in the caption.
% =====================================================================
\begin{figure}[t]
\centering
\resizebox{\textwidth}{!}{%
\begin{tikzpicture}[>=Stealth, line cap=round, line join=round,
    every node/.style={font=\scriptsize},
    blind/.style={densely dashed, draw=black!75},
    promote/.style={->, densely dashed, black!60, semithick},
    sight/.style={->, dotted, black!55}]
% ---- stick-figure pic --------------------------------------------------
\tikzset{stick/.pic={
  \draw[pic actions] (0,0.44) circle (0.085);
  \draw[pic actions] (0,0.355) -- (0,0.10);
  \draw[pic actions] (-0.12,0.20) -- (0,0.29) -- (0.12,0.20);
  \draw[pic actions] (-0.10,-0.10) -- (0,0.10) -- (0.10,-0.10);
}}
% ---- the staircase (Cayley--Dickson ladder) ---------------------------
\draw[black!35, thick]
  (-0.35,-0.55) -- (3.27,-0.55) -- (3.27,0.25) -- (6.89,0.25)
  -- (6.89,1.05) -- (10.51,1.05) -- (10.51,1.85) -- (13.60,1.85)
  -- (13.60,2.65) -- (16.20,2.65);
\node[anchor=north] at (1.25,-0.60) {$\mathbb{C}$ --- one wire};
\node[anchor=north] at (4.87,0.20)  {$\mathbb{H}$ --- two wires};
\node[anchor=north] at (8.49,1.00)  {$\mathbb{O}$ --- three wires};
\node[anchor=north] at (12.11,1.80) {$\mathbb{S}$ --- four wires};
\node[anchor=north, black!55] at (14.45,2.60) {$\mathbb{T},\dots$ --- five wires on};
% =======================================================================
% Panel 1: F's chart of S
% =======================================================================
\draw[rounded corners=2pt, fill=black!5, draw=black!40] (0,0) rectangle (2.5,1.32);
\filldraw[black!70] (1.25,0.72) circle (0.075);
\node[font=\tiny] at (1.25,0.45) {$S$};
\pic[black] at (2.05,1.42) {stick};
\node[anchor=west] at (2.22,1.62) {$F$};
\draw[sight] (1.95,1.60) -- (1.45,1.02);
\draw[blind] (2.05,2.14) circle (0.11);
\node[anchor=east, align=right, font=\tiny, black!75] at (1.82,2.14)
  {his own gauge circle:\\ the phase he cannot\\ see around};
\node[anchor=south west, font=\tiny, black!75] at (0.02,1.34) {fibre $S^{1}$};
\node[anchor=north] at (1.25,-0.06) {$S^{3}\to S^{2}$};
% =======================================================================
% Panel 2: W's chart of {S,F}
% =======================================================================
\draw[rounded corners=2pt, fill=black!5, draw=black!40] (3.62,0.80) rectangle (6.12,2.12);
\filldraw[black!70] (4.07,1.52) circle (0.075);
\node[font=\tiny] at (4.07,1.25) {$S$};
\pic[black!55, scale=0.85] at (4.85,1.16) {stick};
\node[font=\tiny, black!55] at (4.85,0.94) {$F$};
\pic[black] at (5.67,2.22) {stick};
\node[anchor=west] at (5.84,2.42) {$W$};
\draw[sight] (5.57,2.40) -- (5.05,1.82);
\draw[blind] (5.67,2.94) circle (0.11);
\node[anchor=south west, font=\tiny, black!75] at (3.64,2.14) {fibre $S^{3}$};
\node[anchor=north] at (4.87,0.74) {$S^{7}\to S^{4}$};
\draw[promote] (2.17,2.14) to[out=15,in=170] (3.60,2.24);
\node[anchor=south, align=center, font=\tiny, black!60] at (2.85,2.32)
  {gauge below,\\ interferable above};
% =======================================================================
% Panel 3: O_W's chart of {S,F,W}
% =======================================================================
\draw[rounded corners=2pt, fill=black!5, draw=black!40] (7.24,1.60) rectangle (9.74,2.92);
\filldraw[black!70] (7.69,2.32) circle (0.075);
\node[font=\tiny] at (7.69,2.05) {$S$};
\pic[black!55, scale=0.85] at (8.35,1.96) {stick};
\node[font=\tiny, black!55] at (8.35,1.74) {$F$};
\pic[black!55, scale=0.85] at (8.95,1.96) {stick};
\node[font=\tiny, black!55] at (8.95,1.74) {$W$};
\pic[black] at (9.29,3.02) {stick};
\node[anchor=west] at (9.46,3.22) {$O_{W}$};
\draw[sight] (9.19,3.20) -- (8.65,2.62);
\draw[blind] (9.29,3.74) circle (0.11);
\node[anchor=south west, font=\tiny, black!75] at (7.26,2.94) {fibre $S^{7}$};
\node[anchor=north, align=center] at (8.49,1.54) {$S^{15}\to S^{8}$\\[-1pt]
  {\tiny the last fibration}};
\draw[promote] (5.79,2.94) to[out=15,in=170] (7.22,3.04);
% =======================================================================
% Panel 4: the next chart of {S,F,W,O_W} -- sedenions, no fibration
% =======================================================================
\draw[rounded corners=2pt, fill=black!5, draw=black!40] (10.86,2.40) rectangle (13.36,3.72);
\filldraw[black!70] (11.24,3.12) circle (0.075);
\node[font=\tiny] at (11.24,2.85) {$S$};
\pic[black!55, scale=0.80] at (11.78,2.78) {stick};
\node[font=\tiny, black!55] at (11.78,2.56) {$F$};
\pic[black!55, scale=0.80] at (12.28,2.78) {stick};
\node[font=\tiny, black!55] at (12.28,2.56) {$W$};
\pic[black!55, scale=0.80] at (12.78,2.78) {stick};
\node[font=\tiny, black!55] at (12.78,2.56) {$O_{W}$};
\pic[black!45] at (14.35,2.78) {stick};
\draw[sight] (14.18,3.02) -- (13.42,3.12);
\draw[blind] (14.35,3.50) circle (0.11);
\node[anchor=west, align=left, font=\tiny, black!60] at (14.58,3.48)
  {the next observer,\\ and all after:\\ copies only};
\node[anchor=north] at (12.11,2.34) {$\mathbb{C}^{16}=\mathbb{C}\otimes\mathbb{S}$};
% the broken third promotion
\draw[promote] (9.41,3.74) to[out=15,in=185] (10.68,3.86);
\node[font=\scriptsize, black!60] at (10.80,3.92) {$\times$};
\node[anchor=south west, font=\tiny, black!75] at (10.88,3.74)
  {no fourth fibration: $S^{31}\not\to S^{16}$};
% the S3 badge: three dots at 120 degrees with a rotation arrow
\draw[black!75] (11.45,4.30) circle (0.17);
\filldraw[black!75] (11.45,4.47) circle (0.028);
\filldraw[black!75] (11.303,4.215) circle (0.028);
\filldraw[black!75] (11.597,4.215) circle (0.028);
\draw[->, black!75] (11.67,4.36) arc[start angle=20, end angle=100, radius=0.235];
\node[anchor=west, align=left, font=\tiny, black!75] at (11.70,4.36)
  {$S_{3}$: the generation label\\ (a third of a turn)};
\draw[promote, black!70] (10.86,3.96) to[out=60,in=200] (11.24,4.22);
\end{tikzpicture}%
}
\caption{The Wigner's-friend chain mapped onto the Cayley--Dickson ladder. Each
observer holds a chart of the sealed box below him, one rung further up the
ladder: the friend's chart of the system is the complex Hopf bundle, Wigner's
chart of the sealed laboratory the quaternionic, and the chart of the observer
of Wigner --- of the completed triple --- the octonionic, the tower's last
fibration. Every chart-holder carries a dashed circle over his head, his own
gauge phase: the self-referential ignorance he cannot see around. Measurement
promotes each circle into the next chart's fibre, where what was gauge below is
interferable above (\S\ref{sec:wf-phase-promotion}); the third promotion has no
fibration to land in --- the sedenions have zero divisors and no map
$S^{31}\to S^{16}$ --- and what survives of it is discrete: the doubling's
$S_{3}$, a third of a turn rotating the link into its copy, the generation
label (\S\ref{sec:wf-chain-type}, \S\ref{sec:sedenions-generations}). The
regress of circles continues unbounded to the right, but from the fifth wire
onward it adds copies only. The friend drawn as a single qubit is the
redundancy-one idealisation discussed in the text; for the fibrations
themselves, deferred rather than repeated here, see the introduction in
\autocite{szangolies2025standardmodel}.}
\label{fig:wf-hopf-ladder}
\end{figure}

%      % EDITED 2026-10-29 scripts/patch_continuity_tower_premise_2026-10-29T1500.py
% EDITED 2026-09-20 scripts/patch_born_cost_ledger_v2_2026-09-20T0900.py
% EDITED 2026-09-15 scripts/patch_value_continuum_declared_cencov_2026-09-15T1540.py
% EDITED 2026-09-15 scripts/patch_continuity_connected_state_space_2026-09-15T1430.py
% =====================================================================
%  SECTION: Derivation of the quantum postulates (the assembly chapter).
%  AUDITED 2026-07-06 (final pass, sixth stop).
%
%  PLACEMENT (reordered chain, JS 2026-07-06): position 5 --
%      \input{wigners_friend_in_the_hopf_picture_2026-06-13}
%      \input{derivation_of_the_quantum_postulates}
%      \input{division_algebras_section_2026-07-03}
%  Postulate 6 hands off to division-algebras, its immediate successor.
%
%  LABEL RESOLUTIONS (2026-07-06): the old placeholder
%  sec:complex-amplitudes is now minted at its source (self-reference's
%  negativity subsection), so the ref here resolves unchanged.
%  sec:born-rule COLLIDED with the label minted 2026-07-06 on
%  self-reference's Born-from-fibre-privacy subsection; resolution: two
%  routes, two labels -- Postulate 3 here is renamed
%  sec:born-invariance (the U(1)/lowest-degree geometric route) and
%  cross-stitched to sec:born-rule (the record-theoretic route), their
%  agreement flagged as over-determination. sec:local-tomography also
%  minted at source for Postulate 6's operational clause.
%
%  AUDIT RECORD: the section's status discipline was PRESCIENT -- the
%  continuity arrow was already "owed ... enters as a principle" (the
%  exact voice JS mandated for L2), and both of its confessed debts
%  have since been minted as numbered problems; the closing now cites
%  them (op:limit-realisability, op:tsirelson-geometry): "named, and
%  numbered". Postulate 6 now carries BOTH composition duty-holders
%  (Hurwitz multiplicativity + local tomography), matching the
%  continuity section's caveat exactly. Stone's theorem cited
%  (stone1932one, RIS delivered); time-as-operation-count forward
%  flag added; \textquote made portable; opener formalised.
%
%  PROVENANCE NOTE (relocated from the body, 2026-07-06, JS's words):
%  "Using Hopf fibrations to get the Born rule is something Claude AI
%  came up with on its own without my prompting." Preserved here as the
%  historical record. DISCLOSURE STATUS (JS, 2026-07-06): AI use is
%  already acknowledged in the second paragraph of the paper's
%  introduction; the remaining recommendation is only that, at
%  submission, the same fact be echoed in whatever formal
%  Declarations / Acknowledgements slot the target journal's template
%  mandates -- several publishers now require it there specifically,
%  in addition to any in-text mention.
%
%  British English; \autocite; no bullets in prose.
% =====================================================================

% ADDED 2026-07-12: closing paragraph "The postulates in module
% coordinates" -- the Clifford-module reading (skeleton + two
% supplements), with the explanatory arrow guarded per JS: the field
% comes from the Shannon/additivity/commutativity step (self-ignorance),
% and the left-action composability result is recorded as that same
% requirement in operator clothing, not an independent selection of C.
\section{Derivation of the quantum postulates}
\label{sec:postulates}
Enough pieces of the jigsaw are now assembled to derive the postulates of quantum mechanics.

The Dirac--von Neumann postulates are ordinarily laid down as independent
assumptions. In this framework they are not assumptions but a single chain: each
follows from the Inside/Outside Equivalence Principle by one move from the
preceding step. We assemble the chain here and mark the status of each link as
proved, structural, or still owed, so that no step claims more than it has earned.

The chain begins below the postulates. An embedded observer, required by the
equivalence principle to describe a totality that includes itself, must assign a
probability to a proposition about its own future state. The Lawvere--Yanofsky
diagonal \autocite{Yanofsky2003} and Popper's predictor argument
\autocite{Popper1950a, Popper1950b} show that no consistent classical assignment
exists: the self-referential proposition is undecidable to the observer before
the fact. Representing the result of that undecidable self-prediction forces a
signed, non-classical probability, which carries the observer's irreducible
self-ignorance. From this single obstruction the postulates follow.

\paragraph{Postulate 1 (complex Hilbert space).} Two independent steps fix the state space. The
field is forced first: a signed probability, read through the Shannon information measure and
analytically continued, yields a complex amplitude $S=-\ln\lvert p\rvert-\mathrm{i}\theta$,
whose phase $\theta$ is the self-referential ignorance (Section~\ref{sec:complex-amplitudes}).
The metric is forced second, and separately: the principle of reversible continuity
(\S\ref{sec:continuity}) --- that a continuous reversible transformation connects any two pure
states --- selects the $L^{2}$ norm uniquely. For every $\ell^{p}$ with $p\neq 2$ the linear
isometries are generalised permutations \autocite{Banach1932, Lamperti1958}, a group that cannot
connect a basis state to a superposition; only at $p=2$ is the norm-preserving group connected
and transitive on the sphere \autocite{Aaronson2004h}, and Hardy's operational $K=N^{2}$
counting reaches the same quadratic norm by a second route from the same continuity premise
\autocite{Hardy2001a}, and \v{C}encov's theorem with the every-basis lemma reaches it from the
value continuum alone (\S\ref{sec:cencov-route}). The state manifold is therefore the unit
sphere of the Hermitian $2$-norm in $\mathbb{C}^{n}$, taken modulo the global phase
(Section~\ref{sec:l2-metric-fibre}). Placing the metric here, before the Born rule rather than
after it, is deliberate: the Born derivation presupposes the round sphere, so the metric must be
anchored independently to avoid circularity. Two bounds belong with this step. Reversible
continuity fixes the metric, not the field---real and quaternionic quantum theory satisfy it
equally---so the restriction to $\mathbb{C}$ rests on the Shannon step, not on continuity. And
deriving reversible continuity itself from self-reference is the one arrow the chain still owes
--- now priced exactly: one open premise, that every transformation of the regime is a phase of an inside structure,
Open Problem~\ref{op:limit-realisability}, with everything downstream of it proved
(\S\ref{sec:continuity}). Until it is closed, continuity enters as a principle, and the weight
of the $L^{2}$ conclusion is carried jointly with the isometry and counting arguments and with
the statistical route.

\paragraph{Postulate 2 (Hermitian observables).}
A measurement result must be communicable: it must be copyable from one observer
to another, the \textsc{copy} operation of the diagrammatic calculus
\autocite{CoeckeKissinger2017}. The imaginary part of the amplitude is the
private self-referential phase and cannot be broadcast; only the real,
gauge-invariant content can. Observables whose spectra are real and therefore
communicable are exactly the self-adjoint operators. Hermiticity is thus the
condition that a quantity be the kind of thing a measurement can report.

\paragraph{Postulate 3 (Born rule).}
\label{sec:born-invariance}
Outcome probabilities must be invariant under the global phase, the $U(1)$
self-ignorance, realised geometrically as the $S^{1}$ fibre of the complex Hopf
fibration $S^{1}\to S^{3}\to S^{2}$ over the Bloch sphere
\autocite{Hopf1931, Urbantke2003}. There is no polynomial $U(1)$-invariant of degree one,
and the non-polynomial $\lvert\psi\rvert$ is excluded by the analyticity declared at
\S\ref{sec:complex-amplitudes}; the lowest-degree polynomial invariant is quadratic, so the
probability must be $p=\lvert\psi\rvert^{2}$. The squaring is forced
rather than merely permitted, and the argument holds for a single qubit, where
Gleason's theorem is silent. The Born quadratic is not a new structure: it is the
same $L^{2}$ form that fixed the metric in Postulate~1, now read as the
normalisation of the round sphere, and the fibre over which the phase is averaged
has constant length, so the phase is hidden uniformly
\autocite{GluckWarnerZiller1986}. This step is a structural argument rather than a
theorem in the strict sense, and is flagged as such. The record-theoretic derivation of the same
quadratic --- squared magnitudes from fibre privacy, Proposition~\ref{prop:f1-dictionary} --- is
given at \S\ref{sec:born-rule}, and the cost-ledger derivation --- the record as the closed loop of
a history and its return leg, the exponent the number of legs --- at \S\ref{sec:born-cost-ledger}:
three routes to one rule, each holding for a single qubit, the over-determination this framework
treats as evidence.

\paragraph{Postulate 4 (state reduction).}
Measurement is a \textsc{fanout}: the gauge-invariant relational content is copied
into a record, converting a private epistemological amplitude into a public
fact. For the observer who performs it, acquiring that fact is Bayesian updating,
and the projection postulate is its expression. The reduction is perspectival---a
fact relative to the observer in the sense of relational quantum mechanics
\autocite{Rovelli1996a}---and it is graded by redundancy. A single copy is a
relative fact, still reversible by an agent controlling the joint system; redundant
proliferation across many independent fragments makes the fact objective, which is
quantum Darwinism \autocite{zurek_2009} and the stable/relative-fact distinction of
Di~Biagio and Rovelli \autocite{DiBiagioRovelli2021}. ``Collapse'' and unitary
evolution are then the two charts---inside and outside the \textsc{fanout}
boundary---of one process, as the Wigner's-friend analysis shows and \S\ref{sec:fanout-boundary} formalises
(Section~\ref{sec:wigners-friend}); the irreversibility and its $kT\ln 2$ cost
\autocite{Landauer_1961, wigner_1961} accrue with redundancy rather than at the
single act. The reduction postulate is therefore recovered as a chart-dependent
description, not a fundamental dynamical law; the identification of redundancy as
the grading parameter is a synthesis of the framework's parts and is flagged as
such.

\paragraph{Postulate 5 (unitary evolution).}
Conservation of probability is preservation of the $L^{2}$ inner product, and the
inner-product-preserving linear maps are the unitaries $U(n)$---the same
isometry fact that fixed the metric in Postulate~1. Continuity in time --- operation count,
\S\ref{sec:space-of-computations}, a forward reference --- together with Stone's theorem
\autocite{stone1932one} gives $U(t)=\mathrm{e}^{-\mathrm{i}Ht}$, the Schr\"odinger
equation. Reversibility is here a corollary, not a premise: bare reversibility on
the amplitude space yields only the general linear group, and it is preservation
of the quadratic form that carves $U(n)$ out of it.

\paragraph{Postulate 6 (composite systems).}
Systems combine by the tensor product because the division algebra norm is
multiplicative \autocite{hurwitz1898composition}: composition preserves the $L^{2}$ structure
only if states multiply. The operational half closes from the other side: the locality of
tomography derived at \S\ref{sec:local-tomography} says product states determine composites,
which is the tensor product's signature. Cayley--Dickson doubling carries $\mathbb{C}$ to the next
rung, and the ascent of the Hilbert space is the ascent of the Hopf tower, complex
to quaternionic, used in the Wigner's-friend construction
\autocite{mosseri2001hopf}; the ladder itself is the business of the section that follows
(\S\ref{sec:division-algebras}).

\paragraph{Consequence (the Tsirelson bound).}
The postulates fix not only the form of the theory but the strength of its
correlations. In the two-qubit ball the maximal entanglement sits at the centre,
and a Popescu--Rohrlich box \autocite{PopescuRohrlich1994} would require a point
beyond it; the CHSH correlator is capped at $2\sqrt{2}$
\autocite{Cirelson1980QuantumInequality, VerstraeteWolf2002, Held2021}, with the value recovered
at the centre and its general derivation still owed
(Section~\ref{sec:wf-reconciliation}).

\medskip
\noindent The result is that the six postulates are one chain, not six premises.
Of the links, the field and the Born quadratic rest on the framework's signed
probability and $U(1)$-invariance arguments; the $L^{2}$ metric, the unitary group
and the tensor product are over-determined by reversible continuity, the isometry
classification and Hurwitz composition; state reduction is recovered as the
chart-relative face of \textsc{fanout}. The two outstanding debts are the
derivation of reversible continuity from self-reference and the forcing of the Tsirelson value as
a theorem rather than a recovered number --- named, and numbered:
Open Problems~\ref{op:limit-realisability} and~\ref{op:tsirelson-geometry} --- so that a reader
can see exactly where the chain is complete and where it is not.

\paragraph{The postulates in module coordinates.}
One reading of the chain deserves a closing statement, because it positions the derivation
against a mathematical skeleton the reader may know well, and because stating it carelessly would
invert the chain's explanatory arrow. In the coordinates of the division-algebra programmes
(\S\ref{sec:algebraic-dictionary}), complex Hilbert-space quantum mechanics is real Clifford
module theory plus exactly two supplements. The skeleton is Clifford throughout: the operator
algebra of $n$ qubits is the Clifford algebra of $2n$ generators with the state space its own
spinor module, half-integer spin and fermionic statistics are the module's double cover and its
graded tensor, and the classically simulable stabiliser core is the algebra's combinatorics. What
the skeleton does not contain is, first, the choice of complex structure --- the polarisation $J$
that turns a real module into a complex Hilbert space, which standard quantum theory and the
algebraic programmes alike \emph{choose} --- and, second, the Born weights on the fibre. Those two
supplements are precisely what this section has derived, and derived from self-ignorance rather
than from algebra: the field from the Shannon step (information about one's own undecidable state
must still be information --- it must add, addition commutes, and so the amplitude algebra is
commutative, excluding the quaternionic bookkeeping an observer could never balance), and the
quadratic from fibre privacy. The algebraic echo --- that the algebra's own left actions compose
only along a commutative line, the composability result recorded with the dictionary --- is not
an independent selection of $\mathbb{C}$: composition of left actions fails off that line exactly
because the doubled slot carries the \emph{reversed} product, and order-dependence is the same
disease in operator clothing that non-additivity is in information clothing. The mathematics is
the epistemic argument's shadow, welcome as consistency and as a bridge to the module-side
literature (Appendix~\ref{app:winding-number} walks the bridge for charge), but the origin of the
complex amplitude is the observer's ignorance of their own state, and the arrow runs from there
to the algebra, never the reverse.
%      % =====================================================================
%  SECTION: The Division-Algebra Ladder
%  Drafted 2026-07-03. One of the three foundations sections
%  (sec:self-reference, sec:division-algebras, sec:fanout-boundary);
%  this one resolves the dangling \S\ref{sec:division-algebras} in the
%  Hopf-justification section, whose header names it "the
%  Cayley--Dickson / Hurwitz step".
%
%  PLACEMENT in main.tex (REORDERED per JS 2026-07-06): position 6,
%  opening Act II ("why stop at C?") right before its first consumer --
%      \input{self_reference_section_2026-07-04}
%      \input{continuity_section_2026-07-04}
%      \input{L2_metric_global_phase_section_2026_06_07}
%      \input{wigners_friend_in_the_hopf_picture_2026-06-13}
%      \input{derivation_of_the_quantum_postulates}
%      \input{division_algebras_section_2026-07-03}   <-- moved here
%      \input{spacetime_from_non-computability}       <-- first consumer
%                                                          (S^7 exclusion)
%      \input{hopf_justification_section_2026-06-08}
%  The section is position-portable: all adjacency phrasing was
%  retargeted to explicit \S\ref on 2026-07-06. The S^7 hand-off
%  sentence to sec:space-of-computations WRITTEN 2026-07-06 (spacetime
%  file read; the closing now hands first to its immediate consumer,
%  then to hopf-justification). Upload verifications COMPLETE: L2 and
%  quantum-postulates do not consume the ladder; WF's references are
%  forward-voiced or self-carrying.
%
%  DIVISION OF LABOUR (deliberate, to avoid overlap):
%   - sec:self-reference (undrafted): self-reference -> complex
%     amplitudes and the Born rule; why not a real base. This section
%     TAKES that as given and starts at C.
%   - THIS section: the algebra -- the doubling construction, the
%     ledger of losses-as-physics, the internal-algebras-not-a-new-QM
%     scoping (Adler pre-empt), the two appearances of the ladder
%     (register / measurement climb, stated lightly), and Hurwitz as
%     the algebraic termination with the committed zero-divisor
%     reading.
%   - sec:hopf-justification: the geometry (Adams, parallelisability,
%     the fibrations), the join of the two classifications, the
%     entanglement-sensitivity literature, the C-direction honest debt,
%     and the epistemic reading of termination.
%
%  COMMITTED READINGS USED (sources re-read 2026-07-03):
%   - Loss ledger and level dictionary: cayley_dickson_qubits_for_review
%     sec.5 ("each step loses a symmetry, and each lost symmetry is
%     identified with a physical content").
%   - Termination reading: grounding_argument_gauge_group ("the
%     sedenions have zero divisors and cannot support consistent
%     information flow around loops") -- leads; norm-composition-as-
%     probability given as its amplitude-level restatement.
%   - Internal-not-base scoping: cayley_dickson_qubits_for_review Q2 and
%     sec.3.2 caveat (Adler objection; C-qubit treated geometrically).
%
%  HONEST DEBT: ladder proved closed at the top (Hurwitz), derived at
%  the floor. The rung-compulsion revisit-trigger FIRED 2026-07-06:
%  sec:self-reference and op:sigma-functor now exist, so the thesis is
%  minted as op:ladder-compulsion (with the CD-conjugation-as-sigma
%  candidate route), inheriting op:sigma-functor's status.
%
%  AUDITED 2026-07-06 (final pass, third stop): all four citations
%  used; all cross-refs live; CD facts, dims, alternativity, Hurwitz
%  statement verified; Adler subsection delivers everything the
%  register's compressed pointer promises; file pure ASCII; no bare
%  "frame". Fixes: opening rewritten for both predecessors (continuity
%  was inserted between after drafting); "associator's non-vanishing is
%  mass itself" recalibrated to the register's debt-against-selection
%  defusal; the O->S ledger row and Hurwitz subsection reconciled with
%  the fourth-qubit forward remark ("the terminus closes the ladder of
%  forces; what survives past it is, at most, flavour").
%
%  CROSS-REFS: sec:self-reference [sibling dangle -- resolves when that
%  section is drafted], sec:hopf-justification, sec:hopf-information-bound,
%  sec:register, sec:entanglement-classes, sec:generations, sec:masses,
%  sec:wigners-friend. All others exist.
%  CITATIONS: existing keys baez2002octonions, hurwitz1898composition,
%  furey2018three. ONE new key: adler1995quaternionic (RIS in
%  division_algebras_refs_2026-07-03.ris).
%  British English; \autocite; no bullets in prose; the ledger is a
%  table because the comparison is the argument.
% =====================================================================

\section{The Division-Algebra Ladder}
\label{sec:division-algebras}

The foundations behind us end at amplitudes twice over: natively, where the self-referential
argument forced the description of a system into the system's own state space and priced
consistency at a number system whose parts can cancel --- complex amplitudes, and the Born rule
that squares them (\S\ref{sec:self-reference}) --- and independently, where the regime's
continuity delivers the same field through Hardy's theorem (\S\ref{sec:continuity}), the
convergence of the two branches being evidence rather than accident. This section asks the question
both routes leave open. Why
should the number system stop at $\mathbb{C}$? Ordinary quantum mechanics stops there by postulate;
a framework that has forbidden itself postulates must answer, and the answer is that the climb
neither stops at $\mathbb{C}$ nor continues forever. There are exactly two more rungs, each
purchased by surrendering an algebraic law and each paying for itself in physics, and the ladder is
closed at the top by a theorem of 1898. This section sets out the algebra of that ladder: the single
construction that climbs it, the ledger of what each rung costs and buys, and the theorem that ends
it. The geometry of the same ladder --- that each rung acts on a sphere, that exactly those spheres
are parallelisable, and that the algebraic and geometric classifications are one --- is the business
of the geometry section (\S\ref{sec:hopf-justification}).

\subsection{One construction climbs the whole ladder}

The rungs are not a miscellany; a single construction generates them all. Given an algebra with a
conjugation, the Cayley--Dickson construction builds its double on ordered pairs, with a twisted
multiplication in which the conjugate enters and a conjugation $(a,b)^{*} = (a^{*}, -b)$
\autocite{baez2002octonions}. (Conventions differ in the placement of the conjugations; all yield
isomorphic algebras.) Applied repeatedly it produces the sequence
$\mathbb{R} \to \mathbb{C} \to \mathbb{H} \to \mathbb{O} \to \mathbb{S}$, of real dimensions
$1, 2, 4, 8, 16$: the reals, the complex numbers, the quaternions, the octonions, the sedenions.
The ground floor is classical: $\mathbb{R}$ is ordered and positive-definite, the arithmetic of
classical probability. The first doubling is therefore the quantum move itself --- the passage from
probabilities to amplitudes that can interfere --- and it is not made here by choice: it is the
derivation of \S\ref{sec:self-reference}, already behind us. What this section must account for is
the rest of the climb.

\subsection{The ledger}

In mathematics the doublings are usually told as a story of degeneration: each algebra is worse
behaved than the last, and the sequence is studied for where it breaks. This framework reads the
same ledger in the opposite direction. Each rung is bought by surrendering a law, and each
surrendered law is exactly a piece of physics --- the losses are the assets, and the comparison is
sharp enough to be the argument:

\begin{center}
\begin{tabular}{llll}
\hline
doubling & dimension & what is lost & what the loss becomes \\
\hline
$\mathbb{R} \to \mathbb{C}$ & $1 \to 2$ & the order of the reals & phase and interference; amplitudes \\
 & & & that cancel; the Born rule (\S\ref{sec:self-reference}) \\
$\mathbb{C} \to \mathbb{H}$ & $2 \to 4$ & commutativity & non-commuting actions: weak isospin \\
 & & & $SU(2)$, the $B$ register (\S\ref{sec:register}) \\
$\mathbb{H} \to \mathbb{O}$ & $4 \to 8$ & associativity & the associator: colour, chirality, the \\
 & & & generations, and mass as debt \\
 & & & (\S\ref{sec:generations}, \S\ref{sec:masses}) \\
$\mathbb{O} \to \mathbb{S}$ & $8 \to 16$ & division itself & no further force: zero divisors appear \\
 & & (zero divisors) & and the ladder of sectors ends (below) \\
\hline
\end{tabular}
\end{center}

Walking the rows: the order lost at the first doubling is the possibility of interference --- parts
of a probability that can oppose one another --- which is where the derivation of \S\ref{sec:self-reference}
lives. The commutativity lost at the second is the existence of actions whose order matters; the
unit quaternions form the group $SU(2)$, and the framework houses weak isospin there, as the
internal algebra of the $B$ register whose order-sensitive, non-commuting dynamics the interaction sections develop
(\S\ref{sec:register}, \S\ref{sec:entanglement-classes}). The associativity lost at the third ---
the octonions retain the weaker property of alternativity \autocite{baez2002octonions} --- creates
the associator $[a,b,c] = (ab)c - a(bc)$, and the associator is where most of this paper's physics
is transacted: the stabiliser of a preferred complex direction inside the octonions' automorphisms
is the colour group (\S\ref{sec:hopf-justification}); the triality structure and its $K_3$
obstruction generate the three generations and the sterile channel (\S\ref{sec:generations}); and
the associator is the currency of mass --- the debt against the selected direction that
\S\ref{sec:masses} prices, owed by whatever action moves the selection and not by residence at
the rung, which is the defusal \S\ref{sec:register} makes explicit. The algebraic facts
in the third column are theorems; the identifications in the fourth are established in the sections
cited, each carrying its own epistemic flag there, and are not re-litigated here. What this table
asserts on its own account is only the correspondence's shape: one law surrendered per rung, one
sector of physics gained per law.

\subsection{Internal algebras, not a new quantum mechanics}

A reader versed in the quaternionic-quantum-mechanics literature will at this point raise a fair
objection: quantum theory over $\mathbb{H}$ as a foundation is known to be troubled --- the
tensor-product and composite-system difficulties set out by Adler and others
\autocite{adler1995quaternionic} --- and quantum theory over $\mathbb{O}$ founders earlier still,
since a non-associative algebra does not furnish a Hilbert space at all. The framework's position
must therefore be stated exactly, because it does not make the claim the objection targets. The base
of the physical Hilbert space here is $\mathbb{C}$ and remains $\mathbb{C}$; the higher algebras
enter as the \emph{internal} algebras of the register --- $\mathbb{H}$ for the $B$ qubit,
$\mathbb{O}$ for the $C$ qubit (\S\ref{sec:register}) --- in a tensor structure whose outer theory is
ordinary complex quantum mechanics. This is the same move by which the division algebras enter the
constructions of Furey and the related programmes without any proposal to rebase quantum mechanics
itself \autocite{furey2018three}. For the octonionic register the point is sharper still:
$\mathbb{O}^{2}$ is not a Hilbert space, and the framework accordingly treats the $C$ qubit not in
the Dirac--von Neumann manner but geometrically, as the sphere $S^{15}$ with its Hopf structure ---
which is precisely the form in which \S\ref{sec:hopf-justification} needs it. Nothing in this paper therefore
stands or falls with quaternionic or octonionic quantum mechanics as foundations; the ladder is
climbed by the registers' internal algebras, not by the theory's base field.

\subsection{The same ladder, twice}

The ladder appears in this framework in two guises, and it will help the reader to have them
distinguished now. Statically, it is the anatomy of a single relation: the three qubits of the
register take their values one rung apart --- $A$ over $\mathbb{C}$, $B$ over $\mathbb{H}$, $C$ over
$\mathbb{O}$ --- so that every particle carries the whole ladder as its internal structure
(\S\ref{sec:register}). Dynamically, it is the shape of measurement: a single qubit's pure states
fill $S^{3}$, an entangling measurement carries the joint state into $S^{7}$, and a third party into
$S^{15}$ --- the correspondence between qubit entanglement and the Hopf fibrations that
\S\ref{sec:hopf-justification} takes from a mature literature --- so that each act of measurement ascends one rung --- the climb the reader has already made,
narrated in \S\ref{sec:wigners-friend} --- and every completed relation is a terminated climb, which is
the reading \S\ref{sec:hopf-information-bound} gives the tower's endpoint. One ladder, seen twice:
as the anatomy of a relation and as the dynamics of relating. The framework regards this double
appearance not as a coincidence but as the same structure viewed statically and in motion, and the
sections cited carry that identification's weight.

\subsection{Hurwitz closes the ladder}

The construction itself would double forever; what stops it is a classification theorem. By
Hurwitz's theorem, the normed division algebras --- those in which the norm composes,
$\lvert ab\rvert = \lvert a\rvert\,\lvert b\rvert$, and no non-zero elements multiply to zero ---
are exactly $\mathbb{R}$, $\mathbb{C}$, $\mathbb{H}$ and $\mathbb{O}$
\autocite{hurwitz1898composition,baez2002octonions}. The next doubling, the sedenions, has zero
divisors: non-zero elements whose product vanishes. The framework's reading of this failure is the
one that matters physically. Zero divisors cannot support consistent information flow around loops
--- a cycle of compositions can silently annihilate what it carries --- and a framework in which
closed computational paths are load-bearing (\S\ref{sec:hopf-justification},
\S\ref{sec:space-of-computations}) cannot be built on an algebra that permits it. At the level of
amplitudes the same failure reads: the composition law
$\lvert ab\rvert = \lvert a\rvert\,\lvert b\rvert$, which carries probability multiplicatively
through every rung of the ladder, fails beyond the octonions, so that two non-vanishing amplitudes
could compose to nothing and probability would be destroyed by mere composition. The last algebra
fit to carry amplitudes is $\mathbb{O}$. One rung past the terminus is not entirely silent,
however: the sedenions' internal structure --- their automorphism group and their octonion
subalgebras --- is the recorded candidate arena for tagging the generations
(\S\ref{sec:interactions}), and it enters exactly as this section's scoping permits, internally
and never as a base. The terminus closes the ladder of forces; what survives past it is, at most,
flavour. This is the algebraic half of the termination; the
geometric half --- that no sphere above $S^{7}$ is parallelisable --- and the theorem joining the
two halves are the opening business of \S\ref{sec:hopf-justification}, and the epistemic reading of the same endpoint,
as a bound on the information a single relation can carry, is given in
\S\ref{sec:hopf-information-bound}.

\subsection{The honest debt}

As in the geometry section, we close by stating plainly what is proved, what is established, and what is
owed. Proved: the construction, the losses, and the termination are theorems --- the Cayley--Dickson
doubling and its ledger of surrendered laws are standard algebra, and Hurwitz's classification is
among the durable results of the subject. Established: the physical cashing of each rung --- the
Born rule at $\mathbb{C}$, weak isospin at $\mathbb{H}$, colour, generations and mass at
$\mathbb{O}$ --- is carried out in the sections cited above, each under its own epistemic flag. Owed,
and stated without disguise: the ladder as presented here is closed from above and derived at its
floor, but it is not yet derived rung by rung from below. Hurwitz forbids a further rung; the
argument of \S\ref{sec:self-reference} forces the first doubling, from the ordered arithmetic of
classical probability to amplitudes that interfere. That the same self-referential closure, applied
to a system whose amplitudes already take values at one rung, compels the next --- that relational
closure at one level forces precisely the Cayley--Dickson double of that level --- is the
framework's reading of the climb, and it is a thesis, not a theorem --- but the formal apparatus
against which it can be posed now exists, and we pose it.

\begin{openproblem}[The compulsion of the intermediate rungs]\label{op:ladder-compulsion}
Derive, within the categorical apparatus of Open Problem~\ref{op:sigma-functor}, that
self-referential closure at one rung compels the Cayley--Dickson double of that rung, so that the
intermediate steps $\mathbb{C} \to \mathbb{H} \to \mathbb{O}$ follow from the same principle
that forces the floor. The candidate route is the construction's own shape: the Cayley--Dickson
pair with its conjugation, $(a, b)^{*} = (a^{*}, -b)$, is an involution-equipped doubling, and the
route reads it as the algebraic shadow of adjoining the describer --- closure of level-$A$
descriptions under the inside/outside exchange requiring exactly the twisted pair. The problem
inherits the status of Open Problem~\ref{op:sigma-functor}, whose apparatus it presupposes.
\end{openproblem}

If the derivation lands, the whole ladder becomes a consequence of the same self-reference that
opened it. Until then, the honest statement is that the ladder's necessity is established at its
ceiling and at its floor, while the compulsion of the intermediate rungs is inherited from the
physics they are shown to carry rather than derived in advance. What this section hands on is, accordingly, an algebraic ladder with a floor, a ceiling and a
ledger. Its first consumer is immediate: \S\ref{sec:space-of-computations} spends the terminal
fact --- the non-associativity of $S^{7}$'s own multiplication --- as geometry, excluding the
seven-sphere as a ruler and fixing the dimension of space at three.
\S\ref{sec:hopf-justification} then supplies the geometry of the fibrations themselves, and the
join.
   <-- moved here
%      % EDITED 2026-10-07 scripts/patch_infinite_not_flat_2026-10-07T1400.py
% EDITED 2026-09-27 scripts/patch_spacelike_inversion_2026-09-27T1800.py
% EDITED 2026-09-21 scripts/patch_op_sterile_gravity_boundary_2026-09-21T1300.py
% EDITED 2026-09-12 scripts/patch_spacetime_bloch_headline_2026-09-12T1400.py
% EDITED 2026-09-11 scripts/patch_splice_loops_prospects_2026-09-11T0900.py
% EDITED 2026-09-10 scripts/patch_cite_rubino_2026-09-10T0700.py
% EDITED 2026-09-10 scripts/patch_quoting_and_typography_2026-09-10T0620.py
% =====================================================================
%  SECTION: Spacetime from non-computability (sec:space-of-computations).
%  AUDITED 2026-07-06 (final pass, seventh stop).
%
%  PLACEMENT (reordered chain, JS 2026-07-06): position 7 --
%      \input{division_algebras_section_2026-07-03}   <-- hands off here
%      \input{spacetime_from_non-computability}
%      \input{hopf_justification_section_2026-06-08}
%  Division-algebras' closing now carries the S^7 hand-off sentence
%  (written 2026-07-06 after this file was read).
%
%  INPUT DEPENDENCIES (both held in the audited outputs):
%      entropic_gravity_family_paragraph_2026-07-03.tex  (sec:jacobson)
%      gr_nonlocality_revision_2026-06-13.tex            (sec:gr-nonlocality)
%  The 2026-07-04 additions (asymmetry paragraph; Penrose second route;
%  acceleration/Unruh arc) are FULLY PASTED; the separately uploaded
%  additions file is redundant and can be archived.
%
%  AUDIT RECORD: opener formalised (Russell typo; bare ref retargeted
%  to the freshly minted sec:complex-amplitudes); the "rotations of
%  S^7" wording made precise at BOTH occurrences (Moufang loop, not a
%  group; SO(8) disclaimer added -- referee hazard removed); two \\
%  paragraph-break bugs fixed; dark-matter paragraph retargeted to
%  sec:generations and stitched to op:sterile-gravity (the zoo's
%  demotion and this paragraph's priority condition are one question);
%  CP paragraph stitched to the fourth-party reading of
%  sec:interactions; stale mid-file mini-header retired
%  (sec:k3-sterile placeholder dies with it); the section's narrated
%  debts NUMBERED at last -- op:comp-metric and op:three-dimensions
%  minted in the Open-problems subsection, cited from the in-text
%  "named open problems" line and from the closing summary.
%
%  KEYS: adams1960hopf matches corpus. Unique to this section, presumed
%  in the .bib since the paper compiles: Spence-Brown-1969 (suggest
%  normalising the key spelling at leisure), VanRaamsdonk2010,
%  MaldacenaSusskind2013, MuellerMasanes2013, MasanesMuellerEtAl2014,
%  toffoli_2003, Chaitin1998, Penrose1996, HohnMueller2016,
%  GarnerMuellerDahlsten2017, Jacobson1995, Bekenstein1973, tHooft1993,
%  Susskind1995, spekkens2007toy, Andrianov2001, Sakharov1967,
%  Planck2020, Riess1998. Shared keys verified: small_2006, Penrose1994,
%  unruh1976notes (swap-flag stands), baez2002octonions, Wheeler1990,
%  furey2018three, DiBiagioRovelli2021, Rovelli1996a, Spekkens toy.
%  British English; \autocite; no bullets in prose.
% =====================================================================

\section{Spacetime from non-computability}
\label{sec:space-of-computations}

This section is a deliberate diversion: it takes the idea that non-computable steps are measured
by imaginary information (\S\ref{sec:complex-amplitudes}) and asks what geometry it buys.
Spencer-Brown records that Bertrand Russell thought well of the underlying move
\autocite{Spence-Brown-1969}; the reader may judge whether the framework spends it responsibly.

The programmes that derive spacetime from quantum mechanics---Van~Raamsdonk's
\enquote{building spacetime with entanglement} \autocite{VanRaamsdonk2010}, the
ER\,$=$\,EPR conjecture \autocite{MaldacenaSusskind2013}, and M\"uller's
information-theoretic recovery of three-dimensional space from the qubit
\autocite{MuellerMasanes2013}---take entanglement or information as primitive and
grow geometry from it. The route taken here is the same in spirit but inverted in
order: rather than positing spatial dimensions and forbidding superluminal
signalling, we \emph{define} spatial separation as a measure of non-computability
and read the dimension, the signature, and the Lorentz group off the result. The
proposal originates in the author's CASYS'05 paper \autocite{small_2006}, which
first suggested treating spacetime as a space of computations and attempted to
justify three dimensions on this basis; the present account tightens and extends
it. This is the framework's most structural and least formally complete arc; the
three-dimensionality result is a structural conclusion awaiting a formal proof,
and several of the links below are conjectural readings. Each is flagged as such,
so the reader sees exactly where the chain is load-bearing and where it is an
invitation.

\subsection{Timelike and spacelike distance as computation}
\label{sec:comp-distance}

A computation carries a state from axioms to a theorem in a number of steps.
Counting the computable steps between two states gives a real measure---time. A
physical clock counts completed irreversible loops, so proper time literally
counts loop completions and the metric $\mathrm{d}\tau^{2}=g_{\mu\nu}\,
\mathrm{d}x^{\mu}\mathrm{d}x^{\nu}$ is a loop-density tensor, the conversion from
coordinate intervals to loop counts; Toffoli's reading of the Lagrangian as a
count of action quanta is the same statement \autocite{toffoli_2003}.

Spacelike separation is the complementary, non-computable measure. Define the
non-computable distance between two propositions, relative to an axiom system, as
the number of bits that must be added to the axioms to make the second derivable.
A closed timelike curve is a computation whose output is its own input; by the
fixed-point theorem it has a unique consistent solution or none, and the
existence of that solution is one bit not derivable from any axiom system that
does not already contain the loop---one bit per curve, in the manner of Chaitin's
$\Omega$ \autocite{Chaitin1998}. The spacelike distance between two events is then
the number of such curves that would have to be closed to make them causally
connected. Because each added bit is a non-computable proposition, it is carried
by the imaginary part of the complex information measure
$S=-\ln\lvert p\rvert-\mathrm{i}\theta$ derived in
Section~\ref{sec:complex-amplitudes}: computable distance is counted in real
units, non-computable distance in imaginary ones. The Minkowski signature, with
spacelike intervals imaginary relative to timelike ones, is the immediate reading. The convention is not new to this paper: it was proposed in the author's CASYS paper two
decades ago, where counting non-causal steps in imaginary numbers was introduced and its formal
proof flagged, already then, as a key open problem \autocite{small_2006} ---
Open Problem~\ref{op:comp-metric} is that flag's direct descendant, stated now with the machinery
to be answerable.
That the resulting structure is a genuine metric space, and that the signature
follows rather than being fitted, are open problems---at present the
identification is motivated by reproducing what nature does, not proved.

%% VERIFY-SUPPORT 2026-09-27: conway_2003 is cited below for "any two octonions
%% generate an associative subalgebra" (Artin's theorem; a Moufang loop is
%% diassociative) and harvey1982calibrated for "a three-plane in Im O is
%% associative iff the associator of a basis vanishes, and the associative
%% three-planes are the imaginary parts of the quaternionic subalgebras"
%% (Harvey--Lawson, ch. IV). Neither passage was read in the session that wrote
%% this; both statements are checked numerically in Convention A by
%% scripts/antitelephone_network_associativity_2026-09-27T1900.py (PASS).
A reader used to the standard order of ideas should pause here, because the order has been
inverted, and the inversion is used below. In special relativity the tachyonic antitelephone is a
corollary: a signal faster than light, composed with a boost, would arrive before it was sent and
close a timelike loop, therefore there is no such signal \autocite{tolman1917theory,
benfordbooknewcomb1970tachyonic}. The metric comes first and the impossible telephone is a check
on it. Here the telephone comes first. Two observers are spacelike separated when an antitelephone
between them would close a loop; the loops such telephones would close, and carry nothing round,
are what spacelike separation is, their number is the distance, and the metric is read off from
them. Nothing is lost, since the corollary returns as a tautology: an antitelephone that carried
a record would close a causal loop, and its two ends would not have been spacelike after all.

What is gained is a definite meaning for a spatial dimension. Picture many observers, each holding
an antitelephone to each of the others. A message relayed from one observer to another through a
third composes two transports round two loops, and the network is consistent only if every relay
composes to a single transport whatever the bracketing of the relaying: a spatial dimension is a
family of loops whose transports compose associatively, which is a group law on a parallelisable
sphere. That is a condition a family of loops can fail, and \S\ref{sec:three-dimensions} shows
that the loops of the seventh sphere fail it. Two antitelephones always compose, because any two
octonions generate an associative subalgebra \autocite{conway_2003}; three compose only when their
imaginary parts lie in one quaternionic subalgebra, which is what the associative three-planes of
Harvey and Lawson are \autocite{harvey1982calibrated}, and a third link outside that subalgebra is
seen by the network exactly through its component outside it. The largest consistent network of
antitelephones inside $S^{7}$ is therefore an $S^{3}$, and there are three spatial directions.
When the seventh sphere is set aside there as a ruler, which is the octonionic spacetime of the
division-algebra literature, the claim is not that $\mathbb{R}^{8}$ has no metric, nor that the
ten-dimensional Minkowski space of $SL(2,\mathbb{O}) \cong \mathrm{Spin}(9,1)$ fails to exist
\autocite{baez2002octonions}; it is that the octonionic directions are not loops on which
antitelephones compose, and so are algebra rather than distance. In the standard order that
sentence has no content; in this one it is the reason space has three dimensions.

The split just defined explains the most familiar asymmetry in physics: that one can move back and
forth in space but not in time. \S\ref{sec:continuity} identifies the pre-measurement regime with
continuous reversibility --- its transformation group is connected, and every motion in it is
undoable. Spacelike separation, being the non-computable sector, is exactly the part of the
relational structure on which no record has been written; motion across it is motion within the
reversible regime, and it inherits that regime's continuity and its two-sidedness: a place can be
revisited because going and returning writes no record that forbids the return. Timelike separation
is the record chain itself --- proper time counts completed irreversible loops, as above --- and a
moment cannot be revisited because a moment \emph{is} a stratum of records, and returning would
require unwriting them, which is what \textsc{fanout} forbids (\S\ref{sec:fanout-boundary}). The
asymmetry between space and time is therefore not an extra postulate about the manifold: it is the
pre-/post-\textsc{fanout} asymmetry read as geometry --- space is what the reversible group feels
like from inside, time is what the record chain feels like --- and the arrow-of-time discussion
later in this section is the same asymmetry read thermodynamically. The identification is the
framework's structural reading of the split, offered in the same register as the rest of this
subsection.

\subsection{Three dimensions of space}
\label{sec:three-dimensions}

The closed timelike curves that define spacelike distance require information to
be parallel transported consistently around a loop, and consistent parallel
transport on a sphere requires it to be parallelisable. By Adams' theorem the
only parallelisable spheres are $S^{1}$, $S^{3}$ and $S^{7}$
\autocite{adams1960hopf}, the unit spheres of $\mathbb{C}$, $\mathbb{H}$ and
$\mathbb{O}$. Impose now the requirement that information arriving from outside
the causal zone must appear to possess its own causal history---that is, to be the
result of a sequence of information-preserving operations, hence a sequence of
group operations. The spheres $S^{1}$ and $S^{3}$ are groups under their own multiplication --- $U(1)$ and
$SU(2)$ --- and their operations preserve information; $S^{7}$'s multiplication is non-associative,
a Moufang loop rather than a group, so $S^{7}$ cannot serve as a ruler. (The rotation group
$SO(8)$ of the seven-sphere is of course a group; what fails is the sphere's \emph{own}
composition law, the one a ruler would use.) The same non-associativity
forbids a metric built from the sphere's own composition law outright: path
composition $d(A,C)$ through $B$ depends on the bracketing, so the triangle inequality is not merely violated but ill-defined, and
no geodesic structure exists. (What is untouched should be said: the Euclidean metric of
$\mathbb{R}^{8}$, and the ten-dimensional spacetime that $SL(2,\mathbb{O})$ furnishes, are metrics
of the ordinary kind; what fails is a ruler made of antitelephone loops, which is what a spatial
dimension is in this section, as \S\ref{sec:comp-distance} set out. Nor is the failure total:
relays between any two links compose, and among three they compose exactly when the links lie in
one quaternionic subalgebra, so the largest consistent network inside $S^{7}$ is an $S^{3}$
\autocite{conway_2003,harvey1982calibrated}.) Since the abelian $S^{1}$ alone cannot carry
non-abelian structure, spacelike distance must be three-dimensional. Dimension
three is an output, not an input \autocite{small_2006}.

The division of labour is then exact. $S^{1}$ ($\mathbb{C}$) contributes the time
dimension and the $U(1)$ of electromagnetism; $S^{3}$ ($\mathbb{H}$) contributes
the three spatial dimensions and the $SU(2)$ of the weak interaction; $S^{7}$
($\mathbb{O}$) contributes no geometry but, as a filter on what information may be
exchanged, the $SU(3)$ of the strong interaction. A sphere that cannot be a ruler
can still be a filter, which is why the strong force resists geometric
interpretation while electromagnetism and the weak force do not. This is the same
$\mathbb{C}\oplus\mathbb{C}^{3}$ split that yields $G_{\mathrm{SM}}$ in
Section~\ref{sec:entanglement-classes}, now read as the reason spacetime is $(3+1)$-%
dimensional. The conclusion converges with M\"uller and Masanes, who prove from
information-theoretic postulates, assuming no quantum theory, that minimal
direction-carrying systems with continuous reversible dynamics force $d=3$ and
quantum theory on two qubits \autocite{MuellerMasanes2013}, and with the result
that entanglement exists only for a three-dimensional Bloch ball
\autocite{MasanesMuellerEtAl2014}; the framework supplies the reason such systems
exist, Cayley--Dickson doubling from self-reference, while their work supplies
the rigorous dimensional uniqueness. The status of the framework's own route is a
structural result requiring formal proof: the one-bit-per-curve identification and
the precise reduction of the $S^{7}$ filter to $SU(3)$ rather than its full automorphism group
$G_{2}$ are numbered below (Open Problems~\ref{op:comp-metric} and~\ref{op:three-dimensions}).

A scoping remark, and then a second route to the same count. The computational model of spacetime
proposed here has many detailed parallels with the real thing, but a parallel is not an identity,
and the identity is not claimed in this paper: establishing it is work deferred to the second paper.
What can be claimed now is that the parallel is productive in a specific direction --- it suggests a
derivation of the dimension count carried out entirely within general relativity, with no
computational premise. Acceleration and gravitation tip the light cone, and the tipping can be
carried to the extreme: Penrose's discussion runs the cone through a full circle
\autocite[382]{Penrose1994}, at which point the tipped direction closes on itself --- which is to
say, the extreme of tipping is a closed timelike curve, the very object whose closure the
computational definition of \S\ref{sec:comp-distance} counts in bits. The two routes are one route
entered from opposite ends. Now demand of general relativity only this: that observers accelerating
with respect to one another, and with respect to unaccelerated observers, retain a consistent
conception of the distance between them. Each observer's tipping traces a path on a sphere of
directions, and mutual consistency requires the associated rotations to compose without regard to
bracketing --- to form a group. Compositions on $S^{1}$ and $S^{3}$ do, forming $U(1)$ and $SU(2)$ as above; compositions on
$S^{7}$ do not, the multiplication being non-associative, so observers tipped across an
$S^{7}$-worth of directions could not agree on distance even pairwise: the same failure that forbade
$S^{7}$ a metric above forbids it here a shared ruler. The largest sphere admitting a consistent
conception of distance under mutual acceleration is therefore $S^{3}$, and space has at most three
dimensions; with the floor argued above --- the abelian $S^{1}$ alone cannot carry the required
structure --- the count is exactly three. The status matches the subsection's. The tipping and its
extreme are general relativity's own \autocite{Penrose1994}; the consistency demand is
elementary; the associativity filter is the algebraic fact used throughout; the reading of composed
tippings as motion on the parallelisable spheres is the structural step, shared with the
computational route. The convergence of the two routes on the same endpoint is offered as evidence
for the parallel, while the identity it gestures at remains the second paper's burden.

\subsection{Recovering the Lorentz group}
\label{sec:lorentz}

A measurement changes what the observer knows, hence, relationally, the causal
structure: it tips the light cone. This is Penrose's association of cone-tipping
with state reduction \autocite{Penrose1996} run in the opposite causal direction,
reduction driving the tipping rather than the reverse, and the two are then one
operation in two charts---both described by $SL(2,\mathbb{C})$. The algebraic
backbone is the division-algebra correspondence
$SL(2,\mathbb{A})\cong\mathrm{Spin}(1,d{+}1)$ \autocite{baez2002octonions}: $\mathbb{R}$
gives $1{+}1$, $\mathbb{C}$ gives $3{+}1$, $\mathbb{H}$ gives $5{+}1$ and
$\mathbb{O}$ gives $9{+}1$ dimensions. The computability split that singles out a
preferred $\mathbb{C}$ fixes the fundamental state space as $\mathbb{C}^{2}$ and
hence $SL(2,\mathbb{C})=\mathrm{Spin}(1,3)$: the Lorentz group of $3{+}1$
spacetime. The qubit and Minkowski space are not merely both Euclidean-and-three%
-dimensional by coincidence; the Hermitian $2\times2$ matrices are
$\mathbb{R}^{1,3}$, $SL(2,\mathbb{C})$ acts on them as the Lorentz group, and the
qubit density operators fill the future cone---the correspondence M\"uller's group
develops operationally, recovering Lorentz transformations from quantum
communication \autocite{HohnMueller2016} and the complex and quaternionic qubit
from the relativity of simultaneity \autocite{GarnerMuellerDahlsten2017}.
In these coordinates the Bloch ball is the unit-trace slice of the cone and the Bloch sphere its
boundary rays, the celestial sphere on which $SL(2,\mathbb{C})$ acts by M\"obius
transformations; state reduction sends a point of the ball to a point of the sphere, and a boost
sends the sphere to itself, which is what lets the two be one operation in two charts.
The $SL(2,\mathbb{A})$ correspondence is rigorous; its identification with
measurement-as-tipping is the framework's structural reading.

One consequence of the identification has been implied throughout this section and deserves its own
paragraph, because the next paragraph turns on it. If measurement creates a relation between
observer and observed, and if --- relationally --- a new relation alters the causal structure, then
the tipping of the cone is not merely \emph{associated} with the creation of the relation: it is
what the creation of the relation is, seen geometrically. And a tipping of one's cone against the
congruence of one's neighbours already has a name in physics: acceleration. The framework's
statement, made explicit here, is that \emph{creating a new relation between observer and observed
is experienced as acceleration}. The resonance with the equivalence principle is exactly what it
should be: the principle makes acceleration and gravitation locally indistinguishable, Penrose's
association --- run in this section's direction, reduction driving the tipping --- makes state
reduction and cone-tipping one operation, and the framework's identification closes the triangle,
measurement and acceleration meeting where the principle says they must. The claim is structural:
the framework's reading of the cone-tipping established above, not an independent derivation.

Taken literally, the identification has a measurable face, and it is a familiar one: the Unruh
effect, read in this section's own ledger. Recall the definitions of \S\ref{sec:comp-distance}:
timelike distance counts computable steps, spacelike distance counts non-computable bits. What an
observer can compute --- what their \textsc{fanout}s can convert into classical records --- is fixed
by their causal structure, and the previous paragraph has just made that structure
relation-dependent. An accelerated observer's cones are tipped, and with them the boundary of the
recordable: relations that for the inertial observer lie wholly in the reversible, pre-record sector
--- the entanglement structure of the vacuum, non-computable for them, spacelike in the
computational sense --- fall, for the accelerated observer, partly on the recordable side. The mechanism is the one
Chaitin identified for formal systems \autocite{Chaitin1998}: a proposition undecidable relative to
an axiom system becomes decidable when information is added to the axioms --- and
\S\ref{sec:comp-distance} defined spacelike distance as exactly the bits whose addition closes the
gap. Creating a relation is adding information to the observer's axiom base, so timelike and
spacelike distance are converted into one another by the addition and subtraction of information,
and acceleration is the continuous version of that addition. A
relation newly recordable is a detection. Particles appear from nowhere: an accelerated detector in
the inertial vacuum registers a thermal bath at $T = \hbar a / 2\pi c k_{\mathrm{B}}$
\autocite{unruh1976notes}. Two guards keep the reading honest. Nothing here reclassifies the
invariant interval --- timelike and spacelike separation in the emergent metric are Lorentz
invariants, and no event crosses that classification; what is observer-relative is the computational
partition, which relations a given observer's horizon renders recordable. In the ledger of
\S\ref{sec:comp-distance} that partition \emph{is} the timelike/spacelike distinction, so in the
framework's own units timelike distance is \emph{rotated} into spacelike, and spacelike into
timelike, by a change of observer --- the tipping of the cone is literally that rotation --- and the
detector's clicks are the rotation made audible. And the effect is established physics
while the reading is the framework's, flagged structural. What gives the reading its content is a
parallel: the inertial and the accelerated observer stand to the vacuum exactly as Wigner and his
friend stand to the laboratory (\S\ref{sec:wigners-friend}) --- one holds a coherent relation and no
event, the other holds a record, and the facts are relative on both sides in precisely Rovelli's
sense \autocite{DiBiagioRovelli2021}. Since the previous paragraph identifies acceleration with
relation-creation, the parallel is not an analogy but the same structure in gravitational dress, the
horizon standing where the laboratory wall stood: \emph{the Unruh effect is Wigner's friend for
accelerated observers}. And two subsections of this section that seemed independent here meet: the
temperature Jacobson's construction requires across every local Rindler horizon
(\S\ref{sec:jacobson}) is exactly this temperature --- the $T$ of the Clausius relation is the $T$
of the tipped cone, supplied by the very relation-creation with which this subsection began.

\subsection{The relationalist commitment}
\label{sec:relationalism}

General relativity is relational: by the hole argument, spacetime points carry no
intrinsic identity, and the metric is fixed by the network of coincidences and
interactions, not by a pre-existing stage \autocite{Rovelli1996a}. The
computational picture is relational by construction and explains why it must be:
a background metric would require spatial distances to be pre-computable, and
non-computability forbids exactly that. General covariance is the statement that
distances have no pre-assigned values and depend on what information is actually
present. The metric is therefore dynamical not by stipulation but because it
measures a quantity that cannot be fixed in advance.

\subsection{Jacobson, holography, and imaginary information}
\label{sec:jacobson}

Jacobson derives the Einstein equation as an equation of state, from the Clausius
relation $\delta Q=T\,\mathrm{d}S$ across every local Rindler horizon together
with the Bekenstein proportionality of entropy to horizon area
\autocite{Jacobson1995, Bekenstein1973}. In the computational reading this is the
self-consistency condition on the loop-counting network: the Einstein equation is
the requirement that the bookkeeping of computation be internally consistent
across surfaces. The holographic principle, that the information in a region is
bounded by its bounding area \autocite{tHooft1993, Susskind1995}, then acquires a
natural reading here. A surface area is a squared distance; if spacelike distance
is measured in imaginary units of information, an area is measured in the square of
those units, that is in \emph{negative} information---entropy. The holographic
link between area and entropy thus requires the square root of negative
information, and that square root is precisely the imaginary unit of the complex
amplitude, the observer's self-referential phase. Read this way, Jacobson's
account of information crossing surfaces implicitly assigns an imaginary
information value to spacelike distance, and the $\mathrm{i}$ of the Minkowski
signature is the same $\mathrm{i}$ as the $\mathrm{i}$ of the quantum amplitude.
This identification is the most speculative step in the section: it is
structurally suggestive and it makes the holographic form unsurprising, but it is
a conjectural reading rather than a derivation, and is offered as such.

The conjecture has a two-way motivation that this paper's record should carry, because each
direction was on the table before the other was noticed. Forward: if spacelike distance is counted
in imaginary units of non-computable information, then any spacelike area is the product of two
imaginary lengths and is therefore \emph{negative} information --- which is to say entropy, by
Brillouin's identification \autocite{brillouin1956science} --- and a holographic principle arrives
automatically, the sign supplied by $i^{2}$ rather than by interpretation. Converse, and more
striking: in the holographic principle as ordinary physics already states it, surface area is
proportional to entropy, and entropy is negative information; so by implication ordinary physics
has been measuring spacelike separation in imaginary units of information all along, and the
convention of \S\ref{sec:comp-distance} is read off from established results rather than
invented. The same imaginary unit is already employed once in this section --- it is what makes
the signature of the interval Minkowski --- so the sign of the area law is its second job, not a
second assumption. What the two-way motivates it does not prove: the formal statement is exactly
the metric-and-signature half of Open Problem~\ref{op:comp-metric}, together with the
one-bit-per-bit-pair coefficient owed to the second paper. The stakes of the conjecture have
recently risen: Dorau and Much prove, with modular-theoretic rigour, that the semiclassical
Einstein equations follow from quantum relative entropy on a local horizon \emph{once the
entropy--area proportionality is assumed} \autocite{doraumuch2026relative} --- so what this
framework proposes to derive is precisely the remaining input of a derivation already carried
through at that standard.
\input{entropic_gravity_family_paragraph_2026-07-03}

\subsection{Non-locality already present in general relativity}
\label{sec:gr-nonlocality}
\input{gr_nonlocality_revision_2026-06-13}
\medskip
\noindent  The spacetime arc therefore yields, from one definition of spacelike
distance as non-computability: the signature, three spatial dimensions,
the Lorentz group, the dynamical relational metric, a reading of the
holographic principle, the presence of non-locality within general
relativity, the identification of relation-creation with acceleration
and of the Unruh effect with Wigner's friend, and the space--time
mobility asymmetry as the \textsc{fanout} asymmetry read as geometry. Of these,
the $(3+1)$ count and the Lorentz recovery are structural results that converge with independent
rigorous work and await their formal proof --- the count now numbered as
Open Problem~\ref{op:three-dimensions}; the metric-space status, the signature and the
one-bit-per-curve measure (together Open Problem~\ref{op:comp-metric}), the imaginary-information
reading of holography, and the general-relativistic non-locality, the acceleration and
Unruh readings, and the mobility asymmetry are conjectural, and are named here precisely so that the boundary between what is shown and what
is proposed is visible.

\subsection{Why the metric is not to be quantised}\label{sec:no-quantised-gravity}

A consequence has been implicit in everything above and deserves its plain statement. Causal
structure, on this framework's reading, is the bookkeeping of locked relations --- post-\textsc{fanout}
by construction --- while superposition is the signature of the pre-\textsc{fanout} regime. A
``superposition of causal structures'' is therefore not a difficult state but a category error: it
applies the pre-record operation to a record-side object, exactly as ``a superposition of
measurement outcomes'' does. Gravity lives in the post-\textsc{fanout} regime. The quantum layer
beneath it exists and is not in dispute; but it is \emph{pre-geometric} --- it carries amplitudes
over the formation of relations, not amplitudes over formed geometries --- so there is no
fundamental metric to quantise, and the standard programme of superposing spacetimes pursues an
object the ontology does not contain. One guard rail keeps the claim honest: collective
excitations of the locking pattern may perfectly well exist as effective quanta --- gravitons as
the phonons of the relational network --- and the thesis forbids fundamental metric superposition,
not effective field theory; this is Jacobson's equation-of-state moral \autocite{Jacobson1995},
made structural. The relation to gravitational-collapse programmes is an inversion: where
superposed geometries are there read as unstable \autocite{Penrose1996}, here they are read as
unformed --- what decays in those models is what locks in this one.

Three experimental postures follow, each at the status the reading carries. Laboratory
demonstrations of indefinite causal order \autocite{procopio2015experimental, rubino2017indefinite} are, on this
account, demonstrations of the nascent layer: the indefiniteness lives below the locking threshold
and definitises at it, so the framework adopts the programme and adds a boundary --- indefinite
causal order can never be amplified into records while remaining indefinite. Gravitationally
induced entanglement of the kind the Bose--Marletto--Vedral proposals seek
\autocite{bose2017spin, marlettovedral2017gravitationally} is \emph{predicted}, because the
pre-geometric layer is quantum --- while witnessing that layer, not a superposed metric: the
framework splits the experiment's usual dichotomy. And gravitationally driven anomalous collapse
is denied: locking is redundancy and thermodynamics, so decoherence experiments should find
standard rates and nothing gravity-specific. That no consistent quantisation of gravity exists was
the conclusion of the author's CASYS paper by a different route \autocite{small_2006}; the present
framework supplies the mechanism the conclusion then lacked.

\subsection{Open problems and prospects}
\label{sec:spacetime-open-problems}

The spacetime arc is the framework's frontier, and honesty about its edges
matters more here than anywhere else in the paper. We set out the genuine
problems and the prospects together, deferring the deepest to a second paper
rather than papering over them. Two of the arc's debts are well-posed now, and we number them.

\begin{openproblem}[The computational metric]\label{op:comp-metric}
Prove that the paired distance of \S\ref{sec:comp-distance} --- computable steps counted real,
non-computable bits counted imaginary, one bit per closed curve --- defines a genuine
pseudo-metric on the space of computations, with the Minkowski signature derived rather than read
off; equivalently, establish the one-bit-per-curve measure as a measure, and the signature as its
theorem. Until then the identification stands as motivated by reproducing what nature does, which
is the status flagged at its introduction.
\end{openproblem}

\begin{openproblem}[Dimension three as a theorem]\label{op:three-dimensions}
Turn the two convergent routes of \S\ref{sec:three-dimensions} into a proof: formalise the
parallelisability-plus-associativity filter (the computational route) and the mutual-acceleration
consistency demand (the general-relativistic route) so that $d = 3$ follows as a theorem, and
include the reduction clause --- that the $S^{7}$ filter acts through $SU(3)$, the stabiliser of
the selected complex direction, rather than through its full automorphism group $G_{2}$. The
M\"uller--Masanes uniqueness results \autocite{MuellerMasanes2013, MasanesMuellerEtAl2014} are
the rigorous convergents any proof must meet.
\end{openproblem}

\paragraph{The substrate and the generation of events (deferred).}
A computation runs on something; a network of events presupposes that events
occur. Both are the same objection in two forms---the network's hardware and its
clock---and the original CASYS'05 statement of the programme raised the
implementation regress and set it aside \autocite{small_2006}. The framework's
response is that a self-referential structure is the only kind that need not
presuppose a meta-substrate: it generates its own. The natural algebraic emblem
is $E_{8}$, the unique simple Lie group whose adjoint representation is also its
smallest faithful one, so that what it is and what it acts on coincide; the $240$
unit integral octonions that generate it are exactly the $E_{8}$ root system, so
the self-action is welded to the division-algebra spine of
Section~\ref{sec:entanglement-classes}. This is Wheeler's self-excited circuit
\autocite{Wheeler1990} given an algebraic body, and it converges with Furey's use
of an algebra acting on itself to generate particle content
\autocite{furey2018three}. The decisive caveat, which we state rather than hide, is
that the bootstrap is self-consistent but not self-specifying: it does not
eliminate the non-computability identified throughout, it relocates it to a
determined place. The substrate question is, on the framework's own terms,
provably without a computable answer from inside---it is the irreducible residue,
the bump in the carpet that the whole construction has been moving rather than
removing. The leap from ``$E_{8}$ is self-dual'' to ``$E_{8}$ is the substrate''
is at present an emblem awaiting operational content (an $E_{8}$-equivariant
dynamics, or a state space built as a homogeneous space of $E_{8}$), and that
content, together with how a tick is generated at all, is deferred to the second
paper.

\paragraph{Dark matter (deferred).}
If gravity is the bookkeeping of the interaction network and dark matter interacts only
gravitationally, how does it participate in a network it does not exchange information with?
This is a genuine boundary of the programme, at which it may make no prediction, and Open
Problem~\ref{op:sterile-gravity} states it with the criteria that would close it. The prospect, to be developed in the second paper, is
that this is a feature rather than a defect. In the division of labour of
Section~\ref{sec:three-dimensions}, participation in the metric---counting loops
in the $S^{1}$ and $S^{3}$ directions---is universal, carried by anything with a
worldline and a tick-rate, whereas participation in the gauge filters, the
$S^{3}$ of the weak interaction and the $S^{7}$ of colour, is selective, engaged
only by charged matter. A node that counts loops but engages no filter
\emph{is} matter that gravitates without otherwise interacting, and the
framework's Class-4 a-chiral sterile neutrino (\S\ref{sec:generations}) is the candidate that
realises it --- recorded in the particle sections as Open Problem~\ref{op:sterile-gravity}, whose
tension (a gauge-sterile channel gravitating, when gravity is measurement bookkeeping) is exactly
the question named here. The condition for the candidacy to be coherent rather than circular is
that metric participation be established as prior to, and independent of, gauge participation;
making that priority precise is the work deferred, and it is the same work the open problem
requests.

\paragraph{CP violation, the arrow of time, and baryogenesis.}
CP violation appears to demand a time asymmetry owing nothing to the entropic
arrow, which would seem to require a preferred direction built into spacetime
itself. The framework dissolves the appearance rather than accommodating it. CP
violation needs an irreducible complex phase, the CKM phase $\delta$; by Spekkens'
result that negativity and contextuality are the same notion of nonclassicality
\autocite{spekkens2007toy}, a complex phase is a signature of contextuality; so CP
violation should be contextual, sensitive to the near environment of the decaying system --- the
same phase the interaction sections read as anchored by a fourth party
(\S\ref{sec:interactions}) --- which is exactly what Andrianov, Taron and Tarrach established for
neutral kaons in matter, where the environment breaks C, irreversibility breaks
T, and the effective parameters follow a Lindblad equation
\autocite{Andrianov2001}. The framework's reading is that the resulting
informational asymmetry \emph{constitutes} the arrow rather than presupposing it:
CP violation, entropy increase, and the direction of time are one informational
asymmetry seen three ways, not an independent second arrow, so the worry about a
direction built into spacetime does not arise. On this footing the contextual
dependence sets up a positive feedback loop for baryogenesis---an initial
fluctuation makes the environment slightly matter-dominated, contextual
enhancement raises the effective CP violation, more matter is produced, and the
loop runs---with $\delta$ a cosmological fixed point rather than a fundamental
constant. The honest limit of this is that a feedback loop explains amplification
and lock-in, not the sign selection or the baryon-number violation that
Sakharov's conditions still require \autocite{Sakharov1967}; the stronger claim
the framework is positioned to make, that the matter direction and the time
direction are co-selected by the same self-referential asymmetry, needs a source
of baryon-number violation and is left open.

\paragraph{Spatial infinitude and the status of inflation.}
There is no upper bound on non-computability, so the space of computations is
unbounded: the universe is spatially infinite, which leaves its spatial sections flat or of
negative curvature and excludes only the closed case. Infinitude on its own, however, does not
settle flatness, which is a fine-tuning of the Friedmann dynamics the framework has not yet
derived (Open Problem~\ref{op:comp-metric}); and flatness was in any case never inflation's
principal task. The horizon problem is, and an infinite space does nothing for it. The framework's actual lever against
inflation is the relational thesis of Section~\ref{sec:gr-nonlocality}: if
pre-spatial entanglement correlates regions before \textsc{fanout} spreads them
into spatial separation, then the observed large-scale homogeneity is fixed
relationally, before any question of causal contact, and the horizon problem
dissolves without an inflaton. This is a real argument, but it carries a real
bill: inflation also generates the nearly scale-invariant spectrum of primordial
fluctuations seen in the microwave background \autocite{Planck2020}, and any
account that dispenses with inflation owes an alternative origin for that
spectrum. That, not the horizon problem, is the hard empirical debt, and it is
named here rather than discharged.

\paragraph{Expansion and acceleration.}
If spacelike distance is the information separating worldlines, then as worldlines
accumulate distinct interaction histories they accumulate information-distance,
and the universe expands. The picture has the right exclusions built in: bound
systems interact continuously, share information, and so keep low
information-distance---they do not expand---while only worldlines that decouple
diverge, recovering expansion on cosmological scales alone. The feedback in which
greater separation means less light-cone overlap and so faster divergence,
$\mathrm{d}d/\mathrm{d}t \propto d$, yields exponential divergence, a de~Sitter
phase, and hence accelerating expansion without a separate dark-energy term
\autocite{Riess1998}. The limit to be honest about is that pure exponential
divergence is the late-time behaviour only; reproducing the full expansion history,
with its radiation- and matter-dominated eras, and fixing the Hubble rate
quantitatively, remains to be done.

\paragraph{The $S^{7}$ sector as internal structure.}
That $S^{7}$ cannot carry a metric does not make it disappear. Its closed
timelike curves carry the same non-computable information structure as those of
$S^{1}$ and $S^{3}$, but lacking associativity they support algebra---gauge
charge---in place of geometry---distance. This is why gauge symmetry has so often
looked like a set of compactified extra dimensions: the Kaluza--Klein and string
intuition is tracking something real in the $S^{7}$ structure, but misdiagnoses
it as spatial. In the framework these degrees of freedom have no size because
they have no metric; they are internal in the exact sense that they enter the
information-exchange that constitutes gauge charge while contributing nothing to
the loop-counting that constitutes distance. The resemblance to rolled-up
dimensions is therefore genuine and instructive, and the diagnosis that they are
algebraic rather than geometric is what separates this account from theories that
add real spatial dimensions to carry the same structure.

\paragraph{Loops, supertasks, and a second route to parallelisability.}
The argument of \S\ref{sec:three-dimensions} derives the parallelisable-sphere constraint from
closed timelike curves: a computation whose output is its own input requires information to be
transported consistently around a loop, and consistent transport selects the parallelisable spheres
$S^{1}$, $S^{3}$, $S^{7}$ \autocite{adams1960hopf}. There is a suggestion, which we record as a
prospect rather than a result, that the same constraint may govern the loops of perturbative quantum
field theory. The external legs of a Feynman diagram are on-shell classical boundary data; the
quantum content, and the divergences, enter with the internal loops, where the loop momentum is
integrated over all values and the action accumulated inside the loop is in this sense unbounded ---
a supertask, an infinite computation closed into a finite process. If such a loop is, in the
framework's terms, a closed computational path of the same kind as a closed timelike curve --- a
cycle carrying information not derivable from the data outside it --- then the parallelisability
constraint derived for closed timelike curves would apply to it, and the admissible interactions
would be constrained by the same $S^{1}$, $S^{3}$, $S^{7}$ structure that fixes the gauge group in
\S\ref{sec:entanglement-classes}. This would be a second and independent route to the Standard Model
content: not only \enquote{closed paths in spacetime select the parallelisable spheres} but \enquote{closed paths in interactions select them too}, unifying the spacetime arc and the interaction arc under one
principle. The decisive step, which we have not taken, is the identification of a field-theory loop
with a closed computational path in the precise sense the framework requires --- a loop integral is a
closed cycle in momentum space, a closed timelike curve a closed cycle in spacetime, and whether
these are the same kind of self-consistent closed structure is exactly what would have to be
established. It is offered as a direction, not a claim.

\paragraph{Does parallelisability determine renormalisability?}
A sharper and more self-contained question sits alongside the previous one, and we state it in the
hope that it will be taken up by those equipped for it. Renormalisability is the property that the
divergences of a quantum field theory can be consistently absorbed, leaving finite predictions; some
theories have it and some do not, and the standard criterion for which is technical (power-counting
of the coupling's dimension). The framework suggests a more structural reading. Parallelisable
spheres are precisely the spheres on which a global, consistent transport exists where it generically
does not; the non-parallelisable spheres are where the obstruction lives. If the divergence of a loop
is, in the framework's language, an obstruction to closing a computational path consistently, then
parallelisability would be exactly the structure that distinguishes the removable obstruction (a
renormalisable theory) from the irremovable one. It is at least a suggestive coincidence that the
renormalisable interactions of nature are the gauge theories built on $U(1)$, $SU(2)$ and $SU(3)$ ---
the groups the framework ties to the three parallelisable spheres --- and that gravity, which resists
renormalisation, is the sector \S\ref{sec:three-dimensions} ties to the geometry rather than the
algebra. Whether this is structural or accidental we do not know. The question --- whether
parallelisability, in the Hurwitz--Adams sense, is what determines renormalisability --- is posed in
the native language of the division-algebra approach to the Standard Model, and is one we think that
community is best placed to answer. We are careful to claim nothing more than that the question is
well-motivated: this is emphatically not a method for computing renormalised quantities, whose finite
remainders carry the precision agreement with experiment, but a conjecture about which theories admit
such a computation at all. It is, like the previous prospect, a direction with a possible path to an
answer rather than an answer, and it is left open deliberately.


\medskip
\noindent These mark the boundary of the programme as it stands. The substrate,
the generation of events, and the network status of dark matter are deferred to
the second paper; the contextual account of CP violation and baryogenesis, the
relational case against inflation, the information-divergence account of
expansion, and the algebraic reading of the $S^{7}$ sector are offered as
prospects, each with its outstanding debt named so that the line between argument
and aspiration stays visible.       <-- first consumer
%                                                          (S^7 exclusion)
%      % EDITED 2026-09-09 scripts/patch_spekkens_citekey_2026-09-09T1810.py
% =====================================================================
%  SECTION: Why Hopf Fibrations Belong in Particle Physics
%  (sec:hopf-justification). Drafted 2026-06-08; F1 arc added
%  prop:f1-dictionary and op:hopf-frobenius's discharge clause.
%  AUDITED 2026-07-06 (final pass, eighth stop).
%
%  PLACEMENT (reordered chain, JS 2026-07-06): position 8 --
%      \input{spacetime_from_non-computability}
%      \input{hopf_justification_section_2026-06-08}
%      \input{entanglement_classes_and_gauge_groups_2026-06-08}
%  No adjacency phrasing in the prose: position-portable. All
%  cross-refs live -- sec:fanout-direction DRAFTED 2026-07-06 (formerly the one undrafted
%  section, by design).
%
%  AUDIT RECORD: prop:gsm, prop:termination and prop:f1-dictionary
%  verified against every quoting consumer (statements canonical);
%  honest-debt subsection names both the selection assumption and the
%  mereology debt; the epistemic two-climbs disambiguation and the
%  compositional rung-evidence verified; sec:canonical-orbits live;
%  file pure ASCII. Fixes: op:hopf-frobenius's cross-link with
%  op:sigma-functor made TWO-WAY (the "four ways" line now names the
%  op); both remaining functor conditions ANCHORED --- the quaternionic
%  to the established fibre physics of sec:wf-reconciliation (residue:
%  clutching map vs strong complementarity), the octonionic to the
%  proved canonical-orbit images (umbrella label app:ghz-w-proofs
%  minted at the appendix for the purpose). Statuses unchanged.
%  British English; \autocite; no bullets in prose.
% =====================================================================

\section{Why Hopf Fibrations Belong in Particle Physics}
\label{sec:hopf-justification}

Invoking Hopf fibrations in a discussion of the Standard Model invites an immediate and fair
objection: that the move is decorative, an analogy between two pieces of attractive mathematics
with no load-bearing connection between them. The purpose of this section is to dispose of that
objection before the construction is put to work. The connection is not an analogy. It is
carried by two independent and independently published theorems, and the structures it relies
upon are not selected for convenience but compelled by a pair of classification results. We take
as already established the correspondence between qubit entanglement and the Hopf fibrations
themselves --- the subject of a mature literature
\autocite{mosseri2001hopf,bernevigchen2003geometry} --- and do not reproduce it here. What this
section addresses is the prior question of why that correspondence is entitled to carry physical
content at all.

\subsection{Two theorems, joined end to end}

The bridge from qubit entanglement to the Standard Model is built from two distinct results,
neither of which is original to the present framework and neither of which depends on the other.

The first is geometric. The unit sphere in the Hilbert space of $n$ qubits is the total space of
a Hopf fibration: $S^3 \to S^2$ for one qubit, $S^7 \to S^4$ for two, and $S^{15} \to S^8$ for
three. For two and three qubits the fibration is \emph{entanglement-sensitive} --- separability
is precisely the condition under which the base degenerates to a product of single-qubit Bloch
spheres and the additional dimensions associated with the higher imaginary directions collapse
\autocite{mosseri2001hopf,bernevigchen2003geometry}. This is a statement about the geometry of
state spaces and requires no physical interpretation whatever.

The second is algebraic. Singling out a preferred complex structure on the octonions ---
equivalently, fixing the split $\mathbb{O} \cong \mathbb{C} \oplus \mathbb{C}^3$ --- leaves a
residual symmetry, and that residual symmetry is exactly the gauge group of the Standard Model.
This was established through the exceptional Jordan algebra by Dubois-Violette and Todorov
\autocite{duboisviolette2016exceptional,duboisviolettetodorov2019exceptional,todorovduboisviolette2018deducing}
and given its sharpest spinorial form by Krasnov \autocite{krasnov2021so9}, who showed that the
relevant group is the stabiliser of a complex structure inside $\mathrm{Spin}(9)$. We record the
result in the form in which it enters the present work.

\begin{proposition}[Dubois-Violette--Todorov; Krasnov]\label{prop:gsm} Let $J_R$ be the complex
structure on the space $\mathbb{O}^2$ of $\mathrm{Spin}(9)$ spinors determined by a choice of
unit imaginary octonion, equivalently the split
$\mathbb{O} \cong \mathbb{C} \oplus \mathbb{C}^3$. The subgroup of $\mathrm{Spin}(9)$ commuting
with $J_R$ is \[ G_{\mathrm{SM}} \;=\; \frac{SU(3) \times SU(2) \times U(1)}{\mathbb{Z}_6}, \]
acting on the singled-out $\mathbb{C}$ summand as a colour singlet of hypercharge $-1$ and on
the $\mathbb{C}^3$ summand as a colour triplet of hypercharge $\tfrac{1}{3}$, both as $SU(2)$
doublets --- the quantum numbers of one left-handed generation.
\end{proposition}

The two results are independent, and the join between them is the only place a physical theory
must do work. Each result is a theorem; what neither supplies is a reason why nature should fix
one complex direction within $\mathbb{O}$ rather than treating all directions symmetrically. We
return to this --- the genuine load-bearing assumption --- in \S\ref{sec:hopf-honest-debt}, and
resolve it, in candidate form, in \S\ref{sec:fanout-direction}.

\subsection{The ladder is forced, not chosen}

It might still be objected that selecting the octonions, or the Hopf fibrations, is itself a
modelling choice carrying as much latitude as the choice of a grand-unified group. It is not,
and the reason is a pair of classification theorems that close off the alternatives entirely.

By Hurwitz's theorem the only normed division algebras are $\mathbb{R}$, $\mathbb{C}$,
$\mathbb{H}$ and $\mathbb{O}$ \autocite{hurwitz1898composition,baez2002octonions}. By Adams's
resolution of the Hopf-invariant-one problem --- together with the parallelisability results of
Bott--Milnor and Kervaire --- the only parallelisable spheres are $S^0$, $S^1$, $S^3$ and $S^7$,
so that there exist exactly three nontrivial Hopf fibrations, in bijection with $\mathbb{C}$,
$\mathbb{H}$ and $\mathbb{O}$
\autocite{adams1960hopf,bottmilnor1958parallelizability,kervaire1958nonparallelizability}. These
are not heuristics; they are among the more durable theorems in topology and algebra. Their
force here is that, once a framework is committed to amplitudes valued in a normed division
algebra --- a commitment that in the present construction is made far upstream, by the
self-referential argument of \S\ref{sec:self-reference} and \S\ref{sec:division-algebras}, and
not by fiat at this stage --- the ladder of Hopf fibrations is not selected but compelled, and
it terminates of its own accord.

\begin{proposition}[Termination]\label{prop:termination} The Cayley--Dickson tower of normed
division algebras and the tower of parallelisable spheres both terminate at real dimension
eight. Consequently there are exactly three nontrivial Hopf fibrations, the largest being the
octonionic $S^{15} \to S^8$, and the construction admits no extension to a fourth case.
\end{proposition}

The termination is worth dwelling on, because it is unusual for a model of this kind to predict
its own boundary. The next rung of the Cayley--Dickson sequence, the sedenions, possesses zero
divisors; the Hopf map can no longer be globally defined; there is no fourth fibration. This is
not a limitation imposed from outside but a structural fact, and it carries a physical
consequence that we develop later: there is no room in the construction for a fourth gauge force
or a fifth interaction, so that the discovery of either would count as evidence against the
framework rather than as an invitation to extend it. The endpoint is a theorem, not a
stipulation.

\subsection{Independent arrival as corroboration}

The most persuasive evidence that the destination is not an artefact of the route is that it is
reached by routes sharing none of the same machinery. Szangolies
\autocite{szangolies2025standardmodel} reaches $G_{\mathrm{SM}}$ from the emergent-gravity and
Kaluza--Klein programmes, reading the singled-out complex direction as a dimensional reduction
of a $9+1$-dimensional spacetime; the present framework reaches it from self-reference
(\S\ref{sec:self-reference}); and Furey and Gillard--Gresnigt reach the same gauge group, and
the same three-generation structure, through the minimal left ideals of the Dixon and Clifford
algebras, without invoking Hopf fibrations at all
\autocite{furey2018three,gillardgresnigt2019three}; the relation between their minimal ideals
and the entanglement classes used here is taken up in \S\ref{sec:canonical-orbits}. The
octonionic route to particle physics is in any case an old one, going back to the use of
$\mathbb{O}$ to model the strong force \autocite{gunaydingursey1974quark}.

That programmes with disjoint starting points converge on
$SU(3) \times SU(2) \times U(1)/\mathbb{Z}_6$ and on the same quark--lepton division is the
coincidence Szangolies himself remarks upon: the characteristic structure of the Standard Model,
routinely regarded as arbitrary, turns up independently in unrelated settings, and its appearing
twice would then be doubly odd unless the two settings are in fact related. We read the
convergence the same way --- as evidence that the division-algebraic description is singled out
by the mathematics rather than imposed by the modeller, and therefore that
Proposition~\ref{prop:gsm} is structural rather than contrived.

\subsection{The honest debt}
\label{sec:hopf-honest-debt}

It would be misleading to leave the impression that this route is free of assumptions, so we
state plainly the one assumption that all of these programmes share. Proposition~\ref{prop:gsm}
requires that \emph{some} complex direction in $\mathbb{O}$ be singled out, and the uniqueness
of that choice is not supplied by the theorem. Krasnov parametrises the complex structure $J_R$
by an arbitrary unit imaginary octonion, and the residual group is the same for every such
choice because $G_2 = \mathrm{Aut}(\mathbb{O})$ acts transitively on the unit imaginary
octonions. The theorem therefore guarantees the gauge group \emph{given} that a direction is
fixed, while remaining silent on \emph{why} one is fixed. Szangolies treats this as a free
choice and flags it, in his own conclusions, as an open element of the construction --- one to
which several inequivalent answers are available.

This, and not the appeal to Hopf fibrations, is the genuine load-bearing assumption of the whole
approach. The Hopf fibrations are forced (Proposition~\ref{prop:termination}); the singling-out
is not. The present framework offers a candidate resolution that the prior literature lacks: the
act of measurement, modelled as \textsc{fanout}, selects the direction physically, and the
transitivity of $G_2$ then guarantees that the gauge group is independent of which direction the
measurement happens to select. This argument is given in \S\ref{sec:fanout-direction}; we flag
here, at the appropriate epistemic status, that it is structural rather than fully proved, and
that the residual question --- why a measurement selects \emph{exactly one} imaginary direction
rather than merely \emph{permitting} a selection --- remains open. It is now priced as Open
Problem~\ref{op:fanout-selection}.

A second debt is shared with every programme of this character and is not closed by anything
said here: why the abstract quantum state should admit a tensor-product decomposition into
qubits in the first place. This is the quantum factorisation, or mereology, problem. The
entanglement-sensitivity on which the construction turns is meaningful only once such a
decomposition is fixed, and we treat the decomposition as a necessary input rather than a
derived result. Naming it is preferable to letting it pass unremarked.

\subsection{The termination read epistemically}
\label{sec:hopf-information-bound}

The termination of the tower (Proposition~\ref{prop:termination}) has so far been stated
algebraically: the Cayley--Dickson sequence acquires zero divisors at the sedenions and the Hopf
map can no longer be globally defined. There is a second reading of the same boundary, epistemic
rather than algebraic, which we offer as a motivated conjecture rather than a settled result,
because it connects the tower's endpoint to a bound that arises independently in the foundations
of measurement.

A particle, in this framework, is a reified relation between events, and its internal content is
carried by the octonionic $A,B,C$ register --- the $S^{7}$ sector serving as internal degrees of
freedom in the channel that connects the events (\S\ref{sec:space-of-computations}). The
termination of the tower is then a statement about that content: a single such relation can
support genuine entanglement among at most three parties, because there is no fourth fibration
to carry a fourth independent party within the same irreducible relation. Every particle is, on
this reading, a \emph{completed} octonionic-level relation --- topped out at three parties not
by circumstance but because the tower admits no further rung. The information that can be
encoded in one relation-between-events, and hence extracted by one measurement of it, is
correspondingly bounded: it is the information of a completed three-party register, and no
measurement of a single particle can yield more.

This is exactly the kind of bound that the epistemic-restriction programmes reach from the other
direction. Spekkens' toy theory is built on a knowledge-balance principle that caps the
information obtainable about a system \autocite{spekkens2007toy}, and Szangolies derives such an
epistemic horizon not by positing it but from the same self-referential non-computability on
which the present tower rests \autocite{szangolies_j_2018}. That the two accounts --- the
algebraic termination of the Hopf tower and the epistemic bound on measurement information ---
should coincide is, we suggest, not a numerical accident but a shared root: both are
consequences of the self-referential non-computability of \S\ref{sec:self-reference}, wearing an
algebraic face in the one case and an epistemic face in the other. The Hurwitz--Adams ceiling
and the knowledge-balance ceiling would then be one limit seen twice, and Szangolies' programme,
which reaches $G_{\mathrm{SM}}$ by an independent route (\S\ref{sec:hopf-honest-debt}) and the
epistemic horizon from the same non-computable source, shares with this framework not merely a
conclusion but its origin.

The reading must be kept distinct from a different climb, with which it would otherwise be
conflated. The Wigner's-friend regress of \S\ref{sec:wigners-friend} ascends without end: each
observer may be measured by a further observer, each carrying a fresh gauge-private circle, and
that chain has no top --- it is the bump in the carpet, unbounded in length. The termination
bounds something else: not the length of the chain of observers, but the internal structure of
each single relation, each particle, that the chain is built from. The regress is unbounded
precisely because each of its links is a bounded, completed three-party object;
termination-of-each-link and non-termination-of-the-chain are not in tension but complementary,
the former individuating the links of the latter. What the bound caps is the information in one
particle; what the regress leaves uncapped is the number of particles and observers that can be
chained. The precise identification of ``the information of a completed three-party register''
with the formal measures of the epistemic-restriction literature is the technical step that
would turn this conjecture into a theorem, and we flag it as owed; the structural claim --- that
the tower's endpoint and the measurement bound share the non-computable root --- is what we
advance here.

\subsection{The termination read compositionally}
\label{sec:hopf-compositional}

There is a third termination theorem, proved in a community that has so far had no contact with
the two joined above. Coecke and Kissinger showed that the GHZ and W classes --- realised as the
commutative special and the anti-special Frobenius algebras of categorical quantum mechanics ---
are \emph{compositionally complete}: every $N$-qubit entanglement class under SLOCC is generated
by composing these tripartite primitives with bipartite entanglement and product states
\autocite{coeckekissinger2010compositional}. Above three parties nothing new appears; the
entanglement alphabet closes at exactly the point where the fibration tower terminates. The
mutual silence of the two literatures is itself striking: Coecke and Kissinger make no reference
to Hopf fibrations or division algebras, and the fibration literature on entanglement
\autocite{mosseri2001hopf,bernevigchen2003geometry} nowhere cites compositional completeness.
Two classifications, in different mathematics, draw the same line in the same place, unaware of
one another. That silence has, in the months this paper was in draft, just begun to break from
the categorical side, with a projective, base-level reading of the same calculus
\autocite{burtonanwar2026meromorphic}; we engage it below.

The framework's conjecture is that this is one termination wearing a third face: the alphabet
has the letters it has --- Bell, GHZ, W --- because the parallelisable spheres have the
dimensions they have. There is partial evidence at every rung. At one qubit, the commutative
special Frobenius algebras on $\mathbb{C}^{2}$ are classified by orthonormal bases --- points of
$S^{2}$, the base of the first fibration --- so that the copied basis plays the role of a
section of the bundle and the spider theorem expresses its idempotence. At two qubits, the
$SU(2)$ fibre of $S^{7} \to S^{4}$ carries exactly the entanglement that survives local
operations, and L\'{e}vay's result that the bundle's natural connection is the $SU(2)$
Yang--Mills instanton \autocite{levay2004geometry} stands on the geometric side where the
interaction of two Frobenius structures stands on the categorical one. At three qubits, the
split of D\"{u}r, Vidal and Cirac \autocite{durvidalcirac2000three} mirrors the algebra exactly:
GHZ is genuinely tripartite, demanding the full non-associative octonion, while W is
pairwise-robust and housed in the associative subalgebra $\mathbb{H} \subset \mathbb{O}$ --- the
special/anti-special Frobenius distinction in geometric dress.

The first of the four conditions the correspondence requires can be discharged at once, and we
do so here, because a problem whose first step has been walked is a different kind of invitation
from one merely posed.

\begin{proposition}[The one-qubit dictionary]\label{prop:f1-dictionary} The $\dagger$-special
commutative Frobenius algebras on $\mathbb{C}^{2}$ correspond bijectively to
\emph{lifted antipodal skeletons} of the complex Hopf fibration $\pi \colon S^{3} \to S^{2}$:
unordered antipodal pairs $P \subset S^{2}$ together with a section of $\pi$ over $P$. The
copyable states of a structure are exactly the section's image, so the lifts are recorded
faithfully --- the projectively trivial centre of $U(2)$ acts non-trivially on structures ---
and the correspondence is $U(2)$-equivariant, the ordered structures forming a single free and
transitive $U(2)$-orbit. The channel a structure induces on density matrices is dephasing in its
skeleton and is independent of the section: the fibre freedom, a torsor for
$\mathrm{Map}(P, U(1))$, is exactly the data absent from the classical record a structure
induces.
\end{proposition}

\begin{proof}
By the classification of Coecke, Pavlovic and Vicary, such structures correspond to orthonormal
bases, recovered as the sets of copyable states \autocite{coeckepavlovicvicary2013orthogonal}.
An orthonormal basis of $\mathbb{C}^{2}$ is two orthogonal rays --- an antipodal pair on the
Bloch sphere --- together with a unit representative of each, which is precisely a section of
$\pi$ over the pair; and $\delta(e^{i\theta}x) = e^{i\theta}\, x \otimes x$ against
$(e^{i\theta}x) \otimes (e^{i\theta}x) = e^{2i\theta}\, x \otimes x$ pins the phases.
Equivariance, the action of the centre, the orbit structure and the induced channel are direct
computations; an executable verification of every finite check accompanies the code release.
\end{proof}

Two remarks fix the proposition's scope. It corrects a natural first reading: a classical
structure does not project to the base but selects lifts into the total space --- classicality
is section data, and what the copy operation broadcasts is a lift, not a ray. And it claims
nothing topological about cloning: the finiteness of the copyable set is linear, structural
no-cloning residing in non-naturality \autocite{abramskyheunen2012hstar}, so the first place the
bundle's global topology can bite is the quaternionic rung, where the degree-one clutching map
stands ready to be compared with the strong-complementarity law governing an interacting pair of
structures. The dictionary is, finally, the complement of the projective treatment cited above:
Burton and Anwar pass functorially to the base and deliberately discard the fibre
\autocite{burtonanwar2026meromorphic}, and the classical structures turn out to live in exactly
the data their projectivisation forgets --- the two readings compose, and the resulting triangle
is a consistency condition available to constrain the remaining conditions of the
correspondence.

\begin{openproblem}[The Hopf-fibration--Frobenius correspondence]\label{op:hopf-frobenius} Let
$\mathbf{Hopf}$ denote the category of the three fibrations
$S^{2^{k}-1} \hookrightarrow S^{2^{k+1}-1} \to S^{2^{k}}$, $k = 1, 2, 3$, with bundle maps as
morphisms, and let $\mathbf{Frob}$ denote the Frobenius algebras on $\mathbb{C}^{2}$ in
$\mathbf{FdHilb}$, with the commutative-special (GHZ) and anti-special (W) classes
distinguished. Construct a functor $F \colon \mathbf{Hopf} \to \mathbf{Frob}$ --- or prove that
none exists --- under which the complex fibration yields the single-qubit classical structures
--- Proposition~\ref{prop:f1-dictionary} identifies these as the lifted antipodal skeletons,
section data over two base points rather than the points alone; the quaternionic fibration
yields the interaction of two such structures, with the $SU(2)$ fibre carrying the entanglement
that survives the classical quotient; the octonionic fibration yields the GHZ/W pair, GHZ on the
full non-associative fibre and W on the associative $\mathbb{H} \subset \mathbb{O}$; and the
composition by which Coecke--Kissinger generate the $N$-qubit classes corresponds to a
well-defined geometric operation on the bundles. If such an $F$ exists, compositional
completeness follows from Adams's theorem as a topological corollary --- there is no fourth
fibration to supply a fourth primitive --- and if $F$ is faithful enough that the implication
also runs back, the termination of the tower and the completeness of the alphabet are equivalent
expressions of a single classification. The first condition is discharged by
Proposition~\ref{prop:f1-dictionary}; what remains open is the quaternionic and octonionic
conditions and, above all, the correspondence of compositions. Both open conditions are already
constrained: the quaternionic fibre's physical content --- the recoverable relative phase and
its purity ledger --- is established (\S\ref{sec:wf-reconciliation}), leaving the
clutching-map-to-strong-complementarity comparison as the structural residue; and any candidate
functor must reproduce the proved canonical-orbit images, GHZ to the colour triplet and W to the
singlet (\S\ref{sec:canonical-orbits}, Appendix~\ref{app:ghz-w-proofs}). A proof that no such
functor exists, or a correspondence at the level of the classifications without functoriality on
compositions, would each also resolve the problem: the first would make the agreement of two
independent classifications a fact demanding its own explanation; the second would establish the
weaker reading.
\end{openproblem}

A resolution in the affirmative would do three things. It would explain compositional
completeness rather than merely prove it --- a corollary of the division-algebra structure of
quantum mechanics rather than an independent algebraic fact. It would give the diagrammatic
language of this paper's own interaction circuits a topological warrant: the GHZ/W spider
calculus would be the unique language the topology permits, not a convenient notation. And it
would open a two-way bridge, translating problems on compositions of Frobenius structures into
operations on bundles and back. The problem sits at the intersection of categorical quantum
mechanics, the quantum-information geometry of the fibrations, and the division-algebra
approaches to particle physics, and it is posed here so that any of the three communities can
take it up in its own language. It is perhaps worth noticing that both classifications already
answer to the same name --- the fibrations are Heinz Hopf's, and so, through the bialgebra
structure of the diagrammatic calculus, are the Hopf algebras --- a coincidence of eponymy that
the conjecture proposes to promote to a theorem. With the epistemic reading above, an
affirmative resolution would leave the tower's endpoint as one limit seen four ways: algebraic,
geometric, compositional, epistemic --- the last being the fixed-point classification of Open
Problem~\ref{op:sigma-functor}, with which this problem is cross-linked, each constraining the
other.

\subsection{What this sets up}

The entanglement-sensitivity of the Hopf maps is not a side feature but the hinge of the
physical reading that follows. Because separability is exactly the collapse of the base to a
product of Bloch spheres, the structure of the entanglement --- the SLOCC class of a multi-qubit
state, and within it the canonical maximally entangled orbit --- is carried by precisely the
degrees of freedom on which $G_{\mathrm{SM}}$ acts in Proposition~\ref{prop:gsm}. This is what
makes it possible, in \S\ref{sec:entanglement-classes}, to read the two genuinely tripartite
entanglement classes onto the $\mathbb{C} \oplus \mathbb{C}^3$ split through their canonical
representatives: the GHZ orbit confined to the colour-triplet $\mathbb{C}^3$ sector, the W orbit
reaching the $\mathbb{C}$ singlet, recovering the quark--lepton division not as an assignment
but as a consequence of the entanglement structure. The scope of those statements --- canonical
orbits rather than whole classes, and the identification of the canonical orbits with the
physical particles --- is made precise there. The Hopf fibrations are the bridge; the
entanglement classes, treated next, are the traffic that crosses it.

%  The section is position-portable: all adjacency phrasing was
%  retargeted to explicit \S\ref on 2026-07-06. The S^7 hand-off
%  sentence to sec:space-of-computations WRITTEN 2026-07-06 (spacetime
%  file read; the closing now hands first to its immediate consumer,
%  then to hopf-justification). Upload verifications COMPLETE: L2 and
%  quantum-postulates do not consume the ladder; WF's references are
%  forward-voiced or self-carrying.
%
%  DIVISION OF LABOUR (deliberate, to avoid overlap):
%   - sec:self-reference (undrafted): self-reference -> complex
%     amplitudes and the Born rule; why not a real base. This section
%     TAKES that as given and starts at C.
%   - THIS section: the algebra -- the doubling construction, the
%     ledger of losses-as-physics, the internal-algebras-not-a-new-QM
%     scoping (Adler pre-empt), the two appearances of the ladder
%     (register / measurement climb, stated lightly), and Hurwitz as
%     the algebraic termination with the committed zero-divisor
%     reading.
%   - sec:hopf-justification: the geometry (Adams, parallelisability,
%     the fibrations), the join of the two classifications, the
%     entanglement-sensitivity literature, the C-direction honest debt,
%     and the epistemic reading of termination.
%
%  COMMITTED READINGS USED (sources re-read 2026-07-03):
%   - Loss ledger and level dictionary: cayley_dickson_qubits_for_review
%     sec.5 ("each step loses a symmetry, and each lost symmetry is
%     identified with a physical content").
%   - Termination reading: grounding_argument_gauge_group ("the
%     sedenions have zero divisors and cannot support consistent
%     information flow around loops") -- leads; norm-composition-as-
%     probability given as its amplitude-level restatement.
%   - Internal-not-base scoping: cayley_dickson_qubits_for_review Q2 and
%     sec.3.2 caveat (Adler objection; C-qubit treated geometrically).
%
%  HONEST DEBT: ladder proved closed at the top (Hurwitz), derived at
%  the floor. The rung-compulsion revisit-trigger FIRED 2026-07-06:
%  sec:self-reference and op:sigma-functor now exist, so the thesis is
%  minted as op:ladder-compulsion (with the CD-conjugation-as-sigma
%  candidate route), inheriting op:sigma-functor's status.
%
%  AUDITED 2026-07-06 (final pass, third stop): all four citations
%  used; all cross-refs live; CD facts, dims, alternativity, Hurwitz
%  statement verified; Adler subsection delivers everything the
%  register's compressed pointer promises; file pure ASCII; no bare
%  "frame". Fixes: opening rewritten for both predecessors (continuity
%  was inserted between after drafting); "associator's non-vanishing is
%  mass itself" recalibrated to the register's debt-against-selection
%  defusal; the O->S ledger row and Hurwitz subsection reconciled with
%  the fourth-qubit forward remark ("the terminus closes the ladder of
%  forces; what survives past it is, at most, flavour").
%
%  CROSS-REFS: sec:self-reference [sibling dangle -- resolves when that
%  section is drafted], sec:hopf-justification, sec:hopf-information-bound,
%  sec:register, sec:entanglement-classes, sec:generations, sec:masses,
%  sec:wigners-friend. All others exist.
%  CITATIONS: existing keys baez2002octonions, hurwitz1898composition,
%  furey2018three. ONE new key: adler1995quaternionic (RIS in
%  division_algebras_refs_2026-07-03.ris).
%  British English; \autocite; no bullets in prose; the ledger is a
%  table because the comparison is the argument.
% =====================================================================

\section{The Division-Algebra Ladder}
\label{sec:division-algebras}

The foundations behind us end at amplitudes twice over: natively, where the self-referential
argument forced the description of a system into the system's own state space and priced
consistency at a number system whose parts can cancel --- complex amplitudes, and the Born rule
that squares them (\S\ref{sec:self-reference}) --- and independently, where the regime's
continuity delivers the same field through Hardy's theorem (\S\ref{sec:continuity}), the
convergence of the two branches being evidence rather than accident. This section asks the question
both routes leave open. Why
should the number system stop at $\mathbb{C}$? Ordinary quantum mechanics stops there by postulate;
a framework that has forbidden itself postulates must answer, and the answer is that the climb
neither stops at $\mathbb{C}$ nor continues forever. There are exactly two more rungs, each
purchased by surrendering an algebraic law and each paying for itself in physics, and the ladder is
closed at the top by a theorem of 1898. This section sets out the algebra of that ladder: the single
construction that climbs it, the ledger of what each rung costs and buys, and the theorem that ends
it. The geometry of the same ladder --- that each rung acts on a sphere, that exactly those spheres
are parallelisable, and that the algebraic and geometric classifications are one --- is the business
of the geometry section (\S\ref{sec:hopf-justification}).

\subsection{One construction climbs the whole ladder}

The rungs are not a miscellany; a single construction generates them all. Given an algebra with a
conjugation, the Cayley--Dickson construction builds its double on ordered pairs, with a twisted
multiplication in which the conjugate enters and a conjugation $(a,b)^{*} = (a^{*}, -b)$
\autocite{baez2002octonions}. (Conventions differ in the placement of the conjugations; all yield
isomorphic algebras.) Applied repeatedly it produces the sequence
$\mathbb{R} \to \mathbb{C} \to \mathbb{H} \to \mathbb{O} \to \mathbb{S}$, of real dimensions
$1, 2, 4, 8, 16$: the reals, the complex numbers, the quaternions, the octonions, the sedenions.
The ground floor is classical: $\mathbb{R}$ is ordered and positive-definite, the arithmetic of
classical probability. The first doubling is therefore the quantum move itself --- the passage from
probabilities to amplitudes that can interfere --- and it is not made here by choice: it is the
derivation of \S\ref{sec:self-reference}, already behind us. What this section must account for is
the rest of the climb.

\subsection{The ledger}

In mathematics the doublings are usually told as a story of degeneration: each algebra is worse
behaved than the last, and the sequence is studied for where it breaks. This framework reads the
same ledger in the opposite direction. Each rung is bought by surrendering a law, and each
surrendered law is exactly a piece of physics --- the losses are the assets, and the comparison is
sharp enough to be the argument:

\begin{center}
\begin{tabular}{llll}
\hline
doubling & dimension & what is lost & what the loss becomes \\
\hline
$\mathbb{R} \to \mathbb{C}$ & $1 \to 2$ & the order of the reals & phase and interference; amplitudes \\
 & & & that cancel; the Born rule (\S\ref{sec:self-reference}) \\
$\mathbb{C} \to \mathbb{H}$ & $2 \to 4$ & commutativity & non-commuting actions: weak isospin \\
 & & & $SU(2)$, the $B$ register (\S\ref{sec:register}) \\
$\mathbb{H} \to \mathbb{O}$ & $4 \to 8$ & associativity & the associator: colour, chirality, the \\
 & & & generations, and mass as debt \\
 & & & (\S\ref{sec:generations}, \S\ref{sec:masses}) \\
$\mathbb{O} \to \mathbb{S}$ & $8 \to 16$ & division itself & no further force: zero divisors appear \\
 & & (zero divisors) & and the ladder of sectors ends (below) \\
\hline
\end{tabular}
\end{center}

Walking the rows: the order lost at the first doubling is the possibility of interference --- parts
of a probability that can oppose one another --- which is where the derivation of \S\ref{sec:self-reference}
lives. The commutativity lost at the second is the existence of actions whose order matters; the
unit quaternions form the group $SU(2)$, and the framework houses weak isospin there, as the
internal algebra of the $B$ register whose order-sensitive, non-commuting dynamics the interaction sections develop
(\S\ref{sec:register}, \S\ref{sec:entanglement-classes}). The associativity lost at the third ---
the octonions retain the weaker property of alternativity \autocite{baez2002octonions} --- creates
the associator $[a,b,c] = (ab)c - a(bc)$, and the associator is where most of this paper's physics
is transacted: the stabiliser of a preferred complex direction inside the octonions' automorphisms
is the colour group (\S\ref{sec:hopf-justification}); the triality structure and its $K_3$
obstruction generate the three generations and the sterile channel (\S\ref{sec:generations}); and
the associator is the currency of mass --- the debt against the selected direction that
\S\ref{sec:masses} prices, owed by whatever action moves the selection and not by residence at
the rung, which is the defusal \S\ref{sec:register} makes explicit. The algebraic facts
in the third column are theorems; the identifications in the fourth are established in the sections
cited, each carrying its own epistemic flag there, and are not re-litigated here. What this table
asserts on its own account is only the correspondence's shape: one law surrendered per rung, one
sector of physics gained per law.

\subsection{Internal algebras, not a new quantum mechanics}

A reader versed in the quaternionic-quantum-mechanics literature will at this point raise a fair
objection: quantum theory over $\mathbb{H}$ as a foundation is known to be troubled --- the
tensor-product and composite-system difficulties set out by Adler and others
\autocite{adler1995quaternionic} --- and quantum theory over $\mathbb{O}$ founders earlier still,
since a non-associative algebra does not furnish a Hilbert space at all. The framework's position
must therefore be stated exactly, because it does not make the claim the objection targets. The base
of the physical Hilbert space here is $\mathbb{C}$ and remains $\mathbb{C}$; the higher algebras
enter as the \emph{internal} algebras of the register --- $\mathbb{H}$ for the $B$ qubit,
$\mathbb{O}$ for the $C$ qubit (\S\ref{sec:register}) --- in a tensor structure whose outer theory is
ordinary complex quantum mechanics. This is the same move by which the division algebras enter the
constructions of Furey and the related programmes without any proposal to rebase quantum mechanics
itself \autocite{furey2018three}. For the octonionic register the point is sharper still:
$\mathbb{O}^{2}$ is not a Hilbert space, and the framework accordingly treats the $C$ qubit not in
the Dirac--von Neumann manner but geometrically, as the sphere $S^{15}$ with its Hopf structure ---
which is precisely the form in which \S\ref{sec:hopf-justification} needs it. Nothing in this paper therefore
stands or falls with quaternionic or octonionic quantum mechanics as foundations; the ladder is
climbed by the registers' internal algebras, not by the theory's base field.

\subsection{The same ladder, twice}

The ladder appears in this framework in two guises, and it will help the reader to have them
distinguished now. Statically, it is the anatomy of a single relation: the three qubits of the
register take their values one rung apart --- $A$ over $\mathbb{C}$, $B$ over $\mathbb{H}$, $C$ over
$\mathbb{O}$ --- so that every particle carries the whole ladder as its internal structure
(\S\ref{sec:register}). Dynamically, it is the shape of measurement: a single qubit's pure states
fill $S^{3}$, an entangling measurement carries the joint state into $S^{7}$, and a third party into
$S^{15}$ --- the correspondence between qubit entanglement and the Hopf fibrations that
\S\ref{sec:hopf-justification} takes from a mature literature --- so that each act of measurement ascends one rung --- the climb the reader has already made,
narrated in \S\ref{sec:wigners-friend} --- and every completed relation is a terminated climb, which is
the reading \S\ref{sec:hopf-information-bound} gives the tower's endpoint. One ladder, seen twice:
as the anatomy of a relation and as the dynamics of relating. The framework regards this double
appearance not as a coincidence but as the same structure viewed statically and in motion, and the
sections cited carry that identification's weight.

\subsection{Hurwitz closes the ladder}

The construction itself would double forever; what stops it is a classification theorem. By
Hurwitz's theorem, the normed division algebras --- those in which the norm composes,
$\lvert ab\rvert = \lvert a\rvert\,\lvert b\rvert$, and no non-zero elements multiply to zero ---
are exactly $\mathbb{R}$, $\mathbb{C}$, $\mathbb{H}$ and $\mathbb{O}$
\autocite{hurwitz1898composition,baez2002octonions}. The next doubling, the sedenions, has zero
divisors: non-zero elements whose product vanishes. The framework's reading of this failure is the
one that matters physically. Zero divisors cannot support consistent information flow around loops
--- a cycle of compositions can silently annihilate what it carries --- and a framework in which
closed computational paths are load-bearing (\S\ref{sec:hopf-justification},
\S\ref{sec:space-of-computations}) cannot be built on an algebra that permits it. At the level of
amplitudes the same failure reads: the composition law
$\lvert ab\rvert = \lvert a\rvert\,\lvert b\rvert$, which carries probability multiplicatively
through every rung of the ladder, fails beyond the octonions, so that two non-vanishing amplitudes
could compose to nothing and probability would be destroyed by mere composition. The last algebra
fit to carry amplitudes is $\mathbb{O}$. One rung past the terminus is not entirely silent,
however: the sedenions' internal structure --- their automorphism group and their octonion
subalgebras --- is the recorded candidate arena for tagging the generations
(\S\ref{sec:interactions}), and it enters exactly as this section's scoping permits, internally
and never as a base. The terminus closes the ladder of forces; what survives past it is, at most,
flavour. This is the algebraic half of the termination; the
geometric half --- that no sphere above $S^{7}$ is parallelisable --- and the theorem joining the
two halves are the opening business of \S\ref{sec:hopf-justification}, and the epistemic reading of the same endpoint,
as a bound on the information a single relation can carry, is given in
\S\ref{sec:hopf-information-bound}.

\subsection{The honest debt}

As in the geometry section, we close by stating plainly what is proved, what is established, and what is
owed. Proved: the construction, the losses, and the termination are theorems --- the Cayley--Dickson
doubling and its ledger of surrendered laws are standard algebra, and Hurwitz's classification is
among the durable results of the subject. Established: the physical cashing of each rung --- the
Born rule at $\mathbb{C}$, weak isospin at $\mathbb{H}$, colour, generations and mass at
$\mathbb{O}$ --- is carried out in the sections cited above, each under its own epistemic flag. Owed,
and stated without disguise: the ladder as presented here is closed from above and derived at its
floor, but it is not yet derived rung by rung from below. Hurwitz forbids a further rung; the
argument of \S\ref{sec:self-reference} forces the first doubling, from the ordered arithmetic of
classical probability to amplitudes that interfere. That the same self-referential closure, applied
to a system whose amplitudes already take values at one rung, compels the next --- that relational
closure at one level forces precisely the Cayley--Dickson double of that level --- is the
framework's reading of the climb, and it is a thesis, not a theorem --- but the formal apparatus
against which it can be posed now exists, and we pose it.

\begin{openproblem}[The compulsion of the intermediate rungs]\label{op:ladder-compulsion}
Derive, within the categorical apparatus of Open Problem~\ref{op:sigma-functor}, that
self-referential closure at one rung compels the Cayley--Dickson double of that rung, so that the
intermediate steps $\mathbb{C} \to \mathbb{H} \to \mathbb{O}$ follow from the same principle
that forces the floor. The candidate route is the construction's own shape: the Cayley--Dickson
pair with its conjugation, $(a, b)^{*} = (a^{*}, -b)$, is an involution-equipped doubling, and the
route reads it as the algebraic shadow of adjoining the describer --- closure of level-$A$
descriptions under the inside/outside exchange requiring exactly the twisted pair. The problem
inherits the status of Open Problem~\ref{op:sigma-functor}, whose apparatus it presupposes.
\end{openproblem}

If the derivation lands, the whole ladder becomes a consequence of the same self-reference that
opened it. Until then, the honest statement is that the ladder's necessity is established at its
ceiling and at its floor, while the compulsion of the intermediate rungs is inherited from the
physics they are shown to carry rather than derived in advance. What this section hands on is, accordingly, an algebraic ladder with a floor, a ceiling and a
ledger. Its first consumer is immediate: \S\ref{sec:space-of-computations} spends the terminal
fact --- the non-associativity of $S^{7}$'s own multiplication --- as geometry, excluding the
seven-sphere as a ruler and fixing the dimension of space at three.
\S\ref{sec:hopf-justification} then supplies the geometry of the fibrations themselves, and the
join.

%      % EDITED 2026-09-09 scripts/patch_spekkens_citekey_2026-09-09T1810.py
% =====================================================================
%  SECTION: Why Hopf Fibrations Belong in Particle Physics
%  (sec:hopf-justification). Drafted 2026-06-08; F1 arc added
%  prop:f1-dictionary and op:hopf-frobenius's discharge clause.
%  AUDITED 2026-07-06 (final pass, eighth stop).
%
%  PLACEMENT (reordered chain, JS 2026-07-06): position 8 --
%      % EDITED 2026-10-07 scripts/patch_infinite_not_flat_2026-10-07T1400.py
% EDITED 2026-09-27 scripts/patch_spacelike_inversion_2026-09-27T1800.py
% EDITED 2026-09-21 scripts/patch_op_sterile_gravity_boundary_2026-09-21T1300.py
% EDITED 2026-09-12 scripts/patch_spacetime_bloch_headline_2026-09-12T1400.py
% EDITED 2026-09-11 scripts/patch_splice_loops_prospects_2026-09-11T0900.py
% EDITED 2026-09-10 scripts/patch_cite_rubino_2026-09-10T0700.py
% EDITED 2026-09-10 scripts/patch_quoting_and_typography_2026-09-10T0620.py
% =====================================================================
%  SECTION: Spacetime from non-computability (sec:space-of-computations).
%  AUDITED 2026-07-06 (final pass, seventh stop).
%
%  PLACEMENT (reordered chain, JS 2026-07-06): position 7 --
%      \input{division_algebras_section_2026-07-03}   <-- hands off here
%      \input{spacetime_from_non-computability}
%      \input{hopf_justification_section_2026-06-08}
%  Division-algebras' closing now carries the S^7 hand-off sentence
%  (written 2026-07-06 after this file was read).
%
%  INPUT DEPENDENCIES (both held in the audited outputs):
%      entropic_gravity_family_paragraph_2026-07-03.tex  (sec:jacobson)
%      gr_nonlocality_revision_2026-06-13.tex            (sec:gr-nonlocality)
%  The 2026-07-04 additions (asymmetry paragraph; Penrose second route;
%  acceleration/Unruh arc) are FULLY PASTED; the separately uploaded
%  additions file is redundant and can be archived.
%
%  AUDIT RECORD: opener formalised (Russell typo; bare ref retargeted
%  to the freshly minted sec:complex-amplitudes); the "rotations of
%  S^7" wording made precise at BOTH occurrences (Moufang loop, not a
%  group; SO(8) disclaimer added -- referee hazard removed); two \\
%  paragraph-break bugs fixed; dark-matter paragraph retargeted to
%  sec:generations and stitched to op:sterile-gravity (the zoo's
%  demotion and this paragraph's priority condition are one question);
%  CP paragraph stitched to the fourth-party reading of
%  sec:interactions; stale mid-file mini-header retired
%  (sec:k3-sterile placeholder dies with it); the section's narrated
%  debts NUMBERED at last -- op:comp-metric and op:three-dimensions
%  minted in the Open-problems subsection, cited from the in-text
%  "named open problems" line and from the closing summary.
%
%  KEYS: adams1960hopf matches corpus. Unique to this section, presumed
%  in the .bib since the paper compiles: Spence-Brown-1969 (suggest
%  normalising the key spelling at leisure), VanRaamsdonk2010,
%  MaldacenaSusskind2013, MuellerMasanes2013, MasanesMuellerEtAl2014,
%  toffoli_2003, Chaitin1998, Penrose1996, HohnMueller2016,
%  GarnerMuellerDahlsten2017, Jacobson1995, Bekenstein1973, tHooft1993,
%  Susskind1995, spekkens2007toy, Andrianov2001, Sakharov1967,
%  Planck2020, Riess1998. Shared keys verified: small_2006, Penrose1994,
%  unruh1976notes (swap-flag stands), baez2002octonions, Wheeler1990,
%  furey2018three, DiBiagioRovelli2021, Rovelli1996a, Spekkens toy.
%  British English; \autocite; no bullets in prose.
% =====================================================================

\section{Spacetime from non-computability}
\label{sec:space-of-computations}

This section is a deliberate diversion: it takes the idea that non-computable steps are measured
by imaginary information (\S\ref{sec:complex-amplitudes}) and asks what geometry it buys.
Spencer-Brown records that Bertrand Russell thought well of the underlying move
\autocite{Spence-Brown-1969}; the reader may judge whether the framework spends it responsibly.

The programmes that derive spacetime from quantum mechanics---Van~Raamsdonk's
\enquote{building spacetime with entanglement} \autocite{VanRaamsdonk2010}, the
ER\,$=$\,EPR conjecture \autocite{MaldacenaSusskind2013}, and M\"uller's
information-theoretic recovery of three-dimensional space from the qubit
\autocite{MuellerMasanes2013}---take entanglement or information as primitive and
grow geometry from it. The route taken here is the same in spirit but inverted in
order: rather than positing spatial dimensions and forbidding superluminal
signalling, we \emph{define} spatial separation as a measure of non-computability
and read the dimension, the signature, and the Lorentz group off the result. The
proposal originates in the author's CASYS'05 paper \autocite{small_2006}, which
first suggested treating spacetime as a space of computations and attempted to
justify three dimensions on this basis; the present account tightens and extends
it. This is the framework's most structural and least formally complete arc; the
three-dimensionality result is a structural conclusion awaiting a formal proof,
and several of the links below are conjectural readings. Each is flagged as such,
so the reader sees exactly where the chain is load-bearing and where it is an
invitation.

\subsection{Timelike and spacelike distance as computation}
\label{sec:comp-distance}

A computation carries a state from axioms to a theorem in a number of steps.
Counting the computable steps between two states gives a real measure---time. A
physical clock counts completed irreversible loops, so proper time literally
counts loop completions and the metric $\mathrm{d}\tau^{2}=g_{\mu\nu}\,
\mathrm{d}x^{\mu}\mathrm{d}x^{\nu}$ is a loop-density tensor, the conversion from
coordinate intervals to loop counts; Toffoli's reading of the Lagrangian as a
count of action quanta is the same statement \autocite{toffoli_2003}.

Spacelike separation is the complementary, non-computable measure. Define the
non-computable distance between two propositions, relative to an axiom system, as
the number of bits that must be added to the axioms to make the second derivable.
A closed timelike curve is a computation whose output is its own input; by the
fixed-point theorem it has a unique consistent solution or none, and the
existence of that solution is one bit not derivable from any axiom system that
does not already contain the loop---one bit per curve, in the manner of Chaitin's
$\Omega$ \autocite{Chaitin1998}. The spacelike distance between two events is then
the number of such curves that would have to be closed to make them causally
connected. Because each added bit is a non-computable proposition, it is carried
by the imaginary part of the complex information measure
$S=-\ln\lvert p\rvert-\mathrm{i}\theta$ derived in
Section~\ref{sec:complex-amplitudes}: computable distance is counted in real
units, non-computable distance in imaginary ones. The Minkowski signature, with
spacelike intervals imaginary relative to timelike ones, is the immediate reading. The convention is not new to this paper: it was proposed in the author's CASYS paper two
decades ago, where counting non-causal steps in imaginary numbers was introduced and its formal
proof flagged, already then, as a key open problem \autocite{small_2006} ---
Open Problem~\ref{op:comp-metric} is that flag's direct descendant, stated now with the machinery
to be answerable.
That the resulting structure is a genuine metric space, and that the signature
follows rather than being fitted, are open problems---at present the
identification is motivated by reproducing what nature does, not proved.

%% VERIFY-SUPPORT 2026-09-27: conway_2003 is cited below for "any two octonions
%% generate an associative subalgebra" (Artin's theorem; a Moufang loop is
%% diassociative) and harvey1982calibrated for "a three-plane in Im O is
%% associative iff the associator of a basis vanishes, and the associative
%% three-planes are the imaginary parts of the quaternionic subalgebras"
%% (Harvey--Lawson, ch. IV). Neither passage was read in the session that wrote
%% this; both statements are checked numerically in Convention A by
%% scripts/antitelephone_network_associativity_2026-09-27T1900.py (PASS).
A reader used to the standard order of ideas should pause here, because the order has been
inverted, and the inversion is used below. In special relativity the tachyonic antitelephone is a
corollary: a signal faster than light, composed with a boost, would arrive before it was sent and
close a timelike loop, therefore there is no such signal \autocite{tolman1917theory,
benfordbooknewcomb1970tachyonic}. The metric comes first and the impossible telephone is a check
on it. Here the telephone comes first. Two observers are spacelike separated when an antitelephone
between them would close a loop; the loops such telephones would close, and carry nothing round,
are what spacelike separation is, their number is the distance, and the metric is read off from
them. Nothing is lost, since the corollary returns as a tautology: an antitelephone that carried
a record would close a causal loop, and its two ends would not have been spacelike after all.

What is gained is a definite meaning for a spatial dimension. Picture many observers, each holding
an antitelephone to each of the others. A message relayed from one observer to another through a
third composes two transports round two loops, and the network is consistent only if every relay
composes to a single transport whatever the bracketing of the relaying: a spatial dimension is a
family of loops whose transports compose associatively, which is a group law on a parallelisable
sphere. That is a condition a family of loops can fail, and \S\ref{sec:three-dimensions} shows
that the loops of the seventh sphere fail it. Two antitelephones always compose, because any two
octonions generate an associative subalgebra \autocite{conway_2003}; three compose only when their
imaginary parts lie in one quaternionic subalgebra, which is what the associative three-planes of
Harvey and Lawson are \autocite{harvey1982calibrated}, and a third link outside that subalgebra is
seen by the network exactly through its component outside it. The largest consistent network of
antitelephones inside $S^{7}$ is therefore an $S^{3}$, and there are three spatial directions.
When the seventh sphere is set aside there as a ruler, which is the octonionic spacetime of the
division-algebra literature, the claim is not that $\mathbb{R}^{8}$ has no metric, nor that the
ten-dimensional Minkowski space of $SL(2,\mathbb{O}) \cong \mathrm{Spin}(9,1)$ fails to exist
\autocite{baez2002octonions}; it is that the octonionic directions are not loops on which
antitelephones compose, and so are algebra rather than distance. In the standard order that
sentence has no content; in this one it is the reason space has three dimensions.

The split just defined explains the most familiar asymmetry in physics: that one can move back and
forth in space but not in time. \S\ref{sec:continuity} identifies the pre-measurement regime with
continuous reversibility --- its transformation group is connected, and every motion in it is
undoable. Spacelike separation, being the non-computable sector, is exactly the part of the
relational structure on which no record has been written; motion across it is motion within the
reversible regime, and it inherits that regime's continuity and its two-sidedness: a place can be
revisited because going and returning writes no record that forbids the return. Timelike separation
is the record chain itself --- proper time counts completed irreversible loops, as above --- and a
moment cannot be revisited because a moment \emph{is} a stratum of records, and returning would
require unwriting them, which is what \textsc{fanout} forbids (\S\ref{sec:fanout-boundary}). The
asymmetry between space and time is therefore not an extra postulate about the manifold: it is the
pre-/post-\textsc{fanout} asymmetry read as geometry --- space is what the reversible group feels
like from inside, time is what the record chain feels like --- and the arrow-of-time discussion
later in this section is the same asymmetry read thermodynamically. The identification is the
framework's structural reading of the split, offered in the same register as the rest of this
subsection.

\subsection{Three dimensions of space}
\label{sec:three-dimensions}

The closed timelike curves that define spacelike distance require information to
be parallel transported consistently around a loop, and consistent parallel
transport on a sphere requires it to be parallelisable. By Adams' theorem the
only parallelisable spheres are $S^{1}$, $S^{3}$ and $S^{7}$
\autocite{adams1960hopf}, the unit spheres of $\mathbb{C}$, $\mathbb{H}$ and
$\mathbb{O}$. Impose now the requirement that information arriving from outside
the causal zone must appear to possess its own causal history---that is, to be the
result of a sequence of information-preserving operations, hence a sequence of
group operations. The spheres $S^{1}$ and $S^{3}$ are groups under their own multiplication --- $U(1)$ and
$SU(2)$ --- and their operations preserve information; $S^{7}$'s multiplication is non-associative,
a Moufang loop rather than a group, so $S^{7}$ cannot serve as a ruler. (The rotation group
$SO(8)$ of the seven-sphere is of course a group; what fails is the sphere's \emph{own}
composition law, the one a ruler would use.) The same non-associativity
forbids a metric built from the sphere's own composition law outright: path
composition $d(A,C)$ through $B$ depends on the bracketing, so the triangle inequality is not merely violated but ill-defined, and
no geodesic structure exists. (What is untouched should be said: the Euclidean metric of
$\mathbb{R}^{8}$, and the ten-dimensional spacetime that $SL(2,\mathbb{O})$ furnishes, are metrics
of the ordinary kind; what fails is a ruler made of antitelephone loops, which is what a spatial
dimension is in this section, as \S\ref{sec:comp-distance} set out. Nor is the failure total:
relays between any two links compose, and among three they compose exactly when the links lie in
one quaternionic subalgebra, so the largest consistent network inside $S^{7}$ is an $S^{3}$
\autocite{conway_2003,harvey1982calibrated}.) Since the abelian $S^{1}$ alone cannot carry
non-abelian structure, spacelike distance must be three-dimensional. Dimension
three is an output, not an input \autocite{small_2006}.

The division of labour is then exact. $S^{1}$ ($\mathbb{C}$) contributes the time
dimension and the $U(1)$ of electromagnetism; $S^{3}$ ($\mathbb{H}$) contributes
the three spatial dimensions and the $SU(2)$ of the weak interaction; $S^{7}$
($\mathbb{O}$) contributes no geometry but, as a filter on what information may be
exchanged, the $SU(3)$ of the strong interaction. A sphere that cannot be a ruler
can still be a filter, which is why the strong force resists geometric
interpretation while electromagnetism and the weak force do not. This is the same
$\mathbb{C}\oplus\mathbb{C}^{3}$ split that yields $G_{\mathrm{SM}}$ in
Section~\ref{sec:entanglement-classes}, now read as the reason spacetime is $(3+1)$-%
dimensional. The conclusion converges with M\"uller and Masanes, who prove from
information-theoretic postulates, assuming no quantum theory, that minimal
direction-carrying systems with continuous reversible dynamics force $d=3$ and
quantum theory on two qubits \autocite{MuellerMasanes2013}, and with the result
that entanglement exists only for a three-dimensional Bloch ball
\autocite{MasanesMuellerEtAl2014}; the framework supplies the reason such systems
exist, Cayley--Dickson doubling from self-reference, while their work supplies
the rigorous dimensional uniqueness. The status of the framework's own route is a
structural result requiring formal proof: the one-bit-per-curve identification and
the precise reduction of the $S^{7}$ filter to $SU(3)$ rather than its full automorphism group
$G_{2}$ are numbered below (Open Problems~\ref{op:comp-metric} and~\ref{op:three-dimensions}).

A scoping remark, and then a second route to the same count. The computational model of spacetime
proposed here has many detailed parallels with the real thing, but a parallel is not an identity,
and the identity is not claimed in this paper: establishing it is work deferred to the second paper.
What can be claimed now is that the parallel is productive in a specific direction --- it suggests a
derivation of the dimension count carried out entirely within general relativity, with no
computational premise. Acceleration and gravitation tip the light cone, and the tipping can be
carried to the extreme: Penrose's discussion runs the cone through a full circle
\autocite[382]{Penrose1994}, at which point the tipped direction closes on itself --- which is to
say, the extreme of tipping is a closed timelike curve, the very object whose closure the
computational definition of \S\ref{sec:comp-distance} counts in bits. The two routes are one route
entered from opposite ends. Now demand of general relativity only this: that observers accelerating
with respect to one another, and with respect to unaccelerated observers, retain a consistent
conception of the distance between them. Each observer's tipping traces a path on a sphere of
directions, and mutual consistency requires the associated rotations to compose without regard to
bracketing --- to form a group. Compositions on $S^{1}$ and $S^{3}$ do, forming $U(1)$ and $SU(2)$ as above; compositions on
$S^{7}$ do not, the multiplication being non-associative, so observers tipped across an
$S^{7}$-worth of directions could not agree on distance even pairwise: the same failure that forbade
$S^{7}$ a metric above forbids it here a shared ruler. The largest sphere admitting a consistent
conception of distance under mutual acceleration is therefore $S^{3}$, and space has at most three
dimensions; with the floor argued above --- the abelian $S^{1}$ alone cannot carry the required
structure --- the count is exactly three. The status matches the subsection's. The tipping and its
extreme are general relativity's own \autocite{Penrose1994}; the consistency demand is
elementary; the associativity filter is the algebraic fact used throughout; the reading of composed
tippings as motion on the parallelisable spheres is the structural step, shared with the
computational route. The convergence of the two routes on the same endpoint is offered as evidence
for the parallel, while the identity it gestures at remains the second paper's burden.

\subsection{Recovering the Lorentz group}
\label{sec:lorentz}

A measurement changes what the observer knows, hence, relationally, the causal
structure: it tips the light cone. This is Penrose's association of cone-tipping
with state reduction \autocite{Penrose1996} run in the opposite causal direction,
reduction driving the tipping rather than the reverse, and the two are then one
operation in two charts---both described by $SL(2,\mathbb{C})$. The algebraic
backbone is the division-algebra correspondence
$SL(2,\mathbb{A})\cong\mathrm{Spin}(1,d{+}1)$ \autocite{baez2002octonions}: $\mathbb{R}$
gives $1{+}1$, $\mathbb{C}$ gives $3{+}1$, $\mathbb{H}$ gives $5{+}1$ and
$\mathbb{O}$ gives $9{+}1$ dimensions. The computability split that singles out a
preferred $\mathbb{C}$ fixes the fundamental state space as $\mathbb{C}^{2}$ and
hence $SL(2,\mathbb{C})=\mathrm{Spin}(1,3)$: the Lorentz group of $3{+}1$
spacetime. The qubit and Minkowski space are not merely both Euclidean-and-three%
-dimensional by coincidence; the Hermitian $2\times2$ matrices are
$\mathbb{R}^{1,3}$, $SL(2,\mathbb{C})$ acts on them as the Lorentz group, and the
qubit density operators fill the future cone---the correspondence M\"uller's group
develops operationally, recovering Lorentz transformations from quantum
communication \autocite{HohnMueller2016} and the complex and quaternionic qubit
from the relativity of simultaneity \autocite{GarnerMuellerDahlsten2017}.
In these coordinates the Bloch ball is the unit-trace slice of the cone and the Bloch sphere its
boundary rays, the celestial sphere on which $SL(2,\mathbb{C})$ acts by M\"obius
transformations; state reduction sends a point of the ball to a point of the sphere, and a boost
sends the sphere to itself, which is what lets the two be one operation in two charts.
The $SL(2,\mathbb{A})$ correspondence is rigorous; its identification with
measurement-as-tipping is the framework's structural reading.

One consequence of the identification has been implied throughout this section and deserves its own
paragraph, because the next paragraph turns on it. If measurement creates a relation between
observer and observed, and if --- relationally --- a new relation alters the causal structure, then
the tipping of the cone is not merely \emph{associated} with the creation of the relation: it is
what the creation of the relation is, seen geometrically. And a tipping of one's cone against the
congruence of one's neighbours already has a name in physics: acceleration. The framework's
statement, made explicit here, is that \emph{creating a new relation between observer and observed
is experienced as acceleration}. The resonance with the equivalence principle is exactly what it
should be: the principle makes acceleration and gravitation locally indistinguishable, Penrose's
association --- run in this section's direction, reduction driving the tipping --- makes state
reduction and cone-tipping one operation, and the framework's identification closes the triangle,
measurement and acceleration meeting where the principle says they must. The claim is structural:
the framework's reading of the cone-tipping established above, not an independent derivation.

Taken literally, the identification has a measurable face, and it is a familiar one: the Unruh
effect, read in this section's own ledger. Recall the definitions of \S\ref{sec:comp-distance}:
timelike distance counts computable steps, spacelike distance counts non-computable bits. What an
observer can compute --- what their \textsc{fanout}s can convert into classical records --- is fixed
by their causal structure, and the previous paragraph has just made that structure
relation-dependent. An accelerated observer's cones are tipped, and with them the boundary of the
recordable: relations that for the inertial observer lie wholly in the reversible, pre-record sector
--- the entanglement structure of the vacuum, non-computable for them, spacelike in the
computational sense --- fall, for the accelerated observer, partly on the recordable side. The mechanism is the one
Chaitin identified for formal systems \autocite{Chaitin1998}: a proposition undecidable relative to
an axiom system becomes decidable when information is added to the axioms --- and
\S\ref{sec:comp-distance} defined spacelike distance as exactly the bits whose addition closes the
gap. Creating a relation is adding information to the observer's axiom base, so timelike and
spacelike distance are converted into one another by the addition and subtraction of information,
and acceleration is the continuous version of that addition. A
relation newly recordable is a detection. Particles appear from nowhere: an accelerated detector in
the inertial vacuum registers a thermal bath at $T = \hbar a / 2\pi c k_{\mathrm{B}}$
\autocite{unruh1976notes}. Two guards keep the reading honest. Nothing here reclassifies the
invariant interval --- timelike and spacelike separation in the emergent metric are Lorentz
invariants, and no event crosses that classification; what is observer-relative is the computational
partition, which relations a given observer's horizon renders recordable. In the ledger of
\S\ref{sec:comp-distance} that partition \emph{is} the timelike/spacelike distinction, so in the
framework's own units timelike distance is \emph{rotated} into spacelike, and spacelike into
timelike, by a change of observer --- the tipping of the cone is literally that rotation --- and the
detector's clicks are the rotation made audible. And the effect is established physics
while the reading is the framework's, flagged structural. What gives the reading its content is a
parallel: the inertial and the accelerated observer stand to the vacuum exactly as Wigner and his
friend stand to the laboratory (\S\ref{sec:wigners-friend}) --- one holds a coherent relation and no
event, the other holds a record, and the facts are relative on both sides in precisely Rovelli's
sense \autocite{DiBiagioRovelli2021}. Since the previous paragraph identifies acceleration with
relation-creation, the parallel is not an analogy but the same structure in gravitational dress, the
horizon standing where the laboratory wall stood: \emph{the Unruh effect is Wigner's friend for
accelerated observers}. And two subsections of this section that seemed independent here meet: the
temperature Jacobson's construction requires across every local Rindler horizon
(\S\ref{sec:jacobson}) is exactly this temperature --- the $T$ of the Clausius relation is the $T$
of the tipped cone, supplied by the very relation-creation with which this subsection began.

\subsection{The relationalist commitment}
\label{sec:relationalism}

General relativity is relational: by the hole argument, spacetime points carry no
intrinsic identity, and the metric is fixed by the network of coincidences and
interactions, not by a pre-existing stage \autocite{Rovelli1996a}. The
computational picture is relational by construction and explains why it must be:
a background metric would require spatial distances to be pre-computable, and
non-computability forbids exactly that. General covariance is the statement that
distances have no pre-assigned values and depend on what information is actually
present. The metric is therefore dynamical not by stipulation but because it
measures a quantity that cannot be fixed in advance.

\subsection{Jacobson, holography, and imaginary information}
\label{sec:jacobson}

Jacobson derives the Einstein equation as an equation of state, from the Clausius
relation $\delta Q=T\,\mathrm{d}S$ across every local Rindler horizon together
with the Bekenstein proportionality of entropy to horizon area
\autocite{Jacobson1995, Bekenstein1973}. In the computational reading this is the
self-consistency condition on the loop-counting network: the Einstein equation is
the requirement that the bookkeeping of computation be internally consistent
across surfaces. The holographic principle, that the information in a region is
bounded by its bounding area \autocite{tHooft1993, Susskind1995}, then acquires a
natural reading here. A surface area is a squared distance; if spacelike distance
is measured in imaginary units of information, an area is measured in the square of
those units, that is in \emph{negative} information---entropy. The holographic
link between area and entropy thus requires the square root of negative
information, and that square root is precisely the imaginary unit of the complex
amplitude, the observer's self-referential phase. Read this way, Jacobson's
account of information crossing surfaces implicitly assigns an imaginary
information value to spacelike distance, and the $\mathrm{i}$ of the Minkowski
signature is the same $\mathrm{i}$ as the $\mathrm{i}$ of the quantum amplitude.
This identification is the most speculative step in the section: it is
structurally suggestive and it makes the holographic form unsurprising, but it is
a conjectural reading rather than a derivation, and is offered as such.

The conjecture has a two-way motivation that this paper's record should carry, because each
direction was on the table before the other was noticed. Forward: if spacelike distance is counted
in imaginary units of non-computable information, then any spacelike area is the product of two
imaginary lengths and is therefore \emph{negative} information --- which is to say entropy, by
Brillouin's identification \autocite{brillouin1956science} --- and a holographic principle arrives
automatically, the sign supplied by $i^{2}$ rather than by interpretation. Converse, and more
striking: in the holographic principle as ordinary physics already states it, surface area is
proportional to entropy, and entropy is negative information; so by implication ordinary physics
has been measuring spacelike separation in imaginary units of information all along, and the
convention of \S\ref{sec:comp-distance} is read off from established results rather than
invented. The same imaginary unit is already employed once in this section --- it is what makes
the signature of the interval Minkowski --- so the sign of the area law is its second job, not a
second assumption. What the two-way motivates it does not prove: the formal statement is exactly
the metric-and-signature half of Open Problem~\ref{op:comp-metric}, together with the
one-bit-per-bit-pair coefficient owed to the second paper. The stakes of the conjecture have
recently risen: Dorau and Much prove, with modular-theoretic rigour, that the semiclassical
Einstein equations follow from quantum relative entropy on a local horizon \emph{once the
entropy--area proportionality is assumed} \autocite{doraumuch2026relative} --- so what this
framework proposes to derive is precisely the remaining input of a derivation already carried
through at that standard.
\input{entropic_gravity_family_paragraph_2026-07-03}

\subsection{Non-locality already present in general relativity}
\label{sec:gr-nonlocality}
\input{gr_nonlocality_revision_2026-06-13}
\medskip
\noindent  The spacetime arc therefore yields, from one definition of spacelike
distance as non-computability: the signature, three spatial dimensions,
the Lorentz group, the dynamical relational metric, a reading of the
holographic principle, the presence of non-locality within general
relativity, the identification of relation-creation with acceleration
and of the Unruh effect with Wigner's friend, and the space--time
mobility asymmetry as the \textsc{fanout} asymmetry read as geometry. Of these,
the $(3+1)$ count and the Lorentz recovery are structural results that converge with independent
rigorous work and await their formal proof --- the count now numbered as
Open Problem~\ref{op:three-dimensions}; the metric-space status, the signature and the
one-bit-per-curve measure (together Open Problem~\ref{op:comp-metric}), the imaginary-information
reading of holography, and the general-relativistic non-locality, the acceleration and
Unruh readings, and the mobility asymmetry are conjectural, and are named here precisely so that the boundary between what is shown and what
is proposed is visible.

\subsection{Why the metric is not to be quantised}\label{sec:no-quantised-gravity}

A consequence has been implicit in everything above and deserves its plain statement. Causal
structure, on this framework's reading, is the bookkeeping of locked relations --- post-\textsc{fanout}
by construction --- while superposition is the signature of the pre-\textsc{fanout} regime. A
``superposition of causal structures'' is therefore not a difficult state but a category error: it
applies the pre-record operation to a record-side object, exactly as ``a superposition of
measurement outcomes'' does. Gravity lives in the post-\textsc{fanout} regime. The quantum layer
beneath it exists and is not in dispute; but it is \emph{pre-geometric} --- it carries amplitudes
over the formation of relations, not amplitudes over formed geometries --- so there is no
fundamental metric to quantise, and the standard programme of superposing spacetimes pursues an
object the ontology does not contain. One guard rail keeps the claim honest: collective
excitations of the locking pattern may perfectly well exist as effective quanta --- gravitons as
the phonons of the relational network --- and the thesis forbids fundamental metric superposition,
not effective field theory; this is Jacobson's equation-of-state moral \autocite{Jacobson1995},
made structural. The relation to gravitational-collapse programmes is an inversion: where
superposed geometries are there read as unstable \autocite{Penrose1996}, here they are read as
unformed --- what decays in those models is what locks in this one.

Three experimental postures follow, each at the status the reading carries. Laboratory
demonstrations of indefinite causal order \autocite{procopio2015experimental, rubino2017indefinite} are, on this
account, demonstrations of the nascent layer: the indefiniteness lives below the locking threshold
and definitises at it, so the framework adopts the programme and adds a boundary --- indefinite
causal order can never be amplified into records while remaining indefinite. Gravitationally
induced entanglement of the kind the Bose--Marletto--Vedral proposals seek
\autocite{bose2017spin, marlettovedral2017gravitationally} is \emph{predicted}, because the
pre-geometric layer is quantum --- while witnessing that layer, not a superposed metric: the
framework splits the experiment's usual dichotomy. And gravitationally driven anomalous collapse
is denied: locking is redundancy and thermodynamics, so decoherence experiments should find
standard rates and nothing gravity-specific. That no consistent quantisation of gravity exists was
the conclusion of the author's CASYS paper by a different route \autocite{small_2006}; the present
framework supplies the mechanism the conclusion then lacked.

\subsection{Open problems and prospects}
\label{sec:spacetime-open-problems}

The spacetime arc is the framework's frontier, and honesty about its edges
matters more here than anywhere else in the paper. We set out the genuine
problems and the prospects together, deferring the deepest to a second paper
rather than papering over them. Two of the arc's debts are well-posed now, and we number them.

\begin{openproblem}[The computational metric]\label{op:comp-metric}
Prove that the paired distance of \S\ref{sec:comp-distance} --- computable steps counted real,
non-computable bits counted imaginary, one bit per closed curve --- defines a genuine
pseudo-metric on the space of computations, with the Minkowski signature derived rather than read
off; equivalently, establish the one-bit-per-curve measure as a measure, and the signature as its
theorem. Until then the identification stands as motivated by reproducing what nature does, which
is the status flagged at its introduction.
\end{openproblem}

\begin{openproblem}[Dimension three as a theorem]\label{op:three-dimensions}
Turn the two convergent routes of \S\ref{sec:three-dimensions} into a proof: formalise the
parallelisability-plus-associativity filter (the computational route) and the mutual-acceleration
consistency demand (the general-relativistic route) so that $d = 3$ follows as a theorem, and
include the reduction clause --- that the $S^{7}$ filter acts through $SU(3)$, the stabiliser of
the selected complex direction, rather than through its full automorphism group $G_{2}$. The
M\"uller--Masanes uniqueness results \autocite{MuellerMasanes2013, MasanesMuellerEtAl2014} are
the rigorous convergents any proof must meet.
\end{openproblem}

\paragraph{The substrate and the generation of events (deferred).}
A computation runs on something; a network of events presupposes that events
occur. Both are the same objection in two forms---the network's hardware and its
clock---and the original CASYS'05 statement of the programme raised the
implementation regress and set it aside \autocite{small_2006}. The framework's
response is that a self-referential structure is the only kind that need not
presuppose a meta-substrate: it generates its own. The natural algebraic emblem
is $E_{8}$, the unique simple Lie group whose adjoint representation is also its
smallest faithful one, so that what it is and what it acts on coincide; the $240$
unit integral octonions that generate it are exactly the $E_{8}$ root system, so
the self-action is welded to the division-algebra spine of
Section~\ref{sec:entanglement-classes}. This is Wheeler's self-excited circuit
\autocite{Wheeler1990} given an algebraic body, and it converges with Furey's use
of an algebra acting on itself to generate particle content
\autocite{furey2018three}. The decisive caveat, which we state rather than hide, is
that the bootstrap is self-consistent but not self-specifying: it does not
eliminate the non-computability identified throughout, it relocates it to a
determined place. The substrate question is, on the framework's own terms,
provably without a computable answer from inside---it is the irreducible residue,
the bump in the carpet that the whole construction has been moving rather than
removing. The leap from ``$E_{8}$ is self-dual'' to ``$E_{8}$ is the substrate''
is at present an emblem awaiting operational content (an $E_{8}$-equivariant
dynamics, or a state space built as a homogeneous space of $E_{8}$), and that
content, together with how a tick is generated at all, is deferred to the second
paper.

\paragraph{Dark matter (deferred).}
If gravity is the bookkeeping of the interaction network and dark matter interacts only
gravitationally, how does it participate in a network it does not exchange information with?
This is a genuine boundary of the programme, at which it may make no prediction, and Open
Problem~\ref{op:sterile-gravity} states it with the criteria that would close it. The prospect, to be developed in the second paper, is
that this is a feature rather than a defect. In the division of labour of
Section~\ref{sec:three-dimensions}, participation in the metric---counting loops
in the $S^{1}$ and $S^{3}$ directions---is universal, carried by anything with a
worldline and a tick-rate, whereas participation in the gauge filters, the
$S^{3}$ of the weak interaction and the $S^{7}$ of colour, is selective, engaged
only by charged matter. A node that counts loops but engages no filter
\emph{is} matter that gravitates without otherwise interacting, and the
framework's Class-4 a-chiral sterile neutrino (\S\ref{sec:generations}) is the candidate that
realises it --- recorded in the particle sections as Open Problem~\ref{op:sterile-gravity}, whose
tension (a gauge-sterile channel gravitating, when gravity is measurement bookkeeping) is exactly
the question named here. The condition for the candidacy to be coherent rather than circular is
that metric participation be established as prior to, and independent of, gauge participation;
making that priority precise is the work deferred, and it is the same work the open problem
requests.

\paragraph{CP violation, the arrow of time, and baryogenesis.}
CP violation appears to demand a time asymmetry owing nothing to the entropic
arrow, which would seem to require a preferred direction built into spacetime
itself. The framework dissolves the appearance rather than accommodating it. CP
violation needs an irreducible complex phase, the CKM phase $\delta$; by Spekkens'
result that negativity and contextuality are the same notion of nonclassicality
\autocite{spekkens2007toy}, a complex phase is a signature of contextuality; so CP
violation should be contextual, sensitive to the near environment of the decaying system --- the
same phase the interaction sections read as anchored by a fourth party
(\S\ref{sec:interactions}) --- which is exactly what Andrianov, Taron and Tarrach established for
neutral kaons in matter, where the environment breaks C, irreversibility breaks
T, and the effective parameters follow a Lindblad equation
\autocite{Andrianov2001}. The framework's reading is that the resulting
informational asymmetry \emph{constitutes} the arrow rather than presupposing it:
CP violation, entropy increase, and the direction of time are one informational
asymmetry seen three ways, not an independent second arrow, so the worry about a
direction built into spacetime does not arise. On this footing the contextual
dependence sets up a positive feedback loop for baryogenesis---an initial
fluctuation makes the environment slightly matter-dominated, contextual
enhancement raises the effective CP violation, more matter is produced, and the
loop runs---with $\delta$ a cosmological fixed point rather than a fundamental
constant. The honest limit of this is that a feedback loop explains amplification
and lock-in, not the sign selection or the baryon-number violation that
Sakharov's conditions still require \autocite{Sakharov1967}; the stronger claim
the framework is positioned to make, that the matter direction and the time
direction are co-selected by the same self-referential asymmetry, needs a source
of baryon-number violation and is left open.

\paragraph{Spatial infinitude and the status of inflation.}
There is no upper bound on non-computability, so the space of computations is
unbounded: the universe is spatially infinite, which leaves its spatial sections flat or of
negative curvature and excludes only the closed case. Infinitude on its own, however, does not
settle flatness, which is a fine-tuning of the Friedmann dynamics the framework has not yet
derived (Open Problem~\ref{op:comp-metric}); and flatness was in any case never inflation's
principal task. The horizon problem is, and an infinite space does nothing for it. The framework's actual lever against
inflation is the relational thesis of Section~\ref{sec:gr-nonlocality}: if
pre-spatial entanglement correlates regions before \textsc{fanout} spreads them
into spatial separation, then the observed large-scale homogeneity is fixed
relationally, before any question of causal contact, and the horizon problem
dissolves without an inflaton. This is a real argument, but it carries a real
bill: inflation also generates the nearly scale-invariant spectrum of primordial
fluctuations seen in the microwave background \autocite{Planck2020}, and any
account that dispenses with inflation owes an alternative origin for that
spectrum. That, not the horizon problem, is the hard empirical debt, and it is
named here rather than discharged.

\paragraph{Expansion and acceleration.}
If spacelike distance is the information separating worldlines, then as worldlines
accumulate distinct interaction histories they accumulate information-distance,
and the universe expands. The picture has the right exclusions built in: bound
systems interact continuously, share information, and so keep low
information-distance---they do not expand---while only worldlines that decouple
diverge, recovering expansion on cosmological scales alone. The feedback in which
greater separation means less light-cone overlap and so faster divergence,
$\mathrm{d}d/\mathrm{d}t \propto d$, yields exponential divergence, a de~Sitter
phase, and hence accelerating expansion without a separate dark-energy term
\autocite{Riess1998}. The limit to be honest about is that pure exponential
divergence is the late-time behaviour only; reproducing the full expansion history,
with its radiation- and matter-dominated eras, and fixing the Hubble rate
quantitatively, remains to be done.

\paragraph{The $S^{7}$ sector as internal structure.}
That $S^{7}$ cannot carry a metric does not make it disappear. Its closed
timelike curves carry the same non-computable information structure as those of
$S^{1}$ and $S^{3}$, but lacking associativity they support algebra---gauge
charge---in place of geometry---distance. This is why gauge symmetry has so often
looked like a set of compactified extra dimensions: the Kaluza--Klein and string
intuition is tracking something real in the $S^{7}$ structure, but misdiagnoses
it as spatial. In the framework these degrees of freedom have no size because
they have no metric; they are internal in the exact sense that they enter the
information-exchange that constitutes gauge charge while contributing nothing to
the loop-counting that constitutes distance. The resemblance to rolled-up
dimensions is therefore genuine and instructive, and the diagnosis that they are
algebraic rather than geometric is what separates this account from theories that
add real spatial dimensions to carry the same structure.

\paragraph{Loops, supertasks, and a second route to parallelisability.}
The argument of \S\ref{sec:three-dimensions} derives the parallelisable-sphere constraint from
closed timelike curves: a computation whose output is its own input requires information to be
transported consistently around a loop, and consistent transport selects the parallelisable spheres
$S^{1}$, $S^{3}$, $S^{7}$ \autocite{adams1960hopf}. There is a suggestion, which we record as a
prospect rather than a result, that the same constraint may govern the loops of perturbative quantum
field theory. The external legs of a Feynman diagram are on-shell classical boundary data; the
quantum content, and the divergences, enter with the internal loops, where the loop momentum is
integrated over all values and the action accumulated inside the loop is in this sense unbounded ---
a supertask, an infinite computation closed into a finite process. If such a loop is, in the
framework's terms, a closed computational path of the same kind as a closed timelike curve --- a
cycle carrying information not derivable from the data outside it --- then the parallelisability
constraint derived for closed timelike curves would apply to it, and the admissible interactions
would be constrained by the same $S^{1}$, $S^{3}$, $S^{7}$ structure that fixes the gauge group in
\S\ref{sec:entanglement-classes}. This would be a second and independent route to the Standard Model
content: not only \enquote{closed paths in spacetime select the parallelisable spheres} but \enquote{closed paths in interactions select them too}, unifying the spacetime arc and the interaction arc under one
principle. The decisive step, which we have not taken, is the identification of a field-theory loop
with a closed computational path in the precise sense the framework requires --- a loop integral is a
closed cycle in momentum space, a closed timelike curve a closed cycle in spacetime, and whether
these are the same kind of self-consistent closed structure is exactly what would have to be
established. It is offered as a direction, not a claim.

\paragraph{Does parallelisability determine renormalisability?}
A sharper and more self-contained question sits alongside the previous one, and we state it in the
hope that it will be taken up by those equipped for it. Renormalisability is the property that the
divergences of a quantum field theory can be consistently absorbed, leaving finite predictions; some
theories have it and some do not, and the standard criterion for which is technical (power-counting
of the coupling's dimension). The framework suggests a more structural reading. Parallelisable
spheres are precisely the spheres on which a global, consistent transport exists where it generically
does not; the non-parallelisable spheres are where the obstruction lives. If the divergence of a loop
is, in the framework's language, an obstruction to closing a computational path consistently, then
parallelisability would be exactly the structure that distinguishes the removable obstruction (a
renormalisable theory) from the irremovable one. It is at least a suggestive coincidence that the
renormalisable interactions of nature are the gauge theories built on $U(1)$, $SU(2)$ and $SU(3)$ ---
the groups the framework ties to the three parallelisable spheres --- and that gravity, which resists
renormalisation, is the sector \S\ref{sec:three-dimensions} ties to the geometry rather than the
algebra. Whether this is structural or accidental we do not know. The question --- whether
parallelisability, in the Hurwitz--Adams sense, is what determines renormalisability --- is posed in
the native language of the division-algebra approach to the Standard Model, and is one we think that
community is best placed to answer. We are careful to claim nothing more than that the question is
well-motivated: this is emphatically not a method for computing renormalised quantities, whose finite
remainders carry the precision agreement with experiment, but a conjecture about which theories admit
such a computation at all. It is, like the previous prospect, a direction with a possible path to an
answer rather than an answer, and it is left open deliberately.


\medskip
\noindent These mark the boundary of the programme as it stands. The substrate,
the generation of events, and the network status of dark matter are deferred to
the second paper; the contextual account of CP violation and baryogenesis, the
relational case against inflation, the information-divergence account of
expansion, and the algebraic reading of the $S^{7}$ sector are offered as
prospects, each with its outstanding debt named so that the line between argument
and aspiration stays visible.
%      % EDITED 2026-09-09 scripts/patch_spekkens_citekey_2026-09-09T1810.py
% =====================================================================
%  SECTION: Why Hopf Fibrations Belong in Particle Physics
%  (sec:hopf-justification). Drafted 2026-06-08; F1 arc added
%  prop:f1-dictionary and op:hopf-frobenius's discharge clause.
%  AUDITED 2026-07-06 (final pass, eighth stop).
%
%  PLACEMENT (reordered chain, JS 2026-07-06): position 8 --
%      \input{spacetime_from_non-computability}
%      \input{hopf_justification_section_2026-06-08}
%      \input{entanglement_classes_and_gauge_groups_2026-06-08}
%  No adjacency phrasing in the prose: position-portable. All
%  cross-refs live -- sec:fanout-direction DRAFTED 2026-07-06 (formerly the one undrafted
%  section, by design).
%
%  AUDIT RECORD: prop:gsm, prop:termination and prop:f1-dictionary
%  verified against every quoting consumer (statements canonical);
%  honest-debt subsection names both the selection assumption and the
%  mereology debt; the epistemic two-climbs disambiguation and the
%  compositional rung-evidence verified; sec:canonical-orbits live;
%  file pure ASCII. Fixes: op:hopf-frobenius's cross-link with
%  op:sigma-functor made TWO-WAY (the "four ways" line now names the
%  op); both remaining functor conditions ANCHORED --- the quaternionic
%  to the established fibre physics of sec:wf-reconciliation (residue:
%  clutching map vs strong complementarity), the octonionic to the
%  proved canonical-orbit images (umbrella label app:ghz-w-proofs
%  minted at the appendix for the purpose). Statuses unchanged.
%  British English; \autocite; no bullets in prose.
% =====================================================================

\section{Why Hopf Fibrations Belong in Particle Physics}
\label{sec:hopf-justification}

Invoking Hopf fibrations in a discussion of the Standard Model invites an immediate and fair
objection: that the move is decorative, an analogy between two pieces of attractive mathematics
with no load-bearing connection between them. The purpose of this section is to dispose of that
objection before the construction is put to work. The connection is not an analogy. It is
carried by two independent and independently published theorems, and the structures it relies
upon are not selected for convenience but compelled by a pair of classification results. We take
as already established the correspondence between qubit entanglement and the Hopf fibrations
themselves --- the subject of a mature literature
\autocite{mosseri2001hopf,bernevigchen2003geometry} --- and do not reproduce it here. What this
section addresses is the prior question of why that correspondence is entitled to carry physical
content at all.

\subsection{Two theorems, joined end to end}

The bridge from qubit entanglement to the Standard Model is built from two distinct results,
neither of which is original to the present framework and neither of which depends on the other.

The first is geometric. The unit sphere in the Hilbert space of $n$ qubits is the total space of
a Hopf fibration: $S^3 \to S^2$ for one qubit, $S^7 \to S^4$ for two, and $S^{15} \to S^8$ for
three. For two and three qubits the fibration is \emph{entanglement-sensitive} --- separability
is precisely the condition under which the base degenerates to a product of single-qubit Bloch
spheres and the additional dimensions associated with the higher imaginary directions collapse
\autocite{mosseri2001hopf,bernevigchen2003geometry}. This is a statement about the geometry of
state spaces and requires no physical interpretation whatever.

The second is algebraic. Singling out a preferred complex structure on the octonions ---
equivalently, fixing the split $\mathbb{O} \cong \mathbb{C} \oplus \mathbb{C}^3$ --- leaves a
residual symmetry, and that residual symmetry is exactly the gauge group of the Standard Model.
This was established through the exceptional Jordan algebra by Dubois-Violette and Todorov
\autocite{duboisviolette2016exceptional,duboisviolettetodorov2019exceptional,todorovduboisviolette2018deducing}
and given its sharpest spinorial form by Krasnov \autocite{krasnov2021so9}, who showed that the
relevant group is the stabiliser of a complex structure inside $\mathrm{Spin}(9)$. We record the
result in the form in which it enters the present work.

\begin{proposition}[Dubois-Violette--Todorov; Krasnov]\label{prop:gsm} Let $J_R$ be the complex
structure on the space $\mathbb{O}^2$ of $\mathrm{Spin}(9)$ spinors determined by a choice of
unit imaginary octonion, equivalently the split
$\mathbb{O} \cong \mathbb{C} \oplus \mathbb{C}^3$. The subgroup of $\mathrm{Spin}(9)$ commuting
with $J_R$ is \[ G_{\mathrm{SM}} \;=\; \frac{SU(3) \times SU(2) \times U(1)}{\mathbb{Z}_6}, \]
acting on the singled-out $\mathbb{C}$ summand as a colour singlet of hypercharge $-1$ and on
the $\mathbb{C}^3$ summand as a colour triplet of hypercharge $\tfrac{1}{3}$, both as $SU(2)$
doublets --- the quantum numbers of one left-handed generation.
\end{proposition}

The two results are independent, and the join between them is the only place a physical theory
must do work. Each result is a theorem; what neither supplies is a reason why nature should fix
one complex direction within $\mathbb{O}$ rather than treating all directions symmetrically. We
return to this --- the genuine load-bearing assumption --- in \S\ref{sec:hopf-honest-debt}, and
resolve it, in candidate form, in \S\ref{sec:fanout-direction}.

\subsection{The ladder is forced, not chosen}

It might still be objected that selecting the octonions, or the Hopf fibrations, is itself a
modelling choice carrying as much latitude as the choice of a grand-unified group. It is not,
and the reason is a pair of classification theorems that close off the alternatives entirely.

By Hurwitz's theorem the only normed division algebras are $\mathbb{R}$, $\mathbb{C}$,
$\mathbb{H}$ and $\mathbb{O}$ \autocite{hurwitz1898composition,baez2002octonions}. By Adams's
resolution of the Hopf-invariant-one problem --- together with the parallelisability results of
Bott--Milnor and Kervaire --- the only parallelisable spheres are $S^0$, $S^1$, $S^3$ and $S^7$,
so that there exist exactly three nontrivial Hopf fibrations, in bijection with $\mathbb{C}$,
$\mathbb{H}$ and $\mathbb{O}$
\autocite{adams1960hopf,bottmilnor1958parallelizability,kervaire1958nonparallelizability}. These
are not heuristics; they are among the more durable theorems in topology and algebra. Their
force here is that, once a framework is committed to amplitudes valued in a normed division
algebra --- a commitment that in the present construction is made far upstream, by the
self-referential argument of \S\ref{sec:self-reference} and \S\ref{sec:division-algebras}, and
not by fiat at this stage --- the ladder of Hopf fibrations is not selected but compelled, and
it terminates of its own accord.

\begin{proposition}[Termination]\label{prop:termination} The Cayley--Dickson tower of normed
division algebras and the tower of parallelisable spheres both terminate at real dimension
eight. Consequently there are exactly three nontrivial Hopf fibrations, the largest being the
octonionic $S^{15} \to S^8$, and the construction admits no extension to a fourth case.
\end{proposition}

The termination is worth dwelling on, because it is unusual for a model of this kind to predict
its own boundary. The next rung of the Cayley--Dickson sequence, the sedenions, possesses zero
divisors; the Hopf map can no longer be globally defined; there is no fourth fibration. This is
not a limitation imposed from outside but a structural fact, and it carries a physical
consequence that we develop later: there is no room in the construction for a fourth gauge force
or a fifth interaction, so that the discovery of either would count as evidence against the
framework rather than as an invitation to extend it. The endpoint is a theorem, not a
stipulation.

\subsection{Independent arrival as corroboration}

The most persuasive evidence that the destination is not an artefact of the route is that it is
reached by routes sharing none of the same machinery. Szangolies
\autocite{szangolies2025standardmodel} reaches $G_{\mathrm{SM}}$ from the emergent-gravity and
Kaluza--Klein programmes, reading the singled-out complex direction as a dimensional reduction
of a $9+1$-dimensional spacetime; the present framework reaches it from self-reference
(\S\ref{sec:self-reference}); and Furey and Gillard--Gresnigt reach the same gauge group, and
the same three-generation structure, through the minimal left ideals of the Dixon and Clifford
algebras, without invoking Hopf fibrations at all
\autocite{furey2018three,gillardgresnigt2019three}; the relation between their minimal ideals
and the entanglement classes used here is taken up in \S\ref{sec:canonical-orbits}. The
octonionic route to particle physics is in any case an old one, going back to the use of
$\mathbb{O}$ to model the strong force \autocite{gunaydingursey1974quark}.

That programmes with disjoint starting points converge on
$SU(3) \times SU(2) \times U(1)/\mathbb{Z}_6$ and on the same quark--lepton division is the
coincidence Szangolies himself remarks upon: the characteristic structure of the Standard Model,
routinely regarded as arbitrary, turns up independently in unrelated settings, and its appearing
twice would then be doubly odd unless the two settings are in fact related. We read the
convergence the same way --- as evidence that the division-algebraic description is singled out
by the mathematics rather than imposed by the modeller, and therefore that
Proposition~\ref{prop:gsm} is structural rather than contrived.

\subsection{The honest debt}
\label{sec:hopf-honest-debt}

It would be misleading to leave the impression that this route is free of assumptions, so we
state plainly the one assumption that all of these programmes share. Proposition~\ref{prop:gsm}
requires that \emph{some} complex direction in $\mathbb{O}$ be singled out, and the uniqueness
of that choice is not supplied by the theorem. Krasnov parametrises the complex structure $J_R$
by an arbitrary unit imaginary octonion, and the residual group is the same for every such
choice because $G_2 = \mathrm{Aut}(\mathbb{O})$ acts transitively on the unit imaginary
octonions. The theorem therefore guarantees the gauge group \emph{given} that a direction is
fixed, while remaining silent on \emph{why} one is fixed. Szangolies treats this as a free
choice and flags it, in his own conclusions, as an open element of the construction --- one to
which several inequivalent answers are available.

This, and not the appeal to Hopf fibrations, is the genuine load-bearing assumption of the whole
approach. The Hopf fibrations are forced (Proposition~\ref{prop:termination}); the singling-out
is not. The present framework offers a candidate resolution that the prior literature lacks: the
act of measurement, modelled as \textsc{fanout}, selects the direction physically, and the
transitivity of $G_2$ then guarantees that the gauge group is independent of which direction the
measurement happens to select. This argument is given in \S\ref{sec:fanout-direction}; we flag
here, at the appropriate epistemic status, that it is structural rather than fully proved, and
that the residual question --- why a measurement selects \emph{exactly one} imaginary direction
rather than merely \emph{permitting} a selection --- remains open. It is now priced as Open
Problem~\ref{op:fanout-selection}.

A second debt is shared with every programme of this character and is not closed by anything
said here: why the abstract quantum state should admit a tensor-product decomposition into
qubits in the first place. This is the quantum factorisation, or mereology, problem. The
entanglement-sensitivity on which the construction turns is meaningful only once such a
decomposition is fixed, and we treat the decomposition as a necessary input rather than a
derived result. Naming it is preferable to letting it pass unremarked.

\subsection{The termination read epistemically}
\label{sec:hopf-information-bound}

The termination of the tower (Proposition~\ref{prop:termination}) has so far been stated
algebraically: the Cayley--Dickson sequence acquires zero divisors at the sedenions and the Hopf
map can no longer be globally defined. There is a second reading of the same boundary, epistemic
rather than algebraic, which we offer as a motivated conjecture rather than a settled result,
because it connects the tower's endpoint to a bound that arises independently in the foundations
of measurement.

A particle, in this framework, is a reified relation between events, and its internal content is
carried by the octonionic $A,B,C$ register --- the $S^{7}$ sector serving as internal degrees of
freedom in the channel that connects the events (\S\ref{sec:space-of-computations}). The
termination of the tower is then a statement about that content: a single such relation can
support genuine entanglement among at most three parties, because there is no fourth fibration
to carry a fourth independent party within the same irreducible relation. Every particle is, on
this reading, a \emph{completed} octonionic-level relation --- topped out at three parties not
by circumstance but because the tower admits no further rung. The information that can be
encoded in one relation-between-events, and hence extracted by one measurement of it, is
correspondingly bounded: it is the information of a completed three-party register, and no
measurement of a single particle can yield more.

This is exactly the kind of bound that the epistemic-restriction programmes reach from the other
direction. Spekkens' toy theory is built on a knowledge-balance principle that caps the
information obtainable about a system \autocite{spekkens2007toy}, and Szangolies derives such an
epistemic horizon not by positing it but from the same self-referential non-computability on
which the present tower rests \autocite{szangolies_j_2018}. That the two accounts --- the
algebraic termination of the Hopf tower and the epistemic bound on measurement information ---
should coincide is, we suggest, not a numerical accident but a shared root: both are
consequences of the self-referential non-computability of \S\ref{sec:self-reference}, wearing an
algebraic face in the one case and an epistemic face in the other. The Hurwitz--Adams ceiling
and the knowledge-balance ceiling would then be one limit seen twice, and Szangolies' programme,
which reaches $G_{\mathrm{SM}}$ by an independent route (\S\ref{sec:hopf-honest-debt}) and the
epistemic horizon from the same non-computable source, shares with this framework not merely a
conclusion but its origin.

The reading must be kept distinct from a different climb, with which it would otherwise be
conflated. The Wigner's-friend regress of \S\ref{sec:wigners-friend} ascends without end: each
observer may be measured by a further observer, each carrying a fresh gauge-private circle, and
that chain has no top --- it is the bump in the carpet, unbounded in length. The termination
bounds something else: not the length of the chain of observers, but the internal structure of
each single relation, each particle, that the chain is built from. The regress is unbounded
precisely because each of its links is a bounded, completed three-party object;
termination-of-each-link and non-termination-of-the-chain are not in tension but complementary,
the former individuating the links of the latter. What the bound caps is the information in one
particle; what the regress leaves uncapped is the number of particles and observers that can be
chained. The precise identification of ``the information of a completed three-party register''
with the formal measures of the epistemic-restriction literature is the technical step that
would turn this conjecture into a theorem, and we flag it as owed; the structural claim --- that
the tower's endpoint and the measurement bound share the non-computable root --- is what we
advance here.

\subsection{The termination read compositionally}
\label{sec:hopf-compositional}

There is a third termination theorem, proved in a community that has so far had no contact with
the two joined above. Coecke and Kissinger showed that the GHZ and W classes --- realised as the
commutative special and the anti-special Frobenius algebras of categorical quantum mechanics ---
are \emph{compositionally complete}: every $N$-qubit entanglement class under SLOCC is generated
by composing these tripartite primitives with bipartite entanglement and product states
\autocite{coeckekissinger2010compositional}. Above three parties nothing new appears; the
entanglement alphabet closes at exactly the point where the fibration tower terminates. The
mutual silence of the two literatures is itself striking: Coecke and Kissinger make no reference
to Hopf fibrations or division algebras, and the fibration literature on entanglement
\autocite{mosseri2001hopf,bernevigchen2003geometry} nowhere cites compositional completeness.
Two classifications, in different mathematics, draw the same line in the same place, unaware of
one another. That silence has, in the months this paper was in draft, just begun to break from
the categorical side, with a projective, base-level reading of the same calculus
\autocite{burtonanwar2026meromorphic}; we engage it below.

The framework's conjecture is that this is one termination wearing a third face: the alphabet
has the letters it has --- Bell, GHZ, W --- because the parallelisable spheres have the
dimensions they have. There is partial evidence at every rung. At one qubit, the commutative
special Frobenius algebras on $\mathbb{C}^{2}$ are classified by orthonormal bases --- points of
$S^{2}$, the base of the first fibration --- so that the copied basis plays the role of a
section of the bundle and the spider theorem expresses its idempotence. At two qubits, the
$SU(2)$ fibre of $S^{7} \to S^{4}$ carries exactly the entanglement that survives local
operations, and L\'{e}vay's result that the bundle's natural connection is the $SU(2)$
Yang--Mills instanton \autocite{levay2004geometry} stands on the geometric side where the
interaction of two Frobenius structures stands on the categorical one. At three qubits, the
split of D\"{u}r, Vidal and Cirac \autocite{durvidalcirac2000three} mirrors the algebra exactly:
GHZ is genuinely tripartite, demanding the full non-associative octonion, while W is
pairwise-robust and housed in the associative subalgebra $\mathbb{H} \subset \mathbb{O}$ --- the
special/anti-special Frobenius distinction in geometric dress.

The first of the four conditions the correspondence requires can be discharged at once, and we
do so here, because a problem whose first step has been walked is a different kind of invitation
from one merely posed.

\begin{proposition}[The one-qubit dictionary]\label{prop:f1-dictionary} The $\dagger$-special
commutative Frobenius algebras on $\mathbb{C}^{2}$ correspond bijectively to
\emph{lifted antipodal skeletons} of the complex Hopf fibration $\pi \colon S^{3} \to S^{2}$:
unordered antipodal pairs $P \subset S^{2}$ together with a section of $\pi$ over $P$. The
copyable states of a structure are exactly the section's image, so the lifts are recorded
faithfully --- the projectively trivial centre of $U(2)$ acts non-trivially on structures ---
and the correspondence is $U(2)$-equivariant, the ordered structures forming a single free and
transitive $U(2)$-orbit. The channel a structure induces on density matrices is dephasing in its
skeleton and is independent of the section: the fibre freedom, a torsor for
$\mathrm{Map}(P, U(1))$, is exactly the data absent from the classical record a structure
induces.
\end{proposition}

\begin{proof}
By the classification of Coecke, Pavlovic and Vicary, such structures correspond to orthonormal
bases, recovered as the sets of copyable states \autocite{coeckepavlovicvicary2013orthogonal}.
An orthonormal basis of $\mathbb{C}^{2}$ is two orthogonal rays --- an antipodal pair on the
Bloch sphere --- together with a unit representative of each, which is precisely a section of
$\pi$ over the pair; and $\delta(e^{i\theta}x) = e^{i\theta}\, x \otimes x$ against
$(e^{i\theta}x) \otimes (e^{i\theta}x) = e^{2i\theta}\, x \otimes x$ pins the phases.
Equivariance, the action of the centre, the orbit structure and the induced channel are direct
computations; an executable verification of every finite check accompanies the code release.
\end{proof}

Two remarks fix the proposition's scope. It corrects a natural first reading: a classical
structure does not project to the base but selects lifts into the total space --- classicality
is section data, and what the copy operation broadcasts is a lift, not a ray. And it claims
nothing topological about cloning: the finiteness of the copyable set is linear, structural
no-cloning residing in non-naturality \autocite{abramskyheunen2012hstar}, so the first place the
bundle's global topology can bite is the quaternionic rung, where the degree-one clutching map
stands ready to be compared with the strong-complementarity law governing an interacting pair of
structures. The dictionary is, finally, the complement of the projective treatment cited above:
Burton and Anwar pass functorially to the base and deliberately discard the fibre
\autocite{burtonanwar2026meromorphic}, and the classical structures turn out to live in exactly
the data their projectivisation forgets --- the two readings compose, and the resulting triangle
is a consistency condition available to constrain the remaining conditions of the
correspondence.

\begin{openproblem}[The Hopf-fibration--Frobenius correspondence]\label{op:hopf-frobenius} Let
$\mathbf{Hopf}$ denote the category of the three fibrations
$S^{2^{k}-1} \hookrightarrow S^{2^{k+1}-1} \to S^{2^{k}}$, $k = 1, 2, 3$, with bundle maps as
morphisms, and let $\mathbf{Frob}$ denote the Frobenius algebras on $\mathbb{C}^{2}$ in
$\mathbf{FdHilb}$, with the commutative-special (GHZ) and anti-special (W) classes
distinguished. Construct a functor $F \colon \mathbf{Hopf} \to \mathbf{Frob}$ --- or prove that
none exists --- under which the complex fibration yields the single-qubit classical structures
--- Proposition~\ref{prop:f1-dictionary} identifies these as the lifted antipodal skeletons,
section data over two base points rather than the points alone; the quaternionic fibration
yields the interaction of two such structures, with the $SU(2)$ fibre carrying the entanglement
that survives the classical quotient; the octonionic fibration yields the GHZ/W pair, GHZ on the
full non-associative fibre and W on the associative $\mathbb{H} \subset \mathbb{O}$; and the
composition by which Coecke--Kissinger generate the $N$-qubit classes corresponds to a
well-defined geometric operation on the bundles. If such an $F$ exists, compositional
completeness follows from Adams's theorem as a topological corollary --- there is no fourth
fibration to supply a fourth primitive --- and if $F$ is faithful enough that the implication
also runs back, the termination of the tower and the completeness of the alphabet are equivalent
expressions of a single classification. The first condition is discharged by
Proposition~\ref{prop:f1-dictionary}; what remains open is the quaternionic and octonionic
conditions and, above all, the correspondence of compositions. Both open conditions are already
constrained: the quaternionic fibre's physical content --- the recoverable relative phase and
its purity ledger --- is established (\S\ref{sec:wf-reconciliation}), leaving the
clutching-map-to-strong-complementarity comparison as the structural residue; and any candidate
functor must reproduce the proved canonical-orbit images, GHZ to the colour triplet and W to the
singlet (\S\ref{sec:canonical-orbits}, Appendix~\ref{app:ghz-w-proofs}). A proof that no such
functor exists, or a correspondence at the level of the classifications without functoriality on
compositions, would each also resolve the problem: the first would make the agreement of two
independent classifications a fact demanding its own explanation; the second would establish the
weaker reading.
\end{openproblem}

A resolution in the affirmative would do three things. It would explain compositional
completeness rather than merely prove it --- a corollary of the division-algebra structure of
quantum mechanics rather than an independent algebraic fact. It would give the diagrammatic
language of this paper's own interaction circuits a topological warrant: the GHZ/W spider
calculus would be the unique language the topology permits, not a convenient notation. And it
would open a two-way bridge, translating problems on compositions of Frobenius structures into
operations on bundles and back. The problem sits at the intersection of categorical quantum
mechanics, the quantum-information geometry of the fibrations, and the division-algebra
approaches to particle physics, and it is posed here so that any of the three communities can
take it up in its own language. It is perhaps worth noticing that both classifications already
answer to the same name --- the fibrations are Heinz Hopf's, and so, through the bialgebra
structure of the diagrammatic calculus, are the Hopf algebras --- a coincidence of eponymy that
the conjecture proposes to promote to a theorem. With the epistemic reading above, an
affirmative resolution would leave the tower's endpoint as one limit seen four ways: algebraic,
geometric, compositional, epistemic --- the last being the fixed-point classification of Open
Problem~\ref{op:sigma-functor}, with which this problem is cross-linked, each constraining the
other.

\subsection{What this sets up}

The entanglement-sensitivity of the Hopf maps is not a side feature but the hinge of the
physical reading that follows. Because separability is exactly the collapse of the base to a
product of Bloch spheres, the structure of the entanglement --- the SLOCC class of a multi-qubit
state, and within it the canonical maximally entangled orbit --- is carried by precisely the
degrees of freedom on which $G_{\mathrm{SM}}$ acts in Proposition~\ref{prop:gsm}. This is what
makes it possible, in \S\ref{sec:entanglement-classes}, to read the two genuinely tripartite
entanglement classes onto the $\mathbb{C} \oplus \mathbb{C}^3$ split through their canonical
representatives: the GHZ orbit confined to the colour-triplet $\mathbb{C}^3$ sector, the W orbit
reaching the $\mathbb{C}$ singlet, recovering the quark--lepton division not as an assignment
but as a consequence of the entanglement structure. The scope of those statements --- canonical
orbits rather than whole classes, and the identification of the canonical orbits with the
physical particles --- is made precise there. The Hopf fibrations are the bridge; the
entanglement classes, treated next, are the traffic that crosses it.

%      % EDITED 2026-09-23 scripts/patch_prop_w_access_outline_2026-09-23T1100.py
% EDITED 2026-09-10 scripts/patch_quoting_and_typography_2026-09-10T0620.py
% =====================================================================
%  SECTION: Entanglement Classes and Gauge Groups
%  (sec:entanglement-classes). Drafted 2026-06-08; supersedes the
%  earlier same-day draft. AUDITED 2026-07-06 (final pass, ninth stop).
%
%  PLACEMENT (reordered chain, JS 2026-07-06): position 9 --
%      \input{hopf_justification_section_2026-06-08}
%      \input{entanglement_classes_and_gauge_groups_2026-06-08}
%      \input{register_section_2026-06-08}
%  The "previous section" phrases in the prose refer to
%  hopf-justification and are TRUE at this position (the hopf->classes
%  adjacency survived the reorder); the file is position-fragile in
%  that specific pairing, by design.
%
%  APPENDIX LABELS ACQUITTED: app:ghz-proof (line 246 of the appendix)
%  and app:w-proof (line 364) both exist -- the June "written next
%  turn" promise was kept; app:setup and the ceiling lemma live.
%  sec:generations is defined in the particle-zoo section; all forward
%  refs land.
%
%  AUDIT RECORD: lem:cproj, prop:ghz, prop:w-access, prop:w-doublet and
%  the canonical-orbits postulate all verified against every quoting
%  consumer (the frame's (2/3,1/3) quote exact; the register delivers
%  the deferred SU(2)_L/handedness fixing); the e1-alignment point
%  (Fano pairing, not commutativity) sound; loci numbers verified.
%  Fixes: CKW key variant normalised to corpus canon
%  (coffmankunduwootters2000); chirality return-promise given its
%  addresses (sec:frame, sec:generations); absent-gate forward stitch
%  at the DVC sentence; and the TWO-NINTHS IDENTITY drawn out --
%  sup C_frac = r^2 (ceiling lemma) and slack = 1 - 4 det rho = r^2,
%  so access ceiling and budget capacity are one invariant (W: 1/9;
%  GHZ: 0 -- confinement and vanishing C-access are one fact), a
%  partial discharge of op:w-slack's physical-reading clause recorded
%  in the confinement section; script check added. op:ideal-orbit
%  inventoried and stands as posed (well-formed, counterexample-valued).
%  The header's stray Gödel umlaut dies with this rewrite: file now
%  pure ASCII. British English; \autocite; no bullets in prose.
% =====================================================================

% ADDED 2026-07-12: sec:algebraic-dictionary prior-work subsection
% (Furey/Dixon/Krasnov dictionary; two REF-OWED keys; see its header).
\section{Entanglement Classes and Gauge Groups}
I'll quote from Szangolies \autocite{szangolies2025standardmodel} because this is very much his idea. 
\begin{quote}
Our main result is then to show how to use a construction of Dubois-Violette and Todorov \autocite{duboisviolette2016exceptional}, 
as extended and clarified by Krasnov \autocite{krasnov2021so9}, to derive the exact symmetry
and degrees of freedom of (a generation of) Standard Model fermions from the entanglement properties of a three-qubit state.
\end{quote}
I had been looking for a relationship between entanglement and gauge theory for many years and then I saw Szangolies 
give a ten minute session in an online workshop on division algebras in particle physics \autocite{szangolies_j_2021} 
during the presentation he flashed up a screen shot of a paper he was working on relating entanglement to particle physics.
I instantly knew this was what I was looking for and contacted him to ask if the paper was ready for publication. He very kindly let me
see a pre-print and I encouraged him with the phrase "perfection is the enemy of done". So I was very please to see the final thing published. 

A short while after that at a conference on quantum foundations in V\"axj\"o in 2023 I saw a single A4 poster by Filtov describing
a geometrical view of bi-partite entanglement with two Block spheres related by inversion. I already knew that bi-partite entanglement was 
related to SU(2), so the fact that Filtov had so clearly demonstrated the chirality inherent in SU(2) was the most significant 
contribution to the conference. A poster on a single sheet of A4! Which eventually led to \autocite{filatov2024towards}

That got me thinking that relating 3 Bloch spheres by inversion would be impossible and therefore if you could show a relationship between 3 particle entanglement and SU(3)
then the the fact that the colour force is blind to chirality would fall out immediately. And this is where we'll pick up. 

\label{sec:entanglement-classes}

The previous section established that the division-algebra ladder forces exactly three nontrivial
Hopf fibrations, $S^3 \to S^2$, $S^7 \to S^4$ and $S^{15} \to S^8$, with fibres $S^1$, $S^3$ and
$S^7$, and that fixing a complex direction inside the octonions leaves the Standard Model gauge
group as the stabiliser in $\mathrm{Spin}(9)$ (Proposition~\ref{prop:gsm}). This section asks where,
in the entanglement of a three-qubit state, the three gauge groups actually live. The answer is
clean for the first two and forces a genuine question for the third, and it is that question --- and
its resolution through the two classes of tripartite entanglement --- that the section is about.

\subsection{From fibres to groups, and where the pattern breaks}

For the first two rungs the identification needs no interpretation, because the fibre simply
\emph{is} the gauge group. The fibre of the first Hopf fibration is $S^1$, the unit circle in the
complex plane; as a group the unit complex numbers are exactly $U(1)$, the gauge group of
electromagnetism. The fibre of the second is $S^3$, the unit sphere in the quaternions; the unit
quaternions form the group $\mathrm{Sp}(1) \cong SU(2)$, the gauge group of the weak interaction.
So the abelian phase $U(1)$ and the non-abelian $SU(2)$ are not analogues of the Hopf fibres --- they
are those fibres, read as groups. A reader from quantum computation will recognise $S^1$ as the
global phase of a single qubit and $S^3 \cong SU(2)$ as the group of single-qubit unitaries --- the
Bloch-sphere rotations; a reader from gauge theory will recognise the two smallest gauge groups.

The pattern then breaks, and it is important to be clear about why, because the break is what makes
the rest of the section necessary. The fibre of the third, octonionic, Hopf fibration is $S^7$, the
unit sphere in the octonions. But the unit octonions are \emph{not a group}: octonionic
multiplication is non-associative, so $S^7$ does not close into a Lie group at all --- it is a
Moufang loop, not a group. There is therefore no statement ``$S^7 \cong SU(3)$'' to be written down
in parallel with $S^1 \cong U(1)$ and $S^3 \cong SU(2)$. The third gauge group cannot be the fibre.

This is not a defect to be patched but a signpost. The previous section already located the only
candidate: $SU(3)$ is the subgroup of $G_2 = \mathrm{Aut}(\mathbb{O})$ that stabilises a complex
direction, equivalently the split $\mathbb{O} \cong \mathbb{C} \oplus \mathbb{C}^3$
(Proposition~\ref{prop:gsm}). That tells us the group is available; it does not tell us where in the
physics it acts. Because $S^7$ is not itself a group, $SU(3)$ must enter not as the fibre but as a
symmetry of the \emph{structure carried on} the third rung --- and the structure carried there is the
entanglement of the three-qubit state. Two questions follow, and they are the questions this section
answers. First: can the colour $SU(3)$ be located in the entanglement of three qubits, and if so in
which class --- the GHZ class or the W class? Second: does the split between those two classes, the
only two kinds of genuinely tripartite entanglement, map onto the two matter sectors that colour
distinguishes, namely the coloured quarks and the colourless leptons?

These are not idle questions, and they are not new ones. Szangolies has already established part of
the gauge-group link from the representation-theoretic side \autocite{szangolies2025standardmodel};
Filatov and Auzinsh have shown that bipartite entanglement carries an intrinsic chirality of the
kind the weak $SU(2)$ exhibits \autocite{filatov2024towards}; and, for well over a decade, Duff and
collaborators at Imperial College have pursued deep correspondences between multipartite
entanglement and string and brane physics \autocite{duff2007string,borsten2009blackholes}. We take
up the pedigree below and then turn the questions into theorems.

\subsection{Chirality is already present in bipartite entanglement}

Before the tripartite case, a structural observation about the bipartite one, because it bears on
the weak sector that sits on the $S^3 \cong SU(2)$ rung. Filatov and Auzinsh showed that a pure
maximally entangled two-qubit state cannot be drawn on two Bloch spheres of the same handedness:
representing the entanglement faithfully requires the coordinate axes of one sphere to be
right-handed and those of the other left-handed \autocite{filatov2024towards}. Handedness is thus
not an artefact of how one chooses to draw the state; it is forced by the entanglement, and it is a
topological fact about the quaternionic Hopf fibration --- the same fibration whose fibre is the weak
$SU(2)$.

The weak interaction is, famously, the one sector of the Standard Model that is chiral: it couples
to left-handed fermions and not to their right-handed counterparts, and that chirality is usually
treated as an empirical input with no structural explanation. It is therefore at least suggestive
that chirality should appear, unbidden and forced, in the simplest entangled state of the very rung
that carries $SU(2)$. We draw no specific derivation from this here --- the mechanism that turns the
Filatov handedness into the physical left-handedness of the weak force belongs to the later
discussion of the register, the frame, and the generations (\S\ref{sec:register},
\S\ref{sec:frame}, \S\ref{sec:generations}) --- but we record it as a structural encouragement
for a programme that proposes to read gauge content off entanglement, and we return to it there.

\subsection{The two classes of tripartite entanglement}

Three qubits admit exactly two inequivalent classes of genuinely tripartite entanglement under
stochastic local operations and classical communication: the GHZ class and the W class
\autocite{durvidalcirac2000three}. (The operational face of the inequivalence --- the converter
mathematics permits and the vertex dictionary does not stock --- is \S\ref{sec:absent-gate}.) Every natural invariant separates them. The three-tangle is
maximal on GHZ and identically zero on W \autocite{coffmankunduwootters2000}; tracing
out one qubit of a GHZ state leaves the other two separable, while the same operation on a W state
leaves them entangled. Coecke and Kissinger gave the distinction an algebraic form: GHZ-class states
carry a \emph{special} commutative Frobenius algebra and W-class states an \emph{anti-special} one,
the two differing in a purely topological way \autocite{coeckekissinger2010compositional}. Those two
Frobenius structures are the primitives of a calculus expressive enough to build arbitrary $N$-qubit
states, so the entanglement alphabet, like the division-algebra ladder, closes at three: there is no
third primitive to find, just as there is no fourth division algebra (Proposition~\ref{prop:termination}).

We're now at the point where it should seem completely obvious to ask the question "That split between GHZ and W looks suspiciously like the split between 
the quark and lepton sectors, tell me more about it?" That there are exactly two genuinely tripartite classes, and exactly two matter sectors that colour
distinguishes, is the coincidence the rest of the section makes precise.

\subsection{The map into the octonion algebra}

We use the identification of a three-qubit state $\psi \in \mathbb{C}^8$ with a pair of octonions
$(o_1, o_2) \in \mathbb{O}^2$, obtained by splitting on the first qubit under the Cayley--Dickson
correspondence, together with the octonionic Hopf map $\mathbf{h} = o_1\bar{o}_2$. Fixing the
preferred complex direction $e_1$ gives the split $\mathbb{O} \cong \mathbb{C} \oplus \mathbb{C}^3$,
with $\mathbb{C} = \{e_0, e_1\}$ the singlet direction and $\mathbb{C}^3 = \{e_2,\dots,e_7\}$ the
colour-triplet sector carrying the $\mathbf{3} \oplus \bar{\mathbf{3}}$ of the $SU(3)$ of
Proposition~\ref{prop:gsm}. Everything below turns on one identity.

\begin{lemma}[The $\mathbb{C}$-projection is a Hermitian inner product]\label{lem:cproj}
The projection of $\mathbf{h}$ onto the $\mathbb{C} = \{e_0,e_1\}$ summand is
$\pi_{\mathbb{C}}(\mathbf{h}) = \overline{\langle o_1, o_2\rangle_{\mathbb{C}^4}}$, the conjugate of
the standard Hermitian inner product of the two first-qubit branches. Hence
$\pi_{\mathbb{C}}(\mathbf{h}) = 0$ if and only if those branches are $\mathbb{C}^4$-orthogonal.
\end{lemma}

The lemma is proved in Appendix~\ref{app:setup}, where the octonion multiplication table and the
identification of $\mathbb{C}^4$ with $\mathbb{O}$ are fixed explicitly; the identification must be
adapted to the complex structure given by left multiplication by $e_1$, and the sign bookkeeping is
part of the proof. The composition-algebra identity $|\mathbf{h}|^2 = |o_1|^2|o_2|^2$, which holds on
$\mathbb{O}$ and no larger Cayley--Dickson algebra, makes the $\mathbb{C}$-fraction
$C_{\mathrm{frac}} = |\pi_{\mathbb{C}}(\mathbf{h})|^2/|\mathbf{h}|^2$ a well-defined local invariant.
What singles out $e_1$ is not commutativity --- every imaginary direction spans a commutative plane
with $e_0$ --- but alignment: $e_1$ is the unique imaginary direction whose Fano pairing of the basis
units coincides with the pairing of real and imaginary parts in the qubit amplitudes, and the GHZ
cancellation below requires exactly that alignment.

\subsection{GHZ states are confined to the colour sector}

\begin{proposition}[GHZ confinement]\label{prop:ghz}
For every state in the local-unitary orbit of $\lvert\mathrm{GHZ}\rangle$, the $\mathbb{C}$-projection
of the Hopf image vanishes identically, so the image lies entirely in the colour-triplet sector
$\mathbb{C}^3 = \mathbf{3}\oplus\bar{\mathbf{3}}$, symmetrically under each of the three qubit splits.
This property characterises the orbit exactly: the $\mathbb{C}$-projection vanishes across the orbit
under every split if and only if all three single-qubit reduced states are maximally mixed, which
holds if and only if the state is local-unitarily equivalent to $\lvert\mathrm{GHZ}\rangle$.
\end{proposition}

\emph{Proof outline.} By Lemma~\ref{lem:cproj} it is enough that the branches are orthogonal for
every local-unitary image. The bare GHZ state has orthogonal branches of \emph{equal} norm
($\rho_k = \tfrac12\mathbb{I}$ for each $k$); the internal unitaries preserve orthogonality and equal
norms, and the first-qubit unitary mixes the branches with an inner product that evaluates to
$-\bar a\bar b\cdot\tfrac12 + \bar b\bar a\cdot\tfrac12 = 0$, the cancellation holding because
$\mathbb{C}$ is commutative. The full derivation --- the characterisation, and the symmetry across
all three splits --- is given in Appendix~\ref{app:ghz-proof}. \hfill$\square$

The cancellation is exact; over random local-unitary samples the $\mathbb{C}$-fraction sits at
machine zero. Two consequences are proved rather than conjectured. Confinement is frame-independent:
the invariant separating the two canonical orbits is the single-qubit reduced spectrum --- equal
$(\tfrac12,\tfrac12)$ on the GHZ orbit, unequal $(\tfrac23,\tfrac13)$ on the W orbit --- and since no
local unitary can change that spectrum, no change of local frame can carry a colour triplet into the
lepton sector. The scope of the statement is the canonical orbits: SLOCC deformations away from them,
such as $\alpha\lvert000\rangle+\beta\lvert111\rangle$ with $\alpha\neq\beta$, acquire a nonzero
$\mathbb{C}$-projection, and their status is fixed by the identification of
\S\ref{sec:canonical-orbits}. And the QCD vacuum angle vanishes: the GHZ image has
no oriented pairwise linking structure to orient a $\theta\,\mathrm{Tr}(G\wedge G)$ term, so
$\theta_{\mathrm{QCD}} = 0$ is forced and no Peccei--Quinn axion is admitted. The reading is
GHZ $\leftrightarrow$ colour triplet $\leftrightarrow$ quark.

This is the first result to drop out of the identification between GHZ and the colour force, and it was the one 
I was looking for, SU(3) is blind to chirality because you cannot have a geometric construction analogous to the 
2 Bloch sphere case with 3 Bloch spheres. But once a bridge between the gauge groups of the Standard Model
and multi-particle entanglement has been built then a lot of two way traffic becomes possible. One can move hard problems
on one end of the bridge to the other end where they have easy solutions, $\theta_{QCD}=0$ being only the first. This also means
you can express particle interactions in terms of quantum circuits acting on entangled states, which we pick up in \ref{sec:interactions}

\subsection{W states reach the lepton sector and carry the weak doublet}

\begin{proposition}[W accessibility]\label{prop:w-access}
For states in the local-unitary orbit of $\lvert\mathrm{W}\rangle$, the $\mathbb{C}$-projection of the
Hopf image is generically nonzero, vanishing only on a measure-zero subset, with
$\sup C_{\mathrm{frac}} = \tfrac19$.
\end{proposition}

Where the GHZ orbit is confined to the colour sector, the W orbit reaches the lepton sector, and
does so generically: the $\mathbb{C}$-component of its Hopf image is nonzero at every point of the
orbit but a set of measure zero, and it is bounded, never exceeding one ninth of the image's squared
norm. Both facts come from the same one-line computation as Proposition~\ref{prop:ghz}, with one
norm changed.

\emph{Proof outline.} Split on the first qubit. The bare W has branches
$o_1 = \tfrac{1}{\sqrt3}(0,1,1,0)$ and $o_2 = \tfrac{1}{\sqrt3}(1,0,0,0)$, orthogonal but of
\emph{unequal} norm, $\lvert o_1\rvert^2 = \tfrac23$ and $\lvert o_2\rvert^2 = \tfrac13$. Unitaries
on the other two qubits act on both branches alike and change neither their inner product nor
their norms, so $C_{\mathrm{frac}}$ depends only on the first-qubit unitary
$U_1 = \bigl(\begin{smallmatrix} a & b \\ -\bar b & \bar a\end{smallmatrix}\bigr)$. Orthogonality
kills the $\bar a^{2}$ and $\bar b^{2}$ terms of the mixing formula, and the unequal norms leave the
residue $\langle o_1, o_2\rangle = \bar a\bar b\,(\tfrac13 - \tfrac23) = -\tfrac13\bar a\bar b$: the
equal-norm cancellation of the GHZ proof, failing by exactly the norm difference. It vanishes only
when $ab = 0$, $U_1$ diagonal or anti-diagonal, a measure-zero subset of $SU(2)$ on which the bare
representative itself sits. With $p = \lvert a\rvert^2$ the branch norms are
$\lvert o_1\rvert^2 = (1+p)/3$ and $\lvert o_2\rvert^2 = (2-p)/3$, and by Lemma~\ref{lem:cproj}
together with the composition identity,
\[
C_{\mathrm{frac}} \;=\; \frac{p(1-p)}{(1+p)(2-p)},
\]
maximised at $p = \tfrac12$ with value $\tfrac19$: the first-qubit Bloch vector rotated into the
equatorial plane. The reduced spectrum $(\tfrac23,\tfrac13)$ is a local-unitary invariant, so the
statement holds across the orbit, and by the permutation symmetry of $\lvert\mathrm{W}\rangle$ under
each of the three splits. The ceiling is the instance $r = \tfrac13$ of the general identity
$\sup C_{\mathrm{frac}} = r^{2}$, the squared Bloch radius of the marginal (Appendix~\ref{app:setup},
Lemma~\ref{app:lem:ceiling}); the full derivation is in Appendix~\ref{app:w-proof}.
\hfill$\square$

Two remarks. The geometry is sharper than a bound. With $r$ fixed at $\tfrac13$ on the orbit the
colour component of the image has constant norm $\sqrt{8/9}$, so the W orbit's Hopf image lies on
the product of a colour five-sphere of radius $\sqrt{8/9}$ and a Bloch two-sphere of radius
$\tfrac13$, disjoint from the colour equator of unit radius that carries the GHZ image; the missing
ninth of squared norm sits in the $\mathbb{C}$-plane, and the two canonical orbits occupy disjoint
compact loci on the base, which is the geometric form of the statement that no change of local
frame connects them. And the ceiling has been measured: single-wire tomography of the W register on
hardware returned a probe-corrected $r^{2} = 0.1095 \pm 0.0051$ against $\tfrac19 = 0.1111$
(Appendix~\ref{app:hw-ninth}), the number this paper calls the leptonic ninth.

\begin{proposition}[W carries the $SU(2)$ doublet]\label{prop:w-doublet}
Each single-qubit reduced state of a state in the local-unitary orbit of $\lvert\mathrm{W}\rangle$
has the non-degenerate spectrum $(\tfrac23,\tfrac13)$, symmetrically across the three qubits. The
non-degeneracy singles out a distinguished eigenbasis --- a $T_3$ axis with its two doublet members
--- breaking the $SU(2)$ invariance of the marginal down to the $U(1)$ of rotations about that axis.
The degenerate GHZ marginal $\tfrac12\mathbb{1}$ singles out no such direction.
\end{proposition}

\emph{Proof outline.} Immediate from the spectrum. The single-qubit marginal of
$\lvert\mathrm{W}\rangle$ is $\mathrm{diag}(\tfrac23,\tfrac13)$ in the computational basis, the
same on each qubit by the permutation symmetry, and a local-unitary invariant up to conjugation.
Non-degeneracy fixes a unique eigenbasis, a distinguished $T_3$ axis between whose members the
$SU(2)$ raising and lowering operators act, and leaves invariant only the $U(1)$ of rotations about
that axis. The degenerate GHZ marginal $\tfrac12\mathbb{1}$ has every basis as an eigenbasis and
singles out no direction. Full derivations are in Appendix~\ref{app:w-proof}. \hfill$\square$

The contrast with GHZ is entirely in the branch norms: equal norms force the cancellation, unequal
norms leave a residue. The reading is W $\leftrightarrow$ colour singlet $\leftrightarrow$ lepton,
organised into weak-isospin doublets. Which of the two quaternionic actions, $SU(2)_L$ or $SU(2)_R$,
becomes physical weak isospin --- and how the Filatov handedness of \S\ref{sec:entanglement-classes}
becomes physical left-handedness --- is fixed by the structure of the full register, and we defer it
to \S\ref{sec:register}.

\subsection{Canonical orbits are the particles}
\label{sec:canonical-orbits}

The propositions above are scoped to the local-unitary orbits of the standard GHZ and W states ---
the maximally entangled representatives of their classes. The SLOCC classes are larger: each contains
a continuum of deformations, such as $\alpha\lvert000\rangle+\beta\lvert111\rangle$ with
$\alpha\neq\beta$, whose marginal spectra differ and for which the $\mathbb{C}$-projection switches
on. We therefore make the identification explicit, and adopt it as a structural postulate of the
framework: \emph{a physical one-particle state is a canonical maximally entangled orbit; the SLOCC
class is the particle type}. The quantum numbers of \S\ref{sec:charges} and the masses of
\S\ref{sec:masses} are computed on canonical orbits, so the postulate is load-bearing, and we state
it rather than leave it implicit.

The identification is not arbitrary, and the reason lives in machinery the paper already uses ---
and, since \S\ref{sec:superconductor-witness}, in an independent physical witness: the
superconductor, where the Bell and product rows are the two strata of the Anderson--Higgs
mechanism. Under
map--state duality, a maximally entangled state is the state-form of the identity map --- the algebra
acting on itself, written as a state --- and its local-unitary orbit corresponds to the unitaries,
while the SLOCC deformations correspond to invertible but non-unitary maps, excluded as physical
configurations for the same reason non-unitary evolutions are. In the Frobenius language
\autocite{coeckekissinger2010compositional}, the canonical orbits are precisely the \emph{special}
and \emph{anti-special} normalisations of the two tripartite algebras, and dagger-special commutative
Frobenius algebras correspond exactly to orthonormal bases
\autocite{coeckepavlovicvicary2013orthogonal}: the canonical orbit is the Frobenius algebra in its
canonical form, the tilted states its non-special deformations.

This converges with the identification at the heart of Furey's programme, in which particles are
minimal ideals of an algebra acting on itself \autocite{furey2012unified,furey2018three}. The
physical idea is the same --- the particle is the canonical, parameter-free substructure singled out
by the algebra's self-action --- expressed there in module coordinates and here in entanglement
coordinates. The selection principles differ in form (minimality of the ideal; maximal symmetry of
the orbit), and we record their reconciliation as a named open problem.

\begin{openproblem}[Ideals and canonical orbits]\label{op:ideal-orbit}
Prove that the dagger-special condition on the Frobenius side corresponds to the minimal-ideal
condition on the division-algebra side, so that the present identification and Furey's become one
fact described in two coordinate systems. A counterexample separating the two selections on some
sector, or a correspondence holding only on a proper subclass, would each also resolve the problem,
and either outcome would be informative: it would delimit exactly the sectors on which module
coordinates and entanglement coordinates may be used interchangeably.
\end{openproblem}

The status of the tilted, off-canonical states --- whether they admit a reading as
off-shell or perturbed configurations --- is likewise open; no result downstream depends on them.

\subsection{The dictionary}

Collecting the strands gives the correspondence between entanglement structure and gauge content,
the upper rows established in the previous section and the lower two the content of
Propositions~\ref{prop:ghz}--\ref{prop:w-doublet}.

\begin{center}
\begin{tabular}{llll}
\hline
Entanglement & Division algebra & Hopf fibre & Standard Model \\
\hline
separable & $\mathbb{C}$ & $S^1 \cong U(1)$ & electromagnetic phase \\
Bell (two-qubit) & $\mathbb{H}$ & $S^3 \cong SU(2)$ & weak isospin \\
GHZ (special) & $\mathbb{O}$, $\mathbb{C}^3$ sector & $S^7$ & $SU(3)$ colour / quarks (confined) \\
W (anti-special) & $\mathbb{O}$, $\mathbb{C}$ sector & $S^7$ & colour singlets / leptons (free) \\
\hline
\end{tabular}
\end{center}

The first two rows are group identities; the lower two are not, precisely because $S^7$ is not a
group, and that is why they required the theorems above rather than a one-line reading of the fibre.
A clarification of where the content lies: the class names in the left column are labels, and the
physics is in the geometry of the Hopf map they label. The base $S^8$ decomposes as (Bloch vector of
the split qubit) $\oplus$ (colour sector $\mathbb{C}^3$), and the canonical GHZ orbit maps onto the
unit sphere of the colour sector --- the equatorial $S^5\subset S^8$ --- while the canonical W orbit
occupies a disjoint locus of constant colour radius $\sqrt{8/9}$ and Bloch radius $\tfrac13$, with
$\mathbb{C}$-access ceiling $(\tfrac13)^2 = \tfrac19$ (Appendix~\ref{app:setup}). The two octonionic
rows share the same fibre and are separated as disjoint loci on the base: on the colour equator, or
tilted off it into the single-qubit Bloch directions.

One identity deserves to be drawn out, because it joins this section's geometry to the entanglement
budget of \S\ref{sec:budget-converse}. The ceiling lemma of Appendix~\ref{app:setup} gives
$\sup C_{\mathrm{frac}} = r^{2}$, the squared marginal Bloch radius; and the budget slack of a
register qubit is $1 - \tau = 1 - 4\det\rho = r^{2}$ identically. So the $\mathbb{C}$-access
ceiling and the external-correlation capacity are \emph{the same invariant}, $r^{2}$, seen in base
and in budget coordinates: for the W orbit both are exactly $\tfrac19$ --- the leptonic ninth of
Open Problem~\ref{op:w-slack} \emph{is} the lepton's singlet-direction access --- and for the GHZ
orbit both are exactly zero, so confinement and vanishing $\mathbb{C}$-access are one fact, not
two. An executable verification accompanies the code release.

\subsection{Prior work: Szangolies, and the Imperial programme}
\label{sec:entanglement-szangolies}

The overlap with existing work is worth stating precisely, because the convergence is itself
evidence. Szangolies obtains the same lepton/quark split directly from the representation theory of
the $\mathrm{Spin}(9)$ stabiliser: the octonionic spinor decomposes into a leptonic part from the
$\mathbb{C}$ factor, an $SU(3)$ singlet of hypercharge $-1$, and a quark part from $\mathbb{C}^3$, an
$SU(3)$ triplet of hypercharge $\tfrac13$, both $SU(2)$ doublets \autocite{szangolies2025standardmodel}.
This is reached with no reference to entanglement classes; the state is treated as carrying a
generation wholesale, and the split is a fact about the gauge group acting on
$\mathbb{C}\oplus\mathbb{C}^3$. Our result is complementary: where Szangolies fixes the representation
content of each summand, Propositions~\ref{prop:ghz} and~\ref{prop:w-access} fix \emph{which states
populate which summand}, by their SLOCC class. The lepton/quark division he assigns by representation
theory we obtain as a consequence of entanglement type, and his independent arrival from the gauge
side is the strongest external check that the identification is not an artefact of our route.

The longer pedigree matters too, because the idea that multipartite entanglement encodes the
structures of fundamental physics is not new and has been pursued by senior researchers for well over
a decade. Duff and collaborators at Imperial College developed the ``black holes, qubits and
octonions'' programme, beginning from the observation that the entropy of certain stringy black holes
and the tripartite entanglement of three qubits are governed by the same invariant, Cayley's
hyperdeterminant \autocite{duff2007string}, and extending to $N=8$ supergravity, where the
56-dimensional charge space of $E_7$ decomposes into seven tripartite sectors arranged on the lines
of the Fano plane --- the multiplication table of the octonions --- with Cartan's quartic $E_7$
invariant playing the role of the entanglement measure
\autocite{dufferrara2007seven,levay2007fano,borsten2009blackholes}. Read from the present vantage
point, that programme had the octonionic skeleton in hand: the Fano plane, the exceptional groups,
the three-qubit classification. The divergence is in the routing. There, the geometry was carried
into supergravity charge space, and the correspondence is a dictionary between invariants --- an
entanglement measure matched to an entropy formula, number to number, with the mechanism beneath the
dictionary left open. Here, the geometry stays in the state space, where the fibration's stabiliser
supplies the gauge group as a mechanism (Proposition~\ref{prop:gsm}) and the entanglement classes
serve as labels for loci of the Hopf map. The Imperial programme establishes that qubit entanglement,
division algebras and the same three-qubit classification at the foundations of physics is a
direction taken seriously, independently, by others; what we add is the routing of that structure
into the gauge sectors of the Standard Model, as theorems rather than analogies.

\subsection{Prior work: the division-algebra programmes, and the two coordinatisations}
\label{sec:algebraic-dictionary}
% ADDED 2026-07-12. Sibling of the Szangolies/Imperial subsection:
% positions the framework against Furey, Dixon, Gillard--Gresnigt and
% Krasnov as TWO COORDINATISATIONS of one structure. REF-OWED (two new
% keys, .ris cluster to follow): dixon1994division (G. M. Dixon,
% "Division Algebras: Octonions, Quaternions, Complex Numbers and the
% Algebraic Design of Physics", Kluwer 1994); krasnov2021so9
% (K. Krasnov, the Spin(9) / centraliser-of-a-complex-structure
% characterization of the Standard Model gauge group, J. Math. Phys.
% 62, 021703 (2021), DOI 10.1063/5.0039941, arXiv:1912.11282 --
% VERIFIED 2026-07-12); partonpiccinni2018spin9 (Parton & Piccinni,
% The Role of Spin(9) in Octonionic Geometry, arXiv:1810.06288 --
% arXiv verified, journal DOI to confirm at Mendeley import).
% 2026-07-12 (same day, later): charge lemma DISCHARGED --
% appendix_winding_number_operator_2026-07-12.tex, lem:winding-number;
% op:dictionary-lemmas shrunk to the chirality item, label unchanged.
% CUTOVER: add the
% composability sentence with \ref{app:nb-I} and the scan reference
% \ref{app:nb-M} when the notebook record is input.

A third body of prior work runs on the algebraic side of the same division algebras: Dixon's
tensor algebra $\mathbb{R} \otimes \mathbb{C} \otimes \mathbb{H} \otimes \mathbb{O}$
\autocite{dixon1994division}, Furey's particles as minimal ideals with the operator content
generated by left multiplication \autocite{furey2012unified,furey2018three}, the
Gillard--Gresnigt Clifford-algebra routes to three generations
\autocite{gillardgresnigt2019three}, the exceptional-Jordan wing
\autocite{duboisviolettetodorov2019exceptional}, and Krasnov's characterisation of the gauge
group inside $\mathrm{Spin}(9)$ \autocite{krasnov2021so9}. The relationship to the present
framework is closer than parallel effort, and the organising fact is worth stating as a rule:
\emph{the algebra fixes the skeleton; the geometry carries the measure.} Ideal decompositions
and ladder counting yield exact integer data --- charge quantisation in thirds, representation
content, a generation count --- because an algebra has no metric to get wrong. Fibrations yield
the continuous data --- $\tfrac29$, $\tfrac14$, the geodesic masses of \S\ref{sec:masses}, the
Born rule itself --- because spheres do. The two are coordinatisations of one structure, and the
dictionary is short enough to give whole:

\begin{center}
\begin{tabular}{lll}
\hline
object & module coordinates & entanglement coordinates \\
\hline
arena & $\mathbb{C}\otimes\mathbb{O}$; Dixon's $\mathbb{R\,C\,H\,O}$ &
  the register $\mathbb{C}^8$; the nesting $\mathbb{C}\subset\mathbb{H}\subset\mathbb{O}$ \\
operators & $\mathbb{C}l(6)$ from left multiplication &
  the register's left-action envelope \\
selection & nilpotent basis, \emph{chosen} &
  record direction, \emph{derived} (\S\ref{sec:frame}) \\
gauge group & centraliser in $\mathrm{Spin}(9)$ &
  fibration stabiliser (Proposition~\ref{prop:gsm}) \\
multiplets & minimal ideals; $1{+}3{+}3{+}1$ grading &
  canonical orbits; $\mathbb{C}\oplus\mathbb{C}^3$ by class \\
charge & ladder count over three &
  $S^1$ winding through $\mathbb{C}^3$ \\
generations & three ideal prints; $Cl(8)$ copies &
  sedenion $\mathbb{O}_i$ orbit; Fano lines; $K_3$ classes \\
the theory itself & module + $J$ \emph{chosen} + Born posited &
  module + $J$ and Born \emph{derived} (\S\ref{sec:postulates}) \\
\hline
\end{tabular}
\end{center}

Three rows deserve their sentences. The gauge-group row is a composite reading whose ingredients
are separately on record but whose composition we have not found stated: $\mathrm{Spin}(9)$ is
the symmetry group of the octonionic Hopf fibration $S^{15} \to S^8$
\autocite{partonpiccinni2018spin9}, so Krasnov's theorem --- the Standard Model gauge group as the centraliser of a
complex structure inside $\mathrm{Spin}(9)$ --- \emph{is} Proposition~\ref{prop:gsm} in module
coordinates: the symmetries of the top Hopf rung that commute with the selected record direction.
On the algebraic side both inputs are taken as given; the present framework derives the arena
(three qubits, \S\ref{sec:register}) and the selection (the record direction, \S\ref{sec:frame}),
so on this side the characterisation upgrades from a description to a derivation. The selection
row states the general pattern of which that is an instance: what the algebraic programmes
\emph{choose} --- the nilpotent basis whose stabiliser is $SU(3)$, the complex structure, the
preferred unit --- the geometric route \emph{derives}, and this is the single most consequential
thing each side offers the other, in one direction a reason and in the other a calculus. The
multiplet row is Open Problem~\ref{op:ideal-orbit}, already priced above: the exterior-algebra
grading of Furey's ideal, $1{+}3{+}3{+}1$ under $SU(3)$, is the Hopf split
$\mathbb{O} \cong \mathbb{C}\oplus\mathbb{C}^3$ read grade by grade, leptons at the singlet
grades and quarks at the triplet grades exactly as Propositions~\ref{prop:ghz}
and~\ref{prop:w-access} populate them.

The division of labour is then plain. What the module coordinates cannot supply is anything with
a measure in it --- and the last row of the table is this rule at the foundations rung, the
closing paragraph of \S\ref{sec:postulates} stating the reading and guarding its arrow: no $\tfrac29$,
no mass, no probability --- the algebraic programmes have no
Born-rule story at all --- and no compositional statement of confinement, which here is
compactness of the base. What the entanglement coordinates cannot supply without translation is
an associative operator calculus ready for Lagrangians and the parameter-free integer bookkeeping
that makes charge quantisation automatic. The generations row is where the translation is under
active construction: the module-side prints, the sedenion orbit, the Fano lines and the $K_3$
classes are now suspected of being presentations of a single Klein-four-group with its $S_3$ of
automorphisms, and the shadow programme of \S\ref{sec:sedenion-announcement} is that suspicion
made precise.

The charge row has since been converted from correspondence to identity, in the same revision
that tabulated it: Appendix~\ref{app:winding-number} proves that Furey's number operator is the
relative winding of the colour circle's spinor lift --- the engine being the one-line identity
$\mathrm{i}\, c_{x_{k}} c_{y_{k}} = 2N_{k} - 1$, each mode's ladder count \emph{equal} to its
plane's Clifford bilinear --- so that $Q = N/3$ and the fibration's winding are one operator,
the charge thirds are winding integers divided by the three colour planes, and the half-integer
vacuum offset is the spinor double cover doing its usual work
(Lemma~\ref{lem:winding-number}).

\begin{openproblem}[The dictionary's chirality lemma]\label{op:dictionary-lemmas}
One lemma of the dictionary remains at the one-page grade: show that the Filatov
opposite-handedness of \S\ref{sec:generations} is the distinction between left and right
quaternionic multiplication, giving chirality's algebraic address and its geometric proof in one
stroke. A failure of the identification would be equally informative, delimiting where the two
coordinate systems part.
\end{openproblem}

\subsection{What is settled, and what comes next}

What is settled in this section, with the full derivations carried into the appendices of this paper:
GHZ confinement to $\mathbb{C}^3$ (Proposition~\ref{prop:ghz}, Appendix~\ref{app:ghz-proof}); W
accessibility of $\mathbb{C}$ and the $SU(2)$ doublet content of W
(Propositions~\ref{prop:w-access}--\ref{prop:w-doublet}, Appendix~\ref{app:w-proof}); the topological
character of confinement; and $\theta_{\mathrm{QCD}} = 0$. These rest only on the composition-algebra
identity and the local-unitary invariance of single-qubit spectra, with no free parameters. The two
questions posed at the outset are answered: the colour $SU(3)$ is found in the entanglement of the
three-qubit state, in the GHZ class, exactly because $S^7$ is not a group and the gauge content must
be carried structurally; and the GHZ/W split does map onto the quark/lepton division.

What remains is to make the carrier explicit. The dictionary identifies which entanglement class is
which kind of particle, but not yet how a single particle's three gauge charges sit together, how its
electric charge, weak isospin, hypercharge and colour are read off, or how the Filatov handedness
becomes physical chirality. Those are the business of the next sections: the three-qubit register and
the quantum numbers read from it (\S\ref{sec:register}, \S\ref{sec:charges}), the three generations
(\S\ref{sec:generations}), and the masses (\S\ref{sec:masses}). The residual questions that are
properly spacetime-level --- among them the dynamical content of confinement --- are taken up in
Paper~2, alongside the spacetime discussion to which they belong.

%  No adjacency phrasing in the prose: position-portable. All
%  cross-refs live -- sec:fanout-direction DRAFTED 2026-07-06 (formerly the one undrafted
%  section, by design).
%
%  AUDIT RECORD: prop:gsm, prop:termination and prop:f1-dictionary
%  verified against every quoting consumer (statements canonical);
%  honest-debt subsection names both the selection assumption and the
%  mereology debt; the epistemic two-climbs disambiguation and the
%  compositional rung-evidence verified; sec:canonical-orbits live;
%  file pure ASCII. Fixes: op:hopf-frobenius's cross-link with
%  op:sigma-functor made TWO-WAY (the "four ways" line now names the
%  op); both remaining functor conditions ANCHORED --- the quaternionic
%  to the established fibre physics of sec:wf-reconciliation (residue:
%  clutching map vs strong complementarity), the octonionic to the
%  proved canonical-orbit images (umbrella label app:ghz-w-proofs
%  minted at the appendix for the purpose). Statuses unchanged.
%  British English; \autocite; no bullets in prose.
% =====================================================================

\section{Why Hopf Fibrations Belong in Particle Physics}
\label{sec:hopf-justification}

Invoking Hopf fibrations in a discussion of the Standard Model invites an immediate and fair
objection: that the move is decorative, an analogy between two pieces of attractive mathematics
with no load-bearing connection between them. The purpose of this section is to dispose of that
objection before the construction is put to work. The connection is not an analogy. It is
carried by two independent and independently published theorems, and the structures it relies
upon are not selected for convenience but compelled by a pair of classification results. We take
as already established the correspondence between qubit entanglement and the Hopf fibrations
themselves --- the subject of a mature literature
\autocite{mosseri2001hopf,bernevigchen2003geometry} --- and do not reproduce it here. What this
section addresses is the prior question of why that correspondence is entitled to carry physical
content at all.

\subsection{Two theorems, joined end to end}

The bridge from qubit entanglement to the Standard Model is built from two distinct results,
neither of which is original to the present framework and neither of which depends on the other.

The first is geometric. The unit sphere in the Hilbert space of $n$ qubits is the total space of
a Hopf fibration: $S^3 \to S^2$ for one qubit, $S^7 \to S^4$ for two, and $S^{15} \to S^8$ for
three. For two and three qubits the fibration is \emph{entanglement-sensitive} --- separability
is precisely the condition under which the base degenerates to a product of single-qubit Bloch
spheres and the additional dimensions associated with the higher imaginary directions collapse
\autocite{mosseri2001hopf,bernevigchen2003geometry}. This is a statement about the geometry of
state spaces and requires no physical interpretation whatever.

The second is algebraic. Singling out a preferred complex structure on the octonions ---
equivalently, fixing the split $\mathbb{O} \cong \mathbb{C} \oplus \mathbb{C}^3$ --- leaves a
residual symmetry, and that residual symmetry is exactly the gauge group of the Standard Model.
This was established through the exceptional Jordan algebra by Dubois-Violette and Todorov
\autocite{duboisviolette2016exceptional,duboisviolettetodorov2019exceptional,todorovduboisviolette2018deducing}
and given its sharpest spinorial form by Krasnov \autocite{krasnov2021so9}, who showed that the
relevant group is the stabiliser of a complex structure inside $\mathrm{Spin}(9)$. We record the
result in the form in which it enters the present work.

\begin{proposition}[Dubois-Violette--Todorov; Krasnov]\label{prop:gsm} Let $J_R$ be the complex
structure on the space $\mathbb{O}^2$ of $\mathrm{Spin}(9)$ spinors determined by a choice of
unit imaginary octonion, equivalently the split
$\mathbb{O} \cong \mathbb{C} \oplus \mathbb{C}^3$. The subgroup of $\mathrm{Spin}(9)$ commuting
with $J_R$ is \[ G_{\mathrm{SM}} \;=\; \frac{SU(3) \times SU(2) \times U(1)}{\mathbb{Z}_6}, \]
acting on the singled-out $\mathbb{C}$ summand as a colour singlet of hypercharge $-1$ and on
the $\mathbb{C}^3$ summand as a colour triplet of hypercharge $\tfrac{1}{3}$, both as $SU(2)$
doublets --- the quantum numbers of one left-handed generation.
\end{proposition}

The two results are independent, and the join between them is the only place a physical theory
must do work. Each result is a theorem; what neither supplies is a reason why nature should fix
one complex direction within $\mathbb{O}$ rather than treating all directions symmetrically. We
return to this --- the genuine load-bearing assumption --- in \S\ref{sec:hopf-honest-debt}, and
resolve it, in candidate form, in \S\ref{sec:fanout-direction}.

\subsection{The ladder is forced, not chosen}

It might still be objected that selecting the octonions, or the Hopf fibrations, is itself a
modelling choice carrying as much latitude as the choice of a grand-unified group. It is not,
and the reason is a pair of classification theorems that close off the alternatives entirely.

By Hurwitz's theorem the only normed division algebras are $\mathbb{R}$, $\mathbb{C}$,
$\mathbb{H}$ and $\mathbb{O}$ \autocite{hurwitz1898composition,baez2002octonions}. By Adams's
resolution of the Hopf-invariant-one problem --- together with the parallelisability results of
Bott--Milnor and Kervaire --- the only parallelisable spheres are $S^0$, $S^1$, $S^3$ and $S^7$,
so that there exist exactly three nontrivial Hopf fibrations, in bijection with $\mathbb{C}$,
$\mathbb{H}$ and $\mathbb{O}$
\autocite{adams1960hopf,bottmilnor1958parallelizability,kervaire1958nonparallelizability}. These
are not heuristics; they are among the more durable theorems in topology and algebra. Their
force here is that, once a framework is committed to amplitudes valued in a normed division
algebra --- a commitment that in the present construction is made far upstream, by the
self-referential argument of \S\ref{sec:self-reference} and \S\ref{sec:division-algebras}, and
not by fiat at this stage --- the ladder of Hopf fibrations is not selected but compelled, and
it terminates of its own accord.

\begin{proposition}[Termination]\label{prop:termination} The Cayley--Dickson tower of normed
division algebras and the tower of parallelisable spheres both terminate at real dimension
eight. Consequently there are exactly three nontrivial Hopf fibrations, the largest being the
octonionic $S^{15} \to S^8$, and the construction admits no extension to a fourth case.
\end{proposition}

The termination is worth dwelling on, because it is unusual for a model of this kind to predict
its own boundary. The next rung of the Cayley--Dickson sequence, the sedenions, possesses zero
divisors; the Hopf map can no longer be globally defined; there is no fourth fibration. This is
not a limitation imposed from outside but a structural fact, and it carries a physical
consequence that we develop later: there is no room in the construction for a fourth gauge force
or a fifth interaction, so that the discovery of either would count as evidence against the
framework rather than as an invitation to extend it. The endpoint is a theorem, not a
stipulation.

\subsection{Independent arrival as corroboration}

The most persuasive evidence that the destination is not an artefact of the route is that it is
reached by routes sharing none of the same machinery. Szangolies
\autocite{szangolies2025standardmodel} reaches $G_{\mathrm{SM}}$ from the emergent-gravity and
Kaluza--Klein programmes, reading the singled-out complex direction as a dimensional reduction
of a $9+1$-dimensional spacetime; the present framework reaches it from self-reference
(\S\ref{sec:self-reference}); and Furey and Gillard--Gresnigt reach the same gauge group, and
the same three-generation structure, through the minimal left ideals of the Dixon and Clifford
algebras, without invoking Hopf fibrations at all
\autocite{furey2018three,gillardgresnigt2019three}; the relation between their minimal ideals
and the entanglement classes used here is taken up in \S\ref{sec:canonical-orbits}. The
octonionic route to particle physics is in any case an old one, going back to the use of
$\mathbb{O}$ to model the strong force \autocite{gunaydingursey1974quark}.

That programmes with disjoint starting points converge on
$SU(3) \times SU(2) \times U(1)/\mathbb{Z}_6$ and on the same quark--lepton division is the
coincidence Szangolies himself remarks upon: the characteristic structure of the Standard Model,
routinely regarded as arbitrary, turns up independently in unrelated settings, and its appearing
twice would then be doubly odd unless the two settings are in fact related. We read the
convergence the same way --- as evidence that the division-algebraic description is singled out
by the mathematics rather than imposed by the modeller, and therefore that
Proposition~\ref{prop:gsm} is structural rather than contrived.

\subsection{The honest debt}
\label{sec:hopf-honest-debt}

It would be misleading to leave the impression that this route is free of assumptions, so we
state plainly the one assumption that all of these programmes share. Proposition~\ref{prop:gsm}
requires that \emph{some} complex direction in $\mathbb{O}$ be singled out, and the uniqueness
of that choice is not supplied by the theorem. Krasnov parametrises the complex structure $J_R$
by an arbitrary unit imaginary octonion, and the residual group is the same for every such
choice because $G_2 = \mathrm{Aut}(\mathbb{O})$ acts transitively on the unit imaginary
octonions. The theorem therefore guarantees the gauge group \emph{given} that a direction is
fixed, while remaining silent on \emph{why} one is fixed. Szangolies treats this as a free
choice and flags it, in his own conclusions, as an open element of the construction --- one to
which several inequivalent answers are available.

This, and not the appeal to Hopf fibrations, is the genuine load-bearing assumption of the whole
approach. The Hopf fibrations are forced (Proposition~\ref{prop:termination}); the singling-out
is not. The present framework offers a candidate resolution that the prior literature lacks: the
act of measurement, modelled as \textsc{fanout}, selects the direction physically, and the
transitivity of $G_2$ then guarantees that the gauge group is independent of which direction the
measurement happens to select. This argument is given in \S\ref{sec:fanout-direction}; we flag
here, at the appropriate epistemic status, that it is structural rather than fully proved, and
that the residual question --- why a measurement selects \emph{exactly one} imaginary direction
rather than merely \emph{permitting} a selection --- remains open. It is now priced as Open
Problem~\ref{op:fanout-selection}.

A second debt is shared with every programme of this character and is not closed by anything
said here: why the abstract quantum state should admit a tensor-product decomposition into
qubits in the first place. This is the quantum factorisation, or mereology, problem. The
entanglement-sensitivity on which the construction turns is meaningful only once such a
decomposition is fixed, and we treat the decomposition as a necessary input rather than a
derived result. Naming it is preferable to letting it pass unremarked.

\subsection{The termination read epistemically}
\label{sec:hopf-information-bound}

The termination of the tower (Proposition~\ref{prop:termination}) has so far been stated
algebraically: the Cayley--Dickson sequence acquires zero divisors at the sedenions and the Hopf
map can no longer be globally defined. There is a second reading of the same boundary, epistemic
rather than algebraic, which we offer as a motivated conjecture rather than a settled result,
because it connects the tower's endpoint to a bound that arises independently in the foundations
of measurement.

A particle, in this framework, is a reified relation between events, and its internal content is
carried by the octonionic $A,B,C$ register --- the $S^{7}$ sector serving as internal degrees of
freedom in the channel that connects the events (\S\ref{sec:space-of-computations}). The
termination of the tower is then a statement about that content: a single such relation can
support genuine entanglement among at most three parties, because there is no fourth fibration
to carry a fourth independent party within the same irreducible relation. Every particle is, on
this reading, a \emph{completed} octonionic-level relation --- topped out at three parties not
by circumstance but because the tower admits no further rung. The information that can be
encoded in one relation-between-events, and hence extracted by one measurement of it, is
correspondingly bounded: it is the information of a completed three-party register, and no
measurement of a single particle can yield more.

This is exactly the kind of bound that the epistemic-restriction programmes reach from the other
direction. Spekkens' toy theory is built on a knowledge-balance principle that caps the
information obtainable about a system \autocite{spekkens2007toy}, and Szangolies derives such an
epistemic horizon not by positing it but from the same self-referential non-computability on
which the present tower rests \autocite{szangolies_j_2018}. That the two accounts --- the
algebraic termination of the Hopf tower and the epistemic bound on measurement information ---
should coincide is, we suggest, not a numerical accident but a shared root: both are
consequences of the self-referential non-computability of \S\ref{sec:self-reference}, wearing an
algebraic face in the one case and an epistemic face in the other. The Hurwitz--Adams ceiling
and the knowledge-balance ceiling would then be one limit seen twice, and Szangolies' programme,
which reaches $G_{\mathrm{SM}}$ by an independent route (\S\ref{sec:hopf-honest-debt}) and the
epistemic horizon from the same non-computable source, shares with this framework not merely a
conclusion but its origin.

The reading must be kept distinct from a different climb, with which it would otherwise be
conflated. The Wigner's-friend regress of \S\ref{sec:wigners-friend} ascends without end: each
observer may be measured by a further observer, each carrying a fresh gauge-private circle, and
that chain has no top --- it is the bump in the carpet, unbounded in length. The termination
bounds something else: not the length of the chain of observers, but the internal structure of
each single relation, each particle, that the chain is built from. The regress is unbounded
precisely because each of its links is a bounded, completed three-party object;
termination-of-each-link and non-termination-of-the-chain are not in tension but complementary,
the former individuating the links of the latter. What the bound caps is the information in one
particle; what the regress leaves uncapped is the number of particles and observers that can be
chained. The precise identification of ``the information of a completed three-party register''
with the formal measures of the epistemic-restriction literature is the technical step that
would turn this conjecture into a theorem, and we flag it as owed; the structural claim --- that
the tower's endpoint and the measurement bound share the non-computable root --- is what we
advance here.

\subsection{The termination read compositionally}
\label{sec:hopf-compositional}

There is a third termination theorem, proved in a community that has so far had no contact with
the two joined above. Coecke and Kissinger showed that the GHZ and W classes --- realised as the
commutative special and the anti-special Frobenius algebras of categorical quantum mechanics ---
are \emph{compositionally complete}: every $N$-qubit entanglement class under SLOCC is generated
by composing these tripartite primitives with bipartite entanglement and product states
\autocite{coeckekissinger2010compositional}. Above three parties nothing new appears; the
entanglement alphabet closes at exactly the point where the fibration tower terminates. The
mutual silence of the two literatures is itself striking: Coecke and Kissinger make no reference
to Hopf fibrations or division algebras, and the fibration literature on entanglement
\autocite{mosseri2001hopf,bernevigchen2003geometry} nowhere cites compositional completeness.
Two classifications, in different mathematics, draw the same line in the same place, unaware of
one another. That silence has, in the months this paper was in draft, just begun to break from
the categorical side, with a projective, base-level reading of the same calculus
\autocite{burtonanwar2026meromorphic}; we engage it below.

The framework's conjecture is that this is one termination wearing a third face: the alphabet
has the letters it has --- Bell, GHZ, W --- because the parallelisable spheres have the
dimensions they have. There is partial evidence at every rung. At one qubit, the commutative
special Frobenius algebras on $\mathbb{C}^{2}$ are classified by orthonormal bases --- points of
$S^{2}$, the base of the first fibration --- so that the copied basis plays the role of a
section of the bundle and the spider theorem expresses its idempotence. At two qubits, the
$SU(2)$ fibre of $S^{7} \to S^{4}$ carries exactly the entanglement that survives local
operations, and L\'{e}vay's result that the bundle's natural connection is the $SU(2)$
Yang--Mills instanton \autocite{levay2004geometry} stands on the geometric side where the
interaction of two Frobenius structures stands on the categorical one. At three qubits, the
split of D\"{u}r, Vidal and Cirac \autocite{durvidalcirac2000three} mirrors the algebra exactly:
GHZ is genuinely tripartite, demanding the full non-associative octonion, while W is
pairwise-robust and housed in the associative subalgebra $\mathbb{H} \subset \mathbb{O}$ --- the
special/anti-special Frobenius distinction in geometric dress.

The first of the four conditions the correspondence requires can be discharged at once, and we
do so here, because a problem whose first step has been walked is a different kind of invitation
from one merely posed.

\begin{proposition}[The one-qubit dictionary]\label{prop:f1-dictionary} The $\dagger$-special
commutative Frobenius algebras on $\mathbb{C}^{2}$ correspond bijectively to
\emph{lifted antipodal skeletons} of the complex Hopf fibration $\pi \colon S^{3} \to S^{2}$:
unordered antipodal pairs $P \subset S^{2}$ together with a section of $\pi$ over $P$. The
copyable states of a structure are exactly the section's image, so the lifts are recorded
faithfully --- the projectively trivial centre of $U(2)$ acts non-trivially on structures ---
and the correspondence is $U(2)$-equivariant, the ordered structures forming a single free and
transitive $U(2)$-orbit. The channel a structure induces on density matrices is dephasing in its
skeleton and is independent of the section: the fibre freedom, a torsor for
$\mathrm{Map}(P, U(1))$, is exactly the data absent from the classical record a structure
induces.
\end{proposition}

\begin{proof}
By the classification of Coecke, Pavlovic and Vicary, such structures correspond to orthonormal
bases, recovered as the sets of copyable states \autocite{coeckepavlovicvicary2013orthogonal}.
An orthonormal basis of $\mathbb{C}^{2}$ is two orthogonal rays --- an antipodal pair on the
Bloch sphere --- together with a unit representative of each, which is precisely a section of
$\pi$ over the pair; and $\delta(e^{i\theta}x) = e^{i\theta}\, x \otimes x$ against
$(e^{i\theta}x) \otimes (e^{i\theta}x) = e^{2i\theta}\, x \otimes x$ pins the phases.
Equivariance, the action of the centre, the orbit structure and the induced channel are direct
computations; an executable verification of every finite check accompanies the code release.
\end{proof}

Two remarks fix the proposition's scope. It corrects a natural first reading: a classical
structure does not project to the base but selects lifts into the total space --- classicality
is section data, and what the copy operation broadcasts is a lift, not a ray. And it claims
nothing topological about cloning: the finiteness of the copyable set is linear, structural
no-cloning residing in non-naturality \autocite{abramskyheunen2012hstar}, so the first place the
bundle's global topology can bite is the quaternionic rung, where the degree-one clutching map
stands ready to be compared with the strong-complementarity law governing an interacting pair of
structures. The dictionary is, finally, the complement of the projective treatment cited above:
Burton and Anwar pass functorially to the base and deliberately discard the fibre
\autocite{burtonanwar2026meromorphic}, and the classical structures turn out to live in exactly
the data their projectivisation forgets --- the two readings compose, and the resulting triangle
is a consistency condition available to constrain the remaining conditions of the
correspondence.

\begin{openproblem}[The Hopf-fibration--Frobenius correspondence]\label{op:hopf-frobenius} Let
$\mathbf{Hopf}$ denote the category of the three fibrations
$S^{2^{k}-1} \hookrightarrow S^{2^{k+1}-1} \to S^{2^{k}}$, $k = 1, 2, 3$, with bundle maps as
morphisms, and let $\mathbf{Frob}$ denote the Frobenius algebras on $\mathbb{C}^{2}$ in
$\mathbf{FdHilb}$, with the commutative-special (GHZ) and anti-special (W) classes
distinguished. Construct a functor $F \colon \mathbf{Hopf} \to \mathbf{Frob}$ --- or prove that
none exists --- under which the complex fibration yields the single-qubit classical structures
--- Proposition~\ref{prop:f1-dictionary} identifies these as the lifted antipodal skeletons,
section data over two base points rather than the points alone; the quaternionic fibration
yields the interaction of two such structures, with the $SU(2)$ fibre carrying the entanglement
that survives the classical quotient; the octonionic fibration yields the GHZ/W pair, GHZ on the
full non-associative fibre and W on the associative $\mathbb{H} \subset \mathbb{O}$; and the
composition by which Coecke--Kissinger generate the $N$-qubit classes corresponds to a
well-defined geometric operation on the bundles. If such an $F$ exists, compositional
completeness follows from Adams's theorem as a topological corollary --- there is no fourth
fibration to supply a fourth primitive --- and if $F$ is faithful enough that the implication
also runs back, the termination of the tower and the completeness of the alphabet are equivalent
expressions of a single classification. The first condition is discharged by
Proposition~\ref{prop:f1-dictionary}; what remains open is the quaternionic and octonionic
conditions and, above all, the correspondence of compositions. Both open conditions are already
constrained: the quaternionic fibre's physical content --- the recoverable relative phase and
its purity ledger --- is established (\S\ref{sec:wf-reconciliation}), leaving the
clutching-map-to-strong-complementarity comparison as the structural residue; and any candidate
functor must reproduce the proved canonical-orbit images, GHZ to the colour triplet and W to the
singlet (\S\ref{sec:canonical-orbits}, Appendix~\ref{app:ghz-w-proofs}). A proof that no such
functor exists, or a correspondence at the level of the classifications without functoriality on
compositions, would each also resolve the problem: the first would make the agreement of two
independent classifications a fact demanding its own explanation; the second would establish the
weaker reading.
\end{openproblem}

A resolution in the affirmative would do three things. It would explain compositional
completeness rather than merely prove it --- a corollary of the division-algebra structure of
quantum mechanics rather than an independent algebraic fact. It would give the diagrammatic
language of this paper's own interaction circuits a topological warrant: the GHZ/W spider
calculus would be the unique language the topology permits, not a convenient notation. And it
would open a two-way bridge, translating problems on compositions of Frobenius structures into
operations on bundles and back. The problem sits at the intersection of categorical quantum
mechanics, the quantum-information geometry of the fibrations, and the division-algebra
approaches to particle physics, and it is posed here so that any of the three communities can
take it up in its own language. It is perhaps worth noticing that both classifications already
answer to the same name --- the fibrations are Heinz Hopf's, and so, through the bialgebra
structure of the diagrammatic calculus, are the Hopf algebras --- a coincidence of eponymy that
the conjecture proposes to promote to a theorem. With the epistemic reading above, an
affirmative resolution would leave the tower's endpoint as one limit seen four ways: algebraic,
geometric, compositional, epistemic --- the last being the fixed-point classification of Open
Problem~\ref{op:sigma-functor}, with which this problem is cross-linked, each constraining the
other.

\subsection{What this sets up}

The entanglement-sensitivity of the Hopf maps is not a side feature but the hinge of the
physical reading that follows. Because separability is exactly the collapse of the base to a
product of Bloch spheres, the structure of the entanglement --- the SLOCC class of a multi-qubit
state, and within it the canonical maximally entangled orbit --- is carried by precisely the
degrees of freedom on which $G_{\mathrm{SM}}$ acts in Proposition~\ref{prop:gsm}. This is what
makes it possible, in \S\ref{sec:entanglement-classes}, to read the two genuinely tripartite
entanglement classes onto the $\mathbb{C} \oplus \mathbb{C}^3$ split through their canonical
representatives: the GHZ orbit confined to the colour-triplet $\mathbb{C}^3$ sector, the W orbit
reaching the $\mathbb{C}$ singlet, recovering the quark--lepton division not as an assignment
but as a consequence of the entanglement structure. The scope of those statements --- canonical
orbits rather than whole classes, and the identification of the canonical orbits with the
physical particles --- is made precise there. The Hopf fibrations are the bridge; the
entanglement classes, treated next, are the traffic that crosses it.

%
%  SOURCES (re-read 2026-07-04): lawvere_inside_outside_fixed_point_
%  theorem_2026-05-04T1806.md (the spine: the two-principles reduction,
%  the theorem, the sigma candidate, the open items -> op minted per
%  D9); research_summary_neg_prob.md (negativity route; lamp (1+i)/2);
%  dynamics_companion_fanout_unitarity (legs 1a/1c; anti-events; the
%  Brukner-Zeilinger rebuttal per D6); Pure_states_and_mixed_states
%  (local tomography derived; QBism contrast per D7).
%
%  DECISIONS ENCODED (JS 2026-07-04): D1 title as below; D2 Wheeler1990
%  cited here as the title anchor; D3 lineage Popper -> Breuer ->
%  Svozil -> Szangolies (Dalla Chiara omitted); D4 Frauchiger_2018
%  cited here, dissolved in sec:wigners-friend; D5 Kirkwood-Dirac
%  heritage via the Arvidsson-Shukur 2024 review as one-stop; D6
%  Brukner-Zeilinger + the framework rebuttal IN ("Shannon is
%  post-FANOUT"); D7 QBism contrast IN; D8 Aaronson private
%  communication OMITTED (published island paper only); D9 one op
%  minted (op:sigma-functor), cross-linked to op:hopf-frobenius.
%
%  CITATIONS -- existing keys: Wheeler1990, small_2006, Godel1931,
%  Lawvere_1969, Yanofsky2003, CoeckeKissinger2017, Rovelli1996a,
%  Popper1950a, Popper1950b, Breuer1995, svozil_k_2018,
%  szangolies_j_2018, Frauchiger_2018, dirac1942, Feynman_r_82,
%  Feynman1987, Youssef1994 [SWAP-FLAG: intended paper is "Quantum
%  mechanics as complex probability theory", Mod. Phys. Lett. A 9
%  (1994) 2571 -- if that lives under youssef1994quantum, swap],
%  caticha_a_2006, goyal_knuth_skilling_2010, Aaronson2004h, Hardy2001a,
%  muller_p_2011, Thomson_1954, benacerraf1962tasks, cantor1891diagonal,
%  Turing_1937.
%  NEW keys (RIS in self_reference_refs_2026-07-04.ris, all verified
%  this session): tikochinsky1988feynman, renou2021real, chen2022ruling,
%  wu2022refutation, hoffreumonwoods2026real, arvidssonshukur2024kirkwood,
%  bruknerzeilinger2001shannon, fuchsschack2013qbism,
%  shannon1948mathematical.
%
%  CROSS-REFS: forward (fine under two-pass LaTeX): sec:continuity,
%  sec:division-algebras, sec:hopf-justification, prop:f1-dictionary,
%  sec:wigners-friend, op:limit-realisability, op:hopf-frobenius.
%  Still dangling after this section: sec:fanout-boundary (and its
%  sibling label sec:fanout-direction) -- both DRAFTED 2026-07-06.
%
%  AUDITED 2026-07-06 (section-by-section final pass, first stop):
%  proof of prop:inside-outside re-verified by hand; attribution chain
%  clean (Burgin/Khrennikov correctly absent); all header-declared
%  citations confirmed used; lamp arithmetic and antecedents verified;
%  FANOUT smallcaps consistent; no pre-definition use of "frame".
%  Fixes: stray Unicode em-dash (QBism paragraph, the file's only
%  non-ASCII character) -> house ---; trace-blind two-way stitch added
%  to the Born subsection ("one privacy, two duties");
%  op:sigma-functor gains its first standing constraint (F1
%  inheritance via the op:hopf-frobenius discharge) and its previously
%  one-directional local-tomography cargo made explicit.
% =====================================================================

\section{Self-Reference: One Principle}
\label{sec:self-reference-2}

Wheeler asked for it from bit: for physics derived from information, ``every it --- every particle,
every field of force'' obtained from binary answers to yes-or-no questions \autocite{Wheeler1990}.
This paper's route to that programme runs through G\"{o}del: the single principle from which
everything else here is derived, deferred, or priced is \emph{self-reference} --- the fact that any
complete physical description, made from inside the world it describes, must describe the describer.
This section states the principle, proves that the framework's second historical commitment follows
from it as a theorem rather than standing beside it as an axiom, and then extracts the price of the
principle in four instalments: probabilities whose parts can cancel, amplitudes valued in
$\mathbb{C}$, the Born rule, and the local tomography of composite systems. What the section cannot
yet derive is stated as precisely as what it can, in one numbered problem and in the ledger at the
end.

\subsection{Events, not objects}

The framework's ontology is relational: the fundamental entities are events and the relations
between them, and objects are the bookkeeping of relational invariants --- what remains stable
across many relations. The kinship with relational quantum mechanics \autocite{Rovelli1996a} is
genuine and acknowledged: facts are facts relative to the systems that register them. The
difference, developed over the rest of this paper, is that the present framework grounds the
relational structure in a single logical principle and then carries it, through the division
algebras and the Hopf fibrations, all the way to the gauge group and the particle spectrum ---
relations are not only the ontology but the derivation engine. On this reading a measurement is not
the revelation of a pre-existing property; it is the creation of a relation, and the later sections
of this paper read that creation geometrically (\S\ref{sec:hopf-justification}) and even
gravitationally. Here we need only the logical core: whatever describes the world is an event in
the world it describes.

\subsection{One principle, and a theorem}

The framework historically carried two foundational commitments: self-reference, and the
\emph{Inside/Outside Equivalence Principle} --- that the laws of physics must be identical whether
formulated by an embedded participant or by an external observer. Two principles are one too many,
and this subsection reduces them.

The reduction corrects the framework's own origin, and we make the correction in print. The 2005
form of the argument \autocite{small_2006} ran a G\"{o}delian chain: a self-referential process,
being complete, must by the incompleteness theorem \autocite{Godel1931} be inconsistent; an
inconsistent system can negate itself; and the negation of self-reference is the production of
causal sequence. The chain captures the right intuition and is retained here as motivation, but it
is not sound as a proof: classical inconsistency yields \emph{ex falso quodlibet} --- every
proposition provable, every causal sequence an output, no constraint at all. The load-bearing
argument is different, and better: a fixed-point theorem.

Let $\mathcal{C}$ be the category of physical descriptions and admissible processes --- concretely,
the dagger-compact-closed setting with classical structures in which categorical quantum mechanics,
and the whole of this paper, works \autocite{CoeckeKissinger2017}. Lawvere showed that in a
sufficiently rich such setting, any weakly point-surjective self-representation forces every
endomap to have a fixed point, and Yanofsky showed that one diagonal argument of this shape
underlies Cantor's, G\"{o}del's, Turing's and Tarski's theorems alike
\autocite{Lawvere_1969,Yanofsky2003,cantor1891diagonal}. The framework's single commitment is that
$\mathcal{C}$ contains a \emph{self-representing object} --- an object capable of encoding the
processes on itself. That is what ``self-reference is foundational'' means, stated categorically,
and it is the surjectivity hypothesis of the diagonal lemma. Now let
$\sigma \colon \mathcal{C} \to \mathcal{C}$ be the \emph{inside/outside involution}: the operation
exchanging the two complementary presentations of any description --- the inside presentation, as
participant-indexed total space in which the observer is part of the system, and the outside
presentation, as a base-times-fibre decomposition over a classical parameter space. That
$\sigma^{2} = \mathrm{id}$ is not an axiom but a coherence requirement: the outside view of the
inside view of a thing is the thing. The candidate concrete realisation of $\sigma$ is the
Hopf-duality swap --- total space against base-times-fibre for the three fibrations --- which gives
the involution computable content and ties this subsection to \S\ref{sec:hopf-justification}.

\begin{proposition}[Inside/outside fixed point]\label{prop:inside-outside}
Let $\mathcal{C}$ satisfy the richness condition above and contain a self-representing object, and
let $\sigma$ be an involution on $\mathcal{C}$. Then the fixed-point locus
$\mathrm{Fix}(\sigma)$ is non-empty: there exist descriptions invariant under the exchange of
inside and outside perspective.
\end{proposition}

\begin{proof}
Lawvere's construction, applied to $\sigma$ in place of negation: given a weakly point-surjective
$\varphi \colon A \to Y^{A}$, form $g(a) = \sigma\!\left(\varphi(a)(a)\right)$; point-surjectivity
supplies $a^{*}$ with $\varphi(a^{*}) = g$, and then $\varphi(a^{*})(a^{*})$ is a fixed point of
$\sigma$ \autocite{Lawvere_1969,Yanofsky2003}. Where the same construction applied to Boolean
negation yields paradox --- negation has no fixed point on a two-element type, so the conclusion
explodes the premises --- the inside/outside involution is an honest involution on a rich type, and
the conclusion is a constructive fixed point rather than a contradiction.
\end{proof}

\textbf{Definition.} \emph{Physical law is the structure carried by} $\mathrm{Fix}(\sigma)$.

The Inside/Outside Equivalence Principle is thereby demoted from axiom to theorem: physical law is
invariant under the exchange of perspectives \emph{because it is the fixed-point set of that
exchange}, and the fixed-point set is forced to be non-empty by self-reference alone. One
commitment remains foundational; everything the second commitment asserted is now derived. The
reader who accepts categorical quantum mechanics as an existing programme has already accepted the
machinery; the framework adds one premise and collects the rest. And the fixed-point locus is not
empty of structure: it inherits the Frobenius-algebra apparatus of $\mathcal{C}$, so that the
Bell, GHZ and W primitives --- and with them the particle classification of the later sections ---
live inside $\mathrm{Fix}(\sigma)$ from the start. What is proved above is non-emptiness given the
setting; what is assumed is the setting itself and the existence of $\sigma$ as a genuine
involution; what is open is the rigour of both, and it is priced now rather than hidden:

\begin{openproblem}[The involution made functorial]\label{op:sigma-functor}
Fix the category $\mathcal{C}$ rigorously --- the dagger-compact-closed setting with classical
structures is the candidate --- and define $\sigma$ on morphisms, not only on objects, the natural
route being the induced fibre--base exchange on spider diagrams; establish that $\sigma$ is an (at
worst pseudo-)involutive endofunctor; and prove or refute the identification of
$\mathrm{Fix}(\sigma)$ with the three Hopf cross-sections, under which Adams's theorem would become
a corollary of the fixed-point classification. This problem is cross-linked with Open
Problem~\ref{op:hopf-frobenius}: under an affirmative resolution of both, the parallelisable
terminus, the fixed-point classification, and compositional completeness are three statements of
one fact, and progress on either problem constrains the other. One constraint is already in place:
at the complex rung, Proposition~\ref{prop:f1-dictionary} supplies the dictionary between
classical structures and the fibration's lifted skeletons --- the clause of
Open Problem~\ref{op:hopf-frobenius} it discharges --- and any resolution here must reproduce it.
The problem also carries a piece of cargo recorded later in this section: the formalisation, within
the fixed category, of the perspective-accounting that derives local tomography.
\end{openproblem}

\subsection{The lineage}

The obstruction this section builds on is old, and the framework claims its completion, not its
discovery. Popper argued in 1950 that a predictor embedded in the world has irreducibly inferior
knowledge of itself, because it cannot impart information about itself without, by that very act,
rendering the information out of date \autocite{Popper1950a,Popper1950b} --- the passage is quoted
where the framework makes the obstruction structural (\S\ref{sec:wigners-friend}). Breuer proved
the impossibility of accurate state self-measurement within quantum theory \autocite{Breuer1995};
Svozil's monograph surveys the whole territory where physics meets self-reference, undecidability
and Chaitin-style algorithmic limits \autocite{svozil_k_2018}; Szangolies sharpened the modern form
as epistemic horizons \autocite{szangolies_j_2018}; and the Frauchiger--Renner theorem is the
contemporary crisis-form --- nested observers deriving contradictory records when each treats the
others' quantum descriptions as universally valid \autocite{Frauchiger_2018} --- which
\S\ref{sec:wigners-friend} dissolves relationally rather than paradoxically. This framework's
addition to the lineage is twofold. Structurally, Proposition~\ref{prop:f1-dictionary} converts
Popper's temporal obsolescence into algebraic cancellation: the datum a copier cannot put into its
own record is exactly its fibre commitment, which does not arrive stale --- it cancels. And
foundationally, the present section runs the lineage in reverse: where the tradition treats
self-reference as a limit on knowledge \emph{within} physics, the framework takes it as the
generator of physics, and derives the quantum formalism from it.

\subsection{The price is negative probability}\label{sec:complex-amplitudes-2}

What does it cost a probabilistic description to contain its own describer? The first instalment is
the sign. In the regime before any record exists, every event has an anti-event --- the description
must be able to undo what it has provisionally done, or the self-description could never be
revised into consistency --- and a probability calculus closed under undoing cannot be confined to
$[0,1]$: it is a signed calculus, in which the parts of a probability can oppose and cancel. The
framework did not discover negativity as bookkeeping; Dirac admitted negative probabilities as
formal intermediaries \autocite{dirac1942}, Feynman defended them as consistent whenever the final,
recorded probabilities stay positive \autocite{Feynman_r_82,Feynman1987}, and the Kirkwood--Dirac
distribution has carried complex quasiprobability from 1933 to the present, as a recent
comprehensive review documents \autocite{arvidssonshukur2024kirkwood}. What the framework adds is
the forcing: negativity is not a permissible convenience but the first price of self-referential
closure, and it is charged exactly and only in the pre-record regime --- which is why every
negativity in the tradition lives in intermediate, unrecorded quantities and never in a final
record.

The second instalment is complexity. A signed information measure must extend Shannon's
$S = -\ln p$ off the positive axis, and the analytic continuation of the logarithm carries $p$ into
the complex plane: amplitudes, with the phase as the continued information. Here too the extension
has a prior tradition, and the attributions are owed precisely. Tikochinsky derived Feynman's rules
for complex amplitudes from consistency of composition \autocite{tikochinsky1988feynman}; Youssef
reformulated quantum mechanics as complex probability theory outright \autocite{Youssef1994};
Caticha carried objective amplitudes to Bayesian probabilities \autocite{caticha_a_2006}; and
Goyal, Knuth and Skilling derived Feynman's rules from a pair-postulate without assuming
analyticity \autocite{goyal_knuth_skilling_2010}. The extension is prior art. The framework's
claim is the compulsion: the complex extension is not a modelling choice among exotic probability
theories but the unique continuation forced by self-referential closure --- the same forcing that
Aaronson glimpsed when he argued quantum mechanics is what negative probability inevitably becomes
\autocite{Aaronson2004h}.

Why does the continuation stop at $\mathbb{C}$? Downward, the real field is not available: real
quantum theory is experimentally distinguishable from complex quantum theory in network scenarios
and has been disfavoured by experiment
\autocite{renou2021real,chen2022ruling,wu2022refutation}; the recent objection that the exclusion
rests on locally tomographic composition \autocite{hoffreumonwoods2026real} is, for this framework,
not an escape but a confirmation of where the weight sits --- local tomography is \emph{derived}
below, not assumed, so the framework closes from its own resources exactly the door the objection
holds open. Upward, the quaternions and octonions fail as a \emph{base} for the information
continuation: the continued measure is a logarithm, and beyond $\mathbb{C}$ the logarithm of a
product is no longer the sum of logarithms --- non-commutativity severs
$\log(ab) = \log a + \log b$ --- so the additivity that makes it an information measure is lost.
Aaronson's parameter-counting and signalling arguments against a quaternionic base point the same
way \autocite{Aaronson2004h}. The higher algebras are not banished; they return as the
\emph{internal} algebras of the register, which is the business of \S\ref{sec:division-algebras}.
The base is $\mathbb{C}$, for informational reasons, and stays $\mathbb{C}$.

\subsection{The Born rule from fibre privacy}\label{sec:born-rule}

The third instalment is the quadratic rule. An amplitude carries one datum more than any record of
it can: its overall phase. This was asserted throughout the framework's development as the
observer's private, self-referential ignorance; it is now a proved statement, in exactly the form
needed. By Proposition~\ref{prop:f1-dictionary}, a classical structure --- the copy operation any
record requires --- is committed to a specific lift of its skeleton, and that commitment is absent
from every classical record the structure induces: the copier is fibre-committed, the record is
fibre-blind, and the fibre freedom is gauge for every observer. (The computation extends verbatim
from one qubit to the full register, where Proposition~\ref{prop:trace-blind} makes it the ground
of hypercharge's tracelessness --- one privacy, two duties.) Popper's obsolescence, made
cancellation. It follows that anything recordable is invariant under the fibre --- observables and
their statistics must be $U(1)$-invariant functions of the amplitude. The lowest-degree polynomial invariant
of a complex amplitude is the quadratic $|\psi|^{2}$ --- the non-polynomial invariant $|\psi|$,
of degree one, is excluded by the analyticity declared at \S\ref{sec:complex-amplitudes} --- and
normalisation fixes the coefficient:
probability is the minimal record-compatible function of the amplitude,
\[
p \;=\; |\psi|^{2},
\]
the Born rule as the cheapest invariant the fibre permits. The same square is reached from the
cost side in \S\ref{sec:born-cost-ledger}, where a record is the closed loop of a history and its
return leg and the exponent is the number of legs; neither route uses the lattice of projections
that Gleason's theorem needs, and both hold for a single qubit. The derivation is the framework's
own;
its premise --- that the fibre datum is structurally excluded from records --- is
Proposition~\ref{prop:f1-dictionary}, and its geometric reading --- probability as base-point data
under the Hopf projection --- is developed in \S\ref{sec:hopf-justification}.

One standing objection to information-theoretic foundations is answered here, because this
framework needs the answer. Brukner and Zeilinger argued that Shannon information is conceptually
inadequate for quantum measurements: it is basis-dependent, and its operational meaning presupposes
ignorance of a definite but unknown fact, which the quantum case does not supply
\autocite{bruknerzeilinger2001shannon}. The framework's reply is structural, and it is a
demarcation rather than a rebuttal. Shannon's precondition fails exactly where they say it fails
--- in the pre-record regime, where there is no definite unknown to be ignorant of --- and there
the operative measure is the quadratic invariant just derived. Shannon's axioms hold exactly where
the precondition holds --- on records, after \textsc{fanout}, where outcomes are definite and
ignorance is classical --- and there the uniqueness theorem forces $-\log p$
\autocite{shannon1948mathematical}. The basis-dependence dissolves with the same stroke: a basis
\emph{is} selected precisely at \textsc{fanout}, so the basis-dependent measure applies only where
a basis exists. Shannon is not inadequate; Shannon is post-\textsc{fanout}. The framework thereby
keeps both measures, each on its own side of the boundary, and uses the additivity of the
post-record measure --- the very property the previous subsection used to pin the base field ---
with a clear conscience.

\subsection{Thomson's lamp, resolved}

The section owes \S\ref{sec:continuity} a debt incurred there: the lamp's amplitude-side limit.
Thomson's lamp --- toggled at times $1 - 2^{-n}$, its state at $t = 1$ demanded --- has no
classical answer: the task description fixes every finite stage and thereby fixes no limit at all
\autocite{Thomson_1954,benacerraf1962tasks}, which \S\ref{sec:continuity} read as the proof that
the record algebra is not limit-closed. On the amplitude side the limit exists, and the calculus of
this section computes it. The toggle exchanges the two amplitudes; the self-consistent state of the
completed task is the fixed point of the exchange in the complex calculus, and solving it under the
signed continuation gives
\[
P(\mathrm{on}) = \tfrac{1 + i}{2}, \qquad P(\mathrm{off}) = \tfrac{1 - i}{2},
\]
a superposed fixed point whose imaginary part carries the unrecordable datum --- the phase of the
endless alternation, when the lamp \emph{would have} stopped. Two checks certify that the
assignment is calculus rather than decoration. The pair is normalised,
$P(\mathrm{on}) + P(\mathrm{off}) = 1$; and the Born rule of the previous subsection, applied to
the fixed point, returns $|P(\mathrm{on})|^{2} = \tfrac{1}{2}$ --- exactly the symmetric classical
ignorance that is the only sane record-level answer, recovered rather than imposed. The lamp,
which entered the framework at its very beginning as the place where classical computation meets
its boundary, thus exits it as a worked example: undetermined as a record, determinate as an
amplitude, the two sides split precisely along \textsc{fanout} --- and the diagonal construction it
physically instantiates \autocite{cantor1891diagonal,Turing_1937} is the same one that powers
Proposition~\ref{prop:inside-outside}. The status flag is owed here as everywhere: the
undetermination is theorem-level; the specific fixed-point value is the framework's complex
calculus, structural, offered with its two consistency checks shown.

\subsection{Local tomography, derived}\label{sec:local-tomography}

The fourth instalment concerns composites, and it discharges two promises at once. The
Inside/Outside Equivalence Principle --- now a theorem --- requires that an external observer's
characterisation of a composite system be recoverable by the embedded participants, and the
embedded participants have access only to local measurements and their correlations. Equivalence
of the two perspectives therefore entails that the state of a composite is determined by the
statistics of local measurements: \emph{local tomography}, obtained as a consequence rather than
imposed as an axiom. This converts a postulate of the reconstruction programmes
\autocite{Hardy2001a,mullerp2011} into a theorem-shaped consequence of the one principle ---
discharging the caveat recorded in \S\ref{sec:continuity} --- and it closes, from the framework's
own resources, the escape route in the real-number debate noted above: the assumption the
objection of \autocite{hoffreumonwoods2026real} contests is here not an assumption. The status is
flagged plainly: the derivation is structural, one paragraph of perspective-accounting, and its
formalisation within the categorical setting belongs to Open Problem~\ref{op:sigma-functor}.

\subsection{Contrast: QBism, and the ledger}

The nearest neighbour in spirit deserves a paragraph of honest contrast. QBism
\autocite{fuchsschack2013qbism} holds, as this framework does, that quantum probabilities are an
agent's --- that the theory is a calculus of expectation, not a catalogue of properties. What QBism
does not say is why the agent's expectations must be complex-valued and quadratically squared; the
Born rule enters as a normative addition to Bayesian coherence, and the complex amplitude as an
unexplained representation. This section is, in effect, the missing why: the agent's ignorance is
self-referential --- the agent is inside the description --- and self-referential ignorance is
complex, with the quadratic rule as the cheapest record-compatible invariant of it. Where QBism
stops at the agent, the framework derives the agent's calculus.

The ledger, then. Delivered in this section: the reduction of two principles to one
(Proposition~\ref{prop:inside-outside}); the pre-record regime with its anti-events, handed to
\S\ref{sec:continuity} as the five-face regime whose continuity is priced there
(Open Problem~\ref{op:limit-realisability}); the complex base, forced downward by experiment and
upward by the additivity of information; the Born rule from fibre privacy; the lamp's fixed point;
and local tomography. Owed, and stated where they are owed: the functorial rigour of the involution
and the fixed-point classification (Open Problem~\ref{op:sigma-functor}); the remaining premise of continuity (Open Problem~\ref{op:limit-realisability}, in
\S\ref{sec:continuity}); and the
rung-by-rung compulsion of the algebra ladder, posed as Open
Problem~\ref{op:ladder-compulsion} in \S\ref{sec:division-algebras}. The
hand-on is the paper's spine: the $\mathbb{C}$ this section derives natively and the $\mathbb{C}$
that \S\ref{sec:continuity} underwrites through Hardy's theorem are the same field, the floor of
the ladder that \S\ref{sec:division-algebras} climbs and \S\ref{sec:hopf-justification}
geometrises. One principle; everything else is theorem, derivation, or a debt with a stated
invoice.

%      % EDITED 2026-10-29 scripts/patch_continuity_tower_premise_2026-10-29T1500.py
% EDITED 2026-09-15 scripts/patch_value_continuum_declared_cencov_2026-09-15T1540.py
% EDITED 2026-09-15 scripts/patch_continuity_connected_state_space_2026-09-15T1430.py
% =====================================================================
%  SECTION: Continuity Is Not an Axiom
%  Drafted 2026-07-04 from continuity_limit_realisability_2026-07-04.md
%  (now in project documents) and the May record
%  (dynamics_companion_fanout_unitarity_2026-05-20T1430.md).
%
%  PLACEMENT in main.tex: a standalone section between self-reference
%  and division-algebras --
%      % EDITED 2026-10-29 scripts/patch_continuity_tower_premise_2026-10-29T1500.py
% EDITED 2026-09-20 scripts/patch_born_cost_ledger_v2_2026-09-20T0900.py
% =====================================================================
%  SECTION: Self-Reference: One Principle
%  Drafted 2026-07-04. The paper's foundational section; resolves the
%  largest dangling label in the manuscript (sec:self-reference).
%
%  PLACEMENT in main.tex: FIRST of the foundations chain --
%      % EDITED 2026-10-29 scripts/patch_continuity_tower_premise_2026-10-29T1500.py
% EDITED 2026-09-20 scripts/patch_born_cost_ledger_v2_2026-09-20T0900.py
% =====================================================================
%  SECTION: Self-Reference: One Principle
%  Drafted 2026-07-04. The paper's foundational section; resolves the
%  largest dangling label in the manuscript (sec:self-reference).
%
%  PLACEMENT in main.tex: FIRST of the foundations chain --
%      \input{self_reference_section_2026-07-04}
%      \input{continuity_section_2026-07-04}
%      \input{division_algebras_section_2026-07-03}
%      \input{hopf_justification_section_2026-06-08}
%
%  SOURCES (re-read 2026-07-04): lawvere_inside_outside_fixed_point_
%  theorem_2026-05-04T1806.md (the spine: the two-principles reduction,
%  the theorem, the sigma candidate, the open items -> op minted per
%  D9); research_summary_neg_prob.md (negativity route; lamp (1+i)/2);
%  dynamics_companion_fanout_unitarity (legs 1a/1c; anti-events; the
%  Brukner-Zeilinger rebuttal per D6); Pure_states_and_mixed_states
%  (local tomography derived; QBism contrast per D7).
%
%  DECISIONS ENCODED (JS 2026-07-04): D1 title as below; D2 Wheeler1990
%  cited here as the title anchor; D3 lineage Popper -> Breuer ->
%  Svozil -> Szangolies (Dalla Chiara omitted); D4 Frauchiger_2018
%  cited here, dissolved in sec:wigners-friend; D5 Kirkwood-Dirac
%  heritage via the Arvidsson-Shukur 2024 review as one-stop; D6
%  Brukner-Zeilinger + the framework rebuttal IN ("Shannon is
%  post-FANOUT"); D7 QBism contrast IN; D8 Aaronson private
%  communication OMITTED (published island paper only); D9 one op
%  minted (op:sigma-functor), cross-linked to op:hopf-frobenius.
%
%  CITATIONS -- existing keys: Wheeler1990, small_2006, Godel1931,
%  Lawvere_1969, Yanofsky2003, CoeckeKissinger2017, Rovelli1996a,
%  Popper1950a, Popper1950b, Breuer1995, svozil_k_2018,
%  szangolies_j_2018, Frauchiger_2018, dirac1942, Feynman_r_82,
%  Feynman1987, Youssef1994 [SWAP-FLAG: intended paper is "Quantum
%  mechanics as complex probability theory", Mod. Phys. Lett. A 9
%  (1994) 2571 -- if that lives under youssef1994quantum, swap],
%  caticha_a_2006, goyal_knuth_skilling_2010, Aaronson2004h, Hardy2001a,
%  muller_p_2011, Thomson_1954, benacerraf1962tasks, cantor1891diagonal,
%  Turing_1937.
%  NEW keys (RIS in self_reference_refs_2026-07-04.ris, all verified
%  this session): tikochinsky1988feynman, renou2021real, chen2022ruling,
%  wu2022refutation, hoffreumonwoods2026real, arvidssonshukur2024kirkwood,
%  bruknerzeilinger2001shannon, fuchsschack2013qbism,
%  shannon1948mathematical.
%
%  CROSS-REFS: forward (fine under two-pass LaTeX): sec:continuity,
%  sec:division-algebras, sec:hopf-justification, prop:f1-dictionary,
%  sec:wigners-friend, op:limit-realisability, op:hopf-frobenius.
%  Still dangling after this section: sec:fanout-boundary (and its
%  sibling label sec:fanout-direction) -- both DRAFTED 2026-07-06.
%
%  AUDITED 2026-07-06 (section-by-section final pass, first stop):
%  proof of prop:inside-outside re-verified by hand; attribution chain
%  clean (Burgin/Khrennikov correctly absent); all header-declared
%  citations confirmed used; lamp arithmetic and antecedents verified;
%  FANOUT smallcaps consistent; no pre-definition use of "frame".
%  Fixes: stray Unicode em-dash (QBism paragraph, the file's only
%  non-ASCII character) -> house ---; trace-blind two-way stitch added
%  to the Born subsection ("one privacy, two duties");
%  op:sigma-functor gains its first standing constraint (F1
%  inheritance via the op:hopf-frobenius discharge) and its previously
%  one-directional local-tomography cargo made explicit.
% =====================================================================

\section{Self-Reference: One Principle}
\label{sec:self-reference-2}

Wheeler asked for it from bit: for physics derived from information, ``every it --- every particle,
every field of force'' obtained from binary answers to yes-or-no questions \autocite{Wheeler1990}.
This paper's route to that programme runs through G\"{o}del: the single principle from which
everything else here is derived, deferred, or priced is \emph{self-reference} --- the fact that any
complete physical description, made from inside the world it describes, must describe the describer.
This section states the principle, proves that the framework's second historical commitment follows
from it as a theorem rather than standing beside it as an axiom, and then extracts the price of the
principle in four instalments: probabilities whose parts can cancel, amplitudes valued in
$\mathbb{C}$, the Born rule, and the local tomography of composite systems. What the section cannot
yet derive is stated as precisely as what it can, in one numbered problem and in the ledger at the
end.

\subsection{Events, not objects}

The framework's ontology is relational: the fundamental entities are events and the relations
between them, and objects are the bookkeeping of relational invariants --- what remains stable
across many relations. The kinship with relational quantum mechanics \autocite{Rovelli1996a} is
genuine and acknowledged: facts are facts relative to the systems that register them. The
difference, developed over the rest of this paper, is that the present framework grounds the
relational structure in a single logical principle and then carries it, through the division
algebras and the Hopf fibrations, all the way to the gauge group and the particle spectrum ---
relations are not only the ontology but the derivation engine. On this reading a measurement is not
the revelation of a pre-existing property; it is the creation of a relation, and the later sections
of this paper read that creation geometrically (\S\ref{sec:hopf-justification}) and even
gravitationally. Here we need only the logical core: whatever describes the world is an event in
the world it describes.

\subsection{One principle, and a theorem}

The framework historically carried two foundational commitments: self-reference, and the
\emph{Inside/Outside Equivalence Principle} --- that the laws of physics must be identical whether
formulated by an embedded participant or by an external observer. Two principles are one too many,
and this subsection reduces them.

The reduction corrects the framework's own origin, and we make the correction in print. The 2005
form of the argument \autocite{small_2006} ran a G\"{o}delian chain: a self-referential process,
being complete, must by the incompleteness theorem \autocite{Godel1931} be inconsistent; an
inconsistent system can negate itself; and the negation of self-reference is the production of
causal sequence. The chain captures the right intuition and is retained here as motivation, but it
is not sound as a proof: classical inconsistency yields \emph{ex falso quodlibet} --- every
proposition provable, every causal sequence an output, no constraint at all. The load-bearing
argument is different, and better: a fixed-point theorem.

Let $\mathcal{C}$ be the category of physical descriptions and admissible processes --- concretely,
the dagger-compact-closed setting with classical structures in which categorical quantum mechanics,
and the whole of this paper, works \autocite{CoeckeKissinger2017}. Lawvere showed that in a
sufficiently rich such setting, any weakly point-surjective self-representation forces every
endomap to have a fixed point, and Yanofsky showed that one diagonal argument of this shape
underlies Cantor's, G\"{o}del's, Turing's and Tarski's theorems alike
\autocite{Lawvere_1969,Yanofsky2003,cantor1891diagonal}. The framework's single commitment is that
$\mathcal{C}$ contains a \emph{self-representing object} --- an object capable of encoding the
processes on itself. That is what ``self-reference is foundational'' means, stated categorically,
and it is the surjectivity hypothesis of the diagonal lemma. Now let
$\sigma \colon \mathcal{C} \to \mathcal{C}$ be the \emph{inside/outside involution}: the operation
exchanging the two complementary presentations of any description --- the inside presentation, as
participant-indexed total space in which the observer is part of the system, and the outside
presentation, as a base-times-fibre decomposition over a classical parameter space. That
$\sigma^{2} = \mathrm{id}$ is not an axiom but a coherence requirement: the outside view of the
inside view of a thing is the thing. The candidate concrete realisation of $\sigma$ is the
Hopf-duality swap --- total space against base-times-fibre for the three fibrations --- which gives
the involution computable content and ties this subsection to \S\ref{sec:hopf-justification}.

\begin{proposition}[Inside/outside fixed point]\label{prop:inside-outside}
Let $\mathcal{C}$ satisfy the richness condition above and contain a self-representing object, and
let $\sigma$ be an involution on $\mathcal{C}$. Then the fixed-point locus
$\mathrm{Fix}(\sigma)$ is non-empty: there exist descriptions invariant under the exchange of
inside and outside perspective.
\end{proposition}

\begin{proof}
Lawvere's construction, applied to $\sigma$ in place of negation: given a weakly point-surjective
$\varphi \colon A \to Y^{A}$, form $g(a) = \sigma\!\left(\varphi(a)(a)\right)$; point-surjectivity
supplies $a^{*}$ with $\varphi(a^{*}) = g$, and then $\varphi(a^{*})(a^{*})$ is a fixed point of
$\sigma$ \autocite{Lawvere_1969,Yanofsky2003}. Where the same construction applied to Boolean
negation yields paradox --- negation has no fixed point on a two-element type, so the conclusion
explodes the premises --- the inside/outside involution is an honest involution on a rich type, and
the conclusion is a constructive fixed point rather than a contradiction.
\end{proof}

\textbf{Definition.} \emph{Physical law is the structure carried by} $\mathrm{Fix}(\sigma)$.

The Inside/Outside Equivalence Principle is thereby demoted from axiom to theorem: physical law is
invariant under the exchange of perspectives \emph{because it is the fixed-point set of that
exchange}, and the fixed-point set is forced to be non-empty by self-reference alone. One
commitment remains foundational; everything the second commitment asserted is now derived. The
reader who accepts categorical quantum mechanics as an existing programme has already accepted the
machinery; the framework adds one premise and collects the rest. And the fixed-point locus is not
empty of structure: it inherits the Frobenius-algebra apparatus of $\mathcal{C}$, so that the
Bell, GHZ and W primitives --- and with them the particle classification of the later sections ---
live inside $\mathrm{Fix}(\sigma)$ from the start. What is proved above is non-emptiness given the
setting; what is assumed is the setting itself and the existence of $\sigma$ as a genuine
involution; what is open is the rigour of both, and it is priced now rather than hidden:

\begin{openproblem}[The involution made functorial]\label{op:sigma-functor}
Fix the category $\mathcal{C}$ rigorously --- the dagger-compact-closed setting with classical
structures is the candidate --- and define $\sigma$ on morphisms, not only on objects, the natural
route being the induced fibre--base exchange on spider diagrams; establish that $\sigma$ is an (at
worst pseudo-)involutive endofunctor; and prove or refute the identification of
$\mathrm{Fix}(\sigma)$ with the three Hopf cross-sections, under which Adams's theorem would become
a corollary of the fixed-point classification. This problem is cross-linked with Open
Problem~\ref{op:hopf-frobenius}: under an affirmative resolution of both, the parallelisable
terminus, the fixed-point classification, and compositional completeness are three statements of
one fact, and progress on either problem constrains the other. One constraint is already in place:
at the complex rung, Proposition~\ref{prop:f1-dictionary} supplies the dictionary between
classical structures and the fibration's lifted skeletons --- the clause of
Open Problem~\ref{op:hopf-frobenius} it discharges --- and any resolution here must reproduce it.
The problem also carries a piece of cargo recorded later in this section: the formalisation, within
the fixed category, of the perspective-accounting that derives local tomography.
\end{openproblem}

\subsection{The lineage}

The obstruction this section builds on is old, and the framework claims its completion, not its
discovery. Popper argued in 1950 that a predictor embedded in the world has irreducibly inferior
knowledge of itself, because it cannot impart information about itself without, by that very act,
rendering the information out of date \autocite{Popper1950a,Popper1950b} --- the passage is quoted
where the framework makes the obstruction structural (\S\ref{sec:wigners-friend}). Breuer proved
the impossibility of accurate state self-measurement within quantum theory \autocite{Breuer1995};
Svozil's monograph surveys the whole territory where physics meets self-reference, undecidability
and Chaitin-style algorithmic limits \autocite{svozil_k_2018}; Szangolies sharpened the modern form
as epistemic horizons \autocite{szangolies_j_2018}; and the Frauchiger--Renner theorem is the
contemporary crisis-form --- nested observers deriving contradictory records when each treats the
others' quantum descriptions as universally valid \autocite{Frauchiger_2018} --- which
\S\ref{sec:wigners-friend} dissolves relationally rather than paradoxically. This framework's
addition to the lineage is twofold. Structurally, Proposition~\ref{prop:f1-dictionary} converts
Popper's temporal obsolescence into algebraic cancellation: the datum a copier cannot put into its
own record is exactly its fibre commitment, which does not arrive stale --- it cancels. And
foundationally, the present section runs the lineage in reverse: where the tradition treats
self-reference as a limit on knowledge \emph{within} physics, the framework takes it as the
generator of physics, and derives the quantum formalism from it.

\subsection{The price is negative probability}\label{sec:complex-amplitudes-2}

What does it cost a probabilistic description to contain its own describer? The first instalment is
the sign. In the regime before any record exists, every event has an anti-event --- the description
must be able to undo what it has provisionally done, or the self-description could never be
revised into consistency --- and a probability calculus closed under undoing cannot be confined to
$[0,1]$: it is a signed calculus, in which the parts of a probability can oppose and cancel. The
framework did not discover negativity as bookkeeping; Dirac admitted negative probabilities as
formal intermediaries \autocite{dirac1942}, Feynman defended them as consistent whenever the final,
recorded probabilities stay positive \autocite{Feynman_r_82,Feynman1987}, and the Kirkwood--Dirac
distribution has carried complex quasiprobability from 1933 to the present, as a recent
comprehensive review documents \autocite{arvidssonshukur2024kirkwood}. What the framework adds is
the forcing: negativity is not a permissible convenience but the first price of self-referential
closure, and it is charged exactly and only in the pre-record regime --- which is why every
negativity in the tradition lives in intermediate, unrecorded quantities and never in a final
record.

The second instalment is complexity. A signed information measure must extend Shannon's
$S = -\ln p$ off the positive axis, and the analytic continuation of the logarithm carries $p$ into
the complex plane: amplitudes, with the phase as the continued information. Here too the extension
has a prior tradition, and the attributions are owed precisely. Tikochinsky derived Feynman's rules
for complex amplitudes from consistency of composition \autocite{tikochinsky1988feynman}; Youssef
reformulated quantum mechanics as complex probability theory outright \autocite{Youssef1994};
Caticha carried objective amplitudes to Bayesian probabilities \autocite{caticha_a_2006}; and
Goyal, Knuth and Skilling derived Feynman's rules from a pair-postulate without assuming
analyticity \autocite{goyal_knuth_skilling_2010}. The extension is prior art. The framework's
claim is the compulsion: the complex extension is not a modelling choice among exotic probability
theories but the unique continuation forced by self-referential closure --- the same forcing that
Aaronson glimpsed when he argued quantum mechanics is what negative probability inevitably becomes
\autocite{Aaronson2004h}.

Why does the continuation stop at $\mathbb{C}$? Downward, the real field is not available: real
quantum theory is experimentally distinguishable from complex quantum theory in network scenarios
and has been disfavoured by experiment
\autocite{renou2021real,chen2022ruling,wu2022refutation}; the recent objection that the exclusion
rests on locally tomographic composition \autocite{hoffreumonwoods2026real} is, for this framework,
not an escape but a confirmation of where the weight sits --- local tomography is \emph{derived}
below, not assumed, so the framework closes from its own resources exactly the door the objection
holds open. Upward, the quaternions and octonions fail as a \emph{base} for the information
continuation: the continued measure is a logarithm, and beyond $\mathbb{C}$ the logarithm of a
product is no longer the sum of logarithms --- non-commutativity severs
$\log(ab) = \log a + \log b$ --- so the additivity that makes it an information measure is lost.
Aaronson's parameter-counting and signalling arguments against a quaternionic base point the same
way \autocite{Aaronson2004h}. The higher algebras are not banished; they return as the
\emph{internal} algebras of the register, which is the business of \S\ref{sec:division-algebras}.
The base is $\mathbb{C}$, for informational reasons, and stays $\mathbb{C}$.

\subsection{The Born rule from fibre privacy}\label{sec:born-rule}

The third instalment is the quadratic rule. An amplitude carries one datum more than any record of
it can: its overall phase. This was asserted throughout the framework's development as the
observer's private, self-referential ignorance; it is now a proved statement, in exactly the form
needed. By Proposition~\ref{prop:f1-dictionary}, a classical structure --- the copy operation any
record requires --- is committed to a specific lift of its skeleton, and that commitment is absent
from every classical record the structure induces: the copier is fibre-committed, the record is
fibre-blind, and the fibre freedom is gauge for every observer. (The computation extends verbatim
from one qubit to the full register, where Proposition~\ref{prop:trace-blind} makes it the ground
of hypercharge's tracelessness --- one privacy, two duties.) Popper's obsolescence, made
cancellation. It follows that anything recordable is invariant under the fibre --- observables and
their statistics must be $U(1)$-invariant functions of the amplitude. The lowest-degree polynomial invariant
of a complex amplitude is the quadratic $|\psi|^{2}$ --- the non-polynomial invariant $|\psi|$,
of degree one, is excluded by the analyticity declared at \S\ref{sec:complex-amplitudes} --- and
normalisation fixes the coefficient:
probability is the minimal record-compatible function of the amplitude,
\[
p \;=\; |\psi|^{2},
\]
the Born rule as the cheapest invariant the fibre permits. The same square is reached from the
cost side in \S\ref{sec:born-cost-ledger}, where a record is the closed loop of a history and its
return leg and the exponent is the number of legs; neither route uses the lattice of projections
that Gleason's theorem needs, and both hold for a single qubit. The derivation is the framework's
own;
its premise --- that the fibre datum is structurally excluded from records --- is
Proposition~\ref{prop:f1-dictionary}, and its geometric reading --- probability as base-point data
under the Hopf projection --- is developed in \S\ref{sec:hopf-justification}.

One standing objection to information-theoretic foundations is answered here, because this
framework needs the answer. Brukner and Zeilinger argued that Shannon information is conceptually
inadequate for quantum measurements: it is basis-dependent, and its operational meaning presupposes
ignorance of a definite but unknown fact, which the quantum case does not supply
\autocite{bruknerzeilinger2001shannon}. The framework's reply is structural, and it is a
demarcation rather than a rebuttal. Shannon's precondition fails exactly where they say it fails
--- in the pre-record regime, where there is no definite unknown to be ignorant of --- and there
the operative measure is the quadratic invariant just derived. Shannon's axioms hold exactly where
the precondition holds --- on records, after \textsc{fanout}, where outcomes are definite and
ignorance is classical --- and there the uniqueness theorem forces $-\log p$
\autocite{shannon1948mathematical}. The basis-dependence dissolves with the same stroke: a basis
\emph{is} selected precisely at \textsc{fanout}, so the basis-dependent measure applies only where
a basis exists. Shannon is not inadequate; Shannon is post-\textsc{fanout}. The framework thereby
keeps both measures, each on its own side of the boundary, and uses the additivity of the
post-record measure --- the very property the previous subsection used to pin the base field ---
with a clear conscience.

\subsection{Thomson's lamp, resolved}

The section owes \S\ref{sec:continuity} a debt incurred there: the lamp's amplitude-side limit.
Thomson's lamp --- toggled at times $1 - 2^{-n}$, its state at $t = 1$ demanded --- has no
classical answer: the task description fixes every finite stage and thereby fixes no limit at all
\autocite{Thomson_1954,benacerraf1962tasks}, which \S\ref{sec:continuity} read as the proof that
the record algebra is not limit-closed. On the amplitude side the limit exists, and the calculus of
this section computes it. The toggle exchanges the two amplitudes; the self-consistent state of the
completed task is the fixed point of the exchange in the complex calculus, and solving it under the
signed continuation gives
\[
P(\mathrm{on}) = \tfrac{1 + i}{2}, \qquad P(\mathrm{off}) = \tfrac{1 - i}{2},
\]
a superposed fixed point whose imaginary part carries the unrecordable datum --- the phase of the
endless alternation, when the lamp \emph{would have} stopped. Two checks certify that the
assignment is calculus rather than decoration. The pair is normalised,
$P(\mathrm{on}) + P(\mathrm{off}) = 1$; and the Born rule of the previous subsection, applied to
the fixed point, returns $|P(\mathrm{on})|^{2} = \tfrac{1}{2}$ --- exactly the symmetric classical
ignorance that is the only sane record-level answer, recovered rather than imposed. The lamp,
which entered the framework at its very beginning as the place where classical computation meets
its boundary, thus exits it as a worked example: undetermined as a record, determinate as an
amplitude, the two sides split precisely along \textsc{fanout} --- and the diagonal construction it
physically instantiates \autocite{cantor1891diagonal,Turing_1937} is the same one that powers
Proposition~\ref{prop:inside-outside}. The status flag is owed here as everywhere: the
undetermination is theorem-level; the specific fixed-point value is the framework's complex
calculus, structural, offered with its two consistency checks shown.

\subsection{Local tomography, derived}\label{sec:local-tomography}

The fourth instalment concerns composites, and it discharges two promises at once. The
Inside/Outside Equivalence Principle --- now a theorem --- requires that an external observer's
characterisation of a composite system be recoverable by the embedded participants, and the
embedded participants have access only to local measurements and their correlations. Equivalence
of the two perspectives therefore entails that the state of a composite is determined by the
statistics of local measurements: \emph{local tomography}, obtained as a consequence rather than
imposed as an axiom. This converts a postulate of the reconstruction programmes
\autocite{Hardy2001a,mullerp2011} into a theorem-shaped consequence of the one principle ---
discharging the caveat recorded in \S\ref{sec:continuity} --- and it closes, from the framework's
own resources, the escape route in the real-number debate noted above: the assumption the
objection of \autocite{hoffreumonwoods2026real} contests is here not an assumption. The status is
flagged plainly: the derivation is structural, one paragraph of perspective-accounting, and its
formalisation within the categorical setting belongs to Open Problem~\ref{op:sigma-functor}.

\subsection{Contrast: QBism, and the ledger}

The nearest neighbour in spirit deserves a paragraph of honest contrast. QBism
\autocite{fuchsschack2013qbism} holds, as this framework does, that quantum probabilities are an
agent's --- that the theory is a calculus of expectation, not a catalogue of properties. What QBism
does not say is why the agent's expectations must be complex-valued and quadratically squared; the
Born rule enters as a normative addition to Bayesian coherence, and the complex amplitude as an
unexplained representation. This section is, in effect, the missing why: the agent's ignorance is
self-referential --- the agent is inside the description --- and self-referential ignorance is
complex, with the quadratic rule as the cheapest record-compatible invariant of it. Where QBism
stops at the agent, the framework derives the agent's calculus.

The ledger, then. Delivered in this section: the reduction of two principles to one
(Proposition~\ref{prop:inside-outside}); the pre-record regime with its anti-events, handed to
\S\ref{sec:continuity} as the five-face regime whose continuity is priced there
(Open Problem~\ref{op:limit-realisability}); the complex base, forced downward by experiment and
upward by the additivity of information; the Born rule from fibre privacy; the lamp's fixed point;
and local tomography. Owed, and stated where they are owed: the functorial rigour of the involution
and the fixed-point classification (Open Problem~\ref{op:sigma-functor}); the remaining premise of continuity (Open Problem~\ref{op:limit-realisability}, in
\S\ref{sec:continuity}); and the
rung-by-rung compulsion of the algebra ladder, posed as Open
Problem~\ref{op:ladder-compulsion} in \S\ref{sec:division-algebras}. The
hand-on is the paper's spine: the $\mathbb{C}$ this section derives natively and the $\mathbb{C}$
that \S\ref{sec:continuity} underwrites through Hardy's theorem are the same field, the floor of
the ladder that \S\ref{sec:division-algebras} climbs and \S\ref{sec:hopf-justification}
geometrises. One principle; everything else is theorem, derivation, or a debt with a stated
invoice.

%      % EDITED 2026-10-29 scripts/patch_continuity_tower_premise_2026-10-29T1500.py
% EDITED 2026-09-15 scripts/patch_value_continuum_declared_cencov_2026-09-15T1540.py
% EDITED 2026-09-15 scripts/patch_continuity_connected_state_space_2026-09-15T1430.py
% =====================================================================
%  SECTION: Continuity Is Not an Axiom
%  Drafted 2026-07-04 from continuity_limit_realisability_2026-07-04.md
%  (now in project documents) and the May record
%  (dynamics_companion_fanout_unitarity_2026-05-20T1430.md).
%
%  PLACEMENT in main.tex: a standalone section between self-reference
%  and division-algebras --
%      \input{self_reference_section_2026-07-04}
%      \input{continuity_section_2026-07-04}
%      \input{division_algebras_section_2026-07-03}
%  (Architecture call SETTLED by the sync manifest: standalone.)
%
%  CITATIONS -- existing keys (as supplied by JS 2026-07-04):
%    Hardy2001a, Aaronson2004h, Thomson_1954, Lawvere_1969,
%    Turing_1937, muller_p_2011, Brukner2009.
%  NEW keys (RIS in continuity_refs_2026-07-04.ris):
%    cantor1891diagonal, benacerraf1962tasks, broeckertomdieck1985compact.
%  Deliberately excluded per JS: Chiribella--D'Ariano--Perinotti;
%  Bargmann/Wigner's theorem. Yanofsky2003 now paired on the
%  Lawvere-shaped clause (the pending-key excuse was void: the key
%  exists and sec:self-reference uses it).
%
%  CROSS-REFS: sec:self-reference now DRAFTED (2026-07-04, later the
%  same day) -- no longer dangling; only sec:fanout-boundary (and its
%  sibling sec:fanout-direction) -- DRAFTED 2026-07-06.
%
%  AUDITED 2026-07-06 (final pass, second stop): both proofs
%  re-verified by hand (Lagrange divisibility sentence exact; the
%  Baire countable-escape chain sound); caveat-discharge with
%  sec:self-reference confirmed two-way (already present); file pure
%  ASCII; no pre-definition "frame"; all citations used as declared.
%  Fixes: Yanofsky2003 paired on the Lawvere clause; the op's two
%  candidate routes anchored to printed propositions (f1-dictionary
%  for the divisibility premise's referent, inside-outside for the
%  closedness demand) -- statuses unchanged, referents sharpened;
%  header staleness corrected. Deliberate FORWARD references --
%  prop:f1-dictionary and sec:hopf-justification (defined later in the
%  paper; fine under two-pass LaTeX and flagged in prose),
%  sec:division-algebras, sec:space-of-computations.
%
%  DECISIONS ENCODED: op:limit-realisability minted (the owed arrow,
%  priced: exactly two derivations, divisibility and closedness);
%  Lagrange used twice (finite exclusion; component-group triviality);
%  the (1+i)/2 lamp value NOT stated here -- "a superposed fixed point"
%  with the complex-probability treatment deferred to
%  sec:self-reference; the two-branch over-determination stated to
%  pre-empt a circularity reading; Thomson spelling throughout.
% =====================================================================

\section{Continuity Is Not an Axiom}
\label{sec:continuity}

Hardy derives quantum theory from five axioms \autocite{Hardy2001a}. Four of them are satisfied
by classical probability theory; the fifth --- that a \emph{continuous} reversible
transformation connects any two pure states --- is the entire difference, the single postulate
that turns the classical simplex into the complex Hilbert space. Every reconstruction programme
pays an equivalent toll: continuous reversibility in Masanes and M\"{u}ller
\autocite{mullerp2011}, the corresponding postulate in Daki\'{c} and Brukner
\autocite{Brukner2009}. The axiom is load-bearing, and in every programme it is unexplained: it
is where the quantumness is put in by hand. The claim of this section is that in the present
framework it is not an axiom. Continuity is the fifth face of the pre-measurement regime that self-reference forces, and what
stands between the framework and a full derivation is not a postulate about values but one
premise about transformations --- that every transformation the inside performs has a stationary
outside presentation --- derived in outline below from the inside/outside involution of
\S\ref{sec:self-reference} and stated as an open problem at the end of the section with nothing
hidden in it. The continuum itself, the completion that every probability theory adopts, is
named for what it is: the real line, shared with classical probability and carrying no quantum
content.

One structural remark prevents a misreading. The framework reaches complex amplitudes natively in
\S\ref{sec:self-reference}, by the negative-probability route; the present section reaches the same
destination by an independent branch --- the regime's continuity, then Hardy's theorem. The two
branches converging on one structure is over-determination, not circularity, and the convergence is
itself evidence that the destination is forced.

\subsection{Five faces of one regime}

The regime before any measurement --- before any \textsc{fanout} --- can be characterised five
equivalent ways: negative probability is available; anti-events exist, every event being undoable;
no classical record exists; no preferred basis is selected; and a continuous reversible
transformation connects any two pure states. These are not five premises to be combined. They are
one regime under five descriptions, and \textsc{fanout} (\S\ref{sec:fanout-boundary}) is the single
event that fails all five at once: it writes the record, selects the basis, kills the anti-events,
discretises the transformation group, and switches off the negativity. The equivalence of the first
two faces is the anti-event identity of \S\ref{sec:self-reference}; the third is the boundary's
definition. The work of this section is the equivalence of the last two faces, and then the route
from the regime's existence to its continuity. The fourth face has, since
Proposition~\ref{prop:f1-dictionary} (\S\ref{sec:hopf-justification} --- a deliberate forward
reference), an exact form: a basis-selection is the commitment of a classical structure, a lifted
antipodal skeleton, so ``no preferred basis'' means precisely that no copy structure has fired.

\subsection{Basis selection is the collapse of continuity}

\begin{proposition}[Continuity forbids a selected basis]\label{prop:basis-discreteness}
If a finite set of pure states is privileged --- a preferred basis --- then every reversible
transformation respecting the privilege carries the set to itself and acts on it as a permutation;
the respecting transformations therefore realise only a finite group of motions among the
distinguished states, and no continuous family of them connects distinct pure states of the set.
Contrapositively, continuous reversibility entails that no basis is selected.
\end{proposition}

\begin{proof}
A privileged finite set is by definition carried to itself by every structure-preserving reversible
map, whose induced action on the set is a permutation of finitely many points. A continuous path of
such maps induces a continuous path in a finite discrete set, which is constant. Hence a continuous
reversible family moving one distinguished pure state to another respects no such privilege, and if
every pair of pure states is so connected, no finite set is privileged.
\end{proof}

The converse is where the historical difficulty lies: the absence of a privileged finite set does
not by itself force continuity, for the transformation group could in principle be infinite and
discrete. Two structural properties close every such escape at once.

\begin{proposition}[Divisible and closed implies connected]\label{prop:divisible-closed}
Let $G$ be the group of reversible transformations of the pre-measurement regime, acting on the
state body. If $G$ is divisible --- every element has an $n$-th root for every $n$ --- and
topologically closed, then $G$ is a compact connected Lie group; and if $G$ moreover acts
transitively on the pure states, then a continuous reversible transformation connects any two of
them: Hardy's fifth axiom.
\end{proposition}

\begin{proof}
The state body is compact and convex with nonempty interior, so its symmetry group is compact, and
a closed subgroup of a compact matrix group is a compact Lie group with finitely many connected
components \autocite{broeckertomdieck1985compact}. The component group $G/G_{0}$ is a quotient of a
divisible group and hence divisible; but a nontrivial finite group is never divisible --- with
$\lvert H\rvert = m$ every element satisfies $h^{m} = e$, so no element other than the identity
possesses an $m$-th root --- and $G/G_{0}$ is finite, so it is trivial and $G = G_{0}$ is connected.
A connected Lie group is path-connected, and transitivity then supplies the continuous reversible
path between any two pure states.
\end{proof}

The same one-line fact of Lagrange closes the remaining escapes. Finite groups are excluded
directly, being nontrivial and never divisible. Countable groups are excluded because a
countable closed subgroup of a compact group is compact, has an isolated point by the Baire
category theorem, and is therefore discrete and finite. Uncountable disconnected groups are
excluded by the component-group argument of the proposition. Every route to a discrete
transformation group therefore passes through a failure of divisibility or a failure of
closedness. Divisibility, however, asks for more than the conclusion needs, and the next
proposition removes it.

\begin{proposition}[Continuity from a connected state space]\label{prop:connected-transitive}
Let the pure states of a finite-dimensional system form a connected space $P$, and let the group
$G$ of reversible transformations of the pre-measurement regime be topologically closed and act
transitively on $P$. Then the identity component $G_0$ already acts transitively on $P$, and a
continuous reversible transformation connects any two pure states: Hardy's fifth axiom holds,
with no divisibility assumed.
\end{proposition}

\begin{proof}
As in Proposition~\ref{prop:divisible-closed}, closedness makes $G$ a compact Lie group with
finitely many connected components \autocite{broeckertomdieck1985compact}, so $G_0$ is a closed
normal subgroup of finite index and is path-connected. Each $G_0$-orbit in $P$ is the continuous
image of a compact group, hence compact, hence closed. Since $G$ is transitive and $G_0$ is
normal, the $G_0$-orbit of $gx$ is $g$ applied to the $G_0$-orbit of $x$, so the cosets of $G_0$
permute the orbits transitively and there are finitely many of them. A finite partition of $P$
into closed sets is a partition into open sets; $P$ is connected; so there is one orbit and
$G_0$ is transitive. Given pure states $x$ and $y$, choose $g \in G_0$ with $gx = y$ and a
continuous path $g_t$ in $G_0$ from the identity to $g$; then $g_t x$ is a continuous reversible
transformation from $x$ to $y$.
\end{proof}

The proposition relocates the premise. Divisibility was asked for because a transformation group
could be infinite and discrete; but a group acting transitively on a continuum cannot be
countable, and a closed group acting transitively on a connected space cannot be disconnected in
the way that matters, so what the argument needs from the regime is not the availability of
$n$-th roots but the connectedness of the state space itself. That is a statement about values:
that the pure states are all the unit vectors of $\mathbb{C}^{K}$ up to phase, the space
$\mathbb{CP}^{K-1}$, connected for $K \ge 2$, and hence that the amplitudes of
\S\ref{sec:self-reference} fill their continuum. The two standing foils show the premise doing
its work. The stabiliser fragment of quantum theory has complex amplitudes and is contextual,
yet it has finitely many pure states, its outcome probabilities take only the values $0$,
$\tfrac12$ and $1$, and its reversible transformations form a finite group, the Clifford group,
so Hardy's fifth axiom fails in it; Spekkens's toy theory \autocite{spekkens2007toy} is the same
with real amplitudes, its ignorance rationed by a knowledge-balance principle. Neither has a
connected state space, and neither is a model of the regime described here, whose credences take
every value. Divisibility, which both also lack --- the phase gate has no square root in the Clifford group
--- is a symptom of the discreteness, not its cause. What both lack can be said exactly, and the
next subsection says it: each contains a transformation that fixes no state it admits.

\subsection{The premises, and what they mean}

Three premises carry Proposition~\ref{prop:connected-transitive}. Transitivity is the regime's
homogeneity: inside/outside equivalence leaves no pure state privileged, and a proper closed
invariant set of pure states would define a superselection observable, a classical record
predating any \textsc{fanout} (\S\ref{sec:fanout-boundary}). Connectedness of the state space is two statements, and only one of them is quantum. The first
is that the pure states are infinite and dense, with a fixed point for every transformation; the
second is that the state space is complete. Quantum theory over the algebraic numbers ---
amplitudes in the algebraic closure of $\mathbb{Q}$ within $\mathbb{C}$ --- has the first
without the second: it has $K = N^{2}$, the Born rule, superposition, entanglement and Bell
violation at the Tsirelson value, every unitary has an algebraic eigenvector, and it is
countable and totally disconnected. So the whole quantum content of Hardy's axiom lies in the
density, and the completion is the real-number convention of every probability theory. Hardy
adopts it before his fifth axiom is stated, in the first, which makes probabilities real numbers
and gives his classical theory a real simplex of mixtures; this framework adopts it at the
Shannon step, where credences are declared to range over a connected domain so that the
continuation of the logarithm is unique (\S\ref{sec:complex-amplitudes}). Given that convention,
density of the pure states is connectedness of their closure. Divisibility, which an earlier
version of this section asked for, is not needed: it follows from the conclusion, every element
of a compact connected Lie group having roots of every order.

Closedness is the completion, limit-realisability: a Cauchy sequence of available
transformations has its limit available. It cannot be derived, and the reason should be stated
rather than hidden. Any collection of first-order principles about states, transformations and
probabilities that has an infinite model has a countable one, and a countable subset of
$\mathbb{CP}^{K-1}$ is totally disconnected; the founding non-computability does not escape
this, since the arithmetical complex numbers form a countable algebraically closed field
containing non-computable numbers and closed under every limit it can itself describe. Hardy's
axiom and closedness are second-order, quantifying over paths and over sequences, and a
second-order principle is grounded as Dedekind completeness is grounded or not at all. The
grounds are these. The completion of a dense set under the finite-precision metric is unique and
depends on no choice, whereas any countable alternative is a choice of definability level, an
extra law with no ground and no test; and the framework's own result is that no observer's
descriptive level captures the regime, since by \S\ref{sec:self-reference} every observer's
state exceeds its own level and the inside/outside principle makes that a law rather than a
perspective. The burden therefore falls on any restriction, not on the completion, which is the
least committal option and carries no quantum content. What limit-realisability adds to the line
is still worth stating exactly, because it is the corrected form of a natural intuition.
Cantor's diagonal \autocite{cantor1891diagonal} does not build the continuum; its output is an
infinite decimal --- a Cauchy sequence of rationals --- which exists \emph{as a point} only
because the line is complete. The diagonal presupposes limit-realisability and then shows that
the completed line outruns every internal listing. And the character of the completion is the
framework's own: the computable reals are countable \autocite{Turing_1937}, hence of measure
zero, so once the transformation group contains its limits almost all of them are
non-computable. Non-computability is a property of the completed continuum, not the reason for
it. The constructive reading, on which only describable limits exist, remains consistent; it is
the countable theory above, and the two readings differ in nothing an experiment can certify.

Density is where the section's title is earned, and it is derived from the paper's own premise.
Write $(\mathrm{I})$ for the statement that every admissible transformation of the regime is a
phase of some admissible classical structure: diagonal, that is, in the record basis of some
observer the regime contains, a phased spider of that basis in the calculus of
\S\ref{sec:hopf-justification} \autocite{coeckeduncan2011interacting} (VERIFY-SUPPORT). It says
of transformations what the global phase of \S\ref{sec:self-reference} says of states: there are
no absolute ones. It follows from three things the paper already holds. The inside/outside
involution is a genuine involution, so the outside presentation of the inside is a presentation
of the same object with nothing lost: the outside can have complete knowledge of the inside, a
potential record, the joint state lying in a structure the outside carries. Time is operation
count, internal to the system that operates (\S\ref{sec:space-of-computations}), so the outside
does not share the inside's count. And before any \textsc{fanout} that count is not classically
recorded, so the outside's description is pure. Let $g$ be the transformation the inside
performs per operation. The outside's complete knowledge cannot depend on the inside's count,
and it is pure, so it is a pure state of the inside's clock and system that is stationary under
``advance the clock and apply $g$''; the stationary states of that map are clock Fourier modes
tensored with eigenvectors of $g$ and nothing else \autocite{pagewootters1983evolution}
(VERIFY-SUPPORT). So the outside holds a structure in which $g$ acts by phases, which is
$(\mathrm{I})$; and the relative phase in that structure is the outside's ignorance of the
inside's computational count, the same ignorance the global phase records one level down. The
consequences are a theorem. A phase of a structure fixes the structure's copyable states, so
every admissible transformation fixes an admissible pure state; a finite transitive action has
an element moving every point (Jordan), so the pure states are infinite; the transformation
carrying a stabiliser state to the fixed point of the composite of a $Z$-phase and an $X$-phase
turns through $\arccos(1/\sqrt{3})$, an irrational fraction of a turn, so the group has an
element of infinite projective order and its orbits are dense. With the completion,
Proposition~\ref{prop:connected-transitive} fires and Hardy's fifth axiom holds. The two foils
fail $(\mathrm{I})$ exactly where they fail continuity: fourteen of the twenty-four Clifford
operations are phases of no Pauli basis and fix no stabiliser state, their fixed points being
the magic states; the toy theory's three-cycles are phases of no epistemic basis; and in each
the outside's knowledge of the inside under the order-three protocol is either count-dependent
or a mixture over the count, so each makes the inside's time classical before any record exists.
The derivation of $(\mathrm{I})$ is an outline, not a proof, and it is what the open problem
below asks for.

\subsection{Why the cut matters physically}

The regime's continuity is not decorative. Reversibility alone, over complex amplitudes, licenses
the full invertible group $GL(n, \mathbb{C})$, and Aaronson showed what that buys: superluminal
signalling, and computational power at the level of \textsf{PP} \autocite{Aaronson2004h}. The cut
from $GL(n, \mathbb{C})$ down to the unitary group is the conserved quadratic form, and continuity
is what entitles the theory to make it: with the fifth axiom in hand, Hardy's theorem delivers the
complex Hilbert space, the quadratic norm and the unitary group in one stroke \autocite{Hardy2001a}
--- cited here, not re-proved. One honest caveat travels with the citation. Hardy's reconstruction
uses his remaining axioms besides --- composition and subspace structure --- and these must be
derived rather than imported: the composition ledger of \S\ref{sec:division-algebras} and the
locality-of-tomography argument of \S\ref{sec:self-reference} carry that duty, and it is discharged
there, not silently assumed here.

\subsection{Super-tasks at the boundary}

The premise just priced has a fifty-year-old stress test in the philosophy of mathematics, and the
framework passes it in an unexpected way. A super-task --- infinitely many operations completed in
finite time \autocite{Thomson_1954} --- presupposes its terminal moment as an actual point of the
timeline: the accumulation point of a generic sequence of step-times exists only in the completed
line, and the defining question, the state \emph{after all the steps}, is meaningful only under
limit-realisability of time. Super-tasks are downstream consumers of exactly the completeness at
issue in this section. Thomson's lamp then exhibits both sides of the \textsc{fanout} boundary at
once. On the record side, the lamp's description fixes its state at every finite stage and thereby
fixes no limit at all: any terminal value is consistent with the whole task
\autocite{Thomson_1954,benacerraf1962tasks}. That is not a puzzle about lamps; it is the proof that
the classical record algebra is not limit-closed, and cannot be. On the amplitude side the limit
exists --- a superposed fixed point, whose complex-probability treatment belongs to
\S\ref{sec:self-reference} --- because that side \emph{is} limit-closed, which is the fifth face
itself. The boundary is thereby the edge of limit-realisability: complete on the reversible side,
incomplete on the record side, with the lamp standing on both sides at once as the boundary's oldest
exhibit. A final clause dissolves rather than answers the classic question. In this framework time
is operation count (\S\ref{sec:space-of-computations}), and ``after infinitely many operations''
indexes no operational moment: the lamp packs its endless \textsc{fanout}s into the emergent
continuum, and asking for its record at the accumulation point asks the discrete contents for a
value that only the completed container possesses. On or off is not undecided; it is unasked.

\subsection{What remains, priced}

\begin{openproblem}[Every transformation a phase of an inside structure]\label{op:limit-realisability}
Proposition~\ref{prop:connected-transitive} reduces Hardy's fifth axiom to three premises:
connectedness of the pure states, transitivity, and closedness. Transitivity is argued above
from the absence of records. Closedness is the completion, adopted on the grounds stated above
and carrying no quantum content. Connectedness is density plus that completion, and density
follows, by Jordan's theorem and the infinite-order argument, from $(\mathrm{I})$: that every
admissible transformation of the regime is a phase of some admissible classical structure. What
remains open is the rigour of the derivation of $(\mathrm{I})$ from the paper's premise. Derive,
inside the fixed-point locus $\mathrm{Fix}(\sigma)$ of Proposition~\ref{prop:inside-outside},
that every endomap is a phase of a classical structure in $\mathrm{Fix}(\sigma)$, from the
self-representing object, the involution $\sigma$ as a genuine involution, and the reading of
the inside presentation as a sequence indexed by operation count with no record before
\textsc{fanout}. The outline above does this with the stationary-state construction of Page and
Wootters; a derivation passes if it makes that outline a theorem of the calculus, with the
outside's count-independent complete description shown to exist and to be pure from the
involution alone. It fails before any detail is checked if every premise it uses holds in the
stabiliser fragment of quantum theory, in Spekkens's toy theory \autocite{spekkens2007toy}, in
the real qubit, or in quantum theory over the Gaussian rationals, all of which violate
$(\mathrm{I})$ outright and all of which make the inside's operation count classical before any
record exists. It also fails if it imports the completion as a premise about transformations:
the completion is the real line and has no derivation, and a derivation of $(\mathrm{I})$ must
hold in the countable algebraic theory, where $(\mathrm{I})$ is true and continuity is not.
\end{openproblem}

The accounting, in one sentence: the axiom that every reconstruction programme pays is, in this
framework, one derivation posed, of Lawvere's shape --- that every endomap of the fixed-point
locus is a phase of a classical structure in it --- and one convention named, the completion,
which classical probability pays too. And
the hand-on is already visible: the complex amplitudes that this section's continuity underwrites
through Hardy's theorem, and the complex amplitudes that \S\ref{sec:self-reference} derives natively,
are one and the same $\mathbb{C}$ --- the floor of the ladder that \S\ref{sec:division-algebras}
climbs.

%      % =====================================================================
%  SECTION: The Division-Algebra Ladder
%  Drafted 2026-07-03. One of the three foundations sections
%  (sec:self-reference, sec:division-algebras, sec:fanout-boundary);
%  this one resolves the dangling \S\ref{sec:division-algebras} in the
%  Hopf-justification section, whose header names it "the
%  Cayley--Dickson / Hurwitz step".
%
%  PLACEMENT in main.tex (REORDERED per JS 2026-07-06): position 6,
%  opening Act II ("why stop at C?") right before its first consumer --
%      \input{self_reference_section_2026-07-04}
%      \input{continuity_section_2026-07-04}
%      \input{L2_metric_global_phase_section_2026_06_07}
%      \input{wigners_friend_in_the_hopf_picture_2026-06-13}
%      \input{derivation_of_the_quantum_postulates}
%      \input{division_algebras_section_2026-07-03}   <-- moved here
%      \input{spacetime_from_non-computability}       <-- first consumer
%                                                          (S^7 exclusion)
%      \input{hopf_justification_section_2026-06-08}
%  The section is position-portable: all adjacency phrasing was
%  retargeted to explicit \S\ref on 2026-07-06. The S^7 hand-off
%  sentence to sec:space-of-computations WRITTEN 2026-07-06 (spacetime
%  file read; the closing now hands first to its immediate consumer,
%  then to hopf-justification). Upload verifications COMPLETE: L2 and
%  quantum-postulates do not consume the ladder; WF's references are
%  forward-voiced or self-carrying.
%
%  DIVISION OF LABOUR (deliberate, to avoid overlap):
%   - sec:self-reference (undrafted): self-reference -> complex
%     amplitudes and the Born rule; why not a real base. This section
%     TAKES that as given and starts at C.
%   - THIS section: the algebra -- the doubling construction, the
%     ledger of losses-as-physics, the internal-algebras-not-a-new-QM
%     scoping (Adler pre-empt), the two appearances of the ladder
%     (register / measurement climb, stated lightly), and Hurwitz as
%     the algebraic termination with the committed zero-divisor
%     reading.
%   - sec:hopf-justification: the geometry (Adams, parallelisability,
%     the fibrations), the join of the two classifications, the
%     entanglement-sensitivity literature, the C-direction honest debt,
%     and the epistemic reading of termination.
%
%  COMMITTED READINGS USED (sources re-read 2026-07-03):
%   - Loss ledger and level dictionary: cayley_dickson_qubits_for_review
%     sec.5 ("each step loses a symmetry, and each lost symmetry is
%     identified with a physical content").
%   - Termination reading: grounding_argument_gauge_group ("the
%     sedenions have zero divisors and cannot support consistent
%     information flow around loops") -- leads; norm-composition-as-
%     probability given as its amplitude-level restatement.
%   - Internal-not-base scoping: cayley_dickson_qubits_for_review Q2 and
%     sec.3.2 caveat (Adler objection; C-qubit treated geometrically).
%
%  HONEST DEBT: ladder proved closed at the top (Hurwitz), derived at
%  the floor. The rung-compulsion revisit-trigger FIRED 2026-07-06:
%  sec:self-reference and op:sigma-functor now exist, so the thesis is
%  minted as op:ladder-compulsion (with the CD-conjugation-as-sigma
%  candidate route), inheriting op:sigma-functor's status.
%
%  AUDITED 2026-07-06 (final pass, third stop): all four citations
%  used; all cross-refs live; CD facts, dims, alternativity, Hurwitz
%  statement verified; Adler subsection delivers everything the
%  register's compressed pointer promises; file pure ASCII; no bare
%  "frame". Fixes: opening rewritten for both predecessors (continuity
%  was inserted between after drafting); "associator's non-vanishing is
%  mass itself" recalibrated to the register's debt-against-selection
%  defusal; the O->S ledger row and Hurwitz subsection reconciled with
%  the fourth-qubit forward remark ("the terminus closes the ladder of
%  forces; what survives past it is, at most, flavour").
%
%  CROSS-REFS: sec:self-reference [sibling dangle -- resolves when that
%  section is drafted], sec:hopf-justification, sec:hopf-information-bound,
%  sec:register, sec:entanglement-classes, sec:generations, sec:masses,
%  sec:wigners-friend. All others exist.
%  CITATIONS: existing keys baez2002octonions, hurwitz1898composition,
%  furey2018three. ONE new key: adler1995quaternionic (RIS in
%  division_algebras_refs_2026-07-03.ris).
%  British English; \autocite; no bullets in prose; the ledger is a
%  table because the comparison is the argument.
% =====================================================================

\section{The Division-Algebra Ladder}
\label{sec:division-algebras}

The foundations behind us end at amplitudes twice over: natively, where the self-referential
argument forced the description of a system into the system's own state space and priced
consistency at a number system whose parts can cancel --- complex amplitudes, and the Born rule
that squares them (\S\ref{sec:self-reference}) --- and independently, where the regime's
continuity delivers the same field through Hardy's theorem (\S\ref{sec:continuity}), the
convergence of the two branches being evidence rather than accident. This section asks the question
both routes leave open. Why
should the number system stop at $\mathbb{C}$? Ordinary quantum mechanics stops there by postulate;
a framework that has forbidden itself postulates must answer, and the answer is that the climb
neither stops at $\mathbb{C}$ nor continues forever. There are exactly two more rungs, each
purchased by surrendering an algebraic law and each paying for itself in physics, and the ladder is
closed at the top by a theorem of 1898. This section sets out the algebra of that ladder: the single
construction that climbs it, the ledger of what each rung costs and buys, and the theorem that ends
it. The geometry of the same ladder --- that each rung acts on a sphere, that exactly those spheres
are parallelisable, and that the algebraic and geometric classifications are one --- is the business
of the geometry section (\S\ref{sec:hopf-justification}).

\subsection{One construction climbs the whole ladder}

The rungs are not a miscellany; a single construction generates them all. Given an algebra with a
conjugation, the Cayley--Dickson construction builds its double on ordered pairs, with a twisted
multiplication in which the conjugate enters and a conjugation $(a,b)^{*} = (a^{*}, -b)$
\autocite{baez2002octonions}. (Conventions differ in the placement of the conjugations; all yield
isomorphic algebras.) Applied repeatedly it produces the sequence
$\mathbb{R} \to \mathbb{C} \to \mathbb{H} \to \mathbb{O} \to \mathbb{S}$, of real dimensions
$1, 2, 4, 8, 16$: the reals, the complex numbers, the quaternions, the octonions, the sedenions.
The ground floor is classical: $\mathbb{R}$ is ordered and positive-definite, the arithmetic of
classical probability. The first doubling is therefore the quantum move itself --- the passage from
probabilities to amplitudes that can interfere --- and it is not made here by choice: it is the
derivation of \S\ref{sec:self-reference}, already behind us. What this section must account for is
the rest of the climb.

\subsection{The ledger}

In mathematics the doublings are usually told as a story of degeneration: each algebra is worse
behaved than the last, and the sequence is studied for where it breaks. This framework reads the
same ledger in the opposite direction. Each rung is bought by surrendering a law, and each
surrendered law is exactly a piece of physics --- the losses are the assets, and the comparison is
sharp enough to be the argument:

\begin{center}
\begin{tabular}{llll}
\hline
doubling & dimension & what is lost & what the loss becomes \\
\hline
$\mathbb{R} \to \mathbb{C}$ & $1 \to 2$ & the order of the reals & phase and interference; amplitudes \\
 & & & that cancel; the Born rule (\S\ref{sec:self-reference}) \\
$\mathbb{C} \to \mathbb{H}$ & $2 \to 4$ & commutativity & non-commuting actions: weak isospin \\
 & & & $SU(2)$, the $B$ register (\S\ref{sec:register}) \\
$\mathbb{H} \to \mathbb{O}$ & $4 \to 8$ & associativity & the associator: colour, chirality, the \\
 & & & generations, and mass as debt \\
 & & & (\S\ref{sec:generations}, \S\ref{sec:masses}) \\
$\mathbb{O} \to \mathbb{S}$ & $8 \to 16$ & division itself & no further force: zero divisors appear \\
 & & (zero divisors) & and the ladder of sectors ends (below) \\
\hline
\end{tabular}
\end{center}

Walking the rows: the order lost at the first doubling is the possibility of interference --- parts
of a probability that can oppose one another --- which is where the derivation of \S\ref{sec:self-reference}
lives. The commutativity lost at the second is the existence of actions whose order matters; the
unit quaternions form the group $SU(2)$, and the framework houses weak isospin there, as the
internal algebra of the $B$ register whose order-sensitive, non-commuting dynamics the interaction sections develop
(\S\ref{sec:register}, \S\ref{sec:entanglement-classes}). The associativity lost at the third ---
the octonions retain the weaker property of alternativity \autocite{baez2002octonions} --- creates
the associator $[a,b,c] = (ab)c - a(bc)$, and the associator is where most of this paper's physics
is transacted: the stabiliser of a preferred complex direction inside the octonions' automorphisms
is the colour group (\S\ref{sec:hopf-justification}); the triality structure and its $K_3$
obstruction generate the three generations and the sterile channel (\S\ref{sec:generations}); and
the associator is the currency of mass --- the debt against the selected direction that
\S\ref{sec:masses} prices, owed by whatever action moves the selection and not by residence at
the rung, which is the defusal \S\ref{sec:register} makes explicit. The algebraic facts
in the third column are theorems; the identifications in the fourth are established in the sections
cited, each carrying its own epistemic flag there, and are not re-litigated here. What this table
asserts on its own account is only the correspondence's shape: one law surrendered per rung, one
sector of physics gained per law.

\subsection{Internal algebras, not a new quantum mechanics}

A reader versed in the quaternionic-quantum-mechanics literature will at this point raise a fair
objection: quantum theory over $\mathbb{H}$ as a foundation is known to be troubled --- the
tensor-product and composite-system difficulties set out by Adler and others
\autocite{adler1995quaternionic} --- and quantum theory over $\mathbb{O}$ founders earlier still,
since a non-associative algebra does not furnish a Hilbert space at all. The framework's position
must therefore be stated exactly, because it does not make the claim the objection targets. The base
of the physical Hilbert space here is $\mathbb{C}$ and remains $\mathbb{C}$; the higher algebras
enter as the \emph{internal} algebras of the register --- $\mathbb{H}$ for the $B$ qubit,
$\mathbb{O}$ for the $C$ qubit (\S\ref{sec:register}) --- in a tensor structure whose outer theory is
ordinary complex quantum mechanics. This is the same move by which the division algebras enter the
constructions of Furey and the related programmes without any proposal to rebase quantum mechanics
itself \autocite{furey2018three}. For the octonionic register the point is sharper still:
$\mathbb{O}^{2}$ is not a Hilbert space, and the framework accordingly treats the $C$ qubit not in
the Dirac--von Neumann manner but geometrically, as the sphere $S^{15}$ with its Hopf structure ---
which is precisely the form in which \S\ref{sec:hopf-justification} needs it. Nothing in this paper therefore
stands or falls with quaternionic or octonionic quantum mechanics as foundations; the ladder is
climbed by the registers' internal algebras, not by the theory's base field.

\subsection{The same ladder, twice}

The ladder appears in this framework in two guises, and it will help the reader to have them
distinguished now. Statically, it is the anatomy of a single relation: the three qubits of the
register take their values one rung apart --- $A$ over $\mathbb{C}$, $B$ over $\mathbb{H}$, $C$ over
$\mathbb{O}$ --- so that every particle carries the whole ladder as its internal structure
(\S\ref{sec:register}). Dynamically, it is the shape of measurement: a single qubit's pure states
fill $S^{3}$, an entangling measurement carries the joint state into $S^{7}$, and a third party into
$S^{15}$ --- the correspondence between qubit entanglement and the Hopf fibrations that
\S\ref{sec:hopf-justification} takes from a mature literature --- so that each act of measurement ascends one rung --- the climb the reader has already made,
narrated in \S\ref{sec:wigners-friend} --- and every completed relation is a terminated climb, which is
the reading \S\ref{sec:hopf-information-bound} gives the tower's endpoint. One ladder, seen twice:
as the anatomy of a relation and as the dynamics of relating. The framework regards this double
appearance not as a coincidence but as the same structure viewed statically and in motion, and the
sections cited carry that identification's weight.

\subsection{Hurwitz closes the ladder}

The construction itself would double forever; what stops it is a classification theorem. By
Hurwitz's theorem, the normed division algebras --- those in which the norm composes,
$\lvert ab\rvert = \lvert a\rvert\,\lvert b\rvert$, and no non-zero elements multiply to zero ---
are exactly $\mathbb{R}$, $\mathbb{C}$, $\mathbb{H}$ and $\mathbb{O}$
\autocite{hurwitz1898composition,baez2002octonions}. The next doubling, the sedenions, has zero
divisors: non-zero elements whose product vanishes. The framework's reading of this failure is the
one that matters physically. Zero divisors cannot support consistent information flow around loops
--- a cycle of compositions can silently annihilate what it carries --- and a framework in which
closed computational paths are load-bearing (\S\ref{sec:hopf-justification},
\S\ref{sec:space-of-computations}) cannot be built on an algebra that permits it. At the level of
amplitudes the same failure reads: the composition law
$\lvert ab\rvert = \lvert a\rvert\,\lvert b\rvert$, which carries probability multiplicatively
through every rung of the ladder, fails beyond the octonions, so that two non-vanishing amplitudes
could compose to nothing and probability would be destroyed by mere composition. The last algebra
fit to carry amplitudes is $\mathbb{O}$. One rung past the terminus is not entirely silent,
however: the sedenions' internal structure --- their automorphism group and their octonion
subalgebras --- is the recorded candidate arena for tagging the generations
(\S\ref{sec:interactions}), and it enters exactly as this section's scoping permits, internally
and never as a base. The terminus closes the ladder of forces; what survives past it is, at most,
flavour. This is the algebraic half of the termination; the
geometric half --- that no sphere above $S^{7}$ is parallelisable --- and the theorem joining the
two halves are the opening business of \S\ref{sec:hopf-justification}, and the epistemic reading of the same endpoint,
as a bound on the information a single relation can carry, is given in
\S\ref{sec:hopf-information-bound}.

\subsection{The honest debt}

As in the geometry section, we close by stating plainly what is proved, what is established, and what is
owed. Proved: the construction, the losses, and the termination are theorems --- the Cayley--Dickson
doubling and its ledger of surrendered laws are standard algebra, and Hurwitz's classification is
among the durable results of the subject. Established: the physical cashing of each rung --- the
Born rule at $\mathbb{C}$, weak isospin at $\mathbb{H}$, colour, generations and mass at
$\mathbb{O}$ --- is carried out in the sections cited above, each under its own epistemic flag. Owed,
and stated without disguise: the ladder as presented here is closed from above and derived at its
floor, but it is not yet derived rung by rung from below. Hurwitz forbids a further rung; the
argument of \S\ref{sec:self-reference} forces the first doubling, from the ordered arithmetic of
classical probability to amplitudes that interfere. That the same self-referential closure, applied
to a system whose amplitudes already take values at one rung, compels the next --- that relational
closure at one level forces precisely the Cayley--Dickson double of that level --- is the
framework's reading of the climb, and it is a thesis, not a theorem --- but the formal apparatus
against which it can be posed now exists, and we pose it.

\begin{openproblem}[The compulsion of the intermediate rungs]\label{op:ladder-compulsion}
Derive, within the categorical apparatus of Open Problem~\ref{op:sigma-functor}, that
self-referential closure at one rung compels the Cayley--Dickson double of that rung, so that the
intermediate steps $\mathbb{C} \to \mathbb{H} \to \mathbb{O}$ follow from the same principle
that forces the floor. The candidate route is the construction's own shape: the Cayley--Dickson
pair with its conjugation, $(a, b)^{*} = (a^{*}, -b)$, is an involution-equipped doubling, and the
route reads it as the algebraic shadow of adjoining the describer --- closure of level-$A$
descriptions under the inside/outside exchange requiring exactly the twisted pair. The problem
inherits the status of Open Problem~\ref{op:sigma-functor}, whose apparatus it presupposes.
\end{openproblem}

If the derivation lands, the whole ladder becomes a consequence of the same self-reference that
opened it. Until then, the honest statement is that the ladder's necessity is established at its
ceiling and at its floor, while the compulsion of the intermediate rungs is inherited from the
physics they are shown to carry rather than derived in advance. What this section hands on is, accordingly, an algebraic ladder with a floor, a ceiling and a
ledger. Its first consumer is immediate: \S\ref{sec:space-of-computations} spends the terminal
fact --- the non-associativity of $S^{7}$'s own multiplication --- as geometry, excluding the
seven-sphere as a ruler and fixing the dimension of space at three.
\S\ref{sec:hopf-justification} then supplies the geometry of the fibrations themselves, and the
join.

%      % EDITED 2026-09-09 scripts/patch_spekkens_citekey_2026-09-09T1810.py
% =====================================================================
%  SECTION: Why Hopf Fibrations Belong in Particle Physics
%  (sec:hopf-justification). Drafted 2026-06-08; F1 arc added
%  prop:f1-dictionary and op:hopf-frobenius's discharge clause.
%  AUDITED 2026-07-06 (final pass, eighth stop).
%
%  PLACEMENT (reordered chain, JS 2026-07-06): position 8 --
%      \input{spacetime_from_non-computability}
%      \input{hopf_justification_section_2026-06-08}
%      \input{entanglement_classes_and_gauge_groups_2026-06-08}
%  No adjacency phrasing in the prose: position-portable. All
%  cross-refs live -- sec:fanout-direction DRAFTED 2026-07-06 (formerly the one undrafted
%  section, by design).
%
%  AUDIT RECORD: prop:gsm, prop:termination and prop:f1-dictionary
%  verified against every quoting consumer (statements canonical);
%  honest-debt subsection names both the selection assumption and the
%  mereology debt; the epistemic two-climbs disambiguation and the
%  compositional rung-evidence verified; sec:canonical-orbits live;
%  file pure ASCII. Fixes: op:hopf-frobenius's cross-link with
%  op:sigma-functor made TWO-WAY (the "four ways" line now names the
%  op); both remaining functor conditions ANCHORED --- the quaternionic
%  to the established fibre physics of sec:wf-reconciliation (residue:
%  clutching map vs strong complementarity), the octonionic to the
%  proved canonical-orbit images (umbrella label app:ghz-w-proofs
%  minted at the appendix for the purpose). Statuses unchanged.
%  British English; \autocite; no bullets in prose.
% =====================================================================

\section{Why Hopf Fibrations Belong in Particle Physics}
\label{sec:hopf-justification}

Invoking Hopf fibrations in a discussion of the Standard Model invites an immediate and fair
objection: that the move is decorative, an analogy between two pieces of attractive mathematics
with no load-bearing connection between them. The purpose of this section is to dispose of that
objection before the construction is put to work. The connection is not an analogy. It is
carried by two independent and independently published theorems, and the structures it relies
upon are not selected for convenience but compelled by a pair of classification results. We take
as already established the correspondence between qubit entanglement and the Hopf fibrations
themselves --- the subject of a mature literature
\autocite{mosseri2001hopf,bernevigchen2003geometry} --- and do not reproduce it here. What this
section addresses is the prior question of why that correspondence is entitled to carry physical
content at all.

\subsection{Two theorems, joined end to end}

The bridge from qubit entanglement to the Standard Model is built from two distinct results,
neither of which is original to the present framework and neither of which depends on the other.

The first is geometric. The unit sphere in the Hilbert space of $n$ qubits is the total space of
a Hopf fibration: $S^3 \to S^2$ for one qubit, $S^7 \to S^4$ for two, and $S^{15} \to S^8$ for
three. For two and three qubits the fibration is \emph{entanglement-sensitive} --- separability
is precisely the condition under which the base degenerates to a product of single-qubit Bloch
spheres and the additional dimensions associated with the higher imaginary directions collapse
\autocite{mosseri2001hopf,bernevigchen2003geometry}. This is a statement about the geometry of
state spaces and requires no physical interpretation whatever.

The second is algebraic. Singling out a preferred complex structure on the octonions ---
equivalently, fixing the split $\mathbb{O} \cong \mathbb{C} \oplus \mathbb{C}^3$ --- leaves a
residual symmetry, and that residual symmetry is exactly the gauge group of the Standard Model.
This was established through the exceptional Jordan algebra by Dubois-Violette and Todorov
\autocite{duboisviolette2016exceptional,duboisviolettetodorov2019exceptional,todorovduboisviolette2018deducing}
and given its sharpest spinorial form by Krasnov \autocite{krasnov2021so9}, who showed that the
relevant group is the stabiliser of a complex structure inside $\mathrm{Spin}(9)$. We record the
result in the form in which it enters the present work.

\begin{proposition}[Dubois-Violette--Todorov; Krasnov]\label{prop:gsm} Let $J_R$ be the complex
structure on the space $\mathbb{O}^2$ of $\mathrm{Spin}(9)$ spinors determined by a choice of
unit imaginary octonion, equivalently the split
$\mathbb{O} \cong \mathbb{C} \oplus \mathbb{C}^3$. The subgroup of $\mathrm{Spin}(9)$ commuting
with $J_R$ is \[ G_{\mathrm{SM}} \;=\; \frac{SU(3) \times SU(2) \times U(1)}{\mathbb{Z}_6}, \]
acting on the singled-out $\mathbb{C}$ summand as a colour singlet of hypercharge $-1$ and on
the $\mathbb{C}^3$ summand as a colour triplet of hypercharge $\tfrac{1}{3}$, both as $SU(2)$
doublets --- the quantum numbers of one left-handed generation.
\end{proposition}

The two results are independent, and the join between them is the only place a physical theory
must do work. Each result is a theorem; what neither supplies is a reason why nature should fix
one complex direction within $\mathbb{O}$ rather than treating all directions symmetrically. We
return to this --- the genuine load-bearing assumption --- in \S\ref{sec:hopf-honest-debt}, and
resolve it, in candidate form, in \S\ref{sec:fanout-direction}.

\subsection{The ladder is forced, not chosen}

It might still be objected that selecting the octonions, or the Hopf fibrations, is itself a
modelling choice carrying as much latitude as the choice of a grand-unified group. It is not,
and the reason is a pair of classification theorems that close off the alternatives entirely.

By Hurwitz's theorem the only normed division algebras are $\mathbb{R}$, $\mathbb{C}$,
$\mathbb{H}$ and $\mathbb{O}$ \autocite{hurwitz1898composition,baez2002octonions}. By Adams's
resolution of the Hopf-invariant-one problem --- together with the parallelisability results of
Bott--Milnor and Kervaire --- the only parallelisable spheres are $S^0$, $S^1$, $S^3$ and $S^7$,
so that there exist exactly three nontrivial Hopf fibrations, in bijection with $\mathbb{C}$,
$\mathbb{H}$ and $\mathbb{O}$
\autocite{adams1960hopf,bottmilnor1958parallelizability,kervaire1958nonparallelizability}. These
are not heuristics; they are among the more durable theorems in topology and algebra. Their
force here is that, once a framework is committed to amplitudes valued in a normed division
algebra --- a commitment that in the present construction is made far upstream, by the
self-referential argument of \S\ref{sec:self-reference} and \S\ref{sec:division-algebras}, and
not by fiat at this stage --- the ladder of Hopf fibrations is not selected but compelled, and
it terminates of its own accord.

\begin{proposition}[Termination]\label{prop:termination} The Cayley--Dickson tower of normed
division algebras and the tower of parallelisable spheres both terminate at real dimension
eight. Consequently there are exactly three nontrivial Hopf fibrations, the largest being the
octonionic $S^{15} \to S^8$, and the construction admits no extension to a fourth case.
\end{proposition}

The termination is worth dwelling on, because it is unusual for a model of this kind to predict
its own boundary. The next rung of the Cayley--Dickson sequence, the sedenions, possesses zero
divisors; the Hopf map can no longer be globally defined; there is no fourth fibration. This is
not a limitation imposed from outside but a structural fact, and it carries a physical
consequence that we develop later: there is no room in the construction for a fourth gauge force
or a fifth interaction, so that the discovery of either would count as evidence against the
framework rather than as an invitation to extend it. The endpoint is a theorem, not a
stipulation.

\subsection{Independent arrival as corroboration}

The most persuasive evidence that the destination is not an artefact of the route is that it is
reached by routes sharing none of the same machinery. Szangolies
\autocite{szangolies2025standardmodel} reaches $G_{\mathrm{SM}}$ from the emergent-gravity and
Kaluza--Klein programmes, reading the singled-out complex direction as a dimensional reduction
of a $9+1$-dimensional spacetime; the present framework reaches it from self-reference
(\S\ref{sec:self-reference}); and Furey and Gillard--Gresnigt reach the same gauge group, and
the same three-generation structure, through the minimal left ideals of the Dixon and Clifford
algebras, without invoking Hopf fibrations at all
\autocite{furey2018three,gillardgresnigt2019three}; the relation between their minimal ideals
and the entanglement classes used here is taken up in \S\ref{sec:canonical-orbits}. The
octonionic route to particle physics is in any case an old one, going back to the use of
$\mathbb{O}$ to model the strong force \autocite{gunaydingursey1974quark}.

That programmes with disjoint starting points converge on
$SU(3) \times SU(2) \times U(1)/\mathbb{Z}_6$ and on the same quark--lepton division is the
coincidence Szangolies himself remarks upon: the characteristic structure of the Standard Model,
routinely regarded as arbitrary, turns up independently in unrelated settings, and its appearing
twice would then be doubly odd unless the two settings are in fact related. We read the
convergence the same way --- as evidence that the division-algebraic description is singled out
by the mathematics rather than imposed by the modeller, and therefore that
Proposition~\ref{prop:gsm} is structural rather than contrived.

\subsection{The honest debt}
\label{sec:hopf-honest-debt}

It would be misleading to leave the impression that this route is free of assumptions, so we
state plainly the one assumption that all of these programmes share. Proposition~\ref{prop:gsm}
requires that \emph{some} complex direction in $\mathbb{O}$ be singled out, and the uniqueness
of that choice is not supplied by the theorem. Krasnov parametrises the complex structure $J_R$
by an arbitrary unit imaginary octonion, and the residual group is the same for every such
choice because $G_2 = \mathrm{Aut}(\mathbb{O})$ acts transitively on the unit imaginary
octonions. The theorem therefore guarantees the gauge group \emph{given} that a direction is
fixed, while remaining silent on \emph{why} one is fixed. Szangolies treats this as a free
choice and flags it, in his own conclusions, as an open element of the construction --- one to
which several inequivalent answers are available.

This, and not the appeal to Hopf fibrations, is the genuine load-bearing assumption of the whole
approach. The Hopf fibrations are forced (Proposition~\ref{prop:termination}); the singling-out
is not. The present framework offers a candidate resolution that the prior literature lacks: the
act of measurement, modelled as \textsc{fanout}, selects the direction physically, and the
transitivity of $G_2$ then guarantees that the gauge group is independent of which direction the
measurement happens to select. This argument is given in \S\ref{sec:fanout-direction}; we flag
here, at the appropriate epistemic status, that it is structural rather than fully proved, and
that the residual question --- why a measurement selects \emph{exactly one} imaginary direction
rather than merely \emph{permitting} a selection --- remains open. It is now priced as Open
Problem~\ref{op:fanout-selection}.

A second debt is shared with every programme of this character and is not closed by anything
said here: why the abstract quantum state should admit a tensor-product decomposition into
qubits in the first place. This is the quantum factorisation, or mereology, problem. The
entanglement-sensitivity on which the construction turns is meaningful only once such a
decomposition is fixed, and we treat the decomposition as a necessary input rather than a
derived result. Naming it is preferable to letting it pass unremarked.

\subsection{The termination read epistemically}
\label{sec:hopf-information-bound}

The termination of the tower (Proposition~\ref{prop:termination}) has so far been stated
algebraically: the Cayley--Dickson sequence acquires zero divisors at the sedenions and the Hopf
map can no longer be globally defined. There is a second reading of the same boundary, epistemic
rather than algebraic, which we offer as a motivated conjecture rather than a settled result,
because it connects the tower's endpoint to a bound that arises independently in the foundations
of measurement.

A particle, in this framework, is a reified relation between events, and its internal content is
carried by the octonionic $A,B,C$ register --- the $S^{7}$ sector serving as internal degrees of
freedom in the channel that connects the events (\S\ref{sec:space-of-computations}). The
termination of the tower is then a statement about that content: a single such relation can
support genuine entanglement among at most three parties, because there is no fourth fibration
to carry a fourth independent party within the same irreducible relation. Every particle is, on
this reading, a \emph{completed} octonionic-level relation --- topped out at three parties not
by circumstance but because the tower admits no further rung. The information that can be
encoded in one relation-between-events, and hence extracted by one measurement of it, is
correspondingly bounded: it is the information of a completed three-party register, and no
measurement of a single particle can yield more.

This is exactly the kind of bound that the epistemic-restriction programmes reach from the other
direction. Spekkens' toy theory is built on a knowledge-balance principle that caps the
information obtainable about a system \autocite{spekkens2007toy}, and Szangolies derives such an
epistemic horizon not by positing it but from the same self-referential non-computability on
which the present tower rests \autocite{szangolies_j_2018}. That the two accounts --- the
algebraic termination of the Hopf tower and the epistemic bound on measurement information ---
should coincide is, we suggest, not a numerical accident but a shared root: both are
consequences of the self-referential non-computability of \S\ref{sec:self-reference}, wearing an
algebraic face in the one case and an epistemic face in the other. The Hurwitz--Adams ceiling
and the knowledge-balance ceiling would then be one limit seen twice, and Szangolies' programme,
which reaches $G_{\mathrm{SM}}$ by an independent route (\S\ref{sec:hopf-honest-debt}) and the
epistemic horizon from the same non-computable source, shares with this framework not merely a
conclusion but its origin.

The reading must be kept distinct from a different climb, with which it would otherwise be
conflated. The Wigner's-friend regress of \S\ref{sec:wigners-friend} ascends without end: each
observer may be measured by a further observer, each carrying a fresh gauge-private circle, and
that chain has no top --- it is the bump in the carpet, unbounded in length. The termination
bounds something else: not the length of the chain of observers, but the internal structure of
each single relation, each particle, that the chain is built from. The regress is unbounded
precisely because each of its links is a bounded, completed three-party object;
termination-of-each-link and non-termination-of-the-chain are not in tension but complementary,
the former individuating the links of the latter. What the bound caps is the information in one
particle; what the regress leaves uncapped is the number of particles and observers that can be
chained. The precise identification of ``the information of a completed three-party register''
with the formal measures of the epistemic-restriction literature is the technical step that
would turn this conjecture into a theorem, and we flag it as owed; the structural claim --- that
the tower's endpoint and the measurement bound share the non-computable root --- is what we
advance here.

\subsection{The termination read compositionally}
\label{sec:hopf-compositional}

There is a third termination theorem, proved in a community that has so far had no contact with
the two joined above. Coecke and Kissinger showed that the GHZ and W classes --- realised as the
commutative special and the anti-special Frobenius algebras of categorical quantum mechanics ---
are \emph{compositionally complete}: every $N$-qubit entanglement class under SLOCC is generated
by composing these tripartite primitives with bipartite entanglement and product states
\autocite{coeckekissinger2010compositional}. Above three parties nothing new appears; the
entanglement alphabet closes at exactly the point where the fibration tower terminates. The
mutual silence of the two literatures is itself striking: Coecke and Kissinger make no reference
to Hopf fibrations or division algebras, and the fibration literature on entanglement
\autocite{mosseri2001hopf,bernevigchen2003geometry} nowhere cites compositional completeness.
Two classifications, in different mathematics, draw the same line in the same place, unaware of
one another. That silence has, in the months this paper was in draft, just begun to break from
the categorical side, with a projective, base-level reading of the same calculus
\autocite{burtonanwar2026meromorphic}; we engage it below.

The framework's conjecture is that this is one termination wearing a third face: the alphabet
has the letters it has --- Bell, GHZ, W --- because the parallelisable spheres have the
dimensions they have. There is partial evidence at every rung. At one qubit, the commutative
special Frobenius algebras on $\mathbb{C}^{2}$ are classified by orthonormal bases --- points of
$S^{2}$, the base of the first fibration --- so that the copied basis plays the role of a
section of the bundle and the spider theorem expresses its idempotence. At two qubits, the
$SU(2)$ fibre of $S^{7} \to S^{4}$ carries exactly the entanglement that survives local
operations, and L\'{e}vay's result that the bundle's natural connection is the $SU(2)$
Yang--Mills instanton \autocite{levay2004geometry} stands on the geometric side where the
interaction of two Frobenius structures stands on the categorical one. At three qubits, the
split of D\"{u}r, Vidal and Cirac \autocite{durvidalcirac2000three} mirrors the algebra exactly:
GHZ is genuinely tripartite, demanding the full non-associative octonion, while W is
pairwise-robust and housed in the associative subalgebra $\mathbb{H} \subset \mathbb{O}$ --- the
special/anti-special Frobenius distinction in geometric dress.

The first of the four conditions the correspondence requires can be discharged at once, and we
do so here, because a problem whose first step has been walked is a different kind of invitation
from one merely posed.

\begin{proposition}[The one-qubit dictionary]\label{prop:f1-dictionary} The $\dagger$-special
commutative Frobenius algebras on $\mathbb{C}^{2}$ correspond bijectively to
\emph{lifted antipodal skeletons} of the complex Hopf fibration $\pi \colon S^{3} \to S^{2}$:
unordered antipodal pairs $P \subset S^{2}$ together with a section of $\pi$ over $P$. The
copyable states of a structure are exactly the section's image, so the lifts are recorded
faithfully --- the projectively trivial centre of $U(2)$ acts non-trivially on structures ---
and the correspondence is $U(2)$-equivariant, the ordered structures forming a single free and
transitive $U(2)$-orbit. The channel a structure induces on density matrices is dephasing in its
skeleton and is independent of the section: the fibre freedom, a torsor for
$\mathrm{Map}(P, U(1))$, is exactly the data absent from the classical record a structure
induces.
\end{proposition}

\begin{proof}
By the classification of Coecke, Pavlovic and Vicary, such structures correspond to orthonormal
bases, recovered as the sets of copyable states \autocite{coeckepavlovicvicary2013orthogonal}.
An orthonormal basis of $\mathbb{C}^{2}$ is two orthogonal rays --- an antipodal pair on the
Bloch sphere --- together with a unit representative of each, which is precisely a section of
$\pi$ over the pair; and $\delta(e^{i\theta}x) = e^{i\theta}\, x \otimes x$ against
$(e^{i\theta}x) \otimes (e^{i\theta}x) = e^{2i\theta}\, x \otimes x$ pins the phases.
Equivariance, the action of the centre, the orbit structure and the induced channel are direct
computations; an executable verification of every finite check accompanies the code release.
\end{proof}

Two remarks fix the proposition's scope. It corrects a natural first reading: a classical
structure does not project to the base but selects lifts into the total space --- classicality
is section data, and what the copy operation broadcasts is a lift, not a ray. And it claims
nothing topological about cloning: the finiteness of the copyable set is linear, structural
no-cloning residing in non-naturality \autocite{abramskyheunen2012hstar}, so the first place the
bundle's global topology can bite is the quaternionic rung, where the degree-one clutching map
stands ready to be compared with the strong-complementarity law governing an interacting pair of
structures. The dictionary is, finally, the complement of the projective treatment cited above:
Burton and Anwar pass functorially to the base and deliberately discard the fibre
\autocite{burtonanwar2026meromorphic}, and the classical structures turn out to live in exactly
the data their projectivisation forgets --- the two readings compose, and the resulting triangle
is a consistency condition available to constrain the remaining conditions of the
correspondence.

\begin{openproblem}[The Hopf-fibration--Frobenius correspondence]\label{op:hopf-frobenius} Let
$\mathbf{Hopf}$ denote the category of the three fibrations
$S^{2^{k}-1} \hookrightarrow S^{2^{k+1}-1} \to S^{2^{k}}$, $k = 1, 2, 3$, with bundle maps as
morphisms, and let $\mathbf{Frob}$ denote the Frobenius algebras on $\mathbb{C}^{2}$ in
$\mathbf{FdHilb}$, with the commutative-special (GHZ) and anti-special (W) classes
distinguished. Construct a functor $F \colon \mathbf{Hopf} \to \mathbf{Frob}$ --- or prove that
none exists --- under which the complex fibration yields the single-qubit classical structures
--- Proposition~\ref{prop:f1-dictionary} identifies these as the lifted antipodal skeletons,
section data over two base points rather than the points alone; the quaternionic fibration
yields the interaction of two such structures, with the $SU(2)$ fibre carrying the entanglement
that survives the classical quotient; the octonionic fibration yields the GHZ/W pair, GHZ on the
full non-associative fibre and W on the associative $\mathbb{H} \subset \mathbb{O}$; and the
composition by which Coecke--Kissinger generate the $N$-qubit classes corresponds to a
well-defined geometric operation on the bundles. If such an $F$ exists, compositional
completeness follows from Adams's theorem as a topological corollary --- there is no fourth
fibration to supply a fourth primitive --- and if $F$ is faithful enough that the implication
also runs back, the termination of the tower and the completeness of the alphabet are equivalent
expressions of a single classification. The first condition is discharged by
Proposition~\ref{prop:f1-dictionary}; what remains open is the quaternionic and octonionic
conditions and, above all, the correspondence of compositions. Both open conditions are already
constrained: the quaternionic fibre's physical content --- the recoverable relative phase and
its purity ledger --- is established (\S\ref{sec:wf-reconciliation}), leaving the
clutching-map-to-strong-complementarity comparison as the structural residue; and any candidate
functor must reproduce the proved canonical-orbit images, GHZ to the colour triplet and W to the
singlet (\S\ref{sec:canonical-orbits}, Appendix~\ref{app:ghz-w-proofs}). A proof that no such
functor exists, or a correspondence at the level of the classifications without functoriality on
compositions, would each also resolve the problem: the first would make the agreement of two
independent classifications a fact demanding its own explanation; the second would establish the
weaker reading.
\end{openproblem}

A resolution in the affirmative would do three things. It would explain compositional
completeness rather than merely prove it --- a corollary of the division-algebra structure of
quantum mechanics rather than an independent algebraic fact. It would give the diagrammatic
language of this paper's own interaction circuits a topological warrant: the GHZ/W spider
calculus would be the unique language the topology permits, not a convenient notation. And it
would open a two-way bridge, translating problems on compositions of Frobenius structures into
operations on bundles and back. The problem sits at the intersection of categorical quantum
mechanics, the quantum-information geometry of the fibrations, and the division-algebra
approaches to particle physics, and it is posed here so that any of the three communities can
take it up in its own language. It is perhaps worth noticing that both classifications already
answer to the same name --- the fibrations are Heinz Hopf's, and so, through the bialgebra
structure of the diagrammatic calculus, are the Hopf algebras --- a coincidence of eponymy that
the conjecture proposes to promote to a theorem. With the epistemic reading above, an
affirmative resolution would leave the tower's endpoint as one limit seen four ways: algebraic,
geometric, compositional, epistemic --- the last being the fixed-point classification of Open
Problem~\ref{op:sigma-functor}, with which this problem is cross-linked, each constraining the
other.

\subsection{What this sets up}

The entanglement-sensitivity of the Hopf maps is not a side feature but the hinge of the
physical reading that follows. Because separability is exactly the collapse of the base to a
product of Bloch spheres, the structure of the entanglement --- the SLOCC class of a multi-qubit
state, and within it the canonical maximally entangled orbit --- is carried by precisely the
degrees of freedom on which $G_{\mathrm{SM}}$ acts in Proposition~\ref{prop:gsm}. This is what
makes it possible, in \S\ref{sec:entanglement-classes}, to read the two genuinely tripartite
entanglement classes onto the $\mathbb{C} \oplus \mathbb{C}^3$ split through their canonical
representatives: the GHZ orbit confined to the colour-triplet $\mathbb{C}^3$ sector, the W orbit
reaching the $\mathbb{C}$ singlet, recovering the quark--lepton division not as an assignment
but as a consequence of the entanglement structure. The scope of those statements --- canonical
orbits rather than whole classes, and the identification of the canonical orbits with the
physical particles --- is made precise there. The Hopf fibrations are the bridge; the
entanglement classes, treated next, are the traffic that crosses it.

%
%  SOURCES (re-read 2026-07-04): lawvere_inside_outside_fixed_point_
%  theorem_2026-05-04T1806.md (the spine: the two-principles reduction,
%  the theorem, the sigma candidate, the open items -> op minted per
%  D9); research_summary_neg_prob.md (negativity route; lamp (1+i)/2);
%  dynamics_companion_fanout_unitarity (legs 1a/1c; anti-events; the
%  Brukner-Zeilinger rebuttal per D6); Pure_states_and_mixed_states
%  (local tomography derived; QBism contrast per D7).
%
%  DECISIONS ENCODED (JS 2026-07-04): D1 title as below; D2 Wheeler1990
%  cited here as the title anchor; D3 lineage Popper -> Breuer ->
%  Svozil -> Szangolies (Dalla Chiara omitted); D4 Frauchiger_2018
%  cited here, dissolved in sec:wigners-friend; D5 Kirkwood-Dirac
%  heritage via the Arvidsson-Shukur 2024 review as one-stop; D6
%  Brukner-Zeilinger + the framework rebuttal IN ("Shannon is
%  post-FANOUT"); D7 QBism contrast IN; D8 Aaronson private
%  communication OMITTED (published island paper only); D9 one op
%  minted (op:sigma-functor), cross-linked to op:hopf-frobenius.
%
%  CITATIONS -- existing keys: Wheeler1990, small_2006, Godel1931,
%  Lawvere_1969, Yanofsky2003, CoeckeKissinger2017, Rovelli1996a,
%  Popper1950a, Popper1950b, Breuer1995, svozil_k_2018,
%  szangolies_j_2018, Frauchiger_2018, dirac1942, Feynman_r_82,
%  Feynman1987, Youssef1994 [SWAP-FLAG: intended paper is "Quantum
%  mechanics as complex probability theory", Mod. Phys. Lett. A 9
%  (1994) 2571 -- if that lives under youssef1994quantum, swap],
%  caticha_a_2006, goyal_knuth_skilling_2010, Aaronson2004h, Hardy2001a,
%  muller_p_2011, Thomson_1954, benacerraf1962tasks, cantor1891diagonal,
%  Turing_1937.
%  NEW keys (RIS in self_reference_refs_2026-07-04.ris, all verified
%  this session): tikochinsky1988feynman, renou2021real, chen2022ruling,
%  wu2022refutation, hoffreumonwoods2026real, arvidssonshukur2024kirkwood,
%  bruknerzeilinger2001shannon, fuchsschack2013qbism,
%  shannon1948mathematical.
%
%  CROSS-REFS: forward (fine under two-pass LaTeX): sec:continuity,
%  sec:division-algebras, sec:hopf-justification, prop:f1-dictionary,
%  sec:wigners-friend, op:limit-realisability, op:hopf-frobenius.
%  Still dangling after this section: sec:fanout-boundary (and its
%  sibling label sec:fanout-direction) -- both DRAFTED 2026-07-06.
%
%  AUDITED 2026-07-06 (section-by-section final pass, first stop):
%  proof of prop:inside-outside re-verified by hand; attribution chain
%  clean (Burgin/Khrennikov correctly absent); all header-declared
%  citations confirmed used; lamp arithmetic and antecedents verified;
%  FANOUT smallcaps consistent; no pre-definition use of "frame".
%  Fixes: stray Unicode em-dash (QBism paragraph, the file's only
%  non-ASCII character) -> house ---; trace-blind two-way stitch added
%  to the Born subsection ("one privacy, two duties");
%  op:sigma-functor gains its first standing constraint (F1
%  inheritance via the op:hopf-frobenius discharge) and its previously
%  one-directional local-tomography cargo made explicit.
% =====================================================================

\section{Self-Reference: One Principle}
\label{sec:self-reference-2}

Wheeler asked for it from bit: for physics derived from information, ``every it --- every particle,
every field of force'' obtained from binary answers to yes-or-no questions \autocite{Wheeler1990}.
This paper's route to that programme runs through G\"{o}del: the single principle from which
everything else here is derived, deferred, or priced is \emph{self-reference} --- the fact that any
complete physical description, made from inside the world it describes, must describe the describer.
This section states the principle, proves that the framework's second historical commitment follows
from it as a theorem rather than standing beside it as an axiom, and then extracts the price of the
principle in four instalments: probabilities whose parts can cancel, amplitudes valued in
$\mathbb{C}$, the Born rule, and the local tomography of composite systems. What the section cannot
yet derive is stated as precisely as what it can, in one numbered problem and in the ledger at the
end.

\subsection{Events, not objects}

The framework's ontology is relational: the fundamental entities are events and the relations
between them, and objects are the bookkeeping of relational invariants --- what remains stable
across many relations. The kinship with relational quantum mechanics \autocite{Rovelli1996a} is
genuine and acknowledged: facts are facts relative to the systems that register them. The
difference, developed over the rest of this paper, is that the present framework grounds the
relational structure in a single logical principle and then carries it, through the division
algebras and the Hopf fibrations, all the way to the gauge group and the particle spectrum ---
relations are not only the ontology but the derivation engine. On this reading a measurement is not
the revelation of a pre-existing property; it is the creation of a relation, and the later sections
of this paper read that creation geometrically (\S\ref{sec:hopf-justification}) and even
gravitationally. Here we need only the logical core: whatever describes the world is an event in
the world it describes.

\subsection{One principle, and a theorem}

The framework historically carried two foundational commitments: self-reference, and the
\emph{Inside/Outside Equivalence Principle} --- that the laws of physics must be identical whether
formulated by an embedded participant or by an external observer. Two principles are one too many,
and this subsection reduces them.

The reduction corrects the framework's own origin, and we make the correction in print. The 2005
form of the argument \autocite{small_2006} ran a G\"{o}delian chain: a self-referential process,
being complete, must by the incompleteness theorem \autocite{Godel1931} be inconsistent; an
inconsistent system can negate itself; and the negation of self-reference is the production of
causal sequence. The chain captures the right intuition and is retained here as motivation, but it
is not sound as a proof: classical inconsistency yields \emph{ex falso quodlibet} --- every
proposition provable, every causal sequence an output, no constraint at all. The load-bearing
argument is different, and better: a fixed-point theorem.

Let $\mathcal{C}$ be the category of physical descriptions and admissible processes --- concretely,
the dagger-compact-closed setting with classical structures in which categorical quantum mechanics,
and the whole of this paper, works \autocite{CoeckeKissinger2017}. Lawvere showed that in a
sufficiently rich such setting, any weakly point-surjective self-representation forces every
endomap to have a fixed point, and Yanofsky showed that one diagonal argument of this shape
underlies Cantor's, G\"{o}del's, Turing's and Tarski's theorems alike
\autocite{Lawvere_1969,Yanofsky2003,cantor1891diagonal}. The framework's single commitment is that
$\mathcal{C}$ contains a \emph{self-representing object} --- an object capable of encoding the
processes on itself. That is what ``self-reference is foundational'' means, stated categorically,
and it is the surjectivity hypothesis of the diagonal lemma. Now let
$\sigma \colon \mathcal{C} \to \mathcal{C}$ be the \emph{inside/outside involution}: the operation
exchanging the two complementary presentations of any description --- the inside presentation, as
participant-indexed total space in which the observer is part of the system, and the outside
presentation, as a base-times-fibre decomposition over a classical parameter space. That
$\sigma^{2} = \mathrm{id}$ is not an axiom but a coherence requirement: the outside view of the
inside view of a thing is the thing. The candidate concrete realisation of $\sigma$ is the
Hopf-duality swap --- total space against base-times-fibre for the three fibrations --- which gives
the involution computable content and ties this subsection to \S\ref{sec:hopf-justification}.

\begin{proposition}[Inside/outside fixed point]\label{prop:inside-outside}
Let $\mathcal{C}$ satisfy the richness condition above and contain a self-representing object, and
let $\sigma$ be an involution on $\mathcal{C}$. Then the fixed-point locus
$\mathrm{Fix}(\sigma)$ is non-empty: there exist descriptions invariant under the exchange of
inside and outside perspective.
\end{proposition}

\begin{proof}
Lawvere's construction, applied to $\sigma$ in place of negation: given a weakly point-surjective
$\varphi \colon A \to Y^{A}$, form $g(a) = \sigma\!\left(\varphi(a)(a)\right)$; point-surjectivity
supplies $a^{*}$ with $\varphi(a^{*}) = g$, and then $\varphi(a^{*})(a^{*})$ is a fixed point of
$\sigma$ \autocite{Lawvere_1969,Yanofsky2003}. Where the same construction applied to Boolean
negation yields paradox --- negation has no fixed point on a two-element type, so the conclusion
explodes the premises --- the inside/outside involution is an honest involution on a rich type, and
the conclusion is a constructive fixed point rather than a contradiction.
\end{proof}

\textbf{Definition.} \emph{Physical law is the structure carried by} $\mathrm{Fix}(\sigma)$.

The Inside/Outside Equivalence Principle is thereby demoted from axiom to theorem: physical law is
invariant under the exchange of perspectives \emph{because it is the fixed-point set of that
exchange}, and the fixed-point set is forced to be non-empty by self-reference alone. One
commitment remains foundational; everything the second commitment asserted is now derived. The
reader who accepts categorical quantum mechanics as an existing programme has already accepted the
machinery; the framework adds one premise and collects the rest. And the fixed-point locus is not
empty of structure: it inherits the Frobenius-algebra apparatus of $\mathcal{C}$, so that the
Bell, GHZ and W primitives --- and with them the particle classification of the later sections ---
live inside $\mathrm{Fix}(\sigma)$ from the start. What is proved above is non-emptiness given the
setting; what is assumed is the setting itself and the existence of $\sigma$ as a genuine
involution; what is open is the rigour of both, and it is priced now rather than hidden:

\begin{openproblem}[The involution made functorial]\label{op:sigma-functor}
Fix the category $\mathcal{C}$ rigorously --- the dagger-compact-closed setting with classical
structures is the candidate --- and define $\sigma$ on morphisms, not only on objects, the natural
route being the induced fibre--base exchange on spider diagrams; establish that $\sigma$ is an (at
worst pseudo-)involutive endofunctor; and prove or refute the identification of
$\mathrm{Fix}(\sigma)$ with the three Hopf cross-sections, under which Adams's theorem would become
a corollary of the fixed-point classification. This problem is cross-linked with Open
Problem~\ref{op:hopf-frobenius}: under an affirmative resolution of both, the parallelisable
terminus, the fixed-point classification, and compositional completeness are three statements of
one fact, and progress on either problem constrains the other. One constraint is already in place:
at the complex rung, Proposition~\ref{prop:f1-dictionary} supplies the dictionary between
classical structures and the fibration's lifted skeletons --- the clause of
Open Problem~\ref{op:hopf-frobenius} it discharges --- and any resolution here must reproduce it.
The problem also carries a piece of cargo recorded later in this section: the formalisation, within
the fixed category, of the perspective-accounting that derives local tomography.
\end{openproblem}

\subsection{The lineage}

The obstruction this section builds on is old, and the framework claims its completion, not its
discovery. Popper argued in 1950 that a predictor embedded in the world has irreducibly inferior
knowledge of itself, because it cannot impart information about itself without, by that very act,
rendering the information out of date \autocite{Popper1950a,Popper1950b} --- the passage is quoted
where the framework makes the obstruction structural (\S\ref{sec:wigners-friend}). Breuer proved
the impossibility of accurate state self-measurement within quantum theory \autocite{Breuer1995};
Svozil's monograph surveys the whole territory where physics meets self-reference, undecidability
and Chaitin-style algorithmic limits \autocite{svozil_k_2018}; Szangolies sharpened the modern form
as epistemic horizons \autocite{szangolies_j_2018}; and the Frauchiger--Renner theorem is the
contemporary crisis-form --- nested observers deriving contradictory records when each treats the
others' quantum descriptions as universally valid \autocite{Frauchiger_2018} --- which
\S\ref{sec:wigners-friend} dissolves relationally rather than paradoxically. This framework's
addition to the lineage is twofold. Structurally, Proposition~\ref{prop:f1-dictionary} converts
Popper's temporal obsolescence into algebraic cancellation: the datum a copier cannot put into its
own record is exactly its fibre commitment, which does not arrive stale --- it cancels. And
foundationally, the present section runs the lineage in reverse: where the tradition treats
self-reference as a limit on knowledge \emph{within} physics, the framework takes it as the
generator of physics, and derives the quantum formalism from it.

\subsection{The price is negative probability}\label{sec:complex-amplitudes-2}

What does it cost a probabilistic description to contain its own describer? The first instalment is
the sign. In the regime before any record exists, every event has an anti-event --- the description
must be able to undo what it has provisionally done, or the self-description could never be
revised into consistency --- and a probability calculus closed under undoing cannot be confined to
$[0,1]$: it is a signed calculus, in which the parts of a probability can oppose and cancel. The
framework did not discover negativity as bookkeeping; Dirac admitted negative probabilities as
formal intermediaries \autocite{dirac1942}, Feynman defended them as consistent whenever the final,
recorded probabilities stay positive \autocite{Feynman_r_82,Feynman1987}, and the Kirkwood--Dirac
distribution has carried complex quasiprobability from 1933 to the present, as a recent
comprehensive review documents \autocite{arvidssonshukur2024kirkwood}. What the framework adds is
the forcing: negativity is not a permissible convenience but the first price of self-referential
closure, and it is charged exactly and only in the pre-record regime --- which is why every
negativity in the tradition lives in intermediate, unrecorded quantities and never in a final
record.

The second instalment is complexity. A signed information measure must extend Shannon's
$S = -\ln p$ off the positive axis, and the analytic continuation of the logarithm carries $p$ into
the complex plane: amplitudes, with the phase as the continued information. Here too the extension
has a prior tradition, and the attributions are owed precisely. Tikochinsky derived Feynman's rules
for complex amplitudes from consistency of composition \autocite{tikochinsky1988feynman}; Youssef
reformulated quantum mechanics as complex probability theory outright \autocite{Youssef1994};
Caticha carried objective amplitudes to Bayesian probabilities \autocite{caticha_a_2006}; and
Goyal, Knuth and Skilling derived Feynman's rules from a pair-postulate without assuming
analyticity \autocite{goyal_knuth_skilling_2010}. The extension is prior art. The framework's
claim is the compulsion: the complex extension is not a modelling choice among exotic probability
theories but the unique continuation forced by self-referential closure --- the same forcing that
Aaronson glimpsed when he argued quantum mechanics is what negative probability inevitably becomes
\autocite{Aaronson2004h}.

Why does the continuation stop at $\mathbb{C}$? Downward, the real field is not available: real
quantum theory is experimentally distinguishable from complex quantum theory in network scenarios
and has been disfavoured by experiment
\autocite{renou2021real,chen2022ruling,wu2022refutation}; the recent objection that the exclusion
rests on locally tomographic composition \autocite{hoffreumonwoods2026real} is, for this framework,
not an escape but a confirmation of where the weight sits --- local tomography is \emph{derived}
below, not assumed, so the framework closes from its own resources exactly the door the objection
holds open. Upward, the quaternions and octonions fail as a \emph{base} for the information
continuation: the continued measure is a logarithm, and beyond $\mathbb{C}$ the logarithm of a
product is no longer the sum of logarithms --- non-commutativity severs
$\log(ab) = \log a + \log b$ --- so the additivity that makes it an information measure is lost.
Aaronson's parameter-counting and signalling arguments against a quaternionic base point the same
way \autocite{Aaronson2004h}. The higher algebras are not banished; they return as the
\emph{internal} algebras of the register, which is the business of \S\ref{sec:division-algebras}.
The base is $\mathbb{C}$, for informational reasons, and stays $\mathbb{C}$.

\subsection{The Born rule from fibre privacy}\label{sec:born-rule}

The third instalment is the quadratic rule. An amplitude carries one datum more than any record of
it can: its overall phase. This was asserted throughout the framework's development as the
observer's private, self-referential ignorance; it is now a proved statement, in exactly the form
needed. By Proposition~\ref{prop:f1-dictionary}, a classical structure --- the copy operation any
record requires --- is committed to a specific lift of its skeleton, and that commitment is absent
from every classical record the structure induces: the copier is fibre-committed, the record is
fibre-blind, and the fibre freedom is gauge for every observer. (The computation extends verbatim
from one qubit to the full register, where Proposition~\ref{prop:trace-blind} makes it the ground
of hypercharge's tracelessness --- one privacy, two duties.) Popper's obsolescence, made
cancellation. It follows that anything recordable is invariant under the fibre --- observables and
their statistics must be $U(1)$-invariant functions of the amplitude. The lowest-degree polynomial invariant
of a complex amplitude is the quadratic $|\psi|^{2}$ --- the non-polynomial invariant $|\psi|$,
of degree one, is excluded by the analyticity declared at \S\ref{sec:complex-amplitudes} --- and
normalisation fixes the coefficient:
probability is the minimal record-compatible function of the amplitude,
\[
p \;=\; |\psi|^{2},
\]
the Born rule as the cheapest invariant the fibre permits. The same square is reached from the
cost side in \S\ref{sec:born-cost-ledger}, where a record is the closed loop of a history and its
return leg and the exponent is the number of legs; neither route uses the lattice of projections
that Gleason's theorem needs, and both hold for a single qubit. The derivation is the framework's
own;
its premise --- that the fibre datum is structurally excluded from records --- is
Proposition~\ref{prop:f1-dictionary}, and its geometric reading --- probability as base-point data
under the Hopf projection --- is developed in \S\ref{sec:hopf-justification}.

One standing objection to information-theoretic foundations is answered here, because this
framework needs the answer. Brukner and Zeilinger argued that Shannon information is conceptually
inadequate for quantum measurements: it is basis-dependent, and its operational meaning presupposes
ignorance of a definite but unknown fact, which the quantum case does not supply
\autocite{bruknerzeilinger2001shannon}. The framework's reply is structural, and it is a
demarcation rather than a rebuttal. Shannon's precondition fails exactly where they say it fails
--- in the pre-record regime, where there is no definite unknown to be ignorant of --- and there
the operative measure is the quadratic invariant just derived. Shannon's axioms hold exactly where
the precondition holds --- on records, after \textsc{fanout}, where outcomes are definite and
ignorance is classical --- and there the uniqueness theorem forces $-\log p$
\autocite{shannon1948mathematical}. The basis-dependence dissolves with the same stroke: a basis
\emph{is} selected precisely at \textsc{fanout}, so the basis-dependent measure applies only where
a basis exists. Shannon is not inadequate; Shannon is post-\textsc{fanout}. The framework thereby
keeps both measures, each on its own side of the boundary, and uses the additivity of the
post-record measure --- the very property the previous subsection used to pin the base field ---
with a clear conscience.

\subsection{Thomson's lamp, resolved}

The section owes \S\ref{sec:continuity} a debt incurred there: the lamp's amplitude-side limit.
Thomson's lamp --- toggled at times $1 - 2^{-n}$, its state at $t = 1$ demanded --- has no
classical answer: the task description fixes every finite stage and thereby fixes no limit at all
\autocite{Thomson_1954,benacerraf1962tasks}, which \S\ref{sec:continuity} read as the proof that
the record algebra is not limit-closed. On the amplitude side the limit exists, and the calculus of
this section computes it. The toggle exchanges the two amplitudes; the self-consistent state of the
completed task is the fixed point of the exchange in the complex calculus, and solving it under the
signed continuation gives
\[
P(\mathrm{on}) = \tfrac{1 + i}{2}, \qquad P(\mathrm{off}) = \tfrac{1 - i}{2},
\]
a superposed fixed point whose imaginary part carries the unrecordable datum --- the phase of the
endless alternation, when the lamp \emph{would have} stopped. Two checks certify that the
assignment is calculus rather than decoration. The pair is normalised,
$P(\mathrm{on}) + P(\mathrm{off}) = 1$; and the Born rule of the previous subsection, applied to
the fixed point, returns $|P(\mathrm{on})|^{2} = \tfrac{1}{2}$ --- exactly the symmetric classical
ignorance that is the only sane record-level answer, recovered rather than imposed. The lamp,
which entered the framework at its very beginning as the place where classical computation meets
its boundary, thus exits it as a worked example: undetermined as a record, determinate as an
amplitude, the two sides split precisely along \textsc{fanout} --- and the diagonal construction it
physically instantiates \autocite{cantor1891diagonal,Turing_1937} is the same one that powers
Proposition~\ref{prop:inside-outside}. The status flag is owed here as everywhere: the
undetermination is theorem-level; the specific fixed-point value is the framework's complex
calculus, structural, offered with its two consistency checks shown.

\subsection{Local tomography, derived}\label{sec:local-tomography}

The fourth instalment concerns composites, and it discharges two promises at once. The
Inside/Outside Equivalence Principle --- now a theorem --- requires that an external observer's
characterisation of a composite system be recoverable by the embedded participants, and the
embedded participants have access only to local measurements and their correlations. Equivalence
of the two perspectives therefore entails that the state of a composite is determined by the
statistics of local measurements: \emph{local tomography}, obtained as a consequence rather than
imposed as an axiom. This converts a postulate of the reconstruction programmes
\autocite{Hardy2001a,mullerp2011} into a theorem-shaped consequence of the one principle ---
discharging the caveat recorded in \S\ref{sec:continuity} --- and it closes, from the framework's
own resources, the escape route in the real-number debate noted above: the assumption the
objection of \autocite{hoffreumonwoods2026real} contests is here not an assumption. The status is
flagged plainly: the derivation is structural, one paragraph of perspective-accounting, and its
formalisation within the categorical setting belongs to Open Problem~\ref{op:sigma-functor}.

\subsection{Contrast: QBism, and the ledger}

The nearest neighbour in spirit deserves a paragraph of honest contrast. QBism
\autocite{fuchsschack2013qbism} holds, as this framework does, that quantum probabilities are an
agent's --- that the theory is a calculus of expectation, not a catalogue of properties. What QBism
does not say is why the agent's expectations must be complex-valued and quadratically squared; the
Born rule enters as a normative addition to Bayesian coherence, and the complex amplitude as an
unexplained representation. This section is, in effect, the missing why: the agent's ignorance is
self-referential --- the agent is inside the description --- and self-referential ignorance is
complex, with the quadratic rule as the cheapest record-compatible invariant of it. Where QBism
stops at the agent, the framework derives the agent's calculus.

The ledger, then. Delivered in this section: the reduction of two principles to one
(Proposition~\ref{prop:inside-outside}); the pre-record regime with its anti-events, handed to
\S\ref{sec:continuity} as the five-face regime whose continuity is priced there
(Open Problem~\ref{op:limit-realisability}); the complex base, forced downward by experiment and
upward by the additivity of information; the Born rule from fibre privacy; the lamp's fixed point;
and local tomography. Owed, and stated where they are owed: the functorial rigour of the involution
and the fixed-point classification (Open Problem~\ref{op:sigma-functor}); the remaining premise of continuity (Open Problem~\ref{op:limit-realisability}, in
\S\ref{sec:continuity}); and the
rung-by-rung compulsion of the algebra ladder, posed as Open
Problem~\ref{op:ladder-compulsion} in \S\ref{sec:division-algebras}. The
hand-on is the paper's spine: the $\mathbb{C}$ this section derives natively and the $\mathbb{C}$
that \S\ref{sec:continuity} underwrites through Hardy's theorem are the same field, the floor of
the ladder that \S\ref{sec:division-algebras} climbs and \S\ref{sec:hopf-justification}
geometrises. One principle; everything else is theorem, derivation, or a debt with a stated
invoice.

%      % EDITED 2026-10-29 scripts/patch_continuity_tower_premise_2026-10-29T1500.py
% EDITED 2026-09-15 scripts/patch_value_continuum_declared_cencov_2026-09-15T1540.py
% EDITED 2026-09-15 scripts/patch_continuity_connected_state_space_2026-09-15T1430.py
% =====================================================================
%  SECTION: Continuity Is Not an Axiom
%  Drafted 2026-07-04 from continuity_limit_realisability_2026-07-04.md
%  (now in project documents) and the May record
%  (dynamics_companion_fanout_unitarity_2026-05-20T1430.md).
%
%  PLACEMENT in main.tex: a standalone section between self-reference
%  and division-algebras --
%      % EDITED 2026-10-29 scripts/patch_continuity_tower_premise_2026-10-29T1500.py
% EDITED 2026-09-20 scripts/patch_born_cost_ledger_v2_2026-09-20T0900.py
% =====================================================================
%  SECTION: Self-Reference: One Principle
%  Drafted 2026-07-04. The paper's foundational section; resolves the
%  largest dangling label in the manuscript (sec:self-reference).
%
%  PLACEMENT in main.tex: FIRST of the foundations chain --
%      \input{self_reference_section_2026-07-04}
%      \input{continuity_section_2026-07-04}
%      \input{division_algebras_section_2026-07-03}
%      \input{hopf_justification_section_2026-06-08}
%
%  SOURCES (re-read 2026-07-04): lawvere_inside_outside_fixed_point_
%  theorem_2026-05-04T1806.md (the spine: the two-principles reduction,
%  the theorem, the sigma candidate, the open items -> op minted per
%  D9); research_summary_neg_prob.md (negativity route; lamp (1+i)/2);
%  dynamics_companion_fanout_unitarity (legs 1a/1c; anti-events; the
%  Brukner-Zeilinger rebuttal per D6); Pure_states_and_mixed_states
%  (local tomography derived; QBism contrast per D7).
%
%  DECISIONS ENCODED (JS 2026-07-04): D1 title as below; D2 Wheeler1990
%  cited here as the title anchor; D3 lineage Popper -> Breuer ->
%  Svozil -> Szangolies (Dalla Chiara omitted); D4 Frauchiger_2018
%  cited here, dissolved in sec:wigners-friend; D5 Kirkwood-Dirac
%  heritage via the Arvidsson-Shukur 2024 review as one-stop; D6
%  Brukner-Zeilinger + the framework rebuttal IN ("Shannon is
%  post-FANOUT"); D7 QBism contrast IN; D8 Aaronson private
%  communication OMITTED (published island paper only); D9 one op
%  minted (op:sigma-functor), cross-linked to op:hopf-frobenius.
%
%  CITATIONS -- existing keys: Wheeler1990, small_2006, Godel1931,
%  Lawvere_1969, Yanofsky2003, CoeckeKissinger2017, Rovelli1996a,
%  Popper1950a, Popper1950b, Breuer1995, svozil_k_2018,
%  szangolies_j_2018, Frauchiger_2018, dirac1942, Feynman_r_82,
%  Feynman1987, Youssef1994 [SWAP-FLAG: intended paper is "Quantum
%  mechanics as complex probability theory", Mod. Phys. Lett. A 9
%  (1994) 2571 -- if that lives under youssef1994quantum, swap],
%  caticha_a_2006, goyal_knuth_skilling_2010, Aaronson2004h, Hardy2001a,
%  muller_p_2011, Thomson_1954, benacerraf1962tasks, cantor1891diagonal,
%  Turing_1937.
%  NEW keys (RIS in self_reference_refs_2026-07-04.ris, all verified
%  this session): tikochinsky1988feynman, renou2021real, chen2022ruling,
%  wu2022refutation, hoffreumonwoods2026real, arvidssonshukur2024kirkwood,
%  bruknerzeilinger2001shannon, fuchsschack2013qbism,
%  shannon1948mathematical.
%
%  CROSS-REFS: forward (fine under two-pass LaTeX): sec:continuity,
%  sec:division-algebras, sec:hopf-justification, prop:f1-dictionary,
%  sec:wigners-friend, op:limit-realisability, op:hopf-frobenius.
%  Still dangling after this section: sec:fanout-boundary (and its
%  sibling label sec:fanout-direction) -- both DRAFTED 2026-07-06.
%
%  AUDITED 2026-07-06 (section-by-section final pass, first stop):
%  proof of prop:inside-outside re-verified by hand; attribution chain
%  clean (Burgin/Khrennikov correctly absent); all header-declared
%  citations confirmed used; lamp arithmetic and antecedents verified;
%  FANOUT smallcaps consistent; no pre-definition use of "frame".
%  Fixes: stray Unicode em-dash (QBism paragraph, the file's only
%  non-ASCII character) -> house ---; trace-blind two-way stitch added
%  to the Born subsection ("one privacy, two duties");
%  op:sigma-functor gains its first standing constraint (F1
%  inheritance via the op:hopf-frobenius discharge) and its previously
%  one-directional local-tomography cargo made explicit.
% =====================================================================

\section{Self-Reference: One Principle}
\label{sec:self-reference-2}

Wheeler asked for it from bit: for physics derived from information, ``every it --- every particle,
every field of force'' obtained from binary answers to yes-or-no questions \autocite{Wheeler1990}.
This paper's route to that programme runs through G\"{o}del: the single principle from which
everything else here is derived, deferred, or priced is \emph{self-reference} --- the fact that any
complete physical description, made from inside the world it describes, must describe the describer.
This section states the principle, proves that the framework's second historical commitment follows
from it as a theorem rather than standing beside it as an axiom, and then extracts the price of the
principle in four instalments: probabilities whose parts can cancel, amplitudes valued in
$\mathbb{C}$, the Born rule, and the local tomography of composite systems. What the section cannot
yet derive is stated as precisely as what it can, in one numbered problem and in the ledger at the
end.

\subsection{Events, not objects}

The framework's ontology is relational: the fundamental entities are events and the relations
between them, and objects are the bookkeeping of relational invariants --- what remains stable
across many relations. The kinship with relational quantum mechanics \autocite{Rovelli1996a} is
genuine and acknowledged: facts are facts relative to the systems that register them. The
difference, developed over the rest of this paper, is that the present framework grounds the
relational structure in a single logical principle and then carries it, through the division
algebras and the Hopf fibrations, all the way to the gauge group and the particle spectrum ---
relations are not only the ontology but the derivation engine. On this reading a measurement is not
the revelation of a pre-existing property; it is the creation of a relation, and the later sections
of this paper read that creation geometrically (\S\ref{sec:hopf-justification}) and even
gravitationally. Here we need only the logical core: whatever describes the world is an event in
the world it describes.

\subsection{One principle, and a theorem}

The framework historically carried two foundational commitments: self-reference, and the
\emph{Inside/Outside Equivalence Principle} --- that the laws of physics must be identical whether
formulated by an embedded participant or by an external observer. Two principles are one too many,
and this subsection reduces them.

The reduction corrects the framework's own origin, and we make the correction in print. The 2005
form of the argument \autocite{small_2006} ran a G\"{o}delian chain: a self-referential process,
being complete, must by the incompleteness theorem \autocite{Godel1931} be inconsistent; an
inconsistent system can negate itself; and the negation of self-reference is the production of
causal sequence. The chain captures the right intuition and is retained here as motivation, but it
is not sound as a proof: classical inconsistency yields \emph{ex falso quodlibet} --- every
proposition provable, every causal sequence an output, no constraint at all. The load-bearing
argument is different, and better: a fixed-point theorem.

Let $\mathcal{C}$ be the category of physical descriptions and admissible processes --- concretely,
the dagger-compact-closed setting with classical structures in which categorical quantum mechanics,
and the whole of this paper, works \autocite{CoeckeKissinger2017}. Lawvere showed that in a
sufficiently rich such setting, any weakly point-surjective self-representation forces every
endomap to have a fixed point, and Yanofsky showed that one diagonal argument of this shape
underlies Cantor's, G\"{o}del's, Turing's and Tarski's theorems alike
\autocite{Lawvere_1969,Yanofsky2003,cantor1891diagonal}. The framework's single commitment is that
$\mathcal{C}$ contains a \emph{self-representing object} --- an object capable of encoding the
processes on itself. That is what ``self-reference is foundational'' means, stated categorically,
and it is the surjectivity hypothesis of the diagonal lemma. Now let
$\sigma \colon \mathcal{C} \to \mathcal{C}$ be the \emph{inside/outside involution}: the operation
exchanging the two complementary presentations of any description --- the inside presentation, as
participant-indexed total space in which the observer is part of the system, and the outside
presentation, as a base-times-fibre decomposition over a classical parameter space. That
$\sigma^{2} = \mathrm{id}$ is not an axiom but a coherence requirement: the outside view of the
inside view of a thing is the thing. The candidate concrete realisation of $\sigma$ is the
Hopf-duality swap --- total space against base-times-fibre for the three fibrations --- which gives
the involution computable content and ties this subsection to \S\ref{sec:hopf-justification}.

\begin{proposition}[Inside/outside fixed point]\label{prop:inside-outside}
Let $\mathcal{C}$ satisfy the richness condition above and contain a self-representing object, and
let $\sigma$ be an involution on $\mathcal{C}$. Then the fixed-point locus
$\mathrm{Fix}(\sigma)$ is non-empty: there exist descriptions invariant under the exchange of
inside and outside perspective.
\end{proposition}

\begin{proof}
Lawvere's construction, applied to $\sigma$ in place of negation: given a weakly point-surjective
$\varphi \colon A \to Y^{A}$, form $g(a) = \sigma\!\left(\varphi(a)(a)\right)$; point-surjectivity
supplies $a^{*}$ with $\varphi(a^{*}) = g$, and then $\varphi(a^{*})(a^{*})$ is a fixed point of
$\sigma$ \autocite{Lawvere_1969,Yanofsky2003}. Where the same construction applied to Boolean
negation yields paradox --- negation has no fixed point on a two-element type, so the conclusion
explodes the premises --- the inside/outside involution is an honest involution on a rich type, and
the conclusion is a constructive fixed point rather than a contradiction.
\end{proof}

\textbf{Definition.} \emph{Physical law is the structure carried by} $\mathrm{Fix}(\sigma)$.

The Inside/Outside Equivalence Principle is thereby demoted from axiom to theorem: physical law is
invariant under the exchange of perspectives \emph{because it is the fixed-point set of that
exchange}, and the fixed-point set is forced to be non-empty by self-reference alone. One
commitment remains foundational; everything the second commitment asserted is now derived. The
reader who accepts categorical quantum mechanics as an existing programme has already accepted the
machinery; the framework adds one premise and collects the rest. And the fixed-point locus is not
empty of structure: it inherits the Frobenius-algebra apparatus of $\mathcal{C}$, so that the
Bell, GHZ and W primitives --- and with them the particle classification of the later sections ---
live inside $\mathrm{Fix}(\sigma)$ from the start. What is proved above is non-emptiness given the
setting; what is assumed is the setting itself and the existence of $\sigma$ as a genuine
involution; what is open is the rigour of both, and it is priced now rather than hidden:

\begin{openproblem}[The involution made functorial]\label{op:sigma-functor}
Fix the category $\mathcal{C}$ rigorously --- the dagger-compact-closed setting with classical
structures is the candidate --- and define $\sigma$ on morphisms, not only on objects, the natural
route being the induced fibre--base exchange on spider diagrams; establish that $\sigma$ is an (at
worst pseudo-)involutive endofunctor; and prove or refute the identification of
$\mathrm{Fix}(\sigma)$ with the three Hopf cross-sections, under which Adams's theorem would become
a corollary of the fixed-point classification. This problem is cross-linked with Open
Problem~\ref{op:hopf-frobenius}: under an affirmative resolution of both, the parallelisable
terminus, the fixed-point classification, and compositional completeness are three statements of
one fact, and progress on either problem constrains the other. One constraint is already in place:
at the complex rung, Proposition~\ref{prop:f1-dictionary} supplies the dictionary between
classical structures and the fibration's lifted skeletons --- the clause of
Open Problem~\ref{op:hopf-frobenius} it discharges --- and any resolution here must reproduce it.
The problem also carries a piece of cargo recorded later in this section: the formalisation, within
the fixed category, of the perspective-accounting that derives local tomography.
\end{openproblem}

\subsection{The lineage}

The obstruction this section builds on is old, and the framework claims its completion, not its
discovery. Popper argued in 1950 that a predictor embedded in the world has irreducibly inferior
knowledge of itself, because it cannot impart information about itself without, by that very act,
rendering the information out of date \autocite{Popper1950a,Popper1950b} --- the passage is quoted
where the framework makes the obstruction structural (\S\ref{sec:wigners-friend}). Breuer proved
the impossibility of accurate state self-measurement within quantum theory \autocite{Breuer1995};
Svozil's monograph surveys the whole territory where physics meets self-reference, undecidability
and Chaitin-style algorithmic limits \autocite{svozil_k_2018}; Szangolies sharpened the modern form
as epistemic horizons \autocite{szangolies_j_2018}; and the Frauchiger--Renner theorem is the
contemporary crisis-form --- nested observers deriving contradictory records when each treats the
others' quantum descriptions as universally valid \autocite{Frauchiger_2018} --- which
\S\ref{sec:wigners-friend} dissolves relationally rather than paradoxically. This framework's
addition to the lineage is twofold. Structurally, Proposition~\ref{prop:f1-dictionary} converts
Popper's temporal obsolescence into algebraic cancellation: the datum a copier cannot put into its
own record is exactly its fibre commitment, which does not arrive stale --- it cancels. And
foundationally, the present section runs the lineage in reverse: where the tradition treats
self-reference as a limit on knowledge \emph{within} physics, the framework takes it as the
generator of physics, and derives the quantum formalism from it.

\subsection{The price is negative probability}\label{sec:complex-amplitudes-2}

What does it cost a probabilistic description to contain its own describer? The first instalment is
the sign. In the regime before any record exists, every event has an anti-event --- the description
must be able to undo what it has provisionally done, or the self-description could never be
revised into consistency --- and a probability calculus closed under undoing cannot be confined to
$[0,1]$: it is a signed calculus, in which the parts of a probability can oppose and cancel. The
framework did not discover negativity as bookkeeping; Dirac admitted negative probabilities as
formal intermediaries \autocite{dirac1942}, Feynman defended them as consistent whenever the final,
recorded probabilities stay positive \autocite{Feynman_r_82,Feynman1987}, and the Kirkwood--Dirac
distribution has carried complex quasiprobability from 1933 to the present, as a recent
comprehensive review documents \autocite{arvidssonshukur2024kirkwood}. What the framework adds is
the forcing: negativity is not a permissible convenience but the first price of self-referential
closure, and it is charged exactly and only in the pre-record regime --- which is why every
negativity in the tradition lives in intermediate, unrecorded quantities and never in a final
record.

The second instalment is complexity. A signed information measure must extend Shannon's
$S = -\ln p$ off the positive axis, and the analytic continuation of the logarithm carries $p$ into
the complex plane: amplitudes, with the phase as the continued information. Here too the extension
has a prior tradition, and the attributions are owed precisely. Tikochinsky derived Feynman's rules
for complex amplitudes from consistency of composition \autocite{tikochinsky1988feynman}; Youssef
reformulated quantum mechanics as complex probability theory outright \autocite{Youssef1994};
Caticha carried objective amplitudes to Bayesian probabilities \autocite{caticha_a_2006}; and
Goyal, Knuth and Skilling derived Feynman's rules from a pair-postulate without assuming
analyticity \autocite{goyal_knuth_skilling_2010}. The extension is prior art. The framework's
claim is the compulsion: the complex extension is not a modelling choice among exotic probability
theories but the unique continuation forced by self-referential closure --- the same forcing that
Aaronson glimpsed when he argued quantum mechanics is what negative probability inevitably becomes
\autocite{Aaronson2004h}.

Why does the continuation stop at $\mathbb{C}$? Downward, the real field is not available: real
quantum theory is experimentally distinguishable from complex quantum theory in network scenarios
and has been disfavoured by experiment
\autocite{renou2021real,chen2022ruling,wu2022refutation}; the recent objection that the exclusion
rests on locally tomographic composition \autocite{hoffreumonwoods2026real} is, for this framework,
not an escape but a confirmation of where the weight sits --- local tomography is \emph{derived}
below, not assumed, so the framework closes from its own resources exactly the door the objection
holds open. Upward, the quaternions and octonions fail as a \emph{base} for the information
continuation: the continued measure is a logarithm, and beyond $\mathbb{C}$ the logarithm of a
product is no longer the sum of logarithms --- non-commutativity severs
$\log(ab) = \log a + \log b$ --- so the additivity that makes it an information measure is lost.
Aaronson's parameter-counting and signalling arguments against a quaternionic base point the same
way \autocite{Aaronson2004h}. The higher algebras are not banished; they return as the
\emph{internal} algebras of the register, which is the business of \S\ref{sec:division-algebras}.
The base is $\mathbb{C}$, for informational reasons, and stays $\mathbb{C}$.

\subsection{The Born rule from fibre privacy}\label{sec:born-rule}

The third instalment is the quadratic rule. An amplitude carries one datum more than any record of
it can: its overall phase. This was asserted throughout the framework's development as the
observer's private, self-referential ignorance; it is now a proved statement, in exactly the form
needed. By Proposition~\ref{prop:f1-dictionary}, a classical structure --- the copy operation any
record requires --- is committed to a specific lift of its skeleton, and that commitment is absent
from every classical record the structure induces: the copier is fibre-committed, the record is
fibre-blind, and the fibre freedom is gauge for every observer. (The computation extends verbatim
from one qubit to the full register, where Proposition~\ref{prop:trace-blind} makes it the ground
of hypercharge's tracelessness --- one privacy, two duties.) Popper's obsolescence, made
cancellation. It follows that anything recordable is invariant under the fibre --- observables and
their statistics must be $U(1)$-invariant functions of the amplitude. The lowest-degree polynomial invariant
of a complex amplitude is the quadratic $|\psi|^{2}$ --- the non-polynomial invariant $|\psi|$,
of degree one, is excluded by the analyticity declared at \S\ref{sec:complex-amplitudes} --- and
normalisation fixes the coefficient:
probability is the minimal record-compatible function of the amplitude,
\[
p \;=\; |\psi|^{2},
\]
the Born rule as the cheapest invariant the fibre permits. The same square is reached from the
cost side in \S\ref{sec:born-cost-ledger}, where a record is the closed loop of a history and its
return leg and the exponent is the number of legs; neither route uses the lattice of projections
that Gleason's theorem needs, and both hold for a single qubit. The derivation is the framework's
own;
its premise --- that the fibre datum is structurally excluded from records --- is
Proposition~\ref{prop:f1-dictionary}, and its geometric reading --- probability as base-point data
under the Hopf projection --- is developed in \S\ref{sec:hopf-justification}.

One standing objection to information-theoretic foundations is answered here, because this
framework needs the answer. Brukner and Zeilinger argued that Shannon information is conceptually
inadequate for quantum measurements: it is basis-dependent, and its operational meaning presupposes
ignorance of a definite but unknown fact, which the quantum case does not supply
\autocite{bruknerzeilinger2001shannon}. The framework's reply is structural, and it is a
demarcation rather than a rebuttal. Shannon's precondition fails exactly where they say it fails
--- in the pre-record regime, where there is no definite unknown to be ignorant of --- and there
the operative measure is the quadratic invariant just derived. Shannon's axioms hold exactly where
the precondition holds --- on records, after \textsc{fanout}, where outcomes are definite and
ignorance is classical --- and there the uniqueness theorem forces $-\log p$
\autocite{shannon1948mathematical}. The basis-dependence dissolves with the same stroke: a basis
\emph{is} selected precisely at \textsc{fanout}, so the basis-dependent measure applies only where
a basis exists. Shannon is not inadequate; Shannon is post-\textsc{fanout}. The framework thereby
keeps both measures, each on its own side of the boundary, and uses the additivity of the
post-record measure --- the very property the previous subsection used to pin the base field ---
with a clear conscience.

\subsection{Thomson's lamp, resolved}

The section owes \S\ref{sec:continuity} a debt incurred there: the lamp's amplitude-side limit.
Thomson's lamp --- toggled at times $1 - 2^{-n}$, its state at $t = 1$ demanded --- has no
classical answer: the task description fixes every finite stage and thereby fixes no limit at all
\autocite{Thomson_1954,benacerraf1962tasks}, which \S\ref{sec:continuity} read as the proof that
the record algebra is not limit-closed. On the amplitude side the limit exists, and the calculus of
this section computes it. The toggle exchanges the two amplitudes; the self-consistent state of the
completed task is the fixed point of the exchange in the complex calculus, and solving it under the
signed continuation gives
\[
P(\mathrm{on}) = \tfrac{1 + i}{2}, \qquad P(\mathrm{off}) = \tfrac{1 - i}{2},
\]
a superposed fixed point whose imaginary part carries the unrecordable datum --- the phase of the
endless alternation, when the lamp \emph{would have} stopped. Two checks certify that the
assignment is calculus rather than decoration. The pair is normalised,
$P(\mathrm{on}) + P(\mathrm{off}) = 1$; and the Born rule of the previous subsection, applied to
the fixed point, returns $|P(\mathrm{on})|^{2} = \tfrac{1}{2}$ --- exactly the symmetric classical
ignorance that is the only sane record-level answer, recovered rather than imposed. The lamp,
which entered the framework at its very beginning as the place where classical computation meets
its boundary, thus exits it as a worked example: undetermined as a record, determinate as an
amplitude, the two sides split precisely along \textsc{fanout} --- and the diagonal construction it
physically instantiates \autocite{cantor1891diagonal,Turing_1937} is the same one that powers
Proposition~\ref{prop:inside-outside}. The status flag is owed here as everywhere: the
undetermination is theorem-level; the specific fixed-point value is the framework's complex
calculus, structural, offered with its two consistency checks shown.

\subsection{Local tomography, derived}\label{sec:local-tomography}

The fourth instalment concerns composites, and it discharges two promises at once. The
Inside/Outside Equivalence Principle --- now a theorem --- requires that an external observer's
characterisation of a composite system be recoverable by the embedded participants, and the
embedded participants have access only to local measurements and their correlations. Equivalence
of the two perspectives therefore entails that the state of a composite is determined by the
statistics of local measurements: \emph{local tomography}, obtained as a consequence rather than
imposed as an axiom. This converts a postulate of the reconstruction programmes
\autocite{Hardy2001a,mullerp2011} into a theorem-shaped consequence of the one principle ---
discharging the caveat recorded in \S\ref{sec:continuity} --- and it closes, from the framework's
own resources, the escape route in the real-number debate noted above: the assumption the
objection of \autocite{hoffreumonwoods2026real} contests is here not an assumption. The status is
flagged plainly: the derivation is structural, one paragraph of perspective-accounting, and its
formalisation within the categorical setting belongs to Open Problem~\ref{op:sigma-functor}.

\subsection{Contrast: QBism, and the ledger}

The nearest neighbour in spirit deserves a paragraph of honest contrast. QBism
\autocite{fuchsschack2013qbism} holds, as this framework does, that quantum probabilities are an
agent's --- that the theory is a calculus of expectation, not a catalogue of properties. What QBism
does not say is why the agent's expectations must be complex-valued and quadratically squared; the
Born rule enters as a normative addition to Bayesian coherence, and the complex amplitude as an
unexplained representation. This section is, in effect, the missing why: the agent's ignorance is
self-referential --- the agent is inside the description --- and self-referential ignorance is
complex, with the quadratic rule as the cheapest record-compatible invariant of it. Where QBism
stops at the agent, the framework derives the agent's calculus.

The ledger, then. Delivered in this section: the reduction of two principles to one
(Proposition~\ref{prop:inside-outside}); the pre-record regime with its anti-events, handed to
\S\ref{sec:continuity} as the five-face regime whose continuity is priced there
(Open Problem~\ref{op:limit-realisability}); the complex base, forced downward by experiment and
upward by the additivity of information; the Born rule from fibre privacy; the lamp's fixed point;
and local tomography. Owed, and stated where they are owed: the functorial rigour of the involution
and the fixed-point classification (Open Problem~\ref{op:sigma-functor}); the remaining premise of continuity (Open Problem~\ref{op:limit-realisability}, in
\S\ref{sec:continuity}); and the
rung-by-rung compulsion of the algebra ladder, posed as Open
Problem~\ref{op:ladder-compulsion} in \S\ref{sec:division-algebras}. The
hand-on is the paper's spine: the $\mathbb{C}$ this section derives natively and the $\mathbb{C}$
that \S\ref{sec:continuity} underwrites through Hardy's theorem are the same field, the floor of
the ladder that \S\ref{sec:division-algebras} climbs and \S\ref{sec:hopf-justification}
geometrises. One principle; everything else is theorem, derivation, or a debt with a stated
invoice.

%      % EDITED 2026-10-29 scripts/patch_continuity_tower_premise_2026-10-29T1500.py
% EDITED 2026-09-15 scripts/patch_value_continuum_declared_cencov_2026-09-15T1540.py
% EDITED 2026-09-15 scripts/patch_continuity_connected_state_space_2026-09-15T1430.py
% =====================================================================
%  SECTION: Continuity Is Not an Axiom
%  Drafted 2026-07-04 from continuity_limit_realisability_2026-07-04.md
%  (now in project documents) and the May record
%  (dynamics_companion_fanout_unitarity_2026-05-20T1430.md).
%
%  PLACEMENT in main.tex: a standalone section between self-reference
%  and division-algebras --
%      \input{self_reference_section_2026-07-04}
%      \input{continuity_section_2026-07-04}
%      \input{division_algebras_section_2026-07-03}
%  (Architecture call SETTLED by the sync manifest: standalone.)
%
%  CITATIONS -- existing keys (as supplied by JS 2026-07-04):
%    Hardy2001a, Aaronson2004h, Thomson_1954, Lawvere_1969,
%    Turing_1937, muller_p_2011, Brukner2009.
%  NEW keys (RIS in continuity_refs_2026-07-04.ris):
%    cantor1891diagonal, benacerraf1962tasks, broeckertomdieck1985compact.
%  Deliberately excluded per JS: Chiribella--D'Ariano--Perinotti;
%  Bargmann/Wigner's theorem. Yanofsky2003 now paired on the
%  Lawvere-shaped clause (the pending-key excuse was void: the key
%  exists and sec:self-reference uses it).
%
%  CROSS-REFS: sec:self-reference now DRAFTED (2026-07-04, later the
%  same day) -- no longer dangling; only sec:fanout-boundary (and its
%  sibling sec:fanout-direction) -- DRAFTED 2026-07-06.
%
%  AUDITED 2026-07-06 (final pass, second stop): both proofs
%  re-verified by hand (Lagrange divisibility sentence exact; the
%  Baire countable-escape chain sound); caveat-discharge with
%  sec:self-reference confirmed two-way (already present); file pure
%  ASCII; no pre-definition "frame"; all citations used as declared.
%  Fixes: Yanofsky2003 paired on the Lawvere clause; the op's two
%  candidate routes anchored to printed propositions (f1-dictionary
%  for the divisibility premise's referent, inside-outside for the
%  closedness demand) -- statuses unchanged, referents sharpened;
%  header staleness corrected. Deliberate FORWARD references --
%  prop:f1-dictionary and sec:hopf-justification (defined later in the
%  paper; fine under two-pass LaTeX and flagged in prose),
%  sec:division-algebras, sec:space-of-computations.
%
%  DECISIONS ENCODED: op:limit-realisability minted (the owed arrow,
%  priced: exactly two derivations, divisibility and closedness);
%  Lagrange used twice (finite exclusion; component-group triviality);
%  the (1+i)/2 lamp value NOT stated here -- "a superposed fixed point"
%  with the complex-probability treatment deferred to
%  sec:self-reference; the two-branch over-determination stated to
%  pre-empt a circularity reading; Thomson spelling throughout.
% =====================================================================

\section{Continuity Is Not an Axiom}
\label{sec:continuity}

Hardy derives quantum theory from five axioms \autocite{Hardy2001a}. Four of them are satisfied
by classical probability theory; the fifth --- that a \emph{continuous} reversible
transformation connects any two pure states --- is the entire difference, the single postulate
that turns the classical simplex into the complex Hilbert space. Every reconstruction programme
pays an equivalent toll: continuous reversibility in Masanes and M\"{u}ller
\autocite{mullerp2011}, the corresponding postulate in Daki\'{c} and Brukner
\autocite{Brukner2009}. The axiom is load-bearing, and in every programme it is unexplained: it
is where the quantumness is put in by hand. The claim of this section is that in the present
framework it is not an axiom. Continuity is the fifth face of the pre-measurement regime that self-reference forces, and what
stands between the framework and a full derivation is not a postulate about values but one
premise about transformations --- that every transformation the inside performs has a stationary
outside presentation --- derived in outline below from the inside/outside involution of
\S\ref{sec:self-reference} and stated as an open problem at the end of the section with nothing
hidden in it. The continuum itself, the completion that every probability theory adopts, is
named for what it is: the real line, shared with classical probability and carrying no quantum
content.

One structural remark prevents a misreading. The framework reaches complex amplitudes natively in
\S\ref{sec:self-reference}, by the negative-probability route; the present section reaches the same
destination by an independent branch --- the regime's continuity, then Hardy's theorem. The two
branches converging on one structure is over-determination, not circularity, and the convergence is
itself evidence that the destination is forced.

\subsection{Five faces of one regime}

The regime before any measurement --- before any \textsc{fanout} --- can be characterised five
equivalent ways: negative probability is available; anti-events exist, every event being undoable;
no classical record exists; no preferred basis is selected; and a continuous reversible
transformation connects any two pure states. These are not five premises to be combined. They are
one regime under five descriptions, and \textsc{fanout} (\S\ref{sec:fanout-boundary}) is the single
event that fails all five at once: it writes the record, selects the basis, kills the anti-events,
discretises the transformation group, and switches off the negativity. The equivalence of the first
two faces is the anti-event identity of \S\ref{sec:self-reference}; the third is the boundary's
definition. The work of this section is the equivalence of the last two faces, and then the route
from the regime's existence to its continuity. The fourth face has, since
Proposition~\ref{prop:f1-dictionary} (\S\ref{sec:hopf-justification} --- a deliberate forward
reference), an exact form: a basis-selection is the commitment of a classical structure, a lifted
antipodal skeleton, so ``no preferred basis'' means precisely that no copy structure has fired.

\subsection{Basis selection is the collapse of continuity}

\begin{proposition}[Continuity forbids a selected basis]\label{prop:basis-discreteness}
If a finite set of pure states is privileged --- a preferred basis --- then every reversible
transformation respecting the privilege carries the set to itself and acts on it as a permutation;
the respecting transformations therefore realise only a finite group of motions among the
distinguished states, and no continuous family of them connects distinct pure states of the set.
Contrapositively, continuous reversibility entails that no basis is selected.
\end{proposition}

\begin{proof}
A privileged finite set is by definition carried to itself by every structure-preserving reversible
map, whose induced action on the set is a permutation of finitely many points. A continuous path of
such maps induces a continuous path in a finite discrete set, which is constant. Hence a continuous
reversible family moving one distinguished pure state to another respects no such privilege, and if
every pair of pure states is so connected, no finite set is privileged.
\end{proof}

The converse is where the historical difficulty lies: the absence of a privileged finite set does
not by itself force continuity, for the transformation group could in principle be infinite and
discrete. Two structural properties close every such escape at once.

\begin{proposition}[Divisible and closed implies connected]\label{prop:divisible-closed}
Let $G$ be the group of reversible transformations of the pre-measurement regime, acting on the
state body. If $G$ is divisible --- every element has an $n$-th root for every $n$ --- and
topologically closed, then $G$ is a compact connected Lie group; and if $G$ moreover acts
transitively on the pure states, then a continuous reversible transformation connects any two of
them: Hardy's fifth axiom.
\end{proposition}

\begin{proof}
The state body is compact and convex with nonempty interior, so its symmetry group is compact, and
a closed subgroup of a compact matrix group is a compact Lie group with finitely many connected
components \autocite{broeckertomdieck1985compact}. The component group $G/G_{0}$ is a quotient of a
divisible group and hence divisible; but a nontrivial finite group is never divisible --- with
$\lvert H\rvert = m$ every element satisfies $h^{m} = e$, so no element other than the identity
possesses an $m$-th root --- and $G/G_{0}$ is finite, so it is trivial and $G = G_{0}$ is connected.
A connected Lie group is path-connected, and transitivity then supplies the continuous reversible
path between any two pure states.
\end{proof}

The same one-line fact of Lagrange closes the remaining escapes. Finite groups are excluded
directly, being nontrivial and never divisible. Countable groups are excluded because a
countable closed subgroup of a compact group is compact, has an isolated point by the Baire
category theorem, and is therefore discrete and finite. Uncountable disconnected groups are
excluded by the component-group argument of the proposition. Every route to a discrete
transformation group therefore passes through a failure of divisibility or a failure of
closedness. Divisibility, however, asks for more than the conclusion needs, and the next
proposition removes it.

\begin{proposition}[Continuity from a connected state space]\label{prop:connected-transitive}
Let the pure states of a finite-dimensional system form a connected space $P$, and let the group
$G$ of reversible transformations of the pre-measurement regime be topologically closed and act
transitively on $P$. Then the identity component $G_0$ already acts transitively on $P$, and a
continuous reversible transformation connects any two pure states: Hardy's fifth axiom holds,
with no divisibility assumed.
\end{proposition}

\begin{proof}
As in Proposition~\ref{prop:divisible-closed}, closedness makes $G$ a compact Lie group with
finitely many connected components \autocite{broeckertomdieck1985compact}, so $G_0$ is a closed
normal subgroup of finite index and is path-connected. Each $G_0$-orbit in $P$ is the continuous
image of a compact group, hence compact, hence closed. Since $G$ is transitive and $G_0$ is
normal, the $G_0$-orbit of $gx$ is $g$ applied to the $G_0$-orbit of $x$, so the cosets of $G_0$
permute the orbits transitively and there are finitely many of them. A finite partition of $P$
into closed sets is a partition into open sets; $P$ is connected; so there is one orbit and
$G_0$ is transitive. Given pure states $x$ and $y$, choose $g \in G_0$ with $gx = y$ and a
continuous path $g_t$ in $G_0$ from the identity to $g$; then $g_t x$ is a continuous reversible
transformation from $x$ to $y$.
\end{proof}

The proposition relocates the premise. Divisibility was asked for because a transformation group
could be infinite and discrete; but a group acting transitively on a continuum cannot be
countable, and a closed group acting transitively on a connected space cannot be disconnected in
the way that matters, so what the argument needs from the regime is not the availability of
$n$-th roots but the connectedness of the state space itself. That is a statement about values:
that the pure states are all the unit vectors of $\mathbb{C}^{K}$ up to phase, the space
$\mathbb{CP}^{K-1}$, connected for $K \ge 2$, and hence that the amplitudes of
\S\ref{sec:self-reference} fill their continuum. The two standing foils show the premise doing
its work. The stabiliser fragment of quantum theory has complex amplitudes and is contextual,
yet it has finitely many pure states, its outcome probabilities take only the values $0$,
$\tfrac12$ and $1$, and its reversible transformations form a finite group, the Clifford group,
so Hardy's fifth axiom fails in it; Spekkens's toy theory \autocite{spekkens2007toy} is the same
with real amplitudes, its ignorance rationed by a knowledge-balance principle. Neither has a
connected state space, and neither is a model of the regime described here, whose credences take
every value. Divisibility, which both also lack --- the phase gate has no square root in the Clifford group
--- is a symptom of the discreteness, not its cause. What both lack can be said exactly, and the
next subsection says it: each contains a transformation that fixes no state it admits.

\subsection{The premises, and what they mean}

Three premises carry Proposition~\ref{prop:connected-transitive}. Transitivity is the regime's
homogeneity: inside/outside equivalence leaves no pure state privileged, and a proper closed
invariant set of pure states would define a superselection observable, a classical record
predating any \textsc{fanout} (\S\ref{sec:fanout-boundary}). Connectedness of the state space is two statements, and only one of them is quantum. The first
is that the pure states are infinite and dense, with a fixed point for every transformation; the
second is that the state space is complete. Quantum theory over the algebraic numbers ---
amplitudes in the algebraic closure of $\mathbb{Q}$ within $\mathbb{C}$ --- has the first
without the second: it has $K = N^{2}$, the Born rule, superposition, entanglement and Bell
violation at the Tsirelson value, every unitary has an algebraic eigenvector, and it is
countable and totally disconnected. So the whole quantum content of Hardy's axiom lies in the
density, and the completion is the real-number convention of every probability theory. Hardy
adopts it before his fifth axiom is stated, in the first, which makes probabilities real numbers
and gives his classical theory a real simplex of mixtures; this framework adopts it at the
Shannon step, where credences are declared to range over a connected domain so that the
continuation of the logarithm is unique (\S\ref{sec:complex-amplitudes}). Given that convention,
density of the pure states is connectedness of their closure. Divisibility, which an earlier
version of this section asked for, is not needed: it follows from the conclusion, every element
of a compact connected Lie group having roots of every order.

Closedness is the completion, limit-realisability: a Cauchy sequence of available
transformations has its limit available. It cannot be derived, and the reason should be stated
rather than hidden. Any collection of first-order principles about states, transformations and
probabilities that has an infinite model has a countable one, and a countable subset of
$\mathbb{CP}^{K-1}$ is totally disconnected; the founding non-computability does not escape
this, since the arithmetical complex numbers form a countable algebraically closed field
containing non-computable numbers and closed under every limit it can itself describe. Hardy's
axiom and closedness are second-order, quantifying over paths and over sequences, and a
second-order principle is grounded as Dedekind completeness is grounded or not at all. The
grounds are these. The completion of a dense set under the finite-precision metric is unique and
depends on no choice, whereas any countable alternative is a choice of definability level, an
extra law with no ground and no test; and the framework's own result is that no observer's
descriptive level captures the regime, since by \S\ref{sec:self-reference} every observer's
state exceeds its own level and the inside/outside principle makes that a law rather than a
perspective. The burden therefore falls on any restriction, not on the completion, which is the
least committal option and carries no quantum content. What limit-realisability adds to the line
is still worth stating exactly, because it is the corrected form of a natural intuition.
Cantor's diagonal \autocite{cantor1891diagonal} does not build the continuum; its output is an
infinite decimal --- a Cauchy sequence of rationals --- which exists \emph{as a point} only
because the line is complete. The diagonal presupposes limit-realisability and then shows that
the completed line outruns every internal listing. And the character of the completion is the
framework's own: the computable reals are countable \autocite{Turing_1937}, hence of measure
zero, so once the transformation group contains its limits almost all of them are
non-computable. Non-computability is a property of the completed continuum, not the reason for
it. The constructive reading, on which only describable limits exist, remains consistent; it is
the countable theory above, and the two readings differ in nothing an experiment can certify.

Density is where the section's title is earned, and it is derived from the paper's own premise.
Write $(\mathrm{I})$ for the statement that every admissible transformation of the regime is a
phase of some admissible classical structure: diagonal, that is, in the record basis of some
observer the regime contains, a phased spider of that basis in the calculus of
\S\ref{sec:hopf-justification} \autocite{coeckeduncan2011interacting} (VERIFY-SUPPORT). It says
of transformations what the global phase of \S\ref{sec:self-reference} says of states: there are
no absolute ones. It follows from three things the paper already holds. The inside/outside
involution is a genuine involution, so the outside presentation of the inside is a presentation
of the same object with nothing lost: the outside can have complete knowledge of the inside, a
potential record, the joint state lying in a structure the outside carries. Time is operation
count, internal to the system that operates (\S\ref{sec:space-of-computations}), so the outside
does not share the inside's count. And before any \textsc{fanout} that count is not classically
recorded, so the outside's description is pure. Let $g$ be the transformation the inside
performs per operation. The outside's complete knowledge cannot depend on the inside's count,
and it is pure, so it is a pure state of the inside's clock and system that is stationary under
``advance the clock and apply $g$''; the stationary states of that map are clock Fourier modes
tensored with eigenvectors of $g$ and nothing else \autocite{pagewootters1983evolution}
(VERIFY-SUPPORT). So the outside holds a structure in which $g$ acts by phases, which is
$(\mathrm{I})$; and the relative phase in that structure is the outside's ignorance of the
inside's computational count, the same ignorance the global phase records one level down. The
consequences are a theorem. A phase of a structure fixes the structure's copyable states, so
every admissible transformation fixes an admissible pure state; a finite transitive action has
an element moving every point (Jordan), so the pure states are infinite; the transformation
carrying a stabiliser state to the fixed point of the composite of a $Z$-phase and an $X$-phase
turns through $\arccos(1/\sqrt{3})$, an irrational fraction of a turn, so the group has an
element of infinite projective order and its orbits are dense. With the completion,
Proposition~\ref{prop:connected-transitive} fires and Hardy's fifth axiom holds. The two foils
fail $(\mathrm{I})$ exactly where they fail continuity: fourteen of the twenty-four Clifford
operations are phases of no Pauli basis and fix no stabiliser state, their fixed points being
the magic states; the toy theory's three-cycles are phases of no epistemic basis; and in each
the outside's knowledge of the inside under the order-three protocol is either count-dependent
or a mixture over the count, so each makes the inside's time classical before any record exists.
The derivation of $(\mathrm{I})$ is an outline, not a proof, and it is what the open problem
below asks for.

\subsection{Why the cut matters physically}

The regime's continuity is not decorative. Reversibility alone, over complex amplitudes, licenses
the full invertible group $GL(n, \mathbb{C})$, and Aaronson showed what that buys: superluminal
signalling, and computational power at the level of \textsf{PP} \autocite{Aaronson2004h}. The cut
from $GL(n, \mathbb{C})$ down to the unitary group is the conserved quadratic form, and continuity
is what entitles the theory to make it: with the fifth axiom in hand, Hardy's theorem delivers the
complex Hilbert space, the quadratic norm and the unitary group in one stroke \autocite{Hardy2001a}
--- cited here, not re-proved. One honest caveat travels with the citation. Hardy's reconstruction
uses his remaining axioms besides --- composition and subspace structure --- and these must be
derived rather than imported: the composition ledger of \S\ref{sec:division-algebras} and the
locality-of-tomography argument of \S\ref{sec:self-reference} carry that duty, and it is discharged
there, not silently assumed here.

\subsection{Super-tasks at the boundary}

The premise just priced has a fifty-year-old stress test in the philosophy of mathematics, and the
framework passes it in an unexpected way. A super-task --- infinitely many operations completed in
finite time \autocite{Thomson_1954} --- presupposes its terminal moment as an actual point of the
timeline: the accumulation point of a generic sequence of step-times exists only in the completed
line, and the defining question, the state \emph{after all the steps}, is meaningful only under
limit-realisability of time. Super-tasks are downstream consumers of exactly the completeness at
issue in this section. Thomson's lamp then exhibits both sides of the \textsc{fanout} boundary at
once. On the record side, the lamp's description fixes its state at every finite stage and thereby
fixes no limit at all: any terminal value is consistent with the whole task
\autocite{Thomson_1954,benacerraf1962tasks}. That is not a puzzle about lamps; it is the proof that
the classical record algebra is not limit-closed, and cannot be. On the amplitude side the limit
exists --- a superposed fixed point, whose complex-probability treatment belongs to
\S\ref{sec:self-reference} --- because that side \emph{is} limit-closed, which is the fifth face
itself. The boundary is thereby the edge of limit-realisability: complete on the reversible side,
incomplete on the record side, with the lamp standing on both sides at once as the boundary's oldest
exhibit. A final clause dissolves rather than answers the classic question. In this framework time
is operation count (\S\ref{sec:space-of-computations}), and ``after infinitely many operations''
indexes no operational moment: the lamp packs its endless \textsc{fanout}s into the emergent
continuum, and asking for its record at the accumulation point asks the discrete contents for a
value that only the completed container possesses. On or off is not undecided; it is unasked.

\subsection{What remains, priced}

\begin{openproblem}[Every transformation a phase of an inside structure]\label{op:limit-realisability}
Proposition~\ref{prop:connected-transitive} reduces Hardy's fifth axiom to three premises:
connectedness of the pure states, transitivity, and closedness. Transitivity is argued above
from the absence of records. Closedness is the completion, adopted on the grounds stated above
and carrying no quantum content. Connectedness is density plus that completion, and density
follows, by Jordan's theorem and the infinite-order argument, from $(\mathrm{I})$: that every
admissible transformation of the regime is a phase of some admissible classical structure. What
remains open is the rigour of the derivation of $(\mathrm{I})$ from the paper's premise. Derive,
inside the fixed-point locus $\mathrm{Fix}(\sigma)$ of Proposition~\ref{prop:inside-outside},
that every endomap is a phase of a classical structure in $\mathrm{Fix}(\sigma)$, from the
self-representing object, the involution $\sigma$ as a genuine involution, and the reading of
the inside presentation as a sequence indexed by operation count with no record before
\textsc{fanout}. The outline above does this with the stationary-state construction of Page and
Wootters; a derivation passes if it makes that outline a theorem of the calculus, with the
outside's count-independent complete description shown to exist and to be pure from the
involution alone. It fails before any detail is checked if every premise it uses holds in the
stabiliser fragment of quantum theory, in Spekkens's toy theory \autocite{spekkens2007toy}, in
the real qubit, or in quantum theory over the Gaussian rationals, all of which violate
$(\mathrm{I})$ outright and all of which make the inside's operation count classical before any
record exists. It also fails if it imports the completion as a premise about transformations:
the completion is the real line and has no derivation, and a derivation of $(\mathrm{I})$ must
hold in the countable algebraic theory, where $(\mathrm{I})$ is true and continuity is not.
\end{openproblem}

The accounting, in one sentence: the axiom that every reconstruction programme pays is, in this
framework, one derivation posed, of Lawvere's shape --- that every endomap of the fixed-point
locus is a phase of a classical structure in it --- and one convention named, the completion,
which classical probability pays too. And
the hand-on is already visible: the complex amplitudes that this section's continuity underwrites
through Hardy's theorem, and the complex amplitudes that \S\ref{sec:self-reference} derives natively,
are one and the same $\mathbb{C}$ --- the floor of the ladder that \S\ref{sec:division-algebras}
climbs.

%      % =====================================================================
%  SECTION: The Division-Algebra Ladder
%  Drafted 2026-07-03. One of the three foundations sections
%  (sec:self-reference, sec:division-algebras, sec:fanout-boundary);
%  this one resolves the dangling \S\ref{sec:division-algebras} in the
%  Hopf-justification section, whose header names it "the
%  Cayley--Dickson / Hurwitz step".
%
%  PLACEMENT in main.tex (REORDERED per JS 2026-07-06): position 6,
%  opening Act II ("why stop at C?") right before its first consumer --
%      \input{self_reference_section_2026-07-04}
%      \input{continuity_section_2026-07-04}
%      \input{L2_metric_global_phase_section_2026_06_07}
%      \input{wigners_friend_in_the_hopf_picture_2026-06-13}
%      \input{derivation_of_the_quantum_postulates}
%      \input{division_algebras_section_2026-07-03}   <-- moved here
%      \input{spacetime_from_non-computability}       <-- first consumer
%                                                          (S^7 exclusion)
%      \input{hopf_justification_section_2026-06-08}
%  The section is position-portable: all adjacency phrasing was
%  retargeted to explicit \S\ref on 2026-07-06. The S^7 hand-off
%  sentence to sec:space-of-computations WRITTEN 2026-07-06 (spacetime
%  file read; the closing now hands first to its immediate consumer,
%  then to hopf-justification). Upload verifications COMPLETE: L2 and
%  quantum-postulates do not consume the ladder; WF's references are
%  forward-voiced or self-carrying.
%
%  DIVISION OF LABOUR (deliberate, to avoid overlap):
%   - sec:self-reference (undrafted): self-reference -> complex
%     amplitudes and the Born rule; why not a real base. This section
%     TAKES that as given and starts at C.
%   - THIS section: the algebra -- the doubling construction, the
%     ledger of losses-as-physics, the internal-algebras-not-a-new-QM
%     scoping (Adler pre-empt), the two appearances of the ladder
%     (register / measurement climb, stated lightly), and Hurwitz as
%     the algebraic termination with the committed zero-divisor
%     reading.
%   - sec:hopf-justification: the geometry (Adams, parallelisability,
%     the fibrations), the join of the two classifications, the
%     entanglement-sensitivity literature, the C-direction honest debt,
%     and the epistemic reading of termination.
%
%  COMMITTED READINGS USED (sources re-read 2026-07-03):
%   - Loss ledger and level dictionary: cayley_dickson_qubits_for_review
%     sec.5 ("each step loses a symmetry, and each lost symmetry is
%     identified with a physical content").
%   - Termination reading: grounding_argument_gauge_group ("the
%     sedenions have zero divisors and cannot support consistent
%     information flow around loops") -- leads; norm-composition-as-
%     probability given as its amplitude-level restatement.
%   - Internal-not-base scoping: cayley_dickson_qubits_for_review Q2 and
%     sec.3.2 caveat (Adler objection; C-qubit treated geometrically).
%
%  HONEST DEBT: ladder proved closed at the top (Hurwitz), derived at
%  the floor. The rung-compulsion revisit-trigger FIRED 2026-07-06:
%  sec:self-reference and op:sigma-functor now exist, so the thesis is
%  minted as op:ladder-compulsion (with the CD-conjugation-as-sigma
%  candidate route), inheriting op:sigma-functor's status.
%
%  AUDITED 2026-07-06 (final pass, third stop): all four citations
%  used; all cross-refs live; CD facts, dims, alternativity, Hurwitz
%  statement verified; Adler subsection delivers everything the
%  register's compressed pointer promises; file pure ASCII; no bare
%  "frame". Fixes: opening rewritten for both predecessors (continuity
%  was inserted between after drafting); "associator's non-vanishing is
%  mass itself" recalibrated to the register's debt-against-selection
%  defusal; the O->S ledger row and Hurwitz subsection reconciled with
%  the fourth-qubit forward remark ("the terminus closes the ladder of
%  forces; what survives past it is, at most, flavour").
%
%  CROSS-REFS: sec:self-reference [sibling dangle -- resolves when that
%  section is drafted], sec:hopf-justification, sec:hopf-information-bound,
%  sec:register, sec:entanglement-classes, sec:generations, sec:masses,
%  sec:wigners-friend. All others exist.
%  CITATIONS: existing keys baez2002octonions, hurwitz1898composition,
%  furey2018three. ONE new key: adler1995quaternionic (RIS in
%  division_algebras_refs_2026-07-03.ris).
%  British English; \autocite; no bullets in prose; the ledger is a
%  table because the comparison is the argument.
% =====================================================================

\section{The Division-Algebra Ladder}
\label{sec:division-algebras}

The foundations behind us end at amplitudes twice over: natively, where the self-referential
argument forced the description of a system into the system's own state space and priced
consistency at a number system whose parts can cancel --- complex amplitudes, and the Born rule
that squares them (\S\ref{sec:self-reference}) --- and independently, where the regime's
continuity delivers the same field through Hardy's theorem (\S\ref{sec:continuity}), the
convergence of the two branches being evidence rather than accident. This section asks the question
both routes leave open. Why
should the number system stop at $\mathbb{C}$? Ordinary quantum mechanics stops there by postulate;
a framework that has forbidden itself postulates must answer, and the answer is that the climb
neither stops at $\mathbb{C}$ nor continues forever. There are exactly two more rungs, each
purchased by surrendering an algebraic law and each paying for itself in physics, and the ladder is
closed at the top by a theorem of 1898. This section sets out the algebra of that ladder: the single
construction that climbs it, the ledger of what each rung costs and buys, and the theorem that ends
it. The geometry of the same ladder --- that each rung acts on a sphere, that exactly those spheres
are parallelisable, and that the algebraic and geometric classifications are one --- is the business
of the geometry section (\S\ref{sec:hopf-justification}).

\subsection{One construction climbs the whole ladder}

The rungs are not a miscellany; a single construction generates them all. Given an algebra with a
conjugation, the Cayley--Dickson construction builds its double on ordered pairs, with a twisted
multiplication in which the conjugate enters and a conjugation $(a,b)^{*} = (a^{*}, -b)$
\autocite{baez2002octonions}. (Conventions differ in the placement of the conjugations; all yield
isomorphic algebras.) Applied repeatedly it produces the sequence
$\mathbb{R} \to \mathbb{C} \to \mathbb{H} \to \mathbb{O} \to \mathbb{S}$, of real dimensions
$1, 2, 4, 8, 16$: the reals, the complex numbers, the quaternions, the octonions, the sedenions.
The ground floor is classical: $\mathbb{R}$ is ordered and positive-definite, the arithmetic of
classical probability. The first doubling is therefore the quantum move itself --- the passage from
probabilities to amplitudes that can interfere --- and it is not made here by choice: it is the
derivation of \S\ref{sec:self-reference}, already behind us. What this section must account for is
the rest of the climb.

\subsection{The ledger}

In mathematics the doublings are usually told as a story of degeneration: each algebra is worse
behaved than the last, and the sequence is studied for where it breaks. This framework reads the
same ledger in the opposite direction. Each rung is bought by surrendering a law, and each
surrendered law is exactly a piece of physics --- the losses are the assets, and the comparison is
sharp enough to be the argument:

\begin{center}
\begin{tabular}{llll}
\hline
doubling & dimension & what is lost & what the loss becomes \\
\hline
$\mathbb{R} \to \mathbb{C}$ & $1 \to 2$ & the order of the reals & phase and interference; amplitudes \\
 & & & that cancel; the Born rule (\S\ref{sec:self-reference}) \\
$\mathbb{C} \to \mathbb{H}$ & $2 \to 4$ & commutativity & non-commuting actions: weak isospin \\
 & & & $SU(2)$, the $B$ register (\S\ref{sec:register}) \\
$\mathbb{H} \to \mathbb{O}$ & $4 \to 8$ & associativity & the associator: colour, chirality, the \\
 & & & generations, and mass as debt \\
 & & & (\S\ref{sec:generations}, \S\ref{sec:masses}) \\
$\mathbb{O} \to \mathbb{S}$ & $8 \to 16$ & division itself & no further force: zero divisors appear \\
 & & (zero divisors) & and the ladder of sectors ends (below) \\
\hline
\end{tabular}
\end{center}

Walking the rows: the order lost at the first doubling is the possibility of interference --- parts
of a probability that can oppose one another --- which is where the derivation of \S\ref{sec:self-reference}
lives. The commutativity lost at the second is the existence of actions whose order matters; the
unit quaternions form the group $SU(2)$, and the framework houses weak isospin there, as the
internal algebra of the $B$ register whose order-sensitive, non-commuting dynamics the interaction sections develop
(\S\ref{sec:register}, \S\ref{sec:entanglement-classes}). The associativity lost at the third ---
the octonions retain the weaker property of alternativity \autocite{baez2002octonions} --- creates
the associator $[a,b,c] = (ab)c - a(bc)$, and the associator is where most of this paper's physics
is transacted: the stabiliser of a preferred complex direction inside the octonions' automorphisms
is the colour group (\S\ref{sec:hopf-justification}); the triality structure and its $K_3$
obstruction generate the three generations and the sterile channel (\S\ref{sec:generations}); and
the associator is the currency of mass --- the debt against the selected direction that
\S\ref{sec:masses} prices, owed by whatever action moves the selection and not by residence at
the rung, which is the defusal \S\ref{sec:register} makes explicit. The algebraic facts
in the third column are theorems; the identifications in the fourth are established in the sections
cited, each carrying its own epistemic flag there, and are not re-litigated here. What this table
asserts on its own account is only the correspondence's shape: one law surrendered per rung, one
sector of physics gained per law.

\subsection{Internal algebras, not a new quantum mechanics}

A reader versed in the quaternionic-quantum-mechanics literature will at this point raise a fair
objection: quantum theory over $\mathbb{H}$ as a foundation is known to be troubled --- the
tensor-product and composite-system difficulties set out by Adler and others
\autocite{adler1995quaternionic} --- and quantum theory over $\mathbb{O}$ founders earlier still,
since a non-associative algebra does not furnish a Hilbert space at all. The framework's position
must therefore be stated exactly, because it does not make the claim the objection targets. The base
of the physical Hilbert space here is $\mathbb{C}$ and remains $\mathbb{C}$; the higher algebras
enter as the \emph{internal} algebras of the register --- $\mathbb{H}$ for the $B$ qubit,
$\mathbb{O}$ for the $C$ qubit (\S\ref{sec:register}) --- in a tensor structure whose outer theory is
ordinary complex quantum mechanics. This is the same move by which the division algebras enter the
constructions of Furey and the related programmes without any proposal to rebase quantum mechanics
itself \autocite{furey2018three}. For the octonionic register the point is sharper still:
$\mathbb{O}^{2}$ is not a Hilbert space, and the framework accordingly treats the $C$ qubit not in
the Dirac--von Neumann manner but geometrically, as the sphere $S^{15}$ with its Hopf structure ---
which is precisely the form in which \S\ref{sec:hopf-justification} needs it. Nothing in this paper therefore
stands or falls with quaternionic or octonionic quantum mechanics as foundations; the ladder is
climbed by the registers' internal algebras, not by the theory's base field.

\subsection{The same ladder, twice}

The ladder appears in this framework in two guises, and it will help the reader to have them
distinguished now. Statically, it is the anatomy of a single relation: the three qubits of the
register take their values one rung apart --- $A$ over $\mathbb{C}$, $B$ over $\mathbb{H}$, $C$ over
$\mathbb{O}$ --- so that every particle carries the whole ladder as its internal structure
(\S\ref{sec:register}). Dynamically, it is the shape of measurement: a single qubit's pure states
fill $S^{3}$, an entangling measurement carries the joint state into $S^{7}$, and a third party into
$S^{15}$ --- the correspondence between qubit entanglement and the Hopf fibrations that
\S\ref{sec:hopf-justification} takes from a mature literature --- so that each act of measurement ascends one rung --- the climb the reader has already made,
narrated in \S\ref{sec:wigners-friend} --- and every completed relation is a terminated climb, which is
the reading \S\ref{sec:hopf-information-bound} gives the tower's endpoint. One ladder, seen twice:
as the anatomy of a relation and as the dynamics of relating. The framework regards this double
appearance not as a coincidence but as the same structure viewed statically and in motion, and the
sections cited carry that identification's weight.

\subsection{Hurwitz closes the ladder}

The construction itself would double forever; what stops it is a classification theorem. By
Hurwitz's theorem, the normed division algebras --- those in which the norm composes,
$\lvert ab\rvert = \lvert a\rvert\,\lvert b\rvert$, and no non-zero elements multiply to zero ---
are exactly $\mathbb{R}$, $\mathbb{C}$, $\mathbb{H}$ and $\mathbb{O}$
\autocite{hurwitz1898composition,baez2002octonions}. The next doubling, the sedenions, has zero
divisors: non-zero elements whose product vanishes. The framework's reading of this failure is the
one that matters physically. Zero divisors cannot support consistent information flow around loops
--- a cycle of compositions can silently annihilate what it carries --- and a framework in which
closed computational paths are load-bearing (\S\ref{sec:hopf-justification},
\S\ref{sec:space-of-computations}) cannot be built on an algebra that permits it. At the level of
amplitudes the same failure reads: the composition law
$\lvert ab\rvert = \lvert a\rvert\,\lvert b\rvert$, which carries probability multiplicatively
through every rung of the ladder, fails beyond the octonions, so that two non-vanishing amplitudes
could compose to nothing and probability would be destroyed by mere composition. The last algebra
fit to carry amplitudes is $\mathbb{O}$. One rung past the terminus is not entirely silent,
however: the sedenions' internal structure --- their automorphism group and their octonion
subalgebras --- is the recorded candidate arena for tagging the generations
(\S\ref{sec:interactions}), and it enters exactly as this section's scoping permits, internally
and never as a base. The terminus closes the ladder of forces; what survives past it is, at most,
flavour. This is the algebraic half of the termination; the
geometric half --- that no sphere above $S^{7}$ is parallelisable --- and the theorem joining the
two halves are the opening business of \S\ref{sec:hopf-justification}, and the epistemic reading of the same endpoint,
as a bound on the information a single relation can carry, is given in
\S\ref{sec:hopf-information-bound}.

\subsection{The honest debt}

As in the geometry section, we close by stating plainly what is proved, what is established, and what is
owed. Proved: the construction, the losses, and the termination are theorems --- the Cayley--Dickson
doubling and its ledger of surrendered laws are standard algebra, and Hurwitz's classification is
among the durable results of the subject. Established: the physical cashing of each rung --- the
Born rule at $\mathbb{C}$, weak isospin at $\mathbb{H}$, colour, generations and mass at
$\mathbb{O}$ --- is carried out in the sections cited above, each under its own epistemic flag. Owed,
and stated without disguise: the ladder as presented here is closed from above and derived at its
floor, but it is not yet derived rung by rung from below. Hurwitz forbids a further rung; the
argument of \S\ref{sec:self-reference} forces the first doubling, from the ordered arithmetic of
classical probability to amplitudes that interfere. That the same self-referential closure, applied
to a system whose amplitudes already take values at one rung, compels the next --- that relational
closure at one level forces precisely the Cayley--Dickson double of that level --- is the
framework's reading of the climb, and it is a thesis, not a theorem --- but the formal apparatus
against which it can be posed now exists, and we pose it.

\begin{openproblem}[The compulsion of the intermediate rungs]\label{op:ladder-compulsion}
Derive, within the categorical apparatus of Open Problem~\ref{op:sigma-functor}, that
self-referential closure at one rung compels the Cayley--Dickson double of that rung, so that the
intermediate steps $\mathbb{C} \to \mathbb{H} \to \mathbb{O}$ follow from the same principle
that forces the floor. The candidate route is the construction's own shape: the Cayley--Dickson
pair with its conjugation, $(a, b)^{*} = (a^{*}, -b)$, is an involution-equipped doubling, and the
route reads it as the algebraic shadow of adjoining the describer --- closure of level-$A$
descriptions under the inside/outside exchange requiring exactly the twisted pair. The problem
inherits the status of Open Problem~\ref{op:sigma-functor}, whose apparatus it presupposes.
\end{openproblem}

If the derivation lands, the whole ladder becomes a consequence of the same self-reference that
opened it. Until then, the honest statement is that the ladder's necessity is established at its
ceiling and at its floor, while the compulsion of the intermediate rungs is inherited from the
physics they are shown to carry rather than derived in advance. What this section hands on is, accordingly, an algebraic ladder with a floor, a ceiling and a
ledger. Its first consumer is immediate: \S\ref{sec:space-of-computations} spends the terminal
fact --- the non-associativity of $S^{7}$'s own multiplication --- as geometry, excluding the
seven-sphere as a ruler and fixing the dimension of space at three.
\S\ref{sec:hopf-justification} then supplies the geometry of the fibrations themselves, and the
join.

%  (Architecture call SETTLED by the sync manifest: standalone.)
%
%  CITATIONS -- existing keys (as supplied by JS 2026-07-04):
%    Hardy2001a, Aaronson2004h, Thomson_1954, Lawvere_1969,
%    Turing_1937, muller_p_2011, Brukner2009.
%  NEW keys (RIS in continuity_refs_2026-07-04.ris):
%    cantor1891diagonal, benacerraf1962tasks, broeckertomdieck1985compact.
%  Deliberately excluded per JS: Chiribella--D'Ariano--Perinotti;
%  Bargmann/Wigner's theorem. Yanofsky2003 now paired on the
%  Lawvere-shaped clause (the pending-key excuse was void: the key
%  exists and sec:self-reference uses it).
%
%  CROSS-REFS: sec:self-reference now DRAFTED (2026-07-04, later the
%  same day) -- no longer dangling; only sec:fanout-boundary (and its
%  sibling sec:fanout-direction) -- DRAFTED 2026-07-06.
%
%  AUDITED 2026-07-06 (final pass, second stop): both proofs
%  re-verified by hand (Lagrange divisibility sentence exact; the
%  Baire countable-escape chain sound); caveat-discharge with
%  sec:self-reference confirmed two-way (already present); file pure
%  ASCII; no pre-definition "frame"; all citations used as declared.
%  Fixes: Yanofsky2003 paired on the Lawvere clause; the op's two
%  candidate routes anchored to printed propositions (f1-dictionary
%  for the divisibility premise's referent, inside-outside for the
%  closedness demand) -- statuses unchanged, referents sharpened;
%  header staleness corrected. Deliberate FORWARD references --
%  prop:f1-dictionary and sec:hopf-justification (defined later in the
%  paper; fine under two-pass LaTeX and flagged in prose),
%  sec:division-algebras, sec:space-of-computations.
%
%  DECISIONS ENCODED: op:limit-realisability minted (the owed arrow,
%  priced: exactly two derivations, divisibility and closedness);
%  Lagrange used twice (finite exclusion; component-group triviality);
%  the (1+i)/2 lamp value NOT stated here -- "a superposed fixed point"
%  with the complex-probability treatment deferred to
%  sec:self-reference; the two-branch over-determination stated to
%  pre-empt a circularity reading; Thomson spelling throughout.
% =====================================================================

\section{Continuity Is Not an Axiom}
\label{sec:continuity}

Hardy derives quantum theory from five axioms \autocite{Hardy2001a}. Four of them are satisfied
by classical probability theory; the fifth --- that a \emph{continuous} reversible
transformation connects any two pure states --- is the entire difference, the single postulate
that turns the classical simplex into the complex Hilbert space. Every reconstruction programme
pays an equivalent toll: continuous reversibility in Masanes and M\"{u}ller
\autocite{mullerp2011}, the corresponding postulate in Daki\'{c} and Brukner
\autocite{Brukner2009}. The axiom is load-bearing, and in every programme it is unexplained: it
is where the quantumness is put in by hand. The claim of this section is that in the present
framework it is not an axiom. Continuity is the fifth face of the pre-measurement regime that self-reference forces, and what
stands between the framework and a full derivation is not a postulate about values but one
premise about transformations --- that every transformation the inside performs has a stationary
outside presentation --- derived in outline below from the inside/outside involution of
\S\ref{sec:self-reference} and stated as an open problem at the end of the section with nothing
hidden in it. The continuum itself, the completion that every probability theory adopts, is
named for what it is: the real line, shared with classical probability and carrying no quantum
content.

One structural remark prevents a misreading. The framework reaches complex amplitudes natively in
\S\ref{sec:self-reference}, by the negative-probability route; the present section reaches the same
destination by an independent branch --- the regime's continuity, then Hardy's theorem. The two
branches converging on one structure is over-determination, not circularity, and the convergence is
itself evidence that the destination is forced.

\subsection{Five faces of one regime}

The regime before any measurement --- before any \textsc{fanout} --- can be characterised five
equivalent ways: negative probability is available; anti-events exist, every event being undoable;
no classical record exists; no preferred basis is selected; and a continuous reversible
transformation connects any two pure states. These are not five premises to be combined. They are
one regime under five descriptions, and \textsc{fanout} (\S\ref{sec:fanout-boundary}) is the single
event that fails all five at once: it writes the record, selects the basis, kills the anti-events,
discretises the transformation group, and switches off the negativity. The equivalence of the first
two faces is the anti-event identity of \S\ref{sec:self-reference}; the third is the boundary's
definition. The work of this section is the equivalence of the last two faces, and then the route
from the regime's existence to its continuity. The fourth face has, since
Proposition~\ref{prop:f1-dictionary} (\S\ref{sec:hopf-justification} --- a deliberate forward
reference), an exact form: a basis-selection is the commitment of a classical structure, a lifted
antipodal skeleton, so ``no preferred basis'' means precisely that no copy structure has fired.

\subsection{Basis selection is the collapse of continuity}

\begin{proposition}[Continuity forbids a selected basis]\label{prop:basis-discreteness}
If a finite set of pure states is privileged --- a preferred basis --- then every reversible
transformation respecting the privilege carries the set to itself and acts on it as a permutation;
the respecting transformations therefore realise only a finite group of motions among the
distinguished states, and no continuous family of them connects distinct pure states of the set.
Contrapositively, continuous reversibility entails that no basis is selected.
\end{proposition}

\begin{proof}
A privileged finite set is by definition carried to itself by every structure-preserving reversible
map, whose induced action on the set is a permutation of finitely many points. A continuous path of
such maps induces a continuous path in a finite discrete set, which is constant. Hence a continuous
reversible family moving one distinguished pure state to another respects no such privilege, and if
every pair of pure states is so connected, no finite set is privileged.
\end{proof}

The converse is where the historical difficulty lies: the absence of a privileged finite set does
not by itself force continuity, for the transformation group could in principle be infinite and
discrete. Two structural properties close every such escape at once.

\begin{proposition}[Divisible and closed implies connected]\label{prop:divisible-closed}
Let $G$ be the group of reversible transformations of the pre-measurement regime, acting on the
state body. If $G$ is divisible --- every element has an $n$-th root for every $n$ --- and
topologically closed, then $G$ is a compact connected Lie group; and if $G$ moreover acts
transitively on the pure states, then a continuous reversible transformation connects any two of
them: Hardy's fifth axiom.
\end{proposition}

\begin{proof}
The state body is compact and convex with nonempty interior, so its symmetry group is compact, and
a closed subgroup of a compact matrix group is a compact Lie group with finitely many connected
components \autocite{broeckertomdieck1985compact}. The component group $G/G_{0}$ is a quotient of a
divisible group and hence divisible; but a nontrivial finite group is never divisible --- with
$\lvert H\rvert = m$ every element satisfies $h^{m} = e$, so no element other than the identity
possesses an $m$-th root --- and $G/G_{0}$ is finite, so it is trivial and $G = G_{0}$ is connected.
A connected Lie group is path-connected, and transitivity then supplies the continuous reversible
path between any two pure states.
\end{proof}

The same one-line fact of Lagrange closes the remaining escapes. Finite groups are excluded
directly, being nontrivial and never divisible. Countable groups are excluded because a
countable closed subgroup of a compact group is compact, has an isolated point by the Baire
category theorem, and is therefore discrete and finite. Uncountable disconnected groups are
excluded by the component-group argument of the proposition. Every route to a discrete
transformation group therefore passes through a failure of divisibility or a failure of
closedness. Divisibility, however, asks for more than the conclusion needs, and the next
proposition removes it.

\begin{proposition}[Continuity from a connected state space]\label{prop:connected-transitive}
Let the pure states of a finite-dimensional system form a connected space $P$, and let the group
$G$ of reversible transformations of the pre-measurement regime be topologically closed and act
transitively on $P$. Then the identity component $G_0$ already acts transitively on $P$, and a
continuous reversible transformation connects any two pure states: Hardy's fifth axiom holds,
with no divisibility assumed.
\end{proposition}

\begin{proof}
As in Proposition~\ref{prop:divisible-closed}, closedness makes $G$ a compact Lie group with
finitely many connected components \autocite{broeckertomdieck1985compact}, so $G_0$ is a closed
normal subgroup of finite index and is path-connected. Each $G_0$-orbit in $P$ is the continuous
image of a compact group, hence compact, hence closed. Since $G$ is transitive and $G_0$ is
normal, the $G_0$-orbit of $gx$ is $g$ applied to the $G_0$-orbit of $x$, so the cosets of $G_0$
permute the orbits transitively and there are finitely many of them. A finite partition of $P$
into closed sets is a partition into open sets; $P$ is connected; so there is one orbit and
$G_0$ is transitive. Given pure states $x$ and $y$, choose $g \in G_0$ with $gx = y$ and a
continuous path $g_t$ in $G_0$ from the identity to $g$; then $g_t x$ is a continuous reversible
transformation from $x$ to $y$.
\end{proof}

The proposition relocates the premise. Divisibility was asked for because a transformation group
could be infinite and discrete; but a group acting transitively on a continuum cannot be
countable, and a closed group acting transitively on a connected space cannot be disconnected in
the way that matters, so what the argument needs from the regime is not the availability of
$n$-th roots but the connectedness of the state space itself. That is a statement about values:
that the pure states are all the unit vectors of $\mathbb{C}^{K}$ up to phase, the space
$\mathbb{CP}^{K-1}$, connected for $K \ge 2$, and hence that the amplitudes of
\S\ref{sec:self-reference} fill their continuum. The two standing foils show the premise doing
its work. The stabiliser fragment of quantum theory has complex amplitudes and is contextual,
yet it has finitely many pure states, its outcome probabilities take only the values $0$,
$\tfrac12$ and $1$, and its reversible transformations form a finite group, the Clifford group,
so Hardy's fifth axiom fails in it; Spekkens's toy theory \autocite{spekkens2007toy} is the same
with real amplitudes, its ignorance rationed by a knowledge-balance principle. Neither has a
connected state space, and neither is a model of the regime described here, whose credences take
every value. Divisibility, which both also lack --- the phase gate has no square root in the Clifford group
--- is a symptom of the discreteness, not its cause. What both lack can be said exactly, and the
next subsection says it: each contains a transformation that fixes no state it admits.

\subsection{The premises, and what they mean}

Three premises carry Proposition~\ref{prop:connected-transitive}. Transitivity is the regime's
homogeneity: inside/outside equivalence leaves no pure state privileged, and a proper closed
invariant set of pure states would define a superselection observable, a classical record
predating any \textsc{fanout} (\S\ref{sec:fanout-boundary}). Connectedness of the state space is two statements, and only one of them is quantum. The first
is that the pure states are infinite and dense, with a fixed point for every transformation; the
second is that the state space is complete. Quantum theory over the algebraic numbers ---
amplitudes in the algebraic closure of $\mathbb{Q}$ within $\mathbb{C}$ --- has the first
without the second: it has $K = N^{2}$, the Born rule, superposition, entanglement and Bell
violation at the Tsirelson value, every unitary has an algebraic eigenvector, and it is
countable and totally disconnected. So the whole quantum content of Hardy's axiom lies in the
density, and the completion is the real-number convention of every probability theory. Hardy
adopts it before his fifth axiom is stated, in the first, which makes probabilities real numbers
and gives his classical theory a real simplex of mixtures; this framework adopts it at the
Shannon step, where credences are declared to range over a connected domain so that the
continuation of the logarithm is unique (\S\ref{sec:complex-amplitudes}). Given that convention,
density of the pure states is connectedness of their closure. Divisibility, which an earlier
version of this section asked for, is not needed: it follows from the conclusion, every element
of a compact connected Lie group having roots of every order.

Closedness is the completion, limit-realisability: a Cauchy sequence of available
transformations has its limit available. It cannot be derived, and the reason should be stated
rather than hidden. Any collection of first-order principles about states, transformations and
probabilities that has an infinite model has a countable one, and a countable subset of
$\mathbb{CP}^{K-1}$ is totally disconnected; the founding non-computability does not escape
this, since the arithmetical complex numbers form a countable algebraically closed field
containing non-computable numbers and closed under every limit it can itself describe. Hardy's
axiom and closedness are second-order, quantifying over paths and over sequences, and a
second-order principle is grounded as Dedekind completeness is grounded or not at all. The
grounds are these. The completion of a dense set under the finite-precision metric is unique and
depends on no choice, whereas any countable alternative is a choice of definability level, an
extra law with no ground and no test; and the framework's own result is that no observer's
descriptive level captures the regime, since by \S\ref{sec:self-reference} every observer's
state exceeds its own level and the inside/outside principle makes that a law rather than a
perspective. The burden therefore falls on any restriction, not on the completion, which is the
least committal option and carries no quantum content. What limit-realisability adds to the line
is still worth stating exactly, because it is the corrected form of a natural intuition.
Cantor's diagonal \autocite{cantor1891diagonal} does not build the continuum; its output is an
infinite decimal --- a Cauchy sequence of rationals --- which exists \emph{as a point} only
because the line is complete. The diagonal presupposes limit-realisability and then shows that
the completed line outruns every internal listing. And the character of the completion is the
framework's own: the computable reals are countable \autocite{Turing_1937}, hence of measure
zero, so once the transformation group contains its limits almost all of them are
non-computable. Non-computability is a property of the completed continuum, not the reason for
it. The constructive reading, on which only describable limits exist, remains consistent; it is
the countable theory above, and the two readings differ in nothing an experiment can certify.

Density is where the section's title is earned, and it is derived from the paper's own premise.
Write $(\mathrm{I})$ for the statement that every admissible transformation of the regime is a
phase of some admissible classical structure: diagonal, that is, in the record basis of some
observer the regime contains, a phased spider of that basis in the calculus of
\S\ref{sec:hopf-justification} \autocite{coeckeduncan2011interacting} (VERIFY-SUPPORT). It says
of transformations what the global phase of \S\ref{sec:self-reference} says of states: there are
no absolute ones. It follows from three things the paper already holds. The inside/outside
involution is a genuine involution, so the outside presentation of the inside is a presentation
of the same object with nothing lost: the outside can have complete knowledge of the inside, a
potential record, the joint state lying in a structure the outside carries. Time is operation
count, internal to the system that operates (\S\ref{sec:space-of-computations}), so the outside
does not share the inside's count. And before any \textsc{fanout} that count is not classically
recorded, so the outside's description is pure. Let $g$ be the transformation the inside
performs per operation. The outside's complete knowledge cannot depend on the inside's count,
and it is pure, so it is a pure state of the inside's clock and system that is stationary under
``advance the clock and apply $g$''; the stationary states of that map are clock Fourier modes
tensored with eigenvectors of $g$ and nothing else \autocite{pagewootters1983evolution}
(VERIFY-SUPPORT). So the outside holds a structure in which $g$ acts by phases, which is
$(\mathrm{I})$; and the relative phase in that structure is the outside's ignorance of the
inside's computational count, the same ignorance the global phase records one level down. The
consequences are a theorem. A phase of a structure fixes the structure's copyable states, so
every admissible transformation fixes an admissible pure state; a finite transitive action has
an element moving every point (Jordan), so the pure states are infinite; the transformation
carrying a stabiliser state to the fixed point of the composite of a $Z$-phase and an $X$-phase
turns through $\arccos(1/\sqrt{3})$, an irrational fraction of a turn, so the group has an
element of infinite projective order and its orbits are dense. With the completion,
Proposition~\ref{prop:connected-transitive} fires and Hardy's fifth axiom holds. The two foils
fail $(\mathrm{I})$ exactly where they fail continuity: fourteen of the twenty-four Clifford
operations are phases of no Pauli basis and fix no stabiliser state, their fixed points being
the magic states; the toy theory's three-cycles are phases of no epistemic basis; and in each
the outside's knowledge of the inside under the order-three protocol is either count-dependent
or a mixture over the count, so each makes the inside's time classical before any record exists.
The derivation of $(\mathrm{I})$ is an outline, not a proof, and it is what the open problem
below asks for.

\subsection{Why the cut matters physically}

The regime's continuity is not decorative. Reversibility alone, over complex amplitudes, licenses
the full invertible group $GL(n, \mathbb{C})$, and Aaronson showed what that buys: superluminal
signalling, and computational power at the level of \textsf{PP} \autocite{Aaronson2004h}. The cut
from $GL(n, \mathbb{C})$ down to the unitary group is the conserved quadratic form, and continuity
is what entitles the theory to make it: with the fifth axiom in hand, Hardy's theorem delivers the
complex Hilbert space, the quadratic norm and the unitary group in one stroke \autocite{Hardy2001a}
--- cited here, not re-proved. One honest caveat travels with the citation. Hardy's reconstruction
uses his remaining axioms besides --- composition and subspace structure --- and these must be
derived rather than imported: the composition ledger of \S\ref{sec:division-algebras} and the
locality-of-tomography argument of \S\ref{sec:self-reference} carry that duty, and it is discharged
there, not silently assumed here.

\subsection{Super-tasks at the boundary}

The premise just priced has a fifty-year-old stress test in the philosophy of mathematics, and the
framework passes it in an unexpected way. A super-task --- infinitely many operations completed in
finite time \autocite{Thomson_1954} --- presupposes its terminal moment as an actual point of the
timeline: the accumulation point of a generic sequence of step-times exists only in the completed
line, and the defining question, the state \emph{after all the steps}, is meaningful only under
limit-realisability of time. Super-tasks are downstream consumers of exactly the completeness at
issue in this section. Thomson's lamp then exhibits both sides of the \textsc{fanout} boundary at
once. On the record side, the lamp's description fixes its state at every finite stage and thereby
fixes no limit at all: any terminal value is consistent with the whole task
\autocite{Thomson_1954,benacerraf1962tasks}. That is not a puzzle about lamps; it is the proof that
the classical record algebra is not limit-closed, and cannot be. On the amplitude side the limit
exists --- a superposed fixed point, whose complex-probability treatment belongs to
\S\ref{sec:self-reference} --- because that side \emph{is} limit-closed, which is the fifth face
itself. The boundary is thereby the edge of limit-realisability: complete on the reversible side,
incomplete on the record side, with the lamp standing on both sides at once as the boundary's oldest
exhibit. A final clause dissolves rather than answers the classic question. In this framework time
is operation count (\S\ref{sec:space-of-computations}), and ``after infinitely many operations''
indexes no operational moment: the lamp packs its endless \textsc{fanout}s into the emergent
continuum, and asking for its record at the accumulation point asks the discrete contents for a
value that only the completed container possesses. On or off is not undecided; it is unasked.

\subsection{What remains, priced}

\begin{openproblem}[Every transformation a phase of an inside structure]\label{op:limit-realisability}
Proposition~\ref{prop:connected-transitive} reduces Hardy's fifth axiom to three premises:
connectedness of the pure states, transitivity, and closedness. Transitivity is argued above
from the absence of records. Closedness is the completion, adopted on the grounds stated above
and carrying no quantum content. Connectedness is density plus that completion, and density
follows, by Jordan's theorem and the infinite-order argument, from $(\mathrm{I})$: that every
admissible transformation of the regime is a phase of some admissible classical structure. What
remains open is the rigour of the derivation of $(\mathrm{I})$ from the paper's premise. Derive,
inside the fixed-point locus $\mathrm{Fix}(\sigma)$ of Proposition~\ref{prop:inside-outside},
that every endomap is a phase of a classical structure in $\mathrm{Fix}(\sigma)$, from the
self-representing object, the involution $\sigma$ as a genuine involution, and the reading of
the inside presentation as a sequence indexed by operation count with no record before
\textsc{fanout}. The outline above does this with the stationary-state construction of Page and
Wootters; a derivation passes if it makes that outline a theorem of the calculus, with the
outside's count-independent complete description shown to exist and to be pure from the
involution alone. It fails before any detail is checked if every premise it uses holds in the
stabiliser fragment of quantum theory, in Spekkens's toy theory \autocite{spekkens2007toy}, in
the real qubit, or in quantum theory over the Gaussian rationals, all of which violate
$(\mathrm{I})$ outright and all of which make the inside's operation count classical before any
record exists. It also fails if it imports the completion as a premise about transformations:
the completion is the real line and has no derivation, and a derivation of $(\mathrm{I})$ must
hold in the countable algebraic theory, where $(\mathrm{I})$ is true and continuity is not.
\end{openproblem}

The accounting, in one sentence: the axiom that every reconstruction programme pays is, in this
framework, one derivation posed, of Lawvere's shape --- that every endomap of the fixed-point
locus is a phase of a classical structure in it --- and one convention named, the completion,
which classical probability pays too. And
the hand-on is already visible: the complex amplitudes that this section's continuity underwrites
through Hardy's theorem, and the complex amplitudes that \S\ref{sec:self-reference} derives natively,
are one and the same $\mathbb{C}$ --- the floor of the ladder that \S\ref{sec:division-algebras}
climbs.

%      % =====================================================================
%  SECTION: The Division-Algebra Ladder
%  Drafted 2026-07-03. One of the three foundations sections
%  (sec:self-reference, sec:division-algebras, sec:fanout-boundary);
%  this one resolves the dangling \S\ref{sec:division-algebras} in the
%  Hopf-justification section, whose header names it "the
%  Cayley--Dickson / Hurwitz step".
%
%  PLACEMENT in main.tex (REORDERED per JS 2026-07-06): position 6,
%  opening Act II ("why stop at C?") right before its first consumer --
%      % EDITED 2026-10-29 scripts/patch_continuity_tower_premise_2026-10-29T1500.py
% EDITED 2026-09-20 scripts/patch_born_cost_ledger_v2_2026-09-20T0900.py
% =====================================================================
%  SECTION: Self-Reference: One Principle
%  Drafted 2026-07-04. The paper's foundational section; resolves the
%  largest dangling label in the manuscript (sec:self-reference).
%
%  PLACEMENT in main.tex: FIRST of the foundations chain --
%      \input{self_reference_section_2026-07-04}
%      \input{continuity_section_2026-07-04}
%      \input{division_algebras_section_2026-07-03}
%      \input{hopf_justification_section_2026-06-08}
%
%  SOURCES (re-read 2026-07-04): lawvere_inside_outside_fixed_point_
%  theorem_2026-05-04T1806.md (the spine: the two-principles reduction,
%  the theorem, the sigma candidate, the open items -> op minted per
%  D9); research_summary_neg_prob.md (negativity route; lamp (1+i)/2);
%  dynamics_companion_fanout_unitarity (legs 1a/1c; anti-events; the
%  Brukner-Zeilinger rebuttal per D6); Pure_states_and_mixed_states
%  (local tomography derived; QBism contrast per D7).
%
%  DECISIONS ENCODED (JS 2026-07-04): D1 title as below; D2 Wheeler1990
%  cited here as the title anchor; D3 lineage Popper -> Breuer ->
%  Svozil -> Szangolies (Dalla Chiara omitted); D4 Frauchiger_2018
%  cited here, dissolved in sec:wigners-friend; D5 Kirkwood-Dirac
%  heritage via the Arvidsson-Shukur 2024 review as one-stop; D6
%  Brukner-Zeilinger + the framework rebuttal IN ("Shannon is
%  post-FANOUT"); D7 QBism contrast IN; D8 Aaronson private
%  communication OMITTED (published island paper only); D9 one op
%  minted (op:sigma-functor), cross-linked to op:hopf-frobenius.
%
%  CITATIONS -- existing keys: Wheeler1990, small_2006, Godel1931,
%  Lawvere_1969, Yanofsky2003, CoeckeKissinger2017, Rovelli1996a,
%  Popper1950a, Popper1950b, Breuer1995, svozil_k_2018,
%  szangolies_j_2018, Frauchiger_2018, dirac1942, Feynman_r_82,
%  Feynman1987, Youssef1994 [SWAP-FLAG: intended paper is "Quantum
%  mechanics as complex probability theory", Mod. Phys. Lett. A 9
%  (1994) 2571 -- if that lives under youssef1994quantum, swap],
%  caticha_a_2006, goyal_knuth_skilling_2010, Aaronson2004h, Hardy2001a,
%  muller_p_2011, Thomson_1954, benacerraf1962tasks, cantor1891diagonal,
%  Turing_1937.
%  NEW keys (RIS in self_reference_refs_2026-07-04.ris, all verified
%  this session): tikochinsky1988feynman, renou2021real, chen2022ruling,
%  wu2022refutation, hoffreumonwoods2026real, arvidssonshukur2024kirkwood,
%  bruknerzeilinger2001shannon, fuchsschack2013qbism,
%  shannon1948mathematical.
%
%  CROSS-REFS: forward (fine under two-pass LaTeX): sec:continuity,
%  sec:division-algebras, sec:hopf-justification, prop:f1-dictionary,
%  sec:wigners-friend, op:limit-realisability, op:hopf-frobenius.
%  Still dangling after this section: sec:fanout-boundary (and its
%  sibling label sec:fanout-direction) -- both DRAFTED 2026-07-06.
%
%  AUDITED 2026-07-06 (section-by-section final pass, first stop):
%  proof of prop:inside-outside re-verified by hand; attribution chain
%  clean (Burgin/Khrennikov correctly absent); all header-declared
%  citations confirmed used; lamp arithmetic and antecedents verified;
%  FANOUT smallcaps consistent; no pre-definition use of "frame".
%  Fixes: stray Unicode em-dash (QBism paragraph, the file's only
%  non-ASCII character) -> house ---; trace-blind two-way stitch added
%  to the Born subsection ("one privacy, two duties");
%  op:sigma-functor gains its first standing constraint (F1
%  inheritance via the op:hopf-frobenius discharge) and its previously
%  one-directional local-tomography cargo made explicit.
% =====================================================================

\section{Self-Reference: One Principle}
\label{sec:self-reference-2}

Wheeler asked for it from bit: for physics derived from information, ``every it --- every particle,
every field of force'' obtained from binary answers to yes-or-no questions \autocite{Wheeler1990}.
This paper's route to that programme runs through G\"{o}del: the single principle from which
everything else here is derived, deferred, or priced is \emph{self-reference} --- the fact that any
complete physical description, made from inside the world it describes, must describe the describer.
This section states the principle, proves that the framework's second historical commitment follows
from it as a theorem rather than standing beside it as an axiom, and then extracts the price of the
principle in four instalments: probabilities whose parts can cancel, amplitudes valued in
$\mathbb{C}$, the Born rule, and the local tomography of composite systems. What the section cannot
yet derive is stated as precisely as what it can, in one numbered problem and in the ledger at the
end.

\subsection{Events, not objects}

The framework's ontology is relational: the fundamental entities are events and the relations
between them, and objects are the bookkeeping of relational invariants --- what remains stable
across many relations. The kinship with relational quantum mechanics \autocite{Rovelli1996a} is
genuine and acknowledged: facts are facts relative to the systems that register them. The
difference, developed over the rest of this paper, is that the present framework grounds the
relational structure in a single logical principle and then carries it, through the division
algebras and the Hopf fibrations, all the way to the gauge group and the particle spectrum ---
relations are not only the ontology but the derivation engine. On this reading a measurement is not
the revelation of a pre-existing property; it is the creation of a relation, and the later sections
of this paper read that creation geometrically (\S\ref{sec:hopf-justification}) and even
gravitationally. Here we need only the logical core: whatever describes the world is an event in
the world it describes.

\subsection{One principle, and a theorem}

The framework historically carried two foundational commitments: self-reference, and the
\emph{Inside/Outside Equivalence Principle} --- that the laws of physics must be identical whether
formulated by an embedded participant or by an external observer. Two principles are one too many,
and this subsection reduces them.

The reduction corrects the framework's own origin, and we make the correction in print. The 2005
form of the argument \autocite{small_2006} ran a G\"{o}delian chain: a self-referential process,
being complete, must by the incompleteness theorem \autocite{Godel1931} be inconsistent; an
inconsistent system can negate itself; and the negation of self-reference is the production of
causal sequence. The chain captures the right intuition and is retained here as motivation, but it
is not sound as a proof: classical inconsistency yields \emph{ex falso quodlibet} --- every
proposition provable, every causal sequence an output, no constraint at all. The load-bearing
argument is different, and better: a fixed-point theorem.

Let $\mathcal{C}$ be the category of physical descriptions and admissible processes --- concretely,
the dagger-compact-closed setting with classical structures in which categorical quantum mechanics,
and the whole of this paper, works \autocite{CoeckeKissinger2017}. Lawvere showed that in a
sufficiently rich such setting, any weakly point-surjective self-representation forces every
endomap to have a fixed point, and Yanofsky showed that one diagonal argument of this shape
underlies Cantor's, G\"{o}del's, Turing's and Tarski's theorems alike
\autocite{Lawvere_1969,Yanofsky2003,cantor1891diagonal}. The framework's single commitment is that
$\mathcal{C}$ contains a \emph{self-representing object} --- an object capable of encoding the
processes on itself. That is what ``self-reference is foundational'' means, stated categorically,
and it is the surjectivity hypothesis of the diagonal lemma. Now let
$\sigma \colon \mathcal{C} \to \mathcal{C}$ be the \emph{inside/outside involution}: the operation
exchanging the two complementary presentations of any description --- the inside presentation, as
participant-indexed total space in which the observer is part of the system, and the outside
presentation, as a base-times-fibre decomposition over a classical parameter space. That
$\sigma^{2} = \mathrm{id}$ is not an axiom but a coherence requirement: the outside view of the
inside view of a thing is the thing. The candidate concrete realisation of $\sigma$ is the
Hopf-duality swap --- total space against base-times-fibre for the three fibrations --- which gives
the involution computable content and ties this subsection to \S\ref{sec:hopf-justification}.

\begin{proposition}[Inside/outside fixed point]\label{prop:inside-outside}
Let $\mathcal{C}$ satisfy the richness condition above and contain a self-representing object, and
let $\sigma$ be an involution on $\mathcal{C}$. Then the fixed-point locus
$\mathrm{Fix}(\sigma)$ is non-empty: there exist descriptions invariant under the exchange of
inside and outside perspective.
\end{proposition}

\begin{proof}
Lawvere's construction, applied to $\sigma$ in place of negation: given a weakly point-surjective
$\varphi \colon A \to Y^{A}$, form $g(a) = \sigma\!\left(\varphi(a)(a)\right)$; point-surjectivity
supplies $a^{*}$ with $\varphi(a^{*}) = g$, and then $\varphi(a^{*})(a^{*})$ is a fixed point of
$\sigma$ \autocite{Lawvere_1969,Yanofsky2003}. Where the same construction applied to Boolean
negation yields paradox --- negation has no fixed point on a two-element type, so the conclusion
explodes the premises --- the inside/outside involution is an honest involution on a rich type, and
the conclusion is a constructive fixed point rather than a contradiction.
\end{proof}

\textbf{Definition.} \emph{Physical law is the structure carried by} $\mathrm{Fix}(\sigma)$.

The Inside/Outside Equivalence Principle is thereby demoted from axiom to theorem: physical law is
invariant under the exchange of perspectives \emph{because it is the fixed-point set of that
exchange}, and the fixed-point set is forced to be non-empty by self-reference alone. One
commitment remains foundational; everything the second commitment asserted is now derived. The
reader who accepts categorical quantum mechanics as an existing programme has already accepted the
machinery; the framework adds one premise and collects the rest. And the fixed-point locus is not
empty of structure: it inherits the Frobenius-algebra apparatus of $\mathcal{C}$, so that the
Bell, GHZ and W primitives --- and with them the particle classification of the later sections ---
live inside $\mathrm{Fix}(\sigma)$ from the start. What is proved above is non-emptiness given the
setting; what is assumed is the setting itself and the existence of $\sigma$ as a genuine
involution; what is open is the rigour of both, and it is priced now rather than hidden:

\begin{openproblem}[The involution made functorial]\label{op:sigma-functor}
Fix the category $\mathcal{C}$ rigorously --- the dagger-compact-closed setting with classical
structures is the candidate --- and define $\sigma$ on morphisms, not only on objects, the natural
route being the induced fibre--base exchange on spider diagrams; establish that $\sigma$ is an (at
worst pseudo-)involutive endofunctor; and prove or refute the identification of
$\mathrm{Fix}(\sigma)$ with the three Hopf cross-sections, under which Adams's theorem would become
a corollary of the fixed-point classification. This problem is cross-linked with Open
Problem~\ref{op:hopf-frobenius}: under an affirmative resolution of both, the parallelisable
terminus, the fixed-point classification, and compositional completeness are three statements of
one fact, and progress on either problem constrains the other. One constraint is already in place:
at the complex rung, Proposition~\ref{prop:f1-dictionary} supplies the dictionary between
classical structures and the fibration's lifted skeletons --- the clause of
Open Problem~\ref{op:hopf-frobenius} it discharges --- and any resolution here must reproduce it.
The problem also carries a piece of cargo recorded later in this section: the formalisation, within
the fixed category, of the perspective-accounting that derives local tomography.
\end{openproblem}

\subsection{The lineage}

The obstruction this section builds on is old, and the framework claims its completion, not its
discovery. Popper argued in 1950 that a predictor embedded in the world has irreducibly inferior
knowledge of itself, because it cannot impart information about itself without, by that very act,
rendering the information out of date \autocite{Popper1950a,Popper1950b} --- the passage is quoted
where the framework makes the obstruction structural (\S\ref{sec:wigners-friend}). Breuer proved
the impossibility of accurate state self-measurement within quantum theory \autocite{Breuer1995};
Svozil's monograph surveys the whole territory where physics meets self-reference, undecidability
and Chaitin-style algorithmic limits \autocite{svozil_k_2018}; Szangolies sharpened the modern form
as epistemic horizons \autocite{szangolies_j_2018}; and the Frauchiger--Renner theorem is the
contemporary crisis-form --- nested observers deriving contradictory records when each treats the
others' quantum descriptions as universally valid \autocite{Frauchiger_2018} --- which
\S\ref{sec:wigners-friend} dissolves relationally rather than paradoxically. This framework's
addition to the lineage is twofold. Structurally, Proposition~\ref{prop:f1-dictionary} converts
Popper's temporal obsolescence into algebraic cancellation: the datum a copier cannot put into its
own record is exactly its fibre commitment, which does not arrive stale --- it cancels. And
foundationally, the present section runs the lineage in reverse: where the tradition treats
self-reference as a limit on knowledge \emph{within} physics, the framework takes it as the
generator of physics, and derives the quantum formalism from it.

\subsection{The price is negative probability}\label{sec:complex-amplitudes-2}

What does it cost a probabilistic description to contain its own describer? The first instalment is
the sign. In the regime before any record exists, every event has an anti-event --- the description
must be able to undo what it has provisionally done, or the self-description could never be
revised into consistency --- and a probability calculus closed under undoing cannot be confined to
$[0,1]$: it is a signed calculus, in which the parts of a probability can oppose and cancel. The
framework did not discover negativity as bookkeeping; Dirac admitted negative probabilities as
formal intermediaries \autocite{dirac1942}, Feynman defended them as consistent whenever the final,
recorded probabilities stay positive \autocite{Feynman_r_82,Feynman1987}, and the Kirkwood--Dirac
distribution has carried complex quasiprobability from 1933 to the present, as a recent
comprehensive review documents \autocite{arvidssonshukur2024kirkwood}. What the framework adds is
the forcing: negativity is not a permissible convenience but the first price of self-referential
closure, and it is charged exactly and only in the pre-record regime --- which is why every
negativity in the tradition lives in intermediate, unrecorded quantities and never in a final
record.

The second instalment is complexity. A signed information measure must extend Shannon's
$S = -\ln p$ off the positive axis, and the analytic continuation of the logarithm carries $p$ into
the complex plane: amplitudes, with the phase as the continued information. Here too the extension
has a prior tradition, and the attributions are owed precisely. Tikochinsky derived Feynman's rules
for complex amplitudes from consistency of composition \autocite{tikochinsky1988feynman}; Youssef
reformulated quantum mechanics as complex probability theory outright \autocite{Youssef1994};
Caticha carried objective amplitudes to Bayesian probabilities \autocite{caticha_a_2006}; and
Goyal, Knuth and Skilling derived Feynman's rules from a pair-postulate without assuming
analyticity \autocite{goyal_knuth_skilling_2010}. The extension is prior art. The framework's
claim is the compulsion: the complex extension is not a modelling choice among exotic probability
theories but the unique continuation forced by self-referential closure --- the same forcing that
Aaronson glimpsed when he argued quantum mechanics is what negative probability inevitably becomes
\autocite{Aaronson2004h}.

Why does the continuation stop at $\mathbb{C}$? Downward, the real field is not available: real
quantum theory is experimentally distinguishable from complex quantum theory in network scenarios
and has been disfavoured by experiment
\autocite{renou2021real,chen2022ruling,wu2022refutation}; the recent objection that the exclusion
rests on locally tomographic composition \autocite{hoffreumonwoods2026real} is, for this framework,
not an escape but a confirmation of where the weight sits --- local tomography is \emph{derived}
below, not assumed, so the framework closes from its own resources exactly the door the objection
holds open. Upward, the quaternions and octonions fail as a \emph{base} for the information
continuation: the continued measure is a logarithm, and beyond $\mathbb{C}$ the logarithm of a
product is no longer the sum of logarithms --- non-commutativity severs
$\log(ab) = \log a + \log b$ --- so the additivity that makes it an information measure is lost.
Aaronson's parameter-counting and signalling arguments against a quaternionic base point the same
way \autocite{Aaronson2004h}. The higher algebras are not banished; they return as the
\emph{internal} algebras of the register, which is the business of \S\ref{sec:division-algebras}.
The base is $\mathbb{C}$, for informational reasons, and stays $\mathbb{C}$.

\subsection{The Born rule from fibre privacy}\label{sec:born-rule}

The third instalment is the quadratic rule. An amplitude carries one datum more than any record of
it can: its overall phase. This was asserted throughout the framework's development as the
observer's private, self-referential ignorance; it is now a proved statement, in exactly the form
needed. By Proposition~\ref{prop:f1-dictionary}, a classical structure --- the copy operation any
record requires --- is committed to a specific lift of its skeleton, and that commitment is absent
from every classical record the structure induces: the copier is fibre-committed, the record is
fibre-blind, and the fibre freedom is gauge for every observer. (The computation extends verbatim
from one qubit to the full register, where Proposition~\ref{prop:trace-blind} makes it the ground
of hypercharge's tracelessness --- one privacy, two duties.) Popper's obsolescence, made
cancellation. It follows that anything recordable is invariant under the fibre --- observables and
their statistics must be $U(1)$-invariant functions of the amplitude. The lowest-degree polynomial invariant
of a complex amplitude is the quadratic $|\psi|^{2}$ --- the non-polynomial invariant $|\psi|$,
of degree one, is excluded by the analyticity declared at \S\ref{sec:complex-amplitudes} --- and
normalisation fixes the coefficient:
probability is the minimal record-compatible function of the amplitude,
\[
p \;=\; |\psi|^{2},
\]
the Born rule as the cheapest invariant the fibre permits. The same square is reached from the
cost side in \S\ref{sec:born-cost-ledger}, where a record is the closed loop of a history and its
return leg and the exponent is the number of legs; neither route uses the lattice of projections
that Gleason's theorem needs, and both hold for a single qubit. The derivation is the framework's
own;
its premise --- that the fibre datum is structurally excluded from records --- is
Proposition~\ref{prop:f1-dictionary}, and its geometric reading --- probability as base-point data
under the Hopf projection --- is developed in \S\ref{sec:hopf-justification}.

One standing objection to information-theoretic foundations is answered here, because this
framework needs the answer. Brukner and Zeilinger argued that Shannon information is conceptually
inadequate for quantum measurements: it is basis-dependent, and its operational meaning presupposes
ignorance of a definite but unknown fact, which the quantum case does not supply
\autocite{bruknerzeilinger2001shannon}. The framework's reply is structural, and it is a
demarcation rather than a rebuttal. Shannon's precondition fails exactly where they say it fails
--- in the pre-record regime, where there is no definite unknown to be ignorant of --- and there
the operative measure is the quadratic invariant just derived. Shannon's axioms hold exactly where
the precondition holds --- on records, after \textsc{fanout}, where outcomes are definite and
ignorance is classical --- and there the uniqueness theorem forces $-\log p$
\autocite{shannon1948mathematical}. The basis-dependence dissolves with the same stroke: a basis
\emph{is} selected precisely at \textsc{fanout}, so the basis-dependent measure applies only where
a basis exists. Shannon is not inadequate; Shannon is post-\textsc{fanout}. The framework thereby
keeps both measures, each on its own side of the boundary, and uses the additivity of the
post-record measure --- the very property the previous subsection used to pin the base field ---
with a clear conscience.

\subsection{Thomson's lamp, resolved}

The section owes \S\ref{sec:continuity} a debt incurred there: the lamp's amplitude-side limit.
Thomson's lamp --- toggled at times $1 - 2^{-n}$, its state at $t = 1$ demanded --- has no
classical answer: the task description fixes every finite stage and thereby fixes no limit at all
\autocite{Thomson_1954,benacerraf1962tasks}, which \S\ref{sec:continuity} read as the proof that
the record algebra is not limit-closed. On the amplitude side the limit exists, and the calculus of
this section computes it. The toggle exchanges the two amplitudes; the self-consistent state of the
completed task is the fixed point of the exchange in the complex calculus, and solving it under the
signed continuation gives
\[
P(\mathrm{on}) = \tfrac{1 + i}{2}, \qquad P(\mathrm{off}) = \tfrac{1 - i}{2},
\]
a superposed fixed point whose imaginary part carries the unrecordable datum --- the phase of the
endless alternation, when the lamp \emph{would have} stopped. Two checks certify that the
assignment is calculus rather than decoration. The pair is normalised,
$P(\mathrm{on}) + P(\mathrm{off}) = 1$; and the Born rule of the previous subsection, applied to
the fixed point, returns $|P(\mathrm{on})|^{2} = \tfrac{1}{2}$ --- exactly the symmetric classical
ignorance that is the only sane record-level answer, recovered rather than imposed. The lamp,
which entered the framework at its very beginning as the place where classical computation meets
its boundary, thus exits it as a worked example: undetermined as a record, determinate as an
amplitude, the two sides split precisely along \textsc{fanout} --- and the diagonal construction it
physically instantiates \autocite{cantor1891diagonal,Turing_1937} is the same one that powers
Proposition~\ref{prop:inside-outside}. The status flag is owed here as everywhere: the
undetermination is theorem-level; the specific fixed-point value is the framework's complex
calculus, structural, offered with its two consistency checks shown.

\subsection{Local tomography, derived}\label{sec:local-tomography}

The fourth instalment concerns composites, and it discharges two promises at once. The
Inside/Outside Equivalence Principle --- now a theorem --- requires that an external observer's
characterisation of a composite system be recoverable by the embedded participants, and the
embedded participants have access only to local measurements and their correlations. Equivalence
of the two perspectives therefore entails that the state of a composite is determined by the
statistics of local measurements: \emph{local tomography}, obtained as a consequence rather than
imposed as an axiom. This converts a postulate of the reconstruction programmes
\autocite{Hardy2001a,mullerp2011} into a theorem-shaped consequence of the one principle ---
discharging the caveat recorded in \S\ref{sec:continuity} --- and it closes, from the framework's
own resources, the escape route in the real-number debate noted above: the assumption the
objection of \autocite{hoffreumonwoods2026real} contests is here not an assumption. The status is
flagged plainly: the derivation is structural, one paragraph of perspective-accounting, and its
formalisation within the categorical setting belongs to Open Problem~\ref{op:sigma-functor}.

\subsection{Contrast: QBism, and the ledger}

The nearest neighbour in spirit deserves a paragraph of honest contrast. QBism
\autocite{fuchsschack2013qbism} holds, as this framework does, that quantum probabilities are an
agent's --- that the theory is a calculus of expectation, not a catalogue of properties. What QBism
does not say is why the agent's expectations must be complex-valued and quadratically squared; the
Born rule enters as a normative addition to Bayesian coherence, and the complex amplitude as an
unexplained representation. This section is, in effect, the missing why: the agent's ignorance is
self-referential --- the agent is inside the description --- and self-referential ignorance is
complex, with the quadratic rule as the cheapest record-compatible invariant of it. Where QBism
stops at the agent, the framework derives the agent's calculus.

The ledger, then. Delivered in this section: the reduction of two principles to one
(Proposition~\ref{prop:inside-outside}); the pre-record regime with its anti-events, handed to
\S\ref{sec:continuity} as the five-face regime whose continuity is priced there
(Open Problem~\ref{op:limit-realisability}); the complex base, forced downward by experiment and
upward by the additivity of information; the Born rule from fibre privacy; the lamp's fixed point;
and local tomography. Owed, and stated where they are owed: the functorial rigour of the involution
and the fixed-point classification (Open Problem~\ref{op:sigma-functor}); the remaining premise of continuity (Open Problem~\ref{op:limit-realisability}, in
\S\ref{sec:continuity}); and the
rung-by-rung compulsion of the algebra ladder, posed as Open
Problem~\ref{op:ladder-compulsion} in \S\ref{sec:division-algebras}. The
hand-on is the paper's spine: the $\mathbb{C}$ this section derives natively and the $\mathbb{C}$
that \S\ref{sec:continuity} underwrites through Hardy's theorem are the same field, the floor of
the ladder that \S\ref{sec:division-algebras} climbs and \S\ref{sec:hopf-justification}
geometrises. One principle; everything else is theorem, derivation, or a debt with a stated
invoice.

%      % EDITED 2026-10-29 scripts/patch_continuity_tower_premise_2026-10-29T1500.py
% EDITED 2026-09-15 scripts/patch_value_continuum_declared_cencov_2026-09-15T1540.py
% EDITED 2026-09-15 scripts/patch_continuity_connected_state_space_2026-09-15T1430.py
% =====================================================================
%  SECTION: Continuity Is Not an Axiom
%  Drafted 2026-07-04 from continuity_limit_realisability_2026-07-04.md
%  (now in project documents) and the May record
%  (dynamics_companion_fanout_unitarity_2026-05-20T1430.md).
%
%  PLACEMENT in main.tex: a standalone section between self-reference
%  and division-algebras --
%      \input{self_reference_section_2026-07-04}
%      \input{continuity_section_2026-07-04}
%      \input{division_algebras_section_2026-07-03}
%  (Architecture call SETTLED by the sync manifest: standalone.)
%
%  CITATIONS -- existing keys (as supplied by JS 2026-07-04):
%    Hardy2001a, Aaronson2004h, Thomson_1954, Lawvere_1969,
%    Turing_1937, muller_p_2011, Brukner2009.
%  NEW keys (RIS in continuity_refs_2026-07-04.ris):
%    cantor1891diagonal, benacerraf1962tasks, broeckertomdieck1985compact.
%  Deliberately excluded per JS: Chiribella--D'Ariano--Perinotti;
%  Bargmann/Wigner's theorem. Yanofsky2003 now paired on the
%  Lawvere-shaped clause (the pending-key excuse was void: the key
%  exists and sec:self-reference uses it).
%
%  CROSS-REFS: sec:self-reference now DRAFTED (2026-07-04, later the
%  same day) -- no longer dangling; only sec:fanout-boundary (and its
%  sibling sec:fanout-direction) -- DRAFTED 2026-07-06.
%
%  AUDITED 2026-07-06 (final pass, second stop): both proofs
%  re-verified by hand (Lagrange divisibility sentence exact; the
%  Baire countable-escape chain sound); caveat-discharge with
%  sec:self-reference confirmed two-way (already present); file pure
%  ASCII; no pre-definition "frame"; all citations used as declared.
%  Fixes: Yanofsky2003 paired on the Lawvere clause; the op's two
%  candidate routes anchored to printed propositions (f1-dictionary
%  for the divisibility premise's referent, inside-outside for the
%  closedness demand) -- statuses unchanged, referents sharpened;
%  header staleness corrected. Deliberate FORWARD references --
%  prop:f1-dictionary and sec:hopf-justification (defined later in the
%  paper; fine under two-pass LaTeX and flagged in prose),
%  sec:division-algebras, sec:space-of-computations.
%
%  DECISIONS ENCODED: op:limit-realisability minted (the owed arrow,
%  priced: exactly two derivations, divisibility and closedness);
%  Lagrange used twice (finite exclusion; component-group triviality);
%  the (1+i)/2 lamp value NOT stated here -- "a superposed fixed point"
%  with the complex-probability treatment deferred to
%  sec:self-reference; the two-branch over-determination stated to
%  pre-empt a circularity reading; Thomson spelling throughout.
% =====================================================================

\section{Continuity Is Not an Axiom}
\label{sec:continuity}

Hardy derives quantum theory from five axioms \autocite{Hardy2001a}. Four of them are satisfied
by classical probability theory; the fifth --- that a \emph{continuous} reversible
transformation connects any two pure states --- is the entire difference, the single postulate
that turns the classical simplex into the complex Hilbert space. Every reconstruction programme
pays an equivalent toll: continuous reversibility in Masanes and M\"{u}ller
\autocite{mullerp2011}, the corresponding postulate in Daki\'{c} and Brukner
\autocite{Brukner2009}. The axiom is load-bearing, and in every programme it is unexplained: it
is where the quantumness is put in by hand. The claim of this section is that in the present
framework it is not an axiom. Continuity is the fifth face of the pre-measurement regime that self-reference forces, and what
stands between the framework and a full derivation is not a postulate about values but one
premise about transformations --- that every transformation the inside performs has a stationary
outside presentation --- derived in outline below from the inside/outside involution of
\S\ref{sec:self-reference} and stated as an open problem at the end of the section with nothing
hidden in it. The continuum itself, the completion that every probability theory adopts, is
named for what it is: the real line, shared with classical probability and carrying no quantum
content.

One structural remark prevents a misreading. The framework reaches complex amplitudes natively in
\S\ref{sec:self-reference}, by the negative-probability route; the present section reaches the same
destination by an independent branch --- the regime's continuity, then Hardy's theorem. The two
branches converging on one structure is over-determination, not circularity, and the convergence is
itself evidence that the destination is forced.

\subsection{Five faces of one regime}

The regime before any measurement --- before any \textsc{fanout} --- can be characterised five
equivalent ways: negative probability is available; anti-events exist, every event being undoable;
no classical record exists; no preferred basis is selected; and a continuous reversible
transformation connects any two pure states. These are not five premises to be combined. They are
one regime under five descriptions, and \textsc{fanout} (\S\ref{sec:fanout-boundary}) is the single
event that fails all five at once: it writes the record, selects the basis, kills the anti-events,
discretises the transformation group, and switches off the negativity. The equivalence of the first
two faces is the anti-event identity of \S\ref{sec:self-reference}; the third is the boundary's
definition. The work of this section is the equivalence of the last two faces, and then the route
from the regime's existence to its continuity. The fourth face has, since
Proposition~\ref{prop:f1-dictionary} (\S\ref{sec:hopf-justification} --- a deliberate forward
reference), an exact form: a basis-selection is the commitment of a classical structure, a lifted
antipodal skeleton, so ``no preferred basis'' means precisely that no copy structure has fired.

\subsection{Basis selection is the collapse of continuity}

\begin{proposition}[Continuity forbids a selected basis]\label{prop:basis-discreteness}
If a finite set of pure states is privileged --- a preferred basis --- then every reversible
transformation respecting the privilege carries the set to itself and acts on it as a permutation;
the respecting transformations therefore realise only a finite group of motions among the
distinguished states, and no continuous family of them connects distinct pure states of the set.
Contrapositively, continuous reversibility entails that no basis is selected.
\end{proposition}

\begin{proof}
A privileged finite set is by definition carried to itself by every structure-preserving reversible
map, whose induced action on the set is a permutation of finitely many points. A continuous path of
such maps induces a continuous path in a finite discrete set, which is constant. Hence a continuous
reversible family moving one distinguished pure state to another respects no such privilege, and if
every pair of pure states is so connected, no finite set is privileged.
\end{proof}

The converse is where the historical difficulty lies: the absence of a privileged finite set does
not by itself force continuity, for the transformation group could in principle be infinite and
discrete. Two structural properties close every such escape at once.

\begin{proposition}[Divisible and closed implies connected]\label{prop:divisible-closed}
Let $G$ be the group of reversible transformations of the pre-measurement regime, acting on the
state body. If $G$ is divisible --- every element has an $n$-th root for every $n$ --- and
topologically closed, then $G$ is a compact connected Lie group; and if $G$ moreover acts
transitively on the pure states, then a continuous reversible transformation connects any two of
them: Hardy's fifth axiom.
\end{proposition}

\begin{proof}
The state body is compact and convex with nonempty interior, so its symmetry group is compact, and
a closed subgroup of a compact matrix group is a compact Lie group with finitely many connected
components \autocite{broeckertomdieck1985compact}. The component group $G/G_{0}$ is a quotient of a
divisible group and hence divisible; but a nontrivial finite group is never divisible --- with
$\lvert H\rvert = m$ every element satisfies $h^{m} = e$, so no element other than the identity
possesses an $m$-th root --- and $G/G_{0}$ is finite, so it is trivial and $G = G_{0}$ is connected.
A connected Lie group is path-connected, and transitivity then supplies the continuous reversible
path between any two pure states.
\end{proof}

The same one-line fact of Lagrange closes the remaining escapes. Finite groups are excluded
directly, being nontrivial and never divisible. Countable groups are excluded because a
countable closed subgroup of a compact group is compact, has an isolated point by the Baire
category theorem, and is therefore discrete and finite. Uncountable disconnected groups are
excluded by the component-group argument of the proposition. Every route to a discrete
transformation group therefore passes through a failure of divisibility or a failure of
closedness. Divisibility, however, asks for more than the conclusion needs, and the next
proposition removes it.

\begin{proposition}[Continuity from a connected state space]\label{prop:connected-transitive}
Let the pure states of a finite-dimensional system form a connected space $P$, and let the group
$G$ of reversible transformations of the pre-measurement regime be topologically closed and act
transitively on $P$. Then the identity component $G_0$ already acts transitively on $P$, and a
continuous reversible transformation connects any two pure states: Hardy's fifth axiom holds,
with no divisibility assumed.
\end{proposition}

\begin{proof}
As in Proposition~\ref{prop:divisible-closed}, closedness makes $G$ a compact Lie group with
finitely many connected components \autocite{broeckertomdieck1985compact}, so $G_0$ is a closed
normal subgroup of finite index and is path-connected. Each $G_0$-orbit in $P$ is the continuous
image of a compact group, hence compact, hence closed. Since $G$ is transitive and $G_0$ is
normal, the $G_0$-orbit of $gx$ is $g$ applied to the $G_0$-orbit of $x$, so the cosets of $G_0$
permute the orbits transitively and there are finitely many of them. A finite partition of $P$
into closed sets is a partition into open sets; $P$ is connected; so there is one orbit and
$G_0$ is transitive. Given pure states $x$ and $y$, choose $g \in G_0$ with $gx = y$ and a
continuous path $g_t$ in $G_0$ from the identity to $g$; then $g_t x$ is a continuous reversible
transformation from $x$ to $y$.
\end{proof}

The proposition relocates the premise. Divisibility was asked for because a transformation group
could be infinite and discrete; but a group acting transitively on a continuum cannot be
countable, and a closed group acting transitively on a connected space cannot be disconnected in
the way that matters, so what the argument needs from the regime is not the availability of
$n$-th roots but the connectedness of the state space itself. That is a statement about values:
that the pure states are all the unit vectors of $\mathbb{C}^{K}$ up to phase, the space
$\mathbb{CP}^{K-1}$, connected for $K \ge 2$, and hence that the amplitudes of
\S\ref{sec:self-reference} fill their continuum. The two standing foils show the premise doing
its work. The stabiliser fragment of quantum theory has complex amplitudes and is contextual,
yet it has finitely many pure states, its outcome probabilities take only the values $0$,
$\tfrac12$ and $1$, and its reversible transformations form a finite group, the Clifford group,
so Hardy's fifth axiom fails in it; Spekkens's toy theory \autocite{spekkens2007toy} is the same
with real amplitudes, its ignorance rationed by a knowledge-balance principle. Neither has a
connected state space, and neither is a model of the regime described here, whose credences take
every value. Divisibility, which both also lack --- the phase gate has no square root in the Clifford group
--- is a symptom of the discreteness, not its cause. What both lack can be said exactly, and the
next subsection says it: each contains a transformation that fixes no state it admits.

\subsection{The premises, and what they mean}

Three premises carry Proposition~\ref{prop:connected-transitive}. Transitivity is the regime's
homogeneity: inside/outside equivalence leaves no pure state privileged, and a proper closed
invariant set of pure states would define a superselection observable, a classical record
predating any \textsc{fanout} (\S\ref{sec:fanout-boundary}). Connectedness of the state space is two statements, and only one of them is quantum. The first
is that the pure states are infinite and dense, with a fixed point for every transformation; the
second is that the state space is complete. Quantum theory over the algebraic numbers ---
amplitudes in the algebraic closure of $\mathbb{Q}$ within $\mathbb{C}$ --- has the first
without the second: it has $K = N^{2}$, the Born rule, superposition, entanglement and Bell
violation at the Tsirelson value, every unitary has an algebraic eigenvector, and it is
countable and totally disconnected. So the whole quantum content of Hardy's axiom lies in the
density, and the completion is the real-number convention of every probability theory. Hardy
adopts it before his fifth axiom is stated, in the first, which makes probabilities real numbers
and gives his classical theory a real simplex of mixtures; this framework adopts it at the
Shannon step, where credences are declared to range over a connected domain so that the
continuation of the logarithm is unique (\S\ref{sec:complex-amplitudes}). Given that convention,
density of the pure states is connectedness of their closure. Divisibility, which an earlier
version of this section asked for, is not needed: it follows from the conclusion, every element
of a compact connected Lie group having roots of every order.

Closedness is the completion, limit-realisability: a Cauchy sequence of available
transformations has its limit available. It cannot be derived, and the reason should be stated
rather than hidden. Any collection of first-order principles about states, transformations and
probabilities that has an infinite model has a countable one, and a countable subset of
$\mathbb{CP}^{K-1}$ is totally disconnected; the founding non-computability does not escape
this, since the arithmetical complex numbers form a countable algebraically closed field
containing non-computable numbers and closed under every limit it can itself describe. Hardy's
axiom and closedness are second-order, quantifying over paths and over sequences, and a
second-order principle is grounded as Dedekind completeness is grounded or not at all. The
grounds are these. The completion of a dense set under the finite-precision metric is unique and
depends on no choice, whereas any countable alternative is a choice of definability level, an
extra law with no ground and no test; and the framework's own result is that no observer's
descriptive level captures the regime, since by \S\ref{sec:self-reference} every observer's
state exceeds its own level and the inside/outside principle makes that a law rather than a
perspective. The burden therefore falls on any restriction, not on the completion, which is the
least committal option and carries no quantum content. What limit-realisability adds to the line
is still worth stating exactly, because it is the corrected form of a natural intuition.
Cantor's diagonal \autocite{cantor1891diagonal} does not build the continuum; its output is an
infinite decimal --- a Cauchy sequence of rationals --- which exists \emph{as a point} only
because the line is complete. The diagonal presupposes limit-realisability and then shows that
the completed line outruns every internal listing. And the character of the completion is the
framework's own: the computable reals are countable \autocite{Turing_1937}, hence of measure
zero, so once the transformation group contains its limits almost all of them are
non-computable. Non-computability is a property of the completed continuum, not the reason for
it. The constructive reading, on which only describable limits exist, remains consistent; it is
the countable theory above, and the two readings differ in nothing an experiment can certify.

Density is where the section's title is earned, and it is derived from the paper's own premise.
Write $(\mathrm{I})$ for the statement that every admissible transformation of the regime is a
phase of some admissible classical structure: diagonal, that is, in the record basis of some
observer the regime contains, a phased spider of that basis in the calculus of
\S\ref{sec:hopf-justification} \autocite{coeckeduncan2011interacting} (VERIFY-SUPPORT). It says
of transformations what the global phase of \S\ref{sec:self-reference} says of states: there are
no absolute ones. It follows from three things the paper already holds. The inside/outside
involution is a genuine involution, so the outside presentation of the inside is a presentation
of the same object with nothing lost: the outside can have complete knowledge of the inside, a
potential record, the joint state lying in a structure the outside carries. Time is operation
count, internal to the system that operates (\S\ref{sec:space-of-computations}), so the outside
does not share the inside's count. And before any \textsc{fanout} that count is not classically
recorded, so the outside's description is pure. Let $g$ be the transformation the inside
performs per operation. The outside's complete knowledge cannot depend on the inside's count,
and it is pure, so it is a pure state of the inside's clock and system that is stationary under
``advance the clock and apply $g$''; the stationary states of that map are clock Fourier modes
tensored with eigenvectors of $g$ and nothing else \autocite{pagewootters1983evolution}
(VERIFY-SUPPORT). So the outside holds a structure in which $g$ acts by phases, which is
$(\mathrm{I})$; and the relative phase in that structure is the outside's ignorance of the
inside's computational count, the same ignorance the global phase records one level down. The
consequences are a theorem. A phase of a structure fixes the structure's copyable states, so
every admissible transformation fixes an admissible pure state; a finite transitive action has
an element moving every point (Jordan), so the pure states are infinite; the transformation
carrying a stabiliser state to the fixed point of the composite of a $Z$-phase and an $X$-phase
turns through $\arccos(1/\sqrt{3})$, an irrational fraction of a turn, so the group has an
element of infinite projective order and its orbits are dense. With the completion,
Proposition~\ref{prop:connected-transitive} fires and Hardy's fifth axiom holds. The two foils
fail $(\mathrm{I})$ exactly where they fail continuity: fourteen of the twenty-four Clifford
operations are phases of no Pauli basis and fix no stabiliser state, their fixed points being
the magic states; the toy theory's three-cycles are phases of no epistemic basis; and in each
the outside's knowledge of the inside under the order-three protocol is either count-dependent
or a mixture over the count, so each makes the inside's time classical before any record exists.
The derivation of $(\mathrm{I})$ is an outline, not a proof, and it is what the open problem
below asks for.

\subsection{Why the cut matters physically}

The regime's continuity is not decorative. Reversibility alone, over complex amplitudes, licenses
the full invertible group $GL(n, \mathbb{C})$, and Aaronson showed what that buys: superluminal
signalling, and computational power at the level of \textsf{PP} \autocite{Aaronson2004h}. The cut
from $GL(n, \mathbb{C})$ down to the unitary group is the conserved quadratic form, and continuity
is what entitles the theory to make it: with the fifth axiom in hand, Hardy's theorem delivers the
complex Hilbert space, the quadratic norm and the unitary group in one stroke \autocite{Hardy2001a}
--- cited here, not re-proved. One honest caveat travels with the citation. Hardy's reconstruction
uses his remaining axioms besides --- composition and subspace structure --- and these must be
derived rather than imported: the composition ledger of \S\ref{sec:division-algebras} and the
locality-of-tomography argument of \S\ref{sec:self-reference} carry that duty, and it is discharged
there, not silently assumed here.

\subsection{Super-tasks at the boundary}

The premise just priced has a fifty-year-old stress test in the philosophy of mathematics, and the
framework passes it in an unexpected way. A super-task --- infinitely many operations completed in
finite time \autocite{Thomson_1954} --- presupposes its terminal moment as an actual point of the
timeline: the accumulation point of a generic sequence of step-times exists only in the completed
line, and the defining question, the state \emph{after all the steps}, is meaningful only under
limit-realisability of time. Super-tasks are downstream consumers of exactly the completeness at
issue in this section. Thomson's lamp then exhibits both sides of the \textsc{fanout} boundary at
once. On the record side, the lamp's description fixes its state at every finite stage and thereby
fixes no limit at all: any terminal value is consistent with the whole task
\autocite{Thomson_1954,benacerraf1962tasks}. That is not a puzzle about lamps; it is the proof that
the classical record algebra is not limit-closed, and cannot be. On the amplitude side the limit
exists --- a superposed fixed point, whose complex-probability treatment belongs to
\S\ref{sec:self-reference} --- because that side \emph{is} limit-closed, which is the fifth face
itself. The boundary is thereby the edge of limit-realisability: complete on the reversible side,
incomplete on the record side, with the lamp standing on both sides at once as the boundary's oldest
exhibit. A final clause dissolves rather than answers the classic question. In this framework time
is operation count (\S\ref{sec:space-of-computations}), and ``after infinitely many operations''
indexes no operational moment: the lamp packs its endless \textsc{fanout}s into the emergent
continuum, and asking for its record at the accumulation point asks the discrete contents for a
value that only the completed container possesses. On or off is not undecided; it is unasked.

\subsection{What remains, priced}

\begin{openproblem}[Every transformation a phase of an inside structure]\label{op:limit-realisability}
Proposition~\ref{prop:connected-transitive} reduces Hardy's fifth axiom to three premises:
connectedness of the pure states, transitivity, and closedness. Transitivity is argued above
from the absence of records. Closedness is the completion, adopted on the grounds stated above
and carrying no quantum content. Connectedness is density plus that completion, and density
follows, by Jordan's theorem and the infinite-order argument, from $(\mathrm{I})$: that every
admissible transformation of the regime is a phase of some admissible classical structure. What
remains open is the rigour of the derivation of $(\mathrm{I})$ from the paper's premise. Derive,
inside the fixed-point locus $\mathrm{Fix}(\sigma)$ of Proposition~\ref{prop:inside-outside},
that every endomap is a phase of a classical structure in $\mathrm{Fix}(\sigma)$, from the
self-representing object, the involution $\sigma$ as a genuine involution, and the reading of
the inside presentation as a sequence indexed by operation count with no record before
\textsc{fanout}. The outline above does this with the stationary-state construction of Page and
Wootters; a derivation passes if it makes that outline a theorem of the calculus, with the
outside's count-independent complete description shown to exist and to be pure from the
involution alone. It fails before any detail is checked if every premise it uses holds in the
stabiliser fragment of quantum theory, in Spekkens's toy theory \autocite{spekkens2007toy}, in
the real qubit, or in quantum theory over the Gaussian rationals, all of which violate
$(\mathrm{I})$ outright and all of which make the inside's operation count classical before any
record exists. It also fails if it imports the completion as a premise about transformations:
the completion is the real line and has no derivation, and a derivation of $(\mathrm{I})$ must
hold in the countable algebraic theory, where $(\mathrm{I})$ is true and continuity is not.
\end{openproblem}

The accounting, in one sentence: the axiom that every reconstruction programme pays is, in this
framework, one derivation posed, of Lawvere's shape --- that every endomap of the fixed-point
locus is a phase of a classical structure in it --- and one convention named, the completion,
which classical probability pays too. And
the hand-on is already visible: the complex amplitudes that this section's continuity underwrites
through Hardy's theorem, and the complex amplitudes that \S\ref{sec:self-reference} derives natively,
are one and the same $\mathbb{C}$ --- the floor of the ladder that \S\ref{sec:division-algebras}
climbs.

%      % EDITED 2026-10-29 scripts/patch_continuity_tower_premise_2026-10-29T1500.py
% EDITED 2026-09-15 scripts/patch_value_continuum_declared_cencov_2026-09-15T1540.py
% EDITED 2026-09-15 scripts/patch_continuity_connected_state_space_2026-09-15T1430.py
% =====================================================================
%  SECTION: The Euclidean metric on the state sphere, and the global
%  phase as self-referential ignorance
%  Origin: drafted in an earlier conversation (2026-06-07), edited by
%  JS; audited into the current architecture 2026-07-06.
%
%  PLACEMENT in main.tex (reordered chain, JS 2026-07-06): position 3 --
%      \input{self_reference_section_2026-07-04}
%      \input{continuity_section_2026-07-04}
%      \input{L2_metric_global_phase_section_2026_06_07}   <-- here
%      \input{wigners_friend_in_the_hopf_picture_2026-06-13}
%  (WF backward-references sec:l2-metric-fibre: order-consistent.)
%
%  AUDITED 2026-07-06 (final pass): Lamperti/Banach classification,
%  Hopf map components, submersion facts, Hardy K=N^r reading all
%  verified; the recoverable-vs-gauge-private fibre distinction
%  confirmed canonical and untouched; file pure ASCII. Fixes: opening
%  rewritten (Popper C+/C/E forward-voiced to sec:wigners-friend;
%  continuity claim made CONDITIONAL per JS -- "assuming those premises
%  can be proved", op:limit-realisability cited); Aaronson paragraph
%  de-duplicated against sec:continuity by stitch; sec:born-rule was
%  dangling paper-wide -- label minted at its canonical home (the Born
%  subsection of sec:self-reference) and the claim here recalibrated
%  ("forced, not chosen" -> lowest-degree-invariant, geometric face);
%  record-absence upgraded from interpretive to theorem via
%  prop:f1-dictionary while the self-ignorance identification stays
%  interpretive; grammar and dash punctuation repaired; field-not-fixed
%  bound pointed at sec:self-reference.
%
%  CITATIONS: existing keys Hardy2001a, Aaronson2004h, mosseri2001hopf.
%  NEW keys (RIS in L2_metric_refs_2026-07-06.ris): Banach1932,
%  Lamperti1958, Hopf1931, Urbantke2003, GluckWarnerZiller1986,
%  wooters_1981 (key spelling preserved as used; note the single-t).
%  British English; \autocite; no bullets in prose.
% =====================================================================

\section{The Euclidean metric on the state sphere, and the global phase as self-referential ignorance}
\label{sec:l2-metric-fibre}
Popper's systems $C^{+}$, $C$ and $E$ --- the predictor, its target and their shared
environment, whose argument \S\ref{sec:wigners-friend} anchors and quotes --- are, in modern
dress, Wigner, his friend and a qubit $S$ to be measured. The natural arena for analysing them
is the Hopf fibration and the Bloch sphere; but even to draw the Bloch sphere one must first
justify the $L^{2}$ metric, and this section earns that right. The principle doing the work is
the reversible continuity of \S\ref{sec:continuity} --- derived there from self-reference
conditionally, on the premises priced as Open Problem~\ref{op:limit-realisability}, of which one, that every
transformation of the regime is a phase of an inside structure, remains open; the continuum
itself is the completion the Shannon step declares (\S\ref{sec:complex-amplitudes}), the
classical convention, and \S\ref{sec:cencov-route} below reaches the same metric by a
second route that consumes it and nothing else that is open. Assuming that premise can be proved,
everything below inherits the derivation; we carry that status, and do not upgrade it.

We have established that the amplitudes of a single qubit form a complex vector space
$\mathbb{C}^{2}\cong\mathbb{R}^{4}$. We have not yet fixed which norm on this space singles out
the physical state manifold. That choice is not free: the principle of
\emph{reversible continuity} forces the Euclidean ($L^{2}$) norm, and hence the round
three-sphere, before any appeal to the Born rule. The quadratic structure usually \emph{posited}
as the Born rule is therefore already present in the metric, and the global phase emerges as a
geometrically forced fibre rather than an added convention.

\subsection{Reversible continuity selects the 2-norm}
\label{sec:reversible-continuity-l2}

The principle is the one isolated by Hardy as the single axiom distinguishing quantum from
classical probability \autocite{Hardy2001a}:
\emph{there exists a continuous reversible transformation between any two pure states.}
Operationally a reversible transformation of states is a symmetry of the state manifold, hence
norm-preserving; the demand is that such symmetries form a \emph{connected} group acting
\emph{transitively} on the pure states. Linearity is inherited from the complex vector-space
structure already in hand, so the admissible transformations are the complex-linear isometries
of whatever norm we place on $\mathbb{C}^{2}$.

Take the candidate norms to be the $\ell^{p}$ family,
$\lVert x\rVert_{p}=\bigl(\sum_{k}\lvert x_{k}\rvert^{p}\bigr)^{1/p}$, with the state manifold
the unit $\ell^{p}$ sphere. Their linear isometries are completely classified.


\begin{proposition}[Lamperti; after Banach]
\label{prop:lamperti}
For $1\le p<\infty$, $p\neq 2$, every linear isometry of $L^{p}$ (and hence of $\ell^{p}$) is a
generalised permutation: a weighted composition operator which, in the finite-dimensional case,
acts as $(x_{k})\mapsto(\omega_{k}\,x_{\sigma(k)})$ with $\lvert\omega_{k}\rvert=1$ and $\sigma$
a permutation of the coordinates.
\end{proposition}

\noindent Banach gave this for $L^{p}[0,1]$; the general $\sigma$-finite statement is Lamperti's
\autocite{Banach1932, Lamperti1958}. Its consequence for our four real dimensions is decisive.
For every $p\neq 2$ the isometry group is, in the real case, the finite group of signed
permutations (the hyperoctahedral group $B_{4}$), and in the complex case the group
$\mathbb{T}^{n}\rtimes S_{n}$ of diagonal phases composed with permutations. Even the complex
group, though it contains the continuous torus $\mathbb{T}^{n}$, fails to act transitively on
the sphere: the orbit of a basis state $e_{1}$ is only the phased basis vectors
$\{\omega\,e_{k}\}$, a measure-zero union of circles. No continuous reversible transformation
carries a basis state into a generic superposition such as $(e_{1}+e_{2})/\sqrt{2}$. Reversible
continuity therefore \emph{fails} for every $p\neq 2$.

At $p=2$ the isometry group is $O(4)$ in the real picture and $U(2)$ in the complex one; the
identity component acts transitively on the unit sphere, and any two unit vectors are joined by
a continuous one-parameter subgroup. Among the $\ell^{p}$ geometries, reversible continuity is
satisfiable at $p=2$ alone.

Aaronson states the same selection in the quantum-foundations setting and sharpens it into a
physical, rather than merely aesthetic, necessity \autocite{Aaronson2004h}: the $2$-norm is the
unique $p$-norm admitting a nontrivial continuous group of norm-preserving linear maps, every
other $p$ yielding only generalised permutations; and the alternative of keeping linearity and
reversibility while \emph{dropping} norm preservation lands one in the full general linear group
$\mathrm{GL}(n,\mathbb{C})$, which permits superluminal signalling and collapses the complexity
class $\mathsf{PP}$ into polynomial time --- the same fact \S\ref{sec:continuity} uses to price
the classical--quantum cut. Bare reversibility on the amplitude space buys only
$\mathrm{GL}(n,\mathbb{C})$; what carves the unitary group $U(n)$ out of it is preservation of
the quadratic ($\ell^{2}$) form, and reversibility of the resulting maps is then a corollary
rather than a premise.

The conclusion is that the physical state manifold of a single qubit is the unit sphere of the
Euclidean norm,
\begin{equation}
  S^{3}=\Bigl\{\,x\in\mathbb{R}^{4}\;:\;\textstyle\sum_{i=1}^{4}x_{i}^{2}=1\,\Bigr\}
  =\bigl\{\,\psi\in\mathbb{C}^{2}\;:\;\lvert z_{1}\rvert^{2}+\lvert z_{2}\rvert^{2}=1\,\bigr\},
\end{equation}
with $z_{j}=x_{2j-1}+\mathrm{i}\,x_{2j}$, carrying the round metric induced from the ambient
inner product,
\begin{equation}
  \mathrm{d}s^{2}=\sum_{i=1}^{4}\mathrm{d}x_{i}^{2}\,\Big|_{S^{3}}
  =\operatorname{Re}\langle\mathrm{d}\psi,\mathrm{d}\psi\rangle,
  \qquad
  \langle\phi,\psi\rangle=\bar\phi_{1}\psi_{1}+\bar\phi_{2}\psi_{2}.
\end{equation}
One bound deserves explicit statement, in the interest of claiming no more than is proved:
reversible continuity fixes the \emph{metric} as $L^{2}$, but not the underlying field. Real and
quaternionic quantum theory equally satisfy reversible continuity and equally carry Euclidean
state-space metrics; the restriction to $\mathbb{C}$ is pinned elsewhere in the framework ---
the analytic continuation of the information measure and the additivity of phase information,
\S\ref{sec:self-reference} --- not here.

\subsection{Hardy's independent route}
\label{sec:hardy-route}

The same conclusion is reached by Hardy along a route that never invokes
Proposition~\ref{prop:lamperti} \autocite{Hardy2001a}. His reconstruction turns on the counting
relation $K=N^{r}$, where $N$ is the number of states distinguishable in a single measurement
and $K$ the number of fiducial probabilities needed to specify a state. The continuity axiom
excludes the classical case $r=1$: there the pure states are the $N$ vertices of a simplex and
the reversible maps merely permute them --- a discrete group of motions --- so there is no
continuous path between distinct pure states, which is the simplex face of the same $\ell^{1}$
isometry fact. A simplicity assumption then selects $r=2$, yielding complex quantum theory with
its quadratic norm and trace rule. The quadratic structure is not extracted from continuity as a
separate step; it is built into Hardy's first axiom (probabilities linear in the fiducial set)
together with $K=N^{2}$.

We thus have two genuinely independent demonstrations that reversible continuity forces the
$L^{2}$ metric: Hardy's operational $K=N^{2}$ counting, and the Lamperti--Aaronson
classification of the isometries. They converge on the same step by different mathematics, and
should be cited as two roads rather than one road traversed twice. Both roads start from
continuity. A third, in the next subsection, starts one step earlier, from the value continuum
that continuity itself rests on, and reaches the same metric without passing through the
transformation group at all.

\subsection{The statistical route: \v{C}encov's theorem and the every-basis lemma}
\label{sec:cencov-route}

The value continuum that \S\ref{sec:continuity} isolates as the one open premise reaches the
$L^{2}$ metric on its own, by a route that uses no transformation group. \v{C}encov's theorem
states that the only Riemannian metric on the simplex of probability distributions invariant
under coarse-graining --- under the Markov maps that merge or relabel outcomes --- is the
Fisher--Rao metric \autocite{bengtsson2006geometry}, and in the coordinates $x_{k} =
\sqrt{p_{k}}$ the Fisher--Rao metric is the round metric on the positive orthant of the unit
sphere. The premise the theorem consumes is that the simplex is a manifold, which is the value
continuum declared in \S\ref{sec:complex-amplitudes}, and the principle it uses, that
distinguishability does not depend on how outcomes are grouped, is one the Shannon bookkeeper of
that section already grants. What it does not supply is the passage from one orthant to the
whole projective space; that is the following lemma.

\begin{lemma}[The every-basis lemma]\label{lem:every-basis}
Let $g$ be a Riemannian metric on $\mathbb{CP}^{K-1}$ whose restriction to the positive real
orthant of every orthonormal basis --- the states $\sum_{k}\sqrt{p_{k}}\,e_{k}$ with $p$ in the
simplex --- is the Fisher--Rao metric. Then $g$ is the Fubini--Study metric, whose linear lift
is the $L^{2}$ inner product.
\end{lemma}

\begin{proof}
A tangent vector to $\mathbb{CP}^{K-1}$ at a state $\psi$ is represented by a vector $\phi$
orthogonal to $\psi$, and the curve $t \mapsto \cos t\,\psi + \sin t\,\phi$ with real $t$ lies
in the positive orthant of any orthonormal basis that contains $\psi$ and $\phi$. Along that
curve the Fisher--Rao metric is the round metric of the sphere in the coordinates $\sqrt{p}$,
which is the Fubini--Study metric restricted to real amplitudes; so $g$ and the Fubini--Study
metric agree on every tangent vector at every point, the direction $i\phi$ being reached by the
basis that contains $i\phi$ in place of $\phi$.
\end{proof}

The lemma's premise, that every orthonormal basis is a legitimate measurement and none is
privileged, is the fourth face of the pre-measurement regime seen from the measurement side
(\S\ref{sec:continuity}); in the stabiliser fragment of quantum theory it fails in the same way
everything else does, only the Pauli bases being measurements there. So the $L^{2}$ metric now
rests on two legs that share exactly one premise. Through reversible continuity it uses the
value continuum, closure and homogeneity (Proposition~\ref{prop:lamperti}); through \v{C}encov
it uses the value continuum, coarse-graining invariance and the equal standing of bases
(Lemma~\ref{lem:every-basis}). Neither leg needs anything else that is open, and nothing could
avoid the continuum, since a metric is a statement about one. That is the status of everything
downstream of this section, the Hopf fibrations included: derived from the completion every physicist grants, which carries no quantum content, and from
one premise about transformations that the framework derives in outline and has not yet proved
(Open Problem~\ref{op:limit-realisability}).

\subsection{The Hopf fibration and the Bloch sphere}
\label{sec:hopf-bloch}

The round metric is exactly what makes the global phase a clean symmetry. The phase action
$\psi\mapsto\mathrm{e}^{\mathrm{i}\theta}\psi$ is a one-parameter subgroup of
$U(2)\subset SO(4)$ and so acts on $S^{3}$ by \emph{isometries}, a statement that holds only
because the metric is the $L^{2}$ one just derived. Its orbits are great circles $S^{1}$, all of
circumference $2\pi$, and the quotient is the complex projective line,
\begin{equation}
  S^{3}\big/U(1)=\mathbb{CP}^{1}\cong S^{2},
\end{equation}
the Bloch sphere. This quotient is the complex Hopf fibration
\autocite{Hopf1931, Urbantke2003, mosseri2001hopf}
\begin{equation}
  S^{1}\hookrightarrow S^{3}\xrightarrow{\;h\;}S^{2},
\end{equation}
whose projection sends an amplitude to its Bloch vector,
\begin{equation}
  h(\psi)=\bigl(\langle\psi\rvert\sigma_{x}\lvert\psi\rangle,\;
                 \langle\psi\rvert\sigma_{y}\lvert\psi\rangle,\;
                 \langle\psi\rvert\sigma_{z}\lvert\psi\rangle\bigr)
        =\bigl(2\operatorname{Re}\bar z_{1}z_{2},\;
               2\operatorname{Im}\bar z_{1}z_{2},\;
               \lvert z_{1}\rvert^{2}-\lvert z_{2}\rvert^{2}\bigr).
\end{equation}
The Hopf map is quadratic in the amplitudes, and this is the geometric face of the Born rule:
the third component is $p_{0}-p_{1}$ with $p_{k}=\lvert z_{k}\rvert^{2}$, and the squaring is
the lowest-degree $U(1)$-invariant available on the fibre --- the record-theoretic derivation of
the same rule is \S\ref{sec:born-rule}'s. Because the metric is round, the projection is a
Riemannian submersion with totally geodesic fibres of constant length
\autocite{GluckWarnerZiller1986}: the same amount of phase is hidden at every point of the base,
so the phase is hidden \emph{uniformly}, not in a position-dependent way that the metric could
partially betray.

\subsection{The fibre as self-referential ignorance}
\label{sec:fibre-self-ignorance}

The interpretation the framework places on this geometry --- and here we are reading physics
into the mathematics, so the claim is interpretive rather than proved --- is that the $S^{1}$
fibre is the observer's own un-narratable self-ignorance. The global phase is the imaginary part
of the complex information measure $S=-\ln\lvert p\rvert-\mathrm{i}\theta$: information the
observer lacks about their own state. Lacking it, they cannot communicate it; there is no
\textsc{fanout} of information one does not possess. Uncommunicable, it cannot enter any
measurement record --- and this last step is no longer interpretive:
Proposition~\ref{prop:f1-dictionary} proves the global phase absent from every classical record
a copying structure induces --- so every observable is $U(1)$-invariant; and a $U(1)$-invariant
quantity is a function on the quotient. The physical state is therefore the orbit, not the
representative point---the base $S^{2}$---and the $S^{1}$ fibre carries exactly the inaccessible
self-referential phase. It is a null direction of statistical distinguishability: two global
phases are indistinguishable by any measurement whatever \autocite{wooters_1981}, so the fibre
is a gauge direction without physical content. This phase is irrecoverably discarded, and must
remain so: it is the relative phase, living in the $S^{3}$ fibre of the quaternionic Hopf
fibration at two qubits, that is genuinely recoverable, whereas the single-qubit global phase of
the complex fibration is gauge-private without remainder. The self-reference that opened the
derivation thus reappears, at its end, as the circle the observer can never see around.

%      % EDITED 2026-09-12 scripts/patch_revert_sedenion_repoint_2026-09-12T1600.py
% EDITED 2026-09-11 scripts/patch_undefined_section_refs_2026-09-11T2030.py
% =====================================================================
%  SECTION: Wigner's friend in the Hopf picture (v2; supersedes the
%  2026-06-07 draft). AUDITED 2026-07-06 (final pass, fifth stop).
%
%  PLACEMENT (reordered chain, JS 2026-07-06): position 4 --
%      \input{L2_metric_global_phase_section_2026_06_07}
%      \input{wigners_friend_in_the_hopf_picture_2026-06-13}
%      \input{derivation_of_the_quantum_postulates}
%  sec:l2-metric-fibre is REAL (backward, order-consistent);
%  sec:fanout-boundary DRAFTED 2026-07-06 (formerly the one undrafted section);
%  the old sec:monogamy placeholder note was VOID -- no such \ref in
%  the body -- and is removed. prop:f1-dictionary and
%  sec:hopf-justification are FORWARD references (position 8): first
%  F1 use now carries the deliberate-forward flag.
%
%  DEPENDENCY: \input{monogamy_aside_2026-06-07} -- received from JS
%  and audited 2026-07-06; the \input sits inside
%  sec:wf-reconciliation, immediately after op:tsirelson-geometry
%  (relocated 2026-07-06 from end-of-file, where the aside's opening
%  antecedent had been orphaned). A bare replace-all during the earlier
%  patch briefly mangled this header block, leaving a live \input in
%  it; repaired by index-splice, and the lesson recorded: anchors that
%  also occur in comments must be uniqueness-asserted.
%
%  AUDIT RECORD: CHSH arithmetic (2sqrt(1+C^2), PR => r^2 = -2) and the
%  C^2+r^2=1 chain verified; F1 wording canonical; the
%  obsolescence-made-cancellation echo two-way with
%  sec:self-reference; redundancy/Darwinism/DiBiagio-Rovelli/Brukner
%  positioning sound; Landauer placed on amplification, correctly.
%  Fixes: first-person voice removed; \textquote replaced with
%  portable quotes; IOEP given its prop ref; stale Tsirelson
%  mini-header (wrong key name Tsirelson1980) retired;
%  op:tsirelson-geometry MINTED from the honesty paragraph's unpriced
%  owed. Popper page-pin DISCHARGED 2026-07-07: the zip bundles carry
%  OCR .txt sidecars (the earlier "no text layer" verdict was wrong --
%  ZIP-of-JPEGs defeated the PDF tools, not the text); the quoted
%  phrase sits at Part II, p. 187, now pinned in the cite. Bonus from
%  the same dig: "C+" is Popper's own notation (Part II, scan pp.
%  15-17), so the L2 opening's attribution stands as written; "E" not
%  sighted in the sampled pages -- the inline gloss covers it either
%  way. Keys unique to this file (presumed in the .bib,
%  since the paper compiles): VerstraeteWolf2002,
%  Cirelson1980QuantumInequality, PopescuRohrlich1994, Held2021,
%  Brukner2021, zurek_2009, Landauer_1961, wigner_1961.
%  British English; \autocite; no bullets in prose.
% =====================================================================

\section{Wigner's friend in the Hopf picture: one state, two charts, and the redundancy ladder}
\label{sec:wigners-friend}

The framework dissolves the Wigner's-friend paradox \autocite{wigner_1961} not by
adjudicating between the two accounts but by recognising that they are written in
two charts on a single manifold, selected as the problem requires, much as one
moves between cylindrical and spherical coordinates without either being the true
one. The epistemological chart (private, non-clonable, phase-carrying) and the
ontological chart (clonable, public, record-bearing) are related by the
\textsc{fanout} boundary (Section~\ref{sec:fanout-boundary}), and that boundary
is not a sharp wall but a graded one, parametrised by how redundantly a piece of
relational content has been copied. Read this way, the fact that a \textsc{fanout}
has occurred for the friend but not for Wigner is not a contradiction to be
removed; it is the boundary behaving exactly as defined.

Let the system be a qubit $\lvert\psi\rangle_{S}=\alpha\lvert 0\rangle
+\beta\lvert 1\rangle$ and let the friend $F$, idealised, with deliberate
economy, as a single qubit, measure it in the computational basis. The
interaction is the controlled coupling that copies $S$'s value into $F$'s
register: the \textsc{copy} comultiplication of the diagrammatic calculus
\autocite{CoeckeKissinger2017}, and in the circuit picture the very gate that, in
the two-qubit Mosseri--Dandoloff geometry \autocite{mosseri2001hopf},
activates the quaternionic Hopf fibration $S^{3}\to S^{7}\to S^{4}$ and carries
the joint state from the surface of the entanglement ball to its interior. The
friend and Wigner read the two ends of that plunge.

Two idealisations in this set-up should be named at once, because the argument
later depends on seeing past them. The first is the friend as a single qubit,
which we return to as the redundancy-one endpoint of a graded boundary
(Section~\ref{sec:wf-redundancy}). The second, less obvious and more consequential,
is the system $S$ as a pristine qubit with no prior entanglement --- a relation
with no history. No real measured particle is of this kind. A particle is itself a
reified relation between prior events, and by the time the friend couples to it, it
already carries the completed internal structure such a relation possesses: in the
framework's terms, the full octonionic $A,B,C$ register that the $S^{7}$ sector
supplies as internal degrees of freedom (\S\ref{sec:space-of-computations}). The
pristine qubit is a schematic convenience, drawn structureless so that the logic of
a single link is visible; the ``system, friend, Wigner'' triple is likewise the
schematic of \emph{one link} in the chain of observers, not its base. The chain
does not bottom out in a structureless qubit --- there is no such thing --- and we
draw on this when we ask, later, what bounds the information a single such link can
carry.

\subsection{The friend's chart: a \textsc{fanout} of redundancy one}
\label{sec:wf-friend}

Relative to $F$, the coupling is a \textsc{fanout}: $F$ acquires a fact. What is
copied is gauge-invariant, the relative phase, the correlation between $F$ and
$S$, which is a genuine observable. The global phase is never the thing copied:
there is nothing observable there to duplicate, so the $S^{1}$ fibre of the
complex Hopf fibration $S^{1}\to S^{3}\to S^{2}$ of
Section~\ref{sec:l2-metric-fibre} remains $F$'s gauge-private self-ignorance,
untouched. In Hopf terms the division is clean, but it lives at two levels, which
Proposition~\ref{prop:f1-dictionary} (\S\ref{sec:hopf-justification} --- a deliberate forward
reference) separates. The copy operation itself is fibre-committed:
copying is faithful to phase, so what $F$'s \textsc{fanout} duplicates is a lift --- a chosen
fibre-point over the outcome --- never a bare point of the base. Yet everything the operation can
broadcast is base data alone: the record it induces is independent of the lift, so the $S^{1}$
freedom, though pinned in the copier, is absent from every classical record it yields. $F$'s classical record is a
base-point; the structure that produced it committed to a fibre-point; and the gap between the two
is precisely the gauge character of the fibre. The observation has a fifty-year-old
anchor: Popper argued that a predictor's inferior self-knowledge follows from the logical
impossibility of a system's imparting information about itself without thereby rendering it out of
date --- ``making this information obsolete'', in his words
\autocites{Popper1950a}[Part~II, p.~187]{Popper1950b}. Proposition~\ref{prop:f1-dictionary}
makes the impossibility structural rather than temporal, and localises it: the datum the copier
cannot put into its own record is not its state wholesale but exactly its fibre commitment, which
does not arrive stale --- it cancels. Conditioned on the outcome it reads, $S$ is for
$F$ a definite point on the Bloch sphere,
\begin{equation}
  \rho_{S}^{(F)}=\lvert 0\rangle\langle 0\rvert
  \quad\text{(given outcome }0\text{)},
\end{equation}
a pure state of Bloch radius $r=1$ on the surface of the ball, concurrence $C=0$.

This is a \textsc{fanout}, but the minimal one. The fact $F$ holds is a
\emph{relative} fact in the sense of relational quantum mechanics
\autocite{Rovelli1996a}: $S$ definite relative to $F$, not an observer-independent
record. It is one copy, redundancy $R=1$, and as such it remains reversible by an
agent who controls the joint system. The single-qubit friend is therefore not
outside \textsc{fanout}, as a careless reading might suggest; it is the $R=1$
endpoint of the \textsc{fanout} continuum---a \textsc{fanout} that has happened,
and a fact that is real for $F$, while being maximally still reversible from
outside.

\subsection{Wigner's chart: no copy has reached him}
\label{sec:wf-wigner}

Wigner treats the laboratory as closed and no copy of the relational content has
reached him. For him the coupling is the unitary
\begin{equation}
  \lvert\Psi\rangle
  =\alpha\lvert 0\rangle\lvert F_{0}\rangle
  +\beta\lvert 1\rangle\lvert F_{1}\rangle,
\end{equation}
the move up to the larger Hilbert space, $\mathbb{C}^{2}\to\mathbb{C}^{4}$,
$S^{3}\to S^{7}$, complex Hopf to quaternionic Hopf, with the fibre grown from
$S^{1}$ to $S^{3}$. What $F$ underwent as the acquisition of a fact, Wigner
describes as correlation deposited into a joint pure state on $S^{7}$. When he
asks specifically about $S$, he traces $F$ out, and the geometry does the
bookkeeping:
\begin{equation}
  \rho_{S}^{(W)}=\operatorname{Tr}_{F}\lvert\Psi\rangle\langle\Psi\rvert
  =\lvert\alpha\rvert^{2}\lvert 0\rangle\langle 0\rvert
  +\lvert\beta\rvert^{2}\lvert 1\rangle\langle 1\rvert,
  \qquad
  r=\bigl\lvert\,\lvert\alpha\rvert^{2}-\lvert\beta\rvert^{2}\bigr\rvert,
  \quad
  C=2\lvert\alpha\beta\rvert .
\end{equation}
Tracing out discards the $S^{3}$ fibre and carries the pure point on $S^{7}$ to a
mixed point in the interior of the Bloch ball at radius $r=\sqrt{1-C^{2}}$
\autocite{mosseri2001hopf}. The inward displacement is governed by the
concurrence: it is the relational content that has migrated into the $S$--$F$
correlation, stored in the recoverable $S^{3}$ fibre. Nothing has been destroyed;
it has been relocated into the fibre, which is why Wigner can in principle
interfere the branches and reverse the interaction, the content sits at low
redundancy and no further copy has fixed it.

\subsection{One state, two charts: \texorpdfstring{$C^{2}+r^{2}=1$}{C2+r2=1}}
\label{sec:wf-reconciliation}

The system is pure for the friend ($r=1$, on the surface) and mixed for Wigner
($r<1$, in the interior). These are not rival claims about a single fact; they are
the friend at the surface and Wigner at the interior of the \emph{same} state on
$S^{7}$, bound by the Mosseri--Dandoloff constraint
\begin{equation}
  C^{2}+r^{2}=1,
\end{equation}
in which the discrepancy in Bloch radius is exactly the concurrence Wigner
retains and the friend has resolved away. This is the Inside/Outside Equivalence
Principle (Proposition~\ref{prop:inside-outside}) made concrete, the same law written by the embedded participant and by
the external observer, and the relativity of facts of relational quantum
mechanics \autocite{Rovelli1996a} realised in geometry rather than asserted: the
\textsc{fanout} boundary lies \emph{between} $F$ and $W$, so a fact exists
relative to $F$ and not relative to $W$. Neither chart is privileged; one reads in
the epistemological or the ontological chart as the question requires, and their
mutual incompleteness is a theorem of the construction, not a defect of either.

\medskip
\noindent\emph{Aside: the Tsirelson bound and the absence of PR boxes.}
\label{sec:tsirelson}
A ball has a unique centre and no point beyond it. Maximal entanglement sits at
that centre \autocite{mosseri2001hopf}, $r=0$, $C=1$, and because the
Bloch radius cannot be negative, the concurrence cannot exceed unity. For a
two-qubit pure state the maximal CHSH correlator is
\begin{equation}
  \langle\mathrm{CHSH}\rangle_{\max}=2\sqrt{1+C^{2}}=2\sqrt{2-r^{2}}  
\end{equation}
\autocite{VerstraeteWolf2002},
which the centre maximises to $2\sqrt{2}$, the Tsirelson bound
\autocite{Cirelson1980QuantumInequality}. A Popescu--Rohrlich box \autocite{PopescuRohrlich1994},
with CHSH value $4$, would require $C^{2}=3$ and hence $r^{2}=-2$: a point two
units past the centre, which the ball does not contain. The super-quantum
correlation is excluded by the very positivity that places maximal entanglement
at the centre, one cannot trace out to a more-than-maximally-mixed reduced
state. The Tsirelson bound is in this sense the image, under the CHSH map, of the
boundedness of the $L^{2}$ state ball derived in
Section~\ref{sec:l2-metric-fibre}, and the prohibition on PR boxes shares its
origin with the roundness of the sphere: both trace to the reversible-continuity
metric, not to a separate principle.

\noindent This is the two-qubit, pure-state face of the bound, and its status is
worth stating plainly. It exhibits PR correlations as having no representing state
in the Hilbert geometry, and it recovers the value $2\sqrt{2}$ at the centre; but
forcing \emph{exactly} $2\sqrt{2}$ as a theorem of the geometry, rather than
reading it from the known concurrence--CHSH relation, and lifting the argument
from two qubits to the general Tsirelson bound, remain to be made rigorous. The
closest cognate in the literature is Held's derivation of the same bound from the
geometric algebra $G^{3}\cong\mathbb{H}$ \autocite{Held2021}, the same
quaternionic structure in a different language.

\begin{openproblem}[The Tsirelson bound as geometry]\label{op:tsirelson-geometry}
Derive $2\sqrt{2}$ as a theorem of the round state geometry rather than reading it off the
concurrence--CHSH relation, and lift the argument from the two-qubit pure case to the general
Tsirelson bound --- so that the exclusion of super-quantum correlation becomes, like the exclusion
exhibited above, a consequence of the $L^{2}$ ball's positivity (\S\ref{sec:l2-metric-fibre}) and
of nothing else. Held's quaternionic derivation \autocite{Held2021} is the nearest cognate and the
natural first tool.
\end{openproblem}

\input{monogamy_aside_2026-06-07}

\subsection{The redundancy ladder: from relative fact to objective fact}
\label{sec:wf-redundancy}

The friend's fact becomes a fact for Wigner only when Wigner copies it when he
measures $F$. It becomes a fact for a third party only when that party copies it
from Wigner, and so on up the von Neumann chain. The point at which the content
stops being negotiable is therefore not a point but a limit: once it has been
copied into enough mutually independent fragments that many parties can read it
without disturbing it, no one can practically un-copy it, and it is operationally
an objective fact. This is exactly quantum Darwinism \autocite{zurek_2009}:
objectivity is redundancy. It is not an external graft, because the framework's
account of decoherence already \emph{is} this picture---many fragments each take
their own copy of the clonable content, while their private phases, each
fragment's own self-ignorance, fail to add coherently and wash out.

This supplies the grading parameter that ``graded \textsc{fanout}'' had left
unnamed. \textsc{fanout} is graded not because the operation is fuzzy but because
redundancy $R$ runs continuously: $R=1$ is the single relational copy---the
friend's relative fact, reversible by Wigner---and $R\gg 1$ is the objective
fact, effectively irreversible because un-copying it would mean recovering every
fragment. The single-qubit friend is the $R=1$ endpoint; the macroscopic friend
is the same operation at large $R$. One axis, two ends. The relation to
Di~Biagio and Rovelli's distinction \autocite{DiBiagioRovelli2021} is then exact:
their \emph{stable} facts are the high-$R$ limit and their \emph{relative} facts
the low-$R$ regime, and the standard conditional-probability law they use to
define stability is the operational signature of redundancy having grown large.
Brukner's no-go theorem \autocite{Brukner2021}, that qubits are not observers,
sits on the same axis without conflict: it bites only against calling an $R=1$
relative fact an objective one, which the framework never does---the redundancy
ladder is precisely the structure that separates the relational definiteness a
single qubit does carry from the objectivity it does not.

The thermodynamics tracks the same axis. The single relational copy is unitary and
free; the $kT\ln 2$ cost \autocite{Landauer_1961} is incurred not by the friend's
act but by the amplification and stabilisation that drive $R$ upward, when
fragments are irreversibly committed. The relative fact is thermodynamically free;
the objective fact is what costs $kT\ln 2$.

\subsection{Phase promotion, and the bump in the carpet}
\label{sec:wf-phase-promotion}

One point keeps the load-bearing irrecoverability intact while explaining what
Wigner can do that the friend cannot. What Wigner recovers is \emph{not} the
$S^{1}$ self-ignorance; the global phase remains gauge for every observer and that
circle is spent for good. The quantity the friend resolves away is the relative
phase between the $\lvert 0\rangle$ and $\lvert 1\rangle$ components, and moving up
to $\mathbb{C}^{4}$ \emph{promotes} it: the single-qubit relative phase becomes the
relative phase between $\lvert 0\rangle\lvert F_{0}\rangle$ and
$\lvert 1\rangle\lvert F_{1}\rangle$, which lives in the recoverable $S^{3}$ fibre
(itself a circle bundle over $S^{2}$, so the promotion adds a two-sphere of
relative-phase directions to the lone gauge circle). Gauge phase at one level
becomes physical, interferable phase one level up---which is why Wigner can run
the interference experiment the friend cannot. But Wigner now carries his own
fresh global-phase $S^{1}$ at the top of the joint state: his own irreducible
self-ignorance. Ascending the tower does not eliminate the self-reference, only
displaces it one level higher, and the regress up the chain is precisely the
redundancy ladder of the previous subsection. This is the framework's bump in the
carpet: each observer sits atop a Hopf tower with a single gauge circle he can
never see around, and a measurement relocates the next observer's blind circle
into his own recoverable fibre without ever closing the regress.

Two cautions close the section. First, the recoverability Wigner
enjoys is the low-redundancy limit: for a macroscopic, high-$R$ friend the
reversal is exponentially hard and the fact is objective for all practical
purposes, so Wigner's ``no record'' is the $R\!\approx\!1$ idealisation of a
graded boundary, not a free lunch. Second, naming redundancy as the grading
parameter of \textsc{fanout}, and the phase-promotion reading itself, are
syntheses of the framework's parts: each rests on established pieces (the
clonability criterion, quantum Darwinism, the tensor-product promotion of phase)
but is assembled here, and is offered as the cleanest statement of why the two
observers never disagree about anything both can check.
% =====================================================================
%  SPLICE: "Where the chain goes next" -- forward-looking coda for
%  sec:wigners-friend, pointing to sec:sedenions-generations.
%  Created 2026-08-27 per JS.
%
%  THE POINT (JS, 2026-08-27). The Wigner's-friend charts stop at
%  Wigner because the Hopf tower runs out at the octonions, while the
%  text says the chain of observers obviously continues. The sedenions
%  show what the continuation adds: the observer of Wigner (fourth
%  wire, C (x) S) brings a NEW piece of structure -- the finite S_3 of
%  Brown's theorem, the generation label -- and the observer of the
%  observer of Wigner (fifth wire, the trigintaduonions) brings nothing
%  new in kind. This coda records that, so the section's termination at
%  Wigner reads as a fact about the tower rather than an economy of the
%  exposition.
%
%  WHAT THE ALGEBRA ACTUALLY SAYS (verified 2026-08-27, notebook T --
%  notebook_T_trigintaduonion_automorphisms_2026-08-27.py):
%    - Aut(S) = G_2 x S_3 (Brown 1967), DIRECT product; the new S_3 is
%      generated by the doubling involution and Brown's explicit
%      order-three map psi, which fixes 1 and the doubling unit and
%      rotates every (a, a u) plane, a imaginary, by a third of a turn.
%    - psi permutes the three octonion HALVINGS O, psi O, psi^2 O of S
%      (each doubled by the same unit; meeting only in R). It does NOT
%      permute the three copies through the quaternion core -- those are
%      cycled by a lifted G_2 element (the colour-line three-cycle); see
%      the sedenion section's two-triples remark and the corrections
%      ledger of the same date.
%    - Brown: Aut(A_{n-1}) x S_3 = Aut(A_n) with EQUALITY for n = 4, 5, 6
%      (as quoted in gresnigtgourlayvarma2023three, sec. 4.3); Eakin and
%      Sathaye prove Brown's conjecture in general (VERIFY-CITE for the
%      exact general form -- their step is Aut(A_{n+1}) = Aut(A_n) x G,
%      G in {Z/2, S_3}; over R with the standard parameters the branch is
%      S_3 at every rung >= 4, which is what notebook T checks directly
%      at rungs 4, 5, 6). Derivation algebra g_2 from the octonions on
%      (Schafer 1954; notebook T: dim Der = 14 at n = 3, 4, 5).
%  So "nothing new at the fifth wire" is stated PRECISELY: the fifth
%  doubling adds another finite S_3 and nothing else -- no connected
%  symmetry, no fibration; the increment is the same at every rung
%  thereafter. Not "silent": REPETITIVE. The prose says so.
%
%  STATUS: the algebra is ESTABLISHED (Brown; Eakin--Sathaye; Schafer)
%  and directly checked (notebook T). The reading of the repeated finite
%  factor as what the redundancy ladder consumes is a RHYME, flagged.
%  The reading of the fourth rung's discrete orientation as the
%  generation label is the sedenion section's CANDIDATE identification,
%  pointed to, not re-argued here.
%
%  PLACEMENT: a new \subsection at the END of sec:wigners-friend, after
%  the "Two cautions" paragraph. Because that paragraph reads "Two
%  cautions close the section.", the patch also changes that one word:
%  "close the section" -> "close the argument". Both anchors
%  uniqueness-asserted (the July lesson: anchors that also occur in
%  comments must be asserted unique).
%
%  CROSS-REFS OUT (all exist in the synced source except the two
%  MAP-LABELs):
%    sec:hopf-justification   -- the "two climbs" paragraph (length vs
%                                termination); this coda sharpens it
%    sec:wf-redundancy         -- the redundancy ladder (this file)
%    sec:wf-phase-promotion    -- phase promotion (this file)
%    sec:sedenions-generations -- the sedenion section (2026-08-27 rev.)
%    prop:termination          -- Hurwitz/Adams termination
%  MAP-LABELs:
%    app:notebook-record       -- the notebook-record appendix
%                                (appendices/appendix_notebook_record_
%                                2026-07-11.tex); substitute its label.
%  NEW KEYS (RIS delivered: cayley_dickson_automorphism_refs_2026-08-27.ris):
%    eakinsathaye1990automorphisms, schafer1954cayleydickson
%  IN-LIBRARY: brown1967generalized.
%
%  OPTIONAL CONSUMER-POINTER (not patched; one clause for the "two
%  climbs" paragraph of hopf_justification_section_2026-06-08.tex, after
%  "that chain has no top --- it is the bump in the carpet, unbounded in
%  length."):
%      "(unbounded in length, and --- as \S\ref{sec:wf-chain-type}
%      records --- bounded in type: after the fourth observer nothing
%      new in kind arrives)"
%
%  EDITED 2026-08-30 (per JS): one sentence added to the fourth-wire
%  paragraph anchoring the flavour sector's irremovable phase (the
%  CP-phase recognition) on this rung, with the development at
%  sec:sedenions-mixing (2026-08-30 revision). No new keys here.
%  British English; pure ASCII; \autocite; no bullets; spaced dashes.
% =====================================================================
% SPLICE-BEGIN wf-chain-type 2026-08-27

\subsection{Where the chain goes next: one more rung, then copies}
\label{sec:wf-chain-type}

The regress just described is unbounded in length, and nothing here shortens
it: the observer of Wigner carries a gauge circle of his own, and so does the
observer of that observer, as \S\ref{sec:hopf-justification} insists when it
separates the endlessness of the chain from the termination of the tower. But
the algebra says how much of the endless chain adds structure of a new kind,
and the answer is exact: one more step, and then none.

The count runs as follows. The triple system, friend, Wigner is three wires,
and three wires are the three Hopf rungs $\mathbb{C}$, $\mathbb{H}$,
$\mathbb{O}$ --- the completed link of the earlier caveat --- which is why the
charts above stop at Wigner: the chart of whoever looks in on the triple,
$S^{15}\to S^{8}$, is the last fibration the division algebras supply
(Proposition~\ref{prop:termination}). The observer of Wigner is a fourth wire,
and the fourth wire is the sedenions,
$(\mathbb{C}^{2})^{\otimes 4} = \mathbb{C}\otimes\mathbb{S}$
(\S\ref{sec:sedenions-generations}). No fibration comes with it --- the
sedenions have zero divisors and admit no map $S^{31}\to S^{16}$ --- so the
fourth observer gains no recoverable fibre of the tower's kind: the continuous
promotion of relative phase into a larger fibre, the mechanism of the previous
subsection, has no fourth instance. What the doubling adds instead is finite.
Brown's theorem gives $\mathrm{Aut}(\mathbb{S}) = G_{2}\times S_{3}$
\autocite{brown1967generalized}: the identity component is unchanged from the
octonions, and the entire gain is a component group, generated by an
order-three automorphism that fixes the new doubling direction and rotates
every imaginary direction of the completed link into its copy by a third of a
turn. The relative orientation of the link and its copy, which at the Hopf
rungs was a continuous phase promoted into a fibre, survives at the fourth rung
only as a discrete choice among three, and that choice is the generation label
(\S\ref{sec:sedenions-generations}). The fourth observer, then, sees no new
force and no new fibre; he sees the same link three ways, and cannot gauge the
difference, because a component group has no Lie algebra. His rung is also where
the flavour sector's one irremovable phase finds its anchor: three generations
are the threshold at which the mixing matrix first carries a phase no convention
removes, and \S\ref{sec:sedenions-mixing} reads that phase --- through the
paper's negativity--contextuality chain --- as this wire's contextual cargo,
real only from one rung above.

The observer of the observer of Wigner is a fifth wire, the trigintaduonions,
and here the tower has nothing new to say. Brown proved that the next two
doublings each add exactly one further $S_{3}$ and nothing else, Eakin and
Sathaye that Brown's pattern continues at every rung, and the derivation
algebra --- the Lie algebra of the identity component --- is $\mathfrak{g}_{2}$
from the octonions onward
\autocite{brown1967generalized,eakinsathaye1990automorphisms,schafer1954cayleydickson};
the fifth and sixth rungs are also checked directly in the notebook record
(Appendix~\ref{app:notebook-record}). So the fifth observer's doubling repeats
the fourth's gift --- another third-of-a-turn symmetry between the four-wire
structure and its copy, commuting with everything below --- and adds no
connected symmetry and no fibration. Nothing new in kind arrives at the fifth
wire or at any wire after it. The increments of the tower are, in order, a
force, a force, a label, and then the same finite permutation of copies at
every step. That is the precise sense in which the chain ``obviously
continues'': it continues, and after the fourth observer it has nothing further
to announce. The framework reads the repeated factor as what the redundancy
ladder of \S\ref{sec:wf-redundancy} consumes --- a permutation symmetry among
copies, which is the freedom a record has and a type does not --- and records
that reading as a structural resonance rather than a derivation.

The two climbs of \S\ref{sec:hopf-justification} are thereby sharpened. The
chain of observers is unbounded in length and bounded in type: three forces by
the third wire, one family label by the fourth, copies thereafter. The observer
of Wigner is the last observer whose arrival changes what there is.

% SPLICE-END wf-chain-type 2026-08-27

% =====================================================================
%  FIGURE: the Wigner's-friend chain mapped onto the Cayley--Dickson
%  ladder. Created 2026-08-30 per JS: the Hopf-fibration view of
%  Wigner's friend has been load-bearing for months without being drawn;
%  the fibrations themselves are DEFERRED to Szangolies's introduction
%  (szangolies2025standardmodel, in library), and this figure carries
%  only the mapping.
%
%  CONTENT (left to right, one panel per chart, staircase rising):
%    Panel 1: F's chart of S            -- S^3 -> S^2, fibre S^1,  C, one wire
%    Panel 2: W's chart of {S,F}        -- S^7 -> S^4, fibre S^3,  H, two wires
%    Panel 3: O_W's chart of {S,F,W}    -- S^15 -> S^8, fibre S^7, O, three wires
%             (the last fibration)
%    Panel 4: the next observer's chart of {S,F,W,O_W} -- C^16 = C (x) S,
%             no fibration (S^31 -/-> S^16); the doubling's S_3 badge:
%             the generation label, a third of a turn.
%  Each describer perches on his sealed box with a dashed circle over
%  his head (his own gauge circle -- the self-referential ignorance he
%  cannot see around); dashed promotion arrows carry each circle into
%  the next chart's fibre ("gauge here -> interferable one rung up",
%  sec:wf-phase-promotion); the third arrow terminates at "no fourth
%  fibration", with what survives being the discrete third-of-a-turn
%  (sec:wf-chain-type, sec:sedenions-generations). The regress
%  continues to the right: copies only.
%
%  PREAMBLE REQUIREMENT (confirm in main.tex; add if absent):
%      \usepackage{tikz}
%      \usetikzlibrary{arrows.meta}
%  PLACEMENT: \input by patch_wigner_hopf_ladder_fig_2026-08-30.py,
%  after the two-idealisations paragraph of sec:wigners-friend, with
%  the accompanying deferral paragraph (in the patch).
%  CROSS-REFS OUT (all real): sec:wf-phase-promotion,
%  sec:wf-chain-type, sec:sedenions-generations.
%  KEY: szangolies2025standardmodel (in library).
%  British English; pure ASCII; \autocite in the caption.
% =====================================================================
\begin{figure}[t]
\centering
\resizebox{\textwidth}{!}{%
\begin{tikzpicture}[>=Stealth, line cap=round, line join=round,
    every node/.style={font=\scriptsize},
    blind/.style={densely dashed, draw=black!75},
    promote/.style={->, densely dashed, black!60, semithick},
    sight/.style={->, dotted, black!55}]
% ---- stick-figure pic --------------------------------------------------
\tikzset{stick/.pic={
  \draw[pic actions] (0,0.44) circle (0.085);
  \draw[pic actions] (0,0.355) -- (0,0.10);
  \draw[pic actions] (-0.12,0.20) -- (0,0.29) -- (0.12,0.20);
  \draw[pic actions] (-0.10,-0.10) -- (0,0.10) -- (0.10,-0.10);
}}
% ---- the staircase (Cayley--Dickson ladder) ---------------------------
\draw[black!35, thick]
  (-0.35,-0.55) -- (3.27,-0.55) -- (3.27,0.25) -- (6.89,0.25)
  -- (6.89,1.05) -- (10.51,1.05) -- (10.51,1.85) -- (13.60,1.85)
  -- (13.60,2.65) -- (16.20,2.65);
\node[anchor=north] at (1.25,-0.60) {$\mathbb{C}$ --- one wire};
\node[anchor=north] at (4.87,0.20)  {$\mathbb{H}$ --- two wires};
\node[anchor=north] at (8.49,1.00)  {$\mathbb{O}$ --- three wires};
\node[anchor=north] at (12.11,1.80) {$\mathbb{S}$ --- four wires};
\node[anchor=north, black!55] at (14.45,2.60) {$\mathbb{T},\dots$ --- five wires on};
% =======================================================================
% Panel 1: F's chart of S
% =======================================================================
\draw[rounded corners=2pt, fill=black!5, draw=black!40] (0,0) rectangle (2.5,1.32);
\filldraw[black!70] (1.25,0.72) circle (0.075);
\node[font=\tiny] at (1.25,0.45) {$S$};
\pic[black] at (2.05,1.42) {stick};
\node[anchor=west] at (2.22,1.62) {$F$};
\draw[sight] (1.95,1.60) -- (1.45,1.02);
\draw[blind] (2.05,2.14) circle (0.11);
\node[anchor=east, align=right, font=\tiny, black!75] at (1.82,2.14)
  {his own gauge circle:\\ the phase he cannot\\ see around};
\node[anchor=south west, font=\tiny, black!75] at (0.02,1.34) {fibre $S^{1}$};
\node[anchor=north] at (1.25,-0.06) {$S^{3}\to S^{2}$};
% =======================================================================
% Panel 2: W's chart of {S,F}
% =======================================================================
\draw[rounded corners=2pt, fill=black!5, draw=black!40] (3.62,0.80) rectangle (6.12,2.12);
\filldraw[black!70] (4.07,1.52) circle (0.075);
\node[font=\tiny] at (4.07,1.25) {$S$};
\pic[black!55, scale=0.85] at (4.85,1.16) {stick};
\node[font=\tiny, black!55] at (4.85,0.94) {$F$};
\pic[black] at (5.67,2.22) {stick};
\node[anchor=west] at (5.84,2.42) {$W$};
\draw[sight] (5.57,2.40) -- (5.05,1.82);
\draw[blind] (5.67,2.94) circle (0.11);
\node[anchor=south west, font=\tiny, black!75] at (3.64,2.14) {fibre $S^{3}$};
\node[anchor=north] at (4.87,0.74) {$S^{7}\to S^{4}$};
\draw[promote] (2.17,2.14) to[out=15,in=170] (3.60,2.24);
\node[anchor=south, align=center, font=\tiny, black!60] at (2.85,2.32)
  {gauge below,\\ interferable above};
% =======================================================================
% Panel 3: O_W's chart of {S,F,W}
% =======================================================================
\draw[rounded corners=2pt, fill=black!5, draw=black!40] (7.24,1.60) rectangle (9.74,2.92);
\filldraw[black!70] (7.69,2.32) circle (0.075);
\node[font=\tiny] at (7.69,2.05) {$S$};
\pic[black!55, scale=0.85] at (8.35,1.96) {stick};
\node[font=\tiny, black!55] at (8.35,1.74) {$F$};
\pic[black!55, scale=0.85] at (8.95,1.96) {stick};
\node[font=\tiny, black!55] at (8.95,1.74) {$W$};
\pic[black] at (9.29,3.02) {stick};
\node[anchor=west] at (9.46,3.22) {$O_{W}$};
\draw[sight] (9.19,3.20) -- (8.65,2.62);
\draw[blind] (9.29,3.74) circle (0.11);
\node[anchor=south west, font=\tiny, black!75] at (7.26,2.94) {fibre $S^{7}$};
\node[anchor=north, align=center] at (8.49,1.54) {$S^{15}\to S^{8}$\\[-1pt]
  {\tiny the last fibration}};
\draw[promote] (5.79,2.94) to[out=15,in=170] (7.22,3.04);
% =======================================================================
% Panel 4: the next chart of {S,F,W,O_W} -- sedenions, no fibration
% =======================================================================
\draw[rounded corners=2pt, fill=black!5, draw=black!40] (10.86,2.40) rectangle (13.36,3.72);
\filldraw[black!70] (11.24,3.12) circle (0.075);
\node[font=\tiny] at (11.24,2.85) {$S$};
\pic[black!55, scale=0.80] at (11.78,2.78) {stick};
\node[font=\tiny, black!55] at (11.78,2.56) {$F$};
\pic[black!55, scale=0.80] at (12.28,2.78) {stick};
\node[font=\tiny, black!55] at (12.28,2.56) {$W$};
\pic[black!55, scale=0.80] at (12.78,2.78) {stick};
\node[font=\tiny, black!55] at (12.78,2.56) {$O_{W}$};
\pic[black!45] at (14.35,2.78) {stick};
\draw[sight] (14.18,3.02) -- (13.42,3.12);
\draw[blind] (14.35,3.50) circle (0.11);
\node[anchor=west, align=left, font=\tiny, black!60] at (14.58,3.48)
  {the next observer,\\ and all after:\\ copies only};
\node[anchor=north] at (12.11,2.34) {$\mathbb{C}^{16}=\mathbb{C}\otimes\mathbb{S}$};
% the broken third promotion
\draw[promote] (9.41,3.74) to[out=15,in=185] (10.68,3.86);
\node[font=\scriptsize, black!60] at (10.80,3.92) {$\times$};
\node[anchor=south west, font=\tiny, black!75] at (10.88,3.74)
  {no fourth fibration: $S^{31}\not\to S^{16}$};
% the S3 badge: three dots at 120 degrees with a rotation arrow
\draw[black!75] (11.45,4.30) circle (0.17);
\filldraw[black!75] (11.45,4.47) circle (0.028);
\filldraw[black!75] (11.303,4.215) circle (0.028);
\filldraw[black!75] (11.597,4.215) circle (0.028);
\draw[->, black!75] (11.67,4.36) arc[start angle=20, end angle=100, radius=0.235];
\node[anchor=west, align=left, font=\tiny, black!75] at (11.70,4.36)
  {$S_{3}$: the generation label\\ (a third of a turn)};
\draw[promote, black!70] (10.86,3.96) to[out=60,in=200] (11.24,4.22);
\end{tikzpicture}%
}
\caption{The Wigner's-friend chain mapped onto the Cayley--Dickson ladder. Each
observer holds a chart of the sealed box below him, one rung further up the
ladder: the friend's chart of the system is the complex Hopf bundle, Wigner's
chart of the sealed laboratory the quaternionic, and the chart of the observer
of Wigner --- of the completed triple --- the octonionic, the tower's last
fibration. Every chart-holder carries a dashed circle over his head, his own
gauge phase: the self-referential ignorance he cannot see around. Measurement
promotes each circle into the next chart's fibre, where what was gauge below is
interferable above (\S\ref{sec:wf-phase-promotion}); the third promotion has no
fibration to land in --- the sedenions have zero divisors and no map
$S^{31}\to S^{16}$ --- and what survives of it is discrete: the doubling's
$S_{3}$, a third of a turn rotating the link into its copy, the generation
label (\S\ref{sec:wf-chain-type}, \S\ref{sec:sedenions-generations}). The
regress of circles continues unbounded to the right, but from the fifth wire
onward it adds copies only. The friend drawn as a single qubit is the
redundancy-one idealisation discussed in the text; for the fibrations
themselves, deferred rather than repeated here, see the introduction in
\autocite{szangolies2025standardmodel}.}
\label{fig:wf-hopf-ladder}
\end{figure}

%      % EDITED 2026-10-29 scripts/patch_continuity_tower_premise_2026-10-29T1500.py
% EDITED 2026-09-20 scripts/patch_born_cost_ledger_v2_2026-09-20T0900.py
% EDITED 2026-09-15 scripts/patch_value_continuum_declared_cencov_2026-09-15T1540.py
% EDITED 2026-09-15 scripts/patch_continuity_connected_state_space_2026-09-15T1430.py
% =====================================================================
%  SECTION: Derivation of the quantum postulates (the assembly chapter).
%  AUDITED 2026-07-06 (final pass, sixth stop).
%
%  PLACEMENT (reordered chain, JS 2026-07-06): position 5 --
%      \input{wigners_friend_in_the_hopf_picture_2026-06-13}
%      \input{derivation_of_the_quantum_postulates}
%      \input{division_algebras_section_2026-07-03}
%  Postulate 6 hands off to division-algebras, its immediate successor.
%
%  LABEL RESOLUTIONS (2026-07-06): the old placeholder
%  sec:complex-amplitudes is now minted at its source (self-reference's
%  negativity subsection), so the ref here resolves unchanged.
%  sec:born-rule COLLIDED with the label minted 2026-07-06 on
%  self-reference's Born-from-fibre-privacy subsection; resolution: two
%  routes, two labels -- Postulate 3 here is renamed
%  sec:born-invariance (the U(1)/lowest-degree geometric route) and
%  cross-stitched to sec:born-rule (the record-theoretic route), their
%  agreement flagged as over-determination. sec:local-tomography also
%  minted at source for Postulate 6's operational clause.
%
%  AUDIT RECORD: the section's status discipline was PRESCIENT -- the
%  continuity arrow was already "owed ... enters as a principle" (the
%  exact voice JS mandated for L2), and both of its confessed debts
%  have since been minted as numbered problems; the closing now cites
%  them (op:limit-realisability, op:tsirelson-geometry): "named, and
%  numbered". Postulate 6 now carries BOTH composition duty-holders
%  (Hurwitz multiplicativity + local tomography), matching the
%  continuity section's caveat exactly. Stone's theorem cited
%  (stone1932one, RIS delivered); time-as-operation-count forward
%  flag added; \textquote made portable; opener formalised.
%
%  PROVENANCE NOTE (relocated from the body, 2026-07-06, JS's words):
%  "Using Hopf fibrations to get the Born rule is something Claude AI
%  came up with on its own without my prompting." Preserved here as the
%  historical record. DISCLOSURE STATUS (JS, 2026-07-06): AI use is
%  already acknowledged in the second paragraph of the paper's
%  introduction; the remaining recommendation is only that, at
%  submission, the same fact be echoed in whatever formal
%  Declarations / Acknowledgements slot the target journal's template
%  mandates -- several publishers now require it there specifically,
%  in addition to any in-text mention.
%
%  British English; \autocite; no bullets in prose.
% =====================================================================

% ADDED 2026-07-12: closing paragraph "The postulates in module
% coordinates" -- the Clifford-module reading (skeleton + two
% supplements), with the explanatory arrow guarded per JS: the field
% comes from the Shannon/additivity/commutativity step (self-ignorance),
% and the left-action composability result is recorded as that same
% requirement in operator clothing, not an independent selection of C.
\section{Derivation of the quantum postulates}
\label{sec:postulates}
Enough pieces of the jigsaw are now assembled to derive the postulates of quantum mechanics.

The Dirac--von Neumann postulates are ordinarily laid down as independent
assumptions. In this framework they are not assumptions but a single chain: each
follows from the Inside/Outside Equivalence Principle by one move from the
preceding step. We assemble the chain here and mark the status of each link as
proved, structural, or still owed, so that no step claims more than it has earned.

The chain begins below the postulates. An embedded observer, required by the
equivalence principle to describe a totality that includes itself, must assign a
probability to a proposition about its own future state. The Lawvere--Yanofsky
diagonal \autocite{Yanofsky2003} and Popper's predictor argument
\autocite{Popper1950a, Popper1950b} show that no consistent classical assignment
exists: the self-referential proposition is undecidable to the observer before
the fact. Representing the result of that undecidable self-prediction forces a
signed, non-classical probability, which carries the observer's irreducible
self-ignorance. From this single obstruction the postulates follow.

\paragraph{Postulate 1 (complex Hilbert space).} Two independent steps fix the state space. The
field is forced first: a signed probability, read through the Shannon information measure and
analytically continued, yields a complex amplitude $S=-\ln\lvert p\rvert-\mathrm{i}\theta$,
whose phase $\theta$ is the self-referential ignorance (Section~\ref{sec:complex-amplitudes}).
The metric is forced second, and separately: the principle of reversible continuity
(\S\ref{sec:continuity}) --- that a continuous reversible transformation connects any two pure
states --- selects the $L^{2}$ norm uniquely. For every $\ell^{p}$ with $p\neq 2$ the linear
isometries are generalised permutations \autocite{Banach1932, Lamperti1958}, a group that cannot
connect a basis state to a superposition; only at $p=2$ is the norm-preserving group connected
and transitive on the sphere \autocite{Aaronson2004h}, and Hardy's operational $K=N^{2}$
counting reaches the same quadratic norm by a second route from the same continuity premise
\autocite{Hardy2001a}, and \v{C}encov's theorem with the every-basis lemma reaches it from the
value continuum alone (\S\ref{sec:cencov-route}). The state manifold is therefore the unit
sphere of the Hermitian $2$-norm in $\mathbb{C}^{n}$, taken modulo the global phase
(Section~\ref{sec:l2-metric-fibre}). Placing the metric here, before the Born rule rather than
after it, is deliberate: the Born derivation presupposes the round sphere, so the metric must be
anchored independently to avoid circularity. Two bounds belong with this step. Reversible
continuity fixes the metric, not the field---real and quaternionic quantum theory satisfy it
equally---so the restriction to $\mathbb{C}$ rests on the Shannon step, not on continuity. And
deriving reversible continuity itself from self-reference is the one arrow the chain still owes
--- now priced exactly: one open premise, that every transformation of the regime is a phase of an inside structure,
Open Problem~\ref{op:limit-realisability}, with everything downstream of it proved
(\S\ref{sec:continuity}). Until it is closed, continuity enters as a principle, and the weight
of the $L^{2}$ conclusion is carried jointly with the isometry and counting arguments and with
the statistical route.

\paragraph{Postulate 2 (Hermitian observables).}
A measurement result must be communicable: it must be copyable from one observer
to another, the \textsc{copy} operation of the diagrammatic calculus
\autocite{CoeckeKissinger2017}. The imaginary part of the amplitude is the
private self-referential phase and cannot be broadcast; only the real,
gauge-invariant content can. Observables whose spectra are real and therefore
communicable are exactly the self-adjoint operators. Hermiticity is thus the
condition that a quantity be the kind of thing a measurement can report.

\paragraph{Postulate 3 (Born rule).}
\label{sec:born-invariance}
Outcome probabilities must be invariant under the global phase, the $U(1)$
self-ignorance, realised geometrically as the $S^{1}$ fibre of the complex Hopf
fibration $S^{1}\to S^{3}\to S^{2}$ over the Bloch sphere
\autocite{Hopf1931, Urbantke2003}. There is no polynomial $U(1)$-invariant of degree one,
and the non-polynomial $\lvert\psi\rvert$ is excluded by the analyticity declared at
\S\ref{sec:complex-amplitudes}; the lowest-degree polynomial invariant is quadratic, so the
probability must be $p=\lvert\psi\rvert^{2}$. The squaring is forced
rather than merely permitted, and the argument holds for a single qubit, where
Gleason's theorem is silent. The Born quadratic is not a new structure: it is the
same $L^{2}$ form that fixed the metric in Postulate~1, now read as the
normalisation of the round sphere, and the fibre over which the phase is averaged
has constant length, so the phase is hidden uniformly
\autocite{GluckWarnerZiller1986}. This step is a structural argument rather than a
theorem in the strict sense, and is flagged as such. The record-theoretic derivation of the same
quadratic --- squared magnitudes from fibre privacy, Proposition~\ref{prop:f1-dictionary} --- is
given at \S\ref{sec:born-rule}, and the cost-ledger derivation --- the record as the closed loop of
a history and its return leg, the exponent the number of legs --- at \S\ref{sec:born-cost-ledger}:
three routes to one rule, each holding for a single qubit, the over-determination this framework
treats as evidence.

\paragraph{Postulate 4 (state reduction).}
Measurement is a \textsc{fanout}: the gauge-invariant relational content is copied
into a record, converting a private epistemological amplitude into a public
fact. For the observer who performs it, acquiring that fact is Bayesian updating,
and the projection postulate is its expression. The reduction is perspectival---a
fact relative to the observer in the sense of relational quantum mechanics
\autocite{Rovelli1996a}---and it is graded by redundancy. A single copy is a
relative fact, still reversible by an agent controlling the joint system; redundant
proliferation across many independent fragments makes the fact objective, which is
quantum Darwinism \autocite{zurek_2009} and the stable/relative-fact distinction of
Di~Biagio and Rovelli \autocite{DiBiagioRovelli2021}. ``Collapse'' and unitary
evolution are then the two charts---inside and outside the \textsc{fanout}
boundary---of one process, as the Wigner's-friend analysis shows and \S\ref{sec:fanout-boundary} formalises
(Section~\ref{sec:wigners-friend}); the irreversibility and its $kT\ln 2$ cost
\autocite{Landauer_1961, wigner_1961} accrue with redundancy rather than at the
single act. The reduction postulate is therefore recovered as a chart-dependent
description, not a fundamental dynamical law; the identification of redundancy as
the grading parameter is a synthesis of the framework's parts and is flagged as
such.

\paragraph{Postulate 5 (unitary evolution).}
Conservation of probability is preservation of the $L^{2}$ inner product, and the
inner-product-preserving linear maps are the unitaries $U(n)$---the same
isometry fact that fixed the metric in Postulate~1. Continuity in time --- operation count,
\S\ref{sec:space-of-computations}, a forward reference --- together with Stone's theorem
\autocite{stone1932one} gives $U(t)=\mathrm{e}^{-\mathrm{i}Ht}$, the Schr\"odinger
equation. Reversibility is here a corollary, not a premise: bare reversibility on
the amplitude space yields only the general linear group, and it is preservation
of the quadratic form that carves $U(n)$ out of it.

\paragraph{Postulate 6 (composite systems).}
Systems combine by the tensor product because the division algebra norm is
multiplicative \autocite{hurwitz1898composition}: composition preserves the $L^{2}$ structure
only if states multiply. The operational half closes from the other side: the locality of
tomography derived at \S\ref{sec:local-tomography} says product states determine composites,
which is the tensor product's signature. Cayley--Dickson doubling carries $\mathbb{C}$ to the next
rung, and the ascent of the Hilbert space is the ascent of the Hopf tower, complex
to quaternionic, used in the Wigner's-friend construction
\autocite{mosseri2001hopf}; the ladder itself is the business of the section that follows
(\S\ref{sec:division-algebras}).

\paragraph{Consequence (the Tsirelson bound).}
The postulates fix not only the form of the theory but the strength of its
correlations. In the two-qubit ball the maximal entanglement sits at the centre,
and a Popescu--Rohrlich box \autocite{PopescuRohrlich1994} would require a point
beyond it; the CHSH correlator is capped at $2\sqrt{2}$
\autocite{Cirelson1980QuantumInequality, VerstraeteWolf2002, Held2021}, with the value recovered
at the centre and its general derivation still owed
(Section~\ref{sec:wf-reconciliation}).

\medskip
\noindent The result is that the six postulates are one chain, not six premises.
Of the links, the field and the Born quadratic rest on the framework's signed
probability and $U(1)$-invariance arguments; the $L^{2}$ metric, the unitary group
and the tensor product are over-determined by reversible continuity, the isometry
classification and Hurwitz composition; state reduction is recovered as the
chart-relative face of \textsc{fanout}. The two outstanding debts are the
derivation of reversible continuity from self-reference and the forcing of the Tsirelson value as
a theorem rather than a recovered number --- named, and numbered:
Open Problems~\ref{op:limit-realisability} and~\ref{op:tsirelson-geometry} --- so that a reader
can see exactly where the chain is complete and where it is not.

\paragraph{The postulates in module coordinates.}
One reading of the chain deserves a closing statement, because it positions the derivation
against a mathematical skeleton the reader may know well, and because stating it carelessly would
invert the chain's explanatory arrow. In the coordinates of the division-algebra programmes
(\S\ref{sec:algebraic-dictionary}), complex Hilbert-space quantum mechanics is real Clifford
module theory plus exactly two supplements. The skeleton is Clifford throughout: the operator
algebra of $n$ qubits is the Clifford algebra of $2n$ generators with the state space its own
spinor module, half-integer spin and fermionic statistics are the module's double cover and its
graded tensor, and the classically simulable stabiliser core is the algebra's combinatorics. What
the skeleton does not contain is, first, the choice of complex structure --- the polarisation $J$
that turns a real module into a complex Hilbert space, which standard quantum theory and the
algebraic programmes alike \emph{choose} --- and, second, the Born weights on the fibre. Those two
supplements are precisely what this section has derived, and derived from self-ignorance rather
than from algebra: the field from the Shannon step (information about one's own undecidable state
must still be information --- it must add, addition commutes, and so the amplitude algebra is
commutative, excluding the quaternionic bookkeeping an observer could never balance), and the
quadratic from fibre privacy. The algebraic echo --- that the algebra's own left actions compose
only along a commutative line, the composability result recorded with the dictionary --- is not
an independent selection of $\mathbb{C}$: composition of left actions fails off that line exactly
because the doubled slot carries the \emph{reversed} product, and order-dependence is the same
disease in operator clothing that non-additivity is in information clothing. The mathematics is
the epistemic argument's shadow, welcome as consistency and as a bridge to the module-side
literature (Appendix~\ref{app:winding-number} walks the bridge for charge), but the origin of the
complex amplitude is the observer's ignorance of their own state, and the arrow runs from there
to the algebra, never the reverse.
%      % =====================================================================
%  SECTION: The Division-Algebra Ladder
%  Drafted 2026-07-03. One of the three foundations sections
%  (sec:self-reference, sec:division-algebras, sec:fanout-boundary);
%  this one resolves the dangling \S\ref{sec:division-algebras} in the
%  Hopf-justification section, whose header names it "the
%  Cayley--Dickson / Hurwitz step".
%
%  PLACEMENT in main.tex (REORDERED per JS 2026-07-06): position 6,
%  opening Act II ("why stop at C?") right before its first consumer --
%      \input{self_reference_section_2026-07-04}
%      \input{continuity_section_2026-07-04}
%      \input{L2_metric_global_phase_section_2026_06_07}
%      \input{wigners_friend_in_the_hopf_picture_2026-06-13}
%      \input{derivation_of_the_quantum_postulates}
%      \input{division_algebras_section_2026-07-03}   <-- moved here
%      \input{spacetime_from_non-computability}       <-- first consumer
%                                                          (S^7 exclusion)
%      \input{hopf_justification_section_2026-06-08}
%  The section is position-portable: all adjacency phrasing was
%  retargeted to explicit \S\ref on 2026-07-06. The S^7 hand-off
%  sentence to sec:space-of-computations WRITTEN 2026-07-06 (spacetime
%  file read; the closing now hands first to its immediate consumer,
%  then to hopf-justification). Upload verifications COMPLETE: L2 and
%  quantum-postulates do not consume the ladder; WF's references are
%  forward-voiced or self-carrying.
%
%  DIVISION OF LABOUR (deliberate, to avoid overlap):
%   - sec:self-reference (undrafted): self-reference -> complex
%     amplitudes and the Born rule; why not a real base. This section
%     TAKES that as given and starts at C.
%   - THIS section: the algebra -- the doubling construction, the
%     ledger of losses-as-physics, the internal-algebras-not-a-new-QM
%     scoping (Adler pre-empt), the two appearances of the ladder
%     (register / measurement climb, stated lightly), and Hurwitz as
%     the algebraic termination with the committed zero-divisor
%     reading.
%   - sec:hopf-justification: the geometry (Adams, parallelisability,
%     the fibrations), the join of the two classifications, the
%     entanglement-sensitivity literature, the C-direction honest debt,
%     and the epistemic reading of termination.
%
%  COMMITTED READINGS USED (sources re-read 2026-07-03):
%   - Loss ledger and level dictionary: cayley_dickson_qubits_for_review
%     sec.5 ("each step loses a symmetry, and each lost symmetry is
%     identified with a physical content").
%   - Termination reading: grounding_argument_gauge_group ("the
%     sedenions have zero divisors and cannot support consistent
%     information flow around loops") -- leads; norm-composition-as-
%     probability given as its amplitude-level restatement.
%   - Internal-not-base scoping: cayley_dickson_qubits_for_review Q2 and
%     sec.3.2 caveat (Adler objection; C-qubit treated geometrically).
%
%  HONEST DEBT: ladder proved closed at the top (Hurwitz), derived at
%  the floor. The rung-compulsion revisit-trigger FIRED 2026-07-06:
%  sec:self-reference and op:sigma-functor now exist, so the thesis is
%  minted as op:ladder-compulsion (with the CD-conjugation-as-sigma
%  candidate route), inheriting op:sigma-functor's status.
%
%  AUDITED 2026-07-06 (final pass, third stop): all four citations
%  used; all cross-refs live; CD facts, dims, alternativity, Hurwitz
%  statement verified; Adler subsection delivers everything the
%  register's compressed pointer promises; file pure ASCII; no bare
%  "frame". Fixes: opening rewritten for both predecessors (continuity
%  was inserted between after drafting); "associator's non-vanishing is
%  mass itself" recalibrated to the register's debt-against-selection
%  defusal; the O->S ledger row and Hurwitz subsection reconciled with
%  the fourth-qubit forward remark ("the terminus closes the ladder of
%  forces; what survives past it is, at most, flavour").
%
%  CROSS-REFS: sec:self-reference [sibling dangle -- resolves when that
%  section is drafted], sec:hopf-justification, sec:hopf-information-bound,
%  sec:register, sec:entanglement-classes, sec:generations, sec:masses,
%  sec:wigners-friend. All others exist.
%  CITATIONS: existing keys baez2002octonions, hurwitz1898composition,
%  furey2018three. ONE new key: adler1995quaternionic (RIS in
%  division_algebras_refs_2026-07-03.ris).
%  British English; \autocite; no bullets in prose; the ledger is a
%  table because the comparison is the argument.
% =====================================================================

\section{The Division-Algebra Ladder}
\label{sec:division-algebras}

The foundations behind us end at amplitudes twice over: natively, where the self-referential
argument forced the description of a system into the system's own state space and priced
consistency at a number system whose parts can cancel --- complex amplitudes, and the Born rule
that squares them (\S\ref{sec:self-reference}) --- and independently, where the regime's
continuity delivers the same field through Hardy's theorem (\S\ref{sec:continuity}), the
convergence of the two branches being evidence rather than accident. This section asks the question
both routes leave open. Why
should the number system stop at $\mathbb{C}$? Ordinary quantum mechanics stops there by postulate;
a framework that has forbidden itself postulates must answer, and the answer is that the climb
neither stops at $\mathbb{C}$ nor continues forever. There are exactly two more rungs, each
purchased by surrendering an algebraic law and each paying for itself in physics, and the ladder is
closed at the top by a theorem of 1898. This section sets out the algebra of that ladder: the single
construction that climbs it, the ledger of what each rung costs and buys, and the theorem that ends
it. The geometry of the same ladder --- that each rung acts on a sphere, that exactly those spheres
are parallelisable, and that the algebraic and geometric classifications are one --- is the business
of the geometry section (\S\ref{sec:hopf-justification}).

\subsection{One construction climbs the whole ladder}

The rungs are not a miscellany; a single construction generates them all. Given an algebra with a
conjugation, the Cayley--Dickson construction builds its double on ordered pairs, with a twisted
multiplication in which the conjugate enters and a conjugation $(a,b)^{*} = (a^{*}, -b)$
\autocite{baez2002octonions}. (Conventions differ in the placement of the conjugations; all yield
isomorphic algebras.) Applied repeatedly it produces the sequence
$\mathbb{R} \to \mathbb{C} \to \mathbb{H} \to \mathbb{O} \to \mathbb{S}$, of real dimensions
$1, 2, 4, 8, 16$: the reals, the complex numbers, the quaternions, the octonions, the sedenions.
The ground floor is classical: $\mathbb{R}$ is ordered and positive-definite, the arithmetic of
classical probability. The first doubling is therefore the quantum move itself --- the passage from
probabilities to amplitudes that can interfere --- and it is not made here by choice: it is the
derivation of \S\ref{sec:self-reference}, already behind us. What this section must account for is
the rest of the climb.

\subsection{The ledger}

In mathematics the doublings are usually told as a story of degeneration: each algebra is worse
behaved than the last, and the sequence is studied for where it breaks. This framework reads the
same ledger in the opposite direction. Each rung is bought by surrendering a law, and each
surrendered law is exactly a piece of physics --- the losses are the assets, and the comparison is
sharp enough to be the argument:

\begin{center}
\begin{tabular}{llll}
\hline
doubling & dimension & what is lost & what the loss becomes \\
\hline
$\mathbb{R} \to \mathbb{C}$ & $1 \to 2$ & the order of the reals & phase and interference; amplitudes \\
 & & & that cancel; the Born rule (\S\ref{sec:self-reference}) \\
$\mathbb{C} \to \mathbb{H}$ & $2 \to 4$ & commutativity & non-commuting actions: weak isospin \\
 & & & $SU(2)$, the $B$ register (\S\ref{sec:register}) \\
$\mathbb{H} \to \mathbb{O}$ & $4 \to 8$ & associativity & the associator: colour, chirality, the \\
 & & & generations, and mass as debt \\
 & & & (\S\ref{sec:generations}, \S\ref{sec:masses}) \\
$\mathbb{O} \to \mathbb{S}$ & $8 \to 16$ & division itself & no further force: zero divisors appear \\
 & & (zero divisors) & and the ladder of sectors ends (below) \\
\hline
\end{tabular}
\end{center}

Walking the rows: the order lost at the first doubling is the possibility of interference --- parts
of a probability that can oppose one another --- which is where the derivation of \S\ref{sec:self-reference}
lives. The commutativity lost at the second is the existence of actions whose order matters; the
unit quaternions form the group $SU(2)$, and the framework houses weak isospin there, as the
internal algebra of the $B$ register whose order-sensitive, non-commuting dynamics the interaction sections develop
(\S\ref{sec:register}, \S\ref{sec:entanglement-classes}). The associativity lost at the third ---
the octonions retain the weaker property of alternativity \autocite{baez2002octonions} --- creates
the associator $[a,b,c] = (ab)c - a(bc)$, and the associator is where most of this paper's physics
is transacted: the stabiliser of a preferred complex direction inside the octonions' automorphisms
is the colour group (\S\ref{sec:hopf-justification}); the triality structure and its $K_3$
obstruction generate the three generations and the sterile channel (\S\ref{sec:generations}); and
the associator is the currency of mass --- the debt against the selected direction that
\S\ref{sec:masses} prices, owed by whatever action moves the selection and not by residence at
the rung, which is the defusal \S\ref{sec:register} makes explicit. The algebraic facts
in the third column are theorems; the identifications in the fourth are established in the sections
cited, each carrying its own epistemic flag there, and are not re-litigated here. What this table
asserts on its own account is only the correspondence's shape: one law surrendered per rung, one
sector of physics gained per law.

\subsection{Internal algebras, not a new quantum mechanics}

A reader versed in the quaternionic-quantum-mechanics literature will at this point raise a fair
objection: quantum theory over $\mathbb{H}$ as a foundation is known to be troubled --- the
tensor-product and composite-system difficulties set out by Adler and others
\autocite{adler1995quaternionic} --- and quantum theory over $\mathbb{O}$ founders earlier still,
since a non-associative algebra does not furnish a Hilbert space at all. The framework's position
must therefore be stated exactly, because it does not make the claim the objection targets. The base
of the physical Hilbert space here is $\mathbb{C}$ and remains $\mathbb{C}$; the higher algebras
enter as the \emph{internal} algebras of the register --- $\mathbb{H}$ for the $B$ qubit,
$\mathbb{O}$ for the $C$ qubit (\S\ref{sec:register}) --- in a tensor structure whose outer theory is
ordinary complex quantum mechanics. This is the same move by which the division algebras enter the
constructions of Furey and the related programmes without any proposal to rebase quantum mechanics
itself \autocite{furey2018three}. For the octonionic register the point is sharper still:
$\mathbb{O}^{2}$ is not a Hilbert space, and the framework accordingly treats the $C$ qubit not in
the Dirac--von Neumann manner but geometrically, as the sphere $S^{15}$ with its Hopf structure ---
which is precisely the form in which \S\ref{sec:hopf-justification} needs it. Nothing in this paper therefore
stands or falls with quaternionic or octonionic quantum mechanics as foundations; the ladder is
climbed by the registers' internal algebras, not by the theory's base field.

\subsection{The same ladder, twice}

The ladder appears in this framework in two guises, and it will help the reader to have them
distinguished now. Statically, it is the anatomy of a single relation: the three qubits of the
register take their values one rung apart --- $A$ over $\mathbb{C}$, $B$ over $\mathbb{H}$, $C$ over
$\mathbb{O}$ --- so that every particle carries the whole ladder as its internal structure
(\S\ref{sec:register}). Dynamically, it is the shape of measurement: a single qubit's pure states
fill $S^{3}$, an entangling measurement carries the joint state into $S^{7}$, and a third party into
$S^{15}$ --- the correspondence between qubit entanglement and the Hopf fibrations that
\S\ref{sec:hopf-justification} takes from a mature literature --- so that each act of measurement ascends one rung --- the climb the reader has already made,
narrated in \S\ref{sec:wigners-friend} --- and every completed relation is a terminated climb, which is
the reading \S\ref{sec:hopf-information-bound} gives the tower's endpoint. One ladder, seen twice:
as the anatomy of a relation and as the dynamics of relating. The framework regards this double
appearance not as a coincidence but as the same structure viewed statically and in motion, and the
sections cited carry that identification's weight.

\subsection{Hurwitz closes the ladder}

The construction itself would double forever; what stops it is a classification theorem. By
Hurwitz's theorem, the normed division algebras --- those in which the norm composes,
$\lvert ab\rvert = \lvert a\rvert\,\lvert b\rvert$, and no non-zero elements multiply to zero ---
are exactly $\mathbb{R}$, $\mathbb{C}$, $\mathbb{H}$ and $\mathbb{O}$
\autocite{hurwitz1898composition,baez2002octonions}. The next doubling, the sedenions, has zero
divisors: non-zero elements whose product vanishes. The framework's reading of this failure is the
one that matters physically. Zero divisors cannot support consistent information flow around loops
--- a cycle of compositions can silently annihilate what it carries --- and a framework in which
closed computational paths are load-bearing (\S\ref{sec:hopf-justification},
\S\ref{sec:space-of-computations}) cannot be built on an algebra that permits it. At the level of
amplitudes the same failure reads: the composition law
$\lvert ab\rvert = \lvert a\rvert\,\lvert b\rvert$, which carries probability multiplicatively
through every rung of the ladder, fails beyond the octonions, so that two non-vanishing amplitudes
could compose to nothing and probability would be destroyed by mere composition. The last algebra
fit to carry amplitudes is $\mathbb{O}$. One rung past the terminus is not entirely silent,
however: the sedenions' internal structure --- their automorphism group and their octonion
subalgebras --- is the recorded candidate arena for tagging the generations
(\S\ref{sec:interactions}), and it enters exactly as this section's scoping permits, internally
and never as a base. The terminus closes the ladder of forces; what survives past it is, at most,
flavour. This is the algebraic half of the termination; the
geometric half --- that no sphere above $S^{7}$ is parallelisable --- and the theorem joining the
two halves are the opening business of \S\ref{sec:hopf-justification}, and the epistemic reading of the same endpoint,
as a bound on the information a single relation can carry, is given in
\S\ref{sec:hopf-information-bound}.

\subsection{The honest debt}

As in the geometry section, we close by stating plainly what is proved, what is established, and what is
owed. Proved: the construction, the losses, and the termination are theorems --- the Cayley--Dickson
doubling and its ledger of surrendered laws are standard algebra, and Hurwitz's classification is
among the durable results of the subject. Established: the physical cashing of each rung --- the
Born rule at $\mathbb{C}$, weak isospin at $\mathbb{H}$, colour, generations and mass at
$\mathbb{O}$ --- is carried out in the sections cited above, each under its own epistemic flag. Owed,
and stated without disguise: the ladder as presented here is closed from above and derived at its
floor, but it is not yet derived rung by rung from below. Hurwitz forbids a further rung; the
argument of \S\ref{sec:self-reference} forces the first doubling, from the ordered arithmetic of
classical probability to amplitudes that interfere. That the same self-referential closure, applied
to a system whose amplitudes already take values at one rung, compels the next --- that relational
closure at one level forces precisely the Cayley--Dickson double of that level --- is the
framework's reading of the climb, and it is a thesis, not a theorem --- but the formal apparatus
against which it can be posed now exists, and we pose it.

\begin{openproblem}[The compulsion of the intermediate rungs]\label{op:ladder-compulsion}
Derive, within the categorical apparatus of Open Problem~\ref{op:sigma-functor}, that
self-referential closure at one rung compels the Cayley--Dickson double of that rung, so that the
intermediate steps $\mathbb{C} \to \mathbb{H} \to \mathbb{O}$ follow from the same principle
that forces the floor. The candidate route is the construction's own shape: the Cayley--Dickson
pair with its conjugation, $(a, b)^{*} = (a^{*}, -b)$, is an involution-equipped doubling, and the
route reads it as the algebraic shadow of adjoining the describer --- closure of level-$A$
descriptions under the inside/outside exchange requiring exactly the twisted pair. The problem
inherits the status of Open Problem~\ref{op:sigma-functor}, whose apparatus it presupposes.
\end{openproblem}

If the derivation lands, the whole ladder becomes a consequence of the same self-reference that
opened it. Until then, the honest statement is that the ladder's necessity is established at its
ceiling and at its floor, while the compulsion of the intermediate rungs is inherited from the
physics they are shown to carry rather than derived in advance. What this section hands on is, accordingly, an algebraic ladder with a floor, a ceiling and a
ledger. Its first consumer is immediate: \S\ref{sec:space-of-computations} spends the terminal
fact --- the non-associativity of $S^{7}$'s own multiplication --- as geometry, excluding the
seven-sphere as a ruler and fixing the dimension of space at three.
\S\ref{sec:hopf-justification} then supplies the geometry of the fibrations themselves, and the
join.
   <-- moved here
%      % EDITED 2026-10-07 scripts/patch_infinite_not_flat_2026-10-07T1400.py
% EDITED 2026-09-27 scripts/patch_spacelike_inversion_2026-09-27T1800.py
% EDITED 2026-09-21 scripts/patch_op_sterile_gravity_boundary_2026-09-21T1300.py
% EDITED 2026-09-12 scripts/patch_spacetime_bloch_headline_2026-09-12T1400.py
% EDITED 2026-09-11 scripts/patch_splice_loops_prospects_2026-09-11T0900.py
% EDITED 2026-09-10 scripts/patch_cite_rubino_2026-09-10T0700.py
% EDITED 2026-09-10 scripts/patch_quoting_and_typography_2026-09-10T0620.py
% =====================================================================
%  SECTION: Spacetime from non-computability (sec:space-of-computations).
%  AUDITED 2026-07-06 (final pass, seventh stop).
%
%  PLACEMENT (reordered chain, JS 2026-07-06): position 7 --
%      \input{division_algebras_section_2026-07-03}   <-- hands off here
%      \input{spacetime_from_non-computability}
%      \input{hopf_justification_section_2026-06-08}
%  Division-algebras' closing now carries the S^7 hand-off sentence
%  (written 2026-07-06 after this file was read).
%
%  INPUT DEPENDENCIES (both held in the audited outputs):
%      entropic_gravity_family_paragraph_2026-07-03.tex  (sec:jacobson)
%      gr_nonlocality_revision_2026-06-13.tex            (sec:gr-nonlocality)
%  The 2026-07-04 additions (asymmetry paragraph; Penrose second route;
%  acceleration/Unruh arc) are FULLY PASTED; the separately uploaded
%  additions file is redundant and can be archived.
%
%  AUDIT RECORD: opener formalised (Russell typo; bare ref retargeted
%  to the freshly minted sec:complex-amplitudes); the "rotations of
%  S^7" wording made precise at BOTH occurrences (Moufang loop, not a
%  group; SO(8) disclaimer added -- referee hazard removed); two \\
%  paragraph-break bugs fixed; dark-matter paragraph retargeted to
%  sec:generations and stitched to op:sterile-gravity (the zoo's
%  demotion and this paragraph's priority condition are one question);
%  CP paragraph stitched to the fourth-party reading of
%  sec:interactions; stale mid-file mini-header retired
%  (sec:k3-sterile placeholder dies with it); the section's narrated
%  debts NUMBERED at last -- op:comp-metric and op:three-dimensions
%  minted in the Open-problems subsection, cited from the in-text
%  "named open problems" line and from the closing summary.
%
%  KEYS: adams1960hopf matches corpus. Unique to this section, presumed
%  in the .bib since the paper compiles: Spence-Brown-1969 (suggest
%  normalising the key spelling at leisure), VanRaamsdonk2010,
%  MaldacenaSusskind2013, MuellerMasanes2013, MasanesMuellerEtAl2014,
%  toffoli_2003, Chaitin1998, Penrose1996, HohnMueller2016,
%  GarnerMuellerDahlsten2017, Jacobson1995, Bekenstein1973, tHooft1993,
%  Susskind1995, spekkens2007toy, Andrianov2001, Sakharov1967,
%  Planck2020, Riess1998. Shared keys verified: small_2006, Penrose1994,
%  unruh1976notes (swap-flag stands), baez2002octonions, Wheeler1990,
%  furey2018three, DiBiagioRovelli2021, Rovelli1996a, Spekkens toy.
%  British English; \autocite; no bullets in prose.
% =====================================================================

\section{Spacetime from non-computability}
\label{sec:space-of-computations}

This section is a deliberate diversion: it takes the idea that non-computable steps are measured
by imaginary information (\S\ref{sec:complex-amplitudes}) and asks what geometry it buys.
Spencer-Brown records that Bertrand Russell thought well of the underlying move
\autocite{Spence-Brown-1969}; the reader may judge whether the framework spends it responsibly.

The programmes that derive spacetime from quantum mechanics---Van~Raamsdonk's
\enquote{building spacetime with entanglement} \autocite{VanRaamsdonk2010}, the
ER\,$=$\,EPR conjecture \autocite{MaldacenaSusskind2013}, and M\"uller's
information-theoretic recovery of three-dimensional space from the qubit
\autocite{MuellerMasanes2013}---take entanglement or information as primitive and
grow geometry from it. The route taken here is the same in spirit but inverted in
order: rather than positing spatial dimensions and forbidding superluminal
signalling, we \emph{define} spatial separation as a measure of non-computability
and read the dimension, the signature, and the Lorentz group off the result. The
proposal originates in the author's CASYS'05 paper \autocite{small_2006}, which
first suggested treating spacetime as a space of computations and attempted to
justify three dimensions on this basis; the present account tightens and extends
it. This is the framework's most structural and least formally complete arc; the
three-dimensionality result is a structural conclusion awaiting a formal proof,
and several of the links below are conjectural readings. Each is flagged as such,
so the reader sees exactly where the chain is load-bearing and where it is an
invitation.

\subsection{Timelike and spacelike distance as computation}
\label{sec:comp-distance}

A computation carries a state from axioms to a theorem in a number of steps.
Counting the computable steps between two states gives a real measure---time. A
physical clock counts completed irreversible loops, so proper time literally
counts loop completions and the metric $\mathrm{d}\tau^{2}=g_{\mu\nu}\,
\mathrm{d}x^{\mu}\mathrm{d}x^{\nu}$ is a loop-density tensor, the conversion from
coordinate intervals to loop counts; Toffoli's reading of the Lagrangian as a
count of action quanta is the same statement \autocite{toffoli_2003}.

Spacelike separation is the complementary, non-computable measure. Define the
non-computable distance between two propositions, relative to an axiom system, as
the number of bits that must be added to the axioms to make the second derivable.
A closed timelike curve is a computation whose output is its own input; by the
fixed-point theorem it has a unique consistent solution or none, and the
existence of that solution is one bit not derivable from any axiom system that
does not already contain the loop---one bit per curve, in the manner of Chaitin's
$\Omega$ \autocite{Chaitin1998}. The spacelike distance between two events is then
the number of such curves that would have to be closed to make them causally
connected. Because each added bit is a non-computable proposition, it is carried
by the imaginary part of the complex information measure
$S=-\ln\lvert p\rvert-\mathrm{i}\theta$ derived in
Section~\ref{sec:complex-amplitudes}: computable distance is counted in real
units, non-computable distance in imaginary ones. The Minkowski signature, with
spacelike intervals imaginary relative to timelike ones, is the immediate reading. The convention is not new to this paper: it was proposed in the author's CASYS paper two
decades ago, where counting non-causal steps in imaginary numbers was introduced and its formal
proof flagged, already then, as a key open problem \autocite{small_2006} ---
Open Problem~\ref{op:comp-metric} is that flag's direct descendant, stated now with the machinery
to be answerable.
That the resulting structure is a genuine metric space, and that the signature
follows rather than being fitted, are open problems---at present the
identification is motivated by reproducing what nature does, not proved.

%% VERIFY-SUPPORT 2026-09-27: conway_2003 is cited below for "any two octonions
%% generate an associative subalgebra" (Artin's theorem; a Moufang loop is
%% diassociative) and harvey1982calibrated for "a three-plane in Im O is
%% associative iff the associator of a basis vanishes, and the associative
%% three-planes are the imaginary parts of the quaternionic subalgebras"
%% (Harvey--Lawson, ch. IV). Neither passage was read in the session that wrote
%% this; both statements are checked numerically in Convention A by
%% scripts/antitelephone_network_associativity_2026-09-27T1900.py (PASS).
A reader used to the standard order of ideas should pause here, because the order has been
inverted, and the inversion is used below. In special relativity the tachyonic antitelephone is a
corollary: a signal faster than light, composed with a boost, would arrive before it was sent and
close a timelike loop, therefore there is no such signal \autocite{tolman1917theory,
benfordbooknewcomb1970tachyonic}. The metric comes first and the impossible telephone is a check
on it. Here the telephone comes first. Two observers are spacelike separated when an antitelephone
between them would close a loop; the loops such telephones would close, and carry nothing round,
are what spacelike separation is, their number is the distance, and the metric is read off from
them. Nothing is lost, since the corollary returns as a tautology: an antitelephone that carried
a record would close a causal loop, and its two ends would not have been spacelike after all.

What is gained is a definite meaning for a spatial dimension. Picture many observers, each holding
an antitelephone to each of the others. A message relayed from one observer to another through a
third composes two transports round two loops, and the network is consistent only if every relay
composes to a single transport whatever the bracketing of the relaying: a spatial dimension is a
family of loops whose transports compose associatively, which is a group law on a parallelisable
sphere. That is a condition a family of loops can fail, and \S\ref{sec:three-dimensions} shows
that the loops of the seventh sphere fail it. Two antitelephones always compose, because any two
octonions generate an associative subalgebra \autocite{conway_2003}; three compose only when their
imaginary parts lie in one quaternionic subalgebra, which is what the associative three-planes of
Harvey and Lawson are \autocite{harvey1982calibrated}, and a third link outside that subalgebra is
seen by the network exactly through its component outside it. The largest consistent network of
antitelephones inside $S^{7}$ is therefore an $S^{3}$, and there are three spatial directions.
When the seventh sphere is set aside there as a ruler, which is the octonionic spacetime of the
division-algebra literature, the claim is not that $\mathbb{R}^{8}$ has no metric, nor that the
ten-dimensional Minkowski space of $SL(2,\mathbb{O}) \cong \mathrm{Spin}(9,1)$ fails to exist
\autocite{baez2002octonions}; it is that the octonionic directions are not loops on which
antitelephones compose, and so are algebra rather than distance. In the standard order that
sentence has no content; in this one it is the reason space has three dimensions.

The split just defined explains the most familiar asymmetry in physics: that one can move back and
forth in space but not in time. \S\ref{sec:continuity} identifies the pre-measurement regime with
continuous reversibility --- its transformation group is connected, and every motion in it is
undoable. Spacelike separation, being the non-computable sector, is exactly the part of the
relational structure on which no record has been written; motion across it is motion within the
reversible regime, and it inherits that regime's continuity and its two-sidedness: a place can be
revisited because going and returning writes no record that forbids the return. Timelike separation
is the record chain itself --- proper time counts completed irreversible loops, as above --- and a
moment cannot be revisited because a moment \emph{is} a stratum of records, and returning would
require unwriting them, which is what \textsc{fanout} forbids (\S\ref{sec:fanout-boundary}). The
asymmetry between space and time is therefore not an extra postulate about the manifold: it is the
pre-/post-\textsc{fanout} asymmetry read as geometry --- space is what the reversible group feels
like from inside, time is what the record chain feels like --- and the arrow-of-time discussion
later in this section is the same asymmetry read thermodynamically. The identification is the
framework's structural reading of the split, offered in the same register as the rest of this
subsection.

\subsection{Three dimensions of space}
\label{sec:three-dimensions}

The closed timelike curves that define spacelike distance require information to
be parallel transported consistently around a loop, and consistent parallel
transport on a sphere requires it to be parallelisable. By Adams' theorem the
only parallelisable spheres are $S^{1}$, $S^{3}$ and $S^{7}$
\autocite{adams1960hopf}, the unit spheres of $\mathbb{C}$, $\mathbb{H}$ and
$\mathbb{O}$. Impose now the requirement that information arriving from outside
the causal zone must appear to possess its own causal history---that is, to be the
result of a sequence of information-preserving operations, hence a sequence of
group operations. The spheres $S^{1}$ and $S^{3}$ are groups under their own multiplication --- $U(1)$ and
$SU(2)$ --- and their operations preserve information; $S^{7}$'s multiplication is non-associative,
a Moufang loop rather than a group, so $S^{7}$ cannot serve as a ruler. (The rotation group
$SO(8)$ of the seven-sphere is of course a group; what fails is the sphere's \emph{own}
composition law, the one a ruler would use.) The same non-associativity
forbids a metric built from the sphere's own composition law outright: path
composition $d(A,C)$ through $B$ depends on the bracketing, so the triangle inequality is not merely violated but ill-defined, and
no geodesic structure exists. (What is untouched should be said: the Euclidean metric of
$\mathbb{R}^{8}$, and the ten-dimensional spacetime that $SL(2,\mathbb{O})$ furnishes, are metrics
of the ordinary kind; what fails is a ruler made of antitelephone loops, which is what a spatial
dimension is in this section, as \S\ref{sec:comp-distance} set out. Nor is the failure total:
relays between any two links compose, and among three they compose exactly when the links lie in
one quaternionic subalgebra, so the largest consistent network inside $S^{7}$ is an $S^{3}$
\autocite{conway_2003,harvey1982calibrated}.) Since the abelian $S^{1}$ alone cannot carry
non-abelian structure, spacelike distance must be three-dimensional. Dimension
three is an output, not an input \autocite{small_2006}.

The division of labour is then exact. $S^{1}$ ($\mathbb{C}$) contributes the time
dimension and the $U(1)$ of electromagnetism; $S^{3}$ ($\mathbb{H}$) contributes
the three spatial dimensions and the $SU(2)$ of the weak interaction; $S^{7}$
($\mathbb{O}$) contributes no geometry but, as a filter on what information may be
exchanged, the $SU(3)$ of the strong interaction. A sphere that cannot be a ruler
can still be a filter, which is why the strong force resists geometric
interpretation while electromagnetism and the weak force do not. This is the same
$\mathbb{C}\oplus\mathbb{C}^{3}$ split that yields $G_{\mathrm{SM}}$ in
Section~\ref{sec:entanglement-classes}, now read as the reason spacetime is $(3+1)$-%
dimensional. The conclusion converges with M\"uller and Masanes, who prove from
information-theoretic postulates, assuming no quantum theory, that minimal
direction-carrying systems with continuous reversible dynamics force $d=3$ and
quantum theory on two qubits \autocite{MuellerMasanes2013}, and with the result
that entanglement exists only for a three-dimensional Bloch ball
\autocite{MasanesMuellerEtAl2014}; the framework supplies the reason such systems
exist, Cayley--Dickson doubling from self-reference, while their work supplies
the rigorous dimensional uniqueness. The status of the framework's own route is a
structural result requiring formal proof: the one-bit-per-curve identification and
the precise reduction of the $S^{7}$ filter to $SU(3)$ rather than its full automorphism group
$G_{2}$ are numbered below (Open Problems~\ref{op:comp-metric} and~\ref{op:three-dimensions}).

A scoping remark, and then a second route to the same count. The computational model of spacetime
proposed here has many detailed parallels with the real thing, but a parallel is not an identity,
and the identity is not claimed in this paper: establishing it is work deferred to the second paper.
What can be claimed now is that the parallel is productive in a specific direction --- it suggests a
derivation of the dimension count carried out entirely within general relativity, with no
computational premise. Acceleration and gravitation tip the light cone, and the tipping can be
carried to the extreme: Penrose's discussion runs the cone through a full circle
\autocite[382]{Penrose1994}, at which point the tipped direction closes on itself --- which is to
say, the extreme of tipping is a closed timelike curve, the very object whose closure the
computational definition of \S\ref{sec:comp-distance} counts in bits. The two routes are one route
entered from opposite ends. Now demand of general relativity only this: that observers accelerating
with respect to one another, and with respect to unaccelerated observers, retain a consistent
conception of the distance between them. Each observer's tipping traces a path on a sphere of
directions, and mutual consistency requires the associated rotations to compose without regard to
bracketing --- to form a group. Compositions on $S^{1}$ and $S^{3}$ do, forming $U(1)$ and $SU(2)$ as above; compositions on
$S^{7}$ do not, the multiplication being non-associative, so observers tipped across an
$S^{7}$-worth of directions could not agree on distance even pairwise: the same failure that forbade
$S^{7}$ a metric above forbids it here a shared ruler. The largest sphere admitting a consistent
conception of distance under mutual acceleration is therefore $S^{3}$, and space has at most three
dimensions; with the floor argued above --- the abelian $S^{1}$ alone cannot carry the required
structure --- the count is exactly three. The status matches the subsection's. The tipping and its
extreme are general relativity's own \autocite{Penrose1994}; the consistency demand is
elementary; the associativity filter is the algebraic fact used throughout; the reading of composed
tippings as motion on the parallelisable spheres is the structural step, shared with the
computational route. The convergence of the two routes on the same endpoint is offered as evidence
for the parallel, while the identity it gestures at remains the second paper's burden.

\subsection{Recovering the Lorentz group}
\label{sec:lorentz}

A measurement changes what the observer knows, hence, relationally, the causal
structure: it tips the light cone. This is Penrose's association of cone-tipping
with state reduction \autocite{Penrose1996} run in the opposite causal direction,
reduction driving the tipping rather than the reverse, and the two are then one
operation in two charts---both described by $SL(2,\mathbb{C})$. The algebraic
backbone is the division-algebra correspondence
$SL(2,\mathbb{A})\cong\mathrm{Spin}(1,d{+}1)$ \autocite{baez2002octonions}: $\mathbb{R}$
gives $1{+}1$, $\mathbb{C}$ gives $3{+}1$, $\mathbb{H}$ gives $5{+}1$ and
$\mathbb{O}$ gives $9{+}1$ dimensions. The computability split that singles out a
preferred $\mathbb{C}$ fixes the fundamental state space as $\mathbb{C}^{2}$ and
hence $SL(2,\mathbb{C})=\mathrm{Spin}(1,3)$: the Lorentz group of $3{+}1$
spacetime. The qubit and Minkowski space are not merely both Euclidean-and-three%
-dimensional by coincidence; the Hermitian $2\times2$ matrices are
$\mathbb{R}^{1,3}$, $SL(2,\mathbb{C})$ acts on them as the Lorentz group, and the
qubit density operators fill the future cone---the correspondence M\"uller's group
develops operationally, recovering Lorentz transformations from quantum
communication \autocite{HohnMueller2016} and the complex and quaternionic qubit
from the relativity of simultaneity \autocite{GarnerMuellerDahlsten2017}.
In these coordinates the Bloch ball is the unit-trace slice of the cone and the Bloch sphere its
boundary rays, the celestial sphere on which $SL(2,\mathbb{C})$ acts by M\"obius
transformations; state reduction sends a point of the ball to a point of the sphere, and a boost
sends the sphere to itself, which is what lets the two be one operation in two charts.
The $SL(2,\mathbb{A})$ correspondence is rigorous; its identification with
measurement-as-tipping is the framework's structural reading.

One consequence of the identification has been implied throughout this section and deserves its own
paragraph, because the next paragraph turns on it. If measurement creates a relation between
observer and observed, and if --- relationally --- a new relation alters the causal structure, then
the tipping of the cone is not merely \emph{associated} with the creation of the relation: it is
what the creation of the relation is, seen geometrically. And a tipping of one's cone against the
congruence of one's neighbours already has a name in physics: acceleration. The framework's
statement, made explicit here, is that \emph{creating a new relation between observer and observed
is experienced as acceleration}. The resonance with the equivalence principle is exactly what it
should be: the principle makes acceleration and gravitation locally indistinguishable, Penrose's
association --- run in this section's direction, reduction driving the tipping --- makes state
reduction and cone-tipping one operation, and the framework's identification closes the triangle,
measurement and acceleration meeting where the principle says they must. The claim is structural:
the framework's reading of the cone-tipping established above, not an independent derivation.

Taken literally, the identification has a measurable face, and it is a familiar one: the Unruh
effect, read in this section's own ledger. Recall the definitions of \S\ref{sec:comp-distance}:
timelike distance counts computable steps, spacelike distance counts non-computable bits. What an
observer can compute --- what their \textsc{fanout}s can convert into classical records --- is fixed
by their causal structure, and the previous paragraph has just made that structure
relation-dependent. An accelerated observer's cones are tipped, and with them the boundary of the
recordable: relations that for the inertial observer lie wholly in the reversible, pre-record sector
--- the entanglement structure of the vacuum, non-computable for them, spacelike in the
computational sense --- fall, for the accelerated observer, partly on the recordable side. The mechanism is the one
Chaitin identified for formal systems \autocite{Chaitin1998}: a proposition undecidable relative to
an axiom system becomes decidable when information is added to the axioms --- and
\S\ref{sec:comp-distance} defined spacelike distance as exactly the bits whose addition closes the
gap. Creating a relation is adding information to the observer's axiom base, so timelike and
spacelike distance are converted into one another by the addition and subtraction of information,
and acceleration is the continuous version of that addition. A
relation newly recordable is a detection. Particles appear from nowhere: an accelerated detector in
the inertial vacuum registers a thermal bath at $T = \hbar a / 2\pi c k_{\mathrm{B}}$
\autocite{unruh1976notes}. Two guards keep the reading honest. Nothing here reclassifies the
invariant interval --- timelike and spacelike separation in the emergent metric are Lorentz
invariants, and no event crosses that classification; what is observer-relative is the computational
partition, which relations a given observer's horizon renders recordable. In the ledger of
\S\ref{sec:comp-distance} that partition \emph{is} the timelike/spacelike distinction, so in the
framework's own units timelike distance is \emph{rotated} into spacelike, and spacelike into
timelike, by a change of observer --- the tipping of the cone is literally that rotation --- and the
detector's clicks are the rotation made audible. And the effect is established physics
while the reading is the framework's, flagged structural. What gives the reading its content is a
parallel: the inertial and the accelerated observer stand to the vacuum exactly as Wigner and his
friend stand to the laboratory (\S\ref{sec:wigners-friend}) --- one holds a coherent relation and no
event, the other holds a record, and the facts are relative on both sides in precisely Rovelli's
sense \autocite{DiBiagioRovelli2021}. Since the previous paragraph identifies acceleration with
relation-creation, the parallel is not an analogy but the same structure in gravitational dress, the
horizon standing where the laboratory wall stood: \emph{the Unruh effect is Wigner's friend for
accelerated observers}. And two subsections of this section that seemed independent here meet: the
temperature Jacobson's construction requires across every local Rindler horizon
(\S\ref{sec:jacobson}) is exactly this temperature --- the $T$ of the Clausius relation is the $T$
of the tipped cone, supplied by the very relation-creation with which this subsection began.

\subsection{The relationalist commitment}
\label{sec:relationalism}

General relativity is relational: by the hole argument, spacetime points carry no
intrinsic identity, and the metric is fixed by the network of coincidences and
interactions, not by a pre-existing stage \autocite{Rovelli1996a}. The
computational picture is relational by construction and explains why it must be:
a background metric would require spatial distances to be pre-computable, and
non-computability forbids exactly that. General covariance is the statement that
distances have no pre-assigned values and depend on what information is actually
present. The metric is therefore dynamical not by stipulation but because it
measures a quantity that cannot be fixed in advance.

\subsection{Jacobson, holography, and imaginary information}
\label{sec:jacobson}

Jacobson derives the Einstein equation as an equation of state, from the Clausius
relation $\delta Q=T\,\mathrm{d}S$ across every local Rindler horizon together
with the Bekenstein proportionality of entropy to horizon area
\autocite{Jacobson1995, Bekenstein1973}. In the computational reading this is the
self-consistency condition on the loop-counting network: the Einstein equation is
the requirement that the bookkeeping of computation be internally consistent
across surfaces. The holographic principle, that the information in a region is
bounded by its bounding area \autocite{tHooft1993, Susskind1995}, then acquires a
natural reading here. A surface area is a squared distance; if spacelike distance
is measured in imaginary units of information, an area is measured in the square of
those units, that is in \emph{negative} information---entropy. The holographic
link between area and entropy thus requires the square root of negative
information, and that square root is precisely the imaginary unit of the complex
amplitude, the observer's self-referential phase. Read this way, Jacobson's
account of information crossing surfaces implicitly assigns an imaginary
information value to spacelike distance, and the $\mathrm{i}$ of the Minkowski
signature is the same $\mathrm{i}$ as the $\mathrm{i}$ of the quantum amplitude.
This identification is the most speculative step in the section: it is
structurally suggestive and it makes the holographic form unsurprising, but it is
a conjectural reading rather than a derivation, and is offered as such.

The conjecture has a two-way motivation that this paper's record should carry, because each
direction was on the table before the other was noticed. Forward: if spacelike distance is counted
in imaginary units of non-computable information, then any spacelike area is the product of two
imaginary lengths and is therefore \emph{negative} information --- which is to say entropy, by
Brillouin's identification \autocite{brillouin1956science} --- and a holographic principle arrives
automatically, the sign supplied by $i^{2}$ rather than by interpretation. Converse, and more
striking: in the holographic principle as ordinary physics already states it, surface area is
proportional to entropy, and entropy is negative information; so by implication ordinary physics
has been measuring spacelike separation in imaginary units of information all along, and the
convention of \S\ref{sec:comp-distance} is read off from established results rather than
invented. The same imaginary unit is already employed once in this section --- it is what makes
the signature of the interval Minkowski --- so the sign of the area law is its second job, not a
second assumption. What the two-way motivates it does not prove: the formal statement is exactly
the metric-and-signature half of Open Problem~\ref{op:comp-metric}, together with the
one-bit-per-bit-pair coefficient owed to the second paper. The stakes of the conjecture have
recently risen: Dorau and Much prove, with modular-theoretic rigour, that the semiclassical
Einstein equations follow from quantum relative entropy on a local horizon \emph{once the
entropy--area proportionality is assumed} \autocite{doraumuch2026relative} --- so what this
framework proposes to derive is precisely the remaining input of a derivation already carried
through at that standard.
\input{entropic_gravity_family_paragraph_2026-07-03}

\subsection{Non-locality already present in general relativity}
\label{sec:gr-nonlocality}
\input{gr_nonlocality_revision_2026-06-13}
\medskip
\noindent  The spacetime arc therefore yields, from one definition of spacelike
distance as non-computability: the signature, three spatial dimensions,
the Lorentz group, the dynamical relational metric, a reading of the
holographic principle, the presence of non-locality within general
relativity, the identification of relation-creation with acceleration
and of the Unruh effect with Wigner's friend, and the space--time
mobility asymmetry as the \textsc{fanout} asymmetry read as geometry. Of these,
the $(3+1)$ count and the Lorentz recovery are structural results that converge with independent
rigorous work and await their formal proof --- the count now numbered as
Open Problem~\ref{op:three-dimensions}; the metric-space status, the signature and the
one-bit-per-curve measure (together Open Problem~\ref{op:comp-metric}), the imaginary-information
reading of holography, and the general-relativistic non-locality, the acceleration and
Unruh readings, and the mobility asymmetry are conjectural, and are named here precisely so that the boundary between what is shown and what
is proposed is visible.

\subsection{Why the metric is not to be quantised}\label{sec:no-quantised-gravity}

A consequence has been implicit in everything above and deserves its plain statement. Causal
structure, on this framework's reading, is the bookkeeping of locked relations --- post-\textsc{fanout}
by construction --- while superposition is the signature of the pre-\textsc{fanout} regime. A
``superposition of causal structures'' is therefore not a difficult state but a category error: it
applies the pre-record operation to a record-side object, exactly as ``a superposition of
measurement outcomes'' does. Gravity lives in the post-\textsc{fanout} regime. The quantum layer
beneath it exists and is not in dispute; but it is \emph{pre-geometric} --- it carries amplitudes
over the formation of relations, not amplitudes over formed geometries --- so there is no
fundamental metric to quantise, and the standard programme of superposing spacetimes pursues an
object the ontology does not contain. One guard rail keeps the claim honest: collective
excitations of the locking pattern may perfectly well exist as effective quanta --- gravitons as
the phonons of the relational network --- and the thesis forbids fundamental metric superposition,
not effective field theory; this is Jacobson's equation-of-state moral \autocite{Jacobson1995},
made structural. The relation to gravitational-collapse programmes is an inversion: where
superposed geometries are there read as unstable \autocite{Penrose1996}, here they are read as
unformed --- what decays in those models is what locks in this one.

Three experimental postures follow, each at the status the reading carries. Laboratory
demonstrations of indefinite causal order \autocite{procopio2015experimental, rubino2017indefinite} are, on this
account, demonstrations of the nascent layer: the indefiniteness lives below the locking threshold
and definitises at it, so the framework adopts the programme and adds a boundary --- indefinite
causal order can never be amplified into records while remaining indefinite. Gravitationally
induced entanglement of the kind the Bose--Marletto--Vedral proposals seek
\autocite{bose2017spin, marlettovedral2017gravitationally} is \emph{predicted}, because the
pre-geometric layer is quantum --- while witnessing that layer, not a superposed metric: the
framework splits the experiment's usual dichotomy. And gravitationally driven anomalous collapse
is denied: locking is redundancy and thermodynamics, so decoherence experiments should find
standard rates and nothing gravity-specific. That no consistent quantisation of gravity exists was
the conclusion of the author's CASYS paper by a different route \autocite{small_2006}; the present
framework supplies the mechanism the conclusion then lacked.

\subsection{Open problems and prospects}
\label{sec:spacetime-open-problems}

The spacetime arc is the framework's frontier, and honesty about its edges
matters more here than anywhere else in the paper. We set out the genuine
problems and the prospects together, deferring the deepest to a second paper
rather than papering over them. Two of the arc's debts are well-posed now, and we number them.

\begin{openproblem}[The computational metric]\label{op:comp-metric}
Prove that the paired distance of \S\ref{sec:comp-distance} --- computable steps counted real,
non-computable bits counted imaginary, one bit per closed curve --- defines a genuine
pseudo-metric on the space of computations, with the Minkowski signature derived rather than read
off; equivalently, establish the one-bit-per-curve measure as a measure, and the signature as its
theorem. Until then the identification stands as motivated by reproducing what nature does, which
is the status flagged at its introduction.
\end{openproblem}

\begin{openproblem}[Dimension three as a theorem]\label{op:three-dimensions}
Turn the two convergent routes of \S\ref{sec:three-dimensions} into a proof: formalise the
parallelisability-plus-associativity filter (the computational route) and the mutual-acceleration
consistency demand (the general-relativistic route) so that $d = 3$ follows as a theorem, and
include the reduction clause --- that the $S^{7}$ filter acts through $SU(3)$, the stabiliser of
the selected complex direction, rather than through its full automorphism group $G_{2}$. The
M\"uller--Masanes uniqueness results \autocite{MuellerMasanes2013, MasanesMuellerEtAl2014} are
the rigorous convergents any proof must meet.
\end{openproblem}

\paragraph{The substrate and the generation of events (deferred).}
A computation runs on something; a network of events presupposes that events
occur. Both are the same objection in two forms---the network's hardware and its
clock---and the original CASYS'05 statement of the programme raised the
implementation regress and set it aside \autocite{small_2006}. The framework's
response is that a self-referential structure is the only kind that need not
presuppose a meta-substrate: it generates its own. The natural algebraic emblem
is $E_{8}$, the unique simple Lie group whose adjoint representation is also its
smallest faithful one, so that what it is and what it acts on coincide; the $240$
unit integral octonions that generate it are exactly the $E_{8}$ root system, so
the self-action is welded to the division-algebra spine of
Section~\ref{sec:entanglement-classes}. This is Wheeler's self-excited circuit
\autocite{Wheeler1990} given an algebraic body, and it converges with Furey's use
of an algebra acting on itself to generate particle content
\autocite{furey2018three}. The decisive caveat, which we state rather than hide, is
that the bootstrap is self-consistent but not self-specifying: it does not
eliminate the non-computability identified throughout, it relocates it to a
determined place. The substrate question is, on the framework's own terms,
provably without a computable answer from inside---it is the irreducible residue,
the bump in the carpet that the whole construction has been moving rather than
removing. The leap from ``$E_{8}$ is self-dual'' to ``$E_{8}$ is the substrate''
is at present an emblem awaiting operational content (an $E_{8}$-equivariant
dynamics, or a state space built as a homogeneous space of $E_{8}$), and that
content, together with how a tick is generated at all, is deferred to the second
paper.

\paragraph{Dark matter (deferred).}
If gravity is the bookkeeping of the interaction network and dark matter interacts only
gravitationally, how does it participate in a network it does not exchange information with?
This is a genuine boundary of the programme, at which it may make no prediction, and Open
Problem~\ref{op:sterile-gravity} states it with the criteria that would close it. The prospect, to be developed in the second paper, is
that this is a feature rather than a defect. In the division of labour of
Section~\ref{sec:three-dimensions}, participation in the metric---counting loops
in the $S^{1}$ and $S^{3}$ directions---is universal, carried by anything with a
worldline and a tick-rate, whereas participation in the gauge filters, the
$S^{3}$ of the weak interaction and the $S^{7}$ of colour, is selective, engaged
only by charged matter. A node that counts loops but engages no filter
\emph{is} matter that gravitates without otherwise interacting, and the
framework's Class-4 a-chiral sterile neutrino (\S\ref{sec:generations}) is the candidate that
realises it --- recorded in the particle sections as Open Problem~\ref{op:sterile-gravity}, whose
tension (a gauge-sterile channel gravitating, when gravity is measurement bookkeeping) is exactly
the question named here. The condition for the candidacy to be coherent rather than circular is
that metric participation be established as prior to, and independent of, gauge participation;
making that priority precise is the work deferred, and it is the same work the open problem
requests.

\paragraph{CP violation, the arrow of time, and baryogenesis.}
CP violation appears to demand a time asymmetry owing nothing to the entropic
arrow, which would seem to require a preferred direction built into spacetime
itself. The framework dissolves the appearance rather than accommodating it. CP
violation needs an irreducible complex phase, the CKM phase $\delta$; by Spekkens'
result that negativity and contextuality are the same notion of nonclassicality
\autocite{spekkens2007toy}, a complex phase is a signature of contextuality; so CP
violation should be contextual, sensitive to the near environment of the decaying system --- the
same phase the interaction sections read as anchored by a fourth party
(\S\ref{sec:interactions}) --- which is exactly what Andrianov, Taron and Tarrach established for
neutral kaons in matter, where the environment breaks C, irreversibility breaks
T, and the effective parameters follow a Lindblad equation
\autocite{Andrianov2001}. The framework's reading is that the resulting
informational asymmetry \emph{constitutes} the arrow rather than presupposing it:
CP violation, entropy increase, and the direction of time are one informational
asymmetry seen three ways, not an independent second arrow, so the worry about a
direction built into spacetime does not arise. On this footing the contextual
dependence sets up a positive feedback loop for baryogenesis---an initial
fluctuation makes the environment slightly matter-dominated, contextual
enhancement raises the effective CP violation, more matter is produced, and the
loop runs---with $\delta$ a cosmological fixed point rather than a fundamental
constant. The honest limit of this is that a feedback loop explains amplification
and lock-in, not the sign selection or the baryon-number violation that
Sakharov's conditions still require \autocite{Sakharov1967}; the stronger claim
the framework is positioned to make, that the matter direction and the time
direction are co-selected by the same self-referential asymmetry, needs a source
of baryon-number violation and is left open.

\paragraph{Spatial infinitude and the status of inflation.}
There is no upper bound on non-computability, so the space of computations is
unbounded: the universe is spatially infinite, which leaves its spatial sections flat or of
negative curvature and excludes only the closed case. Infinitude on its own, however, does not
settle flatness, which is a fine-tuning of the Friedmann dynamics the framework has not yet
derived (Open Problem~\ref{op:comp-metric}); and flatness was in any case never inflation's
principal task. The horizon problem is, and an infinite space does nothing for it. The framework's actual lever against
inflation is the relational thesis of Section~\ref{sec:gr-nonlocality}: if
pre-spatial entanglement correlates regions before \textsc{fanout} spreads them
into spatial separation, then the observed large-scale homogeneity is fixed
relationally, before any question of causal contact, and the horizon problem
dissolves without an inflaton. This is a real argument, but it carries a real
bill: inflation also generates the nearly scale-invariant spectrum of primordial
fluctuations seen in the microwave background \autocite{Planck2020}, and any
account that dispenses with inflation owes an alternative origin for that
spectrum. That, not the horizon problem, is the hard empirical debt, and it is
named here rather than discharged.

\paragraph{Expansion and acceleration.}
If spacelike distance is the information separating worldlines, then as worldlines
accumulate distinct interaction histories they accumulate information-distance,
and the universe expands. The picture has the right exclusions built in: bound
systems interact continuously, share information, and so keep low
information-distance---they do not expand---while only worldlines that decouple
diverge, recovering expansion on cosmological scales alone. The feedback in which
greater separation means less light-cone overlap and so faster divergence,
$\mathrm{d}d/\mathrm{d}t \propto d$, yields exponential divergence, a de~Sitter
phase, and hence accelerating expansion without a separate dark-energy term
\autocite{Riess1998}. The limit to be honest about is that pure exponential
divergence is the late-time behaviour only; reproducing the full expansion history,
with its radiation- and matter-dominated eras, and fixing the Hubble rate
quantitatively, remains to be done.

\paragraph{The $S^{7}$ sector as internal structure.}
That $S^{7}$ cannot carry a metric does not make it disappear. Its closed
timelike curves carry the same non-computable information structure as those of
$S^{1}$ and $S^{3}$, but lacking associativity they support algebra---gauge
charge---in place of geometry---distance. This is why gauge symmetry has so often
looked like a set of compactified extra dimensions: the Kaluza--Klein and string
intuition is tracking something real in the $S^{7}$ structure, but misdiagnoses
it as spatial. In the framework these degrees of freedom have no size because
they have no metric; they are internal in the exact sense that they enter the
information-exchange that constitutes gauge charge while contributing nothing to
the loop-counting that constitutes distance. The resemblance to rolled-up
dimensions is therefore genuine and instructive, and the diagnosis that they are
algebraic rather than geometric is what separates this account from theories that
add real spatial dimensions to carry the same structure.

\paragraph{Loops, supertasks, and a second route to parallelisability.}
The argument of \S\ref{sec:three-dimensions} derives the parallelisable-sphere constraint from
closed timelike curves: a computation whose output is its own input requires information to be
transported consistently around a loop, and consistent transport selects the parallelisable spheres
$S^{1}$, $S^{3}$, $S^{7}$ \autocite{adams1960hopf}. There is a suggestion, which we record as a
prospect rather than a result, that the same constraint may govern the loops of perturbative quantum
field theory. The external legs of a Feynman diagram are on-shell classical boundary data; the
quantum content, and the divergences, enter with the internal loops, where the loop momentum is
integrated over all values and the action accumulated inside the loop is in this sense unbounded ---
a supertask, an infinite computation closed into a finite process. If such a loop is, in the
framework's terms, a closed computational path of the same kind as a closed timelike curve --- a
cycle carrying information not derivable from the data outside it --- then the parallelisability
constraint derived for closed timelike curves would apply to it, and the admissible interactions
would be constrained by the same $S^{1}$, $S^{3}$, $S^{7}$ structure that fixes the gauge group in
\S\ref{sec:entanglement-classes}. This would be a second and independent route to the Standard Model
content: not only \enquote{closed paths in spacetime select the parallelisable spheres} but \enquote{closed paths in interactions select them too}, unifying the spacetime arc and the interaction arc under one
principle. The decisive step, which we have not taken, is the identification of a field-theory loop
with a closed computational path in the precise sense the framework requires --- a loop integral is a
closed cycle in momentum space, a closed timelike curve a closed cycle in spacetime, and whether
these are the same kind of self-consistent closed structure is exactly what would have to be
established. It is offered as a direction, not a claim.

\paragraph{Does parallelisability determine renormalisability?}
A sharper and more self-contained question sits alongside the previous one, and we state it in the
hope that it will be taken up by those equipped for it. Renormalisability is the property that the
divergences of a quantum field theory can be consistently absorbed, leaving finite predictions; some
theories have it and some do not, and the standard criterion for which is technical (power-counting
of the coupling's dimension). The framework suggests a more structural reading. Parallelisable
spheres are precisely the spheres on which a global, consistent transport exists where it generically
does not; the non-parallelisable spheres are where the obstruction lives. If the divergence of a loop
is, in the framework's language, an obstruction to closing a computational path consistently, then
parallelisability would be exactly the structure that distinguishes the removable obstruction (a
renormalisable theory) from the irremovable one. It is at least a suggestive coincidence that the
renormalisable interactions of nature are the gauge theories built on $U(1)$, $SU(2)$ and $SU(3)$ ---
the groups the framework ties to the three parallelisable spheres --- and that gravity, which resists
renormalisation, is the sector \S\ref{sec:three-dimensions} ties to the geometry rather than the
algebra. Whether this is structural or accidental we do not know. The question --- whether
parallelisability, in the Hurwitz--Adams sense, is what determines renormalisability --- is posed in
the native language of the division-algebra approach to the Standard Model, and is one we think that
community is best placed to answer. We are careful to claim nothing more than that the question is
well-motivated: this is emphatically not a method for computing renormalised quantities, whose finite
remainders carry the precision agreement with experiment, but a conjecture about which theories admit
such a computation at all. It is, like the previous prospect, a direction with a possible path to an
answer rather than an answer, and it is left open deliberately.


\medskip
\noindent These mark the boundary of the programme as it stands. The substrate,
the generation of events, and the network status of dark matter are deferred to
the second paper; the contextual account of CP violation and baryogenesis, the
relational case against inflation, the information-divergence account of
expansion, and the algebraic reading of the $S^{7}$ sector are offered as
prospects, each with its outstanding debt named so that the line between argument
and aspiration stays visible.       <-- first consumer
%                                                          (S^7 exclusion)
%      % EDITED 2026-09-09 scripts/patch_spekkens_citekey_2026-09-09T1810.py
% =====================================================================
%  SECTION: Why Hopf Fibrations Belong in Particle Physics
%  (sec:hopf-justification). Drafted 2026-06-08; F1 arc added
%  prop:f1-dictionary and op:hopf-frobenius's discharge clause.
%  AUDITED 2026-07-06 (final pass, eighth stop).
%
%  PLACEMENT (reordered chain, JS 2026-07-06): position 8 --
%      \input{spacetime_from_non-computability}
%      \input{hopf_justification_section_2026-06-08}
%      \input{entanglement_classes_and_gauge_groups_2026-06-08}
%  No adjacency phrasing in the prose: position-portable. All
%  cross-refs live -- sec:fanout-direction DRAFTED 2026-07-06 (formerly the one undrafted
%  section, by design).
%
%  AUDIT RECORD: prop:gsm, prop:termination and prop:f1-dictionary
%  verified against every quoting consumer (statements canonical);
%  honest-debt subsection names both the selection assumption and the
%  mereology debt; the epistemic two-climbs disambiguation and the
%  compositional rung-evidence verified; sec:canonical-orbits live;
%  file pure ASCII. Fixes: op:hopf-frobenius's cross-link with
%  op:sigma-functor made TWO-WAY (the "four ways" line now names the
%  op); both remaining functor conditions ANCHORED --- the quaternionic
%  to the established fibre physics of sec:wf-reconciliation (residue:
%  clutching map vs strong complementarity), the octonionic to the
%  proved canonical-orbit images (umbrella label app:ghz-w-proofs
%  minted at the appendix for the purpose). Statuses unchanged.
%  British English; \autocite; no bullets in prose.
% =====================================================================

\section{Why Hopf Fibrations Belong in Particle Physics}
\label{sec:hopf-justification}

Invoking Hopf fibrations in a discussion of the Standard Model invites an immediate and fair
objection: that the move is decorative, an analogy between two pieces of attractive mathematics
with no load-bearing connection between them. The purpose of this section is to dispose of that
objection before the construction is put to work. The connection is not an analogy. It is
carried by two independent and independently published theorems, and the structures it relies
upon are not selected for convenience but compelled by a pair of classification results. We take
as already established the correspondence between qubit entanglement and the Hopf fibrations
themselves --- the subject of a mature literature
\autocite{mosseri2001hopf,bernevigchen2003geometry} --- and do not reproduce it here. What this
section addresses is the prior question of why that correspondence is entitled to carry physical
content at all.

\subsection{Two theorems, joined end to end}

The bridge from qubit entanglement to the Standard Model is built from two distinct results,
neither of which is original to the present framework and neither of which depends on the other.

The first is geometric. The unit sphere in the Hilbert space of $n$ qubits is the total space of
a Hopf fibration: $S^3 \to S^2$ for one qubit, $S^7 \to S^4$ for two, and $S^{15} \to S^8$ for
three. For two and three qubits the fibration is \emph{entanglement-sensitive} --- separability
is precisely the condition under which the base degenerates to a product of single-qubit Bloch
spheres and the additional dimensions associated with the higher imaginary directions collapse
\autocite{mosseri2001hopf,bernevigchen2003geometry}. This is a statement about the geometry of
state spaces and requires no physical interpretation whatever.

The second is algebraic. Singling out a preferred complex structure on the octonions ---
equivalently, fixing the split $\mathbb{O} \cong \mathbb{C} \oplus \mathbb{C}^3$ --- leaves a
residual symmetry, and that residual symmetry is exactly the gauge group of the Standard Model.
This was established through the exceptional Jordan algebra by Dubois-Violette and Todorov
\autocite{duboisviolette2016exceptional,duboisviolettetodorov2019exceptional,todorovduboisviolette2018deducing}
and given its sharpest spinorial form by Krasnov \autocite{krasnov2021so9}, who showed that the
relevant group is the stabiliser of a complex structure inside $\mathrm{Spin}(9)$. We record the
result in the form in which it enters the present work.

\begin{proposition}[Dubois-Violette--Todorov; Krasnov]\label{prop:gsm} Let $J_R$ be the complex
structure on the space $\mathbb{O}^2$ of $\mathrm{Spin}(9)$ spinors determined by a choice of
unit imaginary octonion, equivalently the split
$\mathbb{O} \cong \mathbb{C} \oplus \mathbb{C}^3$. The subgroup of $\mathrm{Spin}(9)$ commuting
with $J_R$ is \[ G_{\mathrm{SM}} \;=\; \frac{SU(3) \times SU(2) \times U(1)}{\mathbb{Z}_6}, \]
acting on the singled-out $\mathbb{C}$ summand as a colour singlet of hypercharge $-1$ and on
the $\mathbb{C}^3$ summand as a colour triplet of hypercharge $\tfrac{1}{3}$, both as $SU(2)$
doublets --- the quantum numbers of one left-handed generation.
\end{proposition}

The two results are independent, and the join between them is the only place a physical theory
must do work. Each result is a theorem; what neither supplies is a reason why nature should fix
one complex direction within $\mathbb{O}$ rather than treating all directions symmetrically. We
return to this --- the genuine load-bearing assumption --- in \S\ref{sec:hopf-honest-debt}, and
resolve it, in candidate form, in \S\ref{sec:fanout-direction}.

\subsection{The ladder is forced, not chosen}

It might still be objected that selecting the octonions, or the Hopf fibrations, is itself a
modelling choice carrying as much latitude as the choice of a grand-unified group. It is not,
and the reason is a pair of classification theorems that close off the alternatives entirely.

By Hurwitz's theorem the only normed division algebras are $\mathbb{R}$, $\mathbb{C}$,
$\mathbb{H}$ and $\mathbb{O}$ \autocite{hurwitz1898composition,baez2002octonions}. By Adams's
resolution of the Hopf-invariant-one problem --- together with the parallelisability results of
Bott--Milnor and Kervaire --- the only parallelisable spheres are $S^0$, $S^1$, $S^3$ and $S^7$,
so that there exist exactly three nontrivial Hopf fibrations, in bijection with $\mathbb{C}$,
$\mathbb{H}$ and $\mathbb{O}$
\autocite{adams1960hopf,bottmilnor1958parallelizability,kervaire1958nonparallelizability}. These
are not heuristics; they are among the more durable theorems in topology and algebra. Their
force here is that, once a framework is committed to amplitudes valued in a normed division
algebra --- a commitment that in the present construction is made far upstream, by the
self-referential argument of \S\ref{sec:self-reference} and \S\ref{sec:division-algebras}, and
not by fiat at this stage --- the ladder of Hopf fibrations is not selected but compelled, and
it terminates of its own accord.

\begin{proposition}[Termination]\label{prop:termination} The Cayley--Dickson tower of normed
division algebras and the tower of parallelisable spheres both terminate at real dimension
eight. Consequently there are exactly three nontrivial Hopf fibrations, the largest being the
octonionic $S^{15} \to S^8$, and the construction admits no extension to a fourth case.
\end{proposition}

The termination is worth dwelling on, because it is unusual for a model of this kind to predict
its own boundary. The next rung of the Cayley--Dickson sequence, the sedenions, possesses zero
divisors; the Hopf map can no longer be globally defined; there is no fourth fibration. This is
not a limitation imposed from outside but a structural fact, and it carries a physical
consequence that we develop later: there is no room in the construction for a fourth gauge force
or a fifth interaction, so that the discovery of either would count as evidence against the
framework rather than as an invitation to extend it. The endpoint is a theorem, not a
stipulation.

\subsection{Independent arrival as corroboration}

The most persuasive evidence that the destination is not an artefact of the route is that it is
reached by routes sharing none of the same machinery. Szangolies
\autocite{szangolies2025standardmodel} reaches $G_{\mathrm{SM}}$ from the emergent-gravity and
Kaluza--Klein programmes, reading the singled-out complex direction as a dimensional reduction
of a $9+1$-dimensional spacetime; the present framework reaches it from self-reference
(\S\ref{sec:self-reference}); and Furey and Gillard--Gresnigt reach the same gauge group, and
the same three-generation structure, through the minimal left ideals of the Dixon and Clifford
algebras, without invoking Hopf fibrations at all
\autocite{furey2018three,gillardgresnigt2019three}; the relation between their minimal ideals
and the entanglement classes used here is taken up in \S\ref{sec:canonical-orbits}. The
octonionic route to particle physics is in any case an old one, going back to the use of
$\mathbb{O}$ to model the strong force \autocite{gunaydingursey1974quark}.

That programmes with disjoint starting points converge on
$SU(3) \times SU(2) \times U(1)/\mathbb{Z}_6$ and on the same quark--lepton division is the
coincidence Szangolies himself remarks upon: the characteristic structure of the Standard Model,
routinely regarded as arbitrary, turns up independently in unrelated settings, and its appearing
twice would then be doubly odd unless the two settings are in fact related. We read the
convergence the same way --- as evidence that the division-algebraic description is singled out
by the mathematics rather than imposed by the modeller, and therefore that
Proposition~\ref{prop:gsm} is structural rather than contrived.

\subsection{The honest debt}
\label{sec:hopf-honest-debt}

It would be misleading to leave the impression that this route is free of assumptions, so we
state plainly the one assumption that all of these programmes share. Proposition~\ref{prop:gsm}
requires that \emph{some} complex direction in $\mathbb{O}$ be singled out, and the uniqueness
of that choice is not supplied by the theorem. Krasnov parametrises the complex structure $J_R$
by an arbitrary unit imaginary octonion, and the residual group is the same for every such
choice because $G_2 = \mathrm{Aut}(\mathbb{O})$ acts transitively on the unit imaginary
octonions. The theorem therefore guarantees the gauge group \emph{given} that a direction is
fixed, while remaining silent on \emph{why} one is fixed. Szangolies treats this as a free
choice and flags it, in his own conclusions, as an open element of the construction --- one to
which several inequivalent answers are available.

This, and not the appeal to Hopf fibrations, is the genuine load-bearing assumption of the whole
approach. The Hopf fibrations are forced (Proposition~\ref{prop:termination}); the singling-out
is not. The present framework offers a candidate resolution that the prior literature lacks: the
act of measurement, modelled as \textsc{fanout}, selects the direction physically, and the
transitivity of $G_2$ then guarantees that the gauge group is independent of which direction the
measurement happens to select. This argument is given in \S\ref{sec:fanout-direction}; we flag
here, at the appropriate epistemic status, that it is structural rather than fully proved, and
that the residual question --- why a measurement selects \emph{exactly one} imaginary direction
rather than merely \emph{permitting} a selection --- remains open. It is now priced as Open
Problem~\ref{op:fanout-selection}.

A second debt is shared with every programme of this character and is not closed by anything
said here: why the abstract quantum state should admit a tensor-product decomposition into
qubits in the first place. This is the quantum factorisation, or mereology, problem. The
entanglement-sensitivity on which the construction turns is meaningful only once such a
decomposition is fixed, and we treat the decomposition as a necessary input rather than a
derived result. Naming it is preferable to letting it pass unremarked.

\subsection{The termination read epistemically}
\label{sec:hopf-information-bound}

The termination of the tower (Proposition~\ref{prop:termination}) has so far been stated
algebraically: the Cayley--Dickson sequence acquires zero divisors at the sedenions and the Hopf
map can no longer be globally defined. There is a second reading of the same boundary, epistemic
rather than algebraic, which we offer as a motivated conjecture rather than a settled result,
because it connects the tower's endpoint to a bound that arises independently in the foundations
of measurement.

A particle, in this framework, is a reified relation between events, and its internal content is
carried by the octonionic $A,B,C$ register --- the $S^{7}$ sector serving as internal degrees of
freedom in the channel that connects the events (\S\ref{sec:space-of-computations}). The
termination of the tower is then a statement about that content: a single such relation can
support genuine entanglement among at most three parties, because there is no fourth fibration
to carry a fourth independent party within the same irreducible relation. Every particle is, on
this reading, a \emph{completed} octonionic-level relation --- topped out at three parties not
by circumstance but because the tower admits no further rung. The information that can be
encoded in one relation-between-events, and hence extracted by one measurement of it, is
correspondingly bounded: it is the information of a completed three-party register, and no
measurement of a single particle can yield more.

This is exactly the kind of bound that the epistemic-restriction programmes reach from the other
direction. Spekkens' toy theory is built on a knowledge-balance principle that caps the
information obtainable about a system \autocite{spekkens2007toy}, and Szangolies derives such an
epistemic horizon not by positing it but from the same self-referential non-computability on
which the present tower rests \autocite{szangolies_j_2018}. That the two accounts --- the
algebraic termination of the Hopf tower and the epistemic bound on measurement information ---
should coincide is, we suggest, not a numerical accident but a shared root: both are
consequences of the self-referential non-computability of \S\ref{sec:self-reference}, wearing an
algebraic face in the one case and an epistemic face in the other. The Hurwitz--Adams ceiling
and the knowledge-balance ceiling would then be one limit seen twice, and Szangolies' programme,
which reaches $G_{\mathrm{SM}}$ by an independent route (\S\ref{sec:hopf-honest-debt}) and the
epistemic horizon from the same non-computable source, shares with this framework not merely a
conclusion but its origin.

The reading must be kept distinct from a different climb, with which it would otherwise be
conflated. The Wigner's-friend regress of \S\ref{sec:wigners-friend} ascends without end: each
observer may be measured by a further observer, each carrying a fresh gauge-private circle, and
that chain has no top --- it is the bump in the carpet, unbounded in length. The termination
bounds something else: not the length of the chain of observers, but the internal structure of
each single relation, each particle, that the chain is built from. The regress is unbounded
precisely because each of its links is a bounded, completed three-party object;
termination-of-each-link and non-termination-of-the-chain are not in tension but complementary,
the former individuating the links of the latter. What the bound caps is the information in one
particle; what the regress leaves uncapped is the number of particles and observers that can be
chained. The precise identification of ``the information of a completed three-party register''
with the formal measures of the epistemic-restriction literature is the technical step that
would turn this conjecture into a theorem, and we flag it as owed; the structural claim --- that
the tower's endpoint and the measurement bound share the non-computable root --- is what we
advance here.

\subsection{The termination read compositionally}
\label{sec:hopf-compositional}

There is a third termination theorem, proved in a community that has so far had no contact with
the two joined above. Coecke and Kissinger showed that the GHZ and W classes --- realised as the
commutative special and the anti-special Frobenius algebras of categorical quantum mechanics ---
are \emph{compositionally complete}: every $N$-qubit entanglement class under SLOCC is generated
by composing these tripartite primitives with bipartite entanglement and product states
\autocite{coeckekissinger2010compositional}. Above three parties nothing new appears; the
entanglement alphabet closes at exactly the point where the fibration tower terminates. The
mutual silence of the two literatures is itself striking: Coecke and Kissinger make no reference
to Hopf fibrations or division algebras, and the fibration literature on entanglement
\autocite{mosseri2001hopf,bernevigchen2003geometry} nowhere cites compositional completeness.
Two classifications, in different mathematics, draw the same line in the same place, unaware of
one another. That silence has, in the months this paper was in draft, just begun to break from
the categorical side, with a projective, base-level reading of the same calculus
\autocite{burtonanwar2026meromorphic}; we engage it below.

The framework's conjecture is that this is one termination wearing a third face: the alphabet
has the letters it has --- Bell, GHZ, W --- because the parallelisable spheres have the
dimensions they have. There is partial evidence at every rung. At one qubit, the commutative
special Frobenius algebras on $\mathbb{C}^{2}$ are classified by orthonormal bases --- points of
$S^{2}$, the base of the first fibration --- so that the copied basis plays the role of a
section of the bundle and the spider theorem expresses its idempotence. At two qubits, the
$SU(2)$ fibre of $S^{7} \to S^{4}$ carries exactly the entanglement that survives local
operations, and L\'{e}vay's result that the bundle's natural connection is the $SU(2)$
Yang--Mills instanton \autocite{levay2004geometry} stands on the geometric side where the
interaction of two Frobenius structures stands on the categorical one. At three qubits, the
split of D\"{u}r, Vidal and Cirac \autocite{durvidalcirac2000three} mirrors the algebra exactly:
GHZ is genuinely tripartite, demanding the full non-associative octonion, while W is
pairwise-robust and housed in the associative subalgebra $\mathbb{H} \subset \mathbb{O}$ --- the
special/anti-special Frobenius distinction in geometric dress.

The first of the four conditions the correspondence requires can be discharged at once, and we
do so here, because a problem whose first step has been walked is a different kind of invitation
from one merely posed.

\begin{proposition}[The one-qubit dictionary]\label{prop:f1-dictionary} The $\dagger$-special
commutative Frobenius algebras on $\mathbb{C}^{2}$ correspond bijectively to
\emph{lifted antipodal skeletons} of the complex Hopf fibration $\pi \colon S^{3} \to S^{2}$:
unordered antipodal pairs $P \subset S^{2}$ together with a section of $\pi$ over $P$. The
copyable states of a structure are exactly the section's image, so the lifts are recorded
faithfully --- the projectively trivial centre of $U(2)$ acts non-trivially on structures ---
and the correspondence is $U(2)$-equivariant, the ordered structures forming a single free and
transitive $U(2)$-orbit. The channel a structure induces on density matrices is dephasing in its
skeleton and is independent of the section: the fibre freedom, a torsor for
$\mathrm{Map}(P, U(1))$, is exactly the data absent from the classical record a structure
induces.
\end{proposition}

\begin{proof}
By the classification of Coecke, Pavlovic and Vicary, such structures correspond to orthonormal
bases, recovered as the sets of copyable states \autocite{coeckepavlovicvicary2013orthogonal}.
An orthonormal basis of $\mathbb{C}^{2}$ is two orthogonal rays --- an antipodal pair on the
Bloch sphere --- together with a unit representative of each, which is precisely a section of
$\pi$ over the pair; and $\delta(e^{i\theta}x) = e^{i\theta}\, x \otimes x$ against
$(e^{i\theta}x) \otimes (e^{i\theta}x) = e^{2i\theta}\, x \otimes x$ pins the phases.
Equivariance, the action of the centre, the orbit structure and the induced channel are direct
computations; an executable verification of every finite check accompanies the code release.
\end{proof}

Two remarks fix the proposition's scope. It corrects a natural first reading: a classical
structure does not project to the base but selects lifts into the total space --- classicality
is section data, and what the copy operation broadcasts is a lift, not a ray. And it claims
nothing topological about cloning: the finiteness of the copyable set is linear, structural
no-cloning residing in non-naturality \autocite{abramskyheunen2012hstar}, so the first place the
bundle's global topology can bite is the quaternionic rung, where the degree-one clutching map
stands ready to be compared with the strong-complementarity law governing an interacting pair of
structures. The dictionary is, finally, the complement of the projective treatment cited above:
Burton and Anwar pass functorially to the base and deliberately discard the fibre
\autocite{burtonanwar2026meromorphic}, and the classical structures turn out to live in exactly
the data their projectivisation forgets --- the two readings compose, and the resulting triangle
is a consistency condition available to constrain the remaining conditions of the
correspondence.

\begin{openproblem}[The Hopf-fibration--Frobenius correspondence]\label{op:hopf-frobenius} Let
$\mathbf{Hopf}$ denote the category of the three fibrations
$S^{2^{k}-1} \hookrightarrow S^{2^{k+1}-1} \to S^{2^{k}}$, $k = 1, 2, 3$, with bundle maps as
morphisms, and let $\mathbf{Frob}$ denote the Frobenius algebras on $\mathbb{C}^{2}$ in
$\mathbf{FdHilb}$, with the commutative-special (GHZ) and anti-special (W) classes
distinguished. Construct a functor $F \colon \mathbf{Hopf} \to \mathbf{Frob}$ --- or prove that
none exists --- under which the complex fibration yields the single-qubit classical structures
--- Proposition~\ref{prop:f1-dictionary} identifies these as the lifted antipodal skeletons,
section data over two base points rather than the points alone; the quaternionic fibration
yields the interaction of two such structures, with the $SU(2)$ fibre carrying the entanglement
that survives the classical quotient; the octonionic fibration yields the GHZ/W pair, GHZ on the
full non-associative fibre and W on the associative $\mathbb{H} \subset \mathbb{O}$; and the
composition by which Coecke--Kissinger generate the $N$-qubit classes corresponds to a
well-defined geometric operation on the bundles. If such an $F$ exists, compositional
completeness follows from Adams's theorem as a topological corollary --- there is no fourth
fibration to supply a fourth primitive --- and if $F$ is faithful enough that the implication
also runs back, the termination of the tower and the completeness of the alphabet are equivalent
expressions of a single classification. The first condition is discharged by
Proposition~\ref{prop:f1-dictionary}; what remains open is the quaternionic and octonionic
conditions and, above all, the correspondence of compositions. Both open conditions are already
constrained: the quaternionic fibre's physical content --- the recoverable relative phase and
its purity ledger --- is established (\S\ref{sec:wf-reconciliation}), leaving the
clutching-map-to-strong-complementarity comparison as the structural residue; and any candidate
functor must reproduce the proved canonical-orbit images, GHZ to the colour triplet and W to the
singlet (\S\ref{sec:canonical-orbits}, Appendix~\ref{app:ghz-w-proofs}). A proof that no such
functor exists, or a correspondence at the level of the classifications without functoriality on
compositions, would each also resolve the problem: the first would make the agreement of two
independent classifications a fact demanding its own explanation; the second would establish the
weaker reading.
\end{openproblem}

A resolution in the affirmative would do three things. It would explain compositional
completeness rather than merely prove it --- a corollary of the division-algebra structure of
quantum mechanics rather than an independent algebraic fact. It would give the diagrammatic
language of this paper's own interaction circuits a topological warrant: the GHZ/W spider
calculus would be the unique language the topology permits, not a convenient notation. And it
would open a two-way bridge, translating problems on compositions of Frobenius structures into
operations on bundles and back. The problem sits at the intersection of categorical quantum
mechanics, the quantum-information geometry of the fibrations, and the division-algebra
approaches to particle physics, and it is posed here so that any of the three communities can
take it up in its own language. It is perhaps worth noticing that both classifications already
answer to the same name --- the fibrations are Heinz Hopf's, and so, through the bialgebra
structure of the diagrammatic calculus, are the Hopf algebras --- a coincidence of eponymy that
the conjecture proposes to promote to a theorem. With the epistemic reading above, an
affirmative resolution would leave the tower's endpoint as one limit seen four ways: algebraic,
geometric, compositional, epistemic --- the last being the fixed-point classification of Open
Problem~\ref{op:sigma-functor}, with which this problem is cross-linked, each constraining the
other.

\subsection{What this sets up}

The entanglement-sensitivity of the Hopf maps is not a side feature but the hinge of the
physical reading that follows. Because separability is exactly the collapse of the base to a
product of Bloch spheres, the structure of the entanglement --- the SLOCC class of a multi-qubit
state, and within it the canonical maximally entangled orbit --- is carried by precisely the
degrees of freedom on which $G_{\mathrm{SM}}$ acts in Proposition~\ref{prop:gsm}. This is what
makes it possible, in \S\ref{sec:entanglement-classes}, to read the two genuinely tripartite
entanglement classes onto the $\mathbb{C} \oplus \mathbb{C}^3$ split through their canonical
representatives: the GHZ orbit confined to the colour-triplet $\mathbb{C}^3$ sector, the W orbit
reaching the $\mathbb{C}$ singlet, recovering the quark--lepton division not as an assignment
but as a consequence of the entanglement structure. The scope of those statements --- canonical
orbits rather than whole classes, and the identification of the canonical orbits with the
physical particles --- is made precise there. The Hopf fibrations are the bridge; the
entanglement classes, treated next, are the traffic that crosses it.

%  The section is position-portable: all adjacency phrasing was
%  retargeted to explicit \S\ref on 2026-07-06. The S^7 hand-off
%  sentence to sec:space-of-computations WRITTEN 2026-07-06 (spacetime
%  file read; the closing now hands first to its immediate consumer,
%  then to hopf-justification). Upload verifications COMPLETE: L2 and
%  quantum-postulates do not consume the ladder; WF's references are
%  forward-voiced or self-carrying.
%
%  DIVISION OF LABOUR (deliberate, to avoid overlap):
%   - sec:self-reference (undrafted): self-reference -> complex
%     amplitudes and the Born rule; why not a real base. This section
%     TAKES that as given and starts at C.
%   - THIS section: the algebra -- the doubling construction, the
%     ledger of losses-as-physics, the internal-algebras-not-a-new-QM
%     scoping (Adler pre-empt), the two appearances of the ladder
%     (register / measurement climb, stated lightly), and Hurwitz as
%     the algebraic termination with the committed zero-divisor
%     reading.
%   - sec:hopf-justification: the geometry (Adams, parallelisability,
%     the fibrations), the join of the two classifications, the
%     entanglement-sensitivity literature, the C-direction honest debt,
%     and the epistemic reading of termination.
%
%  COMMITTED READINGS USED (sources re-read 2026-07-03):
%   - Loss ledger and level dictionary: cayley_dickson_qubits_for_review
%     sec.5 ("each step loses a symmetry, and each lost symmetry is
%     identified with a physical content").
%   - Termination reading: grounding_argument_gauge_group ("the
%     sedenions have zero divisors and cannot support consistent
%     information flow around loops") -- leads; norm-composition-as-
%     probability given as its amplitude-level restatement.
%   - Internal-not-base scoping: cayley_dickson_qubits_for_review Q2 and
%     sec.3.2 caveat (Adler objection; C-qubit treated geometrically).
%
%  HONEST DEBT: ladder proved closed at the top (Hurwitz), derived at
%  the floor. The rung-compulsion revisit-trigger FIRED 2026-07-06:
%  sec:self-reference and op:sigma-functor now exist, so the thesis is
%  minted as op:ladder-compulsion (with the CD-conjugation-as-sigma
%  candidate route), inheriting op:sigma-functor's status.
%
%  AUDITED 2026-07-06 (final pass, third stop): all four citations
%  used; all cross-refs live; CD facts, dims, alternativity, Hurwitz
%  statement verified; Adler subsection delivers everything the
%  register's compressed pointer promises; file pure ASCII; no bare
%  "frame". Fixes: opening rewritten for both predecessors (continuity
%  was inserted between after drafting); "associator's non-vanishing is
%  mass itself" recalibrated to the register's debt-against-selection
%  defusal; the O->S ledger row and Hurwitz subsection reconciled with
%  the fourth-qubit forward remark ("the terminus closes the ladder of
%  forces; what survives past it is, at most, flavour").
%
%  CROSS-REFS: sec:self-reference [sibling dangle -- resolves when that
%  section is drafted], sec:hopf-justification, sec:hopf-information-bound,
%  sec:register, sec:entanglement-classes, sec:generations, sec:masses,
%  sec:wigners-friend. All others exist.
%  CITATIONS: existing keys baez2002octonions, hurwitz1898composition,
%  furey2018three. ONE new key: adler1995quaternionic (RIS in
%  division_algebras_refs_2026-07-03.ris).
%  British English; \autocite; no bullets in prose; the ledger is a
%  table because the comparison is the argument.
% =====================================================================

\section{The Division-Algebra Ladder}
\label{sec:division-algebras}

The foundations behind us end at amplitudes twice over: natively, where the self-referential
argument forced the description of a system into the system's own state space and priced
consistency at a number system whose parts can cancel --- complex amplitudes, and the Born rule
that squares them (\S\ref{sec:self-reference}) --- and independently, where the regime's
continuity delivers the same field through Hardy's theorem (\S\ref{sec:continuity}), the
convergence of the two branches being evidence rather than accident. This section asks the question
both routes leave open. Why
should the number system stop at $\mathbb{C}$? Ordinary quantum mechanics stops there by postulate;
a framework that has forbidden itself postulates must answer, and the answer is that the climb
neither stops at $\mathbb{C}$ nor continues forever. There are exactly two more rungs, each
purchased by surrendering an algebraic law and each paying for itself in physics, and the ladder is
closed at the top by a theorem of 1898. This section sets out the algebra of that ladder: the single
construction that climbs it, the ledger of what each rung costs and buys, and the theorem that ends
it. The geometry of the same ladder --- that each rung acts on a sphere, that exactly those spheres
are parallelisable, and that the algebraic and geometric classifications are one --- is the business
of the geometry section (\S\ref{sec:hopf-justification}).

\subsection{One construction climbs the whole ladder}

The rungs are not a miscellany; a single construction generates them all. Given an algebra with a
conjugation, the Cayley--Dickson construction builds its double on ordered pairs, with a twisted
multiplication in which the conjugate enters and a conjugation $(a,b)^{*} = (a^{*}, -b)$
\autocite{baez2002octonions}. (Conventions differ in the placement of the conjugations; all yield
isomorphic algebras.) Applied repeatedly it produces the sequence
$\mathbb{R} \to \mathbb{C} \to \mathbb{H} \to \mathbb{O} \to \mathbb{S}$, of real dimensions
$1, 2, 4, 8, 16$: the reals, the complex numbers, the quaternions, the octonions, the sedenions.
The ground floor is classical: $\mathbb{R}$ is ordered and positive-definite, the arithmetic of
classical probability. The first doubling is therefore the quantum move itself --- the passage from
probabilities to amplitudes that can interfere --- and it is not made here by choice: it is the
derivation of \S\ref{sec:self-reference}, already behind us. What this section must account for is
the rest of the climb.

\subsection{The ledger}

In mathematics the doublings are usually told as a story of degeneration: each algebra is worse
behaved than the last, and the sequence is studied for where it breaks. This framework reads the
same ledger in the opposite direction. Each rung is bought by surrendering a law, and each
surrendered law is exactly a piece of physics --- the losses are the assets, and the comparison is
sharp enough to be the argument:

\begin{center}
\begin{tabular}{llll}
\hline
doubling & dimension & what is lost & what the loss becomes \\
\hline
$\mathbb{R} \to \mathbb{C}$ & $1 \to 2$ & the order of the reals & phase and interference; amplitudes \\
 & & & that cancel; the Born rule (\S\ref{sec:self-reference}) \\
$\mathbb{C} \to \mathbb{H}$ & $2 \to 4$ & commutativity & non-commuting actions: weak isospin \\
 & & & $SU(2)$, the $B$ register (\S\ref{sec:register}) \\
$\mathbb{H} \to \mathbb{O}$ & $4 \to 8$ & associativity & the associator: colour, chirality, the \\
 & & & generations, and mass as debt \\
 & & & (\S\ref{sec:generations}, \S\ref{sec:masses}) \\
$\mathbb{O} \to \mathbb{S}$ & $8 \to 16$ & division itself & no further force: zero divisors appear \\
 & & (zero divisors) & and the ladder of sectors ends (below) \\
\hline
\end{tabular}
\end{center}

Walking the rows: the order lost at the first doubling is the possibility of interference --- parts
of a probability that can oppose one another --- which is where the derivation of \S\ref{sec:self-reference}
lives. The commutativity lost at the second is the existence of actions whose order matters; the
unit quaternions form the group $SU(2)$, and the framework houses weak isospin there, as the
internal algebra of the $B$ register whose order-sensitive, non-commuting dynamics the interaction sections develop
(\S\ref{sec:register}, \S\ref{sec:entanglement-classes}). The associativity lost at the third ---
the octonions retain the weaker property of alternativity \autocite{baez2002octonions} --- creates
the associator $[a,b,c] = (ab)c - a(bc)$, and the associator is where most of this paper's physics
is transacted: the stabiliser of a preferred complex direction inside the octonions' automorphisms
is the colour group (\S\ref{sec:hopf-justification}); the triality structure and its $K_3$
obstruction generate the three generations and the sterile channel (\S\ref{sec:generations}); and
the associator is the currency of mass --- the debt against the selected direction that
\S\ref{sec:masses} prices, owed by whatever action moves the selection and not by residence at
the rung, which is the defusal \S\ref{sec:register} makes explicit. The algebraic facts
in the third column are theorems; the identifications in the fourth are established in the sections
cited, each carrying its own epistemic flag there, and are not re-litigated here. What this table
asserts on its own account is only the correspondence's shape: one law surrendered per rung, one
sector of physics gained per law.

\subsection{Internal algebras, not a new quantum mechanics}

A reader versed in the quaternionic-quantum-mechanics literature will at this point raise a fair
objection: quantum theory over $\mathbb{H}$ as a foundation is known to be troubled --- the
tensor-product and composite-system difficulties set out by Adler and others
\autocite{adler1995quaternionic} --- and quantum theory over $\mathbb{O}$ founders earlier still,
since a non-associative algebra does not furnish a Hilbert space at all. The framework's position
must therefore be stated exactly, because it does not make the claim the objection targets. The base
of the physical Hilbert space here is $\mathbb{C}$ and remains $\mathbb{C}$; the higher algebras
enter as the \emph{internal} algebras of the register --- $\mathbb{H}$ for the $B$ qubit,
$\mathbb{O}$ for the $C$ qubit (\S\ref{sec:register}) --- in a tensor structure whose outer theory is
ordinary complex quantum mechanics. This is the same move by which the division algebras enter the
constructions of Furey and the related programmes without any proposal to rebase quantum mechanics
itself \autocite{furey2018three}. For the octonionic register the point is sharper still:
$\mathbb{O}^{2}$ is not a Hilbert space, and the framework accordingly treats the $C$ qubit not in
the Dirac--von Neumann manner but geometrically, as the sphere $S^{15}$ with its Hopf structure ---
which is precisely the form in which \S\ref{sec:hopf-justification} needs it. Nothing in this paper therefore
stands or falls with quaternionic or octonionic quantum mechanics as foundations; the ladder is
climbed by the registers' internal algebras, not by the theory's base field.

\subsection{The same ladder, twice}

The ladder appears in this framework in two guises, and it will help the reader to have them
distinguished now. Statically, it is the anatomy of a single relation: the three qubits of the
register take their values one rung apart --- $A$ over $\mathbb{C}$, $B$ over $\mathbb{H}$, $C$ over
$\mathbb{O}$ --- so that every particle carries the whole ladder as its internal structure
(\S\ref{sec:register}). Dynamically, it is the shape of measurement: a single qubit's pure states
fill $S^{3}$, an entangling measurement carries the joint state into $S^{7}$, and a third party into
$S^{15}$ --- the correspondence between qubit entanglement and the Hopf fibrations that
\S\ref{sec:hopf-justification} takes from a mature literature --- so that each act of measurement ascends one rung --- the climb the reader has already made,
narrated in \S\ref{sec:wigners-friend} --- and every completed relation is a terminated climb, which is
the reading \S\ref{sec:hopf-information-bound} gives the tower's endpoint. One ladder, seen twice:
as the anatomy of a relation and as the dynamics of relating. The framework regards this double
appearance not as a coincidence but as the same structure viewed statically and in motion, and the
sections cited carry that identification's weight.

\subsection{Hurwitz closes the ladder}

The construction itself would double forever; what stops it is a classification theorem. By
Hurwitz's theorem, the normed division algebras --- those in which the norm composes,
$\lvert ab\rvert = \lvert a\rvert\,\lvert b\rvert$, and no non-zero elements multiply to zero ---
are exactly $\mathbb{R}$, $\mathbb{C}$, $\mathbb{H}$ and $\mathbb{O}$
\autocite{hurwitz1898composition,baez2002octonions}. The next doubling, the sedenions, has zero
divisors: non-zero elements whose product vanishes. The framework's reading of this failure is the
one that matters physically. Zero divisors cannot support consistent information flow around loops
--- a cycle of compositions can silently annihilate what it carries --- and a framework in which
closed computational paths are load-bearing (\S\ref{sec:hopf-justification},
\S\ref{sec:space-of-computations}) cannot be built on an algebra that permits it. At the level of
amplitudes the same failure reads: the composition law
$\lvert ab\rvert = \lvert a\rvert\,\lvert b\rvert$, which carries probability multiplicatively
through every rung of the ladder, fails beyond the octonions, so that two non-vanishing amplitudes
could compose to nothing and probability would be destroyed by mere composition. The last algebra
fit to carry amplitudes is $\mathbb{O}$. One rung past the terminus is not entirely silent,
however: the sedenions' internal structure --- their automorphism group and their octonion
subalgebras --- is the recorded candidate arena for tagging the generations
(\S\ref{sec:interactions}), and it enters exactly as this section's scoping permits, internally
and never as a base. The terminus closes the ladder of forces; what survives past it is, at most,
flavour. This is the algebraic half of the termination; the
geometric half --- that no sphere above $S^{7}$ is parallelisable --- and the theorem joining the
two halves are the opening business of \S\ref{sec:hopf-justification}, and the epistemic reading of the same endpoint,
as a bound on the information a single relation can carry, is given in
\S\ref{sec:hopf-information-bound}.

\subsection{The honest debt}

As in the geometry section, we close by stating plainly what is proved, what is established, and what is
owed. Proved: the construction, the losses, and the termination are theorems --- the Cayley--Dickson
doubling and its ledger of surrendered laws are standard algebra, and Hurwitz's classification is
among the durable results of the subject. Established: the physical cashing of each rung --- the
Born rule at $\mathbb{C}$, weak isospin at $\mathbb{H}$, colour, generations and mass at
$\mathbb{O}$ --- is carried out in the sections cited above, each under its own epistemic flag. Owed,
and stated without disguise: the ladder as presented here is closed from above and derived at its
floor, but it is not yet derived rung by rung from below. Hurwitz forbids a further rung; the
argument of \S\ref{sec:self-reference} forces the first doubling, from the ordered arithmetic of
classical probability to amplitudes that interfere. That the same self-referential closure, applied
to a system whose amplitudes already take values at one rung, compels the next --- that relational
closure at one level forces precisely the Cayley--Dickson double of that level --- is the
framework's reading of the climb, and it is a thesis, not a theorem --- but the formal apparatus
against which it can be posed now exists, and we pose it.

\begin{openproblem}[The compulsion of the intermediate rungs]\label{op:ladder-compulsion}
Derive, within the categorical apparatus of Open Problem~\ref{op:sigma-functor}, that
self-referential closure at one rung compels the Cayley--Dickson double of that rung, so that the
intermediate steps $\mathbb{C} \to \mathbb{H} \to \mathbb{O}$ follow from the same principle
that forces the floor. The candidate route is the construction's own shape: the Cayley--Dickson
pair with its conjugation, $(a, b)^{*} = (a^{*}, -b)$, is an involution-equipped doubling, and the
route reads it as the algebraic shadow of adjoining the describer --- closure of level-$A$
descriptions under the inside/outside exchange requiring exactly the twisted pair. The problem
inherits the status of Open Problem~\ref{op:sigma-functor}, whose apparatus it presupposes.
\end{openproblem}

If the derivation lands, the whole ladder becomes a consequence of the same self-reference that
opened it. Until then, the honest statement is that the ladder's necessity is established at its
ceiling and at its floor, while the compulsion of the intermediate rungs is inherited from the
physics they are shown to carry rather than derived in advance. What this section hands on is, accordingly, an algebraic ladder with a floor, a ceiling and a
ledger. Its first consumer is immediate: \S\ref{sec:space-of-computations} spends the terminal
fact --- the non-associativity of $S^{7}$'s own multiplication --- as geometry, excluding the
seven-sphere as a ruler and fixing the dimension of space at three.
\S\ref{sec:hopf-justification} then supplies the geometry of the fibrations themselves, and the
join.

%  (Architecture call SETTLED by the sync manifest: standalone.)
%
%  CITATIONS -- existing keys (as supplied by JS 2026-07-04):
%    Hardy2001a, Aaronson2004h, Thomson_1954, Lawvere_1969,
%    Turing_1937, muller_p_2011, Brukner2009.
%  NEW keys (RIS in continuity_refs_2026-07-04.ris):
%    cantor1891diagonal, benacerraf1962tasks, broeckertomdieck1985compact.
%  Deliberately excluded per JS: Chiribella--D'Ariano--Perinotti;
%  Bargmann/Wigner's theorem. Yanofsky2003 now paired on the
%  Lawvere-shaped clause (the pending-key excuse was void: the key
%  exists and sec:self-reference uses it).
%
%  CROSS-REFS: sec:self-reference now DRAFTED (2026-07-04, later the
%  same day) -- no longer dangling; only sec:fanout-boundary (and its
%  sibling sec:fanout-direction) -- DRAFTED 2026-07-06.
%
%  AUDITED 2026-07-06 (final pass, second stop): both proofs
%  re-verified by hand (Lagrange divisibility sentence exact; the
%  Baire countable-escape chain sound); caveat-discharge with
%  sec:self-reference confirmed two-way (already present); file pure
%  ASCII; no pre-definition "frame"; all citations used as declared.
%  Fixes: Yanofsky2003 paired on the Lawvere clause; the op's two
%  candidate routes anchored to printed propositions (f1-dictionary
%  for the divisibility premise's referent, inside-outside for the
%  closedness demand) -- statuses unchanged, referents sharpened;
%  header staleness corrected. Deliberate FORWARD references --
%  prop:f1-dictionary and sec:hopf-justification (defined later in the
%  paper; fine under two-pass LaTeX and flagged in prose),
%  sec:division-algebras, sec:space-of-computations.
%
%  DECISIONS ENCODED: op:limit-realisability minted (the owed arrow,
%  priced: exactly two derivations, divisibility and closedness);
%  Lagrange used twice (finite exclusion; component-group triviality);
%  the (1+i)/2 lamp value NOT stated here -- "a superposed fixed point"
%  with the complex-probability treatment deferred to
%  sec:self-reference; the two-branch over-determination stated to
%  pre-empt a circularity reading; Thomson spelling throughout.
% =====================================================================

\section{Continuity Is Not an Axiom}
\label{sec:continuity}

Hardy derives quantum theory from five axioms \autocite{Hardy2001a}. Four of them are satisfied
by classical probability theory; the fifth --- that a \emph{continuous} reversible
transformation connects any two pure states --- is the entire difference, the single postulate
that turns the classical simplex into the complex Hilbert space. Every reconstruction programme
pays an equivalent toll: continuous reversibility in Masanes and M\"{u}ller
\autocite{mullerp2011}, the corresponding postulate in Daki\'{c} and Brukner
\autocite{Brukner2009}. The axiom is load-bearing, and in every programme it is unexplained: it
is where the quantumness is put in by hand. The claim of this section is that in the present
framework it is not an axiom. Continuity is the fifth face of the pre-measurement regime that self-reference forces, and what
stands between the framework and a full derivation is not a postulate about values but one
premise about transformations --- that every transformation the inside performs has a stationary
outside presentation --- derived in outline below from the inside/outside involution of
\S\ref{sec:self-reference} and stated as an open problem at the end of the section with nothing
hidden in it. The continuum itself, the completion that every probability theory adopts, is
named for what it is: the real line, shared with classical probability and carrying no quantum
content.

One structural remark prevents a misreading. The framework reaches complex amplitudes natively in
\S\ref{sec:self-reference}, by the negative-probability route; the present section reaches the same
destination by an independent branch --- the regime's continuity, then Hardy's theorem. The two
branches converging on one structure is over-determination, not circularity, and the convergence is
itself evidence that the destination is forced.

\subsection{Five faces of one regime}

The regime before any measurement --- before any \textsc{fanout} --- can be characterised five
equivalent ways: negative probability is available; anti-events exist, every event being undoable;
no classical record exists; no preferred basis is selected; and a continuous reversible
transformation connects any two pure states. These are not five premises to be combined. They are
one regime under five descriptions, and \textsc{fanout} (\S\ref{sec:fanout-boundary}) is the single
event that fails all five at once: it writes the record, selects the basis, kills the anti-events,
discretises the transformation group, and switches off the negativity. The equivalence of the first
two faces is the anti-event identity of \S\ref{sec:self-reference}; the third is the boundary's
definition. The work of this section is the equivalence of the last two faces, and then the route
from the regime's existence to its continuity. The fourth face has, since
Proposition~\ref{prop:f1-dictionary} (\S\ref{sec:hopf-justification} --- a deliberate forward
reference), an exact form: a basis-selection is the commitment of a classical structure, a lifted
antipodal skeleton, so ``no preferred basis'' means precisely that no copy structure has fired.

\subsection{Basis selection is the collapse of continuity}

\begin{proposition}[Continuity forbids a selected basis]\label{prop:basis-discreteness}
If a finite set of pure states is privileged --- a preferred basis --- then every reversible
transformation respecting the privilege carries the set to itself and acts on it as a permutation;
the respecting transformations therefore realise only a finite group of motions among the
distinguished states, and no continuous family of them connects distinct pure states of the set.
Contrapositively, continuous reversibility entails that no basis is selected.
\end{proposition}

\begin{proof}
A privileged finite set is by definition carried to itself by every structure-preserving reversible
map, whose induced action on the set is a permutation of finitely many points. A continuous path of
such maps induces a continuous path in a finite discrete set, which is constant. Hence a continuous
reversible family moving one distinguished pure state to another respects no such privilege, and if
every pair of pure states is so connected, no finite set is privileged.
\end{proof}

The converse is where the historical difficulty lies: the absence of a privileged finite set does
not by itself force continuity, for the transformation group could in principle be infinite and
discrete. Two structural properties close every such escape at once.

\begin{proposition}[Divisible and closed implies connected]\label{prop:divisible-closed}
Let $G$ be the group of reversible transformations of the pre-measurement regime, acting on the
state body. If $G$ is divisible --- every element has an $n$-th root for every $n$ --- and
topologically closed, then $G$ is a compact connected Lie group; and if $G$ moreover acts
transitively on the pure states, then a continuous reversible transformation connects any two of
them: Hardy's fifth axiom.
\end{proposition}

\begin{proof}
The state body is compact and convex with nonempty interior, so its symmetry group is compact, and
a closed subgroup of a compact matrix group is a compact Lie group with finitely many connected
components \autocite{broeckertomdieck1985compact}. The component group $G/G_{0}$ is a quotient of a
divisible group and hence divisible; but a nontrivial finite group is never divisible --- with
$\lvert H\rvert = m$ every element satisfies $h^{m} = e$, so no element other than the identity
possesses an $m$-th root --- and $G/G_{0}$ is finite, so it is trivial and $G = G_{0}$ is connected.
A connected Lie group is path-connected, and transitivity then supplies the continuous reversible
path between any two pure states.
\end{proof}

The same one-line fact of Lagrange closes the remaining escapes. Finite groups are excluded
directly, being nontrivial and never divisible. Countable groups are excluded because a
countable closed subgroup of a compact group is compact, has an isolated point by the Baire
category theorem, and is therefore discrete and finite. Uncountable disconnected groups are
excluded by the component-group argument of the proposition. Every route to a discrete
transformation group therefore passes through a failure of divisibility or a failure of
closedness. Divisibility, however, asks for more than the conclusion needs, and the next
proposition removes it.

\begin{proposition}[Continuity from a connected state space]\label{prop:connected-transitive}
Let the pure states of a finite-dimensional system form a connected space $P$, and let the group
$G$ of reversible transformations of the pre-measurement regime be topologically closed and act
transitively on $P$. Then the identity component $G_0$ already acts transitively on $P$, and a
continuous reversible transformation connects any two pure states: Hardy's fifth axiom holds,
with no divisibility assumed.
\end{proposition}

\begin{proof}
As in Proposition~\ref{prop:divisible-closed}, closedness makes $G$ a compact Lie group with
finitely many connected components \autocite{broeckertomdieck1985compact}, so $G_0$ is a closed
normal subgroup of finite index and is path-connected. Each $G_0$-orbit in $P$ is the continuous
image of a compact group, hence compact, hence closed. Since $G$ is transitive and $G_0$ is
normal, the $G_0$-orbit of $gx$ is $g$ applied to the $G_0$-orbit of $x$, so the cosets of $G_0$
permute the orbits transitively and there are finitely many of them. A finite partition of $P$
into closed sets is a partition into open sets; $P$ is connected; so there is one orbit and
$G_0$ is transitive. Given pure states $x$ and $y$, choose $g \in G_0$ with $gx = y$ and a
continuous path $g_t$ in $G_0$ from the identity to $g$; then $g_t x$ is a continuous reversible
transformation from $x$ to $y$.
\end{proof}

The proposition relocates the premise. Divisibility was asked for because a transformation group
could be infinite and discrete; but a group acting transitively on a continuum cannot be
countable, and a closed group acting transitively on a connected space cannot be disconnected in
the way that matters, so what the argument needs from the regime is not the availability of
$n$-th roots but the connectedness of the state space itself. That is a statement about values:
that the pure states are all the unit vectors of $\mathbb{C}^{K}$ up to phase, the space
$\mathbb{CP}^{K-1}$, connected for $K \ge 2$, and hence that the amplitudes of
\S\ref{sec:self-reference} fill their continuum. The two standing foils show the premise doing
its work. The stabiliser fragment of quantum theory has complex amplitudes and is contextual,
yet it has finitely many pure states, its outcome probabilities take only the values $0$,
$\tfrac12$ and $1$, and its reversible transformations form a finite group, the Clifford group,
so Hardy's fifth axiom fails in it; Spekkens's toy theory \autocite{spekkens2007toy} is the same
with real amplitudes, its ignorance rationed by a knowledge-balance principle. Neither has a
connected state space, and neither is a model of the regime described here, whose credences take
every value. Divisibility, which both also lack --- the phase gate has no square root in the Clifford group
--- is a symptom of the discreteness, not its cause. What both lack can be said exactly, and the
next subsection says it: each contains a transformation that fixes no state it admits.

\subsection{The premises, and what they mean}

Three premises carry Proposition~\ref{prop:connected-transitive}. Transitivity is the regime's
homogeneity: inside/outside equivalence leaves no pure state privileged, and a proper closed
invariant set of pure states would define a superselection observable, a classical record
predating any \textsc{fanout} (\S\ref{sec:fanout-boundary}). Connectedness of the state space is two statements, and only one of them is quantum. The first
is that the pure states are infinite and dense, with a fixed point for every transformation; the
second is that the state space is complete. Quantum theory over the algebraic numbers ---
amplitudes in the algebraic closure of $\mathbb{Q}$ within $\mathbb{C}$ --- has the first
without the second: it has $K = N^{2}$, the Born rule, superposition, entanglement and Bell
violation at the Tsirelson value, every unitary has an algebraic eigenvector, and it is
countable and totally disconnected. So the whole quantum content of Hardy's axiom lies in the
density, and the completion is the real-number convention of every probability theory. Hardy
adopts it before his fifth axiom is stated, in the first, which makes probabilities real numbers
and gives his classical theory a real simplex of mixtures; this framework adopts it at the
Shannon step, where credences are declared to range over a connected domain so that the
continuation of the logarithm is unique (\S\ref{sec:complex-amplitudes}). Given that convention,
density of the pure states is connectedness of their closure. Divisibility, which an earlier
version of this section asked for, is not needed: it follows from the conclusion, every element
of a compact connected Lie group having roots of every order.

Closedness is the completion, limit-realisability: a Cauchy sequence of available
transformations has its limit available. It cannot be derived, and the reason should be stated
rather than hidden. Any collection of first-order principles about states, transformations and
probabilities that has an infinite model has a countable one, and a countable subset of
$\mathbb{CP}^{K-1}$ is totally disconnected; the founding non-computability does not escape
this, since the arithmetical complex numbers form a countable algebraically closed field
containing non-computable numbers and closed under every limit it can itself describe. Hardy's
axiom and closedness are second-order, quantifying over paths and over sequences, and a
second-order principle is grounded as Dedekind completeness is grounded or not at all. The
grounds are these. The completion of a dense set under the finite-precision metric is unique and
depends on no choice, whereas any countable alternative is a choice of definability level, an
extra law with no ground and no test; and the framework's own result is that no observer's
descriptive level captures the regime, since by \S\ref{sec:self-reference} every observer's
state exceeds its own level and the inside/outside principle makes that a law rather than a
perspective. The burden therefore falls on any restriction, not on the completion, which is the
least committal option and carries no quantum content. What limit-realisability adds to the line
is still worth stating exactly, because it is the corrected form of a natural intuition.
Cantor's diagonal \autocite{cantor1891diagonal} does not build the continuum; its output is an
infinite decimal --- a Cauchy sequence of rationals --- which exists \emph{as a point} only
because the line is complete. The diagonal presupposes limit-realisability and then shows that
the completed line outruns every internal listing. And the character of the completion is the
framework's own: the computable reals are countable \autocite{Turing_1937}, hence of measure
zero, so once the transformation group contains its limits almost all of them are
non-computable. Non-computability is a property of the completed continuum, not the reason for
it. The constructive reading, on which only describable limits exist, remains consistent; it is
the countable theory above, and the two readings differ in nothing an experiment can certify.

Density is where the section's title is earned, and it is derived from the paper's own premise.
Write $(\mathrm{I})$ for the statement that every admissible transformation of the regime is a
phase of some admissible classical structure: diagonal, that is, in the record basis of some
observer the regime contains, a phased spider of that basis in the calculus of
\S\ref{sec:hopf-justification} \autocite{coeckeduncan2011interacting} (VERIFY-SUPPORT). It says
of transformations what the global phase of \S\ref{sec:self-reference} says of states: there are
no absolute ones. It follows from three things the paper already holds. The inside/outside
involution is a genuine involution, so the outside presentation of the inside is a presentation
of the same object with nothing lost: the outside can have complete knowledge of the inside, a
potential record, the joint state lying in a structure the outside carries. Time is operation
count, internal to the system that operates (\S\ref{sec:space-of-computations}), so the outside
does not share the inside's count. And before any \textsc{fanout} that count is not classically
recorded, so the outside's description is pure. Let $g$ be the transformation the inside
performs per operation. The outside's complete knowledge cannot depend on the inside's count,
and it is pure, so it is a pure state of the inside's clock and system that is stationary under
``advance the clock and apply $g$''; the stationary states of that map are clock Fourier modes
tensored with eigenvectors of $g$ and nothing else \autocite{pagewootters1983evolution}
(VERIFY-SUPPORT). So the outside holds a structure in which $g$ acts by phases, which is
$(\mathrm{I})$; and the relative phase in that structure is the outside's ignorance of the
inside's computational count, the same ignorance the global phase records one level down. The
consequences are a theorem. A phase of a structure fixes the structure's copyable states, so
every admissible transformation fixes an admissible pure state; a finite transitive action has
an element moving every point (Jordan), so the pure states are infinite; the transformation
carrying a stabiliser state to the fixed point of the composite of a $Z$-phase and an $X$-phase
turns through $\arccos(1/\sqrt{3})$, an irrational fraction of a turn, so the group has an
element of infinite projective order and its orbits are dense. With the completion,
Proposition~\ref{prop:connected-transitive} fires and Hardy's fifth axiom holds. The two foils
fail $(\mathrm{I})$ exactly where they fail continuity: fourteen of the twenty-four Clifford
operations are phases of no Pauli basis and fix no stabiliser state, their fixed points being
the magic states; the toy theory's three-cycles are phases of no epistemic basis; and in each
the outside's knowledge of the inside under the order-three protocol is either count-dependent
or a mixture over the count, so each makes the inside's time classical before any record exists.
The derivation of $(\mathrm{I})$ is an outline, not a proof, and it is what the open problem
below asks for.

\subsection{Why the cut matters physically}

The regime's continuity is not decorative. Reversibility alone, over complex amplitudes, licenses
the full invertible group $GL(n, \mathbb{C})$, and Aaronson showed what that buys: superluminal
signalling, and computational power at the level of \textsf{PP} \autocite{Aaronson2004h}. The cut
from $GL(n, \mathbb{C})$ down to the unitary group is the conserved quadratic form, and continuity
is what entitles the theory to make it: with the fifth axiom in hand, Hardy's theorem delivers the
complex Hilbert space, the quadratic norm and the unitary group in one stroke \autocite{Hardy2001a}
--- cited here, not re-proved. One honest caveat travels with the citation. Hardy's reconstruction
uses his remaining axioms besides --- composition and subspace structure --- and these must be
derived rather than imported: the composition ledger of \S\ref{sec:division-algebras} and the
locality-of-tomography argument of \S\ref{sec:self-reference} carry that duty, and it is discharged
there, not silently assumed here.

\subsection{Super-tasks at the boundary}

The premise just priced has a fifty-year-old stress test in the philosophy of mathematics, and the
framework passes it in an unexpected way. A super-task --- infinitely many operations completed in
finite time \autocite{Thomson_1954} --- presupposes its terminal moment as an actual point of the
timeline: the accumulation point of a generic sequence of step-times exists only in the completed
line, and the defining question, the state \emph{after all the steps}, is meaningful only under
limit-realisability of time. Super-tasks are downstream consumers of exactly the completeness at
issue in this section. Thomson's lamp then exhibits both sides of the \textsc{fanout} boundary at
once. On the record side, the lamp's description fixes its state at every finite stage and thereby
fixes no limit at all: any terminal value is consistent with the whole task
\autocite{Thomson_1954,benacerraf1962tasks}. That is not a puzzle about lamps; it is the proof that
the classical record algebra is not limit-closed, and cannot be. On the amplitude side the limit
exists --- a superposed fixed point, whose complex-probability treatment belongs to
\S\ref{sec:self-reference} --- because that side \emph{is} limit-closed, which is the fifth face
itself. The boundary is thereby the edge of limit-realisability: complete on the reversible side,
incomplete on the record side, with the lamp standing on both sides at once as the boundary's oldest
exhibit. A final clause dissolves rather than answers the classic question. In this framework time
is operation count (\S\ref{sec:space-of-computations}), and ``after infinitely many operations''
indexes no operational moment: the lamp packs its endless \textsc{fanout}s into the emergent
continuum, and asking for its record at the accumulation point asks the discrete contents for a
value that only the completed container possesses. On or off is not undecided; it is unasked.

\subsection{What remains, priced}

\begin{openproblem}[Every transformation a phase of an inside structure]\label{op:limit-realisability}
Proposition~\ref{prop:connected-transitive} reduces Hardy's fifth axiom to three premises:
connectedness of the pure states, transitivity, and closedness. Transitivity is argued above
from the absence of records. Closedness is the completion, adopted on the grounds stated above
and carrying no quantum content. Connectedness is density plus that completion, and density
follows, by Jordan's theorem and the infinite-order argument, from $(\mathrm{I})$: that every
admissible transformation of the regime is a phase of some admissible classical structure. What
remains open is the rigour of the derivation of $(\mathrm{I})$ from the paper's premise. Derive,
inside the fixed-point locus $\mathrm{Fix}(\sigma)$ of Proposition~\ref{prop:inside-outside},
that every endomap is a phase of a classical structure in $\mathrm{Fix}(\sigma)$, from the
self-representing object, the involution $\sigma$ as a genuine involution, and the reading of
the inside presentation as a sequence indexed by operation count with no record before
\textsc{fanout}. The outline above does this with the stationary-state construction of Page and
Wootters; a derivation passes if it makes that outline a theorem of the calculus, with the
outside's count-independent complete description shown to exist and to be pure from the
involution alone. It fails before any detail is checked if every premise it uses holds in the
stabiliser fragment of quantum theory, in Spekkens's toy theory \autocite{spekkens2007toy}, in
the real qubit, or in quantum theory over the Gaussian rationals, all of which violate
$(\mathrm{I})$ outright and all of which make the inside's operation count classical before any
record exists. It also fails if it imports the completion as a premise about transformations:
the completion is the real line and has no derivation, and a derivation of $(\mathrm{I})$ must
hold in the countable algebraic theory, where $(\mathrm{I})$ is true and continuity is not.
\end{openproblem}

The accounting, in one sentence: the axiom that every reconstruction programme pays is, in this
framework, one derivation posed, of Lawvere's shape --- that every endomap of the fixed-point
locus is a phase of a classical structure in it --- and one convention named, the completion,
which classical probability pays too. And
the hand-on is already visible: the complex amplitudes that this section's continuity underwrites
through Hardy's theorem, and the complex amplitudes that \S\ref{sec:self-reference} derives natively,
are one and the same $\mathbb{C}$ --- the floor of the ladder that \S\ref{sec:division-algebras}
climbs.

%      % =====================================================================
%  SECTION: The Division-Algebra Ladder
%  Drafted 2026-07-03. One of the three foundations sections
%  (sec:self-reference, sec:division-algebras, sec:fanout-boundary);
%  this one resolves the dangling \S\ref{sec:division-algebras} in the
%  Hopf-justification section, whose header names it "the
%  Cayley--Dickson / Hurwitz step".
%
%  PLACEMENT in main.tex (REORDERED per JS 2026-07-06): position 6,
%  opening Act II ("why stop at C?") right before its first consumer --
%      % EDITED 2026-10-29 scripts/patch_continuity_tower_premise_2026-10-29T1500.py
% EDITED 2026-09-20 scripts/patch_born_cost_ledger_v2_2026-09-20T0900.py
% =====================================================================
%  SECTION: Self-Reference: One Principle
%  Drafted 2026-07-04. The paper's foundational section; resolves the
%  largest dangling label in the manuscript (sec:self-reference).
%
%  PLACEMENT in main.tex: FIRST of the foundations chain --
%      % EDITED 2026-10-29 scripts/patch_continuity_tower_premise_2026-10-29T1500.py
% EDITED 2026-09-20 scripts/patch_born_cost_ledger_v2_2026-09-20T0900.py
% =====================================================================
%  SECTION: Self-Reference: One Principle
%  Drafted 2026-07-04. The paper's foundational section; resolves the
%  largest dangling label in the manuscript (sec:self-reference).
%
%  PLACEMENT in main.tex: FIRST of the foundations chain --
%      \input{self_reference_section_2026-07-04}
%      \input{continuity_section_2026-07-04}
%      \input{division_algebras_section_2026-07-03}
%      \input{hopf_justification_section_2026-06-08}
%
%  SOURCES (re-read 2026-07-04): lawvere_inside_outside_fixed_point_
%  theorem_2026-05-04T1806.md (the spine: the two-principles reduction,
%  the theorem, the sigma candidate, the open items -> op minted per
%  D9); research_summary_neg_prob.md (negativity route; lamp (1+i)/2);
%  dynamics_companion_fanout_unitarity (legs 1a/1c; anti-events; the
%  Brukner-Zeilinger rebuttal per D6); Pure_states_and_mixed_states
%  (local tomography derived; QBism contrast per D7).
%
%  DECISIONS ENCODED (JS 2026-07-04): D1 title as below; D2 Wheeler1990
%  cited here as the title anchor; D3 lineage Popper -> Breuer ->
%  Svozil -> Szangolies (Dalla Chiara omitted); D4 Frauchiger_2018
%  cited here, dissolved in sec:wigners-friend; D5 Kirkwood-Dirac
%  heritage via the Arvidsson-Shukur 2024 review as one-stop; D6
%  Brukner-Zeilinger + the framework rebuttal IN ("Shannon is
%  post-FANOUT"); D7 QBism contrast IN; D8 Aaronson private
%  communication OMITTED (published island paper only); D9 one op
%  minted (op:sigma-functor), cross-linked to op:hopf-frobenius.
%
%  CITATIONS -- existing keys: Wheeler1990, small_2006, Godel1931,
%  Lawvere_1969, Yanofsky2003, CoeckeKissinger2017, Rovelli1996a,
%  Popper1950a, Popper1950b, Breuer1995, svozil_k_2018,
%  szangolies_j_2018, Frauchiger_2018, dirac1942, Feynman_r_82,
%  Feynman1987, Youssef1994 [SWAP-FLAG: intended paper is "Quantum
%  mechanics as complex probability theory", Mod. Phys. Lett. A 9
%  (1994) 2571 -- if that lives under youssef1994quantum, swap],
%  caticha_a_2006, goyal_knuth_skilling_2010, Aaronson2004h, Hardy2001a,
%  muller_p_2011, Thomson_1954, benacerraf1962tasks, cantor1891diagonal,
%  Turing_1937.
%  NEW keys (RIS in self_reference_refs_2026-07-04.ris, all verified
%  this session): tikochinsky1988feynman, renou2021real, chen2022ruling,
%  wu2022refutation, hoffreumonwoods2026real, arvidssonshukur2024kirkwood,
%  bruknerzeilinger2001shannon, fuchsschack2013qbism,
%  shannon1948mathematical.
%
%  CROSS-REFS: forward (fine under two-pass LaTeX): sec:continuity,
%  sec:division-algebras, sec:hopf-justification, prop:f1-dictionary,
%  sec:wigners-friend, op:limit-realisability, op:hopf-frobenius.
%  Still dangling after this section: sec:fanout-boundary (and its
%  sibling label sec:fanout-direction) -- both DRAFTED 2026-07-06.
%
%  AUDITED 2026-07-06 (section-by-section final pass, first stop):
%  proof of prop:inside-outside re-verified by hand; attribution chain
%  clean (Burgin/Khrennikov correctly absent); all header-declared
%  citations confirmed used; lamp arithmetic and antecedents verified;
%  FANOUT smallcaps consistent; no pre-definition use of "frame".
%  Fixes: stray Unicode em-dash (QBism paragraph, the file's only
%  non-ASCII character) -> house ---; trace-blind two-way stitch added
%  to the Born subsection ("one privacy, two duties");
%  op:sigma-functor gains its first standing constraint (F1
%  inheritance via the op:hopf-frobenius discharge) and its previously
%  one-directional local-tomography cargo made explicit.
% =====================================================================

\section{Self-Reference: One Principle}
\label{sec:self-reference-2}

Wheeler asked for it from bit: for physics derived from information, ``every it --- every particle,
every field of force'' obtained from binary answers to yes-or-no questions \autocite{Wheeler1990}.
This paper's route to that programme runs through G\"{o}del: the single principle from which
everything else here is derived, deferred, or priced is \emph{self-reference} --- the fact that any
complete physical description, made from inside the world it describes, must describe the describer.
This section states the principle, proves that the framework's second historical commitment follows
from it as a theorem rather than standing beside it as an axiom, and then extracts the price of the
principle in four instalments: probabilities whose parts can cancel, amplitudes valued in
$\mathbb{C}$, the Born rule, and the local tomography of composite systems. What the section cannot
yet derive is stated as precisely as what it can, in one numbered problem and in the ledger at the
end.

\subsection{Events, not objects}

The framework's ontology is relational: the fundamental entities are events and the relations
between them, and objects are the bookkeeping of relational invariants --- what remains stable
across many relations. The kinship with relational quantum mechanics \autocite{Rovelli1996a} is
genuine and acknowledged: facts are facts relative to the systems that register them. The
difference, developed over the rest of this paper, is that the present framework grounds the
relational structure in a single logical principle and then carries it, through the division
algebras and the Hopf fibrations, all the way to the gauge group and the particle spectrum ---
relations are not only the ontology but the derivation engine. On this reading a measurement is not
the revelation of a pre-existing property; it is the creation of a relation, and the later sections
of this paper read that creation geometrically (\S\ref{sec:hopf-justification}) and even
gravitationally. Here we need only the logical core: whatever describes the world is an event in
the world it describes.

\subsection{One principle, and a theorem}

The framework historically carried two foundational commitments: self-reference, and the
\emph{Inside/Outside Equivalence Principle} --- that the laws of physics must be identical whether
formulated by an embedded participant or by an external observer. Two principles are one too many,
and this subsection reduces them.

The reduction corrects the framework's own origin, and we make the correction in print. The 2005
form of the argument \autocite{small_2006} ran a G\"{o}delian chain: a self-referential process,
being complete, must by the incompleteness theorem \autocite{Godel1931} be inconsistent; an
inconsistent system can negate itself; and the negation of self-reference is the production of
causal sequence. The chain captures the right intuition and is retained here as motivation, but it
is not sound as a proof: classical inconsistency yields \emph{ex falso quodlibet} --- every
proposition provable, every causal sequence an output, no constraint at all. The load-bearing
argument is different, and better: a fixed-point theorem.

Let $\mathcal{C}$ be the category of physical descriptions and admissible processes --- concretely,
the dagger-compact-closed setting with classical structures in which categorical quantum mechanics,
and the whole of this paper, works \autocite{CoeckeKissinger2017}. Lawvere showed that in a
sufficiently rich such setting, any weakly point-surjective self-representation forces every
endomap to have a fixed point, and Yanofsky showed that one diagonal argument of this shape
underlies Cantor's, G\"{o}del's, Turing's and Tarski's theorems alike
\autocite{Lawvere_1969,Yanofsky2003,cantor1891diagonal}. The framework's single commitment is that
$\mathcal{C}$ contains a \emph{self-representing object} --- an object capable of encoding the
processes on itself. That is what ``self-reference is foundational'' means, stated categorically,
and it is the surjectivity hypothesis of the diagonal lemma. Now let
$\sigma \colon \mathcal{C} \to \mathcal{C}$ be the \emph{inside/outside involution}: the operation
exchanging the two complementary presentations of any description --- the inside presentation, as
participant-indexed total space in which the observer is part of the system, and the outside
presentation, as a base-times-fibre decomposition over a classical parameter space. That
$\sigma^{2} = \mathrm{id}$ is not an axiom but a coherence requirement: the outside view of the
inside view of a thing is the thing. The candidate concrete realisation of $\sigma$ is the
Hopf-duality swap --- total space against base-times-fibre for the three fibrations --- which gives
the involution computable content and ties this subsection to \S\ref{sec:hopf-justification}.

\begin{proposition}[Inside/outside fixed point]\label{prop:inside-outside}
Let $\mathcal{C}$ satisfy the richness condition above and contain a self-representing object, and
let $\sigma$ be an involution on $\mathcal{C}$. Then the fixed-point locus
$\mathrm{Fix}(\sigma)$ is non-empty: there exist descriptions invariant under the exchange of
inside and outside perspective.
\end{proposition}

\begin{proof}
Lawvere's construction, applied to $\sigma$ in place of negation: given a weakly point-surjective
$\varphi \colon A \to Y^{A}$, form $g(a) = \sigma\!\left(\varphi(a)(a)\right)$; point-surjectivity
supplies $a^{*}$ with $\varphi(a^{*}) = g$, and then $\varphi(a^{*})(a^{*})$ is a fixed point of
$\sigma$ \autocite{Lawvere_1969,Yanofsky2003}. Where the same construction applied to Boolean
negation yields paradox --- negation has no fixed point on a two-element type, so the conclusion
explodes the premises --- the inside/outside involution is an honest involution on a rich type, and
the conclusion is a constructive fixed point rather than a contradiction.
\end{proof}

\textbf{Definition.} \emph{Physical law is the structure carried by} $\mathrm{Fix}(\sigma)$.

The Inside/Outside Equivalence Principle is thereby demoted from axiom to theorem: physical law is
invariant under the exchange of perspectives \emph{because it is the fixed-point set of that
exchange}, and the fixed-point set is forced to be non-empty by self-reference alone. One
commitment remains foundational; everything the second commitment asserted is now derived. The
reader who accepts categorical quantum mechanics as an existing programme has already accepted the
machinery; the framework adds one premise and collects the rest. And the fixed-point locus is not
empty of structure: it inherits the Frobenius-algebra apparatus of $\mathcal{C}$, so that the
Bell, GHZ and W primitives --- and with them the particle classification of the later sections ---
live inside $\mathrm{Fix}(\sigma)$ from the start. What is proved above is non-emptiness given the
setting; what is assumed is the setting itself and the existence of $\sigma$ as a genuine
involution; what is open is the rigour of both, and it is priced now rather than hidden:

\begin{openproblem}[The involution made functorial]\label{op:sigma-functor}
Fix the category $\mathcal{C}$ rigorously --- the dagger-compact-closed setting with classical
structures is the candidate --- and define $\sigma$ on morphisms, not only on objects, the natural
route being the induced fibre--base exchange on spider diagrams; establish that $\sigma$ is an (at
worst pseudo-)involutive endofunctor; and prove or refute the identification of
$\mathrm{Fix}(\sigma)$ with the three Hopf cross-sections, under which Adams's theorem would become
a corollary of the fixed-point classification. This problem is cross-linked with Open
Problem~\ref{op:hopf-frobenius}: under an affirmative resolution of both, the parallelisable
terminus, the fixed-point classification, and compositional completeness are three statements of
one fact, and progress on either problem constrains the other. One constraint is already in place:
at the complex rung, Proposition~\ref{prop:f1-dictionary} supplies the dictionary between
classical structures and the fibration's lifted skeletons --- the clause of
Open Problem~\ref{op:hopf-frobenius} it discharges --- and any resolution here must reproduce it.
The problem also carries a piece of cargo recorded later in this section: the formalisation, within
the fixed category, of the perspective-accounting that derives local tomography.
\end{openproblem}

\subsection{The lineage}

The obstruction this section builds on is old, and the framework claims its completion, not its
discovery. Popper argued in 1950 that a predictor embedded in the world has irreducibly inferior
knowledge of itself, because it cannot impart information about itself without, by that very act,
rendering the information out of date \autocite{Popper1950a,Popper1950b} --- the passage is quoted
where the framework makes the obstruction structural (\S\ref{sec:wigners-friend}). Breuer proved
the impossibility of accurate state self-measurement within quantum theory \autocite{Breuer1995};
Svozil's monograph surveys the whole territory where physics meets self-reference, undecidability
and Chaitin-style algorithmic limits \autocite{svozil_k_2018}; Szangolies sharpened the modern form
as epistemic horizons \autocite{szangolies_j_2018}; and the Frauchiger--Renner theorem is the
contemporary crisis-form --- nested observers deriving contradictory records when each treats the
others' quantum descriptions as universally valid \autocite{Frauchiger_2018} --- which
\S\ref{sec:wigners-friend} dissolves relationally rather than paradoxically. This framework's
addition to the lineage is twofold. Structurally, Proposition~\ref{prop:f1-dictionary} converts
Popper's temporal obsolescence into algebraic cancellation: the datum a copier cannot put into its
own record is exactly its fibre commitment, which does not arrive stale --- it cancels. And
foundationally, the present section runs the lineage in reverse: where the tradition treats
self-reference as a limit on knowledge \emph{within} physics, the framework takes it as the
generator of physics, and derives the quantum formalism from it.

\subsection{The price is negative probability}\label{sec:complex-amplitudes-2}

What does it cost a probabilistic description to contain its own describer? The first instalment is
the sign. In the regime before any record exists, every event has an anti-event --- the description
must be able to undo what it has provisionally done, or the self-description could never be
revised into consistency --- and a probability calculus closed under undoing cannot be confined to
$[0,1]$: it is a signed calculus, in which the parts of a probability can oppose and cancel. The
framework did not discover negativity as bookkeeping; Dirac admitted negative probabilities as
formal intermediaries \autocite{dirac1942}, Feynman defended them as consistent whenever the final,
recorded probabilities stay positive \autocite{Feynman_r_82,Feynman1987}, and the Kirkwood--Dirac
distribution has carried complex quasiprobability from 1933 to the present, as a recent
comprehensive review documents \autocite{arvidssonshukur2024kirkwood}. What the framework adds is
the forcing: negativity is not a permissible convenience but the first price of self-referential
closure, and it is charged exactly and only in the pre-record regime --- which is why every
negativity in the tradition lives in intermediate, unrecorded quantities and never in a final
record.

The second instalment is complexity. A signed information measure must extend Shannon's
$S = -\ln p$ off the positive axis, and the analytic continuation of the logarithm carries $p$ into
the complex plane: amplitudes, with the phase as the continued information. Here too the extension
has a prior tradition, and the attributions are owed precisely. Tikochinsky derived Feynman's rules
for complex amplitudes from consistency of composition \autocite{tikochinsky1988feynman}; Youssef
reformulated quantum mechanics as complex probability theory outright \autocite{Youssef1994};
Caticha carried objective amplitudes to Bayesian probabilities \autocite{caticha_a_2006}; and
Goyal, Knuth and Skilling derived Feynman's rules from a pair-postulate without assuming
analyticity \autocite{goyal_knuth_skilling_2010}. The extension is prior art. The framework's
claim is the compulsion: the complex extension is not a modelling choice among exotic probability
theories but the unique continuation forced by self-referential closure --- the same forcing that
Aaronson glimpsed when he argued quantum mechanics is what negative probability inevitably becomes
\autocite{Aaronson2004h}.

Why does the continuation stop at $\mathbb{C}$? Downward, the real field is not available: real
quantum theory is experimentally distinguishable from complex quantum theory in network scenarios
and has been disfavoured by experiment
\autocite{renou2021real,chen2022ruling,wu2022refutation}; the recent objection that the exclusion
rests on locally tomographic composition \autocite{hoffreumonwoods2026real} is, for this framework,
not an escape but a confirmation of where the weight sits --- local tomography is \emph{derived}
below, not assumed, so the framework closes from its own resources exactly the door the objection
holds open. Upward, the quaternions and octonions fail as a \emph{base} for the information
continuation: the continued measure is a logarithm, and beyond $\mathbb{C}$ the logarithm of a
product is no longer the sum of logarithms --- non-commutativity severs
$\log(ab) = \log a + \log b$ --- so the additivity that makes it an information measure is lost.
Aaronson's parameter-counting and signalling arguments against a quaternionic base point the same
way \autocite{Aaronson2004h}. The higher algebras are not banished; they return as the
\emph{internal} algebras of the register, which is the business of \S\ref{sec:division-algebras}.
The base is $\mathbb{C}$, for informational reasons, and stays $\mathbb{C}$.

\subsection{The Born rule from fibre privacy}\label{sec:born-rule}

The third instalment is the quadratic rule. An amplitude carries one datum more than any record of
it can: its overall phase. This was asserted throughout the framework's development as the
observer's private, self-referential ignorance; it is now a proved statement, in exactly the form
needed. By Proposition~\ref{prop:f1-dictionary}, a classical structure --- the copy operation any
record requires --- is committed to a specific lift of its skeleton, and that commitment is absent
from every classical record the structure induces: the copier is fibre-committed, the record is
fibre-blind, and the fibre freedom is gauge for every observer. (The computation extends verbatim
from one qubit to the full register, where Proposition~\ref{prop:trace-blind} makes it the ground
of hypercharge's tracelessness --- one privacy, two duties.) Popper's obsolescence, made
cancellation. It follows that anything recordable is invariant under the fibre --- observables and
their statistics must be $U(1)$-invariant functions of the amplitude. The lowest-degree polynomial invariant
of a complex amplitude is the quadratic $|\psi|^{2}$ --- the non-polynomial invariant $|\psi|$,
of degree one, is excluded by the analyticity declared at \S\ref{sec:complex-amplitudes} --- and
normalisation fixes the coefficient:
probability is the minimal record-compatible function of the amplitude,
\[
p \;=\; |\psi|^{2},
\]
the Born rule as the cheapest invariant the fibre permits. The same square is reached from the
cost side in \S\ref{sec:born-cost-ledger}, where a record is the closed loop of a history and its
return leg and the exponent is the number of legs; neither route uses the lattice of projections
that Gleason's theorem needs, and both hold for a single qubit. The derivation is the framework's
own;
its premise --- that the fibre datum is structurally excluded from records --- is
Proposition~\ref{prop:f1-dictionary}, and its geometric reading --- probability as base-point data
under the Hopf projection --- is developed in \S\ref{sec:hopf-justification}.

One standing objection to information-theoretic foundations is answered here, because this
framework needs the answer. Brukner and Zeilinger argued that Shannon information is conceptually
inadequate for quantum measurements: it is basis-dependent, and its operational meaning presupposes
ignorance of a definite but unknown fact, which the quantum case does not supply
\autocite{bruknerzeilinger2001shannon}. The framework's reply is structural, and it is a
demarcation rather than a rebuttal. Shannon's precondition fails exactly where they say it fails
--- in the pre-record regime, where there is no definite unknown to be ignorant of --- and there
the operative measure is the quadratic invariant just derived. Shannon's axioms hold exactly where
the precondition holds --- on records, after \textsc{fanout}, where outcomes are definite and
ignorance is classical --- and there the uniqueness theorem forces $-\log p$
\autocite{shannon1948mathematical}. The basis-dependence dissolves with the same stroke: a basis
\emph{is} selected precisely at \textsc{fanout}, so the basis-dependent measure applies only where
a basis exists. Shannon is not inadequate; Shannon is post-\textsc{fanout}. The framework thereby
keeps both measures, each on its own side of the boundary, and uses the additivity of the
post-record measure --- the very property the previous subsection used to pin the base field ---
with a clear conscience.

\subsection{Thomson's lamp, resolved}

The section owes \S\ref{sec:continuity} a debt incurred there: the lamp's amplitude-side limit.
Thomson's lamp --- toggled at times $1 - 2^{-n}$, its state at $t = 1$ demanded --- has no
classical answer: the task description fixes every finite stage and thereby fixes no limit at all
\autocite{Thomson_1954,benacerraf1962tasks}, which \S\ref{sec:continuity} read as the proof that
the record algebra is not limit-closed. On the amplitude side the limit exists, and the calculus of
this section computes it. The toggle exchanges the two amplitudes; the self-consistent state of the
completed task is the fixed point of the exchange in the complex calculus, and solving it under the
signed continuation gives
\[
P(\mathrm{on}) = \tfrac{1 + i}{2}, \qquad P(\mathrm{off}) = \tfrac{1 - i}{2},
\]
a superposed fixed point whose imaginary part carries the unrecordable datum --- the phase of the
endless alternation, when the lamp \emph{would have} stopped. Two checks certify that the
assignment is calculus rather than decoration. The pair is normalised,
$P(\mathrm{on}) + P(\mathrm{off}) = 1$; and the Born rule of the previous subsection, applied to
the fixed point, returns $|P(\mathrm{on})|^{2} = \tfrac{1}{2}$ --- exactly the symmetric classical
ignorance that is the only sane record-level answer, recovered rather than imposed. The lamp,
which entered the framework at its very beginning as the place where classical computation meets
its boundary, thus exits it as a worked example: undetermined as a record, determinate as an
amplitude, the two sides split precisely along \textsc{fanout} --- and the diagonal construction it
physically instantiates \autocite{cantor1891diagonal,Turing_1937} is the same one that powers
Proposition~\ref{prop:inside-outside}. The status flag is owed here as everywhere: the
undetermination is theorem-level; the specific fixed-point value is the framework's complex
calculus, structural, offered with its two consistency checks shown.

\subsection{Local tomography, derived}\label{sec:local-tomography}

The fourth instalment concerns composites, and it discharges two promises at once. The
Inside/Outside Equivalence Principle --- now a theorem --- requires that an external observer's
characterisation of a composite system be recoverable by the embedded participants, and the
embedded participants have access only to local measurements and their correlations. Equivalence
of the two perspectives therefore entails that the state of a composite is determined by the
statistics of local measurements: \emph{local tomography}, obtained as a consequence rather than
imposed as an axiom. This converts a postulate of the reconstruction programmes
\autocite{Hardy2001a,mullerp2011} into a theorem-shaped consequence of the one principle ---
discharging the caveat recorded in \S\ref{sec:continuity} --- and it closes, from the framework's
own resources, the escape route in the real-number debate noted above: the assumption the
objection of \autocite{hoffreumonwoods2026real} contests is here not an assumption. The status is
flagged plainly: the derivation is structural, one paragraph of perspective-accounting, and its
formalisation within the categorical setting belongs to Open Problem~\ref{op:sigma-functor}.

\subsection{Contrast: QBism, and the ledger}

The nearest neighbour in spirit deserves a paragraph of honest contrast. QBism
\autocite{fuchsschack2013qbism} holds, as this framework does, that quantum probabilities are an
agent's --- that the theory is a calculus of expectation, not a catalogue of properties. What QBism
does not say is why the agent's expectations must be complex-valued and quadratically squared; the
Born rule enters as a normative addition to Bayesian coherence, and the complex amplitude as an
unexplained representation. This section is, in effect, the missing why: the agent's ignorance is
self-referential --- the agent is inside the description --- and self-referential ignorance is
complex, with the quadratic rule as the cheapest record-compatible invariant of it. Where QBism
stops at the agent, the framework derives the agent's calculus.

The ledger, then. Delivered in this section: the reduction of two principles to one
(Proposition~\ref{prop:inside-outside}); the pre-record regime with its anti-events, handed to
\S\ref{sec:continuity} as the five-face regime whose continuity is priced there
(Open Problem~\ref{op:limit-realisability}); the complex base, forced downward by experiment and
upward by the additivity of information; the Born rule from fibre privacy; the lamp's fixed point;
and local tomography. Owed, and stated where they are owed: the functorial rigour of the involution
and the fixed-point classification (Open Problem~\ref{op:sigma-functor}); the remaining premise of continuity (Open Problem~\ref{op:limit-realisability}, in
\S\ref{sec:continuity}); and the
rung-by-rung compulsion of the algebra ladder, posed as Open
Problem~\ref{op:ladder-compulsion} in \S\ref{sec:division-algebras}. The
hand-on is the paper's spine: the $\mathbb{C}$ this section derives natively and the $\mathbb{C}$
that \S\ref{sec:continuity} underwrites through Hardy's theorem are the same field, the floor of
the ladder that \S\ref{sec:division-algebras} climbs and \S\ref{sec:hopf-justification}
geometrises. One principle; everything else is theorem, derivation, or a debt with a stated
invoice.

%      % EDITED 2026-10-29 scripts/patch_continuity_tower_premise_2026-10-29T1500.py
% EDITED 2026-09-15 scripts/patch_value_continuum_declared_cencov_2026-09-15T1540.py
% EDITED 2026-09-15 scripts/patch_continuity_connected_state_space_2026-09-15T1430.py
% =====================================================================
%  SECTION: Continuity Is Not an Axiom
%  Drafted 2026-07-04 from continuity_limit_realisability_2026-07-04.md
%  (now in project documents) and the May record
%  (dynamics_companion_fanout_unitarity_2026-05-20T1430.md).
%
%  PLACEMENT in main.tex: a standalone section between self-reference
%  and division-algebras --
%      \input{self_reference_section_2026-07-04}
%      \input{continuity_section_2026-07-04}
%      \input{division_algebras_section_2026-07-03}
%  (Architecture call SETTLED by the sync manifest: standalone.)
%
%  CITATIONS -- existing keys (as supplied by JS 2026-07-04):
%    Hardy2001a, Aaronson2004h, Thomson_1954, Lawvere_1969,
%    Turing_1937, muller_p_2011, Brukner2009.
%  NEW keys (RIS in continuity_refs_2026-07-04.ris):
%    cantor1891diagonal, benacerraf1962tasks, broeckertomdieck1985compact.
%  Deliberately excluded per JS: Chiribella--D'Ariano--Perinotti;
%  Bargmann/Wigner's theorem. Yanofsky2003 now paired on the
%  Lawvere-shaped clause (the pending-key excuse was void: the key
%  exists and sec:self-reference uses it).
%
%  CROSS-REFS: sec:self-reference now DRAFTED (2026-07-04, later the
%  same day) -- no longer dangling; only sec:fanout-boundary (and its
%  sibling sec:fanout-direction) -- DRAFTED 2026-07-06.
%
%  AUDITED 2026-07-06 (final pass, second stop): both proofs
%  re-verified by hand (Lagrange divisibility sentence exact; the
%  Baire countable-escape chain sound); caveat-discharge with
%  sec:self-reference confirmed two-way (already present); file pure
%  ASCII; no pre-definition "frame"; all citations used as declared.
%  Fixes: Yanofsky2003 paired on the Lawvere clause; the op's two
%  candidate routes anchored to printed propositions (f1-dictionary
%  for the divisibility premise's referent, inside-outside for the
%  closedness demand) -- statuses unchanged, referents sharpened;
%  header staleness corrected. Deliberate FORWARD references --
%  prop:f1-dictionary and sec:hopf-justification (defined later in the
%  paper; fine under two-pass LaTeX and flagged in prose),
%  sec:division-algebras, sec:space-of-computations.
%
%  DECISIONS ENCODED: op:limit-realisability minted (the owed arrow,
%  priced: exactly two derivations, divisibility and closedness);
%  Lagrange used twice (finite exclusion; component-group triviality);
%  the (1+i)/2 lamp value NOT stated here -- "a superposed fixed point"
%  with the complex-probability treatment deferred to
%  sec:self-reference; the two-branch over-determination stated to
%  pre-empt a circularity reading; Thomson spelling throughout.
% =====================================================================

\section{Continuity Is Not an Axiom}
\label{sec:continuity}

Hardy derives quantum theory from five axioms \autocite{Hardy2001a}. Four of them are satisfied
by classical probability theory; the fifth --- that a \emph{continuous} reversible
transformation connects any two pure states --- is the entire difference, the single postulate
that turns the classical simplex into the complex Hilbert space. Every reconstruction programme
pays an equivalent toll: continuous reversibility in Masanes and M\"{u}ller
\autocite{mullerp2011}, the corresponding postulate in Daki\'{c} and Brukner
\autocite{Brukner2009}. The axiom is load-bearing, and in every programme it is unexplained: it
is where the quantumness is put in by hand. The claim of this section is that in the present
framework it is not an axiom. Continuity is the fifth face of the pre-measurement regime that self-reference forces, and what
stands between the framework and a full derivation is not a postulate about values but one
premise about transformations --- that every transformation the inside performs has a stationary
outside presentation --- derived in outline below from the inside/outside involution of
\S\ref{sec:self-reference} and stated as an open problem at the end of the section with nothing
hidden in it. The continuum itself, the completion that every probability theory adopts, is
named for what it is: the real line, shared with classical probability and carrying no quantum
content.

One structural remark prevents a misreading. The framework reaches complex amplitudes natively in
\S\ref{sec:self-reference}, by the negative-probability route; the present section reaches the same
destination by an independent branch --- the regime's continuity, then Hardy's theorem. The two
branches converging on one structure is over-determination, not circularity, and the convergence is
itself evidence that the destination is forced.

\subsection{Five faces of one regime}

The regime before any measurement --- before any \textsc{fanout} --- can be characterised five
equivalent ways: negative probability is available; anti-events exist, every event being undoable;
no classical record exists; no preferred basis is selected; and a continuous reversible
transformation connects any two pure states. These are not five premises to be combined. They are
one regime under five descriptions, and \textsc{fanout} (\S\ref{sec:fanout-boundary}) is the single
event that fails all five at once: it writes the record, selects the basis, kills the anti-events,
discretises the transformation group, and switches off the negativity. The equivalence of the first
two faces is the anti-event identity of \S\ref{sec:self-reference}; the third is the boundary's
definition. The work of this section is the equivalence of the last two faces, and then the route
from the regime's existence to its continuity. The fourth face has, since
Proposition~\ref{prop:f1-dictionary} (\S\ref{sec:hopf-justification} --- a deliberate forward
reference), an exact form: a basis-selection is the commitment of a classical structure, a lifted
antipodal skeleton, so ``no preferred basis'' means precisely that no copy structure has fired.

\subsection{Basis selection is the collapse of continuity}

\begin{proposition}[Continuity forbids a selected basis]\label{prop:basis-discreteness}
If a finite set of pure states is privileged --- a preferred basis --- then every reversible
transformation respecting the privilege carries the set to itself and acts on it as a permutation;
the respecting transformations therefore realise only a finite group of motions among the
distinguished states, and no continuous family of them connects distinct pure states of the set.
Contrapositively, continuous reversibility entails that no basis is selected.
\end{proposition}

\begin{proof}
A privileged finite set is by definition carried to itself by every structure-preserving reversible
map, whose induced action on the set is a permutation of finitely many points. A continuous path of
such maps induces a continuous path in a finite discrete set, which is constant. Hence a continuous
reversible family moving one distinguished pure state to another respects no such privilege, and if
every pair of pure states is so connected, no finite set is privileged.
\end{proof}

The converse is where the historical difficulty lies: the absence of a privileged finite set does
not by itself force continuity, for the transformation group could in principle be infinite and
discrete. Two structural properties close every such escape at once.

\begin{proposition}[Divisible and closed implies connected]\label{prop:divisible-closed}
Let $G$ be the group of reversible transformations of the pre-measurement regime, acting on the
state body. If $G$ is divisible --- every element has an $n$-th root for every $n$ --- and
topologically closed, then $G$ is a compact connected Lie group; and if $G$ moreover acts
transitively on the pure states, then a continuous reversible transformation connects any two of
them: Hardy's fifth axiom.
\end{proposition}

\begin{proof}
The state body is compact and convex with nonempty interior, so its symmetry group is compact, and
a closed subgroup of a compact matrix group is a compact Lie group with finitely many connected
components \autocite{broeckertomdieck1985compact}. The component group $G/G_{0}$ is a quotient of a
divisible group and hence divisible; but a nontrivial finite group is never divisible --- with
$\lvert H\rvert = m$ every element satisfies $h^{m} = e$, so no element other than the identity
possesses an $m$-th root --- and $G/G_{0}$ is finite, so it is trivial and $G = G_{0}$ is connected.
A connected Lie group is path-connected, and transitivity then supplies the continuous reversible
path between any two pure states.
\end{proof}

The same one-line fact of Lagrange closes the remaining escapes. Finite groups are excluded
directly, being nontrivial and never divisible. Countable groups are excluded because a
countable closed subgroup of a compact group is compact, has an isolated point by the Baire
category theorem, and is therefore discrete and finite. Uncountable disconnected groups are
excluded by the component-group argument of the proposition. Every route to a discrete
transformation group therefore passes through a failure of divisibility or a failure of
closedness. Divisibility, however, asks for more than the conclusion needs, and the next
proposition removes it.

\begin{proposition}[Continuity from a connected state space]\label{prop:connected-transitive}
Let the pure states of a finite-dimensional system form a connected space $P$, and let the group
$G$ of reversible transformations of the pre-measurement regime be topologically closed and act
transitively on $P$. Then the identity component $G_0$ already acts transitively on $P$, and a
continuous reversible transformation connects any two pure states: Hardy's fifth axiom holds,
with no divisibility assumed.
\end{proposition}

\begin{proof}
As in Proposition~\ref{prop:divisible-closed}, closedness makes $G$ a compact Lie group with
finitely many connected components \autocite{broeckertomdieck1985compact}, so $G_0$ is a closed
normal subgroup of finite index and is path-connected. Each $G_0$-orbit in $P$ is the continuous
image of a compact group, hence compact, hence closed. Since $G$ is transitive and $G_0$ is
normal, the $G_0$-orbit of $gx$ is $g$ applied to the $G_0$-orbit of $x$, so the cosets of $G_0$
permute the orbits transitively and there are finitely many of them. A finite partition of $P$
into closed sets is a partition into open sets; $P$ is connected; so there is one orbit and
$G_0$ is transitive. Given pure states $x$ and $y$, choose $g \in G_0$ with $gx = y$ and a
continuous path $g_t$ in $G_0$ from the identity to $g$; then $g_t x$ is a continuous reversible
transformation from $x$ to $y$.
\end{proof}

The proposition relocates the premise. Divisibility was asked for because a transformation group
could be infinite and discrete; but a group acting transitively on a continuum cannot be
countable, and a closed group acting transitively on a connected space cannot be disconnected in
the way that matters, so what the argument needs from the regime is not the availability of
$n$-th roots but the connectedness of the state space itself. That is a statement about values:
that the pure states are all the unit vectors of $\mathbb{C}^{K}$ up to phase, the space
$\mathbb{CP}^{K-1}$, connected for $K \ge 2$, and hence that the amplitudes of
\S\ref{sec:self-reference} fill their continuum. The two standing foils show the premise doing
its work. The stabiliser fragment of quantum theory has complex amplitudes and is contextual,
yet it has finitely many pure states, its outcome probabilities take only the values $0$,
$\tfrac12$ and $1$, and its reversible transformations form a finite group, the Clifford group,
so Hardy's fifth axiom fails in it; Spekkens's toy theory \autocite{spekkens2007toy} is the same
with real amplitudes, its ignorance rationed by a knowledge-balance principle. Neither has a
connected state space, and neither is a model of the regime described here, whose credences take
every value. Divisibility, which both also lack --- the phase gate has no square root in the Clifford group
--- is a symptom of the discreteness, not its cause. What both lack can be said exactly, and the
next subsection says it: each contains a transformation that fixes no state it admits.

\subsection{The premises, and what they mean}

Three premises carry Proposition~\ref{prop:connected-transitive}. Transitivity is the regime's
homogeneity: inside/outside equivalence leaves no pure state privileged, and a proper closed
invariant set of pure states would define a superselection observable, a classical record
predating any \textsc{fanout} (\S\ref{sec:fanout-boundary}). Connectedness of the state space is two statements, and only one of them is quantum. The first
is that the pure states are infinite and dense, with a fixed point for every transformation; the
second is that the state space is complete. Quantum theory over the algebraic numbers ---
amplitudes in the algebraic closure of $\mathbb{Q}$ within $\mathbb{C}$ --- has the first
without the second: it has $K = N^{2}$, the Born rule, superposition, entanglement and Bell
violation at the Tsirelson value, every unitary has an algebraic eigenvector, and it is
countable and totally disconnected. So the whole quantum content of Hardy's axiom lies in the
density, and the completion is the real-number convention of every probability theory. Hardy
adopts it before his fifth axiom is stated, in the first, which makes probabilities real numbers
and gives his classical theory a real simplex of mixtures; this framework adopts it at the
Shannon step, where credences are declared to range over a connected domain so that the
continuation of the logarithm is unique (\S\ref{sec:complex-amplitudes}). Given that convention,
density of the pure states is connectedness of their closure. Divisibility, which an earlier
version of this section asked for, is not needed: it follows from the conclusion, every element
of a compact connected Lie group having roots of every order.

Closedness is the completion, limit-realisability: a Cauchy sequence of available
transformations has its limit available. It cannot be derived, and the reason should be stated
rather than hidden. Any collection of first-order principles about states, transformations and
probabilities that has an infinite model has a countable one, and a countable subset of
$\mathbb{CP}^{K-1}$ is totally disconnected; the founding non-computability does not escape
this, since the arithmetical complex numbers form a countable algebraically closed field
containing non-computable numbers and closed under every limit it can itself describe. Hardy's
axiom and closedness are second-order, quantifying over paths and over sequences, and a
second-order principle is grounded as Dedekind completeness is grounded or not at all. The
grounds are these. The completion of a dense set under the finite-precision metric is unique and
depends on no choice, whereas any countable alternative is a choice of definability level, an
extra law with no ground and no test; and the framework's own result is that no observer's
descriptive level captures the regime, since by \S\ref{sec:self-reference} every observer's
state exceeds its own level and the inside/outside principle makes that a law rather than a
perspective. The burden therefore falls on any restriction, not on the completion, which is the
least committal option and carries no quantum content. What limit-realisability adds to the line
is still worth stating exactly, because it is the corrected form of a natural intuition.
Cantor's diagonal \autocite{cantor1891diagonal} does not build the continuum; its output is an
infinite decimal --- a Cauchy sequence of rationals --- which exists \emph{as a point} only
because the line is complete. The diagonal presupposes limit-realisability and then shows that
the completed line outruns every internal listing. And the character of the completion is the
framework's own: the computable reals are countable \autocite{Turing_1937}, hence of measure
zero, so once the transformation group contains its limits almost all of them are
non-computable. Non-computability is a property of the completed continuum, not the reason for
it. The constructive reading, on which only describable limits exist, remains consistent; it is
the countable theory above, and the two readings differ in nothing an experiment can certify.

Density is where the section's title is earned, and it is derived from the paper's own premise.
Write $(\mathrm{I})$ for the statement that every admissible transformation of the regime is a
phase of some admissible classical structure: diagonal, that is, in the record basis of some
observer the regime contains, a phased spider of that basis in the calculus of
\S\ref{sec:hopf-justification} \autocite{coeckeduncan2011interacting} (VERIFY-SUPPORT). It says
of transformations what the global phase of \S\ref{sec:self-reference} says of states: there are
no absolute ones. It follows from three things the paper already holds. The inside/outside
involution is a genuine involution, so the outside presentation of the inside is a presentation
of the same object with nothing lost: the outside can have complete knowledge of the inside, a
potential record, the joint state lying in a structure the outside carries. Time is operation
count, internal to the system that operates (\S\ref{sec:space-of-computations}), so the outside
does not share the inside's count. And before any \textsc{fanout} that count is not classically
recorded, so the outside's description is pure. Let $g$ be the transformation the inside
performs per operation. The outside's complete knowledge cannot depend on the inside's count,
and it is pure, so it is a pure state of the inside's clock and system that is stationary under
``advance the clock and apply $g$''; the stationary states of that map are clock Fourier modes
tensored with eigenvectors of $g$ and nothing else \autocite{pagewootters1983evolution}
(VERIFY-SUPPORT). So the outside holds a structure in which $g$ acts by phases, which is
$(\mathrm{I})$; and the relative phase in that structure is the outside's ignorance of the
inside's computational count, the same ignorance the global phase records one level down. The
consequences are a theorem. A phase of a structure fixes the structure's copyable states, so
every admissible transformation fixes an admissible pure state; a finite transitive action has
an element moving every point (Jordan), so the pure states are infinite; the transformation
carrying a stabiliser state to the fixed point of the composite of a $Z$-phase and an $X$-phase
turns through $\arccos(1/\sqrt{3})$, an irrational fraction of a turn, so the group has an
element of infinite projective order and its orbits are dense. With the completion,
Proposition~\ref{prop:connected-transitive} fires and Hardy's fifth axiom holds. The two foils
fail $(\mathrm{I})$ exactly where they fail continuity: fourteen of the twenty-four Clifford
operations are phases of no Pauli basis and fix no stabiliser state, their fixed points being
the magic states; the toy theory's three-cycles are phases of no epistemic basis; and in each
the outside's knowledge of the inside under the order-three protocol is either count-dependent
or a mixture over the count, so each makes the inside's time classical before any record exists.
The derivation of $(\mathrm{I})$ is an outline, not a proof, and it is what the open problem
below asks for.

\subsection{Why the cut matters physically}

The regime's continuity is not decorative. Reversibility alone, over complex amplitudes, licenses
the full invertible group $GL(n, \mathbb{C})$, and Aaronson showed what that buys: superluminal
signalling, and computational power at the level of \textsf{PP} \autocite{Aaronson2004h}. The cut
from $GL(n, \mathbb{C})$ down to the unitary group is the conserved quadratic form, and continuity
is what entitles the theory to make it: with the fifth axiom in hand, Hardy's theorem delivers the
complex Hilbert space, the quadratic norm and the unitary group in one stroke \autocite{Hardy2001a}
--- cited here, not re-proved. One honest caveat travels with the citation. Hardy's reconstruction
uses his remaining axioms besides --- composition and subspace structure --- and these must be
derived rather than imported: the composition ledger of \S\ref{sec:division-algebras} and the
locality-of-tomography argument of \S\ref{sec:self-reference} carry that duty, and it is discharged
there, not silently assumed here.

\subsection{Super-tasks at the boundary}

The premise just priced has a fifty-year-old stress test in the philosophy of mathematics, and the
framework passes it in an unexpected way. A super-task --- infinitely many operations completed in
finite time \autocite{Thomson_1954} --- presupposes its terminal moment as an actual point of the
timeline: the accumulation point of a generic sequence of step-times exists only in the completed
line, and the defining question, the state \emph{after all the steps}, is meaningful only under
limit-realisability of time. Super-tasks are downstream consumers of exactly the completeness at
issue in this section. Thomson's lamp then exhibits both sides of the \textsc{fanout} boundary at
once. On the record side, the lamp's description fixes its state at every finite stage and thereby
fixes no limit at all: any terminal value is consistent with the whole task
\autocite{Thomson_1954,benacerraf1962tasks}. That is not a puzzle about lamps; it is the proof that
the classical record algebra is not limit-closed, and cannot be. On the amplitude side the limit
exists --- a superposed fixed point, whose complex-probability treatment belongs to
\S\ref{sec:self-reference} --- because that side \emph{is} limit-closed, which is the fifth face
itself. The boundary is thereby the edge of limit-realisability: complete on the reversible side,
incomplete on the record side, with the lamp standing on both sides at once as the boundary's oldest
exhibit. A final clause dissolves rather than answers the classic question. In this framework time
is operation count (\S\ref{sec:space-of-computations}), and ``after infinitely many operations''
indexes no operational moment: the lamp packs its endless \textsc{fanout}s into the emergent
continuum, and asking for its record at the accumulation point asks the discrete contents for a
value that only the completed container possesses. On or off is not undecided; it is unasked.

\subsection{What remains, priced}

\begin{openproblem}[Every transformation a phase of an inside structure]\label{op:limit-realisability}
Proposition~\ref{prop:connected-transitive} reduces Hardy's fifth axiom to three premises:
connectedness of the pure states, transitivity, and closedness. Transitivity is argued above
from the absence of records. Closedness is the completion, adopted on the grounds stated above
and carrying no quantum content. Connectedness is density plus that completion, and density
follows, by Jordan's theorem and the infinite-order argument, from $(\mathrm{I})$: that every
admissible transformation of the regime is a phase of some admissible classical structure. What
remains open is the rigour of the derivation of $(\mathrm{I})$ from the paper's premise. Derive,
inside the fixed-point locus $\mathrm{Fix}(\sigma)$ of Proposition~\ref{prop:inside-outside},
that every endomap is a phase of a classical structure in $\mathrm{Fix}(\sigma)$, from the
self-representing object, the involution $\sigma$ as a genuine involution, and the reading of
the inside presentation as a sequence indexed by operation count with no record before
\textsc{fanout}. The outline above does this with the stationary-state construction of Page and
Wootters; a derivation passes if it makes that outline a theorem of the calculus, with the
outside's count-independent complete description shown to exist and to be pure from the
involution alone. It fails before any detail is checked if every premise it uses holds in the
stabiliser fragment of quantum theory, in Spekkens's toy theory \autocite{spekkens2007toy}, in
the real qubit, or in quantum theory over the Gaussian rationals, all of which violate
$(\mathrm{I})$ outright and all of which make the inside's operation count classical before any
record exists. It also fails if it imports the completion as a premise about transformations:
the completion is the real line and has no derivation, and a derivation of $(\mathrm{I})$ must
hold in the countable algebraic theory, where $(\mathrm{I})$ is true and continuity is not.
\end{openproblem}

The accounting, in one sentence: the axiom that every reconstruction programme pays is, in this
framework, one derivation posed, of Lawvere's shape --- that every endomap of the fixed-point
locus is a phase of a classical structure in it --- and one convention named, the completion,
which classical probability pays too. And
the hand-on is already visible: the complex amplitudes that this section's continuity underwrites
through Hardy's theorem, and the complex amplitudes that \S\ref{sec:self-reference} derives natively,
are one and the same $\mathbb{C}$ --- the floor of the ladder that \S\ref{sec:division-algebras}
climbs.

%      % =====================================================================
%  SECTION: The Division-Algebra Ladder
%  Drafted 2026-07-03. One of the three foundations sections
%  (sec:self-reference, sec:division-algebras, sec:fanout-boundary);
%  this one resolves the dangling \S\ref{sec:division-algebras} in the
%  Hopf-justification section, whose header names it "the
%  Cayley--Dickson / Hurwitz step".
%
%  PLACEMENT in main.tex (REORDERED per JS 2026-07-06): position 6,
%  opening Act II ("why stop at C?") right before its first consumer --
%      \input{self_reference_section_2026-07-04}
%      \input{continuity_section_2026-07-04}
%      \input{L2_metric_global_phase_section_2026_06_07}
%      \input{wigners_friend_in_the_hopf_picture_2026-06-13}
%      \input{derivation_of_the_quantum_postulates}
%      \input{division_algebras_section_2026-07-03}   <-- moved here
%      \input{spacetime_from_non-computability}       <-- first consumer
%                                                          (S^7 exclusion)
%      \input{hopf_justification_section_2026-06-08}
%  The section is position-portable: all adjacency phrasing was
%  retargeted to explicit \S\ref on 2026-07-06. The S^7 hand-off
%  sentence to sec:space-of-computations WRITTEN 2026-07-06 (spacetime
%  file read; the closing now hands first to its immediate consumer,
%  then to hopf-justification). Upload verifications COMPLETE: L2 and
%  quantum-postulates do not consume the ladder; WF's references are
%  forward-voiced or self-carrying.
%
%  DIVISION OF LABOUR (deliberate, to avoid overlap):
%   - sec:self-reference (undrafted): self-reference -> complex
%     amplitudes and the Born rule; why not a real base. This section
%     TAKES that as given and starts at C.
%   - THIS section: the algebra -- the doubling construction, the
%     ledger of losses-as-physics, the internal-algebras-not-a-new-QM
%     scoping (Adler pre-empt), the two appearances of the ladder
%     (register / measurement climb, stated lightly), and Hurwitz as
%     the algebraic termination with the committed zero-divisor
%     reading.
%   - sec:hopf-justification: the geometry (Adams, parallelisability,
%     the fibrations), the join of the two classifications, the
%     entanglement-sensitivity literature, the C-direction honest debt,
%     and the epistemic reading of termination.
%
%  COMMITTED READINGS USED (sources re-read 2026-07-03):
%   - Loss ledger and level dictionary: cayley_dickson_qubits_for_review
%     sec.5 ("each step loses a symmetry, and each lost symmetry is
%     identified with a physical content").
%   - Termination reading: grounding_argument_gauge_group ("the
%     sedenions have zero divisors and cannot support consistent
%     information flow around loops") -- leads; norm-composition-as-
%     probability given as its amplitude-level restatement.
%   - Internal-not-base scoping: cayley_dickson_qubits_for_review Q2 and
%     sec.3.2 caveat (Adler objection; C-qubit treated geometrically).
%
%  HONEST DEBT: ladder proved closed at the top (Hurwitz), derived at
%  the floor. The rung-compulsion revisit-trigger FIRED 2026-07-06:
%  sec:self-reference and op:sigma-functor now exist, so the thesis is
%  minted as op:ladder-compulsion (with the CD-conjugation-as-sigma
%  candidate route), inheriting op:sigma-functor's status.
%
%  AUDITED 2026-07-06 (final pass, third stop): all four citations
%  used; all cross-refs live; CD facts, dims, alternativity, Hurwitz
%  statement verified; Adler subsection delivers everything the
%  register's compressed pointer promises; file pure ASCII; no bare
%  "frame". Fixes: opening rewritten for both predecessors (continuity
%  was inserted between after drafting); "associator's non-vanishing is
%  mass itself" recalibrated to the register's debt-against-selection
%  defusal; the O->S ledger row and Hurwitz subsection reconciled with
%  the fourth-qubit forward remark ("the terminus closes the ladder of
%  forces; what survives past it is, at most, flavour").
%
%  CROSS-REFS: sec:self-reference [sibling dangle -- resolves when that
%  section is drafted], sec:hopf-justification, sec:hopf-information-bound,
%  sec:register, sec:entanglement-classes, sec:generations, sec:masses,
%  sec:wigners-friend. All others exist.
%  CITATIONS: existing keys baez2002octonions, hurwitz1898composition,
%  furey2018three. ONE new key: adler1995quaternionic (RIS in
%  division_algebras_refs_2026-07-03.ris).
%  British English; \autocite; no bullets in prose; the ledger is a
%  table because the comparison is the argument.
% =====================================================================

\section{The Division-Algebra Ladder}
\label{sec:division-algebras}

The foundations behind us end at amplitudes twice over: natively, where the self-referential
argument forced the description of a system into the system's own state space and priced
consistency at a number system whose parts can cancel --- complex amplitudes, and the Born rule
that squares them (\S\ref{sec:self-reference}) --- and independently, where the regime's
continuity delivers the same field through Hardy's theorem (\S\ref{sec:continuity}), the
convergence of the two branches being evidence rather than accident. This section asks the question
both routes leave open. Why
should the number system stop at $\mathbb{C}$? Ordinary quantum mechanics stops there by postulate;
a framework that has forbidden itself postulates must answer, and the answer is that the climb
neither stops at $\mathbb{C}$ nor continues forever. There are exactly two more rungs, each
purchased by surrendering an algebraic law and each paying for itself in physics, and the ladder is
closed at the top by a theorem of 1898. This section sets out the algebra of that ladder: the single
construction that climbs it, the ledger of what each rung costs and buys, and the theorem that ends
it. The geometry of the same ladder --- that each rung acts on a sphere, that exactly those spheres
are parallelisable, and that the algebraic and geometric classifications are one --- is the business
of the geometry section (\S\ref{sec:hopf-justification}).

\subsection{One construction climbs the whole ladder}

The rungs are not a miscellany; a single construction generates them all. Given an algebra with a
conjugation, the Cayley--Dickson construction builds its double on ordered pairs, with a twisted
multiplication in which the conjugate enters and a conjugation $(a,b)^{*} = (a^{*}, -b)$
\autocite{baez2002octonions}. (Conventions differ in the placement of the conjugations; all yield
isomorphic algebras.) Applied repeatedly it produces the sequence
$\mathbb{R} \to \mathbb{C} \to \mathbb{H} \to \mathbb{O} \to \mathbb{S}$, of real dimensions
$1, 2, 4, 8, 16$: the reals, the complex numbers, the quaternions, the octonions, the sedenions.
The ground floor is classical: $\mathbb{R}$ is ordered and positive-definite, the arithmetic of
classical probability. The first doubling is therefore the quantum move itself --- the passage from
probabilities to amplitudes that can interfere --- and it is not made here by choice: it is the
derivation of \S\ref{sec:self-reference}, already behind us. What this section must account for is
the rest of the climb.

\subsection{The ledger}

In mathematics the doublings are usually told as a story of degeneration: each algebra is worse
behaved than the last, and the sequence is studied for where it breaks. This framework reads the
same ledger in the opposite direction. Each rung is bought by surrendering a law, and each
surrendered law is exactly a piece of physics --- the losses are the assets, and the comparison is
sharp enough to be the argument:

\begin{center}
\begin{tabular}{llll}
\hline
doubling & dimension & what is lost & what the loss becomes \\
\hline
$\mathbb{R} \to \mathbb{C}$ & $1 \to 2$ & the order of the reals & phase and interference; amplitudes \\
 & & & that cancel; the Born rule (\S\ref{sec:self-reference}) \\
$\mathbb{C} \to \mathbb{H}$ & $2 \to 4$ & commutativity & non-commuting actions: weak isospin \\
 & & & $SU(2)$, the $B$ register (\S\ref{sec:register}) \\
$\mathbb{H} \to \mathbb{O}$ & $4 \to 8$ & associativity & the associator: colour, chirality, the \\
 & & & generations, and mass as debt \\
 & & & (\S\ref{sec:generations}, \S\ref{sec:masses}) \\
$\mathbb{O} \to \mathbb{S}$ & $8 \to 16$ & division itself & no further force: zero divisors appear \\
 & & (zero divisors) & and the ladder of sectors ends (below) \\
\hline
\end{tabular}
\end{center}

Walking the rows: the order lost at the first doubling is the possibility of interference --- parts
of a probability that can oppose one another --- which is where the derivation of \S\ref{sec:self-reference}
lives. The commutativity lost at the second is the existence of actions whose order matters; the
unit quaternions form the group $SU(2)$, and the framework houses weak isospin there, as the
internal algebra of the $B$ register whose order-sensitive, non-commuting dynamics the interaction sections develop
(\S\ref{sec:register}, \S\ref{sec:entanglement-classes}). The associativity lost at the third ---
the octonions retain the weaker property of alternativity \autocite{baez2002octonions} --- creates
the associator $[a,b,c] = (ab)c - a(bc)$, and the associator is where most of this paper's physics
is transacted: the stabiliser of a preferred complex direction inside the octonions' automorphisms
is the colour group (\S\ref{sec:hopf-justification}); the triality structure and its $K_3$
obstruction generate the three generations and the sterile channel (\S\ref{sec:generations}); and
the associator is the currency of mass --- the debt against the selected direction that
\S\ref{sec:masses} prices, owed by whatever action moves the selection and not by residence at
the rung, which is the defusal \S\ref{sec:register} makes explicit. The algebraic facts
in the third column are theorems; the identifications in the fourth are established in the sections
cited, each carrying its own epistemic flag there, and are not re-litigated here. What this table
asserts on its own account is only the correspondence's shape: one law surrendered per rung, one
sector of physics gained per law.

\subsection{Internal algebras, not a new quantum mechanics}

A reader versed in the quaternionic-quantum-mechanics literature will at this point raise a fair
objection: quantum theory over $\mathbb{H}$ as a foundation is known to be troubled --- the
tensor-product and composite-system difficulties set out by Adler and others
\autocite{adler1995quaternionic} --- and quantum theory over $\mathbb{O}$ founders earlier still,
since a non-associative algebra does not furnish a Hilbert space at all. The framework's position
must therefore be stated exactly, because it does not make the claim the objection targets. The base
of the physical Hilbert space here is $\mathbb{C}$ and remains $\mathbb{C}$; the higher algebras
enter as the \emph{internal} algebras of the register --- $\mathbb{H}$ for the $B$ qubit,
$\mathbb{O}$ for the $C$ qubit (\S\ref{sec:register}) --- in a tensor structure whose outer theory is
ordinary complex quantum mechanics. This is the same move by which the division algebras enter the
constructions of Furey and the related programmes without any proposal to rebase quantum mechanics
itself \autocite{furey2018three}. For the octonionic register the point is sharper still:
$\mathbb{O}^{2}$ is not a Hilbert space, and the framework accordingly treats the $C$ qubit not in
the Dirac--von Neumann manner but geometrically, as the sphere $S^{15}$ with its Hopf structure ---
which is precisely the form in which \S\ref{sec:hopf-justification} needs it. Nothing in this paper therefore
stands or falls with quaternionic or octonionic quantum mechanics as foundations; the ladder is
climbed by the registers' internal algebras, not by the theory's base field.

\subsection{The same ladder, twice}

The ladder appears in this framework in two guises, and it will help the reader to have them
distinguished now. Statically, it is the anatomy of a single relation: the three qubits of the
register take their values one rung apart --- $A$ over $\mathbb{C}$, $B$ over $\mathbb{H}$, $C$ over
$\mathbb{O}$ --- so that every particle carries the whole ladder as its internal structure
(\S\ref{sec:register}). Dynamically, it is the shape of measurement: a single qubit's pure states
fill $S^{3}$, an entangling measurement carries the joint state into $S^{7}$, and a third party into
$S^{15}$ --- the correspondence between qubit entanglement and the Hopf fibrations that
\S\ref{sec:hopf-justification} takes from a mature literature --- so that each act of measurement ascends one rung --- the climb the reader has already made,
narrated in \S\ref{sec:wigners-friend} --- and every completed relation is a terminated climb, which is
the reading \S\ref{sec:hopf-information-bound} gives the tower's endpoint. One ladder, seen twice:
as the anatomy of a relation and as the dynamics of relating. The framework regards this double
appearance not as a coincidence but as the same structure viewed statically and in motion, and the
sections cited carry that identification's weight.

\subsection{Hurwitz closes the ladder}

The construction itself would double forever; what stops it is a classification theorem. By
Hurwitz's theorem, the normed division algebras --- those in which the norm composes,
$\lvert ab\rvert = \lvert a\rvert\,\lvert b\rvert$, and no non-zero elements multiply to zero ---
are exactly $\mathbb{R}$, $\mathbb{C}$, $\mathbb{H}$ and $\mathbb{O}$
\autocite{hurwitz1898composition,baez2002octonions}. The next doubling, the sedenions, has zero
divisors: non-zero elements whose product vanishes. The framework's reading of this failure is the
one that matters physically. Zero divisors cannot support consistent information flow around loops
--- a cycle of compositions can silently annihilate what it carries --- and a framework in which
closed computational paths are load-bearing (\S\ref{sec:hopf-justification},
\S\ref{sec:space-of-computations}) cannot be built on an algebra that permits it. At the level of
amplitudes the same failure reads: the composition law
$\lvert ab\rvert = \lvert a\rvert\,\lvert b\rvert$, which carries probability multiplicatively
through every rung of the ladder, fails beyond the octonions, so that two non-vanishing amplitudes
could compose to nothing and probability would be destroyed by mere composition. The last algebra
fit to carry amplitudes is $\mathbb{O}$. One rung past the terminus is not entirely silent,
however: the sedenions' internal structure --- their automorphism group and their octonion
subalgebras --- is the recorded candidate arena for tagging the generations
(\S\ref{sec:interactions}), and it enters exactly as this section's scoping permits, internally
and never as a base. The terminus closes the ladder of forces; what survives past it is, at most,
flavour. This is the algebraic half of the termination; the
geometric half --- that no sphere above $S^{7}$ is parallelisable --- and the theorem joining the
two halves are the opening business of \S\ref{sec:hopf-justification}, and the epistemic reading of the same endpoint,
as a bound on the information a single relation can carry, is given in
\S\ref{sec:hopf-information-bound}.

\subsection{The honest debt}

As in the geometry section, we close by stating plainly what is proved, what is established, and what is
owed. Proved: the construction, the losses, and the termination are theorems --- the Cayley--Dickson
doubling and its ledger of surrendered laws are standard algebra, and Hurwitz's classification is
among the durable results of the subject. Established: the physical cashing of each rung --- the
Born rule at $\mathbb{C}$, weak isospin at $\mathbb{H}$, colour, generations and mass at
$\mathbb{O}$ --- is carried out in the sections cited above, each under its own epistemic flag. Owed,
and stated without disguise: the ladder as presented here is closed from above and derived at its
floor, but it is not yet derived rung by rung from below. Hurwitz forbids a further rung; the
argument of \S\ref{sec:self-reference} forces the first doubling, from the ordered arithmetic of
classical probability to amplitudes that interfere. That the same self-referential closure, applied
to a system whose amplitudes already take values at one rung, compels the next --- that relational
closure at one level forces precisely the Cayley--Dickson double of that level --- is the
framework's reading of the climb, and it is a thesis, not a theorem --- but the formal apparatus
against which it can be posed now exists, and we pose it.

\begin{openproblem}[The compulsion of the intermediate rungs]\label{op:ladder-compulsion}
Derive, within the categorical apparatus of Open Problem~\ref{op:sigma-functor}, that
self-referential closure at one rung compels the Cayley--Dickson double of that rung, so that the
intermediate steps $\mathbb{C} \to \mathbb{H} \to \mathbb{O}$ follow from the same principle
that forces the floor. The candidate route is the construction's own shape: the Cayley--Dickson
pair with its conjugation, $(a, b)^{*} = (a^{*}, -b)$, is an involution-equipped doubling, and the
route reads it as the algebraic shadow of adjoining the describer --- closure of level-$A$
descriptions under the inside/outside exchange requiring exactly the twisted pair. The problem
inherits the status of Open Problem~\ref{op:sigma-functor}, whose apparatus it presupposes.
\end{openproblem}

If the derivation lands, the whole ladder becomes a consequence of the same self-reference that
opened it. Until then, the honest statement is that the ladder's necessity is established at its
ceiling and at its floor, while the compulsion of the intermediate rungs is inherited from the
physics they are shown to carry rather than derived in advance. What this section hands on is, accordingly, an algebraic ladder with a floor, a ceiling and a
ledger. Its first consumer is immediate: \S\ref{sec:space-of-computations} spends the terminal
fact --- the non-associativity of $S^{7}$'s own multiplication --- as geometry, excluding the
seven-sphere as a ruler and fixing the dimension of space at three.
\S\ref{sec:hopf-justification} then supplies the geometry of the fibrations themselves, and the
join.

%      % EDITED 2026-09-09 scripts/patch_spekkens_citekey_2026-09-09T1810.py
% =====================================================================
%  SECTION: Why Hopf Fibrations Belong in Particle Physics
%  (sec:hopf-justification). Drafted 2026-06-08; F1 arc added
%  prop:f1-dictionary and op:hopf-frobenius's discharge clause.
%  AUDITED 2026-07-06 (final pass, eighth stop).
%
%  PLACEMENT (reordered chain, JS 2026-07-06): position 8 --
%      \input{spacetime_from_non-computability}
%      \input{hopf_justification_section_2026-06-08}
%      \input{entanglement_classes_and_gauge_groups_2026-06-08}
%  No adjacency phrasing in the prose: position-portable. All
%  cross-refs live -- sec:fanout-direction DRAFTED 2026-07-06 (formerly the one undrafted
%  section, by design).
%
%  AUDIT RECORD: prop:gsm, prop:termination and prop:f1-dictionary
%  verified against every quoting consumer (statements canonical);
%  honest-debt subsection names both the selection assumption and the
%  mereology debt; the epistemic two-climbs disambiguation and the
%  compositional rung-evidence verified; sec:canonical-orbits live;
%  file pure ASCII. Fixes: op:hopf-frobenius's cross-link with
%  op:sigma-functor made TWO-WAY (the "four ways" line now names the
%  op); both remaining functor conditions ANCHORED --- the quaternionic
%  to the established fibre physics of sec:wf-reconciliation (residue:
%  clutching map vs strong complementarity), the octonionic to the
%  proved canonical-orbit images (umbrella label app:ghz-w-proofs
%  minted at the appendix for the purpose). Statuses unchanged.
%  British English; \autocite; no bullets in prose.
% =====================================================================

\section{Why Hopf Fibrations Belong in Particle Physics}
\label{sec:hopf-justification}

Invoking Hopf fibrations in a discussion of the Standard Model invites an immediate and fair
objection: that the move is decorative, an analogy between two pieces of attractive mathematics
with no load-bearing connection between them. The purpose of this section is to dispose of that
objection before the construction is put to work. The connection is not an analogy. It is
carried by two independent and independently published theorems, and the structures it relies
upon are not selected for convenience but compelled by a pair of classification results. We take
as already established the correspondence between qubit entanglement and the Hopf fibrations
themselves --- the subject of a mature literature
\autocite{mosseri2001hopf,bernevigchen2003geometry} --- and do not reproduce it here. What this
section addresses is the prior question of why that correspondence is entitled to carry physical
content at all.

\subsection{Two theorems, joined end to end}

The bridge from qubit entanglement to the Standard Model is built from two distinct results,
neither of which is original to the present framework and neither of which depends on the other.

The first is geometric. The unit sphere in the Hilbert space of $n$ qubits is the total space of
a Hopf fibration: $S^3 \to S^2$ for one qubit, $S^7 \to S^4$ for two, and $S^{15} \to S^8$ for
three. For two and three qubits the fibration is \emph{entanglement-sensitive} --- separability
is precisely the condition under which the base degenerates to a product of single-qubit Bloch
spheres and the additional dimensions associated with the higher imaginary directions collapse
\autocite{mosseri2001hopf,bernevigchen2003geometry}. This is a statement about the geometry of
state spaces and requires no physical interpretation whatever.

The second is algebraic. Singling out a preferred complex structure on the octonions ---
equivalently, fixing the split $\mathbb{O} \cong \mathbb{C} \oplus \mathbb{C}^3$ --- leaves a
residual symmetry, and that residual symmetry is exactly the gauge group of the Standard Model.
This was established through the exceptional Jordan algebra by Dubois-Violette and Todorov
\autocite{duboisviolette2016exceptional,duboisviolettetodorov2019exceptional,todorovduboisviolette2018deducing}
and given its sharpest spinorial form by Krasnov \autocite{krasnov2021so9}, who showed that the
relevant group is the stabiliser of a complex structure inside $\mathrm{Spin}(9)$. We record the
result in the form in which it enters the present work.

\begin{proposition}[Dubois-Violette--Todorov; Krasnov]\label{prop:gsm} Let $J_R$ be the complex
structure on the space $\mathbb{O}^2$ of $\mathrm{Spin}(9)$ spinors determined by a choice of
unit imaginary octonion, equivalently the split
$\mathbb{O} \cong \mathbb{C} \oplus \mathbb{C}^3$. The subgroup of $\mathrm{Spin}(9)$ commuting
with $J_R$ is \[ G_{\mathrm{SM}} \;=\; \frac{SU(3) \times SU(2) \times U(1)}{\mathbb{Z}_6}, \]
acting on the singled-out $\mathbb{C}$ summand as a colour singlet of hypercharge $-1$ and on
the $\mathbb{C}^3$ summand as a colour triplet of hypercharge $\tfrac{1}{3}$, both as $SU(2)$
doublets --- the quantum numbers of one left-handed generation.
\end{proposition}

The two results are independent, and the join between them is the only place a physical theory
must do work. Each result is a theorem; what neither supplies is a reason why nature should fix
one complex direction within $\mathbb{O}$ rather than treating all directions symmetrically. We
return to this --- the genuine load-bearing assumption --- in \S\ref{sec:hopf-honest-debt}, and
resolve it, in candidate form, in \S\ref{sec:fanout-direction}.

\subsection{The ladder is forced, not chosen}

It might still be objected that selecting the octonions, or the Hopf fibrations, is itself a
modelling choice carrying as much latitude as the choice of a grand-unified group. It is not,
and the reason is a pair of classification theorems that close off the alternatives entirely.

By Hurwitz's theorem the only normed division algebras are $\mathbb{R}$, $\mathbb{C}$,
$\mathbb{H}$ and $\mathbb{O}$ \autocite{hurwitz1898composition,baez2002octonions}. By Adams's
resolution of the Hopf-invariant-one problem --- together with the parallelisability results of
Bott--Milnor and Kervaire --- the only parallelisable spheres are $S^0$, $S^1$, $S^3$ and $S^7$,
so that there exist exactly three nontrivial Hopf fibrations, in bijection with $\mathbb{C}$,
$\mathbb{H}$ and $\mathbb{O}$
\autocite{adams1960hopf,bottmilnor1958parallelizability,kervaire1958nonparallelizability}. These
are not heuristics; they are among the more durable theorems in topology and algebra. Their
force here is that, once a framework is committed to amplitudes valued in a normed division
algebra --- a commitment that in the present construction is made far upstream, by the
self-referential argument of \S\ref{sec:self-reference} and \S\ref{sec:division-algebras}, and
not by fiat at this stage --- the ladder of Hopf fibrations is not selected but compelled, and
it terminates of its own accord.

\begin{proposition}[Termination]\label{prop:termination} The Cayley--Dickson tower of normed
division algebras and the tower of parallelisable spheres both terminate at real dimension
eight. Consequently there are exactly three nontrivial Hopf fibrations, the largest being the
octonionic $S^{15} \to S^8$, and the construction admits no extension to a fourth case.
\end{proposition}

The termination is worth dwelling on, because it is unusual for a model of this kind to predict
its own boundary. The next rung of the Cayley--Dickson sequence, the sedenions, possesses zero
divisors; the Hopf map can no longer be globally defined; there is no fourth fibration. This is
not a limitation imposed from outside but a structural fact, and it carries a physical
consequence that we develop later: there is no room in the construction for a fourth gauge force
or a fifth interaction, so that the discovery of either would count as evidence against the
framework rather than as an invitation to extend it. The endpoint is a theorem, not a
stipulation.

\subsection{Independent arrival as corroboration}

The most persuasive evidence that the destination is not an artefact of the route is that it is
reached by routes sharing none of the same machinery. Szangolies
\autocite{szangolies2025standardmodel} reaches $G_{\mathrm{SM}}$ from the emergent-gravity and
Kaluza--Klein programmes, reading the singled-out complex direction as a dimensional reduction
of a $9+1$-dimensional spacetime; the present framework reaches it from self-reference
(\S\ref{sec:self-reference}); and Furey and Gillard--Gresnigt reach the same gauge group, and
the same three-generation structure, through the minimal left ideals of the Dixon and Clifford
algebras, without invoking Hopf fibrations at all
\autocite{furey2018three,gillardgresnigt2019three}; the relation between their minimal ideals
and the entanglement classes used here is taken up in \S\ref{sec:canonical-orbits}. The
octonionic route to particle physics is in any case an old one, going back to the use of
$\mathbb{O}$ to model the strong force \autocite{gunaydingursey1974quark}.

That programmes with disjoint starting points converge on
$SU(3) \times SU(2) \times U(1)/\mathbb{Z}_6$ and on the same quark--lepton division is the
coincidence Szangolies himself remarks upon: the characteristic structure of the Standard Model,
routinely regarded as arbitrary, turns up independently in unrelated settings, and its appearing
twice would then be doubly odd unless the two settings are in fact related. We read the
convergence the same way --- as evidence that the division-algebraic description is singled out
by the mathematics rather than imposed by the modeller, and therefore that
Proposition~\ref{prop:gsm} is structural rather than contrived.

\subsection{The honest debt}
\label{sec:hopf-honest-debt}

It would be misleading to leave the impression that this route is free of assumptions, so we
state plainly the one assumption that all of these programmes share. Proposition~\ref{prop:gsm}
requires that \emph{some} complex direction in $\mathbb{O}$ be singled out, and the uniqueness
of that choice is not supplied by the theorem. Krasnov parametrises the complex structure $J_R$
by an arbitrary unit imaginary octonion, and the residual group is the same for every such
choice because $G_2 = \mathrm{Aut}(\mathbb{O})$ acts transitively on the unit imaginary
octonions. The theorem therefore guarantees the gauge group \emph{given} that a direction is
fixed, while remaining silent on \emph{why} one is fixed. Szangolies treats this as a free
choice and flags it, in his own conclusions, as an open element of the construction --- one to
which several inequivalent answers are available.

This, and not the appeal to Hopf fibrations, is the genuine load-bearing assumption of the whole
approach. The Hopf fibrations are forced (Proposition~\ref{prop:termination}); the singling-out
is not. The present framework offers a candidate resolution that the prior literature lacks: the
act of measurement, modelled as \textsc{fanout}, selects the direction physically, and the
transitivity of $G_2$ then guarantees that the gauge group is independent of which direction the
measurement happens to select. This argument is given in \S\ref{sec:fanout-direction}; we flag
here, at the appropriate epistemic status, that it is structural rather than fully proved, and
that the residual question --- why a measurement selects \emph{exactly one} imaginary direction
rather than merely \emph{permitting} a selection --- remains open. It is now priced as Open
Problem~\ref{op:fanout-selection}.

A second debt is shared with every programme of this character and is not closed by anything
said here: why the abstract quantum state should admit a tensor-product decomposition into
qubits in the first place. This is the quantum factorisation, or mereology, problem. The
entanglement-sensitivity on which the construction turns is meaningful only once such a
decomposition is fixed, and we treat the decomposition as a necessary input rather than a
derived result. Naming it is preferable to letting it pass unremarked.

\subsection{The termination read epistemically}
\label{sec:hopf-information-bound}

The termination of the tower (Proposition~\ref{prop:termination}) has so far been stated
algebraically: the Cayley--Dickson sequence acquires zero divisors at the sedenions and the Hopf
map can no longer be globally defined. There is a second reading of the same boundary, epistemic
rather than algebraic, which we offer as a motivated conjecture rather than a settled result,
because it connects the tower's endpoint to a bound that arises independently in the foundations
of measurement.

A particle, in this framework, is a reified relation between events, and its internal content is
carried by the octonionic $A,B,C$ register --- the $S^{7}$ sector serving as internal degrees of
freedom in the channel that connects the events (\S\ref{sec:space-of-computations}). The
termination of the tower is then a statement about that content: a single such relation can
support genuine entanglement among at most three parties, because there is no fourth fibration
to carry a fourth independent party within the same irreducible relation. Every particle is, on
this reading, a \emph{completed} octonionic-level relation --- topped out at three parties not
by circumstance but because the tower admits no further rung. The information that can be
encoded in one relation-between-events, and hence extracted by one measurement of it, is
correspondingly bounded: it is the information of a completed three-party register, and no
measurement of a single particle can yield more.

This is exactly the kind of bound that the epistemic-restriction programmes reach from the other
direction. Spekkens' toy theory is built on a knowledge-balance principle that caps the
information obtainable about a system \autocite{spekkens2007toy}, and Szangolies derives such an
epistemic horizon not by positing it but from the same self-referential non-computability on
which the present tower rests \autocite{szangolies_j_2018}. That the two accounts --- the
algebraic termination of the Hopf tower and the epistemic bound on measurement information ---
should coincide is, we suggest, not a numerical accident but a shared root: both are
consequences of the self-referential non-computability of \S\ref{sec:self-reference}, wearing an
algebraic face in the one case and an epistemic face in the other. The Hurwitz--Adams ceiling
and the knowledge-balance ceiling would then be one limit seen twice, and Szangolies' programme,
which reaches $G_{\mathrm{SM}}$ by an independent route (\S\ref{sec:hopf-honest-debt}) and the
epistemic horizon from the same non-computable source, shares with this framework not merely a
conclusion but its origin.

The reading must be kept distinct from a different climb, with which it would otherwise be
conflated. The Wigner's-friend regress of \S\ref{sec:wigners-friend} ascends without end: each
observer may be measured by a further observer, each carrying a fresh gauge-private circle, and
that chain has no top --- it is the bump in the carpet, unbounded in length. The termination
bounds something else: not the length of the chain of observers, but the internal structure of
each single relation, each particle, that the chain is built from. The regress is unbounded
precisely because each of its links is a bounded, completed three-party object;
termination-of-each-link and non-termination-of-the-chain are not in tension but complementary,
the former individuating the links of the latter. What the bound caps is the information in one
particle; what the regress leaves uncapped is the number of particles and observers that can be
chained. The precise identification of ``the information of a completed three-party register''
with the formal measures of the epistemic-restriction literature is the technical step that
would turn this conjecture into a theorem, and we flag it as owed; the structural claim --- that
the tower's endpoint and the measurement bound share the non-computable root --- is what we
advance here.

\subsection{The termination read compositionally}
\label{sec:hopf-compositional}

There is a third termination theorem, proved in a community that has so far had no contact with
the two joined above. Coecke and Kissinger showed that the GHZ and W classes --- realised as the
commutative special and the anti-special Frobenius algebras of categorical quantum mechanics ---
are \emph{compositionally complete}: every $N$-qubit entanglement class under SLOCC is generated
by composing these tripartite primitives with bipartite entanglement and product states
\autocite{coeckekissinger2010compositional}. Above three parties nothing new appears; the
entanglement alphabet closes at exactly the point where the fibration tower terminates. The
mutual silence of the two literatures is itself striking: Coecke and Kissinger make no reference
to Hopf fibrations or division algebras, and the fibration literature on entanglement
\autocite{mosseri2001hopf,bernevigchen2003geometry} nowhere cites compositional completeness.
Two classifications, in different mathematics, draw the same line in the same place, unaware of
one another. That silence has, in the months this paper was in draft, just begun to break from
the categorical side, with a projective, base-level reading of the same calculus
\autocite{burtonanwar2026meromorphic}; we engage it below.

The framework's conjecture is that this is one termination wearing a third face: the alphabet
has the letters it has --- Bell, GHZ, W --- because the parallelisable spheres have the
dimensions they have. There is partial evidence at every rung. At one qubit, the commutative
special Frobenius algebras on $\mathbb{C}^{2}$ are classified by orthonormal bases --- points of
$S^{2}$, the base of the first fibration --- so that the copied basis plays the role of a
section of the bundle and the spider theorem expresses its idempotence. At two qubits, the
$SU(2)$ fibre of $S^{7} \to S^{4}$ carries exactly the entanglement that survives local
operations, and L\'{e}vay's result that the bundle's natural connection is the $SU(2)$
Yang--Mills instanton \autocite{levay2004geometry} stands on the geometric side where the
interaction of two Frobenius structures stands on the categorical one. At three qubits, the
split of D\"{u}r, Vidal and Cirac \autocite{durvidalcirac2000three} mirrors the algebra exactly:
GHZ is genuinely tripartite, demanding the full non-associative octonion, while W is
pairwise-robust and housed in the associative subalgebra $\mathbb{H} \subset \mathbb{O}$ --- the
special/anti-special Frobenius distinction in geometric dress.

The first of the four conditions the correspondence requires can be discharged at once, and we
do so here, because a problem whose first step has been walked is a different kind of invitation
from one merely posed.

\begin{proposition}[The one-qubit dictionary]\label{prop:f1-dictionary} The $\dagger$-special
commutative Frobenius algebras on $\mathbb{C}^{2}$ correspond bijectively to
\emph{lifted antipodal skeletons} of the complex Hopf fibration $\pi \colon S^{3} \to S^{2}$:
unordered antipodal pairs $P \subset S^{2}$ together with a section of $\pi$ over $P$. The
copyable states of a structure are exactly the section's image, so the lifts are recorded
faithfully --- the projectively trivial centre of $U(2)$ acts non-trivially on structures ---
and the correspondence is $U(2)$-equivariant, the ordered structures forming a single free and
transitive $U(2)$-orbit. The channel a structure induces on density matrices is dephasing in its
skeleton and is independent of the section: the fibre freedom, a torsor for
$\mathrm{Map}(P, U(1))$, is exactly the data absent from the classical record a structure
induces.
\end{proposition}

\begin{proof}
By the classification of Coecke, Pavlovic and Vicary, such structures correspond to orthonormal
bases, recovered as the sets of copyable states \autocite{coeckepavlovicvicary2013orthogonal}.
An orthonormal basis of $\mathbb{C}^{2}$ is two orthogonal rays --- an antipodal pair on the
Bloch sphere --- together with a unit representative of each, which is precisely a section of
$\pi$ over the pair; and $\delta(e^{i\theta}x) = e^{i\theta}\, x \otimes x$ against
$(e^{i\theta}x) \otimes (e^{i\theta}x) = e^{2i\theta}\, x \otimes x$ pins the phases.
Equivariance, the action of the centre, the orbit structure and the induced channel are direct
computations; an executable verification of every finite check accompanies the code release.
\end{proof}

Two remarks fix the proposition's scope. It corrects a natural first reading: a classical
structure does not project to the base but selects lifts into the total space --- classicality
is section data, and what the copy operation broadcasts is a lift, not a ray. And it claims
nothing topological about cloning: the finiteness of the copyable set is linear, structural
no-cloning residing in non-naturality \autocite{abramskyheunen2012hstar}, so the first place the
bundle's global topology can bite is the quaternionic rung, where the degree-one clutching map
stands ready to be compared with the strong-complementarity law governing an interacting pair of
structures. The dictionary is, finally, the complement of the projective treatment cited above:
Burton and Anwar pass functorially to the base and deliberately discard the fibre
\autocite{burtonanwar2026meromorphic}, and the classical structures turn out to live in exactly
the data their projectivisation forgets --- the two readings compose, and the resulting triangle
is a consistency condition available to constrain the remaining conditions of the
correspondence.

\begin{openproblem}[The Hopf-fibration--Frobenius correspondence]\label{op:hopf-frobenius} Let
$\mathbf{Hopf}$ denote the category of the three fibrations
$S^{2^{k}-1} \hookrightarrow S^{2^{k+1}-1} \to S^{2^{k}}$, $k = 1, 2, 3$, with bundle maps as
morphisms, and let $\mathbf{Frob}$ denote the Frobenius algebras on $\mathbb{C}^{2}$ in
$\mathbf{FdHilb}$, with the commutative-special (GHZ) and anti-special (W) classes
distinguished. Construct a functor $F \colon \mathbf{Hopf} \to \mathbf{Frob}$ --- or prove that
none exists --- under which the complex fibration yields the single-qubit classical structures
--- Proposition~\ref{prop:f1-dictionary} identifies these as the lifted antipodal skeletons,
section data over two base points rather than the points alone; the quaternionic fibration
yields the interaction of two such structures, with the $SU(2)$ fibre carrying the entanglement
that survives the classical quotient; the octonionic fibration yields the GHZ/W pair, GHZ on the
full non-associative fibre and W on the associative $\mathbb{H} \subset \mathbb{O}$; and the
composition by which Coecke--Kissinger generate the $N$-qubit classes corresponds to a
well-defined geometric operation on the bundles. If such an $F$ exists, compositional
completeness follows from Adams's theorem as a topological corollary --- there is no fourth
fibration to supply a fourth primitive --- and if $F$ is faithful enough that the implication
also runs back, the termination of the tower and the completeness of the alphabet are equivalent
expressions of a single classification. The first condition is discharged by
Proposition~\ref{prop:f1-dictionary}; what remains open is the quaternionic and octonionic
conditions and, above all, the correspondence of compositions. Both open conditions are already
constrained: the quaternionic fibre's physical content --- the recoverable relative phase and
its purity ledger --- is established (\S\ref{sec:wf-reconciliation}), leaving the
clutching-map-to-strong-complementarity comparison as the structural residue; and any candidate
functor must reproduce the proved canonical-orbit images, GHZ to the colour triplet and W to the
singlet (\S\ref{sec:canonical-orbits}, Appendix~\ref{app:ghz-w-proofs}). A proof that no such
functor exists, or a correspondence at the level of the classifications without functoriality on
compositions, would each also resolve the problem: the first would make the agreement of two
independent classifications a fact demanding its own explanation; the second would establish the
weaker reading.
\end{openproblem}

A resolution in the affirmative would do three things. It would explain compositional
completeness rather than merely prove it --- a corollary of the division-algebra structure of
quantum mechanics rather than an independent algebraic fact. It would give the diagrammatic
language of this paper's own interaction circuits a topological warrant: the GHZ/W spider
calculus would be the unique language the topology permits, not a convenient notation. And it
would open a two-way bridge, translating problems on compositions of Frobenius structures into
operations on bundles and back. The problem sits at the intersection of categorical quantum
mechanics, the quantum-information geometry of the fibrations, and the division-algebra
approaches to particle physics, and it is posed here so that any of the three communities can
take it up in its own language. It is perhaps worth noticing that both classifications already
answer to the same name --- the fibrations are Heinz Hopf's, and so, through the bialgebra
structure of the diagrammatic calculus, are the Hopf algebras --- a coincidence of eponymy that
the conjecture proposes to promote to a theorem. With the epistemic reading above, an
affirmative resolution would leave the tower's endpoint as one limit seen four ways: algebraic,
geometric, compositional, epistemic --- the last being the fixed-point classification of Open
Problem~\ref{op:sigma-functor}, with which this problem is cross-linked, each constraining the
other.

\subsection{What this sets up}

The entanglement-sensitivity of the Hopf maps is not a side feature but the hinge of the
physical reading that follows. Because separability is exactly the collapse of the base to a
product of Bloch spheres, the structure of the entanglement --- the SLOCC class of a multi-qubit
state, and within it the canonical maximally entangled orbit --- is carried by precisely the
degrees of freedom on which $G_{\mathrm{SM}}$ acts in Proposition~\ref{prop:gsm}. This is what
makes it possible, in \S\ref{sec:entanglement-classes}, to read the two genuinely tripartite
entanglement classes onto the $\mathbb{C} \oplus \mathbb{C}^3$ split through their canonical
representatives: the GHZ orbit confined to the colour-triplet $\mathbb{C}^3$ sector, the W orbit
reaching the $\mathbb{C}$ singlet, recovering the quark--lepton division not as an assignment
but as a consequence of the entanglement structure. The scope of those statements --- canonical
orbits rather than whole classes, and the identification of the canonical orbits with the
physical particles --- is made precise there. The Hopf fibrations are the bridge; the
entanglement classes, treated next, are the traffic that crosses it.

%
%  SOURCES (re-read 2026-07-04): lawvere_inside_outside_fixed_point_
%  theorem_2026-05-04T1806.md (the spine: the two-principles reduction,
%  the theorem, the sigma candidate, the open items -> op minted per
%  D9); research_summary_neg_prob.md (negativity route; lamp (1+i)/2);
%  dynamics_companion_fanout_unitarity (legs 1a/1c; anti-events; the
%  Brukner-Zeilinger rebuttal per D6); Pure_states_and_mixed_states
%  (local tomography derived; QBism contrast per D7).
%
%  DECISIONS ENCODED (JS 2026-07-04): D1 title as below; D2 Wheeler1990
%  cited here as the title anchor; D3 lineage Popper -> Breuer ->
%  Svozil -> Szangolies (Dalla Chiara omitted); D4 Frauchiger_2018
%  cited here, dissolved in sec:wigners-friend; D5 Kirkwood-Dirac
%  heritage via the Arvidsson-Shukur 2024 review as one-stop; D6
%  Brukner-Zeilinger + the framework rebuttal IN ("Shannon is
%  post-FANOUT"); D7 QBism contrast IN; D8 Aaronson private
%  communication OMITTED (published island paper only); D9 one op
%  minted (op:sigma-functor), cross-linked to op:hopf-frobenius.
%
%  CITATIONS -- existing keys: Wheeler1990, small_2006, Godel1931,
%  Lawvere_1969, Yanofsky2003, CoeckeKissinger2017, Rovelli1996a,
%  Popper1950a, Popper1950b, Breuer1995, svozil_k_2018,
%  szangolies_j_2018, Frauchiger_2018, dirac1942, Feynman_r_82,
%  Feynman1987, Youssef1994 [SWAP-FLAG: intended paper is "Quantum
%  mechanics as complex probability theory", Mod. Phys. Lett. A 9
%  (1994) 2571 -- if that lives under youssef1994quantum, swap],
%  caticha_a_2006, goyal_knuth_skilling_2010, Aaronson2004h, Hardy2001a,
%  muller_p_2011, Thomson_1954, benacerraf1962tasks, cantor1891diagonal,
%  Turing_1937.
%  NEW keys (RIS in self_reference_refs_2026-07-04.ris, all verified
%  this session): tikochinsky1988feynman, renou2021real, chen2022ruling,
%  wu2022refutation, hoffreumonwoods2026real, arvidssonshukur2024kirkwood,
%  bruknerzeilinger2001shannon, fuchsschack2013qbism,
%  shannon1948mathematical.
%
%  CROSS-REFS: forward (fine under two-pass LaTeX): sec:continuity,
%  sec:division-algebras, sec:hopf-justification, prop:f1-dictionary,
%  sec:wigners-friend, op:limit-realisability, op:hopf-frobenius.
%  Still dangling after this section: sec:fanout-boundary (and its
%  sibling label sec:fanout-direction) -- both DRAFTED 2026-07-06.
%
%  AUDITED 2026-07-06 (section-by-section final pass, first stop):
%  proof of prop:inside-outside re-verified by hand; attribution chain
%  clean (Burgin/Khrennikov correctly absent); all header-declared
%  citations confirmed used; lamp arithmetic and antecedents verified;
%  FANOUT smallcaps consistent; no pre-definition use of "frame".
%  Fixes: stray Unicode em-dash (QBism paragraph, the file's only
%  non-ASCII character) -> house ---; trace-blind two-way stitch added
%  to the Born subsection ("one privacy, two duties");
%  op:sigma-functor gains its first standing constraint (F1
%  inheritance via the op:hopf-frobenius discharge) and its previously
%  one-directional local-tomography cargo made explicit.
% =====================================================================

\section{Self-Reference: One Principle}
\label{sec:self-reference-2}

Wheeler asked for it from bit: for physics derived from information, ``every it --- every particle,
every field of force'' obtained from binary answers to yes-or-no questions \autocite{Wheeler1990}.
This paper's route to that programme runs through G\"{o}del: the single principle from which
everything else here is derived, deferred, or priced is \emph{self-reference} --- the fact that any
complete physical description, made from inside the world it describes, must describe the describer.
This section states the principle, proves that the framework's second historical commitment follows
from it as a theorem rather than standing beside it as an axiom, and then extracts the price of the
principle in four instalments: probabilities whose parts can cancel, amplitudes valued in
$\mathbb{C}$, the Born rule, and the local tomography of composite systems. What the section cannot
yet derive is stated as precisely as what it can, in one numbered problem and in the ledger at the
end.

\subsection{Events, not objects}

The framework's ontology is relational: the fundamental entities are events and the relations
between them, and objects are the bookkeeping of relational invariants --- what remains stable
across many relations. The kinship with relational quantum mechanics \autocite{Rovelli1996a} is
genuine and acknowledged: facts are facts relative to the systems that register them. The
difference, developed over the rest of this paper, is that the present framework grounds the
relational structure in a single logical principle and then carries it, through the division
algebras and the Hopf fibrations, all the way to the gauge group and the particle spectrum ---
relations are not only the ontology but the derivation engine. On this reading a measurement is not
the revelation of a pre-existing property; it is the creation of a relation, and the later sections
of this paper read that creation geometrically (\S\ref{sec:hopf-justification}) and even
gravitationally. Here we need only the logical core: whatever describes the world is an event in
the world it describes.

\subsection{One principle, and a theorem}

The framework historically carried two foundational commitments: self-reference, and the
\emph{Inside/Outside Equivalence Principle} --- that the laws of physics must be identical whether
formulated by an embedded participant or by an external observer. Two principles are one too many,
and this subsection reduces them.

The reduction corrects the framework's own origin, and we make the correction in print. The 2005
form of the argument \autocite{small_2006} ran a G\"{o}delian chain: a self-referential process,
being complete, must by the incompleteness theorem \autocite{Godel1931} be inconsistent; an
inconsistent system can negate itself; and the negation of self-reference is the production of
causal sequence. The chain captures the right intuition and is retained here as motivation, but it
is not sound as a proof: classical inconsistency yields \emph{ex falso quodlibet} --- every
proposition provable, every causal sequence an output, no constraint at all. The load-bearing
argument is different, and better: a fixed-point theorem.

Let $\mathcal{C}$ be the category of physical descriptions and admissible processes --- concretely,
the dagger-compact-closed setting with classical structures in which categorical quantum mechanics,
and the whole of this paper, works \autocite{CoeckeKissinger2017}. Lawvere showed that in a
sufficiently rich such setting, any weakly point-surjective self-representation forces every
endomap to have a fixed point, and Yanofsky showed that one diagonal argument of this shape
underlies Cantor's, G\"{o}del's, Turing's and Tarski's theorems alike
\autocite{Lawvere_1969,Yanofsky2003,cantor1891diagonal}. The framework's single commitment is that
$\mathcal{C}$ contains a \emph{self-representing object} --- an object capable of encoding the
processes on itself. That is what ``self-reference is foundational'' means, stated categorically,
and it is the surjectivity hypothesis of the diagonal lemma. Now let
$\sigma \colon \mathcal{C} \to \mathcal{C}$ be the \emph{inside/outside involution}: the operation
exchanging the two complementary presentations of any description --- the inside presentation, as
participant-indexed total space in which the observer is part of the system, and the outside
presentation, as a base-times-fibre decomposition over a classical parameter space. That
$\sigma^{2} = \mathrm{id}$ is not an axiom but a coherence requirement: the outside view of the
inside view of a thing is the thing. The candidate concrete realisation of $\sigma$ is the
Hopf-duality swap --- total space against base-times-fibre for the three fibrations --- which gives
the involution computable content and ties this subsection to \S\ref{sec:hopf-justification}.

\begin{proposition}[Inside/outside fixed point]\label{prop:inside-outside}
Let $\mathcal{C}$ satisfy the richness condition above and contain a self-representing object, and
let $\sigma$ be an involution on $\mathcal{C}$. Then the fixed-point locus
$\mathrm{Fix}(\sigma)$ is non-empty: there exist descriptions invariant under the exchange of
inside and outside perspective.
\end{proposition}

\begin{proof}
Lawvere's construction, applied to $\sigma$ in place of negation: given a weakly point-surjective
$\varphi \colon A \to Y^{A}$, form $g(a) = \sigma\!\left(\varphi(a)(a)\right)$; point-surjectivity
supplies $a^{*}$ with $\varphi(a^{*}) = g$, and then $\varphi(a^{*})(a^{*})$ is a fixed point of
$\sigma$ \autocite{Lawvere_1969,Yanofsky2003}. Where the same construction applied to Boolean
negation yields paradox --- negation has no fixed point on a two-element type, so the conclusion
explodes the premises --- the inside/outside involution is an honest involution on a rich type, and
the conclusion is a constructive fixed point rather than a contradiction.
\end{proof}

\textbf{Definition.} \emph{Physical law is the structure carried by} $\mathrm{Fix}(\sigma)$.

The Inside/Outside Equivalence Principle is thereby demoted from axiom to theorem: physical law is
invariant under the exchange of perspectives \emph{because it is the fixed-point set of that
exchange}, and the fixed-point set is forced to be non-empty by self-reference alone. One
commitment remains foundational; everything the second commitment asserted is now derived. The
reader who accepts categorical quantum mechanics as an existing programme has already accepted the
machinery; the framework adds one premise and collects the rest. And the fixed-point locus is not
empty of structure: it inherits the Frobenius-algebra apparatus of $\mathcal{C}$, so that the
Bell, GHZ and W primitives --- and with them the particle classification of the later sections ---
live inside $\mathrm{Fix}(\sigma)$ from the start. What is proved above is non-emptiness given the
setting; what is assumed is the setting itself and the existence of $\sigma$ as a genuine
involution; what is open is the rigour of both, and it is priced now rather than hidden:

\begin{openproblem}[The involution made functorial]\label{op:sigma-functor}
Fix the category $\mathcal{C}$ rigorously --- the dagger-compact-closed setting with classical
structures is the candidate --- and define $\sigma$ on morphisms, not only on objects, the natural
route being the induced fibre--base exchange on spider diagrams; establish that $\sigma$ is an (at
worst pseudo-)involutive endofunctor; and prove or refute the identification of
$\mathrm{Fix}(\sigma)$ with the three Hopf cross-sections, under which Adams's theorem would become
a corollary of the fixed-point classification. This problem is cross-linked with Open
Problem~\ref{op:hopf-frobenius}: under an affirmative resolution of both, the parallelisable
terminus, the fixed-point classification, and compositional completeness are three statements of
one fact, and progress on either problem constrains the other. One constraint is already in place:
at the complex rung, Proposition~\ref{prop:f1-dictionary} supplies the dictionary between
classical structures and the fibration's lifted skeletons --- the clause of
Open Problem~\ref{op:hopf-frobenius} it discharges --- and any resolution here must reproduce it.
The problem also carries a piece of cargo recorded later in this section: the formalisation, within
the fixed category, of the perspective-accounting that derives local tomography.
\end{openproblem}

\subsection{The lineage}

The obstruction this section builds on is old, and the framework claims its completion, not its
discovery. Popper argued in 1950 that a predictor embedded in the world has irreducibly inferior
knowledge of itself, because it cannot impart information about itself without, by that very act,
rendering the information out of date \autocite{Popper1950a,Popper1950b} --- the passage is quoted
where the framework makes the obstruction structural (\S\ref{sec:wigners-friend}). Breuer proved
the impossibility of accurate state self-measurement within quantum theory \autocite{Breuer1995};
Svozil's monograph surveys the whole territory where physics meets self-reference, undecidability
and Chaitin-style algorithmic limits \autocite{svozil_k_2018}; Szangolies sharpened the modern form
as epistemic horizons \autocite{szangolies_j_2018}; and the Frauchiger--Renner theorem is the
contemporary crisis-form --- nested observers deriving contradictory records when each treats the
others' quantum descriptions as universally valid \autocite{Frauchiger_2018} --- which
\S\ref{sec:wigners-friend} dissolves relationally rather than paradoxically. This framework's
addition to the lineage is twofold. Structurally, Proposition~\ref{prop:f1-dictionary} converts
Popper's temporal obsolescence into algebraic cancellation: the datum a copier cannot put into its
own record is exactly its fibre commitment, which does not arrive stale --- it cancels. And
foundationally, the present section runs the lineage in reverse: where the tradition treats
self-reference as a limit on knowledge \emph{within} physics, the framework takes it as the
generator of physics, and derives the quantum formalism from it.

\subsection{The price is negative probability}\label{sec:complex-amplitudes-2}

What does it cost a probabilistic description to contain its own describer? The first instalment is
the sign. In the regime before any record exists, every event has an anti-event --- the description
must be able to undo what it has provisionally done, or the self-description could never be
revised into consistency --- and a probability calculus closed under undoing cannot be confined to
$[0,1]$: it is a signed calculus, in which the parts of a probability can oppose and cancel. The
framework did not discover negativity as bookkeeping; Dirac admitted negative probabilities as
formal intermediaries \autocite{dirac1942}, Feynman defended them as consistent whenever the final,
recorded probabilities stay positive \autocite{Feynman_r_82,Feynman1987}, and the Kirkwood--Dirac
distribution has carried complex quasiprobability from 1933 to the present, as a recent
comprehensive review documents \autocite{arvidssonshukur2024kirkwood}. What the framework adds is
the forcing: negativity is not a permissible convenience but the first price of self-referential
closure, and it is charged exactly and only in the pre-record regime --- which is why every
negativity in the tradition lives in intermediate, unrecorded quantities and never in a final
record.

The second instalment is complexity. A signed information measure must extend Shannon's
$S = -\ln p$ off the positive axis, and the analytic continuation of the logarithm carries $p$ into
the complex plane: amplitudes, with the phase as the continued information. Here too the extension
has a prior tradition, and the attributions are owed precisely. Tikochinsky derived Feynman's rules
for complex amplitudes from consistency of composition \autocite{tikochinsky1988feynman}; Youssef
reformulated quantum mechanics as complex probability theory outright \autocite{Youssef1994};
Caticha carried objective amplitudes to Bayesian probabilities \autocite{caticha_a_2006}; and
Goyal, Knuth and Skilling derived Feynman's rules from a pair-postulate without assuming
analyticity \autocite{goyal_knuth_skilling_2010}. The extension is prior art. The framework's
claim is the compulsion: the complex extension is not a modelling choice among exotic probability
theories but the unique continuation forced by self-referential closure --- the same forcing that
Aaronson glimpsed when he argued quantum mechanics is what negative probability inevitably becomes
\autocite{Aaronson2004h}.

Why does the continuation stop at $\mathbb{C}$? Downward, the real field is not available: real
quantum theory is experimentally distinguishable from complex quantum theory in network scenarios
and has been disfavoured by experiment
\autocite{renou2021real,chen2022ruling,wu2022refutation}; the recent objection that the exclusion
rests on locally tomographic composition \autocite{hoffreumonwoods2026real} is, for this framework,
not an escape but a confirmation of where the weight sits --- local tomography is \emph{derived}
below, not assumed, so the framework closes from its own resources exactly the door the objection
holds open. Upward, the quaternions and octonions fail as a \emph{base} for the information
continuation: the continued measure is a logarithm, and beyond $\mathbb{C}$ the logarithm of a
product is no longer the sum of logarithms --- non-commutativity severs
$\log(ab) = \log a + \log b$ --- so the additivity that makes it an information measure is lost.
Aaronson's parameter-counting and signalling arguments against a quaternionic base point the same
way \autocite{Aaronson2004h}. The higher algebras are not banished; they return as the
\emph{internal} algebras of the register, which is the business of \S\ref{sec:division-algebras}.
The base is $\mathbb{C}$, for informational reasons, and stays $\mathbb{C}$.

\subsection{The Born rule from fibre privacy}\label{sec:born-rule}

The third instalment is the quadratic rule. An amplitude carries one datum more than any record of
it can: its overall phase. This was asserted throughout the framework's development as the
observer's private, self-referential ignorance; it is now a proved statement, in exactly the form
needed. By Proposition~\ref{prop:f1-dictionary}, a classical structure --- the copy operation any
record requires --- is committed to a specific lift of its skeleton, and that commitment is absent
from every classical record the structure induces: the copier is fibre-committed, the record is
fibre-blind, and the fibre freedom is gauge for every observer. (The computation extends verbatim
from one qubit to the full register, where Proposition~\ref{prop:trace-blind} makes it the ground
of hypercharge's tracelessness --- one privacy, two duties.) Popper's obsolescence, made
cancellation. It follows that anything recordable is invariant under the fibre --- observables and
their statistics must be $U(1)$-invariant functions of the amplitude. The lowest-degree polynomial invariant
of a complex amplitude is the quadratic $|\psi|^{2}$ --- the non-polynomial invariant $|\psi|$,
of degree one, is excluded by the analyticity declared at \S\ref{sec:complex-amplitudes} --- and
normalisation fixes the coefficient:
probability is the minimal record-compatible function of the amplitude,
\[
p \;=\; |\psi|^{2},
\]
the Born rule as the cheapest invariant the fibre permits. The same square is reached from the
cost side in \S\ref{sec:born-cost-ledger}, where a record is the closed loop of a history and its
return leg and the exponent is the number of legs; neither route uses the lattice of projections
that Gleason's theorem needs, and both hold for a single qubit. The derivation is the framework's
own;
its premise --- that the fibre datum is structurally excluded from records --- is
Proposition~\ref{prop:f1-dictionary}, and its geometric reading --- probability as base-point data
under the Hopf projection --- is developed in \S\ref{sec:hopf-justification}.

One standing objection to information-theoretic foundations is answered here, because this
framework needs the answer. Brukner and Zeilinger argued that Shannon information is conceptually
inadequate for quantum measurements: it is basis-dependent, and its operational meaning presupposes
ignorance of a definite but unknown fact, which the quantum case does not supply
\autocite{bruknerzeilinger2001shannon}. The framework's reply is structural, and it is a
demarcation rather than a rebuttal. Shannon's precondition fails exactly where they say it fails
--- in the pre-record regime, where there is no definite unknown to be ignorant of --- and there
the operative measure is the quadratic invariant just derived. Shannon's axioms hold exactly where
the precondition holds --- on records, after \textsc{fanout}, where outcomes are definite and
ignorance is classical --- and there the uniqueness theorem forces $-\log p$
\autocite{shannon1948mathematical}. The basis-dependence dissolves with the same stroke: a basis
\emph{is} selected precisely at \textsc{fanout}, so the basis-dependent measure applies only where
a basis exists. Shannon is not inadequate; Shannon is post-\textsc{fanout}. The framework thereby
keeps both measures, each on its own side of the boundary, and uses the additivity of the
post-record measure --- the very property the previous subsection used to pin the base field ---
with a clear conscience.

\subsection{Thomson's lamp, resolved}

The section owes \S\ref{sec:continuity} a debt incurred there: the lamp's amplitude-side limit.
Thomson's lamp --- toggled at times $1 - 2^{-n}$, its state at $t = 1$ demanded --- has no
classical answer: the task description fixes every finite stage and thereby fixes no limit at all
\autocite{Thomson_1954,benacerraf1962tasks}, which \S\ref{sec:continuity} read as the proof that
the record algebra is not limit-closed. On the amplitude side the limit exists, and the calculus of
this section computes it. The toggle exchanges the two amplitudes; the self-consistent state of the
completed task is the fixed point of the exchange in the complex calculus, and solving it under the
signed continuation gives
\[
P(\mathrm{on}) = \tfrac{1 + i}{2}, \qquad P(\mathrm{off}) = \tfrac{1 - i}{2},
\]
a superposed fixed point whose imaginary part carries the unrecordable datum --- the phase of the
endless alternation, when the lamp \emph{would have} stopped. Two checks certify that the
assignment is calculus rather than decoration. The pair is normalised,
$P(\mathrm{on}) + P(\mathrm{off}) = 1$; and the Born rule of the previous subsection, applied to
the fixed point, returns $|P(\mathrm{on})|^{2} = \tfrac{1}{2}$ --- exactly the symmetric classical
ignorance that is the only sane record-level answer, recovered rather than imposed. The lamp,
which entered the framework at its very beginning as the place where classical computation meets
its boundary, thus exits it as a worked example: undetermined as a record, determinate as an
amplitude, the two sides split precisely along \textsc{fanout} --- and the diagonal construction it
physically instantiates \autocite{cantor1891diagonal,Turing_1937} is the same one that powers
Proposition~\ref{prop:inside-outside}. The status flag is owed here as everywhere: the
undetermination is theorem-level; the specific fixed-point value is the framework's complex
calculus, structural, offered with its two consistency checks shown.

\subsection{Local tomography, derived}\label{sec:local-tomography}

The fourth instalment concerns composites, and it discharges two promises at once. The
Inside/Outside Equivalence Principle --- now a theorem --- requires that an external observer's
characterisation of a composite system be recoverable by the embedded participants, and the
embedded participants have access only to local measurements and their correlations. Equivalence
of the two perspectives therefore entails that the state of a composite is determined by the
statistics of local measurements: \emph{local tomography}, obtained as a consequence rather than
imposed as an axiom. This converts a postulate of the reconstruction programmes
\autocite{Hardy2001a,mullerp2011} into a theorem-shaped consequence of the one principle ---
discharging the caveat recorded in \S\ref{sec:continuity} --- and it closes, from the framework's
own resources, the escape route in the real-number debate noted above: the assumption the
objection of \autocite{hoffreumonwoods2026real} contests is here not an assumption. The status is
flagged plainly: the derivation is structural, one paragraph of perspective-accounting, and its
formalisation within the categorical setting belongs to Open Problem~\ref{op:sigma-functor}.

\subsection{Contrast: QBism, and the ledger}

The nearest neighbour in spirit deserves a paragraph of honest contrast. QBism
\autocite{fuchsschack2013qbism} holds, as this framework does, that quantum probabilities are an
agent's --- that the theory is a calculus of expectation, not a catalogue of properties. What QBism
does not say is why the agent's expectations must be complex-valued and quadratically squared; the
Born rule enters as a normative addition to Bayesian coherence, and the complex amplitude as an
unexplained representation. This section is, in effect, the missing why: the agent's ignorance is
self-referential --- the agent is inside the description --- and self-referential ignorance is
complex, with the quadratic rule as the cheapest record-compatible invariant of it. Where QBism
stops at the agent, the framework derives the agent's calculus.

The ledger, then. Delivered in this section: the reduction of two principles to one
(Proposition~\ref{prop:inside-outside}); the pre-record regime with its anti-events, handed to
\S\ref{sec:continuity} as the five-face regime whose continuity is priced there
(Open Problem~\ref{op:limit-realisability}); the complex base, forced downward by experiment and
upward by the additivity of information; the Born rule from fibre privacy; the lamp's fixed point;
and local tomography. Owed, and stated where they are owed: the functorial rigour of the involution
and the fixed-point classification (Open Problem~\ref{op:sigma-functor}); the remaining premise of continuity (Open Problem~\ref{op:limit-realisability}, in
\S\ref{sec:continuity}); and the
rung-by-rung compulsion of the algebra ladder, posed as Open
Problem~\ref{op:ladder-compulsion} in \S\ref{sec:division-algebras}. The
hand-on is the paper's spine: the $\mathbb{C}$ this section derives natively and the $\mathbb{C}$
that \S\ref{sec:continuity} underwrites through Hardy's theorem are the same field, the floor of
the ladder that \S\ref{sec:division-algebras} climbs and \S\ref{sec:hopf-justification}
geometrises. One principle; everything else is theorem, derivation, or a debt with a stated
invoice.

%      % EDITED 2026-10-29 scripts/patch_continuity_tower_premise_2026-10-29T1500.py
% EDITED 2026-09-15 scripts/patch_value_continuum_declared_cencov_2026-09-15T1540.py
% EDITED 2026-09-15 scripts/patch_continuity_connected_state_space_2026-09-15T1430.py
% =====================================================================
%  SECTION: Continuity Is Not an Axiom
%  Drafted 2026-07-04 from continuity_limit_realisability_2026-07-04.md
%  (now in project documents) and the May record
%  (dynamics_companion_fanout_unitarity_2026-05-20T1430.md).
%
%  PLACEMENT in main.tex: a standalone section between self-reference
%  and division-algebras --
%      % EDITED 2026-10-29 scripts/patch_continuity_tower_premise_2026-10-29T1500.py
% EDITED 2026-09-20 scripts/patch_born_cost_ledger_v2_2026-09-20T0900.py
% =====================================================================
%  SECTION: Self-Reference: One Principle
%  Drafted 2026-07-04. The paper's foundational section; resolves the
%  largest dangling label in the manuscript (sec:self-reference).
%
%  PLACEMENT in main.tex: FIRST of the foundations chain --
%      \input{self_reference_section_2026-07-04}
%      \input{continuity_section_2026-07-04}
%      \input{division_algebras_section_2026-07-03}
%      \input{hopf_justification_section_2026-06-08}
%
%  SOURCES (re-read 2026-07-04): lawvere_inside_outside_fixed_point_
%  theorem_2026-05-04T1806.md (the spine: the two-principles reduction,
%  the theorem, the sigma candidate, the open items -> op minted per
%  D9); research_summary_neg_prob.md (negativity route; lamp (1+i)/2);
%  dynamics_companion_fanout_unitarity (legs 1a/1c; anti-events; the
%  Brukner-Zeilinger rebuttal per D6); Pure_states_and_mixed_states
%  (local tomography derived; QBism contrast per D7).
%
%  DECISIONS ENCODED (JS 2026-07-04): D1 title as below; D2 Wheeler1990
%  cited here as the title anchor; D3 lineage Popper -> Breuer ->
%  Svozil -> Szangolies (Dalla Chiara omitted); D4 Frauchiger_2018
%  cited here, dissolved in sec:wigners-friend; D5 Kirkwood-Dirac
%  heritage via the Arvidsson-Shukur 2024 review as one-stop; D6
%  Brukner-Zeilinger + the framework rebuttal IN ("Shannon is
%  post-FANOUT"); D7 QBism contrast IN; D8 Aaronson private
%  communication OMITTED (published island paper only); D9 one op
%  minted (op:sigma-functor), cross-linked to op:hopf-frobenius.
%
%  CITATIONS -- existing keys: Wheeler1990, small_2006, Godel1931,
%  Lawvere_1969, Yanofsky2003, CoeckeKissinger2017, Rovelli1996a,
%  Popper1950a, Popper1950b, Breuer1995, svozil_k_2018,
%  szangolies_j_2018, Frauchiger_2018, dirac1942, Feynman_r_82,
%  Feynman1987, Youssef1994 [SWAP-FLAG: intended paper is "Quantum
%  mechanics as complex probability theory", Mod. Phys. Lett. A 9
%  (1994) 2571 -- if that lives under youssef1994quantum, swap],
%  caticha_a_2006, goyal_knuth_skilling_2010, Aaronson2004h, Hardy2001a,
%  muller_p_2011, Thomson_1954, benacerraf1962tasks, cantor1891diagonal,
%  Turing_1937.
%  NEW keys (RIS in self_reference_refs_2026-07-04.ris, all verified
%  this session): tikochinsky1988feynman, renou2021real, chen2022ruling,
%  wu2022refutation, hoffreumonwoods2026real, arvidssonshukur2024kirkwood,
%  bruknerzeilinger2001shannon, fuchsschack2013qbism,
%  shannon1948mathematical.
%
%  CROSS-REFS: forward (fine under two-pass LaTeX): sec:continuity,
%  sec:division-algebras, sec:hopf-justification, prop:f1-dictionary,
%  sec:wigners-friend, op:limit-realisability, op:hopf-frobenius.
%  Still dangling after this section: sec:fanout-boundary (and its
%  sibling label sec:fanout-direction) -- both DRAFTED 2026-07-06.
%
%  AUDITED 2026-07-06 (section-by-section final pass, first stop):
%  proof of prop:inside-outside re-verified by hand; attribution chain
%  clean (Burgin/Khrennikov correctly absent); all header-declared
%  citations confirmed used; lamp arithmetic and antecedents verified;
%  FANOUT smallcaps consistent; no pre-definition use of "frame".
%  Fixes: stray Unicode em-dash (QBism paragraph, the file's only
%  non-ASCII character) -> house ---; trace-blind two-way stitch added
%  to the Born subsection ("one privacy, two duties");
%  op:sigma-functor gains its first standing constraint (F1
%  inheritance via the op:hopf-frobenius discharge) and its previously
%  one-directional local-tomography cargo made explicit.
% =====================================================================

\section{Self-Reference: One Principle}
\label{sec:self-reference-2}

Wheeler asked for it from bit: for physics derived from information, ``every it --- every particle,
every field of force'' obtained from binary answers to yes-or-no questions \autocite{Wheeler1990}.
This paper's route to that programme runs through G\"{o}del: the single principle from which
everything else here is derived, deferred, or priced is \emph{self-reference} --- the fact that any
complete physical description, made from inside the world it describes, must describe the describer.
This section states the principle, proves that the framework's second historical commitment follows
from it as a theorem rather than standing beside it as an axiom, and then extracts the price of the
principle in four instalments: probabilities whose parts can cancel, amplitudes valued in
$\mathbb{C}$, the Born rule, and the local tomography of composite systems. What the section cannot
yet derive is stated as precisely as what it can, in one numbered problem and in the ledger at the
end.

\subsection{Events, not objects}

The framework's ontology is relational: the fundamental entities are events and the relations
between them, and objects are the bookkeeping of relational invariants --- what remains stable
across many relations. The kinship with relational quantum mechanics \autocite{Rovelli1996a} is
genuine and acknowledged: facts are facts relative to the systems that register them. The
difference, developed over the rest of this paper, is that the present framework grounds the
relational structure in a single logical principle and then carries it, through the division
algebras and the Hopf fibrations, all the way to the gauge group and the particle spectrum ---
relations are not only the ontology but the derivation engine. On this reading a measurement is not
the revelation of a pre-existing property; it is the creation of a relation, and the later sections
of this paper read that creation geometrically (\S\ref{sec:hopf-justification}) and even
gravitationally. Here we need only the logical core: whatever describes the world is an event in
the world it describes.

\subsection{One principle, and a theorem}

The framework historically carried two foundational commitments: self-reference, and the
\emph{Inside/Outside Equivalence Principle} --- that the laws of physics must be identical whether
formulated by an embedded participant or by an external observer. Two principles are one too many,
and this subsection reduces them.

The reduction corrects the framework's own origin, and we make the correction in print. The 2005
form of the argument \autocite{small_2006} ran a G\"{o}delian chain: a self-referential process,
being complete, must by the incompleteness theorem \autocite{Godel1931} be inconsistent; an
inconsistent system can negate itself; and the negation of self-reference is the production of
causal sequence. The chain captures the right intuition and is retained here as motivation, but it
is not sound as a proof: classical inconsistency yields \emph{ex falso quodlibet} --- every
proposition provable, every causal sequence an output, no constraint at all. The load-bearing
argument is different, and better: a fixed-point theorem.

Let $\mathcal{C}$ be the category of physical descriptions and admissible processes --- concretely,
the dagger-compact-closed setting with classical structures in which categorical quantum mechanics,
and the whole of this paper, works \autocite{CoeckeKissinger2017}. Lawvere showed that in a
sufficiently rich such setting, any weakly point-surjective self-representation forces every
endomap to have a fixed point, and Yanofsky showed that one diagonal argument of this shape
underlies Cantor's, G\"{o}del's, Turing's and Tarski's theorems alike
\autocite{Lawvere_1969,Yanofsky2003,cantor1891diagonal}. The framework's single commitment is that
$\mathcal{C}$ contains a \emph{self-representing object} --- an object capable of encoding the
processes on itself. That is what ``self-reference is foundational'' means, stated categorically,
and it is the surjectivity hypothesis of the diagonal lemma. Now let
$\sigma \colon \mathcal{C} \to \mathcal{C}$ be the \emph{inside/outside involution}: the operation
exchanging the two complementary presentations of any description --- the inside presentation, as
participant-indexed total space in which the observer is part of the system, and the outside
presentation, as a base-times-fibre decomposition over a classical parameter space. That
$\sigma^{2} = \mathrm{id}$ is not an axiom but a coherence requirement: the outside view of the
inside view of a thing is the thing. The candidate concrete realisation of $\sigma$ is the
Hopf-duality swap --- total space against base-times-fibre for the three fibrations --- which gives
the involution computable content and ties this subsection to \S\ref{sec:hopf-justification}.

\begin{proposition}[Inside/outside fixed point]\label{prop:inside-outside}
Let $\mathcal{C}$ satisfy the richness condition above and contain a self-representing object, and
let $\sigma$ be an involution on $\mathcal{C}$. Then the fixed-point locus
$\mathrm{Fix}(\sigma)$ is non-empty: there exist descriptions invariant under the exchange of
inside and outside perspective.
\end{proposition}

\begin{proof}
Lawvere's construction, applied to $\sigma$ in place of negation: given a weakly point-surjective
$\varphi \colon A \to Y^{A}$, form $g(a) = \sigma\!\left(\varphi(a)(a)\right)$; point-surjectivity
supplies $a^{*}$ with $\varphi(a^{*}) = g$, and then $\varphi(a^{*})(a^{*})$ is a fixed point of
$\sigma$ \autocite{Lawvere_1969,Yanofsky2003}. Where the same construction applied to Boolean
negation yields paradox --- negation has no fixed point on a two-element type, so the conclusion
explodes the premises --- the inside/outside involution is an honest involution on a rich type, and
the conclusion is a constructive fixed point rather than a contradiction.
\end{proof}

\textbf{Definition.} \emph{Physical law is the structure carried by} $\mathrm{Fix}(\sigma)$.

The Inside/Outside Equivalence Principle is thereby demoted from axiom to theorem: physical law is
invariant under the exchange of perspectives \emph{because it is the fixed-point set of that
exchange}, and the fixed-point set is forced to be non-empty by self-reference alone. One
commitment remains foundational; everything the second commitment asserted is now derived. The
reader who accepts categorical quantum mechanics as an existing programme has already accepted the
machinery; the framework adds one premise and collects the rest. And the fixed-point locus is not
empty of structure: it inherits the Frobenius-algebra apparatus of $\mathcal{C}$, so that the
Bell, GHZ and W primitives --- and with them the particle classification of the later sections ---
live inside $\mathrm{Fix}(\sigma)$ from the start. What is proved above is non-emptiness given the
setting; what is assumed is the setting itself and the existence of $\sigma$ as a genuine
involution; what is open is the rigour of both, and it is priced now rather than hidden:

\begin{openproblem}[The involution made functorial]\label{op:sigma-functor}
Fix the category $\mathcal{C}$ rigorously --- the dagger-compact-closed setting with classical
structures is the candidate --- and define $\sigma$ on morphisms, not only on objects, the natural
route being the induced fibre--base exchange on spider diagrams; establish that $\sigma$ is an (at
worst pseudo-)involutive endofunctor; and prove or refute the identification of
$\mathrm{Fix}(\sigma)$ with the three Hopf cross-sections, under which Adams's theorem would become
a corollary of the fixed-point classification. This problem is cross-linked with Open
Problem~\ref{op:hopf-frobenius}: under an affirmative resolution of both, the parallelisable
terminus, the fixed-point classification, and compositional completeness are three statements of
one fact, and progress on either problem constrains the other. One constraint is already in place:
at the complex rung, Proposition~\ref{prop:f1-dictionary} supplies the dictionary between
classical structures and the fibration's lifted skeletons --- the clause of
Open Problem~\ref{op:hopf-frobenius} it discharges --- and any resolution here must reproduce it.
The problem also carries a piece of cargo recorded later in this section: the formalisation, within
the fixed category, of the perspective-accounting that derives local tomography.
\end{openproblem}

\subsection{The lineage}

The obstruction this section builds on is old, and the framework claims its completion, not its
discovery. Popper argued in 1950 that a predictor embedded in the world has irreducibly inferior
knowledge of itself, because it cannot impart information about itself without, by that very act,
rendering the information out of date \autocite{Popper1950a,Popper1950b} --- the passage is quoted
where the framework makes the obstruction structural (\S\ref{sec:wigners-friend}). Breuer proved
the impossibility of accurate state self-measurement within quantum theory \autocite{Breuer1995};
Svozil's monograph surveys the whole territory where physics meets self-reference, undecidability
and Chaitin-style algorithmic limits \autocite{svozil_k_2018}; Szangolies sharpened the modern form
as epistemic horizons \autocite{szangolies_j_2018}; and the Frauchiger--Renner theorem is the
contemporary crisis-form --- nested observers deriving contradictory records when each treats the
others' quantum descriptions as universally valid \autocite{Frauchiger_2018} --- which
\S\ref{sec:wigners-friend} dissolves relationally rather than paradoxically. This framework's
addition to the lineage is twofold. Structurally, Proposition~\ref{prop:f1-dictionary} converts
Popper's temporal obsolescence into algebraic cancellation: the datum a copier cannot put into its
own record is exactly its fibre commitment, which does not arrive stale --- it cancels. And
foundationally, the present section runs the lineage in reverse: where the tradition treats
self-reference as a limit on knowledge \emph{within} physics, the framework takes it as the
generator of physics, and derives the quantum formalism from it.

\subsection{The price is negative probability}\label{sec:complex-amplitudes-2}

What does it cost a probabilistic description to contain its own describer? The first instalment is
the sign. In the regime before any record exists, every event has an anti-event --- the description
must be able to undo what it has provisionally done, or the self-description could never be
revised into consistency --- and a probability calculus closed under undoing cannot be confined to
$[0,1]$: it is a signed calculus, in which the parts of a probability can oppose and cancel. The
framework did not discover negativity as bookkeeping; Dirac admitted negative probabilities as
formal intermediaries \autocite{dirac1942}, Feynman defended them as consistent whenever the final,
recorded probabilities stay positive \autocite{Feynman_r_82,Feynman1987}, and the Kirkwood--Dirac
distribution has carried complex quasiprobability from 1933 to the present, as a recent
comprehensive review documents \autocite{arvidssonshukur2024kirkwood}. What the framework adds is
the forcing: negativity is not a permissible convenience but the first price of self-referential
closure, and it is charged exactly and only in the pre-record regime --- which is why every
negativity in the tradition lives in intermediate, unrecorded quantities and never in a final
record.

The second instalment is complexity. A signed information measure must extend Shannon's
$S = -\ln p$ off the positive axis, and the analytic continuation of the logarithm carries $p$ into
the complex plane: amplitudes, with the phase as the continued information. Here too the extension
has a prior tradition, and the attributions are owed precisely. Tikochinsky derived Feynman's rules
for complex amplitudes from consistency of composition \autocite{tikochinsky1988feynman}; Youssef
reformulated quantum mechanics as complex probability theory outright \autocite{Youssef1994};
Caticha carried objective amplitudes to Bayesian probabilities \autocite{caticha_a_2006}; and
Goyal, Knuth and Skilling derived Feynman's rules from a pair-postulate without assuming
analyticity \autocite{goyal_knuth_skilling_2010}. The extension is prior art. The framework's
claim is the compulsion: the complex extension is not a modelling choice among exotic probability
theories but the unique continuation forced by self-referential closure --- the same forcing that
Aaronson glimpsed when he argued quantum mechanics is what negative probability inevitably becomes
\autocite{Aaronson2004h}.

Why does the continuation stop at $\mathbb{C}$? Downward, the real field is not available: real
quantum theory is experimentally distinguishable from complex quantum theory in network scenarios
and has been disfavoured by experiment
\autocite{renou2021real,chen2022ruling,wu2022refutation}; the recent objection that the exclusion
rests on locally tomographic composition \autocite{hoffreumonwoods2026real} is, for this framework,
not an escape but a confirmation of where the weight sits --- local tomography is \emph{derived}
below, not assumed, so the framework closes from its own resources exactly the door the objection
holds open. Upward, the quaternions and octonions fail as a \emph{base} for the information
continuation: the continued measure is a logarithm, and beyond $\mathbb{C}$ the logarithm of a
product is no longer the sum of logarithms --- non-commutativity severs
$\log(ab) = \log a + \log b$ --- so the additivity that makes it an information measure is lost.
Aaronson's parameter-counting and signalling arguments against a quaternionic base point the same
way \autocite{Aaronson2004h}. The higher algebras are not banished; they return as the
\emph{internal} algebras of the register, which is the business of \S\ref{sec:division-algebras}.
The base is $\mathbb{C}$, for informational reasons, and stays $\mathbb{C}$.

\subsection{The Born rule from fibre privacy}\label{sec:born-rule}

The third instalment is the quadratic rule. An amplitude carries one datum more than any record of
it can: its overall phase. This was asserted throughout the framework's development as the
observer's private, self-referential ignorance; it is now a proved statement, in exactly the form
needed. By Proposition~\ref{prop:f1-dictionary}, a classical structure --- the copy operation any
record requires --- is committed to a specific lift of its skeleton, and that commitment is absent
from every classical record the structure induces: the copier is fibre-committed, the record is
fibre-blind, and the fibre freedom is gauge for every observer. (The computation extends verbatim
from one qubit to the full register, where Proposition~\ref{prop:trace-blind} makes it the ground
of hypercharge's tracelessness --- one privacy, two duties.) Popper's obsolescence, made
cancellation. It follows that anything recordable is invariant under the fibre --- observables and
their statistics must be $U(1)$-invariant functions of the amplitude. The lowest-degree polynomial invariant
of a complex amplitude is the quadratic $|\psi|^{2}$ --- the non-polynomial invariant $|\psi|$,
of degree one, is excluded by the analyticity declared at \S\ref{sec:complex-amplitudes} --- and
normalisation fixes the coefficient:
probability is the minimal record-compatible function of the amplitude,
\[
p \;=\; |\psi|^{2},
\]
the Born rule as the cheapest invariant the fibre permits. The same square is reached from the
cost side in \S\ref{sec:born-cost-ledger}, where a record is the closed loop of a history and its
return leg and the exponent is the number of legs; neither route uses the lattice of projections
that Gleason's theorem needs, and both hold for a single qubit. The derivation is the framework's
own;
its premise --- that the fibre datum is structurally excluded from records --- is
Proposition~\ref{prop:f1-dictionary}, and its geometric reading --- probability as base-point data
under the Hopf projection --- is developed in \S\ref{sec:hopf-justification}.

One standing objection to information-theoretic foundations is answered here, because this
framework needs the answer. Brukner and Zeilinger argued that Shannon information is conceptually
inadequate for quantum measurements: it is basis-dependent, and its operational meaning presupposes
ignorance of a definite but unknown fact, which the quantum case does not supply
\autocite{bruknerzeilinger2001shannon}. The framework's reply is structural, and it is a
demarcation rather than a rebuttal. Shannon's precondition fails exactly where they say it fails
--- in the pre-record regime, where there is no definite unknown to be ignorant of --- and there
the operative measure is the quadratic invariant just derived. Shannon's axioms hold exactly where
the precondition holds --- on records, after \textsc{fanout}, where outcomes are definite and
ignorance is classical --- and there the uniqueness theorem forces $-\log p$
\autocite{shannon1948mathematical}. The basis-dependence dissolves with the same stroke: a basis
\emph{is} selected precisely at \textsc{fanout}, so the basis-dependent measure applies only where
a basis exists. Shannon is not inadequate; Shannon is post-\textsc{fanout}. The framework thereby
keeps both measures, each on its own side of the boundary, and uses the additivity of the
post-record measure --- the very property the previous subsection used to pin the base field ---
with a clear conscience.

\subsection{Thomson's lamp, resolved}

The section owes \S\ref{sec:continuity} a debt incurred there: the lamp's amplitude-side limit.
Thomson's lamp --- toggled at times $1 - 2^{-n}$, its state at $t = 1$ demanded --- has no
classical answer: the task description fixes every finite stage and thereby fixes no limit at all
\autocite{Thomson_1954,benacerraf1962tasks}, which \S\ref{sec:continuity} read as the proof that
the record algebra is not limit-closed. On the amplitude side the limit exists, and the calculus of
this section computes it. The toggle exchanges the two amplitudes; the self-consistent state of the
completed task is the fixed point of the exchange in the complex calculus, and solving it under the
signed continuation gives
\[
P(\mathrm{on}) = \tfrac{1 + i}{2}, \qquad P(\mathrm{off}) = \tfrac{1 - i}{2},
\]
a superposed fixed point whose imaginary part carries the unrecordable datum --- the phase of the
endless alternation, when the lamp \emph{would have} stopped. Two checks certify that the
assignment is calculus rather than decoration. The pair is normalised,
$P(\mathrm{on}) + P(\mathrm{off}) = 1$; and the Born rule of the previous subsection, applied to
the fixed point, returns $|P(\mathrm{on})|^{2} = \tfrac{1}{2}$ --- exactly the symmetric classical
ignorance that is the only sane record-level answer, recovered rather than imposed. The lamp,
which entered the framework at its very beginning as the place where classical computation meets
its boundary, thus exits it as a worked example: undetermined as a record, determinate as an
amplitude, the two sides split precisely along \textsc{fanout} --- and the diagonal construction it
physically instantiates \autocite{cantor1891diagonal,Turing_1937} is the same one that powers
Proposition~\ref{prop:inside-outside}. The status flag is owed here as everywhere: the
undetermination is theorem-level; the specific fixed-point value is the framework's complex
calculus, structural, offered with its two consistency checks shown.

\subsection{Local tomography, derived}\label{sec:local-tomography}

The fourth instalment concerns composites, and it discharges two promises at once. The
Inside/Outside Equivalence Principle --- now a theorem --- requires that an external observer's
characterisation of a composite system be recoverable by the embedded participants, and the
embedded participants have access only to local measurements and their correlations. Equivalence
of the two perspectives therefore entails that the state of a composite is determined by the
statistics of local measurements: \emph{local tomography}, obtained as a consequence rather than
imposed as an axiom. This converts a postulate of the reconstruction programmes
\autocite{Hardy2001a,mullerp2011} into a theorem-shaped consequence of the one principle ---
discharging the caveat recorded in \S\ref{sec:continuity} --- and it closes, from the framework's
own resources, the escape route in the real-number debate noted above: the assumption the
objection of \autocite{hoffreumonwoods2026real} contests is here not an assumption. The status is
flagged plainly: the derivation is structural, one paragraph of perspective-accounting, and its
formalisation within the categorical setting belongs to Open Problem~\ref{op:sigma-functor}.

\subsection{Contrast: QBism, and the ledger}

The nearest neighbour in spirit deserves a paragraph of honest contrast. QBism
\autocite{fuchsschack2013qbism} holds, as this framework does, that quantum probabilities are an
agent's --- that the theory is a calculus of expectation, not a catalogue of properties. What QBism
does not say is why the agent's expectations must be complex-valued and quadratically squared; the
Born rule enters as a normative addition to Bayesian coherence, and the complex amplitude as an
unexplained representation. This section is, in effect, the missing why: the agent's ignorance is
self-referential --- the agent is inside the description --- and self-referential ignorance is
complex, with the quadratic rule as the cheapest record-compatible invariant of it. Where QBism
stops at the agent, the framework derives the agent's calculus.

The ledger, then. Delivered in this section: the reduction of two principles to one
(Proposition~\ref{prop:inside-outside}); the pre-record regime with its anti-events, handed to
\S\ref{sec:continuity} as the five-face regime whose continuity is priced there
(Open Problem~\ref{op:limit-realisability}); the complex base, forced downward by experiment and
upward by the additivity of information; the Born rule from fibre privacy; the lamp's fixed point;
and local tomography. Owed, and stated where they are owed: the functorial rigour of the involution
and the fixed-point classification (Open Problem~\ref{op:sigma-functor}); the remaining premise of continuity (Open Problem~\ref{op:limit-realisability}, in
\S\ref{sec:continuity}); and the
rung-by-rung compulsion of the algebra ladder, posed as Open
Problem~\ref{op:ladder-compulsion} in \S\ref{sec:division-algebras}. The
hand-on is the paper's spine: the $\mathbb{C}$ this section derives natively and the $\mathbb{C}$
that \S\ref{sec:continuity} underwrites through Hardy's theorem are the same field, the floor of
the ladder that \S\ref{sec:division-algebras} climbs and \S\ref{sec:hopf-justification}
geometrises. One principle; everything else is theorem, derivation, or a debt with a stated
invoice.

%      % EDITED 2026-10-29 scripts/patch_continuity_tower_premise_2026-10-29T1500.py
% EDITED 2026-09-15 scripts/patch_value_continuum_declared_cencov_2026-09-15T1540.py
% EDITED 2026-09-15 scripts/patch_continuity_connected_state_space_2026-09-15T1430.py
% =====================================================================
%  SECTION: Continuity Is Not an Axiom
%  Drafted 2026-07-04 from continuity_limit_realisability_2026-07-04.md
%  (now in project documents) and the May record
%  (dynamics_companion_fanout_unitarity_2026-05-20T1430.md).
%
%  PLACEMENT in main.tex: a standalone section between self-reference
%  and division-algebras --
%      \input{self_reference_section_2026-07-04}
%      \input{continuity_section_2026-07-04}
%      \input{division_algebras_section_2026-07-03}
%  (Architecture call SETTLED by the sync manifest: standalone.)
%
%  CITATIONS -- existing keys (as supplied by JS 2026-07-04):
%    Hardy2001a, Aaronson2004h, Thomson_1954, Lawvere_1969,
%    Turing_1937, muller_p_2011, Brukner2009.
%  NEW keys (RIS in continuity_refs_2026-07-04.ris):
%    cantor1891diagonal, benacerraf1962tasks, broeckertomdieck1985compact.
%  Deliberately excluded per JS: Chiribella--D'Ariano--Perinotti;
%  Bargmann/Wigner's theorem. Yanofsky2003 now paired on the
%  Lawvere-shaped clause (the pending-key excuse was void: the key
%  exists and sec:self-reference uses it).
%
%  CROSS-REFS: sec:self-reference now DRAFTED (2026-07-04, later the
%  same day) -- no longer dangling; only sec:fanout-boundary (and its
%  sibling sec:fanout-direction) -- DRAFTED 2026-07-06.
%
%  AUDITED 2026-07-06 (final pass, second stop): both proofs
%  re-verified by hand (Lagrange divisibility sentence exact; the
%  Baire countable-escape chain sound); caveat-discharge with
%  sec:self-reference confirmed two-way (already present); file pure
%  ASCII; no pre-definition "frame"; all citations used as declared.
%  Fixes: Yanofsky2003 paired on the Lawvere clause; the op's two
%  candidate routes anchored to printed propositions (f1-dictionary
%  for the divisibility premise's referent, inside-outside for the
%  closedness demand) -- statuses unchanged, referents sharpened;
%  header staleness corrected. Deliberate FORWARD references --
%  prop:f1-dictionary and sec:hopf-justification (defined later in the
%  paper; fine under two-pass LaTeX and flagged in prose),
%  sec:division-algebras, sec:space-of-computations.
%
%  DECISIONS ENCODED: op:limit-realisability minted (the owed arrow,
%  priced: exactly two derivations, divisibility and closedness);
%  Lagrange used twice (finite exclusion; component-group triviality);
%  the (1+i)/2 lamp value NOT stated here -- "a superposed fixed point"
%  with the complex-probability treatment deferred to
%  sec:self-reference; the two-branch over-determination stated to
%  pre-empt a circularity reading; Thomson spelling throughout.
% =====================================================================

\section{Continuity Is Not an Axiom}
\label{sec:continuity}

Hardy derives quantum theory from five axioms \autocite{Hardy2001a}. Four of them are satisfied
by classical probability theory; the fifth --- that a \emph{continuous} reversible
transformation connects any two pure states --- is the entire difference, the single postulate
that turns the classical simplex into the complex Hilbert space. Every reconstruction programme
pays an equivalent toll: continuous reversibility in Masanes and M\"{u}ller
\autocite{mullerp2011}, the corresponding postulate in Daki\'{c} and Brukner
\autocite{Brukner2009}. The axiom is load-bearing, and in every programme it is unexplained: it
is where the quantumness is put in by hand. The claim of this section is that in the present
framework it is not an axiom. Continuity is the fifth face of the pre-measurement regime that self-reference forces, and what
stands between the framework and a full derivation is not a postulate about values but one
premise about transformations --- that every transformation the inside performs has a stationary
outside presentation --- derived in outline below from the inside/outside involution of
\S\ref{sec:self-reference} and stated as an open problem at the end of the section with nothing
hidden in it. The continuum itself, the completion that every probability theory adopts, is
named for what it is: the real line, shared with classical probability and carrying no quantum
content.

One structural remark prevents a misreading. The framework reaches complex amplitudes natively in
\S\ref{sec:self-reference}, by the negative-probability route; the present section reaches the same
destination by an independent branch --- the regime's continuity, then Hardy's theorem. The two
branches converging on one structure is over-determination, not circularity, and the convergence is
itself evidence that the destination is forced.

\subsection{Five faces of one regime}

The regime before any measurement --- before any \textsc{fanout} --- can be characterised five
equivalent ways: negative probability is available; anti-events exist, every event being undoable;
no classical record exists; no preferred basis is selected; and a continuous reversible
transformation connects any two pure states. These are not five premises to be combined. They are
one regime under five descriptions, and \textsc{fanout} (\S\ref{sec:fanout-boundary}) is the single
event that fails all five at once: it writes the record, selects the basis, kills the anti-events,
discretises the transformation group, and switches off the negativity. The equivalence of the first
two faces is the anti-event identity of \S\ref{sec:self-reference}; the third is the boundary's
definition. The work of this section is the equivalence of the last two faces, and then the route
from the regime's existence to its continuity. The fourth face has, since
Proposition~\ref{prop:f1-dictionary} (\S\ref{sec:hopf-justification} --- a deliberate forward
reference), an exact form: a basis-selection is the commitment of a classical structure, a lifted
antipodal skeleton, so ``no preferred basis'' means precisely that no copy structure has fired.

\subsection{Basis selection is the collapse of continuity}

\begin{proposition}[Continuity forbids a selected basis]\label{prop:basis-discreteness}
If a finite set of pure states is privileged --- a preferred basis --- then every reversible
transformation respecting the privilege carries the set to itself and acts on it as a permutation;
the respecting transformations therefore realise only a finite group of motions among the
distinguished states, and no continuous family of them connects distinct pure states of the set.
Contrapositively, continuous reversibility entails that no basis is selected.
\end{proposition}

\begin{proof}
A privileged finite set is by definition carried to itself by every structure-preserving reversible
map, whose induced action on the set is a permutation of finitely many points. A continuous path of
such maps induces a continuous path in a finite discrete set, which is constant. Hence a continuous
reversible family moving one distinguished pure state to another respects no such privilege, and if
every pair of pure states is so connected, no finite set is privileged.
\end{proof}

The converse is where the historical difficulty lies: the absence of a privileged finite set does
not by itself force continuity, for the transformation group could in principle be infinite and
discrete. Two structural properties close every such escape at once.

\begin{proposition}[Divisible and closed implies connected]\label{prop:divisible-closed}
Let $G$ be the group of reversible transformations of the pre-measurement regime, acting on the
state body. If $G$ is divisible --- every element has an $n$-th root for every $n$ --- and
topologically closed, then $G$ is a compact connected Lie group; and if $G$ moreover acts
transitively on the pure states, then a continuous reversible transformation connects any two of
them: Hardy's fifth axiom.
\end{proposition}

\begin{proof}
The state body is compact and convex with nonempty interior, so its symmetry group is compact, and
a closed subgroup of a compact matrix group is a compact Lie group with finitely many connected
components \autocite{broeckertomdieck1985compact}. The component group $G/G_{0}$ is a quotient of a
divisible group and hence divisible; but a nontrivial finite group is never divisible --- with
$\lvert H\rvert = m$ every element satisfies $h^{m} = e$, so no element other than the identity
possesses an $m$-th root --- and $G/G_{0}$ is finite, so it is trivial and $G = G_{0}$ is connected.
A connected Lie group is path-connected, and transitivity then supplies the continuous reversible
path between any two pure states.
\end{proof}

The same one-line fact of Lagrange closes the remaining escapes. Finite groups are excluded
directly, being nontrivial and never divisible. Countable groups are excluded because a
countable closed subgroup of a compact group is compact, has an isolated point by the Baire
category theorem, and is therefore discrete and finite. Uncountable disconnected groups are
excluded by the component-group argument of the proposition. Every route to a discrete
transformation group therefore passes through a failure of divisibility or a failure of
closedness. Divisibility, however, asks for more than the conclusion needs, and the next
proposition removes it.

\begin{proposition}[Continuity from a connected state space]\label{prop:connected-transitive}
Let the pure states of a finite-dimensional system form a connected space $P$, and let the group
$G$ of reversible transformations of the pre-measurement regime be topologically closed and act
transitively on $P$. Then the identity component $G_0$ already acts transitively on $P$, and a
continuous reversible transformation connects any two pure states: Hardy's fifth axiom holds,
with no divisibility assumed.
\end{proposition}

\begin{proof}
As in Proposition~\ref{prop:divisible-closed}, closedness makes $G$ a compact Lie group with
finitely many connected components \autocite{broeckertomdieck1985compact}, so $G_0$ is a closed
normal subgroup of finite index and is path-connected. Each $G_0$-orbit in $P$ is the continuous
image of a compact group, hence compact, hence closed. Since $G$ is transitive and $G_0$ is
normal, the $G_0$-orbit of $gx$ is $g$ applied to the $G_0$-orbit of $x$, so the cosets of $G_0$
permute the orbits transitively and there are finitely many of them. A finite partition of $P$
into closed sets is a partition into open sets; $P$ is connected; so there is one orbit and
$G_0$ is transitive. Given pure states $x$ and $y$, choose $g \in G_0$ with $gx = y$ and a
continuous path $g_t$ in $G_0$ from the identity to $g$; then $g_t x$ is a continuous reversible
transformation from $x$ to $y$.
\end{proof}

The proposition relocates the premise. Divisibility was asked for because a transformation group
could be infinite and discrete; but a group acting transitively on a continuum cannot be
countable, and a closed group acting transitively on a connected space cannot be disconnected in
the way that matters, so what the argument needs from the regime is not the availability of
$n$-th roots but the connectedness of the state space itself. That is a statement about values:
that the pure states are all the unit vectors of $\mathbb{C}^{K}$ up to phase, the space
$\mathbb{CP}^{K-1}$, connected for $K \ge 2$, and hence that the amplitudes of
\S\ref{sec:self-reference} fill their continuum. The two standing foils show the premise doing
its work. The stabiliser fragment of quantum theory has complex amplitudes and is contextual,
yet it has finitely many pure states, its outcome probabilities take only the values $0$,
$\tfrac12$ and $1$, and its reversible transformations form a finite group, the Clifford group,
so Hardy's fifth axiom fails in it; Spekkens's toy theory \autocite{spekkens2007toy} is the same
with real amplitudes, its ignorance rationed by a knowledge-balance principle. Neither has a
connected state space, and neither is a model of the regime described here, whose credences take
every value. Divisibility, which both also lack --- the phase gate has no square root in the Clifford group
--- is a symptom of the discreteness, not its cause. What both lack can be said exactly, and the
next subsection says it: each contains a transformation that fixes no state it admits.

\subsection{The premises, and what they mean}

Three premises carry Proposition~\ref{prop:connected-transitive}. Transitivity is the regime's
homogeneity: inside/outside equivalence leaves no pure state privileged, and a proper closed
invariant set of pure states would define a superselection observable, a classical record
predating any \textsc{fanout} (\S\ref{sec:fanout-boundary}). Connectedness of the state space is two statements, and only one of them is quantum. The first
is that the pure states are infinite and dense, with a fixed point for every transformation; the
second is that the state space is complete. Quantum theory over the algebraic numbers ---
amplitudes in the algebraic closure of $\mathbb{Q}$ within $\mathbb{C}$ --- has the first
without the second: it has $K = N^{2}$, the Born rule, superposition, entanglement and Bell
violation at the Tsirelson value, every unitary has an algebraic eigenvector, and it is
countable and totally disconnected. So the whole quantum content of Hardy's axiom lies in the
density, and the completion is the real-number convention of every probability theory. Hardy
adopts it before his fifth axiom is stated, in the first, which makes probabilities real numbers
and gives his classical theory a real simplex of mixtures; this framework adopts it at the
Shannon step, where credences are declared to range over a connected domain so that the
continuation of the logarithm is unique (\S\ref{sec:complex-amplitudes}). Given that convention,
density of the pure states is connectedness of their closure. Divisibility, which an earlier
version of this section asked for, is not needed: it follows from the conclusion, every element
of a compact connected Lie group having roots of every order.

Closedness is the completion, limit-realisability: a Cauchy sequence of available
transformations has its limit available. It cannot be derived, and the reason should be stated
rather than hidden. Any collection of first-order principles about states, transformations and
probabilities that has an infinite model has a countable one, and a countable subset of
$\mathbb{CP}^{K-1}$ is totally disconnected; the founding non-computability does not escape
this, since the arithmetical complex numbers form a countable algebraically closed field
containing non-computable numbers and closed under every limit it can itself describe. Hardy's
axiom and closedness are second-order, quantifying over paths and over sequences, and a
second-order principle is grounded as Dedekind completeness is grounded or not at all. The
grounds are these. The completion of a dense set under the finite-precision metric is unique and
depends on no choice, whereas any countable alternative is a choice of definability level, an
extra law with no ground and no test; and the framework's own result is that no observer's
descriptive level captures the regime, since by \S\ref{sec:self-reference} every observer's
state exceeds its own level and the inside/outside principle makes that a law rather than a
perspective. The burden therefore falls on any restriction, not on the completion, which is the
least committal option and carries no quantum content. What limit-realisability adds to the line
is still worth stating exactly, because it is the corrected form of a natural intuition.
Cantor's diagonal \autocite{cantor1891diagonal} does not build the continuum; its output is an
infinite decimal --- a Cauchy sequence of rationals --- which exists \emph{as a point} only
because the line is complete. The diagonal presupposes limit-realisability and then shows that
the completed line outruns every internal listing. And the character of the completion is the
framework's own: the computable reals are countable \autocite{Turing_1937}, hence of measure
zero, so once the transformation group contains its limits almost all of them are
non-computable. Non-computability is a property of the completed continuum, not the reason for
it. The constructive reading, on which only describable limits exist, remains consistent; it is
the countable theory above, and the two readings differ in nothing an experiment can certify.

Density is where the section's title is earned, and it is derived from the paper's own premise.
Write $(\mathrm{I})$ for the statement that every admissible transformation of the regime is a
phase of some admissible classical structure: diagonal, that is, in the record basis of some
observer the regime contains, a phased spider of that basis in the calculus of
\S\ref{sec:hopf-justification} \autocite{coeckeduncan2011interacting} (VERIFY-SUPPORT). It says
of transformations what the global phase of \S\ref{sec:self-reference} says of states: there are
no absolute ones. It follows from three things the paper already holds. The inside/outside
involution is a genuine involution, so the outside presentation of the inside is a presentation
of the same object with nothing lost: the outside can have complete knowledge of the inside, a
potential record, the joint state lying in a structure the outside carries. Time is operation
count, internal to the system that operates (\S\ref{sec:space-of-computations}), so the outside
does not share the inside's count. And before any \textsc{fanout} that count is not classically
recorded, so the outside's description is pure. Let $g$ be the transformation the inside
performs per operation. The outside's complete knowledge cannot depend on the inside's count,
and it is pure, so it is a pure state of the inside's clock and system that is stationary under
``advance the clock and apply $g$''; the stationary states of that map are clock Fourier modes
tensored with eigenvectors of $g$ and nothing else \autocite{pagewootters1983evolution}
(VERIFY-SUPPORT). So the outside holds a structure in which $g$ acts by phases, which is
$(\mathrm{I})$; and the relative phase in that structure is the outside's ignorance of the
inside's computational count, the same ignorance the global phase records one level down. The
consequences are a theorem. A phase of a structure fixes the structure's copyable states, so
every admissible transformation fixes an admissible pure state; a finite transitive action has
an element moving every point (Jordan), so the pure states are infinite; the transformation
carrying a stabiliser state to the fixed point of the composite of a $Z$-phase and an $X$-phase
turns through $\arccos(1/\sqrt{3})$, an irrational fraction of a turn, so the group has an
element of infinite projective order and its orbits are dense. With the completion,
Proposition~\ref{prop:connected-transitive} fires and Hardy's fifth axiom holds. The two foils
fail $(\mathrm{I})$ exactly where they fail continuity: fourteen of the twenty-four Clifford
operations are phases of no Pauli basis and fix no stabiliser state, their fixed points being
the magic states; the toy theory's three-cycles are phases of no epistemic basis; and in each
the outside's knowledge of the inside under the order-three protocol is either count-dependent
or a mixture over the count, so each makes the inside's time classical before any record exists.
The derivation of $(\mathrm{I})$ is an outline, not a proof, and it is what the open problem
below asks for.

\subsection{Why the cut matters physically}

The regime's continuity is not decorative. Reversibility alone, over complex amplitudes, licenses
the full invertible group $GL(n, \mathbb{C})$, and Aaronson showed what that buys: superluminal
signalling, and computational power at the level of \textsf{PP} \autocite{Aaronson2004h}. The cut
from $GL(n, \mathbb{C})$ down to the unitary group is the conserved quadratic form, and continuity
is what entitles the theory to make it: with the fifth axiom in hand, Hardy's theorem delivers the
complex Hilbert space, the quadratic norm and the unitary group in one stroke \autocite{Hardy2001a}
--- cited here, not re-proved. One honest caveat travels with the citation. Hardy's reconstruction
uses his remaining axioms besides --- composition and subspace structure --- and these must be
derived rather than imported: the composition ledger of \S\ref{sec:division-algebras} and the
locality-of-tomography argument of \S\ref{sec:self-reference} carry that duty, and it is discharged
there, not silently assumed here.

\subsection{Super-tasks at the boundary}

The premise just priced has a fifty-year-old stress test in the philosophy of mathematics, and the
framework passes it in an unexpected way. A super-task --- infinitely many operations completed in
finite time \autocite{Thomson_1954} --- presupposes its terminal moment as an actual point of the
timeline: the accumulation point of a generic sequence of step-times exists only in the completed
line, and the defining question, the state \emph{after all the steps}, is meaningful only under
limit-realisability of time. Super-tasks are downstream consumers of exactly the completeness at
issue in this section. Thomson's lamp then exhibits both sides of the \textsc{fanout} boundary at
once. On the record side, the lamp's description fixes its state at every finite stage and thereby
fixes no limit at all: any terminal value is consistent with the whole task
\autocite{Thomson_1954,benacerraf1962tasks}. That is not a puzzle about lamps; it is the proof that
the classical record algebra is not limit-closed, and cannot be. On the amplitude side the limit
exists --- a superposed fixed point, whose complex-probability treatment belongs to
\S\ref{sec:self-reference} --- because that side \emph{is} limit-closed, which is the fifth face
itself. The boundary is thereby the edge of limit-realisability: complete on the reversible side,
incomplete on the record side, with the lamp standing on both sides at once as the boundary's oldest
exhibit. A final clause dissolves rather than answers the classic question. In this framework time
is operation count (\S\ref{sec:space-of-computations}), and ``after infinitely many operations''
indexes no operational moment: the lamp packs its endless \textsc{fanout}s into the emergent
continuum, and asking for its record at the accumulation point asks the discrete contents for a
value that only the completed container possesses. On or off is not undecided; it is unasked.

\subsection{What remains, priced}

\begin{openproblem}[Every transformation a phase of an inside structure]\label{op:limit-realisability}
Proposition~\ref{prop:connected-transitive} reduces Hardy's fifth axiom to three premises:
connectedness of the pure states, transitivity, and closedness. Transitivity is argued above
from the absence of records. Closedness is the completion, adopted on the grounds stated above
and carrying no quantum content. Connectedness is density plus that completion, and density
follows, by Jordan's theorem and the infinite-order argument, from $(\mathrm{I})$: that every
admissible transformation of the regime is a phase of some admissible classical structure. What
remains open is the rigour of the derivation of $(\mathrm{I})$ from the paper's premise. Derive,
inside the fixed-point locus $\mathrm{Fix}(\sigma)$ of Proposition~\ref{prop:inside-outside},
that every endomap is a phase of a classical structure in $\mathrm{Fix}(\sigma)$, from the
self-representing object, the involution $\sigma$ as a genuine involution, and the reading of
the inside presentation as a sequence indexed by operation count with no record before
\textsc{fanout}. The outline above does this with the stationary-state construction of Page and
Wootters; a derivation passes if it makes that outline a theorem of the calculus, with the
outside's count-independent complete description shown to exist and to be pure from the
involution alone. It fails before any detail is checked if every premise it uses holds in the
stabiliser fragment of quantum theory, in Spekkens's toy theory \autocite{spekkens2007toy}, in
the real qubit, or in quantum theory over the Gaussian rationals, all of which violate
$(\mathrm{I})$ outright and all of which make the inside's operation count classical before any
record exists. It also fails if it imports the completion as a premise about transformations:
the completion is the real line and has no derivation, and a derivation of $(\mathrm{I})$ must
hold in the countable algebraic theory, where $(\mathrm{I})$ is true and continuity is not.
\end{openproblem}

The accounting, in one sentence: the axiom that every reconstruction programme pays is, in this
framework, one derivation posed, of Lawvere's shape --- that every endomap of the fixed-point
locus is a phase of a classical structure in it --- and one convention named, the completion,
which classical probability pays too. And
the hand-on is already visible: the complex amplitudes that this section's continuity underwrites
through Hardy's theorem, and the complex amplitudes that \S\ref{sec:self-reference} derives natively,
are one and the same $\mathbb{C}$ --- the floor of the ladder that \S\ref{sec:division-algebras}
climbs.

%      % =====================================================================
%  SECTION: The Division-Algebra Ladder
%  Drafted 2026-07-03. One of the three foundations sections
%  (sec:self-reference, sec:division-algebras, sec:fanout-boundary);
%  this one resolves the dangling \S\ref{sec:division-algebras} in the
%  Hopf-justification section, whose header names it "the
%  Cayley--Dickson / Hurwitz step".
%
%  PLACEMENT in main.tex (REORDERED per JS 2026-07-06): position 6,
%  opening Act II ("why stop at C?") right before its first consumer --
%      \input{self_reference_section_2026-07-04}
%      \input{continuity_section_2026-07-04}
%      \input{L2_metric_global_phase_section_2026_06_07}
%      \input{wigners_friend_in_the_hopf_picture_2026-06-13}
%      \input{derivation_of_the_quantum_postulates}
%      \input{division_algebras_section_2026-07-03}   <-- moved here
%      \input{spacetime_from_non-computability}       <-- first consumer
%                                                          (S^7 exclusion)
%      \input{hopf_justification_section_2026-06-08}
%  The section is position-portable: all adjacency phrasing was
%  retargeted to explicit \S\ref on 2026-07-06. The S^7 hand-off
%  sentence to sec:space-of-computations WRITTEN 2026-07-06 (spacetime
%  file read; the closing now hands first to its immediate consumer,
%  then to hopf-justification). Upload verifications COMPLETE: L2 and
%  quantum-postulates do not consume the ladder; WF's references are
%  forward-voiced or self-carrying.
%
%  DIVISION OF LABOUR (deliberate, to avoid overlap):
%   - sec:self-reference (undrafted): self-reference -> complex
%     amplitudes and the Born rule; why not a real base. This section
%     TAKES that as given and starts at C.
%   - THIS section: the algebra -- the doubling construction, the
%     ledger of losses-as-physics, the internal-algebras-not-a-new-QM
%     scoping (Adler pre-empt), the two appearances of the ladder
%     (register / measurement climb, stated lightly), and Hurwitz as
%     the algebraic termination with the committed zero-divisor
%     reading.
%   - sec:hopf-justification: the geometry (Adams, parallelisability,
%     the fibrations), the join of the two classifications, the
%     entanglement-sensitivity literature, the C-direction honest debt,
%     and the epistemic reading of termination.
%
%  COMMITTED READINGS USED (sources re-read 2026-07-03):
%   - Loss ledger and level dictionary: cayley_dickson_qubits_for_review
%     sec.5 ("each step loses a symmetry, and each lost symmetry is
%     identified with a physical content").
%   - Termination reading: grounding_argument_gauge_group ("the
%     sedenions have zero divisors and cannot support consistent
%     information flow around loops") -- leads; norm-composition-as-
%     probability given as its amplitude-level restatement.
%   - Internal-not-base scoping: cayley_dickson_qubits_for_review Q2 and
%     sec.3.2 caveat (Adler objection; C-qubit treated geometrically).
%
%  HONEST DEBT: ladder proved closed at the top (Hurwitz), derived at
%  the floor. The rung-compulsion revisit-trigger FIRED 2026-07-06:
%  sec:self-reference and op:sigma-functor now exist, so the thesis is
%  minted as op:ladder-compulsion (with the CD-conjugation-as-sigma
%  candidate route), inheriting op:sigma-functor's status.
%
%  AUDITED 2026-07-06 (final pass, third stop): all four citations
%  used; all cross-refs live; CD facts, dims, alternativity, Hurwitz
%  statement verified; Adler subsection delivers everything the
%  register's compressed pointer promises; file pure ASCII; no bare
%  "frame". Fixes: opening rewritten for both predecessors (continuity
%  was inserted between after drafting); "associator's non-vanishing is
%  mass itself" recalibrated to the register's debt-against-selection
%  defusal; the O->S ledger row and Hurwitz subsection reconciled with
%  the fourth-qubit forward remark ("the terminus closes the ladder of
%  forces; what survives past it is, at most, flavour").
%
%  CROSS-REFS: sec:self-reference [sibling dangle -- resolves when that
%  section is drafted], sec:hopf-justification, sec:hopf-information-bound,
%  sec:register, sec:entanglement-classes, sec:generations, sec:masses,
%  sec:wigners-friend. All others exist.
%  CITATIONS: existing keys baez2002octonions, hurwitz1898composition,
%  furey2018three. ONE new key: adler1995quaternionic (RIS in
%  division_algebras_refs_2026-07-03.ris).
%  British English; \autocite; no bullets in prose; the ledger is a
%  table because the comparison is the argument.
% =====================================================================

\section{The Division-Algebra Ladder}
\label{sec:division-algebras}

The foundations behind us end at amplitudes twice over: natively, where the self-referential
argument forced the description of a system into the system's own state space and priced
consistency at a number system whose parts can cancel --- complex amplitudes, and the Born rule
that squares them (\S\ref{sec:self-reference}) --- and independently, where the regime's
continuity delivers the same field through Hardy's theorem (\S\ref{sec:continuity}), the
convergence of the two branches being evidence rather than accident. This section asks the question
both routes leave open. Why
should the number system stop at $\mathbb{C}$? Ordinary quantum mechanics stops there by postulate;
a framework that has forbidden itself postulates must answer, and the answer is that the climb
neither stops at $\mathbb{C}$ nor continues forever. There are exactly two more rungs, each
purchased by surrendering an algebraic law and each paying for itself in physics, and the ladder is
closed at the top by a theorem of 1898. This section sets out the algebra of that ladder: the single
construction that climbs it, the ledger of what each rung costs and buys, and the theorem that ends
it. The geometry of the same ladder --- that each rung acts on a sphere, that exactly those spheres
are parallelisable, and that the algebraic and geometric classifications are one --- is the business
of the geometry section (\S\ref{sec:hopf-justification}).

\subsection{One construction climbs the whole ladder}

The rungs are not a miscellany; a single construction generates them all. Given an algebra with a
conjugation, the Cayley--Dickson construction builds its double on ordered pairs, with a twisted
multiplication in which the conjugate enters and a conjugation $(a,b)^{*} = (a^{*}, -b)$
\autocite{baez2002octonions}. (Conventions differ in the placement of the conjugations; all yield
isomorphic algebras.) Applied repeatedly it produces the sequence
$\mathbb{R} \to \mathbb{C} \to \mathbb{H} \to \mathbb{O} \to \mathbb{S}$, of real dimensions
$1, 2, 4, 8, 16$: the reals, the complex numbers, the quaternions, the octonions, the sedenions.
The ground floor is classical: $\mathbb{R}$ is ordered and positive-definite, the arithmetic of
classical probability. The first doubling is therefore the quantum move itself --- the passage from
probabilities to amplitudes that can interfere --- and it is not made here by choice: it is the
derivation of \S\ref{sec:self-reference}, already behind us. What this section must account for is
the rest of the climb.

\subsection{The ledger}

In mathematics the doublings are usually told as a story of degeneration: each algebra is worse
behaved than the last, and the sequence is studied for where it breaks. This framework reads the
same ledger in the opposite direction. Each rung is bought by surrendering a law, and each
surrendered law is exactly a piece of physics --- the losses are the assets, and the comparison is
sharp enough to be the argument:

\begin{center}
\begin{tabular}{llll}
\hline
doubling & dimension & what is lost & what the loss becomes \\
\hline
$\mathbb{R} \to \mathbb{C}$ & $1 \to 2$ & the order of the reals & phase and interference; amplitudes \\
 & & & that cancel; the Born rule (\S\ref{sec:self-reference}) \\
$\mathbb{C} \to \mathbb{H}$ & $2 \to 4$ & commutativity & non-commuting actions: weak isospin \\
 & & & $SU(2)$, the $B$ register (\S\ref{sec:register}) \\
$\mathbb{H} \to \mathbb{O}$ & $4 \to 8$ & associativity & the associator: colour, chirality, the \\
 & & & generations, and mass as debt \\
 & & & (\S\ref{sec:generations}, \S\ref{sec:masses}) \\
$\mathbb{O} \to \mathbb{S}$ & $8 \to 16$ & division itself & no further force: zero divisors appear \\
 & & (zero divisors) & and the ladder of sectors ends (below) \\
\hline
\end{tabular}
\end{center}

Walking the rows: the order lost at the first doubling is the possibility of interference --- parts
of a probability that can oppose one another --- which is where the derivation of \S\ref{sec:self-reference}
lives. The commutativity lost at the second is the existence of actions whose order matters; the
unit quaternions form the group $SU(2)$, and the framework houses weak isospin there, as the
internal algebra of the $B$ register whose order-sensitive, non-commuting dynamics the interaction sections develop
(\S\ref{sec:register}, \S\ref{sec:entanglement-classes}). The associativity lost at the third ---
the octonions retain the weaker property of alternativity \autocite{baez2002octonions} --- creates
the associator $[a,b,c] = (ab)c - a(bc)$, and the associator is where most of this paper's physics
is transacted: the stabiliser of a preferred complex direction inside the octonions' automorphisms
is the colour group (\S\ref{sec:hopf-justification}); the triality structure and its $K_3$
obstruction generate the three generations and the sterile channel (\S\ref{sec:generations}); and
the associator is the currency of mass --- the debt against the selected direction that
\S\ref{sec:masses} prices, owed by whatever action moves the selection and not by residence at
the rung, which is the defusal \S\ref{sec:register} makes explicit. The algebraic facts
in the third column are theorems; the identifications in the fourth are established in the sections
cited, each carrying its own epistemic flag there, and are not re-litigated here. What this table
asserts on its own account is only the correspondence's shape: one law surrendered per rung, one
sector of physics gained per law.

\subsection{Internal algebras, not a new quantum mechanics}

A reader versed in the quaternionic-quantum-mechanics literature will at this point raise a fair
objection: quantum theory over $\mathbb{H}$ as a foundation is known to be troubled --- the
tensor-product and composite-system difficulties set out by Adler and others
\autocite{adler1995quaternionic} --- and quantum theory over $\mathbb{O}$ founders earlier still,
since a non-associative algebra does not furnish a Hilbert space at all. The framework's position
must therefore be stated exactly, because it does not make the claim the objection targets. The base
of the physical Hilbert space here is $\mathbb{C}$ and remains $\mathbb{C}$; the higher algebras
enter as the \emph{internal} algebras of the register --- $\mathbb{H}$ for the $B$ qubit,
$\mathbb{O}$ for the $C$ qubit (\S\ref{sec:register}) --- in a tensor structure whose outer theory is
ordinary complex quantum mechanics. This is the same move by which the division algebras enter the
constructions of Furey and the related programmes without any proposal to rebase quantum mechanics
itself \autocite{furey2018three}. For the octonionic register the point is sharper still:
$\mathbb{O}^{2}$ is not a Hilbert space, and the framework accordingly treats the $C$ qubit not in
the Dirac--von Neumann manner but geometrically, as the sphere $S^{15}$ with its Hopf structure ---
which is precisely the form in which \S\ref{sec:hopf-justification} needs it. Nothing in this paper therefore
stands or falls with quaternionic or octonionic quantum mechanics as foundations; the ladder is
climbed by the registers' internal algebras, not by the theory's base field.

\subsection{The same ladder, twice}

The ladder appears in this framework in two guises, and it will help the reader to have them
distinguished now. Statically, it is the anatomy of a single relation: the three qubits of the
register take their values one rung apart --- $A$ over $\mathbb{C}$, $B$ over $\mathbb{H}$, $C$ over
$\mathbb{O}$ --- so that every particle carries the whole ladder as its internal structure
(\S\ref{sec:register}). Dynamically, it is the shape of measurement: a single qubit's pure states
fill $S^{3}$, an entangling measurement carries the joint state into $S^{7}$, and a third party into
$S^{15}$ --- the correspondence between qubit entanglement and the Hopf fibrations that
\S\ref{sec:hopf-justification} takes from a mature literature --- so that each act of measurement ascends one rung --- the climb the reader has already made,
narrated in \S\ref{sec:wigners-friend} --- and every completed relation is a terminated climb, which is
the reading \S\ref{sec:hopf-information-bound} gives the tower's endpoint. One ladder, seen twice:
as the anatomy of a relation and as the dynamics of relating. The framework regards this double
appearance not as a coincidence but as the same structure viewed statically and in motion, and the
sections cited carry that identification's weight.

\subsection{Hurwitz closes the ladder}

The construction itself would double forever; what stops it is a classification theorem. By
Hurwitz's theorem, the normed division algebras --- those in which the norm composes,
$\lvert ab\rvert = \lvert a\rvert\,\lvert b\rvert$, and no non-zero elements multiply to zero ---
are exactly $\mathbb{R}$, $\mathbb{C}$, $\mathbb{H}$ and $\mathbb{O}$
\autocite{hurwitz1898composition,baez2002octonions}. The next doubling, the sedenions, has zero
divisors: non-zero elements whose product vanishes. The framework's reading of this failure is the
one that matters physically. Zero divisors cannot support consistent information flow around loops
--- a cycle of compositions can silently annihilate what it carries --- and a framework in which
closed computational paths are load-bearing (\S\ref{sec:hopf-justification},
\S\ref{sec:space-of-computations}) cannot be built on an algebra that permits it. At the level of
amplitudes the same failure reads: the composition law
$\lvert ab\rvert = \lvert a\rvert\,\lvert b\rvert$, which carries probability multiplicatively
through every rung of the ladder, fails beyond the octonions, so that two non-vanishing amplitudes
could compose to nothing and probability would be destroyed by mere composition. The last algebra
fit to carry amplitudes is $\mathbb{O}$. One rung past the terminus is not entirely silent,
however: the sedenions' internal structure --- their automorphism group and their octonion
subalgebras --- is the recorded candidate arena for tagging the generations
(\S\ref{sec:interactions}), and it enters exactly as this section's scoping permits, internally
and never as a base. The terminus closes the ladder of forces; what survives past it is, at most,
flavour. This is the algebraic half of the termination; the
geometric half --- that no sphere above $S^{7}$ is parallelisable --- and the theorem joining the
two halves are the opening business of \S\ref{sec:hopf-justification}, and the epistemic reading of the same endpoint,
as a bound on the information a single relation can carry, is given in
\S\ref{sec:hopf-information-bound}.

\subsection{The honest debt}

As in the geometry section, we close by stating plainly what is proved, what is established, and what is
owed. Proved: the construction, the losses, and the termination are theorems --- the Cayley--Dickson
doubling and its ledger of surrendered laws are standard algebra, and Hurwitz's classification is
among the durable results of the subject. Established: the physical cashing of each rung --- the
Born rule at $\mathbb{C}$, weak isospin at $\mathbb{H}$, colour, generations and mass at
$\mathbb{O}$ --- is carried out in the sections cited above, each under its own epistemic flag. Owed,
and stated without disguise: the ladder as presented here is closed from above and derived at its
floor, but it is not yet derived rung by rung from below. Hurwitz forbids a further rung; the
argument of \S\ref{sec:self-reference} forces the first doubling, from the ordered arithmetic of
classical probability to amplitudes that interfere. That the same self-referential closure, applied
to a system whose amplitudes already take values at one rung, compels the next --- that relational
closure at one level forces precisely the Cayley--Dickson double of that level --- is the
framework's reading of the climb, and it is a thesis, not a theorem --- but the formal apparatus
against which it can be posed now exists, and we pose it.

\begin{openproblem}[The compulsion of the intermediate rungs]\label{op:ladder-compulsion}
Derive, within the categorical apparatus of Open Problem~\ref{op:sigma-functor}, that
self-referential closure at one rung compels the Cayley--Dickson double of that rung, so that the
intermediate steps $\mathbb{C} \to \mathbb{H} \to \mathbb{O}$ follow from the same principle
that forces the floor. The candidate route is the construction's own shape: the Cayley--Dickson
pair with its conjugation, $(a, b)^{*} = (a^{*}, -b)$, is an involution-equipped doubling, and the
route reads it as the algebraic shadow of adjoining the describer --- closure of level-$A$
descriptions under the inside/outside exchange requiring exactly the twisted pair. The problem
inherits the status of Open Problem~\ref{op:sigma-functor}, whose apparatus it presupposes.
\end{openproblem}

If the derivation lands, the whole ladder becomes a consequence of the same self-reference that
opened it. Until then, the honest statement is that the ladder's necessity is established at its
ceiling and at its floor, while the compulsion of the intermediate rungs is inherited from the
physics they are shown to carry rather than derived in advance. What this section hands on is, accordingly, an algebraic ladder with a floor, a ceiling and a
ledger. Its first consumer is immediate: \S\ref{sec:space-of-computations} spends the terminal
fact --- the non-associativity of $S^{7}$'s own multiplication --- as geometry, excluding the
seven-sphere as a ruler and fixing the dimension of space at three.
\S\ref{sec:hopf-justification} then supplies the geometry of the fibrations themselves, and the
join.

%  (Architecture call SETTLED by the sync manifest: standalone.)
%
%  CITATIONS -- existing keys (as supplied by JS 2026-07-04):
%    Hardy2001a, Aaronson2004h, Thomson_1954, Lawvere_1969,
%    Turing_1937, muller_p_2011, Brukner2009.
%  NEW keys (RIS in continuity_refs_2026-07-04.ris):
%    cantor1891diagonal, benacerraf1962tasks, broeckertomdieck1985compact.
%  Deliberately excluded per JS: Chiribella--D'Ariano--Perinotti;
%  Bargmann/Wigner's theorem. Yanofsky2003 now paired on the
%  Lawvere-shaped clause (the pending-key excuse was void: the key
%  exists and sec:self-reference uses it).
%
%  CROSS-REFS: sec:self-reference now DRAFTED (2026-07-04, later the
%  same day) -- no longer dangling; only sec:fanout-boundary (and its
%  sibling sec:fanout-direction) -- DRAFTED 2026-07-06.
%
%  AUDITED 2026-07-06 (final pass, second stop): both proofs
%  re-verified by hand (Lagrange divisibility sentence exact; the
%  Baire countable-escape chain sound); caveat-discharge with
%  sec:self-reference confirmed two-way (already present); file pure
%  ASCII; no pre-definition "frame"; all citations used as declared.
%  Fixes: Yanofsky2003 paired on the Lawvere clause; the op's two
%  candidate routes anchored to printed propositions (f1-dictionary
%  for the divisibility premise's referent, inside-outside for the
%  closedness demand) -- statuses unchanged, referents sharpened;
%  header staleness corrected. Deliberate FORWARD references --
%  prop:f1-dictionary and sec:hopf-justification (defined later in the
%  paper; fine under two-pass LaTeX and flagged in prose),
%  sec:division-algebras, sec:space-of-computations.
%
%  DECISIONS ENCODED: op:limit-realisability minted (the owed arrow,
%  priced: exactly two derivations, divisibility and closedness);
%  Lagrange used twice (finite exclusion; component-group triviality);
%  the (1+i)/2 lamp value NOT stated here -- "a superposed fixed point"
%  with the complex-probability treatment deferred to
%  sec:self-reference; the two-branch over-determination stated to
%  pre-empt a circularity reading; Thomson spelling throughout.
% =====================================================================

\section{Continuity Is Not an Axiom}
\label{sec:continuity}

Hardy derives quantum theory from five axioms \autocite{Hardy2001a}. Four of them are satisfied
by classical probability theory; the fifth --- that a \emph{continuous} reversible
transformation connects any two pure states --- is the entire difference, the single postulate
that turns the classical simplex into the complex Hilbert space. Every reconstruction programme
pays an equivalent toll: continuous reversibility in Masanes and M\"{u}ller
\autocite{mullerp2011}, the corresponding postulate in Daki\'{c} and Brukner
\autocite{Brukner2009}. The axiom is load-bearing, and in every programme it is unexplained: it
is where the quantumness is put in by hand. The claim of this section is that in the present
framework it is not an axiom. Continuity is the fifth face of the pre-measurement regime that self-reference forces, and what
stands between the framework and a full derivation is not a postulate about values but one
premise about transformations --- that every transformation the inside performs has a stationary
outside presentation --- derived in outline below from the inside/outside involution of
\S\ref{sec:self-reference} and stated as an open problem at the end of the section with nothing
hidden in it. The continuum itself, the completion that every probability theory adopts, is
named for what it is: the real line, shared with classical probability and carrying no quantum
content.

One structural remark prevents a misreading. The framework reaches complex amplitudes natively in
\S\ref{sec:self-reference}, by the negative-probability route; the present section reaches the same
destination by an independent branch --- the regime's continuity, then Hardy's theorem. The two
branches converging on one structure is over-determination, not circularity, and the convergence is
itself evidence that the destination is forced.

\subsection{Five faces of one regime}

The regime before any measurement --- before any \textsc{fanout} --- can be characterised five
equivalent ways: negative probability is available; anti-events exist, every event being undoable;
no classical record exists; no preferred basis is selected; and a continuous reversible
transformation connects any two pure states. These are not five premises to be combined. They are
one regime under five descriptions, and \textsc{fanout} (\S\ref{sec:fanout-boundary}) is the single
event that fails all five at once: it writes the record, selects the basis, kills the anti-events,
discretises the transformation group, and switches off the negativity. The equivalence of the first
two faces is the anti-event identity of \S\ref{sec:self-reference}; the third is the boundary's
definition. The work of this section is the equivalence of the last two faces, and then the route
from the regime's existence to its continuity. The fourth face has, since
Proposition~\ref{prop:f1-dictionary} (\S\ref{sec:hopf-justification} --- a deliberate forward
reference), an exact form: a basis-selection is the commitment of a classical structure, a lifted
antipodal skeleton, so ``no preferred basis'' means precisely that no copy structure has fired.

\subsection{Basis selection is the collapse of continuity}

\begin{proposition}[Continuity forbids a selected basis]\label{prop:basis-discreteness}
If a finite set of pure states is privileged --- a preferred basis --- then every reversible
transformation respecting the privilege carries the set to itself and acts on it as a permutation;
the respecting transformations therefore realise only a finite group of motions among the
distinguished states, and no continuous family of them connects distinct pure states of the set.
Contrapositively, continuous reversibility entails that no basis is selected.
\end{proposition}

\begin{proof}
A privileged finite set is by definition carried to itself by every structure-preserving reversible
map, whose induced action on the set is a permutation of finitely many points. A continuous path of
such maps induces a continuous path in a finite discrete set, which is constant. Hence a continuous
reversible family moving one distinguished pure state to another respects no such privilege, and if
every pair of pure states is so connected, no finite set is privileged.
\end{proof}

The converse is where the historical difficulty lies: the absence of a privileged finite set does
not by itself force continuity, for the transformation group could in principle be infinite and
discrete. Two structural properties close every such escape at once.

\begin{proposition}[Divisible and closed implies connected]\label{prop:divisible-closed}
Let $G$ be the group of reversible transformations of the pre-measurement regime, acting on the
state body. If $G$ is divisible --- every element has an $n$-th root for every $n$ --- and
topologically closed, then $G$ is a compact connected Lie group; and if $G$ moreover acts
transitively on the pure states, then a continuous reversible transformation connects any two of
them: Hardy's fifth axiom.
\end{proposition}

\begin{proof}
The state body is compact and convex with nonempty interior, so its symmetry group is compact, and
a closed subgroup of a compact matrix group is a compact Lie group with finitely many connected
components \autocite{broeckertomdieck1985compact}. The component group $G/G_{0}$ is a quotient of a
divisible group and hence divisible; but a nontrivial finite group is never divisible --- with
$\lvert H\rvert = m$ every element satisfies $h^{m} = e$, so no element other than the identity
possesses an $m$-th root --- and $G/G_{0}$ is finite, so it is trivial and $G = G_{0}$ is connected.
A connected Lie group is path-connected, and transitivity then supplies the continuous reversible
path between any two pure states.
\end{proof}

The same one-line fact of Lagrange closes the remaining escapes. Finite groups are excluded
directly, being nontrivial and never divisible. Countable groups are excluded because a
countable closed subgroup of a compact group is compact, has an isolated point by the Baire
category theorem, and is therefore discrete and finite. Uncountable disconnected groups are
excluded by the component-group argument of the proposition. Every route to a discrete
transformation group therefore passes through a failure of divisibility or a failure of
closedness. Divisibility, however, asks for more than the conclusion needs, and the next
proposition removes it.

\begin{proposition}[Continuity from a connected state space]\label{prop:connected-transitive}
Let the pure states of a finite-dimensional system form a connected space $P$, and let the group
$G$ of reversible transformations of the pre-measurement regime be topologically closed and act
transitively on $P$. Then the identity component $G_0$ already acts transitively on $P$, and a
continuous reversible transformation connects any two pure states: Hardy's fifth axiom holds,
with no divisibility assumed.
\end{proposition}

\begin{proof}
As in Proposition~\ref{prop:divisible-closed}, closedness makes $G$ a compact Lie group with
finitely many connected components \autocite{broeckertomdieck1985compact}, so $G_0$ is a closed
normal subgroup of finite index and is path-connected. Each $G_0$-orbit in $P$ is the continuous
image of a compact group, hence compact, hence closed. Since $G$ is transitive and $G_0$ is
normal, the $G_0$-orbit of $gx$ is $g$ applied to the $G_0$-orbit of $x$, so the cosets of $G_0$
permute the orbits transitively and there are finitely many of them. A finite partition of $P$
into closed sets is a partition into open sets; $P$ is connected; so there is one orbit and
$G_0$ is transitive. Given pure states $x$ and $y$, choose $g \in G_0$ with $gx = y$ and a
continuous path $g_t$ in $G_0$ from the identity to $g$; then $g_t x$ is a continuous reversible
transformation from $x$ to $y$.
\end{proof}

The proposition relocates the premise. Divisibility was asked for because a transformation group
could be infinite and discrete; but a group acting transitively on a continuum cannot be
countable, and a closed group acting transitively on a connected space cannot be disconnected in
the way that matters, so what the argument needs from the regime is not the availability of
$n$-th roots but the connectedness of the state space itself. That is a statement about values:
that the pure states are all the unit vectors of $\mathbb{C}^{K}$ up to phase, the space
$\mathbb{CP}^{K-1}$, connected for $K \ge 2$, and hence that the amplitudes of
\S\ref{sec:self-reference} fill their continuum. The two standing foils show the premise doing
its work. The stabiliser fragment of quantum theory has complex amplitudes and is contextual,
yet it has finitely many pure states, its outcome probabilities take only the values $0$,
$\tfrac12$ and $1$, and its reversible transformations form a finite group, the Clifford group,
so Hardy's fifth axiom fails in it; Spekkens's toy theory \autocite{spekkens2007toy} is the same
with real amplitudes, its ignorance rationed by a knowledge-balance principle. Neither has a
connected state space, and neither is a model of the regime described here, whose credences take
every value. Divisibility, which both also lack --- the phase gate has no square root in the Clifford group
--- is a symptom of the discreteness, not its cause. What both lack can be said exactly, and the
next subsection says it: each contains a transformation that fixes no state it admits.

\subsection{The premises, and what they mean}

Three premises carry Proposition~\ref{prop:connected-transitive}. Transitivity is the regime's
homogeneity: inside/outside equivalence leaves no pure state privileged, and a proper closed
invariant set of pure states would define a superselection observable, a classical record
predating any \textsc{fanout} (\S\ref{sec:fanout-boundary}). Connectedness of the state space is two statements, and only one of them is quantum. The first
is that the pure states are infinite and dense, with a fixed point for every transformation; the
second is that the state space is complete. Quantum theory over the algebraic numbers ---
amplitudes in the algebraic closure of $\mathbb{Q}$ within $\mathbb{C}$ --- has the first
without the second: it has $K = N^{2}$, the Born rule, superposition, entanglement and Bell
violation at the Tsirelson value, every unitary has an algebraic eigenvector, and it is
countable and totally disconnected. So the whole quantum content of Hardy's axiom lies in the
density, and the completion is the real-number convention of every probability theory. Hardy
adopts it before his fifth axiom is stated, in the first, which makes probabilities real numbers
and gives his classical theory a real simplex of mixtures; this framework adopts it at the
Shannon step, where credences are declared to range over a connected domain so that the
continuation of the logarithm is unique (\S\ref{sec:complex-amplitudes}). Given that convention,
density of the pure states is connectedness of their closure. Divisibility, which an earlier
version of this section asked for, is not needed: it follows from the conclusion, every element
of a compact connected Lie group having roots of every order.

Closedness is the completion, limit-realisability: a Cauchy sequence of available
transformations has its limit available. It cannot be derived, and the reason should be stated
rather than hidden. Any collection of first-order principles about states, transformations and
probabilities that has an infinite model has a countable one, and a countable subset of
$\mathbb{CP}^{K-1}$ is totally disconnected; the founding non-computability does not escape
this, since the arithmetical complex numbers form a countable algebraically closed field
containing non-computable numbers and closed under every limit it can itself describe. Hardy's
axiom and closedness are second-order, quantifying over paths and over sequences, and a
second-order principle is grounded as Dedekind completeness is grounded or not at all. The
grounds are these. The completion of a dense set under the finite-precision metric is unique and
depends on no choice, whereas any countable alternative is a choice of definability level, an
extra law with no ground and no test; and the framework's own result is that no observer's
descriptive level captures the regime, since by \S\ref{sec:self-reference} every observer's
state exceeds its own level and the inside/outside principle makes that a law rather than a
perspective. The burden therefore falls on any restriction, not on the completion, which is the
least committal option and carries no quantum content. What limit-realisability adds to the line
is still worth stating exactly, because it is the corrected form of a natural intuition.
Cantor's diagonal \autocite{cantor1891diagonal} does not build the continuum; its output is an
infinite decimal --- a Cauchy sequence of rationals --- which exists \emph{as a point} only
because the line is complete. The diagonal presupposes limit-realisability and then shows that
the completed line outruns every internal listing. And the character of the completion is the
framework's own: the computable reals are countable \autocite{Turing_1937}, hence of measure
zero, so once the transformation group contains its limits almost all of them are
non-computable. Non-computability is a property of the completed continuum, not the reason for
it. The constructive reading, on which only describable limits exist, remains consistent; it is
the countable theory above, and the two readings differ in nothing an experiment can certify.

Density is where the section's title is earned, and it is derived from the paper's own premise.
Write $(\mathrm{I})$ for the statement that every admissible transformation of the regime is a
phase of some admissible classical structure: diagonal, that is, in the record basis of some
observer the regime contains, a phased spider of that basis in the calculus of
\S\ref{sec:hopf-justification} \autocite{coeckeduncan2011interacting} (VERIFY-SUPPORT). It says
of transformations what the global phase of \S\ref{sec:self-reference} says of states: there are
no absolute ones. It follows from three things the paper already holds. The inside/outside
involution is a genuine involution, so the outside presentation of the inside is a presentation
of the same object with nothing lost: the outside can have complete knowledge of the inside, a
potential record, the joint state lying in a structure the outside carries. Time is operation
count, internal to the system that operates (\S\ref{sec:space-of-computations}), so the outside
does not share the inside's count. And before any \textsc{fanout} that count is not classically
recorded, so the outside's description is pure. Let $g$ be the transformation the inside
performs per operation. The outside's complete knowledge cannot depend on the inside's count,
and it is pure, so it is a pure state of the inside's clock and system that is stationary under
``advance the clock and apply $g$''; the stationary states of that map are clock Fourier modes
tensored with eigenvectors of $g$ and nothing else \autocite{pagewootters1983evolution}
(VERIFY-SUPPORT). So the outside holds a structure in which $g$ acts by phases, which is
$(\mathrm{I})$; and the relative phase in that structure is the outside's ignorance of the
inside's computational count, the same ignorance the global phase records one level down. The
consequences are a theorem. A phase of a structure fixes the structure's copyable states, so
every admissible transformation fixes an admissible pure state; a finite transitive action has
an element moving every point (Jordan), so the pure states are infinite; the transformation
carrying a stabiliser state to the fixed point of the composite of a $Z$-phase and an $X$-phase
turns through $\arccos(1/\sqrt{3})$, an irrational fraction of a turn, so the group has an
element of infinite projective order and its orbits are dense. With the completion,
Proposition~\ref{prop:connected-transitive} fires and Hardy's fifth axiom holds. The two foils
fail $(\mathrm{I})$ exactly where they fail continuity: fourteen of the twenty-four Clifford
operations are phases of no Pauli basis and fix no stabiliser state, their fixed points being
the magic states; the toy theory's three-cycles are phases of no epistemic basis; and in each
the outside's knowledge of the inside under the order-three protocol is either count-dependent
or a mixture over the count, so each makes the inside's time classical before any record exists.
The derivation of $(\mathrm{I})$ is an outline, not a proof, and it is what the open problem
below asks for.

\subsection{Why the cut matters physically}

The regime's continuity is not decorative. Reversibility alone, over complex amplitudes, licenses
the full invertible group $GL(n, \mathbb{C})$, and Aaronson showed what that buys: superluminal
signalling, and computational power at the level of \textsf{PP} \autocite{Aaronson2004h}. The cut
from $GL(n, \mathbb{C})$ down to the unitary group is the conserved quadratic form, and continuity
is what entitles the theory to make it: with the fifth axiom in hand, Hardy's theorem delivers the
complex Hilbert space, the quadratic norm and the unitary group in one stroke \autocite{Hardy2001a}
--- cited here, not re-proved. One honest caveat travels with the citation. Hardy's reconstruction
uses his remaining axioms besides --- composition and subspace structure --- and these must be
derived rather than imported: the composition ledger of \S\ref{sec:division-algebras} and the
locality-of-tomography argument of \S\ref{sec:self-reference} carry that duty, and it is discharged
there, not silently assumed here.

\subsection{Super-tasks at the boundary}

The premise just priced has a fifty-year-old stress test in the philosophy of mathematics, and the
framework passes it in an unexpected way. A super-task --- infinitely many operations completed in
finite time \autocite{Thomson_1954} --- presupposes its terminal moment as an actual point of the
timeline: the accumulation point of a generic sequence of step-times exists only in the completed
line, and the defining question, the state \emph{after all the steps}, is meaningful only under
limit-realisability of time. Super-tasks are downstream consumers of exactly the completeness at
issue in this section. Thomson's lamp then exhibits both sides of the \textsc{fanout} boundary at
once. On the record side, the lamp's description fixes its state at every finite stage and thereby
fixes no limit at all: any terminal value is consistent with the whole task
\autocite{Thomson_1954,benacerraf1962tasks}. That is not a puzzle about lamps; it is the proof that
the classical record algebra is not limit-closed, and cannot be. On the amplitude side the limit
exists --- a superposed fixed point, whose complex-probability treatment belongs to
\S\ref{sec:self-reference} --- because that side \emph{is} limit-closed, which is the fifth face
itself. The boundary is thereby the edge of limit-realisability: complete on the reversible side,
incomplete on the record side, with the lamp standing on both sides at once as the boundary's oldest
exhibit. A final clause dissolves rather than answers the classic question. In this framework time
is operation count (\S\ref{sec:space-of-computations}), and ``after infinitely many operations''
indexes no operational moment: the lamp packs its endless \textsc{fanout}s into the emergent
continuum, and asking for its record at the accumulation point asks the discrete contents for a
value that only the completed container possesses. On or off is not undecided; it is unasked.

\subsection{What remains, priced}

\begin{openproblem}[Every transformation a phase of an inside structure]\label{op:limit-realisability}
Proposition~\ref{prop:connected-transitive} reduces Hardy's fifth axiom to three premises:
connectedness of the pure states, transitivity, and closedness. Transitivity is argued above
from the absence of records. Closedness is the completion, adopted on the grounds stated above
and carrying no quantum content. Connectedness is density plus that completion, and density
follows, by Jordan's theorem and the infinite-order argument, from $(\mathrm{I})$: that every
admissible transformation of the regime is a phase of some admissible classical structure. What
remains open is the rigour of the derivation of $(\mathrm{I})$ from the paper's premise. Derive,
inside the fixed-point locus $\mathrm{Fix}(\sigma)$ of Proposition~\ref{prop:inside-outside},
that every endomap is a phase of a classical structure in $\mathrm{Fix}(\sigma)$, from the
self-representing object, the involution $\sigma$ as a genuine involution, and the reading of
the inside presentation as a sequence indexed by operation count with no record before
\textsc{fanout}. The outline above does this with the stationary-state construction of Page and
Wootters; a derivation passes if it makes that outline a theorem of the calculus, with the
outside's count-independent complete description shown to exist and to be pure from the
involution alone. It fails before any detail is checked if every premise it uses holds in the
stabiliser fragment of quantum theory, in Spekkens's toy theory \autocite{spekkens2007toy}, in
the real qubit, or in quantum theory over the Gaussian rationals, all of which violate
$(\mathrm{I})$ outright and all of which make the inside's operation count classical before any
record exists. It also fails if it imports the completion as a premise about transformations:
the completion is the real line and has no derivation, and a derivation of $(\mathrm{I})$ must
hold in the countable algebraic theory, where $(\mathrm{I})$ is true and continuity is not.
\end{openproblem}

The accounting, in one sentence: the axiom that every reconstruction programme pays is, in this
framework, one derivation posed, of Lawvere's shape --- that every endomap of the fixed-point
locus is a phase of a classical structure in it --- and one convention named, the completion,
which classical probability pays too. And
the hand-on is already visible: the complex amplitudes that this section's continuity underwrites
through Hardy's theorem, and the complex amplitudes that \S\ref{sec:self-reference} derives natively,
are one and the same $\mathbb{C}$ --- the floor of the ladder that \S\ref{sec:division-algebras}
climbs.

%      % EDITED 2026-10-29 scripts/patch_continuity_tower_premise_2026-10-29T1500.py
% EDITED 2026-09-15 scripts/patch_value_continuum_declared_cencov_2026-09-15T1540.py
% EDITED 2026-09-15 scripts/patch_continuity_connected_state_space_2026-09-15T1430.py
% =====================================================================
%  SECTION: The Euclidean metric on the state sphere, and the global
%  phase as self-referential ignorance
%  Origin: drafted in an earlier conversation (2026-06-07), edited by
%  JS; audited into the current architecture 2026-07-06.
%
%  PLACEMENT in main.tex (reordered chain, JS 2026-07-06): position 3 --
%      % EDITED 2026-10-29 scripts/patch_continuity_tower_premise_2026-10-29T1500.py
% EDITED 2026-09-20 scripts/patch_born_cost_ledger_v2_2026-09-20T0900.py
% =====================================================================
%  SECTION: Self-Reference: One Principle
%  Drafted 2026-07-04. The paper's foundational section; resolves the
%  largest dangling label in the manuscript (sec:self-reference).
%
%  PLACEMENT in main.tex: FIRST of the foundations chain --
%      \input{self_reference_section_2026-07-04}
%      \input{continuity_section_2026-07-04}
%      \input{division_algebras_section_2026-07-03}
%      \input{hopf_justification_section_2026-06-08}
%
%  SOURCES (re-read 2026-07-04): lawvere_inside_outside_fixed_point_
%  theorem_2026-05-04T1806.md (the spine: the two-principles reduction,
%  the theorem, the sigma candidate, the open items -> op minted per
%  D9); research_summary_neg_prob.md (negativity route; lamp (1+i)/2);
%  dynamics_companion_fanout_unitarity (legs 1a/1c; anti-events; the
%  Brukner-Zeilinger rebuttal per D6); Pure_states_and_mixed_states
%  (local tomography derived; QBism contrast per D7).
%
%  DECISIONS ENCODED (JS 2026-07-04): D1 title as below; D2 Wheeler1990
%  cited here as the title anchor; D3 lineage Popper -> Breuer ->
%  Svozil -> Szangolies (Dalla Chiara omitted); D4 Frauchiger_2018
%  cited here, dissolved in sec:wigners-friend; D5 Kirkwood-Dirac
%  heritage via the Arvidsson-Shukur 2024 review as one-stop; D6
%  Brukner-Zeilinger + the framework rebuttal IN ("Shannon is
%  post-FANOUT"); D7 QBism contrast IN; D8 Aaronson private
%  communication OMITTED (published island paper only); D9 one op
%  minted (op:sigma-functor), cross-linked to op:hopf-frobenius.
%
%  CITATIONS -- existing keys: Wheeler1990, small_2006, Godel1931,
%  Lawvere_1969, Yanofsky2003, CoeckeKissinger2017, Rovelli1996a,
%  Popper1950a, Popper1950b, Breuer1995, svozil_k_2018,
%  szangolies_j_2018, Frauchiger_2018, dirac1942, Feynman_r_82,
%  Feynman1987, Youssef1994 [SWAP-FLAG: intended paper is "Quantum
%  mechanics as complex probability theory", Mod. Phys. Lett. A 9
%  (1994) 2571 -- if that lives under youssef1994quantum, swap],
%  caticha_a_2006, goyal_knuth_skilling_2010, Aaronson2004h, Hardy2001a,
%  muller_p_2011, Thomson_1954, benacerraf1962tasks, cantor1891diagonal,
%  Turing_1937.
%  NEW keys (RIS in self_reference_refs_2026-07-04.ris, all verified
%  this session): tikochinsky1988feynman, renou2021real, chen2022ruling,
%  wu2022refutation, hoffreumonwoods2026real, arvidssonshukur2024kirkwood,
%  bruknerzeilinger2001shannon, fuchsschack2013qbism,
%  shannon1948mathematical.
%
%  CROSS-REFS: forward (fine under two-pass LaTeX): sec:continuity,
%  sec:division-algebras, sec:hopf-justification, prop:f1-dictionary,
%  sec:wigners-friend, op:limit-realisability, op:hopf-frobenius.
%  Still dangling after this section: sec:fanout-boundary (and its
%  sibling label sec:fanout-direction) -- both DRAFTED 2026-07-06.
%
%  AUDITED 2026-07-06 (section-by-section final pass, first stop):
%  proof of prop:inside-outside re-verified by hand; attribution chain
%  clean (Burgin/Khrennikov correctly absent); all header-declared
%  citations confirmed used; lamp arithmetic and antecedents verified;
%  FANOUT smallcaps consistent; no pre-definition use of "frame".
%  Fixes: stray Unicode em-dash (QBism paragraph, the file's only
%  non-ASCII character) -> house ---; trace-blind two-way stitch added
%  to the Born subsection ("one privacy, two duties");
%  op:sigma-functor gains its first standing constraint (F1
%  inheritance via the op:hopf-frobenius discharge) and its previously
%  one-directional local-tomography cargo made explicit.
% =====================================================================

\section{Self-Reference: One Principle}
\label{sec:self-reference-2}

Wheeler asked for it from bit: for physics derived from information, ``every it --- every particle,
every field of force'' obtained from binary answers to yes-or-no questions \autocite{Wheeler1990}.
This paper's route to that programme runs through G\"{o}del: the single principle from which
everything else here is derived, deferred, or priced is \emph{self-reference} --- the fact that any
complete physical description, made from inside the world it describes, must describe the describer.
This section states the principle, proves that the framework's second historical commitment follows
from it as a theorem rather than standing beside it as an axiom, and then extracts the price of the
principle in four instalments: probabilities whose parts can cancel, amplitudes valued in
$\mathbb{C}$, the Born rule, and the local tomography of composite systems. What the section cannot
yet derive is stated as precisely as what it can, in one numbered problem and in the ledger at the
end.

\subsection{Events, not objects}

The framework's ontology is relational: the fundamental entities are events and the relations
between them, and objects are the bookkeeping of relational invariants --- what remains stable
across many relations. The kinship with relational quantum mechanics \autocite{Rovelli1996a} is
genuine and acknowledged: facts are facts relative to the systems that register them. The
difference, developed over the rest of this paper, is that the present framework grounds the
relational structure in a single logical principle and then carries it, through the division
algebras and the Hopf fibrations, all the way to the gauge group and the particle spectrum ---
relations are not only the ontology but the derivation engine. On this reading a measurement is not
the revelation of a pre-existing property; it is the creation of a relation, and the later sections
of this paper read that creation geometrically (\S\ref{sec:hopf-justification}) and even
gravitationally. Here we need only the logical core: whatever describes the world is an event in
the world it describes.

\subsection{One principle, and a theorem}

The framework historically carried two foundational commitments: self-reference, and the
\emph{Inside/Outside Equivalence Principle} --- that the laws of physics must be identical whether
formulated by an embedded participant or by an external observer. Two principles are one too many,
and this subsection reduces them.

The reduction corrects the framework's own origin, and we make the correction in print. The 2005
form of the argument \autocite{small_2006} ran a G\"{o}delian chain: a self-referential process,
being complete, must by the incompleteness theorem \autocite{Godel1931} be inconsistent; an
inconsistent system can negate itself; and the negation of self-reference is the production of
causal sequence. The chain captures the right intuition and is retained here as motivation, but it
is not sound as a proof: classical inconsistency yields \emph{ex falso quodlibet} --- every
proposition provable, every causal sequence an output, no constraint at all. The load-bearing
argument is different, and better: a fixed-point theorem.

Let $\mathcal{C}$ be the category of physical descriptions and admissible processes --- concretely,
the dagger-compact-closed setting with classical structures in which categorical quantum mechanics,
and the whole of this paper, works \autocite{CoeckeKissinger2017}. Lawvere showed that in a
sufficiently rich such setting, any weakly point-surjective self-representation forces every
endomap to have a fixed point, and Yanofsky showed that one diagonal argument of this shape
underlies Cantor's, G\"{o}del's, Turing's and Tarski's theorems alike
\autocite{Lawvere_1969,Yanofsky2003,cantor1891diagonal}. The framework's single commitment is that
$\mathcal{C}$ contains a \emph{self-representing object} --- an object capable of encoding the
processes on itself. That is what ``self-reference is foundational'' means, stated categorically,
and it is the surjectivity hypothesis of the diagonal lemma. Now let
$\sigma \colon \mathcal{C} \to \mathcal{C}$ be the \emph{inside/outside involution}: the operation
exchanging the two complementary presentations of any description --- the inside presentation, as
participant-indexed total space in which the observer is part of the system, and the outside
presentation, as a base-times-fibre decomposition over a classical parameter space. That
$\sigma^{2} = \mathrm{id}$ is not an axiom but a coherence requirement: the outside view of the
inside view of a thing is the thing. The candidate concrete realisation of $\sigma$ is the
Hopf-duality swap --- total space against base-times-fibre for the three fibrations --- which gives
the involution computable content and ties this subsection to \S\ref{sec:hopf-justification}.

\begin{proposition}[Inside/outside fixed point]\label{prop:inside-outside}
Let $\mathcal{C}$ satisfy the richness condition above and contain a self-representing object, and
let $\sigma$ be an involution on $\mathcal{C}$. Then the fixed-point locus
$\mathrm{Fix}(\sigma)$ is non-empty: there exist descriptions invariant under the exchange of
inside and outside perspective.
\end{proposition}

\begin{proof}
Lawvere's construction, applied to $\sigma$ in place of negation: given a weakly point-surjective
$\varphi \colon A \to Y^{A}$, form $g(a) = \sigma\!\left(\varphi(a)(a)\right)$; point-surjectivity
supplies $a^{*}$ with $\varphi(a^{*}) = g$, and then $\varphi(a^{*})(a^{*})$ is a fixed point of
$\sigma$ \autocite{Lawvere_1969,Yanofsky2003}. Where the same construction applied to Boolean
negation yields paradox --- negation has no fixed point on a two-element type, so the conclusion
explodes the premises --- the inside/outside involution is an honest involution on a rich type, and
the conclusion is a constructive fixed point rather than a contradiction.
\end{proof}

\textbf{Definition.} \emph{Physical law is the structure carried by} $\mathrm{Fix}(\sigma)$.

The Inside/Outside Equivalence Principle is thereby demoted from axiom to theorem: physical law is
invariant under the exchange of perspectives \emph{because it is the fixed-point set of that
exchange}, and the fixed-point set is forced to be non-empty by self-reference alone. One
commitment remains foundational; everything the second commitment asserted is now derived. The
reader who accepts categorical quantum mechanics as an existing programme has already accepted the
machinery; the framework adds one premise and collects the rest. And the fixed-point locus is not
empty of structure: it inherits the Frobenius-algebra apparatus of $\mathcal{C}$, so that the
Bell, GHZ and W primitives --- and with them the particle classification of the later sections ---
live inside $\mathrm{Fix}(\sigma)$ from the start. What is proved above is non-emptiness given the
setting; what is assumed is the setting itself and the existence of $\sigma$ as a genuine
involution; what is open is the rigour of both, and it is priced now rather than hidden:

\begin{openproblem}[The involution made functorial]\label{op:sigma-functor}
Fix the category $\mathcal{C}$ rigorously --- the dagger-compact-closed setting with classical
structures is the candidate --- and define $\sigma$ on morphisms, not only on objects, the natural
route being the induced fibre--base exchange on spider diagrams; establish that $\sigma$ is an (at
worst pseudo-)involutive endofunctor; and prove or refute the identification of
$\mathrm{Fix}(\sigma)$ with the three Hopf cross-sections, under which Adams's theorem would become
a corollary of the fixed-point classification. This problem is cross-linked with Open
Problem~\ref{op:hopf-frobenius}: under an affirmative resolution of both, the parallelisable
terminus, the fixed-point classification, and compositional completeness are three statements of
one fact, and progress on either problem constrains the other. One constraint is already in place:
at the complex rung, Proposition~\ref{prop:f1-dictionary} supplies the dictionary between
classical structures and the fibration's lifted skeletons --- the clause of
Open Problem~\ref{op:hopf-frobenius} it discharges --- and any resolution here must reproduce it.
The problem also carries a piece of cargo recorded later in this section: the formalisation, within
the fixed category, of the perspective-accounting that derives local tomography.
\end{openproblem}

\subsection{The lineage}

The obstruction this section builds on is old, and the framework claims its completion, not its
discovery. Popper argued in 1950 that a predictor embedded in the world has irreducibly inferior
knowledge of itself, because it cannot impart information about itself without, by that very act,
rendering the information out of date \autocite{Popper1950a,Popper1950b} --- the passage is quoted
where the framework makes the obstruction structural (\S\ref{sec:wigners-friend}). Breuer proved
the impossibility of accurate state self-measurement within quantum theory \autocite{Breuer1995};
Svozil's monograph surveys the whole territory where physics meets self-reference, undecidability
and Chaitin-style algorithmic limits \autocite{svozil_k_2018}; Szangolies sharpened the modern form
as epistemic horizons \autocite{szangolies_j_2018}; and the Frauchiger--Renner theorem is the
contemporary crisis-form --- nested observers deriving contradictory records when each treats the
others' quantum descriptions as universally valid \autocite{Frauchiger_2018} --- which
\S\ref{sec:wigners-friend} dissolves relationally rather than paradoxically. This framework's
addition to the lineage is twofold. Structurally, Proposition~\ref{prop:f1-dictionary} converts
Popper's temporal obsolescence into algebraic cancellation: the datum a copier cannot put into its
own record is exactly its fibre commitment, which does not arrive stale --- it cancels. And
foundationally, the present section runs the lineage in reverse: where the tradition treats
self-reference as a limit on knowledge \emph{within} physics, the framework takes it as the
generator of physics, and derives the quantum formalism from it.

\subsection{The price is negative probability}\label{sec:complex-amplitudes-2}

What does it cost a probabilistic description to contain its own describer? The first instalment is
the sign. In the regime before any record exists, every event has an anti-event --- the description
must be able to undo what it has provisionally done, or the self-description could never be
revised into consistency --- and a probability calculus closed under undoing cannot be confined to
$[0,1]$: it is a signed calculus, in which the parts of a probability can oppose and cancel. The
framework did not discover negativity as bookkeeping; Dirac admitted negative probabilities as
formal intermediaries \autocite{dirac1942}, Feynman defended them as consistent whenever the final,
recorded probabilities stay positive \autocite{Feynman_r_82,Feynman1987}, and the Kirkwood--Dirac
distribution has carried complex quasiprobability from 1933 to the present, as a recent
comprehensive review documents \autocite{arvidssonshukur2024kirkwood}. What the framework adds is
the forcing: negativity is not a permissible convenience but the first price of self-referential
closure, and it is charged exactly and only in the pre-record regime --- which is why every
negativity in the tradition lives in intermediate, unrecorded quantities and never in a final
record.

The second instalment is complexity. A signed information measure must extend Shannon's
$S = -\ln p$ off the positive axis, and the analytic continuation of the logarithm carries $p$ into
the complex plane: amplitudes, with the phase as the continued information. Here too the extension
has a prior tradition, and the attributions are owed precisely. Tikochinsky derived Feynman's rules
for complex amplitudes from consistency of composition \autocite{tikochinsky1988feynman}; Youssef
reformulated quantum mechanics as complex probability theory outright \autocite{Youssef1994};
Caticha carried objective amplitudes to Bayesian probabilities \autocite{caticha_a_2006}; and
Goyal, Knuth and Skilling derived Feynman's rules from a pair-postulate without assuming
analyticity \autocite{goyal_knuth_skilling_2010}. The extension is prior art. The framework's
claim is the compulsion: the complex extension is not a modelling choice among exotic probability
theories but the unique continuation forced by self-referential closure --- the same forcing that
Aaronson glimpsed when he argued quantum mechanics is what negative probability inevitably becomes
\autocite{Aaronson2004h}.

Why does the continuation stop at $\mathbb{C}$? Downward, the real field is not available: real
quantum theory is experimentally distinguishable from complex quantum theory in network scenarios
and has been disfavoured by experiment
\autocite{renou2021real,chen2022ruling,wu2022refutation}; the recent objection that the exclusion
rests on locally tomographic composition \autocite{hoffreumonwoods2026real} is, for this framework,
not an escape but a confirmation of where the weight sits --- local tomography is \emph{derived}
below, not assumed, so the framework closes from its own resources exactly the door the objection
holds open. Upward, the quaternions and octonions fail as a \emph{base} for the information
continuation: the continued measure is a logarithm, and beyond $\mathbb{C}$ the logarithm of a
product is no longer the sum of logarithms --- non-commutativity severs
$\log(ab) = \log a + \log b$ --- so the additivity that makes it an information measure is lost.
Aaronson's parameter-counting and signalling arguments against a quaternionic base point the same
way \autocite{Aaronson2004h}. The higher algebras are not banished; they return as the
\emph{internal} algebras of the register, which is the business of \S\ref{sec:division-algebras}.
The base is $\mathbb{C}$, for informational reasons, and stays $\mathbb{C}$.

\subsection{The Born rule from fibre privacy}\label{sec:born-rule}

The third instalment is the quadratic rule. An amplitude carries one datum more than any record of
it can: its overall phase. This was asserted throughout the framework's development as the
observer's private, self-referential ignorance; it is now a proved statement, in exactly the form
needed. By Proposition~\ref{prop:f1-dictionary}, a classical structure --- the copy operation any
record requires --- is committed to a specific lift of its skeleton, and that commitment is absent
from every classical record the structure induces: the copier is fibre-committed, the record is
fibre-blind, and the fibre freedom is gauge for every observer. (The computation extends verbatim
from one qubit to the full register, where Proposition~\ref{prop:trace-blind} makes it the ground
of hypercharge's tracelessness --- one privacy, two duties.) Popper's obsolescence, made
cancellation. It follows that anything recordable is invariant under the fibre --- observables and
their statistics must be $U(1)$-invariant functions of the amplitude. The lowest-degree polynomial invariant
of a complex amplitude is the quadratic $|\psi|^{2}$ --- the non-polynomial invariant $|\psi|$,
of degree one, is excluded by the analyticity declared at \S\ref{sec:complex-amplitudes} --- and
normalisation fixes the coefficient:
probability is the minimal record-compatible function of the amplitude,
\[
p \;=\; |\psi|^{2},
\]
the Born rule as the cheapest invariant the fibre permits. The same square is reached from the
cost side in \S\ref{sec:born-cost-ledger}, where a record is the closed loop of a history and its
return leg and the exponent is the number of legs; neither route uses the lattice of projections
that Gleason's theorem needs, and both hold for a single qubit. The derivation is the framework's
own;
its premise --- that the fibre datum is structurally excluded from records --- is
Proposition~\ref{prop:f1-dictionary}, and its geometric reading --- probability as base-point data
under the Hopf projection --- is developed in \S\ref{sec:hopf-justification}.

One standing objection to information-theoretic foundations is answered here, because this
framework needs the answer. Brukner and Zeilinger argued that Shannon information is conceptually
inadequate for quantum measurements: it is basis-dependent, and its operational meaning presupposes
ignorance of a definite but unknown fact, which the quantum case does not supply
\autocite{bruknerzeilinger2001shannon}. The framework's reply is structural, and it is a
demarcation rather than a rebuttal. Shannon's precondition fails exactly where they say it fails
--- in the pre-record regime, where there is no definite unknown to be ignorant of --- and there
the operative measure is the quadratic invariant just derived. Shannon's axioms hold exactly where
the precondition holds --- on records, after \textsc{fanout}, where outcomes are definite and
ignorance is classical --- and there the uniqueness theorem forces $-\log p$
\autocite{shannon1948mathematical}. The basis-dependence dissolves with the same stroke: a basis
\emph{is} selected precisely at \textsc{fanout}, so the basis-dependent measure applies only where
a basis exists. Shannon is not inadequate; Shannon is post-\textsc{fanout}. The framework thereby
keeps both measures, each on its own side of the boundary, and uses the additivity of the
post-record measure --- the very property the previous subsection used to pin the base field ---
with a clear conscience.

\subsection{Thomson's lamp, resolved}

The section owes \S\ref{sec:continuity} a debt incurred there: the lamp's amplitude-side limit.
Thomson's lamp --- toggled at times $1 - 2^{-n}$, its state at $t = 1$ demanded --- has no
classical answer: the task description fixes every finite stage and thereby fixes no limit at all
\autocite{Thomson_1954,benacerraf1962tasks}, which \S\ref{sec:continuity} read as the proof that
the record algebra is not limit-closed. On the amplitude side the limit exists, and the calculus of
this section computes it. The toggle exchanges the two amplitudes; the self-consistent state of the
completed task is the fixed point of the exchange in the complex calculus, and solving it under the
signed continuation gives
\[
P(\mathrm{on}) = \tfrac{1 + i}{2}, \qquad P(\mathrm{off}) = \tfrac{1 - i}{2},
\]
a superposed fixed point whose imaginary part carries the unrecordable datum --- the phase of the
endless alternation, when the lamp \emph{would have} stopped. Two checks certify that the
assignment is calculus rather than decoration. The pair is normalised,
$P(\mathrm{on}) + P(\mathrm{off}) = 1$; and the Born rule of the previous subsection, applied to
the fixed point, returns $|P(\mathrm{on})|^{2} = \tfrac{1}{2}$ --- exactly the symmetric classical
ignorance that is the only sane record-level answer, recovered rather than imposed. The lamp,
which entered the framework at its very beginning as the place where classical computation meets
its boundary, thus exits it as a worked example: undetermined as a record, determinate as an
amplitude, the two sides split precisely along \textsc{fanout} --- and the diagonal construction it
physically instantiates \autocite{cantor1891diagonal,Turing_1937} is the same one that powers
Proposition~\ref{prop:inside-outside}. The status flag is owed here as everywhere: the
undetermination is theorem-level; the specific fixed-point value is the framework's complex
calculus, structural, offered with its two consistency checks shown.

\subsection{Local tomography, derived}\label{sec:local-tomography}

The fourth instalment concerns composites, and it discharges two promises at once. The
Inside/Outside Equivalence Principle --- now a theorem --- requires that an external observer's
characterisation of a composite system be recoverable by the embedded participants, and the
embedded participants have access only to local measurements and their correlations. Equivalence
of the two perspectives therefore entails that the state of a composite is determined by the
statistics of local measurements: \emph{local tomography}, obtained as a consequence rather than
imposed as an axiom. This converts a postulate of the reconstruction programmes
\autocite{Hardy2001a,mullerp2011} into a theorem-shaped consequence of the one principle ---
discharging the caveat recorded in \S\ref{sec:continuity} --- and it closes, from the framework's
own resources, the escape route in the real-number debate noted above: the assumption the
objection of \autocite{hoffreumonwoods2026real} contests is here not an assumption. The status is
flagged plainly: the derivation is structural, one paragraph of perspective-accounting, and its
formalisation within the categorical setting belongs to Open Problem~\ref{op:sigma-functor}.

\subsection{Contrast: QBism, and the ledger}

The nearest neighbour in spirit deserves a paragraph of honest contrast. QBism
\autocite{fuchsschack2013qbism} holds, as this framework does, that quantum probabilities are an
agent's --- that the theory is a calculus of expectation, not a catalogue of properties. What QBism
does not say is why the agent's expectations must be complex-valued and quadratically squared; the
Born rule enters as a normative addition to Bayesian coherence, and the complex amplitude as an
unexplained representation. This section is, in effect, the missing why: the agent's ignorance is
self-referential --- the agent is inside the description --- and self-referential ignorance is
complex, with the quadratic rule as the cheapest record-compatible invariant of it. Where QBism
stops at the agent, the framework derives the agent's calculus.

The ledger, then. Delivered in this section: the reduction of two principles to one
(Proposition~\ref{prop:inside-outside}); the pre-record regime with its anti-events, handed to
\S\ref{sec:continuity} as the five-face regime whose continuity is priced there
(Open Problem~\ref{op:limit-realisability}); the complex base, forced downward by experiment and
upward by the additivity of information; the Born rule from fibre privacy; the lamp's fixed point;
and local tomography. Owed, and stated where they are owed: the functorial rigour of the involution
and the fixed-point classification (Open Problem~\ref{op:sigma-functor}); the remaining premise of continuity (Open Problem~\ref{op:limit-realisability}, in
\S\ref{sec:continuity}); and the
rung-by-rung compulsion of the algebra ladder, posed as Open
Problem~\ref{op:ladder-compulsion} in \S\ref{sec:division-algebras}. The
hand-on is the paper's spine: the $\mathbb{C}$ this section derives natively and the $\mathbb{C}$
that \S\ref{sec:continuity} underwrites through Hardy's theorem are the same field, the floor of
the ladder that \S\ref{sec:division-algebras} climbs and \S\ref{sec:hopf-justification}
geometrises. One principle; everything else is theorem, derivation, or a debt with a stated
invoice.

%      % EDITED 2026-10-29 scripts/patch_continuity_tower_premise_2026-10-29T1500.py
% EDITED 2026-09-15 scripts/patch_value_continuum_declared_cencov_2026-09-15T1540.py
% EDITED 2026-09-15 scripts/patch_continuity_connected_state_space_2026-09-15T1430.py
% =====================================================================
%  SECTION: Continuity Is Not an Axiom
%  Drafted 2026-07-04 from continuity_limit_realisability_2026-07-04.md
%  (now in project documents) and the May record
%  (dynamics_companion_fanout_unitarity_2026-05-20T1430.md).
%
%  PLACEMENT in main.tex: a standalone section between self-reference
%  and division-algebras --
%      \input{self_reference_section_2026-07-04}
%      \input{continuity_section_2026-07-04}
%      \input{division_algebras_section_2026-07-03}
%  (Architecture call SETTLED by the sync manifest: standalone.)
%
%  CITATIONS -- existing keys (as supplied by JS 2026-07-04):
%    Hardy2001a, Aaronson2004h, Thomson_1954, Lawvere_1969,
%    Turing_1937, muller_p_2011, Brukner2009.
%  NEW keys (RIS in continuity_refs_2026-07-04.ris):
%    cantor1891diagonal, benacerraf1962tasks, broeckertomdieck1985compact.
%  Deliberately excluded per JS: Chiribella--D'Ariano--Perinotti;
%  Bargmann/Wigner's theorem. Yanofsky2003 now paired on the
%  Lawvere-shaped clause (the pending-key excuse was void: the key
%  exists and sec:self-reference uses it).
%
%  CROSS-REFS: sec:self-reference now DRAFTED (2026-07-04, later the
%  same day) -- no longer dangling; only sec:fanout-boundary (and its
%  sibling sec:fanout-direction) -- DRAFTED 2026-07-06.
%
%  AUDITED 2026-07-06 (final pass, second stop): both proofs
%  re-verified by hand (Lagrange divisibility sentence exact; the
%  Baire countable-escape chain sound); caveat-discharge with
%  sec:self-reference confirmed two-way (already present); file pure
%  ASCII; no pre-definition "frame"; all citations used as declared.
%  Fixes: Yanofsky2003 paired on the Lawvere clause; the op's two
%  candidate routes anchored to printed propositions (f1-dictionary
%  for the divisibility premise's referent, inside-outside for the
%  closedness demand) -- statuses unchanged, referents sharpened;
%  header staleness corrected. Deliberate FORWARD references --
%  prop:f1-dictionary and sec:hopf-justification (defined later in the
%  paper; fine under two-pass LaTeX and flagged in prose),
%  sec:division-algebras, sec:space-of-computations.
%
%  DECISIONS ENCODED: op:limit-realisability minted (the owed arrow,
%  priced: exactly two derivations, divisibility and closedness);
%  Lagrange used twice (finite exclusion; component-group triviality);
%  the (1+i)/2 lamp value NOT stated here -- "a superposed fixed point"
%  with the complex-probability treatment deferred to
%  sec:self-reference; the two-branch over-determination stated to
%  pre-empt a circularity reading; Thomson spelling throughout.
% =====================================================================

\section{Continuity Is Not an Axiom}
\label{sec:continuity}

Hardy derives quantum theory from five axioms \autocite{Hardy2001a}. Four of them are satisfied
by classical probability theory; the fifth --- that a \emph{continuous} reversible
transformation connects any two pure states --- is the entire difference, the single postulate
that turns the classical simplex into the complex Hilbert space. Every reconstruction programme
pays an equivalent toll: continuous reversibility in Masanes and M\"{u}ller
\autocite{mullerp2011}, the corresponding postulate in Daki\'{c} and Brukner
\autocite{Brukner2009}. The axiom is load-bearing, and in every programme it is unexplained: it
is where the quantumness is put in by hand. The claim of this section is that in the present
framework it is not an axiom. Continuity is the fifth face of the pre-measurement regime that self-reference forces, and what
stands between the framework and a full derivation is not a postulate about values but one
premise about transformations --- that every transformation the inside performs has a stationary
outside presentation --- derived in outline below from the inside/outside involution of
\S\ref{sec:self-reference} and stated as an open problem at the end of the section with nothing
hidden in it. The continuum itself, the completion that every probability theory adopts, is
named for what it is: the real line, shared with classical probability and carrying no quantum
content.

One structural remark prevents a misreading. The framework reaches complex amplitudes natively in
\S\ref{sec:self-reference}, by the negative-probability route; the present section reaches the same
destination by an independent branch --- the regime's continuity, then Hardy's theorem. The two
branches converging on one structure is over-determination, not circularity, and the convergence is
itself evidence that the destination is forced.

\subsection{Five faces of one regime}

The regime before any measurement --- before any \textsc{fanout} --- can be characterised five
equivalent ways: negative probability is available; anti-events exist, every event being undoable;
no classical record exists; no preferred basis is selected; and a continuous reversible
transformation connects any two pure states. These are not five premises to be combined. They are
one regime under five descriptions, and \textsc{fanout} (\S\ref{sec:fanout-boundary}) is the single
event that fails all five at once: it writes the record, selects the basis, kills the anti-events,
discretises the transformation group, and switches off the negativity. The equivalence of the first
two faces is the anti-event identity of \S\ref{sec:self-reference}; the third is the boundary's
definition. The work of this section is the equivalence of the last two faces, and then the route
from the regime's existence to its continuity. The fourth face has, since
Proposition~\ref{prop:f1-dictionary} (\S\ref{sec:hopf-justification} --- a deliberate forward
reference), an exact form: a basis-selection is the commitment of a classical structure, a lifted
antipodal skeleton, so ``no preferred basis'' means precisely that no copy structure has fired.

\subsection{Basis selection is the collapse of continuity}

\begin{proposition}[Continuity forbids a selected basis]\label{prop:basis-discreteness}
If a finite set of pure states is privileged --- a preferred basis --- then every reversible
transformation respecting the privilege carries the set to itself and acts on it as a permutation;
the respecting transformations therefore realise only a finite group of motions among the
distinguished states, and no continuous family of them connects distinct pure states of the set.
Contrapositively, continuous reversibility entails that no basis is selected.
\end{proposition}

\begin{proof}
A privileged finite set is by definition carried to itself by every structure-preserving reversible
map, whose induced action on the set is a permutation of finitely many points. A continuous path of
such maps induces a continuous path in a finite discrete set, which is constant. Hence a continuous
reversible family moving one distinguished pure state to another respects no such privilege, and if
every pair of pure states is so connected, no finite set is privileged.
\end{proof}

The converse is where the historical difficulty lies: the absence of a privileged finite set does
not by itself force continuity, for the transformation group could in principle be infinite and
discrete. Two structural properties close every such escape at once.

\begin{proposition}[Divisible and closed implies connected]\label{prop:divisible-closed}
Let $G$ be the group of reversible transformations of the pre-measurement regime, acting on the
state body. If $G$ is divisible --- every element has an $n$-th root for every $n$ --- and
topologically closed, then $G$ is a compact connected Lie group; and if $G$ moreover acts
transitively on the pure states, then a continuous reversible transformation connects any two of
them: Hardy's fifth axiom.
\end{proposition}

\begin{proof}
The state body is compact and convex with nonempty interior, so its symmetry group is compact, and
a closed subgroup of a compact matrix group is a compact Lie group with finitely many connected
components \autocite{broeckertomdieck1985compact}. The component group $G/G_{0}$ is a quotient of a
divisible group and hence divisible; but a nontrivial finite group is never divisible --- with
$\lvert H\rvert = m$ every element satisfies $h^{m} = e$, so no element other than the identity
possesses an $m$-th root --- and $G/G_{0}$ is finite, so it is trivial and $G = G_{0}$ is connected.
A connected Lie group is path-connected, and transitivity then supplies the continuous reversible
path between any two pure states.
\end{proof}

The same one-line fact of Lagrange closes the remaining escapes. Finite groups are excluded
directly, being nontrivial and never divisible. Countable groups are excluded because a
countable closed subgroup of a compact group is compact, has an isolated point by the Baire
category theorem, and is therefore discrete and finite. Uncountable disconnected groups are
excluded by the component-group argument of the proposition. Every route to a discrete
transformation group therefore passes through a failure of divisibility or a failure of
closedness. Divisibility, however, asks for more than the conclusion needs, and the next
proposition removes it.

\begin{proposition}[Continuity from a connected state space]\label{prop:connected-transitive}
Let the pure states of a finite-dimensional system form a connected space $P$, and let the group
$G$ of reversible transformations of the pre-measurement regime be topologically closed and act
transitively on $P$. Then the identity component $G_0$ already acts transitively on $P$, and a
continuous reversible transformation connects any two pure states: Hardy's fifth axiom holds,
with no divisibility assumed.
\end{proposition}

\begin{proof}
As in Proposition~\ref{prop:divisible-closed}, closedness makes $G$ a compact Lie group with
finitely many connected components \autocite{broeckertomdieck1985compact}, so $G_0$ is a closed
normal subgroup of finite index and is path-connected. Each $G_0$-orbit in $P$ is the continuous
image of a compact group, hence compact, hence closed. Since $G$ is transitive and $G_0$ is
normal, the $G_0$-orbit of $gx$ is $g$ applied to the $G_0$-orbit of $x$, so the cosets of $G_0$
permute the orbits transitively and there are finitely many of them. A finite partition of $P$
into closed sets is a partition into open sets; $P$ is connected; so there is one orbit and
$G_0$ is transitive. Given pure states $x$ and $y$, choose $g \in G_0$ with $gx = y$ and a
continuous path $g_t$ in $G_0$ from the identity to $g$; then $g_t x$ is a continuous reversible
transformation from $x$ to $y$.
\end{proof}

The proposition relocates the premise. Divisibility was asked for because a transformation group
could be infinite and discrete; but a group acting transitively on a continuum cannot be
countable, and a closed group acting transitively on a connected space cannot be disconnected in
the way that matters, so what the argument needs from the regime is not the availability of
$n$-th roots but the connectedness of the state space itself. That is a statement about values:
that the pure states are all the unit vectors of $\mathbb{C}^{K}$ up to phase, the space
$\mathbb{CP}^{K-1}$, connected for $K \ge 2$, and hence that the amplitudes of
\S\ref{sec:self-reference} fill their continuum. The two standing foils show the premise doing
its work. The stabiliser fragment of quantum theory has complex amplitudes and is contextual,
yet it has finitely many pure states, its outcome probabilities take only the values $0$,
$\tfrac12$ and $1$, and its reversible transformations form a finite group, the Clifford group,
so Hardy's fifth axiom fails in it; Spekkens's toy theory \autocite{spekkens2007toy} is the same
with real amplitudes, its ignorance rationed by a knowledge-balance principle. Neither has a
connected state space, and neither is a model of the regime described here, whose credences take
every value. Divisibility, which both also lack --- the phase gate has no square root in the Clifford group
--- is a symptom of the discreteness, not its cause. What both lack can be said exactly, and the
next subsection says it: each contains a transformation that fixes no state it admits.

\subsection{The premises, and what they mean}

Three premises carry Proposition~\ref{prop:connected-transitive}. Transitivity is the regime's
homogeneity: inside/outside equivalence leaves no pure state privileged, and a proper closed
invariant set of pure states would define a superselection observable, a classical record
predating any \textsc{fanout} (\S\ref{sec:fanout-boundary}). Connectedness of the state space is two statements, and only one of them is quantum. The first
is that the pure states are infinite and dense, with a fixed point for every transformation; the
second is that the state space is complete. Quantum theory over the algebraic numbers ---
amplitudes in the algebraic closure of $\mathbb{Q}$ within $\mathbb{C}$ --- has the first
without the second: it has $K = N^{2}$, the Born rule, superposition, entanglement and Bell
violation at the Tsirelson value, every unitary has an algebraic eigenvector, and it is
countable and totally disconnected. So the whole quantum content of Hardy's axiom lies in the
density, and the completion is the real-number convention of every probability theory. Hardy
adopts it before his fifth axiom is stated, in the first, which makes probabilities real numbers
and gives his classical theory a real simplex of mixtures; this framework adopts it at the
Shannon step, where credences are declared to range over a connected domain so that the
continuation of the logarithm is unique (\S\ref{sec:complex-amplitudes}). Given that convention,
density of the pure states is connectedness of their closure. Divisibility, which an earlier
version of this section asked for, is not needed: it follows from the conclusion, every element
of a compact connected Lie group having roots of every order.

Closedness is the completion, limit-realisability: a Cauchy sequence of available
transformations has its limit available. It cannot be derived, and the reason should be stated
rather than hidden. Any collection of first-order principles about states, transformations and
probabilities that has an infinite model has a countable one, and a countable subset of
$\mathbb{CP}^{K-1}$ is totally disconnected; the founding non-computability does not escape
this, since the arithmetical complex numbers form a countable algebraically closed field
containing non-computable numbers and closed under every limit it can itself describe. Hardy's
axiom and closedness are second-order, quantifying over paths and over sequences, and a
second-order principle is grounded as Dedekind completeness is grounded or not at all. The
grounds are these. The completion of a dense set under the finite-precision metric is unique and
depends on no choice, whereas any countable alternative is a choice of definability level, an
extra law with no ground and no test; and the framework's own result is that no observer's
descriptive level captures the regime, since by \S\ref{sec:self-reference} every observer's
state exceeds its own level and the inside/outside principle makes that a law rather than a
perspective. The burden therefore falls on any restriction, not on the completion, which is the
least committal option and carries no quantum content. What limit-realisability adds to the line
is still worth stating exactly, because it is the corrected form of a natural intuition.
Cantor's diagonal \autocite{cantor1891diagonal} does not build the continuum; its output is an
infinite decimal --- a Cauchy sequence of rationals --- which exists \emph{as a point} only
because the line is complete. The diagonal presupposes limit-realisability and then shows that
the completed line outruns every internal listing. And the character of the completion is the
framework's own: the computable reals are countable \autocite{Turing_1937}, hence of measure
zero, so once the transformation group contains its limits almost all of them are
non-computable. Non-computability is a property of the completed continuum, not the reason for
it. The constructive reading, on which only describable limits exist, remains consistent; it is
the countable theory above, and the two readings differ in nothing an experiment can certify.

Density is where the section's title is earned, and it is derived from the paper's own premise.
Write $(\mathrm{I})$ for the statement that every admissible transformation of the regime is a
phase of some admissible classical structure: diagonal, that is, in the record basis of some
observer the regime contains, a phased spider of that basis in the calculus of
\S\ref{sec:hopf-justification} \autocite{coeckeduncan2011interacting} (VERIFY-SUPPORT). It says
of transformations what the global phase of \S\ref{sec:self-reference} says of states: there are
no absolute ones. It follows from three things the paper already holds. The inside/outside
involution is a genuine involution, so the outside presentation of the inside is a presentation
of the same object with nothing lost: the outside can have complete knowledge of the inside, a
potential record, the joint state lying in a structure the outside carries. Time is operation
count, internal to the system that operates (\S\ref{sec:space-of-computations}), so the outside
does not share the inside's count. And before any \textsc{fanout} that count is not classically
recorded, so the outside's description is pure. Let $g$ be the transformation the inside
performs per operation. The outside's complete knowledge cannot depend on the inside's count,
and it is pure, so it is a pure state of the inside's clock and system that is stationary under
``advance the clock and apply $g$''; the stationary states of that map are clock Fourier modes
tensored with eigenvectors of $g$ and nothing else \autocite{pagewootters1983evolution}
(VERIFY-SUPPORT). So the outside holds a structure in which $g$ acts by phases, which is
$(\mathrm{I})$; and the relative phase in that structure is the outside's ignorance of the
inside's computational count, the same ignorance the global phase records one level down. The
consequences are a theorem. A phase of a structure fixes the structure's copyable states, so
every admissible transformation fixes an admissible pure state; a finite transitive action has
an element moving every point (Jordan), so the pure states are infinite; the transformation
carrying a stabiliser state to the fixed point of the composite of a $Z$-phase and an $X$-phase
turns through $\arccos(1/\sqrt{3})$, an irrational fraction of a turn, so the group has an
element of infinite projective order and its orbits are dense. With the completion,
Proposition~\ref{prop:connected-transitive} fires and Hardy's fifth axiom holds. The two foils
fail $(\mathrm{I})$ exactly where they fail continuity: fourteen of the twenty-four Clifford
operations are phases of no Pauli basis and fix no stabiliser state, their fixed points being
the magic states; the toy theory's three-cycles are phases of no epistemic basis; and in each
the outside's knowledge of the inside under the order-three protocol is either count-dependent
or a mixture over the count, so each makes the inside's time classical before any record exists.
The derivation of $(\mathrm{I})$ is an outline, not a proof, and it is what the open problem
below asks for.

\subsection{Why the cut matters physically}

The regime's continuity is not decorative. Reversibility alone, over complex amplitudes, licenses
the full invertible group $GL(n, \mathbb{C})$, and Aaronson showed what that buys: superluminal
signalling, and computational power at the level of \textsf{PP} \autocite{Aaronson2004h}. The cut
from $GL(n, \mathbb{C})$ down to the unitary group is the conserved quadratic form, and continuity
is what entitles the theory to make it: with the fifth axiom in hand, Hardy's theorem delivers the
complex Hilbert space, the quadratic norm and the unitary group in one stroke \autocite{Hardy2001a}
--- cited here, not re-proved. One honest caveat travels with the citation. Hardy's reconstruction
uses his remaining axioms besides --- composition and subspace structure --- and these must be
derived rather than imported: the composition ledger of \S\ref{sec:division-algebras} and the
locality-of-tomography argument of \S\ref{sec:self-reference} carry that duty, and it is discharged
there, not silently assumed here.

\subsection{Super-tasks at the boundary}

The premise just priced has a fifty-year-old stress test in the philosophy of mathematics, and the
framework passes it in an unexpected way. A super-task --- infinitely many operations completed in
finite time \autocite{Thomson_1954} --- presupposes its terminal moment as an actual point of the
timeline: the accumulation point of a generic sequence of step-times exists only in the completed
line, and the defining question, the state \emph{after all the steps}, is meaningful only under
limit-realisability of time. Super-tasks are downstream consumers of exactly the completeness at
issue in this section. Thomson's lamp then exhibits both sides of the \textsc{fanout} boundary at
once. On the record side, the lamp's description fixes its state at every finite stage and thereby
fixes no limit at all: any terminal value is consistent with the whole task
\autocite{Thomson_1954,benacerraf1962tasks}. That is not a puzzle about lamps; it is the proof that
the classical record algebra is not limit-closed, and cannot be. On the amplitude side the limit
exists --- a superposed fixed point, whose complex-probability treatment belongs to
\S\ref{sec:self-reference} --- because that side \emph{is} limit-closed, which is the fifth face
itself. The boundary is thereby the edge of limit-realisability: complete on the reversible side,
incomplete on the record side, with the lamp standing on both sides at once as the boundary's oldest
exhibit. A final clause dissolves rather than answers the classic question. In this framework time
is operation count (\S\ref{sec:space-of-computations}), and ``after infinitely many operations''
indexes no operational moment: the lamp packs its endless \textsc{fanout}s into the emergent
continuum, and asking for its record at the accumulation point asks the discrete contents for a
value that only the completed container possesses. On or off is not undecided; it is unasked.

\subsection{What remains, priced}

\begin{openproblem}[Every transformation a phase of an inside structure]\label{op:limit-realisability}
Proposition~\ref{prop:connected-transitive} reduces Hardy's fifth axiom to three premises:
connectedness of the pure states, transitivity, and closedness. Transitivity is argued above
from the absence of records. Closedness is the completion, adopted on the grounds stated above
and carrying no quantum content. Connectedness is density plus that completion, and density
follows, by Jordan's theorem and the infinite-order argument, from $(\mathrm{I})$: that every
admissible transformation of the regime is a phase of some admissible classical structure. What
remains open is the rigour of the derivation of $(\mathrm{I})$ from the paper's premise. Derive,
inside the fixed-point locus $\mathrm{Fix}(\sigma)$ of Proposition~\ref{prop:inside-outside},
that every endomap is a phase of a classical structure in $\mathrm{Fix}(\sigma)$, from the
self-representing object, the involution $\sigma$ as a genuine involution, and the reading of
the inside presentation as a sequence indexed by operation count with no record before
\textsc{fanout}. The outline above does this with the stationary-state construction of Page and
Wootters; a derivation passes if it makes that outline a theorem of the calculus, with the
outside's count-independent complete description shown to exist and to be pure from the
involution alone. It fails before any detail is checked if every premise it uses holds in the
stabiliser fragment of quantum theory, in Spekkens's toy theory \autocite{spekkens2007toy}, in
the real qubit, or in quantum theory over the Gaussian rationals, all of which violate
$(\mathrm{I})$ outright and all of which make the inside's operation count classical before any
record exists. It also fails if it imports the completion as a premise about transformations:
the completion is the real line and has no derivation, and a derivation of $(\mathrm{I})$ must
hold in the countable algebraic theory, where $(\mathrm{I})$ is true and continuity is not.
\end{openproblem}

The accounting, in one sentence: the axiom that every reconstruction programme pays is, in this
framework, one derivation posed, of Lawvere's shape --- that every endomap of the fixed-point
locus is a phase of a classical structure in it --- and one convention named, the completion,
which classical probability pays too. And
the hand-on is already visible: the complex amplitudes that this section's continuity underwrites
through Hardy's theorem, and the complex amplitudes that \S\ref{sec:self-reference} derives natively,
are one and the same $\mathbb{C}$ --- the floor of the ladder that \S\ref{sec:division-algebras}
climbs.

%      % EDITED 2026-10-29 scripts/patch_continuity_tower_premise_2026-10-29T1500.py
% EDITED 2026-09-15 scripts/patch_value_continuum_declared_cencov_2026-09-15T1540.py
% EDITED 2026-09-15 scripts/patch_continuity_connected_state_space_2026-09-15T1430.py
% =====================================================================
%  SECTION: The Euclidean metric on the state sphere, and the global
%  phase as self-referential ignorance
%  Origin: drafted in an earlier conversation (2026-06-07), edited by
%  JS; audited into the current architecture 2026-07-06.
%
%  PLACEMENT in main.tex (reordered chain, JS 2026-07-06): position 3 --
%      \input{self_reference_section_2026-07-04}
%      \input{continuity_section_2026-07-04}
%      \input{L2_metric_global_phase_section_2026_06_07}   <-- here
%      \input{wigners_friend_in_the_hopf_picture_2026-06-13}
%  (WF backward-references sec:l2-metric-fibre: order-consistent.)
%
%  AUDITED 2026-07-06 (final pass): Lamperti/Banach classification,
%  Hopf map components, submersion facts, Hardy K=N^r reading all
%  verified; the recoverable-vs-gauge-private fibre distinction
%  confirmed canonical and untouched; file pure ASCII. Fixes: opening
%  rewritten (Popper C+/C/E forward-voiced to sec:wigners-friend;
%  continuity claim made CONDITIONAL per JS -- "assuming those premises
%  can be proved", op:limit-realisability cited); Aaronson paragraph
%  de-duplicated against sec:continuity by stitch; sec:born-rule was
%  dangling paper-wide -- label minted at its canonical home (the Born
%  subsection of sec:self-reference) and the claim here recalibrated
%  ("forced, not chosen" -> lowest-degree-invariant, geometric face);
%  record-absence upgraded from interpretive to theorem via
%  prop:f1-dictionary while the self-ignorance identification stays
%  interpretive; grammar and dash punctuation repaired; field-not-fixed
%  bound pointed at sec:self-reference.
%
%  CITATIONS: existing keys Hardy2001a, Aaronson2004h, mosseri2001hopf.
%  NEW keys (RIS in L2_metric_refs_2026-07-06.ris): Banach1932,
%  Lamperti1958, Hopf1931, Urbantke2003, GluckWarnerZiller1986,
%  wooters_1981 (key spelling preserved as used; note the single-t).
%  British English; \autocite; no bullets in prose.
% =====================================================================

\section{The Euclidean metric on the state sphere, and the global phase as self-referential ignorance}
\label{sec:l2-metric-fibre}
Popper's systems $C^{+}$, $C$ and $E$ --- the predictor, its target and their shared
environment, whose argument \S\ref{sec:wigners-friend} anchors and quotes --- are, in modern
dress, Wigner, his friend and a qubit $S$ to be measured. The natural arena for analysing them
is the Hopf fibration and the Bloch sphere; but even to draw the Bloch sphere one must first
justify the $L^{2}$ metric, and this section earns that right. The principle doing the work is
the reversible continuity of \S\ref{sec:continuity} --- derived there from self-reference
conditionally, on the premises priced as Open Problem~\ref{op:limit-realisability}, of which one, that every
transformation of the regime is a phase of an inside structure, remains open; the continuum
itself is the completion the Shannon step declares (\S\ref{sec:complex-amplitudes}), the
classical convention, and \S\ref{sec:cencov-route} below reaches the same metric by a
second route that consumes it and nothing else that is open. Assuming that premise can be proved,
everything below inherits the derivation; we carry that status, and do not upgrade it.

We have established that the amplitudes of a single qubit form a complex vector space
$\mathbb{C}^{2}\cong\mathbb{R}^{4}$. We have not yet fixed which norm on this space singles out
the physical state manifold. That choice is not free: the principle of
\emph{reversible continuity} forces the Euclidean ($L^{2}$) norm, and hence the round
three-sphere, before any appeal to the Born rule. The quadratic structure usually \emph{posited}
as the Born rule is therefore already present in the metric, and the global phase emerges as a
geometrically forced fibre rather than an added convention.

\subsection{Reversible continuity selects the 2-norm}
\label{sec:reversible-continuity-l2}

The principle is the one isolated by Hardy as the single axiom distinguishing quantum from
classical probability \autocite{Hardy2001a}:
\emph{there exists a continuous reversible transformation between any two pure states.}
Operationally a reversible transformation of states is a symmetry of the state manifold, hence
norm-preserving; the demand is that such symmetries form a \emph{connected} group acting
\emph{transitively} on the pure states. Linearity is inherited from the complex vector-space
structure already in hand, so the admissible transformations are the complex-linear isometries
of whatever norm we place on $\mathbb{C}^{2}$.

Take the candidate norms to be the $\ell^{p}$ family,
$\lVert x\rVert_{p}=\bigl(\sum_{k}\lvert x_{k}\rvert^{p}\bigr)^{1/p}$, with the state manifold
the unit $\ell^{p}$ sphere. Their linear isometries are completely classified.


\begin{proposition}[Lamperti; after Banach]
\label{prop:lamperti}
For $1\le p<\infty$, $p\neq 2$, every linear isometry of $L^{p}$ (and hence of $\ell^{p}$) is a
generalised permutation: a weighted composition operator which, in the finite-dimensional case,
acts as $(x_{k})\mapsto(\omega_{k}\,x_{\sigma(k)})$ with $\lvert\omega_{k}\rvert=1$ and $\sigma$
a permutation of the coordinates.
\end{proposition}

\noindent Banach gave this for $L^{p}[0,1]$; the general $\sigma$-finite statement is Lamperti's
\autocite{Banach1932, Lamperti1958}. Its consequence for our four real dimensions is decisive.
For every $p\neq 2$ the isometry group is, in the real case, the finite group of signed
permutations (the hyperoctahedral group $B_{4}$), and in the complex case the group
$\mathbb{T}^{n}\rtimes S_{n}$ of diagonal phases composed with permutations. Even the complex
group, though it contains the continuous torus $\mathbb{T}^{n}$, fails to act transitively on
the sphere: the orbit of a basis state $e_{1}$ is only the phased basis vectors
$\{\omega\,e_{k}\}$, a measure-zero union of circles. No continuous reversible transformation
carries a basis state into a generic superposition such as $(e_{1}+e_{2})/\sqrt{2}$. Reversible
continuity therefore \emph{fails} for every $p\neq 2$.

At $p=2$ the isometry group is $O(4)$ in the real picture and $U(2)$ in the complex one; the
identity component acts transitively on the unit sphere, and any two unit vectors are joined by
a continuous one-parameter subgroup. Among the $\ell^{p}$ geometries, reversible continuity is
satisfiable at $p=2$ alone.

Aaronson states the same selection in the quantum-foundations setting and sharpens it into a
physical, rather than merely aesthetic, necessity \autocite{Aaronson2004h}: the $2$-norm is the
unique $p$-norm admitting a nontrivial continuous group of norm-preserving linear maps, every
other $p$ yielding only generalised permutations; and the alternative of keeping linearity and
reversibility while \emph{dropping} norm preservation lands one in the full general linear group
$\mathrm{GL}(n,\mathbb{C})$, which permits superluminal signalling and collapses the complexity
class $\mathsf{PP}$ into polynomial time --- the same fact \S\ref{sec:continuity} uses to price
the classical--quantum cut. Bare reversibility on the amplitude space buys only
$\mathrm{GL}(n,\mathbb{C})$; what carves the unitary group $U(n)$ out of it is preservation of
the quadratic ($\ell^{2}$) form, and reversibility of the resulting maps is then a corollary
rather than a premise.

The conclusion is that the physical state manifold of a single qubit is the unit sphere of the
Euclidean norm,
\begin{equation}
  S^{3}=\Bigl\{\,x\in\mathbb{R}^{4}\;:\;\textstyle\sum_{i=1}^{4}x_{i}^{2}=1\,\Bigr\}
  =\bigl\{\,\psi\in\mathbb{C}^{2}\;:\;\lvert z_{1}\rvert^{2}+\lvert z_{2}\rvert^{2}=1\,\bigr\},
\end{equation}
with $z_{j}=x_{2j-1}+\mathrm{i}\,x_{2j}$, carrying the round metric induced from the ambient
inner product,
\begin{equation}
  \mathrm{d}s^{2}=\sum_{i=1}^{4}\mathrm{d}x_{i}^{2}\,\Big|_{S^{3}}
  =\operatorname{Re}\langle\mathrm{d}\psi,\mathrm{d}\psi\rangle,
  \qquad
  \langle\phi,\psi\rangle=\bar\phi_{1}\psi_{1}+\bar\phi_{2}\psi_{2}.
\end{equation}
One bound deserves explicit statement, in the interest of claiming no more than is proved:
reversible continuity fixes the \emph{metric} as $L^{2}$, but not the underlying field. Real and
quaternionic quantum theory equally satisfy reversible continuity and equally carry Euclidean
state-space metrics; the restriction to $\mathbb{C}$ is pinned elsewhere in the framework ---
the analytic continuation of the information measure and the additivity of phase information,
\S\ref{sec:self-reference} --- not here.

\subsection{Hardy's independent route}
\label{sec:hardy-route}

The same conclusion is reached by Hardy along a route that never invokes
Proposition~\ref{prop:lamperti} \autocite{Hardy2001a}. His reconstruction turns on the counting
relation $K=N^{r}$, where $N$ is the number of states distinguishable in a single measurement
and $K$ the number of fiducial probabilities needed to specify a state. The continuity axiom
excludes the classical case $r=1$: there the pure states are the $N$ vertices of a simplex and
the reversible maps merely permute them --- a discrete group of motions --- so there is no
continuous path between distinct pure states, which is the simplex face of the same $\ell^{1}$
isometry fact. A simplicity assumption then selects $r=2$, yielding complex quantum theory with
its quadratic norm and trace rule. The quadratic structure is not extracted from continuity as a
separate step; it is built into Hardy's first axiom (probabilities linear in the fiducial set)
together with $K=N^{2}$.

We thus have two genuinely independent demonstrations that reversible continuity forces the
$L^{2}$ metric: Hardy's operational $K=N^{2}$ counting, and the Lamperti--Aaronson
classification of the isometries. They converge on the same step by different mathematics, and
should be cited as two roads rather than one road traversed twice. Both roads start from
continuity. A third, in the next subsection, starts one step earlier, from the value continuum
that continuity itself rests on, and reaches the same metric without passing through the
transformation group at all.

\subsection{The statistical route: \v{C}encov's theorem and the every-basis lemma}
\label{sec:cencov-route}

The value continuum that \S\ref{sec:continuity} isolates as the one open premise reaches the
$L^{2}$ metric on its own, by a route that uses no transformation group. \v{C}encov's theorem
states that the only Riemannian metric on the simplex of probability distributions invariant
under coarse-graining --- under the Markov maps that merge or relabel outcomes --- is the
Fisher--Rao metric \autocite{bengtsson2006geometry}, and in the coordinates $x_{k} =
\sqrt{p_{k}}$ the Fisher--Rao metric is the round metric on the positive orthant of the unit
sphere. The premise the theorem consumes is that the simplex is a manifold, which is the value
continuum declared in \S\ref{sec:complex-amplitudes}, and the principle it uses, that
distinguishability does not depend on how outcomes are grouped, is one the Shannon bookkeeper of
that section already grants. What it does not supply is the passage from one orthant to the
whole projective space; that is the following lemma.

\begin{lemma}[The every-basis lemma]\label{lem:every-basis}
Let $g$ be a Riemannian metric on $\mathbb{CP}^{K-1}$ whose restriction to the positive real
orthant of every orthonormal basis --- the states $\sum_{k}\sqrt{p_{k}}\,e_{k}$ with $p$ in the
simplex --- is the Fisher--Rao metric. Then $g$ is the Fubini--Study metric, whose linear lift
is the $L^{2}$ inner product.
\end{lemma}

\begin{proof}
A tangent vector to $\mathbb{CP}^{K-1}$ at a state $\psi$ is represented by a vector $\phi$
orthogonal to $\psi$, and the curve $t \mapsto \cos t\,\psi + \sin t\,\phi$ with real $t$ lies
in the positive orthant of any orthonormal basis that contains $\psi$ and $\phi$. Along that
curve the Fisher--Rao metric is the round metric of the sphere in the coordinates $\sqrt{p}$,
which is the Fubini--Study metric restricted to real amplitudes; so $g$ and the Fubini--Study
metric agree on every tangent vector at every point, the direction $i\phi$ being reached by the
basis that contains $i\phi$ in place of $\phi$.
\end{proof}

The lemma's premise, that every orthonormal basis is a legitimate measurement and none is
privileged, is the fourth face of the pre-measurement regime seen from the measurement side
(\S\ref{sec:continuity}); in the stabiliser fragment of quantum theory it fails in the same way
everything else does, only the Pauli bases being measurements there. So the $L^{2}$ metric now
rests on two legs that share exactly one premise. Through reversible continuity it uses the
value continuum, closure and homogeneity (Proposition~\ref{prop:lamperti}); through \v{C}encov
it uses the value continuum, coarse-graining invariance and the equal standing of bases
(Lemma~\ref{lem:every-basis}). Neither leg needs anything else that is open, and nothing could
avoid the continuum, since a metric is a statement about one. That is the status of everything
downstream of this section, the Hopf fibrations included: derived from the completion every physicist grants, which carries no quantum content, and from
one premise about transformations that the framework derives in outline and has not yet proved
(Open Problem~\ref{op:limit-realisability}).

\subsection{The Hopf fibration and the Bloch sphere}
\label{sec:hopf-bloch}

The round metric is exactly what makes the global phase a clean symmetry. The phase action
$\psi\mapsto\mathrm{e}^{\mathrm{i}\theta}\psi$ is a one-parameter subgroup of
$U(2)\subset SO(4)$ and so acts on $S^{3}$ by \emph{isometries}, a statement that holds only
because the metric is the $L^{2}$ one just derived. Its orbits are great circles $S^{1}$, all of
circumference $2\pi$, and the quotient is the complex projective line,
\begin{equation}
  S^{3}\big/U(1)=\mathbb{CP}^{1}\cong S^{2},
\end{equation}
the Bloch sphere. This quotient is the complex Hopf fibration
\autocite{Hopf1931, Urbantke2003, mosseri2001hopf}
\begin{equation}
  S^{1}\hookrightarrow S^{3}\xrightarrow{\;h\;}S^{2},
\end{equation}
whose projection sends an amplitude to its Bloch vector,
\begin{equation}
  h(\psi)=\bigl(\langle\psi\rvert\sigma_{x}\lvert\psi\rangle,\;
                 \langle\psi\rvert\sigma_{y}\lvert\psi\rangle,\;
                 \langle\psi\rvert\sigma_{z}\lvert\psi\rangle\bigr)
        =\bigl(2\operatorname{Re}\bar z_{1}z_{2},\;
               2\operatorname{Im}\bar z_{1}z_{2},\;
               \lvert z_{1}\rvert^{2}-\lvert z_{2}\rvert^{2}\bigr).
\end{equation}
The Hopf map is quadratic in the amplitudes, and this is the geometric face of the Born rule:
the third component is $p_{0}-p_{1}$ with $p_{k}=\lvert z_{k}\rvert^{2}$, and the squaring is
the lowest-degree $U(1)$-invariant available on the fibre --- the record-theoretic derivation of
the same rule is \S\ref{sec:born-rule}'s. Because the metric is round, the projection is a
Riemannian submersion with totally geodesic fibres of constant length
\autocite{GluckWarnerZiller1986}: the same amount of phase is hidden at every point of the base,
so the phase is hidden \emph{uniformly}, not in a position-dependent way that the metric could
partially betray.

\subsection{The fibre as self-referential ignorance}
\label{sec:fibre-self-ignorance}

The interpretation the framework places on this geometry --- and here we are reading physics
into the mathematics, so the claim is interpretive rather than proved --- is that the $S^{1}$
fibre is the observer's own un-narratable self-ignorance. The global phase is the imaginary part
of the complex information measure $S=-\ln\lvert p\rvert-\mathrm{i}\theta$: information the
observer lacks about their own state. Lacking it, they cannot communicate it; there is no
\textsc{fanout} of information one does not possess. Uncommunicable, it cannot enter any
measurement record --- and this last step is no longer interpretive:
Proposition~\ref{prop:f1-dictionary} proves the global phase absent from every classical record
a copying structure induces --- so every observable is $U(1)$-invariant; and a $U(1)$-invariant
quantity is a function on the quotient. The physical state is therefore the orbit, not the
representative point---the base $S^{2}$---and the $S^{1}$ fibre carries exactly the inaccessible
self-referential phase. It is a null direction of statistical distinguishability: two global
phases are indistinguishable by any measurement whatever \autocite{wooters_1981}, so the fibre
is a gauge direction without physical content. This phase is irrecoverably discarded, and must
remain so: it is the relative phase, living in the $S^{3}$ fibre of the quaternionic Hopf
fibration at two qubits, that is genuinely recoverable, whereas the single-qubit global phase of
the complex fibration is gauge-private without remainder. The self-reference that opened the
derivation thus reappears, at its end, as the circle the observer can never see around.
   <-- here
%      % EDITED 2026-09-12 scripts/patch_revert_sedenion_repoint_2026-09-12T1600.py
% EDITED 2026-09-11 scripts/patch_undefined_section_refs_2026-09-11T2030.py
% =====================================================================
%  SECTION: Wigner's friend in the Hopf picture (v2; supersedes the
%  2026-06-07 draft). AUDITED 2026-07-06 (final pass, fifth stop).
%
%  PLACEMENT (reordered chain, JS 2026-07-06): position 4 --
%      \input{L2_metric_global_phase_section_2026_06_07}
%      \input{wigners_friend_in_the_hopf_picture_2026-06-13}
%      \input{derivation_of_the_quantum_postulates}
%  sec:l2-metric-fibre is REAL (backward, order-consistent);
%  sec:fanout-boundary DRAFTED 2026-07-06 (formerly the one undrafted section);
%  the old sec:monogamy placeholder note was VOID -- no such \ref in
%  the body -- and is removed. prop:f1-dictionary and
%  sec:hopf-justification are FORWARD references (position 8): first
%  F1 use now carries the deliberate-forward flag.
%
%  DEPENDENCY: \input{monogamy_aside_2026-06-07} -- received from JS
%  and audited 2026-07-06; the \input sits inside
%  sec:wf-reconciliation, immediately after op:tsirelson-geometry
%  (relocated 2026-07-06 from end-of-file, where the aside's opening
%  antecedent had been orphaned). A bare replace-all during the earlier
%  patch briefly mangled this header block, leaving a live \input in
%  it; repaired by index-splice, and the lesson recorded: anchors that
%  also occur in comments must be uniqueness-asserted.
%
%  AUDIT RECORD: CHSH arithmetic (2sqrt(1+C^2), PR => r^2 = -2) and the
%  C^2+r^2=1 chain verified; F1 wording canonical; the
%  obsolescence-made-cancellation echo two-way with
%  sec:self-reference; redundancy/Darwinism/DiBiagio-Rovelli/Brukner
%  positioning sound; Landauer placed on amplification, correctly.
%  Fixes: first-person voice removed; \textquote replaced with
%  portable quotes; IOEP given its prop ref; stale Tsirelson
%  mini-header (wrong key name Tsirelson1980) retired;
%  op:tsirelson-geometry MINTED from the honesty paragraph's unpriced
%  owed. Popper page-pin DISCHARGED 2026-07-07: the zip bundles carry
%  OCR .txt sidecars (the earlier "no text layer" verdict was wrong --
%  ZIP-of-JPEGs defeated the PDF tools, not the text); the quoted
%  phrase sits at Part II, p. 187, now pinned in the cite. Bonus from
%  the same dig: "C+" is Popper's own notation (Part II, scan pp.
%  15-17), so the L2 opening's attribution stands as written; "E" not
%  sighted in the sampled pages -- the inline gloss covers it either
%  way. Keys unique to this file (presumed in the .bib,
%  since the paper compiles): VerstraeteWolf2002,
%  Cirelson1980QuantumInequality, PopescuRohrlich1994, Held2021,
%  Brukner2021, zurek_2009, Landauer_1961, wigner_1961.
%  British English; \autocite; no bullets in prose.
% =====================================================================

\section{Wigner's friend in the Hopf picture: one state, two charts, and the redundancy ladder}
\label{sec:wigners-friend}

The framework dissolves the Wigner's-friend paradox \autocite{wigner_1961} not by
adjudicating between the two accounts but by recognising that they are written in
two charts on a single manifold, selected as the problem requires, much as one
moves between cylindrical and spherical coordinates without either being the true
one. The epistemological chart (private, non-clonable, phase-carrying) and the
ontological chart (clonable, public, record-bearing) are related by the
\textsc{fanout} boundary (Section~\ref{sec:fanout-boundary}), and that boundary
is not a sharp wall but a graded one, parametrised by how redundantly a piece of
relational content has been copied. Read this way, the fact that a \textsc{fanout}
has occurred for the friend but not for Wigner is not a contradiction to be
removed; it is the boundary behaving exactly as defined.

Let the system be a qubit $\lvert\psi\rangle_{S}=\alpha\lvert 0\rangle
+\beta\lvert 1\rangle$ and let the friend $F$, idealised, with deliberate
economy, as a single qubit, measure it in the computational basis. The
interaction is the controlled coupling that copies $S$'s value into $F$'s
register: the \textsc{copy} comultiplication of the diagrammatic calculus
\autocite{CoeckeKissinger2017}, and in the circuit picture the very gate that, in
the two-qubit Mosseri--Dandoloff geometry \autocite{mosseri2001hopf},
activates the quaternionic Hopf fibration $S^{3}\to S^{7}\to S^{4}$ and carries
the joint state from the surface of the entanglement ball to its interior. The
friend and Wigner read the two ends of that plunge.

Two idealisations in this set-up should be named at once, because the argument
later depends on seeing past them. The first is the friend as a single qubit,
which we return to as the redundancy-one endpoint of a graded boundary
(Section~\ref{sec:wf-redundancy}). The second, less obvious and more consequential,
is the system $S$ as a pristine qubit with no prior entanglement --- a relation
with no history. No real measured particle is of this kind. A particle is itself a
reified relation between prior events, and by the time the friend couples to it, it
already carries the completed internal structure such a relation possesses: in the
framework's terms, the full octonionic $A,B,C$ register that the $S^{7}$ sector
supplies as internal degrees of freedom (\S\ref{sec:space-of-computations}). The
pristine qubit is a schematic convenience, drawn structureless so that the logic of
a single link is visible; the ``system, friend, Wigner'' triple is likewise the
schematic of \emph{one link} in the chain of observers, not its base. The chain
does not bottom out in a structureless qubit --- there is no such thing --- and we
draw on this when we ask, later, what bounds the information a single such link can
carry.

\subsection{The friend's chart: a \textsc{fanout} of redundancy one}
\label{sec:wf-friend}

Relative to $F$, the coupling is a \textsc{fanout}: $F$ acquires a fact. What is
copied is gauge-invariant, the relative phase, the correlation between $F$ and
$S$, which is a genuine observable. The global phase is never the thing copied:
there is nothing observable there to duplicate, so the $S^{1}$ fibre of the
complex Hopf fibration $S^{1}\to S^{3}\to S^{2}$ of
Section~\ref{sec:l2-metric-fibre} remains $F$'s gauge-private self-ignorance,
untouched. In Hopf terms the division is clean, but it lives at two levels, which
Proposition~\ref{prop:f1-dictionary} (\S\ref{sec:hopf-justification} --- a deliberate forward
reference) separates. The copy operation itself is fibre-committed:
copying is faithful to phase, so what $F$'s \textsc{fanout} duplicates is a lift --- a chosen
fibre-point over the outcome --- never a bare point of the base. Yet everything the operation can
broadcast is base data alone: the record it induces is independent of the lift, so the $S^{1}$
freedom, though pinned in the copier, is absent from every classical record it yields. $F$'s classical record is a
base-point; the structure that produced it committed to a fibre-point; and the gap between the two
is precisely the gauge character of the fibre. The observation has a fifty-year-old
anchor: Popper argued that a predictor's inferior self-knowledge follows from the logical
impossibility of a system's imparting information about itself without thereby rendering it out of
date --- ``making this information obsolete'', in his words
\autocites{Popper1950a}[Part~II, p.~187]{Popper1950b}. Proposition~\ref{prop:f1-dictionary}
makes the impossibility structural rather than temporal, and localises it: the datum the copier
cannot put into its own record is not its state wholesale but exactly its fibre commitment, which
does not arrive stale --- it cancels. Conditioned on the outcome it reads, $S$ is for
$F$ a definite point on the Bloch sphere,
\begin{equation}
  \rho_{S}^{(F)}=\lvert 0\rangle\langle 0\rvert
  \quad\text{(given outcome }0\text{)},
\end{equation}
a pure state of Bloch radius $r=1$ on the surface of the ball, concurrence $C=0$.

This is a \textsc{fanout}, but the minimal one. The fact $F$ holds is a
\emph{relative} fact in the sense of relational quantum mechanics
\autocite{Rovelli1996a}: $S$ definite relative to $F$, not an observer-independent
record. It is one copy, redundancy $R=1$, and as such it remains reversible by an
agent who controls the joint system. The single-qubit friend is therefore not
outside \textsc{fanout}, as a careless reading might suggest; it is the $R=1$
endpoint of the \textsc{fanout} continuum---a \textsc{fanout} that has happened,
and a fact that is real for $F$, while being maximally still reversible from
outside.

\subsection{Wigner's chart: no copy has reached him}
\label{sec:wf-wigner}

Wigner treats the laboratory as closed and no copy of the relational content has
reached him. For him the coupling is the unitary
\begin{equation}
  \lvert\Psi\rangle
  =\alpha\lvert 0\rangle\lvert F_{0}\rangle
  +\beta\lvert 1\rangle\lvert F_{1}\rangle,
\end{equation}
the move up to the larger Hilbert space, $\mathbb{C}^{2}\to\mathbb{C}^{4}$,
$S^{3}\to S^{7}$, complex Hopf to quaternionic Hopf, with the fibre grown from
$S^{1}$ to $S^{3}$. What $F$ underwent as the acquisition of a fact, Wigner
describes as correlation deposited into a joint pure state on $S^{7}$. When he
asks specifically about $S$, he traces $F$ out, and the geometry does the
bookkeeping:
\begin{equation}
  \rho_{S}^{(W)}=\operatorname{Tr}_{F}\lvert\Psi\rangle\langle\Psi\rvert
  =\lvert\alpha\rvert^{2}\lvert 0\rangle\langle 0\rvert
  +\lvert\beta\rvert^{2}\lvert 1\rangle\langle 1\rvert,
  \qquad
  r=\bigl\lvert\,\lvert\alpha\rvert^{2}-\lvert\beta\rvert^{2}\bigr\rvert,
  \quad
  C=2\lvert\alpha\beta\rvert .
\end{equation}
Tracing out discards the $S^{3}$ fibre and carries the pure point on $S^{7}$ to a
mixed point in the interior of the Bloch ball at radius $r=\sqrt{1-C^{2}}$
\autocite{mosseri2001hopf}. The inward displacement is governed by the
concurrence: it is the relational content that has migrated into the $S$--$F$
correlation, stored in the recoverable $S^{3}$ fibre. Nothing has been destroyed;
it has been relocated into the fibre, which is why Wigner can in principle
interfere the branches and reverse the interaction, the content sits at low
redundancy and no further copy has fixed it.

\subsection{One state, two charts: \texorpdfstring{$C^{2}+r^{2}=1$}{C2+r2=1}}
\label{sec:wf-reconciliation}

The system is pure for the friend ($r=1$, on the surface) and mixed for Wigner
($r<1$, in the interior). These are not rival claims about a single fact; they are
the friend at the surface and Wigner at the interior of the \emph{same} state on
$S^{7}$, bound by the Mosseri--Dandoloff constraint
\begin{equation}
  C^{2}+r^{2}=1,
\end{equation}
in which the discrepancy in Bloch radius is exactly the concurrence Wigner
retains and the friend has resolved away. This is the Inside/Outside Equivalence
Principle (Proposition~\ref{prop:inside-outside}) made concrete, the same law written by the embedded participant and by
the external observer, and the relativity of facts of relational quantum
mechanics \autocite{Rovelli1996a} realised in geometry rather than asserted: the
\textsc{fanout} boundary lies \emph{between} $F$ and $W$, so a fact exists
relative to $F$ and not relative to $W$. Neither chart is privileged; one reads in
the epistemological or the ontological chart as the question requires, and their
mutual incompleteness is a theorem of the construction, not a defect of either.

\medskip
\noindent\emph{Aside: the Tsirelson bound and the absence of PR boxes.}
\label{sec:tsirelson}
A ball has a unique centre and no point beyond it. Maximal entanglement sits at
that centre \autocite{mosseri2001hopf}, $r=0$, $C=1$, and because the
Bloch radius cannot be negative, the concurrence cannot exceed unity. For a
two-qubit pure state the maximal CHSH correlator is
\begin{equation}
  \langle\mathrm{CHSH}\rangle_{\max}=2\sqrt{1+C^{2}}=2\sqrt{2-r^{2}}  
\end{equation}
\autocite{VerstraeteWolf2002},
which the centre maximises to $2\sqrt{2}$, the Tsirelson bound
\autocite{Cirelson1980QuantumInequality}. A Popescu--Rohrlich box \autocite{PopescuRohrlich1994},
with CHSH value $4$, would require $C^{2}=3$ and hence $r^{2}=-2$: a point two
units past the centre, which the ball does not contain. The super-quantum
correlation is excluded by the very positivity that places maximal entanglement
at the centre, one cannot trace out to a more-than-maximally-mixed reduced
state. The Tsirelson bound is in this sense the image, under the CHSH map, of the
boundedness of the $L^{2}$ state ball derived in
Section~\ref{sec:l2-metric-fibre}, and the prohibition on PR boxes shares its
origin with the roundness of the sphere: both trace to the reversible-continuity
metric, not to a separate principle.

\noindent This is the two-qubit, pure-state face of the bound, and its status is
worth stating plainly. It exhibits PR correlations as having no representing state
in the Hilbert geometry, and it recovers the value $2\sqrt{2}$ at the centre; but
forcing \emph{exactly} $2\sqrt{2}$ as a theorem of the geometry, rather than
reading it from the known concurrence--CHSH relation, and lifting the argument
from two qubits to the general Tsirelson bound, remain to be made rigorous. The
closest cognate in the literature is Held's derivation of the same bound from the
geometric algebra $G^{3}\cong\mathbb{H}$ \autocite{Held2021}, the same
quaternionic structure in a different language.

\begin{openproblem}[The Tsirelson bound as geometry]\label{op:tsirelson-geometry}
Derive $2\sqrt{2}$ as a theorem of the round state geometry rather than reading it off the
concurrence--CHSH relation, and lift the argument from the two-qubit pure case to the general
Tsirelson bound --- so that the exclusion of super-quantum correlation becomes, like the exclusion
exhibited above, a consequence of the $L^{2}$ ball's positivity (\S\ref{sec:l2-metric-fibre}) and
of nothing else. Held's quaternionic derivation \autocite{Held2021} is the nearest cognate and the
natural first tool.
\end{openproblem}

\input{monogamy_aside_2026-06-07}

\subsection{The redundancy ladder: from relative fact to objective fact}
\label{sec:wf-redundancy}

The friend's fact becomes a fact for Wigner only when Wigner copies it when he
measures $F$. It becomes a fact for a third party only when that party copies it
from Wigner, and so on up the von Neumann chain. The point at which the content
stops being negotiable is therefore not a point but a limit: once it has been
copied into enough mutually independent fragments that many parties can read it
without disturbing it, no one can practically un-copy it, and it is operationally
an objective fact. This is exactly quantum Darwinism \autocite{zurek_2009}:
objectivity is redundancy. It is not an external graft, because the framework's
account of decoherence already \emph{is} this picture---many fragments each take
their own copy of the clonable content, while their private phases, each
fragment's own self-ignorance, fail to add coherently and wash out.

This supplies the grading parameter that ``graded \textsc{fanout}'' had left
unnamed. \textsc{fanout} is graded not because the operation is fuzzy but because
redundancy $R$ runs continuously: $R=1$ is the single relational copy---the
friend's relative fact, reversible by Wigner---and $R\gg 1$ is the objective
fact, effectively irreversible because un-copying it would mean recovering every
fragment. The single-qubit friend is the $R=1$ endpoint; the macroscopic friend
is the same operation at large $R$. One axis, two ends. The relation to
Di~Biagio and Rovelli's distinction \autocite{DiBiagioRovelli2021} is then exact:
their \emph{stable} facts are the high-$R$ limit and their \emph{relative} facts
the low-$R$ regime, and the standard conditional-probability law they use to
define stability is the operational signature of redundancy having grown large.
Brukner's no-go theorem \autocite{Brukner2021}, that qubits are not observers,
sits on the same axis without conflict: it bites only against calling an $R=1$
relative fact an objective one, which the framework never does---the redundancy
ladder is precisely the structure that separates the relational definiteness a
single qubit does carry from the objectivity it does not.

The thermodynamics tracks the same axis. The single relational copy is unitary and
free; the $kT\ln 2$ cost \autocite{Landauer_1961} is incurred not by the friend's
act but by the amplification and stabilisation that drive $R$ upward, when
fragments are irreversibly committed. The relative fact is thermodynamically free;
the objective fact is what costs $kT\ln 2$.

\subsection{Phase promotion, and the bump in the carpet}
\label{sec:wf-phase-promotion}

One point keeps the load-bearing irrecoverability intact while explaining what
Wigner can do that the friend cannot. What Wigner recovers is \emph{not} the
$S^{1}$ self-ignorance; the global phase remains gauge for every observer and that
circle is spent for good. The quantity the friend resolves away is the relative
phase between the $\lvert 0\rangle$ and $\lvert 1\rangle$ components, and moving up
to $\mathbb{C}^{4}$ \emph{promotes} it: the single-qubit relative phase becomes the
relative phase between $\lvert 0\rangle\lvert F_{0}\rangle$ and
$\lvert 1\rangle\lvert F_{1}\rangle$, which lives in the recoverable $S^{3}$ fibre
(itself a circle bundle over $S^{2}$, so the promotion adds a two-sphere of
relative-phase directions to the lone gauge circle). Gauge phase at one level
becomes physical, interferable phase one level up---which is why Wigner can run
the interference experiment the friend cannot. But Wigner now carries his own
fresh global-phase $S^{1}$ at the top of the joint state: his own irreducible
self-ignorance. Ascending the tower does not eliminate the self-reference, only
displaces it one level higher, and the regress up the chain is precisely the
redundancy ladder of the previous subsection. This is the framework's bump in the
carpet: each observer sits atop a Hopf tower with a single gauge circle he can
never see around, and a measurement relocates the next observer's blind circle
into his own recoverable fibre without ever closing the regress.

Two cautions close the section. First, the recoverability Wigner
enjoys is the low-redundancy limit: for a macroscopic, high-$R$ friend the
reversal is exponentially hard and the fact is objective for all practical
purposes, so Wigner's ``no record'' is the $R\!\approx\!1$ idealisation of a
graded boundary, not a free lunch. Second, naming redundancy as the grading
parameter of \textsc{fanout}, and the phase-promotion reading itself, are
syntheses of the framework's parts: each rests on established pieces (the
clonability criterion, quantum Darwinism, the tensor-product promotion of phase)
but is assembled here, and is offered as the cleanest statement of why the two
observers never disagree about anything both can check.
% =====================================================================
%  SPLICE: "Where the chain goes next" -- forward-looking coda for
%  sec:wigners-friend, pointing to sec:sedenions-generations.
%  Created 2026-08-27 per JS.
%
%  THE POINT (JS, 2026-08-27). The Wigner's-friend charts stop at
%  Wigner because the Hopf tower runs out at the octonions, while the
%  text says the chain of observers obviously continues. The sedenions
%  show what the continuation adds: the observer of Wigner (fourth
%  wire, C (x) S) brings a NEW piece of structure -- the finite S_3 of
%  Brown's theorem, the generation label -- and the observer of the
%  observer of Wigner (fifth wire, the trigintaduonions) brings nothing
%  new in kind. This coda records that, so the section's termination at
%  Wigner reads as a fact about the tower rather than an economy of the
%  exposition.
%
%  WHAT THE ALGEBRA ACTUALLY SAYS (verified 2026-08-27, notebook T --
%  notebook_T_trigintaduonion_automorphisms_2026-08-27.py):
%    - Aut(S) = G_2 x S_3 (Brown 1967), DIRECT product; the new S_3 is
%      generated by the doubling involution and Brown's explicit
%      order-three map psi, which fixes 1 and the doubling unit and
%      rotates every (a, a u) plane, a imaginary, by a third of a turn.
%    - psi permutes the three octonion HALVINGS O, psi O, psi^2 O of S
%      (each doubled by the same unit; meeting only in R). It does NOT
%      permute the three copies through the quaternion core -- those are
%      cycled by a lifted G_2 element (the colour-line three-cycle); see
%      the sedenion section's two-triples remark and the corrections
%      ledger of the same date.
%    - Brown: Aut(A_{n-1}) x S_3 = Aut(A_n) with EQUALITY for n = 4, 5, 6
%      (as quoted in gresnigtgourlayvarma2023three, sec. 4.3); Eakin and
%      Sathaye prove Brown's conjecture in general (VERIFY-CITE for the
%      exact general form -- their step is Aut(A_{n+1}) = Aut(A_n) x G,
%      G in {Z/2, S_3}; over R with the standard parameters the branch is
%      S_3 at every rung >= 4, which is what notebook T checks directly
%      at rungs 4, 5, 6). Derivation algebra g_2 from the octonions on
%      (Schafer 1954; notebook T: dim Der = 14 at n = 3, 4, 5).
%  So "nothing new at the fifth wire" is stated PRECISELY: the fifth
%  doubling adds another finite S_3 and nothing else -- no connected
%  symmetry, no fibration; the increment is the same at every rung
%  thereafter. Not "silent": REPETITIVE. The prose says so.
%
%  STATUS: the algebra is ESTABLISHED (Brown; Eakin--Sathaye; Schafer)
%  and directly checked (notebook T). The reading of the repeated finite
%  factor as what the redundancy ladder consumes is a RHYME, flagged.
%  The reading of the fourth rung's discrete orientation as the
%  generation label is the sedenion section's CANDIDATE identification,
%  pointed to, not re-argued here.
%
%  PLACEMENT: a new \subsection at the END of sec:wigners-friend, after
%  the "Two cautions" paragraph. Because that paragraph reads "Two
%  cautions close the section.", the patch also changes that one word:
%  "close the section" -> "close the argument". Both anchors
%  uniqueness-asserted (the July lesson: anchors that also occur in
%  comments must be asserted unique).
%
%  CROSS-REFS OUT (all exist in the synced source except the two
%  MAP-LABELs):
%    sec:hopf-justification   -- the "two climbs" paragraph (length vs
%                                termination); this coda sharpens it
%    sec:wf-redundancy         -- the redundancy ladder (this file)
%    sec:wf-phase-promotion    -- phase promotion (this file)
%    sec:sedenions-generations -- the sedenion section (2026-08-27 rev.)
%    prop:termination          -- Hurwitz/Adams termination
%  MAP-LABELs:
%    app:notebook-record       -- the notebook-record appendix
%                                (appendices/appendix_notebook_record_
%                                2026-07-11.tex); substitute its label.
%  NEW KEYS (RIS delivered: cayley_dickson_automorphism_refs_2026-08-27.ris):
%    eakinsathaye1990automorphisms, schafer1954cayleydickson
%  IN-LIBRARY: brown1967generalized.
%
%  OPTIONAL CONSUMER-POINTER (not patched; one clause for the "two
%  climbs" paragraph of hopf_justification_section_2026-06-08.tex, after
%  "that chain has no top --- it is the bump in the carpet, unbounded in
%  length."):
%      "(unbounded in length, and --- as \S\ref{sec:wf-chain-type}
%      records --- bounded in type: after the fourth observer nothing
%      new in kind arrives)"
%
%  EDITED 2026-08-30 (per JS): one sentence added to the fourth-wire
%  paragraph anchoring the flavour sector's irremovable phase (the
%  CP-phase recognition) on this rung, with the development at
%  sec:sedenions-mixing (2026-08-30 revision). No new keys here.
%  British English; pure ASCII; \autocite; no bullets; spaced dashes.
% =====================================================================
% SPLICE-BEGIN wf-chain-type 2026-08-27

\subsection{Where the chain goes next: one more rung, then copies}
\label{sec:wf-chain-type}

The regress just described is unbounded in length, and nothing here shortens
it: the observer of Wigner carries a gauge circle of his own, and so does the
observer of that observer, as \S\ref{sec:hopf-justification} insists when it
separates the endlessness of the chain from the termination of the tower. But
the algebra says how much of the endless chain adds structure of a new kind,
and the answer is exact: one more step, and then none.

The count runs as follows. The triple system, friend, Wigner is three wires,
and three wires are the three Hopf rungs $\mathbb{C}$, $\mathbb{H}$,
$\mathbb{O}$ --- the completed link of the earlier caveat --- which is why the
charts above stop at Wigner: the chart of whoever looks in on the triple,
$S^{15}\to S^{8}$, is the last fibration the division algebras supply
(Proposition~\ref{prop:termination}). The observer of Wigner is a fourth wire,
and the fourth wire is the sedenions,
$(\mathbb{C}^{2})^{\otimes 4} = \mathbb{C}\otimes\mathbb{S}$
(\S\ref{sec:sedenions-generations}). No fibration comes with it --- the
sedenions have zero divisors and admit no map $S^{31}\to S^{16}$ --- so the
fourth observer gains no recoverable fibre of the tower's kind: the continuous
promotion of relative phase into a larger fibre, the mechanism of the previous
subsection, has no fourth instance. What the doubling adds instead is finite.
Brown's theorem gives $\mathrm{Aut}(\mathbb{S}) = G_{2}\times S_{3}$
\autocite{brown1967generalized}: the identity component is unchanged from the
octonions, and the entire gain is a component group, generated by an
order-three automorphism that fixes the new doubling direction and rotates
every imaginary direction of the completed link into its copy by a third of a
turn. The relative orientation of the link and its copy, which at the Hopf
rungs was a continuous phase promoted into a fibre, survives at the fourth rung
only as a discrete choice among three, and that choice is the generation label
(\S\ref{sec:sedenions-generations}). The fourth observer, then, sees no new
force and no new fibre; he sees the same link three ways, and cannot gauge the
difference, because a component group has no Lie algebra. His rung is also where
the flavour sector's one irremovable phase finds its anchor: three generations
are the threshold at which the mixing matrix first carries a phase no convention
removes, and \S\ref{sec:sedenions-mixing} reads that phase --- through the
paper's negativity--contextuality chain --- as this wire's contextual cargo,
real only from one rung above.

The observer of the observer of Wigner is a fifth wire, the trigintaduonions,
and here the tower has nothing new to say. Brown proved that the next two
doublings each add exactly one further $S_{3}$ and nothing else, Eakin and
Sathaye that Brown's pattern continues at every rung, and the derivation
algebra --- the Lie algebra of the identity component --- is $\mathfrak{g}_{2}$
from the octonions onward
\autocite{brown1967generalized,eakinsathaye1990automorphisms,schafer1954cayleydickson};
the fifth and sixth rungs are also checked directly in the notebook record
(Appendix~\ref{app:notebook-record}). So the fifth observer's doubling repeats
the fourth's gift --- another third-of-a-turn symmetry between the four-wire
structure and its copy, commuting with everything below --- and adds no
connected symmetry and no fibration. Nothing new in kind arrives at the fifth
wire or at any wire after it. The increments of the tower are, in order, a
force, a force, a label, and then the same finite permutation of copies at
every step. That is the precise sense in which the chain ``obviously
continues'': it continues, and after the fourth observer it has nothing further
to announce. The framework reads the repeated factor as what the redundancy
ladder of \S\ref{sec:wf-redundancy} consumes --- a permutation symmetry among
copies, which is the freedom a record has and a type does not --- and records
that reading as a structural resonance rather than a derivation.

The two climbs of \S\ref{sec:hopf-justification} are thereby sharpened. The
chain of observers is unbounded in length and bounded in type: three forces by
the third wire, one family label by the fourth, copies thereafter. The observer
of Wigner is the last observer whose arrival changes what there is.

% SPLICE-END wf-chain-type 2026-08-27

% =====================================================================
%  FIGURE: the Wigner's-friend chain mapped onto the Cayley--Dickson
%  ladder. Created 2026-08-30 per JS: the Hopf-fibration view of
%  Wigner's friend has been load-bearing for months without being drawn;
%  the fibrations themselves are DEFERRED to Szangolies's introduction
%  (szangolies2025standardmodel, in library), and this figure carries
%  only the mapping.
%
%  CONTENT (left to right, one panel per chart, staircase rising):
%    Panel 1: F's chart of S            -- S^3 -> S^2, fibre S^1,  C, one wire
%    Panel 2: W's chart of {S,F}        -- S^7 -> S^4, fibre S^3,  H, two wires
%    Panel 3: O_W's chart of {S,F,W}    -- S^15 -> S^8, fibre S^7, O, three wires
%             (the last fibration)
%    Panel 4: the next observer's chart of {S,F,W,O_W} -- C^16 = C (x) S,
%             no fibration (S^31 -/-> S^16); the doubling's S_3 badge:
%             the generation label, a third of a turn.
%  Each describer perches on his sealed box with a dashed circle over
%  his head (his own gauge circle -- the self-referential ignorance he
%  cannot see around); dashed promotion arrows carry each circle into
%  the next chart's fibre ("gauge here -> interferable one rung up",
%  sec:wf-phase-promotion); the third arrow terminates at "no fourth
%  fibration", with what survives being the discrete third-of-a-turn
%  (sec:wf-chain-type, sec:sedenions-generations). The regress
%  continues to the right: copies only.
%
%  PREAMBLE REQUIREMENT (confirm in main.tex; add if absent):
%      \usepackage{tikz}
%      \usetikzlibrary{arrows.meta}
%  PLACEMENT: \input by patch_wigner_hopf_ladder_fig_2026-08-30.py,
%  after the two-idealisations paragraph of sec:wigners-friend, with
%  the accompanying deferral paragraph (in the patch).
%  CROSS-REFS OUT (all real): sec:wf-phase-promotion,
%  sec:wf-chain-type, sec:sedenions-generations.
%  KEY: szangolies2025standardmodel (in library).
%  British English; pure ASCII; \autocite in the caption.
% =====================================================================
\begin{figure}[t]
\centering
\resizebox{\textwidth}{!}{%
\begin{tikzpicture}[>=Stealth, line cap=round, line join=round,
    every node/.style={font=\scriptsize},
    blind/.style={densely dashed, draw=black!75},
    promote/.style={->, densely dashed, black!60, semithick},
    sight/.style={->, dotted, black!55}]
% ---- stick-figure pic --------------------------------------------------
\tikzset{stick/.pic={
  \draw[pic actions] (0,0.44) circle (0.085);
  \draw[pic actions] (0,0.355) -- (0,0.10);
  \draw[pic actions] (-0.12,0.20) -- (0,0.29) -- (0.12,0.20);
  \draw[pic actions] (-0.10,-0.10) -- (0,0.10) -- (0.10,-0.10);
}}
% ---- the staircase (Cayley--Dickson ladder) ---------------------------
\draw[black!35, thick]
  (-0.35,-0.55) -- (3.27,-0.55) -- (3.27,0.25) -- (6.89,0.25)
  -- (6.89,1.05) -- (10.51,1.05) -- (10.51,1.85) -- (13.60,1.85)
  -- (13.60,2.65) -- (16.20,2.65);
\node[anchor=north] at (1.25,-0.60) {$\mathbb{C}$ --- one wire};
\node[anchor=north] at (4.87,0.20)  {$\mathbb{H}$ --- two wires};
\node[anchor=north] at (8.49,1.00)  {$\mathbb{O}$ --- three wires};
\node[anchor=north] at (12.11,1.80) {$\mathbb{S}$ --- four wires};
\node[anchor=north, black!55] at (14.45,2.60) {$\mathbb{T},\dots$ --- five wires on};
% =======================================================================
% Panel 1: F's chart of S
% =======================================================================
\draw[rounded corners=2pt, fill=black!5, draw=black!40] (0,0) rectangle (2.5,1.32);
\filldraw[black!70] (1.25,0.72) circle (0.075);
\node[font=\tiny] at (1.25,0.45) {$S$};
\pic[black] at (2.05,1.42) {stick};
\node[anchor=west] at (2.22,1.62) {$F$};
\draw[sight] (1.95,1.60) -- (1.45,1.02);
\draw[blind] (2.05,2.14) circle (0.11);
\node[anchor=east, align=right, font=\tiny, black!75] at (1.82,2.14)
  {his own gauge circle:\\ the phase he cannot\\ see around};
\node[anchor=south west, font=\tiny, black!75] at (0.02,1.34) {fibre $S^{1}$};
\node[anchor=north] at (1.25,-0.06) {$S^{3}\to S^{2}$};
% =======================================================================
% Panel 2: W's chart of {S,F}
% =======================================================================
\draw[rounded corners=2pt, fill=black!5, draw=black!40] (3.62,0.80) rectangle (6.12,2.12);
\filldraw[black!70] (4.07,1.52) circle (0.075);
\node[font=\tiny] at (4.07,1.25) {$S$};
\pic[black!55, scale=0.85] at (4.85,1.16) {stick};
\node[font=\tiny, black!55] at (4.85,0.94) {$F$};
\pic[black] at (5.67,2.22) {stick};
\node[anchor=west] at (5.84,2.42) {$W$};
\draw[sight] (5.57,2.40) -- (5.05,1.82);
\draw[blind] (5.67,2.94) circle (0.11);
\node[anchor=south west, font=\tiny, black!75] at (3.64,2.14) {fibre $S^{3}$};
\node[anchor=north] at (4.87,0.74) {$S^{7}\to S^{4}$};
\draw[promote] (2.17,2.14) to[out=15,in=170] (3.60,2.24);
\node[anchor=south, align=center, font=\tiny, black!60] at (2.85,2.32)
  {gauge below,\\ interferable above};
% =======================================================================
% Panel 3: O_W's chart of {S,F,W}
% =======================================================================
\draw[rounded corners=2pt, fill=black!5, draw=black!40] (7.24,1.60) rectangle (9.74,2.92);
\filldraw[black!70] (7.69,2.32) circle (0.075);
\node[font=\tiny] at (7.69,2.05) {$S$};
\pic[black!55, scale=0.85] at (8.35,1.96) {stick};
\node[font=\tiny, black!55] at (8.35,1.74) {$F$};
\pic[black!55, scale=0.85] at (8.95,1.96) {stick};
\node[font=\tiny, black!55] at (8.95,1.74) {$W$};
\pic[black] at (9.29,3.02) {stick};
\node[anchor=west] at (9.46,3.22) {$O_{W}$};
\draw[sight] (9.19,3.20) -- (8.65,2.62);
\draw[blind] (9.29,3.74) circle (0.11);
\node[anchor=south west, font=\tiny, black!75] at (7.26,2.94) {fibre $S^{7}$};
\node[anchor=north, align=center] at (8.49,1.54) {$S^{15}\to S^{8}$\\[-1pt]
  {\tiny the last fibration}};
\draw[promote] (5.79,2.94) to[out=15,in=170] (7.22,3.04);
% =======================================================================
% Panel 4: the next chart of {S,F,W,O_W} -- sedenions, no fibration
% =======================================================================
\draw[rounded corners=2pt, fill=black!5, draw=black!40] (10.86,2.40) rectangle (13.36,3.72);
\filldraw[black!70] (11.24,3.12) circle (0.075);
\node[font=\tiny] at (11.24,2.85) {$S$};
\pic[black!55, scale=0.80] at (11.78,2.78) {stick};
\node[font=\tiny, black!55] at (11.78,2.56) {$F$};
\pic[black!55, scale=0.80] at (12.28,2.78) {stick};
\node[font=\tiny, black!55] at (12.28,2.56) {$W$};
\pic[black!55, scale=0.80] at (12.78,2.78) {stick};
\node[font=\tiny, black!55] at (12.78,2.56) {$O_{W}$};
\pic[black!45] at (14.35,2.78) {stick};
\draw[sight] (14.18,3.02) -- (13.42,3.12);
\draw[blind] (14.35,3.50) circle (0.11);
\node[anchor=west, align=left, font=\tiny, black!60] at (14.58,3.48)
  {the next observer,\\ and all after:\\ copies only};
\node[anchor=north] at (12.11,2.34) {$\mathbb{C}^{16}=\mathbb{C}\otimes\mathbb{S}$};
% the broken third promotion
\draw[promote] (9.41,3.74) to[out=15,in=185] (10.68,3.86);
\node[font=\scriptsize, black!60] at (10.80,3.92) {$\times$};
\node[anchor=south west, font=\tiny, black!75] at (10.88,3.74)
  {no fourth fibration: $S^{31}\not\to S^{16}$};
% the S3 badge: three dots at 120 degrees with a rotation arrow
\draw[black!75] (11.45,4.30) circle (0.17);
\filldraw[black!75] (11.45,4.47) circle (0.028);
\filldraw[black!75] (11.303,4.215) circle (0.028);
\filldraw[black!75] (11.597,4.215) circle (0.028);
\draw[->, black!75] (11.67,4.36) arc[start angle=20, end angle=100, radius=0.235];
\node[anchor=west, align=left, font=\tiny, black!75] at (11.70,4.36)
  {$S_{3}$: the generation label\\ (a third of a turn)};
\draw[promote, black!70] (10.86,3.96) to[out=60,in=200] (11.24,4.22);
\end{tikzpicture}%
}
\caption{The Wigner's-friend chain mapped onto the Cayley--Dickson ladder. Each
observer holds a chart of the sealed box below him, one rung further up the
ladder: the friend's chart of the system is the complex Hopf bundle, Wigner's
chart of the sealed laboratory the quaternionic, and the chart of the observer
of Wigner --- of the completed triple --- the octonionic, the tower's last
fibration. Every chart-holder carries a dashed circle over his head, his own
gauge phase: the self-referential ignorance he cannot see around. Measurement
promotes each circle into the next chart's fibre, where what was gauge below is
interferable above (\S\ref{sec:wf-phase-promotion}); the third promotion has no
fibration to land in --- the sedenions have zero divisors and no map
$S^{31}\to S^{16}$ --- and what survives of it is discrete: the doubling's
$S_{3}$, a third of a turn rotating the link into its copy, the generation
label (\S\ref{sec:wf-chain-type}, \S\ref{sec:sedenions-generations}). The
regress of circles continues unbounded to the right, but from the fifth wire
onward it adds copies only. The friend drawn as a single qubit is the
redundancy-one idealisation discussed in the text; for the fibrations
themselves, deferred rather than repeated here, see the introduction in
\autocite{szangolies2025standardmodel}.}
\label{fig:wf-hopf-ladder}
\end{figure}

%  (WF backward-references sec:l2-metric-fibre: order-consistent.)
%
%  AUDITED 2026-07-06 (final pass): Lamperti/Banach classification,
%  Hopf map components, submersion facts, Hardy K=N^r reading all
%  verified; the recoverable-vs-gauge-private fibre distinction
%  confirmed canonical and untouched; file pure ASCII. Fixes: opening
%  rewritten (Popper C+/C/E forward-voiced to sec:wigners-friend;
%  continuity claim made CONDITIONAL per JS -- "assuming those premises
%  can be proved", op:limit-realisability cited); Aaronson paragraph
%  de-duplicated against sec:continuity by stitch; sec:born-rule was
%  dangling paper-wide -- label minted at its canonical home (the Born
%  subsection of sec:self-reference) and the claim here recalibrated
%  ("forced, not chosen" -> lowest-degree-invariant, geometric face);
%  record-absence upgraded from interpretive to theorem via
%  prop:f1-dictionary while the self-ignorance identification stays
%  interpretive; grammar and dash punctuation repaired; field-not-fixed
%  bound pointed at sec:self-reference.
%
%  CITATIONS: existing keys Hardy2001a, Aaronson2004h, mosseri2001hopf.
%  NEW keys (RIS in L2_metric_refs_2026-07-06.ris): Banach1932,
%  Lamperti1958, Hopf1931, Urbantke2003, GluckWarnerZiller1986,
%  wooters_1981 (key spelling preserved as used; note the single-t).
%  British English; \autocite; no bullets in prose.
% =====================================================================

\section{The Euclidean metric on the state sphere, and the global phase as self-referential ignorance}
\label{sec:l2-metric-fibre}
Popper's systems $C^{+}$, $C$ and $E$ --- the predictor, its target and their shared
environment, whose argument \S\ref{sec:wigners-friend} anchors and quotes --- are, in modern
dress, Wigner, his friend and a qubit $S$ to be measured. The natural arena for analysing them
is the Hopf fibration and the Bloch sphere; but even to draw the Bloch sphere one must first
justify the $L^{2}$ metric, and this section earns that right. The principle doing the work is
the reversible continuity of \S\ref{sec:continuity} --- derived there from self-reference
conditionally, on the premises priced as Open Problem~\ref{op:limit-realisability}, of which one, that every
transformation of the regime is a phase of an inside structure, remains open; the continuum
itself is the completion the Shannon step declares (\S\ref{sec:complex-amplitudes}), the
classical convention, and \S\ref{sec:cencov-route} below reaches the same metric by a
second route that consumes it and nothing else that is open. Assuming that premise can be proved,
everything below inherits the derivation; we carry that status, and do not upgrade it.

We have established that the amplitudes of a single qubit form a complex vector space
$\mathbb{C}^{2}\cong\mathbb{R}^{4}$. We have not yet fixed which norm on this space singles out
the physical state manifold. That choice is not free: the principle of
\emph{reversible continuity} forces the Euclidean ($L^{2}$) norm, and hence the round
three-sphere, before any appeal to the Born rule. The quadratic structure usually \emph{posited}
as the Born rule is therefore already present in the metric, and the global phase emerges as a
geometrically forced fibre rather than an added convention.

\subsection{Reversible continuity selects the 2-norm}
\label{sec:reversible-continuity-l2}

The principle is the one isolated by Hardy as the single axiom distinguishing quantum from
classical probability \autocite{Hardy2001a}:
\emph{there exists a continuous reversible transformation between any two pure states.}
Operationally a reversible transformation of states is a symmetry of the state manifold, hence
norm-preserving; the demand is that such symmetries form a \emph{connected} group acting
\emph{transitively} on the pure states. Linearity is inherited from the complex vector-space
structure already in hand, so the admissible transformations are the complex-linear isometries
of whatever norm we place on $\mathbb{C}^{2}$.

Take the candidate norms to be the $\ell^{p}$ family,
$\lVert x\rVert_{p}=\bigl(\sum_{k}\lvert x_{k}\rvert^{p}\bigr)^{1/p}$, with the state manifold
the unit $\ell^{p}$ sphere. Their linear isometries are completely classified.


\begin{proposition}[Lamperti; after Banach]
\label{prop:lamperti}
For $1\le p<\infty$, $p\neq 2$, every linear isometry of $L^{p}$ (and hence of $\ell^{p}$) is a
generalised permutation: a weighted composition operator which, in the finite-dimensional case,
acts as $(x_{k})\mapsto(\omega_{k}\,x_{\sigma(k)})$ with $\lvert\omega_{k}\rvert=1$ and $\sigma$
a permutation of the coordinates.
\end{proposition}

\noindent Banach gave this for $L^{p}[0,1]$; the general $\sigma$-finite statement is Lamperti's
\autocite{Banach1932, Lamperti1958}. Its consequence for our four real dimensions is decisive.
For every $p\neq 2$ the isometry group is, in the real case, the finite group of signed
permutations (the hyperoctahedral group $B_{4}$), and in the complex case the group
$\mathbb{T}^{n}\rtimes S_{n}$ of diagonal phases composed with permutations. Even the complex
group, though it contains the continuous torus $\mathbb{T}^{n}$, fails to act transitively on
the sphere: the orbit of a basis state $e_{1}$ is only the phased basis vectors
$\{\omega\,e_{k}\}$, a measure-zero union of circles. No continuous reversible transformation
carries a basis state into a generic superposition such as $(e_{1}+e_{2})/\sqrt{2}$. Reversible
continuity therefore \emph{fails} for every $p\neq 2$.

At $p=2$ the isometry group is $O(4)$ in the real picture and $U(2)$ in the complex one; the
identity component acts transitively on the unit sphere, and any two unit vectors are joined by
a continuous one-parameter subgroup. Among the $\ell^{p}$ geometries, reversible continuity is
satisfiable at $p=2$ alone.

Aaronson states the same selection in the quantum-foundations setting and sharpens it into a
physical, rather than merely aesthetic, necessity \autocite{Aaronson2004h}: the $2$-norm is the
unique $p$-norm admitting a nontrivial continuous group of norm-preserving linear maps, every
other $p$ yielding only generalised permutations; and the alternative of keeping linearity and
reversibility while \emph{dropping} norm preservation lands one in the full general linear group
$\mathrm{GL}(n,\mathbb{C})$, which permits superluminal signalling and collapses the complexity
class $\mathsf{PP}$ into polynomial time --- the same fact \S\ref{sec:continuity} uses to price
the classical--quantum cut. Bare reversibility on the amplitude space buys only
$\mathrm{GL}(n,\mathbb{C})$; what carves the unitary group $U(n)$ out of it is preservation of
the quadratic ($\ell^{2}$) form, and reversibility of the resulting maps is then a corollary
rather than a premise.

The conclusion is that the physical state manifold of a single qubit is the unit sphere of the
Euclidean norm,
\begin{equation}
  S^{3}=\Bigl\{\,x\in\mathbb{R}^{4}\;:\;\textstyle\sum_{i=1}^{4}x_{i}^{2}=1\,\Bigr\}
  =\bigl\{\,\psi\in\mathbb{C}^{2}\;:\;\lvert z_{1}\rvert^{2}+\lvert z_{2}\rvert^{2}=1\,\bigr\},
\end{equation}
with $z_{j}=x_{2j-1}+\mathrm{i}\,x_{2j}$, carrying the round metric induced from the ambient
inner product,
\begin{equation}
  \mathrm{d}s^{2}=\sum_{i=1}^{4}\mathrm{d}x_{i}^{2}\,\Big|_{S^{3}}
  =\operatorname{Re}\langle\mathrm{d}\psi,\mathrm{d}\psi\rangle,
  \qquad
  \langle\phi,\psi\rangle=\bar\phi_{1}\psi_{1}+\bar\phi_{2}\psi_{2}.
\end{equation}
One bound deserves explicit statement, in the interest of claiming no more than is proved:
reversible continuity fixes the \emph{metric} as $L^{2}$, but not the underlying field. Real and
quaternionic quantum theory equally satisfy reversible continuity and equally carry Euclidean
state-space metrics; the restriction to $\mathbb{C}$ is pinned elsewhere in the framework ---
the analytic continuation of the information measure and the additivity of phase information,
\S\ref{sec:self-reference} --- not here.

\subsection{Hardy's independent route}
\label{sec:hardy-route}

The same conclusion is reached by Hardy along a route that never invokes
Proposition~\ref{prop:lamperti} \autocite{Hardy2001a}. His reconstruction turns on the counting
relation $K=N^{r}$, where $N$ is the number of states distinguishable in a single measurement
and $K$ the number of fiducial probabilities needed to specify a state. The continuity axiom
excludes the classical case $r=1$: there the pure states are the $N$ vertices of a simplex and
the reversible maps merely permute them --- a discrete group of motions --- so there is no
continuous path between distinct pure states, which is the simplex face of the same $\ell^{1}$
isometry fact. A simplicity assumption then selects $r=2$, yielding complex quantum theory with
its quadratic norm and trace rule. The quadratic structure is not extracted from continuity as a
separate step; it is built into Hardy's first axiom (probabilities linear in the fiducial set)
together with $K=N^{2}$.

We thus have two genuinely independent demonstrations that reversible continuity forces the
$L^{2}$ metric: Hardy's operational $K=N^{2}$ counting, and the Lamperti--Aaronson
classification of the isometries. They converge on the same step by different mathematics, and
should be cited as two roads rather than one road traversed twice. Both roads start from
continuity. A third, in the next subsection, starts one step earlier, from the value continuum
that continuity itself rests on, and reaches the same metric without passing through the
transformation group at all.

\subsection{The statistical route: \v{C}encov's theorem and the every-basis lemma}
\label{sec:cencov-route}

The value continuum that \S\ref{sec:continuity} isolates as the one open premise reaches the
$L^{2}$ metric on its own, by a route that uses no transformation group. \v{C}encov's theorem
states that the only Riemannian metric on the simplex of probability distributions invariant
under coarse-graining --- under the Markov maps that merge or relabel outcomes --- is the
Fisher--Rao metric \autocite{bengtsson2006geometry}, and in the coordinates $x_{k} =
\sqrt{p_{k}}$ the Fisher--Rao metric is the round metric on the positive orthant of the unit
sphere. The premise the theorem consumes is that the simplex is a manifold, which is the value
continuum declared in \S\ref{sec:complex-amplitudes}, and the principle it uses, that
distinguishability does not depend on how outcomes are grouped, is one the Shannon bookkeeper of
that section already grants. What it does not supply is the passage from one orthant to the
whole projective space; that is the following lemma.

\begin{lemma}[The every-basis lemma]\label{lem:every-basis}
Let $g$ be a Riemannian metric on $\mathbb{CP}^{K-1}$ whose restriction to the positive real
orthant of every orthonormal basis --- the states $\sum_{k}\sqrt{p_{k}}\,e_{k}$ with $p$ in the
simplex --- is the Fisher--Rao metric. Then $g$ is the Fubini--Study metric, whose linear lift
is the $L^{2}$ inner product.
\end{lemma}

\begin{proof}
A tangent vector to $\mathbb{CP}^{K-1}$ at a state $\psi$ is represented by a vector $\phi$
orthogonal to $\psi$, and the curve $t \mapsto \cos t\,\psi + \sin t\,\phi$ with real $t$ lies
in the positive orthant of any orthonormal basis that contains $\psi$ and $\phi$. Along that
curve the Fisher--Rao metric is the round metric of the sphere in the coordinates $\sqrt{p}$,
which is the Fubini--Study metric restricted to real amplitudes; so $g$ and the Fubini--Study
metric agree on every tangent vector at every point, the direction $i\phi$ being reached by the
basis that contains $i\phi$ in place of $\phi$.
\end{proof}

The lemma's premise, that every orthonormal basis is a legitimate measurement and none is
privileged, is the fourth face of the pre-measurement regime seen from the measurement side
(\S\ref{sec:continuity}); in the stabiliser fragment of quantum theory it fails in the same way
everything else does, only the Pauli bases being measurements there. So the $L^{2}$ metric now
rests on two legs that share exactly one premise. Through reversible continuity it uses the
value continuum, closure and homogeneity (Proposition~\ref{prop:lamperti}); through \v{C}encov
it uses the value continuum, coarse-graining invariance and the equal standing of bases
(Lemma~\ref{lem:every-basis}). Neither leg needs anything else that is open, and nothing could
avoid the continuum, since a metric is a statement about one. That is the status of everything
downstream of this section, the Hopf fibrations included: derived from the completion every physicist grants, which carries no quantum content, and from
one premise about transformations that the framework derives in outline and has not yet proved
(Open Problem~\ref{op:limit-realisability}).

\subsection{The Hopf fibration and the Bloch sphere}
\label{sec:hopf-bloch}

The round metric is exactly what makes the global phase a clean symmetry. The phase action
$\psi\mapsto\mathrm{e}^{\mathrm{i}\theta}\psi$ is a one-parameter subgroup of
$U(2)\subset SO(4)$ and so acts on $S^{3}$ by \emph{isometries}, a statement that holds only
because the metric is the $L^{2}$ one just derived. Its orbits are great circles $S^{1}$, all of
circumference $2\pi$, and the quotient is the complex projective line,
\begin{equation}
  S^{3}\big/U(1)=\mathbb{CP}^{1}\cong S^{2},
\end{equation}
the Bloch sphere. This quotient is the complex Hopf fibration
\autocite{Hopf1931, Urbantke2003, mosseri2001hopf}
\begin{equation}
  S^{1}\hookrightarrow S^{3}\xrightarrow{\;h\;}S^{2},
\end{equation}
whose projection sends an amplitude to its Bloch vector,
\begin{equation}
  h(\psi)=\bigl(\langle\psi\rvert\sigma_{x}\lvert\psi\rangle,\;
                 \langle\psi\rvert\sigma_{y}\lvert\psi\rangle,\;
                 \langle\psi\rvert\sigma_{z}\lvert\psi\rangle\bigr)
        =\bigl(2\operatorname{Re}\bar z_{1}z_{2},\;
               2\operatorname{Im}\bar z_{1}z_{2},\;
               \lvert z_{1}\rvert^{2}-\lvert z_{2}\rvert^{2}\bigr).
\end{equation}
The Hopf map is quadratic in the amplitudes, and this is the geometric face of the Born rule:
the third component is $p_{0}-p_{1}$ with $p_{k}=\lvert z_{k}\rvert^{2}$, and the squaring is
the lowest-degree $U(1)$-invariant available on the fibre --- the record-theoretic derivation of
the same rule is \S\ref{sec:born-rule}'s. Because the metric is round, the projection is a
Riemannian submersion with totally geodesic fibres of constant length
\autocite{GluckWarnerZiller1986}: the same amount of phase is hidden at every point of the base,
so the phase is hidden \emph{uniformly}, not in a position-dependent way that the metric could
partially betray.

\subsection{The fibre as self-referential ignorance}
\label{sec:fibre-self-ignorance}

The interpretation the framework places on this geometry --- and here we are reading physics
into the mathematics, so the claim is interpretive rather than proved --- is that the $S^{1}$
fibre is the observer's own un-narratable self-ignorance. The global phase is the imaginary part
of the complex information measure $S=-\ln\lvert p\rvert-\mathrm{i}\theta$: information the
observer lacks about their own state. Lacking it, they cannot communicate it; there is no
\textsc{fanout} of information one does not possess. Uncommunicable, it cannot enter any
measurement record --- and this last step is no longer interpretive:
Proposition~\ref{prop:f1-dictionary} proves the global phase absent from every classical record
a copying structure induces --- so every observable is $U(1)$-invariant; and a $U(1)$-invariant
quantity is a function on the quotient. The physical state is therefore the orbit, not the
representative point---the base $S^{2}$---and the $S^{1}$ fibre carries exactly the inaccessible
self-referential phase. It is a null direction of statistical distinguishability: two global
phases are indistinguishable by any measurement whatever \autocite{wooters_1981}, so the fibre
is a gauge direction without physical content. This phase is irrecoverably discarded, and must
remain so: it is the relative phase, living in the $S^{3}$ fibre of the quaternionic Hopf
fibration at two qubits, that is genuinely recoverable, whereas the single-qubit global phase of
the complex fibration is gauge-private without remainder. The self-reference that opened the
derivation thus reappears, at its end, as the circle the observer can never see around.

%      % EDITED 2026-09-12 scripts/patch_revert_sedenion_repoint_2026-09-12T1600.py
% EDITED 2026-09-11 scripts/patch_undefined_section_refs_2026-09-11T2030.py
% =====================================================================
%  SECTION: Wigner's friend in the Hopf picture (v2; supersedes the
%  2026-06-07 draft). AUDITED 2026-07-06 (final pass, fifth stop).
%
%  PLACEMENT (reordered chain, JS 2026-07-06): position 4 --
%      % EDITED 2026-10-29 scripts/patch_continuity_tower_premise_2026-10-29T1500.py
% EDITED 2026-09-15 scripts/patch_value_continuum_declared_cencov_2026-09-15T1540.py
% EDITED 2026-09-15 scripts/patch_continuity_connected_state_space_2026-09-15T1430.py
% =====================================================================
%  SECTION: The Euclidean metric on the state sphere, and the global
%  phase as self-referential ignorance
%  Origin: drafted in an earlier conversation (2026-06-07), edited by
%  JS; audited into the current architecture 2026-07-06.
%
%  PLACEMENT in main.tex (reordered chain, JS 2026-07-06): position 3 --
%      \input{self_reference_section_2026-07-04}
%      \input{continuity_section_2026-07-04}
%      \input{L2_metric_global_phase_section_2026_06_07}   <-- here
%      \input{wigners_friend_in_the_hopf_picture_2026-06-13}
%  (WF backward-references sec:l2-metric-fibre: order-consistent.)
%
%  AUDITED 2026-07-06 (final pass): Lamperti/Banach classification,
%  Hopf map components, submersion facts, Hardy K=N^r reading all
%  verified; the recoverable-vs-gauge-private fibre distinction
%  confirmed canonical and untouched; file pure ASCII. Fixes: opening
%  rewritten (Popper C+/C/E forward-voiced to sec:wigners-friend;
%  continuity claim made CONDITIONAL per JS -- "assuming those premises
%  can be proved", op:limit-realisability cited); Aaronson paragraph
%  de-duplicated against sec:continuity by stitch; sec:born-rule was
%  dangling paper-wide -- label minted at its canonical home (the Born
%  subsection of sec:self-reference) and the claim here recalibrated
%  ("forced, not chosen" -> lowest-degree-invariant, geometric face);
%  record-absence upgraded from interpretive to theorem via
%  prop:f1-dictionary while the self-ignorance identification stays
%  interpretive; grammar and dash punctuation repaired; field-not-fixed
%  bound pointed at sec:self-reference.
%
%  CITATIONS: existing keys Hardy2001a, Aaronson2004h, mosseri2001hopf.
%  NEW keys (RIS in L2_metric_refs_2026-07-06.ris): Banach1932,
%  Lamperti1958, Hopf1931, Urbantke2003, GluckWarnerZiller1986,
%  wooters_1981 (key spelling preserved as used; note the single-t).
%  British English; \autocite; no bullets in prose.
% =====================================================================

\section{The Euclidean metric on the state sphere, and the global phase as self-referential ignorance}
\label{sec:l2-metric-fibre}
Popper's systems $C^{+}$, $C$ and $E$ --- the predictor, its target and their shared
environment, whose argument \S\ref{sec:wigners-friend} anchors and quotes --- are, in modern
dress, Wigner, his friend and a qubit $S$ to be measured. The natural arena for analysing them
is the Hopf fibration and the Bloch sphere; but even to draw the Bloch sphere one must first
justify the $L^{2}$ metric, and this section earns that right. The principle doing the work is
the reversible continuity of \S\ref{sec:continuity} --- derived there from self-reference
conditionally, on the premises priced as Open Problem~\ref{op:limit-realisability}, of which one, that every
transformation of the regime is a phase of an inside structure, remains open; the continuum
itself is the completion the Shannon step declares (\S\ref{sec:complex-amplitudes}), the
classical convention, and \S\ref{sec:cencov-route} below reaches the same metric by a
second route that consumes it and nothing else that is open. Assuming that premise can be proved,
everything below inherits the derivation; we carry that status, and do not upgrade it.

We have established that the amplitudes of a single qubit form a complex vector space
$\mathbb{C}^{2}\cong\mathbb{R}^{4}$. We have not yet fixed which norm on this space singles out
the physical state manifold. That choice is not free: the principle of
\emph{reversible continuity} forces the Euclidean ($L^{2}$) norm, and hence the round
three-sphere, before any appeal to the Born rule. The quadratic structure usually \emph{posited}
as the Born rule is therefore already present in the metric, and the global phase emerges as a
geometrically forced fibre rather than an added convention.

\subsection{Reversible continuity selects the 2-norm}
\label{sec:reversible-continuity-l2}

The principle is the one isolated by Hardy as the single axiom distinguishing quantum from
classical probability \autocite{Hardy2001a}:
\emph{there exists a continuous reversible transformation between any two pure states.}
Operationally a reversible transformation of states is a symmetry of the state manifold, hence
norm-preserving; the demand is that such symmetries form a \emph{connected} group acting
\emph{transitively} on the pure states. Linearity is inherited from the complex vector-space
structure already in hand, so the admissible transformations are the complex-linear isometries
of whatever norm we place on $\mathbb{C}^{2}$.

Take the candidate norms to be the $\ell^{p}$ family,
$\lVert x\rVert_{p}=\bigl(\sum_{k}\lvert x_{k}\rvert^{p}\bigr)^{1/p}$, with the state manifold
the unit $\ell^{p}$ sphere. Their linear isometries are completely classified.


\begin{proposition}[Lamperti; after Banach]
\label{prop:lamperti}
For $1\le p<\infty$, $p\neq 2$, every linear isometry of $L^{p}$ (and hence of $\ell^{p}$) is a
generalised permutation: a weighted composition operator which, in the finite-dimensional case,
acts as $(x_{k})\mapsto(\omega_{k}\,x_{\sigma(k)})$ with $\lvert\omega_{k}\rvert=1$ and $\sigma$
a permutation of the coordinates.
\end{proposition}

\noindent Banach gave this for $L^{p}[0,1]$; the general $\sigma$-finite statement is Lamperti's
\autocite{Banach1932, Lamperti1958}. Its consequence for our four real dimensions is decisive.
For every $p\neq 2$ the isometry group is, in the real case, the finite group of signed
permutations (the hyperoctahedral group $B_{4}$), and in the complex case the group
$\mathbb{T}^{n}\rtimes S_{n}$ of diagonal phases composed with permutations. Even the complex
group, though it contains the continuous torus $\mathbb{T}^{n}$, fails to act transitively on
the sphere: the orbit of a basis state $e_{1}$ is only the phased basis vectors
$\{\omega\,e_{k}\}$, a measure-zero union of circles. No continuous reversible transformation
carries a basis state into a generic superposition such as $(e_{1}+e_{2})/\sqrt{2}$. Reversible
continuity therefore \emph{fails} for every $p\neq 2$.

At $p=2$ the isometry group is $O(4)$ in the real picture and $U(2)$ in the complex one; the
identity component acts transitively on the unit sphere, and any two unit vectors are joined by
a continuous one-parameter subgroup. Among the $\ell^{p}$ geometries, reversible continuity is
satisfiable at $p=2$ alone.

Aaronson states the same selection in the quantum-foundations setting and sharpens it into a
physical, rather than merely aesthetic, necessity \autocite{Aaronson2004h}: the $2$-norm is the
unique $p$-norm admitting a nontrivial continuous group of norm-preserving linear maps, every
other $p$ yielding only generalised permutations; and the alternative of keeping linearity and
reversibility while \emph{dropping} norm preservation lands one in the full general linear group
$\mathrm{GL}(n,\mathbb{C})$, which permits superluminal signalling and collapses the complexity
class $\mathsf{PP}$ into polynomial time --- the same fact \S\ref{sec:continuity} uses to price
the classical--quantum cut. Bare reversibility on the amplitude space buys only
$\mathrm{GL}(n,\mathbb{C})$; what carves the unitary group $U(n)$ out of it is preservation of
the quadratic ($\ell^{2}$) form, and reversibility of the resulting maps is then a corollary
rather than a premise.

The conclusion is that the physical state manifold of a single qubit is the unit sphere of the
Euclidean norm,
\begin{equation}
  S^{3}=\Bigl\{\,x\in\mathbb{R}^{4}\;:\;\textstyle\sum_{i=1}^{4}x_{i}^{2}=1\,\Bigr\}
  =\bigl\{\,\psi\in\mathbb{C}^{2}\;:\;\lvert z_{1}\rvert^{2}+\lvert z_{2}\rvert^{2}=1\,\bigr\},
\end{equation}
with $z_{j}=x_{2j-1}+\mathrm{i}\,x_{2j}$, carrying the round metric induced from the ambient
inner product,
\begin{equation}
  \mathrm{d}s^{2}=\sum_{i=1}^{4}\mathrm{d}x_{i}^{2}\,\Big|_{S^{3}}
  =\operatorname{Re}\langle\mathrm{d}\psi,\mathrm{d}\psi\rangle,
  \qquad
  \langle\phi,\psi\rangle=\bar\phi_{1}\psi_{1}+\bar\phi_{2}\psi_{2}.
\end{equation}
One bound deserves explicit statement, in the interest of claiming no more than is proved:
reversible continuity fixes the \emph{metric} as $L^{2}$, but not the underlying field. Real and
quaternionic quantum theory equally satisfy reversible continuity and equally carry Euclidean
state-space metrics; the restriction to $\mathbb{C}$ is pinned elsewhere in the framework ---
the analytic continuation of the information measure and the additivity of phase information,
\S\ref{sec:self-reference} --- not here.

\subsection{Hardy's independent route}
\label{sec:hardy-route}

The same conclusion is reached by Hardy along a route that never invokes
Proposition~\ref{prop:lamperti} \autocite{Hardy2001a}. His reconstruction turns on the counting
relation $K=N^{r}$, where $N$ is the number of states distinguishable in a single measurement
and $K$ the number of fiducial probabilities needed to specify a state. The continuity axiom
excludes the classical case $r=1$: there the pure states are the $N$ vertices of a simplex and
the reversible maps merely permute them --- a discrete group of motions --- so there is no
continuous path between distinct pure states, which is the simplex face of the same $\ell^{1}$
isometry fact. A simplicity assumption then selects $r=2$, yielding complex quantum theory with
its quadratic norm and trace rule. The quadratic structure is not extracted from continuity as a
separate step; it is built into Hardy's first axiom (probabilities linear in the fiducial set)
together with $K=N^{2}$.

We thus have two genuinely independent demonstrations that reversible continuity forces the
$L^{2}$ metric: Hardy's operational $K=N^{2}$ counting, and the Lamperti--Aaronson
classification of the isometries. They converge on the same step by different mathematics, and
should be cited as two roads rather than one road traversed twice. Both roads start from
continuity. A third, in the next subsection, starts one step earlier, from the value continuum
that continuity itself rests on, and reaches the same metric without passing through the
transformation group at all.

\subsection{The statistical route: \v{C}encov's theorem and the every-basis lemma}
\label{sec:cencov-route}

The value continuum that \S\ref{sec:continuity} isolates as the one open premise reaches the
$L^{2}$ metric on its own, by a route that uses no transformation group. \v{C}encov's theorem
states that the only Riemannian metric on the simplex of probability distributions invariant
under coarse-graining --- under the Markov maps that merge or relabel outcomes --- is the
Fisher--Rao metric \autocite{bengtsson2006geometry}, and in the coordinates $x_{k} =
\sqrt{p_{k}}$ the Fisher--Rao metric is the round metric on the positive orthant of the unit
sphere. The premise the theorem consumes is that the simplex is a manifold, which is the value
continuum declared in \S\ref{sec:complex-amplitudes}, and the principle it uses, that
distinguishability does not depend on how outcomes are grouped, is one the Shannon bookkeeper of
that section already grants. What it does not supply is the passage from one orthant to the
whole projective space; that is the following lemma.

\begin{lemma}[The every-basis lemma]\label{lem:every-basis}
Let $g$ be a Riemannian metric on $\mathbb{CP}^{K-1}$ whose restriction to the positive real
orthant of every orthonormal basis --- the states $\sum_{k}\sqrt{p_{k}}\,e_{k}$ with $p$ in the
simplex --- is the Fisher--Rao metric. Then $g$ is the Fubini--Study metric, whose linear lift
is the $L^{2}$ inner product.
\end{lemma}

\begin{proof}
A tangent vector to $\mathbb{CP}^{K-1}$ at a state $\psi$ is represented by a vector $\phi$
orthogonal to $\psi$, and the curve $t \mapsto \cos t\,\psi + \sin t\,\phi$ with real $t$ lies
in the positive orthant of any orthonormal basis that contains $\psi$ and $\phi$. Along that
curve the Fisher--Rao metric is the round metric of the sphere in the coordinates $\sqrt{p}$,
which is the Fubini--Study metric restricted to real amplitudes; so $g$ and the Fubini--Study
metric agree on every tangent vector at every point, the direction $i\phi$ being reached by the
basis that contains $i\phi$ in place of $\phi$.
\end{proof}

The lemma's premise, that every orthonormal basis is a legitimate measurement and none is
privileged, is the fourth face of the pre-measurement regime seen from the measurement side
(\S\ref{sec:continuity}); in the stabiliser fragment of quantum theory it fails in the same way
everything else does, only the Pauli bases being measurements there. So the $L^{2}$ metric now
rests on two legs that share exactly one premise. Through reversible continuity it uses the
value continuum, closure and homogeneity (Proposition~\ref{prop:lamperti}); through \v{C}encov
it uses the value continuum, coarse-graining invariance and the equal standing of bases
(Lemma~\ref{lem:every-basis}). Neither leg needs anything else that is open, and nothing could
avoid the continuum, since a metric is a statement about one. That is the status of everything
downstream of this section, the Hopf fibrations included: derived from the completion every physicist grants, which carries no quantum content, and from
one premise about transformations that the framework derives in outline and has not yet proved
(Open Problem~\ref{op:limit-realisability}).

\subsection{The Hopf fibration and the Bloch sphere}
\label{sec:hopf-bloch}

The round metric is exactly what makes the global phase a clean symmetry. The phase action
$\psi\mapsto\mathrm{e}^{\mathrm{i}\theta}\psi$ is a one-parameter subgroup of
$U(2)\subset SO(4)$ and so acts on $S^{3}$ by \emph{isometries}, a statement that holds only
because the metric is the $L^{2}$ one just derived. Its orbits are great circles $S^{1}$, all of
circumference $2\pi$, and the quotient is the complex projective line,
\begin{equation}
  S^{3}\big/U(1)=\mathbb{CP}^{1}\cong S^{2},
\end{equation}
the Bloch sphere. This quotient is the complex Hopf fibration
\autocite{Hopf1931, Urbantke2003, mosseri2001hopf}
\begin{equation}
  S^{1}\hookrightarrow S^{3}\xrightarrow{\;h\;}S^{2},
\end{equation}
whose projection sends an amplitude to its Bloch vector,
\begin{equation}
  h(\psi)=\bigl(\langle\psi\rvert\sigma_{x}\lvert\psi\rangle,\;
                 \langle\psi\rvert\sigma_{y}\lvert\psi\rangle,\;
                 \langle\psi\rvert\sigma_{z}\lvert\psi\rangle\bigr)
        =\bigl(2\operatorname{Re}\bar z_{1}z_{2},\;
               2\operatorname{Im}\bar z_{1}z_{2},\;
               \lvert z_{1}\rvert^{2}-\lvert z_{2}\rvert^{2}\bigr).
\end{equation}
The Hopf map is quadratic in the amplitudes, and this is the geometric face of the Born rule:
the third component is $p_{0}-p_{1}$ with $p_{k}=\lvert z_{k}\rvert^{2}$, and the squaring is
the lowest-degree $U(1)$-invariant available on the fibre --- the record-theoretic derivation of
the same rule is \S\ref{sec:born-rule}'s. Because the metric is round, the projection is a
Riemannian submersion with totally geodesic fibres of constant length
\autocite{GluckWarnerZiller1986}: the same amount of phase is hidden at every point of the base,
so the phase is hidden \emph{uniformly}, not in a position-dependent way that the metric could
partially betray.

\subsection{The fibre as self-referential ignorance}
\label{sec:fibre-self-ignorance}

The interpretation the framework places on this geometry --- and here we are reading physics
into the mathematics, so the claim is interpretive rather than proved --- is that the $S^{1}$
fibre is the observer's own un-narratable self-ignorance. The global phase is the imaginary part
of the complex information measure $S=-\ln\lvert p\rvert-\mathrm{i}\theta$: information the
observer lacks about their own state. Lacking it, they cannot communicate it; there is no
\textsc{fanout} of information one does not possess. Uncommunicable, it cannot enter any
measurement record --- and this last step is no longer interpretive:
Proposition~\ref{prop:f1-dictionary} proves the global phase absent from every classical record
a copying structure induces --- so every observable is $U(1)$-invariant; and a $U(1)$-invariant
quantity is a function on the quotient. The physical state is therefore the orbit, not the
representative point---the base $S^{2}$---and the $S^{1}$ fibre carries exactly the inaccessible
self-referential phase. It is a null direction of statistical distinguishability: two global
phases are indistinguishable by any measurement whatever \autocite{wooters_1981}, so the fibre
is a gauge direction without physical content. This phase is irrecoverably discarded, and must
remain so: it is the relative phase, living in the $S^{3}$ fibre of the quaternionic Hopf
fibration at two qubits, that is genuinely recoverable, whereas the single-qubit global phase of
the complex fibration is gauge-private without remainder. The self-reference that opened the
derivation thus reappears, at its end, as the circle the observer can never see around.

%      % EDITED 2026-09-12 scripts/patch_revert_sedenion_repoint_2026-09-12T1600.py
% EDITED 2026-09-11 scripts/patch_undefined_section_refs_2026-09-11T2030.py
% =====================================================================
%  SECTION: Wigner's friend in the Hopf picture (v2; supersedes the
%  2026-06-07 draft). AUDITED 2026-07-06 (final pass, fifth stop).
%
%  PLACEMENT (reordered chain, JS 2026-07-06): position 4 --
%      \input{L2_metric_global_phase_section_2026_06_07}
%      \input{wigners_friend_in_the_hopf_picture_2026-06-13}
%      \input{derivation_of_the_quantum_postulates}
%  sec:l2-metric-fibre is REAL (backward, order-consistent);
%  sec:fanout-boundary DRAFTED 2026-07-06 (formerly the one undrafted section);
%  the old sec:monogamy placeholder note was VOID -- no such \ref in
%  the body -- and is removed. prop:f1-dictionary and
%  sec:hopf-justification are FORWARD references (position 8): first
%  F1 use now carries the deliberate-forward flag.
%
%  DEPENDENCY: \input{monogamy_aside_2026-06-07} -- received from JS
%  and audited 2026-07-06; the \input sits inside
%  sec:wf-reconciliation, immediately after op:tsirelson-geometry
%  (relocated 2026-07-06 from end-of-file, where the aside's opening
%  antecedent had been orphaned). A bare replace-all during the earlier
%  patch briefly mangled this header block, leaving a live \input in
%  it; repaired by index-splice, and the lesson recorded: anchors that
%  also occur in comments must be uniqueness-asserted.
%
%  AUDIT RECORD: CHSH arithmetic (2sqrt(1+C^2), PR => r^2 = -2) and the
%  C^2+r^2=1 chain verified; F1 wording canonical; the
%  obsolescence-made-cancellation echo two-way with
%  sec:self-reference; redundancy/Darwinism/DiBiagio-Rovelli/Brukner
%  positioning sound; Landauer placed on amplification, correctly.
%  Fixes: first-person voice removed; \textquote replaced with
%  portable quotes; IOEP given its prop ref; stale Tsirelson
%  mini-header (wrong key name Tsirelson1980) retired;
%  op:tsirelson-geometry MINTED from the honesty paragraph's unpriced
%  owed. Popper page-pin DISCHARGED 2026-07-07: the zip bundles carry
%  OCR .txt sidecars (the earlier "no text layer" verdict was wrong --
%  ZIP-of-JPEGs defeated the PDF tools, not the text); the quoted
%  phrase sits at Part II, p. 187, now pinned in the cite. Bonus from
%  the same dig: "C+" is Popper's own notation (Part II, scan pp.
%  15-17), so the L2 opening's attribution stands as written; "E" not
%  sighted in the sampled pages -- the inline gloss covers it either
%  way. Keys unique to this file (presumed in the .bib,
%  since the paper compiles): VerstraeteWolf2002,
%  Cirelson1980QuantumInequality, PopescuRohrlich1994, Held2021,
%  Brukner2021, zurek_2009, Landauer_1961, wigner_1961.
%  British English; \autocite; no bullets in prose.
% =====================================================================

\section{Wigner's friend in the Hopf picture: one state, two charts, and the redundancy ladder}
\label{sec:wigners-friend}

The framework dissolves the Wigner's-friend paradox \autocite{wigner_1961} not by
adjudicating between the two accounts but by recognising that they are written in
two charts on a single manifold, selected as the problem requires, much as one
moves between cylindrical and spherical coordinates without either being the true
one. The epistemological chart (private, non-clonable, phase-carrying) and the
ontological chart (clonable, public, record-bearing) are related by the
\textsc{fanout} boundary (Section~\ref{sec:fanout-boundary}), and that boundary
is not a sharp wall but a graded one, parametrised by how redundantly a piece of
relational content has been copied. Read this way, the fact that a \textsc{fanout}
has occurred for the friend but not for Wigner is not a contradiction to be
removed; it is the boundary behaving exactly as defined.

Let the system be a qubit $\lvert\psi\rangle_{S}=\alpha\lvert 0\rangle
+\beta\lvert 1\rangle$ and let the friend $F$, idealised, with deliberate
economy, as a single qubit, measure it in the computational basis. The
interaction is the controlled coupling that copies $S$'s value into $F$'s
register: the \textsc{copy} comultiplication of the diagrammatic calculus
\autocite{CoeckeKissinger2017}, and in the circuit picture the very gate that, in
the two-qubit Mosseri--Dandoloff geometry \autocite{mosseri2001hopf},
activates the quaternionic Hopf fibration $S^{3}\to S^{7}\to S^{4}$ and carries
the joint state from the surface of the entanglement ball to its interior. The
friend and Wigner read the two ends of that plunge.

Two idealisations in this set-up should be named at once, because the argument
later depends on seeing past them. The first is the friend as a single qubit,
which we return to as the redundancy-one endpoint of a graded boundary
(Section~\ref{sec:wf-redundancy}). The second, less obvious and more consequential,
is the system $S$ as a pristine qubit with no prior entanglement --- a relation
with no history. No real measured particle is of this kind. A particle is itself a
reified relation between prior events, and by the time the friend couples to it, it
already carries the completed internal structure such a relation possesses: in the
framework's terms, the full octonionic $A,B,C$ register that the $S^{7}$ sector
supplies as internal degrees of freedom (\S\ref{sec:space-of-computations}). The
pristine qubit is a schematic convenience, drawn structureless so that the logic of
a single link is visible; the ``system, friend, Wigner'' triple is likewise the
schematic of \emph{one link} in the chain of observers, not its base. The chain
does not bottom out in a structureless qubit --- there is no such thing --- and we
draw on this when we ask, later, what bounds the information a single such link can
carry.

\subsection{The friend's chart: a \textsc{fanout} of redundancy one}
\label{sec:wf-friend}

Relative to $F$, the coupling is a \textsc{fanout}: $F$ acquires a fact. What is
copied is gauge-invariant, the relative phase, the correlation between $F$ and
$S$, which is a genuine observable. The global phase is never the thing copied:
there is nothing observable there to duplicate, so the $S^{1}$ fibre of the
complex Hopf fibration $S^{1}\to S^{3}\to S^{2}$ of
Section~\ref{sec:l2-metric-fibre} remains $F$'s gauge-private self-ignorance,
untouched. In Hopf terms the division is clean, but it lives at two levels, which
Proposition~\ref{prop:f1-dictionary} (\S\ref{sec:hopf-justification} --- a deliberate forward
reference) separates. The copy operation itself is fibre-committed:
copying is faithful to phase, so what $F$'s \textsc{fanout} duplicates is a lift --- a chosen
fibre-point over the outcome --- never a bare point of the base. Yet everything the operation can
broadcast is base data alone: the record it induces is independent of the lift, so the $S^{1}$
freedom, though pinned in the copier, is absent from every classical record it yields. $F$'s classical record is a
base-point; the structure that produced it committed to a fibre-point; and the gap between the two
is precisely the gauge character of the fibre. The observation has a fifty-year-old
anchor: Popper argued that a predictor's inferior self-knowledge follows from the logical
impossibility of a system's imparting information about itself without thereby rendering it out of
date --- ``making this information obsolete'', in his words
\autocites{Popper1950a}[Part~II, p.~187]{Popper1950b}. Proposition~\ref{prop:f1-dictionary}
makes the impossibility structural rather than temporal, and localises it: the datum the copier
cannot put into its own record is not its state wholesale but exactly its fibre commitment, which
does not arrive stale --- it cancels. Conditioned on the outcome it reads, $S$ is for
$F$ a definite point on the Bloch sphere,
\begin{equation}
  \rho_{S}^{(F)}=\lvert 0\rangle\langle 0\rvert
  \quad\text{(given outcome }0\text{)},
\end{equation}
a pure state of Bloch radius $r=1$ on the surface of the ball, concurrence $C=0$.

This is a \textsc{fanout}, but the minimal one. The fact $F$ holds is a
\emph{relative} fact in the sense of relational quantum mechanics
\autocite{Rovelli1996a}: $S$ definite relative to $F$, not an observer-independent
record. It is one copy, redundancy $R=1$, and as such it remains reversible by an
agent who controls the joint system. The single-qubit friend is therefore not
outside \textsc{fanout}, as a careless reading might suggest; it is the $R=1$
endpoint of the \textsc{fanout} continuum---a \textsc{fanout} that has happened,
and a fact that is real for $F$, while being maximally still reversible from
outside.

\subsection{Wigner's chart: no copy has reached him}
\label{sec:wf-wigner}

Wigner treats the laboratory as closed and no copy of the relational content has
reached him. For him the coupling is the unitary
\begin{equation}
  \lvert\Psi\rangle
  =\alpha\lvert 0\rangle\lvert F_{0}\rangle
  +\beta\lvert 1\rangle\lvert F_{1}\rangle,
\end{equation}
the move up to the larger Hilbert space, $\mathbb{C}^{2}\to\mathbb{C}^{4}$,
$S^{3}\to S^{7}$, complex Hopf to quaternionic Hopf, with the fibre grown from
$S^{1}$ to $S^{3}$. What $F$ underwent as the acquisition of a fact, Wigner
describes as correlation deposited into a joint pure state on $S^{7}$. When he
asks specifically about $S$, he traces $F$ out, and the geometry does the
bookkeeping:
\begin{equation}
  \rho_{S}^{(W)}=\operatorname{Tr}_{F}\lvert\Psi\rangle\langle\Psi\rvert
  =\lvert\alpha\rvert^{2}\lvert 0\rangle\langle 0\rvert
  +\lvert\beta\rvert^{2}\lvert 1\rangle\langle 1\rvert,
  \qquad
  r=\bigl\lvert\,\lvert\alpha\rvert^{2}-\lvert\beta\rvert^{2}\bigr\rvert,
  \quad
  C=2\lvert\alpha\beta\rvert .
\end{equation}
Tracing out discards the $S^{3}$ fibre and carries the pure point on $S^{7}$ to a
mixed point in the interior of the Bloch ball at radius $r=\sqrt{1-C^{2}}$
\autocite{mosseri2001hopf}. The inward displacement is governed by the
concurrence: it is the relational content that has migrated into the $S$--$F$
correlation, stored in the recoverable $S^{3}$ fibre. Nothing has been destroyed;
it has been relocated into the fibre, which is why Wigner can in principle
interfere the branches and reverse the interaction, the content sits at low
redundancy and no further copy has fixed it.

\subsection{One state, two charts: \texorpdfstring{$C^{2}+r^{2}=1$}{C2+r2=1}}
\label{sec:wf-reconciliation}

The system is pure for the friend ($r=1$, on the surface) and mixed for Wigner
($r<1$, in the interior). These are not rival claims about a single fact; they are
the friend at the surface and Wigner at the interior of the \emph{same} state on
$S^{7}$, bound by the Mosseri--Dandoloff constraint
\begin{equation}
  C^{2}+r^{2}=1,
\end{equation}
in which the discrepancy in Bloch radius is exactly the concurrence Wigner
retains and the friend has resolved away. This is the Inside/Outside Equivalence
Principle (Proposition~\ref{prop:inside-outside}) made concrete, the same law written by the embedded participant and by
the external observer, and the relativity of facts of relational quantum
mechanics \autocite{Rovelli1996a} realised in geometry rather than asserted: the
\textsc{fanout} boundary lies \emph{between} $F$ and $W$, so a fact exists
relative to $F$ and not relative to $W$. Neither chart is privileged; one reads in
the epistemological or the ontological chart as the question requires, and their
mutual incompleteness is a theorem of the construction, not a defect of either.

\medskip
\noindent\emph{Aside: the Tsirelson bound and the absence of PR boxes.}
\label{sec:tsirelson}
A ball has a unique centre and no point beyond it. Maximal entanglement sits at
that centre \autocite{mosseri2001hopf}, $r=0$, $C=1$, and because the
Bloch radius cannot be negative, the concurrence cannot exceed unity. For a
two-qubit pure state the maximal CHSH correlator is
\begin{equation}
  \langle\mathrm{CHSH}\rangle_{\max}=2\sqrt{1+C^{2}}=2\sqrt{2-r^{2}}  
\end{equation}
\autocite{VerstraeteWolf2002},
which the centre maximises to $2\sqrt{2}$, the Tsirelson bound
\autocite{Cirelson1980QuantumInequality}. A Popescu--Rohrlich box \autocite{PopescuRohrlich1994},
with CHSH value $4$, would require $C^{2}=3$ and hence $r^{2}=-2$: a point two
units past the centre, which the ball does not contain. The super-quantum
correlation is excluded by the very positivity that places maximal entanglement
at the centre, one cannot trace out to a more-than-maximally-mixed reduced
state. The Tsirelson bound is in this sense the image, under the CHSH map, of the
boundedness of the $L^{2}$ state ball derived in
Section~\ref{sec:l2-metric-fibre}, and the prohibition on PR boxes shares its
origin with the roundness of the sphere: both trace to the reversible-continuity
metric, not to a separate principle.

\noindent This is the two-qubit, pure-state face of the bound, and its status is
worth stating plainly. It exhibits PR correlations as having no representing state
in the Hilbert geometry, and it recovers the value $2\sqrt{2}$ at the centre; but
forcing \emph{exactly} $2\sqrt{2}$ as a theorem of the geometry, rather than
reading it from the known concurrence--CHSH relation, and lifting the argument
from two qubits to the general Tsirelson bound, remain to be made rigorous. The
closest cognate in the literature is Held's derivation of the same bound from the
geometric algebra $G^{3}\cong\mathbb{H}$ \autocite{Held2021}, the same
quaternionic structure in a different language.

\begin{openproblem}[The Tsirelson bound as geometry]\label{op:tsirelson-geometry}
Derive $2\sqrt{2}$ as a theorem of the round state geometry rather than reading it off the
concurrence--CHSH relation, and lift the argument from the two-qubit pure case to the general
Tsirelson bound --- so that the exclusion of super-quantum correlation becomes, like the exclusion
exhibited above, a consequence of the $L^{2}$ ball's positivity (\S\ref{sec:l2-metric-fibre}) and
of nothing else. Held's quaternionic derivation \autocite{Held2021} is the nearest cognate and the
natural first tool.
\end{openproblem}

\input{monogamy_aside_2026-06-07}

\subsection{The redundancy ladder: from relative fact to objective fact}
\label{sec:wf-redundancy}

The friend's fact becomes a fact for Wigner only when Wigner copies it when he
measures $F$. It becomes a fact for a third party only when that party copies it
from Wigner, and so on up the von Neumann chain. The point at which the content
stops being negotiable is therefore not a point but a limit: once it has been
copied into enough mutually independent fragments that many parties can read it
without disturbing it, no one can practically un-copy it, and it is operationally
an objective fact. This is exactly quantum Darwinism \autocite{zurek_2009}:
objectivity is redundancy. It is not an external graft, because the framework's
account of decoherence already \emph{is} this picture---many fragments each take
their own copy of the clonable content, while their private phases, each
fragment's own self-ignorance, fail to add coherently and wash out.

This supplies the grading parameter that ``graded \textsc{fanout}'' had left
unnamed. \textsc{fanout} is graded not because the operation is fuzzy but because
redundancy $R$ runs continuously: $R=1$ is the single relational copy---the
friend's relative fact, reversible by Wigner---and $R\gg 1$ is the objective
fact, effectively irreversible because un-copying it would mean recovering every
fragment. The single-qubit friend is the $R=1$ endpoint; the macroscopic friend
is the same operation at large $R$. One axis, two ends. The relation to
Di~Biagio and Rovelli's distinction \autocite{DiBiagioRovelli2021} is then exact:
their \emph{stable} facts are the high-$R$ limit and their \emph{relative} facts
the low-$R$ regime, and the standard conditional-probability law they use to
define stability is the operational signature of redundancy having grown large.
Brukner's no-go theorem \autocite{Brukner2021}, that qubits are not observers,
sits on the same axis without conflict: it bites only against calling an $R=1$
relative fact an objective one, which the framework never does---the redundancy
ladder is precisely the structure that separates the relational definiteness a
single qubit does carry from the objectivity it does not.

The thermodynamics tracks the same axis. The single relational copy is unitary and
free; the $kT\ln 2$ cost \autocite{Landauer_1961} is incurred not by the friend's
act but by the amplification and stabilisation that drive $R$ upward, when
fragments are irreversibly committed. The relative fact is thermodynamically free;
the objective fact is what costs $kT\ln 2$.

\subsection{Phase promotion, and the bump in the carpet}
\label{sec:wf-phase-promotion}

One point keeps the load-bearing irrecoverability intact while explaining what
Wigner can do that the friend cannot. What Wigner recovers is \emph{not} the
$S^{1}$ self-ignorance; the global phase remains gauge for every observer and that
circle is spent for good. The quantity the friend resolves away is the relative
phase between the $\lvert 0\rangle$ and $\lvert 1\rangle$ components, and moving up
to $\mathbb{C}^{4}$ \emph{promotes} it: the single-qubit relative phase becomes the
relative phase between $\lvert 0\rangle\lvert F_{0}\rangle$ and
$\lvert 1\rangle\lvert F_{1}\rangle$, which lives in the recoverable $S^{3}$ fibre
(itself a circle bundle over $S^{2}$, so the promotion adds a two-sphere of
relative-phase directions to the lone gauge circle). Gauge phase at one level
becomes physical, interferable phase one level up---which is why Wigner can run
the interference experiment the friend cannot. But Wigner now carries his own
fresh global-phase $S^{1}$ at the top of the joint state: his own irreducible
self-ignorance. Ascending the tower does not eliminate the self-reference, only
displaces it one level higher, and the regress up the chain is precisely the
redundancy ladder of the previous subsection. This is the framework's bump in the
carpet: each observer sits atop a Hopf tower with a single gauge circle he can
never see around, and a measurement relocates the next observer's blind circle
into his own recoverable fibre without ever closing the regress.

Two cautions close the section. First, the recoverability Wigner
enjoys is the low-redundancy limit: for a macroscopic, high-$R$ friend the
reversal is exponentially hard and the fact is objective for all practical
purposes, so Wigner's ``no record'' is the $R\!\approx\!1$ idealisation of a
graded boundary, not a free lunch. Second, naming redundancy as the grading
parameter of \textsc{fanout}, and the phase-promotion reading itself, are
syntheses of the framework's parts: each rests on established pieces (the
clonability criterion, quantum Darwinism, the tensor-product promotion of phase)
but is assembled here, and is offered as the cleanest statement of why the two
observers never disagree about anything both can check.
% =====================================================================
%  SPLICE: "Where the chain goes next" -- forward-looking coda for
%  sec:wigners-friend, pointing to sec:sedenions-generations.
%  Created 2026-08-27 per JS.
%
%  THE POINT (JS, 2026-08-27). The Wigner's-friend charts stop at
%  Wigner because the Hopf tower runs out at the octonions, while the
%  text says the chain of observers obviously continues. The sedenions
%  show what the continuation adds: the observer of Wigner (fourth
%  wire, C (x) S) brings a NEW piece of structure -- the finite S_3 of
%  Brown's theorem, the generation label -- and the observer of the
%  observer of Wigner (fifth wire, the trigintaduonions) brings nothing
%  new in kind. This coda records that, so the section's termination at
%  Wigner reads as a fact about the tower rather than an economy of the
%  exposition.
%
%  WHAT THE ALGEBRA ACTUALLY SAYS (verified 2026-08-27, notebook T --
%  notebook_T_trigintaduonion_automorphisms_2026-08-27.py):
%    - Aut(S) = G_2 x S_3 (Brown 1967), DIRECT product; the new S_3 is
%      generated by the doubling involution and Brown's explicit
%      order-three map psi, which fixes 1 and the doubling unit and
%      rotates every (a, a u) plane, a imaginary, by a third of a turn.
%    - psi permutes the three octonion HALVINGS O, psi O, psi^2 O of S
%      (each doubled by the same unit; meeting only in R). It does NOT
%      permute the three copies through the quaternion core -- those are
%      cycled by a lifted G_2 element (the colour-line three-cycle); see
%      the sedenion section's two-triples remark and the corrections
%      ledger of the same date.
%    - Brown: Aut(A_{n-1}) x S_3 = Aut(A_n) with EQUALITY for n = 4, 5, 6
%      (as quoted in gresnigtgourlayvarma2023three, sec. 4.3); Eakin and
%      Sathaye prove Brown's conjecture in general (VERIFY-CITE for the
%      exact general form -- their step is Aut(A_{n+1}) = Aut(A_n) x G,
%      G in {Z/2, S_3}; over R with the standard parameters the branch is
%      S_3 at every rung >= 4, which is what notebook T checks directly
%      at rungs 4, 5, 6). Derivation algebra g_2 from the octonions on
%      (Schafer 1954; notebook T: dim Der = 14 at n = 3, 4, 5).
%  So "nothing new at the fifth wire" is stated PRECISELY: the fifth
%  doubling adds another finite S_3 and nothing else -- no connected
%  symmetry, no fibration; the increment is the same at every rung
%  thereafter. Not "silent": REPETITIVE. The prose says so.
%
%  STATUS: the algebra is ESTABLISHED (Brown; Eakin--Sathaye; Schafer)
%  and directly checked (notebook T). The reading of the repeated finite
%  factor as what the redundancy ladder consumes is a RHYME, flagged.
%  The reading of the fourth rung's discrete orientation as the
%  generation label is the sedenion section's CANDIDATE identification,
%  pointed to, not re-argued here.
%
%  PLACEMENT: a new \subsection at the END of sec:wigners-friend, after
%  the "Two cautions" paragraph. Because that paragraph reads "Two
%  cautions close the section.", the patch also changes that one word:
%  "close the section" -> "close the argument". Both anchors
%  uniqueness-asserted (the July lesson: anchors that also occur in
%  comments must be asserted unique).
%
%  CROSS-REFS OUT (all exist in the synced source except the two
%  MAP-LABELs):
%    sec:hopf-justification   -- the "two climbs" paragraph (length vs
%                                termination); this coda sharpens it
%    sec:wf-redundancy         -- the redundancy ladder (this file)
%    sec:wf-phase-promotion    -- phase promotion (this file)
%    sec:sedenions-generations -- the sedenion section (2026-08-27 rev.)
%    prop:termination          -- Hurwitz/Adams termination
%  MAP-LABELs:
%    app:notebook-record       -- the notebook-record appendix
%                                (appendices/appendix_notebook_record_
%                                2026-07-11.tex); substitute its label.
%  NEW KEYS (RIS delivered: cayley_dickson_automorphism_refs_2026-08-27.ris):
%    eakinsathaye1990automorphisms, schafer1954cayleydickson
%  IN-LIBRARY: brown1967generalized.
%
%  OPTIONAL CONSUMER-POINTER (not patched; one clause for the "two
%  climbs" paragraph of hopf_justification_section_2026-06-08.tex, after
%  "that chain has no top --- it is the bump in the carpet, unbounded in
%  length."):
%      "(unbounded in length, and --- as \S\ref{sec:wf-chain-type}
%      records --- bounded in type: after the fourth observer nothing
%      new in kind arrives)"
%
%  EDITED 2026-08-30 (per JS): one sentence added to the fourth-wire
%  paragraph anchoring the flavour sector's irremovable phase (the
%  CP-phase recognition) on this rung, with the development at
%  sec:sedenions-mixing (2026-08-30 revision). No new keys here.
%  British English; pure ASCII; \autocite; no bullets; spaced dashes.
% =====================================================================
% SPLICE-BEGIN wf-chain-type 2026-08-27

\subsection{Where the chain goes next: one more rung, then copies}
\label{sec:wf-chain-type}

The regress just described is unbounded in length, and nothing here shortens
it: the observer of Wigner carries a gauge circle of his own, and so does the
observer of that observer, as \S\ref{sec:hopf-justification} insists when it
separates the endlessness of the chain from the termination of the tower. But
the algebra says how much of the endless chain adds structure of a new kind,
and the answer is exact: one more step, and then none.

The count runs as follows. The triple system, friend, Wigner is three wires,
and three wires are the three Hopf rungs $\mathbb{C}$, $\mathbb{H}$,
$\mathbb{O}$ --- the completed link of the earlier caveat --- which is why the
charts above stop at Wigner: the chart of whoever looks in on the triple,
$S^{15}\to S^{8}$, is the last fibration the division algebras supply
(Proposition~\ref{prop:termination}). The observer of Wigner is a fourth wire,
and the fourth wire is the sedenions,
$(\mathbb{C}^{2})^{\otimes 4} = \mathbb{C}\otimes\mathbb{S}$
(\S\ref{sec:sedenions-generations}). No fibration comes with it --- the
sedenions have zero divisors and admit no map $S^{31}\to S^{16}$ --- so the
fourth observer gains no recoverable fibre of the tower's kind: the continuous
promotion of relative phase into a larger fibre, the mechanism of the previous
subsection, has no fourth instance. What the doubling adds instead is finite.
Brown's theorem gives $\mathrm{Aut}(\mathbb{S}) = G_{2}\times S_{3}$
\autocite{brown1967generalized}: the identity component is unchanged from the
octonions, and the entire gain is a component group, generated by an
order-three automorphism that fixes the new doubling direction and rotates
every imaginary direction of the completed link into its copy by a third of a
turn. The relative orientation of the link and its copy, which at the Hopf
rungs was a continuous phase promoted into a fibre, survives at the fourth rung
only as a discrete choice among three, and that choice is the generation label
(\S\ref{sec:sedenions-generations}). The fourth observer, then, sees no new
force and no new fibre; he sees the same link three ways, and cannot gauge the
difference, because a component group has no Lie algebra. His rung is also where
the flavour sector's one irremovable phase finds its anchor: three generations
are the threshold at which the mixing matrix first carries a phase no convention
removes, and \S\ref{sec:sedenions-mixing} reads that phase --- through the
paper's negativity--contextuality chain --- as this wire's contextual cargo,
real only from one rung above.

The observer of the observer of Wigner is a fifth wire, the trigintaduonions,
and here the tower has nothing new to say. Brown proved that the next two
doublings each add exactly one further $S_{3}$ and nothing else, Eakin and
Sathaye that Brown's pattern continues at every rung, and the derivation
algebra --- the Lie algebra of the identity component --- is $\mathfrak{g}_{2}$
from the octonions onward
\autocite{brown1967generalized,eakinsathaye1990automorphisms,schafer1954cayleydickson};
the fifth and sixth rungs are also checked directly in the notebook record
(Appendix~\ref{app:notebook-record}). So the fifth observer's doubling repeats
the fourth's gift --- another third-of-a-turn symmetry between the four-wire
structure and its copy, commuting with everything below --- and adds no
connected symmetry and no fibration. Nothing new in kind arrives at the fifth
wire or at any wire after it. The increments of the tower are, in order, a
force, a force, a label, and then the same finite permutation of copies at
every step. That is the precise sense in which the chain ``obviously
continues'': it continues, and after the fourth observer it has nothing further
to announce. The framework reads the repeated factor as what the redundancy
ladder of \S\ref{sec:wf-redundancy} consumes --- a permutation symmetry among
copies, which is the freedom a record has and a type does not --- and records
that reading as a structural resonance rather than a derivation.

The two climbs of \S\ref{sec:hopf-justification} are thereby sharpened. The
chain of observers is unbounded in length and bounded in type: three forces by
the third wire, one family label by the fourth, copies thereafter. The observer
of Wigner is the last observer whose arrival changes what there is.

% SPLICE-END wf-chain-type 2026-08-27

% =====================================================================
%  FIGURE: the Wigner's-friend chain mapped onto the Cayley--Dickson
%  ladder. Created 2026-08-30 per JS: the Hopf-fibration view of
%  Wigner's friend has been load-bearing for months without being drawn;
%  the fibrations themselves are DEFERRED to Szangolies's introduction
%  (szangolies2025standardmodel, in library), and this figure carries
%  only the mapping.
%
%  CONTENT (left to right, one panel per chart, staircase rising):
%    Panel 1: F's chart of S            -- S^3 -> S^2, fibre S^1,  C, one wire
%    Panel 2: W's chart of {S,F}        -- S^7 -> S^4, fibre S^3,  H, two wires
%    Panel 3: O_W's chart of {S,F,W}    -- S^15 -> S^8, fibre S^7, O, three wires
%             (the last fibration)
%    Panel 4: the next observer's chart of {S,F,W,O_W} -- C^16 = C (x) S,
%             no fibration (S^31 -/-> S^16); the doubling's S_3 badge:
%             the generation label, a third of a turn.
%  Each describer perches on his sealed box with a dashed circle over
%  his head (his own gauge circle -- the self-referential ignorance he
%  cannot see around); dashed promotion arrows carry each circle into
%  the next chart's fibre ("gauge here -> interferable one rung up",
%  sec:wf-phase-promotion); the third arrow terminates at "no fourth
%  fibration", with what survives being the discrete third-of-a-turn
%  (sec:wf-chain-type, sec:sedenions-generations). The regress
%  continues to the right: copies only.
%
%  PREAMBLE REQUIREMENT (confirm in main.tex; add if absent):
%      \usepackage{tikz}
%      \usetikzlibrary{arrows.meta}
%  PLACEMENT: \input by patch_wigner_hopf_ladder_fig_2026-08-30.py,
%  after the two-idealisations paragraph of sec:wigners-friend, with
%  the accompanying deferral paragraph (in the patch).
%  CROSS-REFS OUT (all real): sec:wf-phase-promotion,
%  sec:wf-chain-type, sec:sedenions-generations.
%  KEY: szangolies2025standardmodel (in library).
%  British English; pure ASCII; \autocite in the caption.
% =====================================================================
\begin{figure}[t]
\centering
\resizebox{\textwidth}{!}{%
\begin{tikzpicture}[>=Stealth, line cap=round, line join=round,
    every node/.style={font=\scriptsize},
    blind/.style={densely dashed, draw=black!75},
    promote/.style={->, densely dashed, black!60, semithick},
    sight/.style={->, dotted, black!55}]
% ---- stick-figure pic --------------------------------------------------
\tikzset{stick/.pic={
  \draw[pic actions] (0,0.44) circle (0.085);
  \draw[pic actions] (0,0.355) -- (0,0.10);
  \draw[pic actions] (-0.12,0.20) -- (0,0.29) -- (0.12,0.20);
  \draw[pic actions] (-0.10,-0.10) -- (0,0.10) -- (0.10,-0.10);
}}
% ---- the staircase (Cayley--Dickson ladder) ---------------------------
\draw[black!35, thick]
  (-0.35,-0.55) -- (3.27,-0.55) -- (3.27,0.25) -- (6.89,0.25)
  -- (6.89,1.05) -- (10.51,1.05) -- (10.51,1.85) -- (13.60,1.85)
  -- (13.60,2.65) -- (16.20,2.65);
\node[anchor=north] at (1.25,-0.60) {$\mathbb{C}$ --- one wire};
\node[anchor=north] at (4.87,0.20)  {$\mathbb{H}$ --- two wires};
\node[anchor=north] at (8.49,1.00)  {$\mathbb{O}$ --- three wires};
\node[anchor=north] at (12.11,1.80) {$\mathbb{S}$ --- four wires};
\node[anchor=north, black!55] at (14.45,2.60) {$\mathbb{T},\dots$ --- five wires on};
% =======================================================================
% Panel 1: F's chart of S
% =======================================================================
\draw[rounded corners=2pt, fill=black!5, draw=black!40] (0,0) rectangle (2.5,1.32);
\filldraw[black!70] (1.25,0.72) circle (0.075);
\node[font=\tiny] at (1.25,0.45) {$S$};
\pic[black] at (2.05,1.42) {stick};
\node[anchor=west] at (2.22,1.62) {$F$};
\draw[sight] (1.95,1.60) -- (1.45,1.02);
\draw[blind] (2.05,2.14) circle (0.11);
\node[anchor=east, align=right, font=\tiny, black!75] at (1.82,2.14)
  {his own gauge circle:\\ the phase he cannot\\ see around};
\node[anchor=south west, font=\tiny, black!75] at (0.02,1.34) {fibre $S^{1}$};
\node[anchor=north] at (1.25,-0.06) {$S^{3}\to S^{2}$};
% =======================================================================
% Panel 2: W's chart of {S,F}
% =======================================================================
\draw[rounded corners=2pt, fill=black!5, draw=black!40] (3.62,0.80) rectangle (6.12,2.12);
\filldraw[black!70] (4.07,1.52) circle (0.075);
\node[font=\tiny] at (4.07,1.25) {$S$};
\pic[black!55, scale=0.85] at (4.85,1.16) {stick};
\node[font=\tiny, black!55] at (4.85,0.94) {$F$};
\pic[black] at (5.67,2.22) {stick};
\node[anchor=west] at (5.84,2.42) {$W$};
\draw[sight] (5.57,2.40) -- (5.05,1.82);
\draw[blind] (5.67,2.94) circle (0.11);
\node[anchor=south west, font=\tiny, black!75] at (3.64,2.14) {fibre $S^{3}$};
\node[anchor=north] at (4.87,0.74) {$S^{7}\to S^{4}$};
\draw[promote] (2.17,2.14) to[out=15,in=170] (3.60,2.24);
\node[anchor=south, align=center, font=\tiny, black!60] at (2.85,2.32)
  {gauge below,\\ interferable above};
% =======================================================================
% Panel 3: O_W's chart of {S,F,W}
% =======================================================================
\draw[rounded corners=2pt, fill=black!5, draw=black!40] (7.24,1.60) rectangle (9.74,2.92);
\filldraw[black!70] (7.69,2.32) circle (0.075);
\node[font=\tiny] at (7.69,2.05) {$S$};
\pic[black!55, scale=0.85] at (8.35,1.96) {stick};
\node[font=\tiny, black!55] at (8.35,1.74) {$F$};
\pic[black!55, scale=0.85] at (8.95,1.96) {stick};
\node[font=\tiny, black!55] at (8.95,1.74) {$W$};
\pic[black] at (9.29,3.02) {stick};
\node[anchor=west] at (9.46,3.22) {$O_{W}$};
\draw[sight] (9.19,3.20) -- (8.65,2.62);
\draw[blind] (9.29,3.74) circle (0.11);
\node[anchor=south west, font=\tiny, black!75] at (7.26,2.94) {fibre $S^{7}$};
\node[anchor=north, align=center] at (8.49,1.54) {$S^{15}\to S^{8}$\\[-1pt]
  {\tiny the last fibration}};
\draw[promote] (5.79,2.94) to[out=15,in=170] (7.22,3.04);
% =======================================================================
% Panel 4: the next chart of {S,F,W,O_W} -- sedenions, no fibration
% =======================================================================
\draw[rounded corners=2pt, fill=black!5, draw=black!40] (10.86,2.40) rectangle (13.36,3.72);
\filldraw[black!70] (11.24,3.12) circle (0.075);
\node[font=\tiny] at (11.24,2.85) {$S$};
\pic[black!55, scale=0.80] at (11.78,2.78) {stick};
\node[font=\tiny, black!55] at (11.78,2.56) {$F$};
\pic[black!55, scale=0.80] at (12.28,2.78) {stick};
\node[font=\tiny, black!55] at (12.28,2.56) {$W$};
\pic[black!55, scale=0.80] at (12.78,2.78) {stick};
\node[font=\tiny, black!55] at (12.78,2.56) {$O_{W}$};
\pic[black!45] at (14.35,2.78) {stick};
\draw[sight] (14.18,3.02) -- (13.42,3.12);
\draw[blind] (14.35,3.50) circle (0.11);
\node[anchor=west, align=left, font=\tiny, black!60] at (14.58,3.48)
  {the next observer,\\ and all after:\\ copies only};
\node[anchor=north] at (12.11,2.34) {$\mathbb{C}^{16}=\mathbb{C}\otimes\mathbb{S}$};
% the broken third promotion
\draw[promote] (9.41,3.74) to[out=15,in=185] (10.68,3.86);
\node[font=\scriptsize, black!60] at (10.80,3.92) {$\times$};
\node[anchor=south west, font=\tiny, black!75] at (10.88,3.74)
  {no fourth fibration: $S^{31}\not\to S^{16}$};
% the S3 badge: three dots at 120 degrees with a rotation arrow
\draw[black!75] (11.45,4.30) circle (0.17);
\filldraw[black!75] (11.45,4.47) circle (0.028);
\filldraw[black!75] (11.303,4.215) circle (0.028);
\filldraw[black!75] (11.597,4.215) circle (0.028);
\draw[->, black!75] (11.67,4.36) arc[start angle=20, end angle=100, radius=0.235];
\node[anchor=west, align=left, font=\tiny, black!75] at (11.70,4.36)
  {$S_{3}$: the generation label\\ (a third of a turn)};
\draw[promote, black!70] (10.86,3.96) to[out=60,in=200] (11.24,4.22);
\end{tikzpicture}%
}
\caption{The Wigner's-friend chain mapped onto the Cayley--Dickson ladder. Each
observer holds a chart of the sealed box below him, one rung further up the
ladder: the friend's chart of the system is the complex Hopf bundle, Wigner's
chart of the sealed laboratory the quaternionic, and the chart of the observer
of Wigner --- of the completed triple --- the octonionic, the tower's last
fibration. Every chart-holder carries a dashed circle over his head, his own
gauge phase: the self-referential ignorance he cannot see around. Measurement
promotes each circle into the next chart's fibre, where what was gauge below is
interferable above (\S\ref{sec:wf-phase-promotion}); the third promotion has no
fibration to land in --- the sedenions have zero divisors and no map
$S^{31}\to S^{16}$ --- and what survives of it is discrete: the doubling's
$S_{3}$, a third of a turn rotating the link into its copy, the generation
label (\S\ref{sec:wf-chain-type}, \S\ref{sec:sedenions-generations}). The
regress of circles continues unbounded to the right, but from the fifth wire
onward it adds copies only. The friend drawn as a single qubit is the
redundancy-one idealisation discussed in the text; for the fibrations
themselves, deferred rather than repeated here, see the introduction in
\autocite{szangolies2025standardmodel}.}
\label{fig:wf-hopf-ladder}
\end{figure}

%      % EDITED 2026-10-29 scripts/patch_continuity_tower_premise_2026-10-29T1500.py
% EDITED 2026-09-20 scripts/patch_born_cost_ledger_v2_2026-09-20T0900.py
% EDITED 2026-09-15 scripts/patch_value_continuum_declared_cencov_2026-09-15T1540.py
% EDITED 2026-09-15 scripts/patch_continuity_connected_state_space_2026-09-15T1430.py
% =====================================================================
%  SECTION: Derivation of the quantum postulates (the assembly chapter).
%  AUDITED 2026-07-06 (final pass, sixth stop).
%
%  PLACEMENT (reordered chain, JS 2026-07-06): position 5 --
%      \input{wigners_friend_in_the_hopf_picture_2026-06-13}
%      \input{derivation_of_the_quantum_postulates}
%      \input{division_algebras_section_2026-07-03}
%  Postulate 6 hands off to division-algebras, its immediate successor.
%
%  LABEL RESOLUTIONS (2026-07-06): the old placeholder
%  sec:complex-amplitudes is now minted at its source (self-reference's
%  negativity subsection), so the ref here resolves unchanged.
%  sec:born-rule COLLIDED with the label minted 2026-07-06 on
%  self-reference's Born-from-fibre-privacy subsection; resolution: two
%  routes, two labels -- Postulate 3 here is renamed
%  sec:born-invariance (the U(1)/lowest-degree geometric route) and
%  cross-stitched to sec:born-rule (the record-theoretic route), their
%  agreement flagged as over-determination. sec:local-tomography also
%  minted at source for Postulate 6's operational clause.
%
%  AUDIT RECORD: the section's status discipline was PRESCIENT -- the
%  continuity arrow was already "owed ... enters as a principle" (the
%  exact voice JS mandated for L2), and both of its confessed debts
%  have since been minted as numbered problems; the closing now cites
%  them (op:limit-realisability, op:tsirelson-geometry): "named, and
%  numbered". Postulate 6 now carries BOTH composition duty-holders
%  (Hurwitz multiplicativity + local tomography), matching the
%  continuity section's caveat exactly. Stone's theorem cited
%  (stone1932one, RIS delivered); time-as-operation-count forward
%  flag added; \textquote made portable; opener formalised.
%
%  PROVENANCE NOTE (relocated from the body, 2026-07-06, JS's words):
%  "Using Hopf fibrations to get the Born rule is something Claude AI
%  came up with on its own without my prompting." Preserved here as the
%  historical record. DISCLOSURE STATUS (JS, 2026-07-06): AI use is
%  already acknowledged in the second paragraph of the paper's
%  introduction; the remaining recommendation is only that, at
%  submission, the same fact be echoed in whatever formal
%  Declarations / Acknowledgements slot the target journal's template
%  mandates -- several publishers now require it there specifically,
%  in addition to any in-text mention.
%
%  British English; \autocite; no bullets in prose.
% =====================================================================

% ADDED 2026-07-12: closing paragraph "The postulates in module
% coordinates" -- the Clifford-module reading (skeleton + two
% supplements), with the explanatory arrow guarded per JS: the field
% comes from the Shannon/additivity/commutativity step (self-ignorance),
% and the left-action composability result is recorded as that same
% requirement in operator clothing, not an independent selection of C.
\section{Derivation of the quantum postulates}
\label{sec:postulates}
Enough pieces of the jigsaw are now assembled to derive the postulates of quantum mechanics.

The Dirac--von Neumann postulates are ordinarily laid down as independent
assumptions. In this framework they are not assumptions but a single chain: each
follows from the Inside/Outside Equivalence Principle by one move from the
preceding step. We assemble the chain here and mark the status of each link as
proved, structural, or still owed, so that no step claims more than it has earned.

The chain begins below the postulates. An embedded observer, required by the
equivalence principle to describe a totality that includes itself, must assign a
probability to a proposition about its own future state. The Lawvere--Yanofsky
diagonal \autocite{Yanofsky2003} and Popper's predictor argument
\autocite{Popper1950a, Popper1950b} show that no consistent classical assignment
exists: the self-referential proposition is undecidable to the observer before
the fact. Representing the result of that undecidable self-prediction forces a
signed, non-classical probability, which carries the observer's irreducible
self-ignorance. From this single obstruction the postulates follow.

\paragraph{Postulate 1 (complex Hilbert space).} Two independent steps fix the state space. The
field is forced first: a signed probability, read through the Shannon information measure and
analytically continued, yields a complex amplitude $S=-\ln\lvert p\rvert-\mathrm{i}\theta$,
whose phase $\theta$ is the self-referential ignorance (Section~\ref{sec:complex-amplitudes}).
The metric is forced second, and separately: the principle of reversible continuity
(\S\ref{sec:continuity}) --- that a continuous reversible transformation connects any two pure
states --- selects the $L^{2}$ norm uniquely. For every $\ell^{p}$ with $p\neq 2$ the linear
isometries are generalised permutations \autocite{Banach1932, Lamperti1958}, a group that cannot
connect a basis state to a superposition; only at $p=2$ is the norm-preserving group connected
and transitive on the sphere \autocite{Aaronson2004h}, and Hardy's operational $K=N^{2}$
counting reaches the same quadratic norm by a second route from the same continuity premise
\autocite{Hardy2001a}, and \v{C}encov's theorem with the every-basis lemma reaches it from the
value continuum alone (\S\ref{sec:cencov-route}). The state manifold is therefore the unit
sphere of the Hermitian $2$-norm in $\mathbb{C}^{n}$, taken modulo the global phase
(Section~\ref{sec:l2-metric-fibre}). Placing the metric here, before the Born rule rather than
after it, is deliberate: the Born derivation presupposes the round sphere, so the metric must be
anchored independently to avoid circularity. Two bounds belong with this step. Reversible
continuity fixes the metric, not the field---real and quaternionic quantum theory satisfy it
equally---so the restriction to $\mathbb{C}$ rests on the Shannon step, not on continuity. And
deriving reversible continuity itself from self-reference is the one arrow the chain still owes
--- now priced exactly: one open premise, that every transformation of the regime is a phase of an inside structure,
Open Problem~\ref{op:limit-realisability}, with everything downstream of it proved
(\S\ref{sec:continuity}). Until it is closed, continuity enters as a principle, and the weight
of the $L^{2}$ conclusion is carried jointly with the isometry and counting arguments and with
the statistical route.

\paragraph{Postulate 2 (Hermitian observables).}
A measurement result must be communicable: it must be copyable from one observer
to another, the \textsc{copy} operation of the diagrammatic calculus
\autocite{CoeckeKissinger2017}. The imaginary part of the amplitude is the
private self-referential phase and cannot be broadcast; only the real,
gauge-invariant content can. Observables whose spectra are real and therefore
communicable are exactly the self-adjoint operators. Hermiticity is thus the
condition that a quantity be the kind of thing a measurement can report.

\paragraph{Postulate 3 (Born rule).}
\label{sec:born-invariance}
Outcome probabilities must be invariant under the global phase, the $U(1)$
self-ignorance, realised geometrically as the $S^{1}$ fibre of the complex Hopf
fibration $S^{1}\to S^{3}\to S^{2}$ over the Bloch sphere
\autocite{Hopf1931, Urbantke2003}. There is no polynomial $U(1)$-invariant of degree one,
and the non-polynomial $\lvert\psi\rvert$ is excluded by the analyticity declared at
\S\ref{sec:complex-amplitudes}; the lowest-degree polynomial invariant is quadratic, so the
probability must be $p=\lvert\psi\rvert^{2}$. The squaring is forced
rather than merely permitted, and the argument holds for a single qubit, where
Gleason's theorem is silent. The Born quadratic is not a new structure: it is the
same $L^{2}$ form that fixed the metric in Postulate~1, now read as the
normalisation of the round sphere, and the fibre over which the phase is averaged
has constant length, so the phase is hidden uniformly
\autocite{GluckWarnerZiller1986}. This step is a structural argument rather than a
theorem in the strict sense, and is flagged as such. The record-theoretic derivation of the same
quadratic --- squared magnitudes from fibre privacy, Proposition~\ref{prop:f1-dictionary} --- is
given at \S\ref{sec:born-rule}, and the cost-ledger derivation --- the record as the closed loop of
a history and its return leg, the exponent the number of legs --- at \S\ref{sec:born-cost-ledger}:
three routes to one rule, each holding for a single qubit, the over-determination this framework
treats as evidence.

\paragraph{Postulate 4 (state reduction).}
Measurement is a \textsc{fanout}: the gauge-invariant relational content is copied
into a record, converting a private epistemological amplitude into a public
fact. For the observer who performs it, acquiring that fact is Bayesian updating,
and the projection postulate is its expression. The reduction is perspectival---a
fact relative to the observer in the sense of relational quantum mechanics
\autocite{Rovelli1996a}---and it is graded by redundancy. A single copy is a
relative fact, still reversible by an agent controlling the joint system; redundant
proliferation across many independent fragments makes the fact objective, which is
quantum Darwinism \autocite{zurek_2009} and the stable/relative-fact distinction of
Di~Biagio and Rovelli \autocite{DiBiagioRovelli2021}. ``Collapse'' and unitary
evolution are then the two charts---inside and outside the \textsc{fanout}
boundary---of one process, as the Wigner's-friend analysis shows and \S\ref{sec:fanout-boundary} formalises
(Section~\ref{sec:wigners-friend}); the irreversibility and its $kT\ln 2$ cost
\autocite{Landauer_1961, wigner_1961} accrue with redundancy rather than at the
single act. The reduction postulate is therefore recovered as a chart-dependent
description, not a fundamental dynamical law; the identification of redundancy as
the grading parameter is a synthesis of the framework's parts and is flagged as
such.

\paragraph{Postulate 5 (unitary evolution).}
Conservation of probability is preservation of the $L^{2}$ inner product, and the
inner-product-preserving linear maps are the unitaries $U(n)$---the same
isometry fact that fixed the metric in Postulate~1. Continuity in time --- operation count,
\S\ref{sec:space-of-computations}, a forward reference --- together with Stone's theorem
\autocite{stone1932one} gives $U(t)=\mathrm{e}^{-\mathrm{i}Ht}$, the Schr\"odinger
equation. Reversibility is here a corollary, not a premise: bare reversibility on
the amplitude space yields only the general linear group, and it is preservation
of the quadratic form that carves $U(n)$ out of it.

\paragraph{Postulate 6 (composite systems).}
Systems combine by the tensor product because the division algebra norm is
multiplicative \autocite{hurwitz1898composition}: composition preserves the $L^{2}$ structure
only if states multiply. The operational half closes from the other side: the locality of
tomography derived at \S\ref{sec:local-tomography} says product states determine composites,
which is the tensor product's signature. Cayley--Dickson doubling carries $\mathbb{C}$ to the next
rung, and the ascent of the Hilbert space is the ascent of the Hopf tower, complex
to quaternionic, used in the Wigner's-friend construction
\autocite{mosseri2001hopf}; the ladder itself is the business of the section that follows
(\S\ref{sec:division-algebras}).

\paragraph{Consequence (the Tsirelson bound).}
The postulates fix not only the form of the theory but the strength of its
correlations. In the two-qubit ball the maximal entanglement sits at the centre,
and a Popescu--Rohrlich box \autocite{PopescuRohrlich1994} would require a point
beyond it; the CHSH correlator is capped at $2\sqrt{2}$
\autocite{Cirelson1980QuantumInequality, VerstraeteWolf2002, Held2021}, with the value recovered
at the centre and its general derivation still owed
(Section~\ref{sec:wf-reconciliation}).

\medskip
\noindent The result is that the six postulates are one chain, not six premises.
Of the links, the field and the Born quadratic rest on the framework's signed
probability and $U(1)$-invariance arguments; the $L^{2}$ metric, the unitary group
and the tensor product are over-determined by reversible continuity, the isometry
classification and Hurwitz composition; state reduction is recovered as the
chart-relative face of \textsc{fanout}. The two outstanding debts are the
derivation of reversible continuity from self-reference and the forcing of the Tsirelson value as
a theorem rather than a recovered number --- named, and numbered:
Open Problems~\ref{op:limit-realisability} and~\ref{op:tsirelson-geometry} --- so that a reader
can see exactly where the chain is complete and where it is not.

\paragraph{The postulates in module coordinates.}
One reading of the chain deserves a closing statement, because it positions the derivation
against a mathematical skeleton the reader may know well, and because stating it carelessly would
invert the chain's explanatory arrow. In the coordinates of the division-algebra programmes
(\S\ref{sec:algebraic-dictionary}), complex Hilbert-space quantum mechanics is real Clifford
module theory plus exactly two supplements. The skeleton is Clifford throughout: the operator
algebra of $n$ qubits is the Clifford algebra of $2n$ generators with the state space its own
spinor module, half-integer spin and fermionic statistics are the module's double cover and its
graded tensor, and the classically simulable stabiliser core is the algebra's combinatorics. What
the skeleton does not contain is, first, the choice of complex structure --- the polarisation $J$
that turns a real module into a complex Hilbert space, which standard quantum theory and the
algebraic programmes alike \emph{choose} --- and, second, the Born weights on the fibre. Those two
supplements are precisely what this section has derived, and derived from self-ignorance rather
than from algebra: the field from the Shannon step (information about one's own undecidable state
must still be information --- it must add, addition commutes, and so the amplitude algebra is
commutative, excluding the quaternionic bookkeeping an observer could never balance), and the
quadratic from fibre privacy. The algebraic echo --- that the algebra's own left actions compose
only along a commutative line, the composability result recorded with the dictionary --- is not
an independent selection of $\mathbb{C}$: composition of left actions fails off that line exactly
because the doubled slot carries the \emph{reversed} product, and order-dependence is the same
disease in operator clothing that non-additivity is in information clothing. The mathematics is
the epistemic argument's shadow, welcome as consistency and as a bridge to the module-side
literature (Appendix~\ref{app:winding-number} walks the bridge for charge), but the origin of the
complex amplitude is the observer's ignorance of their own state, and the arrow runs from there
to the algebra, never the reverse.
%  sec:l2-metric-fibre is REAL (backward, order-consistent);
%  sec:fanout-boundary DRAFTED 2026-07-06 (formerly the one undrafted section);
%  the old sec:monogamy placeholder note was VOID -- no such \ref in
%  the body -- and is removed. prop:f1-dictionary and
%  sec:hopf-justification are FORWARD references (position 8): first
%  F1 use now carries the deliberate-forward flag.
%
%  DEPENDENCY: % =====================================================================
%  ASIDE: entanglement monogamy as the budget reading of the ball.
%  Input by wigners_friend_in_the_hopf_picture_2026-06-13.tex inside
%  sec:wf-reconciliation, immediately after op:tsirelson-geometry (the
%  \input was relocated there 2026-07-06 from end-of-file, where the
%  opening antecedent "the same boundedness" had been orphaned).
%
%  AUDITED 2026-07-06 (final pass, with the WF section): the identity
%  tau = 1 - r^2 = 4 det rho verified by hand; the CKW derivation from
%  the one-tangle budget is sound, with C^2 + r^2 = 1 correctly the
%  bipartite-pure special case (no monogamy/complementarity
%  conflation). Fixes: orphan label sec:monogamy deleted (referenced
%  nowhere); keys retargeted to corpus canon (CKW2000 ->
%  coffmankunduwootters2000; OsborneVerstraete2006 ->
%  osborneverstraete2006 -- the June-7 "reconcile with Mendeley" debt
%  coming due); the closing confinement pointer retargeted from
%  sec:entanglement-classes to its true grown home, sec:confinement,
%  with the converse-theorem stitch to sec:budget-converse added.
%  This paragraph is the 2026-06-07 seed of the confinement budget
%  arc; it now points at its tree.
% =====================================================================

\noindent
\label{sec:monogamy}
The same boundedness carries entanglement monogamy, and by the same
move. A qubit $A$ has a single reduced state $\rho_{A}$, hence a single Bloch
radius $r_{A}$ and a single position in its ball, fixed once and for all by the
global state. The inward depth from the surface,
\begin{equation}
  \tau_{A} \;=\; 1 - r_{A}^{2} \;=\; 4\det\rho_{A},
\end{equation}
is $A$'s entire entanglement budget---one number, the one-tangle, equal in the
bipartite pure case to $C^{2}(A,B)$ through the Mosseri--Dandoloff relation
\autocite{mosseri2001hopf}. Every partner $A$ entangles with must be paid
for out of this same $r_{A}$: because $C^{2}(A,B)$ and $C^{2}(A,C)$ are both read
off the one reduced state, the squared concurrences sum within the budget,
\begin{equation}
  C^{2}(A,B) + C^{2}(A,C) \;\le\; \tau_{A} \;=\; 1 - r_{A}^{2},
\end{equation}
which is the Coffman--Kundu--Wootters monogamy inequality
\autocite{coffmankunduwootters2000}, proved for arbitrarily many qubits by Osborne and
Verstraete \autocite{osborneverstraete2006}. Driving $A$ to maximal entanglement with $B$
sends it to the centre, $r_{A}=0$, which spends the whole budget and leaves
nothing for $C$: entanglement is necessarily monogamous. Geometrically, monogamy
and the Tsirelson bound are one fact seen twice. A ball has a single centre and
one cannot pass it; the single radial depth toward that centre is a fixed budget
shared among all of $A$'s partners. The Tsirelson bound is that depth read
through the CHSH correlator---no point past the centre, so no violation past
$2\sqrt{2}$; monogamy is that depth partitioned among partners---one budget, so
maximal sharing with one excludes the rest. The inequality itself is standard;
what the Hopf picture adds is the reading that grounds it in the compactness and
unique centre of the base space, and that makes the same bounded radius do both
jobs at once. Pushed to its extreme the budget picture also yields confinement:
three GHZ-correlated quarks saturate the tripartite budget and retain none for partners outside
the triple --- developed in full at \S\ref{sec:confinement}, whose ledger
\S\ref{sec:budget-converse} completes with the converse reading, the W class's unspent ninth.
 -- received from JS
%  and audited 2026-07-06; the \input sits inside
%  sec:wf-reconciliation, immediately after op:tsirelson-geometry
%  (relocated 2026-07-06 from end-of-file, where the aside's opening
%  antecedent had been orphaned). A bare replace-all during the earlier
%  patch briefly mangled this header block, leaving a live \input in
%  it; repaired by index-splice, and the lesson recorded: anchors that
%  also occur in comments must be uniqueness-asserted.
%
%  AUDIT RECORD: CHSH arithmetic (2sqrt(1+C^2), PR => r^2 = -2) and the
%  C^2+r^2=1 chain verified; F1 wording canonical; the
%  obsolescence-made-cancellation echo two-way with
%  sec:self-reference; redundancy/Darwinism/DiBiagio-Rovelli/Brukner
%  positioning sound; Landauer placed on amplification, correctly.
%  Fixes: first-person voice removed; \textquote replaced with
%  portable quotes; IOEP given its prop ref; stale Tsirelson
%  mini-header (wrong key name Tsirelson1980) retired;
%  op:tsirelson-geometry MINTED from the honesty paragraph's unpriced
%  owed. Popper page-pin DISCHARGED 2026-07-07: the zip bundles carry
%  OCR .txt sidecars (the earlier "no text layer" verdict was wrong --
%  ZIP-of-JPEGs defeated the PDF tools, not the text); the quoted
%  phrase sits at Part II, p. 187, now pinned in the cite. Bonus from
%  the same dig: "C+" is Popper's own notation (Part II, scan pp.
%  15-17), so the L2 opening's attribution stands as written; "E" not
%  sighted in the sampled pages -- the inline gloss covers it either
%  way. Keys unique to this file (presumed in the .bib,
%  since the paper compiles): VerstraeteWolf2002,
%  Cirelson1980QuantumInequality, PopescuRohrlich1994, Held2021,
%  Brukner2021, zurek_2009, Landauer_1961, wigner_1961.
%  British English; \autocite; no bullets in prose.
% =====================================================================

\section{Wigner's friend in the Hopf picture: one state, two charts, and the redundancy ladder}
\label{sec:wigners-friend}

The framework dissolves the Wigner's-friend paradox \autocite{wigner_1961} not by
adjudicating between the two accounts but by recognising that they are written in
two charts on a single manifold, selected as the problem requires, much as one
moves between cylindrical and spherical coordinates without either being the true
one. The epistemological chart (private, non-clonable, phase-carrying) and the
ontological chart (clonable, public, record-bearing) are related by the
\textsc{fanout} boundary (Section~\ref{sec:fanout-boundary}), and that boundary
is not a sharp wall but a graded one, parametrised by how redundantly a piece of
relational content has been copied. Read this way, the fact that a \textsc{fanout}
has occurred for the friend but not for Wigner is not a contradiction to be
removed; it is the boundary behaving exactly as defined.

Let the system be a qubit $\lvert\psi\rangle_{S}=\alpha\lvert 0\rangle
+\beta\lvert 1\rangle$ and let the friend $F$, idealised, with deliberate
economy, as a single qubit, measure it in the computational basis. The
interaction is the controlled coupling that copies $S$'s value into $F$'s
register: the \textsc{copy} comultiplication of the diagrammatic calculus
\autocite{CoeckeKissinger2017}, and in the circuit picture the very gate that, in
the two-qubit Mosseri--Dandoloff geometry \autocite{mosseri2001hopf},
activates the quaternionic Hopf fibration $S^{3}\to S^{7}\to S^{4}$ and carries
the joint state from the surface of the entanglement ball to its interior. The
friend and Wigner read the two ends of that plunge.

Two idealisations in this set-up should be named at once, because the argument
later depends on seeing past them. The first is the friend as a single qubit,
which we return to as the redundancy-one endpoint of a graded boundary
(Section~\ref{sec:wf-redundancy}). The second, less obvious and more consequential,
is the system $S$ as a pristine qubit with no prior entanglement --- a relation
with no history. No real measured particle is of this kind. A particle is itself a
reified relation between prior events, and by the time the friend couples to it, it
already carries the completed internal structure such a relation possesses: in the
framework's terms, the full octonionic $A,B,C$ register that the $S^{7}$ sector
supplies as internal degrees of freedom (\S\ref{sec:space-of-computations}). The
pristine qubit is a schematic convenience, drawn structureless so that the logic of
a single link is visible; the ``system, friend, Wigner'' triple is likewise the
schematic of \emph{one link} in the chain of observers, not its base. The chain
does not bottom out in a structureless qubit --- there is no such thing --- and we
draw on this when we ask, later, what bounds the information a single such link can
carry.

\subsection{The friend's chart: a \textsc{fanout} of redundancy one}
\label{sec:wf-friend}

Relative to $F$, the coupling is a \textsc{fanout}: $F$ acquires a fact. What is
copied is gauge-invariant, the relative phase, the correlation between $F$ and
$S$, which is a genuine observable. The global phase is never the thing copied:
there is nothing observable there to duplicate, so the $S^{1}$ fibre of the
complex Hopf fibration $S^{1}\to S^{3}\to S^{2}$ of
Section~\ref{sec:l2-metric-fibre} remains $F$'s gauge-private self-ignorance,
untouched. In Hopf terms the division is clean, but it lives at two levels, which
Proposition~\ref{prop:f1-dictionary} (\S\ref{sec:hopf-justification} --- a deliberate forward
reference) separates. The copy operation itself is fibre-committed:
copying is faithful to phase, so what $F$'s \textsc{fanout} duplicates is a lift --- a chosen
fibre-point over the outcome --- never a bare point of the base. Yet everything the operation can
broadcast is base data alone: the record it induces is independent of the lift, so the $S^{1}$
freedom, though pinned in the copier, is absent from every classical record it yields. $F$'s classical record is a
base-point; the structure that produced it committed to a fibre-point; and the gap between the two
is precisely the gauge character of the fibre. The observation has a fifty-year-old
anchor: Popper argued that a predictor's inferior self-knowledge follows from the logical
impossibility of a system's imparting information about itself without thereby rendering it out of
date --- ``making this information obsolete'', in his words
\autocites{Popper1950a}[Part~II, p.~187]{Popper1950b}. Proposition~\ref{prop:f1-dictionary}
makes the impossibility structural rather than temporal, and localises it: the datum the copier
cannot put into its own record is not its state wholesale but exactly its fibre commitment, which
does not arrive stale --- it cancels. Conditioned on the outcome it reads, $S$ is for
$F$ a definite point on the Bloch sphere,
\begin{equation}
  \rho_{S}^{(F)}=\lvert 0\rangle\langle 0\rvert
  \quad\text{(given outcome }0\text{)},
\end{equation}
a pure state of Bloch radius $r=1$ on the surface of the ball, concurrence $C=0$.

This is a \textsc{fanout}, but the minimal one. The fact $F$ holds is a
\emph{relative} fact in the sense of relational quantum mechanics
\autocite{Rovelli1996a}: $S$ definite relative to $F$, not an observer-independent
record. It is one copy, redundancy $R=1$, and as such it remains reversible by an
agent who controls the joint system. The single-qubit friend is therefore not
outside \textsc{fanout}, as a careless reading might suggest; it is the $R=1$
endpoint of the \textsc{fanout} continuum---a \textsc{fanout} that has happened,
and a fact that is real for $F$, while being maximally still reversible from
outside.

\subsection{Wigner's chart: no copy has reached him}
\label{sec:wf-wigner}

Wigner treats the laboratory as closed and no copy of the relational content has
reached him. For him the coupling is the unitary
\begin{equation}
  \lvert\Psi\rangle
  =\alpha\lvert 0\rangle\lvert F_{0}\rangle
  +\beta\lvert 1\rangle\lvert F_{1}\rangle,
\end{equation}
the move up to the larger Hilbert space, $\mathbb{C}^{2}\to\mathbb{C}^{4}$,
$S^{3}\to S^{7}$, complex Hopf to quaternionic Hopf, with the fibre grown from
$S^{1}$ to $S^{3}$. What $F$ underwent as the acquisition of a fact, Wigner
describes as correlation deposited into a joint pure state on $S^{7}$. When he
asks specifically about $S$, he traces $F$ out, and the geometry does the
bookkeeping:
\begin{equation}
  \rho_{S}^{(W)}=\operatorname{Tr}_{F}\lvert\Psi\rangle\langle\Psi\rvert
  =\lvert\alpha\rvert^{2}\lvert 0\rangle\langle 0\rvert
  +\lvert\beta\rvert^{2}\lvert 1\rangle\langle 1\rvert,
  \qquad
  r=\bigl\lvert\,\lvert\alpha\rvert^{2}-\lvert\beta\rvert^{2}\bigr\rvert,
  \quad
  C=2\lvert\alpha\beta\rvert .
\end{equation}
Tracing out discards the $S^{3}$ fibre and carries the pure point on $S^{7}$ to a
mixed point in the interior of the Bloch ball at radius $r=\sqrt{1-C^{2}}$
\autocite{mosseri2001hopf}. The inward displacement is governed by the
concurrence: it is the relational content that has migrated into the $S$--$F$
correlation, stored in the recoverable $S^{3}$ fibre. Nothing has been destroyed;
it has been relocated into the fibre, which is why Wigner can in principle
interfere the branches and reverse the interaction, the content sits at low
redundancy and no further copy has fixed it.

\subsection{One state, two charts: \texorpdfstring{$C^{2}+r^{2}=1$}{C2+r2=1}}
\label{sec:wf-reconciliation}

The system is pure for the friend ($r=1$, on the surface) and mixed for Wigner
($r<1$, in the interior). These are not rival claims about a single fact; they are
the friend at the surface and Wigner at the interior of the \emph{same} state on
$S^{7}$, bound by the Mosseri--Dandoloff constraint
\begin{equation}
  C^{2}+r^{2}=1,
\end{equation}
in which the discrepancy in Bloch radius is exactly the concurrence Wigner
retains and the friend has resolved away. This is the Inside/Outside Equivalence
Principle (Proposition~\ref{prop:inside-outside}) made concrete, the same law written by the embedded participant and by
the external observer, and the relativity of facts of relational quantum
mechanics \autocite{Rovelli1996a} realised in geometry rather than asserted: the
\textsc{fanout} boundary lies \emph{between} $F$ and $W$, so a fact exists
relative to $F$ and not relative to $W$. Neither chart is privileged; one reads in
the epistemological or the ontological chart as the question requires, and their
mutual incompleteness is a theorem of the construction, not a defect of either.

\medskip
\noindent\emph{Aside: the Tsirelson bound and the absence of PR boxes.}
\label{sec:tsirelson}
A ball has a unique centre and no point beyond it. Maximal entanglement sits at
that centre \autocite{mosseri2001hopf}, $r=0$, $C=1$, and because the
Bloch radius cannot be negative, the concurrence cannot exceed unity. For a
two-qubit pure state the maximal CHSH correlator is
\begin{equation}
  \langle\mathrm{CHSH}\rangle_{\max}=2\sqrt{1+C^{2}}=2\sqrt{2-r^{2}}  
\end{equation}
\autocite{VerstraeteWolf2002},
which the centre maximises to $2\sqrt{2}$, the Tsirelson bound
\autocite{Cirelson1980QuantumInequality}. A Popescu--Rohrlich box \autocite{PopescuRohrlich1994},
with CHSH value $4$, would require $C^{2}=3$ and hence $r^{2}=-2$: a point two
units past the centre, which the ball does not contain. The super-quantum
correlation is excluded by the very positivity that places maximal entanglement
at the centre, one cannot trace out to a more-than-maximally-mixed reduced
state. The Tsirelson bound is in this sense the image, under the CHSH map, of the
boundedness of the $L^{2}$ state ball derived in
Section~\ref{sec:l2-metric-fibre}, and the prohibition on PR boxes shares its
origin with the roundness of the sphere: both trace to the reversible-continuity
metric, not to a separate principle.

\noindent This is the two-qubit, pure-state face of the bound, and its status is
worth stating plainly. It exhibits PR correlations as having no representing state
in the Hilbert geometry, and it recovers the value $2\sqrt{2}$ at the centre; but
forcing \emph{exactly} $2\sqrt{2}$ as a theorem of the geometry, rather than
reading it from the known concurrence--CHSH relation, and lifting the argument
from two qubits to the general Tsirelson bound, remain to be made rigorous. The
closest cognate in the literature is Held's derivation of the same bound from the
geometric algebra $G^{3}\cong\mathbb{H}$ \autocite{Held2021}, the same
quaternionic structure in a different language.

\begin{openproblem}[The Tsirelson bound as geometry]\label{op:tsirelson-geometry}
Derive $2\sqrt{2}$ as a theorem of the round state geometry rather than reading it off the
concurrence--CHSH relation, and lift the argument from the two-qubit pure case to the general
Tsirelson bound --- so that the exclusion of super-quantum correlation becomes, like the exclusion
exhibited above, a consequence of the $L^{2}$ ball's positivity (\S\ref{sec:l2-metric-fibre}) and
of nothing else. Held's quaternionic derivation \autocite{Held2021} is the nearest cognate and the
natural first tool.
\end{openproblem}

% =====================================================================
%  ASIDE: entanglement monogamy as the budget reading of the ball.
%  Input by wigners_friend_in_the_hopf_picture_2026-06-13.tex inside
%  sec:wf-reconciliation, immediately after op:tsirelson-geometry (the
%  \input was relocated there 2026-07-06 from end-of-file, where the
%  opening antecedent "the same boundedness" had been orphaned).
%
%  AUDITED 2026-07-06 (final pass, with the WF section): the identity
%  tau = 1 - r^2 = 4 det rho verified by hand; the CKW derivation from
%  the one-tangle budget is sound, with C^2 + r^2 = 1 correctly the
%  bipartite-pure special case (no monogamy/complementarity
%  conflation). Fixes: orphan label sec:monogamy deleted (referenced
%  nowhere); keys retargeted to corpus canon (CKW2000 ->
%  coffmankunduwootters2000; OsborneVerstraete2006 ->
%  osborneverstraete2006 -- the June-7 "reconcile with Mendeley" debt
%  coming due); the closing confinement pointer retargeted from
%  sec:entanglement-classes to its true grown home, sec:confinement,
%  with the converse-theorem stitch to sec:budget-converse added.
%  This paragraph is the 2026-06-07 seed of the confinement budget
%  arc; it now points at its tree.
% =====================================================================

\noindent
\label{sec:monogamy}
The same boundedness carries entanglement monogamy, and by the same
move. A qubit $A$ has a single reduced state $\rho_{A}$, hence a single Bloch
radius $r_{A}$ and a single position in its ball, fixed once and for all by the
global state. The inward depth from the surface,
\begin{equation}
  \tau_{A} \;=\; 1 - r_{A}^{2} \;=\; 4\det\rho_{A},
\end{equation}
is $A$'s entire entanglement budget---one number, the one-tangle, equal in the
bipartite pure case to $C^{2}(A,B)$ through the Mosseri--Dandoloff relation
\autocite{mosseri2001hopf}. Every partner $A$ entangles with must be paid
for out of this same $r_{A}$: because $C^{2}(A,B)$ and $C^{2}(A,C)$ are both read
off the one reduced state, the squared concurrences sum within the budget,
\begin{equation}
  C^{2}(A,B) + C^{2}(A,C) \;\le\; \tau_{A} \;=\; 1 - r_{A}^{2},
\end{equation}
which is the Coffman--Kundu--Wootters monogamy inequality
\autocite{coffmankunduwootters2000}, proved for arbitrarily many qubits by Osborne and
Verstraete \autocite{osborneverstraete2006}. Driving $A$ to maximal entanglement with $B$
sends it to the centre, $r_{A}=0$, which spends the whole budget and leaves
nothing for $C$: entanglement is necessarily monogamous. Geometrically, monogamy
and the Tsirelson bound are one fact seen twice. A ball has a single centre and
one cannot pass it; the single radial depth toward that centre is a fixed budget
shared among all of $A$'s partners. The Tsirelson bound is that depth read
through the CHSH correlator---no point past the centre, so no violation past
$2\sqrt{2}$; monogamy is that depth partitioned among partners---one budget, so
maximal sharing with one excludes the rest. The inequality itself is standard;
what the Hopf picture adds is the reading that grounds it in the compactness and
unique centre of the base space, and that makes the same bounded radius do both
jobs at once. Pushed to its extreme the budget picture also yields confinement:
three GHZ-correlated quarks saturate the tripartite budget and retain none for partners outside
the triple --- developed in full at \S\ref{sec:confinement}, whose ledger
\S\ref{sec:budget-converse} completes with the converse reading, the W class's unspent ninth.


\subsection{The redundancy ladder: from relative fact to objective fact}
\label{sec:wf-redundancy}

The friend's fact becomes a fact for Wigner only when Wigner copies it when he
measures $F$. It becomes a fact for a third party only when that party copies it
from Wigner, and so on up the von Neumann chain. The point at which the content
stops being negotiable is therefore not a point but a limit: once it has been
copied into enough mutually independent fragments that many parties can read it
without disturbing it, no one can practically un-copy it, and it is operationally
an objective fact. This is exactly quantum Darwinism \autocite{zurek_2009}:
objectivity is redundancy. It is not an external graft, because the framework's
account of decoherence already \emph{is} this picture---many fragments each take
their own copy of the clonable content, while their private phases, each
fragment's own self-ignorance, fail to add coherently and wash out.

This supplies the grading parameter that ``graded \textsc{fanout}'' had left
unnamed. \textsc{fanout} is graded not because the operation is fuzzy but because
redundancy $R$ runs continuously: $R=1$ is the single relational copy---the
friend's relative fact, reversible by Wigner---and $R\gg 1$ is the objective
fact, effectively irreversible because un-copying it would mean recovering every
fragment. The single-qubit friend is the $R=1$ endpoint; the macroscopic friend
is the same operation at large $R$. One axis, two ends. The relation to
Di~Biagio and Rovelli's distinction \autocite{DiBiagioRovelli2021} is then exact:
their \emph{stable} facts are the high-$R$ limit and their \emph{relative} facts
the low-$R$ regime, and the standard conditional-probability law they use to
define stability is the operational signature of redundancy having grown large.
Brukner's no-go theorem \autocite{Brukner2021}, that qubits are not observers,
sits on the same axis without conflict: it bites only against calling an $R=1$
relative fact an objective one, which the framework never does---the redundancy
ladder is precisely the structure that separates the relational definiteness a
single qubit does carry from the objectivity it does not.

The thermodynamics tracks the same axis. The single relational copy is unitary and
free; the $kT\ln 2$ cost \autocite{Landauer_1961} is incurred not by the friend's
act but by the amplification and stabilisation that drive $R$ upward, when
fragments are irreversibly committed. The relative fact is thermodynamically free;
the objective fact is what costs $kT\ln 2$.

\subsection{Phase promotion, and the bump in the carpet}
\label{sec:wf-phase-promotion}

One point keeps the load-bearing irrecoverability intact while explaining what
Wigner can do that the friend cannot. What Wigner recovers is \emph{not} the
$S^{1}$ self-ignorance; the global phase remains gauge for every observer and that
circle is spent for good. The quantity the friend resolves away is the relative
phase between the $\lvert 0\rangle$ and $\lvert 1\rangle$ components, and moving up
to $\mathbb{C}^{4}$ \emph{promotes} it: the single-qubit relative phase becomes the
relative phase between $\lvert 0\rangle\lvert F_{0}\rangle$ and
$\lvert 1\rangle\lvert F_{1}\rangle$, which lives in the recoverable $S^{3}$ fibre
(itself a circle bundle over $S^{2}$, so the promotion adds a two-sphere of
relative-phase directions to the lone gauge circle). Gauge phase at one level
becomes physical, interferable phase one level up---which is why Wigner can run
the interference experiment the friend cannot. But Wigner now carries his own
fresh global-phase $S^{1}$ at the top of the joint state: his own irreducible
self-ignorance. Ascending the tower does not eliminate the self-reference, only
displaces it one level higher, and the regress up the chain is precisely the
redundancy ladder of the previous subsection. This is the framework's bump in the
carpet: each observer sits atop a Hopf tower with a single gauge circle he can
never see around, and a measurement relocates the next observer's blind circle
into his own recoverable fibre without ever closing the regress.

Two cautions close the section. First, the recoverability Wigner
enjoys is the low-redundancy limit: for a macroscopic, high-$R$ friend the
reversal is exponentially hard and the fact is objective for all practical
purposes, so Wigner's ``no record'' is the $R\!\approx\!1$ idealisation of a
graded boundary, not a free lunch. Second, naming redundancy as the grading
parameter of \textsc{fanout}, and the phase-promotion reading itself, are
syntheses of the framework's parts: each rests on established pieces (the
clonability criterion, quantum Darwinism, the tensor-product promotion of phase)
but is assembled here, and is offered as the cleanest statement of why the two
observers never disagree about anything both can check.
% =====================================================================
%  SPLICE: "Where the chain goes next" -- forward-looking coda for
%  sec:wigners-friend, pointing to sec:sedenions-generations.
%  Created 2026-08-27 per JS.
%
%  THE POINT (JS, 2026-08-27). The Wigner's-friend charts stop at
%  Wigner because the Hopf tower runs out at the octonions, while the
%  text says the chain of observers obviously continues. The sedenions
%  show what the continuation adds: the observer of Wigner (fourth
%  wire, C (x) S) brings a NEW piece of structure -- the finite S_3 of
%  Brown's theorem, the generation label -- and the observer of the
%  observer of Wigner (fifth wire, the trigintaduonions) brings nothing
%  new in kind. This coda records that, so the section's termination at
%  Wigner reads as a fact about the tower rather than an economy of the
%  exposition.
%
%  WHAT THE ALGEBRA ACTUALLY SAYS (verified 2026-08-27, notebook T --
%  notebook_T_trigintaduonion_automorphisms_2026-08-27.py):
%    - Aut(S) = G_2 x S_3 (Brown 1967), DIRECT product; the new S_3 is
%      generated by the doubling involution and Brown's explicit
%      order-three map psi, which fixes 1 and the doubling unit and
%      rotates every (a, a u) plane, a imaginary, by a third of a turn.
%    - psi permutes the three octonion HALVINGS O, psi O, psi^2 O of S
%      (each doubled by the same unit; meeting only in R). It does NOT
%      permute the three copies through the quaternion core -- those are
%      cycled by a lifted G_2 element (the colour-line three-cycle); see
%      the sedenion section's two-triples remark and the corrections
%      ledger of the same date.
%    - Brown: Aut(A_{n-1}) x S_3 = Aut(A_n) with EQUALITY for n = 4, 5, 6
%      (as quoted in gresnigtgourlayvarma2023three, sec. 4.3); Eakin and
%      Sathaye prove Brown's conjecture in general (VERIFY-CITE for the
%      exact general form -- their step is Aut(A_{n+1}) = Aut(A_n) x G,
%      G in {Z/2, S_3}; over R with the standard parameters the branch is
%      S_3 at every rung >= 4, which is what notebook T checks directly
%      at rungs 4, 5, 6). Derivation algebra g_2 from the octonions on
%      (Schafer 1954; notebook T: dim Der = 14 at n = 3, 4, 5).
%  So "nothing new at the fifth wire" is stated PRECISELY: the fifth
%  doubling adds another finite S_3 and nothing else -- no connected
%  symmetry, no fibration; the increment is the same at every rung
%  thereafter. Not "silent": REPETITIVE. The prose says so.
%
%  STATUS: the algebra is ESTABLISHED (Brown; Eakin--Sathaye; Schafer)
%  and directly checked (notebook T). The reading of the repeated finite
%  factor as what the redundancy ladder consumes is a RHYME, flagged.
%  The reading of the fourth rung's discrete orientation as the
%  generation label is the sedenion section's CANDIDATE identification,
%  pointed to, not re-argued here.
%
%  PLACEMENT: a new \subsection at the END of sec:wigners-friend, after
%  the "Two cautions" paragraph. Because that paragraph reads "Two
%  cautions close the section.", the patch also changes that one word:
%  "close the section" -> "close the argument". Both anchors
%  uniqueness-asserted (the July lesson: anchors that also occur in
%  comments must be asserted unique).
%
%  CROSS-REFS OUT (all exist in the synced source except the two
%  MAP-LABELs):
%    sec:hopf-justification   -- the "two climbs" paragraph (length vs
%                                termination); this coda sharpens it
%    sec:wf-redundancy         -- the redundancy ladder (this file)
%    sec:wf-phase-promotion    -- phase promotion (this file)
%    sec:sedenions-generations -- the sedenion section (2026-08-27 rev.)
%    prop:termination          -- Hurwitz/Adams termination
%  MAP-LABELs:
%    app:notebook-record       -- the notebook-record appendix
%                                (appendices/appendix_notebook_record_
%                                2026-07-11.tex); substitute its label.
%  NEW KEYS (RIS delivered: cayley_dickson_automorphism_refs_2026-08-27.ris):
%    eakinsathaye1990automorphisms, schafer1954cayleydickson
%  IN-LIBRARY: brown1967generalized.
%
%  OPTIONAL CONSUMER-POINTER (not patched; one clause for the "two
%  climbs" paragraph of hopf_justification_section_2026-06-08.tex, after
%  "that chain has no top --- it is the bump in the carpet, unbounded in
%  length."):
%      "(unbounded in length, and --- as \S\ref{sec:wf-chain-type}
%      records --- bounded in type: after the fourth observer nothing
%      new in kind arrives)"
%
%  EDITED 2026-08-30 (per JS): one sentence added to the fourth-wire
%  paragraph anchoring the flavour sector's irremovable phase (the
%  CP-phase recognition) on this rung, with the development at
%  sec:sedenions-mixing (2026-08-30 revision). No new keys here.
%  British English; pure ASCII; \autocite; no bullets; spaced dashes.
% =====================================================================
% SPLICE-BEGIN wf-chain-type 2026-08-27

\subsection{Where the chain goes next: one more rung, then copies}
\label{sec:wf-chain-type}

The regress just described is unbounded in length, and nothing here shortens
it: the observer of Wigner carries a gauge circle of his own, and so does the
observer of that observer, as \S\ref{sec:hopf-justification} insists when it
separates the endlessness of the chain from the termination of the tower. But
the algebra says how much of the endless chain adds structure of a new kind,
and the answer is exact: one more step, and then none.

The count runs as follows. The triple system, friend, Wigner is three wires,
and three wires are the three Hopf rungs $\mathbb{C}$, $\mathbb{H}$,
$\mathbb{O}$ --- the completed link of the earlier caveat --- which is why the
charts above stop at Wigner: the chart of whoever looks in on the triple,
$S^{15}\to S^{8}$, is the last fibration the division algebras supply
(Proposition~\ref{prop:termination}). The observer of Wigner is a fourth wire,
and the fourth wire is the sedenions,
$(\mathbb{C}^{2})^{\otimes 4} = \mathbb{C}\otimes\mathbb{S}$
(\S\ref{sec:sedenions-generations}). No fibration comes with it --- the
sedenions have zero divisors and admit no map $S^{31}\to S^{16}$ --- so the
fourth observer gains no recoverable fibre of the tower's kind: the continuous
promotion of relative phase into a larger fibre, the mechanism of the previous
subsection, has no fourth instance. What the doubling adds instead is finite.
Brown's theorem gives $\mathrm{Aut}(\mathbb{S}) = G_{2}\times S_{3}$
\autocite{brown1967generalized}: the identity component is unchanged from the
octonions, and the entire gain is a component group, generated by an
order-three automorphism that fixes the new doubling direction and rotates
every imaginary direction of the completed link into its copy by a third of a
turn. The relative orientation of the link and its copy, which at the Hopf
rungs was a continuous phase promoted into a fibre, survives at the fourth rung
only as a discrete choice among three, and that choice is the generation label
(\S\ref{sec:sedenions-generations}). The fourth observer, then, sees no new
force and no new fibre; he sees the same link three ways, and cannot gauge the
difference, because a component group has no Lie algebra. His rung is also where
the flavour sector's one irremovable phase finds its anchor: three generations
are the threshold at which the mixing matrix first carries a phase no convention
removes, and \S\ref{sec:sedenions-mixing} reads that phase --- through the
paper's negativity--contextuality chain --- as this wire's contextual cargo,
real only from one rung above.

The observer of the observer of Wigner is a fifth wire, the trigintaduonions,
and here the tower has nothing new to say. Brown proved that the next two
doublings each add exactly one further $S_{3}$ and nothing else, Eakin and
Sathaye that Brown's pattern continues at every rung, and the derivation
algebra --- the Lie algebra of the identity component --- is $\mathfrak{g}_{2}$
from the octonions onward
\autocite{brown1967generalized,eakinsathaye1990automorphisms,schafer1954cayleydickson};
the fifth and sixth rungs are also checked directly in the notebook record
(Appendix~\ref{app:notebook-record}). So the fifth observer's doubling repeats
the fourth's gift --- another third-of-a-turn symmetry between the four-wire
structure and its copy, commuting with everything below --- and adds no
connected symmetry and no fibration. Nothing new in kind arrives at the fifth
wire or at any wire after it. The increments of the tower are, in order, a
force, a force, a label, and then the same finite permutation of copies at
every step. That is the precise sense in which the chain ``obviously
continues'': it continues, and after the fourth observer it has nothing further
to announce. The framework reads the repeated factor as what the redundancy
ladder of \S\ref{sec:wf-redundancy} consumes --- a permutation symmetry among
copies, which is the freedom a record has and a type does not --- and records
that reading as a structural resonance rather than a derivation.

The two climbs of \S\ref{sec:hopf-justification} are thereby sharpened. The
chain of observers is unbounded in length and bounded in type: three forces by
the third wire, one family label by the fourth, copies thereafter. The observer
of Wigner is the last observer whose arrival changes what there is.

% SPLICE-END wf-chain-type 2026-08-27

% =====================================================================
%  FIGURE: the Wigner's-friend chain mapped onto the Cayley--Dickson
%  ladder. Created 2026-08-30 per JS: the Hopf-fibration view of
%  Wigner's friend has been load-bearing for months without being drawn;
%  the fibrations themselves are DEFERRED to Szangolies's introduction
%  (szangolies2025standardmodel, in library), and this figure carries
%  only the mapping.
%
%  CONTENT (left to right, one panel per chart, staircase rising):
%    Panel 1: F's chart of S            -- S^3 -> S^2, fibre S^1,  C, one wire
%    Panel 2: W's chart of {S,F}        -- S^7 -> S^4, fibre S^3,  H, two wires
%    Panel 3: O_W's chart of {S,F,W}    -- S^15 -> S^8, fibre S^7, O, three wires
%             (the last fibration)
%    Panel 4: the next observer's chart of {S,F,W,O_W} -- C^16 = C (x) S,
%             no fibration (S^31 -/-> S^16); the doubling's S_3 badge:
%             the generation label, a third of a turn.
%  Each describer perches on his sealed box with a dashed circle over
%  his head (his own gauge circle -- the self-referential ignorance he
%  cannot see around); dashed promotion arrows carry each circle into
%  the next chart's fibre ("gauge here -> interferable one rung up",
%  sec:wf-phase-promotion); the third arrow terminates at "no fourth
%  fibration", with what survives being the discrete third-of-a-turn
%  (sec:wf-chain-type, sec:sedenions-generations). The regress
%  continues to the right: copies only.
%
%  PREAMBLE REQUIREMENT (confirm in main.tex; add if absent):
%      \usepackage{tikz}
%      \usetikzlibrary{arrows.meta}
%  PLACEMENT: \input by patch_wigner_hopf_ladder_fig_2026-08-30.py,
%  after the two-idealisations paragraph of sec:wigners-friend, with
%  the accompanying deferral paragraph (in the patch).
%  CROSS-REFS OUT (all real): sec:wf-phase-promotion,
%  sec:wf-chain-type, sec:sedenions-generations.
%  KEY: szangolies2025standardmodel (in library).
%  British English; pure ASCII; \autocite in the caption.
% =====================================================================
\begin{figure}[t]
\centering
\resizebox{\textwidth}{!}{%
\begin{tikzpicture}[>=Stealth, line cap=round, line join=round,
    every node/.style={font=\scriptsize},
    blind/.style={densely dashed, draw=black!75},
    promote/.style={->, densely dashed, black!60, semithick},
    sight/.style={->, dotted, black!55}]
% ---- stick-figure pic --------------------------------------------------
\tikzset{stick/.pic={
  \draw[pic actions] (0,0.44) circle (0.085);
  \draw[pic actions] (0,0.355) -- (0,0.10);
  \draw[pic actions] (-0.12,0.20) -- (0,0.29) -- (0.12,0.20);
  \draw[pic actions] (-0.10,-0.10) -- (0,0.10) -- (0.10,-0.10);
}}
% ---- the staircase (Cayley--Dickson ladder) ---------------------------
\draw[black!35, thick]
  (-0.35,-0.55) -- (3.27,-0.55) -- (3.27,0.25) -- (6.89,0.25)
  -- (6.89,1.05) -- (10.51,1.05) -- (10.51,1.85) -- (13.60,1.85)
  -- (13.60,2.65) -- (16.20,2.65);
\node[anchor=north] at (1.25,-0.60) {$\mathbb{C}$ --- one wire};
\node[anchor=north] at (4.87,0.20)  {$\mathbb{H}$ --- two wires};
\node[anchor=north] at (8.49,1.00)  {$\mathbb{O}$ --- three wires};
\node[anchor=north] at (12.11,1.80) {$\mathbb{S}$ --- four wires};
\node[anchor=north, black!55] at (14.45,2.60) {$\mathbb{T},\dots$ --- five wires on};
% =======================================================================
% Panel 1: F's chart of S
% =======================================================================
\draw[rounded corners=2pt, fill=black!5, draw=black!40] (0,0) rectangle (2.5,1.32);
\filldraw[black!70] (1.25,0.72) circle (0.075);
\node[font=\tiny] at (1.25,0.45) {$S$};
\pic[black] at (2.05,1.42) {stick};
\node[anchor=west] at (2.22,1.62) {$F$};
\draw[sight] (1.95,1.60) -- (1.45,1.02);
\draw[blind] (2.05,2.14) circle (0.11);
\node[anchor=east, align=right, font=\tiny, black!75] at (1.82,2.14)
  {his own gauge circle:\\ the phase he cannot\\ see around};
\node[anchor=south west, font=\tiny, black!75] at (0.02,1.34) {fibre $S^{1}$};
\node[anchor=north] at (1.25,-0.06) {$S^{3}\to S^{2}$};
% =======================================================================
% Panel 2: W's chart of {S,F}
% =======================================================================
\draw[rounded corners=2pt, fill=black!5, draw=black!40] (3.62,0.80) rectangle (6.12,2.12);
\filldraw[black!70] (4.07,1.52) circle (0.075);
\node[font=\tiny] at (4.07,1.25) {$S$};
\pic[black!55, scale=0.85] at (4.85,1.16) {stick};
\node[font=\tiny, black!55] at (4.85,0.94) {$F$};
\pic[black] at (5.67,2.22) {stick};
\node[anchor=west] at (5.84,2.42) {$W$};
\draw[sight] (5.57,2.40) -- (5.05,1.82);
\draw[blind] (5.67,2.94) circle (0.11);
\node[anchor=south west, font=\tiny, black!75] at (3.64,2.14) {fibre $S^{3}$};
\node[anchor=north] at (4.87,0.74) {$S^{7}\to S^{4}$};
\draw[promote] (2.17,2.14) to[out=15,in=170] (3.60,2.24);
\node[anchor=south, align=center, font=\tiny, black!60] at (2.85,2.32)
  {gauge below,\\ interferable above};
% =======================================================================
% Panel 3: O_W's chart of {S,F,W}
% =======================================================================
\draw[rounded corners=2pt, fill=black!5, draw=black!40] (7.24,1.60) rectangle (9.74,2.92);
\filldraw[black!70] (7.69,2.32) circle (0.075);
\node[font=\tiny] at (7.69,2.05) {$S$};
\pic[black!55, scale=0.85] at (8.35,1.96) {stick};
\node[font=\tiny, black!55] at (8.35,1.74) {$F$};
\pic[black!55, scale=0.85] at (8.95,1.96) {stick};
\node[font=\tiny, black!55] at (8.95,1.74) {$W$};
\pic[black] at (9.29,3.02) {stick};
\node[anchor=west] at (9.46,3.22) {$O_{W}$};
\draw[sight] (9.19,3.20) -- (8.65,2.62);
\draw[blind] (9.29,3.74) circle (0.11);
\node[anchor=south west, font=\tiny, black!75] at (7.26,2.94) {fibre $S^{7}$};
\node[anchor=north, align=center] at (8.49,1.54) {$S^{15}\to S^{8}$\\[-1pt]
  {\tiny the last fibration}};
\draw[promote] (5.79,2.94) to[out=15,in=170] (7.22,3.04);
% =======================================================================
% Panel 4: the next chart of {S,F,W,O_W} -- sedenions, no fibration
% =======================================================================
\draw[rounded corners=2pt, fill=black!5, draw=black!40] (10.86,2.40) rectangle (13.36,3.72);
\filldraw[black!70] (11.24,3.12) circle (0.075);
\node[font=\tiny] at (11.24,2.85) {$S$};
\pic[black!55, scale=0.80] at (11.78,2.78) {stick};
\node[font=\tiny, black!55] at (11.78,2.56) {$F$};
\pic[black!55, scale=0.80] at (12.28,2.78) {stick};
\node[font=\tiny, black!55] at (12.28,2.56) {$W$};
\pic[black!55, scale=0.80] at (12.78,2.78) {stick};
\node[font=\tiny, black!55] at (12.78,2.56) {$O_{W}$};
\pic[black!45] at (14.35,2.78) {stick};
\draw[sight] (14.18,3.02) -- (13.42,3.12);
\draw[blind] (14.35,3.50) circle (0.11);
\node[anchor=west, align=left, font=\tiny, black!60] at (14.58,3.48)
  {the next observer,\\ and all after:\\ copies only};
\node[anchor=north] at (12.11,2.34) {$\mathbb{C}^{16}=\mathbb{C}\otimes\mathbb{S}$};
% the broken third promotion
\draw[promote] (9.41,3.74) to[out=15,in=185] (10.68,3.86);
\node[font=\scriptsize, black!60] at (10.80,3.92) {$\times$};
\node[anchor=south west, font=\tiny, black!75] at (10.88,3.74)
  {no fourth fibration: $S^{31}\not\to S^{16}$};
% the S3 badge: three dots at 120 degrees with a rotation arrow
\draw[black!75] (11.45,4.30) circle (0.17);
\filldraw[black!75] (11.45,4.47) circle (0.028);
\filldraw[black!75] (11.303,4.215) circle (0.028);
\filldraw[black!75] (11.597,4.215) circle (0.028);
\draw[->, black!75] (11.67,4.36) arc[start angle=20, end angle=100, radius=0.235];
\node[anchor=west, align=left, font=\tiny, black!75] at (11.70,4.36)
  {$S_{3}$: the generation label\\ (a third of a turn)};
\draw[promote, black!70] (10.86,3.96) to[out=60,in=200] (11.24,4.22);
\end{tikzpicture}%
}
\caption{The Wigner's-friend chain mapped onto the Cayley--Dickson ladder. Each
observer holds a chart of the sealed box below him, one rung further up the
ladder: the friend's chart of the system is the complex Hopf bundle, Wigner's
chart of the sealed laboratory the quaternionic, and the chart of the observer
of Wigner --- of the completed triple --- the octonionic, the tower's last
fibration. Every chart-holder carries a dashed circle over his head, his own
gauge phase: the self-referential ignorance he cannot see around. Measurement
promotes each circle into the next chart's fibre, where what was gauge below is
interferable above (\S\ref{sec:wf-phase-promotion}); the third promotion has no
fibration to land in --- the sedenions have zero divisors and no map
$S^{31}\to S^{16}$ --- and what survives of it is discrete: the doubling's
$S_{3}$, a third of a turn rotating the link into its copy, the generation
label (\S\ref{sec:wf-chain-type}, \S\ref{sec:sedenions-generations}). The
regress of circles continues unbounded to the right, but from the fifth wire
onward it adds copies only. The friend drawn as a single qubit is the
redundancy-one idealisation discussed in the text; for the fibrations
themselves, deferred rather than repeated here, see the introduction in
\autocite{szangolies2025standardmodel}.}
\label{fig:wf-hopf-ladder}
\end{figure}

%      % EDITED 2026-10-29 scripts/patch_continuity_tower_premise_2026-10-29T1500.py
% EDITED 2026-09-20 scripts/patch_born_cost_ledger_v2_2026-09-20T0900.py
% EDITED 2026-09-15 scripts/patch_value_continuum_declared_cencov_2026-09-15T1540.py
% EDITED 2026-09-15 scripts/patch_continuity_connected_state_space_2026-09-15T1430.py
% =====================================================================
%  SECTION: Derivation of the quantum postulates (the assembly chapter).
%  AUDITED 2026-07-06 (final pass, sixth stop).
%
%  PLACEMENT (reordered chain, JS 2026-07-06): position 5 --
%      % EDITED 2026-09-12 scripts/patch_revert_sedenion_repoint_2026-09-12T1600.py
% EDITED 2026-09-11 scripts/patch_undefined_section_refs_2026-09-11T2030.py
% =====================================================================
%  SECTION: Wigner's friend in the Hopf picture (v2; supersedes the
%  2026-06-07 draft). AUDITED 2026-07-06 (final pass, fifth stop).
%
%  PLACEMENT (reordered chain, JS 2026-07-06): position 4 --
%      \input{L2_metric_global_phase_section_2026_06_07}
%      \input{wigners_friend_in_the_hopf_picture_2026-06-13}
%      \input{derivation_of_the_quantum_postulates}
%  sec:l2-metric-fibre is REAL (backward, order-consistent);
%  sec:fanout-boundary DRAFTED 2026-07-06 (formerly the one undrafted section);
%  the old sec:monogamy placeholder note was VOID -- no such \ref in
%  the body -- and is removed. prop:f1-dictionary and
%  sec:hopf-justification are FORWARD references (position 8): first
%  F1 use now carries the deliberate-forward flag.
%
%  DEPENDENCY: \input{monogamy_aside_2026-06-07} -- received from JS
%  and audited 2026-07-06; the \input sits inside
%  sec:wf-reconciliation, immediately after op:tsirelson-geometry
%  (relocated 2026-07-06 from end-of-file, where the aside's opening
%  antecedent had been orphaned). A bare replace-all during the earlier
%  patch briefly mangled this header block, leaving a live \input in
%  it; repaired by index-splice, and the lesson recorded: anchors that
%  also occur in comments must be uniqueness-asserted.
%
%  AUDIT RECORD: CHSH arithmetic (2sqrt(1+C^2), PR => r^2 = -2) and the
%  C^2+r^2=1 chain verified; F1 wording canonical; the
%  obsolescence-made-cancellation echo two-way with
%  sec:self-reference; redundancy/Darwinism/DiBiagio-Rovelli/Brukner
%  positioning sound; Landauer placed on amplification, correctly.
%  Fixes: first-person voice removed; \textquote replaced with
%  portable quotes; IOEP given its prop ref; stale Tsirelson
%  mini-header (wrong key name Tsirelson1980) retired;
%  op:tsirelson-geometry MINTED from the honesty paragraph's unpriced
%  owed. Popper page-pin DISCHARGED 2026-07-07: the zip bundles carry
%  OCR .txt sidecars (the earlier "no text layer" verdict was wrong --
%  ZIP-of-JPEGs defeated the PDF tools, not the text); the quoted
%  phrase sits at Part II, p. 187, now pinned in the cite. Bonus from
%  the same dig: "C+" is Popper's own notation (Part II, scan pp.
%  15-17), so the L2 opening's attribution stands as written; "E" not
%  sighted in the sampled pages -- the inline gloss covers it either
%  way. Keys unique to this file (presumed in the .bib,
%  since the paper compiles): VerstraeteWolf2002,
%  Cirelson1980QuantumInequality, PopescuRohrlich1994, Held2021,
%  Brukner2021, zurek_2009, Landauer_1961, wigner_1961.
%  British English; \autocite; no bullets in prose.
% =====================================================================

\section{Wigner's friend in the Hopf picture: one state, two charts, and the redundancy ladder}
\label{sec:wigners-friend}

The framework dissolves the Wigner's-friend paradox \autocite{wigner_1961} not by
adjudicating between the two accounts but by recognising that they are written in
two charts on a single manifold, selected as the problem requires, much as one
moves between cylindrical and spherical coordinates without either being the true
one. The epistemological chart (private, non-clonable, phase-carrying) and the
ontological chart (clonable, public, record-bearing) are related by the
\textsc{fanout} boundary (Section~\ref{sec:fanout-boundary}), and that boundary
is not a sharp wall but a graded one, parametrised by how redundantly a piece of
relational content has been copied. Read this way, the fact that a \textsc{fanout}
has occurred for the friend but not for Wigner is not a contradiction to be
removed; it is the boundary behaving exactly as defined.

Let the system be a qubit $\lvert\psi\rangle_{S}=\alpha\lvert 0\rangle
+\beta\lvert 1\rangle$ and let the friend $F$, idealised, with deliberate
economy, as a single qubit, measure it in the computational basis. The
interaction is the controlled coupling that copies $S$'s value into $F$'s
register: the \textsc{copy} comultiplication of the diagrammatic calculus
\autocite{CoeckeKissinger2017}, and in the circuit picture the very gate that, in
the two-qubit Mosseri--Dandoloff geometry \autocite{mosseri2001hopf},
activates the quaternionic Hopf fibration $S^{3}\to S^{7}\to S^{4}$ and carries
the joint state from the surface of the entanglement ball to its interior. The
friend and Wigner read the two ends of that plunge.

Two idealisations in this set-up should be named at once, because the argument
later depends on seeing past them. The first is the friend as a single qubit,
which we return to as the redundancy-one endpoint of a graded boundary
(Section~\ref{sec:wf-redundancy}). The second, less obvious and more consequential,
is the system $S$ as a pristine qubit with no prior entanglement --- a relation
with no history. No real measured particle is of this kind. A particle is itself a
reified relation between prior events, and by the time the friend couples to it, it
already carries the completed internal structure such a relation possesses: in the
framework's terms, the full octonionic $A,B,C$ register that the $S^{7}$ sector
supplies as internal degrees of freedom (\S\ref{sec:space-of-computations}). The
pristine qubit is a schematic convenience, drawn structureless so that the logic of
a single link is visible; the ``system, friend, Wigner'' triple is likewise the
schematic of \emph{one link} in the chain of observers, not its base. The chain
does not bottom out in a structureless qubit --- there is no such thing --- and we
draw on this when we ask, later, what bounds the information a single such link can
carry.

\subsection{The friend's chart: a \textsc{fanout} of redundancy one}
\label{sec:wf-friend}

Relative to $F$, the coupling is a \textsc{fanout}: $F$ acquires a fact. What is
copied is gauge-invariant, the relative phase, the correlation between $F$ and
$S$, which is a genuine observable. The global phase is never the thing copied:
there is nothing observable there to duplicate, so the $S^{1}$ fibre of the
complex Hopf fibration $S^{1}\to S^{3}\to S^{2}$ of
Section~\ref{sec:l2-metric-fibre} remains $F$'s gauge-private self-ignorance,
untouched. In Hopf terms the division is clean, but it lives at two levels, which
Proposition~\ref{prop:f1-dictionary} (\S\ref{sec:hopf-justification} --- a deliberate forward
reference) separates. The copy operation itself is fibre-committed:
copying is faithful to phase, so what $F$'s \textsc{fanout} duplicates is a lift --- a chosen
fibre-point over the outcome --- never a bare point of the base. Yet everything the operation can
broadcast is base data alone: the record it induces is independent of the lift, so the $S^{1}$
freedom, though pinned in the copier, is absent from every classical record it yields. $F$'s classical record is a
base-point; the structure that produced it committed to a fibre-point; and the gap between the two
is precisely the gauge character of the fibre. The observation has a fifty-year-old
anchor: Popper argued that a predictor's inferior self-knowledge follows from the logical
impossibility of a system's imparting information about itself without thereby rendering it out of
date --- ``making this information obsolete'', in his words
\autocites{Popper1950a}[Part~II, p.~187]{Popper1950b}. Proposition~\ref{prop:f1-dictionary}
makes the impossibility structural rather than temporal, and localises it: the datum the copier
cannot put into its own record is not its state wholesale but exactly its fibre commitment, which
does not arrive stale --- it cancels. Conditioned on the outcome it reads, $S$ is for
$F$ a definite point on the Bloch sphere,
\begin{equation}
  \rho_{S}^{(F)}=\lvert 0\rangle\langle 0\rvert
  \quad\text{(given outcome }0\text{)},
\end{equation}
a pure state of Bloch radius $r=1$ on the surface of the ball, concurrence $C=0$.

This is a \textsc{fanout}, but the minimal one. The fact $F$ holds is a
\emph{relative} fact in the sense of relational quantum mechanics
\autocite{Rovelli1996a}: $S$ definite relative to $F$, not an observer-independent
record. It is one copy, redundancy $R=1$, and as such it remains reversible by an
agent who controls the joint system. The single-qubit friend is therefore not
outside \textsc{fanout}, as a careless reading might suggest; it is the $R=1$
endpoint of the \textsc{fanout} continuum---a \textsc{fanout} that has happened,
and a fact that is real for $F$, while being maximally still reversible from
outside.

\subsection{Wigner's chart: no copy has reached him}
\label{sec:wf-wigner}

Wigner treats the laboratory as closed and no copy of the relational content has
reached him. For him the coupling is the unitary
\begin{equation}
  \lvert\Psi\rangle
  =\alpha\lvert 0\rangle\lvert F_{0}\rangle
  +\beta\lvert 1\rangle\lvert F_{1}\rangle,
\end{equation}
the move up to the larger Hilbert space, $\mathbb{C}^{2}\to\mathbb{C}^{4}$,
$S^{3}\to S^{7}$, complex Hopf to quaternionic Hopf, with the fibre grown from
$S^{1}$ to $S^{3}$. What $F$ underwent as the acquisition of a fact, Wigner
describes as correlation deposited into a joint pure state on $S^{7}$. When he
asks specifically about $S$, he traces $F$ out, and the geometry does the
bookkeeping:
\begin{equation}
  \rho_{S}^{(W)}=\operatorname{Tr}_{F}\lvert\Psi\rangle\langle\Psi\rvert
  =\lvert\alpha\rvert^{2}\lvert 0\rangle\langle 0\rvert
  +\lvert\beta\rvert^{2}\lvert 1\rangle\langle 1\rvert,
  \qquad
  r=\bigl\lvert\,\lvert\alpha\rvert^{2}-\lvert\beta\rvert^{2}\bigr\rvert,
  \quad
  C=2\lvert\alpha\beta\rvert .
\end{equation}
Tracing out discards the $S^{3}$ fibre and carries the pure point on $S^{7}$ to a
mixed point in the interior of the Bloch ball at radius $r=\sqrt{1-C^{2}}$
\autocite{mosseri2001hopf}. The inward displacement is governed by the
concurrence: it is the relational content that has migrated into the $S$--$F$
correlation, stored in the recoverable $S^{3}$ fibre. Nothing has been destroyed;
it has been relocated into the fibre, which is why Wigner can in principle
interfere the branches and reverse the interaction, the content sits at low
redundancy and no further copy has fixed it.

\subsection{One state, two charts: \texorpdfstring{$C^{2}+r^{2}=1$}{C2+r2=1}}
\label{sec:wf-reconciliation}

The system is pure for the friend ($r=1$, on the surface) and mixed for Wigner
($r<1$, in the interior). These are not rival claims about a single fact; they are
the friend at the surface and Wigner at the interior of the \emph{same} state on
$S^{7}$, bound by the Mosseri--Dandoloff constraint
\begin{equation}
  C^{2}+r^{2}=1,
\end{equation}
in which the discrepancy in Bloch radius is exactly the concurrence Wigner
retains and the friend has resolved away. This is the Inside/Outside Equivalence
Principle (Proposition~\ref{prop:inside-outside}) made concrete, the same law written by the embedded participant and by
the external observer, and the relativity of facts of relational quantum
mechanics \autocite{Rovelli1996a} realised in geometry rather than asserted: the
\textsc{fanout} boundary lies \emph{between} $F$ and $W$, so a fact exists
relative to $F$ and not relative to $W$. Neither chart is privileged; one reads in
the epistemological or the ontological chart as the question requires, and their
mutual incompleteness is a theorem of the construction, not a defect of either.

\medskip
\noindent\emph{Aside: the Tsirelson bound and the absence of PR boxes.}
\label{sec:tsirelson}
A ball has a unique centre and no point beyond it. Maximal entanglement sits at
that centre \autocite{mosseri2001hopf}, $r=0$, $C=1$, and because the
Bloch radius cannot be negative, the concurrence cannot exceed unity. For a
two-qubit pure state the maximal CHSH correlator is
\begin{equation}
  \langle\mathrm{CHSH}\rangle_{\max}=2\sqrt{1+C^{2}}=2\sqrt{2-r^{2}}  
\end{equation}
\autocite{VerstraeteWolf2002},
which the centre maximises to $2\sqrt{2}$, the Tsirelson bound
\autocite{Cirelson1980QuantumInequality}. A Popescu--Rohrlich box \autocite{PopescuRohrlich1994},
with CHSH value $4$, would require $C^{2}=3$ and hence $r^{2}=-2$: a point two
units past the centre, which the ball does not contain. The super-quantum
correlation is excluded by the very positivity that places maximal entanglement
at the centre, one cannot trace out to a more-than-maximally-mixed reduced
state. The Tsirelson bound is in this sense the image, under the CHSH map, of the
boundedness of the $L^{2}$ state ball derived in
Section~\ref{sec:l2-metric-fibre}, and the prohibition on PR boxes shares its
origin with the roundness of the sphere: both trace to the reversible-continuity
metric, not to a separate principle.

\noindent This is the two-qubit, pure-state face of the bound, and its status is
worth stating plainly. It exhibits PR correlations as having no representing state
in the Hilbert geometry, and it recovers the value $2\sqrt{2}$ at the centre; but
forcing \emph{exactly} $2\sqrt{2}$ as a theorem of the geometry, rather than
reading it from the known concurrence--CHSH relation, and lifting the argument
from two qubits to the general Tsirelson bound, remain to be made rigorous. The
closest cognate in the literature is Held's derivation of the same bound from the
geometric algebra $G^{3}\cong\mathbb{H}$ \autocite{Held2021}, the same
quaternionic structure in a different language.

\begin{openproblem}[The Tsirelson bound as geometry]\label{op:tsirelson-geometry}
Derive $2\sqrt{2}$ as a theorem of the round state geometry rather than reading it off the
concurrence--CHSH relation, and lift the argument from the two-qubit pure case to the general
Tsirelson bound --- so that the exclusion of super-quantum correlation becomes, like the exclusion
exhibited above, a consequence of the $L^{2}$ ball's positivity (\S\ref{sec:l2-metric-fibre}) and
of nothing else. Held's quaternionic derivation \autocite{Held2021} is the nearest cognate and the
natural first tool.
\end{openproblem}

\input{monogamy_aside_2026-06-07}

\subsection{The redundancy ladder: from relative fact to objective fact}
\label{sec:wf-redundancy}

The friend's fact becomes a fact for Wigner only when Wigner copies it when he
measures $F$. It becomes a fact for a third party only when that party copies it
from Wigner, and so on up the von Neumann chain. The point at which the content
stops being negotiable is therefore not a point but a limit: once it has been
copied into enough mutually independent fragments that many parties can read it
without disturbing it, no one can practically un-copy it, and it is operationally
an objective fact. This is exactly quantum Darwinism \autocite{zurek_2009}:
objectivity is redundancy. It is not an external graft, because the framework's
account of decoherence already \emph{is} this picture---many fragments each take
their own copy of the clonable content, while their private phases, each
fragment's own self-ignorance, fail to add coherently and wash out.

This supplies the grading parameter that ``graded \textsc{fanout}'' had left
unnamed. \textsc{fanout} is graded not because the operation is fuzzy but because
redundancy $R$ runs continuously: $R=1$ is the single relational copy---the
friend's relative fact, reversible by Wigner---and $R\gg 1$ is the objective
fact, effectively irreversible because un-copying it would mean recovering every
fragment. The single-qubit friend is the $R=1$ endpoint; the macroscopic friend
is the same operation at large $R$. One axis, two ends. The relation to
Di~Biagio and Rovelli's distinction \autocite{DiBiagioRovelli2021} is then exact:
their \emph{stable} facts are the high-$R$ limit and their \emph{relative} facts
the low-$R$ regime, and the standard conditional-probability law they use to
define stability is the operational signature of redundancy having grown large.
Brukner's no-go theorem \autocite{Brukner2021}, that qubits are not observers,
sits on the same axis without conflict: it bites only against calling an $R=1$
relative fact an objective one, which the framework never does---the redundancy
ladder is precisely the structure that separates the relational definiteness a
single qubit does carry from the objectivity it does not.

The thermodynamics tracks the same axis. The single relational copy is unitary and
free; the $kT\ln 2$ cost \autocite{Landauer_1961} is incurred not by the friend's
act but by the amplification and stabilisation that drive $R$ upward, when
fragments are irreversibly committed. The relative fact is thermodynamically free;
the objective fact is what costs $kT\ln 2$.

\subsection{Phase promotion, and the bump in the carpet}
\label{sec:wf-phase-promotion}

One point keeps the load-bearing irrecoverability intact while explaining what
Wigner can do that the friend cannot. What Wigner recovers is \emph{not} the
$S^{1}$ self-ignorance; the global phase remains gauge for every observer and that
circle is spent for good. The quantity the friend resolves away is the relative
phase between the $\lvert 0\rangle$ and $\lvert 1\rangle$ components, and moving up
to $\mathbb{C}^{4}$ \emph{promotes} it: the single-qubit relative phase becomes the
relative phase between $\lvert 0\rangle\lvert F_{0}\rangle$ and
$\lvert 1\rangle\lvert F_{1}\rangle$, which lives in the recoverable $S^{3}$ fibre
(itself a circle bundle over $S^{2}$, so the promotion adds a two-sphere of
relative-phase directions to the lone gauge circle). Gauge phase at one level
becomes physical, interferable phase one level up---which is why Wigner can run
the interference experiment the friend cannot. But Wigner now carries his own
fresh global-phase $S^{1}$ at the top of the joint state: his own irreducible
self-ignorance. Ascending the tower does not eliminate the self-reference, only
displaces it one level higher, and the regress up the chain is precisely the
redundancy ladder of the previous subsection. This is the framework's bump in the
carpet: each observer sits atop a Hopf tower with a single gauge circle he can
never see around, and a measurement relocates the next observer's blind circle
into his own recoverable fibre without ever closing the regress.

Two cautions close the section. First, the recoverability Wigner
enjoys is the low-redundancy limit: for a macroscopic, high-$R$ friend the
reversal is exponentially hard and the fact is objective for all practical
purposes, so Wigner's ``no record'' is the $R\!\approx\!1$ idealisation of a
graded boundary, not a free lunch. Second, naming redundancy as the grading
parameter of \textsc{fanout}, and the phase-promotion reading itself, are
syntheses of the framework's parts: each rests on established pieces (the
clonability criterion, quantum Darwinism, the tensor-product promotion of phase)
but is assembled here, and is offered as the cleanest statement of why the two
observers never disagree about anything both can check.
% =====================================================================
%  SPLICE: "Where the chain goes next" -- forward-looking coda for
%  sec:wigners-friend, pointing to sec:sedenions-generations.
%  Created 2026-08-27 per JS.
%
%  THE POINT (JS, 2026-08-27). The Wigner's-friend charts stop at
%  Wigner because the Hopf tower runs out at the octonions, while the
%  text says the chain of observers obviously continues. The sedenions
%  show what the continuation adds: the observer of Wigner (fourth
%  wire, C (x) S) brings a NEW piece of structure -- the finite S_3 of
%  Brown's theorem, the generation label -- and the observer of the
%  observer of Wigner (fifth wire, the trigintaduonions) brings nothing
%  new in kind. This coda records that, so the section's termination at
%  Wigner reads as a fact about the tower rather than an economy of the
%  exposition.
%
%  WHAT THE ALGEBRA ACTUALLY SAYS (verified 2026-08-27, notebook T --
%  notebook_T_trigintaduonion_automorphisms_2026-08-27.py):
%    - Aut(S) = G_2 x S_3 (Brown 1967), DIRECT product; the new S_3 is
%      generated by the doubling involution and Brown's explicit
%      order-three map psi, which fixes 1 and the doubling unit and
%      rotates every (a, a u) plane, a imaginary, by a third of a turn.
%    - psi permutes the three octonion HALVINGS O, psi O, psi^2 O of S
%      (each doubled by the same unit; meeting only in R). It does NOT
%      permute the three copies through the quaternion core -- those are
%      cycled by a lifted G_2 element (the colour-line three-cycle); see
%      the sedenion section's two-triples remark and the corrections
%      ledger of the same date.
%    - Brown: Aut(A_{n-1}) x S_3 = Aut(A_n) with EQUALITY for n = 4, 5, 6
%      (as quoted in gresnigtgourlayvarma2023three, sec. 4.3); Eakin and
%      Sathaye prove Brown's conjecture in general (VERIFY-CITE for the
%      exact general form -- their step is Aut(A_{n+1}) = Aut(A_n) x G,
%      G in {Z/2, S_3}; over R with the standard parameters the branch is
%      S_3 at every rung >= 4, which is what notebook T checks directly
%      at rungs 4, 5, 6). Derivation algebra g_2 from the octonions on
%      (Schafer 1954; notebook T: dim Der = 14 at n = 3, 4, 5).
%  So "nothing new at the fifth wire" is stated PRECISELY: the fifth
%  doubling adds another finite S_3 and nothing else -- no connected
%  symmetry, no fibration; the increment is the same at every rung
%  thereafter. Not "silent": REPETITIVE. The prose says so.
%
%  STATUS: the algebra is ESTABLISHED (Brown; Eakin--Sathaye; Schafer)
%  and directly checked (notebook T). The reading of the repeated finite
%  factor as what the redundancy ladder consumes is a RHYME, flagged.
%  The reading of the fourth rung's discrete orientation as the
%  generation label is the sedenion section's CANDIDATE identification,
%  pointed to, not re-argued here.
%
%  PLACEMENT: a new \subsection at the END of sec:wigners-friend, after
%  the "Two cautions" paragraph. Because that paragraph reads "Two
%  cautions close the section.", the patch also changes that one word:
%  "close the section" -> "close the argument". Both anchors
%  uniqueness-asserted (the July lesson: anchors that also occur in
%  comments must be asserted unique).
%
%  CROSS-REFS OUT (all exist in the synced source except the two
%  MAP-LABELs):
%    sec:hopf-justification   -- the "two climbs" paragraph (length vs
%                                termination); this coda sharpens it
%    sec:wf-redundancy         -- the redundancy ladder (this file)
%    sec:wf-phase-promotion    -- phase promotion (this file)
%    sec:sedenions-generations -- the sedenion section (2026-08-27 rev.)
%    prop:termination          -- Hurwitz/Adams termination
%  MAP-LABELs:
%    app:notebook-record       -- the notebook-record appendix
%                                (appendices/appendix_notebook_record_
%                                2026-07-11.tex); substitute its label.
%  NEW KEYS (RIS delivered: cayley_dickson_automorphism_refs_2026-08-27.ris):
%    eakinsathaye1990automorphisms, schafer1954cayleydickson
%  IN-LIBRARY: brown1967generalized.
%
%  OPTIONAL CONSUMER-POINTER (not patched; one clause for the "two
%  climbs" paragraph of hopf_justification_section_2026-06-08.tex, after
%  "that chain has no top --- it is the bump in the carpet, unbounded in
%  length."):
%      "(unbounded in length, and --- as \S\ref{sec:wf-chain-type}
%      records --- bounded in type: after the fourth observer nothing
%      new in kind arrives)"
%
%  EDITED 2026-08-30 (per JS): one sentence added to the fourth-wire
%  paragraph anchoring the flavour sector's irremovable phase (the
%  CP-phase recognition) on this rung, with the development at
%  sec:sedenions-mixing (2026-08-30 revision). No new keys here.
%  British English; pure ASCII; \autocite; no bullets; spaced dashes.
% =====================================================================
% SPLICE-BEGIN wf-chain-type 2026-08-27

\subsection{Where the chain goes next: one more rung, then copies}
\label{sec:wf-chain-type}

The regress just described is unbounded in length, and nothing here shortens
it: the observer of Wigner carries a gauge circle of his own, and so does the
observer of that observer, as \S\ref{sec:hopf-justification} insists when it
separates the endlessness of the chain from the termination of the tower. But
the algebra says how much of the endless chain adds structure of a new kind,
and the answer is exact: one more step, and then none.

The count runs as follows. The triple system, friend, Wigner is three wires,
and three wires are the three Hopf rungs $\mathbb{C}$, $\mathbb{H}$,
$\mathbb{O}$ --- the completed link of the earlier caveat --- which is why the
charts above stop at Wigner: the chart of whoever looks in on the triple,
$S^{15}\to S^{8}$, is the last fibration the division algebras supply
(Proposition~\ref{prop:termination}). The observer of Wigner is a fourth wire,
and the fourth wire is the sedenions,
$(\mathbb{C}^{2})^{\otimes 4} = \mathbb{C}\otimes\mathbb{S}$
(\S\ref{sec:sedenions-generations}). No fibration comes with it --- the
sedenions have zero divisors and admit no map $S^{31}\to S^{16}$ --- so the
fourth observer gains no recoverable fibre of the tower's kind: the continuous
promotion of relative phase into a larger fibre, the mechanism of the previous
subsection, has no fourth instance. What the doubling adds instead is finite.
Brown's theorem gives $\mathrm{Aut}(\mathbb{S}) = G_{2}\times S_{3}$
\autocite{brown1967generalized}: the identity component is unchanged from the
octonions, and the entire gain is a component group, generated by an
order-three automorphism that fixes the new doubling direction and rotates
every imaginary direction of the completed link into its copy by a third of a
turn. The relative orientation of the link and its copy, which at the Hopf
rungs was a continuous phase promoted into a fibre, survives at the fourth rung
only as a discrete choice among three, and that choice is the generation label
(\S\ref{sec:sedenions-generations}). The fourth observer, then, sees no new
force and no new fibre; he sees the same link three ways, and cannot gauge the
difference, because a component group has no Lie algebra. His rung is also where
the flavour sector's one irremovable phase finds its anchor: three generations
are the threshold at which the mixing matrix first carries a phase no convention
removes, and \S\ref{sec:sedenions-mixing} reads that phase --- through the
paper's negativity--contextuality chain --- as this wire's contextual cargo,
real only from one rung above.

The observer of the observer of Wigner is a fifth wire, the trigintaduonions,
and here the tower has nothing new to say. Brown proved that the next two
doublings each add exactly one further $S_{3}$ and nothing else, Eakin and
Sathaye that Brown's pattern continues at every rung, and the derivation
algebra --- the Lie algebra of the identity component --- is $\mathfrak{g}_{2}$
from the octonions onward
\autocite{brown1967generalized,eakinsathaye1990automorphisms,schafer1954cayleydickson};
the fifth and sixth rungs are also checked directly in the notebook record
(Appendix~\ref{app:notebook-record}). So the fifth observer's doubling repeats
the fourth's gift --- another third-of-a-turn symmetry between the four-wire
structure and its copy, commuting with everything below --- and adds no
connected symmetry and no fibration. Nothing new in kind arrives at the fifth
wire or at any wire after it. The increments of the tower are, in order, a
force, a force, a label, and then the same finite permutation of copies at
every step. That is the precise sense in which the chain ``obviously
continues'': it continues, and after the fourth observer it has nothing further
to announce. The framework reads the repeated factor as what the redundancy
ladder of \S\ref{sec:wf-redundancy} consumes --- a permutation symmetry among
copies, which is the freedom a record has and a type does not --- and records
that reading as a structural resonance rather than a derivation.

The two climbs of \S\ref{sec:hopf-justification} are thereby sharpened. The
chain of observers is unbounded in length and bounded in type: three forces by
the third wire, one family label by the fourth, copies thereafter. The observer
of Wigner is the last observer whose arrival changes what there is.

% SPLICE-END wf-chain-type 2026-08-27

% =====================================================================
%  FIGURE: the Wigner's-friend chain mapped onto the Cayley--Dickson
%  ladder. Created 2026-08-30 per JS: the Hopf-fibration view of
%  Wigner's friend has been load-bearing for months without being drawn;
%  the fibrations themselves are DEFERRED to Szangolies's introduction
%  (szangolies2025standardmodel, in library), and this figure carries
%  only the mapping.
%
%  CONTENT (left to right, one panel per chart, staircase rising):
%    Panel 1: F's chart of S            -- S^3 -> S^2, fibre S^1,  C, one wire
%    Panel 2: W's chart of {S,F}        -- S^7 -> S^4, fibre S^3,  H, two wires
%    Panel 3: O_W's chart of {S,F,W}    -- S^15 -> S^8, fibre S^7, O, three wires
%             (the last fibration)
%    Panel 4: the next observer's chart of {S,F,W,O_W} -- C^16 = C (x) S,
%             no fibration (S^31 -/-> S^16); the doubling's S_3 badge:
%             the generation label, a third of a turn.
%  Each describer perches on his sealed box with a dashed circle over
%  his head (his own gauge circle -- the self-referential ignorance he
%  cannot see around); dashed promotion arrows carry each circle into
%  the next chart's fibre ("gauge here -> interferable one rung up",
%  sec:wf-phase-promotion); the third arrow terminates at "no fourth
%  fibration", with what survives being the discrete third-of-a-turn
%  (sec:wf-chain-type, sec:sedenions-generations). The regress
%  continues to the right: copies only.
%
%  PREAMBLE REQUIREMENT (confirm in main.tex; add if absent):
%      \usepackage{tikz}
%      \usetikzlibrary{arrows.meta}
%  PLACEMENT: \input by patch_wigner_hopf_ladder_fig_2026-08-30.py,
%  after the two-idealisations paragraph of sec:wigners-friend, with
%  the accompanying deferral paragraph (in the patch).
%  CROSS-REFS OUT (all real): sec:wf-phase-promotion,
%  sec:wf-chain-type, sec:sedenions-generations.
%  KEY: szangolies2025standardmodel (in library).
%  British English; pure ASCII; \autocite in the caption.
% =====================================================================
\begin{figure}[t]
\centering
\resizebox{\textwidth}{!}{%
\begin{tikzpicture}[>=Stealth, line cap=round, line join=round,
    every node/.style={font=\scriptsize},
    blind/.style={densely dashed, draw=black!75},
    promote/.style={->, densely dashed, black!60, semithick},
    sight/.style={->, dotted, black!55}]
% ---- stick-figure pic --------------------------------------------------
\tikzset{stick/.pic={
  \draw[pic actions] (0,0.44) circle (0.085);
  \draw[pic actions] (0,0.355) -- (0,0.10);
  \draw[pic actions] (-0.12,0.20) -- (0,0.29) -- (0.12,0.20);
  \draw[pic actions] (-0.10,-0.10) -- (0,0.10) -- (0.10,-0.10);
}}
% ---- the staircase (Cayley--Dickson ladder) ---------------------------
\draw[black!35, thick]
  (-0.35,-0.55) -- (3.27,-0.55) -- (3.27,0.25) -- (6.89,0.25)
  -- (6.89,1.05) -- (10.51,1.05) -- (10.51,1.85) -- (13.60,1.85)
  -- (13.60,2.65) -- (16.20,2.65);
\node[anchor=north] at (1.25,-0.60) {$\mathbb{C}$ --- one wire};
\node[anchor=north] at (4.87,0.20)  {$\mathbb{H}$ --- two wires};
\node[anchor=north] at (8.49,1.00)  {$\mathbb{O}$ --- three wires};
\node[anchor=north] at (12.11,1.80) {$\mathbb{S}$ --- four wires};
\node[anchor=north, black!55] at (14.45,2.60) {$\mathbb{T},\dots$ --- five wires on};
% =======================================================================
% Panel 1: F's chart of S
% =======================================================================
\draw[rounded corners=2pt, fill=black!5, draw=black!40] (0,0) rectangle (2.5,1.32);
\filldraw[black!70] (1.25,0.72) circle (0.075);
\node[font=\tiny] at (1.25,0.45) {$S$};
\pic[black] at (2.05,1.42) {stick};
\node[anchor=west] at (2.22,1.62) {$F$};
\draw[sight] (1.95,1.60) -- (1.45,1.02);
\draw[blind] (2.05,2.14) circle (0.11);
\node[anchor=east, align=right, font=\tiny, black!75] at (1.82,2.14)
  {his own gauge circle:\\ the phase he cannot\\ see around};
\node[anchor=south west, font=\tiny, black!75] at (0.02,1.34) {fibre $S^{1}$};
\node[anchor=north] at (1.25,-0.06) {$S^{3}\to S^{2}$};
% =======================================================================
% Panel 2: W's chart of {S,F}
% =======================================================================
\draw[rounded corners=2pt, fill=black!5, draw=black!40] (3.62,0.80) rectangle (6.12,2.12);
\filldraw[black!70] (4.07,1.52) circle (0.075);
\node[font=\tiny] at (4.07,1.25) {$S$};
\pic[black!55, scale=0.85] at (4.85,1.16) {stick};
\node[font=\tiny, black!55] at (4.85,0.94) {$F$};
\pic[black] at (5.67,2.22) {stick};
\node[anchor=west] at (5.84,2.42) {$W$};
\draw[sight] (5.57,2.40) -- (5.05,1.82);
\draw[blind] (5.67,2.94) circle (0.11);
\node[anchor=south west, font=\tiny, black!75] at (3.64,2.14) {fibre $S^{3}$};
\node[anchor=north] at (4.87,0.74) {$S^{7}\to S^{4}$};
\draw[promote] (2.17,2.14) to[out=15,in=170] (3.60,2.24);
\node[anchor=south, align=center, font=\tiny, black!60] at (2.85,2.32)
  {gauge below,\\ interferable above};
% =======================================================================
% Panel 3: O_W's chart of {S,F,W}
% =======================================================================
\draw[rounded corners=2pt, fill=black!5, draw=black!40] (7.24,1.60) rectangle (9.74,2.92);
\filldraw[black!70] (7.69,2.32) circle (0.075);
\node[font=\tiny] at (7.69,2.05) {$S$};
\pic[black!55, scale=0.85] at (8.35,1.96) {stick};
\node[font=\tiny, black!55] at (8.35,1.74) {$F$};
\pic[black!55, scale=0.85] at (8.95,1.96) {stick};
\node[font=\tiny, black!55] at (8.95,1.74) {$W$};
\pic[black] at (9.29,3.02) {stick};
\node[anchor=west] at (9.46,3.22) {$O_{W}$};
\draw[sight] (9.19,3.20) -- (8.65,2.62);
\draw[blind] (9.29,3.74) circle (0.11);
\node[anchor=south west, font=\tiny, black!75] at (7.26,2.94) {fibre $S^{7}$};
\node[anchor=north, align=center] at (8.49,1.54) {$S^{15}\to S^{8}$\\[-1pt]
  {\tiny the last fibration}};
\draw[promote] (5.79,2.94) to[out=15,in=170] (7.22,3.04);
% =======================================================================
% Panel 4: the next chart of {S,F,W,O_W} -- sedenions, no fibration
% =======================================================================
\draw[rounded corners=2pt, fill=black!5, draw=black!40] (10.86,2.40) rectangle (13.36,3.72);
\filldraw[black!70] (11.24,3.12) circle (0.075);
\node[font=\tiny] at (11.24,2.85) {$S$};
\pic[black!55, scale=0.80] at (11.78,2.78) {stick};
\node[font=\tiny, black!55] at (11.78,2.56) {$F$};
\pic[black!55, scale=0.80] at (12.28,2.78) {stick};
\node[font=\tiny, black!55] at (12.28,2.56) {$W$};
\pic[black!55, scale=0.80] at (12.78,2.78) {stick};
\node[font=\tiny, black!55] at (12.78,2.56) {$O_{W}$};
\pic[black!45] at (14.35,2.78) {stick};
\draw[sight] (14.18,3.02) -- (13.42,3.12);
\draw[blind] (14.35,3.50) circle (0.11);
\node[anchor=west, align=left, font=\tiny, black!60] at (14.58,3.48)
  {the next observer,\\ and all after:\\ copies only};
\node[anchor=north] at (12.11,2.34) {$\mathbb{C}^{16}=\mathbb{C}\otimes\mathbb{S}$};
% the broken third promotion
\draw[promote] (9.41,3.74) to[out=15,in=185] (10.68,3.86);
\node[font=\scriptsize, black!60] at (10.80,3.92) {$\times$};
\node[anchor=south west, font=\tiny, black!75] at (10.88,3.74)
  {no fourth fibration: $S^{31}\not\to S^{16}$};
% the S3 badge: three dots at 120 degrees with a rotation arrow
\draw[black!75] (11.45,4.30) circle (0.17);
\filldraw[black!75] (11.45,4.47) circle (0.028);
\filldraw[black!75] (11.303,4.215) circle (0.028);
\filldraw[black!75] (11.597,4.215) circle (0.028);
\draw[->, black!75] (11.67,4.36) arc[start angle=20, end angle=100, radius=0.235];
\node[anchor=west, align=left, font=\tiny, black!75] at (11.70,4.36)
  {$S_{3}$: the generation label\\ (a third of a turn)};
\draw[promote, black!70] (10.86,3.96) to[out=60,in=200] (11.24,4.22);
\end{tikzpicture}%
}
\caption{The Wigner's-friend chain mapped onto the Cayley--Dickson ladder. Each
observer holds a chart of the sealed box below him, one rung further up the
ladder: the friend's chart of the system is the complex Hopf bundle, Wigner's
chart of the sealed laboratory the quaternionic, and the chart of the observer
of Wigner --- of the completed triple --- the octonionic, the tower's last
fibration. Every chart-holder carries a dashed circle over his head, his own
gauge phase: the self-referential ignorance he cannot see around. Measurement
promotes each circle into the next chart's fibre, where what was gauge below is
interferable above (\S\ref{sec:wf-phase-promotion}); the third promotion has no
fibration to land in --- the sedenions have zero divisors and no map
$S^{31}\to S^{16}$ --- and what survives of it is discrete: the doubling's
$S_{3}$, a third of a turn rotating the link into its copy, the generation
label (\S\ref{sec:wf-chain-type}, \S\ref{sec:sedenions-generations}). The
regress of circles continues unbounded to the right, but from the fifth wire
onward it adds copies only. The friend drawn as a single qubit is the
redundancy-one idealisation discussed in the text; for the fibrations
themselves, deferred rather than repeated here, see the introduction in
\autocite{szangolies2025standardmodel}.}
\label{fig:wf-hopf-ladder}
\end{figure}

%      % EDITED 2026-10-29 scripts/patch_continuity_tower_premise_2026-10-29T1500.py
% EDITED 2026-09-20 scripts/patch_born_cost_ledger_v2_2026-09-20T0900.py
% EDITED 2026-09-15 scripts/patch_value_continuum_declared_cencov_2026-09-15T1540.py
% EDITED 2026-09-15 scripts/patch_continuity_connected_state_space_2026-09-15T1430.py
% =====================================================================
%  SECTION: Derivation of the quantum postulates (the assembly chapter).
%  AUDITED 2026-07-06 (final pass, sixth stop).
%
%  PLACEMENT (reordered chain, JS 2026-07-06): position 5 --
%      \input{wigners_friend_in_the_hopf_picture_2026-06-13}
%      \input{derivation_of_the_quantum_postulates}
%      \input{division_algebras_section_2026-07-03}
%  Postulate 6 hands off to division-algebras, its immediate successor.
%
%  LABEL RESOLUTIONS (2026-07-06): the old placeholder
%  sec:complex-amplitudes is now minted at its source (self-reference's
%  negativity subsection), so the ref here resolves unchanged.
%  sec:born-rule COLLIDED with the label minted 2026-07-06 on
%  self-reference's Born-from-fibre-privacy subsection; resolution: two
%  routes, two labels -- Postulate 3 here is renamed
%  sec:born-invariance (the U(1)/lowest-degree geometric route) and
%  cross-stitched to sec:born-rule (the record-theoretic route), their
%  agreement flagged as over-determination. sec:local-tomography also
%  minted at source for Postulate 6's operational clause.
%
%  AUDIT RECORD: the section's status discipline was PRESCIENT -- the
%  continuity arrow was already "owed ... enters as a principle" (the
%  exact voice JS mandated for L2), and both of its confessed debts
%  have since been minted as numbered problems; the closing now cites
%  them (op:limit-realisability, op:tsirelson-geometry): "named, and
%  numbered". Postulate 6 now carries BOTH composition duty-holders
%  (Hurwitz multiplicativity + local tomography), matching the
%  continuity section's caveat exactly. Stone's theorem cited
%  (stone1932one, RIS delivered); time-as-operation-count forward
%  flag added; \textquote made portable; opener formalised.
%
%  PROVENANCE NOTE (relocated from the body, 2026-07-06, JS's words):
%  "Using Hopf fibrations to get the Born rule is something Claude AI
%  came up with on its own without my prompting." Preserved here as the
%  historical record. DISCLOSURE STATUS (JS, 2026-07-06): AI use is
%  already acknowledged in the second paragraph of the paper's
%  introduction; the remaining recommendation is only that, at
%  submission, the same fact be echoed in whatever formal
%  Declarations / Acknowledgements slot the target journal's template
%  mandates -- several publishers now require it there specifically,
%  in addition to any in-text mention.
%
%  British English; \autocite; no bullets in prose.
% =====================================================================

% ADDED 2026-07-12: closing paragraph "The postulates in module
% coordinates" -- the Clifford-module reading (skeleton + two
% supplements), with the explanatory arrow guarded per JS: the field
% comes from the Shannon/additivity/commutativity step (self-ignorance),
% and the left-action composability result is recorded as that same
% requirement in operator clothing, not an independent selection of C.
\section{Derivation of the quantum postulates}
\label{sec:postulates}
Enough pieces of the jigsaw are now assembled to derive the postulates of quantum mechanics.

The Dirac--von Neumann postulates are ordinarily laid down as independent
assumptions. In this framework they are not assumptions but a single chain: each
follows from the Inside/Outside Equivalence Principle by one move from the
preceding step. We assemble the chain here and mark the status of each link as
proved, structural, or still owed, so that no step claims more than it has earned.

The chain begins below the postulates. An embedded observer, required by the
equivalence principle to describe a totality that includes itself, must assign a
probability to a proposition about its own future state. The Lawvere--Yanofsky
diagonal \autocite{Yanofsky2003} and Popper's predictor argument
\autocite{Popper1950a, Popper1950b} show that no consistent classical assignment
exists: the self-referential proposition is undecidable to the observer before
the fact. Representing the result of that undecidable self-prediction forces a
signed, non-classical probability, which carries the observer's irreducible
self-ignorance. From this single obstruction the postulates follow.

\paragraph{Postulate 1 (complex Hilbert space).} Two independent steps fix the state space. The
field is forced first: a signed probability, read through the Shannon information measure and
analytically continued, yields a complex amplitude $S=-\ln\lvert p\rvert-\mathrm{i}\theta$,
whose phase $\theta$ is the self-referential ignorance (Section~\ref{sec:complex-amplitudes}).
The metric is forced second, and separately: the principle of reversible continuity
(\S\ref{sec:continuity}) --- that a continuous reversible transformation connects any two pure
states --- selects the $L^{2}$ norm uniquely. For every $\ell^{p}$ with $p\neq 2$ the linear
isometries are generalised permutations \autocite{Banach1932, Lamperti1958}, a group that cannot
connect a basis state to a superposition; only at $p=2$ is the norm-preserving group connected
and transitive on the sphere \autocite{Aaronson2004h}, and Hardy's operational $K=N^{2}$
counting reaches the same quadratic norm by a second route from the same continuity premise
\autocite{Hardy2001a}, and \v{C}encov's theorem with the every-basis lemma reaches it from the
value continuum alone (\S\ref{sec:cencov-route}). The state manifold is therefore the unit
sphere of the Hermitian $2$-norm in $\mathbb{C}^{n}$, taken modulo the global phase
(Section~\ref{sec:l2-metric-fibre}). Placing the metric here, before the Born rule rather than
after it, is deliberate: the Born derivation presupposes the round sphere, so the metric must be
anchored independently to avoid circularity. Two bounds belong with this step. Reversible
continuity fixes the metric, not the field---real and quaternionic quantum theory satisfy it
equally---so the restriction to $\mathbb{C}$ rests on the Shannon step, not on continuity. And
deriving reversible continuity itself from self-reference is the one arrow the chain still owes
--- now priced exactly: one open premise, that every transformation of the regime is a phase of an inside structure,
Open Problem~\ref{op:limit-realisability}, with everything downstream of it proved
(\S\ref{sec:continuity}). Until it is closed, continuity enters as a principle, and the weight
of the $L^{2}$ conclusion is carried jointly with the isometry and counting arguments and with
the statistical route.

\paragraph{Postulate 2 (Hermitian observables).}
A measurement result must be communicable: it must be copyable from one observer
to another, the \textsc{copy} operation of the diagrammatic calculus
\autocite{CoeckeKissinger2017}. The imaginary part of the amplitude is the
private self-referential phase and cannot be broadcast; only the real,
gauge-invariant content can. Observables whose spectra are real and therefore
communicable are exactly the self-adjoint operators. Hermiticity is thus the
condition that a quantity be the kind of thing a measurement can report.

\paragraph{Postulate 3 (Born rule).}
\label{sec:born-invariance}
Outcome probabilities must be invariant under the global phase, the $U(1)$
self-ignorance, realised geometrically as the $S^{1}$ fibre of the complex Hopf
fibration $S^{1}\to S^{3}\to S^{2}$ over the Bloch sphere
\autocite{Hopf1931, Urbantke2003}. There is no polynomial $U(1)$-invariant of degree one,
and the non-polynomial $\lvert\psi\rvert$ is excluded by the analyticity declared at
\S\ref{sec:complex-amplitudes}; the lowest-degree polynomial invariant is quadratic, so the
probability must be $p=\lvert\psi\rvert^{2}$. The squaring is forced
rather than merely permitted, and the argument holds for a single qubit, where
Gleason's theorem is silent. The Born quadratic is not a new structure: it is the
same $L^{2}$ form that fixed the metric in Postulate~1, now read as the
normalisation of the round sphere, and the fibre over which the phase is averaged
has constant length, so the phase is hidden uniformly
\autocite{GluckWarnerZiller1986}. This step is a structural argument rather than a
theorem in the strict sense, and is flagged as such. The record-theoretic derivation of the same
quadratic --- squared magnitudes from fibre privacy, Proposition~\ref{prop:f1-dictionary} --- is
given at \S\ref{sec:born-rule}, and the cost-ledger derivation --- the record as the closed loop of
a history and its return leg, the exponent the number of legs --- at \S\ref{sec:born-cost-ledger}:
three routes to one rule, each holding for a single qubit, the over-determination this framework
treats as evidence.

\paragraph{Postulate 4 (state reduction).}
Measurement is a \textsc{fanout}: the gauge-invariant relational content is copied
into a record, converting a private epistemological amplitude into a public
fact. For the observer who performs it, acquiring that fact is Bayesian updating,
and the projection postulate is its expression. The reduction is perspectival---a
fact relative to the observer in the sense of relational quantum mechanics
\autocite{Rovelli1996a}---and it is graded by redundancy. A single copy is a
relative fact, still reversible by an agent controlling the joint system; redundant
proliferation across many independent fragments makes the fact objective, which is
quantum Darwinism \autocite{zurek_2009} and the stable/relative-fact distinction of
Di~Biagio and Rovelli \autocite{DiBiagioRovelli2021}. ``Collapse'' and unitary
evolution are then the two charts---inside and outside the \textsc{fanout}
boundary---of one process, as the Wigner's-friend analysis shows and \S\ref{sec:fanout-boundary} formalises
(Section~\ref{sec:wigners-friend}); the irreversibility and its $kT\ln 2$ cost
\autocite{Landauer_1961, wigner_1961} accrue with redundancy rather than at the
single act. The reduction postulate is therefore recovered as a chart-dependent
description, not a fundamental dynamical law; the identification of redundancy as
the grading parameter is a synthesis of the framework's parts and is flagged as
such.

\paragraph{Postulate 5 (unitary evolution).}
Conservation of probability is preservation of the $L^{2}$ inner product, and the
inner-product-preserving linear maps are the unitaries $U(n)$---the same
isometry fact that fixed the metric in Postulate~1. Continuity in time --- operation count,
\S\ref{sec:space-of-computations}, a forward reference --- together with Stone's theorem
\autocite{stone1932one} gives $U(t)=\mathrm{e}^{-\mathrm{i}Ht}$, the Schr\"odinger
equation. Reversibility is here a corollary, not a premise: bare reversibility on
the amplitude space yields only the general linear group, and it is preservation
of the quadratic form that carves $U(n)$ out of it.

\paragraph{Postulate 6 (composite systems).}
Systems combine by the tensor product because the division algebra norm is
multiplicative \autocite{hurwitz1898composition}: composition preserves the $L^{2}$ structure
only if states multiply. The operational half closes from the other side: the locality of
tomography derived at \S\ref{sec:local-tomography} says product states determine composites,
which is the tensor product's signature. Cayley--Dickson doubling carries $\mathbb{C}$ to the next
rung, and the ascent of the Hilbert space is the ascent of the Hopf tower, complex
to quaternionic, used in the Wigner's-friend construction
\autocite{mosseri2001hopf}; the ladder itself is the business of the section that follows
(\S\ref{sec:division-algebras}).

\paragraph{Consequence (the Tsirelson bound).}
The postulates fix not only the form of the theory but the strength of its
correlations. In the two-qubit ball the maximal entanglement sits at the centre,
and a Popescu--Rohrlich box \autocite{PopescuRohrlich1994} would require a point
beyond it; the CHSH correlator is capped at $2\sqrt{2}$
\autocite{Cirelson1980QuantumInequality, VerstraeteWolf2002, Held2021}, with the value recovered
at the centre and its general derivation still owed
(Section~\ref{sec:wf-reconciliation}).

\medskip
\noindent The result is that the six postulates are one chain, not six premises.
Of the links, the field and the Born quadratic rest on the framework's signed
probability and $U(1)$-invariance arguments; the $L^{2}$ metric, the unitary group
and the tensor product are over-determined by reversible continuity, the isometry
classification and Hurwitz composition; state reduction is recovered as the
chart-relative face of \textsc{fanout}. The two outstanding debts are the
derivation of reversible continuity from self-reference and the forcing of the Tsirelson value as
a theorem rather than a recovered number --- named, and numbered:
Open Problems~\ref{op:limit-realisability} and~\ref{op:tsirelson-geometry} --- so that a reader
can see exactly where the chain is complete and where it is not.

\paragraph{The postulates in module coordinates.}
One reading of the chain deserves a closing statement, because it positions the derivation
against a mathematical skeleton the reader may know well, and because stating it carelessly would
invert the chain's explanatory arrow. In the coordinates of the division-algebra programmes
(\S\ref{sec:algebraic-dictionary}), complex Hilbert-space quantum mechanics is real Clifford
module theory plus exactly two supplements. The skeleton is Clifford throughout: the operator
algebra of $n$ qubits is the Clifford algebra of $2n$ generators with the state space its own
spinor module, half-integer spin and fermionic statistics are the module's double cover and its
graded tensor, and the classically simulable stabiliser core is the algebra's combinatorics. What
the skeleton does not contain is, first, the choice of complex structure --- the polarisation $J$
that turns a real module into a complex Hilbert space, which standard quantum theory and the
algebraic programmes alike \emph{choose} --- and, second, the Born weights on the fibre. Those two
supplements are precisely what this section has derived, and derived from self-ignorance rather
than from algebra: the field from the Shannon step (information about one's own undecidable state
must still be information --- it must add, addition commutes, and so the amplitude algebra is
commutative, excluding the quaternionic bookkeeping an observer could never balance), and the
quadratic from fibre privacy. The algebraic echo --- that the algebra's own left actions compose
only along a commutative line, the composability result recorded with the dictionary --- is not
an independent selection of $\mathbb{C}$: composition of left actions fails off that line exactly
because the doubled slot carries the \emph{reversed} product, and order-dependence is the same
disease in operator clothing that non-additivity is in information clothing. The mathematics is
the epistemic argument's shadow, welcome as consistency and as a bridge to the module-side
literature (Appendix~\ref{app:winding-number} walks the bridge for charge), but the origin of the
complex amplitude is the observer's ignorance of their own state, and the arrow runs from there
to the algebra, never the reverse.
%      % =====================================================================
%  SECTION: The Division-Algebra Ladder
%  Drafted 2026-07-03. One of the three foundations sections
%  (sec:self-reference, sec:division-algebras, sec:fanout-boundary);
%  this one resolves the dangling \S\ref{sec:division-algebras} in the
%  Hopf-justification section, whose header names it "the
%  Cayley--Dickson / Hurwitz step".
%
%  PLACEMENT in main.tex (REORDERED per JS 2026-07-06): position 6,
%  opening Act II ("why stop at C?") right before its first consumer --
%      \input{self_reference_section_2026-07-04}
%      \input{continuity_section_2026-07-04}
%      \input{L2_metric_global_phase_section_2026_06_07}
%      \input{wigners_friend_in_the_hopf_picture_2026-06-13}
%      \input{derivation_of_the_quantum_postulates}
%      \input{division_algebras_section_2026-07-03}   <-- moved here
%      \input{spacetime_from_non-computability}       <-- first consumer
%                                                          (S^7 exclusion)
%      \input{hopf_justification_section_2026-06-08}
%  The section is position-portable: all adjacency phrasing was
%  retargeted to explicit \S\ref on 2026-07-06. The S^7 hand-off
%  sentence to sec:space-of-computations WRITTEN 2026-07-06 (spacetime
%  file read; the closing now hands first to its immediate consumer,
%  then to hopf-justification). Upload verifications COMPLETE: L2 and
%  quantum-postulates do not consume the ladder; WF's references are
%  forward-voiced or self-carrying.
%
%  DIVISION OF LABOUR (deliberate, to avoid overlap):
%   - sec:self-reference (undrafted): self-reference -> complex
%     amplitudes and the Born rule; why not a real base. This section
%     TAKES that as given and starts at C.
%   - THIS section: the algebra -- the doubling construction, the
%     ledger of losses-as-physics, the internal-algebras-not-a-new-QM
%     scoping (Adler pre-empt), the two appearances of the ladder
%     (register / measurement climb, stated lightly), and Hurwitz as
%     the algebraic termination with the committed zero-divisor
%     reading.
%   - sec:hopf-justification: the geometry (Adams, parallelisability,
%     the fibrations), the join of the two classifications, the
%     entanglement-sensitivity literature, the C-direction honest debt,
%     and the epistemic reading of termination.
%
%  COMMITTED READINGS USED (sources re-read 2026-07-03):
%   - Loss ledger and level dictionary: cayley_dickson_qubits_for_review
%     sec.5 ("each step loses a symmetry, and each lost symmetry is
%     identified with a physical content").
%   - Termination reading: grounding_argument_gauge_group ("the
%     sedenions have zero divisors and cannot support consistent
%     information flow around loops") -- leads; norm-composition-as-
%     probability given as its amplitude-level restatement.
%   - Internal-not-base scoping: cayley_dickson_qubits_for_review Q2 and
%     sec.3.2 caveat (Adler objection; C-qubit treated geometrically).
%
%  HONEST DEBT: ladder proved closed at the top (Hurwitz), derived at
%  the floor. The rung-compulsion revisit-trigger FIRED 2026-07-06:
%  sec:self-reference and op:sigma-functor now exist, so the thesis is
%  minted as op:ladder-compulsion (with the CD-conjugation-as-sigma
%  candidate route), inheriting op:sigma-functor's status.
%
%  AUDITED 2026-07-06 (final pass, third stop): all four citations
%  used; all cross-refs live; CD facts, dims, alternativity, Hurwitz
%  statement verified; Adler subsection delivers everything the
%  register's compressed pointer promises; file pure ASCII; no bare
%  "frame". Fixes: opening rewritten for both predecessors (continuity
%  was inserted between after drafting); "associator's non-vanishing is
%  mass itself" recalibrated to the register's debt-against-selection
%  defusal; the O->S ledger row and Hurwitz subsection reconciled with
%  the fourth-qubit forward remark ("the terminus closes the ladder of
%  forces; what survives past it is, at most, flavour").
%
%  CROSS-REFS: sec:self-reference [sibling dangle -- resolves when that
%  section is drafted], sec:hopf-justification, sec:hopf-information-bound,
%  sec:register, sec:entanglement-classes, sec:generations, sec:masses,
%  sec:wigners-friend. All others exist.
%  CITATIONS: existing keys baez2002octonions, hurwitz1898composition,
%  furey2018three. ONE new key: adler1995quaternionic (RIS in
%  division_algebras_refs_2026-07-03.ris).
%  British English; \autocite; no bullets in prose; the ledger is a
%  table because the comparison is the argument.
% =====================================================================

\section{The Division-Algebra Ladder}
\label{sec:division-algebras}

The foundations behind us end at amplitudes twice over: natively, where the self-referential
argument forced the description of a system into the system's own state space and priced
consistency at a number system whose parts can cancel --- complex amplitudes, and the Born rule
that squares them (\S\ref{sec:self-reference}) --- and independently, where the regime's
continuity delivers the same field through Hardy's theorem (\S\ref{sec:continuity}), the
convergence of the two branches being evidence rather than accident. This section asks the question
both routes leave open. Why
should the number system stop at $\mathbb{C}$? Ordinary quantum mechanics stops there by postulate;
a framework that has forbidden itself postulates must answer, and the answer is that the climb
neither stops at $\mathbb{C}$ nor continues forever. There are exactly two more rungs, each
purchased by surrendering an algebraic law and each paying for itself in physics, and the ladder is
closed at the top by a theorem of 1898. This section sets out the algebra of that ladder: the single
construction that climbs it, the ledger of what each rung costs and buys, and the theorem that ends
it. The geometry of the same ladder --- that each rung acts on a sphere, that exactly those spheres
are parallelisable, and that the algebraic and geometric classifications are one --- is the business
of the geometry section (\S\ref{sec:hopf-justification}).

\subsection{One construction climbs the whole ladder}

The rungs are not a miscellany; a single construction generates them all. Given an algebra with a
conjugation, the Cayley--Dickson construction builds its double on ordered pairs, with a twisted
multiplication in which the conjugate enters and a conjugation $(a,b)^{*} = (a^{*}, -b)$
\autocite{baez2002octonions}. (Conventions differ in the placement of the conjugations; all yield
isomorphic algebras.) Applied repeatedly it produces the sequence
$\mathbb{R} \to \mathbb{C} \to \mathbb{H} \to \mathbb{O} \to \mathbb{S}$, of real dimensions
$1, 2, 4, 8, 16$: the reals, the complex numbers, the quaternions, the octonions, the sedenions.
The ground floor is classical: $\mathbb{R}$ is ordered and positive-definite, the arithmetic of
classical probability. The first doubling is therefore the quantum move itself --- the passage from
probabilities to amplitudes that can interfere --- and it is not made here by choice: it is the
derivation of \S\ref{sec:self-reference}, already behind us. What this section must account for is
the rest of the climb.

\subsection{The ledger}

In mathematics the doublings are usually told as a story of degeneration: each algebra is worse
behaved than the last, and the sequence is studied for where it breaks. This framework reads the
same ledger in the opposite direction. Each rung is bought by surrendering a law, and each
surrendered law is exactly a piece of physics --- the losses are the assets, and the comparison is
sharp enough to be the argument:

\begin{center}
\begin{tabular}{llll}
\hline
doubling & dimension & what is lost & what the loss becomes \\
\hline
$\mathbb{R} \to \mathbb{C}$ & $1 \to 2$ & the order of the reals & phase and interference; amplitudes \\
 & & & that cancel; the Born rule (\S\ref{sec:self-reference}) \\
$\mathbb{C} \to \mathbb{H}$ & $2 \to 4$ & commutativity & non-commuting actions: weak isospin \\
 & & & $SU(2)$, the $B$ register (\S\ref{sec:register}) \\
$\mathbb{H} \to \mathbb{O}$ & $4 \to 8$ & associativity & the associator: colour, chirality, the \\
 & & & generations, and mass as debt \\
 & & & (\S\ref{sec:generations}, \S\ref{sec:masses}) \\
$\mathbb{O} \to \mathbb{S}$ & $8 \to 16$ & division itself & no further force: zero divisors appear \\
 & & (zero divisors) & and the ladder of sectors ends (below) \\
\hline
\end{tabular}
\end{center}

Walking the rows: the order lost at the first doubling is the possibility of interference --- parts
of a probability that can oppose one another --- which is where the derivation of \S\ref{sec:self-reference}
lives. The commutativity lost at the second is the existence of actions whose order matters; the
unit quaternions form the group $SU(2)$, and the framework houses weak isospin there, as the
internal algebra of the $B$ register whose order-sensitive, non-commuting dynamics the interaction sections develop
(\S\ref{sec:register}, \S\ref{sec:entanglement-classes}). The associativity lost at the third ---
the octonions retain the weaker property of alternativity \autocite{baez2002octonions} --- creates
the associator $[a,b,c] = (ab)c - a(bc)$, and the associator is where most of this paper's physics
is transacted: the stabiliser of a preferred complex direction inside the octonions' automorphisms
is the colour group (\S\ref{sec:hopf-justification}); the triality structure and its $K_3$
obstruction generate the three generations and the sterile channel (\S\ref{sec:generations}); and
the associator is the currency of mass --- the debt against the selected direction that
\S\ref{sec:masses} prices, owed by whatever action moves the selection and not by residence at
the rung, which is the defusal \S\ref{sec:register} makes explicit. The algebraic facts
in the third column are theorems; the identifications in the fourth are established in the sections
cited, each carrying its own epistemic flag there, and are not re-litigated here. What this table
asserts on its own account is only the correspondence's shape: one law surrendered per rung, one
sector of physics gained per law.

\subsection{Internal algebras, not a new quantum mechanics}

A reader versed in the quaternionic-quantum-mechanics literature will at this point raise a fair
objection: quantum theory over $\mathbb{H}$ as a foundation is known to be troubled --- the
tensor-product and composite-system difficulties set out by Adler and others
\autocite{adler1995quaternionic} --- and quantum theory over $\mathbb{O}$ founders earlier still,
since a non-associative algebra does not furnish a Hilbert space at all. The framework's position
must therefore be stated exactly, because it does not make the claim the objection targets. The base
of the physical Hilbert space here is $\mathbb{C}$ and remains $\mathbb{C}$; the higher algebras
enter as the \emph{internal} algebras of the register --- $\mathbb{H}$ for the $B$ qubit,
$\mathbb{O}$ for the $C$ qubit (\S\ref{sec:register}) --- in a tensor structure whose outer theory is
ordinary complex quantum mechanics. This is the same move by which the division algebras enter the
constructions of Furey and the related programmes without any proposal to rebase quantum mechanics
itself \autocite{furey2018three}. For the octonionic register the point is sharper still:
$\mathbb{O}^{2}$ is not a Hilbert space, and the framework accordingly treats the $C$ qubit not in
the Dirac--von Neumann manner but geometrically, as the sphere $S^{15}$ with its Hopf structure ---
which is precisely the form in which \S\ref{sec:hopf-justification} needs it. Nothing in this paper therefore
stands or falls with quaternionic or octonionic quantum mechanics as foundations; the ladder is
climbed by the registers' internal algebras, not by the theory's base field.

\subsection{The same ladder, twice}

The ladder appears in this framework in two guises, and it will help the reader to have them
distinguished now. Statically, it is the anatomy of a single relation: the three qubits of the
register take their values one rung apart --- $A$ over $\mathbb{C}$, $B$ over $\mathbb{H}$, $C$ over
$\mathbb{O}$ --- so that every particle carries the whole ladder as its internal structure
(\S\ref{sec:register}). Dynamically, it is the shape of measurement: a single qubit's pure states
fill $S^{3}$, an entangling measurement carries the joint state into $S^{7}$, and a third party into
$S^{15}$ --- the correspondence between qubit entanglement and the Hopf fibrations that
\S\ref{sec:hopf-justification} takes from a mature literature --- so that each act of measurement ascends one rung --- the climb the reader has already made,
narrated in \S\ref{sec:wigners-friend} --- and every completed relation is a terminated climb, which is
the reading \S\ref{sec:hopf-information-bound} gives the tower's endpoint. One ladder, seen twice:
as the anatomy of a relation and as the dynamics of relating. The framework regards this double
appearance not as a coincidence but as the same structure viewed statically and in motion, and the
sections cited carry that identification's weight.

\subsection{Hurwitz closes the ladder}

The construction itself would double forever; what stops it is a classification theorem. By
Hurwitz's theorem, the normed division algebras --- those in which the norm composes,
$\lvert ab\rvert = \lvert a\rvert\,\lvert b\rvert$, and no non-zero elements multiply to zero ---
are exactly $\mathbb{R}$, $\mathbb{C}$, $\mathbb{H}$ and $\mathbb{O}$
\autocite{hurwitz1898composition,baez2002octonions}. The next doubling, the sedenions, has zero
divisors: non-zero elements whose product vanishes. The framework's reading of this failure is the
one that matters physically. Zero divisors cannot support consistent information flow around loops
--- a cycle of compositions can silently annihilate what it carries --- and a framework in which
closed computational paths are load-bearing (\S\ref{sec:hopf-justification},
\S\ref{sec:space-of-computations}) cannot be built on an algebra that permits it. At the level of
amplitudes the same failure reads: the composition law
$\lvert ab\rvert = \lvert a\rvert\,\lvert b\rvert$, which carries probability multiplicatively
through every rung of the ladder, fails beyond the octonions, so that two non-vanishing amplitudes
could compose to nothing and probability would be destroyed by mere composition. The last algebra
fit to carry amplitudes is $\mathbb{O}$. One rung past the terminus is not entirely silent,
however: the sedenions' internal structure --- their automorphism group and their octonion
subalgebras --- is the recorded candidate arena for tagging the generations
(\S\ref{sec:interactions}), and it enters exactly as this section's scoping permits, internally
and never as a base. The terminus closes the ladder of forces; what survives past it is, at most,
flavour. This is the algebraic half of the termination; the
geometric half --- that no sphere above $S^{7}$ is parallelisable --- and the theorem joining the
two halves are the opening business of \S\ref{sec:hopf-justification}, and the epistemic reading of the same endpoint,
as a bound on the information a single relation can carry, is given in
\S\ref{sec:hopf-information-bound}.

\subsection{The honest debt}

As in the geometry section, we close by stating plainly what is proved, what is established, and what is
owed. Proved: the construction, the losses, and the termination are theorems --- the Cayley--Dickson
doubling and its ledger of surrendered laws are standard algebra, and Hurwitz's classification is
among the durable results of the subject. Established: the physical cashing of each rung --- the
Born rule at $\mathbb{C}$, weak isospin at $\mathbb{H}$, colour, generations and mass at
$\mathbb{O}$ --- is carried out in the sections cited above, each under its own epistemic flag. Owed,
and stated without disguise: the ladder as presented here is closed from above and derived at its
floor, but it is not yet derived rung by rung from below. Hurwitz forbids a further rung; the
argument of \S\ref{sec:self-reference} forces the first doubling, from the ordered arithmetic of
classical probability to amplitudes that interfere. That the same self-referential closure, applied
to a system whose amplitudes already take values at one rung, compels the next --- that relational
closure at one level forces precisely the Cayley--Dickson double of that level --- is the
framework's reading of the climb, and it is a thesis, not a theorem --- but the formal apparatus
against which it can be posed now exists, and we pose it.

\begin{openproblem}[The compulsion of the intermediate rungs]\label{op:ladder-compulsion}
Derive, within the categorical apparatus of Open Problem~\ref{op:sigma-functor}, that
self-referential closure at one rung compels the Cayley--Dickson double of that rung, so that the
intermediate steps $\mathbb{C} \to \mathbb{H} \to \mathbb{O}$ follow from the same principle
that forces the floor. The candidate route is the construction's own shape: the Cayley--Dickson
pair with its conjugation, $(a, b)^{*} = (a^{*}, -b)$, is an involution-equipped doubling, and the
route reads it as the algebraic shadow of adjoining the describer --- closure of level-$A$
descriptions under the inside/outside exchange requiring exactly the twisted pair. The problem
inherits the status of Open Problem~\ref{op:sigma-functor}, whose apparatus it presupposes.
\end{openproblem}

If the derivation lands, the whole ladder becomes a consequence of the same self-reference that
opened it. Until then, the honest statement is that the ladder's necessity is established at its
ceiling and at its floor, while the compulsion of the intermediate rungs is inherited from the
physics they are shown to carry rather than derived in advance. What this section hands on is, accordingly, an algebraic ladder with a floor, a ceiling and a
ledger. Its first consumer is immediate: \S\ref{sec:space-of-computations} spends the terminal
fact --- the non-associativity of $S^{7}$'s own multiplication --- as geometry, excluding the
seven-sphere as a ruler and fixing the dimension of space at three.
\S\ref{sec:hopf-justification} then supplies the geometry of the fibrations themselves, and the
join.

%  Postulate 6 hands off to division-algebras, its immediate successor.
%
%  LABEL RESOLUTIONS (2026-07-06): the old placeholder
%  sec:complex-amplitudes is now minted at its source (self-reference's
%  negativity subsection), so the ref here resolves unchanged.
%  sec:born-rule COLLIDED with the label minted 2026-07-06 on
%  self-reference's Born-from-fibre-privacy subsection; resolution: two
%  routes, two labels -- Postulate 3 here is renamed
%  sec:born-invariance (the U(1)/lowest-degree geometric route) and
%  cross-stitched to sec:born-rule (the record-theoretic route), their
%  agreement flagged as over-determination. sec:local-tomography also
%  minted at source for Postulate 6's operational clause.
%
%  AUDIT RECORD: the section's status discipline was PRESCIENT -- the
%  continuity arrow was already "owed ... enters as a principle" (the
%  exact voice JS mandated for L2), and both of its confessed debts
%  have since been minted as numbered problems; the closing now cites
%  them (op:limit-realisability, op:tsirelson-geometry): "named, and
%  numbered". Postulate 6 now carries BOTH composition duty-holders
%  (Hurwitz multiplicativity + local tomography), matching the
%  continuity section's caveat exactly. Stone's theorem cited
%  (stone1932one, RIS delivered); time-as-operation-count forward
%  flag added; \textquote made portable; opener formalised.
%
%  PROVENANCE NOTE (relocated from the body, 2026-07-06, JS's words):
%  "Using Hopf fibrations to get the Born rule is something Claude AI
%  came up with on its own without my prompting." Preserved here as the
%  historical record. DISCLOSURE STATUS (JS, 2026-07-06): AI use is
%  already acknowledged in the second paragraph of the paper's
%  introduction; the remaining recommendation is only that, at
%  submission, the same fact be echoed in whatever formal
%  Declarations / Acknowledgements slot the target journal's template
%  mandates -- several publishers now require it there specifically,
%  in addition to any in-text mention.
%
%  British English; \autocite; no bullets in prose.
% =====================================================================

% ADDED 2026-07-12: closing paragraph "The postulates in module
% coordinates" -- the Clifford-module reading (skeleton + two
% supplements), with the explanatory arrow guarded per JS: the field
% comes from the Shannon/additivity/commutativity step (self-ignorance),
% and the left-action composability result is recorded as that same
% requirement in operator clothing, not an independent selection of C.
\section{Derivation of the quantum postulates}
\label{sec:postulates}
Enough pieces of the jigsaw are now assembled to derive the postulates of quantum mechanics.

The Dirac--von Neumann postulates are ordinarily laid down as independent
assumptions. In this framework they are not assumptions but a single chain: each
follows from the Inside/Outside Equivalence Principle by one move from the
preceding step. We assemble the chain here and mark the status of each link as
proved, structural, or still owed, so that no step claims more than it has earned.

The chain begins below the postulates. An embedded observer, required by the
equivalence principle to describe a totality that includes itself, must assign a
probability to a proposition about its own future state. The Lawvere--Yanofsky
diagonal \autocite{Yanofsky2003} and Popper's predictor argument
\autocite{Popper1950a, Popper1950b} show that no consistent classical assignment
exists: the self-referential proposition is undecidable to the observer before
the fact. Representing the result of that undecidable self-prediction forces a
signed, non-classical probability, which carries the observer's irreducible
self-ignorance. From this single obstruction the postulates follow.

\paragraph{Postulate 1 (complex Hilbert space).} Two independent steps fix the state space. The
field is forced first: a signed probability, read through the Shannon information measure and
analytically continued, yields a complex amplitude $S=-\ln\lvert p\rvert-\mathrm{i}\theta$,
whose phase $\theta$ is the self-referential ignorance (Section~\ref{sec:complex-amplitudes}).
The metric is forced second, and separately: the principle of reversible continuity
(\S\ref{sec:continuity}) --- that a continuous reversible transformation connects any two pure
states --- selects the $L^{2}$ norm uniquely. For every $\ell^{p}$ with $p\neq 2$ the linear
isometries are generalised permutations \autocite{Banach1932, Lamperti1958}, a group that cannot
connect a basis state to a superposition; only at $p=2$ is the norm-preserving group connected
and transitive on the sphere \autocite{Aaronson2004h}, and Hardy's operational $K=N^{2}$
counting reaches the same quadratic norm by a second route from the same continuity premise
\autocite{Hardy2001a}, and \v{C}encov's theorem with the every-basis lemma reaches it from the
value continuum alone (\S\ref{sec:cencov-route}). The state manifold is therefore the unit
sphere of the Hermitian $2$-norm in $\mathbb{C}^{n}$, taken modulo the global phase
(Section~\ref{sec:l2-metric-fibre}). Placing the metric here, before the Born rule rather than
after it, is deliberate: the Born derivation presupposes the round sphere, so the metric must be
anchored independently to avoid circularity. Two bounds belong with this step. Reversible
continuity fixes the metric, not the field---real and quaternionic quantum theory satisfy it
equally---so the restriction to $\mathbb{C}$ rests on the Shannon step, not on continuity. And
deriving reversible continuity itself from self-reference is the one arrow the chain still owes
--- now priced exactly: one open premise, that every transformation of the regime is a phase of an inside structure,
Open Problem~\ref{op:limit-realisability}, with everything downstream of it proved
(\S\ref{sec:continuity}). Until it is closed, continuity enters as a principle, and the weight
of the $L^{2}$ conclusion is carried jointly with the isometry and counting arguments and with
the statistical route.

\paragraph{Postulate 2 (Hermitian observables).}
A measurement result must be communicable: it must be copyable from one observer
to another, the \textsc{copy} operation of the diagrammatic calculus
\autocite{CoeckeKissinger2017}. The imaginary part of the amplitude is the
private self-referential phase and cannot be broadcast; only the real,
gauge-invariant content can. Observables whose spectra are real and therefore
communicable are exactly the self-adjoint operators. Hermiticity is thus the
condition that a quantity be the kind of thing a measurement can report.

\paragraph{Postulate 3 (Born rule).}
\label{sec:born-invariance}
Outcome probabilities must be invariant under the global phase, the $U(1)$
self-ignorance, realised geometrically as the $S^{1}$ fibre of the complex Hopf
fibration $S^{1}\to S^{3}\to S^{2}$ over the Bloch sphere
\autocite{Hopf1931, Urbantke2003}. There is no polynomial $U(1)$-invariant of degree one,
and the non-polynomial $\lvert\psi\rvert$ is excluded by the analyticity declared at
\S\ref{sec:complex-amplitudes}; the lowest-degree polynomial invariant is quadratic, so the
probability must be $p=\lvert\psi\rvert^{2}$. The squaring is forced
rather than merely permitted, and the argument holds for a single qubit, where
Gleason's theorem is silent. The Born quadratic is not a new structure: it is the
same $L^{2}$ form that fixed the metric in Postulate~1, now read as the
normalisation of the round sphere, and the fibre over which the phase is averaged
has constant length, so the phase is hidden uniformly
\autocite{GluckWarnerZiller1986}. This step is a structural argument rather than a
theorem in the strict sense, and is flagged as such. The record-theoretic derivation of the same
quadratic --- squared magnitudes from fibre privacy, Proposition~\ref{prop:f1-dictionary} --- is
given at \S\ref{sec:born-rule}, and the cost-ledger derivation --- the record as the closed loop of
a history and its return leg, the exponent the number of legs --- at \S\ref{sec:born-cost-ledger}:
three routes to one rule, each holding for a single qubit, the over-determination this framework
treats as evidence.

\paragraph{Postulate 4 (state reduction).}
Measurement is a \textsc{fanout}: the gauge-invariant relational content is copied
into a record, converting a private epistemological amplitude into a public
fact. For the observer who performs it, acquiring that fact is Bayesian updating,
and the projection postulate is its expression. The reduction is perspectival---a
fact relative to the observer in the sense of relational quantum mechanics
\autocite{Rovelli1996a}---and it is graded by redundancy. A single copy is a
relative fact, still reversible by an agent controlling the joint system; redundant
proliferation across many independent fragments makes the fact objective, which is
quantum Darwinism \autocite{zurek_2009} and the stable/relative-fact distinction of
Di~Biagio and Rovelli \autocite{DiBiagioRovelli2021}. ``Collapse'' and unitary
evolution are then the two charts---inside and outside the \textsc{fanout}
boundary---of one process, as the Wigner's-friend analysis shows and \S\ref{sec:fanout-boundary} formalises
(Section~\ref{sec:wigners-friend}); the irreversibility and its $kT\ln 2$ cost
\autocite{Landauer_1961, wigner_1961} accrue with redundancy rather than at the
single act. The reduction postulate is therefore recovered as a chart-dependent
description, not a fundamental dynamical law; the identification of redundancy as
the grading parameter is a synthesis of the framework's parts and is flagged as
such.

\paragraph{Postulate 5 (unitary evolution).}
Conservation of probability is preservation of the $L^{2}$ inner product, and the
inner-product-preserving linear maps are the unitaries $U(n)$---the same
isometry fact that fixed the metric in Postulate~1. Continuity in time --- operation count,
\S\ref{sec:space-of-computations}, a forward reference --- together with Stone's theorem
\autocite{stone1932one} gives $U(t)=\mathrm{e}^{-\mathrm{i}Ht}$, the Schr\"odinger
equation. Reversibility is here a corollary, not a premise: bare reversibility on
the amplitude space yields only the general linear group, and it is preservation
of the quadratic form that carves $U(n)$ out of it.

\paragraph{Postulate 6 (composite systems).}
Systems combine by the tensor product because the division algebra norm is
multiplicative \autocite{hurwitz1898composition}: composition preserves the $L^{2}$ structure
only if states multiply. The operational half closes from the other side: the locality of
tomography derived at \S\ref{sec:local-tomography} says product states determine composites,
which is the tensor product's signature. Cayley--Dickson doubling carries $\mathbb{C}$ to the next
rung, and the ascent of the Hilbert space is the ascent of the Hopf tower, complex
to quaternionic, used in the Wigner's-friend construction
\autocite{mosseri2001hopf}; the ladder itself is the business of the section that follows
(\S\ref{sec:division-algebras}).

\paragraph{Consequence (the Tsirelson bound).}
The postulates fix not only the form of the theory but the strength of its
correlations. In the two-qubit ball the maximal entanglement sits at the centre,
and a Popescu--Rohrlich box \autocite{PopescuRohrlich1994} would require a point
beyond it; the CHSH correlator is capped at $2\sqrt{2}$
\autocite{Cirelson1980QuantumInequality, VerstraeteWolf2002, Held2021}, with the value recovered
at the centre and its general derivation still owed
(Section~\ref{sec:wf-reconciliation}).

\medskip
\noindent The result is that the six postulates are one chain, not six premises.
Of the links, the field and the Born quadratic rest on the framework's signed
probability and $U(1)$-invariance arguments; the $L^{2}$ metric, the unitary group
and the tensor product are over-determined by reversible continuity, the isometry
classification and Hurwitz composition; state reduction is recovered as the
chart-relative face of \textsc{fanout}. The two outstanding debts are the
derivation of reversible continuity from self-reference and the forcing of the Tsirelson value as
a theorem rather than a recovered number --- named, and numbered:
Open Problems~\ref{op:limit-realisability} and~\ref{op:tsirelson-geometry} --- so that a reader
can see exactly where the chain is complete and where it is not.

\paragraph{The postulates in module coordinates.}
One reading of the chain deserves a closing statement, because it positions the derivation
against a mathematical skeleton the reader may know well, and because stating it carelessly would
invert the chain's explanatory arrow. In the coordinates of the division-algebra programmes
(\S\ref{sec:algebraic-dictionary}), complex Hilbert-space quantum mechanics is real Clifford
module theory plus exactly two supplements. The skeleton is Clifford throughout: the operator
algebra of $n$ qubits is the Clifford algebra of $2n$ generators with the state space its own
spinor module, half-integer spin and fermionic statistics are the module's double cover and its
graded tensor, and the classically simulable stabiliser core is the algebra's combinatorics. What
the skeleton does not contain is, first, the choice of complex structure --- the polarisation $J$
that turns a real module into a complex Hilbert space, which standard quantum theory and the
algebraic programmes alike \emph{choose} --- and, second, the Born weights on the fibre. Those two
supplements are precisely what this section has derived, and derived from self-ignorance rather
than from algebra: the field from the Shannon step (information about one's own undecidable state
must still be information --- it must add, addition commutes, and so the amplitude algebra is
commutative, excluding the quaternionic bookkeeping an observer could never balance), and the
quadratic from fibre privacy. The algebraic echo --- that the algebra's own left actions compose
only along a commutative line, the composability result recorded with the dictionary --- is not
an independent selection of $\mathbb{C}$: composition of left actions fails off that line exactly
because the doubled slot carries the \emph{reversed} product, and order-dependence is the same
disease in operator clothing that non-additivity is in information clothing. The mathematics is
the epistemic argument's shadow, welcome as consistency and as a bridge to the module-side
literature (Appendix~\ref{app:winding-number} walks the bridge for charge), but the origin of the
complex amplitude is the observer's ignorance of their own state, and the arrow runs from there
to the algebra, never the reverse.
%      % =====================================================================
%  SECTION: The Division-Algebra Ladder
%  Drafted 2026-07-03. One of the three foundations sections
%  (sec:self-reference, sec:division-algebras, sec:fanout-boundary);
%  this one resolves the dangling \S\ref{sec:division-algebras} in the
%  Hopf-justification section, whose header names it "the
%  Cayley--Dickson / Hurwitz step".
%
%  PLACEMENT in main.tex (REORDERED per JS 2026-07-06): position 6,
%  opening Act II ("why stop at C?") right before its first consumer --
%      % EDITED 2026-10-29 scripts/patch_continuity_tower_premise_2026-10-29T1500.py
% EDITED 2026-09-20 scripts/patch_born_cost_ledger_v2_2026-09-20T0900.py
% =====================================================================
%  SECTION: Self-Reference: One Principle
%  Drafted 2026-07-04. The paper's foundational section; resolves the
%  largest dangling label in the manuscript (sec:self-reference).
%
%  PLACEMENT in main.tex: FIRST of the foundations chain --
%      \input{self_reference_section_2026-07-04}
%      \input{continuity_section_2026-07-04}
%      \input{division_algebras_section_2026-07-03}
%      \input{hopf_justification_section_2026-06-08}
%
%  SOURCES (re-read 2026-07-04): lawvere_inside_outside_fixed_point_
%  theorem_2026-05-04T1806.md (the spine: the two-principles reduction,
%  the theorem, the sigma candidate, the open items -> op minted per
%  D9); research_summary_neg_prob.md (negativity route; lamp (1+i)/2);
%  dynamics_companion_fanout_unitarity (legs 1a/1c; anti-events; the
%  Brukner-Zeilinger rebuttal per D6); Pure_states_and_mixed_states
%  (local tomography derived; QBism contrast per D7).
%
%  DECISIONS ENCODED (JS 2026-07-04): D1 title as below; D2 Wheeler1990
%  cited here as the title anchor; D3 lineage Popper -> Breuer ->
%  Svozil -> Szangolies (Dalla Chiara omitted); D4 Frauchiger_2018
%  cited here, dissolved in sec:wigners-friend; D5 Kirkwood-Dirac
%  heritage via the Arvidsson-Shukur 2024 review as one-stop; D6
%  Brukner-Zeilinger + the framework rebuttal IN ("Shannon is
%  post-FANOUT"); D7 QBism contrast IN; D8 Aaronson private
%  communication OMITTED (published island paper only); D9 one op
%  minted (op:sigma-functor), cross-linked to op:hopf-frobenius.
%
%  CITATIONS -- existing keys: Wheeler1990, small_2006, Godel1931,
%  Lawvere_1969, Yanofsky2003, CoeckeKissinger2017, Rovelli1996a,
%  Popper1950a, Popper1950b, Breuer1995, svozil_k_2018,
%  szangolies_j_2018, Frauchiger_2018, dirac1942, Feynman_r_82,
%  Feynman1987, Youssef1994 [SWAP-FLAG: intended paper is "Quantum
%  mechanics as complex probability theory", Mod. Phys. Lett. A 9
%  (1994) 2571 -- if that lives under youssef1994quantum, swap],
%  caticha_a_2006, goyal_knuth_skilling_2010, Aaronson2004h, Hardy2001a,
%  muller_p_2011, Thomson_1954, benacerraf1962tasks, cantor1891diagonal,
%  Turing_1937.
%  NEW keys (RIS in self_reference_refs_2026-07-04.ris, all verified
%  this session): tikochinsky1988feynman, renou2021real, chen2022ruling,
%  wu2022refutation, hoffreumonwoods2026real, arvidssonshukur2024kirkwood,
%  bruknerzeilinger2001shannon, fuchsschack2013qbism,
%  shannon1948mathematical.
%
%  CROSS-REFS: forward (fine under two-pass LaTeX): sec:continuity,
%  sec:division-algebras, sec:hopf-justification, prop:f1-dictionary,
%  sec:wigners-friend, op:limit-realisability, op:hopf-frobenius.
%  Still dangling after this section: sec:fanout-boundary (and its
%  sibling label sec:fanout-direction) -- both DRAFTED 2026-07-06.
%
%  AUDITED 2026-07-06 (section-by-section final pass, first stop):
%  proof of prop:inside-outside re-verified by hand; attribution chain
%  clean (Burgin/Khrennikov correctly absent); all header-declared
%  citations confirmed used; lamp arithmetic and antecedents verified;
%  FANOUT smallcaps consistent; no pre-definition use of "frame".
%  Fixes: stray Unicode em-dash (QBism paragraph, the file's only
%  non-ASCII character) -> house ---; trace-blind two-way stitch added
%  to the Born subsection ("one privacy, two duties");
%  op:sigma-functor gains its first standing constraint (F1
%  inheritance via the op:hopf-frobenius discharge) and its previously
%  one-directional local-tomography cargo made explicit.
% =====================================================================

\section{Self-Reference: One Principle}
\label{sec:self-reference-2}

Wheeler asked for it from bit: for physics derived from information, ``every it --- every particle,
every field of force'' obtained from binary answers to yes-or-no questions \autocite{Wheeler1990}.
This paper's route to that programme runs through G\"{o}del: the single principle from which
everything else here is derived, deferred, or priced is \emph{self-reference} --- the fact that any
complete physical description, made from inside the world it describes, must describe the describer.
This section states the principle, proves that the framework's second historical commitment follows
from it as a theorem rather than standing beside it as an axiom, and then extracts the price of the
principle in four instalments: probabilities whose parts can cancel, amplitudes valued in
$\mathbb{C}$, the Born rule, and the local tomography of composite systems. What the section cannot
yet derive is stated as precisely as what it can, in one numbered problem and in the ledger at the
end.

\subsection{Events, not objects}

The framework's ontology is relational: the fundamental entities are events and the relations
between them, and objects are the bookkeeping of relational invariants --- what remains stable
across many relations. The kinship with relational quantum mechanics \autocite{Rovelli1996a} is
genuine and acknowledged: facts are facts relative to the systems that register them. The
difference, developed over the rest of this paper, is that the present framework grounds the
relational structure in a single logical principle and then carries it, through the division
algebras and the Hopf fibrations, all the way to the gauge group and the particle spectrum ---
relations are not only the ontology but the derivation engine. On this reading a measurement is not
the revelation of a pre-existing property; it is the creation of a relation, and the later sections
of this paper read that creation geometrically (\S\ref{sec:hopf-justification}) and even
gravitationally. Here we need only the logical core: whatever describes the world is an event in
the world it describes.

\subsection{One principle, and a theorem}

The framework historically carried two foundational commitments: self-reference, and the
\emph{Inside/Outside Equivalence Principle} --- that the laws of physics must be identical whether
formulated by an embedded participant or by an external observer. Two principles are one too many,
and this subsection reduces them.

The reduction corrects the framework's own origin, and we make the correction in print. The 2005
form of the argument \autocite{small_2006} ran a G\"{o}delian chain: a self-referential process,
being complete, must by the incompleteness theorem \autocite{Godel1931} be inconsistent; an
inconsistent system can negate itself; and the negation of self-reference is the production of
causal sequence. The chain captures the right intuition and is retained here as motivation, but it
is not sound as a proof: classical inconsistency yields \emph{ex falso quodlibet} --- every
proposition provable, every causal sequence an output, no constraint at all. The load-bearing
argument is different, and better: a fixed-point theorem.

Let $\mathcal{C}$ be the category of physical descriptions and admissible processes --- concretely,
the dagger-compact-closed setting with classical structures in which categorical quantum mechanics,
and the whole of this paper, works \autocite{CoeckeKissinger2017}. Lawvere showed that in a
sufficiently rich such setting, any weakly point-surjective self-representation forces every
endomap to have a fixed point, and Yanofsky showed that one diagonal argument of this shape
underlies Cantor's, G\"{o}del's, Turing's and Tarski's theorems alike
\autocite{Lawvere_1969,Yanofsky2003,cantor1891diagonal}. The framework's single commitment is that
$\mathcal{C}$ contains a \emph{self-representing object} --- an object capable of encoding the
processes on itself. That is what ``self-reference is foundational'' means, stated categorically,
and it is the surjectivity hypothesis of the diagonal lemma. Now let
$\sigma \colon \mathcal{C} \to \mathcal{C}$ be the \emph{inside/outside involution}: the operation
exchanging the two complementary presentations of any description --- the inside presentation, as
participant-indexed total space in which the observer is part of the system, and the outside
presentation, as a base-times-fibre decomposition over a classical parameter space. That
$\sigma^{2} = \mathrm{id}$ is not an axiom but a coherence requirement: the outside view of the
inside view of a thing is the thing. The candidate concrete realisation of $\sigma$ is the
Hopf-duality swap --- total space against base-times-fibre for the three fibrations --- which gives
the involution computable content and ties this subsection to \S\ref{sec:hopf-justification}.

\begin{proposition}[Inside/outside fixed point]\label{prop:inside-outside}
Let $\mathcal{C}$ satisfy the richness condition above and contain a self-representing object, and
let $\sigma$ be an involution on $\mathcal{C}$. Then the fixed-point locus
$\mathrm{Fix}(\sigma)$ is non-empty: there exist descriptions invariant under the exchange of
inside and outside perspective.
\end{proposition}

\begin{proof}
Lawvere's construction, applied to $\sigma$ in place of negation: given a weakly point-surjective
$\varphi \colon A \to Y^{A}$, form $g(a) = \sigma\!\left(\varphi(a)(a)\right)$; point-surjectivity
supplies $a^{*}$ with $\varphi(a^{*}) = g$, and then $\varphi(a^{*})(a^{*})$ is a fixed point of
$\sigma$ \autocite{Lawvere_1969,Yanofsky2003}. Where the same construction applied to Boolean
negation yields paradox --- negation has no fixed point on a two-element type, so the conclusion
explodes the premises --- the inside/outside involution is an honest involution on a rich type, and
the conclusion is a constructive fixed point rather than a contradiction.
\end{proof}

\textbf{Definition.} \emph{Physical law is the structure carried by} $\mathrm{Fix}(\sigma)$.

The Inside/Outside Equivalence Principle is thereby demoted from axiom to theorem: physical law is
invariant under the exchange of perspectives \emph{because it is the fixed-point set of that
exchange}, and the fixed-point set is forced to be non-empty by self-reference alone. One
commitment remains foundational; everything the second commitment asserted is now derived. The
reader who accepts categorical quantum mechanics as an existing programme has already accepted the
machinery; the framework adds one premise and collects the rest. And the fixed-point locus is not
empty of structure: it inherits the Frobenius-algebra apparatus of $\mathcal{C}$, so that the
Bell, GHZ and W primitives --- and with them the particle classification of the later sections ---
live inside $\mathrm{Fix}(\sigma)$ from the start. What is proved above is non-emptiness given the
setting; what is assumed is the setting itself and the existence of $\sigma$ as a genuine
involution; what is open is the rigour of both, and it is priced now rather than hidden:

\begin{openproblem}[The involution made functorial]\label{op:sigma-functor}
Fix the category $\mathcal{C}$ rigorously --- the dagger-compact-closed setting with classical
structures is the candidate --- and define $\sigma$ on morphisms, not only on objects, the natural
route being the induced fibre--base exchange on spider diagrams; establish that $\sigma$ is an (at
worst pseudo-)involutive endofunctor; and prove or refute the identification of
$\mathrm{Fix}(\sigma)$ with the three Hopf cross-sections, under which Adams's theorem would become
a corollary of the fixed-point classification. This problem is cross-linked with Open
Problem~\ref{op:hopf-frobenius}: under an affirmative resolution of both, the parallelisable
terminus, the fixed-point classification, and compositional completeness are three statements of
one fact, and progress on either problem constrains the other. One constraint is already in place:
at the complex rung, Proposition~\ref{prop:f1-dictionary} supplies the dictionary between
classical structures and the fibration's lifted skeletons --- the clause of
Open Problem~\ref{op:hopf-frobenius} it discharges --- and any resolution here must reproduce it.
The problem also carries a piece of cargo recorded later in this section: the formalisation, within
the fixed category, of the perspective-accounting that derives local tomography.
\end{openproblem}

\subsection{The lineage}

The obstruction this section builds on is old, and the framework claims its completion, not its
discovery. Popper argued in 1950 that a predictor embedded in the world has irreducibly inferior
knowledge of itself, because it cannot impart information about itself without, by that very act,
rendering the information out of date \autocite{Popper1950a,Popper1950b} --- the passage is quoted
where the framework makes the obstruction structural (\S\ref{sec:wigners-friend}). Breuer proved
the impossibility of accurate state self-measurement within quantum theory \autocite{Breuer1995};
Svozil's monograph surveys the whole territory where physics meets self-reference, undecidability
and Chaitin-style algorithmic limits \autocite{svozil_k_2018}; Szangolies sharpened the modern form
as epistemic horizons \autocite{szangolies_j_2018}; and the Frauchiger--Renner theorem is the
contemporary crisis-form --- nested observers deriving contradictory records when each treats the
others' quantum descriptions as universally valid \autocite{Frauchiger_2018} --- which
\S\ref{sec:wigners-friend} dissolves relationally rather than paradoxically. This framework's
addition to the lineage is twofold. Structurally, Proposition~\ref{prop:f1-dictionary} converts
Popper's temporal obsolescence into algebraic cancellation: the datum a copier cannot put into its
own record is exactly its fibre commitment, which does not arrive stale --- it cancels. And
foundationally, the present section runs the lineage in reverse: where the tradition treats
self-reference as a limit on knowledge \emph{within} physics, the framework takes it as the
generator of physics, and derives the quantum formalism from it.

\subsection{The price is negative probability}\label{sec:complex-amplitudes-2}

What does it cost a probabilistic description to contain its own describer? The first instalment is
the sign. In the regime before any record exists, every event has an anti-event --- the description
must be able to undo what it has provisionally done, or the self-description could never be
revised into consistency --- and a probability calculus closed under undoing cannot be confined to
$[0,1]$: it is a signed calculus, in which the parts of a probability can oppose and cancel. The
framework did not discover negativity as bookkeeping; Dirac admitted negative probabilities as
formal intermediaries \autocite{dirac1942}, Feynman defended them as consistent whenever the final,
recorded probabilities stay positive \autocite{Feynman_r_82,Feynman1987}, and the Kirkwood--Dirac
distribution has carried complex quasiprobability from 1933 to the present, as a recent
comprehensive review documents \autocite{arvidssonshukur2024kirkwood}. What the framework adds is
the forcing: negativity is not a permissible convenience but the first price of self-referential
closure, and it is charged exactly and only in the pre-record regime --- which is why every
negativity in the tradition lives in intermediate, unrecorded quantities and never in a final
record.

The second instalment is complexity. A signed information measure must extend Shannon's
$S = -\ln p$ off the positive axis, and the analytic continuation of the logarithm carries $p$ into
the complex plane: amplitudes, with the phase as the continued information. Here too the extension
has a prior tradition, and the attributions are owed precisely. Tikochinsky derived Feynman's rules
for complex amplitudes from consistency of composition \autocite{tikochinsky1988feynman}; Youssef
reformulated quantum mechanics as complex probability theory outright \autocite{Youssef1994};
Caticha carried objective amplitudes to Bayesian probabilities \autocite{caticha_a_2006}; and
Goyal, Knuth and Skilling derived Feynman's rules from a pair-postulate without assuming
analyticity \autocite{goyal_knuth_skilling_2010}. The extension is prior art. The framework's
claim is the compulsion: the complex extension is not a modelling choice among exotic probability
theories but the unique continuation forced by self-referential closure --- the same forcing that
Aaronson glimpsed when he argued quantum mechanics is what negative probability inevitably becomes
\autocite{Aaronson2004h}.

Why does the continuation stop at $\mathbb{C}$? Downward, the real field is not available: real
quantum theory is experimentally distinguishable from complex quantum theory in network scenarios
and has been disfavoured by experiment
\autocite{renou2021real,chen2022ruling,wu2022refutation}; the recent objection that the exclusion
rests on locally tomographic composition \autocite{hoffreumonwoods2026real} is, for this framework,
not an escape but a confirmation of where the weight sits --- local tomography is \emph{derived}
below, not assumed, so the framework closes from its own resources exactly the door the objection
holds open. Upward, the quaternions and octonions fail as a \emph{base} for the information
continuation: the continued measure is a logarithm, and beyond $\mathbb{C}$ the logarithm of a
product is no longer the sum of logarithms --- non-commutativity severs
$\log(ab) = \log a + \log b$ --- so the additivity that makes it an information measure is lost.
Aaronson's parameter-counting and signalling arguments against a quaternionic base point the same
way \autocite{Aaronson2004h}. The higher algebras are not banished; they return as the
\emph{internal} algebras of the register, which is the business of \S\ref{sec:division-algebras}.
The base is $\mathbb{C}$, for informational reasons, and stays $\mathbb{C}$.

\subsection{The Born rule from fibre privacy}\label{sec:born-rule}

The third instalment is the quadratic rule. An amplitude carries one datum more than any record of
it can: its overall phase. This was asserted throughout the framework's development as the
observer's private, self-referential ignorance; it is now a proved statement, in exactly the form
needed. By Proposition~\ref{prop:f1-dictionary}, a classical structure --- the copy operation any
record requires --- is committed to a specific lift of its skeleton, and that commitment is absent
from every classical record the structure induces: the copier is fibre-committed, the record is
fibre-blind, and the fibre freedom is gauge for every observer. (The computation extends verbatim
from one qubit to the full register, where Proposition~\ref{prop:trace-blind} makes it the ground
of hypercharge's tracelessness --- one privacy, two duties.) Popper's obsolescence, made
cancellation. It follows that anything recordable is invariant under the fibre --- observables and
their statistics must be $U(1)$-invariant functions of the amplitude. The lowest-degree polynomial invariant
of a complex amplitude is the quadratic $|\psi|^{2}$ --- the non-polynomial invariant $|\psi|$,
of degree one, is excluded by the analyticity declared at \S\ref{sec:complex-amplitudes} --- and
normalisation fixes the coefficient:
probability is the minimal record-compatible function of the amplitude,
\[
p \;=\; |\psi|^{2},
\]
the Born rule as the cheapest invariant the fibre permits. The same square is reached from the
cost side in \S\ref{sec:born-cost-ledger}, where a record is the closed loop of a history and its
return leg and the exponent is the number of legs; neither route uses the lattice of projections
that Gleason's theorem needs, and both hold for a single qubit. The derivation is the framework's
own;
its premise --- that the fibre datum is structurally excluded from records --- is
Proposition~\ref{prop:f1-dictionary}, and its geometric reading --- probability as base-point data
under the Hopf projection --- is developed in \S\ref{sec:hopf-justification}.

One standing objection to information-theoretic foundations is answered here, because this
framework needs the answer. Brukner and Zeilinger argued that Shannon information is conceptually
inadequate for quantum measurements: it is basis-dependent, and its operational meaning presupposes
ignorance of a definite but unknown fact, which the quantum case does not supply
\autocite{bruknerzeilinger2001shannon}. The framework's reply is structural, and it is a
demarcation rather than a rebuttal. Shannon's precondition fails exactly where they say it fails
--- in the pre-record regime, where there is no definite unknown to be ignorant of --- and there
the operative measure is the quadratic invariant just derived. Shannon's axioms hold exactly where
the precondition holds --- on records, after \textsc{fanout}, where outcomes are definite and
ignorance is classical --- and there the uniqueness theorem forces $-\log p$
\autocite{shannon1948mathematical}. The basis-dependence dissolves with the same stroke: a basis
\emph{is} selected precisely at \textsc{fanout}, so the basis-dependent measure applies only where
a basis exists. Shannon is not inadequate; Shannon is post-\textsc{fanout}. The framework thereby
keeps both measures, each on its own side of the boundary, and uses the additivity of the
post-record measure --- the very property the previous subsection used to pin the base field ---
with a clear conscience.

\subsection{Thomson's lamp, resolved}

The section owes \S\ref{sec:continuity} a debt incurred there: the lamp's amplitude-side limit.
Thomson's lamp --- toggled at times $1 - 2^{-n}$, its state at $t = 1$ demanded --- has no
classical answer: the task description fixes every finite stage and thereby fixes no limit at all
\autocite{Thomson_1954,benacerraf1962tasks}, which \S\ref{sec:continuity} read as the proof that
the record algebra is not limit-closed. On the amplitude side the limit exists, and the calculus of
this section computes it. The toggle exchanges the two amplitudes; the self-consistent state of the
completed task is the fixed point of the exchange in the complex calculus, and solving it under the
signed continuation gives
\[
P(\mathrm{on}) = \tfrac{1 + i}{2}, \qquad P(\mathrm{off}) = \tfrac{1 - i}{2},
\]
a superposed fixed point whose imaginary part carries the unrecordable datum --- the phase of the
endless alternation, when the lamp \emph{would have} stopped. Two checks certify that the
assignment is calculus rather than decoration. The pair is normalised,
$P(\mathrm{on}) + P(\mathrm{off}) = 1$; and the Born rule of the previous subsection, applied to
the fixed point, returns $|P(\mathrm{on})|^{2} = \tfrac{1}{2}$ --- exactly the symmetric classical
ignorance that is the only sane record-level answer, recovered rather than imposed. The lamp,
which entered the framework at its very beginning as the place where classical computation meets
its boundary, thus exits it as a worked example: undetermined as a record, determinate as an
amplitude, the two sides split precisely along \textsc{fanout} --- and the diagonal construction it
physically instantiates \autocite{cantor1891diagonal,Turing_1937} is the same one that powers
Proposition~\ref{prop:inside-outside}. The status flag is owed here as everywhere: the
undetermination is theorem-level; the specific fixed-point value is the framework's complex
calculus, structural, offered with its two consistency checks shown.

\subsection{Local tomography, derived}\label{sec:local-tomography}

The fourth instalment concerns composites, and it discharges two promises at once. The
Inside/Outside Equivalence Principle --- now a theorem --- requires that an external observer's
characterisation of a composite system be recoverable by the embedded participants, and the
embedded participants have access only to local measurements and their correlations. Equivalence
of the two perspectives therefore entails that the state of a composite is determined by the
statistics of local measurements: \emph{local tomography}, obtained as a consequence rather than
imposed as an axiom. This converts a postulate of the reconstruction programmes
\autocite{Hardy2001a,mullerp2011} into a theorem-shaped consequence of the one principle ---
discharging the caveat recorded in \S\ref{sec:continuity} --- and it closes, from the framework's
own resources, the escape route in the real-number debate noted above: the assumption the
objection of \autocite{hoffreumonwoods2026real} contests is here not an assumption. The status is
flagged plainly: the derivation is structural, one paragraph of perspective-accounting, and its
formalisation within the categorical setting belongs to Open Problem~\ref{op:sigma-functor}.

\subsection{Contrast: QBism, and the ledger}

The nearest neighbour in spirit deserves a paragraph of honest contrast. QBism
\autocite{fuchsschack2013qbism} holds, as this framework does, that quantum probabilities are an
agent's --- that the theory is a calculus of expectation, not a catalogue of properties. What QBism
does not say is why the agent's expectations must be complex-valued and quadratically squared; the
Born rule enters as a normative addition to Bayesian coherence, and the complex amplitude as an
unexplained representation. This section is, in effect, the missing why: the agent's ignorance is
self-referential --- the agent is inside the description --- and self-referential ignorance is
complex, with the quadratic rule as the cheapest record-compatible invariant of it. Where QBism
stops at the agent, the framework derives the agent's calculus.

The ledger, then. Delivered in this section: the reduction of two principles to one
(Proposition~\ref{prop:inside-outside}); the pre-record regime with its anti-events, handed to
\S\ref{sec:continuity} as the five-face regime whose continuity is priced there
(Open Problem~\ref{op:limit-realisability}); the complex base, forced downward by experiment and
upward by the additivity of information; the Born rule from fibre privacy; the lamp's fixed point;
and local tomography. Owed, and stated where they are owed: the functorial rigour of the involution
and the fixed-point classification (Open Problem~\ref{op:sigma-functor}); the remaining premise of continuity (Open Problem~\ref{op:limit-realisability}, in
\S\ref{sec:continuity}); and the
rung-by-rung compulsion of the algebra ladder, posed as Open
Problem~\ref{op:ladder-compulsion} in \S\ref{sec:division-algebras}. The
hand-on is the paper's spine: the $\mathbb{C}$ this section derives natively and the $\mathbb{C}$
that \S\ref{sec:continuity} underwrites through Hardy's theorem are the same field, the floor of
the ladder that \S\ref{sec:division-algebras} climbs and \S\ref{sec:hopf-justification}
geometrises. One principle; everything else is theorem, derivation, or a debt with a stated
invoice.

%      % EDITED 2026-10-29 scripts/patch_continuity_tower_premise_2026-10-29T1500.py
% EDITED 2026-09-15 scripts/patch_value_continuum_declared_cencov_2026-09-15T1540.py
% EDITED 2026-09-15 scripts/patch_continuity_connected_state_space_2026-09-15T1430.py
% =====================================================================
%  SECTION: Continuity Is Not an Axiom
%  Drafted 2026-07-04 from continuity_limit_realisability_2026-07-04.md
%  (now in project documents) and the May record
%  (dynamics_companion_fanout_unitarity_2026-05-20T1430.md).
%
%  PLACEMENT in main.tex: a standalone section between self-reference
%  and division-algebras --
%      \input{self_reference_section_2026-07-04}
%      \input{continuity_section_2026-07-04}
%      \input{division_algebras_section_2026-07-03}
%  (Architecture call SETTLED by the sync manifest: standalone.)
%
%  CITATIONS -- existing keys (as supplied by JS 2026-07-04):
%    Hardy2001a, Aaronson2004h, Thomson_1954, Lawvere_1969,
%    Turing_1937, muller_p_2011, Brukner2009.
%  NEW keys (RIS in continuity_refs_2026-07-04.ris):
%    cantor1891diagonal, benacerraf1962tasks, broeckertomdieck1985compact.
%  Deliberately excluded per JS: Chiribella--D'Ariano--Perinotti;
%  Bargmann/Wigner's theorem. Yanofsky2003 now paired on the
%  Lawvere-shaped clause (the pending-key excuse was void: the key
%  exists and sec:self-reference uses it).
%
%  CROSS-REFS: sec:self-reference now DRAFTED (2026-07-04, later the
%  same day) -- no longer dangling; only sec:fanout-boundary (and its
%  sibling sec:fanout-direction) -- DRAFTED 2026-07-06.
%
%  AUDITED 2026-07-06 (final pass, second stop): both proofs
%  re-verified by hand (Lagrange divisibility sentence exact; the
%  Baire countable-escape chain sound); caveat-discharge with
%  sec:self-reference confirmed two-way (already present); file pure
%  ASCII; no pre-definition "frame"; all citations used as declared.
%  Fixes: Yanofsky2003 paired on the Lawvere clause; the op's two
%  candidate routes anchored to printed propositions (f1-dictionary
%  for the divisibility premise's referent, inside-outside for the
%  closedness demand) -- statuses unchanged, referents sharpened;
%  header staleness corrected. Deliberate FORWARD references --
%  prop:f1-dictionary and sec:hopf-justification (defined later in the
%  paper; fine under two-pass LaTeX and flagged in prose),
%  sec:division-algebras, sec:space-of-computations.
%
%  DECISIONS ENCODED: op:limit-realisability minted (the owed arrow,
%  priced: exactly two derivations, divisibility and closedness);
%  Lagrange used twice (finite exclusion; component-group triviality);
%  the (1+i)/2 lamp value NOT stated here -- "a superposed fixed point"
%  with the complex-probability treatment deferred to
%  sec:self-reference; the two-branch over-determination stated to
%  pre-empt a circularity reading; Thomson spelling throughout.
% =====================================================================

\section{Continuity Is Not an Axiom}
\label{sec:continuity}

Hardy derives quantum theory from five axioms \autocite{Hardy2001a}. Four of them are satisfied
by classical probability theory; the fifth --- that a \emph{continuous} reversible
transformation connects any two pure states --- is the entire difference, the single postulate
that turns the classical simplex into the complex Hilbert space. Every reconstruction programme
pays an equivalent toll: continuous reversibility in Masanes and M\"{u}ller
\autocite{mullerp2011}, the corresponding postulate in Daki\'{c} and Brukner
\autocite{Brukner2009}. The axiom is load-bearing, and in every programme it is unexplained: it
is where the quantumness is put in by hand. The claim of this section is that in the present
framework it is not an axiom. Continuity is the fifth face of the pre-measurement regime that self-reference forces, and what
stands between the framework and a full derivation is not a postulate about values but one
premise about transformations --- that every transformation the inside performs has a stationary
outside presentation --- derived in outline below from the inside/outside involution of
\S\ref{sec:self-reference} and stated as an open problem at the end of the section with nothing
hidden in it. The continuum itself, the completion that every probability theory adopts, is
named for what it is: the real line, shared with classical probability and carrying no quantum
content.

One structural remark prevents a misreading. The framework reaches complex amplitudes natively in
\S\ref{sec:self-reference}, by the negative-probability route; the present section reaches the same
destination by an independent branch --- the regime's continuity, then Hardy's theorem. The two
branches converging on one structure is over-determination, not circularity, and the convergence is
itself evidence that the destination is forced.

\subsection{Five faces of one regime}

The regime before any measurement --- before any \textsc{fanout} --- can be characterised five
equivalent ways: negative probability is available; anti-events exist, every event being undoable;
no classical record exists; no preferred basis is selected; and a continuous reversible
transformation connects any two pure states. These are not five premises to be combined. They are
one regime under five descriptions, and \textsc{fanout} (\S\ref{sec:fanout-boundary}) is the single
event that fails all five at once: it writes the record, selects the basis, kills the anti-events,
discretises the transformation group, and switches off the negativity. The equivalence of the first
two faces is the anti-event identity of \S\ref{sec:self-reference}; the third is the boundary's
definition. The work of this section is the equivalence of the last two faces, and then the route
from the regime's existence to its continuity. The fourth face has, since
Proposition~\ref{prop:f1-dictionary} (\S\ref{sec:hopf-justification} --- a deliberate forward
reference), an exact form: a basis-selection is the commitment of a classical structure, a lifted
antipodal skeleton, so ``no preferred basis'' means precisely that no copy structure has fired.

\subsection{Basis selection is the collapse of continuity}

\begin{proposition}[Continuity forbids a selected basis]\label{prop:basis-discreteness}
If a finite set of pure states is privileged --- a preferred basis --- then every reversible
transformation respecting the privilege carries the set to itself and acts on it as a permutation;
the respecting transformations therefore realise only a finite group of motions among the
distinguished states, and no continuous family of them connects distinct pure states of the set.
Contrapositively, continuous reversibility entails that no basis is selected.
\end{proposition}

\begin{proof}
A privileged finite set is by definition carried to itself by every structure-preserving reversible
map, whose induced action on the set is a permutation of finitely many points. A continuous path of
such maps induces a continuous path in a finite discrete set, which is constant. Hence a continuous
reversible family moving one distinguished pure state to another respects no such privilege, and if
every pair of pure states is so connected, no finite set is privileged.
\end{proof}

The converse is where the historical difficulty lies: the absence of a privileged finite set does
not by itself force continuity, for the transformation group could in principle be infinite and
discrete. Two structural properties close every such escape at once.

\begin{proposition}[Divisible and closed implies connected]\label{prop:divisible-closed}
Let $G$ be the group of reversible transformations of the pre-measurement regime, acting on the
state body. If $G$ is divisible --- every element has an $n$-th root for every $n$ --- and
topologically closed, then $G$ is a compact connected Lie group; and if $G$ moreover acts
transitively on the pure states, then a continuous reversible transformation connects any two of
them: Hardy's fifth axiom.
\end{proposition}

\begin{proof}
The state body is compact and convex with nonempty interior, so its symmetry group is compact, and
a closed subgroup of a compact matrix group is a compact Lie group with finitely many connected
components \autocite{broeckertomdieck1985compact}. The component group $G/G_{0}$ is a quotient of a
divisible group and hence divisible; but a nontrivial finite group is never divisible --- with
$\lvert H\rvert = m$ every element satisfies $h^{m} = e$, so no element other than the identity
possesses an $m$-th root --- and $G/G_{0}$ is finite, so it is trivial and $G = G_{0}$ is connected.
A connected Lie group is path-connected, and transitivity then supplies the continuous reversible
path between any two pure states.
\end{proof}

The same one-line fact of Lagrange closes the remaining escapes. Finite groups are excluded
directly, being nontrivial and never divisible. Countable groups are excluded because a
countable closed subgroup of a compact group is compact, has an isolated point by the Baire
category theorem, and is therefore discrete and finite. Uncountable disconnected groups are
excluded by the component-group argument of the proposition. Every route to a discrete
transformation group therefore passes through a failure of divisibility or a failure of
closedness. Divisibility, however, asks for more than the conclusion needs, and the next
proposition removes it.

\begin{proposition}[Continuity from a connected state space]\label{prop:connected-transitive}
Let the pure states of a finite-dimensional system form a connected space $P$, and let the group
$G$ of reversible transformations of the pre-measurement regime be topologically closed and act
transitively on $P$. Then the identity component $G_0$ already acts transitively on $P$, and a
continuous reversible transformation connects any two pure states: Hardy's fifth axiom holds,
with no divisibility assumed.
\end{proposition}

\begin{proof}
As in Proposition~\ref{prop:divisible-closed}, closedness makes $G$ a compact Lie group with
finitely many connected components \autocite{broeckertomdieck1985compact}, so $G_0$ is a closed
normal subgroup of finite index and is path-connected. Each $G_0$-orbit in $P$ is the continuous
image of a compact group, hence compact, hence closed. Since $G$ is transitive and $G_0$ is
normal, the $G_0$-orbit of $gx$ is $g$ applied to the $G_0$-orbit of $x$, so the cosets of $G_0$
permute the orbits transitively and there are finitely many of them. A finite partition of $P$
into closed sets is a partition into open sets; $P$ is connected; so there is one orbit and
$G_0$ is transitive. Given pure states $x$ and $y$, choose $g \in G_0$ with $gx = y$ and a
continuous path $g_t$ in $G_0$ from the identity to $g$; then $g_t x$ is a continuous reversible
transformation from $x$ to $y$.
\end{proof}

The proposition relocates the premise. Divisibility was asked for because a transformation group
could be infinite and discrete; but a group acting transitively on a continuum cannot be
countable, and a closed group acting transitively on a connected space cannot be disconnected in
the way that matters, so what the argument needs from the regime is not the availability of
$n$-th roots but the connectedness of the state space itself. That is a statement about values:
that the pure states are all the unit vectors of $\mathbb{C}^{K}$ up to phase, the space
$\mathbb{CP}^{K-1}$, connected for $K \ge 2$, and hence that the amplitudes of
\S\ref{sec:self-reference} fill their continuum. The two standing foils show the premise doing
its work. The stabiliser fragment of quantum theory has complex amplitudes and is contextual,
yet it has finitely many pure states, its outcome probabilities take only the values $0$,
$\tfrac12$ and $1$, and its reversible transformations form a finite group, the Clifford group,
so Hardy's fifth axiom fails in it; Spekkens's toy theory \autocite{spekkens2007toy} is the same
with real amplitudes, its ignorance rationed by a knowledge-balance principle. Neither has a
connected state space, and neither is a model of the regime described here, whose credences take
every value. Divisibility, which both also lack --- the phase gate has no square root in the Clifford group
--- is a symptom of the discreteness, not its cause. What both lack can be said exactly, and the
next subsection says it: each contains a transformation that fixes no state it admits.

\subsection{The premises, and what they mean}

Three premises carry Proposition~\ref{prop:connected-transitive}. Transitivity is the regime's
homogeneity: inside/outside equivalence leaves no pure state privileged, and a proper closed
invariant set of pure states would define a superselection observable, a classical record
predating any \textsc{fanout} (\S\ref{sec:fanout-boundary}). Connectedness of the state space is two statements, and only one of them is quantum. The first
is that the pure states are infinite and dense, with a fixed point for every transformation; the
second is that the state space is complete. Quantum theory over the algebraic numbers ---
amplitudes in the algebraic closure of $\mathbb{Q}$ within $\mathbb{C}$ --- has the first
without the second: it has $K = N^{2}$, the Born rule, superposition, entanglement and Bell
violation at the Tsirelson value, every unitary has an algebraic eigenvector, and it is
countable and totally disconnected. So the whole quantum content of Hardy's axiom lies in the
density, and the completion is the real-number convention of every probability theory. Hardy
adopts it before his fifth axiom is stated, in the first, which makes probabilities real numbers
and gives his classical theory a real simplex of mixtures; this framework adopts it at the
Shannon step, where credences are declared to range over a connected domain so that the
continuation of the logarithm is unique (\S\ref{sec:complex-amplitudes}). Given that convention,
density of the pure states is connectedness of their closure. Divisibility, which an earlier
version of this section asked for, is not needed: it follows from the conclusion, every element
of a compact connected Lie group having roots of every order.

Closedness is the completion, limit-realisability: a Cauchy sequence of available
transformations has its limit available. It cannot be derived, and the reason should be stated
rather than hidden. Any collection of first-order principles about states, transformations and
probabilities that has an infinite model has a countable one, and a countable subset of
$\mathbb{CP}^{K-1}$ is totally disconnected; the founding non-computability does not escape
this, since the arithmetical complex numbers form a countable algebraically closed field
containing non-computable numbers and closed under every limit it can itself describe. Hardy's
axiom and closedness are second-order, quantifying over paths and over sequences, and a
second-order principle is grounded as Dedekind completeness is grounded or not at all. The
grounds are these. The completion of a dense set under the finite-precision metric is unique and
depends on no choice, whereas any countable alternative is a choice of definability level, an
extra law with no ground and no test; and the framework's own result is that no observer's
descriptive level captures the regime, since by \S\ref{sec:self-reference} every observer's
state exceeds its own level and the inside/outside principle makes that a law rather than a
perspective. The burden therefore falls on any restriction, not on the completion, which is the
least committal option and carries no quantum content. What limit-realisability adds to the line
is still worth stating exactly, because it is the corrected form of a natural intuition.
Cantor's diagonal \autocite{cantor1891diagonal} does not build the continuum; its output is an
infinite decimal --- a Cauchy sequence of rationals --- which exists \emph{as a point} only
because the line is complete. The diagonal presupposes limit-realisability and then shows that
the completed line outruns every internal listing. And the character of the completion is the
framework's own: the computable reals are countable \autocite{Turing_1937}, hence of measure
zero, so once the transformation group contains its limits almost all of them are
non-computable. Non-computability is a property of the completed continuum, not the reason for
it. The constructive reading, on which only describable limits exist, remains consistent; it is
the countable theory above, and the two readings differ in nothing an experiment can certify.

Density is where the section's title is earned, and it is derived from the paper's own premise.
Write $(\mathrm{I})$ for the statement that every admissible transformation of the regime is a
phase of some admissible classical structure: diagonal, that is, in the record basis of some
observer the regime contains, a phased spider of that basis in the calculus of
\S\ref{sec:hopf-justification} \autocite{coeckeduncan2011interacting} (VERIFY-SUPPORT). It says
of transformations what the global phase of \S\ref{sec:self-reference} says of states: there are
no absolute ones. It follows from three things the paper already holds. The inside/outside
involution is a genuine involution, so the outside presentation of the inside is a presentation
of the same object with nothing lost: the outside can have complete knowledge of the inside, a
potential record, the joint state lying in a structure the outside carries. Time is operation
count, internal to the system that operates (\S\ref{sec:space-of-computations}), so the outside
does not share the inside's count. And before any \textsc{fanout} that count is not classically
recorded, so the outside's description is pure. Let $g$ be the transformation the inside
performs per operation. The outside's complete knowledge cannot depend on the inside's count,
and it is pure, so it is a pure state of the inside's clock and system that is stationary under
``advance the clock and apply $g$''; the stationary states of that map are clock Fourier modes
tensored with eigenvectors of $g$ and nothing else \autocite{pagewootters1983evolution}
(VERIFY-SUPPORT). So the outside holds a structure in which $g$ acts by phases, which is
$(\mathrm{I})$; and the relative phase in that structure is the outside's ignorance of the
inside's computational count, the same ignorance the global phase records one level down. The
consequences are a theorem. A phase of a structure fixes the structure's copyable states, so
every admissible transformation fixes an admissible pure state; a finite transitive action has
an element moving every point (Jordan), so the pure states are infinite; the transformation
carrying a stabiliser state to the fixed point of the composite of a $Z$-phase and an $X$-phase
turns through $\arccos(1/\sqrt{3})$, an irrational fraction of a turn, so the group has an
element of infinite projective order and its orbits are dense. With the completion,
Proposition~\ref{prop:connected-transitive} fires and Hardy's fifth axiom holds. The two foils
fail $(\mathrm{I})$ exactly where they fail continuity: fourteen of the twenty-four Clifford
operations are phases of no Pauli basis and fix no stabiliser state, their fixed points being
the magic states; the toy theory's three-cycles are phases of no epistemic basis; and in each
the outside's knowledge of the inside under the order-three protocol is either count-dependent
or a mixture over the count, so each makes the inside's time classical before any record exists.
The derivation of $(\mathrm{I})$ is an outline, not a proof, and it is what the open problem
below asks for.

\subsection{Why the cut matters physically}

The regime's continuity is not decorative. Reversibility alone, over complex amplitudes, licenses
the full invertible group $GL(n, \mathbb{C})$, and Aaronson showed what that buys: superluminal
signalling, and computational power at the level of \textsf{PP} \autocite{Aaronson2004h}. The cut
from $GL(n, \mathbb{C})$ down to the unitary group is the conserved quadratic form, and continuity
is what entitles the theory to make it: with the fifth axiom in hand, Hardy's theorem delivers the
complex Hilbert space, the quadratic norm and the unitary group in one stroke \autocite{Hardy2001a}
--- cited here, not re-proved. One honest caveat travels with the citation. Hardy's reconstruction
uses his remaining axioms besides --- composition and subspace structure --- and these must be
derived rather than imported: the composition ledger of \S\ref{sec:division-algebras} and the
locality-of-tomography argument of \S\ref{sec:self-reference} carry that duty, and it is discharged
there, not silently assumed here.

\subsection{Super-tasks at the boundary}

The premise just priced has a fifty-year-old stress test in the philosophy of mathematics, and the
framework passes it in an unexpected way. A super-task --- infinitely many operations completed in
finite time \autocite{Thomson_1954} --- presupposes its terminal moment as an actual point of the
timeline: the accumulation point of a generic sequence of step-times exists only in the completed
line, and the defining question, the state \emph{after all the steps}, is meaningful only under
limit-realisability of time. Super-tasks are downstream consumers of exactly the completeness at
issue in this section. Thomson's lamp then exhibits both sides of the \textsc{fanout} boundary at
once. On the record side, the lamp's description fixes its state at every finite stage and thereby
fixes no limit at all: any terminal value is consistent with the whole task
\autocite{Thomson_1954,benacerraf1962tasks}. That is not a puzzle about lamps; it is the proof that
the classical record algebra is not limit-closed, and cannot be. On the amplitude side the limit
exists --- a superposed fixed point, whose complex-probability treatment belongs to
\S\ref{sec:self-reference} --- because that side \emph{is} limit-closed, which is the fifth face
itself. The boundary is thereby the edge of limit-realisability: complete on the reversible side,
incomplete on the record side, with the lamp standing on both sides at once as the boundary's oldest
exhibit. A final clause dissolves rather than answers the classic question. In this framework time
is operation count (\S\ref{sec:space-of-computations}), and ``after infinitely many operations''
indexes no operational moment: the lamp packs its endless \textsc{fanout}s into the emergent
continuum, and asking for its record at the accumulation point asks the discrete contents for a
value that only the completed container possesses. On or off is not undecided; it is unasked.

\subsection{What remains, priced}

\begin{openproblem}[Every transformation a phase of an inside structure]\label{op:limit-realisability}
Proposition~\ref{prop:connected-transitive} reduces Hardy's fifth axiom to three premises:
connectedness of the pure states, transitivity, and closedness. Transitivity is argued above
from the absence of records. Closedness is the completion, adopted on the grounds stated above
and carrying no quantum content. Connectedness is density plus that completion, and density
follows, by Jordan's theorem and the infinite-order argument, from $(\mathrm{I})$: that every
admissible transformation of the regime is a phase of some admissible classical structure. What
remains open is the rigour of the derivation of $(\mathrm{I})$ from the paper's premise. Derive,
inside the fixed-point locus $\mathrm{Fix}(\sigma)$ of Proposition~\ref{prop:inside-outside},
that every endomap is a phase of a classical structure in $\mathrm{Fix}(\sigma)$, from the
self-representing object, the involution $\sigma$ as a genuine involution, and the reading of
the inside presentation as a sequence indexed by operation count with no record before
\textsc{fanout}. The outline above does this with the stationary-state construction of Page and
Wootters; a derivation passes if it makes that outline a theorem of the calculus, with the
outside's count-independent complete description shown to exist and to be pure from the
involution alone. It fails before any detail is checked if every premise it uses holds in the
stabiliser fragment of quantum theory, in Spekkens's toy theory \autocite{spekkens2007toy}, in
the real qubit, or in quantum theory over the Gaussian rationals, all of which violate
$(\mathrm{I})$ outright and all of which make the inside's operation count classical before any
record exists. It also fails if it imports the completion as a premise about transformations:
the completion is the real line and has no derivation, and a derivation of $(\mathrm{I})$ must
hold in the countable algebraic theory, where $(\mathrm{I})$ is true and continuity is not.
\end{openproblem}

The accounting, in one sentence: the axiom that every reconstruction programme pays is, in this
framework, one derivation posed, of Lawvere's shape --- that every endomap of the fixed-point
locus is a phase of a classical structure in it --- and one convention named, the completion,
which classical probability pays too. And
the hand-on is already visible: the complex amplitudes that this section's continuity underwrites
through Hardy's theorem, and the complex amplitudes that \S\ref{sec:self-reference} derives natively,
are one and the same $\mathbb{C}$ --- the floor of the ladder that \S\ref{sec:division-algebras}
climbs.

%      % EDITED 2026-10-29 scripts/patch_continuity_tower_premise_2026-10-29T1500.py
% EDITED 2026-09-15 scripts/patch_value_continuum_declared_cencov_2026-09-15T1540.py
% EDITED 2026-09-15 scripts/patch_continuity_connected_state_space_2026-09-15T1430.py
% =====================================================================
%  SECTION: The Euclidean metric on the state sphere, and the global
%  phase as self-referential ignorance
%  Origin: drafted in an earlier conversation (2026-06-07), edited by
%  JS; audited into the current architecture 2026-07-06.
%
%  PLACEMENT in main.tex (reordered chain, JS 2026-07-06): position 3 --
%      \input{self_reference_section_2026-07-04}
%      \input{continuity_section_2026-07-04}
%      \input{L2_metric_global_phase_section_2026_06_07}   <-- here
%      \input{wigners_friend_in_the_hopf_picture_2026-06-13}
%  (WF backward-references sec:l2-metric-fibre: order-consistent.)
%
%  AUDITED 2026-07-06 (final pass): Lamperti/Banach classification,
%  Hopf map components, submersion facts, Hardy K=N^r reading all
%  verified; the recoverable-vs-gauge-private fibre distinction
%  confirmed canonical and untouched; file pure ASCII. Fixes: opening
%  rewritten (Popper C+/C/E forward-voiced to sec:wigners-friend;
%  continuity claim made CONDITIONAL per JS -- "assuming those premises
%  can be proved", op:limit-realisability cited); Aaronson paragraph
%  de-duplicated against sec:continuity by stitch; sec:born-rule was
%  dangling paper-wide -- label minted at its canonical home (the Born
%  subsection of sec:self-reference) and the claim here recalibrated
%  ("forced, not chosen" -> lowest-degree-invariant, geometric face);
%  record-absence upgraded from interpretive to theorem via
%  prop:f1-dictionary while the self-ignorance identification stays
%  interpretive; grammar and dash punctuation repaired; field-not-fixed
%  bound pointed at sec:self-reference.
%
%  CITATIONS: existing keys Hardy2001a, Aaronson2004h, mosseri2001hopf.
%  NEW keys (RIS in L2_metric_refs_2026-07-06.ris): Banach1932,
%  Lamperti1958, Hopf1931, Urbantke2003, GluckWarnerZiller1986,
%  wooters_1981 (key spelling preserved as used; note the single-t).
%  British English; \autocite; no bullets in prose.
% =====================================================================

\section{The Euclidean metric on the state sphere, and the global phase as self-referential ignorance}
\label{sec:l2-metric-fibre}
Popper's systems $C^{+}$, $C$ and $E$ --- the predictor, its target and their shared
environment, whose argument \S\ref{sec:wigners-friend} anchors and quotes --- are, in modern
dress, Wigner, his friend and a qubit $S$ to be measured. The natural arena for analysing them
is the Hopf fibration and the Bloch sphere; but even to draw the Bloch sphere one must first
justify the $L^{2}$ metric, and this section earns that right. The principle doing the work is
the reversible continuity of \S\ref{sec:continuity} --- derived there from self-reference
conditionally, on the premises priced as Open Problem~\ref{op:limit-realisability}, of which one, that every
transformation of the regime is a phase of an inside structure, remains open; the continuum
itself is the completion the Shannon step declares (\S\ref{sec:complex-amplitudes}), the
classical convention, and \S\ref{sec:cencov-route} below reaches the same metric by a
second route that consumes it and nothing else that is open. Assuming that premise can be proved,
everything below inherits the derivation; we carry that status, and do not upgrade it.

We have established that the amplitudes of a single qubit form a complex vector space
$\mathbb{C}^{2}\cong\mathbb{R}^{4}$. We have not yet fixed which norm on this space singles out
the physical state manifold. That choice is not free: the principle of
\emph{reversible continuity} forces the Euclidean ($L^{2}$) norm, and hence the round
three-sphere, before any appeal to the Born rule. The quadratic structure usually \emph{posited}
as the Born rule is therefore already present in the metric, and the global phase emerges as a
geometrically forced fibre rather than an added convention.

\subsection{Reversible continuity selects the 2-norm}
\label{sec:reversible-continuity-l2}

The principle is the one isolated by Hardy as the single axiom distinguishing quantum from
classical probability \autocite{Hardy2001a}:
\emph{there exists a continuous reversible transformation between any two pure states.}
Operationally a reversible transformation of states is a symmetry of the state manifold, hence
norm-preserving; the demand is that such symmetries form a \emph{connected} group acting
\emph{transitively} on the pure states. Linearity is inherited from the complex vector-space
structure already in hand, so the admissible transformations are the complex-linear isometries
of whatever norm we place on $\mathbb{C}^{2}$.

Take the candidate norms to be the $\ell^{p}$ family,
$\lVert x\rVert_{p}=\bigl(\sum_{k}\lvert x_{k}\rvert^{p}\bigr)^{1/p}$, with the state manifold
the unit $\ell^{p}$ sphere. Their linear isometries are completely classified.


\begin{proposition}[Lamperti; after Banach]
\label{prop:lamperti}
For $1\le p<\infty$, $p\neq 2$, every linear isometry of $L^{p}$ (and hence of $\ell^{p}$) is a
generalised permutation: a weighted composition operator which, in the finite-dimensional case,
acts as $(x_{k})\mapsto(\omega_{k}\,x_{\sigma(k)})$ with $\lvert\omega_{k}\rvert=1$ and $\sigma$
a permutation of the coordinates.
\end{proposition}

\noindent Banach gave this for $L^{p}[0,1]$; the general $\sigma$-finite statement is Lamperti's
\autocite{Banach1932, Lamperti1958}. Its consequence for our four real dimensions is decisive.
For every $p\neq 2$ the isometry group is, in the real case, the finite group of signed
permutations (the hyperoctahedral group $B_{4}$), and in the complex case the group
$\mathbb{T}^{n}\rtimes S_{n}$ of diagonal phases composed with permutations. Even the complex
group, though it contains the continuous torus $\mathbb{T}^{n}$, fails to act transitively on
the sphere: the orbit of a basis state $e_{1}$ is only the phased basis vectors
$\{\omega\,e_{k}\}$, a measure-zero union of circles. No continuous reversible transformation
carries a basis state into a generic superposition such as $(e_{1}+e_{2})/\sqrt{2}$. Reversible
continuity therefore \emph{fails} for every $p\neq 2$.

At $p=2$ the isometry group is $O(4)$ in the real picture and $U(2)$ in the complex one; the
identity component acts transitively on the unit sphere, and any two unit vectors are joined by
a continuous one-parameter subgroup. Among the $\ell^{p}$ geometries, reversible continuity is
satisfiable at $p=2$ alone.

Aaronson states the same selection in the quantum-foundations setting and sharpens it into a
physical, rather than merely aesthetic, necessity \autocite{Aaronson2004h}: the $2$-norm is the
unique $p$-norm admitting a nontrivial continuous group of norm-preserving linear maps, every
other $p$ yielding only generalised permutations; and the alternative of keeping linearity and
reversibility while \emph{dropping} norm preservation lands one in the full general linear group
$\mathrm{GL}(n,\mathbb{C})$, which permits superluminal signalling and collapses the complexity
class $\mathsf{PP}$ into polynomial time --- the same fact \S\ref{sec:continuity} uses to price
the classical--quantum cut. Bare reversibility on the amplitude space buys only
$\mathrm{GL}(n,\mathbb{C})$; what carves the unitary group $U(n)$ out of it is preservation of
the quadratic ($\ell^{2}$) form, and reversibility of the resulting maps is then a corollary
rather than a premise.

The conclusion is that the physical state manifold of a single qubit is the unit sphere of the
Euclidean norm,
\begin{equation}
  S^{3}=\Bigl\{\,x\in\mathbb{R}^{4}\;:\;\textstyle\sum_{i=1}^{4}x_{i}^{2}=1\,\Bigr\}
  =\bigl\{\,\psi\in\mathbb{C}^{2}\;:\;\lvert z_{1}\rvert^{2}+\lvert z_{2}\rvert^{2}=1\,\bigr\},
\end{equation}
with $z_{j}=x_{2j-1}+\mathrm{i}\,x_{2j}$, carrying the round metric induced from the ambient
inner product,
\begin{equation}
  \mathrm{d}s^{2}=\sum_{i=1}^{4}\mathrm{d}x_{i}^{2}\,\Big|_{S^{3}}
  =\operatorname{Re}\langle\mathrm{d}\psi,\mathrm{d}\psi\rangle,
  \qquad
  \langle\phi,\psi\rangle=\bar\phi_{1}\psi_{1}+\bar\phi_{2}\psi_{2}.
\end{equation}
One bound deserves explicit statement, in the interest of claiming no more than is proved:
reversible continuity fixes the \emph{metric} as $L^{2}$, but not the underlying field. Real and
quaternionic quantum theory equally satisfy reversible continuity and equally carry Euclidean
state-space metrics; the restriction to $\mathbb{C}$ is pinned elsewhere in the framework ---
the analytic continuation of the information measure and the additivity of phase information,
\S\ref{sec:self-reference} --- not here.

\subsection{Hardy's independent route}
\label{sec:hardy-route}

The same conclusion is reached by Hardy along a route that never invokes
Proposition~\ref{prop:lamperti} \autocite{Hardy2001a}. His reconstruction turns on the counting
relation $K=N^{r}$, where $N$ is the number of states distinguishable in a single measurement
and $K$ the number of fiducial probabilities needed to specify a state. The continuity axiom
excludes the classical case $r=1$: there the pure states are the $N$ vertices of a simplex and
the reversible maps merely permute them --- a discrete group of motions --- so there is no
continuous path between distinct pure states, which is the simplex face of the same $\ell^{1}$
isometry fact. A simplicity assumption then selects $r=2$, yielding complex quantum theory with
its quadratic norm and trace rule. The quadratic structure is not extracted from continuity as a
separate step; it is built into Hardy's first axiom (probabilities linear in the fiducial set)
together with $K=N^{2}$.

We thus have two genuinely independent demonstrations that reversible continuity forces the
$L^{2}$ metric: Hardy's operational $K=N^{2}$ counting, and the Lamperti--Aaronson
classification of the isometries. They converge on the same step by different mathematics, and
should be cited as two roads rather than one road traversed twice. Both roads start from
continuity. A third, in the next subsection, starts one step earlier, from the value continuum
that continuity itself rests on, and reaches the same metric without passing through the
transformation group at all.

\subsection{The statistical route: \v{C}encov's theorem and the every-basis lemma}
\label{sec:cencov-route}

The value continuum that \S\ref{sec:continuity} isolates as the one open premise reaches the
$L^{2}$ metric on its own, by a route that uses no transformation group. \v{C}encov's theorem
states that the only Riemannian metric on the simplex of probability distributions invariant
under coarse-graining --- under the Markov maps that merge or relabel outcomes --- is the
Fisher--Rao metric \autocite{bengtsson2006geometry}, and in the coordinates $x_{k} =
\sqrt{p_{k}}$ the Fisher--Rao metric is the round metric on the positive orthant of the unit
sphere. The premise the theorem consumes is that the simplex is a manifold, which is the value
continuum declared in \S\ref{sec:complex-amplitudes}, and the principle it uses, that
distinguishability does not depend on how outcomes are grouped, is one the Shannon bookkeeper of
that section already grants. What it does not supply is the passage from one orthant to the
whole projective space; that is the following lemma.

\begin{lemma}[The every-basis lemma]\label{lem:every-basis}
Let $g$ be a Riemannian metric on $\mathbb{CP}^{K-1}$ whose restriction to the positive real
orthant of every orthonormal basis --- the states $\sum_{k}\sqrt{p_{k}}\,e_{k}$ with $p$ in the
simplex --- is the Fisher--Rao metric. Then $g$ is the Fubini--Study metric, whose linear lift
is the $L^{2}$ inner product.
\end{lemma}

\begin{proof}
A tangent vector to $\mathbb{CP}^{K-1}$ at a state $\psi$ is represented by a vector $\phi$
orthogonal to $\psi$, and the curve $t \mapsto \cos t\,\psi + \sin t\,\phi$ with real $t$ lies
in the positive orthant of any orthonormal basis that contains $\psi$ and $\phi$. Along that
curve the Fisher--Rao metric is the round metric of the sphere in the coordinates $\sqrt{p}$,
which is the Fubini--Study metric restricted to real amplitudes; so $g$ and the Fubini--Study
metric agree on every tangent vector at every point, the direction $i\phi$ being reached by the
basis that contains $i\phi$ in place of $\phi$.
\end{proof}

The lemma's premise, that every orthonormal basis is a legitimate measurement and none is
privileged, is the fourth face of the pre-measurement regime seen from the measurement side
(\S\ref{sec:continuity}); in the stabiliser fragment of quantum theory it fails in the same way
everything else does, only the Pauli bases being measurements there. So the $L^{2}$ metric now
rests on two legs that share exactly one premise. Through reversible continuity it uses the
value continuum, closure and homogeneity (Proposition~\ref{prop:lamperti}); through \v{C}encov
it uses the value continuum, coarse-graining invariance and the equal standing of bases
(Lemma~\ref{lem:every-basis}). Neither leg needs anything else that is open, and nothing could
avoid the continuum, since a metric is a statement about one. That is the status of everything
downstream of this section, the Hopf fibrations included: derived from the completion every physicist grants, which carries no quantum content, and from
one premise about transformations that the framework derives in outline and has not yet proved
(Open Problem~\ref{op:limit-realisability}).

\subsection{The Hopf fibration and the Bloch sphere}
\label{sec:hopf-bloch}

The round metric is exactly what makes the global phase a clean symmetry. The phase action
$\psi\mapsto\mathrm{e}^{\mathrm{i}\theta}\psi$ is a one-parameter subgroup of
$U(2)\subset SO(4)$ and so acts on $S^{3}$ by \emph{isometries}, a statement that holds only
because the metric is the $L^{2}$ one just derived. Its orbits are great circles $S^{1}$, all of
circumference $2\pi$, and the quotient is the complex projective line,
\begin{equation}
  S^{3}\big/U(1)=\mathbb{CP}^{1}\cong S^{2},
\end{equation}
the Bloch sphere. This quotient is the complex Hopf fibration
\autocite{Hopf1931, Urbantke2003, mosseri2001hopf}
\begin{equation}
  S^{1}\hookrightarrow S^{3}\xrightarrow{\;h\;}S^{2},
\end{equation}
whose projection sends an amplitude to its Bloch vector,
\begin{equation}
  h(\psi)=\bigl(\langle\psi\rvert\sigma_{x}\lvert\psi\rangle,\;
                 \langle\psi\rvert\sigma_{y}\lvert\psi\rangle,\;
                 \langle\psi\rvert\sigma_{z}\lvert\psi\rangle\bigr)
        =\bigl(2\operatorname{Re}\bar z_{1}z_{2},\;
               2\operatorname{Im}\bar z_{1}z_{2},\;
               \lvert z_{1}\rvert^{2}-\lvert z_{2}\rvert^{2}\bigr).
\end{equation}
The Hopf map is quadratic in the amplitudes, and this is the geometric face of the Born rule:
the third component is $p_{0}-p_{1}$ with $p_{k}=\lvert z_{k}\rvert^{2}$, and the squaring is
the lowest-degree $U(1)$-invariant available on the fibre --- the record-theoretic derivation of
the same rule is \S\ref{sec:born-rule}'s. Because the metric is round, the projection is a
Riemannian submersion with totally geodesic fibres of constant length
\autocite{GluckWarnerZiller1986}: the same amount of phase is hidden at every point of the base,
so the phase is hidden \emph{uniformly}, not in a position-dependent way that the metric could
partially betray.

\subsection{The fibre as self-referential ignorance}
\label{sec:fibre-self-ignorance}

The interpretation the framework places on this geometry --- and here we are reading physics
into the mathematics, so the claim is interpretive rather than proved --- is that the $S^{1}$
fibre is the observer's own un-narratable self-ignorance. The global phase is the imaginary part
of the complex information measure $S=-\ln\lvert p\rvert-\mathrm{i}\theta$: information the
observer lacks about their own state. Lacking it, they cannot communicate it; there is no
\textsc{fanout} of information one does not possess. Uncommunicable, it cannot enter any
measurement record --- and this last step is no longer interpretive:
Proposition~\ref{prop:f1-dictionary} proves the global phase absent from every classical record
a copying structure induces --- so every observable is $U(1)$-invariant; and a $U(1)$-invariant
quantity is a function on the quotient. The physical state is therefore the orbit, not the
representative point---the base $S^{2}$---and the $S^{1}$ fibre carries exactly the inaccessible
self-referential phase. It is a null direction of statistical distinguishability: two global
phases are indistinguishable by any measurement whatever \autocite{wooters_1981}, so the fibre
is a gauge direction without physical content. This phase is irrecoverably discarded, and must
remain so: it is the relative phase, living in the $S^{3}$ fibre of the quaternionic Hopf
fibration at two qubits, that is genuinely recoverable, whereas the single-qubit global phase of
the complex fibration is gauge-private without remainder. The self-reference that opened the
derivation thus reappears, at its end, as the circle the observer can never see around.

%      % EDITED 2026-09-12 scripts/patch_revert_sedenion_repoint_2026-09-12T1600.py
% EDITED 2026-09-11 scripts/patch_undefined_section_refs_2026-09-11T2030.py
% =====================================================================
%  SECTION: Wigner's friend in the Hopf picture (v2; supersedes the
%  2026-06-07 draft). AUDITED 2026-07-06 (final pass, fifth stop).
%
%  PLACEMENT (reordered chain, JS 2026-07-06): position 4 --
%      \input{L2_metric_global_phase_section_2026_06_07}
%      \input{wigners_friend_in_the_hopf_picture_2026-06-13}
%      \input{derivation_of_the_quantum_postulates}
%  sec:l2-metric-fibre is REAL (backward, order-consistent);
%  sec:fanout-boundary DRAFTED 2026-07-06 (formerly the one undrafted section);
%  the old sec:monogamy placeholder note was VOID -- no such \ref in
%  the body -- and is removed. prop:f1-dictionary and
%  sec:hopf-justification are FORWARD references (position 8): first
%  F1 use now carries the deliberate-forward flag.
%
%  DEPENDENCY: \input{monogamy_aside_2026-06-07} -- received from JS
%  and audited 2026-07-06; the \input sits inside
%  sec:wf-reconciliation, immediately after op:tsirelson-geometry
%  (relocated 2026-07-06 from end-of-file, where the aside's opening
%  antecedent had been orphaned). A bare replace-all during the earlier
%  patch briefly mangled this header block, leaving a live \input in
%  it; repaired by index-splice, and the lesson recorded: anchors that
%  also occur in comments must be uniqueness-asserted.
%
%  AUDIT RECORD: CHSH arithmetic (2sqrt(1+C^2), PR => r^2 = -2) and the
%  C^2+r^2=1 chain verified; F1 wording canonical; the
%  obsolescence-made-cancellation echo two-way with
%  sec:self-reference; redundancy/Darwinism/DiBiagio-Rovelli/Brukner
%  positioning sound; Landauer placed on amplification, correctly.
%  Fixes: first-person voice removed; \textquote replaced with
%  portable quotes; IOEP given its prop ref; stale Tsirelson
%  mini-header (wrong key name Tsirelson1980) retired;
%  op:tsirelson-geometry MINTED from the honesty paragraph's unpriced
%  owed. Popper page-pin DISCHARGED 2026-07-07: the zip bundles carry
%  OCR .txt sidecars (the earlier "no text layer" verdict was wrong --
%  ZIP-of-JPEGs defeated the PDF tools, not the text); the quoted
%  phrase sits at Part II, p. 187, now pinned in the cite. Bonus from
%  the same dig: "C+" is Popper's own notation (Part II, scan pp.
%  15-17), so the L2 opening's attribution stands as written; "E" not
%  sighted in the sampled pages -- the inline gloss covers it either
%  way. Keys unique to this file (presumed in the .bib,
%  since the paper compiles): VerstraeteWolf2002,
%  Cirelson1980QuantumInequality, PopescuRohrlich1994, Held2021,
%  Brukner2021, zurek_2009, Landauer_1961, wigner_1961.
%  British English; \autocite; no bullets in prose.
% =====================================================================

\section{Wigner's friend in the Hopf picture: one state, two charts, and the redundancy ladder}
\label{sec:wigners-friend}

The framework dissolves the Wigner's-friend paradox \autocite{wigner_1961} not by
adjudicating between the two accounts but by recognising that they are written in
two charts on a single manifold, selected as the problem requires, much as one
moves between cylindrical and spherical coordinates without either being the true
one. The epistemological chart (private, non-clonable, phase-carrying) and the
ontological chart (clonable, public, record-bearing) are related by the
\textsc{fanout} boundary (Section~\ref{sec:fanout-boundary}), and that boundary
is not a sharp wall but a graded one, parametrised by how redundantly a piece of
relational content has been copied. Read this way, the fact that a \textsc{fanout}
has occurred for the friend but not for Wigner is not a contradiction to be
removed; it is the boundary behaving exactly as defined.

Let the system be a qubit $\lvert\psi\rangle_{S}=\alpha\lvert 0\rangle
+\beta\lvert 1\rangle$ and let the friend $F$, idealised, with deliberate
economy, as a single qubit, measure it in the computational basis. The
interaction is the controlled coupling that copies $S$'s value into $F$'s
register: the \textsc{copy} comultiplication of the diagrammatic calculus
\autocite{CoeckeKissinger2017}, and in the circuit picture the very gate that, in
the two-qubit Mosseri--Dandoloff geometry \autocite{mosseri2001hopf},
activates the quaternionic Hopf fibration $S^{3}\to S^{7}\to S^{4}$ and carries
the joint state from the surface of the entanglement ball to its interior. The
friend and Wigner read the two ends of that plunge.

Two idealisations in this set-up should be named at once, because the argument
later depends on seeing past them. The first is the friend as a single qubit,
which we return to as the redundancy-one endpoint of a graded boundary
(Section~\ref{sec:wf-redundancy}). The second, less obvious and more consequential,
is the system $S$ as a pristine qubit with no prior entanglement --- a relation
with no history. No real measured particle is of this kind. A particle is itself a
reified relation between prior events, and by the time the friend couples to it, it
already carries the completed internal structure such a relation possesses: in the
framework's terms, the full octonionic $A,B,C$ register that the $S^{7}$ sector
supplies as internal degrees of freedom (\S\ref{sec:space-of-computations}). The
pristine qubit is a schematic convenience, drawn structureless so that the logic of
a single link is visible; the ``system, friend, Wigner'' triple is likewise the
schematic of \emph{one link} in the chain of observers, not its base. The chain
does not bottom out in a structureless qubit --- there is no such thing --- and we
draw on this when we ask, later, what bounds the information a single such link can
carry.

\subsection{The friend's chart: a \textsc{fanout} of redundancy one}
\label{sec:wf-friend}

Relative to $F$, the coupling is a \textsc{fanout}: $F$ acquires a fact. What is
copied is gauge-invariant, the relative phase, the correlation between $F$ and
$S$, which is a genuine observable. The global phase is never the thing copied:
there is nothing observable there to duplicate, so the $S^{1}$ fibre of the
complex Hopf fibration $S^{1}\to S^{3}\to S^{2}$ of
Section~\ref{sec:l2-metric-fibre} remains $F$'s gauge-private self-ignorance,
untouched. In Hopf terms the division is clean, but it lives at two levels, which
Proposition~\ref{prop:f1-dictionary} (\S\ref{sec:hopf-justification} --- a deliberate forward
reference) separates. The copy operation itself is fibre-committed:
copying is faithful to phase, so what $F$'s \textsc{fanout} duplicates is a lift --- a chosen
fibre-point over the outcome --- never a bare point of the base. Yet everything the operation can
broadcast is base data alone: the record it induces is independent of the lift, so the $S^{1}$
freedom, though pinned in the copier, is absent from every classical record it yields. $F$'s classical record is a
base-point; the structure that produced it committed to a fibre-point; and the gap between the two
is precisely the gauge character of the fibre. The observation has a fifty-year-old
anchor: Popper argued that a predictor's inferior self-knowledge follows from the logical
impossibility of a system's imparting information about itself without thereby rendering it out of
date --- ``making this information obsolete'', in his words
\autocites{Popper1950a}[Part~II, p.~187]{Popper1950b}. Proposition~\ref{prop:f1-dictionary}
makes the impossibility structural rather than temporal, and localises it: the datum the copier
cannot put into its own record is not its state wholesale but exactly its fibre commitment, which
does not arrive stale --- it cancels. Conditioned on the outcome it reads, $S$ is for
$F$ a definite point on the Bloch sphere,
\begin{equation}
  \rho_{S}^{(F)}=\lvert 0\rangle\langle 0\rvert
  \quad\text{(given outcome }0\text{)},
\end{equation}
a pure state of Bloch radius $r=1$ on the surface of the ball, concurrence $C=0$.

This is a \textsc{fanout}, but the minimal one. The fact $F$ holds is a
\emph{relative} fact in the sense of relational quantum mechanics
\autocite{Rovelli1996a}: $S$ definite relative to $F$, not an observer-independent
record. It is one copy, redundancy $R=1$, and as such it remains reversible by an
agent who controls the joint system. The single-qubit friend is therefore not
outside \textsc{fanout}, as a careless reading might suggest; it is the $R=1$
endpoint of the \textsc{fanout} continuum---a \textsc{fanout} that has happened,
and a fact that is real for $F$, while being maximally still reversible from
outside.

\subsection{Wigner's chart: no copy has reached him}
\label{sec:wf-wigner}

Wigner treats the laboratory as closed and no copy of the relational content has
reached him. For him the coupling is the unitary
\begin{equation}
  \lvert\Psi\rangle
  =\alpha\lvert 0\rangle\lvert F_{0}\rangle
  +\beta\lvert 1\rangle\lvert F_{1}\rangle,
\end{equation}
the move up to the larger Hilbert space, $\mathbb{C}^{2}\to\mathbb{C}^{4}$,
$S^{3}\to S^{7}$, complex Hopf to quaternionic Hopf, with the fibre grown from
$S^{1}$ to $S^{3}$. What $F$ underwent as the acquisition of a fact, Wigner
describes as correlation deposited into a joint pure state on $S^{7}$. When he
asks specifically about $S$, he traces $F$ out, and the geometry does the
bookkeeping:
\begin{equation}
  \rho_{S}^{(W)}=\operatorname{Tr}_{F}\lvert\Psi\rangle\langle\Psi\rvert
  =\lvert\alpha\rvert^{2}\lvert 0\rangle\langle 0\rvert
  +\lvert\beta\rvert^{2}\lvert 1\rangle\langle 1\rvert,
  \qquad
  r=\bigl\lvert\,\lvert\alpha\rvert^{2}-\lvert\beta\rvert^{2}\bigr\rvert,
  \quad
  C=2\lvert\alpha\beta\rvert .
\end{equation}
Tracing out discards the $S^{3}$ fibre and carries the pure point on $S^{7}$ to a
mixed point in the interior of the Bloch ball at radius $r=\sqrt{1-C^{2}}$
\autocite{mosseri2001hopf}. The inward displacement is governed by the
concurrence: it is the relational content that has migrated into the $S$--$F$
correlation, stored in the recoverable $S^{3}$ fibre. Nothing has been destroyed;
it has been relocated into the fibre, which is why Wigner can in principle
interfere the branches and reverse the interaction, the content sits at low
redundancy and no further copy has fixed it.

\subsection{One state, two charts: \texorpdfstring{$C^{2}+r^{2}=1$}{C2+r2=1}}
\label{sec:wf-reconciliation}

The system is pure for the friend ($r=1$, on the surface) and mixed for Wigner
($r<1$, in the interior). These are not rival claims about a single fact; they are
the friend at the surface and Wigner at the interior of the \emph{same} state on
$S^{7}$, bound by the Mosseri--Dandoloff constraint
\begin{equation}
  C^{2}+r^{2}=1,
\end{equation}
in which the discrepancy in Bloch radius is exactly the concurrence Wigner
retains and the friend has resolved away. This is the Inside/Outside Equivalence
Principle (Proposition~\ref{prop:inside-outside}) made concrete, the same law written by the embedded participant and by
the external observer, and the relativity of facts of relational quantum
mechanics \autocite{Rovelli1996a} realised in geometry rather than asserted: the
\textsc{fanout} boundary lies \emph{between} $F$ and $W$, so a fact exists
relative to $F$ and not relative to $W$. Neither chart is privileged; one reads in
the epistemological or the ontological chart as the question requires, and their
mutual incompleteness is a theorem of the construction, not a defect of either.

\medskip
\noindent\emph{Aside: the Tsirelson bound and the absence of PR boxes.}
\label{sec:tsirelson}
A ball has a unique centre and no point beyond it. Maximal entanglement sits at
that centre \autocite{mosseri2001hopf}, $r=0$, $C=1$, and because the
Bloch radius cannot be negative, the concurrence cannot exceed unity. For a
two-qubit pure state the maximal CHSH correlator is
\begin{equation}
  \langle\mathrm{CHSH}\rangle_{\max}=2\sqrt{1+C^{2}}=2\sqrt{2-r^{2}}  
\end{equation}
\autocite{VerstraeteWolf2002},
which the centre maximises to $2\sqrt{2}$, the Tsirelson bound
\autocite{Cirelson1980QuantumInequality}. A Popescu--Rohrlich box \autocite{PopescuRohrlich1994},
with CHSH value $4$, would require $C^{2}=3$ and hence $r^{2}=-2$: a point two
units past the centre, which the ball does not contain. The super-quantum
correlation is excluded by the very positivity that places maximal entanglement
at the centre, one cannot trace out to a more-than-maximally-mixed reduced
state. The Tsirelson bound is in this sense the image, under the CHSH map, of the
boundedness of the $L^{2}$ state ball derived in
Section~\ref{sec:l2-metric-fibre}, and the prohibition on PR boxes shares its
origin with the roundness of the sphere: both trace to the reversible-continuity
metric, not to a separate principle.

\noindent This is the two-qubit, pure-state face of the bound, and its status is
worth stating plainly. It exhibits PR correlations as having no representing state
in the Hilbert geometry, and it recovers the value $2\sqrt{2}$ at the centre; but
forcing \emph{exactly} $2\sqrt{2}$ as a theorem of the geometry, rather than
reading it from the known concurrence--CHSH relation, and lifting the argument
from two qubits to the general Tsirelson bound, remain to be made rigorous. The
closest cognate in the literature is Held's derivation of the same bound from the
geometric algebra $G^{3}\cong\mathbb{H}$ \autocite{Held2021}, the same
quaternionic structure in a different language.

\begin{openproblem}[The Tsirelson bound as geometry]\label{op:tsirelson-geometry}
Derive $2\sqrt{2}$ as a theorem of the round state geometry rather than reading it off the
concurrence--CHSH relation, and lift the argument from the two-qubit pure case to the general
Tsirelson bound --- so that the exclusion of super-quantum correlation becomes, like the exclusion
exhibited above, a consequence of the $L^{2}$ ball's positivity (\S\ref{sec:l2-metric-fibre}) and
of nothing else. Held's quaternionic derivation \autocite{Held2021} is the nearest cognate and the
natural first tool.
\end{openproblem}

\input{monogamy_aside_2026-06-07}

\subsection{The redundancy ladder: from relative fact to objective fact}
\label{sec:wf-redundancy}

The friend's fact becomes a fact for Wigner only when Wigner copies it when he
measures $F$. It becomes a fact for a third party only when that party copies it
from Wigner, and so on up the von Neumann chain. The point at which the content
stops being negotiable is therefore not a point but a limit: once it has been
copied into enough mutually independent fragments that many parties can read it
without disturbing it, no one can practically un-copy it, and it is operationally
an objective fact. This is exactly quantum Darwinism \autocite{zurek_2009}:
objectivity is redundancy. It is not an external graft, because the framework's
account of decoherence already \emph{is} this picture---many fragments each take
their own copy of the clonable content, while their private phases, each
fragment's own self-ignorance, fail to add coherently and wash out.

This supplies the grading parameter that ``graded \textsc{fanout}'' had left
unnamed. \textsc{fanout} is graded not because the operation is fuzzy but because
redundancy $R$ runs continuously: $R=1$ is the single relational copy---the
friend's relative fact, reversible by Wigner---and $R\gg 1$ is the objective
fact, effectively irreversible because un-copying it would mean recovering every
fragment. The single-qubit friend is the $R=1$ endpoint; the macroscopic friend
is the same operation at large $R$. One axis, two ends. The relation to
Di~Biagio and Rovelli's distinction \autocite{DiBiagioRovelli2021} is then exact:
their \emph{stable} facts are the high-$R$ limit and their \emph{relative} facts
the low-$R$ regime, and the standard conditional-probability law they use to
define stability is the operational signature of redundancy having grown large.
Brukner's no-go theorem \autocite{Brukner2021}, that qubits are not observers,
sits on the same axis without conflict: it bites only against calling an $R=1$
relative fact an objective one, which the framework never does---the redundancy
ladder is precisely the structure that separates the relational definiteness a
single qubit does carry from the objectivity it does not.

The thermodynamics tracks the same axis. The single relational copy is unitary and
free; the $kT\ln 2$ cost \autocite{Landauer_1961} is incurred not by the friend's
act but by the amplification and stabilisation that drive $R$ upward, when
fragments are irreversibly committed. The relative fact is thermodynamically free;
the objective fact is what costs $kT\ln 2$.

\subsection{Phase promotion, and the bump in the carpet}
\label{sec:wf-phase-promotion}

One point keeps the load-bearing irrecoverability intact while explaining what
Wigner can do that the friend cannot. What Wigner recovers is \emph{not} the
$S^{1}$ self-ignorance; the global phase remains gauge for every observer and that
circle is spent for good. The quantity the friend resolves away is the relative
phase between the $\lvert 0\rangle$ and $\lvert 1\rangle$ components, and moving up
to $\mathbb{C}^{4}$ \emph{promotes} it: the single-qubit relative phase becomes the
relative phase between $\lvert 0\rangle\lvert F_{0}\rangle$ and
$\lvert 1\rangle\lvert F_{1}\rangle$, which lives in the recoverable $S^{3}$ fibre
(itself a circle bundle over $S^{2}$, so the promotion adds a two-sphere of
relative-phase directions to the lone gauge circle). Gauge phase at one level
becomes physical, interferable phase one level up---which is why Wigner can run
the interference experiment the friend cannot. But Wigner now carries his own
fresh global-phase $S^{1}$ at the top of the joint state: his own irreducible
self-ignorance. Ascending the tower does not eliminate the self-reference, only
displaces it one level higher, and the regress up the chain is precisely the
redundancy ladder of the previous subsection. This is the framework's bump in the
carpet: each observer sits atop a Hopf tower with a single gauge circle he can
never see around, and a measurement relocates the next observer's blind circle
into his own recoverable fibre without ever closing the regress.

Two cautions close the section. First, the recoverability Wigner
enjoys is the low-redundancy limit: for a macroscopic, high-$R$ friend the
reversal is exponentially hard and the fact is objective for all practical
purposes, so Wigner's ``no record'' is the $R\!\approx\!1$ idealisation of a
graded boundary, not a free lunch. Second, naming redundancy as the grading
parameter of \textsc{fanout}, and the phase-promotion reading itself, are
syntheses of the framework's parts: each rests on established pieces (the
clonability criterion, quantum Darwinism, the tensor-product promotion of phase)
but is assembled here, and is offered as the cleanest statement of why the two
observers never disagree about anything both can check.
% =====================================================================
%  SPLICE: "Where the chain goes next" -- forward-looking coda for
%  sec:wigners-friend, pointing to sec:sedenions-generations.
%  Created 2026-08-27 per JS.
%
%  THE POINT (JS, 2026-08-27). The Wigner's-friend charts stop at
%  Wigner because the Hopf tower runs out at the octonions, while the
%  text says the chain of observers obviously continues. The sedenions
%  show what the continuation adds: the observer of Wigner (fourth
%  wire, C (x) S) brings a NEW piece of structure -- the finite S_3 of
%  Brown's theorem, the generation label -- and the observer of the
%  observer of Wigner (fifth wire, the trigintaduonions) brings nothing
%  new in kind. This coda records that, so the section's termination at
%  Wigner reads as a fact about the tower rather than an economy of the
%  exposition.
%
%  WHAT THE ALGEBRA ACTUALLY SAYS (verified 2026-08-27, notebook T --
%  notebook_T_trigintaduonion_automorphisms_2026-08-27.py):
%    - Aut(S) = G_2 x S_3 (Brown 1967), DIRECT product; the new S_3 is
%      generated by the doubling involution and Brown's explicit
%      order-three map psi, which fixes 1 and the doubling unit and
%      rotates every (a, a u) plane, a imaginary, by a third of a turn.
%    - psi permutes the three octonion HALVINGS O, psi O, psi^2 O of S
%      (each doubled by the same unit; meeting only in R). It does NOT
%      permute the three copies through the quaternion core -- those are
%      cycled by a lifted G_2 element (the colour-line three-cycle); see
%      the sedenion section's two-triples remark and the corrections
%      ledger of the same date.
%    - Brown: Aut(A_{n-1}) x S_3 = Aut(A_n) with EQUALITY for n = 4, 5, 6
%      (as quoted in gresnigtgourlayvarma2023three, sec. 4.3); Eakin and
%      Sathaye prove Brown's conjecture in general (VERIFY-CITE for the
%      exact general form -- their step is Aut(A_{n+1}) = Aut(A_n) x G,
%      G in {Z/2, S_3}; over R with the standard parameters the branch is
%      S_3 at every rung >= 4, which is what notebook T checks directly
%      at rungs 4, 5, 6). Derivation algebra g_2 from the octonions on
%      (Schafer 1954; notebook T: dim Der = 14 at n = 3, 4, 5).
%  So "nothing new at the fifth wire" is stated PRECISELY: the fifth
%  doubling adds another finite S_3 and nothing else -- no connected
%  symmetry, no fibration; the increment is the same at every rung
%  thereafter. Not "silent": REPETITIVE. The prose says so.
%
%  STATUS: the algebra is ESTABLISHED (Brown; Eakin--Sathaye; Schafer)
%  and directly checked (notebook T). The reading of the repeated finite
%  factor as what the redundancy ladder consumes is a RHYME, flagged.
%  The reading of the fourth rung's discrete orientation as the
%  generation label is the sedenion section's CANDIDATE identification,
%  pointed to, not re-argued here.
%
%  PLACEMENT: a new \subsection at the END of sec:wigners-friend, after
%  the "Two cautions" paragraph. Because that paragraph reads "Two
%  cautions close the section.", the patch also changes that one word:
%  "close the section" -> "close the argument". Both anchors
%  uniqueness-asserted (the July lesson: anchors that also occur in
%  comments must be asserted unique).
%
%  CROSS-REFS OUT (all exist in the synced source except the two
%  MAP-LABELs):
%    sec:hopf-justification   -- the "two climbs" paragraph (length vs
%                                termination); this coda sharpens it
%    sec:wf-redundancy         -- the redundancy ladder (this file)
%    sec:wf-phase-promotion    -- phase promotion (this file)
%    sec:sedenions-generations -- the sedenion section (2026-08-27 rev.)
%    prop:termination          -- Hurwitz/Adams termination
%  MAP-LABELs:
%    app:notebook-record       -- the notebook-record appendix
%                                (appendices/appendix_notebook_record_
%                                2026-07-11.tex); substitute its label.
%  NEW KEYS (RIS delivered: cayley_dickson_automorphism_refs_2026-08-27.ris):
%    eakinsathaye1990automorphisms, schafer1954cayleydickson
%  IN-LIBRARY: brown1967generalized.
%
%  OPTIONAL CONSUMER-POINTER (not patched; one clause for the "two
%  climbs" paragraph of hopf_justification_section_2026-06-08.tex, after
%  "that chain has no top --- it is the bump in the carpet, unbounded in
%  length."):
%      "(unbounded in length, and --- as \S\ref{sec:wf-chain-type}
%      records --- bounded in type: after the fourth observer nothing
%      new in kind arrives)"
%
%  EDITED 2026-08-30 (per JS): one sentence added to the fourth-wire
%  paragraph anchoring the flavour sector's irremovable phase (the
%  CP-phase recognition) on this rung, with the development at
%  sec:sedenions-mixing (2026-08-30 revision). No new keys here.
%  British English; pure ASCII; \autocite; no bullets; spaced dashes.
% =====================================================================
% SPLICE-BEGIN wf-chain-type 2026-08-27

\subsection{Where the chain goes next: one more rung, then copies}
\label{sec:wf-chain-type}

The regress just described is unbounded in length, and nothing here shortens
it: the observer of Wigner carries a gauge circle of his own, and so does the
observer of that observer, as \S\ref{sec:hopf-justification} insists when it
separates the endlessness of the chain from the termination of the tower. But
the algebra says how much of the endless chain adds structure of a new kind,
and the answer is exact: one more step, and then none.

The count runs as follows. The triple system, friend, Wigner is three wires,
and three wires are the three Hopf rungs $\mathbb{C}$, $\mathbb{H}$,
$\mathbb{O}$ --- the completed link of the earlier caveat --- which is why the
charts above stop at Wigner: the chart of whoever looks in on the triple,
$S^{15}\to S^{8}$, is the last fibration the division algebras supply
(Proposition~\ref{prop:termination}). The observer of Wigner is a fourth wire,
and the fourth wire is the sedenions,
$(\mathbb{C}^{2})^{\otimes 4} = \mathbb{C}\otimes\mathbb{S}$
(\S\ref{sec:sedenions-generations}). No fibration comes with it --- the
sedenions have zero divisors and admit no map $S^{31}\to S^{16}$ --- so the
fourth observer gains no recoverable fibre of the tower's kind: the continuous
promotion of relative phase into a larger fibre, the mechanism of the previous
subsection, has no fourth instance. What the doubling adds instead is finite.
Brown's theorem gives $\mathrm{Aut}(\mathbb{S}) = G_{2}\times S_{3}$
\autocite{brown1967generalized}: the identity component is unchanged from the
octonions, and the entire gain is a component group, generated by an
order-three automorphism that fixes the new doubling direction and rotates
every imaginary direction of the completed link into its copy by a third of a
turn. The relative orientation of the link and its copy, which at the Hopf
rungs was a continuous phase promoted into a fibre, survives at the fourth rung
only as a discrete choice among three, and that choice is the generation label
(\S\ref{sec:sedenions-generations}). The fourth observer, then, sees no new
force and no new fibre; he sees the same link three ways, and cannot gauge the
difference, because a component group has no Lie algebra. His rung is also where
the flavour sector's one irremovable phase finds its anchor: three generations
are the threshold at which the mixing matrix first carries a phase no convention
removes, and \S\ref{sec:sedenions-mixing} reads that phase --- through the
paper's negativity--contextuality chain --- as this wire's contextual cargo,
real only from one rung above.

The observer of the observer of Wigner is a fifth wire, the trigintaduonions,
and here the tower has nothing new to say. Brown proved that the next two
doublings each add exactly one further $S_{3}$ and nothing else, Eakin and
Sathaye that Brown's pattern continues at every rung, and the derivation
algebra --- the Lie algebra of the identity component --- is $\mathfrak{g}_{2}$
from the octonions onward
\autocite{brown1967generalized,eakinsathaye1990automorphisms,schafer1954cayleydickson};
the fifth and sixth rungs are also checked directly in the notebook record
(Appendix~\ref{app:notebook-record}). So the fifth observer's doubling repeats
the fourth's gift --- another third-of-a-turn symmetry between the four-wire
structure and its copy, commuting with everything below --- and adds no
connected symmetry and no fibration. Nothing new in kind arrives at the fifth
wire or at any wire after it. The increments of the tower are, in order, a
force, a force, a label, and then the same finite permutation of copies at
every step. That is the precise sense in which the chain ``obviously
continues'': it continues, and after the fourth observer it has nothing further
to announce. The framework reads the repeated factor as what the redundancy
ladder of \S\ref{sec:wf-redundancy} consumes --- a permutation symmetry among
copies, which is the freedom a record has and a type does not --- and records
that reading as a structural resonance rather than a derivation.

The two climbs of \S\ref{sec:hopf-justification} are thereby sharpened. The
chain of observers is unbounded in length and bounded in type: three forces by
the third wire, one family label by the fourth, copies thereafter. The observer
of Wigner is the last observer whose arrival changes what there is.

% SPLICE-END wf-chain-type 2026-08-27

% =====================================================================
%  FIGURE: the Wigner's-friend chain mapped onto the Cayley--Dickson
%  ladder. Created 2026-08-30 per JS: the Hopf-fibration view of
%  Wigner's friend has been load-bearing for months without being drawn;
%  the fibrations themselves are DEFERRED to Szangolies's introduction
%  (szangolies2025standardmodel, in library), and this figure carries
%  only the mapping.
%
%  CONTENT (left to right, one panel per chart, staircase rising):
%    Panel 1: F's chart of S            -- S^3 -> S^2, fibre S^1,  C, one wire
%    Panel 2: W's chart of {S,F}        -- S^7 -> S^4, fibre S^3,  H, two wires
%    Panel 3: O_W's chart of {S,F,W}    -- S^15 -> S^8, fibre S^7, O, three wires
%             (the last fibration)
%    Panel 4: the next observer's chart of {S,F,W,O_W} -- C^16 = C (x) S,
%             no fibration (S^31 -/-> S^16); the doubling's S_3 badge:
%             the generation label, a third of a turn.
%  Each describer perches on his sealed box with a dashed circle over
%  his head (his own gauge circle -- the self-referential ignorance he
%  cannot see around); dashed promotion arrows carry each circle into
%  the next chart's fibre ("gauge here -> interferable one rung up",
%  sec:wf-phase-promotion); the third arrow terminates at "no fourth
%  fibration", with what survives being the discrete third-of-a-turn
%  (sec:wf-chain-type, sec:sedenions-generations). The regress
%  continues to the right: copies only.
%
%  PREAMBLE REQUIREMENT (confirm in main.tex; add if absent):
%      \usepackage{tikz}
%      \usetikzlibrary{arrows.meta}
%  PLACEMENT: \input by patch_wigner_hopf_ladder_fig_2026-08-30.py,
%  after the two-idealisations paragraph of sec:wigners-friend, with
%  the accompanying deferral paragraph (in the patch).
%  CROSS-REFS OUT (all real): sec:wf-phase-promotion,
%  sec:wf-chain-type, sec:sedenions-generations.
%  KEY: szangolies2025standardmodel (in library).
%  British English; pure ASCII; \autocite in the caption.
% =====================================================================
\begin{figure}[t]
\centering
\resizebox{\textwidth}{!}{%
\begin{tikzpicture}[>=Stealth, line cap=round, line join=round,
    every node/.style={font=\scriptsize},
    blind/.style={densely dashed, draw=black!75},
    promote/.style={->, densely dashed, black!60, semithick},
    sight/.style={->, dotted, black!55}]
% ---- stick-figure pic --------------------------------------------------
\tikzset{stick/.pic={
  \draw[pic actions] (0,0.44) circle (0.085);
  \draw[pic actions] (0,0.355) -- (0,0.10);
  \draw[pic actions] (-0.12,0.20) -- (0,0.29) -- (0.12,0.20);
  \draw[pic actions] (-0.10,-0.10) -- (0,0.10) -- (0.10,-0.10);
}}
% ---- the staircase (Cayley--Dickson ladder) ---------------------------
\draw[black!35, thick]
  (-0.35,-0.55) -- (3.27,-0.55) -- (3.27,0.25) -- (6.89,0.25)
  -- (6.89,1.05) -- (10.51,1.05) -- (10.51,1.85) -- (13.60,1.85)
  -- (13.60,2.65) -- (16.20,2.65);
\node[anchor=north] at (1.25,-0.60) {$\mathbb{C}$ --- one wire};
\node[anchor=north] at (4.87,0.20)  {$\mathbb{H}$ --- two wires};
\node[anchor=north] at (8.49,1.00)  {$\mathbb{O}$ --- three wires};
\node[anchor=north] at (12.11,1.80) {$\mathbb{S}$ --- four wires};
\node[anchor=north, black!55] at (14.45,2.60) {$\mathbb{T},\dots$ --- five wires on};
% =======================================================================
% Panel 1: F's chart of S
% =======================================================================
\draw[rounded corners=2pt, fill=black!5, draw=black!40] (0,0) rectangle (2.5,1.32);
\filldraw[black!70] (1.25,0.72) circle (0.075);
\node[font=\tiny] at (1.25,0.45) {$S$};
\pic[black] at (2.05,1.42) {stick};
\node[anchor=west] at (2.22,1.62) {$F$};
\draw[sight] (1.95,1.60) -- (1.45,1.02);
\draw[blind] (2.05,2.14) circle (0.11);
\node[anchor=east, align=right, font=\tiny, black!75] at (1.82,2.14)
  {his own gauge circle:\\ the phase he cannot\\ see around};
\node[anchor=south west, font=\tiny, black!75] at (0.02,1.34) {fibre $S^{1}$};
\node[anchor=north] at (1.25,-0.06) {$S^{3}\to S^{2}$};
% =======================================================================
% Panel 2: W's chart of {S,F}
% =======================================================================
\draw[rounded corners=2pt, fill=black!5, draw=black!40] (3.62,0.80) rectangle (6.12,2.12);
\filldraw[black!70] (4.07,1.52) circle (0.075);
\node[font=\tiny] at (4.07,1.25) {$S$};
\pic[black!55, scale=0.85] at (4.85,1.16) {stick};
\node[font=\tiny, black!55] at (4.85,0.94) {$F$};
\pic[black] at (5.67,2.22) {stick};
\node[anchor=west] at (5.84,2.42) {$W$};
\draw[sight] (5.57,2.40) -- (5.05,1.82);
\draw[blind] (5.67,2.94) circle (0.11);
\node[anchor=south west, font=\tiny, black!75] at (3.64,2.14) {fibre $S^{3}$};
\node[anchor=north] at (4.87,0.74) {$S^{7}\to S^{4}$};
\draw[promote] (2.17,2.14) to[out=15,in=170] (3.60,2.24);
\node[anchor=south, align=center, font=\tiny, black!60] at (2.85,2.32)
  {gauge below,\\ interferable above};
% =======================================================================
% Panel 3: O_W's chart of {S,F,W}
% =======================================================================
\draw[rounded corners=2pt, fill=black!5, draw=black!40] (7.24,1.60) rectangle (9.74,2.92);
\filldraw[black!70] (7.69,2.32) circle (0.075);
\node[font=\tiny] at (7.69,2.05) {$S$};
\pic[black!55, scale=0.85] at (8.35,1.96) {stick};
\node[font=\tiny, black!55] at (8.35,1.74) {$F$};
\pic[black!55, scale=0.85] at (8.95,1.96) {stick};
\node[font=\tiny, black!55] at (8.95,1.74) {$W$};
\pic[black] at (9.29,3.02) {stick};
\node[anchor=west] at (9.46,3.22) {$O_{W}$};
\draw[sight] (9.19,3.20) -- (8.65,2.62);
\draw[blind] (9.29,3.74) circle (0.11);
\node[anchor=south west, font=\tiny, black!75] at (7.26,2.94) {fibre $S^{7}$};
\node[anchor=north, align=center] at (8.49,1.54) {$S^{15}\to S^{8}$\\[-1pt]
  {\tiny the last fibration}};
\draw[promote] (5.79,2.94) to[out=15,in=170] (7.22,3.04);
% =======================================================================
% Panel 4: the next chart of {S,F,W,O_W} -- sedenions, no fibration
% =======================================================================
\draw[rounded corners=2pt, fill=black!5, draw=black!40] (10.86,2.40) rectangle (13.36,3.72);
\filldraw[black!70] (11.24,3.12) circle (0.075);
\node[font=\tiny] at (11.24,2.85) {$S$};
\pic[black!55, scale=0.80] at (11.78,2.78) {stick};
\node[font=\tiny, black!55] at (11.78,2.56) {$F$};
\pic[black!55, scale=0.80] at (12.28,2.78) {stick};
\node[font=\tiny, black!55] at (12.28,2.56) {$W$};
\pic[black!55, scale=0.80] at (12.78,2.78) {stick};
\node[font=\tiny, black!55] at (12.78,2.56) {$O_{W}$};
\pic[black!45] at (14.35,2.78) {stick};
\draw[sight] (14.18,3.02) -- (13.42,3.12);
\draw[blind] (14.35,3.50) circle (0.11);
\node[anchor=west, align=left, font=\tiny, black!60] at (14.58,3.48)
  {the next observer,\\ and all after:\\ copies only};
\node[anchor=north] at (12.11,2.34) {$\mathbb{C}^{16}=\mathbb{C}\otimes\mathbb{S}$};
% the broken third promotion
\draw[promote] (9.41,3.74) to[out=15,in=185] (10.68,3.86);
\node[font=\scriptsize, black!60] at (10.80,3.92) {$\times$};
\node[anchor=south west, font=\tiny, black!75] at (10.88,3.74)
  {no fourth fibration: $S^{31}\not\to S^{16}$};
% the S3 badge: three dots at 120 degrees with a rotation arrow
\draw[black!75] (11.45,4.30) circle (0.17);
\filldraw[black!75] (11.45,4.47) circle (0.028);
\filldraw[black!75] (11.303,4.215) circle (0.028);
\filldraw[black!75] (11.597,4.215) circle (0.028);
\draw[->, black!75] (11.67,4.36) arc[start angle=20, end angle=100, radius=0.235];
\node[anchor=west, align=left, font=\tiny, black!75] at (11.70,4.36)
  {$S_{3}$: the generation label\\ (a third of a turn)};
\draw[promote, black!70] (10.86,3.96) to[out=60,in=200] (11.24,4.22);
\end{tikzpicture}%
}
\caption{The Wigner's-friend chain mapped onto the Cayley--Dickson ladder. Each
observer holds a chart of the sealed box below him, one rung further up the
ladder: the friend's chart of the system is the complex Hopf bundle, Wigner's
chart of the sealed laboratory the quaternionic, and the chart of the observer
of Wigner --- of the completed triple --- the octonionic, the tower's last
fibration. Every chart-holder carries a dashed circle over his head, his own
gauge phase: the self-referential ignorance he cannot see around. Measurement
promotes each circle into the next chart's fibre, where what was gauge below is
interferable above (\S\ref{sec:wf-phase-promotion}); the third promotion has no
fibration to land in --- the sedenions have zero divisors and no map
$S^{31}\to S^{16}$ --- and what survives of it is discrete: the doubling's
$S_{3}$, a third of a turn rotating the link into its copy, the generation
label (\S\ref{sec:wf-chain-type}, \S\ref{sec:sedenions-generations}). The
regress of circles continues unbounded to the right, but from the fifth wire
onward it adds copies only. The friend drawn as a single qubit is the
redundancy-one idealisation discussed in the text; for the fibrations
themselves, deferred rather than repeated here, see the introduction in
\autocite{szangolies2025standardmodel}.}
\label{fig:wf-hopf-ladder}
\end{figure}

%      % EDITED 2026-10-29 scripts/patch_continuity_tower_premise_2026-10-29T1500.py
% EDITED 2026-09-20 scripts/patch_born_cost_ledger_v2_2026-09-20T0900.py
% EDITED 2026-09-15 scripts/patch_value_continuum_declared_cencov_2026-09-15T1540.py
% EDITED 2026-09-15 scripts/patch_continuity_connected_state_space_2026-09-15T1430.py
% =====================================================================
%  SECTION: Derivation of the quantum postulates (the assembly chapter).
%  AUDITED 2026-07-06 (final pass, sixth stop).
%
%  PLACEMENT (reordered chain, JS 2026-07-06): position 5 --
%      \input{wigners_friend_in_the_hopf_picture_2026-06-13}
%      \input{derivation_of_the_quantum_postulates}
%      \input{division_algebras_section_2026-07-03}
%  Postulate 6 hands off to division-algebras, its immediate successor.
%
%  LABEL RESOLUTIONS (2026-07-06): the old placeholder
%  sec:complex-amplitudes is now minted at its source (self-reference's
%  negativity subsection), so the ref here resolves unchanged.
%  sec:born-rule COLLIDED with the label minted 2026-07-06 on
%  self-reference's Born-from-fibre-privacy subsection; resolution: two
%  routes, two labels -- Postulate 3 here is renamed
%  sec:born-invariance (the U(1)/lowest-degree geometric route) and
%  cross-stitched to sec:born-rule (the record-theoretic route), their
%  agreement flagged as over-determination. sec:local-tomography also
%  minted at source for Postulate 6's operational clause.
%
%  AUDIT RECORD: the section's status discipline was PRESCIENT -- the
%  continuity arrow was already "owed ... enters as a principle" (the
%  exact voice JS mandated for L2), and both of its confessed debts
%  have since been minted as numbered problems; the closing now cites
%  them (op:limit-realisability, op:tsirelson-geometry): "named, and
%  numbered". Postulate 6 now carries BOTH composition duty-holders
%  (Hurwitz multiplicativity + local tomography), matching the
%  continuity section's caveat exactly. Stone's theorem cited
%  (stone1932one, RIS delivered); time-as-operation-count forward
%  flag added; \textquote made portable; opener formalised.
%
%  PROVENANCE NOTE (relocated from the body, 2026-07-06, JS's words):
%  "Using Hopf fibrations to get the Born rule is something Claude AI
%  came up with on its own without my prompting." Preserved here as the
%  historical record. DISCLOSURE STATUS (JS, 2026-07-06): AI use is
%  already acknowledged in the second paragraph of the paper's
%  introduction; the remaining recommendation is only that, at
%  submission, the same fact be echoed in whatever formal
%  Declarations / Acknowledgements slot the target journal's template
%  mandates -- several publishers now require it there specifically,
%  in addition to any in-text mention.
%
%  British English; \autocite; no bullets in prose.
% =====================================================================

% ADDED 2026-07-12: closing paragraph "The postulates in module
% coordinates" -- the Clifford-module reading (skeleton + two
% supplements), with the explanatory arrow guarded per JS: the field
% comes from the Shannon/additivity/commutativity step (self-ignorance),
% and the left-action composability result is recorded as that same
% requirement in operator clothing, not an independent selection of C.
\section{Derivation of the quantum postulates}
\label{sec:postulates}
Enough pieces of the jigsaw are now assembled to derive the postulates of quantum mechanics.

The Dirac--von Neumann postulates are ordinarily laid down as independent
assumptions. In this framework they are not assumptions but a single chain: each
follows from the Inside/Outside Equivalence Principle by one move from the
preceding step. We assemble the chain here and mark the status of each link as
proved, structural, or still owed, so that no step claims more than it has earned.

The chain begins below the postulates. An embedded observer, required by the
equivalence principle to describe a totality that includes itself, must assign a
probability to a proposition about its own future state. The Lawvere--Yanofsky
diagonal \autocite{Yanofsky2003} and Popper's predictor argument
\autocite{Popper1950a, Popper1950b} show that no consistent classical assignment
exists: the self-referential proposition is undecidable to the observer before
the fact. Representing the result of that undecidable self-prediction forces a
signed, non-classical probability, which carries the observer's irreducible
self-ignorance. From this single obstruction the postulates follow.

\paragraph{Postulate 1 (complex Hilbert space).} Two independent steps fix the state space. The
field is forced first: a signed probability, read through the Shannon information measure and
analytically continued, yields a complex amplitude $S=-\ln\lvert p\rvert-\mathrm{i}\theta$,
whose phase $\theta$ is the self-referential ignorance (Section~\ref{sec:complex-amplitudes}).
The metric is forced second, and separately: the principle of reversible continuity
(\S\ref{sec:continuity}) --- that a continuous reversible transformation connects any two pure
states --- selects the $L^{2}$ norm uniquely. For every $\ell^{p}$ with $p\neq 2$ the linear
isometries are generalised permutations \autocite{Banach1932, Lamperti1958}, a group that cannot
connect a basis state to a superposition; only at $p=2$ is the norm-preserving group connected
and transitive on the sphere \autocite{Aaronson2004h}, and Hardy's operational $K=N^{2}$
counting reaches the same quadratic norm by a second route from the same continuity premise
\autocite{Hardy2001a}, and \v{C}encov's theorem with the every-basis lemma reaches it from the
value continuum alone (\S\ref{sec:cencov-route}). The state manifold is therefore the unit
sphere of the Hermitian $2$-norm in $\mathbb{C}^{n}$, taken modulo the global phase
(Section~\ref{sec:l2-metric-fibre}). Placing the metric here, before the Born rule rather than
after it, is deliberate: the Born derivation presupposes the round sphere, so the metric must be
anchored independently to avoid circularity. Two bounds belong with this step. Reversible
continuity fixes the metric, not the field---real and quaternionic quantum theory satisfy it
equally---so the restriction to $\mathbb{C}$ rests on the Shannon step, not on continuity. And
deriving reversible continuity itself from self-reference is the one arrow the chain still owes
--- now priced exactly: one open premise, that every transformation of the regime is a phase of an inside structure,
Open Problem~\ref{op:limit-realisability}, with everything downstream of it proved
(\S\ref{sec:continuity}). Until it is closed, continuity enters as a principle, and the weight
of the $L^{2}$ conclusion is carried jointly with the isometry and counting arguments and with
the statistical route.

\paragraph{Postulate 2 (Hermitian observables).}
A measurement result must be communicable: it must be copyable from one observer
to another, the \textsc{copy} operation of the diagrammatic calculus
\autocite{CoeckeKissinger2017}. The imaginary part of the amplitude is the
private self-referential phase and cannot be broadcast; only the real,
gauge-invariant content can. Observables whose spectra are real and therefore
communicable are exactly the self-adjoint operators. Hermiticity is thus the
condition that a quantity be the kind of thing a measurement can report.

\paragraph{Postulate 3 (Born rule).}
\label{sec:born-invariance}
Outcome probabilities must be invariant under the global phase, the $U(1)$
self-ignorance, realised geometrically as the $S^{1}$ fibre of the complex Hopf
fibration $S^{1}\to S^{3}\to S^{2}$ over the Bloch sphere
\autocite{Hopf1931, Urbantke2003}. There is no polynomial $U(1)$-invariant of degree one,
and the non-polynomial $\lvert\psi\rvert$ is excluded by the analyticity declared at
\S\ref{sec:complex-amplitudes}; the lowest-degree polynomial invariant is quadratic, so the
probability must be $p=\lvert\psi\rvert^{2}$. The squaring is forced
rather than merely permitted, and the argument holds for a single qubit, where
Gleason's theorem is silent. The Born quadratic is not a new structure: it is the
same $L^{2}$ form that fixed the metric in Postulate~1, now read as the
normalisation of the round sphere, and the fibre over which the phase is averaged
has constant length, so the phase is hidden uniformly
\autocite{GluckWarnerZiller1986}. This step is a structural argument rather than a
theorem in the strict sense, and is flagged as such. The record-theoretic derivation of the same
quadratic --- squared magnitudes from fibre privacy, Proposition~\ref{prop:f1-dictionary} --- is
given at \S\ref{sec:born-rule}, and the cost-ledger derivation --- the record as the closed loop of
a history and its return leg, the exponent the number of legs --- at \S\ref{sec:born-cost-ledger}:
three routes to one rule, each holding for a single qubit, the over-determination this framework
treats as evidence.

\paragraph{Postulate 4 (state reduction).}
Measurement is a \textsc{fanout}: the gauge-invariant relational content is copied
into a record, converting a private epistemological amplitude into a public
fact. For the observer who performs it, acquiring that fact is Bayesian updating,
and the projection postulate is its expression. The reduction is perspectival---a
fact relative to the observer in the sense of relational quantum mechanics
\autocite{Rovelli1996a}---and it is graded by redundancy. A single copy is a
relative fact, still reversible by an agent controlling the joint system; redundant
proliferation across many independent fragments makes the fact objective, which is
quantum Darwinism \autocite{zurek_2009} and the stable/relative-fact distinction of
Di~Biagio and Rovelli \autocite{DiBiagioRovelli2021}. ``Collapse'' and unitary
evolution are then the two charts---inside and outside the \textsc{fanout}
boundary---of one process, as the Wigner's-friend analysis shows and \S\ref{sec:fanout-boundary} formalises
(Section~\ref{sec:wigners-friend}); the irreversibility and its $kT\ln 2$ cost
\autocite{Landauer_1961, wigner_1961} accrue with redundancy rather than at the
single act. The reduction postulate is therefore recovered as a chart-dependent
description, not a fundamental dynamical law; the identification of redundancy as
the grading parameter is a synthesis of the framework's parts and is flagged as
such.

\paragraph{Postulate 5 (unitary evolution).}
Conservation of probability is preservation of the $L^{2}$ inner product, and the
inner-product-preserving linear maps are the unitaries $U(n)$---the same
isometry fact that fixed the metric in Postulate~1. Continuity in time --- operation count,
\S\ref{sec:space-of-computations}, a forward reference --- together with Stone's theorem
\autocite{stone1932one} gives $U(t)=\mathrm{e}^{-\mathrm{i}Ht}$, the Schr\"odinger
equation. Reversibility is here a corollary, not a premise: bare reversibility on
the amplitude space yields only the general linear group, and it is preservation
of the quadratic form that carves $U(n)$ out of it.

\paragraph{Postulate 6 (composite systems).}
Systems combine by the tensor product because the division algebra norm is
multiplicative \autocite{hurwitz1898composition}: composition preserves the $L^{2}$ structure
only if states multiply. The operational half closes from the other side: the locality of
tomography derived at \S\ref{sec:local-tomography} says product states determine composites,
which is the tensor product's signature. Cayley--Dickson doubling carries $\mathbb{C}$ to the next
rung, and the ascent of the Hilbert space is the ascent of the Hopf tower, complex
to quaternionic, used in the Wigner's-friend construction
\autocite{mosseri2001hopf}; the ladder itself is the business of the section that follows
(\S\ref{sec:division-algebras}).

\paragraph{Consequence (the Tsirelson bound).}
The postulates fix not only the form of the theory but the strength of its
correlations. In the two-qubit ball the maximal entanglement sits at the centre,
and a Popescu--Rohrlich box \autocite{PopescuRohrlich1994} would require a point
beyond it; the CHSH correlator is capped at $2\sqrt{2}$
\autocite{Cirelson1980QuantumInequality, VerstraeteWolf2002, Held2021}, with the value recovered
at the centre and its general derivation still owed
(Section~\ref{sec:wf-reconciliation}).

\medskip
\noindent The result is that the six postulates are one chain, not six premises.
Of the links, the field and the Born quadratic rest on the framework's signed
probability and $U(1)$-invariance arguments; the $L^{2}$ metric, the unitary group
and the tensor product are over-determined by reversible continuity, the isometry
classification and Hurwitz composition; state reduction is recovered as the
chart-relative face of \textsc{fanout}. The two outstanding debts are the
derivation of reversible continuity from self-reference and the forcing of the Tsirelson value as
a theorem rather than a recovered number --- named, and numbered:
Open Problems~\ref{op:limit-realisability} and~\ref{op:tsirelson-geometry} --- so that a reader
can see exactly where the chain is complete and where it is not.

\paragraph{The postulates in module coordinates.}
One reading of the chain deserves a closing statement, because it positions the derivation
against a mathematical skeleton the reader may know well, and because stating it carelessly would
invert the chain's explanatory arrow. In the coordinates of the division-algebra programmes
(\S\ref{sec:algebraic-dictionary}), complex Hilbert-space quantum mechanics is real Clifford
module theory plus exactly two supplements. The skeleton is Clifford throughout: the operator
algebra of $n$ qubits is the Clifford algebra of $2n$ generators with the state space its own
spinor module, half-integer spin and fermionic statistics are the module's double cover and its
graded tensor, and the classically simulable stabiliser core is the algebra's combinatorics. What
the skeleton does not contain is, first, the choice of complex structure --- the polarisation $J$
that turns a real module into a complex Hilbert space, which standard quantum theory and the
algebraic programmes alike \emph{choose} --- and, second, the Born weights on the fibre. Those two
supplements are precisely what this section has derived, and derived from self-ignorance rather
than from algebra: the field from the Shannon step (information about one's own undecidable state
must still be information --- it must add, addition commutes, and so the amplitude algebra is
commutative, excluding the quaternionic bookkeeping an observer could never balance), and the
quadratic from fibre privacy. The algebraic echo --- that the algebra's own left actions compose
only along a commutative line, the composability result recorded with the dictionary --- is not
an independent selection of $\mathbb{C}$: composition of left actions fails off that line exactly
because the doubled slot carries the \emph{reversed} product, and order-dependence is the same
disease in operator clothing that non-additivity is in information clothing. The mathematics is
the epistemic argument's shadow, welcome as consistency and as a bridge to the module-side
literature (Appendix~\ref{app:winding-number} walks the bridge for charge), but the origin of the
complex amplitude is the observer's ignorance of their own state, and the arrow runs from there
to the algebra, never the reverse.
%      % =====================================================================
%  SECTION: The Division-Algebra Ladder
%  Drafted 2026-07-03. One of the three foundations sections
%  (sec:self-reference, sec:division-algebras, sec:fanout-boundary);
%  this one resolves the dangling \S\ref{sec:division-algebras} in the
%  Hopf-justification section, whose header names it "the
%  Cayley--Dickson / Hurwitz step".
%
%  PLACEMENT in main.tex (REORDERED per JS 2026-07-06): position 6,
%  opening Act II ("why stop at C?") right before its first consumer --
%      \input{self_reference_section_2026-07-04}
%      \input{continuity_section_2026-07-04}
%      \input{L2_metric_global_phase_section_2026_06_07}
%      \input{wigners_friend_in_the_hopf_picture_2026-06-13}
%      \input{derivation_of_the_quantum_postulates}
%      \input{division_algebras_section_2026-07-03}   <-- moved here
%      \input{spacetime_from_non-computability}       <-- first consumer
%                                                          (S^7 exclusion)
%      \input{hopf_justification_section_2026-06-08}
%  The section is position-portable: all adjacency phrasing was
%  retargeted to explicit \S\ref on 2026-07-06. The S^7 hand-off
%  sentence to sec:space-of-computations WRITTEN 2026-07-06 (spacetime
%  file read; the closing now hands first to its immediate consumer,
%  then to hopf-justification). Upload verifications COMPLETE: L2 and
%  quantum-postulates do not consume the ladder; WF's references are
%  forward-voiced or self-carrying.
%
%  DIVISION OF LABOUR (deliberate, to avoid overlap):
%   - sec:self-reference (undrafted): self-reference -> complex
%     amplitudes and the Born rule; why not a real base. This section
%     TAKES that as given and starts at C.
%   - THIS section: the algebra -- the doubling construction, the
%     ledger of losses-as-physics, the internal-algebras-not-a-new-QM
%     scoping (Adler pre-empt), the two appearances of the ladder
%     (register / measurement climb, stated lightly), and Hurwitz as
%     the algebraic termination with the committed zero-divisor
%     reading.
%   - sec:hopf-justification: the geometry (Adams, parallelisability,
%     the fibrations), the join of the two classifications, the
%     entanglement-sensitivity literature, the C-direction honest debt,
%     and the epistemic reading of termination.
%
%  COMMITTED READINGS USED (sources re-read 2026-07-03):
%   - Loss ledger and level dictionary: cayley_dickson_qubits_for_review
%     sec.5 ("each step loses a symmetry, and each lost symmetry is
%     identified with a physical content").
%   - Termination reading: grounding_argument_gauge_group ("the
%     sedenions have zero divisors and cannot support consistent
%     information flow around loops") -- leads; norm-composition-as-
%     probability given as its amplitude-level restatement.
%   - Internal-not-base scoping: cayley_dickson_qubits_for_review Q2 and
%     sec.3.2 caveat (Adler objection; C-qubit treated geometrically).
%
%  HONEST DEBT: ladder proved closed at the top (Hurwitz), derived at
%  the floor. The rung-compulsion revisit-trigger FIRED 2026-07-06:
%  sec:self-reference and op:sigma-functor now exist, so the thesis is
%  minted as op:ladder-compulsion (with the CD-conjugation-as-sigma
%  candidate route), inheriting op:sigma-functor's status.
%
%  AUDITED 2026-07-06 (final pass, third stop): all four citations
%  used; all cross-refs live; CD facts, dims, alternativity, Hurwitz
%  statement verified; Adler subsection delivers everything the
%  register's compressed pointer promises; file pure ASCII; no bare
%  "frame". Fixes: opening rewritten for both predecessors (continuity
%  was inserted between after drafting); "associator's non-vanishing is
%  mass itself" recalibrated to the register's debt-against-selection
%  defusal; the O->S ledger row and Hurwitz subsection reconciled with
%  the fourth-qubit forward remark ("the terminus closes the ladder of
%  forces; what survives past it is, at most, flavour").
%
%  CROSS-REFS: sec:self-reference [sibling dangle -- resolves when that
%  section is drafted], sec:hopf-justification, sec:hopf-information-bound,
%  sec:register, sec:entanglement-classes, sec:generations, sec:masses,
%  sec:wigners-friend. All others exist.
%  CITATIONS: existing keys baez2002octonions, hurwitz1898composition,
%  furey2018three. ONE new key: adler1995quaternionic (RIS in
%  division_algebras_refs_2026-07-03.ris).
%  British English; \autocite; no bullets in prose; the ledger is a
%  table because the comparison is the argument.
% =====================================================================

\section{The Division-Algebra Ladder}
\label{sec:division-algebras}

The foundations behind us end at amplitudes twice over: natively, where the self-referential
argument forced the description of a system into the system's own state space and priced
consistency at a number system whose parts can cancel --- complex amplitudes, and the Born rule
that squares them (\S\ref{sec:self-reference}) --- and independently, where the regime's
continuity delivers the same field through Hardy's theorem (\S\ref{sec:continuity}), the
convergence of the two branches being evidence rather than accident. This section asks the question
both routes leave open. Why
should the number system stop at $\mathbb{C}$? Ordinary quantum mechanics stops there by postulate;
a framework that has forbidden itself postulates must answer, and the answer is that the climb
neither stops at $\mathbb{C}$ nor continues forever. There are exactly two more rungs, each
purchased by surrendering an algebraic law and each paying for itself in physics, and the ladder is
closed at the top by a theorem of 1898. This section sets out the algebra of that ladder: the single
construction that climbs it, the ledger of what each rung costs and buys, and the theorem that ends
it. The geometry of the same ladder --- that each rung acts on a sphere, that exactly those spheres
are parallelisable, and that the algebraic and geometric classifications are one --- is the business
of the geometry section (\S\ref{sec:hopf-justification}).

\subsection{One construction climbs the whole ladder}

The rungs are not a miscellany; a single construction generates them all. Given an algebra with a
conjugation, the Cayley--Dickson construction builds its double on ordered pairs, with a twisted
multiplication in which the conjugate enters and a conjugation $(a,b)^{*} = (a^{*}, -b)$
\autocite{baez2002octonions}. (Conventions differ in the placement of the conjugations; all yield
isomorphic algebras.) Applied repeatedly it produces the sequence
$\mathbb{R} \to \mathbb{C} \to \mathbb{H} \to \mathbb{O} \to \mathbb{S}$, of real dimensions
$1, 2, 4, 8, 16$: the reals, the complex numbers, the quaternions, the octonions, the sedenions.
The ground floor is classical: $\mathbb{R}$ is ordered and positive-definite, the arithmetic of
classical probability. The first doubling is therefore the quantum move itself --- the passage from
probabilities to amplitudes that can interfere --- and it is not made here by choice: it is the
derivation of \S\ref{sec:self-reference}, already behind us. What this section must account for is
the rest of the climb.

\subsection{The ledger}

In mathematics the doublings are usually told as a story of degeneration: each algebra is worse
behaved than the last, and the sequence is studied for where it breaks. This framework reads the
same ledger in the opposite direction. Each rung is bought by surrendering a law, and each
surrendered law is exactly a piece of physics --- the losses are the assets, and the comparison is
sharp enough to be the argument:

\begin{center}
\begin{tabular}{llll}
\hline
doubling & dimension & what is lost & what the loss becomes \\
\hline
$\mathbb{R} \to \mathbb{C}$ & $1 \to 2$ & the order of the reals & phase and interference; amplitudes \\
 & & & that cancel; the Born rule (\S\ref{sec:self-reference}) \\
$\mathbb{C} \to \mathbb{H}$ & $2 \to 4$ & commutativity & non-commuting actions: weak isospin \\
 & & & $SU(2)$, the $B$ register (\S\ref{sec:register}) \\
$\mathbb{H} \to \mathbb{O}$ & $4 \to 8$ & associativity & the associator: colour, chirality, the \\
 & & & generations, and mass as debt \\
 & & & (\S\ref{sec:generations}, \S\ref{sec:masses}) \\
$\mathbb{O} \to \mathbb{S}$ & $8 \to 16$ & division itself & no further force: zero divisors appear \\
 & & (zero divisors) & and the ladder of sectors ends (below) \\
\hline
\end{tabular}
\end{center}

Walking the rows: the order lost at the first doubling is the possibility of interference --- parts
of a probability that can oppose one another --- which is where the derivation of \S\ref{sec:self-reference}
lives. The commutativity lost at the second is the existence of actions whose order matters; the
unit quaternions form the group $SU(2)$, and the framework houses weak isospin there, as the
internal algebra of the $B$ register whose order-sensitive, non-commuting dynamics the interaction sections develop
(\S\ref{sec:register}, \S\ref{sec:entanglement-classes}). The associativity lost at the third ---
the octonions retain the weaker property of alternativity \autocite{baez2002octonions} --- creates
the associator $[a,b,c] = (ab)c - a(bc)$, and the associator is where most of this paper's physics
is transacted: the stabiliser of a preferred complex direction inside the octonions' automorphisms
is the colour group (\S\ref{sec:hopf-justification}); the triality structure and its $K_3$
obstruction generate the three generations and the sterile channel (\S\ref{sec:generations}); and
the associator is the currency of mass --- the debt against the selected direction that
\S\ref{sec:masses} prices, owed by whatever action moves the selection and not by residence at
the rung, which is the defusal \S\ref{sec:register} makes explicit. The algebraic facts
in the third column are theorems; the identifications in the fourth are established in the sections
cited, each carrying its own epistemic flag there, and are not re-litigated here. What this table
asserts on its own account is only the correspondence's shape: one law surrendered per rung, one
sector of physics gained per law.

\subsection{Internal algebras, not a new quantum mechanics}

A reader versed in the quaternionic-quantum-mechanics literature will at this point raise a fair
objection: quantum theory over $\mathbb{H}$ as a foundation is known to be troubled --- the
tensor-product and composite-system difficulties set out by Adler and others
\autocite{adler1995quaternionic} --- and quantum theory over $\mathbb{O}$ founders earlier still,
since a non-associative algebra does not furnish a Hilbert space at all. The framework's position
must therefore be stated exactly, because it does not make the claim the objection targets. The base
of the physical Hilbert space here is $\mathbb{C}$ and remains $\mathbb{C}$; the higher algebras
enter as the \emph{internal} algebras of the register --- $\mathbb{H}$ for the $B$ qubit,
$\mathbb{O}$ for the $C$ qubit (\S\ref{sec:register}) --- in a tensor structure whose outer theory is
ordinary complex quantum mechanics. This is the same move by which the division algebras enter the
constructions of Furey and the related programmes without any proposal to rebase quantum mechanics
itself \autocite{furey2018three}. For the octonionic register the point is sharper still:
$\mathbb{O}^{2}$ is not a Hilbert space, and the framework accordingly treats the $C$ qubit not in
the Dirac--von Neumann manner but geometrically, as the sphere $S^{15}$ with its Hopf structure ---
which is precisely the form in which \S\ref{sec:hopf-justification} needs it. Nothing in this paper therefore
stands or falls with quaternionic or octonionic quantum mechanics as foundations; the ladder is
climbed by the registers' internal algebras, not by the theory's base field.

\subsection{The same ladder, twice}

The ladder appears in this framework in two guises, and it will help the reader to have them
distinguished now. Statically, it is the anatomy of a single relation: the three qubits of the
register take their values one rung apart --- $A$ over $\mathbb{C}$, $B$ over $\mathbb{H}$, $C$ over
$\mathbb{O}$ --- so that every particle carries the whole ladder as its internal structure
(\S\ref{sec:register}). Dynamically, it is the shape of measurement: a single qubit's pure states
fill $S^{3}$, an entangling measurement carries the joint state into $S^{7}$, and a third party into
$S^{15}$ --- the correspondence between qubit entanglement and the Hopf fibrations that
\S\ref{sec:hopf-justification} takes from a mature literature --- so that each act of measurement ascends one rung --- the climb the reader has already made,
narrated in \S\ref{sec:wigners-friend} --- and every completed relation is a terminated climb, which is
the reading \S\ref{sec:hopf-information-bound} gives the tower's endpoint. One ladder, seen twice:
as the anatomy of a relation and as the dynamics of relating. The framework regards this double
appearance not as a coincidence but as the same structure viewed statically and in motion, and the
sections cited carry that identification's weight.

\subsection{Hurwitz closes the ladder}

The construction itself would double forever; what stops it is a classification theorem. By
Hurwitz's theorem, the normed division algebras --- those in which the norm composes,
$\lvert ab\rvert = \lvert a\rvert\,\lvert b\rvert$, and no non-zero elements multiply to zero ---
are exactly $\mathbb{R}$, $\mathbb{C}$, $\mathbb{H}$ and $\mathbb{O}$
\autocite{hurwitz1898composition,baez2002octonions}. The next doubling, the sedenions, has zero
divisors: non-zero elements whose product vanishes. The framework's reading of this failure is the
one that matters physically. Zero divisors cannot support consistent information flow around loops
--- a cycle of compositions can silently annihilate what it carries --- and a framework in which
closed computational paths are load-bearing (\S\ref{sec:hopf-justification},
\S\ref{sec:space-of-computations}) cannot be built on an algebra that permits it. At the level of
amplitudes the same failure reads: the composition law
$\lvert ab\rvert = \lvert a\rvert\,\lvert b\rvert$, which carries probability multiplicatively
through every rung of the ladder, fails beyond the octonions, so that two non-vanishing amplitudes
could compose to nothing and probability would be destroyed by mere composition. The last algebra
fit to carry amplitudes is $\mathbb{O}$. One rung past the terminus is not entirely silent,
however: the sedenions' internal structure --- their automorphism group and their octonion
subalgebras --- is the recorded candidate arena for tagging the generations
(\S\ref{sec:interactions}), and it enters exactly as this section's scoping permits, internally
and never as a base. The terminus closes the ladder of forces; what survives past it is, at most,
flavour. This is the algebraic half of the termination; the
geometric half --- that no sphere above $S^{7}$ is parallelisable --- and the theorem joining the
two halves are the opening business of \S\ref{sec:hopf-justification}, and the epistemic reading of the same endpoint,
as a bound on the information a single relation can carry, is given in
\S\ref{sec:hopf-information-bound}.

\subsection{The honest debt}

As in the geometry section, we close by stating plainly what is proved, what is established, and what is
owed. Proved: the construction, the losses, and the termination are theorems --- the Cayley--Dickson
doubling and its ledger of surrendered laws are standard algebra, and Hurwitz's classification is
among the durable results of the subject. Established: the physical cashing of each rung --- the
Born rule at $\mathbb{C}$, weak isospin at $\mathbb{H}$, colour, generations and mass at
$\mathbb{O}$ --- is carried out in the sections cited above, each under its own epistemic flag. Owed,
and stated without disguise: the ladder as presented here is closed from above and derived at its
floor, but it is not yet derived rung by rung from below. Hurwitz forbids a further rung; the
argument of \S\ref{sec:self-reference} forces the first doubling, from the ordered arithmetic of
classical probability to amplitudes that interfere. That the same self-referential closure, applied
to a system whose amplitudes already take values at one rung, compels the next --- that relational
closure at one level forces precisely the Cayley--Dickson double of that level --- is the
framework's reading of the climb, and it is a thesis, not a theorem --- but the formal apparatus
against which it can be posed now exists, and we pose it.

\begin{openproblem}[The compulsion of the intermediate rungs]\label{op:ladder-compulsion}
Derive, within the categorical apparatus of Open Problem~\ref{op:sigma-functor}, that
self-referential closure at one rung compels the Cayley--Dickson double of that rung, so that the
intermediate steps $\mathbb{C} \to \mathbb{H} \to \mathbb{O}$ follow from the same principle
that forces the floor. The candidate route is the construction's own shape: the Cayley--Dickson
pair with its conjugation, $(a, b)^{*} = (a^{*}, -b)$, is an involution-equipped doubling, and the
route reads it as the algebraic shadow of adjoining the describer --- closure of level-$A$
descriptions under the inside/outside exchange requiring exactly the twisted pair. The problem
inherits the status of Open Problem~\ref{op:sigma-functor}, whose apparatus it presupposes.
\end{openproblem}

If the derivation lands, the whole ladder becomes a consequence of the same self-reference that
opened it. Until then, the honest statement is that the ladder's necessity is established at its
ceiling and at its floor, while the compulsion of the intermediate rungs is inherited from the
physics they are shown to carry rather than derived in advance. What this section hands on is, accordingly, an algebraic ladder with a floor, a ceiling and a
ledger. Its first consumer is immediate: \S\ref{sec:space-of-computations} spends the terminal
fact --- the non-associativity of $S^{7}$'s own multiplication --- as geometry, excluding the
seven-sphere as a ruler and fixing the dimension of space at three.
\S\ref{sec:hopf-justification} then supplies the geometry of the fibrations themselves, and the
join.
   <-- moved here
%      % EDITED 2026-10-07 scripts/patch_infinite_not_flat_2026-10-07T1400.py
% EDITED 2026-09-27 scripts/patch_spacelike_inversion_2026-09-27T1800.py
% EDITED 2026-09-21 scripts/patch_op_sterile_gravity_boundary_2026-09-21T1300.py
% EDITED 2026-09-12 scripts/patch_spacetime_bloch_headline_2026-09-12T1400.py
% EDITED 2026-09-11 scripts/patch_splice_loops_prospects_2026-09-11T0900.py
% EDITED 2026-09-10 scripts/patch_cite_rubino_2026-09-10T0700.py
% EDITED 2026-09-10 scripts/patch_quoting_and_typography_2026-09-10T0620.py
% =====================================================================
%  SECTION: Spacetime from non-computability (sec:space-of-computations).
%  AUDITED 2026-07-06 (final pass, seventh stop).
%
%  PLACEMENT (reordered chain, JS 2026-07-06): position 7 --
%      \input{division_algebras_section_2026-07-03}   <-- hands off here
%      \input{spacetime_from_non-computability}
%      \input{hopf_justification_section_2026-06-08}
%  Division-algebras' closing now carries the S^7 hand-off sentence
%  (written 2026-07-06 after this file was read).
%
%  INPUT DEPENDENCIES (both held in the audited outputs):
%      entropic_gravity_family_paragraph_2026-07-03.tex  (sec:jacobson)
%      gr_nonlocality_revision_2026-06-13.tex            (sec:gr-nonlocality)
%  The 2026-07-04 additions (asymmetry paragraph; Penrose second route;
%  acceleration/Unruh arc) are FULLY PASTED; the separately uploaded
%  additions file is redundant and can be archived.
%
%  AUDIT RECORD: opener formalised (Russell typo; bare ref retargeted
%  to the freshly minted sec:complex-amplitudes); the "rotations of
%  S^7" wording made precise at BOTH occurrences (Moufang loop, not a
%  group; SO(8) disclaimer added -- referee hazard removed); two \\
%  paragraph-break bugs fixed; dark-matter paragraph retargeted to
%  sec:generations and stitched to op:sterile-gravity (the zoo's
%  demotion and this paragraph's priority condition are one question);
%  CP paragraph stitched to the fourth-party reading of
%  sec:interactions; stale mid-file mini-header retired
%  (sec:k3-sterile placeholder dies with it); the section's narrated
%  debts NUMBERED at last -- op:comp-metric and op:three-dimensions
%  minted in the Open-problems subsection, cited from the in-text
%  "named open problems" line and from the closing summary.
%
%  KEYS: adams1960hopf matches corpus. Unique to this section, presumed
%  in the .bib since the paper compiles: Spence-Brown-1969 (suggest
%  normalising the key spelling at leisure), VanRaamsdonk2010,
%  MaldacenaSusskind2013, MuellerMasanes2013, MasanesMuellerEtAl2014,
%  toffoli_2003, Chaitin1998, Penrose1996, HohnMueller2016,
%  GarnerMuellerDahlsten2017, Jacobson1995, Bekenstein1973, tHooft1993,
%  Susskind1995, spekkens2007toy, Andrianov2001, Sakharov1967,
%  Planck2020, Riess1998. Shared keys verified: small_2006, Penrose1994,
%  unruh1976notes (swap-flag stands), baez2002octonions, Wheeler1990,
%  furey2018three, DiBiagioRovelli2021, Rovelli1996a, Spekkens toy.
%  British English; \autocite; no bullets in prose.
% =====================================================================

\section{Spacetime from non-computability}
\label{sec:space-of-computations}

This section is a deliberate diversion: it takes the idea that non-computable steps are measured
by imaginary information (\S\ref{sec:complex-amplitudes}) and asks what geometry it buys.
Spencer-Brown records that Bertrand Russell thought well of the underlying move
\autocite{Spence-Brown-1969}; the reader may judge whether the framework spends it responsibly.

The programmes that derive spacetime from quantum mechanics---Van~Raamsdonk's
\enquote{building spacetime with entanglement} \autocite{VanRaamsdonk2010}, the
ER\,$=$\,EPR conjecture \autocite{MaldacenaSusskind2013}, and M\"uller's
information-theoretic recovery of three-dimensional space from the qubit
\autocite{MuellerMasanes2013}---take entanglement or information as primitive and
grow geometry from it. The route taken here is the same in spirit but inverted in
order: rather than positing spatial dimensions and forbidding superluminal
signalling, we \emph{define} spatial separation as a measure of non-computability
and read the dimension, the signature, and the Lorentz group off the result. The
proposal originates in the author's CASYS'05 paper \autocite{small_2006}, which
first suggested treating spacetime as a space of computations and attempted to
justify three dimensions on this basis; the present account tightens and extends
it. This is the framework's most structural and least formally complete arc; the
three-dimensionality result is a structural conclusion awaiting a formal proof,
and several of the links below are conjectural readings. Each is flagged as such,
so the reader sees exactly where the chain is load-bearing and where it is an
invitation.

\subsection{Timelike and spacelike distance as computation}
\label{sec:comp-distance}

A computation carries a state from axioms to a theorem in a number of steps.
Counting the computable steps between two states gives a real measure---time. A
physical clock counts completed irreversible loops, so proper time literally
counts loop completions and the metric $\mathrm{d}\tau^{2}=g_{\mu\nu}\,
\mathrm{d}x^{\mu}\mathrm{d}x^{\nu}$ is a loop-density tensor, the conversion from
coordinate intervals to loop counts; Toffoli's reading of the Lagrangian as a
count of action quanta is the same statement \autocite{toffoli_2003}.

Spacelike separation is the complementary, non-computable measure. Define the
non-computable distance between two propositions, relative to an axiom system, as
the number of bits that must be added to the axioms to make the second derivable.
A closed timelike curve is a computation whose output is its own input; by the
fixed-point theorem it has a unique consistent solution or none, and the
existence of that solution is one bit not derivable from any axiom system that
does not already contain the loop---one bit per curve, in the manner of Chaitin's
$\Omega$ \autocite{Chaitin1998}. The spacelike distance between two events is then
the number of such curves that would have to be closed to make them causally
connected. Because each added bit is a non-computable proposition, it is carried
by the imaginary part of the complex information measure
$S=-\ln\lvert p\rvert-\mathrm{i}\theta$ derived in
Section~\ref{sec:complex-amplitudes}: computable distance is counted in real
units, non-computable distance in imaginary ones. The Minkowski signature, with
spacelike intervals imaginary relative to timelike ones, is the immediate reading. The convention is not new to this paper: it was proposed in the author's CASYS paper two
decades ago, where counting non-causal steps in imaginary numbers was introduced and its formal
proof flagged, already then, as a key open problem \autocite{small_2006} ---
Open Problem~\ref{op:comp-metric} is that flag's direct descendant, stated now with the machinery
to be answerable.
That the resulting structure is a genuine metric space, and that the signature
follows rather than being fitted, are open problems---at present the
identification is motivated by reproducing what nature does, not proved.

%% VERIFY-SUPPORT 2026-09-27: conway_2003 is cited below for "any two octonions
%% generate an associative subalgebra" (Artin's theorem; a Moufang loop is
%% diassociative) and harvey1982calibrated for "a three-plane in Im O is
%% associative iff the associator of a basis vanishes, and the associative
%% three-planes are the imaginary parts of the quaternionic subalgebras"
%% (Harvey--Lawson, ch. IV). Neither passage was read in the session that wrote
%% this; both statements are checked numerically in Convention A by
%% scripts/antitelephone_network_associativity_2026-09-27T1900.py (PASS).
A reader used to the standard order of ideas should pause here, because the order has been
inverted, and the inversion is used below. In special relativity the tachyonic antitelephone is a
corollary: a signal faster than light, composed with a boost, would arrive before it was sent and
close a timelike loop, therefore there is no such signal \autocite{tolman1917theory,
benfordbooknewcomb1970tachyonic}. The metric comes first and the impossible telephone is a check
on it. Here the telephone comes first. Two observers are spacelike separated when an antitelephone
between them would close a loop; the loops such telephones would close, and carry nothing round,
are what spacelike separation is, their number is the distance, and the metric is read off from
them. Nothing is lost, since the corollary returns as a tautology: an antitelephone that carried
a record would close a causal loop, and its two ends would not have been spacelike after all.

What is gained is a definite meaning for a spatial dimension. Picture many observers, each holding
an antitelephone to each of the others. A message relayed from one observer to another through a
third composes two transports round two loops, and the network is consistent only if every relay
composes to a single transport whatever the bracketing of the relaying: a spatial dimension is a
family of loops whose transports compose associatively, which is a group law on a parallelisable
sphere. That is a condition a family of loops can fail, and \S\ref{sec:three-dimensions} shows
that the loops of the seventh sphere fail it. Two antitelephones always compose, because any two
octonions generate an associative subalgebra \autocite{conway_2003}; three compose only when their
imaginary parts lie in one quaternionic subalgebra, which is what the associative three-planes of
Harvey and Lawson are \autocite{harvey1982calibrated}, and a third link outside that subalgebra is
seen by the network exactly through its component outside it. The largest consistent network of
antitelephones inside $S^{7}$ is therefore an $S^{3}$, and there are three spatial directions.
When the seventh sphere is set aside there as a ruler, which is the octonionic spacetime of the
division-algebra literature, the claim is not that $\mathbb{R}^{8}$ has no metric, nor that the
ten-dimensional Minkowski space of $SL(2,\mathbb{O}) \cong \mathrm{Spin}(9,1)$ fails to exist
\autocite{baez2002octonions}; it is that the octonionic directions are not loops on which
antitelephones compose, and so are algebra rather than distance. In the standard order that
sentence has no content; in this one it is the reason space has three dimensions.

The split just defined explains the most familiar asymmetry in physics: that one can move back and
forth in space but not in time. \S\ref{sec:continuity} identifies the pre-measurement regime with
continuous reversibility --- its transformation group is connected, and every motion in it is
undoable. Spacelike separation, being the non-computable sector, is exactly the part of the
relational structure on which no record has been written; motion across it is motion within the
reversible regime, and it inherits that regime's continuity and its two-sidedness: a place can be
revisited because going and returning writes no record that forbids the return. Timelike separation
is the record chain itself --- proper time counts completed irreversible loops, as above --- and a
moment cannot be revisited because a moment \emph{is} a stratum of records, and returning would
require unwriting them, which is what \textsc{fanout} forbids (\S\ref{sec:fanout-boundary}). The
asymmetry between space and time is therefore not an extra postulate about the manifold: it is the
pre-/post-\textsc{fanout} asymmetry read as geometry --- space is what the reversible group feels
like from inside, time is what the record chain feels like --- and the arrow-of-time discussion
later in this section is the same asymmetry read thermodynamically. The identification is the
framework's structural reading of the split, offered in the same register as the rest of this
subsection.

\subsection{Three dimensions of space}
\label{sec:three-dimensions}

The closed timelike curves that define spacelike distance require information to
be parallel transported consistently around a loop, and consistent parallel
transport on a sphere requires it to be parallelisable. By Adams' theorem the
only parallelisable spheres are $S^{1}$, $S^{3}$ and $S^{7}$
\autocite{adams1960hopf}, the unit spheres of $\mathbb{C}$, $\mathbb{H}$ and
$\mathbb{O}$. Impose now the requirement that information arriving from outside
the causal zone must appear to possess its own causal history---that is, to be the
result of a sequence of information-preserving operations, hence a sequence of
group operations. The spheres $S^{1}$ and $S^{3}$ are groups under their own multiplication --- $U(1)$ and
$SU(2)$ --- and their operations preserve information; $S^{7}$'s multiplication is non-associative,
a Moufang loop rather than a group, so $S^{7}$ cannot serve as a ruler. (The rotation group
$SO(8)$ of the seven-sphere is of course a group; what fails is the sphere's \emph{own}
composition law, the one a ruler would use.) The same non-associativity
forbids a metric built from the sphere's own composition law outright: path
composition $d(A,C)$ through $B$ depends on the bracketing, so the triangle inequality is not merely violated but ill-defined, and
no geodesic structure exists. (What is untouched should be said: the Euclidean metric of
$\mathbb{R}^{8}$, and the ten-dimensional spacetime that $SL(2,\mathbb{O})$ furnishes, are metrics
of the ordinary kind; what fails is a ruler made of antitelephone loops, which is what a spatial
dimension is in this section, as \S\ref{sec:comp-distance} set out. Nor is the failure total:
relays between any two links compose, and among three they compose exactly when the links lie in
one quaternionic subalgebra, so the largest consistent network inside $S^{7}$ is an $S^{3}$
\autocite{conway_2003,harvey1982calibrated}.) Since the abelian $S^{1}$ alone cannot carry
non-abelian structure, spacelike distance must be three-dimensional. Dimension
three is an output, not an input \autocite{small_2006}.

The division of labour is then exact. $S^{1}$ ($\mathbb{C}$) contributes the time
dimension and the $U(1)$ of electromagnetism; $S^{3}$ ($\mathbb{H}$) contributes
the three spatial dimensions and the $SU(2)$ of the weak interaction; $S^{7}$
($\mathbb{O}$) contributes no geometry but, as a filter on what information may be
exchanged, the $SU(3)$ of the strong interaction. A sphere that cannot be a ruler
can still be a filter, which is why the strong force resists geometric
interpretation while electromagnetism and the weak force do not. This is the same
$\mathbb{C}\oplus\mathbb{C}^{3}$ split that yields $G_{\mathrm{SM}}$ in
Section~\ref{sec:entanglement-classes}, now read as the reason spacetime is $(3+1)$-%
dimensional. The conclusion converges with M\"uller and Masanes, who prove from
information-theoretic postulates, assuming no quantum theory, that minimal
direction-carrying systems with continuous reversible dynamics force $d=3$ and
quantum theory on two qubits \autocite{MuellerMasanes2013}, and with the result
that entanglement exists only for a three-dimensional Bloch ball
\autocite{MasanesMuellerEtAl2014}; the framework supplies the reason such systems
exist, Cayley--Dickson doubling from self-reference, while their work supplies
the rigorous dimensional uniqueness. The status of the framework's own route is a
structural result requiring formal proof: the one-bit-per-curve identification and
the precise reduction of the $S^{7}$ filter to $SU(3)$ rather than its full automorphism group
$G_{2}$ are numbered below (Open Problems~\ref{op:comp-metric} and~\ref{op:three-dimensions}).

A scoping remark, and then a second route to the same count. The computational model of spacetime
proposed here has many detailed parallels with the real thing, but a parallel is not an identity,
and the identity is not claimed in this paper: establishing it is work deferred to the second paper.
What can be claimed now is that the parallel is productive in a specific direction --- it suggests a
derivation of the dimension count carried out entirely within general relativity, with no
computational premise. Acceleration and gravitation tip the light cone, and the tipping can be
carried to the extreme: Penrose's discussion runs the cone through a full circle
\autocite[382]{Penrose1994}, at which point the tipped direction closes on itself --- which is to
say, the extreme of tipping is a closed timelike curve, the very object whose closure the
computational definition of \S\ref{sec:comp-distance} counts in bits. The two routes are one route
entered from opposite ends. Now demand of general relativity only this: that observers accelerating
with respect to one another, and with respect to unaccelerated observers, retain a consistent
conception of the distance between them. Each observer's tipping traces a path on a sphere of
directions, and mutual consistency requires the associated rotations to compose without regard to
bracketing --- to form a group. Compositions on $S^{1}$ and $S^{3}$ do, forming $U(1)$ and $SU(2)$ as above; compositions on
$S^{7}$ do not, the multiplication being non-associative, so observers tipped across an
$S^{7}$-worth of directions could not agree on distance even pairwise: the same failure that forbade
$S^{7}$ a metric above forbids it here a shared ruler. The largest sphere admitting a consistent
conception of distance under mutual acceleration is therefore $S^{3}$, and space has at most three
dimensions; with the floor argued above --- the abelian $S^{1}$ alone cannot carry the required
structure --- the count is exactly three. The status matches the subsection's. The tipping and its
extreme are general relativity's own \autocite{Penrose1994}; the consistency demand is
elementary; the associativity filter is the algebraic fact used throughout; the reading of composed
tippings as motion on the parallelisable spheres is the structural step, shared with the
computational route. The convergence of the two routes on the same endpoint is offered as evidence
for the parallel, while the identity it gestures at remains the second paper's burden.

\subsection{Recovering the Lorentz group}
\label{sec:lorentz}

A measurement changes what the observer knows, hence, relationally, the causal
structure: it tips the light cone. This is Penrose's association of cone-tipping
with state reduction \autocite{Penrose1996} run in the opposite causal direction,
reduction driving the tipping rather than the reverse, and the two are then one
operation in two charts---both described by $SL(2,\mathbb{C})$. The algebraic
backbone is the division-algebra correspondence
$SL(2,\mathbb{A})\cong\mathrm{Spin}(1,d{+}1)$ \autocite{baez2002octonions}: $\mathbb{R}$
gives $1{+}1$, $\mathbb{C}$ gives $3{+}1$, $\mathbb{H}$ gives $5{+}1$ and
$\mathbb{O}$ gives $9{+}1$ dimensions. The computability split that singles out a
preferred $\mathbb{C}$ fixes the fundamental state space as $\mathbb{C}^{2}$ and
hence $SL(2,\mathbb{C})=\mathrm{Spin}(1,3)$: the Lorentz group of $3{+}1$
spacetime. The qubit and Minkowski space are not merely both Euclidean-and-three%
-dimensional by coincidence; the Hermitian $2\times2$ matrices are
$\mathbb{R}^{1,3}$, $SL(2,\mathbb{C})$ acts on them as the Lorentz group, and the
qubit density operators fill the future cone---the correspondence M\"uller's group
develops operationally, recovering Lorentz transformations from quantum
communication \autocite{HohnMueller2016} and the complex and quaternionic qubit
from the relativity of simultaneity \autocite{GarnerMuellerDahlsten2017}.
In these coordinates the Bloch ball is the unit-trace slice of the cone and the Bloch sphere its
boundary rays, the celestial sphere on which $SL(2,\mathbb{C})$ acts by M\"obius
transformations; state reduction sends a point of the ball to a point of the sphere, and a boost
sends the sphere to itself, which is what lets the two be one operation in two charts.
The $SL(2,\mathbb{A})$ correspondence is rigorous; its identification with
measurement-as-tipping is the framework's structural reading.

One consequence of the identification has been implied throughout this section and deserves its own
paragraph, because the next paragraph turns on it. If measurement creates a relation between
observer and observed, and if --- relationally --- a new relation alters the causal structure, then
the tipping of the cone is not merely \emph{associated} with the creation of the relation: it is
what the creation of the relation is, seen geometrically. And a tipping of one's cone against the
congruence of one's neighbours already has a name in physics: acceleration. The framework's
statement, made explicit here, is that \emph{creating a new relation between observer and observed
is experienced as acceleration}. The resonance with the equivalence principle is exactly what it
should be: the principle makes acceleration and gravitation locally indistinguishable, Penrose's
association --- run in this section's direction, reduction driving the tipping --- makes state
reduction and cone-tipping one operation, and the framework's identification closes the triangle,
measurement and acceleration meeting where the principle says they must. The claim is structural:
the framework's reading of the cone-tipping established above, not an independent derivation.

Taken literally, the identification has a measurable face, and it is a familiar one: the Unruh
effect, read in this section's own ledger. Recall the definitions of \S\ref{sec:comp-distance}:
timelike distance counts computable steps, spacelike distance counts non-computable bits. What an
observer can compute --- what their \textsc{fanout}s can convert into classical records --- is fixed
by their causal structure, and the previous paragraph has just made that structure
relation-dependent. An accelerated observer's cones are tipped, and with them the boundary of the
recordable: relations that for the inertial observer lie wholly in the reversible, pre-record sector
--- the entanglement structure of the vacuum, non-computable for them, spacelike in the
computational sense --- fall, for the accelerated observer, partly on the recordable side. The mechanism is the one
Chaitin identified for formal systems \autocite{Chaitin1998}: a proposition undecidable relative to
an axiom system becomes decidable when information is added to the axioms --- and
\S\ref{sec:comp-distance} defined spacelike distance as exactly the bits whose addition closes the
gap. Creating a relation is adding information to the observer's axiom base, so timelike and
spacelike distance are converted into one another by the addition and subtraction of information,
and acceleration is the continuous version of that addition. A
relation newly recordable is a detection. Particles appear from nowhere: an accelerated detector in
the inertial vacuum registers a thermal bath at $T = \hbar a / 2\pi c k_{\mathrm{B}}$
\autocite{unruh1976notes}. Two guards keep the reading honest. Nothing here reclassifies the
invariant interval --- timelike and spacelike separation in the emergent metric are Lorentz
invariants, and no event crosses that classification; what is observer-relative is the computational
partition, which relations a given observer's horizon renders recordable. In the ledger of
\S\ref{sec:comp-distance} that partition \emph{is} the timelike/spacelike distinction, so in the
framework's own units timelike distance is \emph{rotated} into spacelike, and spacelike into
timelike, by a change of observer --- the tipping of the cone is literally that rotation --- and the
detector's clicks are the rotation made audible. And the effect is established physics
while the reading is the framework's, flagged structural. What gives the reading its content is a
parallel: the inertial and the accelerated observer stand to the vacuum exactly as Wigner and his
friend stand to the laboratory (\S\ref{sec:wigners-friend}) --- one holds a coherent relation and no
event, the other holds a record, and the facts are relative on both sides in precisely Rovelli's
sense \autocite{DiBiagioRovelli2021}. Since the previous paragraph identifies acceleration with
relation-creation, the parallel is not an analogy but the same structure in gravitational dress, the
horizon standing where the laboratory wall stood: \emph{the Unruh effect is Wigner's friend for
accelerated observers}. And two subsections of this section that seemed independent here meet: the
temperature Jacobson's construction requires across every local Rindler horizon
(\S\ref{sec:jacobson}) is exactly this temperature --- the $T$ of the Clausius relation is the $T$
of the tipped cone, supplied by the very relation-creation with which this subsection began.

\subsection{The relationalist commitment}
\label{sec:relationalism}

General relativity is relational: by the hole argument, spacetime points carry no
intrinsic identity, and the metric is fixed by the network of coincidences and
interactions, not by a pre-existing stage \autocite{Rovelli1996a}. The
computational picture is relational by construction and explains why it must be:
a background metric would require spatial distances to be pre-computable, and
non-computability forbids exactly that. General covariance is the statement that
distances have no pre-assigned values and depend on what information is actually
present. The metric is therefore dynamical not by stipulation but because it
measures a quantity that cannot be fixed in advance.

\subsection{Jacobson, holography, and imaginary information}
\label{sec:jacobson}

Jacobson derives the Einstein equation as an equation of state, from the Clausius
relation $\delta Q=T\,\mathrm{d}S$ across every local Rindler horizon together
with the Bekenstein proportionality of entropy to horizon area
\autocite{Jacobson1995, Bekenstein1973}. In the computational reading this is the
self-consistency condition on the loop-counting network: the Einstein equation is
the requirement that the bookkeeping of computation be internally consistent
across surfaces. The holographic principle, that the information in a region is
bounded by its bounding area \autocite{tHooft1993, Susskind1995}, then acquires a
natural reading here. A surface area is a squared distance; if spacelike distance
is measured in imaginary units of information, an area is measured in the square of
those units, that is in \emph{negative} information---entropy. The holographic
link between area and entropy thus requires the square root of negative
information, and that square root is precisely the imaginary unit of the complex
amplitude, the observer's self-referential phase. Read this way, Jacobson's
account of information crossing surfaces implicitly assigns an imaginary
information value to spacelike distance, and the $\mathrm{i}$ of the Minkowski
signature is the same $\mathrm{i}$ as the $\mathrm{i}$ of the quantum amplitude.
This identification is the most speculative step in the section: it is
structurally suggestive and it makes the holographic form unsurprising, but it is
a conjectural reading rather than a derivation, and is offered as such.

The conjecture has a two-way motivation that this paper's record should carry, because each
direction was on the table before the other was noticed. Forward: if spacelike distance is counted
in imaginary units of non-computable information, then any spacelike area is the product of two
imaginary lengths and is therefore \emph{negative} information --- which is to say entropy, by
Brillouin's identification \autocite{brillouin1956science} --- and a holographic principle arrives
automatically, the sign supplied by $i^{2}$ rather than by interpretation. Converse, and more
striking: in the holographic principle as ordinary physics already states it, surface area is
proportional to entropy, and entropy is negative information; so by implication ordinary physics
has been measuring spacelike separation in imaginary units of information all along, and the
convention of \S\ref{sec:comp-distance} is read off from established results rather than
invented. The same imaginary unit is already employed once in this section --- it is what makes
the signature of the interval Minkowski --- so the sign of the area law is its second job, not a
second assumption. What the two-way motivates it does not prove: the formal statement is exactly
the metric-and-signature half of Open Problem~\ref{op:comp-metric}, together with the
one-bit-per-bit-pair coefficient owed to the second paper. The stakes of the conjecture have
recently risen: Dorau and Much prove, with modular-theoretic rigour, that the semiclassical
Einstein equations follow from quantum relative entropy on a local horizon \emph{once the
entropy--area proportionality is assumed} \autocite{doraumuch2026relative} --- so what this
framework proposes to derive is precisely the remaining input of a derivation already carried
through at that standard.
\input{entropic_gravity_family_paragraph_2026-07-03}

\subsection{Non-locality already present in general relativity}
\label{sec:gr-nonlocality}
\input{gr_nonlocality_revision_2026-06-13}
\medskip
\noindent  The spacetime arc therefore yields, from one definition of spacelike
distance as non-computability: the signature, three spatial dimensions,
the Lorentz group, the dynamical relational metric, a reading of the
holographic principle, the presence of non-locality within general
relativity, the identification of relation-creation with acceleration
and of the Unruh effect with Wigner's friend, and the space--time
mobility asymmetry as the \textsc{fanout} asymmetry read as geometry. Of these,
the $(3+1)$ count and the Lorentz recovery are structural results that converge with independent
rigorous work and await their formal proof --- the count now numbered as
Open Problem~\ref{op:three-dimensions}; the metric-space status, the signature and the
one-bit-per-curve measure (together Open Problem~\ref{op:comp-metric}), the imaginary-information
reading of holography, and the general-relativistic non-locality, the acceleration and
Unruh readings, and the mobility asymmetry are conjectural, and are named here precisely so that the boundary between what is shown and what
is proposed is visible.

\subsection{Why the metric is not to be quantised}\label{sec:no-quantised-gravity}

A consequence has been implicit in everything above and deserves its plain statement. Causal
structure, on this framework's reading, is the bookkeeping of locked relations --- post-\textsc{fanout}
by construction --- while superposition is the signature of the pre-\textsc{fanout} regime. A
``superposition of causal structures'' is therefore not a difficult state but a category error: it
applies the pre-record operation to a record-side object, exactly as ``a superposition of
measurement outcomes'' does. Gravity lives in the post-\textsc{fanout} regime. The quantum layer
beneath it exists and is not in dispute; but it is \emph{pre-geometric} --- it carries amplitudes
over the formation of relations, not amplitudes over formed geometries --- so there is no
fundamental metric to quantise, and the standard programme of superposing spacetimes pursues an
object the ontology does not contain. One guard rail keeps the claim honest: collective
excitations of the locking pattern may perfectly well exist as effective quanta --- gravitons as
the phonons of the relational network --- and the thesis forbids fundamental metric superposition,
not effective field theory; this is Jacobson's equation-of-state moral \autocite{Jacobson1995},
made structural. The relation to gravitational-collapse programmes is an inversion: where
superposed geometries are there read as unstable \autocite{Penrose1996}, here they are read as
unformed --- what decays in those models is what locks in this one.

Three experimental postures follow, each at the status the reading carries. Laboratory
demonstrations of indefinite causal order \autocite{procopio2015experimental, rubino2017indefinite} are, on this
account, demonstrations of the nascent layer: the indefiniteness lives below the locking threshold
and definitises at it, so the framework adopts the programme and adds a boundary --- indefinite
causal order can never be amplified into records while remaining indefinite. Gravitationally
induced entanglement of the kind the Bose--Marletto--Vedral proposals seek
\autocite{bose2017spin, marlettovedral2017gravitationally} is \emph{predicted}, because the
pre-geometric layer is quantum --- while witnessing that layer, not a superposed metric: the
framework splits the experiment's usual dichotomy. And gravitationally driven anomalous collapse
is denied: locking is redundancy and thermodynamics, so decoherence experiments should find
standard rates and nothing gravity-specific. That no consistent quantisation of gravity exists was
the conclusion of the author's CASYS paper by a different route \autocite{small_2006}; the present
framework supplies the mechanism the conclusion then lacked.

\subsection{Open problems and prospects}
\label{sec:spacetime-open-problems}

The spacetime arc is the framework's frontier, and honesty about its edges
matters more here than anywhere else in the paper. We set out the genuine
problems and the prospects together, deferring the deepest to a second paper
rather than papering over them. Two of the arc's debts are well-posed now, and we number them.

\begin{openproblem}[The computational metric]\label{op:comp-metric}
Prove that the paired distance of \S\ref{sec:comp-distance} --- computable steps counted real,
non-computable bits counted imaginary, one bit per closed curve --- defines a genuine
pseudo-metric on the space of computations, with the Minkowski signature derived rather than read
off; equivalently, establish the one-bit-per-curve measure as a measure, and the signature as its
theorem. Until then the identification stands as motivated by reproducing what nature does, which
is the status flagged at its introduction.
\end{openproblem}

\begin{openproblem}[Dimension three as a theorem]\label{op:three-dimensions}
Turn the two convergent routes of \S\ref{sec:three-dimensions} into a proof: formalise the
parallelisability-plus-associativity filter (the computational route) and the mutual-acceleration
consistency demand (the general-relativistic route) so that $d = 3$ follows as a theorem, and
include the reduction clause --- that the $S^{7}$ filter acts through $SU(3)$, the stabiliser of
the selected complex direction, rather than through its full automorphism group $G_{2}$. The
M\"uller--Masanes uniqueness results \autocite{MuellerMasanes2013, MasanesMuellerEtAl2014} are
the rigorous convergents any proof must meet.
\end{openproblem}

\paragraph{The substrate and the generation of events (deferred).}
A computation runs on something; a network of events presupposes that events
occur. Both are the same objection in two forms---the network's hardware and its
clock---and the original CASYS'05 statement of the programme raised the
implementation regress and set it aside \autocite{small_2006}. The framework's
response is that a self-referential structure is the only kind that need not
presuppose a meta-substrate: it generates its own. The natural algebraic emblem
is $E_{8}$, the unique simple Lie group whose adjoint representation is also its
smallest faithful one, so that what it is and what it acts on coincide; the $240$
unit integral octonions that generate it are exactly the $E_{8}$ root system, so
the self-action is welded to the division-algebra spine of
Section~\ref{sec:entanglement-classes}. This is Wheeler's self-excited circuit
\autocite{Wheeler1990} given an algebraic body, and it converges with Furey's use
of an algebra acting on itself to generate particle content
\autocite{furey2018three}. The decisive caveat, which we state rather than hide, is
that the bootstrap is self-consistent but not self-specifying: it does not
eliminate the non-computability identified throughout, it relocates it to a
determined place. The substrate question is, on the framework's own terms,
provably without a computable answer from inside---it is the irreducible residue,
the bump in the carpet that the whole construction has been moving rather than
removing. The leap from ``$E_{8}$ is self-dual'' to ``$E_{8}$ is the substrate''
is at present an emblem awaiting operational content (an $E_{8}$-equivariant
dynamics, or a state space built as a homogeneous space of $E_{8}$), and that
content, together with how a tick is generated at all, is deferred to the second
paper.

\paragraph{Dark matter (deferred).}
If gravity is the bookkeeping of the interaction network and dark matter interacts only
gravitationally, how does it participate in a network it does not exchange information with?
This is a genuine boundary of the programme, at which it may make no prediction, and Open
Problem~\ref{op:sterile-gravity} states it with the criteria that would close it. The prospect, to be developed in the second paper, is
that this is a feature rather than a defect. In the division of labour of
Section~\ref{sec:three-dimensions}, participation in the metric---counting loops
in the $S^{1}$ and $S^{3}$ directions---is universal, carried by anything with a
worldline and a tick-rate, whereas participation in the gauge filters, the
$S^{3}$ of the weak interaction and the $S^{7}$ of colour, is selective, engaged
only by charged matter. A node that counts loops but engages no filter
\emph{is} matter that gravitates without otherwise interacting, and the
framework's Class-4 a-chiral sterile neutrino (\S\ref{sec:generations}) is the candidate that
realises it --- recorded in the particle sections as Open Problem~\ref{op:sterile-gravity}, whose
tension (a gauge-sterile channel gravitating, when gravity is measurement bookkeeping) is exactly
the question named here. The condition for the candidacy to be coherent rather than circular is
that metric participation be established as prior to, and independent of, gauge participation;
making that priority precise is the work deferred, and it is the same work the open problem
requests.

\paragraph{CP violation, the arrow of time, and baryogenesis.}
CP violation appears to demand a time asymmetry owing nothing to the entropic
arrow, which would seem to require a preferred direction built into spacetime
itself. The framework dissolves the appearance rather than accommodating it. CP
violation needs an irreducible complex phase, the CKM phase $\delta$; by Spekkens'
result that negativity and contextuality are the same notion of nonclassicality
\autocite{spekkens2007toy}, a complex phase is a signature of contextuality; so CP
violation should be contextual, sensitive to the near environment of the decaying system --- the
same phase the interaction sections read as anchored by a fourth party
(\S\ref{sec:interactions}) --- which is exactly what Andrianov, Taron and Tarrach established for
neutral kaons in matter, where the environment breaks C, irreversibility breaks
T, and the effective parameters follow a Lindblad equation
\autocite{Andrianov2001}. The framework's reading is that the resulting
informational asymmetry \emph{constitutes} the arrow rather than presupposing it:
CP violation, entropy increase, and the direction of time are one informational
asymmetry seen three ways, not an independent second arrow, so the worry about a
direction built into spacetime does not arise. On this footing the contextual
dependence sets up a positive feedback loop for baryogenesis---an initial
fluctuation makes the environment slightly matter-dominated, contextual
enhancement raises the effective CP violation, more matter is produced, and the
loop runs---with $\delta$ a cosmological fixed point rather than a fundamental
constant. The honest limit of this is that a feedback loop explains amplification
and lock-in, not the sign selection or the baryon-number violation that
Sakharov's conditions still require \autocite{Sakharov1967}; the stronger claim
the framework is positioned to make, that the matter direction and the time
direction are co-selected by the same self-referential asymmetry, needs a source
of baryon-number violation and is left open.

\paragraph{Spatial infinitude and the status of inflation.}
There is no upper bound on non-computability, so the space of computations is
unbounded: the universe is spatially infinite, which leaves its spatial sections flat or of
negative curvature and excludes only the closed case. Infinitude on its own, however, does not
settle flatness, which is a fine-tuning of the Friedmann dynamics the framework has not yet
derived (Open Problem~\ref{op:comp-metric}); and flatness was in any case never inflation's
principal task. The horizon problem is, and an infinite space does nothing for it. The framework's actual lever against
inflation is the relational thesis of Section~\ref{sec:gr-nonlocality}: if
pre-spatial entanglement correlates regions before \textsc{fanout} spreads them
into spatial separation, then the observed large-scale homogeneity is fixed
relationally, before any question of causal contact, and the horizon problem
dissolves without an inflaton. This is a real argument, but it carries a real
bill: inflation also generates the nearly scale-invariant spectrum of primordial
fluctuations seen in the microwave background \autocite{Planck2020}, and any
account that dispenses with inflation owes an alternative origin for that
spectrum. That, not the horizon problem, is the hard empirical debt, and it is
named here rather than discharged.

\paragraph{Expansion and acceleration.}
If spacelike distance is the information separating worldlines, then as worldlines
accumulate distinct interaction histories they accumulate information-distance,
and the universe expands. The picture has the right exclusions built in: bound
systems interact continuously, share information, and so keep low
information-distance---they do not expand---while only worldlines that decouple
diverge, recovering expansion on cosmological scales alone. The feedback in which
greater separation means less light-cone overlap and so faster divergence,
$\mathrm{d}d/\mathrm{d}t \propto d$, yields exponential divergence, a de~Sitter
phase, and hence accelerating expansion without a separate dark-energy term
\autocite{Riess1998}. The limit to be honest about is that pure exponential
divergence is the late-time behaviour only; reproducing the full expansion history,
with its radiation- and matter-dominated eras, and fixing the Hubble rate
quantitatively, remains to be done.

\paragraph{The $S^{7}$ sector as internal structure.}
That $S^{7}$ cannot carry a metric does not make it disappear. Its closed
timelike curves carry the same non-computable information structure as those of
$S^{1}$ and $S^{3}$, but lacking associativity they support algebra---gauge
charge---in place of geometry---distance. This is why gauge symmetry has so often
looked like a set of compactified extra dimensions: the Kaluza--Klein and string
intuition is tracking something real in the $S^{7}$ structure, but misdiagnoses
it as spatial. In the framework these degrees of freedom have no size because
they have no metric; they are internal in the exact sense that they enter the
information-exchange that constitutes gauge charge while contributing nothing to
the loop-counting that constitutes distance. The resemblance to rolled-up
dimensions is therefore genuine and instructive, and the diagnosis that they are
algebraic rather than geometric is what separates this account from theories that
add real spatial dimensions to carry the same structure.

\paragraph{Loops, supertasks, and a second route to parallelisability.}
The argument of \S\ref{sec:three-dimensions} derives the parallelisable-sphere constraint from
closed timelike curves: a computation whose output is its own input requires information to be
transported consistently around a loop, and consistent transport selects the parallelisable spheres
$S^{1}$, $S^{3}$, $S^{7}$ \autocite{adams1960hopf}. There is a suggestion, which we record as a
prospect rather than a result, that the same constraint may govern the loops of perturbative quantum
field theory. The external legs of a Feynman diagram are on-shell classical boundary data; the
quantum content, and the divergences, enter with the internal loops, where the loop momentum is
integrated over all values and the action accumulated inside the loop is in this sense unbounded ---
a supertask, an infinite computation closed into a finite process. If such a loop is, in the
framework's terms, a closed computational path of the same kind as a closed timelike curve --- a
cycle carrying information not derivable from the data outside it --- then the parallelisability
constraint derived for closed timelike curves would apply to it, and the admissible interactions
would be constrained by the same $S^{1}$, $S^{3}$, $S^{7}$ structure that fixes the gauge group in
\S\ref{sec:entanglement-classes}. This would be a second and independent route to the Standard Model
content: not only \enquote{closed paths in spacetime select the parallelisable spheres} but \enquote{closed paths in interactions select them too}, unifying the spacetime arc and the interaction arc under one
principle. The decisive step, which we have not taken, is the identification of a field-theory loop
with a closed computational path in the precise sense the framework requires --- a loop integral is a
closed cycle in momentum space, a closed timelike curve a closed cycle in spacetime, and whether
these are the same kind of self-consistent closed structure is exactly what would have to be
established. It is offered as a direction, not a claim.

\paragraph{Does parallelisability determine renormalisability?}
A sharper and more self-contained question sits alongside the previous one, and we state it in the
hope that it will be taken up by those equipped for it. Renormalisability is the property that the
divergences of a quantum field theory can be consistently absorbed, leaving finite predictions; some
theories have it and some do not, and the standard criterion for which is technical (power-counting
of the coupling's dimension). The framework suggests a more structural reading. Parallelisable
spheres are precisely the spheres on which a global, consistent transport exists where it generically
does not; the non-parallelisable spheres are where the obstruction lives. If the divergence of a loop
is, in the framework's language, an obstruction to closing a computational path consistently, then
parallelisability would be exactly the structure that distinguishes the removable obstruction (a
renormalisable theory) from the irremovable one. It is at least a suggestive coincidence that the
renormalisable interactions of nature are the gauge theories built on $U(1)$, $SU(2)$ and $SU(3)$ ---
the groups the framework ties to the three parallelisable spheres --- and that gravity, which resists
renormalisation, is the sector \S\ref{sec:three-dimensions} ties to the geometry rather than the
algebra. Whether this is structural or accidental we do not know. The question --- whether
parallelisability, in the Hurwitz--Adams sense, is what determines renormalisability --- is posed in
the native language of the division-algebra approach to the Standard Model, and is one we think that
community is best placed to answer. We are careful to claim nothing more than that the question is
well-motivated: this is emphatically not a method for computing renormalised quantities, whose finite
remainders carry the precision agreement with experiment, but a conjecture about which theories admit
such a computation at all. It is, like the previous prospect, a direction with a possible path to an
answer rather than an answer, and it is left open deliberately.


\medskip
\noindent These mark the boundary of the programme as it stands. The substrate,
the generation of events, and the network status of dark matter are deferred to
the second paper; the contextual account of CP violation and baryogenesis, the
relational case against inflation, the information-divergence account of
expansion, and the algebraic reading of the $S^{7}$ sector are offered as
prospects, each with its outstanding debt named so that the line between argument
and aspiration stays visible.       <-- first consumer
%                                                          (S^7 exclusion)
%      % EDITED 2026-09-09 scripts/patch_spekkens_citekey_2026-09-09T1810.py
% =====================================================================
%  SECTION: Why Hopf Fibrations Belong in Particle Physics
%  (sec:hopf-justification). Drafted 2026-06-08; F1 arc added
%  prop:f1-dictionary and op:hopf-frobenius's discharge clause.
%  AUDITED 2026-07-06 (final pass, eighth stop).
%
%  PLACEMENT (reordered chain, JS 2026-07-06): position 8 --
%      \input{spacetime_from_non-computability}
%      \input{hopf_justification_section_2026-06-08}
%      \input{entanglement_classes_and_gauge_groups_2026-06-08}
%  No adjacency phrasing in the prose: position-portable. All
%  cross-refs live -- sec:fanout-direction DRAFTED 2026-07-06 (formerly the one undrafted
%  section, by design).
%
%  AUDIT RECORD: prop:gsm, prop:termination and prop:f1-dictionary
%  verified against every quoting consumer (statements canonical);
%  honest-debt subsection names both the selection assumption and the
%  mereology debt; the epistemic two-climbs disambiguation and the
%  compositional rung-evidence verified; sec:canonical-orbits live;
%  file pure ASCII. Fixes: op:hopf-frobenius's cross-link with
%  op:sigma-functor made TWO-WAY (the "four ways" line now names the
%  op); both remaining functor conditions ANCHORED --- the quaternionic
%  to the established fibre physics of sec:wf-reconciliation (residue:
%  clutching map vs strong complementarity), the octonionic to the
%  proved canonical-orbit images (umbrella label app:ghz-w-proofs
%  minted at the appendix for the purpose). Statuses unchanged.
%  British English; \autocite; no bullets in prose.
% =====================================================================

\section{Why Hopf Fibrations Belong in Particle Physics}
\label{sec:hopf-justification}

Invoking Hopf fibrations in a discussion of the Standard Model invites an immediate and fair
objection: that the move is decorative, an analogy between two pieces of attractive mathematics
with no load-bearing connection between them. The purpose of this section is to dispose of that
objection before the construction is put to work. The connection is not an analogy. It is
carried by two independent and independently published theorems, and the structures it relies
upon are not selected for convenience but compelled by a pair of classification results. We take
as already established the correspondence between qubit entanglement and the Hopf fibrations
themselves --- the subject of a mature literature
\autocite{mosseri2001hopf,bernevigchen2003geometry} --- and do not reproduce it here. What this
section addresses is the prior question of why that correspondence is entitled to carry physical
content at all.

\subsection{Two theorems, joined end to end}

The bridge from qubit entanglement to the Standard Model is built from two distinct results,
neither of which is original to the present framework and neither of which depends on the other.

The first is geometric. The unit sphere in the Hilbert space of $n$ qubits is the total space of
a Hopf fibration: $S^3 \to S^2$ for one qubit, $S^7 \to S^4$ for two, and $S^{15} \to S^8$ for
three. For two and three qubits the fibration is \emph{entanglement-sensitive} --- separability
is precisely the condition under which the base degenerates to a product of single-qubit Bloch
spheres and the additional dimensions associated with the higher imaginary directions collapse
\autocite{mosseri2001hopf,bernevigchen2003geometry}. This is a statement about the geometry of
state spaces and requires no physical interpretation whatever.

The second is algebraic. Singling out a preferred complex structure on the octonions ---
equivalently, fixing the split $\mathbb{O} \cong \mathbb{C} \oplus \mathbb{C}^3$ --- leaves a
residual symmetry, and that residual symmetry is exactly the gauge group of the Standard Model.
This was established through the exceptional Jordan algebra by Dubois-Violette and Todorov
\autocite{duboisviolette2016exceptional,duboisviolettetodorov2019exceptional,todorovduboisviolette2018deducing}
and given its sharpest spinorial form by Krasnov \autocite{krasnov2021so9}, who showed that the
relevant group is the stabiliser of a complex structure inside $\mathrm{Spin}(9)$. We record the
result in the form in which it enters the present work.

\begin{proposition}[Dubois-Violette--Todorov; Krasnov]\label{prop:gsm} Let $J_R$ be the complex
structure on the space $\mathbb{O}^2$ of $\mathrm{Spin}(9)$ spinors determined by a choice of
unit imaginary octonion, equivalently the split
$\mathbb{O} \cong \mathbb{C} \oplus \mathbb{C}^3$. The subgroup of $\mathrm{Spin}(9)$ commuting
with $J_R$ is \[ G_{\mathrm{SM}} \;=\; \frac{SU(3) \times SU(2) \times U(1)}{\mathbb{Z}_6}, \]
acting on the singled-out $\mathbb{C}$ summand as a colour singlet of hypercharge $-1$ and on
the $\mathbb{C}^3$ summand as a colour triplet of hypercharge $\tfrac{1}{3}$, both as $SU(2)$
doublets --- the quantum numbers of one left-handed generation.
\end{proposition}

The two results are independent, and the join between them is the only place a physical theory
must do work. Each result is a theorem; what neither supplies is a reason why nature should fix
one complex direction within $\mathbb{O}$ rather than treating all directions symmetrically. We
return to this --- the genuine load-bearing assumption --- in \S\ref{sec:hopf-honest-debt}, and
resolve it, in candidate form, in \S\ref{sec:fanout-direction}.

\subsection{The ladder is forced, not chosen}

It might still be objected that selecting the octonions, or the Hopf fibrations, is itself a
modelling choice carrying as much latitude as the choice of a grand-unified group. It is not,
and the reason is a pair of classification theorems that close off the alternatives entirely.

By Hurwitz's theorem the only normed division algebras are $\mathbb{R}$, $\mathbb{C}$,
$\mathbb{H}$ and $\mathbb{O}$ \autocite{hurwitz1898composition,baez2002octonions}. By Adams's
resolution of the Hopf-invariant-one problem --- together with the parallelisability results of
Bott--Milnor and Kervaire --- the only parallelisable spheres are $S^0$, $S^1$, $S^3$ and $S^7$,
so that there exist exactly three nontrivial Hopf fibrations, in bijection with $\mathbb{C}$,
$\mathbb{H}$ and $\mathbb{O}$
\autocite{adams1960hopf,bottmilnor1958parallelizability,kervaire1958nonparallelizability}. These
are not heuristics; they are among the more durable theorems in topology and algebra. Their
force here is that, once a framework is committed to amplitudes valued in a normed division
algebra --- a commitment that in the present construction is made far upstream, by the
self-referential argument of \S\ref{sec:self-reference} and \S\ref{sec:division-algebras}, and
not by fiat at this stage --- the ladder of Hopf fibrations is not selected but compelled, and
it terminates of its own accord.

\begin{proposition}[Termination]\label{prop:termination} The Cayley--Dickson tower of normed
division algebras and the tower of parallelisable spheres both terminate at real dimension
eight. Consequently there are exactly three nontrivial Hopf fibrations, the largest being the
octonionic $S^{15} \to S^8$, and the construction admits no extension to a fourth case.
\end{proposition}

The termination is worth dwelling on, because it is unusual for a model of this kind to predict
its own boundary. The next rung of the Cayley--Dickson sequence, the sedenions, possesses zero
divisors; the Hopf map can no longer be globally defined; there is no fourth fibration. This is
not a limitation imposed from outside but a structural fact, and it carries a physical
consequence that we develop later: there is no room in the construction for a fourth gauge force
or a fifth interaction, so that the discovery of either would count as evidence against the
framework rather than as an invitation to extend it. The endpoint is a theorem, not a
stipulation.

\subsection{Independent arrival as corroboration}

The most persuasive evidence that the destination is not an artefact of the route is that it is
reached by routes sharing none of the same machinery. Szangolies
\autocite{szangolies2025standardmodel} reaches $G_{\mathrm{SM}}$ from the emergent-gravity and
Kaluza--Klein programmes, reading the singled-out complex direction as a dimensional reduction
of a $9+1$-dimensional spacetime; the present framework reaches it from self-reference
(\S\ref{sec:self-reference}); and Furey and Gillard--Gresnigt reach the same gauge group, and
the same three-generation structure, through the minimal left ideals of the Dixon and Clifford
algebras, without invoking Hopf fibrations at all
\autocite{furey2018three,gillardgresnigt2019three}; the relation between their minimal ideals
and the entanglement classes used here is taken up in \S\ref{sec:canonical-orbits}. The
octonionic route to particle physics is in any case an old one, going back to the use of
$\mathbb{O}$ to model the strong force \autocite{gunaydingursey1974quark}.

That programmes with disjoint starting points converge on
$SU(3) \times SU(2) \times U(1)/\mathbb{Z}_6$ and on the same quark--lepton division is the
coincidence Szangolies himself remarks upon: the characteristic structure of the Standard Model,
routinely regarded as arbitrary, turns up independently in unrelated settings, and its appearing
twice would then be doubly odd unless the two settings are in fact related. We read the
convergence the same way --- as evidence that the division-algebraic description is singled out
by the mathematics rather than imposed by the modeller, and therefore that
Proposition~\ref{prop:gsm} is structural rather than contrived.

\subsection{The honest debt}
\label{sec:hopf-honest-debt}

It would be misleading to leave the impression that this route is free of assumptions, so we
state plainly the one assumption that all of these programmes share. Proposition~\ref{prop:gsm}
requires that \emph{some} complex direction in $\mathbb{O}$ be singled out, and the uniqueness
of that choice is not supplied by the theorem. Krasnov parametrises the complex structure $J_R$
by an arbitrary unit imaginary octonion, and the residual group is the same for every such
choice because $G_2 = \mathrm{Aut}(\mathbb{O})$ acts transitively on the unit imaginary
octonions. The theorem therefore guarantees the gauge group \emph{given} that a direction is
fixed, while remaining silent on \emph{why} one is fixed. Szangolies treats this as a free
choice and flags it, in his own conclusions, as an open element of the construction --- one to
which several inequivalent answers are available.

This, and not the appeal to Hopf fibrations, is the genuine load-bearing assumption of the whole
approach. The Hopf fibrations are forced (Proposition~\ref{prop:termination}); the singling-out
is not. The present framework offers a candidate resolution that the prior literature lacks: the
act of measurement, modelled as \textsc{fanout}, selects the direction physically, and the
transitivity of $G_2$ then guarantees that the gauge group is independent of which direction the
measurement happens to select. This argument is given in \S\ref{sec:fanout-direction}; we flag
here, at the appropriate epistemic status, that it is structural rather than fully proved, and
that the residual question --- why a measurement selects \emph{exactly one} imaginary direction
rather than merely \emph{permitting} a selection --- remains open. It is now priced as Open
Problem~\ref{op:fanout-selection}.

A second debt is shared with every programme of this character and is not closed by anything
said here: why the abstract quantum state should admit a tensor-product decomposition into
qubits in the first place. This is the quantum factorisation, or mereology, problem. The
entanglement-sensitivity on which the construction turns is meaningful only once such a
decomposition is fixed, and we treat the decomposition as a necessary input rather than a
derived result. Naming it is preferable to letting it pass unremarked.

\subsection{The termination read epistemically}
\label{sec:hopf-information-bound}

The termination of the tower (Proposition~\ref{prop:termination}) has so far been stated
algebraically: the Cayley--Dickson sequence acquires zero divisors at the sedenions and the Hopf
map can no longer be globally defined. There is a second reading of the same boundary, epistemic
rather than algebraic, which we offer as a motivated conjecture rather than a settled result,
because it connects the tower's endpoint to a bound that arises independently in the foundations
of measurement.

A particle, in this framework, is a reified relation between events, and its internal content is
carried by the octonionic $A,B,C$ register --- the $S^{7}$ sector serving as internal degrees of
freedom in the channel that connects the events (\S\ref{sec:space-of-computations}). The
termination of the tower is then a statement about that content: a single such relation can
support genuine entanglement among at most three parties, because there is no fourth fibration
to carry a fourth independent party within the same irreducible relation. Every particle is, on
this reading, a \emph{completed} octonionic-level relation --- topped out at three parties not
by circumstance but because the tower admits no further rung. The information that can be
encoded in one relation-between-events, and hence extracted by one measurement of it, is
correspondingly bounded: it is the information of a completed three-party register, and no
measurement of a single particle can yield more.

This is exactly the kind of bound that the epistemic-restriction programmes reach from the other
direction. Spekkens' toy theory is built on a knowledge-balance principle that caps the
information obtainable about a system \autocite{spekkens2007toy}, and Szangolies derives such an
epistemic horizon not by positing it but from the same self-referential non-computability on
which the present tower rests \autocite{szangolies_j_2018}. That the two accounts --- the
algebraic termination of the Hopf tower and the epistemic bound on measurement information ---
should coincide is, we suggest, not a numerical accident but a shared root: both are
consequences of the self-referential non-computability of \S\ref{sec:self-reference}, wearing an
algebraic face in the one case and an epistemic face in the other. The Hurwitz--Adams ceiling
and the knowledge-balance ceiling would then be one limit seen twice, and Szangolies' programme,
which reaches $G_{\mathrm{SM}}$ by an independent route (\S\ref{sec:hopf-honest-debt}) and the
epistemic horizon from the same non-computable source, shares with this framework not merely a
conclusion but its origin.

The reading must be kept distinct from a different climb, with which it would otherwise be
conflated. The Wigner's-friend regress of \S\ref{sec:wigners-friend} ascends without end: each
observer may be measured by a further observer, each carrying a fresh gauge-private circle, and
that chain has no top --- it is the bump in the carpet, unbounded in length. The termination
bounds something else: not the length of the chain of observers, but the internal structure of
each single relation, each particle, that the chain is built from. The regress is unbounded
precisely because each of its links is a bounded, completed three-party object;
termination-of-each-link and non-termination-of-the-chain are not in tension but complementary,
the former individuating the links of the latter. What the bound caps is the information in one
particle; what the regress leaves uncapped is the number of particles and observers that can be
chained. The precise identification of ``the information of a completed three-party register''
with the formal measures of the epistemic-restriction literature is the technical step that
would turn this conjecture into a theorem, and we flag it as owed; the structural claim --- that
the tower's endpoint and the measurement bound share the non-computable root --- is what we
advance here.

\subsection{The termination read compositionally}
\label{sec:hopf-compositional}

There is a third termination theorem, proved in a community that has so far had no contact with
the two joined above. Coecke and Kissinger showed that the GHZ and W classes --- realised as the
commutative special and the anti-special Frobenius algebras of categorical quantum mechanics ---
are \emph{compositionally complete}: every $N$-qubit entanglement class under SLOCC is generated
by composing these tripartite primitives with bipartite entanglement and product states
\autocite{coeckekissinger2010compositional}. Above three parties nothing new appears; the
entanglement alphabet closes at exactly the point where the fibration tower terminates. The
mutual silence of the two literatures is itself striking: Coecke and Kissinger make no reference
to Hopf fibrations or division algebras, and the fibration literature on entanglement
\autocite{mosseri2001hopf,bernevigchen2003geometry} nowhere cites compositional completeness.
Two classifications, in different mathematics, draw the same line in the same place, unaware of
one another. That silence has, in the months this paper was in draft, just begun to break from
the categorical side, with a projective, base-level reading of the same calculus
\autocite{burtonanwar2026meromorphic}; we engage it below.

The framework's conjecture is that this is one termination wearing a third face: the alphabet
has the letters it has --- Bell, GHZ, W --- because the parallelisable spheres have the
dimensions they have. There is partial evidence at every rung. At one qubit, the commutative
special Frobenius algebras on $\mathbb{C}^{2}$ are classified by orthonormal bases --- points of
$S^{2}$, the base of the first fibration --- so that the copied basis plays the role of a
section of the bundle and the spider theorem expresses its idempotence. At two qubits, the
$SU(2)$ fibre of $S^{7} \to S^{4}$ carries exactly the entanglement that survives local
operations, and L\'{e}vay's result that the bundle's natural connection is the $SU(2)$
Yang--Mills instanton \autocite{levay2004geometry} stands on the geometric side where the
interaction of two Frobenius structures stands on the categorical one. At three qubits, the
split of D\"{u}r, Vidal and Cirac \autocite{durvidalcirac2000three} mirrors the algebra exactly:
GHZ is genuinely tripartite, demanding the full non-associative octonion, while W is
pairwise-robust and housed in the associative subalgebra $\mathbb{H} \subset \mathbb{O}$ --- the
special/anti-special Frobenius distinction in geometric dress.

The first of the four conditions the correspondence requires can be discharged at once, and we
do so here, because a problem whose first step has been walked is a different kind of invitation
from one merely posed.

\begin{proposition}[The one-qubit dictionary]\label{prop:f1-dictionary} The $\dagger$-special
commutative Frobenius algebras on $\mathbb{C}^{2}$ correspond bijectively to
\emph{lifted antipodal skeletons} of the complex Hopf fibration $\pi \colon S^{3} \to S^{2}$:
unordered antipodal pairs $P \subset S^{2}$ together with a section of $\pi$ over $P$. The
copyable states of a structure are exactly the section's image, so the lifts are recorded
faithfully --- the projectively trivial centre of $U(2)$ acts non-trivially on structures ---
and the correspondence is $U(2)$-equivariant, the ordered structures forming a single free and
transitive $U(2)$-orbit. The channel a structure induces on density matrices is dephasing in its
skeleton and is independent of the section: the fibre freedom, a torsor for
$\mathrm{Map}(P, U(1))$, is exactly the data absent from the classical record a structure
induces.
\end{proposition}

\begin{proof}
By the classification of Coecke, Pavlovic and Vicary, such structures correspond to orthonormal
bases, recovered as the sets of copyable states \autocite{coeckepavlovicvicary2013orthogonal}.
An orthonormal basis of $\mathbb{C}^{2}$ is two orthogonal rays --- an antipodal pair on the
Bloch sphere --- together with a unit representative of each, which is precisely a section of
$\pi$ over the pair; and $\delta(e^{i\theta}x) = e^{i\theta}\, x \otimes x$ against
$(e^{i\theta}x) \otimes (e^{i\theta}x) = e^{2i\theta}\, x \otimes x$ pins the phases.
Equivariance, the action of the centre, the orbit structure and the induced channel are direct
computations; an executable verification of every finite check accompanies the code release.
\end{proof}

Two remarks fix the proposition's scope. It corrects a natural first reading: a classical
structure does not project to the base but selects lifts into the total space --- classicality
is section data, and what the copy operation broadcasts is a lift, not a ray. And it claims
nothing topological about cloning: the finiteness of the copyable set is linear, structural
no-cloning residing in non-naturality \autocite{abramskyheunen2012hstar}, so the first place the
bundle's global topology can bite is the quaternionic rung, where the degree-one clutching map
stands ready to be compared with the strong-complementarity law governing an interacting pair of
structures. The dictionary is, finally, the complement of the projective treatment cited above:
Burton and Anwar pass functorially to the base and deliberately discard the fibre
\autocite{burtonanwar2026meromorphic}, and the classical structures turn out to live in exactly
the data their projectivisation forgets --- the two readings compose, and the resulting triangle
is a consistency condition available to constrain the remaining conditions of the
correspondence.

\begin{openproblem}[The Hopf-fibration--Frobenius correspondence]\label{op:hopf-frobenius} Let
$\mathbf{Hopf}$ denote the category of the three fibrations
$S^{2^{k}-1} \hookrightarrow S^{2^{k+1}-1} \to S^{2^{k}}$, $k = 1, 2, 3$, with bundle maps as
morphisms, and let $\mathbf{Frob}$ denote the Frobenius algebras on $\mathbb{C}^{2}$ in
$\mathbf{FdHilb}$, with the commutative-special (GHZ) and anti-special (W) classes
distinguished. Construct a functor $F \colon \mathbf{Hopf} \to \mathbf{Frob}$ --- or prove that
none exists --- under which the complex fibration yields the single-qubit classical structures
--- Proposition~\ref{prop:f1-dictionary} identifies these as the lifted antipodal skeletons,
section data over two base points rather than the points alone; the quaternionic fibration
yields the interaction of two such structures, with the $SU(2)$ fibre carrying the entanglement
that survives the classical quotient; the octonionic fibration yields the GHZ/W pair, GHZ on the
full non-associative fibre and W on the associative $\mathbb{H} \subset \mathbb{O}$; and the
composition by which Coecke--Kissinger generate the $N$-qubit classes corresponds to a
well-defined geometric operation on the bundles. If such an $F$ exists, compositional
completeness follows from Adams's theorem as a topological corollary --- there is no fourth
fibration to supply a fourth primitive --- and if $F$ is faithful enough that the implication
also runs back, the termination of the tower and the completeness of the alphabet are equivalent
expressions of a single classification. The first condition is discharged by
Proposition~\ref{prop:f1-dictionary}; what remains open is the quaternionic and octonionic
conditions and, above all, the correspondence of compositions. Both open conditions are already
constrained: the quaternionic fibre's physical content --- the recoverable relative phase and
its purity ledger --- is established (\S\ref{sec:wf-reconciliation}), leaving the
clutching-map-to-strong-complementarity comparison as the structural residue; and any candidate
functor must reproduce the proved canonical-orbit images, GHZ to the colour triplet and W to the
singlet (\S\ref{sec:canonical-orbits}, Appendix~\ref{app:ghz-w-proofs}). A proof that no such
functor exists, or a correspondence at the level of the classifications without functoriality on
compositions, would each also resolve the problem: the first would make the agreement of two
independent classifications a fact demanding its own explanation; the second would establish the
weaker reading.
\end{openproblem}

A resolution in the affirmative would do three things. It would explain compositional
completeness rather than merely prove it --- a corollary of the division-algebra structure of
quantum mechanics rather than an independent algebraic fact. It would give the diagrammatic
language of this paper's own interaction circuits a topological warrant: the GHZ/W spider
calculus would be the unique language the topology permits, not a convenient notation. And it
would open a two-way bridge, translating problems on compositions of Frobenius structures into
operations on bundles and back. The problem sits at the intersection of categorical quantum
mechanics, the quantum-information geometry of the fibrations, and the division-algebra
approaches to particle physics, and it is posed here so that any of the three communities can
take it up in its own language. It is perhaps worth noticing that both classifications already
answer to the same name --- the fibrations are Heinz Hopf's, and so, through the bialgebra
structure of the diagrammatic calculus, are the Hopf algebras --- a coincidence of eponymy that
the conjecture proposes to promote to a theorem. With the epistemic reading above, an
affirmative resolution would leave the tower's endpoint as one limit seen four ways: algebraic,
geometric, compositional, epistemic --- the last being the fixed-point classification of Open
Problem~\ref{op:sigma-functor}, with which this problem is cross-linked, each constraining the
other.

\subsection{What this sets up}

The entanglement-sensitivity of the Hopf maps is not a side feature but the hinge of the
physical reading that follows. Because separability is exactly the collapse of the base to a
product of Bloch spheres, the structure of the entanglement --- the SLOCC class of a multi-qubit
state, and within it the canonical maximally entangled orbit --- is carried by precisely the
degrees of freedom on which $G_{\mathrm{SM}}$ acts in Proposition~\ref{prop:gsm}. This is what
makes it possible, in \S\ref{sec:entanglement-classes}, to read the two genuinely tripartite
entanglement classes onto the $\mathbb{C} \oplus \mathbb{C}^3$ split through their canonical
representatives: the GHZ orbit confined to the colour-triplet $\mathbb{C}^3$ sector, the W orbit
reaching the $\mathbb{C}$ singlet, recovering the quark--lepton division not as an assignment
but as a consequence of the entanglement structure. The scope of those statements --- canonical
orbits rather than whole classes, and the identification of the canonical orbits with the
physical particles --- is made precise there. The Hopf fibrations are the bridge; the
entanglement classes, treated next, are the traffic that crosses it.

%  The section is position-portable: all adjacency phrasing was
%  retargeted to explicit \S\ref on 2026-07-06. The S^7 hand-off
%  sentence to sec:space-of-computations WRITTEN 2026-07-06 (spacetime
%  file read; the closing now hands first to its immediate consumer,
%  then to hopf-justification). Upload verifications COMPLETE: L2 and
%  quantum-postulates do not consume the ladder; WF's references are
%  forward-voiced or self-carrying.
%
%  DIVISION OF LABOUR (deliberate, to avoid overlap):
%   - sec:self-reference (undrafted): self-reference -> complex
%     amplitudes and the Born rule; why not a real base. This section
%     TAKES that as given and starts at C.
%   - THIS section: the algebra -- the doubling construction, the
%     ledger of losses-as-physics, the internal-algebras-not-a-new-QM
%     scoping (Adler pre-empt), the two appearances of the ladder
%     (register / measurement climb, stated lightly), and Hurwitz as
%     the algebraic termination with the committed zero-divisor
%     reading.
%   - sec:hopf-justification: the geometry (Adams, parallelisability,
%     the fibrations), the join of the two classifications, the
%     entanglement-sensitivity literature, the C-direction honest debt,
%     and the epistemic reading of termination.
%
%  COMMITTED READINGS USED (sources re-read 2026-07-03):
%   - Loss ledger and level dictionary: cayley_dickson_qubits_for_review
%     sec.5 ("each step loses a symmetry, and each lost symmetry is
%     identified with a physical content").
%   - Termination reading: grounding_argument_gauge_group ("the
%     sedenions have zero divisors and cannot support consistent
%     information flow around loops") -- leads; norm-composition-as-
%     probability given as its amplitude-level restatement.
%   - Internal-not-base scoping: cayley_dickson_qubits_for_review Q2 and
%     sec.3.2 caveat (Adler objection; C-qubit treated geometrically).
%
%  HONEST DEBT: ladder proved closed at the top (Hurwitz), derived at
%  the floor. The rung-compulsion revisit-trigger FIRED 2026-07-06:
%  sec:self-reference and op:sigma-functor now exist, so the thesis is
%  minted as op:ladder-compulsion (with the CD-conjugation-as-sigma
%  candidate route), inheriting op:sigma-functor's status.
%
%  AUDITED 2026-07-06 (final pass, third stop): all four citations
%  used; all cross-refs live; CD facts, dims, alternativity, Hurwitz
%  statement verified; Adler subsection delivers everything the
%  register's compressed pointer promises; file pure ASCII; no bare
%  "frame". Fixes: opening rewritten for both predecessors (continuity
%  was inserted between after drafting); "associator's non-vanishing is
%  mass itself" recalibrated to the register's debt-against-selection
%  defusal; the O->S ledger row and Hurwitz subsection reconciled with
%  the fourth-qubit forward remark ("the terminus closes the ladder of
%  forces; what survives past it is, at most, flavour").
%
%  CROSS-REFS: sec:self-reference [sibling dangle -- resolves when that
%  section is drafted], sec:hopf-justification, sec:hopf-information-bound,
%  sec:register, sec:entanglement-classes, sec:generations, sec:masses,
%  sec:wigners-friend. All others exist.
%  CITATIONS: existing keys baez2002octonions, hurwitz1898composition,
%  furey2018three. ONE new key: adler1995quaternionic (RIS in
%  division_algebras_refs_2026-07-03.ris).
%  British English; \autocite; no bullets in prose; the ledger is a
%  table because the comparison is the argument.
% =====================================================================

\section{The Division-Algebra Ladder}
\label{sec:division-algebras}

The foundations behind us end at amplitudes twice over: natively, where the self-referential
argument forced the description of a system into the system's own state space and priced
consistency at a number system whose parts can cancel --- complex amplitudes, and the Born rule
that squares them (\S\ref{sec:self-reference}) --- and independently, where the regime's
continuity delivers the same field through Hardy's theorem (\S\ref{sec:continuity}), the
convergence of the two branches being evidence rather than accident. This section asks the question
both routes leave open. Why
should the number system stop at $\mathbb{C}$? Ordinary quantum mechanics stops there by postulate;
a framework that has forbidden itself postulates must answer, and the answer is that the climb
neither stops at $\mathbb{C}$ nor continues forever. There are exactly two more rungs, each
purchased by surrendering an algebraic law and each paying for itself in physics, and the ladder is
closed at the top by a theorem of 1898. This section sets out the algebra of that ladder: the single
construction that climbs it, the ledger of what each rung costs and buys, and the theorem that ends
it. The geometry of the same ladder --- that each rung acts on a sphere, that exactly those spheres
are parallelisable, and that the algebraic and geometric classifications are one --- is the business
of the geometry section (\S\ref{sec:hopf-justification}).

\subsection{One construction climbs the whole ladder}

The rungs are not a miscellany; a single construction generates them all. Given an algebra with a
conjugation, the Cayley--Dickson construction builds its double on ordered pairs, with a twisted
multiplication in which the conjugate enters and a conjugation $(a,b)^{*} = (a^{*}, -b)$
\autocite{baez2002octonions}. (Conventions differ in the placement of the conjugations; all yield
isomorphic algebras.) Applied repeatedly it produces the sequence
$\mathbb{R} \to \mathbb{C} \to \mathbb{H} \to \mathbb{O} \to \mathbb{S}$, of real dimensions
$1, 2, 4, 8, 16$: the reals, the complex numbers, the quaternions, the octonions, the sedenions.
The ground floor is classical: $\mathbb{R}$ is ordered and positive-definite, the arithmetic of
classical probability. The first doubling is therefore the quantum move itself --- the passage from
probabilities to amplitudes that can interfere --- and it is not made here by choice: it is the
derivation of \S\ref{sec:self-reference}, already behind us. What this section must account for is
the rest of the climb.

\subsection{The ledger}

In mathematics the doublings are usually told as a story of degeneration: each algebra is worse
behaved than the last, and the sequence is studied for where it breaks. This framework reads the
same ledger in the opposite direction. Each rung is bought by surrendering a law, and each
surrendered law is exactly a piece of physics --- the losses are the assets, and the comparison is
sharp enough to be the argument:

\begin{center}
\begin{tabular}{llll}
\hline
doubling & dimension & what is lost & what the loss becomes \\
\hline
$\mathbb{R} \to \mathbb{C}$ & $1 \to 2$ & the order of the reals & phase and interference; amplitudes \\
 & & & that cancel; the Born rule (\S\ref{sec:self-reference}) \\
$\mathbb{C} \to \mathbb{H}$ & $2 \to 4$ & commutativity & non-commuting actions: weak isospin \\
 & & & $SU(2)$, the $B$ register (\S\ref{sec:register}) \\
$\mathbb{H} \to \mathbb{O}$ & $4 \to 8$ & associativity & the associator: colour, chirality, the \\
 & & & generations, and mass as debt \\
 & & & (\S\ref{sec:generations}, \S\ref{sec:masses}) \\
$\mathbb{O} \to \mathbb{S}$ & $8 \to 16$ & division itself & no further force: zero divisors appear \\
 & & (zero divisors) & and the ladder of sectors ends (below) \\
\hline
\end{tabular}
\end{center}

Walking the rows: the order lost at the first doubling is the possibility of interference --- parts
of a probability that can oppose one another --- which is where the derivation of \S\ref{sec:self-reference}
lives. The commutativity lost at the second is the existence of actions whose order matters; the
unit quaternions form the group $SU(2)$, and the framework houses weak isospin there, as the
internal algebra of the $B$ register whose order-sensitive, non-commuting dynamics the interaction sections develop
(\S\ref{sec:register}, \S\ref{sec:entanglement-classes}). The associativity lost at the third ---
the octonions retain the weaker property of alternativity \autocite{baez2002octonions} --- creates
the associator $[a,b,c] = (ab)c - a(bc)$, and the associator is where most of this paper's physics
is transacted: the stabiliser of a preferred complex direction inside the octonions' automorphisms
is the colour group (\S\ref{sec:hopf-justification}); the triality structure and its $K_3$
obstruction generate the three generations and the sterile channel (\S\ref{sec:generations}); and
the associator is the currency of mass --- the debt against the selected direction that
\S\ref{sec:masses} prices, owed by whatever action moves the selection and not by residence at
the rung, which is the defusal \S\ref{sec:register} makes explicit. The algebraic facts
in the third column are theorems; the identifications in the fourth are established in the sections
cited, each carrying its own epistemic flag there, and are not re-litigated here. What this table
asserts on its own account is only the correspondence's shape: one law surrendered per rung, one
sector of physics gained per law.

\subsection{Internal algebras, not a new quantum mechanics}

A reader versed in the quaternionic-quantum-mechanics literature will at this point raise a fair
objection: quantum theory over $\mathbb{H}$ as a foundation is known to be troubled --- the
tensor-product and composite-system difficulties set out by Adler and others
\autocite{adler1995quaternionic} --- and quantum theory over $\mathbb{O}$ founders earlier still,
since a non-associative algebra does not furnish a Hilbert space at all. The framework's position
must therefore be stated exactly, because it does not make the claim the objection targets. The base
of the physical Hilbert space here is $\mathbb{C}$ and remains $\mathbb{C}$; the higher algebras
enter as the \emph{internal} algebras of the register --- $\mathbb{H}$ for the $B$ qubit,
$\mathbb{O}$ for the $C$ qubit (\S\ref{sec:register}) --- in a tensor structure whose outer theory is
ordinary complex quantum mechanics. This is the same move by which the division algebras enter the
constructions of Furey and the related programmes without any proposal to rebase quantum mechanics
itself \autocite{furey2018three}. For the octonionic register the point is sharper still:
$\mathbb{O}^{2}$ is not a Hilbert space, and the framework accordingly treats the $C$ qubit not in
the Dirac--von Neumann manner but geometrically, as the sphere $S^{15}$ with its Hopf structure ---
which is precisely the form in which \S\ref{sec:hopf-justification} needs it. Nothing in this paper therefore
stands or falls with quaternionic or octonionic quantum mechanics as foundations; the ladder is
climbed by the registers' internal algebras, not by the theory's base field.

\subsection{The same ladder, twice}

The ladder appears in this framework in two guises, and it will help the reader to have them
distinguished now. Statically, it is the anatomy of a single relation: the three qubits of the
register take their values one rung apart --- $A$ over $\mathbb{C}$, $B$ over $\mathbb{H}$, $C$ over
$\mathbb{O}$ --- so that every particle carries the whole ladder as its internal structure
(\S\ref{sec:register}). Dynamically, it is the shape of measurement: a single qubit's pure states
fill $S^{3}$, an entangling measurement carries the joint state into $S^{7}$, and a third party into
$S^{15}$ --- the correspondence between qubit entanglement and the Hopf fibrations that
\S\ref{sec:hopf-justification} takes from a mature literature --- so that each act of measurement ascends one rung --- the climb the reader has already made,
narrated in \S\ref{sec:wigners-friend} --- and every completed relation is a terminated climb, which is
the reading \S\ref{sec:hopf-information-bound} gives the tower's endpoint. One ladder, seen twice:
as the anatomy of a relation and as the dynamics of relating. The framework regards this double
appearance not as a coincidence but as the same structure viewed statically and in motion, and the
sections cited carry that identification's weight.

\subsection{Hurwitz closes the ladder}

The construction itself would double forever; what stops it is a classification theorem. By
Hurwitz's theorem, the normed division algebras --- those in which the norm composes,
$\lvert ab\rvert = \lvert a\rvert\,\lvert b\rvert$, and no non-zero elements multiply to zero ---
are exactly $\mathbb{R}$, $\mathbb{C}$, $\mathbb{H}$ and $\mathbb{O}$
\autocite{hurwitz1898composition,baez2002octonions}. The next doubling, the sedenions, has zero
divisors: non-zero elements whose product vanishes. The framework's reading of this failure is the
one that matters physically. Zero divisors cannot support consistent information flow around loops
--- a cycle of compositions can silently annihilate what it carries --- and a framework in which
closed computational paths are load-bearing (\S\ref{sec:hopf-justification},
\S\ref{sec:space-of-computations}) cannot be built on an algebra that permits it. At the level of
amplitudes the same failure reads: the composition law
$\lvert ab\rvert = \lvert a\rvert\,\lvert b\rvert$, which carries probability multiplicatively
through every rung of the ladder, fails beyond the octonions, so that two non-vanishing amplitudes
could compose to nothing and probability would be destroyed by mere composition. The last algebra
fit to carry amplitudes is $\mathbb{O}$. One rung past the terminus is not entirely silent,
however: the sedenions' internal structure --- their automorphism group and their octonion
subalgebras --- is the recorded candidate arena for tagging the generations
(\S\ref{sec:interactions}), and it enters exactly as this section's scoping permits, internally
and never as a base. The terminus closes the ladder of forces; what survives past it is, at most,
flavour. This is the algebraic half of the termination; the
geometric half --- that no sphere above $S^{7}$ is parallelisable --- and the theorem joining the
two halves are the opening business of \S\ref{sec:hopf-justification}, and the epistemic reading of the same endpoint,
as a bound on the information a single relation can carry, is given in
\S\ref{sec:hopf-information-bound}.

\subsection{The honest debt}

As in the geometry section, we close by stating plainly what is proved, what is established, and what is
owed. Proved: the construction, the losses, and the termination are theorems --- the Cayley--Dickson
doubling and its ledger of surrendered laws are standard algebra, and Hurwitz's classification is
among the durable results of the subject. Established: the physical cashing of each rung --- the
Born rule at $\mathbb{C}$, weak isospin at $\mathbb{H}$, colour, generations and mass at
$\mathbb{O}$ --- is carried out in the sections cited above, each under its own epistemic flag. Owed,
and stated without disguise: the ladder as presented here is closed from above and derived at its
floor, but it is not yet derived rung by rung from below. Hurwitz forbids a further rung; the
argument of \S\ref{sec:self-reference} forces the first doubling, from the ordered arithmetic of
classical probability to amplitudes that interfere. That the same self-referential closure, applied
to a system whose amplitudes already take values at one rung, compels the next --- that relational
closure at one level forces precisely the Cayley--Dickson double of that level --- is the
framework's reading of the climb, and it is a thesis, not a theorem --- but the formal apparatus
against which it can be posed now exists, and we pose it.

\begin{openproblem}[The compulsion of the intermediate rungs]\label{op:ladder-compulsion}
Derive, within the categorical apparatus of Open Problem~\ref{op:sigma-functor}, that
self-referential closure at one rung compels the Cayley--Dickson double of that rung, so that the
intermediate steps $\mathbb{C} \to \mathbb{H} \to \mathbb{O}$ follow from the same principle
that forces the floor. The candidate route is the construction's own shape: the Cayley--Dickson
pair with its conjugation, $(a, b)^{*} = (a^{*}, -b)$, is an involution-equipped doubling, and the
route reads it as the algebraic shadow of adjoining the describer --- closure of level-$A$
descriptions under the inside/outside exchange requiring exactly the twisted pair. The problem
inherits the status of Open Problem~\ref{op:sigma-functor}, whose apparatus it presupposes.
\end{openproblem}

If the derivation lands, the whole ladder becomes a consequence of the same self-reference that
opened it. Until then, the honest statement is that the ladder's necessity is established at its
ceiling and at its floor, while the compulsion of the intermediate rungs is inherited from the
physics they are shown to carry rather than derived in advance. What this section hands on is, accordingly, an algebraic ladder with a floor, a ceiling and a
ledger. Its first consumer is immediate: \S\ref{sec:space-of-computations} spends the terminal
fact --- the non-associativity of $S^{7}$'s own multiplication --- as geometry, excluding the
seven-sphere as a ruler and fixing the dimension of space at three.
\S\ref{sec:hopf-justification} then supplies the geometry of the fibrations themselves, and the
join.
   <-- moved here
%      % EDITED 2026-10-07 scripts/patch_infinite_not_flat_2026-10-07T1400.py
% EDITED 2026-09-27 scripts/patch_spacelike_inversion_2026-09-27T1800.py
% EDITED 2026-09-21 scripts/patch_op_sterile_gravity_boundary_2026-09-21T1300.py
% EDITED 2026-09-12 scripts/patch_spacetime_bloch_headline_2026-09-12T1400.py
% EDITED 2026-09-11 scripts/patch_splice_loops_prospects_2026-09-11T0900.py
% EDITED 2026-09-10 scripts/patch_cite_rubino_2026-09-10T0700.py
% EDITED 2026-09-10 scripts/patch_quoting_and_typography_2026-09-10T0620.py
% =====================================================================
%  SECTION: Spacetime from non-computability (sec:space-of-computations).
%  AUDITED 2026-07-06 (final pass, seventh stop).
%
%  PLACEMENT (reordered chain, JS 2026-07-06): position 7 --
%      % =====================================================================
%  SECTION: The Division-Algebra Ladder
%  Drafted 2026-07-03. One of the three foundations sections
%  (sec:self-reference, sec:division-algebras, sec:fanout-boundary);
%  this one resolves the dangling \S\ref{sec:division-algebras} in the
%  Hopf-justification section, whose header names it "the
%  Cayley--Dickson / Hurwitz step".
%
%  PLACEMENT in main.tex (REORDERED per JS 2026-07-06): position 6,
%  opening Act II ("why stop at C?") right before its first consumer --
%      \input{self_reference_section_2026-07-04}
%      \input{continuity_section_2026-07-04}
%      \input{L2_metric_global_phase_section_2026_06_07}
%      \input{wigners_friend_in_the_hopf_picture_2026-06-13}
%      \input{derivation_of_the_quantum_postulates}
%      \input{division_algebras_section_2026-07-03}   <-- moved here
%      \input{spacetime_from_non-computability}       <-- first consumer
%                                                          (S^7 exclusion)
%      \input{hopf_justification_section_2026-06-08}
%  The section is position-portable: all adjacency phrasing was
%  retargeted to explicit \S\ref on 2026-07-06. The S^7 hand-off
%  sentence to sec:space-of-computations WRITTEN 2026-07-06 (spacetime
%  file read; the closing now hands first to its immediate consumer,
%  then to hopf-justification). Upload verifications COMPLETE: L2 and
%  quantum-postulates do not consume the ladder; WF's references are
%  forward-voiced or self-carrying.
%
%  DIVISION OF LABOUR (deliberate, to avoid overlap):
%   - sec:self-reference (undrafted): self-reference -> complex
%     amplitudes and the Born rule; why not a real base. This section
%     TAKES that as given and starts at C.
%   - THIS section: the algebra -- the doubling construction, the
%     ledger of losses-as-physics, the internal-algebras-not-a-new-QM
%     scoping (Adler pre-empt), the two appearances of the ladder
%     (register / measurement climb, stated lightly), and Hurwitz as
%     the algebraic termination with the committed zero-divisor
%     reading.
%   - sec:hopf-justification: the geometry (Adams, parallelisability,
%     the fibrations), the join of the two classifications, the
%     entanglement-sensitivity literature, the C-direction honest debt,
%     and the epistemic reading of termination.
%
%  COMMITTED READINGS USED (sources re-read 2026-07-03):
%   - Loss ledger and level dictionary: cayley_dickson_qubits_for_review
%     sec.5 ("each step loses a symmetry, and each lost symmetry is
%     identified with a physical content").
%   - Termination reading: grounding_argument_gauge_group ("the
%     sedenions have zero divisors and cannot support consistent
%     information flow around loops") -- leads; norm-composition-as-
%     probability given as its amplitude-level restatement.
%   - Internal-not-base scoping: cayley_dickson_qubits_for_review Q2 and
%     sec.3.2 caveat (Adler objection; C-qubit treated geometrically).
%
%  HONEST DEBT: ladder proved closed at the top (Hurwitz), derived at
%  the floor. The rung-compulsion revisit-trigger FIRED 2026-07-06:
%  sec:self-reference and op:sigma-functor now exist, so the thesis is
%  minted as op:ladder-compulsion (with the CD-conjugation-as-sigma
%  candidate route), inheriting op:sigma-functor's status.
%
%  AUDITED 2026-07-06 (final pass, third stop): all four citations
%  used; all cross-refs live; CD facts, dims, alternativity, Hurwitz
%  statement verified; Adler subsection delivers everything the
%  register's compressed pointer promises; file pure ASCII; no bare
%  "frame". Fixes: opening rewritten for both predecessors (continuity
%  was inserted between after drafting); "associator's non-vanishing is
%  mass itself" recalibrated to the register's debt-against-selection
%  defusal; the O->S ledger row and Hurwitz subsection reconciled with
%  the fourth-qubit forward remark ("the terminus closes the ladder of
%  forces; what survives past it is, at most, flavour").
%
%  CROSS-REFS: sec:self-reference [sibling dangle -- resolves when that
%  section is drafted], sec:hopf-justification, sec:hopf-information-bound,
%  sec:register, sec:entanglement-classes, sec:generations, sec:masses,
%  sec:wigners-friend. All others exist.
%  CITATIONS: existing keys baez2002octonions, hurwitz1898composition,
%  furey2018three. ONE new key: adler1995quaternionic (RIS in
%  division_algebras_refs_2026-07-03.ris).
%  British English; \autocite; no bullets in prose; the ledger is a
%  table because the comparison is the argument.
% =====================================================================

\section{The Division-Algebra Ladder}
\label{sec:division-algebras}

The foundations behind us end at amplitudes twice over: natively, where the self-referential
argument forced the description of a system into the system's own state space and priced
consistency at a number system whose parts can cancel --- complex amplitudes, and the Born rule
that squares them (\S\ref{sec:self-reference}) --- and independently, where the regime's
continuity delivers the same field through Hardy's theorem (\S\ref{sec:continuity}), the
convergence of the two branches being evidence rather than accident. This section asks the question
both routes leave open. Why
should the number system stop at $\mathbb{C}$? Ordinary quantum mechanics stops there by postulate;
a framework that has forbidden itself postulates must answer, and the answer is that the climb
neither stops at $\mathbb{C}$ nor continues forever. There are exactly two more rungs, each
purchased by surrendering an algebraic law and each paying for itself in physics, and the ladder is
closed at the top by a theorem of 1898. This section sets out the algebra of that ladder: the single
construction that climbs it, the ledger of what each rung costs and buys, and the theorem that ends
it. The geometry of the same ladder --- that each rung acts on a sphere, that exactly those spheres
are parallelisable, and that the algebraic and geometric classifications are one --- is the business
of the geometry section (\S\ref{sec:hopf-justification}).

\subsection{One construction climbs the whole ladder}

The rungs are not a miscellany; a single construction generates them all. Given an algebra with a
conjugation, the Cayley--Dickson construction builds its double on ordered pairs, with a twisted
multiplication in which the conjugate enters and a conjugation $(a,b)^{*} = (a^{*}, -b)$
\autocite{baez2002octonions}. (Conventions differ in the placement of the conjugations; all yield
isomorphic algebras.) Applied repeatedly it produces the sequence
$\mathbb{R} \to \mathbb{C} \to \mathbb{H} \to \mathbb{O} \to \mathbb{S}$, of real dimensions
$1, 2, 4, 8, 16$: the reals, the complex numbers, the quaternions, the octonions, the sedenions.
The ground floor is classical: $\mathbb{R}$ is ordered and positive-definite, the arithmetic of
classical probability. The first doubling is therefore the quantum move itself --- the passage from
probabilities to amplitudes that can interfere --- and it is not made here by choice: it is the
derivation of \S\ref{sec:self-reference}, already behind us. What this section must account for is
the rest of the climb.

\subsection{The ledger}

In mathematics the doublings are usually told as a story of degeneration: each algebra is worse
behaved than the last, and the sequence is studied for where it breaks. This framework reads the
same ledger in the opposite direction. Each rung is bought by surrendering a law, and each
surrendered law is exactly a piece of physics --- the losses are the assets, and the comparison is
sharp enough to be the argument:

\begin{center}
\begin{tabular}{llll}
\hline
doubling & dimension & what is lost & what the loss becomes \\
\hline
$\mathbb{R} \to \mathbb{C}$ & $1 \to 2$ & the order of the reals & phase and interference; amplitudes \\
 & & & that cancel; the Born rule (\S\ref{sec:self-reference}) \\
$\mathbb{C} \to \mathbb{H}$ & $2 \to 4$ & commutativity & non-commuting actions: weak isospin \\
 & & & $SU(2)$, the $B$ register (\S\ref{sec:register}) \\
$\mathbb{H} \to \mathbb{O}$ & $4 \to 8$ & associativity & the associator: colour, chirality, the \\
 & & & generations, and mass as debt \\
 & & & (\S\ref{sec:generations}, \S\ref{sec:masses}) \\
$\mathbb{O} \to \mathbb{S}$ & $8 \to 16$ & division itself & no further force: zero divisors appear \\
 & & (zero divisors) & and the ladder of sectors ends (below) \\
\hline
\end{tabular}
\end{center}

Walking the rows: the order lost at the first doubling is the possibility of interference --- parts
of a probability that can oppose one another --- which is where the derivation of \S\ref{sec:self-reference}
lives. The commutativity lost at the second is the existence of actions whose order matters; the
unit quaternions form the group $SU(2)$, and the framework houses weak isospin there, as the
internal algebra of the $B$ register whose order-sensitive, non-commuting dynamics the interaction sections develop
(\S\ref{sec:register}, \S\ref{sec:entanglement-classes}). The associativity lost at the third ---
the octonions retain the weaker property of alternativity \autocite{baez2002octonions} --- creates
the associator $[a,b,c] = (ab)c - a(bc)$, and the associator is where most of this paper's physics
is transacted: the stabiliser of a preferred complex direction inside the octonions' automorphisms
is the colour group (\S\ref{sec:hopf-justification}); the triality structure and its $K_3$
obstruction generate the three generations and the sterile channel (\S\ref{sec:generations}); and
the associator is the currency of mass --- the debt against the selected direction that
\S\ref{sec:masses} prices, owed by whatever action moves the selection and not by residence at
the rung, which is the defusal \S\ref{sec:register} makes explicit. The algebraic facts
in the third column are theorems; the identifications in the fourth are established in the sections
cited, each carrying its own epistemic flag there, and are not re-litigated here. What this table
asserts on its own account is only the correspondence's shape: one law surrendered per rung, one
sector of physics gained per law.

\subsection{Internal algebras, not a new quantum mechanics}

A reader versed in the quaternionic-quantum-mechanics literature will at this point raise a fair
objection: quantum theory over $\mathbb{H}$ as a foundation is known to be troubled --- the
tensor-product and composite-system difficulties set out by Adler and others
\autocite{adler1995quaternionic} --- and quantum theory over $\mathbb{O}$ founders earlier still,
since a non-associative algebra does not furnish a Hilbert space at all. The framework's position
must therefore be stated exactly, because it does not make the claim the objection targets. The base
of the physical Hilbert space here is $\mathbb{C}$ and remains $\mathbb{C}$; the higher algebras
enter as the \emph{internal} algebras of the register --- $\mathbb{H}$ for the $B$ qubit,
$\mathbb{O}$ for the $C$ qubit (\S\ref{sec:register}) --- in a tensor structure whose outer theory is
ordinary complex quantum mechanics. This is the same move by which the division algebras enter the
constructions of Furey and the related programmes without any proposal to rebase quantum mechanics
itself \autocite{furey2018three}. For the octonionic register the point is sharper still:
$\mathbb{O}^{2}$ is not a Hilbert space, and the framework accordingly treats the $C$ qubit not in
the Dirac--von Neumann manner but geometrically, as the sphere $S^{15}$ with its Hopf structure ---
which is precisely the form in which \S\ref{sec:hopf-justification} needs it. Nothing in this paper therefore
stands or falls with quaternionic or octonionic quantum mechanics as foundations; the ladder is
climbed by the registers' internal algebras, not by the theory's base field.

\subsection{The same ladder, twice}

The ladder appears in this framework in two guises, and it will help the reader to have them
distinguished now. Statically, it is the anatomy of a single relation: the three qubits of the
register take their values one rung apart --- $A$ over $\mathbb{C}$, $B$ over $\mathbb{H}$, $C$ over
$\mathbb{O}$ --- so that every particle carries the whole ladder as its internal structure
(\S\ref{sec:register}). Dynamically, it is the shape of measurement: a single qubit's pure states
fill $S^{3}$, an entangling measurement carries the joint state into $S^{7}$, and a third party into
$S^{15}$ --- the correspondence between qubit entanglement and the Hopf fibrations that
\S\ref{sec:hopf-justification} takes from a mature literature --- so that each act of measurement ascends one rung --- the climb the reader has already made,
narrated in \S\ref{sec:wigners-friend} --- and every completed relation is a terminated climb, which is
the reading \S\ref{sec:hopf-information-bound} gives the tower's endpoint. One ladder, seen twice:
as the anatomy of a relation and as the dynamics of relating. The framework regards this double
appearance not as a coincidence but as the same structure viewed statically and in motion, and the
sections cited carry that identification's weight.

\subsection{Hurwitz closes the ladder}

The construction itself would double forever; what stops it is a classification theorem. By
Hurwitz's theorem, the normed division algebras --- those in which the norm composes,
$\lvert ab\rvert = \lvert a\rvert\,\lvert b\rvert$, and no non-zero elements multiply to zero ---
are exactly $\mathbb{R}$, $\mathbb{C}$, $\mathbb{H}$ and $\mathbb{O}$
\autocite{hurwitz1898composition,baez2002octonions}. The next doubling, the sedenions, has zero
divisors: non-zero elements whose product vanishes. The framework's reading of this failure is the
one that matters physically. Zero divisors cannot support consistent information flow around loops
--- a cycle of compositions can silently annihilate what it carries --- and a framework in which
closed computational paths are load-bearing (\S\ref{sec:hopf-justification},
\S\ref{sec:space-of-computations}) cannot be built on an algebra that permits it. At the level of
amplitudes the same failure reads: the composition law
$\lvert ab\rvert = \lvert a\rvert\,\lvert b\rvert$, which carries probability multiplicatively
through every rung of the ladder, fails beyond the octonions, so that two non-vanishing amplitudes
could compose to nothing and probability would be destroyed by mere composition. The last algebra
fit to carry amplitudes is $\mathbb{O}$. One rung past the terminus is not entirely silent,
however: the sedenions' internal structure --- their automorphism group and their octonion
subalgebras --- is the recorded candidate arena for tagging the generations
(\S\ref{sec:interactions}), and it enters exactly as this section's scoping permits, internally
and never as a base. The terminus closes the ladder of forces; what survives past it is, at most,
flavour. This is the algebraic half of the termination; the
geometric half --- that no sphere above $S^{7}$ is parallelisable --- and the theorem joining the
two halves are the opening business of \S\ref{sec:hopf-justification}, and the epistemic reading of the same endpoint,
as a bound on the information a single relation can carry, is given in
\S\ref{sec:hopf-information-bound}.

\subsection{The honest debt}

As in the geometry section, we close by stating plainly what is proved, what is established, and what is
owed. Proved: the construction, the losses, and the termination are theorems --- the Cayley--Dickson
doubling and its ledger of surrendered laws are standard algebra, and Hurwitz's classification is
among the durable results of the subject. Established: the physical cashing of each rung --- the
Born rule at $\mathbb{C}$, weak isospin at $\mathbb{H}$, colour, generations and mass at
$\mathbb{O}$ --- is carried out in the sections cited above, each under its own epistemic flag. Owed,
and stated without disguise: the ladder as presented here is closed from above and derived at its
floor, but it is not yet derived rung by rung from below. Hurwitz forbids a further rung; the
argument of \S\ref{sec:self-reference} forces the first doubling, from the ordered arithmetic of
classical probability to amplitudes that interfere. That the same self-referential closure, applied
to a system whose amplitudes already take values at one rung, compels the next --- that relational
closure at one level forces precisely the Cayley--Dickson double of that level --- is the
framework's reading of the climb, and it is a thesis, not a theorem --- but the formal apparatus
against which it can be posed now exists, and we pose it.

\begin{openproblem}[The compulsion of the intermediate rungs]\label{op:ladder-compulsion}
Derive, within the categorical apparatus of Open Problem~\ref{op:sigma-functor}, that
self-referential closure at one rung compels the Cayley--Dickson double of that rung, so that the
intermediate steps $\mathbb{C} \to \mathbb{H} \to \mathbb{O}$ follow from the same principle
that forces the floor. The candidate route is the construction's own shape: the Cayley--Dickson
pair with its conjugation, $(a, b)^{*} = (a^{*}, -b)$, is an involution-equipped doubling, and the
route reads it as the algebraic shadow of adjoining the describer --- closure of level-$A$
descriptions under the inside/outside exchange requiring exactly the twisted pair. The problem
inherits the status of Open Problem~\ref{op:sigma-functor}, whose apparatus it presupposes.
\end{openproblem}

If the derivation lands, the whole ladder becomes a consequence of the same self-reference that
opened it. Until then, the honest statement is that the ladder's necessity is established at its
ceiling and at its floor, while the compulsion of the intermediate rungs is inherited from the
physics they are shown to carry rather than derived in advance. What this section hands on is, accordingly, an algebraic ladder with a floor, a ceiling and a
ledger. Its first consumer is immediate: \S\ref{sec:space-of-computations} spends the terminal
fact --- the non-associativity of $S^{7}$'s own multiplication --- as geometry, excluding the
seven-sphere as a ruler and fixing the dimension of space at three.
\S\ref{sec:hopf-justification} then supplies the geometry of the fibrations themselves, and the
join.
   <-- hands off here
%      % EDITED 2026-10-07 scripts/patch_infinite_not_flat_2026-10-07T1400.py
% EDITED 2026-09-27 scripts/patch_spacelike_inversion_2026-09-27T1800.py
% EDITED 2026-09-21 scripts/patch_op_sterile_gravity_boundary_2026-09-21T1300.py
% EDITED 2026-09-12 scripts/patch_spacetime_bloch_headline_2026-09-12T1400.py
% EDITED 2026-09-11 scripts/patch_splice_loops_prospects_2026-09-11T0900.py
% EDITED 2026-09-10 scripts/patch_cite_rubino_2026-09-10T0700.py
% EDITED 2026-09-10 scripts/patch_quoting_and_typography_2026-09-10T0620.py
% =====================================================================
%  SECTION: Spacetime from non-computability (sec:space-of-computations).
%  AUDITED 2026-07-06 (final pass, seventh stop).
%
%  PLACEMENT (reordered chain, JS 2026-07-06): position 7 --
%      \input{division_algebras_section_2026-07-03}   <-- hands off here
%      \input{spacetime_from_non-computability}
%      \input{hopf_justification_section_2026-06-08}
%  Division-algebras' closing now carries the S^7 hand-off sentence
%  (written 2026-07-06 after this file was read).
%
%  INPUT DEPENDENCIES (both held in the audited outputs):
%      entropic_gravity_family_paragraph_2026-07-03.tex  (sec:jacobson)
%      gr_nonlocality_revision_2026-06-13.tex            (sec:gr-nonlocality)
%  The 2026-07-04 additions (asymmetry paragraph; Penrose second route;
%  acceleration/Unruh arc) are FULLY PASTED; the separately uploaded
%  additions file is redundant and can be archived.
%
%  AUDIT RECORD: opener formalised (Russell typo; bare ref retargeted
%  to the freshly minted sec:complex-amplitudes); the "rotations of
%  S^7" wording made precise at BOTH occurrences (Moufang loop, not a
%  group; SO(8) disclaimer added -- referee hazard removed); two \\
%  paragraph-break bugs fixed; dark-matter paragraph retargeted to
%  sec:generations and stitched to op:sterile-gravity (the zoo's
%  demotion and this paragraph's priority condition are one question);
%  CP paragraph stitched to the fourth-party reading of
%  sec:interactions; stale mid-file mini-header retired
%  (sec:k3-sterile placeholder dies with it); the section's narrated
%  debts NUMBERED at last -- op:comp-metric and op:three-dimensions
%  minted in the Open-problems subsection, cited from the in-text
%  "named open problems" line and from the closing summary.
%
%  KEYS: adams1960hopf matches corpus. Unique to this section, presumed
%  in the .bib since the paper compiles: Spence-Brown-1969 (suggest
%  normalising the key spelling at leisure), VanRaamsdonk2010,
%  MaldacenaSusskind2013, MuellerMasanes2013, MasanesMuellerEtAl2014,
%  toffoli_2003, Chaitin1998, Penrose1996, HohnMueller2016,
%  GarnerMuellerDahlsten2017, Jacobson1995, Bekenstein1973, tHooft1993,
%  Susskind1995, spekkens2007toy, Andrianov2001, Sakharov1967,
%  Planck2020, Riess1998. Shared keys verified: small_2006, Penrose1994,
%  unruh1976notes (swap-flag stands), baez2002octonions, Wheeler1990,
%  furey2018three, DiBiagioRovelli2021, Rovelli1996a, Spekkens toy.
%  British English; \autocite; no bullets in prose.
% =====================================================================

\section{Spacetime from non-computability}
\label{sec:space-of-computations}

This section is a deliberate diversion: it takes the idea that non-computable steps are measured
by imaginary information (\S\ref{sec:complex-amplitudes}) and asks what geometry it buys.
Spencer-Brown records that Bertrand Russell thought well of the underlying move
\autocite{Spence-Brown-1969}; the reader may judge whether the framework spends it responsibly.

The programmes that derive spacetime from quantum mechanics---Van~Raamsdonk's
\enquote{building spacetime with entanglement} \autocite{VanRaamsdonk2010}, the
ER\,$=$\,EPR conjecture \autocite{MaldacenaSusskind2013}, and M\"uller's
information-theoretic recovery of three-dimensional space from the qubit
\autocite{MuellerMasanes2013}---take entanglement or information as primitive and
grow geometry from it. The route taken here is the same in spirit but inverted in
order: rather than positing spatial dimensions and forbidding superluminal
signalling, we \emph{define} spatial separation as a measure of non-computability
and read the dimension, the signature, and the Lorentz group off the result. The
proposal originates in the author's CASYS'05 paper \autocite{small_2006}, which
first suggested treating spacetime as a space of computations and attempted to
justify three dimensions on this basis; the present account tightens and extends
it. This is the framework's most structural and least formally complete arc; the
three-dimensionality result is a structural conclusion awaiting a formal proof,
and several of the links below are conjectural readings. Each is flagged as such,
so the reader sees exactly where the chain is load-bearing and where it is an
invitation.

\subsection{Timelike and spacelike distance as computation}
\label{sec:comp-distance}

A computation carries a state from axioms to a theorem in a number of steps.
Counting the computable steps between two states gives a real measure---time. A
physical clock counts completed irreversible loops, so proper time literally
counts loop completions and the metric $\mathrm{d}\tau^{2}=g_{\mu\nu}\,
\mathrm{d}x^{\mu}\mathrm{d}x^{\nu}$ is a loop-density tensor, the conversion from
coordinate intervals to loop counts; Toffoli's reading of the Lagrangian as a
count of action quanta is the same statement \autocite{toffoli_2003}.

Spacelike separation is the complementary, non-computable measure. Define the
non-computable distance between two propositions, relative to an axiom system, as
the number of bits that must be added to the axioms to make the second derivable.
A closed timelike curve is a computation whose output is its own input; by the
fixed-point theorem it has a unique consistent solution or none, and the
existence of that solution is one bit not derivable from any axiom system that
does not already contain the loop---one bit per curve, in the manner of Chaitin's
$\Omega$ \autocite{Chaitin1998}. The spacelike distance between two events is then
the number of such curves that would have to be closed to make them causally
connected. Because each added bit is a non-computable proposition, it is carried
by the imaginary part of the complex information measure
$S=-\ln\lvert p\rvert-\mathrm{i}\theta$ derived in
Section~\ref{sec:complex-amplitudes}: computable distance is counted in real
units, non-computable distance in imaginary ones. The Minkowski signature, with
spacelike intervals imaginary relative to timelike ones, is the immediate reading. The convention is not new to this paper: it was proposed in the author's CASYS paper two
decades ago, where counting non-causal steps in imaginary numbers was introduced and its formal
proof flagged, already then, as a key open problem \autocite{small_2006} ---
Open Problem~\ref{op:comp-metric} is that flag's direct descendant, stated now with the machinery
to be answerable.
That the resulting structure is a genuine metric space, and that the signature
follows rather than being fitted, are open problems---at present the
identification is motivated by reproducing what nature does, not proved.

%% VERIFY-SUPPORT 2026-09-27: conway_2003 is cited below for "any two octonions
%% generate an associative subalgebra" (Artin's theorem; a Moufang loop is
%% diassociative) and harvey1982calibrated for "a three-plane in Im O is
%% associative iff the associator of a basis vanishes, and the associative
%% three-planes are the imaginary parts of the quaternionic subalgebras"
%% (Harvey--Lawson, ch. IV). Neither passage was read in the session that wrote
%% this; both statements are checked numerically in Convention A by
%% scripts/antitelephone_network_associativity_2026-09-27T1900.py (PASS).
A reader used to the standard order of ideas should pause here, because the order has been
inverted, and the inversion is used below. In special relativity the tachyonic antitelephone is a
corollary: a signal faster than light, composed with a boost, would arrive before it was sent and
close a timelike loop, therefore there is no such signal \autocite{tolman1917theory,
benfordbooknewcomb1970tachyonic}. The metric comes first and the impossible telephone is a check
on it. Here the telephone comes first. Two observers are spacelike separated when an antitelephone
between them would close a loop; the loops such telephones would close, and carry nothing round,
are what spacelike separation is, their number is the distance, and the metric is read off from
them. Nothing is lost, since the corollary returns as a tautology: an antitelephone that carried
a record would close a causal loop, and its two ends would not have been spacelike after all.

What is gained is a definite meaning for a spatial dimension. Picture many observers, each holding
an antitelephone to each of the others. A message relayed from one observer to another through a
third composes two transports round two loops, and the network is consistent only if every relay
composes to a single transport whatever the bracketing of the relaying: a spatial dimension is a
family of loops whose transports compose associatively, which is a group law on a parallelisable
sphere. That is a condition a family of loops can fail, and \S\ref{sec:three-dimensions} shows
that the loops of the seventh sphere fail it. Two antitelephones always compose, because any two
octonions generate an associative subalgebra \autocite{conway_2003}; three compose only when their
imaginary parts lie in one quaternionic subalgebra, which is what the associative three-planes of
Harvey and Lawson are \autocite{harvey1982calibrated}, and a third link outside that subalgebra is
seen by the network exactly through its component outside it. The largest consistent network of
antitelephones inside $S^{7}$ is therefore an $S^{3}$, and there are three spatial directions.
When the seventh sphere is set aside there as a ruler, which is the octonionic spacetime of the
division-algebra literature, the claim is not that $\mathbb{R}^{8}$ has no metric, nor that the
ten-dimensional Minkowski space of $SL(2,\mathbb{O}) \cong \mathrm{Spin}(9,1)$ fails to exist
\autocite{baez2002octonions}; it is that the octonionic directions are not loops on which
antitelephones compose, and so are algebra rather than distance. In the standard order that
sentence has no content; in this one it is the reason space has three dimensions.

The split just defined explains the most familiar asymmetry in physics: that one can move back and
forth in space but not in time. \S\ref{sec:continuity} identifies the pre-measurement regime with
continuous reversibility --- its transformation group is connected, and every motion in it is
undoable. Spacelike separation, being the non-computable sector, is exactly the part of the
relational structure on which no record has been written; motion across it is motion within the
reversible regime, and it inherits that regime's continuity and its two-sidedness: a place can be
revisited because going and returning writes no record that forbids the return. Timelike separation
is the record chain itself --- proper time counts completed irreversible loops, as above --- and a
moment cannot be revisited because a moment \emph{is} a stratum of records, and returning would
require unwriting them, which is what \textsc{fanout} forbids (\S\ref{sec:fanout-boundary}). The
asymmetry between space and time is therefore not an extra postulate about the manifold: it is the
pre-/post-\textsc{fanout} asymmetry read as geometry --- space is what the reversible group feels
like from inside, time is what the record chain feels like --- and the arrow-of-time discussion
later in this section is the same asymmetry read thermodynamically. The identification is the
framework's structural reading of the split, offered in the same register as the rest of this
subsection.

\subsection{Three dimensions of space}
\label{sec:three-dimensions}

The closed timelike curves that define spacelike distance require information to
be parallel transported consistently around a loop, and consistent parallel
transport on a sphere requires it to be parallelisable. By Adams' theorem the
only parallelisable spheres are $S^{1}$, $S^{3}$ and $S^{7}$
\autocite{adams1960hopf}, the unit spheres of $\mathbb{C}$, $\mathbb{H}$ and
$\mathbb{O}$. Impose now the requirement that information arriving from outside
the causal zone must appear to possess its own causal history---that is, to be the
result of a sequence of information-preserving operations, hence a sequence of
group operations. The spheres $S^{1}$ and $S^{3}$ are groups under their own multiplication --- $U(1)$ and
$SU(2)$ --- and their operations preserve information; $S^{7}$'s multiplication is non-associative,
a Moufang loop rather than a group, so $S^{7}$ cannot serve as a ruler. (The rotation group
$SO(8)$ of the seven-sphere is of course a group; what fails is the sphere's \emph{own}
composition law, the one a ruler would use.) The same non-associativity
forbids a metric built from the sphere's own composition law outright: path
composition $d(A,C)$ through $B$ depends on the bracketing, so the triangle inequality is not merely violated but ill-defined, and
no geodesic structure exists. (What is untouched should be said: the Euclidean metric of
$\mathbb{R}^{8}$, and the ten-dimensional spacetime that $SL(2,\mathbb{O})$ furnishes, are metrics
of the ordinary kind; what fails is a ruler made of antitelephone loops, which is what a spatial
dimension is in this section, as \S\ref{sec:comp-distance} set out. Nor is the failure total:
relays between any two links compose, and among three they compose exactly when the links lie in
one quaternionic subalgebra, so the largest consistent network inside $S^{7}$ is an $S^{3}$
\autocite{conway_2003,harvey1982calibrated}.) Since the abelian $S^{1}$ alone cannot carry
non-abelian structure, spacelike distance must be three-dimensional. Dimension
three is an output, not an input \autocite{small_2006}.

The division of labour is then exact. $S^{1}$ ($\mathbb{C}$) contributes the time
dimension and the $U(1)$ of electromagnetism; $S^{3}$ ($\mathbb{H}$) contributes
the three spatial dimensions and the $SU(2)$ of the weak interaction; $S^{7}$
($\mathbb{O}$) contributes no geometry but, as a filter on what information may be
exchanged, the $SU(3)$ of the strong interaction. A sphere that cannot be a ruler
can still be a filter, which is why the strong force resists geometric
interpretation while electromagnetism and the weak force do not. This is the same
$\mathbb{C}\oplus\mathbb{C}^{3}$ split that yields $G_{\mathrm{SM}}$ in
Section~\ref{sec:entanglement-classes}, now read as the reason spacetime is $(3+1)$-%
dimensional. The conclusion converges with M\"uller and Masanes, who prove from
information-theoretic postulates, assuming no quantum theory, that minimal
direction-carrying systems with continuous reversible dynamics force $d=3$ and
quantum theory on two qubits \autocite{MuellerMasanes2013}, and with the result
that entanglement exists only for a three-dimensional Bloch ball
\autocite{MasanesMuellerEtAl2014}; the framework supplies the reason such systems
exist, Cayley--Dickson doubling from self-reference, while their work supplies
the rigorous dimensional uniqueness. The status of the framework's own route is a
structural result requiring formal proof: the one-bit-per-curve identification and
the precise reduction of the $S^{7}$ filter to $SU(3)$ rather than its full automorphism group
$G_{2}$ are numbered below (Open Problems~\ref{op:comp-metric} and~\ref{op:three-dimensions}).

A scoping remark, and then a second route to the same count. The computational model of spacetime
proposed here has many detailed parallels with the real thing, but a parallel is not an identity,
and the identity is not claimed in this paper: establishing it is work deferred to the second paper.
What can be claimed now is that the parallel is productive in a specific direction --- it suggests a
derivation of the dimension count carried out entirely within general relativity, with no
computational premise. Acceleration and gravitation tip the light cone, and the tipping can be
carried to the extreme: Penrose's discussion runs the cone through a full circle
\autocite[382]{Penrose1994}, at which point the tipped direction closes on itself --- which is to
say, the extreme of tipping is a closed timelike curve, the very object whose closure the
computational definition of \S\ref{sec:comp-distance} counts in bits. The two routes are one route
entered from opposite ends. Now demand of general relativity only this: that observers accelerating
with respect to one another, and with respect to unaccelerated observers, retain a consistent
conception of the distance between them. Each observer's tipping traces a path on a sphere of
directions, and mutual consistency requires the associated rotations to compose without regard to
bracketing --- to form a group. Compositions on $S^{1}$ and $S^{3}$ do, forming $U(1)$ and $SU(2)$ as above; compositions on
$S^{7}$ do not, the multiplication being non-associative, so observers tipped across an
$S^{7}$-worth of directions could not agree on distance even pairwise: the same failure that forbade
$S^{7}$ a metric above forbids it here a shared ruler. The largest sphere admitting a consistent
conception of distance under mutual acceleration is therefore $S^{3}$, and space has at most three
dimensions; with the floor argued above --- the abelian $S^{1}$ alone cannot carry the required
structure --- the count is exactly three. The status matches the subsection's. The tipping and its
extreme are general relativity's own \autocite{Penrose1994}; the consistency demand is
elementary; the associativity filter is the algebraic fact used throughout; the reading of composed
tippings as motion on the parallelisable spheres is the structural step, shared with the
computational route. The convergence of the two routes on the same endpoint is offered as evidence
for the parallel, while the identity it gestures at remains the second paper's burden.

\subsection{Recovering the Lorentz group}
\label{sec:lorentz}

A measurement changes what the observer knows, hence, relationally, the causal
structure: it tips the light cone. This is Penrose's association of cone-tipping
with state reduction \autocite{Penrose1996} run in the opposite causal direction,
reduction driving the tipping rather than the reverse, and the two are then one
operation in two charts---both described by $SL(2,\mathbb{C})$. The algebraic
backbone is the division-algebra correspondence
$SL(2,\mathbb{A})\cong\mathrm{Spin}(1,d{+}1)$ \autocite{baez2002octonions}: $\mathbb{R}$
gives $1{+}1$, $\mathbb{C}$ gives $3{+}1$, $\mathbb{H}$ gives $5{+}1$ and
$\mathbb{O}$ gives $9{+}1$ dimensions. The computability split that singles out a
preferred $\mathbb{C}$ fixes the fundamental state space as $\mathbb{C}^{2}$ and
hence $SL(2,\mathbb{C})=\mathrm{Spin}(1,3)$: the Lorentz group of $3{+}1$
spacetime. The qubit and Minkowski space are not merely both Euclidean-and-three%
-dimensional by coincidence; the Hermitian $2\times2$ matrices are
$\mathbb{R}^{1,3}$, $SL(2,\mathbb{C})$ acts on them as the Lorentz group, and the
qubit density operators fill the future cone---the correspondence M\"uller's group
develops operationally, recovering Lorentz transformations from quantum
communication \autocite{HohnMueller2016} and the complex and quaternionic qubit
from the relativity of simultaneity \autocite{GarnerMuellerDahlsten2017}.
In these coordinates the Bloch ball is the unit-trace slice of the cone and the Bloch sphere its
boundary rays, the celestial sphere on which $SL(2,\mathbb{C})$ acts by M\"obius
transformations; state reduction sends a point of the ball to a point of the sphere, and a boost
sends the sphere to itself, which is what lets the two be one operation in two charts.
The $SL(2,\mathbb{A})$ correspondence is rigorous; its identification with
measurement-as-tipping is the framework's structural reading.

One consequence of the identification has been implied throughout this section and deserves its own
paragraph, because the next paragraph turns on it. If measurement creates a relation between
observer and observed, and if --- relationally --- a new relation alters the causal structure, then
the tipping of the cone is not merely \emph{associated} with the creation of the relation: it is
what the creation of the relation is, seen geometrically. And a tipping of one's cone against the
congruence of one's neighbours already has a name in physics: acceleration. The framework's
statement, made explicit here, is that \emph{creating a new relation between observer and observed
is experienced as acceleration}. The resonance with the equivalence principle is exactly what it
should be: the principle makes acceleration and gravitation locally indistinguishable, Penrose's
association --- run in this section's direction, reduction driving the tipping --- makes state
reduction and cone-tipping one operation, and the framework's identification closes the triangle,
measurement and acceleration meeting where the principle says they must. The claim is structural:
the framework's reading of the cone-tipping established above, not an independent derivation.

Taken literally, the identification has a measurable face, and it is a familiar one: the Unruh
effect, read in this section's own ledger. Recall the definitions of \S\ref{sec:comp-distance}:
timelike distance counts computable steps, spacelike distance counts non-computable bits. What an
observer can compute --- what their \textsc{fanout}s can convert into classical records --- is fixed
by their causal structure, and the previous paragraph has just made that structure
relation-dependent. An accelerated observer's cones are tipped, and with them the boundary of the
recordable: relations that for the inertial observer lie wholly in the reversible, pre-record sector
--- the entanglement structure of the vacuum, non-computable for them, spacelike in the
computational sense --- fall, for the accelerated observer, partly on the recordable side. The mechanism is the one
Chaitin identified for formal systems \autocite{Chaitin1998}: a proposition undecidable relative to
an axiom system becomes decidable when information is added to the axioms --- and
\S\ref{sec:comp-distance} defined spacelike distance as exactly the bits whose addition closes the
gap. Creating a relation is adding information to the observer's axiom base, so timelike and
spacelike distance are converted into one another by the addition and subtraction of information,
and acceleration is the continuous version of that addition. A
relation newly recordable is a detection. Particles appear from nowhere: an accelerated detector in
the inertial vacuum registers a thermal bath at $T = \hbar a / 2\pi c k_{\mathrm{B}}$
\autocite{unruh1976notes}. Two guards keep the reading honest. Nothing here reclassifies the
invariant interval --- timelike and spacelike separation in the emergent metric are Lorentz
invariants, and no event crosses that classification; what is observer-relative is the computational
partition, which relations a given observer's horizon renders recordable. In the ledger of
\S\ref{sec:comp-distance} that partition \emph{is} the timelike/spacelike distinction, so in the
framework's own units timelike distance is \emph{rotated} into spacelike, and spacelike into
timelike, by a change of observer --- the tipping of the cone is literally that rotation --- and the
detector's clicks are the rotation made audible. And the effect is established physics
while the reading is the framework's, flagged structural. What gives the reading its content is a
parallel: the inertial and the accelerated observer stand to the vacuum exactly as Wigner and his
friend stand to the laboratory (\S\ref{sec:wigners-friend}) --- one holds a coherent relation and no
event, the other holds a record, and the facts are relative on both sides in precisely Rovelli's
sense \autocite{DiBiagioRovelli2021}. Since the previous paragraph identifies acceleration with
relation-creation, the parallel is not an analogy but the same structure in gravitational dress, the
horizon standing where the laboratory wall stood: \emph{the Unruh effect is Wigner's friend for
accelerated observers}. And two subsections of this section that seemed independent here meet: the
temperature Jacobson's construction requires across every local Rindler horizon
(\S\ref{sec:jacobson}) is exactly this temperature --- the $T$ of the Clausius relation is the $T$
of the tipped cone, supplied by the very relation-creation with which this subsection began.

\subsection{The relationalist commitment}
\label{sec:relationalism}

General relativity is relational: by the hole argument, spacetime points carry no
intrinsic identity, and the metric is fixed by the network of coincidences and
interactions, not by a pre-existing stage \autocite{Rovelli1996a}. The
computational picture is relational by construction and explains why it must be:
a background metric would require spatial distances to be pre-computable, and
non-computability forbids exactly that. General covariance is the statement that
distances have no pre-assigned values and depend on what information is actually
present. The metric is therefore dynamical not by stipulation but because it
measures a quantity that cannot be fixed in advance.

\subsection{Jacobson, holography, and imaginary information}
\label{sec:jacobson}

Jacobson derives the Einstein equation as an equation of state, from the Clausius
relation $\delta Q=T\,\mathrm{d}S$ across every local Rindler horizon together
with the Bekenstein proportionality of entropy to horizon area
\autocite{Jacobson1995, Bekenstein1973}. In the computational reading this is the
self-consistency condition on the loop-counting network: the Einstein equation is
the requirement that the bookkeeping of computation be internally consistent
across surfaces. The holographic principle, that the information in a region is
bounded by its bounding area \autocite{tHooft1993, Susskind1995}, then acquires a
natural reading here. A surface area is a squared distance; if spacelike distance
is measured in imaginary units of information, an area is measured in the square of
those units, that is in \emph{negative} information---entropy. The holographic
link between area and entropy thus requires the square root of negative
information, and that square root is precisely the imaginary unit of the complex
amplitude, the observer's self-referential phase. Read this way, Jacobson's
account of information crossing surfaces implicitly assigns an imaginary
information value to spacelike distance, and the $\mathrm{i}$ of the Minkowski
signature is the same $\mathrm{i}$ as the $\mathrm{i}$ of the quantum amplitude.
This identification is the most speculative step in the section: it is
structurally suggestive and it makes the holographic form unsurprising, but it is
a conjectural reading rather than a derivation, and is offered as such.

The conjecture has a two-way motivation that this paper's record should carry, because each
direction was on the table before the other was noticed. Forward: if spacelike distance is counted
in imaginary units of non-computable information, then any spacelike area is the product of two
imaginary lengths and is therefore \emph{negative} information --- which is to say entropy, by
Brillouin's identification \autocite{brillouin1956science} --- and a holographic principle arrives
automatically, the sign supplied by $i^{2}$ rather than by interpretation. Converse, and more
striking: in the holographic principle as ordinary physics already states it, surface area is
proportional to entropy, and entropy is negative information; so by implication ordinary physics
has been measuring spacelike separation in imaginary units of information all along, and the
convention of \S\ref{sec:comp-distance} is read off from established results rather than
invented. The same imaginary unit is already employed once in this section --- it is what makes
the signature of the interval Minkowski --- so the sign of the area law is its second job, not a
second assumption. What the two-way motivates it does not prove: the formal statement is exactly
the metric-and-signature half of Open Problem~\ref{op:comp-metric}, together with the
one-bit-per-bit-pair coefficient owed to the second paper. The stakes of the conjecture have
recently risen: Dorau and Much prove, with modular-theoretic rigour, that the semiclassical
Einstein equations follow from quantum relative entropy on a local horizon \emph{once the
entropy--area proportionality is assumed} \autocite{doraumuch2026relative} --- so what this
framework proposes to derive is precisely the remaining input of a derivation already carried
through at that standard.
\input{entropic_gravity_family_paragraph_2026-07-03}

\subsection{Non-locality already present in general relativity}
\label{sec:gr-nonlocality}
\input{gr_nonlocality_revision_2026-06-13}
\medskip
\noindent  The spacetime arc therefore yields, from one definition of spacelike
distance as non-computability: the signature, three spatial dimensions,
the Lorentz group, the dynamical relational metric, a reading of the
holographic principle, the presence of non-locality within general
relativity, the identification of relation-creation with acceleration
and of the Unruh effect with Wigner's friend, and the space--time
mobility asymmetry as the \textsc{fanout} asymmetry read as geometry. Of these,
the $(3+1)$ count and the Lorentz recovery are structural results that converge with independent
rigorous work and await their formal proof --- the count now numbered as
Open Problem~\ref{op:three-dimensions}; the metric-space status, the signature and the
one-bit-per-curve measure (together Open Problem~\ref{op:comp-metric}), the imaginary-information
reading of holography, and the general-relativistic non-locality, the acceleration and
Unruh readings, and the mobility asymmetry are conjectural, and are named here precisely so that the boundary between what is shown and what
is proposed is visible.

\subsection{Why the metric is not to be quantised}\label{sec:no-quantised-gravity}

A consequence has been implicit in everything above and deserves its plain statement. Causal
structure, on this framework's reading, is the bookkeeping of locked relations --- post-\textsc{fanout}
by construction --- while superposition is the signature of the pre-\textsc{fanout} regime. A
``superposition of causal structures'' is therefore not a difficult state but a category error: it
applies the pre-record operation to a record-side object, exactly as ``a superposition of
measurement outcomes'' does. Gravity lives in the post-\textsc{fanout} regime. The quantum layer
beneath it exists and is not in dispute; but it is \emph{pre-geometric} --- it carries amplitudes
over the formation of relations, not amplitudes over formed geometries --- so there is no
fundamental metric to quantise, and the standard programme of superposing spacetimes pursues an
object the ontology does not contain. One guard rail keeps the claim honest: collective
excitations of the locking pattern may perfectly well exist as effective quanta --- gravitons as
the phonons of the relational network --- and the thesis forbids fundamental metric superposition,
not effective field theory; this is Jacobson's equation-of-state moral \autocite{Jacobson1995},
made structural. The relation to gravitational-collapse programmes is an inversion: where
superposed geometries are there read as unstable \autocite{Penrose1996}, here they are read as
unformed --- what decays in those models is what locks in this one.

Three experimental postures follow, each at the status the reading carries. Laboratory
demonstrations of indefinite causal order \autocite{procopio2015experimental, rubino2017indefinite} are, on this
account, demonstrations of the nascent layer: the indefiniteness lives below the locking threshold
and definitises at it, so the framework adopts the programme and adds a boundary --- indefinite
causal order can never be amplified into records while remaining indefinite. Gravitationally
induced entanglement of the kind the Bose--Marletto--Vedral proposals seek
\autocite{bose2017spin, marlettovedral2017gravitationally} is \emph{predicted}, because the
pre-geometric layer is quantum --- while witnessing that layer, not a superposed metric: the
framework splits the experiment's usual dichotomy. And gravitationally driven anomalous collapse
is denied: locking is redundancy and thermodynamics, so decoherence experiments should find
standard rates and nothing gravity-specific. That no consistent quantisation of gravity exists was
the conclusion of the author's CASYS paper by a different route \autocite{small_2006}; the present
framework supplies the mechanism the conclusion then lacked.

\subsection{Open problems and prospects}
\label{sec:spacetime-open-problems}

The spacetime arc is the framework's frontier, and honesty about its edges
matters more here than anywhere else in the paper. We set out the genuine
problems and the prospects together, deferring the deepest to a second paper
rather than papering over them. Two of the arc's debts are well-posed now, and we number them.

\begin{openproblem}[The computational metric]\label{op:comp-metric}
Prove that the paired distance of \S\ref{sec:comp-distance} --- computable steps counted real,
non-computable bits counted imaginary, one bit per closed curve --- defines a genuine
pseudo-metric on the space of computations, with the Minkowski signature derived rather than read
off; equivalently, establish the one-bit-per-curve measure as a measure, and the signature as its
theorem. Until then the identification stands as motivated by reproducing what nature does, which
is the status flagged at its introduction.
\end{openproblem}

\begin{openproblem}[Dimension three as a theorem]\label{op:three-dimensions}
Turn the two convergent routes of \S\ref{sec:three-dimensions} into a proof: formalise the
parallelisability-plus-associativity filter (the computational route) and the mutual-acceleration
consistency demand (the general-relativistic route) so that $d = 3$ follows as a theorem, and
include the reduction clause --- that the $S^{7}$ filter acts through $SU(3)$, the stabiliser of
the selected complex direction, rather than through its full automorphism group $G_{2}$. The
M\"uller--Masanes uniqueness results \autocite{MuellerMasanes2013, MasanesMuellerEtAl2014} are
the rigorous convergents any proof must meet.
\end{openproblem}

\paragraph{The substrate and the generation of events (deferred).}
A computation runs on something; a network of events presupposes that events
occur. Both are the same objection in two forms---the network's hardware and its
clock---and the original CASYS'05 statement of the programme raised the
implementation regress and set it aside \autocite{small_2006}. The framework's
response is that a self-referential structure is the only kind that need not
presuppose a meta-substrate: it generates its own. The natural algebraic emblem
is $E_{8}$, the unique simple Lie group whose adjoint representation is also its
smallest faithful one, so that what it is and what it acts on coincide; the $240$
unit integral octonions that generate it are exactly the $E_{8}$ root system, so
the self-action is welded to the division-algebra spine of
Section~\ref{sec:entanglement-classes}. This is Wheeler's self-excited circuit
\autocite{Wheeler1990} given an algebraic body, and it converges with Furey's use
of an algebra acting on itself to generate particle content
\autocite{furey2018three}. The decisive caveat, which we state rather than hide, is
that the bootstrap is self-consistent but not self-specifying: it does not
eliminate the non-computability identified throughout, it relocates it to a
determined place. The substrate question is, on the framework's own terms,
provably without a computable answer from inside---it is the irreducible residue,
the bump in the carpet that the whole construction has been moving rather than
removing. The leap from ``$E_{8}$ is self-dual'' to ``$E_{8}$ is the substrate''
is at present an emblem awaiting operational content (an $E_{8}$-equivariant
dynamics, or a state space built as a homogeneous space of $E_{8}$), and that
content, together with how a tick is generated at all, is deferred to the second
paper.

\paragraph{Dark matter (deferred).}
If gravity is the bookkeeping of the interaction network and dark matter interacts only
gravitationally, how does it participate in a network it does not exchange information with?
This is a genuine boundary of the programme, at which it may make no prediction, and Open
Problem~\ref{op:sterile-gravity} states it with the criteria that would close it. The prospect, to be developed in the second paper, is
that this is a feature rather than a defect. In the division of labour of
Section~\ref{sec:three-dimensions}, participation in the metric---counting loops
in the $S^{1}$ and $S^{3}$ directions---is universal, carried by anything with a
worldline and a tick-rate, whereas participation in the gauge filters, the
$S^{3}$ of the weak interaction and the $S^{7}$ of colour, is selective, engaged
only by charged matter. A node that counts loops but engages no filter
\emph{is} matter that gravitates without otherwise interacting, and the
framework's Class-4 a-chiral sterile neutrino (\S\ref{sec:generations}) is the candidate that
realises it --- recorded in the particle sections as Open Problem~\ref{op:sterile-gravity}, whose
tension (a gauge-sterile channel gravitating, when gravity is measurement bookkeeping) is exactly
the question named here. The condition for the candidacy to be coherent rather than circular is
that metric participation be established as prior to, and independent of, gauge participation;
making that priority precise is the work deferred, and it is the same work the open problem
requests.

\paragraph{CP violation, the arrow of time, and baryogenesis.}
CP violation appears to demand a time asymmetry owing nothing to the entropic
arrow, which would seem to require a preferred direction built into spacetime
itself. The framework dissolves the appearance rather than accommodating it. CP
violation needs an irreducible complex phase, the CKM phase $\delta$; by Spekkens'
result that negativity and contextuality are the same notion of nonclassicality
\autocite{spekkens2007toy}, a complex phase is a signature of contextuality; so CP
violation should be contextual, sensitive to the near environment of the decaying system --- the
same phase the interaction sections read as anchored by a fourth party
(\S\ref{sec:interactions}) --- which is exactly what Andrianov, Taron and Tarrach established for
neutral kaons in matter, where the environment breaks C, irreversibility breaks
T, and the effective parameters follow a Lindblad equation
\autocite{Andrianov2001}. The framework's reading is that the resulting
informational asymmetry \emph{constitutes} the arrow rather than presupposing it:
CP violation, entropy increase, and the direction of time are one informational
asymmetry seen three ways, not an independent second arrow, so the worry about a
direction built into spacetime does not arise. On this footing the contextual
dependence sets up a positive feedback loop for baryogenesis---an initial
fluctuation makes the environment slightly matter-dominated, contextual
enhancement raises the effective CP violation, more matter is produced, and the
loop runs---with $\delta$ a cosmological fixed point rather than a fundamental
constant. The honest limit of this is that a feedback loop explains amplification
and lock-in, not the sign selection or the baryon-number violation that
Sakharov's conditions still require \autocite{Sakharov1967}; the stronger claim
the framework is positioned to make, that the matter direction and the time
direction are co-selected by the same self-referential asymmetry, needs a source
of baryon-number violation and is left open.

\paragraph{Spatial infinitude and the status of inflation.}
There is no upper bound on non-computability, so the space of computations is
unbounded: the universe is spatially infinite, which leaves its spatial sections flat or of
negative curvature and excludes only the closed case. Infinitude on its own, however, does not
settle flatness, which is a fine-tuning of the Friedmann dynamics the framework has not yet
derived (Open Problem~\ref{op:comp-metric}); and flatness was in any case never inflation's
principal task. The horizon problem is, and an infinite space does nothing for it. The framework's actual lever against
inflation is the relational thesis of Section~\ref{sec:gr-nonlocality}: if
pre-spatial entanglement correlates regions before \textsc{fanout} spreads them
into spatial separation, then the observed large-scale homogeneity is fixed
relationally, before any question of causal contact, and the horizon problem
dissolves without an inflaton. This is a real argument, but it carries a real
bill: inflation also generates the nearly scale-invariant spectrum of primordial
fluctuations seen in the microwave background \autocite{Planck2020}, and any
account that dispenses with inflation owes an alternative origin for that
spectrum. That, not the horizon problem, is the hard empirical debt, and it is
named here rather than discharged.

\paragraph{Expansion and acceleration.}
If spacelike distance is the information separating worldlines, then as worldlines
accumulate distinct interaction histories they accumulate information-distance,
and the universe expands. The picture has the right exclusions built in: bound
systems interact continuously, share information, and so keep low
information-distance---they do not expand---while only worldlines that decouple
diverge, recovering expansion on cosmological scales alone. The feedback in which
greater separation means less light-cone overlap and so faster divergence,
$\mathrm{d}d/\mathrm{d}t \propto d$, yields exponential divergence, a de~Sitter
phase, and hence accelerating expansion without a separate dark-energy term
\autocite{Riess1998}. The limit to be honest about is that pure exponential
divergence is the late-time behaviour only; reproducing the full expansion history,
with its radiation- and matter-dominated eras, and fixing the Hubble rate
quantitatively, remains to be done.

\paragraph{The $S^{7}$ sector as internal structure.}
That $S^{7}$ cannot carry a metric does not make it disappear. Its closed
timelike curves carry the same non-computable information structure as those of
$S^{1}$ and $S^{3}$, but lacking associativity they support algebra---gauge
charge---in place of geometry---distance. This is why gauge symmetry has so often
looked like a set of compactified extra dimensions: the Kaluza--Klein and string
intuition is tracking something real in the $S^{7}$ structure, but misdiagnoses
it as spatial. In the framework these degrees of freedom have no size because
they have no metric; they are internal in the exact sense that they enter the
information-exchange that constitutes gauge charge while contributing nothing to
the loop-counting that constitutes distance. The resemblance to rolled-up
dimensions is therefore genuine and instructive, and the diagnosis that they are
algebraic rather than geometric is what separates this account from theories that
add real spatial dimensions to carry the same structure.

\paragraph{Loops, supertasks, and a second route to parallelisability.}
The argument of \S\ref{sec:three-dimensions} derives the parallelisable-sphere constraint from
closed timelike curves: a computation whose output is its own input requires information to be
transported consistently around a loop, and consistent transport selects the parallelisable spheres
$S^{1}$, $S^{3}$, $S^{7}$ \autocite{adams1960hopf}. There is a suggestion, which we record as a
prospect rather than a result, that the same constraint may govern the loops of perturbative quantum
field theory. The external legs of a Feynman diagram are on-shell classical boundary data; the
quantum content, and the divergences, enter with the internal loops, where the loop momentum is
integrated over all values and the action accumulated inside the loop is in this sense unbounded ---
a supertask, an infinite computation closed into a finite process. If such a loop is, in the
framework's terms, a closed computational path of the same kind as a closed timelike curve --- a
cycle carrying information not derivable from the data outside it --- then the parallelisability
constraint derived for closed timelike curves would apply to it, and the admissible interactions
would be constrained by the same $S^{1}$, $S^{3}$, $S^{7}$ structure that fixes the gauge group in
\S\ref{sec:entanglement-classes}. This would be a second and independent route to the Standard Model
content: not only \enquote{closed paths in spacetime select the parallelisable spheres} but \enquote{closed paths in interactions select them too}, unifying the spacetime arc and the interaction arc under one
principle. The decisive step, which we have not taken, is the identification of a field-theory loop
with a closed computational path in the precise sense the framework requires --- a loop integral is a
closed cycle in momentum space, a closed timelike curve a closed cycle in spacetime, and whether
these are the same kind of self-consistent closed structure is exactly what would have to be
established. It is offered as a direction, not a claim.

\paragraph{Does parallelisability determine renormalisability?}
A sharper and more self-contained question sits alongside the previous one, and we state it in the
hope that it will be taken up by those equipped for it. Renormalisability is the property that the
divergences of a quantum field theory can be consistently absorbed, leaving finite predictions; some
theories have it and some do not, and the standard criterion for which is technical (power-counting
of the coupling's dimension). The framework suggests a more structural reading. Parallelisable
spheres are precisely the spheres on which a global, consistent transport exists where it generically
does not; the non-parallelisable spheres are where the obstruction lives. If the divergence of a loop
is, in the framework's language, an obstruction to closing a computational path consistently, then
parallelisability would be exactly the structure that distinguishes the removable obstruction (a
renormalisable theory) from the irremovable one. It is at least a suggestive coincidence that the
renormalisable interactions of nature are the gauge theories built on $U(1)$, $SU(2)$ and $SU(3)$ ---
the groups the framework ties to the three parallelisable spheres --- and that gravity, which resists
renormalisation, is the sector \S\ref{sec:three-dimensions} ties to the geometry rather than the
algebra. Whether this is structural or accidental we do not know. The question --- whether
parallelisability, in the Hurwitz--Adams sense, is what determines renormalisability --- is posed in
the native language of the division-algebra approach to the Standard Model, and is one we think that
community is best placed to answer. We are careful to claim nothing more than that the question is
well-motivated: this is emphatically not a method for computing renormalised quantities, whose finite
remainders carry the precision agreement with experiment, but a conjecture about which theories admit
such a computation at all. It is, like the previous prospect, a direction with a possible path to an
answer rather than an answer, and it is left open deliberately.


\medskip
\noindent These mark the boundary of the programme as it stands. The substrate,
the generation of events, and the network status of dark matter are deferred to
the second paper; the contextual account of CP violation and baryogenesis, the
relational case against inflation, the information-divergence account of
expansion, and the algebraic reading of the $S^{7}$ sector are offered as
prospects, each with its outstanding debt named so that the line between argument
and aspiration stays visible.
%      % EDITED 2026-09-09 scripts/patch_spekkens_citekey_2026-09-09T1810.py
% =====================================================================
%  SECTION: Why Hopf Fibrations Belong in Particle Physics
%  (sec:hopf-justification). Drafted 2026-06-08; F1 arc added
%  prop:f1-dictionary and op:hopf-frobenius's discharge clause.
%  AUDITED 2026-07-06 (final pass, eighth stop).
%
%  PLACEMENT (reordered chain, JS 2026-07-06): position 8 --
%      \input{spacetime_from_non-computability}
%      \input{hopf_justification_section_2026-06-08}
%      \input{entanglement_classes_and_gauge_groups_2026-06-08}
%  No adjacency phrasing in the prose: position-portable. All
%  cross-refs live -- sec:fanout-direction DRAFTED 2026-07-06 (formerly the one undrafted
%  section, by design).
%
%  AUDIT RECORD: prop:gsm, prop:termination and prop:f1-dictionary
%  verified against every quoting consumer (statements canonical);
%  honest-debt subsection names both the selection assumption and the
%  mereology debt; the epistemic two-climbs disambiguation and the
%  compositional rung-evidence verified; sec:canonical-orbits live;
%  file pure ASCII. Fixes: op:hopf-frobenius's cross-link with
%  op:sigma-functor made TWO-WAY (the "four ways" line now names the
%  op); both remaining functor conditions ANCHORED --- the quaternionic
%  to the established fibre physics of sec:wf-reconciliation (residue:
%  clutching map vs strong complementarity), the octonionic to the
%  proved canonical-orbit images (umbrella label app:ghz-w-proofs
%  minted at the appendix for the purpose). Statuses unchanged.
%  British English; \autocite; no bullets in prose.
% =====================================================================

\section{Why Hopf Fibrations Belong in Particle Physics}
\label{sec:hopf-justification}

Invoking Hopf fibrations in a discussion of the Standard Model invites an immediate and fair
objection: that the move is decorative, an analogy between two pieces of attractive mathematics
with no load-bearing connection between them. The purpose of this section is to dispose of that
objection before the construction is put to work. The connection is not an analogy. It is
carried by two independent and independently published theorems, and the structures it relies
upon are not selected for convenience but compelled by a pair of classification results. We take
as already established the correspondence between qubit entanglement and the Hopf fibrations
themselves --- the subject of a mature literature
\autocite{mosseri2001hopf,bernevigchen2003geometry} --- and do not reproduce it here. What this
section addresses is the prior question of why that correspondence is entitled to carry physical
content at all.

\subsection{Two theorems, joined end to end}

The bridge from qubit entanglement to the Standard Model is built from two distinct results,
neither of which is original to the present framework and neither of which depends on the other.

The first is geometric. The unit sphere in the Hilbert space of $n$ qubits is the total space of
a Hopf fibration: $S^3 \to S^2$ for one qubit, $S^7 \to S^4$ for two, and $S^{15} \to S^8$ for
three. For two and three qubits the fibration is \emph{entanglement-sensitive} --- separability
is precisely the condition under which the base degenerates to a product of single-qubit Bloch
spheres and the additional dimensions associated with the higher imaginary directions collapse
\autocite{mosseri2001hopf,bernevigchen2003geometry}. This is a statement about the geometry of
state spaces and requires no physical interpretation whatever.

The second is algebraic. Singling out a preferred complex structure on the octonions ---
equivalently, fixing the split $\mathbb{O} \cong \mathbb{C} \oplus \mathbb{C}^3$ --- leaves a
residual symmetry, and that residual symmetry is exactly the gauge group of the Standard Model.
This was established through the exceptional Jordan algebra by Dubois-Violette and Todorov
\autocite{duboisviolette2016exceptional,duboisviolettetodorov2019exceptional,todorovduboisviolette2018deducing}
and given its sharpest spinorial form by Krasnov \autocite{krasnov2021so9}, who showed that the
relevant group is the stabiliser of a complex structure inside $\mathrm{Spin}(9)$. We record the
result in the form in which it enters the present work.

\begin{proposition}[Dubois-Violette--Todorov; Krasnov]\label{prop:gsm} Let $J_R$ be the complex
structure on the space $\mathbb{O}^2$ of $\mathrm{Spin}(9)$ spinors determined by a choice of
unit imaginary octonion, equivalently the split
$\mathbb{O} \cong \mathbb{C} \oplus \mathbb{C}^3$. The subgroup of $\mathrm{Spin}(9)$ commuting
with $J_R$ is \[ G_{\mathrm{SM}} \;=\; \frac{SU(3) \times SU(2) \times U(1)}{\mathbb{Z}_6}, \]
acting on the singled-out $\mathbb{C}$ summand as a colour singlet of hypercharge $-1$ and on
the $\mathbb{C}^3$ summand as a colour triplet of hypercharge $\tfrac{1}{3}$, both as $SU(2)$
doublets --- the quantum numbers of one left-handed generation.
\end{proposition}

The two results are independent, and the join between them is the only place a physical theory
must do work. Each result is a theorem; what neither supplies is a reason why nature should fix
one complex direction within $\mathbb{O}$ rather than treating all directions symmetrically. We
return to this --- the genuine load-bearing assumption --- in \S\ref{sec:hopf-honest-debt}, and
resolve it, in candidate form, in \S\ref{sec:fanout-direction}.

\subsection{The ladder is forced, not chosen}

It might still be objected that selecting the octonions, or the Hopf fibrations, is itself a
modelling choice carrying as much latitude as the choice of a grand-unified group. It is not,
and the reason is a pair of classification theorems that close off the alternatives entirely.

By Hurwitz's theorem the only normed division algebras are $\mathbb{R}$, $\mathbb{C}$,
$\mathbb{H}$ and $\mathbb{O}$ \autocite{hurwitz1898composition,baez2002octonions}. By Adams's
resolution of the Hopf-invariant-one problem --- together with the parallelisability results of
Bott--Milnor and Kervaire --- the only parallelisable spheres are $S^0$, $S^1$, $S^3$ and $S^7$,
so that there exist exactly three nontrivial Hopf fibrations, in bijection with $\mathbb{C}$,
$\mathbb{H}$ and $\mathbb{O}$
\autocite{adams1960hopf,bottmilnor1958parallelizability,kervaire1958nonparallelizability}. These
are not heuristics; they are among the more durable theorems in topology and algebra. Their
force here is that, once a framework is committed to amplitudes valued in a normed division
algebra --- a commitment that in the present construction is made far upstream, by the
self-referential argument of \S\ref{sec:self-reference} and \S\ref{sec:division-algebras}, and
not by fiat at this stage --- the ladder of Hopf fibrations is not selected but compelled, and
it terminates of its own accord.

\begin{proposition}[Termination]\label{prop:termination} The Cayley--Dickson tower of normed
division algebras and the tower of parallelisable spheres both terminate at real dimension
eight. Consequently there are exactly three nontrivial Hopf fibrations, the largest being the
octonionic $S^{15} \to S^8$, and the construction admits no extension to a fourth case.
\end{proposition}

The termination is worth dwelling on, because it is unusual for a model of this kind to predict
its own boundary. The next rung of the Cayley--Dickson sequence, the sedenions, possesses zero
divisors; the Hopf map can no longer be globally defined; there is no fourth fibration. This is
not a limitation imposed from outside but a structural fact, and it carries a physical
consequence that we develop later: there is no room in the construction for a fourth gauge force
or a fifth interaction, so that the discovery of either would count as evidence against the
framework rather than as an invitation to extend it. The endpoint is a theorem, not a
stipulation.

\subsection{Independent arrival as corroboration}

The most persuasive evidence that the destination is not an artefact of the route is that it is
reached by routes sharing none of the same machinery. Szangolies
\autocite{szangolies2025standardmodel} reaches $G_{\mathrm{SM}}$ from the emergent-gravity and
Kaluza--Klein programmes, reading the singled-out complex direction as a dimensional reduction
of a $9+1$-dimensional spacetime; the present framework reaches it from self-reference
(\S\ref{sec:self-reference}); and Furey and Gillard--Gresnigt reach the same gauge group, and
the same three-generation structure, through the minimal left ideals of the Dixon and Clifford
algebras, without invoking Hopf fibrations at all
\autocite{furey2018three,gillardgresnigt2019three}; the relation between their minimal ideals
and the entanglement classes used here is taken up in \S\ref{sec:canonical-orbits}. The
octonionic route to particle physics is in any case an old one, going back to the use of
$\mathbb{O}$ to model the strong force \autocite{gunaydingursey1974quark}.

That programmes with disjoint starting points converge on
$SU(3) \times SU(2) \times U(1)/\mathbb{Z}_6$ and on the same quark--lepton division is the
coincidence Szangolies himself remarks upon: the characteristic structure of the Standard Model,
routinely regarded as arbitrary, turns up independently in unrelated settings, and its appearing
twice would then be doubly odd unless the two settings are in fact related. We read the
convergence the same way --- as evidence that the division-algebraic description is singled out
by the mathematics rather than imposed by the modeller, and therefore that
Proposition~\ref{prop:gsm} is structural rather than contrived.

\subsection{The honest debt}
\label{sec:hopf-honest-debt}

It would be misleading to leave the impression that this route is free of assumptions, so we
state plainly the one assumption that all of these programmes share. Proposition~\ref{prop:gsm}
requires that \emph{some} complex direction in $\mathbb{O}$ be singled out, and the uniqueness
of that choice is not supplied by the theorem. Krasnov parametrises the complex structure $J_R$
by an arbitrary unit imaginary octonion, and the residual group is the same for every such
choice because $G_2 = \mathrm{Aut}(\mathbb{O})$ acts transitively on the unit imaginary
octonions. The theorem therefore guarantees the gauge group \emph{given} that a direction is
fixed, while remaining silent on \emph{why} one is fixed. Szangolies treats this as a free
choice and flags it, in his own conclusions, as an open element of the construction --- one to
which several inequivalent answers are available.

This, and not the appeal to Hopf fibrations, is the genuine load-bearing assumption of the whole
approach. The Hopf fibrations are forced (Proposition~\ref{prop:termination}); the singling-out
is not. The present framework offers a candidate resolution that the prior literature lacks: the
act of measurement, modelled as \textsc{fanout}, selects the direction physically, and the
transitivity of $G_2$ then guarantees that the gauge group is independent of which direction the
measurement happens to select. This argument is given in \S\ref{sec:fanout-direction}; we flag
here, at the appropriate epistemic status, that it is structural rather than fully proved, and
that the residual question --- why a measurement selects \emph{exactly one} imaginary direction
rather than merely \emph{permitting} a selection --- remains open. It is now priced as Open
Problem~\ref{op:fanout-selection}.

A second debt is shared with every programme of this character and is not closed by anything
said here: why the abstract quantum state should admit a tensor-product decomposition into
qubits in the first place. This is the quantum factorisation, or mereology, problem. The
entanglement-sensitivity on which the construction turns is meaningful only once such a
decomposition is fixed, and we treat the decomposition as a necessary input rather than a
derived result. Naming it is preferable to letting it pass unremarked.

\subsection{The termination read epistemically}
\label{sec:hopf-information-bound}

The termination of the tower (Proposition~\ref{prop:termination}) has so far been stated
algebraically: the Cayley--Dickson sequence acquires zero divisors at the sedenions and the Hopf
map can no longer be globally defined. There is a second reading of the same boundary, epistemic
rather than algebraic, which we offer as a motivated conjecture rather than a settled result,
because it connects the tower's endpoint to a bound that arises independently in the foundations
of measurement.

A particle, in this framework, is a reified relation between events, and its internal content is
carried by the octonionic $A,B,C$ register --- the $S^{7}$ sector serving as internal degrees of
freedom in the channel that connects the events (\S\ref{sec:space-of-computations}). The
termination of the tower is then a statement about that content: a single such relation can
support genuine entanglement among at most three parties, because there is no fourth fibration
to carry a fourth independent party within the same irreducible relation. Every particle is, on
this reading, a \emph{completed} octonionic-level relation --- topped out at three parties not
by circumstance but because the tower admits no further rung. The information that can be
encoded in one relation-between-events, and hence extracted by one measurement of it, is
correspondingly bounded: it is the information of a completed three-party register, and no
measurement of a single particle can yield more.

This is exactly the kind of bound that the epistemic-restriction programmes reach from the other
direction. Spekkens' toy theory is built on a knowledge-balance principle that caps the
information obtainable about a system \autocite{spekkens2007toy}, and Szangolies derives such an
epistemic horizon not by positing it but from the same self-referential non-computability on
which the present tower rests \autocite{szangolies_j_2018}. That the two accounts --- the
algebraic termination of the Hopf tower and the epistemic bound on measurement information ---
should coincide is, we suggest, not a numerical accident but a shared root: both are
consequences of the self-referential non-computability of \S\ref{sec:self-reference}, wearing an
algebraic face in the one case and an epistemic face in the other. The Hurwitz--Adams ceiling
and the knowledge-balance ceiling would then be one limit seen twice, and Szangolies' programme,
which reaches $G_{\mathrm{SM}}$ by an independent route (\S\ref{sec:hopf-honest-debt}) and the
epistemic horizon from the same non-computable source, shares with this framework not merely a
conclusion but its origin.

The reading must be kept distinct from a different climb, with which it would otherwise be
conflated. The Wigner's-friend regress of \S\ref{sec:wigners-friend} ascends without end: each
observer may be measured by a further observer, each carrying a fresh gauge-private circle, and
that chain has no top --- it is the bump in the carpet, unbounded in length. The termination
bounds something else: not the length of the chain of observers, but the internal structure of
each single relation, each particle, that the chain is built from. The regress is unbounded
precisely because each of its links is a bounded, completed three-party object;
termination-of-each-link and non-termination-of-the-chain are not in tension but complementary,
the former individuating the links of the latter. What the bound caps is the information in one
particle; what the regress leaves uncapped is the number of particles and observers that can be
chained. The precise identification of ``the information of a completed three-party register''
with the formal measures of the epistemic-restriction literature is the technical step that
would turn this conjecture into a theorem, and we flag it as owed; the structural claim --- that
the tower's endpoint and the measurement bound share the non-computable root --- is what we
advance here.

\subsection{The termination read compositionally}
\label{sec:hopf-compositional}

There is a third termination theorem, proved in a community that has so far had no contact with
the two joined above. Coecke and Kissinger showed that the GHZ and W classes --- realised as the
commutative special and the anti-special Frobenius algebras of categorical quantum mechanics ---
are \emph{compositionally complete}: every $N$-qubit entanglement class under SLOCC is generated
by composing these tripartite primitives with bipartite entanglement and product states
\autocite{coeckekissinger2010compositional}. Above three parties nothing new appears; the
entanglement alphabet closes at exactly the point where the fibration tower terminates. The
mutual silence of the two literatures is itself striking: Coecke and Kissinger make no reference
to Hopf fibrations or division algebras, and the fibration literature on entanglement
\autocite{mosseri2001hopf,bernevigchen2003geometry} nowhere cites compositional completeness.
Two classifications, in different mathematics, draw the same line in the same place, unaware of
one another. That silence has, in the months this paper was in draft, just begun to break from
the categorical side, with a projective, base-level reading of the same calculus
\autocite{burtonanwar2026meromorphic}; we engage it below.

The framework's conjecture is that this is one termination wearing a third face: the alphabet
has the letters it has --- Bell, GHZ, W --- because the parallelisable spheres have the
dimensions they have. There is partial evidence at every rung. At one qubit, the commutative
special Frobenius algebras on $\mathbb{C}^{2}$ are classified by orthonormal bases --- points of
$S^{2}$, the base of the first fibration --- so that the copied basis plays the role of a
section of the bundle and the spider theorem expresses its idempotence. At two qubits, the
$SU(2)$ fibre of $S^{7} \to S^{4}$ carries exactly the entanglement that survives local
operations, and L\'{e}vay's result that the bundle's natural connection is the $SU(2)$
Yang--Mills instanton \autocite{levay2004geometry} stands on the geometric side where the
interaction of two Frobenius structures stands on the categorical one. At three qubits, the
split of D\"{u}r, Vidal and Cirac \autocite{durvidalcirac2000three} mirrors the algebra exactly:
GHZ is genuinely tripartite, demanding the full non-associative octonion, while W is
pairwise-robust and housed in the associative subalgebra $\mathbb{H} \subset \mathbb{O}$ --- the
special/anti-special Frobenius distinction in geometric dress.

The first of the four conditions the correspondence requires can be discharged at once, and we
do so here, because a problem whose first step has been walked is a different kind of invitation
from one merely posed.

\begin{proposition}[The one-qubit dictionary]\label{prop:f1-dictionary} The $\dagger$-special
commutative Frobenius algebras on $\mathbb{C}^{2}$ correspond bijectively to
\emph{lifted antipodal skeletons} of the complex Hopf fibration $\pi \colon S^{3} \to S^{2}$:
unordered antipodal pairs $P \subset S^{2}$ together with a section of $\pi$ over $P$. The
copyable states of a structure are exactly the section's image, so the lifts are recorded
faithfully --- the projectively trivial centre of $U(2)$ acts non-trivially on structures ---
and the correspondence is $U(2)$-equivariant, the ordered structures forming a single free and
transitive $U(2)$-orbit. The channel a structure induces on density matrices is dephasing in its
skeleton and is independent of the section: the fibre freedom, a torsor for
$\mathrm{Map}(P, U(1))$, is exactly the data absent from the classical record a structure
induces.
\end{proposition}

\begin{proof}
By the classification of Coecke, Pavlovic and Vicary, such structures correspond to orthonormal
bases, recovered as the sets of copyable states \autocite{coeckepavlovicvicary2013orthogonal}.
An orthonormal basis of $\mathbb{C}^{2}$ is two orthogonal rays --- an antipodal pair on the
Bloch sphere --- together with a unit representative of each, which is precisely a section of
$\pi$ over the pair; and $\delta(e^{i\theta}x) = e^{i\theta}\, x \otimes x$ against
$(e^{i\theta}x) \otimes (e^{i\theta}x) = e^{2i\theta}\, x \otimes x$ pins the phases.
Equivariance, the action of the centre, the orbit structure and the induced channel are direct
computations; an executable verification of every finite check accompanies the code release.
\end{proof}

Two remarks fix the proposition's scope. It corrects a natural first reading: a classical
structure does not project to the base but selects lifts into the total space --- classicality
is section data, and what the copy operation broadcasts is a lift, not a ray. And it claims
nothing topological about cloning: the finiteness of the copyable set is linear, structural
no-cloning residing in non-naturality \autocite{abramskyheunen2012hstar}, so the first place the
bundle's global topology can bite is the quaternionic rung, where the degree-one clutching map
stands ready to be compared with the strong-complementarity law governing an interacting pair of
structures. The dictionary is, finally, the complement of the projective treatment cited above:
Burton and Anwar pass functorially to the base and deliberately discard the fibre
\autocite{burtonanwar2026meromorphic}, and the classical structures turn out to live in exactly
the data their projectivisation forgets --- the two readings compose, and the resulting triangle
is a consistency condition available to constrain the remaining conditions of the
correspondence.

\begin{openproblem}[The Hopf-fibration--Frobenius correspondence]\label{op:hopf-frobenius} Let
$\mathbf{Hopf}$ denote the category of the three fibrations
$S^{2^{k}-1} \hookrightarrow S^{2^{k+1}-1} \to S^{2^{k}}$, $k = 1, 2, 3$, with bundle maps as
morphisms, and let $\mathbf{Frob}$ denote the Frobenius algebras on $\mathbb{C}^{2}$ in
$\mathbf{FdHilb}$, with the commutative-special (GHZ) and anti-special (W) classes
distinguished. Construct a functor $F \colon \mathbf{Hopf} \to \mathbf{Frob}$ --- or prove that
none exists --- under which the complex fibration yields the single-qubit classical structures
--- Proposition~\ref{prop:f1-dictionary} identifies these as the lifted antipodal skeletons,
section data over two base points rather than the points alone; the quaternionic fibration
yields the interaction of two such structures, with the $SU(2)$ fibre carrying the entanglement
that survives the classical quotient; the octonionic fibration yields the GHZ/W pair, GHZ on the
full non-associative fibre and W on the associative $\mathbb{H} \subset \mathbb{O}$; and the
composition by which Coecke--Kissinger generate the $N$-qubit classes corresponds to a
well-defined geometric operation on the bundles. If such an $F$ exists, compositional
completeness follows from Adams's theorem as a topological corollary --- there is no fourth
fibration to supply a fourth primitive --- and if $F$ is faithful enough that the implication
also runs back, the termination of the tower and the completeness of the alphabet are equivalent
expressions of a single classification. The first condition is discharged by
Proposition~\ref{prop:f1-dictionary}; what remains open is the quaternionic and octonionic
conditions and, above all, the correspondence of compositions. Both open conditions are already
constrained: the quaternionic fibre's physical content --- the recoverable relative phase and
its purity ledger --- is established (\S\ref{sec:wf-reconciliation}), leaving the
clutching-map-to-strong-complementarity comparison as the structural residue; and any candidate
functor must reproduce the proved canonical-orbit images, GHZ to the colour triplet and W to the
singlet (\S\ref{sec:canonical-orbits}, Appendix~\ref{app:ghz-w-proofs}). A proof that no such
functor exists, or a correspondence at the level of the classifications without functoriality on
compositions, would each also resolve the problem: the first would make the agreement of two
independent classifications a fact demanding its own explanation; the second would establish the
weaker reading.
\end{openproblem}

A resolution in the affirmative would do three things. It would explain compositional
completeness rather than merely prove it --- a corollary of the division-algebra structure of
quantum mechanics rather than an independent algebraic fact. It would give the diagrammatic
language of this paper's own interaction circuits a topological warrant: the GHZ/W spider
calculus would be the unique language the topology permits, not a convenient notation. And it
would open a two-way bridge, translating problems on compositions of Frobenius structures into
operations on bundles and back. The problem sits at the intersection of categorical quantum
mechanics, the quantum-information geometry of the fibrations, and the division-algebra
approaches to particle physics, and it is posed here so that any of the three communities can
take it up in its own language. It is perhaps worth noticing that both classifications already
answer to the same name --- the fibrations are Heinz Hopf's, and so, through the bialgebra
structure of the diagrammatic calculus, are the Hopf algebras --- a coincidence of eponymy that
the conjecture proposes to promote to a theorem. With the epistemic reading above, an
affirmative resolution would leave the tower's endpoint as one limit seen four ways: algebraic,
geometric, compositional, epistemic --- the last being the fixed-point classification of Open
Problem~\ref{op:sigma-functor}, with which this problem is cross-linked, each constraining the
other.

\subsection{What this sets up}

The entanglement-sensitivity of the Hopf maps is not a side feature but the hinge of the
physical reading that follows. Because separability is exactly the collapse of the base to a
product of Bloch spheres, the structure of the entanglement --- the SLOCC class of a multi-qubit
state, and within it the canonical maximally entangled orbit --- is carried by precisely the
degrees of freedom on which $G_{\mathrm{SM}}$ acts in Proposition~\ref{prop:gsm}. This is what
makes it possible, in \S\ref{sec:entanglement-classes}, to read the two genuinely tripartite
entanglement classes onto the $\mathbb{C} \oplus \mathbb{C}^3$ split through their canonical
representatives: the GHZ orbit confined to the colour-triplet $\mathbb{C}^3$ sector, the W orbit
reaching the $\mathbb{C}$ singlet, recovering the quark--lepton division not as an assignment
but as a consequence of the entanglement structure. The scope of those statements --- canonical
orbits rather than whole classes, and the identification of the canonical orbits with the
physical particles --- is made precise there. The Hopf fibrations are the bridge; the
entanglement classes, treated next, are the traffic that crosses it.

%  Division-algebras' closing now carries the S^7 hand-off sentence
%  (written 2026-07-06 after this file was read).
%
%  INPUT DEPENDENCIES (both held in the audited outputs):
%      entropic_gravity_family_paragraph_2026-07-03.tex  (sec:jacobson)
%      gr_nonlocality_revision_2026-06-13.tex            (sec:gr-nonlocality)
%  The 2026-07-04 additions (asymmetry paragraph; Penrose second route;
%  acceleration/Unruh arc) are FULLY PASTED; the separately uploaded
%  additions file is redundant and can be archived.
%
%  AUDIT RECORD: opener formalised (Russell typo; bare ref retargeted
%  to the freshly minted sec:complex-amplitudes); the "rotations of
%  S^7" wording made precise at BOTH occurrences (Moufang loop, not a
%  group; SO(8) disclaimer added -- referee hazard removed); two \\
%  paragraph-break bugs fixed; dark-matter paragraph retargeted to
%  sec:generations and stitched to op:sterile-gravity (the zoo's
%  demotion and this paragraph's priority condition are one question);
%  CP paragraph stitched to the fourth-party reading of
%  sec:interactions; stale mid-file mini-header retired
%  (sec:k3-sterile placeholder dies with it); the section's narrated
%  debts NUMBERED at last -- op:comp-metric and op:three-dimensions
%  minted in the Open-problems subsection, cited from the in-text
%  "named open problems" line and from the closing summary.
%
%  KEYS: adams1960hopf matches corpus. Unique to this section, presumed
%  in the .bib since the paper compiles: Spence-Brown-1969 (suggest
%  normalising the key spelling at leisure), VanRaamsdonk2010,
%  MaldacenaSusskind2013, MuellerMasanes2013, MasanesMuellerEtAl2014,
%  toffoli_2003, Chaitin1998, Penrose1996, HohnMueller2016,
%  GarnerMuellerDahlsten2017, Jacobson1995, Bekenstein1973, tHooft1993,
%  Susskind1995, spekkens2007toy, Andrianov2001, Sakharov1967,
%  Planck2020, Riess1998. Shared keys verified: small_2006, Penrose1994,
%  unruh1976notes (swap-flag stands), baez2002octonions, Wheeler1990,
%  furey2018three, DiBiagioRovelli2021, Rovelli1996a, Spekkens toy.
%  British English; \autocite; no bullets in prose.
% =====================================================================

\section{Spacetime from non-computability}
\label{sec:space-of-computations}

This section is a deliberate diversion: it takes the idea that non-computable steps are measured
by imaginary information (\S\ref{sec:complex-amplitudes}) and asks what geometry it buys.
Spencer-Brown records that Bertrand Russell thought well of the underlying move
\autocite{Spence-Brown-1969}; the reader may judge whether the framework spends it responsibly.

The programmes that derive spacetime from quantum mechanics---Van~Raamsdonk's
\enquote{building spacetime with entanglement} \autocite{VanRaamsdonk2010}, the
ER\,$=$\,EPR conjecture \autocite{MaldacenaSusskind2013}, and M\"uller's
information-theoretic recovery of three-dimensional space from the qubit
\autocite{MuellerMasanes2013}---take entanglement or information as primitive and
grow geometry from it. The route taken here is the same in spirit but inverted in
order: rather than positing spatial dimensions and forbidding superluminal
signalling, we \emph{define} spatial separation as a measure of non-computability
and read the dimension, the signature, and the Lorentz group off the result. The
proposal originates in the author's CASYS'05 paper \autocite{small_2006}, which
first suggested treating spacetime as a space of computations and attempted to
justify three dimensions on this basis; the present account tightens and extends
it. This is the framework's most structural and least formally complete arc; the
three-dimensionality result is a structural conclusion awaiting a formal proof,
and several of the links below are conjectural readings. Each is flagged as such,
so the reader sees exactly where the chain is load-bearing and where it is an
invitation.

\subsection{Timelike and spacelike distance as computation}
\label{sec:comp-distance}

A computation carries a state from axioms to a theorem in a number of steps.
Counting the computable steps between two states gives a real measure---time. A
physical clock counts completed irreversible loops, so proper time literally
counts loop completions and the metric $\mathrm{d}\tau^{2}=g_{\mu\nu}\,
\mathrm{d}x^{\mu}\mathrm{d}x^{\nu}$ is a loop-density tensor, the conversion from
coordinate intervals to loop counts; Toffoli's reading of the Lagrangian as a
count of action quanta is the same statement \autocite{toffoli_2003}.

Spacelike separation is the complementary, non-computable measure. Define the
non-computable distance between two propositions, relative to an axiom system, as
the number of bits that must be added to the axioms to make the second derivable.
A closed timelike curve is a computation whose output is its own input; by the
fixed-point theorem it has a unique consistent solution or none, and the
existence of that solution is one bit not derivable from any axiom system that
does not already contain the loop---one bit per curve, in the manner of Chaitin's
$\Omega$ \autocite{Chaitin1998}. The spacelike distance between two events is then
the number of such curves that would have to be closed to make them causally
connected. Because each added bit is a non-computable proposition, it is carried
by the imaginary part of the complex information measure
$S=-\ln\lvert p\rvert-\mathrm{i}\theta$ derived in
Section~\ref{sec:complex-amplitudes}: computable distance is counted in real
units, non-computable distance in imaginary ones. The Minkowski signature, with
spacelike intervals imaginary relative to timelike ones, is the immediate reading. The convention is not new to this paper: it was proposed in the author's CASYS paper two
decades ago, where counting non-causal steps in imaginary numbers was introduced and its formal
proof flagged, already then, as a key open problem \autocite{small_2006} ---
Open Problem~\ref{op:comp-metric} is that flag's direct descendant, stated now with the machinery
to be answerable.
That the resulting structure is a genuine metric space, and that the signature
follows rather than being fitted, are open problems---at present the
identification is motivated by reproducing what nature does, not proved.

%% VERIFY-SUPPORT 2026-09-27: conway_2003 is cited below for "any two octonions
%% generate an associative subalgebra" (Artin's theorem; a Moufang loop is
%% diassociative) and harvey1982calibrated for "a three-plane in Im O is
%% associative iff the associator of a basis vanishes, and the associative
%% three-planes are the imaginary parts of the quaternionic subalgebras"
%% (Harvey--Lawson, ch. IV). Neither passage was read in the session that wrote
%% this; both statements are checked numerically in Convention A by
%% scripts/antitelephone_network_associativity_2026-09-27T1900.py (PASS).
A reader used to the standard order of ideas should pause here, because the order has been
inverted, and the inversion is used below. In special relativity the tachyonic antitelephone is a
corollary: a signal faster than light, composed with a boost, would arrive before it was sent and
close a timelike loop, therefore there is no such signal \autocite{tolman1917theory,
benfordbooknewcomb1970tachyonic}. The metric comes first and the impossible telephone is a check
on it. Here the telephone comes first. Two observers are spacelike separated when an antitelephone
between them would close a loop; the loops such telephones would close, and carry nothing round,
are what spacelike separation is, their number is the distance, and the metric is read off from
them. Nothing is lost, since the corollary returns as a tautology: an antitelephone that carried
a record would close a causal loop, and its two ends would not have been spacelike after all.

What is gained is a definite meaning for a spatial dimension. Picture many observers, each holding
an antitelephone to each of the others. A message relayed from one observer to another through a
third composes two transports round two loops, and the network is consistent only if every relay
composes to a single transport whatever the bracketing of the relaying: a spatial dimension is a
family of loops whose transports compose associatively, which is a group law on a parallelisable
sphere. That is a condition a family of loops can fail, and \S\ref{sec:three-dimensions} shows
that the loops of the seventh sphere fail it. Two antitelephones always compose, because any two
octonions generate an associative subalgebra \autocite{conway_2003}; three compose only when their
imaginary parts lie in one quaternionic subalgebra, which is what the associative three-planes of
Harvey and Lawson are \autocite{harvey1982calibrated}, and a third link outside that subalgebra is
seen by the network exactly through its component outside it. The largest consistent network of
antitelephones inside $S^{7}$ is therefore an $S^{3}$, and there are three spatial directions.
When the seventh sphere is set aside there as a ruler, which is the octonionic spacetime of the
division-algebra literature, the claim is not that $\mathbb{R}^{8}$ has no metric, nor that the
ten-dimensional Minkowski space of $SL(2,\mathbb{O}) \cong \mathrm{Spin}(9,1)$ fails to exist
\autocite{baez2002octonions}; it is that the octonionic directions are not loops on which
antitelephones compose, and so are algebra rather than distance. In the standard order that
sentence has no content; in this one it is the reason space has three dimensions.

The split just defined explains the most familiar asymmetry in physics: that one can move back and
forth in space but not in time. \S\ref{sec:continuity} identifies the pre-measurement regime with
continuous reversibility --- its transformation group is connected, and every motion in it is
undoable. Spacelike separation, being the non-computable sector, is exactly the part of the
relational structure on which no record has been written; motion across it is motion within the
reversible regime, and it inherits that regime's continuity and its two-sidedness: a place can be
revisited because going and returning writes no record that forbids the return. Timelike separation
is the record chain itself --- proper time counts completed irreversible loops, as above --- and a
moment cannot be revisited because a moment \emph{is} a stratum of records, and returning would
require unwriting them, which is what \textsc{fanout} forbids (\S\ref{sec:fanout-boundary}). The
asymmetry between space and time is therefore not an extra postulate about the manifold: it is the
pre-/post-\textsc{fanout} asymmetry read as geometry --- space is what the reversible group feels
like from inside, time is what the record chain feels like --- and the arrow-of-time discussion
later in this section is the same asymmetry read thermodynamically. The identification is the
framework's structural reading of the split, offered in the same register as the rest of this
subsection.

\subsection{Three dimensions of space}
\label{sec:three-dimensions}

The closed timelike curves that define spacelike distance require information to
be parallel transported consistently around a loop, and consistent parallel
transport on a sphere requires it to be parallelisable. By Adams' theorem the
only parallelisable spheres are $S^{1}$, $S^{3}$ and $S^{7}$
\autocite{adams1960hopf}, the unit spheres of $\mathbb{C}$, $\mathbb{H}$ and
$\mathbb{O}$. Impose now the requirement that information arriving from outside
the causal zone must appear to possess its own causal history---that is, to be the
result of a sequence of information-preserving operations, hence a sequence of
group operations. The spheres $S^{1}$ and $S^{3}$ are groups under their own multiplication --- $U(1)$ and
$SU(2)$ --- and their operations preserve information; $S^{7}$'s multiplication is non-associative,
a Moufang loop rather than a group, so $S^{7}$ cannot serve as a ruler. (The rotation group
$SO(8)$ of the seven-sphere is of course a group; what fails is the sphere's \emph{own}
composition law, the one a ruler would use.) The same non-associativity
forbids a metric built from the sphere's own composition law outright: path
composition $d(A,C)$ through $B$ depends on the bracketing, so the triangle inequality is not merely violated but ill-defined, and
no geodesic structure exists. (What is untouched should be said: the Euclidean metric of
$\mathbb{R}^{8}$, and the ten-dimensional spacetime that $SL(2,\mathbb{O})$ furnishes, are metrics
of the ordinary kind; what fails is a ruler made of antitelephone loops, which is what a spatial
dimension is in this section, as \S\ref{sec:comp-distance} set out. Nor is the failure total:
relays between any two links compose, and among three they compose exactly when the links lie in
one quaternionic subalgebra, so the largest consistent network inside $S^{7}$ is an $S^{3}$
\autocite{conway_2003,harvey1982calibrated}.) Since the abelian $S^{1}$ alone cannot carry
non-abelian structure, spacelike distance must be three-dimensional. Dimension
three is an output, not an input \autocite{small_2006}.

The division of labour is then exact. $S^{1}$ ($\mathbb{C}$) contributes the time
dimension and the $U(1)$ of electromagnetism; $S^{3}$ ($\mathbb{H}$) contributes
the three spatial dimensions and the $SU(2)$ of the weak interaction; $S^{7}$
($\mathbb{O}$) contributes no geometry but, as a filter on what information may be
exchanged, the $SU(3)$ of the strong interaction. A sphere that cannot be a ruler
can still be a filter, which is why the strong force resists geometric
interpretation while electromagnetism and the weak force do not. This is the same
$\mathbb{C}\oplus\mathbb{C}^{3}$ split that yields $G_{\mathrm{SM}}$ in
Section~\ref{sec:entanglement-classes}, now read as the reason spacetime is $(3+1)$-%
dimensional. The conclusion converges with M\"uller and Masanes, who prove from
information-theoretic postulates, assuming no quantum theory, that minimal
direction-carrying systems with continuous reversible dynamics force $d=3$ and
quantum theory on two qubits \autocite{MuellerMasanes2013}, and with the result
that entanglement exists only for a three-dimensional Bloch ball
\autocite{MasanesMuellerEtAl2014}; the framework supplies the reason such systems
exist, Cayley--Dickson doubling from self-reference, while their work supplies
the rigorous dimensional uniqueness. The status of the framework's own route is a
structural result requiring formal proof: the one-bit-per-curve identification and
the precise reduction of the $S^{7}$ filter to $SU(3)$ rather than its full automorphism group
$G_{2}$ are numbered below (Open Problems~\ref{op:comp-metric} and~\ref{op:three-dimensions}).

A scoping remark, and then a second route to the same count. The computational model of spacetime
proposed here has many detailed parallels with the real thing, but a parallel is not an identity,
and the identity is not claimed in this paper: establishing it is work deferred to the second paper.
What can be claimed now is that the parallel is productive in a specific direction --- it suggests a
derivation of the dimension count carried out entirely within general relativity, with no
computational premise. Acceleration and gravitation tip the light cone, and the tipping can be
carried to the extreme: Penrose's discussion runs the cone through a full circle
\autocite[382]{Penrose1994}, at which point the tipped direction closes on itself --- which is to
say, the extreme of tipping is a closed timelike curve, the very object whose closure the
computational definition of \S\ref{sec:comp-distance} counts in bits. The two routes are one route
entered from opposite ends. Now demand of general relativity only this: that observers accelerating
with respect to one another, and with respect to unaccelerated observers, retain a consistent
conception of the distance between them. Each observer's tipping traces a path on a sphere of
directions, and mutual consistency requires the associated rotations to compose without regard to
bracketing --- to form a group. Compositions on $S^{1}$ and $S^{3}$ do, forming $U(1)$ and $SU(2)$ as above; compositions on
$S^{7}$ do not, the multiplication being non-associative, so observers tipped across an
$S^{7}$-worth of directions could not agree on distance even pairwise: the same failure that forbade
$S^{7}$ a metric above forbids it here a shared ruler. The largest sphere admitting a consistent
conception of distance under mutual acceleration is therefore $S^{3}$, and space has at most three
dimensions; with the floor argued above --- the abelian $S^{1}$ alone cannot carry the required
structure --- the count is exactly three. The status matches the subsection's. The tipping and its
extreme are general relativity's own \autocite{Penrose1994}; the consistency demand is
elementary; the associativity filter is the algebraic fact used throughout; the reading of composed
tippings as motion on the parallelisable spheres is the structural step, shared with the
computational route. The convergence of the two routes on the same endpoint is offered as evidence
for the parallel, while the identity it gestures at remains the second paper's burden.

\subsection{Recovering the Lorentz group}
\label{sec:lorentz}

A measurement changes what the observer knows, hence, relationally, the causal
structure: it tips the light cone. This is Penrose's association of cone-tipping
with state reduction \autocite{Penrose1996} run in the opposite causal direction,
reduction driving the tipping rather than the reverse, and the two are then one
operation in two charts---both described by $SL(2,\mathbb{C})$. The algebraic
backbone is the division-algebra correspondence
$SL(2,\mathbb{A})\cong\mathrm{Spin}(1,d{+}1)$ \autocite{baez2002octonions}: $\mathbb{R}$
gives $1{+}1$, $\mathbb{C}$ gives $3{+}1$, $\mathbb{H}$ gives $5{+}1$ and
$\mathbb{O}$ gives $9{+}1$ dimensions. The computability split that singles out a
preferred $\mathbb{C}$ fixes the fundamental state space as $\mathbb{C}^{2}$ and
hence $SL(2,\mathbb{C})=\mathrm{Spin}(1,3)$: the Lorentz group of $3{+}1$
spacetime. The qubit and Minkowski space are not merely both Euclidean-and-three%
-dimensional by coincidence; the Hermitian $2\times2$ matrices are
$\mathbb{R}^{1,3}$, $SL(2,\mathbb{C})$ acts on them as the Lorentz group, and the
qubit density operators fill the future cone---the correspondence M\"uller's group
develops operationally, recovering Lorentz transformations from quantum
communication \autocite{HohnMueller2016} and the complex and quaternionic qubit
from the relativity of simultaneity \autocite{GarnerMuellerDahlsten2017}.
In these coordinates the Bloch ball is the unit-trace slice of the cone and the Bloch sphere its
boundary rays, the celestial sphere on which $SL(2,\mathbb{C})$ acts by M\"obius
transformations; state reduction sends a point of the ball to a point of the sphere, and a boost
sends the sphere to itself, which is what lets the two be one operation in two charts.
The $SL(2,\mathbb{A})$ correspondence is rigorous; its identification with
measurement-as-tipping is the framework's structural reading.

One consequence of the identification has been implied throughout this section and deserves its own
paragraph, because the next paragraph turns on it. If measurement creates a relation between
observer and observed, and if --- relationally --- a new relation alters the causal structure, then
the tipping of the cone is not merely \emph{associated} with the creation of the relation: it is
what the creation of the relation is, seen geometrically. And a tipping of one's cone against the
congruence of one's neighbours already has a name in physics: acceleration. The framework's
statement, made explicit here, is that \emph{creating a new relation between observer and observed
is experienced as acceleration}. The resonance with the equivalence principle is exactly what it
should be: the principle makes acceleration and gravitation locally indistinguishable, Penrose's
association --- run in this section's direction, reduction driving the tipping --- makes state
reduction and cone-tipping one operation, and the framework's identification closes the triangle,
measurement and acceleration meeting where the principle says they must. The claim is structural:
the framework's reading of the cone-tipping established above, not an independent derivation.

Taken literally, the identification has a measurable face, and it is a familiar one: the Unruh
effect, read in this section's own ledger. Recall the definitions of \S\ref{sec:comp-distance}:
timelike distance counts computable steps, spacelike distance counts non-computable bits. What an
observer can compute --- what their \textsc{fanout}s can convert into classical records --- is fixed
by their causal structure, and the previous paragraph has just made that structure
relation-dependent. An accelerated observer's cones are tipped, and with them the boundary of the
recordable: relations that for the inertial observer lie wholly in the reversible, pre-record sector
--- the entanglement structure of the vacuum, non-computable for them, spacelike in the
computational sense --- fall, for the accelerated observer, partly on the recordable side. The mechanism is the one
Chaitin identified for formal systems \autocite{Chaitin1998}: a proposition undecidable relative to
an axiom system becomes decidable when information is added to the axioms --- and
\S\ref{sec:comp-distance} defined spacelike distance as exactly the bits whose addition closes the
gap. Creating a relation is adding information to the observer's axiom base, so timelike and
spacelike distance are converted into one another by the addition and subtraction of information,
and acceleration is the continuous version of that addition. A
relation newly recordable is a detection. Particles appear from nowhere: an accelerated detector in
the inertial vacuum registers a thermal bath at $T = \hbar a / 2\pi c k_{\mathrm{B}}$
\autocite{unruh1976notes}. Two guards keep the reading honest. Nothing here reclassifies the
invariant interval --- timelike and spacelike separation in the emergent metric are Lorentz
invariants, and no event crosses that classification; what is observer-relative is the computational
partition, which relations a given observer's horizon renders recordable. In the ledger of
\S\ref{sec:comp-distance} that partition \emph{is} the timelike/spacelike distinction, so in the
framework's own units timelike distance is \emph{rotated} into spacelike, and spacelike into
timelike, by a change of observer --- the tipping of the cone is literally that rotation --- and the
detector's clicks are the rotation made audible. And the effect is established physics
while the reading is the framework's, flagged structural. What gives the reading its content is a
parallel: the inertial and the accelerated observer stand to the vacuum exactly as Wigner and his
friend stand to the laboratory (\S\ref{sec:wigners-friend}) --- one holds a coherent relation and no
event, the other holds a record, and the facts are relative on both sides in precisely Rovelli's
sense \autocite{DiBiagioRovelli2021}. Since the previous paragraph identifies acceleration with
relation-creation, the parallel is not an analogy but the same structure in gravitational dress, the
horizon standing where the laboratory wall stood: \emph{the Unruh effect is Wigner's friend for
accelerated observers}. And two subsections of this section that seemed independent here meet: the
temperature Jacobson's construction requires across every local Rindler horizon
(\S\ref{sec:jacobson}) is exactly this temperature --- the $T$ of the Clausius relation is the $T$
of the tipped cone, supplied by the very relation-creation with which this subsection began.

\subsection{The relationalist commitment}
\label{sec:relationalism}

General relativity is relational: by the hole argument, spacetime points carry no
intrinsic identity, and the metric is fixed by the network of coincidences and
interactions, not by a pre-existing stage \autocite{Rovelli1996a}. The
computational picture is relational by construction and explains why it must be:
a background metric would require spatial distances to be pre-computable, and
non-computability forbids exactly that. General covariance is the statement that
distances have no pre-assigned values and depend on what information is actually
present. The metric is therefore dynamical not by stipulation but because it
measures a quantity that cannot be fixed in advance.

\subsection{Jacobson, holography, and imaginary information}
\label{sec:jacobson}

Jacobson derives the Einstein equation as an equation of state, from the Clausius
relation $\delta Q=T\,\mathrm{d}S$ across every local Rindler horizon together
with the Bekenstein proportionality of entropy to horizon area
\autocite{Jacobson1995, Bekenstein1973}. In the computational reading this is the
self-consistency condition on the loop-counting network: the Einstein equation is
the requirement that the bookkeeping of computation be internally consistent
across surfaces. The holographic principle, that the information in a region is
bounded by its bounding area \autocite{tHooft1993, Susskind1995}, then acquires a
natural reading here. A surface area is a squared distance; if spacelike distance
is measured in imaginary units of information, an area is measured in the square of
those units, that is in \emph{negative} information---entropy. The holographic
link between area and entropy thus requires the square root of negative
information, and that square root is precisely the imaginary unit of the complex
amplitude, the observer's self-referential phase. Read this way, Jacobson's
account of information crossing surfaces implicitly assigns an imaginary
information value to spacelike distance, and the $\mathrm{i}$ of the Minkowski
signature is the same $\mathrm{i}$ as the $\mathrm{i}$ of the quantum amplitude.
This identification is the most speculative step in the section: it is
structurally suggestive and it makes the holographic form unsurprising, but it is
a conjectural reading rather than a derivation, and is offered as such.

The conjecture has a two-way motivation that this paper's record should carry, because each
direction was on the table before the other was noticed. Forward: if spacelike distance is counted
in imaginary units of non-computable information, then any spacelike area is the product of two
imaginary lengths and is therefore \emph{negative} information --- which is to say entropy, by
Brillouin's identification \autocite{brillouin1956science} --- and a holographic principle arrives
automatically, the sign supplied by $i^{2}$ rather than by interpretation. Converse, and more
striking: in the holographic principle as ordinary physics already states it, surface area is
proportional to entropy, and entropy is negative information; so by implication ordinary physics
has been measuring spacelike separation in imaginary units of information all along, and the
convention of \S\ref{sec:comp-distance} is read off from established results rather than
invented. The same imaginary unit is already employed once in this section --- it is what makes
the signature of the interval Minkowski --- so the sign of the area law is its second job, not a
second assumption. What the two-way motivates it does not prove: the formal statement is exactly
the metric-and-signature half of Open Problem~\ref{op:comp-metric}, together with the
one-bit-per-bit-pair coefficient owed to the second paper. The stakes of the conjecture have
recently risen: Dorau and Much prove, with modular-theoretic rigour, that the semiclassical
Einstein equations follow from quantum relative entropy on a local horizon \emph{once the
entropy--area proportionality is assumed} \autocite{doraumuch2026relative} --- so what this
framework proposes to derive is precisely the remaining input of a derivation already carried
through at that standard.
% =====================================================================
%  DROP-IN: short closing passage for sec:jacobson ("Jacobson,
%  holography, and imaginary information") in
%  spacetime_from_non-computability.tex.  Revised 2026-07-03 per JS:
%  mention in passing only -- "other people are thinking in terms of
%  computation and spacetime, refer to this work, and move on."
%
%  PLACEMENT: append at the end of sec:jacobson, after "...a conjectural
%  reading rather than a derivation, and is offered as such." and before
%  \subsection{Non-locality already present in general relativity}.
%
%  DELIBERATE DECISIONS (recorded so the drafting history is clear):
%  - No argument against Vopson in the paper. The critical assessment
%    (fixed lattice of labelled cells = the "atoms of spacetime"
%    Einstein's hole-argument resolution rules out; the claim that
%    physical entropy falls under gravitational clustering contradicts
%    gravitational thermodynamics) lives in the conversation record and
%    in the Mendeley N1 notes, NOT in the manuscript.
%  - One neutral classificatory clause only (pre-existing structure vs
%    relational emergence), so the citation cannot be read as alignment.
%  - Verlinde cited alongside since he is named and was previously
%    uncited in this section.
%
%  CROSS-REFS: sec:relationalism.  Citation keys: verlinde2011origin,
%  vopson2025gravity (RIS in entropic_gravity_refs_2026-07-03.ris).
% =====================================================================

The question this section asks --- whether gravitation is a signature of computation --- is now
asked more widely: Verlinde derives Newtonian gravity as an entropic force associated with
information on holographic screens \autocite{verlinde2011origin}, and Vopson has recently read
gravitational clustering as data compression in a discrete, information-bearing space, offering it
as evidence of a computational universe \autocite{vopson2025gravity}. Those routes take space to be
a pre-existing information-bearing structure, where the present account takes space itself to be
relational and emergent (\S\ref{sec:relationalism}); the resemblance is at the level of the
question rather than the answer.


\subsection{Non-locality already present in general relativity}
\label{sec:gr-nonlocality}
% =====================================================================
%  REVISION of sec:gr-nonlocality in spacetime_from_non-computability.tex
%  Drafted 2026-06-13. Three reinforcing changes, all within this one
%  subsection:
%    (1) ER=EPR / Van Raamsdonk recast from "same direction of travel"
%        (convergence) to genuine CONTRAST (near-opposite mechanism):
%        they build a geometric bridge; we observe there was never a gap.
%    (2) The Vaxjo 2025 poster stated explicitly as the PRIORITY claim --
%        first public recognition that relational GR already accommodates
%        non-locality. Anchored to \autocite{Small2025} (existing key).
%    (3) The knowledge-space / co-location MECHANISM expanded (what JS
%        asked for): entangled systems share one knowledge-space location
%        and acquire distinct locations only by interacting with already-
%        located things. Includes the deeper-than-emergent-space subtlety
%        (co-location =/= coincidence) and the forward link to
%        sec:confinement.
%
%  HOW TO APPLY: this REPLACES the passage in sec:gr-nonlocality from
%  "The upshot is that general relativity, taken relationally and
%  seriously, already contains quantum non-locality..." (the sentence at
%  ~line 186) through the end of that paragraph at ~line 197
%  ("...flagged as a thesis rather than a theorem."). Everything BEFORE
%  that sentence (the sealed-bottle electron, the pre-spatial setup) is
%  retained unchanged; this revision picks up from the "upshot" and
%  extends it. The \medskip summary paragraph after (line 199+) is also
%  retained unchanged.
%
%  No new citations required (Small2025, VanRaamsdonk2010,
%  MaldacenaSusskind2013 all already present; fanout-boundary and Zurek
%  already cited elsewhere).
% =====================================================================

The upshot is a thesis about general relativity that the wider literature has not
drawn, and which the present author first stated publicly in the 2025 V\"axj\"o
poster \autocite{Small2025}: relational general relativity, taken seriously,
\emph{already} contains quantum non-locality. Entanglement is the pre-spatial
relational structure, and what looks like spooky action between distant points is
the resolution of a relation that was never spatially separated to begin with. The
claim is one of recognition rather than construction. Nothing need be added to
general relativity to accommodate the non-locality of quantum mechanics; one need
only notice that a theory in which spatial separation is fixed by the network of
interactions, and not by a pre-existing stage, has never committed to the
separation of two systems that have not yet interacted through that network. The
non-locality is already there, unrecognised, in the relational reading of the
theory, and the poster's contribution---advanced earlier in the author's CASYS'05
paper \autocite{small_2006} and developed here---was to point this out.

It is worth being exact about how this differs from the programmes it most
resembles, because the resemblance is real at the level of the phenomenon and
misleading at the level of the mechanism. Building spacetime from entanglement
\autocite{VanRaamsdonk2010} and the ER\,$=$\,EPR conjecture
\autocite{MaldacenaSusskind2013} share with the present account the conviction
that the spatial separation of entangled partners is not fundamental. They differ
from it---and from each other less than they differ from it---in what they take to
follow. Those programmes \emph{construct} a geometric connection: a wormhole, an
Einstein--Rosen bridge, a throat of new spatial geometry that realises the
entanglement as a connection through space. Entanglement, on that reading, builds
or is geometric structure linking the partners. The present account does the
opposite. It builds nothing and connects nothing, because there was never a gap to
span: the partners were never at two places, so no bridge between two places is
required. Where ER\,$=$\,EPR adds a geometric object to explain the correlation,
the relational reading subtracts an assumption---that the partners were
separated---and the puzzle dissolves rather than being mechanised. One approach
erects a structure across a gap; the other observes that relational general
relativity never opened the gap. This is not the same argument reached by a
different route, and it is not a corollary of the holographic results; it is a
distinct, and in its mechanism nearly opposite, account of the same fact, and it
is offered as the framework's interpretive thesis, flagged as a thesis rather than
a theorem.

The mechanism is made precise by reading the space of computations from the
observer's side, as a \emph{knowledge space}: a position in it is a state of
information, not a point of a manifold, and what fixes a system's position is what
is known of it through its interactions. Entanglement is co-location in this
space. Two systems in an entangled relation do not occupy two knowledge-space
positions awaiting a connection; they occupy \emph{one}, because the relation is a
single computational object that has not been resolved into parts. A system
acquires a \emph{distinct} location only by interacting with other systems that
already hold locations in the network---location is inherited relationally, a
place obtained by being related to things that have places. An entangled pair
shares one knowledge-space location until one partner is fanned out
(\S\ref{sec:fanout-boundary}) into the already-located network; that interaction
is what assigns the partner a distinct location, and the other partner's location
follows from the now-resolved relation. Spatial separation is in this sense
\emph{produced} by interaction, not revealed by it: the separation did not
pre-exist the measurement that established it.

One distinction must be drawn sharply, because without it the reading invites a
false consequence. Co-location in knowledge space is not coincidence in emergent
space. Knowledge space is the deeper, pre-spatial layer from which emergent
spacetime is built (\S\ref{sec:comp-distance}, \S\ref{sec:three-dimensions}); two
systems may share a knowledge-space location while the composite they belong to,
having interacted with the rest of the network, holds an emergent-space location
of finite extent. The point has a direct realisation in the strong sector,
developed in \S\ref{sec:confinement}. The three quarks of a colour singlet are
maximally entangled, so they share a single knowledge-space location; the
colour-singlet constraint, which forbids them the external entanglement that
location-acquisition requires, is exactly what prevents any one of them from
acquiring a distinct location of its own. Confinement is, on this reading,
pre-spatial co-location that cannot be resolved: a quark cannot be drawn to a
place apart because, in the knowledge space that is fundamental, it has no place
apart to be drawn to, and the singlet will not let it acquire one. That the proton
nonetheless has a finite radius in emergent space is no contradiction---the
emergent extent is the structure of the bound state as a whole, not a separation
of the singlet---and relating the knowledge-space co-location to the
emergent-space potential between the charges is precisely the open problem of
\S\ref{sec:confinement}.

The graded process by which new interactions create new knowledge-space locations,
and so create and sharpen emergent spacetime structure, is itself graded in the
manner of \S\ref{sec:fanout-boundary}---a location is as definite as the
interactions that fix it are redundant---and its dynamics is the subject of the
second paper. The present paper states what knowledge-space co-location is; how
locations proliferate, and how emergent distance is built from their proliferation,
is deferred.

\medskip
\noindent  The spacetime arc therefore yields, from one definition of spacelike
distance as non-computability: the signature, three spatial dimensions,
the Lorentz group, the dynamical relational metric, a reading of the
holographic principle, the presence of non-locality within general
relativity, the identification of relation-creation with acceleration
and of the Unruh effect with Wigner's friend, and the space--time
mobility asymmetry as the \textsc{fanout} asymmetry read as geometry. Of these,
the $(3+1)$ count and the Lorentz recovery are structural results that converge with independent
rigorous work and await their formal proof --- the count now numbered as
Open Problem~\ref{op:three-dimensions}; the metric-space status, the signature and the
one-bit-per-curve measure (together Open Problem~\ref{op:comp-metric}), the imaginary-information
reading of holography, and the general-relativistic non-locality, the acceleration and
Unruh readings, and the mobility asymmetry are conjectural, and are named here precisely so that the boundary between what is shown and what
is proposed is visible.

\subsection{Why the metric is not to be quantised}\label{sec:no-quantised-gravity}

A consequence has been implicit in everything above and deserves its plain statement. Causal
structure, on this framework's reading, is the bookkeeping of locked relations --- post-\textsc{fanout}
by construction --- while superposition is the signature of the pre-\textsc{fanout} regime. A
``superposition of causal structures'' is therefore not a difficult state but a category error: it
applies the pre-record operation to a record-side object, exactly as ``a superposition of
measurement outcomes'' does. Gravity lives in the post-\textsc{fanout} regime. The quantum layer
beneath it exists and is not in dispute; but it is \emph{pre-geometric} --- it carries amplitudes
over the formation of relations, not amplitudes over formed geometries --- so there is no
fundamental metric to quantise, and the standard programme of superposing spacetimes pursues an
object the ontology does not contain. One guard rail keeps the claim honest: collective
excitations of the locking pattern may perfectly well exist as effective quanta --- gravitons as
the phonons of the relational network --- and the thesis forbids fundamental metric superposition,
not effective field theory; this is Jacobson's equation-of-state moral \autocite{Jacobson1995},
made structural. The relation to gravitational-collapse programmes is an inversion: where
superposed geometries are there read as unstable \autocite{Penrose1996}, here they are read as
unformed --- what decays in those models is what locks in this one.

Three experimental postures follow, each at the status the reading carries. Laboratory
demonstrations of indefinite causal order \autocite{procopio2015experimental, rubino2017indefinite} are, on this
account, demonstrations of the nascent layer: the indefiniteness lives below the locking threshold
and definitises at it, so the framework adopts the programme and adds a boundary --- indefinite
causal order can never be amplified into records while remaining indefinite. Gravitationally
induced entanglement of the kind the Bose--Marletto--Vedral proposals seek
\autocite{bose2017spin, marlettovedral2017gravitationally} is \emph{predicted}, because the
pre-geometric layer is quantum --- while witnessing that layer, not a superposed metric: the
framework splits the experiment's usual dichotomy. And gravitationally driven anomalous collapse
is denied: locking is redundancy and thermodynamics, so decoherence experiments should find
standard rates and nothing gravity-specific. That no consistent quantisation of gravity exists was
the conclusion of the author's CASYS paper by a different route \autocite{small_2006}; the present
framework supplies the mechanism the conclusion then lacked.

\subsection{Open problems and prospects}
\label{sec:spacetime-open-problems}

The spacetime arc is the framework's frontier, and honesty about its edges
matters more here than anywhere else in the paper. We set out the genuine
problems and the prospects together, deferring the deepest to a second paper
rather than papering over them. Two of the arc's debts are well-posed now, and we number them.

\begin{openproblem}[The computational metric]\label{op:comp-metric}
Prove that the paired distance of \S\ref{sec:comp-distance} --- computable steps counted real,
non-computable bits counted imaginary, one bit per closed curve --- defines a genuine
pseudo-metric on the space of computations, with the Minkowski signature derived rather than read
off; equivalently, establish the one-bit-per-curve measure as a measure, and the signature as its
theorem. Until then the identification stands as motivated by reproducing what nature does, which
is the status flagged at its introduction.
\end{openproblem}

\begin{openproblem}[Dimension three as a theorem]\label{op:three-dimensions}
Turn the two convergent routes of \S\ref{sec:three-dimensions} into a proof: formalise the
parallelisability-plus-associativity filter (the computational route) and the mutual-acceleration
consistency demand (the general-relativistic route) so that $d = 3$ follows as a theorem, and
include the reduction clause --- that the $S^{7}$ filter acts through $SU(3)$, the stabiliser of
the selected complex direction, rather than through its full automorphism group $G_{2}$. The
M\"uller--Masanes uniqueness results \autocite{MuellerMasanes2013, MasanesMuellerEtAl2014} are
the rigorous convergents any proof must meet.
\end{openproblem}

\paragraph{The substrate and the generation of events (deferred).}
A computation runs on something; a network of events presupposes that events
occur. Both are the same objection in two forms---the network's hardware and its
clock---and the original CASYS'05 statement of the programme raised the
implementation regress and set it aside \autocite{small_2006}. The framework's
response is that a self-referential structure is the only kind that need not
presuppose a meta-substrate: it generates its own. The natural algebraic emblem
is $E_{8}$, the unique simple Lie group whose adjoint representation is also its
smallest faithful one, so that what it is and what it acts on coincide; the $240$
unit integral octonions that generate it are exactly the $E_{8}$ root system, so
the self-action is welded to the division-algebra spine of
Section~\ref{sec:entanglement-classes}. This is Wheeler's self-excited circuit
\autocite{Wheeler1990} given an algebraic body, and it converges with Furey's use
of an algebra acting on itself to generate particle content
\autocite{furey2018three}. The decisive caveat, which we state rather than hide, is
that the bootstrap is self-consistent but not self-specifying: it does not
eliminate the non-computability identified throughout, it relocates it to a
determined place. The substrate question is, on the framework's own terms,
provably without a computable answer from inside---it is the irreducible residue,
the bump in the carpet that the whole construction has been moving rather than
removing. The leap from ``$E_{8}$ is self-dual'' to ``$E_{8}$ is the substrate''
is at present an emblem awaiting operational content (an $E_{8}$-equivariant
dynamics, or a state space built as a homogeneous space of $E_{8}$), and that
content, together with how a tick is generated at all, is deferred to the second
paper.

\paragraph{Dark matter (deferred).}
If gravity is the bookkeeping of the interaction network and dark matter interacts only
gravitationally, how does it participate in a network it does not exchange information with?
This is a genuine boundary of the programme, at which it may make no prediction, and Open
Problem~\ref{op:sterile-gravity} states it with the criteria that would close it. The prospect, to be developed in the second paper, is
that this is a feature rather than a defect. In the division of labour of
Section~\ref{sec:three-dimensions}, participation in the metric---counting loops
in the $S^{1}$ and $S^{3}$ directions---is universal, carried by anything with a
worldline and a tick-rate, whereas participation in the gauge filters, the
$S^{3}$ of the weak interaction and the $S^{7}$ of colour, is selective, engaged
only by charged matter. A node that counts loops but engages no filter
\emph{is} matter that gravitates without otherwise interacting, and the
framework's Class-4 a-chiral sterile neutrino (\S\ref{sec:generations}) is the candidate that
realises it --- recorded in the particle sections as Open Problem~\ref{op:sterile-gravity}, whose
tension (a gauge-sterile channel gravitating, when gravity is measurement bookkeeping) is exactly
the question named here. The condition for the candidacy to be coherent rather than circular is
that metric participation be established as prior to, and independent of, gauge participation;
making that priority precise is the work deferred, and it is the same work the open problem
requests.

\paragraph{CP violation, the arrow of time, and baryogenesis.}
CP violation appears to demand a time asymmetry owing nothing to the entropic
arrow, which would seem to require a preferred direction built into spacetime
itself. The framework dissolves the appearance rather than accommodating it. CP
violation needs an irreducible complex phase, the CKM phase $\delta$; by Spekkens'
result that negativity and contextuality are the same notion of nonclassicality
\autocite{spekkens2007toy}, a complex phase is a signature of contextuality; so CP
violation should be contextual, sensitive to the near environment of the decaying system --- the
same phase the interaction sections read as anchored by a fourth party
(\S\ref{sec:interactions}) --- which is exactly what Andrianov, Taron and Tarrach established for
neutral kaons in matter, where the environment breaks C, irreversibility breaks
T, and the effective parameters follow a Lindblad equation
\autocite{Andrianov2001}. The framework's reading is that the resulting
informational asymmetry \emph{constitutes} the arrow rather than presupposing it:
CP violation, entropy increase, and the direction of time are one informational
asymmetry seen three ways, not an independent second arrow, so the worry about a
direction built into spacetime does not arise. On this footing the contextual
dependence sets up a positive feedback loop for baryogenesis---an initial
fluctuation makes the environment slightly matter-dominated, contextual
enhancement raises the effective CP violation, more matter is produced, and the
loop runs---with $\delta$ a cosmological fixed point rather than a fundamental
constant. The honest limit of this is that a feedback loop explains amplification
and lock-in, not the sign selection or the baryon-number violation that
Sakharov's conditions still require \autocite{Sakharov1967}; the stronger claim
the framework is positioned to make, that the matter direction and the time
direction are co-selected by the same self-referential asymmetry, needs a source
of baryon-number violation and is left open.

\paragraph{Spatial infinitude and the status of inflation.}
There is no upper bound on non-computability, so the space of computations is
unbounded: the universe is spatially infinite, which leaves its spatial sections flat or of
negative curvature and excludes only the closed case. Infinitude on its own, however, does not
settle flatness, which is a fine-tuning of the Friedmann dynamics the framework has not yet
derived (Open Problem~\ref{op:comp-metric}); and flatness was in any case never inflation's
principal task. The horizon problem is, and an infinite space does nothing for it. The framework's actual lever against
inflation is the relational thesis of Section~\ref{sec:gr-nonlocality}: if
pre-spatial entanglement correlates regions before \textsc{fanout} spreads them
into spatial separation, then the observed large-scale homogeneity is fixed
relationally, before any question of causal contact, and the horizon problem
dissolves without an inflaton. This is a real argument, but it carries a real
bill: inflation also generates the nearly scale-invariant spectrum of primordial
fluctuations seen in the microwave background \autocite{Planck2020}, and any
account that dispenses with inflation owes an alternative origin for that
spectrum. That, not the horizon problem, is the hard empirical debt, and it is
named here rather than discharged.

\paragraph{Expansion and acceleration.}
If spacelike distance is the information separating worldlines, then as worldlines
accumulate distinct interaction histories they accumulate information-distance,
and the universe expands. The picture has the right exclusions built in: bound
systems interact continuously, share information, and so keep low
information-distance---they do not expand---while only worldlines that decouple
diverge, recovering expansion on cosmological scales alone. The feedback in which
greater separation means less light-cone overlap and so faster divergence,
$\mathrm{d}d/\mathrm{d}t \propto d$, yields exponential divergence, a de~Sitter
phase, and hence accelerating expansion without a separate dark-energy term
\autocite{Riess1998}. The limit to be honest about is that pure exponential
divergence is the late-time behaviour only; reproducing the full expansion history,
with its radiation- and matter-dominated eras, and fixing the Hubble rate
quantitatively, remains to be done.

\paragraph{The $S^{7}$ sector as internal structure.}
That $S^{7}$ cannot carry a metric does not make it disappear. Its closed
timelike curves carry the same non-computable information structure as those of
$S^{1}$ and $S^{3}$, but lacking associativity they support algebra---gauge
charge---in place of geometry---distance. This is why gauge symmetry has so often
looked like a set of compactified extra dimensions: the Kaluza--Klein and string
intuition is tracking something real in the $S^{7}$ structure, but misdiagnoses
it as spatial. In the framework these degrees of freedom have no size because
they have no metric; they are internal in the exact sense that they enter the
information-exchange that constitutes gauge charge while contributing nothing to
the loop-counting that constitutes distance. The resemblance to rolled-up
dimensions is therefore genuine and instructive, and the diagnosis that they are
algebraic rather than geometric is what separates this account from theories that
add real spatial dimensions to carry the same structure.

\paragraph{Loops, supertasks, and a second route to parallelisability.}
The argument of \S\ref{sec:three-dimensions} derives the parallelisable-sphere constraint from
closed timelike curves: a computation whose output is its own input requires information to be
transported consistently around a loop, and consistent transport selects the parallelisable spheres
$S^{1}$, $S^{3}$, $S^{7}$ \autocite{adams1960hopf}. There is a suggestion, which we record as a
prospect rather than a result, that the same constraint may govern the loops of perturbative quantum
field theory. The external legs of a Feynman diagram are on-shell classical boundary data; the
quantum content, and the divergences, enter with the internal loops, where the loop momentum is
integrated over all values and the action accumulated inside the loop is in this sense unbounded ---
a supertask, an infinite computation closed into a finite process. If such a loop is, in the
framework's terms, a closed computational path of the same kind as a closed timelike curve --- a
cycle carrying information not derivable from the data outside it --- then the parallelisability
constraint derived for closed timelike curves would apply to it, and the admissible interactions
would be constrained by the same $S^{1}$, $S^{3}$, $S^{7}$ structure that fixes the gauge group in
\S\ref{sec:entanglement-classes}. This would be a second and independent route to the Standard Model
content: not only \enquote{closed paths in spacetime select the parallelisable spheres} but \enquote{closed paths in interactions select them too}, unifying the spacetime arc and the interaction arc under one
principle. The decisive step, which we have not taken, is the identification of a field-theory loop
with a closed computational path in the precise sense the framework requires --- a loop integral is a
closed cycle in momentum space, a closed timelike curve a closed cycle in spacetime, and whether
these are the same kind of self-consistent closed structure is exactly what would have to be
established. It is offered as a direction, not a claim.

\paragraph{Does parallelisability determine renormalisability?}
A sharper and more self-contained question sits alongside the previous one, and we state it in the
hope that it will be taken up by those equipped for it. Renormalisability is the property that the
divergences of a quantum field theory can be consistently absorbed, leaving finite predictions; some
theories have it and some do not, and the standard criterion for which is technical (power-counting
of the coupling's dimension). The framework suggests a more structural reading. Parallelisable
spheres are precisely the spheres on which a global, consistent transport exists where it generically
does not; the non-parallelisable spheres are where the obstruction lives. If the divergence of a loop
is, in the framework's language, an obstruction to closing a computational path consistently, then
parallelisability would be exactly the structure that distinguishes the removable obstruction (a
renormalisable theory) from the irremovable one. It is at least a suggestive coincidence that the
renormalisable interactions of nature are the gauge theories built on $U(1)$, $SU(2)$ and $SU(3)$ ---
the groups the framework ties to the three parallelisable spheres --- and that gravity, which resists
renormalisation, is the sector \S\ref{sec:three-dimensions} ties to the geometry rather than the
algebra. Whether this is structural or accidental we do not know. The question --- whether
parallelisability, in the Hurwitz--Adams sense, is what determines renormalisability --- is posed in
the native language of the division-algebra approach to the Standard Model, and is one we think that
community is best placed to answer. We are careful to claim nothing more than that the question is
well-motivated: this is emphatically not a method for computing renormalised quantities, whose finite
remainders carry the precision agreement with experiment, but a conjecture about which theories admit
such a computation at all. It is, like the previous prospect, a direction with a possible path to an
answer rather than an answer, and it is left open deliberately.


\medskip
\noindent These mark the boundary of the programme as it stands. The substrate,
the generation of events, and the network status of dark matter are deferred to
the second paper; the contextual account of CP violation and baryogenesis, the
relational case against inflation, the information-divergence account of
expansion, and the algebraic reading of the $S^{7}$ sector are offered as
prospects, each with its outstanding debt named so that the line between argument
and aspiration stays visible.       <-- first consumer
%                                                          (S^7 exclusion)
%      % EDITED 2026-09-09 scripts/patch_spekkens_citekey_2026-09-09T1810.py
% =====================================================================
%  SECTION: Why Hopf Fibrations Belong in Particle Physics
%  (sec:hopf-justification). Drafted 2026-06-08; F1 arc added
%  prop:f1-dictionary and op:hopf-frobenius's discharge clause.
%  AUDITED 2026-07-06 (final pass, eighth stop).
%
%  PLACEMENT (reordered chain, JS 2026-07-06): position 8 --
%      % EDITED 2026-10-07 scripts/patch_infinite_not_flat_2026-10-07T1400.py
% EDITED 2026-09-27 scripts/patch_spacelike_inversion_2026-09-27T1800.py
% EDITED 2026-09-21 scripts/patch_op_sterile_gravity_boundary_2026-09-21T1300.py
% EDITED 2026-09-12 scripts/patch_spacetime_bloch_headline_2026-09-12T1400.py
% EDITED 2026-09-11 scripts/patch_splice_loops_prospects_2026-09-11T0900.py
% EDITED 2026-09-10 scripts/patch_cite_rubino_2026-09-10T0700.py
% EDITED 2026-09-10 scripts/patch_quoting_and_typography_2026-09-10T0620.py
% =====================================================================
%  SECTION: Spacetime from non-computability (sec:space-of-computations).
%  AUDITED 2026-07-06 (final pass, seventh stop).
%
%  PLACEMENT (reordered chain, JS 2026-07-06): position 7 --
%      \input{division_algebras_section_2026-07-03}   <-- hands off here
%      \input{spacetime_from_non-computability}
%      \input{hopf_justification_section_2026-06-08}
%  Division-algebras' closing now carries the S^7 hand-off sentence
%  (written 2026-07-06 after this file was read).
%
%  INPUT DEPENDENCIES (both held in the audited outputs):
%      entropic_gravity_family_paragraph_2026-07-03.tex  (sec:jacobson)
%      gr_nonlocality_revision_2026-06-13.tex            (sec:gr-nonlocality)
%  The 2026-07-04 additions (asymmetry paragraph; Penrose second route;
%  acceleration/Unruh arc) are FULLY PASTED; the separately uploaded
%  additions file is redundant and can be archived.
%
%  AUDIT RECORD: opener formalised (Russell typo; bare ref retargeted
%  to the freshly minted sec:complex-amplitudes); the "rotations of
%  S^7" wording made precise at BOTH occurrences (Moufang loop, not a
%  group; SO(8) disclaimer added -- referee hazard removed); two \\
%  paragraph-break bugs fixed; dark-matter paragraph retargeted to
%  sec:generations and stitched to op:sterile-gravity (the zoo's
%  demotion and this paragraph's priority condition are one question);
%  CP paragraph stitched to the fourth-party reading of
%  sec:interactions; stale mid-file mini-header retired
%  (sec:k3-sterile placeholder dies with it); the section's narrated
%  debts NUMBERED at last -- op:comp-metric and op:three-dimensions
%  minted in the Open-problems subsection, cited from the in-text
%  "named open problems" line and from the closing summary.
%
%  KEYS: adams1960hopf matches corpus. Unique to this section, presumed
%  in the .bib since the paper compiles: Spence-Brown-1969 (suggest
%  normalising the key spelling at leisure), VanRaamsdonk2010,
%  MaldacenaSusskind2013, MuellerMasanes2013, MasanesMuellerEtAl2014,
%  toffoli_2003, Chaitin1998, Penrose1996, HohnMueller2016,
%  GarnerMuellerDahlsten2017, Jacobson1995, Bekenstein1973, tHooft1993,
%  Susskind1995, spekkens2007toy, Andrianov2001, Sakharov1967,
%  Planck2020, Riess1998. Shared keys verified: small_2006, Penrose1994,
%  unruh1976notes (swap-flag stands), baez2002octonions, Wheeler1990,
%  furey2018three, DiBiagioRovelli2021, Rovelli1996a, Spekkens toy.
%  British English; \autocite; no bullets in prose.
% =====================================================================

\section{Spacetime from non-computability}
\label{sec:space-of-computations}

This section is a deliberate diversion: it takes the idea that non-computable steps are measured
by imaginary information (\S\ref{sec:complex-amplitudes}) and asks what geometry it buys.
Spencer-Brown records that Bertrand Russell thought well of the underlying move
\autocite{Spence-Brown-1969}; the reader may judge whether the framework spends it responsibly.

The programmes that derive spacetime from quantum mechanics---Van~Raamsdonk's
\enquote{building spacetime with entanglement} \autocite{VanRaamsdonk2010}, the
ER\,$=$\,EPR conjecture \autocite{MaldacenaSusskind2013}, and M\"uller's
information-theoretic recovery of three-dimensional space from the qubit
\autocite{MuellerMasanes2013}---take entanglement or information as primitive and
grow geometry from it. The route taken here is the same in spirit but inverted in
order: rather than positing spatial dimensions and forbidding superluminal
signalling, we \emph{define} spatial separation as a measure of non-computability
and read the dimension, the signature, and the Lorentz group off the result. The
proposal originates in the author's CASYS'05 paper \autocite{small_2006}, which
first suggested treating spacetime as a space of computations and attempted to
justify three dimensions on this basis; the present account tightens and extends
it. This is the framework's most structural and least formally complete arc; the
three-dimensionality result is a structural conclusion awaiting a formal proof,
and several of the links below are conjectural readings. Each is flagged as such,
so the reader sees exactly where the chain is load-bearing and where it is an
invitation.

\subsection{Timelike and spacelike distance as computation}
\label{sec:comp-distance}

A computation carries a state from axioms to a theorem in a number of steps.
Counting the computable steps between two states gives a real measure---time. A
physical clock counts completed irreversible loops, so proper time literally
counts loop completions and the metric $\mathrm{d}\tau^{2}=g_{\mu\nu}\,
\mathrm{d}x^{\mu}\mathrm{d}x^{\nu}$ is a loop-density tensor, the conversion from
coordinate intervals to loop counts; Toffoli's reading of the Lagrangian as a
count of action quanta is the same statement \autocite{toffoli_2003}.

Spacelike separation is the complementary, non-computable measure. Define the
non-computable distance between two propositions, relative to an axiom system, as
the number of bits that must be added to the axioms to make the second derivable.
A closed timelike curve is a computation whose output is its own input; by the
fixed-point theorem it has a unique consistent solution or none, and the
existence of that solution is one bit not derivable from any axiom system that
does not already contain the loop---one bit per curve, in the manner of Chaitin's
$\Omega$ \autocite{Chaitin1998}. The spacelike distance between two events is then
the number of such curves that would have to be closed to make them causally
connected. Because each added bit is a non-computable proposition, it is carried
by the imaginary part of the complex information measure
$S=-\ln\lvert p\rvert-\mathrm{i}\theta$ derived in
Section~\ref{sec:complex-amplitudes}: computable distance is counted in real
units, non-computable distance in imaginary ones. The Minkowski signature, with
spacelike intervals imaginary relative to timelike ones, is the immediate reading. The convention is not new to this paper: it was proposed in the author's CASYS paper two
decades ago, where counting non-causal steps in imaginary numbers was introduced and its formal
proof flagged, already then, as a key open problem \autocite{small_2006} ---
Open Problem~\ref{op:comp-metric} is that flag's direct descendant, stated now with the machinery
to be answerable.
That the resulting structure is a genuine metric space, and that the signature
follows rather than being fitted, are open problems---at present the
identification is motivated by reproducing what nature does, not proved.

%% VERIFY-SUPPORT 2026-09-27: conway_2003 is cited below for "any two octonions
%% generate an associative subalgebra" (Artin's theorem; a Moufang loop is
%% diassociative) and harvey1982calibrated for "a three-plane in Im O is
%% associative iff the associator of a basis vanishes, and the associative
%% three-planes are the imaginary parts of the quaternionic subalgebras"
%% (Harvey--Lawson, ch. IV). Neither passage was read in the session that wrote
%% this; both statements are checked numerically in Convention A by
%% scripts/antitelephone_network_associativity_2026-09-27T1900.py (PASS).
A reader used to the standard order of ideas should pause here, because the order has been
inverted, and the inversion is used below. In special relativity the tachyonic antitelephone is a
corollary: a signal faster than light, composed with a boost, would arrive before it was sent and
close a timelike loop, therefore there is no such signal \autocite{tolman1917theory,
benfordbooknewcomb1970tachyonic}. The metric comes first and the impossible telephone is a check
on it. Here the telephone comes first. Two observers are spacelike separated when an antitelephone
between them would close a loop; the loops such telephones would close, and carry nothing round,
are what spacelike separation is, their number is the distance, and the metric is read off from
them. Nothing is lost, since the corollary returns as a tautology: an antitelephone that carried
a record would close a causal loop, and its two ends would not have been spacelike after all.

What is gained is a definite meaning for a spatial dimension. Picture many observers, each holding
an antitelephone to each of the others. A message relayed from one observer to another through a
third composes two transports round two loops, and the network is consistent only if every relay
composes to a single transport whatever the bracketing of the relaying: a spatial dimension is a
family of loops whose transports compose associatively, which is a group law on a parallelisable
sphere. That is a condition a family of loops can fail, and \S\ref{sec:three-dimensions} shows
that the loops of the seventh sphere fail it. Two antitelephones always compose, because any two
octonions generate an associative subalgebra \autocite{conway_2003}; three compose only when their
imaginary parts lie in one quaternionic subalgebra, which is what the associative three-planes of
Harvey and Lawson are \autocite{harvey1982calibrated}, and a third link outside that subalgebra is
seen by the network exactly through its component outside it. The largest consistent network of
antitelephones inside $S^{7}$ is therefore an $S^{3}$, and there are three spatial directions.
When the seventh sphere is set aside there as a ruler, which is the octonionic spacetime of the
division-algebra literature, the claim is not that $\mathbb{R}^{8}$ has no metric, nor that the
ten-dimensional Minkowski space of $SL(2,\mathbb{O}) \cong \mathrm{Spin}(9,1)$ fails to exist
\autocite{baez2002octonions}; it is that the octonionic directions are not loops on which
antitelephones compose, and so are algebra rather than distance. In the standard order that
sentence has no content; in this one it is the reason space has three dimensions.

The split just defined explains the most familiar asymmetry in physics: that one can move back and
forth in space but not in time. \S\ref{sec:continuity} identifies the pre-measurement regime with
continuous reversibility --- its transformation group is connected, and every motion in it is
undoable. Spacelike separation, being the non-computable sector, is exactly the part of the
relational structure on which no record has been written; motion across it is motion within the
reversible regime, and it inherits that regime's continuity and its two-sidedness: a place can be
revisited because going and returning writes no record that forbids the return. Timelike separation
is the record chain itself --- proper time counts completed irreversible loops, as above --- and a
moment cannot be revisited because a moment \emph{is} a stratum of records, and returning would
require unwriting them, which is what \textsc{fanout} forbids (\S\ref{sec:fanout-boundary}). The
asymmetry between space and time is therefore not an extra postulate about the manifold: it is the
pre-/post-\textsc{fanout} asymmetry read as geometry --- space is what the reversible group feels
like from inside, time is what the record chain feels like --- and the arrow-of-time discussion
later in this section is the same asymmetry read thermodynamically. The identification is the
framework's structural reading of the split, offered in the same register as the rest of this
subsection.

\subsection{Three dimensions of space}
\label{sec:three-dimensions}

The closed timelike curves that define spacelike distance require information to
be parallel transported consistently around a loop, and consistent parallel
transport on a sphere requires it to be parallelisable. By Adams' theorem the
only parallelisable spheres are $S^{1}$, $S^{3}$ and $S^{7}$
\autocite{adams1960hopf}, the unit spheres of $\mathbb{C}$, $\mathbb{H}$ and
$\mathbb{O}$. Impose now the requirement that information arriving from outside
the causal zone must appear to possess its own causal history---that is, to be the
result of a sequence of information-preserving operations, hence a sequence of
group operations. The spheres $S^{1}$ and $S^{3}$ are groups under their own multiplication --- $U(1)$ and
$SU(2)$ --- and their operations preserve information; $S^{7}$'s multiplication is non-associative,
a Moufang loop rather than a group, so $S^{7}$ cannot serve as a ruler. (The rotation group
$SO(8)$ of the seven-sphere is of course a group; what fails is the sphere's \emph{own}
composition law, the one a ruler would use.) The same non-associativity
forbids a metric built from the sphere's own composition law outright: path
composition $d(A,C)$ through $B$ depends on the bracketing, so the triangle inequality is not merely violated but ill-defined, and
no geodesic structure exists. (What is untouched should be said: the Euclidean metric of
$\mathbb{R}^{8}$, and the ten-dimensional spacetime that $SL(2,\mathbb{O})$ furnishes, are metrics
of the ordinary kind; what fails is a ruler made of antitelephone loops, which is what a spatial
dimension is in this section, as \S\ref{sec:comp-distance} set out. Nor is the failure total:
relays between any two links compose, and among three they compose exactly when the links lie in
one quaternionic subalgebra, so the largest consistent network inside $S^{7}$ is an $S^{3}$
\autocite{conway_2003,harvey1982calibrated}.) Since the abelian $S^{1}$ alone cannot carry
non-abelian structure, spacelike distance must be three-dimensional. Dimension
three is an output, not an input \autocite{small_2006}.

The division of labour is then exact. $S^{1}$ ($\mathbb{C}$) contributes the time
dimension and the $U(1)$ of electromagnetism; $S^{3}$ ($\mathbb{H}$) contributes
the three spatial dimensions and the $SU(2)$ of the weak interaction; $S^{7}$
($\mathbb{O}$) contributes no geometry but, as a filter on what information may be
exchanged, the $SU(3)$ of the strong interaction. A sphere that cannot be a ruler
can still be a filter, which is why the strong force resists geometric
interpretation while electromagnetism and the weak force do not. This is the same
$\mathbb{C}\oplus\mathbb{C}^{3}$ split that yields $G_{\mathrm{SM}}$ in
Section~\ref{sec:entanglement-classes}, now read as the reason spacetime is $(3+1)$-%
dimensional. The conclusion converges with M\"uller and Masanes, who prove from
information-theoretic postulates, assuming no quantum theory, that minimal
direction-carrying systems with continuous reversible dynamics force $d=3$ and
quantum theory on two qubits \autocite{MuellerMasanes2013}, and with the result
that entanglement exists only for a three-dimensional Bloch ball
\autocite{MasanesMuellerEtAl2014}; the framework supplies the reason such systems
exist, Cayley--Dickson doubling from self-reference, while their work supplies
the rigorous dimensional uniqueness. The status of the framework's own route is a
structural result requiring formal proof: the one-bit-per-curve identification and
the precise reduction of the $S^{7}$ filter to $SU(3)$ rather than its full automorphism group
$G_{2}$ are numbered below (Open Problems~\ref{op:comp-metric} and~\ref{op:three-dimensions}).

A scoping remark, and then a second route to the same count. The computational model of spacetime
proposed here has many detailed parallels with the real thing, but a parallel is not an identity,
and the identity is not claimed in this paper: establishing it is work deferred to the second paper.
What can be claimed now is that the parallel is productive in a specific direction --- it suggests a
derivation of the dimension count carried out entirely within general relativity, with no
computational premise. Acceleration and gravitation tip the light cone, and the tipping can be
carried to the extreme: Penrose's discussion runs the cone through a full circle
\autocite[382]{Penrose1994}, at which point the tipped direction closes on itself --- which is to
say, the extreme of tipping is a closed timelike curve, the very object whose closure the
computational definition of \S\ref{sec:comp-distance} counts in bits. The two routes are one route
entered from opposite ends. Now demand of general relativity only this: that observers accelerating
with respect to one another, and with respect to unaccelerated observers, retain a consistent
conception of the distance between them. Each observer's tipping traces a path on a sphere of
directions, and mutual consistency requires the associated rotations to compose without regard to
bracketing --- to form a group. Compositions on $S^{1}$ and $S^{3}$ do, forming $U(1)$ and $SU(2)$ as above; compositions on
$S^{7}$ do not, the multiplication being non-associative, so observers tipped across an
$S^{7}$-worth of directions could not agree on distance even pairwise: the same failure that forbade
$S^{7}$ a metric above forbids it here a shared ruler. The largest sphere admitting a consistent
conception of distance under mutual acceleration is therefore $S^{3}$, and space has at most three
dimensions; with the floor argued above --- the abelian $S^{1}$ alone cannot carry the required
structure --- the count is exactly three. The status matches the subsection's. The tipping and its
extreme are general relativity's own \autocite{Penrose1994}; the consistency demand is
elementary; the associativity filter is the algebraic fact used throughout; the reading of composed
tippings as motion on the parallelisable spheres is the structural step, shared with the
computational route. The convergence of the two routes on the same endpoint is offered as evidence
for the parallel, while the identity it gestures at remains the second paper's burden.

\subsection{Recovering the Lorentz group}
\label{sec:lorentz}

A measurement changes what the observer knows, hence, relationally, the causal
structure: it tips the light cone. This is Penrose's association of cone-tipping
with state reduction \autocite{Penrose1996} run in the opposite causal direction,
reduction driving the tipping rather than the reverse, and the two are then one
operation in two charts---both described by $SL(2,\mathbb{C})$. The algebraic
backbone is the division-algebra correspondence
$SL(2,\mathbb{A})\cong\mathrm{Spin}(1,d{+}1)$ \autocite{baez2002octonions}: $\mathbb{R}$
gives $1{+}1$, $\mathbb{C}$ gives $3{+}1$, $\mathbb{H}$ gives $5{+}1$ and
$\mathbb{O}$ gives $9{+}1$ dimensions. The computability split that singles out a
preferred $\mathbb{C}$ fixes the fundamental state space as $\mathbb{C}^{2}$ and
hence $SL(2,\mathbb{C})=\mathrm{Spin}(1,3)$: the Lorentz group of $3{+}1$
spacetime. The qubit and Minkowski space are not merely both Euclidean-and-three%
-dimensional by coincidence; the Hermitian $2\times2$ matrices are
$\mathbb{R}^{1,3}$, $SL(2,\mathbb{C})$ acts on them as the Lorentz group, and the
qubit density operators fill the future cone---the correspondence M\"uller's group
develops operationally, recovering Lorentz transformations from quantum
communication \autocite{HohnMueller2016} and the complex and quaternionic qubit
from the relativity of simultaneity \autocite{GarnerMuellerDahlsten2017}.
In these coordinates the Bloch ball is the unit-trace slice of the cone and the Bloch sphere its
boundary rays, the celestial sphere on which $SL(2,\mathbb{C})$ acts by M\"obius
transformations; state reduction sends a point of the ball to a point of the sphere, and a boost
sends the sphere to itself, which is what lets the two be one operation in two charts.
The $SL(2,\mathbb{A})$ correspondence is rigorous; its identification with
measurement-as-tipping is the framework's structural reading.

One consequence of the identification has been implied throughout this section and deserves its own
paragraph, because the next paragraph turns on it. If measurement creates a relation between
observer and observed, and if --- relationally --- a new relation alters the causal structure, then
the tipping of the cone is not merely \emph{associated} with the creation of the relation: it is
what the creation of the relation is, seen geometrically. And a tipping of one's cone against the
congruence of one's neighbours already has a name in physics: acceleration. The framework's
statement, made explicit here, is that \emph{creating a new relation between observer and observed
is experienced as acceleration}. The resonance with the equivalence principle is exactly what it
should be: the principle makes acceleration and gravitation locally indistinguishable, Penrose's
association --- run in this section's direction, reduction driving the tipping --- makes state
reduction and cone-tipping one operation, and the framework's identification closes the triangle,
measurement and acceleration meeting where the principle says they must. The claim is structural:
the framework's reading of the cone-tipping established above, not an independent derivation.

Taken literally, the identification has a measurable face, and it is a familiar one: the Unruh
effect, read in this section's own ledger. Recall the definitions of \S\ref{sec:comp-distance}:
timelike distance counts computable steps, spacelike distance counts non-computable bits. What an
observer can compute --- what their \textsc{fanout}s can convert into classical records --- is fixed
by their causal structure, and the previous paragraph has just made that structure
relation-dependent. An accelerated observer's cones are tipped, and with them the boundary of the
recordable: relations that for the inertial observer lie wholly in the reversible, pre-record sector
--- the entanglement structure of the vacuum, non-computable for them, spacelike in the
computational sense --- fall, for the accelerated observer, partly on the recordable side. The mechanism is the one
Chaitin identified for formal systems \autocite{Chaitin1998}: a proposition undecidable relative to
an axiom system becomes decidable when information is added to the axioms --- and
\S\ref{sec:comp-distance} defined spacelike distance as exactly the bits whose addition closes the
gap. Creating a relation is adding information to the observer's axiom base, so timelike and
spacelike distance are converted into one another by the addition and subtraction of information,
and acceleration is the continuous version of that addition. A
relation newly recordable is a detection. Particles appear from nowhere: an accelerated detector in
the inertial vacuum registers a thermal bath at $T = \hbar a / 2\pi c k_{\mathrm{B}}$
\autocite{unruh1976notes}. Two guards keep the reading honest. Nothing here reclassifies the
invariant interval --- timelike and spacelike separation in the emergent metric are Lorentz
invariants, and no event crosses that classification; what is observer-relative is the computational
partition, which relations a given observer's horizon renders recordable. In the ledger of
\S\ref{sec:comp-distance} that partition \emph{is} the timelike/spacelike distinction, so in the
framework's own units timelike distance is \emph{rotated} into spacelike, and spacelike into
timelike, by a change of observer --- the tipping of the cone is literally that rotation --- and the
detector's clicks are the rotation made audible. And the effect is established physics
while the reading is the framework's, flagged structural. What gives the reading its content is a
parallel: the inertial and the accelerated observer stand to the vacuum exactly as Wigner and his
friend stand to the laboratory (\S\ref{sec:wigners-friend}) --- one holds a coherent relation and no
event, the other holds a record, and the facts are relative on both sides in precisely Rovelli's
sense \autocite{DiBiagioRovelli2021}. Since the previous paragraph identifies acceleration with
relation-creation, the parallel is not an analogy but the same structure in gravitational dress, the
horizon standing where the laboratory wall stood: \emph{the Unruh effect is Wigner's friend for
accelerated observers}. And two subsections of this section that seemed independent here meet: the
temperature Jacobson's construction requires across every local Rindler horizon
(\S\ref{sec:jacobson}) is exactly this temperature --- the $T$ of the Clausius relation is the $T$
of the tipped cone, supplied by the very relation-creation with which this subsection began.

\subsection{The relationalist commitment}
\label{sec:relationalism}

General relativity is relational: by the hole argument, spacetime points carry no
intrinsic identity, and the metric is fixed by the network of coincidences and
interactions, not by a pre-existing stage \autocite{Rovelli1996a}. The
computational picture is relational by construction and explains why it must be:
a background metric would require spatial distances to be pre-computable, and
non-computability forbids exactly that. General covariance is the statement that
distances have no pre-assigned values and depend on what information is actually
present. The metric is therefore dynamical not by stipulation but because it
measures a quantity that cannot be fixed in advance.

\subsection{Jacobson, holography, and imaginary information}
\label{sec:jacobson}

Jacobson derives the Einstein equation as an equation of state, from the Clausius
relation $\delta Q=T\,\mathrm{d}S$ across every local Rindler horizon together
with the Bekenstein proportionality of entropy to horizon area
\autocite{Jacobson1995, Bekenstein1973}. In the computational reading this is the
self-consistency condition on the loop-counting network: the Einstein equation is
the requirement that the bookkeeping of computation be internally consistent
across surfaces. The holographic principle, that the information in a region is
bounded by its bounding area \autocite{tHooft1993, Susskind1995}, then acquires a
natural reading here. A surface area is a squared distance; if spacelike distance
is measured in imaginary units of information, an area is measured in the square of
those units, that is in \emph{negative} information---entropy. The holographic
link between area and entropy thus requires the square root of negative
information, and that square root is precisely the imaginary unit of the complex
amplitude, the observer's self-referential phase. Read this way, Jacobson's
account of information crossing surfaces implicitly assigns an imaginary
information value to spacelike distance, and the $\mathrm{i}$ of the Minkowski
signature is the same $\mathrm{i}$ as the $\mathrm{i}$ of the quantum amplitude.
This identification is the most speculative step in the section: it is
structurally suggestive and it makes the holographic form unsurprising, but it is
a conjectural reading rather than a derivation, and is offered as such.

The conjecture has a two-way motivation that this paper's record should carry, because each
direction was on the table before the other was noticed. Forward: if spacelike distance is counted
in imaginary units of non-computable information, then any spacelike area is the product of two
imaginary lengths and is therefore \emph{negative} information --- which is to say entropy, by
Brillouin's identification \autocite{brillouin1956science} --- and a holographic principle arrives
automatically, the sign supplied by $i^{2}$ rather than by interpretation. Converse, and more
striking: in the holographic principle as ordinary physics already states it, surface area is
proportional to entropy, and entropy is negative information; so by implication ordinary physics
has been measuring spacelike separation in imaginary units of information all along, and the
convention of \S\ref{sec:comp-distance} is read off from established results rather than
invented. The same imaginary unit is already employed once in this section --- it is what makes
the signature of the interval Minkowski --- so the sign of the area law is its second job, not a
second assumption. What the two-way motivates it does not prove: the formal statement is exactly
the metric-and-signature half of Open Problem~\ref{op:comp-metric}, together with the
one-bit-per-bit-pair coefficient owed to the second paper. The stakes of the conjecture have
recently risen: Dorau and Much prove, with modular-theoretic rigour, that the semiclassical
Einstein equations follow from quantum relative entropy on a local horizon \emph{once the
entropy--area proportionality is assumed} \autocite{doraumuch2026relative} --- so what this
framework proposes to derive is precisely the remaining input of a derivation already carried
through at that standard.
\input{entropic_gravity_family_paragraph_2026-07-03}

\subsection{Non-locality already present in general relativity}
\label{sec:gr-nonlocality}
\input{gr_nonlocality_revision_2026-06-13}
\medskip
\noindent  The spacetime arc therefore yields, from one definition of spacelike
distance as non-computability: the signature, three spatial dimensions,
the Lorentz group, the dynamical relational metric, a reading of the
holographic principle, the presence of non-locality within general
relativity, the identification of relation-creation with acceleration
and of the Unruh effect with Wigner's friend, and the space--time
mobility asymmetry as the \textsc{fanout} asymmetry read as geometry. Of these,
the $(3+1)$ count and the Lorentz recovery are structural results that converge with independent
rigorous work and await their formal proof --- the count now numbered as
Open Problem~\ref{op:three-dimensions}; the metric-space status, the signature and the
one-bit-per-curve measure (together Open Problem~\ref{op:comp-metric}), the imaginary-information
reading of holography, and the general-relativistic non-locality, the acceleration and
Unruh readings, and the mobility asymmetry are conjectural, and are named here precisely so that the boundary between what is shown and what
is proposed is visible.

\subsection{Why the metric is not to be quantised}\label{sec:no-quantised-gravity}

A consequence has been implicit in everything above and deserves its plain statement. Causal
structure, on this framework's reading, is the bookkeeping of locked relations --- post-\textsc{fanout}
by construction --- while superposition is the signature of the pre-\textsc{fanout} regime. A
``superposition of causal structures'' is therefore not a difficult state but a category error: it
applies the pre-record operation to a record-side object, exactly as ``a superposition of
measurement outcomes'' does. Gravity lives in the post-\textsc{fanout} regime. The quantum layer
beneath it exists and is not in dispute; but it is \emph{pre-geometric} --- it carries amplitudes
over the formation of relations, not amplitudes over formed geometries --- so there is no
fundamental metric to quantise, and the standard programme of superposing spacetimes pursues an
object the ontology does not contain. One guard rail keeps the claim honest: collective
excitations of the locking pattern may perfectly well exist as effective quanta --- gravitons as
the phonons of the relational network --- and the thesis forbids fundamental metric superposition,
not effective field theory; this is Jacobson's equation-of-state moral \autocite{Jacobson1995},
made structural. The relation to gravitational-collapse programmes is an inversion: where
superposed geometries are there read as unstable \autocite{Penrose1996}, here they are read as
unformed --- what decays in those models is what locks in this one.

Three experimental postures follow, each at the status the reading carries. Laboratory
demonstrations of indefinite causal order \autocite{procopio2015experimental, rubino2017indefinite} are, on this
account, demonstrations of the nascent layer: the indefiniteness lives below the locking threshold
and definitises at it, so the framework adopts the programme and adds a boundary --- indefinite
causal order can never be amplified into records while remaining indefinite. Gravitationally
induced entanglement of the kind the Bose--Marletto--Vedral proposals seek
\autocite{bose2017spin, marlettovedral2017gravitationally} is \emph{predicted}, because the
pre-geometric layer is quantum --- while witnessing that layer, not a superposed metric: the
framework splits the experiment's usual dichotomy. And gravitationally driven anomalous collapse
is denied: locking is redundancy and thermodynamics, so decoherence experiments should find
standard rates and nothing gravity-specific. That no consistent quantisation of gravity exists was
the conclusion of the author's CASYS paper by a different route \autocite{small_2006}; the present
framework supplies the mechanism the conclusion then lacked.

\subsection{Open problems and prospects}
\label{sec:spacetime-open-problems}

The spacetime arc is the framework's frontier, and honesty about its edges
matters more here than anywhere else in the paper. We set out the genuine
problems and the prospects together, deferring the deepest to a second paper
rather than papering over them. Two of the arc's debts are well-posed now, and we number them.

\begin{openproblem}[The computational metric]\label{op:comp-metric}
Prove that the paired distance of \S\ref{sec:comp-distance} --- computable steps counted real,
non-computable bits counted imaginary, one bit per closed curve --- defines a genuine
pseudo-metric on the space of computations, with the Minkowski signature derived rather than read
off; equivalently, establish the one-bit-per-curve measure as a measure, and the signature as its
theorem. Until then the identification stands as motivated by reproducing what nature does, which
is the status flagged at its introduction.
\end{openproblem}

\begin{openproblem}[Dimension three as a theorem]\label{op:three-dimensions}
Turn the two convergent routes of \S\ref{sec:three-dimensions} into a proof: formalise the
parallelisability-plus-associativity filter (the computational route) and the mutual-acceleration
consistency demand (the general-relativistic route) so that $d = 3$ follows as a theorem, and
include the reduction clause --- that the $S^{7}$ filter acts through $SU(3)$, the stabiliser of
the selected complex direction, rather than through its full automorphism group $G_{2}$. The
M\"uller--Masanes uniqueness results \autocite{MuellerMasanes2013, MasanesMuellerEtAl2014} are
the rigorous convergents any proof must meet.
\end{openproblem}

\paragraph{The substrate and the generation of events (deferred).}
A computation runs on something; a network of events presupposes that events
occur. Both are the same objection in two forms---the network's hardware and its
clock---and the original CASYS'05 statement of the programme raised the
implementation regress and set it aside \autocite{small_2006}. The framework's
response is that a self-referential structure is the only kind that need not
presuppose a meta-substrate: it generates its own. The natural algebraic emblem
is $E_{8}$, the unique simple Lie group whose adjoint representation is also its
smallest faithful one, so that what it is and what it acts on coincide; the $240$
unit integral octonions that generate it are exactly the $E_{8}$ root system, so
the self-action is welded to the division-algebra spine of
Section~\ref{sec:entanglement-classes}. This is Wheeler's self-excited circuit
\autocite{Wheeler1990} given an algebraic body, and it converges with Furey's use
of an algebra acting on itself to generate particle content
\autocite{furey2018three}. The decisive caveat, which we state rather than hide, is
that the bootstrap is self-consistent but not self-specifying: it does not
eliminate the non-computability identified throughout, it relocates it to a
determined place. The substrate question is, on the framework's own terms,
provably without a computable answer from inside---it is the irreducible residue,
the bump in the carpet that the whole construction has been moving rather than
removing. The leap from ``$E_{8}$ is self-dual'' to ``$E_{8}$ is the substrate''
is at present an emblem awaiting operational content (an $E_{8}$-equivariant
dynamics, or a state space built as a homogeneous space of $E_{8}$), and that
content, together with how a tick is generated at all, is deferred to the second
paper.

\paragraph{Dark matter (deferred).}
If gravity is the bookkeeping of the interaction network and dark matter interacts only
gravitationally, how does it participate in a network it does not exchange information with?
This is a genuine boundary of the programme, at which it may make no prediction, and Open
Problem~\ref{op:sterile-gravity} states it with the criteria that would close it. The prospect, to be developed in the second paper, is
that this is a feature rather than a defect. In the division of labour of
Section~\ref{sec:three-dimensions}, participation in the metric---counting loops
in the $S^{1}$ and $S^{3}$ directions---is universal, carried by anything with a
worldline and a tick-rate, whereas participation in the gauge filters, the
$S^{3}$ of the weak interaction and the $S^{7}$ of colour, is selective, engaged
only by charged matter. A node that counts loops but engages no filter
\emph{is} matter that gravitates without otherwise interacting, and the
framework's Class-4 a-chiral sterile neutrino (\S\ref{sec:generations}) is the candidate that
realises it --- recorded in the particle sections as Open Problem~\ref{op:sterile-gravity}, whose
tension (a gauge-sterile channel gravitating, when gravity is measurement bookkeeping) is exactly
the question named here. The condition for the candidacy to be coherent rather than circular is
that metric participation be established as prior to, and independent of, gauge participation;
making that priority precise is the work deferred, and it is the same work the open problem
requests.

\paragraph{CP violation, the arrow of time, and baryogenesis.}
CP violation appears to demand a time asymmetry owing nothing to the entropic
arrow, which would seem to require a preferred direction built into spacetime
itself. The framework dissolves the appearance rather than accommodating it. CP
violation needs an irreducible complex phase, the CKM phase $\delta$; by Spekkens'
result that negativity and contextuality are the same notion of nonclassicality
\autocite{spekkens2007toy}, a complex phase is a signature of contextuality; so CP
violation should be contextual, sensitive to the near environment of the decaying system --- the
same phase the interaction sections read as anchored by a fourth party
(\S\ref{sec:interactions}) --- which is exactly what Andrianov, Taron and Tarrach established for
neutral kaons in matter, where the environment breaks C, irreversibility breaks
T, and the effective parameters follow a Lindblad equation
\autocite{Andrianov2001}. The framework's reading is that the resulting
informational asymmetry \emph{constitutes} the arrow rather than presupposing it:
CP violation, entropy increase, and the direction of time are one informational
asymmetry seen three ways, not an independent second arrow, so the worry about a
direction built into spacetime does not arise. On this footing the contextual
dependence sets up a positive feedback loop for baryogenesis---an initial
fluctuation makes the environment slightly matter-dominated, contextual
enhancement raises the effective CP violation, more matter is produced, and the
loop runs---with $\delta$ a cosmological fixed point rather than a fundamental
constant. The honest limit of this is that a feedback loop explains amplification
and lock-in, not the sign selection or the baryon-number violation that
Sakharov's conditions still require \autocite{Sakharov1967}; the stronger claim
the framework is positioned to make, that the matter direction and the time
direction are co-selected by the same self-referential asymmetry, needs a source
of baryon-number violation and is left open.

\paragraph{Spatial infinitude and the status of inflation.}
There is no upper bound on non-computability, so the space of computations is
unbounded: the universe is spatially infinite, which leaves its spatial sections flat or of
negative curvature and excludes only the closed case. Infinitude on its own, however, does not
settle flatness, which is a fine-tuning of the Friedmann dynamics the framework has not yet
derived (Open Problem~\ref{op:comp-metric}); and flatness was in any case never inflation's
principal task. The horizon problem is, and an infinite space does nothing for it. The framework's actual lever against
inflation is the relational thesis of Section~\ref{sec:gr-nonlocality}: if
pre-spatial entanglement correlates regions before \textsc{fanout} spreads them
into spatial separation, then the observed large-scale homogeneity is fixed
relationally, before any question of causal contact, and the horizon problem
dissolves without an inflaton. This is a real argument, but it carries a real
bill: inflation also generates the nearly scale-invariant spectrum of primordial
fluctuations seen in the microwave background \autocite{Planck2020}, and any
account that dispenses with inflation owes an alternative origin for that
spectrum. That, not the horizon problem, is the hard empirical debt, and it is
named here rather than discharged.

\paragraph{Expansion and acceleration.}
If spacelike distance is the information separating worldlines, then as worldlines
accumulate distinct interaction histories they accumulate information-distance,
and the universe expands. The picture has the right exclusions built in: bound
systems interact continuously, share information, and so keep low
information-distance---they do not expand---while only worldlines that decouple
diverge, recovering expansion on cosmological scales alone. The feedback in which
greater separation means less light-cone overlap and so faster divergence,
$\mathrm{d}d/\mathrm{d}t \propto d$, yields exponential divergence, a de~Sitter
phase, and hence accelerating expansion without a separate dark-energy term
\autocite{Riess1998}. The limit to be honest about is that pure exponential
divergence is the late-time behaviour only; reproducing the full expansion history,
with its radiation- and matter-dominated eras, and fixing the Hubble rate
quantitatively, remains to be done.

\paragraph{The $S^{7}$ sector as internal structure.}
That $S^{7}$ cannot carry a metric does not make it disappear. Its closed
timelike curves carry the same non-computable information structure as those of
$S^{1}$ and $S^{3}$, but lacking associativity they support algebra---gauge
charge---in place of geometry---distance. This is why gauge symmetry has so often
looked like a set of compactified extra dimensions: the Kaluza--Klein and string
intuition is tracking something real in the $S^{7}$ structure, but misdiagnoses
it as spatial. In the framework these degrees of freedom have no size because
they have no metric; they are internal in the exact sense that they enter the
information-exchange that constitutes gauge charge while contributing nothing to
the loop-counting that constitutes distance. The resemblance to rolled-up
dimensions is therefore genuine and instructive, and the diagnosis that they are
algebraic rather than geometric is what separates this account from theories that
add real spatial dimensions to carry the same structure.

\paragraph{Loops, supertasks, and a second route to parallelisability.}
The argument of \S\ref{sec:three-dimensions} derives the parallelisable-sphere constraint from
closed timelike curves: a computation whose output is its own input requires information to be
transported consistently around a loop, and consistent transport selects the parallelisable spheres
$S^{1}$, $S^{3}$, $S^{7}$ \autocite{adams1960hopf}. There is a suggestion, which we record as a
prospect rather than a result, that the same constraint may govern the loops of perturbative quantum
field theory. The external legs of a Feynman diagram are on-shell classical boundary data; the
quantum content, and the divergences, enter with the internal loops, where the loop momentum is
integrated over all values and the action accumulated inside the loop is in this sense unbounded ---
a supertask, an infinite computation closed into a finite process. If such a loop is, in the
framework's terms, a closed computational path of the same kind as a closed timelike curve --- a
cycle carrying information not derivable from the data outside it --- then the parallelisability
constraint derived for closed timelike curves would apply to it, and the admissible interactions
would be constrained by the same $S^{1}$, $S^{3}$, $S^{7}$ structure that fixes the gauge group in
\S\ref{sec:entanglement-classes}. This would be a second and independent route to the Standard Model
content: not only \enquote{closed paths in spacetime select the parallelisable spheres} but \enquote{closed paths in interactions select them too}, unifying the spacetime arc and the interaction arc under one
principle. The decisive step, which we have not taken, is the identification of a field-theory loop
with a closed computational path in the precise sense the framework requires --- a loop integral is a
closed cycle in momentum space, a closed timelike curve a closed cycle in spacetime, and whether
these are the same kind of self-consistent closed structure is exactly what would have to be
established. It is offered as a direction, not a claim.

\paragraph{Does parallelisability determine renormalisability?}
A sharper and more self-contained question sits alongside the previous one, and we state it in the
hope that it will be taken up by those equipped for it. Renormalisability is the property that the
divergences of a quantum field theory can be consistently absorbed, leaving finite predictions; some
theories have it and some do not, and the standard criterion for which is technical (power-counting
of the coupling's dimension). The framework suggests a more structural reading. Parallelisable
spheres are precisely the spheres on which a global, consistent transport exists where it generically
does not; the non-parallelisable spheres are where the obstruction lives. If the divergence of a loop
is, in the framework's language, an obstruction to closing a computational path consistently, then
parallelisability would be exactly the structure that distinguishes the removable obstruction (a
renormalisable theory) from the irremovable one. It is at least a suggestive coincidence that the
renormalisable interactions of nature are the gauge theories built on $U(1)$, $SU(2)$ and $SU(3)$ ---
the groups the framework ties to the three parallelisable spheres --- and that gravity, which resists
renormalisation, is the sector \S\ref{sec:three-dimensions} ties to the geometry rather than the
algebra. Whether this is structural or accidental we do not know. The question --- whether
parallelisability, in the Hurwitz--Adams sense, is what determines renormalisability --- is posed in
the native language of the division-algebra approach to the Standard Model, and is one we think that
community is best placed to answer. We are careful to claim nothing more than that the question is
well-motivated: this is emphatically not a method for computing renormalised quantities, whose finite
remainders carry the precision agreement with experiment, but a conjecture about which theories admit
such a computation at all. It is, like the previous prospect, a direction with a possible path to an
answer rather than an answer, and it is left open deliberately.


\medskip
\noindent These mark the boundary of the programme as it stands. The substrate,
the generation of events, and the network status of dark matter are deferred to
the second paper; the contextual account of CP violation and baryogenesis, the
relational case against inflation, the information-divergence account of
expansion, and the algebraic reading of the $S^{7}$ sector are offered as
prospects, each with its outstanding debt named so that the line between argument
and aspiration stays visible.
%      % EDITED 2026-09-09 scripts/patch_spekkens_citekey_2026-09-09T1810.py
% =====================================================================
%  SECTION: Why Hopf Fibrations Belong in Particle Physics
%  (sec:hopf-justification). Drafted 2026-06-08; F1 arc added
%  prop:f1-dictionary and op:hopf-frobenius's discharge clause.
%  AUDITED 2026-07-06 (final pass, eighth stop).
%
%  PLACEMENT (reordered chain, JS 2026-07-06): position 8 --
%      \input{spacetime_from_non-computability}
%      \input{hopf_justification_section_2026-06-08}
%      \input{entanglement_classes_and_gauge_groups_2026-06-08}
%  No adjacency phrasing in the prose: position-portable. All
%  cross-refs live -- sec:fanout-direction DRAFTED 2026-07-06 (formerly the one undrafted
%  section, by design).
%
%  AUDIT RECORD: prop:gsm, prop:termination and prop:f1-dictionary
%  verified against every quoting consumer (statements canonical);
%  honest-debt subsection names both the selection assumption and the
%  mereology debt; the epistemic two-climbs disambiguation and the
%  compositional rung-evidence verified; sec:canonical-orbits live;
%  file pure ASCII. Fixes: op:hopf-frobenius's cross-link with
%  op:sigma-functor made TWO-WAY (the "four ways" line now names the
%  op); both remaining functor conditions ANCHORED --- the quaternionic
%  to the established fibre physics of sec:wf-reconciliation (residue:
%  clutching map vs strong complementarity), the octonionic to the
%  proved canonical-orbit images (umbrella label app:ghz-w-proofs
%  minted at the appendix for the purpose). Statuses unchanged.
%  British English; \autocite; no bullets in prose.
% =====================================================================

\section{Why Hopf Fibrations Belong in Particle Physics}
\label{sec:hopf-justification}

Invoking Hopf fibrations in a discussion of the Standard Model invites an immediate and fair
objection: that the move is decorative, an analogy between two pieces of attractive mathematics
with no load-bearing connection between them. The purpose of this section is to dispose of that
objection before the construction is put to work. The connection is not an analogy. It is
carried by two independent and independently published theorems, and the structures it relies
upon are not selected for convenience but compelled by a pair of classification results. We take
as already established the correspondence between qubit entanglement and the Hopf fibrations
themselves --- the subject of a mature literature
\autocite{mosseri2001hopf,bernevigchen2003geometry} --- and do not reproduce it here. What this
section addresses is the prior question of why that correspondence is entitled to carry physical
content at all.

\subsection{Two theorems, joined end to end}

The bridge from qubit entanglement to the Standard Model is built from two distinct results,
neither of which is original to the present framework and neither of which depends on the other.

The first is geometric. The unit sphere in the Hilbert space of $n$ qubits is the total space of
a Hopf fibration: $S^3 \to S^2$ for one qubit, $S^7 \to S^4$ for two, and $S^{15} \to S^8$ for
three. For two and three qubits the fibration is \emph{entanglement-sensitive} --- separability
is precisely the condition under which the base degenerates to a product of single-qubit Bloch
spheres and the additional dimensions associated with the higher imaginary directions collapse
\autocite{mosseri2001hopf,bernevigchen2003geometry}. This is a statement about the geometry of
state spaces and requires no physical interpretation whatever.

The second is algebraic. Singling out a preferred complex structure on the octonions ---
equivalently, fixing the split $\mathbb{O} \cong \mathbb{C} \oplus \mathbb{C}^3$ --- leaves a
residual symmetry, and that residual symmetry is exactly the gauge group of the Standard Model.
This was established through the exceptional Jordan algebra by Dubois-Violette and Todorov
\autocite{duboisviolette2016exceptional,duboisviolettetodorov2019exceptional,todorovduboisviolette2018deducing}
and given its sharpest spinorial form by Krasnov \autocite{krasnov2021so9}, who showed that the
relevant group is the stabiliser of a complex structure inside $\mathrm{Spin}(9)$. We record the
result in the form in which it enters the present work.

\begin{proposition}[Dubois-Violette--Todorov; Krasnov]\label{prop:gsm} Let $J_R$ be the complex
structure on the space $\mathbb{O}^2$ of $\mathrm{Spin}(9)$ spinors determined by a choice of
unit imaginary octonion, equivalently the split
$\mathbb{O} \cong \mathbb{C} \oplus \mathbb{C}^3$. The subgroup of $\mathrm{Spin}(9)$ commuting
with $J_R$ is \[ G_{\mathrm{SM}} \;=\; \frac{SU(3) \times SU(2) \times U(1)}{\mathbb{Z}_6}, \]
acting on the singled-out $\mathbb{C}$ summand as a colour singlet of hypercharge $-1$ and on
the $\mathbb{C}^3$ summand as a colour triplet of hypercharge $\tfrac{1}{3}$, both as $SU(2)$
doublets --- the quantum numbers of one left-handed generation.
\end{proposition}

The two results are independent, and the join between them is the only place a physical theory
must do work. Each result is a theorem; what neither supplies is a reason why nature should fix
one complex direction within $\mathbb{O}$ rather than treating all directions symmetrically. We
return to this --- the genuine load-bearing assumption --- in \S\ref{sec:hopf-honest-debt}, and
resolve it, in candidate form, in \S\ref{sec:fanout-direction}.

\subsection{The ladder is forced, not chosen}

It might still be objected that selecting the octonions, or the Hopf fibrations, is itself a
modelling choice carrying as much latitude as the choice of a grand-unified group. It is not,
and the reason is a pair of classification theorems that close off the alternatives entirely.

By Hurwitz's theorem the only normed division algebras are $\mathbb{R}$, $\mathbb{C}$,
$\mathbb{H}$ and $\mathbb{O}$ \autocite{hurwitz1898composition,baez2002octonions}. By Adams's
resolution of the Hopf-invariant-one problem --- together with the parallelisability results of
Bott--Milnor and Kervaire --- the only parallelisable spheres are $S^0$, $S^1$, $S^3$ and $S^7$,
so that there exist exactly three nontrivial Hopf fibrations, in bijection with $\mathbb{C}$,
$\mathbb{H}$ and $\mathbb{O}$
\autocite{adams1960hopf,bottmilnor1958parallelizability,kervaire1958nonparallelizability}. These
are not heuristics; they are among the more durable theorems in topology and algebra. Their
force here is that, once a framework is committed to amplitudes valued in a normed division
algebra --- a commitment that in the present construction is made far upstream, by the
self-referential argument of \S\ref{sec:self-reference} and \S\ref{sec:division-algebras}, and
not by fiat at this stage --- the ladder of Hopf fibrations is not selected but compelled, and
it terminates of its own accord.

\begin{proposition}[Termination]\label{prop:termination} The Cayley--Dickson tower of normed
division algebras and the tower of parallelisable spheres both terminate at real dimension
eight. Consequently there are exactly three nontrivial Hopf fibrations, the largest being the
octonionic $S^{15} \to S^8$, and the construction admits no extension to a fourth case.
\end{proposition}

The termination is worth dwelling on, because it is unusual for a model of this kind to predict
its own boundary. The next rung of the Cayley--Dickson sequence, the sedenions, possesses zero
divisors; the Hopf map can no longer be globally defined; there is no fourth fibration. This is
not a limitation imposed from outside but a structural fact, and it carries a physical
consequence that we develop later: there is no room in the construction for a fourth gauge force
or a fifth interaction, so that the discovery of either would count as evidence against the
framework rather than as an invitation to extend it. The endpoint is a theorem, not a
stipulation.

\subsection{Independent arrival as corroboration}

The most persuasive evidence that the destination is not an artefact of the route is that it is
reached by routes sharing none of the same machinery. Szangolies
\autocite{szangolies2025standardmodel} reaches $G_{\mathrm{SM}}$ from the emergent-gravity and
Kaluza--Klein programmes, reading the singled-out complex direction as a dimensional reduction
of a $9+1$-dimensional spacetime; the present framework reaches it from self-reference
(\S\ref{sec:self-reference}); and Furey and Gillard--Gresnigt reach the same gauge group, and
the same three-generation structure, through the minimal left ideals of the Dixon and Clifford
algebras, without invoking Hopf fibrations at all
\autocite{furey2018three,gillardgresnigt2019three}; the relation between their minimal ideals
and the entanglement classes used here is taken up in \S\ref{sec:canonical-orbits}. The
octonionic route to particle physics is in any case an old one, going back to the use of
$\mathbb{O}$ to model the strong force \autocite{gunaydingursey1974quark}.

That programmes with disjoint starting points converge on
$SU(3) \times SU(2) \times U(1)/\mathbb{Z}_6$ and on the same quark--lepton division is the
coincidence Szangolies himself remarks upon: the characteristic structure of the Standard Model,
routinely regarded as arbitrary, turns up independently in unrelated settings, and its appearing
twice would then be doubly odd unless the two settings are in fact related. We read the
convergence the same way --- as evidence that the division-algebraic description is singled out
by the mathematics rather than imposed by the modeller, and therefore that
Proposition~\ref{prop:gsm} is structural rather than contrived.

\subsection{The honest debt}
\label{sec:hopf-honest-debt}

It would be misleading to leave the impression that this route is free of assumptions, so we
state plainly the one assumption that all of these programmes share. Proposition~\ref{prop:gsm}
requires that \emph{some} complex direction in $\mathbb{O}$ be singled out, and the uniqueness
of that choice is not supplied by the theorem. Krasnov parametrises the complex structure $J_R$
by an arbitrary unit imaginary octonion, and the residual group is the same for every such
choice because $G_2 = \mathrm{Aut}(\mathbb{O})$ acts transitively on the unit imaginary
octonions. The theorem therefore guarantees the gauge group \emph{given} that a direction is
fixed, while remaining silent on \emph{why} one is fixed. Szangolies treats this as a free
choice and flags it, in his own conclusions, as an open element of the construction --- one to
which several inequivalent answers are available.

This, and not the appeal to Hopf fibrations, is the genuine load-bearing assumption of the whole
approach. The Hopf fibrations are forced (Proposition~\ref{prop:termination}); the singling-out
is not. The present framework offers a candidate resolution that the prior literature lacks: the
act of measurement, modelled as \textsc{fanout}, selects the direction physically, and the
transitivity of $G_2$ then guarantees that the gauge group is independent of which direction the
measurement happens to select. This argument is given in \S\ref{sec:fanout-direction}; we flag
here, at the appropriate epistemic status, that it is structural rather than fully proved, and
that the residual question --- why a measurement selects \emph{exactly one} imaginary direction
rather than merely \emph{permitting} a selection --- remains open. It is now priced as Open
Problem~\ref{op:fanout-selection}.

A second debt is shared with every programme of this character and is not closed by anything
said here: why the abstract quantum state should admit a tensor-product decomposition into
qubits in the first place. This is the quantum factorisation, or mereology, problem. The
entanglement-sensitivity on which the construction turns is meaningful only once such a
decomposition is fixed, and we treat the decomposition as a necessary input rather than a
derived result. Naming it is preferable to letting it pass unremarked.

\subsection{The termination read epistemically}
\label{sec:hopf-information-bound}

The termination of the tower (Proposition~\ref{prop:termination}) has so far been stated
algebraically: the Cayley--Dickson sequence acquires zero divisors at the sedenions and the Hopf
map can no longer be globally defined. There is a second reading of the same boundary, epistemic
rather than algebraic, which we offer as a motivated conjecture rather than a settled result,
because it connects the tower's endpoint to a bound that arises independently in the foundations
of measurement.

A particle, in this framework, is a reified relation between events, and its internal content is
carried by the octonionic $A,B,C$ register --- the $S^{7}$ sector serving as internal degrees of
freedom in the channel that connects the events (\S\ref{sec:space-of-computations}). The
termination of the tower is then a statement about that content: a single such relation can
support genuine entanglement among at most three parties, because there is no fourth fibration
to carry a fourth independent party within the same irreducible relation. Every particle is, on
this reading, a \emph{completed} octonionic-level relation --- topped out at three parties not
by circumstance but because the tower admits no further rung. The information that can be
encoded in one relation-between-events, and hence extracted by one measurement of it, is
correspondingly bounded: it is the information of a completed three-party register, and no
measurement of a single particle can yield more.

This is exactly the kind of bound that the epistemic-restriction programmes reach from the other
direction. Spekkens' toy theory is built on a knowledge-balance principle that caps the
information obtainable about a system \autocite{spekkens2007toy}, and Szangolies derives such an
epistemic horizon not by positing it but from the same self-referential non-computability on
which the present tower rests \autocite{szangolies_j_2018}. That the two accounts --- the
algebraic termination of the Hopf tower and the epistemic bound on measurement information ---
should coincide is, we suggest, not a numerical accident but a shared root: both are
consequences of the self-referential non-computability of \S\ref{sec:self-reference}, wearing an
algebraic face in the one case and an epistemic face in the other. The Hurwitz--Adams ceiling
and the knowledge-balance ceiling would then be one limit seen twice, and Szangolies' programme,
which reaches $G_{\mathrm{SM}}$ by an independent route (\S\ref{sec:hopf-honest-debt}) and the
epistemic horizon from the same non-computable source, shares with this framework not merely a
conclusion but its origin.

The reading must be kept distinct from a different climb, with which it would otherwise be
conflated. The Wigner's-friend regress of \S\ref{sec:wigners-friend} ascends without end: each
observer may be measured by a further observer, each carrying a fresh gauge-private circle, and
that chain has no top --- it is the bump in the carpet, unbounded in length. The termination
bounds something else: not the length of the chain of observers, but the internal structure of
each single relation, each particle, that the chain is built from. The regress is unbounded
precisely because each of its links is a bounded, completed three-party object;
termination-of-each-link and non-termination-of-the-chain are not in tension but complementary,
the former individuating the links of the latter. What the bound caps is the information in one
particle; what the regress leaves uncapped is the number of particles and observers that can be
chained. The precise identification of ``the information of a completed three-party register''
with the formal measures of the epistemic-restriction literature is the technical step that
would turn this conjecture into a theorem, and we flag it as owed; the structural claim --- that
the tower's endpoint and the measurement bound share the non-computable root --- is what we
advance here.

\subsection{The termination read compositionally}
\label{sec:hopf-compositional}

There is a third termination theorem, proved in a community that has so far had no contact with
the two joined above. Coecke and Kissinger showed that the GHZ and W classes --- realised as the
commutative special and the anti-special Frobenius algebras of categorical quantum mechanics ---
are \emph{compositionally complete}: every $N$-qubit entanglement class under SLOCC is generated
by composing these tripartite primitives with bipartite entanglement and product states
\autocite{coeckekissinger2010compositional}. Above three parties nothing new appears; the
entanglement alphabet closes at exactly the point where the fibration tower terminates. The
mutual silence of the two literatures is itself striking: Coecke and Kissinger make no reference
to Hopf fibrations or division algebras, and the fibration literature on entanglement
\autocite{mosseri2001hopf,bernevigchen2003geometry} nowhere cites compositional completeness.
Two classifications, in different mathematics, draw the same line in the same place, unaware of
one another. That silence has, in the months this paper was in draft, just begun to break from
the categorical side, with a projective, base-level reading of the same calculus
\autocite{burtonanwar2026meromorphic}; we engage it below.

The framework's conjecture is that this is one termination wearing a third face: the alphabet
has the letters it has --- Bell, GHZ, W --- because the parallelisable spheres have the
dimensions they have. There is partial evidence at every rung. At one qubit, the commutative
special Frobenius algebras on $\mathbb{C}^{2}$ are classified by orthonormal bases --- points of
$S^{2}$, the base of the first fibration --- so that the copied basis plays the role of a
section of the bundle and the spider theorem expresses its idempotence. At two qubits, the
$SU(2)$ fibre of $S^{7} \to S^{4}$ carries exactly the entanglement that survives local
operations, and L\'{e}vay's result that the bundle's natural connection is the $SU(2)$
Yang--Mills instanton \autocite{levay2004geometry} stands on the geometric side where the
interaction of two Frobenius structures stands on the categorical one. At three qubits, the
split of D\"{u}r, Vidal and Cirac \autocite{durvidalcirac2000three} mirrors the algebra exactly:
GHZ is genuinely tripartite, demanding the full non-associative octonion, while W is
pairwise-robust and housed in the associative subalgebra $\mathbb{H} \subset \mathbb{O}$ --- the
special/anti-special Frobenius distinction in geometric dress.

The first of the four conditions the correspondence requires can be discharged at once, and we
do so here, because a problem whose first step has been walked is a different kind of invitation
from one merely posed.

\begin{proposition}[The one-qubit dictionary]\label{prop:f1-dictionary} The $\dagger$-special
commutative Frobenius algebras on $\mathbb{C}^{2}$ correspond bijectively to
\emph{lifted antipodal skeletons} of the complex Hopf fibration $\pi \colon S^{3} \to S^{2}$:
unordered antipodal pairs $P \subset S^{2}$ together with a section of $\pi$ over $P$. The
copyable states of a structure are exactly the section's image, so the lifts are recorded
faithfully --- the projectively trivial centre of $U(2)$ acts non-trivially on structures ---
and the correspondence is $U(2)$-equivariant, the ordered structures forming a single free and
transitive $U(2)$-orbit. The channel a structure induces on density matrices is dephasing in its
skeleton and is independent of the section: the fibre freedom, a torsor for
$\mathrm{Map}(P, U(1))$, is exactly the data absent from the classical record a structure
induces.
\end{proposition}

\begin{proof}
By the classification of Coecke, Pavlovic and Vicary, such structures correspond to orthonormal
bases, recovered as the sets of copyable states \autocite{coeckepavlovicvicary2013orthogonal}.
An orthonormal basis of $\mathbb{C}^{2}$ is two orthogonal rays --- an antipodal pair on the
Bloch sphere --- together with a unit representative of each, which is precisely a section of
$\pi$ over the pair; and $\delta(e^{i\theta}x) = e^{i\theta}\, x \otimes x$ against
$(e^{i\theta}x) \otimes (e^{i\theta}x) = e^{2i\theta}\, x \otimes x$ pins the phases.
Equivariance, the action of the centre, the orbit structure and the induced channel are direct
computations; an executable verification of every finite check accompanies the code release.
\end{proof}

Two remarks fix the proposition's scope. It corrects a natural first reading: a classical
structure does not project to the base but selects lifts into the total space --- classicality
is section data, and what the copy operation broadcasts is a lift, not a ray. And it claims
nothing topological about cloning: the finiteness of the copyable set is linear, structural
no-cloning residing in non-naturality \autocite{abramskyheunen2012hstar}, so the first place the
bundle's global topology can bite is the quaternionic rung, where the degree-one clutching map
stands ready to be compared with the strong-complementarity law governing an interacting pair of
structures. The dictionary is, finally, the complement of the projective treatment cited above:
Burton and Anwar pass functorially to the base and deliberately discard the fibre
\autocite{burtonanwar2026meromorphic}, and the classical structures turn out to live in exactly
the data their projectivisation forgets --- the two readings compose, and the resulting triangle
is a consistency condition available to constrain the remaining conditions of the
correspondence.

\begin{openproblem}[The Hopf-fibration--Frobenius correspondence]\label{op:hopf-frobenius} Let
$\mathbf{Hopf}$ denote the category of the three fibrations
$S^{2^{k}-1} \hookrightarrow S^{2^{k+1}-1} \to S^{2^{k}}$, $k = 1, 2, 3$, with bundle maps as
morphisms, and let $\mathbf{Frob}$ denote the Frobenius algebras on $\mathbb{C}^{2}$ in
$\mathbf{FdHilb}$, with the commutative-special (GHZ) and anti-special (W) classes
distinguished. Construct a functor $F \colon \mathbf{Hopf} \to \mathbf{Frob}$ --- or prove that
none exists --- under which the complex fibration yields the single-qubit classical structures
--- Proposition~\ref{prop:f1-dictionary} identifies these as the lifted antipodal skeletons,
section data over two base points rather than the points alone; the quaternionic fibration
yields the interaction of two such structures, with the $SU(2)$ fibre carrying the entanglement
that survives the classical quotient; the octonionic fibration yields the GHZ/W pair, GHZ on the
full non-associative fibre and W on the associative $\mathbb{H} \subset \mathbb{O}$; and the
composition by which Coecke--Kissinger generate the $N$-qubit classes corresponds to a
well-defined geometric operation on the bundles. If such an $F$ exists, compositional
completeness follows from Adams's theorem as a topological corollary --- there is no fourth
fibration to supply a fourth primitive --- and if $F$ is faithful enough that the implication
also runs back, the termination of the tower and the completeness of the alphabet are equivalent
expressions of a single classification. The first condition is discharged by
Proposition~\ref{prop:f1-dictionary}; what remains open is the quaternionic and octonionic
conditions and, above all, the correspondence of compositions. Both open conditions are already
constrained: the quaternionic fibre's physical content --- the recoverable relative phase and
its purity ledger --- is established (\S\ref{sec:wf-reconciliation}), leaving the
clutching-map-to-strong-complementarity comparison as the structural residue; and any candidate
functor must reproduce the proved canonical-orbit images, GHZ to the colour triplet and W to the
singlet (\S\ref{sec:canonical-orbits}, Appendix~\ref{app:ghz-w-proofs}). A proof that no such
functor exists, or a correspondence at the level of the classifications without functoriality on
compositions, would each also resolve the problem: the first would make the agreement of two
independent classifications a fact demanding its own explanation; the second would establish the
weaker reading.
\end{openproblem}

A resolution in the affirmative would do three things. It would explain compositional
completeness rather than merely prove it --- a corollary of the division-algebra structure of
quantum mechanics rather than an independent algebraic fact. It would give the diagrammatic
language of this paper's own interaction circuits a topological warrant: the GHZ/W spider
calculus would be the unique language the topology permits, not a convenient notation. And it
would open a two-way bridge, translating problems on compositions of Frobenius structures into
operations on bundles and back. The problem sits at the intersection of categorical quantum
mechanics, the quantum-information geometry of the fibrations, and the division-algebra
approaches to particle physics, and it is posed here so that any of the three communities can
take it up in its own language. It is perhaps worth noticing that both classifications already
answer to the same name --- the fibrations are Heinz Hopf's, and so, through the bialgebra
structure of the diagrammatic calculus, are the Hopf algebras --- a coincidence of eponymy that
the conjecture proposes to promote to a theorem. With the epistemic reading above, an
affirmative resolution would leave the tower's endpoint as one limit seen four ways: algebraic,
geometric, compositional, epistemic --- the last being the fixed-point classification of Open
Problem~\ref{op:sigma-functor}, with which this problem is cross-linked, each constraining the
other.

\subsection{What this sets up}

The entanglement-sensitivity of the Hopf maps is not a side feature but the hinge of the
physical reading that follows. Because separability is exactly the collapse of the base to a
product of Bloch spheres, the structure of the entanglement --- the SLOCC class of a multi-qubit
state, and within it the canonical maximally entangled orbit --- is carried by precisely the
degrees of freedom on which $G_{\mathrm{SM}}$ acts in Proposition~\ref{prop:gsm}. This is what
makes it possible, in \S\ref{sec:entanglement-classes}, to read the two genuinely tripartite
entanglement classes onto the $\mathbb{C} \oplus \mathbb{C}^3$ split through their canonical
representatives: the GHZ orbit confined to the colour-triplet $\mathbb{C}^3$ sector, the W orbit
reaching the $\mathbb{C}$ singlet, recovering the quark--lepton division not as an assignment
but as a consequence of the entanglement structure. The scope of those statements --- canonical
orbits rather than whole classes, and the identification of the canonical orbits with the
physical particles --- is made precise there. The Hopf fibrations are the bridge; the
entanglement classes, treated next, are the traffic that crosses it.

%      % EDITED 2026-09-23 scripts/patch_prop_w_access_outline_2026-09-23T1100.py
% EDITED 2026-09-10 scripts/patch_quoting_and_typography_2026-09-10T0620.py
% =====================================================================
%  SECTION: Entanglement Classes and Gauge Groups
%  (sec:entanglement-classes). Drafted 2026-06-08; supersedes the
%  earlier same-day draft. AUDITED 2026-07-06 (final pass, ninth stop).
%
%  PLACEMENT (reordered chain, JS 2026-07-06): position 9 --
%      \input{hopf_justification_section_2026-06-08}
%      \input{entanglement_classes_and_gauge_groups_2026-06-08}
%      \input{register_section_2026-06-08}
%  The "previous section" phrases in the prose refer to
%  hopf-justification and are TRUE at this position (the hopf->classes
%  adjacency survived the reorder); the file is position-fragile in
%  that specific pairing, by design.
%
%  APPENDIX LABELS ACQUITTED: app:ghz-proof (line 246 of the appendix)
%  and app:w-proof (line 364) both exist -- the June "written next
%  turn" promise was kept; app:setup and the ceiling lemma live.
%  sec:generations is defined in the particle-zoo section; all forward
%  refs land.
%
%  AUDIT RECORD: lem:cproj, prop:ghz, prop:w-access, prop:w-doublet and
%  the canonical-orbits postulate all verified against every quoting
%  consumer (the frame's (2/3,1/3) quote exact; the register delivers
%  the deferred SU(2)_L/handedness fixing); the e1-alignment point
%  (Fano pairing, not commutativity) sound; loci numbers verified.
%  Fixes: CKW key variant normalised to corpus canon
%  (coffmankunduwootters2000); chirality return-promise given its
%  addresses (sec:frame, sec:generations); absent-gate forward stitch
%  at the DVC sentence; and the TWO-NINTHS IDENTITY drawn out --
%  sup C_frac = r^2 (ceiling lemma) and slack = 1 - 4 det rho = r^2,
%  so access ceiling and budget capacity are one invariant (W: 1/9;
%  GHZ: 0 -- confinement and vanishing C-access are one fact), a
%  partial discharge of op:w-slack's physical-reading clause recorded
%  in the confinement section; script check added. op:ideal-orbit
%  inventoried and stands as posed (well-formed, counterexample-valued).
%  The header's stray Gödel umlaut dies with this rewrite: file now
%  pure ASCII. British English; \autocite; no bullets in prose.
% =====================================================================

% ADDED 2026-07-12: sec:algebraic-dictionary prior-work subsection
% (Furey/Dixon/Krasnov dictionary; two REF-OWED keys; see its header).
\section{Entanglement Classes and Gauge Groups}
I'll quote from Szangolies \autocite{szangolies2025standardmodel} because this is very much his idea. 
\begin{quote}
Our main result is then to show how to use a construction of Dubois-Violette and Todorov \autocite{duboisviolette2016exceptional}, 
as extended and clarified by Krasnov \autocite{krasnov2021so9}, to derive the exact symmetry
and degrees of freedom of (a generation of) Standard Model fermions from the entanglement properties of a three-qubit state.
\end{quote}
I had been looking for a relationship between entanglement and gauge theory for many years and then I saw Szangolies 
give a ten minute session in an online workshop on division algebras in particle physics \autocite{szangolies_j_2021} 
during the presentation he flashed up a screen shot of a paper he was working on relating entanglement to particle physics.
I instantly knew this was what I was looking for and contacted him to ask if the paper was ready for publication. He very kindly let me
see a pre-print and I encouraged him with the phrase "perfection is the enemy of done". So I was very please to see the final thing published. 

A short while after that at a conference on quantum foundations in V\"axj\"o in 2023 I saw a single A4 poster by Filtov describing
a geometrical view of bi-partite entanglement with two Block spheres related by inversion. I already knew that bi-partite entanglement was 
related to SU(2), so the fact that Filtov had so clearly demonstrated the chirality inherent in SU(2) was the most significant 
contribution to the conference. A poster on a single sheet of A4! Which eventually led to \autocite{filatov2024towards}

That got me thinking that relating 3 Bloch spheres by inversion would be impossible and therefore if you could show a relationship between 3 particle entanglement and SU(3)
then the the fact that the colour force is blind to chirality would fall out immediately. And this is where we'll pick up. 

\label{sec:entanglement-classes}

The previous section established that the division-algebra ladder forces exactly three nontrivial
Hopf fibrations, $S^3 \to S^2$, $S^7 \to S^4$ and $S^{15} \to S^8$, with fibres $S^1$, $S^3$ and
$S^7$, and that fixing a complex direction inside the octonions leaves the Standard Model gauge
group as the stabiliser in $\mathrm{Spin}(9)$ (Proposition~\ref{prop:gsm}). This section asks where,
in the entanglement of a three-qubit state, the three gauge groups actually live. The answer is
clean for the first two and forces a genuine question for the third, and it is that question --- and
its resolution through the two classes of tripartite entanglement --- that the section is about.

\subsection{From fibres to groups, and where the pattern breaks}

For the first two rungs the identification needs no interpretation, because the fibre simply
\emph{is} the gauge group. The fibre of the first Hopf fibration is $S^1$, the unit circle in the
complex plane; as a group the unit complex numbers are exactly $U(1)$, the gauge group of
electromagnetism. The fibre of the second is $S^3$, the unit sphere in the quaternions; the unit
quaternions form the group $\mathrm{Sp}(1) \cong SU(2)$, the gauge group of the weak interaction.
So the abelian phase $U(1)$ and the non-abelian $SU(2)$ are not analogues of the Hopf fibres --- they
are those fibres, read as groups. A reader from quantum computation will recognise $S^1$ as the
global phase of a single qubit and $S^3 \cong SU(2)$ as the group of single-qubit unitaries --- the
Bloch-sphere rotations; a reader from gauge theory will recognise the two smallest gauge groups.

The pattern then breaks, and it is important to be clear about why, because the break is what makes
the rest of the section necessary. The fibre of the third, octonionic, Hopf fibration is $S^7$, the
unit sphere in the octonions. But the unit octonions are \emph{not a group}: octonionic
multiplication is non-associative, so $S^7$ does not close into a Lie group at all --- it is a
Moufang loop, not a group. There is therefore no statement ``$S^7 \cong SU(3)$'' to be written down
in parallel with $S^1 \cong U(1)$ and $S^3 \cong SU(2)$. The third gauge group cannot be the fibre.

This is not a defect to be patched but a signpost. The previous section already located the only
candidate: $SU(3)$ is the subgroup of $G_2 = \mathrm{Aut}(\mathbb{O})$ that stabilises a complex
direction, equivalently the split $\mathbb{O} \cong \mathbb{C} \oplus \mathbb{C}^3$
(Proposition~\ref{prop:gsm}). That tells us the group is available; it does not tell us where in the
physics it acts. Because $S^7$ is not itself a group, $SU(3)$ must enter not as the fibre but as a
symmetry of the \emph{structure carried on} the third rung --- and the structure carried there is the
entanglement of the three-qubit state. Two questions follow, and they are the questions this section
answers. First: can the colour $SU(3)$ be located in the entanglement of three qubits, and if so in
which class --- the GHZ class or the W class? Second: does the split between those two classes, the
only two kinds of genuinely tripartite entanglement, map onto the two matter sectors that colour
distinguishes, namely the coloured quarks and the colourless leptons?

These are not idle questions, and they are not new ones. Szangolies has already established part of
the gauge-group link from the representation-theoretic side \autocite{szangolies2025standardmodel};
Filatov and Auzinsh have shown that bipartite entanglement carries an intrinsic chirality of the
kind the weak $SU(2)$ exhibits \autocite{filatov2024towards}; and, for well over a decade, Duff and
collaborators at Imperial College have pursued deep correspondences between multipartite
entanglement and string and brane physics \autocite{duff2007string,borsten2009blackholes}. We take
up the pedigree below and then turn the questions into theorems.

\subsection{Chirality is already present in bipartite entanglement}

Before the tripartite case, a structural observation about the bipartite one, because it bears on
the weak sector that sits on the $S^3 \cong SU(2)$ rung. Filatov and Auzinsh showed that a pure
maximally entangled two-qubit state cannot be drawn on two Bloch spheres of the same handedness:
representing the entanglement faithfully requires the coordinate axes of one sphere to be
right-handed and those of the other left-handed \autocite{filatov2024towards}. Handedness is thus
not an artefact of how one chooses to draw the state; it is forced by the entanglement, and it is a
topological fact about the quaternionic Hopf fibration --- the same fibration whose fibre is the weak
$SU(2)$.

The weak interaction is, famously, the one sector of the Standard Model that is chiral: it couples
to left-handed fermions and not to their right-handed counterparts, and that chirality is usually
treated as an empirical input with no structural explanation. It is therefore at least suggestive
that chirality should appear, unbidden and forced, in the simplest entangled state of the very rung
that carries $SU(2)$. We draw no specific derivation from this here --- the mechanism that turns the
Filatov handedness into the physical left-handedness of the weak force belongs to the later
discussion of the register, the frame, and the generations (\S\ref{sec:register},
\S\ref{sec:frame}, \S\ref{sec:generations}) --- but we record it as a structural encouragement
for a programme that proposes to read gauge content off entanglement, and we return to it there.

\subsection{The two classes of tripartite entanglement}

Three qubits admit exactly two inequivalent classes of genuinely tripartite entanglement under
stochastic local operations and classical communication: the GHZ class and the W class
\autocite{durvidalcirac2000three}. (The operational face of the inequivalence --- the converter
mathematics permits and the vertex dictionary does not stock --- is \S\ref{sec:absent-gate}.) Every natural invariant separates them. The three-tangle is
maximal on GHZ and identically zero on W \autocite{coffmankunduwootters2000}; tracing
out one qubit of a GHZ state leaves the other two separable, while the same operation on a W state
leaves them entangled. Coecke and Kissinger gave the distinction an algebraic form: GHZ-class states
carry a \emph{special} commutative Frobenius algebra and W-class states an \emph{anti-special} one,
the two differing in a purely topological way \autocite{coeckekissinger2010compositional}. Those two
Frobenius structures are the primitives of a calculus expressive enough to build arbitrary $N$-qubit
states, so the entanglement alphabet, like the division-algebra ladder, closes at three: there is no
third primitive to find, just as there is no fourth division algebra (Proposition~\ref{prop:termination}).

We're now at the point where it should seem completely obvious to ask the question "That split between GHZ and W looks suspiciously like the split between 
the quark and lepton sectors, tell me more about it?" That there are exactly two genuinely tripartite classes, and exactly two matter sectors that colour
distinguishes, is the coincidence the rest of the section makes precise.

\subsection{The map into the octonion algebra}

We use the identification of a three-qubit state $\psi \in \mathbb{C}^8$ with a pair of octonions
$(o_1, o_2) \in \mathbb{O}^2$, obtained by splitting on the first qubit under the Cayley--Dickson
correspondence, together with the octonionic Hopf map $\mathbf{h} = o_1\bar{o}_2$. Fixing the
preferred complex direction $e_1$ gives the split $\mathbb{O} \cong \mathbb{C} \oplus \mathbb{C}^3$,
with $\mathbb{C} = \{e_0, e_1\}$ the singlet direction and $\mathbb{C}^3 = \{e_2,\dots,e_7\}$ the
colour-triplet sector carrying the $\mathbf{3} \oplus \bar{\mathbf{3}}$ of the $SU(3)$ of
Proposition~\ref{prop:gsm}. Everything below turns on one identity.

\begin{lemma}[The $\mathbb{C}$-projection is a Hermitian inner product]\label{lem:cproj}
The projection of $\mathbf{h}$ onto the $\mathbb{C} = \{e_0,e_1\}$ summand is
$\pi_{\mathbb{C}}(\mathbf{h}) = \overline{\langle o_1, o_2\rangle_{\mathbb{C}^4}}$, the conjugate of
the standard Hermitian inner product of the two first-qubit branches. Hence
$\pi_{\mathbb{C}}(\mathbf{h}) = 0$ if and only if those branches are $\mathbb{C}^4$-orthogonal.
\end{lemma}

The lemma is proved in Appendix~\ref{app:setup}, where the octonion multiplication table and the
identification of $\mathbb{C}^4$ with $\mathbb{O}$ are fixed explicitly; the identification must be
adapted to the complex structure given by left multiplication by $e_1$, and the sign bookkeeping is
part of the proof. The composition-algebra identity $|\mathbf{h}|^2 = |o_1|^2|o_2|^2$, which holds on
$\mathbb{O}$ and no larger Cayley--Dickson algebra, makes the $\mathbb{C}$-fraction
$C_{\mathrm{frac}} = |\pi_{\mathbb{C}}(\mathbf{h})|^2/|\mathbf{h}|^2$ a well-defined local invariant.
What singles out $e_1$ is not commutativity --- every imaginary direction spans a commutative plane
with $e_0$ --- but alignment: $e_1$ is the unique imaginary direction whose Fano pairing of the basis
units coincides with the pairing of real and imaginary parts in the qubit amplitudes, and the GHZ
cancellation below requires exactly that alignment.

\subsection{GHZ states are confined to the colour sector}

\begin{proposition}[GHZ confinement]\label{prop:ghz}
For every state in the local-unitary orbit of $\lvert\mathrm{GHZ}\rangle$, the $\mathbb{C}$-projection
of the Hopf image vanishes identically, so the image lies entirely in the colour-triplet sector
$\mathbb{C}^3 = \mathbf{3}\oplus\bar{\mathbf{3}}$, symmetrically under each of the three qubit splits.
This property characterises the orbit exactly: the $\mathbb{C}$-projection vanishes across the orbit
under every split if and only if all three single-qubit reduced states are maximally mixed, which
holds if and only if the state is local-unitarily equivalent to $\lvert\mathrm{GHZ}\rangle$.
\end{proposition}

\emph{Proof outline.} By Lemma~\ref{lem:cproj} it is enough that the branches are orthogonal for
every local-unitary image. The bare GHZ state has orthogonal branches of \emph{equal} norm
($\rho_k = \tfrac12\mathbb{I}$ for each $k$); the internal unitaries preserve orthogonality and equal
norms, and the first-qubit unitary mixes the branches with an inner product that evaluates to
$-\bar a\bar b\cdot\tfrac12 + \bar b\bar a\cdot\tfrac12 = 0$, the cancellation holding because
$\mathbb{C}$ is commutative. The full derivation --- the characterisation, and the symmetry across
all three splits --- is given in Appendix~\ref{app:ghz-proof}. \hfill$\square$

The cancellation is exact; over random local-unitary samples the $\mathbb{C}$-fraction sits at
machine zero. Two consequences are proved rather than conjectured. Confinement is frame-independent:
the invariant separating the two canonical orbits is the single-qubit reduced spectrum --- equal
$(\tfrac12,\tfrac12)$ on the GHZ orbit, unequal $(\tfrac23,\tfrac13)$ on the W orbit --- and since no
local unitary can change that spectrum, no change of local frame can carry a colour triplet into the
lepton sector. The scope of the statement is the canonical orbits: SLOCC deformations away from them,
such as $\alpha\lvert000\rangle+\beta\lvert111\rangle$ with $\alpha\neq\beta$, acquire a nonzero
$\mathbb{C}$-projection, and their status is fixed by the identification of
\S\ref{sec:canonical-orbits}. And the QCD vacuum angle vanishes: the GHZ image has
no oriented pairwise linking structure to orient a $\theta\,\mathrm{Tr}(G\wedge G)$ term, so
$\theta_{\mathrm{QCD}} = 0$ is forced and no Peccei--Quinn axion is admitted. The reading is
GHZ $\leftrightarrow$ colour triplet $\leftrightarrow$ quark.

This is the first result to drop out of the identification between GHZ and the colour force, and it was the one 
I was looking for, SU(3) is blind to chirality because you cannot have a geometric construction analogous to the 
2 Bloch sphere case with 3 Bloch spheres. But once a bridge between the gauge groups of the Standard Model
and multi-particle entanglement has been built then a lot of two way traffic becomes possible. One can move hard problems
on one end of the bridge to the other end where they have easy solutions, $\theta_{QCD}=0$ being only the first. This also means
you can express particle interactions in terms of quantum circuits acting on entangled states, which we pick up in \ref{sec:interactions}

\subsection{W states reach the lepton sector and carry the weak doublet}

\begin{proposition}[W accessibility]\label{prop:w-access}
For states in the local-unitary orbit of $\lvert\mathrm{W}\rangle$, the $\mathbb{C}$-projection of the
Hopf image is generically nonzero, vanishing only on a measure-zero subset, with
$\sup C_{\mathrm{frac}} = \tfrac19$.
\end{proposition}

Where the GHZ orbit is confined to the colour sector, the W orbit reaches the lepton sector, and
does so generically: the $\mathbb{C}$-component of its Hopf image is nonzero at every point of the
orbit but a set of measure zero, and it is bounded, never exceeding one ninth of the image's squared
norm. Both facts come from the same one-line computation as Proposition~\ref{prop:ghz}, with one
norm changed.

\emph{Proof outline.} Split on the first qubit. The bare W has branches
$o_1 = \tfrac{1}{\sqrt3}(0,1,1,0)$ and $o_2 = \tfrac{1}{\sqrt3}(1,0,0,0)$, orthogonal but of
\emph{unequal} norm, $\lvert o_1\rvert^2 = \tfrac23$ and $\lvert o_2\rvert^2 = \tfrac13$. Unitaries
on the other two qubits act on both branches alike and change neither their inner product nor
their norms, so $C_{\mathrm{frac}}$ depends only on the first-qubit unitary
$U_1 = \bigl(\begin{smallmatrix} a & b \\ -\bar b & \bar a\end{smallmatrix}\bigr)$. Orthogonality
kills the $\bar a^{2}$ and $\bar b^{2}$ terms of the mixing formula, and the unequal norms leave the
residue $\langle o_1, o_2\rangle = \bar a\bar b\,(\tfrac13 - \tfrac23) = -\tfrac13\bar a\bar b$: the
equal-norm cancellation of the GHZ proof, failing by exactly the norm difference. It vanishes only
when $ab = 0$, $U_1$ diagonal or anti-diagonal, a measure-zero subset of $SU(2)$ on which the bare
representative itself sits. With $p = \lvert a\rvert^2$ the branch norms are
$\lvert o_1\rvert^2 = (1+p)/3$ and $\lvert o_2\rvert^2 = (2-p)/3$, and by Lemma~\ref{lem:cproj}
together with the composition identity,
\[
C_{\mathrm{frac}} \;=\; \frac{p(1-p)}{(1+p)(2-p)},
\]
maximised at $p = \tfrac12$ with value $\tfrac19$: the first-qubit Bloch vector rotated into the
equatorial plane. The reduced spectrum $(\tfrac23,\tfrac13)$ is a local-unitary invariant, so the
statement holds across the orbit, and by the permutation symmetry of $\lvert\mathrm{W}\rangle$ under
each of the three splits. The ceiling is the instance $r = \tfrac13$ of the general identity
$\sup C_{\mathrm{frac}} = r^{2}$, the squared Bloch radius of the marginal (Appendix~\ref{app:setup},
Lemma~\ref{app:lem:ceiling}); the full derivation is in Appendix~\ref{app:w-proof}.
\hfill$\square$

Two remarks. The geometry is sharper than a bound. With $r$ fixed at $\tfrac13$ on the orbit the
colour component of the image has constant norm $\sqrt{8/9}$, so the W orbit's Hopf image lies on
the product of a colour five-sphere of radius $\sqrt{8/9}$ and a Bloch two-sphere of radius
$\tfrac13$, disjoint from the colour equator of unit radius that carries the GHZ image; the missing
ninth of squared norm sits in the $\mathbb{C}$-plane, and the two canonical orbits occupy disjoint
compact loci on the base, which is the geometric form of the statement that no change of local
frame connects them. And the ceiling has been measured: single-wire tomography of the W register on
hardware returned a probe-corrected $r^{2} = 0.1095 \pm 0.0051$ against $\tfrac19 = 0.1111$
(Appendix~\ref{app:hw-ninth}), the number this paper calls the leptonic ninth.

\begin{proposition}[W carries the $SU(2)$ doublet]\label{prop:w-doublet}
Each single-qubit reduced state of a state in the local-unitary orbit of $\lvert\mathrm{W}\rangle$
has the non-degenerate spectrum $(\tfrac23,\tfrac13)$, symmetrically across the three qubits. The
non-degeneracy singles out a distinguished eigenbasis --- a $T_3$ axis with its two doublet members
--- breaking the $SU(2)$ invariance of the marginal down to the $U(1)$ of rotations about that axis.
The degenerate GHZ marginal $\tfrac12\mathbb{1}$ singles out no such direction.
\end{proposition}

\emph{Proof outline.} Immediate from the spectrum. The single-qubit marginal of
$\lvert\mathrm{W}\rangle$ is $\mathrm{diag}(\tfrac23,\tfrac13)$ in the computational basis, the
same on each qubit by the permutation symmetry, and a local-unitary invariant up to conjugation.
Non-degeneracy fixes a unique eigenbasis, a distinguished $T_3$ axis between whose members the
$SU(2)$ raising and lowering operators act, and leaves invariant only the $U(1)$ of rotations about
that axis. The degenerate GHZ marginal $\tfrac12\mathbb{1}$ has every basis as an eigenbasis and
singles out no direction. Full derivations are in Appendix~\ref{app:w-proof}. \hfill$\square$

The contrast with GHZ is entirely in the branch norms: equal norms force the cancellation, unequal
norms leave a residue. The reading is W $\leftrightarrow$ colour singlet $\leftrightarrow$ lepton,
organised into weak-isospin doublets. Which of the two quaternionic actions, $SU(2)_L$ or $SU(2)_R$,
becomes physical weak isospin --- and how the Filatov handedness of \S\ref{sec:entanglement-classes}
becomes physical left-handedness --- is fixed by the structure of the full register, and we defer it
to \S\ref{sec:register}.

\subsection{Canonical orbits are the particles}
\label{sec:canonical-orbits}

The propositions above are scoped to the local-unitary orbits of the standard GHZ and W states ---
the maximally entangled representatives of their classes. The SLOCC classes are larger: each contains
a continuum of deformations, such as $\alpha\lvert000\rangle+\beta\lvert111\rangle$ with
$\alpha\neq\beta$, whose marginal spectra differ and for which the $\mathbb{C}$-projection switches
on. We therefore make the identification explicit, and adopt it as a structural postulate of the
framework: \emph{a physical one-particle state is a canonical maximally entangled orbit; the SLOCC
class is the particle type}. The quantum numbers of \S\ref{sec:charges} and the masses of
\S\ref{sec:masses} are computed on canonical orbits, so the postulate is load-bearing, and we state
it rather than leave it implicit.

The identification is not arbitrary, and the reason lives in machinery the paper already uses ---
and, since \S\ref{sec:superconductor-witness}, in an independent physical witness: the
superconductor, where the Bell and product rows are the two strata of the Anderson--Higgs
mechanism. Under
map--state duality, a maximally entangled state is the state-form of the identity map --- the algebra
acting on itself, written as a state --- and its local-unitary orbit corresponds to the unitaries,
while the SLOCC deformations correspond to invertible but non-unitary maps, excluded as physical
configurations for the same reason non-unitary evolutions are. In the Frobenius language
\autocite{coeckekissinger2010compositional}, the canonical orbits are precisely the \emph{special}
and \emph{anti-special} normalisations of the two tripartite algebras, and dagger-special commutative
Frobenius algebras correspond exactly to orthonormal bases
\autocite{coeckepavlovicvicary2013orthogonal}: the canonical orbit is the Frobenius algebra in its
canonical form, the tilted states its non-special deformations.

This converges with the identification at the heart of Furey's programme, in which particles are
minimal ideals of an algebra acting on itself \autocite{furey2012unified,furey2018three}. The
physical idea is the same --- the particle is the canonical, parameter-free substructure singled out
by the algebra's self-action --- expressed there in module coordinates and here in entanglement
coordinates. The selection principles differ in form (minimality of the ideal; maximal symmetry of
the orbit), and we record their reconciliation as a named open problem.

\begin{openproblem}[Ideals and canonical orbits]\label{op:ideal-orbit}
Prove that the dagger-special condition on the Frobenius side corresponds to the minimal-ideal
condition on the division-algebra side, so that the present identification and Furey's become one
fact described in two coordinate systems. A counterexample separating the two selections on some
sector, or a correspondence holding only on a proper subclass, would each also resolve the problem,
and either outcome would be informative: it would delimit exactly the sectors on which module
coordinates and entanglement coordinates may be used interchangeably.
\end{openproblem}

The status of the tilted, off-canonical states --- whether they admit a reading as
off-shell or perturbed configurations --- is likewise open; no result downstream depends on them.

\subsection{The dictionary}

Collecting the strands gives the correspondence between entanglement structure and gauge content,
the upper rows established in the previous section and the lower two the content of
Propositions~\ref{prop:ghz}--\ref{prop:w-doublet}.

\begin{center}
\begin{tabular}{llll}
\hline
Entanglement & Division algebra & Hopf fibre & Standard Model \\
\hline
separable & $\mathbb{C}$ & $S^1 \cong U(1)$ & electromagnetic phase \\
Bell (two-qubit) & $\mathbb{H}$ & $S^3 \cong SU(2)$ & weak isospin \\
GHZ (special) & $\mathbb{O}$, $\mathbb{C}^3$ sector & $S^7$ & $SU(3)$ colour / quarks (confined) \\
W (anti-special) & $\mathbb{O}$, $\mathbb{C}$ sector & $S^7$ & colour singlets / leptons (free) \\
\hline
\end{tabular}
\end{center}

The first two rows are group identities; the lower two are not, precisely because $S^7$ is not a
group, and that is why they required the theorems above rather than a one-line reading of the fibre.
A clarification of where the content lies: the class names in the left column are labels, and the
physics is in the geometry of the Hopf map they label. The base $S^8$ decomposes as (Bloch vector of
the split qubit) $\oplus$ (colour sector $\mathbb{C}^3$), and the canonical GHZ orbit maps onto the
unit sphere of the colour sector --- the equatorial $S^5\subset S^8$ --- while the canonical W orbit
occupies a disjoint locus of constant colour radius $\sqrt{8/9}$ and Bloch radius $\tfrac13$, with
$\mathbb{C}$-access ceiling $(\tfrac13)^2 = \tfrac19$ (Appendix~\ref{app:setup}). The two octonionic
rows share the same fibre and are separated as disjoint loci on the base: on the colour equator, or
tilted off it into the single-qubit Bloch directions.

One identity deserves to be drawn out, because it joins this section's geometry to the entanglement
budget of \S\ref{sec:budget-converse}. The ceiling lemma of Appendix~\ref{app:setup} gives
$\sup C_{\mathrm{frac}} = r^{2}$, the squared marginal Bloch radius; and the budget slack of a
register qubit is $1 - \tau = 1 - 4\det\rho = r^{2}$ identically. So the $\mathbb{C}$-access
ceiling and the external-correlation capacity are \emph{the same invariant}, $r^{2}$, seen in base
and in budget coordinates: for the W orbit both are exactly $\tfrac19$ --- the leptonic ninth of
Open Problem~\ref{op:w-slack} \emph{is} the lepton's singlet-direction access --- and for the GHZ
orbit both are exactly zero, so confinement and vanishing $\mathbb{C}$-access are one fact, not
two. An executable verification accompanies the code release.

\subsection{Prior work: Szangolies, and the Imperial programme}
\label{sec:entanglement-szangolies}

The overlap with existing work is worth stating precisely, because the convergence is itself
evidence. Szangolies obtains the same lepton/quark split directly from the representation theory of
the $\mathrm{Spin}(9)$ stabiliser: the octonionic spinor decomposes into a leptonic part from the
$\mathbb{C}$ factor, an $SU(3)$ singlet of hypercharge $-1$, and a quark part from $\mathbb{C}^3$, an
$SU(3)$ triplet of hypercharge $\tfrac13$, both $SU(2)$ doublets \autocite{szangolies2025standardmodel}.
This is reached with no reference to entanglement classes; the state is treated as carrying a
generation wholesale, and the split is a fact about the gauge group acting on
$\mathbb{C}\oplus\mathbb{C}^3$. Our result is complementary: where Szangolies fixes the representation
content of each summand, Propositions~\ref{prop:ghz} and~\ref{prop:w-access} fix \emph{which states
populate which summand}, by their SLOCC class. The lepton/quark division he assigns by representation
theory we obtain as a consequence of entanglement type, and his independent arrival from the gauge
side is the strongest external check that the identification is not an artefact of our route.

The longer pedigree matters too, because the idea that multipartite entanglement encodes the
structures of fundamental physics is not new and has been pursued by senior researchers for well over
a decade. Duff and collaborators at Imperial College developed the ``black holes, qubits and
octonions'' programme, beginning from the observation that the entropy of certain stringy black holes
and the tripartite entanglement of three qubits are governed by the same invariant, Cayley's
hyperdeterminant \autocite{duff2007string}, and extending to $N=8$ supergravity, where the
56-dimensional charge space of $E_7$ decomposes into seven tripartite sectors arranged on the lines
of the Fano plane --- the multiplication table of the octonions --- with Cartan's quartic $E_7$
invariant playing the role of the entanglement measure
\autocite{dufferrara2007seven,levay2007fano,borsten2009blackholes}. Read from the present vantage
point, that programme had the octonionic skeleton in hand: the Fano plane, the exceptional groups,
the three-qubit classification. The divergence is in the routing. There, the geometry was carried
into supergravity charge space, and the correspondence is a dictionary between invariants --- an
entanglement measure matched to an entropy formula, number to number, with the mechanism beneath the
dictionary left open. Here, the geometry stays in the state space, where the fibration's stabiliser
supplies the gauge group as a mechanism (Proposition~\ref{prop:gsm}) and the entanglement classes
serve as labels for loci of the Hopf map. The Imperial programme establishes that qubit entanglement,
division algebras and the same three-qubit classification at the foundations of physics is a
direction taken seriously, independently, by others; what we add is the routing of that structure
into the gauge sectors of the Standard Model, as theorems rather than analogies.

\subsection{Prior work: the division-algebra programmes, and the two coordinatisations}
\label{sec:algebraic-dictionary}
% ADDED 2026-07-12. Sibling of the Szangolies/Imperial subsection:
% positions the framework against Furey, Dixon, Gillard--Gresnigt and
% Krasnov as TWO COORDINATISATIONS of one structure. REF-OWED (two new
% keys, .ris cluster to follow): dixon1994division (G. M. Dixon,
% "Division Algebras: Octonions, Quaternions, Complex Numbers and the
% Algebraic Design of Physics", Kluwer 1994); krasnov2021so9
% (K. Krasnov, the Spin(9) / centraliser-of-a-complex-structure
% characterization of the Standard Model gauge group, J. Math. Phys.
% 62, 021703 (2021), DOI 10.1063/5.0039941, arXiv:1912.11282 --
% VERIFIED 2026-07-12); partonpiccinni2018spin9 (Parton & Piccinni,
% The Role of Spin(9) in Octonionic Geometry, arXiv:1810.06288 --
% arXiv verified, journal DOI to confirm at Mendeley import).
% 2026-07-12 (same day, later): charge lemma DISCHARGED --
% appendix_winding_number_operator_2026-07-12.tex, lem:winding-number;
% op:dictionary-lemmas shrunk to the chirality item, label unchanged.
% CUTOVER: add the
% composability sentence with \ref{app:nb-I} and the scan reference
% \ref{app:nb-M} when the notebook record is input.

A third body of prior work runs on the algebraic side of the same division algebras: Dixon's
tensor algebra $\mathbb{R} \otimes \mathbb{C} \otimes \mathbb{H} \otimes \mathbb{O}$
\autocite{dixon1994division}, Furey's particles as minimal ideals with the operator content
generated by left multiplication \autocite{furey2012unified,furey2018three}, the
Gillard--Gresnigt Clifford-algebra routes to three generations
\autocite{gillardgresnigt2019three}, the exceptional-Jordan wing
\autocite{duboisviolettetodorov2019exceptional}, and Krasnov's characterisation of the gauge
group inside $\mathrm{Spin}(9)$ \autocite{krasnov2021so9}. The relationship to the present
framework is closer than parallel effort, and the organising fact is worth stating as a rule:
\emph{the algebra fixes the skeleton; the geometry carries the measure.} Ideal decompositions
and ladder counting yield exact integer data --- charge quantisation in thirds, representation
content, a generation count --- because an algebra has no metric to get wrong. Fibrations yield
the continuous data --- $\tfrac29$, $\tfrac14$, the geodesic masses of \S\ref{sec:masses}, the
Born rule itself --- because spheres do. The two are coordinatisations of one structure, and the
dictionary is short enough to give whole:

\begin{center}
\begin{tabular}{lll}
\hline
object & module coordinates & entanglement coordinates \\
\hline
arena & $\mathbb{C}\otimes\mathbb{O}$; Dixon's $\mathbb{R\,C\,H\,O}$ &
  the register $\mathbb{C}^8$; the nesting $\mathbb{C}\subset\mathbb{H}\subset\mathbb{O}$ \\
operators & $\mathbb{C}l(6)$ from left multiplication &
  the register's left-action envelope \\
selection & nilpotent basis, \emph{chosen} &
  record direction, \emph{derived} (\S\ref{sec:frame}) \\
gauge group & centraliser in $\mathrm{Spin}(9)$ &
  fibration stabiliser (Proposition~\ref{prop:gsm}) \\
multiplets & minimal ideals; $1{+}3{+}3{+}1$ grading &
  canonical orbits; $\mathbb{C}\oplus\mathbb{C}^3$ by class \\
charge & ladder count over three &
  $S^1$ winding through $\mathbb{C}^3$ \\
generations & three ideal prints; $Cl(8)$ copies &
  sedenion $\mathbb{O}_i$ orbit; Fano lines; $K_3$ classes \\
the theory itself & module + $J$ \emph{chosen} + Born posited &
  module + $J$ and Born \emph{derived} (\S\ref{sec:postulates}) \\
\hline
\end{tabular}
\end{center}

Three rows deserve their sentences. The gauge-group row is a composite reading whose ingredients
are separately on record but whose composition we have not found stated: $\mathrm{Spin}(9)$ is
the symmetry group of the octonionic Hopf fibration $S^{15} \to S^8$
\autocite{partonpiccinni2018spin9}, so Krasnov's theorem --- the Standard Model gauge group as the centraliser of a
complex structure inside $\mathrm{Spin}(9)$ --- \emph{is} Proposition~\ref{prop:gsm} in module
coordinates: the symmetries of the top Hopf rung that commute with the selected record direction.
On the algebraic side both inputs are taken as given; the present framework derives the arena
(three qubits, \S\ref{sec:register}) and the selection (the record direction, \S\ref{sec:frame}),
so on this side the characterisation upgrades from a description to a derivation. The selection
row states the general pattern of which that is an instance: what the algebraic programmes
\emph{choose} --- the nilpotent basis whose stabiliser is $SU(3)$, the complex structure, the
preferred unit --- the geometric route \emph{derives}, and this is the single most consequential
thing each side offers the other, in one direction a reason and in the other a calculus. The
multiplet row is Open Problem~\ref{op:ideal-orbit}, already priced above: the exterior-algebra
grading of Furey's ideal, $1{+}3{+}3{+}1$ under $SU(3)$, is the Hopf split
$\mathbb{O} \cong \mathbb{C}\oplus\mathbb{C}^3$ read grade by grade, leptons at the singlet
grades and quarks at the triplet grades exactly as Propositions~\ref{prop:ghz}
and~\ref{prop:w-access} populate them.

The division of labour is then plain. What the module coordinates cannot supply is anything with
a measure in it --- and the last row of the table is this rule at the foundations rung, the
closing paragraph of \S\ref{sec:postulates} stating the reading and guarding its arrow: no $\tfrac29$,
no mass, no probability --- the algebraic programmes have no
Born-rule story at all --- and no compositional statement of confinement, which here is
compactness of the base. What the entanglement coordinates cannot supply without translation is
an associative operator calculus ready for Lagrangians and the parameter-free integer bookkeeping
that makes charge quantisation automatic. The generations row is where the translation is under
active construction: the module-side prints, the sedenion orbit, the Fano lines and the $K_3$
classes are now suspected of being presentations of a single Klein-four-group with its $S_3$ of
automorphisms, and the shadow programme of \S\ref{sec:sedenion-announcement} is that suspicion
made precise.

The charge row has since been converted from correspondence to identity, in the same revision
that tabulated it: Appendix~\ref{app:winding-number} proves that Furey's number operator is the
relative winding of the colour circle's spinor lift --- the engine being the one-line identity
$\mathrm{i}\, c_{x_{k}} c_{y_{k}} = 2N_{k} - 1$, each mode's ladder count \emph{equal} to its
plane's Clifford bilinear --- so that $Q = N/3$ and the fibration's winding are one operator,
the charge thirds are winding integers divided by the three colour planes, and the half-integer
vacuum offset is the spinor double cover doing its usual work
(Lemma~\ref{lem:winding-number}).

\begin{openproblem}[The dictionary's chirality lemma]\label{op:dictionary-lemmas}
One lemma of the dictionary remains at the one-page grade: show that the Filatov
opposite-handedness of \S\ref{sec:generations} is the distinction between left and right
quaternionic multiplication, giving chirality's algebraic address and its geometric proof in one
stroke. A failure of the identification would be equally informative, delimiting where the two
coordinate systems part.
\end{openproblem}

\subsection{What is settled, and what comes next}

What is settled in this section, with the full derivations carried into the appendices of this paper:
GHZ confinement to $\mathbb{C}^3$ (Proposition~\ref{prop:ghz}, Appendix~\ref{app:ghz-proof}); W
accessibility of $\mathbb{C}$ and the $SU(2)$ doublet content of W
(Propositions~\ref{prop:w-access}--\ref{prop:w-doublet}, Appendix~\ref{app:w-proof}); the topological
character of confinement; and $\theta_{\mathrm{QCD}} = 0$. These rest only on the composition-algebra
identity and the local-unitary invariance of single-qubit spectra, with no free parameters. The two
questions posed at the outset are answered: the colour $SU(3)$ is found in the entanglement of the
three-qubit state, in the GHZ class, exactly because $S^7$ is not a group and the gauge content must
be carried structurally; and the GHZ/W split does map onto the quark/lepton division.

What remains is to make the carrier explicit. The dictionary identifies which entanglement class is
which kind of particle, but not yet how a single particle's three gauge charges sit together, how its
electric charge, weak isospin, hypercharge and colour are read off, or how the Filatov handedness
becomes physical chirality. Those are the business of the next sections: the three-qubit register and
the quantum numbers read from it (\S\ref{sec:register}, \S\ref{sec:charges}), the three generations
(\S\ref{sec:generations}), and the masses (\S\ref{sec:masses}). The residual questions that are
properly spacetime-level --- among them the dynamical content of confinement --- are taken up in
Paper~2, alongside the spacetime discussion to which they belong.

%  No adjacency phrasing in the prose: position-portable. All
%  cross-refs live -- sec:fanout-direction DRAFTED 2026-07-06 (formerly the one undrafted
%  section, by design).
%
%  AUDIT RECORD: prop:gsm, prop:termination and prop:f1-dictionary
%  verified against every quoting consumer (statements canonical);
%  honest-debt subsection names both the selection assumption and the
%  mereology debt; the epistemic two-climbs disambiguation and the
%  compositional rung-evidence verified; sec:canonical-orbits live;
%  file pure ASCII. Fixes: op:hopf-frobenius's cross-link with
%  op:sigma-functor made TWO-WAY (the "four ways" line now names the
%  op); both remaining functor conditions ANCHORED --- the quaternionic
%  to the established fibre physics of sec:wf-reconciliation (residue:
%  clutching map vs strong complementarity), the octonionic to the
%  proved canonical-orbit images (umbrella label app:ghz-w-proofs
%  minted at the appendix for the purpose). Statuses unchanged.
%  British English; \autocite; no bullets in prose.
% =====================================================================

\section{Why Hopf Fibrations Belong in Particle Physics}
\label{sec:hopf-justification}

Invoking Hopf fibrations in a discussion of the Standard Model invites an immediate and fair
objection: that the move is decorative, an analogy between two pieces of attractive mathematics
with no load-bearing connection between them. The purpose of this section is to dispose of that
objection before the construction is put to work. The connection is not an analogy. It is
carried by two independent and independently published theorems, and the structures it relies
upon are not selected for convenience but compelled by a pair of classification results. We take
as already established the correspondence between qubit entanglement and the Hopf fibrations
themselves --- the subject of a mature literature
\autocite{mosseri2001hopf,bernevigchen2003geometry} --- and do not reproduce it here. What this
section addresses is the prior question of why that correspondence is entitled to carry physical
content at all.

\subsection{Two theorems, joined end to end}

The bridge from qubit entanglement to the Standard Model is built from two distinct results,
neither of which is original to the present framework and neither of which depends on the other.

The first is geometric. The unit sphere in the Hilbert space of $n$ qubits is the total space of
a Hopf fibration: $S^3 \to S^2$ for one qubit, $S^7 \to S^4$ for two, and $S^{15} \to S^8$ for
three. For two and three qubits the fibration is \emph{entanglement-sensitive} --- separability
is precisely the condition under which the base degenerates to a product of single-qubit Bloch
spheres and the additional dimensions associated with the higher imaginary directions collapse
\autocite{mosseri2001hopf,bernevigchen2003geometry}. This is a statement about the geometry of
state spaces and requires no physical interpretation whatever.

The second is algebraic. Singling out a preferred complex structure on the octonions ---
equivalently, fixing the split $\mathbb{O} \cong \mathbb{C} \oplus \mathbb{C}^3$ --- leaves a
residual symmetry, and that residual symmetry is exactly the gauge group of the Standard Model.
This was established through the exceptional Jordan algebra by Dubois-Violette and Todorov
\autocite{duboisviolette2016exceptional,duboisviolettetodorov2019exceptional,todorovduboisviolette2018deducing}
and given its sharpest spinorial form by Krasnov \autocite{krasnov2021so9}, who showed that the
relevant group is the stabiliser of a complex structure inside $\mathrm{Spin}(9)$. We record the
result in the form in which it enters the present work.

\begin{proposition}[Dubois-Violette--Todorov; Krasnov]\label{prop:gsm} Let $J_R$ be the complex
structure on the space $\mathbb{O}^2$ of $\mathrm{Spin}(9)$ spinors determined by a choice of
unit imaginary octonion, equivalently the split
$\mathbb{O} \cong \mathbb{C} \oplus \mathbb{C}^3$. The subgroup of $\mathrm{Spin}(9)$ commuting
with $J_R$ is \[ G_{\mathrm{SM}} \;=\; \frac{SU(3) \times SU(2) \times U(1)}{\mathbb{Z}_6}, \]
acting on the singled-out $\mathbb{C}$ summand as a colour singlet of hypercharge $-1$ and on
the $\mathbb{C}^3$ summand as a colour triplet of hypercharge $\tfrac{1}{3}$, both as $SU(2)$
doublets --- the quantum numbers of one left-handed generation.
\end{proposition}

The two results are independent, and the join between them is the only place a physical theory
must do work. Each result is a theorem; what neither supplies is a reason why nature should fix
one complex direction within $\mathbb{O}$ rather than treating all directions symmetrically. We
return to this --- the genuine load-bearing assumption --- in \S\ref{sec:hopf-honest-debt}, and
resolve it, in candidate form, in \S\ref{sec:fanout-direction}.

\subsection{The ladder is forced, not chosen}

It might still be objected that selecting the octonions, or the Hopf fibrations, is itself a
modelling choice carrying as much latitude as the choice of a grand-unified group. It is not,
and the reason is a pair of classification theorems that close off the alternatives entirely.

By Hurwitz's theorem the only normed division algebras are $\mathbb{R}$, $\mathbb{C}$,
$\mathbb{H}$ and $\mathbb{O}$ \autocite{hurwitz1898composition,baez2002octonions}. By Adams's
resolution of the Hopf-invariant-one problem --- together with the parallelisability results of
Bott--Milnor and Kervaire --- the only parallelisable spheres are $S^0$, $S^1$, $S^3$ and $S^7$,
so that there exist exactly three nontrivial Hopf fibrations, in bijection with $\mathbb{C}$,
$\mathbb{H}$ and $\mathbb{O}$
\autocite{adams1960hopf,bottmilnor1958parallelizability,kervaire1958nonparallelizability}. These
are not heuristics; they are among the more durable theorems in topology and algebra. Their
force here is that, once a framework is committed to amplitudes valued in a normed division
algebra --- a commitment that in the present construction is made far upstream, by the
self-referential argument of \S\ref{sec:self-reference} and \S\ref{sec:division-algebras}, and
not by fiat at this stage --- the ladder of Hopf fibrations is not selected but compelled, and
it terminates of its own accord.

\begin{proposition}[Termination]\label{prop:termination} The Cayley--Dickson tower of normed
division algebras and the tower of parallelisable spheres both terminate at real dimension
eight. Consequently there are exactly three nontrivial Hopf fibrations, the largest being the
octonionic $S^{15} \to S^8$, and the construction admits no extension to a fourth case.
\end{proposition}

The termination is worth dwelling on, because it is unusual for a model of this kind to predict
its own boundary. The next rung of the Cayley--Dickson sequence, the sedenions, possesses zero
divisors; the Hopf map can no longer be globally defined; there is no fourth fibration. This is
not a limitation imposed from outside but a structural fact, and it carries a physical
consequence that we develop later: there is no room in the construction for a fourth gauge force
or a fifth interaction, so that the discovery of either would count as evidence against the
framework rather than as an invitation to extend it. The endpoint is a theorem, not a
stipulation.

\subsection{Independent arrival as corroboration}

The most persuasive evidence that the destination is not an artefact of the route is that it is
reached by routes sharing none of the same machinery. Szangolies
\autocite{szangolies2025standardmodel} reaches $G_{\mathrm{SM}}$ from the emergent-gravity and
Kaluza--Klein programmes, reading the singled-out complex direction as a dimensional reduction
of a $9+1$-dimensional spacetime; the present framework reaches it from self-reference
(\S\ref{sec:self-reference}); and Furey and Gillard--Gresnigt reach the same gauge group, and
the same three-generation structure, through the minimal left ideals of the Dixon and Clifford
algebras, without invoking Hopf fibrations at all
\autocite{furey2018three,gillardgresnigt2019three}; the relation between their minimal ideals
and the entanglement classes used here is taken up in \S\ref{sec:canonical-orbits}. The
octonionic route to particle physics is in any case an old one, going back to the use of
$\mathbb{O}$ to model the strong force \autocite{gunaydingursey1974quark}.

That programmes with disjoint starting points converge on
$SU(3) \times SU(2) \times U(1)/\mathbb{Z}_6$ and on the same quark--lepton division is the
coincidence Szangolies himself remarks upon: the characteristic structure of the Standard Model,
routinely regarded as arbitrary, turns up independently in unrelated settings, and its appearing
twice would then be doubly odd unless the two settings are in fact related. We read the
convergence the same way --- as evidence that the division-algebraic description is singled out
by the mathematics rather than imposed by the modeller, and therefore that
Proposition~\ref{prop:gsm} is structural rather than contrived.

\subsection{The honest debt}
\label{sec:hopf-honest-debt}

It would be misleading to leave the impression that this route is free of assumptions, so we
state plainly the one assumption that all of these programmes share. Proposition~\ref{prop:gsm}
requires that \emph{some} complex direction in $\mathbb{O}$ be singled out, and the uniqueness
of that choice is not supplied by the theorem. Krasnov parametrises the complex structure $J_R$
by an arbitrary unit imaginary octonion, and the residual group is the same for every such
choice because $G_2 = \mathrm{Aut}(\mathbb{O})$ acts transitively on the unit imaginary
octonions. The theorem therefore guarantees the gauge group \emph{given} that a direction is
fixed, while remaining silent on \emph{why} one is fixed. Szangolies treats this as a free
choice and flags it, in his own conclusions, as an open element of the construction --- one to
which several inequivalent answers are available.

This, and not the appeal to Hopf fibrations, is the genuine load-bearing assumption of the whole
approach. The Hopf fibrations are forced (Proposition~\ref{prop:termination}); the singling-out
is not. The present framework offers a candidate resolution that the prior literature lacks: the
act of measurement, modelled as \textsc{fanout}, selects the direction physically, and the
transitivity of $G_2$ then guarantees that the gauge group is independent of which direction the
measurement happens to select. This argument is given in \S\ref{sec:fanout-direction}; we flag
here, at the appropriate epistemic status, that it is structural rather than fully proved, and
that the residual question --- why a measurement selects \emph{exactly one} imaginary direction
rather than merely \emph{permitting} a selection --- remains open. It is now priced as Open
Problem~\ref{op:fanout-selection}.

A second debt is shared with every programme of this character and is not closed by anything
said here: why the abstract quantum state should admit a tensor-product decomposition into
qubits in the first place. This is the quantum factorisation, or mereology, problem. The
entanglement-sensitivity on which the construction turns is meaningful only once such a
decomposition is fixed, and we treat the decomposition as a necessary input rather than a
derived result. Naming it is preferable to letting it pass unremarked.

\subsection{The termination read epistemically}
\label{sec:hopf-information-bound}

The termination of the tower (Proposition~\ref{prop:termination}) has so far been stated
algebraically: the Cayley--Dickson sequence acquires zero divisors at the sedenions and the Hopf
map can no longer be globally defined. There is a second reading of the same boundary, epistemic
rather than algebraic, which we offer as a motivated conjecture rather than a settled result,
because it connects the tower's endpoint to a bound that arises independently in the foundations
of measurement.

A particle, in this framework, is a reified relation between events, and its internal content is
carried by the octonionic $A,B,C$ register --- the $S^{7}$ sector serving as internal degrees of
freedom in the channel that connects the events (\S\ref{sec:space-of-computations}). The
termination of the tower is then a statement about that content: a single such relation can
support genuine entanglement among at most three parties, because there is no fourth fibration
to carry a fourth independent party within the same irreducible relation. Every particle is, on
this reading, a \emph{completed} octonionic-level relation --- topped out at three parties not
by circumstance but because the tower admits no further rung. The information that can be
encoded in one relation-between-events, and hence extracted by one measurement of it, is
correspondingly bounded: it is the information of a completed three-party register, and no
measurement of a single particle can yield more.

This is exactly the kind of bound that the epistemic-restriction programmes reach from the other
direction. Spekkens' toy theory is built on a knowledge-balance principle that caps the
information obtainable about a system \autocite{spekkens2007toy}, and Szangolies derives such an
epistemic horizon not by positing it but from the same self-referential non-computability on
which the present tower rests \autocite{szangolies_j_2018}. That the two accounts --- the
algebraic termination of the Hopf tower and the epistemic bound on measurement information ---
should coincide is, we suggest, not a numerical accident but a shared root: both are
consequences of the self-referential non-computability of \S\ref{sec:self-reference}, wearing an
algebraic face in the one case and an epistemic face in the other. The Hurwitz--Adams ceiling
and the knowledge-balance ceiling would then be one limit seen twice, and Szangolies' programme,
which reaches $G_{\mathrm{SM}}$ by an independent route (\S\ref{sec:hopf-honest-debt}) and the
epistemic horizon from the same non-computable source, shares with this framework not merely a
conclusion but its origin.

The reading must be kept distinct from a different climb, with which it would otherwise be
conflated. The Wigner's-friend regress of \S\ref{sec:wigners-friend} ascends without end: each
observer may be measured by a further observer, each carrying a fresh gauge-private circle, and
that chain has no top --- it is the bump in the carpet, unbounded in length. The termination
bounds something else: not the length of the chain of observers, but the internal structure of
each single relation, each particle, that the chain is built from. The regress is unbounded
precisely because each of its links is a bounded, completed three-party object;
termination-of-each-link and non-termination-of-the-chain are not in tension but complementary,
the former individuating the links of the latter. What the bound caps is the information in one
particle; what the regress leaves uncapped is the number of particles and observers that can be
chained. The precise identification of ``the information of a completed three-party register''
with the formal measures of the epistemic-restriction literature is the technical step that
would turn this conjecture into a theorem, and we flag it as owed; the structural claim --- that
the tower's endpoint and the measurement bound share the non-computable root --- is what we
advance here.

\subsection{The termination read compositionally}
\label{sec:hopf-compositional}

There is a third termination theorem, proved in a community that has so far had no contact with
the two joined above. Coecke and Kissinger showed that the GHZ and W classes --- realised as the
commutative special and the anti-special Frobenius algebras of categorical quantum mechanics ---
are \emph{compositionally complete}: every $N$-qubit entanglement class under SLOCC is generated
by composing these tripartite primitives with bipartite entanglement and product states
\autocite{coeckekissinger2010compositional}. Above three parties nothing new appears; the
entanglement alphabet closes at exactly the point where the fibration tower terminates. The
mutual silence of the two literatures is itself striking: Coecke and Kissinger make no reference
to Hopf fibrations or division algebras, and the fibration literature on entanglement
\autocite{mosseri2001hopf,bernevigchen2003geometry} nowhere cites compositional completeness.
Two classifications, in different mathematics, draw the same line in the same place, unaware of
one another. That silence has, in the months this paper was in draft, just begun to break from
the categorical side, with a projective, base-level reading of the same calculus
\autocite{burtonanwar2026meromorphic}; we engage it below.

The framework's conjecture is that this is one termination wearing a third face: the alphabet
has the letters it has --- Bell, GHZ, W --- because the parallelisable spheres have the
dimensions they have. There is partial evidence at every rung. At one qubit, the commutative
special Frobenius algebras on $\mathbb{C}^{2}$ are classified by orthonormal bases --- points of
$S^{2}$, the base of the first fibration --- so that the copied basis plays the role of a
section of the bundle and the spider theorem expresses its idempotence. At two qubits, the
$SU(2)$ fibre of $S^{7} \to S^{4}$ carries exactly the entanglement that survives local
operations, and L\'{e}vay's result that the bundle's natural connection is the $SU(2)$
Yang--Mills instanton \autocite{levay2004geometry} stands on the geometric side where the
interaction of two Frobenius structures stands on the categorical one. At three qubits, the
split of D\"{u}r, Vidal and Cirac \autocite{durvidalcirac2000three} mirrors the algebra exactly:
GHZ is genuinely tripartite, demanding the full non-associative octonion, while W is
pairwise-robust and housed in the associative subalgebra $\mathbb{H} \subset \mathbb{O}$ --- the
special/anti-special Frobenius distinction in geometric dress.

The first of the four conditions the correspondence requires can be discharged at once, and we
do so here, because a problem whose first step has been walked is a different kind of invitation
from one merely posed.

\begin{proposition}[The one-qubit dictionary]\label{prop:f1-dictionary} The $\dagger$-special
commutative Frobenius algebras on $\mathbb{C}^{2}$ correspond bijectively to
\emph{lifted antipodal skeletons} of the complex Hopf fibration $\pi \colon S^{3} \to S^{2}$:
unordered antipodal pairs $P \subset S^{2}$ together with a section of $\pi$ over $P$. The
copyable states of a structure are exactly the section's image, so the lifts are recorded
faithfully --- the projectively trivial centre of $U(2)$ acts non-trivially on structures ---
and the correspondence is $U(2)$-equivariant, the ordered structures forming a single free and
transitive $U(2)$-orbit. The channel a structure induces on density matrices is dephasing in its
skeleton and is independent of the section: the fibre freedom, a torsor for
$\mathrm{Map}(P, U(1))$, is exactly the data absent from the classical record a structure
induces.
\end{proposition}

\begin{proof}
By the classification of Coecke, Pavlovic and Vicary, such structures correspond to orthonormal
bases, recovered as the sets of copyable states \autocite{coeckepavlovicvicary2013orthogonal}.
An orthonormal basis of $\mathbb{C}^{2}$ is two orthogonal rays --- an antipodal pair on the
Bloch sphere --- together with a unit representative of each, which is precisely a section of
$\pi$ over the pair; and $\delta(e^{i\theta}x) = e^{i\theta}\, x \otimes x$ against
$(e^{i\theta}x) \otimes (e^{i\theta}x) = e^{2i\theta}\, x \otimes x$ pins the phases.
Equivariance, the action of the centre, the orbit structure and the induced channel are direct
computations; an executable verification of every finite check accompanies the code release.
\end{proof}

Two remarks fix the proposition's scope. It corrects a natural first reading: a classical
structure does not project to the base but selects lifts into the total space --- classicality
is section data, and what the copy operation broadcasts is a lift, not a ray. And it claims
nothing topological about cloning: the finiteness of the copyable set is linear, structural
no-cloning residing in non-naturality \autocite{abramskyheunen2012hstar}, so the first place the
bundle's global topology can bite is the quaternionic rung, where the degree-one clutching map
stands ready to be compared with the strong-complementarity law governing an interacting pair of
structures. The dictionary is, finally, the complement of the projective treatment cited above:
Burton and Anwar pass functorially to the base and deliberately discard the fibre
\autocite{burtonanwar2026meromorphic}, and the classical structures turn out to live in exactly
the data their projectivisation forgets --- the two readings compose, and the resulting triangle
is a consistency condition available to constrain the remaining conditions of the
correspondence.

\begin{openproblem}[The Hopf-fibration--Frobenius correspondence]\label{op:hopf-frobenius} Let
$\mathbf{Hopf}$ denote the category of the three fibrations
$S^{2^{k}-1} \hookrightarrow S^{2^{k+1}-1} \to S^{2^{k}}$, $k = 1, 2, 3$, with bundle maps as
morphisms, and let $\mathbf{Frob}$ denote the Frobenius algebras on $\mathbb{C}^{2}$ in
$\mathbf{FdHilb}$, with the commutative-special (GHZ) and anti-special (W) classes
distinguished. Construct a functor $F \colon \mathbf{Hopf} \to \mathbf{Frob}$ --- or prove that
none exists --- under which the complex fibration yields the single-qubit classical structures
--- Proposition~\ref{prop:f1-dictionary} identifies these as the lifted antipodal skeletons,
section data over two base points rather than the points alone; the quaternionic fibration
yields the interaction of two such structures, with the $SU(2)$ fibre carrying the entanglement
that survives the classical quotient; the octonionic fibration yields the GHZ/W pair, GHZ on the
full non-associative fibre and W on the associative $\mathbb{H} \subset \mathbb{O}$; and the
composition by which Coecke--Kissinger generate the $N$-qubit classes corresponds to a
well-defined geometric operation on the bundles. If such an $F$ exists, compositional
completeness follows from Adams's theorem as a topological corollary --- there is no fourth
fibration to supply a fourth primitive --- and if $F$ is faithful enough that the implication
also runs back, the termination of the tower and the completeness of the alphabet are equivalent
expressions of a single classification. The first condition is discharged by
Proposition~\ref{prop:f1-dictionary}; what remains open is the quaternionic and octonionic
conditions and, above all, the correspondence of compositions. Both open conditions are already
constrained: the quaternionic fibre's physical content --- the recoverable relative phase and
its purity ledger --- is established (\S\ref{sec:wf-reconciliation}), leaving the
clutching-map-to-strong-complementarity comparison as the structural residue; and any candidate
functor must reproduce the proved canonical-orbit images, GHZ to the colour triplet and W to the
singlet (\S\ref{sec:canonical-orbits}, Appendix~\ref{app:ghz-w-proofs}). A proof that no such
functor exists, or a correspondence at the level of the classifications without functoriality on
compositions, would each also resolve the problem: the first would make the agreement of two
independent classifications a fact demanding its own explanation; the second would establish the
weaker reading.
\end{openproblem}

A resolution in the affirmative would do three things. It would explain compositional
completeness rather than merely prove it --- a corollary of the division-algebra structure of
quantum mechanics rather than an independent algebraic fact. It would give the diagrammatic
language of this paper's own interaction circuits a topological warrant: the GHZ/W spider
calculus would be the unique language the topology permits, not a convenient notation. And it
would open a two-way bridge, translating problems on compositions of Frobenius structures into
operations on bundles and back. The problem sits at the intersection of categorical quantum
mechanics, the quantum-information geometry of the fibrations, and the division-algebra
approaches to particle physics, and it is posed here so that any of the three communities can
take it up in its own language. It is perhaps worth noticing that both classifications already
answer to the same name --- the fibrations are Heinz Hopf's, and so, through the bialgebra
structure of the diagrammatic calculus, are the Hopf algebras --- a coincidence of eponymy that
the conjecture proposes to promote to a theorem. With the epistemic reading above, an
affirmative resolution would leave the tower's endpoint as one limit seen four ways: algebraic,
geometric, compositional, epistemic --- the last being the fixed-point classification of Open
Problem~\ref{op:sigma-functor}, with which this problem is cross-linked, each constraining the
other.

\subsection{What this sets up}

The entanglement-sensitivity of the Hopf maps is not a side feature but the hinge of the
physical reading that follows. Because separability is exactly the collapse of the base to a
product of Bloch spheres, the structure of the entanglement --- the SLOCC class of a multi-qubit
state, and within it the canonical maximally entangled orbit --- is carried by precisely the
degrees of freedom on which $G_{\mathrm{SM}}$ acts in Proposition~\ref{prop:gsm}. This is what
makes it possible, in \S\ref{sec:entanglement-classes}, to read the two genuinely tripartite
entanglement classes onto the $\mathbb{C} \oplus \mathbb{C}^3$ split through their canonical
representatives: the GHZ orbit confined to the colour-triplet $\mathbb{C}^3$ sector, the W orbit
reaching the $\mathbb{C}$ singlet, recovering the quark--lepton division not as an assignment
but as a consequence of the entanglement structure. The scope of those statements --- canonical
orbits rather than whole classes, and the identification of the canonical orbits with the
physical particles --- is made precise there. The Hopf fibrations are the bridge; the
entanglement classes, treated next, are the traffic that crosses it.

%  The section is position-portable: all adjacency phrasing was
%  retargeted to explicit \S\ref on 2026-07-06. The S^7 hand-off
%  sentence to sec:space-of-computations WRITTEN 2026-07-06 (spacetime
%  file read; the closing now hands first to its immediate consumer,
%  then to hopf-justification). Upload verifications COMPLETE: L2 and
%  quantum-postulates do not consume the ladder; WF's references are
%  forward-voiced or self-carrying.
%
%  DIVISION OF LABOUR (deliberate, to avoid overlap):
%   - sec:self-reference (undrafted): self-reference -> complex
%     amplitudes and the Born rule; why not a real base. This section
%     TAKES that as given and starts at C.
%   - THIS section: the algebra -- the doubling construction, the
%     ledger of losses-as-physics, the internal-algebras-not-a-new-QM
%     scoping (Adler pre-empt), the two appearances of the ladder
%     (register / measurement climb, stated lightly), and Hurwitz as
%     the algebraic termination with the committed zero-divisor
%     reading.
%   - sec:hopf-justification: the geometry (Adams, parallelisability,
%     the fibrations), the join of the two classifications, the
%     entanglement-sensitivity literature, the C-direction honest debt,
%     and the epistemic reading of termination.
%
%  COMMITTED READINGS USED (sources re-read 2026-07-03):
%   - Loss ledger and level dictionary: cayley_dickson_qubits_for_review
%     sec.5 ("each step loses a symmetry, and each lost symmetry is
%     identified with a physical content").
%   - Termination reading: grounding_argument_gauge_group ("the
%     sedenions have zero divisors and cannot support consistent
%     information flow around loops") -- leads; norm-composition-as-
%     probability given as its amplitude-level restatement.
%   - Internal-not-base scoping: cayley_dickson_qubits_for_review Q2 and
%     sec.3.2 caveat (Adler objection; C-qubit treated geometrically).
%
%  HONEST DEBT: ladder proved closed at the top (Hurwitz), derived at
%  the floor. The rung-compulsion revisit-trigger FIRED 2026-07-06:
%  sec:self-reference and op:sigma-functor now exist, so the thesis is
%  minted as op:ladder-compulsion (with the CD-conjugation-as-sigma
%  candidate route), inheriting op:sigma-functor's status.
%
%  AUDITED 2026-07-06 (final pass, third stop): all four citations
%  used; all cross-refs live; CD facts, dims, alternativity, Hurwitz
%  statement verified; Adler subsection delivers everything the
%  register's compressed pointer promises; file pure ASCII; no bare
%  "frame". Fixes: opening rewritten for both predecessors (continuity
%  was inserted between after drafting); "associator's non-vanishing is
%  mass itself" recalibrated to the register's debt-against-selection
%  defusal; the O->S ledger row and Hurwitz subsection reconciled with
%  the fourth-qubit forward remark ("the terminus closes the ladder of
%  forces; what survives past it is, at most, flavour").
%
%  CROSS-REFS: sec:self-reference [sibling dangle -- resolves when that
%  section is drafted], sec:hopf-justification, sec:hopf-information-bound,
%  sec:register, sec:entanglement-classes, sec:generations, sec:masses,
%  sec:wigners-friend. All others exist.
%  CITATIONS: existing keys baez2002octonions, hurwitz1898composition,
%  furey2018three. ONE new key: adler1995quaternionic (RIS in
%  division_algebras_refs_2026-07-03.ris).
%  British English; \autocite; no bullets in prose; the ledger is a
%  table because the comparison is the argument.
% =====================================================================

\section{The Division-Algebra Ladder}
\label{sec:division-algebras}

The foundations behind us end at amplitudes twice over: natively, where the self-referential
argument forced the description of a system into the system's own state space and priced
consistency at a number system whose parts can cancel --- complex amplitudes, and the Born rule
that squares them (\S\ref{sec:self-reference}) --- and independently, where the regime's
continuity delivers the same field through Hardy's theorem (\S\ref{sec:continuity}), the
convergence of the two branches being evidence rather than accident. This section asks the question
both routes leave open. Why
should the number system stop at $\mathbb{C}$? Ordinary quantum mechanics stops there by postulate;
a framework that has forbidden itself postulates must answer, and the answer is that the climb
neither stops at $\mathbb{C}$ nor continues forever. There are exactly two more rungs, each
purchased by surrendering an algebraic law and each paying for itself in physics, and the ladder is
closed at the top by a theorem of 1898. This section sets out the algebra of that ladder: the single
construction that climbs it, the ledger of what each rung costs and buys, and the theorem that ends
it. The geometry of the same ladder --- that each rung acts on a sphere, that exactly those spheres
are parallelisable, and that the algebraic and geometric classifications are one --- is the business
of the geometry section (\S\ref{sec:hopf-justification}).

\subsection{One construction climbs the whole ladder}

The rungs are not a miscellany; a single construction generates them all. Given an algebra with a
conjugation, the Cayley--Dickson construction builds its double on ordered pairs, with a twisted
multiplication in which the conjugate enters and a conjugation $(a,b)^{*} = (a^{*}, -b)$
\autocite{baez2002octonions}. (Conventions differ in the placement of the conjugations; all yield
isomorphic algebras.) Applied repeatedly it produces the sequence
$\mathbb{R} \to \mathbb{C} \to \mathbb{H} \to \mathbb{O} \to \mathbb{S}$, of real dimensions
$1, 2, 4, 8, 16$: the reals, the complex numbers, the quaternions, the octonions, the sedenions.
The ground floor is classical: $\mathbb{R}$ is ordered and positive-definite, the arithmetic of
classical probability. The first doubling is therefore the quantum move itself --- the passage from
probabilities to amplitudes that can interfere --- and it is not made here by choice: it is the
derivation of \S\ref{sec:self-reference}, already behind us. What this section must account for is
the rest of the climb.

\subsection{The ledger}

In mathematics the doublings are usually told as a story of degeneration: each algebra is worse
behaved than the last, and the sequence is studied for where it breaks. This framework reads the
same ledger in the opposite direction. Each rung is bought by surrendering a law, and each
surrendered law is exactly a piece of physics --- the losses are the assets, and the comparison is
sharp enough to be the argument:

\begin{center}
\begin{tabular}{llll}
\hline
doubling & dimension & what is lost & what the loss becomes \\
\hline
$\mathbb{R} \to \mathbb{C}$ & $1 \to 2$ & the order of the reals & phase and interference; amplitudes \\
 & & & that cancel; the Born rule (\S\ref{sec:self-reference}) \\
$\mathbb{C} \to \mathbb{H}$ & $2 \to 4$ & commutativity & non-commuting actions: weak isospin \\
 & & & $SU(2)$, the $B$ register (\S\ref{sec:register}) \\
$\mathbb{H} \to \mathbb{O}$ & $4 \to 8$ & associativity & the associator: colour, chirality, the \\
 & & & generations, and mass as debt \\
 & & & (\S\ref{sec:generations}, \S\ref{sec:masses}) \\
$\mathbb{O} \to \mathbb{S}$ & $8 \to 16$ & division itself & no further force: zero divisors appear \\
 & & (zero divisors) & and the ladder of sectors ends (below) \\
\hline
\end{tabular}
\end{center}

Walking the rows: the order lost at the first doubling is the possibility of interference --- parts
of a probability that can oppose one another --- which is where the derivation of \S\ref{sec:self-reference}
lives. The commutativity lost at the second is the existence of actions whose order matters; the
unit quaternions form the group $SU(2)$, and the framework houses weak isospin there, as the
internal algebra of the $B$ register whose order-sensitive, non-commuting dynamics the interaction sections develop
(\S\ref{sec:register}, \S\ref{sec:entanglement-classes}). The associativity lost at the third ---
the octonions retain the weaker property of alternativity \autocite{baez2002octonions} --- creates
the associator $[a,b,c] = (ab)c - a(bc)$, and the associator is where most of this paper's physics
is transacted: the stabiliser of a preferred complex direction inside the octonions' automorphisms
is the colour group (\S\ref{sec:hopf-justification}); the triality structure and its $K_3$
obstruction generate the three generations and the sterile channel (\S\ref{sec:generations}); and
the associator is the currency of mass --- the debt against the selected direction that
\S\ref{sec:masses} prices, owed by whatever action moves the selection and not by residence at
the rung, which is the defusal \S\ref{sec:register} makes explicit. The algebraic facts
in the third column are theorems; the identifications in the fourth are established in the sections
cited, each carrying its own epistemic flag there, and are not re-litigated here. What this table
asserts on its own account is only the correspondence's shape: one law surrendered per rung, one
sector of physics gained per law.

\subsection{Internal algebras, not a new quantum mechanics}

A reader versed in the quaternionic-quantum-mechanics literature will at this point raise a fair
objection: quantum theory over $\mathbb{H}$ as a foundation is known to be troubled --- the
tensor-product and composite-system difficulties set out by Adler and others
\autocite{adler1995quaternionic} --- and quantum theory over $\mathbb{O}$ founders earlier still,
since a non-associative algebra does not furnish a Hilbert space at all. The framework's position
must therefore be stated exactly, because it does not make the claim the objection targets. The base
of the physical Hilbert space here is $\mathbb{C}$ and remains $\mathbb{C}$; the higher algebras
enter as the \emph{internal} algebras of the register --- $\mathbb{H}$ for the $B$ qubit,
$\mathbb{O}$ for the $C$ qubit (\S\ref{sec:register}) --- in a tensor structure whose outer theory is
ordinary complex quantum mechanics. This is the same move by which the division algebras enter the
constructions of Furey and the related programmes without any proposal to rebase quantum mechanics
itself \autocite{furey2018three}. For the octonionic register the point is sharper still:
$\mathbb{O}^{2}$ is not a Hilbert space, and the framework accordingly treats the $C$ qubit not in
the Dirac--von Neumann manner but geometrically, as the sphere $S^{15}$ with its Hopf structure ---
which is precisely the form in which \S\ref{sec:hopf-justification} needs it. Nothing in this paper therefore
stands or falls with quaternionic or octonionic quantum mechanics as foundations; the ladder is
climbed by the registers' internal algebras, not by the theory's base field.

\subsection{The same ladder, twice}

The ladder appears in this framework in two guises, and it will help the reader to have them
distinguished now. Statically, it is the anatomy of a single relation: the three qubits of the
register take their values one rung apart --- $A$ over $\mathbb{C}$, $B$ over $\mathbb{H}$, $C$ over
$\mathbb{O}$ --- so that every particle carries the whole ladder as its internal structure
(\S\ref{sec:register}). Dynamically, it is the shape of measurement: a single qubit's pure states
fill $S^{3}$, an entangling measurement carries the joint state into $S^{7}$, and a third party into
$S^{15}$ --- the correspondence between qubit entanglement and the Hopf fibrations that
\S\ref{sec:hopf-justification} takes from a mature literature --- so that each act of measurement ascends one rung --- the climb the reader has already made,
narrated in \S\ref{sec:wigners-friend} --- and every completed relation is a terminated climb, which is
the reading \S\ref{sec:hopf-information-bound} gives the tower's endpoint. One ladder, seen twice:
as the anatomy of a relation and as the dynamics of relating. The framework regards this double
appearance not as a coincidence but as the same structure viewed statically and in motion, and the
sections cited carry that identification's weight.

\subsection{Hurwitz closes the ladder}

The construction itself would double forever; what stops it is a classification theorem. By
Hurwitz's theorem, the normed division algebras --- those in which the norm composes,
$\lvert ab\rvert = \lvert a\rvert\,\lvert b\rvert$, and no non-zero elements multiply to zero ---
are exactly $\mathbb{R}$, $\mathbb{C}$, $\mathbb{H}$ and $\mathbb{O}$
\autocite{hurwitz1898composition,baez2002octonions}. The next doubling, the sedenions, has zero
divisors: non-zero elements whose product vanishes. The framework's reading of this failure is the
one that matters physically. Zero divisors cannot support consistent information flow around loops
--- a cycle of compositions can silently annihilate what it carries --- and a framework in which
closed computational paths are load-bearing (\S\ref{sec:hopf-justification},
\S\ref{sec:space-of-computations}) cannot be built on an algebra that permits it. At the level of
amplitudes the same failure reads: the composition law
$\lvert ab\rvert = \lvert a\rvert\,\lvert b\rvert$, which carries probability multiplicatively
through every rung of the ladder, fails beyond the octonions, so that two non-vanishing amplitudes
could compose to nothing and probability would be destroyed by mere composition. The last algebra
fit to carry amplitudes is $\mathbb{O}$. One rung past the terminus is not entirely silent,
however: the sedenions' internal structure --- their automorphism group and their octonion
subalgebras --- is the recorded candidate arena for tagging the generations
(\S\ref{sec:interactions}), and it enters exactly as this section's scoping permits, internally
and never as a base. The terminus closes the ladder of forces; what survives past it is, at most,
flavour. This is the algebraic half of the termination; the
geometric half --- that no sphere above $S^{7}$ is parallelisable --- and the theorem joining the
two halves are the opening business of \S\ref{sec:hopf-justification}, and the epistemic reading of the same endpoint,
as a bound on the information a single relation can carry, is given in
\S\ref{sec:hopf-information-bound}.

\subsection{The honest debt}

As in the geometry section, we close by stating plainly what is proved, what is established, and what is
owed. Proved: the construction, the losses, and the termination are theorems --- the Cayley--Dickson
doubling and its ledger of surrendered laws are standard algebra, and Hurwitz's classification is
among the durable results of the subject. Established: the physical cashing of each rung --- the
Born rule at $\mathbb{C}$, weak isospin at $\mathbb{H}$, colour, generations and mass at
$\mathbb{O}$ --- is carried out in the sections cited above, each under its own epistemic flag. Owed,
and stated without disguise: the ladder as presented here is closed from above and derived at its
floor, but it is not yet derived rung by rung from below. Hurwitz forbids a further rung; the
argument of \S\ref{sec:self-reference} forces the first doubling, from the ordered arithmetic of
classical probability to amplitudes that interfere. That the same self-referential closure, applied
to a system whose amplitudes already take values at one rung, compels the next --- that relational
closure at one level forces precisely the Cayley--Dickson double of that level --- is the
framework's reading of the climb, and it is a thesis, not a theorem --- but the formal apparatus
against which it can be posed now exists, and we pose it.

\begin{openproblem}[The compulsion of the intermediate rungs]\label{op:ladder-compulsion}
Derive, within the categorical apparatus of Open Problem~\ref{op:sigma-functor}, that
self-referential closure at one rung compels the Cayley--Dickson double of that rung, so that the
intermediate steps $\mathbb{C} \to \mathbb{H} \to \mathbb{O}$ follow from the same principle
that forces the floor. The candidate route is the construction's own shape: the Cayley--Dickson
pair with its conjugation, $(a, b)^{*} = (a^{*}, -b)$, is an involution-equipped doubling, and the
route reads it as the algebraic shadow of adjoining the describer --- closure of level-$A$
descriptions under the inside/outside exchange requiring exactly the twisted pair. The problem
inherits the status of Open Problem~\ref{op:sigma-functor}, whose apparatus it presupposes.
\end{openproblem}

If the derivation lands, the whole ladder becomes a consequence of the same self-reference that
opened it. Until then, the honest statement is that the ladder's necessity is established at its
ceiling and at its floor, while the compulsion of the intermediate rungs is inherited from the
physics they are shown to carry rather than derived in advance. What this section hands on is, accordingly, an algebraic ladder with a floor, a ceiling and a
ledger. Its first consumer is immediate: \S\ref{sec:space-of-computations} spends the terminal
fact --- the non-associativity of $S^{7}$'s own multiplication --- as geometry, excluding the
seven-sphere as a ruler and fixing the dimension of space at three.
\S\ref{sec:hopf-justification} then supplies the geometry of the fibrations themselves, and the
join.

%  (Architecture call SETTLED by the sync manifest: standalone.)
%
%  CITATIONS -- existing keys (as supplied by JS 2026-07-04):
%    Hardy2001a, Aaronson2004h, Thomson_1954, Lawvere_1969,
%    Turing_1937, muller_p_2011, Brukner2009.
%  NEW keys (RIS in continuity_refs_2026-07-04.ris):
%    cantor1891diagonal, benacerraf1962tasks, broeckertomdieck1985compact.
%  Deliberately excluded per JS: Chiribella--D'Ariano--Perinotti;
%  Bargmann/Wigner's theorem. Yanofsky2003 now paired on the
%  Lawvere-shaped clause (the pending-key excuse was void: the key
%  exists and sec:self-reference uses it).
%
%  CROSS-REFS: sec:self-reference now DRAFTED (2026-07-04, later the
%  same day) -- no longer dangling; only sec:fanout-boundary (and its
%  sibling sec:fanout-direction) -- DRAFTED 2026-07-06.
%
%  AUDITED 2026-07-06 (final pass, second stop): both proofs
%  re-verified by hand (Lagrange divisibility sentence exact; the
%  Baire countable-escape chain sound); caveat-discharge with
%  sec:self-reference confirmed two-way (already present); file pure
%  ASCII; no pre-definition "frame"; all citations used as declared.
%  Fixes: Yanofsky2003 paired on the Lawvere clause; the op's two
%  candidate routes anchored to printed propositions (f1-dictionary
%  for the divisibility premise's referent, inside-outside for the
%  closedness demand) -- statuses unchanged, referents sharpened;
%  header staleness corrected. Deliberate FORWARD references --
%  prop:f1-dictionary and sec:hopf-justification (defined later in the
%  paper; fine under two-pass LaTeX and flagged in prose),
%  sec:division-algebras, sec:space-of-computations.
%
%  DECISIONS ENCODED: op:limit-realisability minted (the owed arrow,
%  priced: exactly two derivations, divisibility and closedness);
%  Lagrange used twice (finite exclusion; component-group triviality);
%  the (1+i)/2 lamp value NOT stated here -- "a superposed fixed point"
%  with the complex-probability treatment deferred to
%  sec:self-reference; the two-branch over-determination stated to
%  pre-empt a circularity reading; Thomson spelling throughout.
% =====================================================================

\section{Continuity Is Not an Axiom}
\label{sec:continuity}

Hardy derives quantum theory from five axioms \autocite{Hardy2001a}. Four of them are satisfied
by classical probability theory; the fifth --- that a \emph{continuous} reversible
transformation connects any two pure states --- is the entire difference, the single postulate
that turns the classical simplex into the complex Hilbert space. Every reconstruction programme
pays an equivalent toll: continuous reversibility in Masanes and M\"{u}ller
\autocite{mullerp2011}, the corresponding postulate in Daki\'{c} and Brukner
\autocite{Brukner2009}. The axiom is load-bearing, and in every programme it is unexplained: it
is where the quantumness is put in by hand. The claim of this section is that in the present
framework it is not an axiom. Continuity is the fifth face of the pre-measurement regime that self-reference forces, and what
stands between the framework and a full derivation is not a postulate about values but one
premise about transformations --- that every transformation the inside performs has a stationary
outside presentation --- derived in outline below from the inside/outside involution of
\S\ref{sec:self-reference} and stated as an open problem at the end of the section with nothing
hidden in it. The continuum itself, the completion that every probability theory adopts, is
named for what it is: the real line, shared with classical probability and carrying no quantum
content.

One structural remark prevents a misreading. The framework reaches complex amplitudes natively in
\S\ref{sec:self-reference}, by the negative-probability route; the present section reaches the same
destination by an independent branch --- the regime's continuity, then Hardy's theorem. The two
branches converging on one structure is over-determination, not circularity, and the convergence is
itself evidence that the destination is forced.

\subsection{Five faces of one regime}

The regime before any measurement --- before any \textsc{fanout} --- can be characterised five
equivalent ways: negative probability is available; anti-events exist, every event being undoable;
no classical record exists; no preferred basis is selected; and a continuous reversible
transformation connects any two pure states. These are not five premises to be combined. They are
one regime under five descriptions, and \textsc{fanout} (\S\ref{sec:fanout-boundary}) is the single
event that fails all five at once: it writes the record, selects the basis, kills the anti-events,
discretises the transformation group, and switches off the negativity. The equivalence of the first
two faces is the anti-event identity of \S\ref{sec:self-reference}; the third is the boundary's
definition. The work of this section is the equivalence of the last two faces, and then the route
from the regime's existence to its continuity. The fourth face has, since
Proposition~\ref{prop:f1-dictionary} (\S\ref{sec:hopf-justification} --- a deliberate forward
reference), an exact form: a basis-selection is the commitment of a classical structure, a lifted
antipodal skeleton, so ``no preferred basis'' means precisely that no copy structure has fired.

\subsection{Basis selection is the collapse of continuity}

\begin{proposition}[Continuity forbids a selected basis]\label{prop:basis-discreteness}
If a finite set of pure states is privileged --- a preferred basis --- then every reversible
transformation respecting the privilege carries the set to itself and acts on it as a permutation;
the respecting transformations therefore realise only a finite group of motions among the
distinguished states, and no continuous family of them connects distinct pure states of the set.
Contrapositively, continuous reversibility entails that no basis is selected.
\end{proposition}

\begin{proof}
A privileged finite set is by definition carried to itself by every structure-preserving reversible
map, whose induced action on the set is a permutation of finitely many points. A continuous path of
such maps induces a continuous path in a finite discrete set, which is constant. Hence a continuous
reversible family moving one distinguished pure state to another respects no such privilege, and if
every pair of pure states is so connected, no finite set is privileged.
\end{proof}

The converse is where the historical difficulty lies: the absence of a privileged finite set does
not by itself force continuity, for the transformation group could in principle be infinite and
discrete. Two structural properties close every such escape at once.

\begin{proposition}[Divisible and closed implies connected]\label{prop:divisible-closed}
Let $G$ be the group of reversible transformations of the pre-measurement regime, acting on the
state body. If $G$ is divisible --- every element has an $n$-th root for every $n$ --- and
topologically closed, then $G$ is a compact connected Lie group; and if $G$ moreover acts
transitively on the pure states, then a continuous reversible transformation connects any two of
them: Hardy's fifth axiom.
\end{proposition}

\begin{proof}
The state body is compact and convex with nonempty interior, so its symmetry group is compact, and
a closed subgroup of a compact matrix group is a compact Lie group with finitely many connected
components \autocite{broeckertomdieck1985compact}. The component group $G/G_{0}$ is a quotient of a
divisible group and hence divisible; but a nontrivial finite group is never divisible --- with
$\lvert H\rvert = m$ every element satisfies $h^{m} = e$, so no element other than the identity
possesses an $m$-th root --- and $G/G_{0}$ is finite, so it is trivial and $G = G_{0}$ is connected.
A connected Lie group is path-connected, and transitivity then supplies the continuous reversible
path between any two pure states.
\end{proof}

The same one-line fact of Lagrange closes the remaining escapes. Finite groups are excluded
directly, being nontrivial and never divisible. Countable groups are excluded because a
countable closed subgroup of a compact group is compact, has an isolated point by the Baire
category theorem, and is therefore discrete and finite. Uncountable disconnected groups are
excluded by the component-group argument of the proposition. Every route to a discrete
transformation group therefore passes through a failure of divisibility or a failure of
closedness. Divisibility, however, asks for more than the conclusion needs, and the next
proposition removes it.

\begin{proposition}[Continuity from a connected state space]\label{prop:connected-transitive}
Let the pure states of a finite-dimensional system form a connected space $P$, and let the group
$G$ of reversible transformations of the pre-measurement regime be topologically closed and act
transitively on $P$. Then the identity component $G_0$ already acts transitively on $P$, and a
continuous reversible transformation connects any two pure states: Hardy's fifth axiom holds,
with no divisibility assumed.
\end{proposition}

\begin{proof}
As in Proposition~\ref{prop:divisible-closed}, closedness makes $G$ a compact Lie group with
finitely many connected components \autocite{broeckertomdieck1985compact}, so $G_0$ is a closed
normal subgroup of finite index and is path-connected. Each $G_0$-orbit in $P$ is the continuous
image of a compact group, hence compact, hence closed. Since $G$ is transitive and $G_0$ is
normal, the $G_0$-orbit of $gx$ is $g$ applied to the $G_0$-orbit of $x$, so the cosets of $G_0$
permute the orbits transitively and there are finitely many of them. A finite partition of $P$
into closed sets is a partition into open sets; $P$ is connected; so there is one orbit and
$G_0$ is transitive. Given pure states $x$ and $y$, choose $g \in G_0$ with $gx = y$ and a
continuous path $g_t$ in $G_0$ from the identity to $g$; then $g_t x$ is a continuous reversible
transformation from $x$ to $y$.
\end{proof}

The proposition relocates the premise. Divisibility was asked for because a transformation group
could be infinite and discrete; but a group acting transitively on a continuum cannot be
countable, and a closed group acting transitively on a connected space cannot be disconnected in
the way that matters, so what the argument needs from the regime is not the availability of
$n$-th roots but the connectedness of the state space itself. That is a statement about values:
that the pure states are all the unit vectors of $\mathbb{C}^{K}$ up to phase, the space
$\mathbb{CP}^{K-1}$, connected for $K \ge 2$, and hence that the amplitudes of
\S\ref{sec:self-reference} fill their continuum. The two standing foils show the premise doing
its work. The stabiliser fragment of quantum theory has complex amplitudes and is contextual,
yet it has finitely many pure states, its outcome probabilities take only the values $0$,
$\tfrac12$ and $1$, and its reversible transformations form a finite group, the Clifford group,
so Hardy's fifth axiom fails in it; Spekkens's toy theory \autocite{spekkens2007toy} is the same
with real amplitudes, its ignorance rationed by a knowledge-balance principle. Neither has a
connected state space, and neither is a model of the regime described here, whose credences take
every value. Divisibility, which both also lack --- the phase gate has no square root in the Clifford group
--- is a symptom of the discreteness, not its cause. What both lack can be said exactly, and the
next subsection says it: each contains a transformation that fixes no state it admits.

\subsection{The premises, and what they mean}

Three premises carry Proposition~\ref{prop:connected-transitive}. Transitivity is the regime's
homogeneity: inside/outside equivalence leaves no pure state privileged, and a proper closed
invariant set of pure states would define a superselection observable, a classical record
predating any \textsc{fanout} (\S\ref{sec:fanout-boundary}). Connectedness of the state space is two statements, and only one of them is quantum. The first
is that the pure states are infinite and dense, with a fixed point for every transformation; the
second is that the state space is complete. Quantum theory over the algebraic numbers ---
amplitudes in the algebraic closure of $\mathbb{Q}$ within $\mathbb{C}$ --- has the first
without the second: it has $K = N^{2}$, the Born rule, superposition, entanglement and Bell
violation at the Tsirelson value, every unitary has an algebraic eigenvector, and it is
countable and totally disconnected. So the whole quantum content of Hardy's axiom lies in the
density, and the completion is the real-number convention of every probability theory. Hardy
adopts it before his fifth axiom is stated, in the first, which makes probabilities real numbers
and gives his classical theory a real simplex of mixtures; this framework adopts it at the
Shannon step, where credences are declared to range over a connected domain so that the
continuation of the logarithm is unique (\S\ref{sec:complex-amplitudes}). Given that convention,
density of the pure states is connectedness of their closure. Divisibility, which an earlier
version of this section asked for, is not needed: it follows from the conclusion, every element
of a compact connected Lie group having roots of every order.

Closedness is the completion, limit-realisability: a Cauchy sequence of available
transformations has its limit available. It cannot be derived, and the reason should be stated
rather than hidden. Any collection of first-order principles about states, transformations and
probabilities that has an infinite model has a countable one, and a countable subset of
$\mathbb{CP}^{K-1}$ is totally disconnected; the founding non-computability does not escape
this, since the arithmetical complex numbers form a countable algebraically closed field
containing non-computable numbers and closed under every limit it can itself describe. Hardy's
axiom and closedness are second-order, quantifying over paths and over sequences, and a
second-order principle is grounded as Dedekind completeness is grounded or not at all. The
grounds are these. The completion of a dense set under the finite-precision metric is unique and
depends on no choice, whereas any countable alternative is a choice of definability level, an
extra law with no ground and no test; and the framework's own result is that no observer's
descriptive level captures the regime, since by \S\ref{sec:self-reference} every observer's
state exceeds its own level and the inside/outside principle makes that a law rather than a
perspective. The burden therefore falls on any restriction, not on the completion, which is the
least committal option and carries no quantum content. What limit-realisability adds to the line
is still worth stating exactly, because it is the corrected form of a natural intuition.
Cantor's diagonal \autocite{cantor1891diagonal} does not build the continuum; its output is an
infinite decimal --- a Cauchy sequence of rationals --- which exists \emph{as a point} only
because the line is complete. The diagonal presupposes limit-realisability and then shows that
the completed line outruns every internal listing. And the character of the completion is the
framework's own: the computable reals are countable \autocite{Turing_1937}, hence of measure
zero, so once the transformation group contains its limits almost all of them are
non-computable. Non-computability is a property of the completed continuum, not the reason for
it. The constructive reading, on which only describable limits exist, remains consistent; it is
the countable theory above, and the two readings differ in nothing an experiment can certify.

Density is where the section's title is earned, and it is derived from the paper's own premise.
Write $(\mathrm{I})$ for the statement that every admissible transformation of the regime is a
phase of some admissible classical structure: diagonal, that is, in the record basis of some
observer the regime contains, a phased spider of that basis in the calculus of
\S\ref{sec:hopf-justification} \autocite{coeckeduncan2011interacting} (VERIFY-SUPPORT). It says
of transformations what the global phase of \S\ref{sec:self-reference} says of states: there are
no absolute ones. It follows from three things the paper already holds. The inside/outside
involution is a genuine involution, so the outside presentation of the inside is a presentation
of the same object with nothing lost: the outside can have complete knowledge of the inside, a
potential record, the joint state lying in a structure the outside carries. Time is operation
count, internal to the system that operates (\S\ref{sec:space-of-computations}), so the outside
does not share the inside's count. And before any \textsc{fanout} that count is not classically
recorded, so the outside's description is pure. Let $g$ be the transformation the inside
performs per operation. The outside's complete knowledge cannot depend on the inside's count,
and it is pure, so it is a pure state of the inside's clock and system that is stationary under
``advance the clock and apply $g$''; the stationary states of that map are clock Fourier modes
tensored with eigenvectors of $g$ and nothing else \autocite{pagewootters1983evolution}
(VERIFY-SUPPORT). So the outside holds a structure in which $g$ acts by phases, which is
$(\mathrm{I})$; and the relative phase in that structure is the outside's ignorance of the
inside's computational count, the same ignorance the global phase records one level down. The
consequences are a theorem. A phase of a structure fixes the structure's copyable states, so
every admissible transformation fixes an admissible pure state; a finite transitive action has
an element moving every point (Jordan), so the pure states are infinite; the transformation
carrying a stabiliser state to the fixed point of the composite of a $Z$-phase and an $X$-phase
turns through $\arccos(1/\sqrt{3})$, an irrational fraction of a turn, so the group has an
element of infinite projective order and its orbits are dense. With the completion,
Proposition~\ref{prop:connected-transitive} fires and Hardy's fifth axiom holds. The two foils
fail $(\mathrm{I})$ exactly where they fail continuity: fourteen of the twenty-four Clifford
operations are phases of no Pauli basis and fix no stabiliser state, their fixed points being
the magic states; the toy theory's three-cycles are phases of no epistemic basis; and in each
the outside's knowledge of the inside under the order-three protocol is either count-dependent
or a mixture over the count, so each makes the inside's time classical before any record exists.
The derivation of $(\mathrm{I})$ is an outline, not a proof, and it is what the open problem
below asks for.

\subsection{Why the cut matters physically}

The regime's continuity is not decorative. Reversibility alone, over complex amplitudes, licenses
the full invertible group $GL(n, \mathbb{C})$, and Aaronson showed what that buys: superluminal
signalling, and computational power at the level of \textsf{PP} \autocite{Aaronson2004h}. The cut
from $GL(n, \mathbb{C})$ down to the unitary group is the conserved quadratic form, and continuity
is what entitles the theory to make it: with the fifth axiom in hand, Hardy's theorem delivers the
complex Hilbert space, the quadratic norm and the unitary group in one stroke \autocite{Hardy2001a}
--- cited here, not re-proved. One honest caveat travels with the citation. Hardy's reconstruction
uses his remaining axioms besides --- composition and subspace structure --- and these must be
derived rather than imported: the composition ledger of \S\ref{sec:division-algebras} and the
locality-of-tomography argument of \S\ref{sec:self-reference} carry that duty, and it is discharged
there, not silently assumed here.

\subsection{Super-tasks at the boundary}

The premise just priced has a fifty-year-old stress test in the philosophy of mathematics, and the
framework passes it in an unexpected way. A super-task --- infinitely many operations completed in
finite time \autocite{Thomson_1954} --- presupposes its terminal moment as an actual point of the
timeline: the accumulation point of a generic sequence of step-times exists only in the completed
line, and the defining question, the state \emph{after all the steps}, is meaningful only under
limit-realisability of time. Super-tasks are downstream consumers of exactly the completeness at
issue in this section. Thomson's lamp then exhibits both sides of the \textsc{fanout} boundary at
once. On the record side, the lamp's description fixes its state at every finite stage and thereby
fixes no limit at all: any terminal value is consistent with the whole task
\autocite{Thomson_1954,benacerraf1962tasks}. That is not a puzzle about lamps; it is the proof that
the classical record algebra is not limit-closed, and cannot be. On the amplitude side the limit
exists --- a superposed fixed point, whose complex-probability treatment belongs to
\S\ref{sec:self-reference} --- because that side \emph{is} limit-closed, which is the fifth face
itself. The boundary is thereby the edge of limit-realisability: complete on the reversible side,
incomplete on the record side, with the lamp standing on both sides at once as the boundary's oldest
exhibit. A final clause dissolves rather than answers the classic question. In this framework time
is operation count (\S\ref{sec:space-of-computations}), and ``after infinitely many operations''
indexes no operational moment: the lamp packs its endless \textsc{fanout}s into the emergent
continuum, and asking for its record at the accumulation point asks the discrete contents for a
value that only the completed container possesses. On or off is not undecided; it is unasked.

\subsection{What remains, priced}

\begin{openproblem}[Every transformation a phase of an inside structure]\label{op:limit-realisability}
Proposition~\ref{prop:connected-transitive} reduces Hardy's fifth axiom to three premises:
connectedness of the pure states, transitivity, and closedness. Transitivity is argued above
from the absence of records. Closedness is the completion, adopted on the grounds stated above
and carrying no quantum content. Connectedness is density plus that completion, and density
follows, by Jordan's theorem and the infinite-order argument, from $(\mathrm{I})$: that every
admissible transformation of the regime is a phase of some admissible classical structure. What
remains open is the rigour of the derivation of $(\mathrm{I})$ from the paper's premise. Derive,
inside the fixed-point locus $\mathrm{Fix}(\sigma)$ of Proposition~\ref{prop:inside-outside},
that every endomap is a phase of a classical structure in $\mathrm{Fix}(\sigma)$, from the
self-representing object, the involution $\sigma$ as a genuine involution, and the reading of
the inside presentation as a sequence indexed by operation count with no record before
\textsc{fanout}. The outline above does this with the stationary-state construction of Page and
Wootters; a derivation passes if it makes that outline a theorem of the calculus, with the
outside's count-independent complete description shown to exist and to be pure from the
involution alone. It fails before any detail is checked if every premise it uses holds in the
stabiliser fragment of quantum theory, in Spekkens's toy theory \autocite{spekkens2007toy}, in
the real qubit, or in quantum theory over the Gaussian rationals, all of which violate
$(\mathrm{I})$ outright and all of which make the inside's operation count classical before any
record exists. It also fails if it imports the completion as a premise about transformations:
the completion is the real line and has no derivation, and a derivation of $(\mathrm{I})$ must
hold in the countable algebraic theory, where $(\mathrm{I})$ is true and continuity is not.
\end{openproblem}

The accounting, in one sentence: the axiom that every reconstruction programme pays is, in this
framework, one derivation posed, of Lawvere's shape --- that every endomap of the fixed-point
locus is a phase of a classical structure in it --- and one convention named, the completion,
which classical probability pays too. And
the hand-on is already visible: the complex amplitudes that this section's continuity underwrites
through Hardy's theorem, and the complex amplitudes that \S\ref{sec:self-reference} derives natively,
are one and the same $\mathbb{C}$ --- the floor of the ladder that \S\ref{sec:division-algebras}
climbs.
